\documentclass[11pt]{article}
\usepackage[T1]{fontenc}
\usepackage[utf8]{inputenc}
\usepackage{lmodern}
\usepackage[margin=1in]{geometry}
\usepackage{amsmath,amssymb,mathtools}
\usepackage{booktabs,longtable,array}
\usepackage{fancyvrb}
\usepackage{graphicx}
\usepackage[hidelinks]{hyperref}
\usepackage{microtype}
\DeclareUnicodeCharacter{00B0}{\ensuremath{^{\circ}}}
\DeclareUnicodeCharacter{00B1}{\ensuremath{\pm}}
\DeclareUnicodeCharacter{00D7}{\ensuremath{\times}}
\DeclareUnicodeCharacter{03A8}{\ensuremath{\Psi}}
\DeclareUnicodeCharacter{03B1}{\ensuremath{\alpha}}
\DeclareUnicodeCharacter{03BC}{\ensuremath{\mu}}
\DeclareUnicodeCharacter{03C3}{\ensuremath{\sigma}}
\DeclareUnicodeCharacter{03C4}{\ensuremath{\tau}}
\DeclareUnicodeCharacter{03C9}{\ensuremath{\omega}}
\pdfobjcompresslevel=0
\pdfinfoomitdate=1
\pdftrailerid{}
\pdfsuppressptexinfo=15
\setlength{\parindent}{0pt}
\setlength{\parskip}{0.65em}
\setlength{\emergencystretch}{3em}
\title{dirac16complex: a textbook for students\\[0.5em]\large From school algebra to a 16-component spinor field in 4+4 dimensions: formalism, field equations, Kohn--Sham approximations, pairing theorems, matter and antimatter, and open problems}
\author{dirac16complex project (with Claude Opus 5.5)}
\date{September 2026}
\begin{document}
\maketitle
\tableofcontents
\newpage
\setcounter{section}{-1}
\begin{abstract}
This book teaches, starting from school algebra and the calculus of one variable, everything that the dirac16complex project has set up, derived, proved and computed, and it says with equal care what the project has not established. The object of study is a field with sixteen complex components on an eight-dimensional space with four space-like and four time-like directions (signature (4,4), coordinates $x_0,\dots,x_7$, time $x_4$), in two versions: dirac16complex, whose components are anticommuting (Grassmann) quantities, and dirac16complex00, whose components are ordinary commuting complex numbers. Part I builds the mathematics: vectors, matrices, complex numbers, indices, derivatives in several variables, metrics and signatures, Clifford algebras and spinors, split octonions, and curved space with its spin connection. Part II builds the field theory: Lagrangians, field equations, the energy--momentum tensor, kinetic and potential energy, pressure and equations of state in an arbitrary gravitational field, for both fields, and canonical quantization, which in signature (4,4) leads to an indefinite (Krein) state space and to unstable modes along the extra times. Part III treats the primordial gravitational field of the author's notebook, the errors found in the notebook and their corrections, the numerical methods, and the numerical experiments on dark matter and dark energy. Part IV introduces density functional theory from zero and derives the Kohn--Sham approximation for both fields in the static primordial field, whose exchange terms differ only in sign. The ground and first excited states of dirac16complex are computed by a Rust solver and compared with an independent reference solver in 63 checks, of which 62 pass; the failed one is a deviation of up to $2.2\times10^{-6}$ (about twice the tolerance) in two deep levels of one state, whose cause is open. Those of dirac16complex00 are computed in a single configuration. Part V proves three exact pairing theorems that relate a universe of mass $+M$ to one of mass $-M$, and then examines two questions without overstatement. Does the big bang create universes in pairs? A pair of classical fields related by the chirality map has the total energy, momentum and charge of the empty state, so the conservation of the totals does not forbid its creation; but the charge of each member is conserved on its own, no cancellation between two independently quantized universes follows, and no creation process, rate or amplitude is derived, so the statement remains the notebook's hypothesis. Does the theory explain the excess of matter over antimatter? It does not: its charge is exactly conserved, and each field has an exact symmetry that reverses the charge (C for dirac16complex00 in every gravitational field; CP for dirac16complex, which has no exact C at nonzero mass, in flat space and wherever the reflection is an isometry), so the first two of Sakharov's conditions fail; no departure from equilibrium is computed, and the theory contains no baryons of the Standard Model. Every statement carries one of five labels (proved, computed, assumed, hypothesis, open) in an honesty ledger, every number is copied from a committed report of the repository, and every proof is paired with the exact machine check that verifies it where one exists. The book records the state of the project on 2026-09-30, when the Stage-4 gate had not been run and the numerics, the documents and the gate of Stage 5 were not finished; what remains open is collected as a list of problems for students.

\end{abstract}

\section{How to read this book}

This chapter is the map. It says what the book is about, which request it answers, how each statement is labelled so that you always know whether it is proved, computed, assumed, a hypothesis or an open question, which conventions hold on every page, and how you can check any statement yourself. It ends with the honesty ledger: a table of the main statements of the project with their status and the place where each one is verified. The ledger records the state of the repository on 2026-09-30. Stages 1 to 3 of the project are finished; Stages 4 and 5 are not, and every row that depends on them says which part is finished and which is not (Section 0.5).

\subsection{What this book is about}

Physics describes the world with \textbf{fields}. A field is a rule that attaches numbers to every point of space and time: the temperature in a room attaches one number to each point, the wind attaches three (its components along three directions). The field of this book attaches sixteen \textbf{complex numbers} to each point (complex numbers are introduced in Chapter 1), and the space it lives on has eight directions instead of the four of everyday space and time. Four of the eight directions behave like space and four behave like time. One of the four time-like directions, called $x_4$, is chosen as the time in which everything evolves; the other three are called the extra times.

The field is called \textbf{dirac16complex}. Its sixteen components are not ordinary numbers but anticommuting quantities, called Grassmann numbers (Chapter 5), because it is meant to describe fermions, particles of the kind of electrons and quarks. A second version, \textbf{dirac16complex00}, has the same sixteen components but they are ordinary commuting complex numbers; it is the analogue of the four-component wave function of Dirac's 1928 equation for the electron, taken as a classical field (Chapter 6).

The project that this book explains started from a Mathematica notebook by Patrick L. Nash, the file \texttt{Pair\_Cre\allowbreak{}ation\_of\allowbreak{}\_Univers\allowbreak{}es\_WaveF\allowbreak{}unctionO\allowbreak{}fUnivers\allowbreak{}e-4+4-Ei\allowbreak{}nstein-L\allowbreak{}ovelock-\allowbreak{}Nash.nb} in the repository root, which we call \textbf{the notebook}. The notebook proposes a gravitational field in eight dimensions in which ordinary 3-space inflates while the three extra times deflate, and it states the hypothesis that at time $x_4=0$ universes of masses $+M$ and $-M$ are created in pairs. The project wrote down the field theory of dirac16complex in any gravitational field, then in the notebook's field. Many of its exact statements were checked by computer algebra, most of them in two independent programs (Wolfram Language and Python) and some, among them the reading of the notebook's stored output, in Wolfram Language alone; a few rest on a derivation that no program checks. The ledger of Section 0.9 names the checks and marks the statements that rest on a derivation only. The project found and corrected errors in the notebook, computed numerical solutions, built a density-functional (Kohn--Sham) approximation for many particles, and formulated exact pairing theorems that relate solutions of mass $+m$ to solutions of mass $-m$ (read, in the notebook's picture, as universes of masses $+M$ and $-M$). The first of these results, the chirality map (ledger row L16), was proved in Stage 1 of the project; the complete set, the theorems T1 to T3, was proved in Stage 5, whose exact theory is complete while its numerical part and its documents are not (Section 0.5). A separate exact analysis examined what the theory says about matter and antimatter. The ledger of Section 0.9 records the status of each result.

You need school algebra and the calculus of one variable (derivatives and integrals of functions of one variable). Everything else is built from zero in Chapter 1 and the chapters after it.

\subsection{The request and what can honestly be delivered}

The book was written in answer to this request of the author (2026-09-30): "create a new teaching 'textbook' (.md, .tex, and .pdf files) for ignorant students that clearly and correctly describes, in entirety, each and every step in setting up the formalism, deriving the governing field equations, introducing and deriving 'DFT' type approximations, solving all equations (to the extent possible, given current knowledge), proving that the 'big bang' creates universes in pairs, this theory solves matter anti-matter mysteries."

Four parts of the request are delivered as far as this project has carried them; Chapter 18 lists what remains open and why. The formalism is set up in Chapters 1 to 5 and 8. The governing field equations are derived in Chapters 6 and 7. The density-functional (DFT) approximations are introduced from zero in Chapter 12 and derived for the two fields in Chapters 13 and 14. The equations are solved exactly where that is possible and numerically elsewhere, in Chapters 9 to 11 and 13 to 15 (for dirac16complex00 the Kohn--Sham states were computed in one configuration only, Chapter 14), and the places where they are not solved are listed as open problems in Chapter 18. Some of these are limits of present knowledge; others, such as the Einstein--Lovelock terms of orders 2 and 3, which are computable but were not computed (row L25), are limits of this project.

Two parts of the request are not established by this project, and a textbook that wrote them as established would teach something false. This book therefore follows one rule above every other, the \textbf{honesty rule}: it teaches exactly what the repository proves and computes, together with every assumption, and it never writes "proved" for a statement that is not proved.

The first of the two parts is "the big bang creates universes in pairs". The notebook states it as a hypothesis, and its own words make that clear. Cell 6 says: "HYPOTHESIS: If, employing the Einstein eqs (or Einstein-Lovelock eqs), superluminal inflation/deflation exists, then at time x4 = 0 ... a pair of universes with MASSES ± M is created". Cell 7 asks: "Are Universe (s) of masses ± M created in pairs at time x4 = 0 (before the particles of the standard model exist) ?". Cell 17 records the task: "TODO: prove Universe(s) of masses ±M are created in pairs!" (the three cells are quoted from the survey \texttt{handoff/\allowbreak{}surveys/\allowbreak{}survey\_n\allowbreak{}otebook-\allowbreak{}physics.\allowbreak{}md}). What the project proves are exact \textbf{pairing theorems} (Chapter 15): a map of the field (the chirality map $\gamma^8$ of Chapter 2, proved in Stage 1) that turns every solution with mass $+m$ and self-coupling $\lambda$ into a solution with mass $-m$ and self-coupling $-\lambda$ (the self-coupling $\lambda$ is the strength with which the field interacts with itself, Section 0.6 and Chapter 6); a mirror map across the brane of the primordial field (the surface $y=0$ of the hidden coordinate $y$, along which two mirror copies of the gravitational field are glued, Section 9.16), which relates mass $+m$ to mass $-m$ with the same self-coupling; and the corresponding statement for the Kohn--Sham states. The last two, and the consequence stated next, are the work of Stage 5: their exact proofs are complete and machine-checked (rows L37 to L39), while the numerical demonstration of the Kohn--Sham statement is partial (row L39). Their consequences are strong but limited. As classical fields, a pair related by $\gamma^8$ (the second member with mass $-m$ and self-coupling $-\lambda$) carries zero total energy, momentum and charge in any gravitational field (Theorem T1 of Chapter 15). The creation of such a pair from nothing is therefore not forbidden by the conservation of the totals or by the gravitational constraint equations (the parts of the field equations of gravity that restrict the state at a single instant, Chapter 16): with the pair as its only source, the Einstein (or Einstein--Lovelock) equations become the source-free equations. In the theory as built, however, the two members do not interact, so the charge of each member is conserved on its own, and a pair whose members carry charge cannot appear from the empty state (Proposition 16.5 of Chapter 16). Whether a given gravitational field can have such a pair as its only source is a separate question, since the field must then be a source-free solution; the notebook's primordial field, for example, is not a source-free solution of eight-dimensional Einstein gravity (row L21), and for Einstein--Lovelock gravity the question is open because the Lovelock terms were not computed (row L25). At the quantum level the vanishing of the totals is an identity between the field $\Psi$ and its own image $\gamma^8\Psi$: the image is the same quantum system (the same operators and the same state), described with the Krein form of Chapter 8 taken with the opposite sign, and the totals vanish as identities of the form $X+(-X)=0$. It is not a cancellation between two independent universes: a universe of mass $-M$ that is quantized on its own has positive energies, which add to ours, and its charge cancels ours only under an additional assumption about its state (\texttt{provenan\allowbreak{}ce/DIRAC\allowbreak{}16COMPLE\allowbreak{}X\_MATTER\allowbreak{}\_ANTIMAT\allowbreak{}TER.md}, §7.3, and \texttt{artifact\allowbreak{}s/dirac1\allowbreak{}6complex\allowbreak{}/matter-\allowbreak{}antimatt\allowbreak{}er/matte\allowbreak{}r-antima\allowbreak{}tter-the\allowbreak{}ory.json}, key \texttt{M5.krein\allowbreak{}LevelCav\allowbreak{}eat}; Chapters 15 and 17). And no creation process, no rate and no probability amplitude is derived, and no dynamical big bang is computed. Stage 1 of the project said it in one sentence, which this book repeats: the pairing of the masses $\pm m$ is "a structural property of the equations, not a claim that universes of masses $\pm M$ are created in pairs" (Stage-1 document \texttt{provenan\allowbreak{}ce/DIRAC\allowbreak{}16COMPLE\allowbreak{}X\_ARBITR\allowbreak{}ARY\_FIEL\allowbreak{}D.md}, §1.3, item 6). Chapter 16 states exactly what is proved, what is merely consistent and what is not derived; the missing step is an open problem.

The second part is "this theory solves matter anti-matter mysteries". The observed universe contains matter but almost no antimatter, and explaining this requires, according to the conditions found by Sakharov in 1967, interactions that change the number of baryons, a violation of the symmetries called C and CP, and a period out of thermal equilibrium (Chapter 17 explains all of this from zero). The theory as built fails the first two conditions exactly and does not address the third. For both fields, every self-interaction and every gravitational field, its Lagrangian does not change when the field is multiplied by a constant phase, $\Psi\to e^{i\alpha}\Psi$ (an exact U(1) symmetry), so its charge $Q$ is exactly conserved (when no charge flows through the boundary of space) and no process of the theory can create a net charge inside one universe: Sakharov's first condition fails (Theorem M1 of the matter--antimatter analysis, row L44). Each of the two fields also has an exact symmetry that reverses $Q$ without reversing the time $x_4$. dirac16complex00 has an exact charge conjugation C, $\Psi\to\Psi^\ast$, in every gravitational field. dirac16complex has no exact C when $m\ne0$, but an exact CP, a charge conjugation combined with the reflection of an odd number of the space-like directions $x_0,\dots,x_3$, in flat space and in every gravitational field in which that reflection is an isometry (a generic gravitational field breaks this CP, but since $Q$ stays conserved this opens no way to an asymmetry). So Sakharov's second condition fails as well (Theorem M2, row L45). No departure from equilibrium is computed, and the theory contains no baryons of the Standard Model of particle physics (\texttt{provenan\allowbreak{}ce/DIRAC\allowbreak{}16COMPLE\allowbreak{}X\_MATTER\allowbreak{}\_ANTIMAT\allowbreak{}TER.md}, §1 and §9). It therefore does not solve the matter--antimatter problem. What it does contribute is an exact symmetry: the chirality map $\gamma^8$ turns every solution with mass $+m$ and self-coupling $\lambda$ into a solution with mass $-m$, self-coupling $-\lambda$ and the opposite charge (in the notebook's picture, a universe of mass $+M$ into one of mass $-M$); this holds for the classical fields and, as identities between operators, for the quantized field. If a second classical field is in the configuration $\gamma^8\Psi$ of the first (an assumption, not derived), the two carry zero total charge; for two independently quantized universes there is no such cancellation (Chapter 17). This places the theory in the family of "universe and anti-universe" ideas; Chapter 17 describes that family, lists what would be needed to turn such an idea into an explanation, and labels every scenario as a hypothesis.

\subsection{Five words for the status of a statement}

Every statement of this book that matters has one of five labels. They are defined here once, and the ledger of Section 0.9 uses them.

\textbf{PROVED} means: the statement follows from its stated definitions and assumptions by a complete derivation, and this book gives that derivation step by step. Where the repository contains an exact machine check (a computation with exact integers, fractions or symbols, in Wolfram Language and, for most checks, independently in Python), the check is named: the report file and the name of the check inside it. One distinction matters. When a statement is about a finite list of definite objects, for example "these eight matrices of integers satisfy these 64 equations", an exact computation of all cases is itself a complete proof. When a statement is about infinitely many objects, for example "for every gravitational field", the proof is the general derivation; an exact check at a few test points then confirms the derivation (and would have caught many errors), but it is not the proof. A third case lies between the two: when the computer keeps the unknown quantities as symbols (a free function $a_4$, a mass $m$, a temperature $T$), the exact symbolic computation holds for all their values and is itself a proof, like a derivation done by hand. The ledger says which of the three applies.

\textbf{COMPUTED} means: the statement is a number, a table or a curve produced by a numerical program, with floating-point arithmetic, finite grids and error tolerances. Computed results are reproducible from the repository and are compared with independent programs, but they are not proofs; each carries an accuracy, and the book states it.

\textbf{ASSUMED} means: the statement is a starting point that is not derived. There are three kinds: conventions (counting from 0, the sign pattern of the metric, the sign of the energy--momentum tensor), physical inputs (for example, that the gravitational field of the notebook is given and not computed from its source), and approximations (for example, that the sixteen field components can be replaced by average numbers, the mean-field approximation). A proved statement is often proved under assumptions; the ledger names the important ones.

\textbf{HYPOTHESIS} means: a claim, physical or about how a result was produced, that someone has proposed and that the book states, but that is neither proved nor computed here. The notebook's pair-creation hypothesis is the main example; row L20 of the ledger is a hypothesis about how the notebook produced one of its outputs.

\textbf{OPEN} means: a question that this project has not answered. Chapter 18 collects the open questions and suggests how a student could attack them.

A small example shows how the labels work. The statement "the matrix $A=\begin{pmatrix}0&1\\1&0\end{pmatrix}$ satisfies $A^2=I$" is PROVED by multiplying out (Chapter 1 shows how). The statement "in the numerical experiment EXP-4 a gas of dirac16complex quanta ends with the equation-of-state parameter $w=0.0359$" is COMPUTED ($w$ is the pressure divided by the energy density, Section 5.8; the number is the key \texttt{thermal.\allowbreak{}wAtAEnd} of \texttt{artifact\allowbreak{}s/dirac1\allowbreak{}6complex\allowbreak{}/numeric\allowbreak{}s/exp4/s\allowbreak{}ummary.j\allowbreak{}son}, rounded, and it holds on the truncated momentum grid of the experiment, Chapter 11). The statement "we count from 0" is ASSUMED. The statement "the big bang creates universes in pairs" is a HYPOTHESIS. The question "what process creates such a pair, and with what probability?" is OPEN.

\subsection{The map of the book}

The book has five parts and two appendices. Each chapter is written to use only what earlier chapters have built; where a chapter needs a later result, it says so.

Part I, Foundations.

\begin{itemize}
\item Chapter 0 (this chapter): how to read the book, and the honesty ledger.
\item Chapter 1: the mathematical toolkit from zero: vectors, matrices, complex numbers, indices and the summation convention, derivatives in several variables, metrics and signatures, the (4,4) world of this book.
\item Chapter 2: Clifford algebras and spinors from zero, from the Pauli matrices to the sixteen-dimensional representation of Pin(4,4), its chirality and its two eight-dimensional halves.
\item Chapter 3: split octonions and the notebook's construction of the gamma matrices, and how it is related to the other standard construction.
\item Chapter 4: curved space from zero: metric, Christoffel symbols, curvature, the Einstein and Einstein--Lovelock equations, the vielbein, the spin connection, and why the notebook's spin-connection contraction is wrong.
\end{itemize}

Part II, Field theory.

\begin{itemize}
\item Chapter 5: classical field theory from zero: action, Euler--Lagrange equations, symmetries and conserved currents, the energy--momentum tensor, Grassmann numbers, and why the notebook's Lagrangian gives no wave equation for a real Grassmann field.
\item Chapter 6: the two fields dirac16complex and dirac16complex00 and their Lagrangians with explicit mass terms and coupling to gravity.
\item Chapter 7: field equations, energy--momentum tensor, kinetic and potential energy, pressure, energy density and equations of state in an arbitrary gravitational field.
\item Chapter 8: canonical quantization in 4+4 dimensions: the Hamiltonian, anticommutators, the Krein space, the good sector and the unstable extra-time modes.
\end{itemize}

Part III, The primordial universe.

\begin{itemize}
\item Chapter 9: the primordial gravitational field of the notebook, what it requires as a source, the notebook errata, the static warped form and the mirror brane.
\item Chapter 10: solving differential equations on a computer from zero.
\item Chapter 11: the dark-sector experiments EXP-1 to EXP-5 and what they show.
\end{itemize}

Part IV, Density functional theory.

\begin{itemize}
\item Chapter 12: many-body quantum mechanics and density functional theory from zero: Hartree, Hartree--Fock, Hohenberg--Kohn, Kohn--Sham, the local density approximation, finite temperature, and excited states.
\item Chapter 13: the Kohn--Sham approximation for dirac16complex in the primordial field, with its ground and first excited states.
\item Chapter 14: the Kohn--Sham approximation for dirac16complex00.
\end{itemize}

Part V, Pairs of universes, matter and antimatter.

\begin{itemize}
\item Chapter 15: the pairing theorems T1 to T3 with complete proofs, and their Kohn--Sham demonstration.
\item Chapter 16: does the big bang create universes in pairs? What is proved, what is consistent, what is not derived.
\item Chapter 17: matter and antimatter: the observed asymmetry, the Sakharov conditions, universe and anti-universe ideas, and what this theory does and does not contribute.
\item Chapter 18: open problems and how a student could attack them.
\end{itemize}

Appendices.

\begin{itemize}
\item Chapter 19: reproducing everything: the commands, the expected outputs and the verification gates.
\item Chapter 20: glossary and index of the verifier checks.
\end{itemize}

If you are in a hurry, read Chapters 0, 1, 2, 4, 5, 6 and 7 for the field theory (Chapter 3 can be skipped on a first reading), add Chapters 8, 9, 12 and 13 for the Kohn--Sham approximation, and read Chapters 15 to 17 for the questions about pairs of universes and antimatter. If you want to understand every step, read the chapters in order and do the exercises.

\subsection{The stages of the project}

The project was carried out in five stages, and a separate exact analysis of matter and antimatter was added to them. A finished stage has a document in the folder \texttt{provenan\allowbreak{}ce/} (a Markdown source, a LaTeX file generated from it, and a PDF), machine-readable reports under \texttt{artifact\allowbreak{}s/dirac1\allowbreak{}6complex\allowbreak{}/}, and a \textbf{gate}, a script that reruns every program of the stage and compares the results with the committed reports (Section 0.8). The book is built on these documents and reports, and it quotes numbers only from the reports.

\begingroup
\small
\begin{longtable}{@{}>{\raggedright\arraybackslash}p{0.200\linewidth}>{\raggedright\arraybackslash}p{0.200\linewidth}>{\raggedright\arraybackslash}p{0.200\linewidth}>{\raggedright\arraybackslash}p{0.200\linewidth}@{}}
\toprule
\textbf{Stage} & \textbf{Question} & \textbf{Document} & \textbf{State on 2026-09-30} \\
\midrule
\endfirsthead
\toprule
\textbf{Stage} & \textbf{Question} & \textbf{Document} & \textbf{State on 2026-09-30} \\
\midrule
\endhead
1 & the field in an arbitrary gravitational field & \texttt{provenan\allowbreak{}ce/DIRAC\allowbreak{}16COMPLE\allowbreak{}X\_ARBITR\allowbreak{}ARY\_FIEL\allowbreak{}D.md} & complete: 153 of 153 exact checks; the gate verified from a fresh public clone \\
2 & the field in the primordial field of the notebook & \texttt{provenan\allowbreak{}ce/DIRAC\allowbreak{}16COMPLE\allowbreak{}X\_PRIMOR\allowbreak{}DIAL\_FIE\allowbreak{}LD.md} & complete: 126 of 126 Wolfram and 16 of 16 Python checks; the gate verified from a fresh public clone \\
3 & pressure, energy density, dark matter and dark energy from numerical solutions & \texttt{provenan\allowbreak{}ce/DIRAC\allowbreak{}16COMPLE\allowbreak{}X\_DARK\_S\allowbreak{}ECTOR\_NU\allowbreak{}MERICS.m\allowbreak{}d} and \texttt{provenan\allowbreak{}ce/DIRAC\allowbreak{}16COMPLE\allowbreak{}X\_STUDEN\allowbreak{}T\_GUIDE.\allowbreak{}md} & complete; the gate verified from a fresh public clone \\
4 & Kohn--Sham ground and first excited states in the primordial field & not finished: only a Markdown draft, \texttt{provenan\allowbreak{}ce/DIRAC\allowbreak{}16COMPLE\allowbreak{}X\_KOHN\_S\allowbreak{}HAM\_PRIM\allowbreak{}ORDIAL.m\allowbreak{}d} & exact theory complete (125 of 125 Wolfram and 157 of 157 Python checks); the Rust solver complete (397 of 397 of its own checks); 56 runs of the independent reference solver; the final cross-check of the two solvers has 63 checks, of which 1 failed, \texttt{canonica\allowbreak{}l\_eigenv\allowbreak{}alues} (row L34); the Stage-4 gate not yet run \\
5 & dirac16complex00 and the pairing of universes of masses $+M$ and $-M$ & not written & exact theory complete (46 of 46 and 141 of 141 Wolfram, 49 of 49 and 172 of 172 Python checks); numerics partial; no Stage-5 gate yet \\
--- & matter and antimatter in the theory & \texttt{provenan\allowbreak{}ce/DIRAC\allowbreak{}16COMPLE\allowbreak{}X\_MATTER\allowbreak{}\_ANTIMAT\allowbreak{}TER.md} & complete: 44 of 44 Wolfram and 77 of 77 Python checks; its quantum-level statements taken from Stage 5 are cited as provisional \\
\bottomrule
\end{longtable}
\endgroup

The check counts are those of the following reports, in the order of the table: Stage 1 in the first line, Stage 2 in the next two, the exact theory of Stage 4 in the next two, the Stage-4 cross-check in the next one, Stage 5 in the next four, and the matter--antimatter analysis in the last two. The 397 checks of the Stage-4 Rust solver are those of the five files \texttt{summary.\allowbreak{}json} in the subfolders \texttt{\detokenize{spectrum}}, \texttt{\detokenize{scf}}, \texttt{\detokenize{excited}}, \texttt{\detokenize{thermo}} and \texttt{\detokenize{emt}} of \texttt{artifact\allowbreak{}s/dirac1\allowbreak{}6complex\allowbreak{}/kohn-sh\allowbreak{}am/rust/} (33, 137, 65, 102 and 60 checks).

\begin{Verbatim}[fontsize=\small]
artifacts/dirac16complex/arbitrary-field/stage1-summary.json
artifacts/dirac16complex/primordial-field/wolfram-primordial-report.json
artifacts/dirac16complex/primordial-field/python-primordial-report.json
artifacts/dirac16complex/kohn-sham/wolfram-kohn-sham-report.json
artifacts/dirac16complex/kohn-sham/python-theory-report.json
artifacts/dirac16complex/kohn-sham/python-check-report.json
artifacts/dirac16complex/pair-creation/wolfram-dirac16complex00-report.json
artifacts/dirac16complex/pair-creation/wolfram-pairing-report.json
artifacts/dirac16complex/pair-creation/python-dirac16complex00-report.json
artifacts/dirac16complex/pair-creation/python-pairing-report.json
artifacts/dirac16complex/matter-antimatter/wolfram-matter-antimatter-report.json
artifacts/dirac16complex/matter-antimatter/python-matter-antimatter-report.json
\end{Verbatim}

The state of each stage is recorded in \texttt{README.m\allowbreak{}d} and, in more detail, in \texttt{HANDOFF.\allowbreak{}md} section 2. Where the book uses results of an unfinished stage, it says so at the place where it uses them.

\subsection{Conventions used everywhere}

These conventions hold on every page of the book. They are choices, not results (status ASSUMED), and they are the conventions of every document of the project.

\begin{itemize}
\item \textbf{Counting from 0.} Coordinates are $x=(x_0,x_1,\dots,x_7)$, the components of the field are $\Psi_0,\dots,\Psi_{15}$, frame indices run over $0,\dots,7$, and the first entry of every list is entry 0. The exceptions are the numbers of the sections, exercises and ledger rows of this book (which start at 1 within each chapter or table), the numbering of the sections and results of the project documents, the cell numbers of the notebook, and the lists of Mathematica (which start at 1); where such a 1-based number could be confused with a 0-based one, the text says so.
\item \textbf{The roles of the coordinates.} $x_0$ is a hidden space direction, $x_1,x_2,x_3$ are ordinary 3-space, $x_4$ is the time in which everything evolves, and $x_5,x_6,x_7$ are the three extra times.
\item \textbf{The flat metric.} $\eta=\mathrm{diag}(+1,+1,+1,+1,-1,-1,-1,-1)$: the directions 0 to 3 are space-like and the directions 4 to 7 are time-like (Chapter 1 explains what this means). The signature is (4,4).
\item \textbf{Index letters.} Greek letters $\mu,\nu,\rho,\sigma,\lambda$ run over the coordinates $0,\dots,7$ (curved indices, Section 1.8 and Chapter 4), Latin letters $a,b,c$ over the frame directions $0,\dots,7$ (the directions of the vielbein, Section 4.10), and spinor indices over $0,\dots,15$. In the primordial field, $i$ runs over $1,2,3$ and $j$ over $5,6,7$; the Kohn--Sham chapters are the exception (Chapter 13, Section 13.1, and the passages of other chapters that use its blocks, for example in Chapters 14 and 15): there $j=\pm1$ is the type of a $2\times2$ block (Section 13.4), and $l$ runs over the six warped directions $1,2,3,5,6,7$, with $l\le3$ the 3-space directions. Three of these letters also have fixed meanings as symbols, not indices: $\lambda$ is the strength of the self-interaction $U(S)=\tfrac\lambda2S^2$ of the field, where $S=\bar\Psi\Psi$ is a product of the field with itself (Chapter 6; in Chapters 1, 4 and 8 it also denotes an eigenvalue or a curve parameter, as the text says there), $\rho$ is an energy density, and $\sigma$ with a label, such as $\sigma_{16}$, is a fixed matrix (Chapters 2 and 3). A letter written as a subscript or superscript of another symbol is an index; otherwise it is the fixed symbol.
\item \textbf{Units.} Planck's constant and the speed of light are set to 1 ($\hbar=c=1$). Numerical work sets one more scale to 1 (for example $H=1$, where $H$ is the notebook's single inverse length, Section 9.2, or $m=1$) and says which.
\item \textbf{Names.} "The notebook" is the author's Mathematica notebook named in Section 0.1. "The Stage-N document" is the document of Stage N listed in Section 0.5. A file is always named by its path relative to the repository root, written in typewriter type, for example \texttt{README.m\allowbreak{}d}.
\end{itemize}

\subsection{How each chapter is built and how to study it}

Every chapter has the same parts. It begins with the motivation in plain words: what problem the chapter solves and why it is needed. Then come the definitions, each introduced before it is used. Then the derivations, step by step, with no step left to "it can be shown"; where a standard mathematical fact is used without proof, the text says so and the "What we proved and what we assumed" section of the chapter lists it. Then worked examples with small matrices and small numbers, which you should redo with pencil and paper. The chapter ends with the "What we proved and what we assumed" section, exercises, and complete answers to the exercises.

References inside the book are written "Chapter 1" or "Section 0.3". References to files of the repository are written in typewriter type. A reference to a section of one of the project documents names the document and writes the section with the sign §, as in "Stage-1 document, §7.7", so that it cannot be mistaken for a section of this book.

Three pieces of advice. First, do the exercises before you read the answers; mathematics is learned by doing it. Second, when a symbol is unfamiliar, look it up in the glossary of Chapter 20 or in the chapter where it was defined. Third, when a statement surprises you, find it in the ledger of Section 0.9 and check it with the method of Section 0.8.

\subsection{How to check a statement yourself}

You can check a statement at three levels.

The first level is the derivation in the book. Follow it line by line. If a line does not follow from the lines before it, you have found either a gap in your understanding or an error in the book; both are worth finding.

The second level is the report of the machine check. Every exact check writes its result into a JSON file under \texttt{artifact\allowbreak{}s/dirac1\allowbreak{}6complex\allowbreak{}/}. JSON is a plain-text format for lists and tables that Python reads directly. Each report has an entry \texttt{\detokenize{checks}}, which maps the name of every check to \texttt{\detokenize{true}} or \texttt{\detokenize{false}}. For example, to count the checks of Stage 1, start Python in the repository root (type \texttt{\detokenize{python}} in PowerShell or in Git Bash, \texttt{\detokenize{python3}} in the Terminal of macOS or Linux; if Git Bash shows no prompt, type \texttt{winpty p\allowbreak{}ython} instead). Python shows its prompt, three greater-than signs, at the start of each input line; type the text that follows the prompt in the lines below:

\begin{Verbatim}[fontsize=\small]
>>> import json
>>> path = "artifacts/dirac16complex/arbitrary-field/stage1-summary.json"
>>> checks = json.load(open(path, encoding="utf-8"))["checks"]
>>> sum(checks.values()), len(checks)
(153, 153)
\end{Verbatim}

The last line is Python's answer: 153 checks are true out of 153. In Python, \texttt{\detokenize{True}} counts as 1 and \texttt{\detokenize{False}} as 0 when added, so \texttt{sum(chec\allowbreak{}ks.value\allowbreak{}s())} is the number of true checks. To look at one named check, for example the check \texttt{ALG\_gamm\allowbreak{}a8Map} of the chirality map in the Wolfram algebra report:

\begin{Verbatim}[fontsize=\small]
>>> path = "artifacts/dirac16complex/arbitrary-field/wolfram-algebra-report.json"
>>> json.load(open(path, encoding="utf-8"))["checks"]["ALG_gamma8Map"]
True
\end{Verbatim}

Type \texttt{\detokenize{exit()}} to leave Python. Each report of an exact verifier (for example \texttt{wolfram-\allowbreak{}algebra-\allowbreak{}report.j\allowbreak{}son}) also has an entry \texttt{measurem\allowbreak{}ents} with the numbers that the checks measured, and an entry \texttt{\detokenize{producer}} that names the program that wrote it. The summary \texttt{stage1-s\allowbreak{}ummary.j\allowbreak{}son} has no \texttt{measurem\allowbreak{}ents}; it re-reads the five Stage-1 reports and records the measurements compared between the two programs under \texttt{agreedMe\allowbreak{}asuremen\allowbreak{}ts}.

The third level is to run the programs again. Each finished stage has a gate, a script that reruns every verifier of the stage and compares the results with the committed reports; it ends with a line such as \texttt{stage2\_p\allowbreak{}rimordia\allowbreak{}l\_field\_\allowbreak{}verifica\allowbreak{}tion=OK}. (The gate of Stage 4 exists but has not been run; Stage 5 has no gate yet, Section 0.5.) The gates need Python, WolframScript, a TeX distribution for the PDFs and, for the numerical stages, the Rust toolchain; some run for minutes, some for hours. Chapter 19 lists all commands for PowerShell and for Bash with their expected outputs.

\subsection{The honesty ledger}

The ledger lists the main statements of the project, each with its status (Section 0.3) and the place where it is verified. A report named without a folder lies in the folder named at the start of its group. "Derivation only" means that the statement is proved by a derivation that is not itself a machine check. A check name ending in \texttt{\detokenize{_*}} stands for the whole family of checks with that beginning. Every symbol in the ledger is defined in the chapter named in the last column; on a first reading you may skip the formulas and read only the words and the status. The ledger describes the state of 2026-09-30 (Section 0.5): the rows on Stages 1 to 3 are final; the rows on Stages 4 and 5 state which parts are finished, and a result of an unfinished part carries that limitation in its row.

\textbf{Algebra and representations} (reports in \texttt{artifact\allowbreak{}s/dirac1\allowbreak{}6complex\allowbreak{}/arbitra\allowbreak{}ry-field\allowbreak{}/}).

\begingroup
\small
\begin{longtable}{@{}>{\raggedright\arraybackslash}p{0.267\linewidth}>{\raggedright\arraybackslash}p{0.267\linewidth}>{\raggedright\arraybackslash}p{0.267\linewidth}@{}}
\toprule
\textbf{Statement} & \textbf{Status} & \textbf{Where verified} \\
\midrule
\endfirsthead
\toprule
\textbf{Statement} & \textbf{Status} & \textbf{Where verified} \\
\midrule
\endhead
L1. The eight $16\times16$ integer matrices $\gamma^0,\dots,\gamma^7$ of the notebook satisfy $\{\gamma^a,\gamma^b\}=2\eta^{ab}I_{16}$. & PROVED (finite exact computation) & \texttt{ALG\_clif\allowbreak{}ford} in \texttt{wolfram-\allowbreak{}algebra-\allowbreak{}report.j\allowbreak{}son} and \texttt{python-a\allowbreak{}lgebra-r\allowbreak{}eport.js\allowbreak{}on}; Chapter 2 \\
L2. $C=\sigma_{16}=\gamma^0\gamma^1\gamma^2\gamma^3$ is symmetric with $C^2=I_{16}$, and every $C\gamma^a$ is antisymmetric (expression [1] of Section 2.9). & PROVED (finite exact computation) & \texttt{ALG\_char\allowbreak{}geMatrix}, \texttt{ALG\_expr\allowbreak{}ession1} in both algebra reports; Chapter 2 \\
L3. The chirality $\gamma^8=\gamma^0\gamma^1\cdots\gamma^7=\mathrm{diag}(-I_8,I_8)$ anticommutes with every $\gamma^a$. & PROVED (finite exact computation) & \texttt{ALG\_chir\allowbreak{}ality} in both algebra reports; Chapter 2 \\
L4. $\mathbb C^{16}$ is irreducible under Pin(4,4); under Spin(4,4) it is the sum of two inequivalent irreducible 8-dimensional parts. & PROVED (exact commutant dimensions 1 and 2, cross-intertwiner dimension 0) & \texttt{ALG\_pinI\allowbreak{}rreducib\allowbreak{}leComple\allowbreak{}x}, \texttt{ALG\_spin\allowbreak{}Decompos\allowbreak{}ition} in both algebra reports; Chapter 2 \\
L5. The notebook's split-octonion gammas and the tensor-product (Clifford) gammas describe one module; the intertwiners K\_clifford and K\_octonion are unique up to a factor. & PROVED; the comparison with the published dirac-main files needs the private folder \texttt{dirac-ma\allowbreak{}in/} & \texttt{ALG\_clif\allowbreak{}fordPict\allowbreak{}ureInter\allowbreak{}twiner}, \texttt{ALG\_octo\allowbreak{}nionPict\allowbreak{}ureInter\allowbreak{}twiner} in both algebra reports; Chapter 3 \\
\bottomrule
\end{longtable}
\endgroup

\textbf{Curved space and field theory in an arbitrary gravitational field} (reports in \texttt{artifact\allowbreak{}s/dirac1\allowbreak{}6complex\allowbreak{}/arbitra\allowbreak{}ry-field\allowbreak{}/}; G1, G2 and G3 are the exact test geometries of the Stage-1 document).

\begingroup
\small
\begin{longtable}{@{}>{\raggedright\arraybackslash}p{0.267\linewidth}>{\raggedright\arraybackslash}p{0.267\linewidth}>{\raggedright\arraybackslash}p{0.267\linewidth}@{}}
\toprule
\textbf{Statement} & \textbf{Status} & \textbf{Where verified} \\
\midrule
\endfirsthead
\toprule
\textbf{Statement} & \textbf{Status} & \textbf{Where verified} \\
\midrule
\endhead
L6. With the canonical spin connection $\omega_{\mu ab}=\eta_{ac}\,\omega_\mu{}^c{}_b$ the curved gamma matrices are covariantly constant. & PROVED in general by derivation; exact checks at test points & \texttt{GEO\_gamm\allowbreak{}aCovaria\allowbreak{}ntConsta\allowbreak{}ncy\_G1}, \texttt{GEO\_gamm\allowbreak{}aCovaria\allowbreak{}ntConsta\allowbreak{}ncy\_G2} in \texttt{wolfram-\allowbreak{}geometry\allowbreak{}-report.\allowbreak{}json} and \texttt{python-g\allowbreak{}eometry-\allowbreak{}report.j\allowbreak{}son}; Chapter 4 \\
L7. The notebook's contraction of $\omega_\mu{}^a{}_b$ with $S^{ab}$ lacks one metric factor: with it the gammas are not covariantly constant. & PROVED (a counterexample proves a "not") & \texttt{GEO\_note\allowbreak{}bookCont\allowbreak{}ractionF\allowbreak{}ails\_G1}, \texttt{GEO\_note\allowbreak{}bookCont\allowbreak{}ractionF\allowbreak{}ails\_G2} in both geometry reports; Chapter 4 \\
L8. For a real anticommuting field, the notebook's Lagrangian Lg[] with the canonical spin connection is a total derivative, so its field equations are empty ($0=0$). With the notebook's own contraction its field equations are the algebraic equations $X\Psi=0$, which force $\Psi=0$ wherever the $16\times16$ matrix $X$ has rank 16, as it has at all six test points where it was computed. Either way Lg[] gives no wave equation. & PROVED in general by derivation (Stage-1 document, Theorem 6.1); exact checks at test points; the rank 16 of $X$ is computed at the six test points only & \texttt{LAG\_note\allowbreak{}bookLgGr\allowbreak{}assmannT\allowbreak{}rivial\_G\allowbreak{}1}, \texttt{LAG\_note\allowbreak{}bookLgGr\allowbreak{}assmannT\allowbreak{}rivial\_G\allowbreak{}2} in \texttt{wolfram-\allowbreak{}geometry\allowbreak{}-report.\allowbreak{}json}; \texttt{GR\_noteb\allowbreak{}ookLgPur\allowbreak{}eDiverge\allowbreak{}nce} in \texttt{grassman\allowbreak{}n-demo-r\allowbreak{}eport.js\allowbreak{}on}; Chapter 5 \\
L9. The Euler--Lagrange equation of the dirac16complex Lagrangian is $\gamma^\mu D_\mu\Psi=(m+U'(S))\Psi$ with $S=\bar\Psi\Psi$. & PROVED in general by derivation; exact checks at test points & \texttt{LAG\_eule\allowbreak{}rLagrang\allowbreak{}ePsibar\_\allowbreak{}G1}, \texttt{LAG\_eule\allowbreak{}rLagrang\allowbreak{}ePsi\_G1} (and the G2 checks) in both geometry reports; \texttt{GR\_compl\allowbreak{}exQuarti\allowbreak{}cEL} in \texttt{grassman\allowbreak{}n-demo-r\allowbreak{}eport.js\allowbreak{}on}; Chapter 7 \\
L10. The energy--momentum tensor is symmetric, Hermitian and conserved when the field equations hold; when they hold, its trace is $-mS+7SU'-8U$. & PROVED by derivation; conservation checked with exact solutions at points, not as an identity & \texttt{EMT\_symm\allowbreak{}etricHer\allowbreak{}mitian\_G\allowbreak{}1} (and G2) in \texttt{wolfram-\allowbreak{}geometry\allowbreak{}-report.\allowbreak{}json}; \texttt{EMT\_cons\allowbreak{}ervation\allowbreak{}\_G1}, \texttt{EMT\_trac\allowbreak{}e\_G1} (and G2) in both geometry reports; Chapter 7 \\
L11. In homogeneous diagonal backgrounds the energy density is $\rho=mS+U$ and, when the field equations hold, the pressure is $p=SU'-U$ in all seven directions other than $x_4$. & PROVED, with the field products read as average numbers (mean field, ASSUMED) & \texttt{EMT\_homo\allowbreak{}geneousR\allowbreak{}eduction\allowbreak{}\_G3} in \texttt{wolfram-\allowbreak{}geometry\allowbreak{}-report.\allowbreak{}json}; \texttt{EMT\_homo\allowbreak{}geneousR\allowbreak{}eduction} in \texttt{python-g\allowbreak{}eometry-\allowbreak{}report.j\allowbreak{}son}; Chapter 7 \\
L12. Canonical quantization with the time $x_4$ gives a state space with an indefinite inner product (a Krein space), and in signature (4,4) every Spin(4,4)-invariant charge density is indefinite (computed classification of the invariant forms). & PROVED (finite exact computation plus derivation) & \texttt{ALG\_inva\allowbreak{}riantFor\allowbreak{}ms}, \texttt{QNT\_krei\allowbreak{}nSignatu\allowbreak{}re} in both algebra reports; Chapter 8 \\
L13. Modes with momentum along $x_5,x_6,x_7$ grow exponentially when $k_5^2+k_6^2+k_7^2>m^2+k_0^2+\dots+k_3^2$. & PROVED (exact dispersion relation); also COMPUTED in EXP-5 & \texttt{QNT\_flat\allowbreak{}ModeHami\allowbreak{}ltonian} in both algebra reports; \texttt{artifact\allowbreak{}s/dirac1\allowbreak{}6complex\allowbreak{}/numeric\allowbreak{}s/exp5/}; Chapters 8 and 11 \\
L14. In flat space, the reduced theory of fields that do not depend on the extra times $x_5,x_6,x_7$ (a (4+1)-dimensional theory, not a subspace of the full theory) has an ordinary positive Fock space. & PROVED (derivation only) & Stage-1 document, §10.6 and §10.7; Chapter 8 \\
L15. An interacting quantum field theory of dirac16complex (interacting states, renormalization). & OPEN & Chapter 18 \\
L16. The chirality map $\Psi\mapsto\gamma^8\Psi$ gives $\mathcal L_{m,\lambda}[\gamma^8\Psi]=-\mathcal L_{-m,-\lambda}[\Psi]$ and turns every solution with $(m,\lambda)$ into one with $(-m,-\lambda)$. & PROVED by derivation from finite matrix facts checked exactly & \texttt{ALG\_gamm\allowbreak{}a8Map} in both algebra reports; Chapters 6 and 15 \\
\bottomrule
\end{longtable}
\endgroup

\textbf{The primordial field of the notebook} (Stage-2 reports in \texttt{artifact\allowbreak{}s/dirac1\allowbreak{}6complex\allowbreak{}/primord\allowbreak{}ial-fiel\allowbreak{}d/}, Stage-4 geometry in \texttt{artifact\allowbreak{}s/dirac1\allowbreak{}6complex\allowbreak{}/kohn-sh\allowbreak{}am/}).

\begingroup
\small
\begin{longtable}{@{}>{\raggedright\arraybackslash}p{0.267\linewidth}>{\raggedright\arraybackslash}p{0.267\linewidth}>{\raggedright\arraybackslash}p{0.267\linewidth}@{}}
\toprule
\textbf{Statement} & \textbf{Status} & \textbf{Where verified} \\
\midrule
\endfirsthead
\toprule
\textbf{Statement} & \textbf{Status} & \textbf{Where verified} \\
\midrule
\endhead
L17. For the notebook's metric $\det g=+\cos^2z$, and the spin connection enters the Dirac operator only through $\gamma^\mu\Omega_\mu=3H\gamma^0$, which does not depend on the free function $a_4$. & PROVED (exact symbolic computation) & \texttt{P\_metric\allowbreak{}\_detG\_eq\allowbreak{}uals\_plu\allowbreak{}s\_cos2z}, \texttt{P\_Omega\_\allowbreak{}gammaSla\allowbreak{}sh3Hgamm\allowbreak{}a0} in \texttt{wolfram-\allowbreak{}primordi\allowbreak{}al-repor\allowbreak{}t.json}; the group checks \texttt{\detokenize{P_metric}}, \texttt{\detokenize{P_Omega}} in \texttt{python-p\allowbreak{}rimordia\allowbreak{}l-report\allowbreak{}.json}; Chapter 9 \\
L18. The substitution rule of notebook cell 1058 is wrong: $1/\sqrt{\sin^{1/3}z/e^{2a_4}}$ equals $e^{a_4}/\sin^{1/6}z$, not $1/(e^{a_4}\sin^{1/6}z)$. & PROVED (school algebra, derivation only; Exercise 1.3 of Chapter 1) & Chapter 1; Chapter 9 \\
L19. The notebook's stored cell-1137 equations agree with the correct equations except for a term $q=Q_1\sinh(a_4)\,a_4'\,e^{-a_4}$. & PROVED (exact reproduction of all 16 stored equations, in Wolfram Language only) & \texttt{P\_notebo\allowbreak{}okCompar\allowbreak{}e\_cell11\allowbreak{}37Reprod\allowbreak{}uced16of\allowbreak{}16}, \texttt{P\_notebo\allowbreak{}okCompar\allowbreak{}e\_correc\allowbreak{}tVsStore\allowbreak{}dDifferO\allowbreak{}nlyByQ} in \texttt{wolfram-\allowbreak{}primordi\allowbreak{}al-repor\allowbreak{}t.json}; Chapter 9 \\
L20. The $q$ term was produced by the wrong rule of cell 1058. & HYPOTHESIS (a reconstruction that reproduces the stored output exactly; a literal re-execution of cell 1058 gives no $q$ term) & \texttt{P\_notebo\allowbreak{}okCompar\allowbreak{}e\_recons\allowbreak{}truction\allowbreak{}Reproduc\allowbreak{}esStored\allowbreak{}Ela} in \texttt{wolfram-\allowbreak{}primordi\allowbreak{}al-repor\allowbreak{}t.json}; Stage-2 document, §11.3; Chapter 9 \\
L21. Eight-dimensional Einstein gravity needs the negative energy density $\rho_{\mathrm{req}}=-3H^2(7+a_4'^2)/\kappa$ to produce this field. & PROVED (exact symbolic computation) & \texttt{P\_einste\allowbreak{}in\_requi\allowbreak{}redSourc\allowbreak{}e}, \texttt{P\_einste\allowbreak{}in\_rhoRe\allowbreak{}quiredNe\allowbreak{}gative} in \texttt{wolfram-\allowbreak{}primordi\allowbreak{}al-repor\allowbreak{}t.json}; the group check \texttt{P\_einste\allowbreak{}in} in \texttt{python-p\allowbreak{}rimordia\allowbreak{}l-report\allowbreak{}.json}; Chapter 9 \\
L22. For linear $a_4$ ($a_4=ct$ plus a constant) an $x_0$-independent dirac16complex condensate is an exact source if and only if its three-gamma bilinears vanish, $mS=-36H^2/\kappa$ and $\lambda S^2=2H^2(15-3c^2)/\kappa$; such condensates exist, and their energy density is negative. For $a_4''\ne0$ no state that depends on $x_0$ and $x_4$ only is a source. & PROVED, with the field products read as numbers (mean field, ASSUMED) & \texttt{P\_source\allowbreak{}\_x0Indep\allowbreak{}endentSo\allowbreak{}urceCond\allowbreak{}itions}, \texttt{P\_source\allowbreak{}\_x0Indep\allowbreak{}endentEx\allowbreak{}actExamp\allowbreak{}les}, \texttt{P\_source\allowbreak{}\_einstei\allowbreak{}nTransve\allowbreak{}rseDiffe\allowbreak{}renceIs2\allowbreak{}H2a4pp} in \texttt{wolfram-\allowbreak{}primordi\allowbreak{}al-repor\allowbreak{}t.json}; the group check \texttt{\detokenize{P_source}} in \texttt{python-p\allowbreak{}rimordia\allowbreak{}l-report\allowbreak{}.json} (the tensor of the state and the case $a_4''\ne0$); Chapter 9 \\
L23. The static member of the field ($a_4$ constant), written in the warped coordinate $y$: $R=-42H^2$, $\rho_{\mathrm{req}}=-21H^2/\kappa$, $p_{\mathrm{req}}=+15H^2/\kappa$; the mirror brane at $y=0$ carries the stress $S^\mu{}_\mu=-10H/\kappa$ (no sum) for $\mu=1,2,3,5,6,7$ and $S^4{}_4=-12H/\kappa$. & PROVED (exact symbolic computation) & \texttt{KS\_geome\allowbreak{}try\_ricc\allowbreak{}iScalarM\allowbreak{}inus42H2}, \texttt{KS\_geome\allowbreak{}try\_requ\allowbreak{}iredSour\allowbreak{}ce}, \texttt{KS\_geome\allowbreak{}try\_isra\allowbreak{}elStress} in \texttt{wolfram-\allowbreak{}kohn-sha\allowbreak{}m-report\allowbreak{}.json} and \texttt{python-t\allowbreak{}heory-re\allowbreak{}port.jso\allowbreak{}n}; Chapter 9 \\
L24. The continuation of the field across $y=0$ as a mirror copy (the $Z_2$ extension, the notebook's "pair of universes" picture). & ASSUMED (a construction chosen to model the hypothesis) & Chapter 9 \\
L25. The primordial field solves Einstein--Lovelock gravity with a physically acceptable source; what determines $a_4$. & OPEN (the Lovelock terms are not computed) & Chapter 18 \\
\bottomrule
\end{longtable}
\endgroup

\textbf{Numerical experiments of Stage 3} (outputs in \texttt{artifact\allowbreak{}s/dirac1\allowbreak{}6complex\allowbreak{}/numeric\allowbreak{}s/}).

\begingroup
\small
\begin{longtable}{@{}>{\raggedright\arraybackslash}p{0.267\linewidth}>{\raggedright\arraybackslash}p{0.267\linewidth}>{\raggedright\arraybackslash}p{0.267\linewidth}@{}}
\toprule
\textbf{Statement} & \textbf{Status} & \textbf{Where verified} \\
\midrule
\endfirsthead
\toprule
\textbf{Statement} & \textbf{Status} & \textbf{Where verified} \\
\midrule
\endhead
L26. A thermal gas of dirac16complex quanta has $w=0.3329$ at $T=10m$ and $w=0.0359$ after the scale factor of 3-space has grown a hundredfold (on the momentum grid $k\le12T_i$, with the initial temperature $T_i=10m$; 0.0361 without the cut). & COMPUTED & \texttt{exp4/sum\allowbreak{}mary.jso\allowbreak{}n}, keys \texttt{thermal.\allowbreak{}wAtA1} and \texttt{thermal.\allowbreak{}wAtAEnd}; Chapter 11 \\
L27. Dark matter: the gas has the equation of state of dark matter and is created by the expansion for $m>0$, but no dark-matter model is established (abundance, darkness and the stabilization of the extra dimensions are not computed). & COMPUTED (the behaviour); OPEN (a model) & \texttt{exp4/sum\allowbreak{}mary.jso\allowbreak{}n}, \texttt{exp4/pyt\allowbreak{}hon-chec\allowbreak{}k-report\allowbreak{}.json}; Chapter 11 \\
L28. Dark energy: tuned to $w_0=-0.861$ today, the condensate has $w_a=-4.81$; within everything computed the framework gives no viable dark energy. & COMPUTED & \texttt{exp3/fit\allowbreak{}s.json}, key \texttt{scan.tan\allowbreak{}gentWaAt\allowbreak{}UniteW0}; Chapter 11 \\
L29. The five experiments pass 69 self-checks of the Rust program, 167 checks of five independent Python checkers, 71 assertions of the Jupyter notebook and 49 checks of the Mathematica notebook. & COMPUTED & the \texttt{summary.\allowbreak{}json} and \texttt{python-c\allowbreak{}heck-rep\allowbreak{}ort.json} of \texttt{\detokenize{exp1}} to \texttt{\detokenize{exp5}}, \texttt{notebook\allowbreak{}-report.\allowbreak{}json}, \texttt{mathemat\allowbreak{}ica-repo\allowbreak{}rt.json}; Chapter 11 \\
\bottomrule
\end{longtable}
\endgroup

\textbf{Kohn--Sham approximation, Stage 4} (reports in \texttt{artifact\allowbreak{}s/dirac1\allowbreak{}6complex\allowbreak{}/kohn-sh\allowbreak{}am/}: \texttt{wolfram-\allowbreak{}kohn-sha\allowbreak{}m-report\allowbreak{}.json} with 125 of 125 checks, \texttt{python-t\allowbreak{}heory-re\allowbreak{}port.jso\allowbreak{}n} with 157 of 157, the Rust outputs under \texttt{\detokenize{rust/}}, and the cross-check \texttt{python-c\allowbreak{}heck-rep\allowbreak{}ort.json} with 62 of 63).

\begingroup
\small
\begin{longtable}{@{}>{\raggedright\arraybackslash}p{0.267\linewidth}>{\raggedright\arraybackslash}p{0.267\linewidth}>{\raggedright\arraybackslash}p{0.267\linewidth}@{}}
\toprule
\textbf{Statement} & \textbf{Status} & \textbf{Where verified} \\
\midrule
\endfirsthead
\toprule
\textbf{Statement} & \textbf{Status} & \textbf{Where verified} \\
\midrule
\endhead
L30. In the static primordial field the Kohn--Sham equation reduces exactly to eight $2\times2$ first-order systems; four of the blocks carry the negated spectrum of $h-v$ of the other four (the negated spectrum of $h$ when the exchange potential $v$ vanishes). & PROVED (exact symbolic computation) & \texttt{KS\_reduc\allowbreak{}tion\_blo\allowbreak{}ckODEMat\allowbreak{}rix}, \texttt{KS\_reduc\allowbreak{}tion\_hMi\allowbreak{}nusEqual\allowbreak{}sMinusHP\allowbreak{}lus} in both Stage-4 reports; Chapter 13 \\
L31. For the contact interaction the Fock (exchange) energy is exactly local: in the static primordial field, for every quasi-free state built from the eight $2\times2$ blocks, $E_x=-\tfrac\lambda4\int dV\sum_\beta(n_\beta^2+s_\beta^2-t_\beta^2-c_\beta^2)$; for the uniform gas with isotropic occupations it reduces to $e_x=-(\lambda/32)(n^2+S^2)$ at every temperature. & PROVED (exact symbolic computation) & \texttt{KS\_excha\allowbreak{}nge\_bloc\allowbreak{}kFormOfF\allowbreak{}ockTerm}, \texttt{KS\_excha\allowbreak{}nge\_unif\allowbreak{}ormGasCl\allowbreak{}osedForm}, \texttt{KS\_excha\allowbreak{}nge\_angu\allowbreak{}larAvera\allowbreak{}geOfPdot\allowbreak{}QVanishe\allowbreak{}s} in both Stage-4 reports; Chapter 13 \\
L32. The Kohn--Sham model: a static mean field of finitely many quanta in the fixed field, with the exchange of the uniform gas written as a local functional (exact for the uniform gas, an approximation otherwise), no correlation term and no back-reaction on the metric. & ASSUMED (an approximation) & \texttt{README.m\allowbreak{}d}, "Scientific boundary"; Chapters 12 and 13 \\
L33. Kohn--Sham spectra, densities, ground and first excited states, gaps, thermodynamics and energy--momentum tensor of dirac16complex. & COMPUTED by the Rust solver: all 397 of its own checks pass (33 in \texttt{\detokenize{spectrum}}, 137 in \texttt{\detokenize{scf}}, 65 in \texttt{\detokenize{excited}}, 102 in \texttt{\detokenize{thermo}}, 60 in \texttt{\detokenize{emt}}); a repeated run is byte-identical in all 330 files compared; a run with tightened tolerances agrees to $5.96\times10^{-8}$ relative in the energies (65 runs) and to $2.8\times10^{-8}$ in 364829 eigenvalues & the five \texttt{rust/*/s\allowbreak{}ummary.j\allowbreak{}son} and \texttt{rust/det\allowbreak{}erminism\allowbreak{}-report.\allowbreak{}json}; Chapter 13 (Sections 13.10 to 13.15) \\
L34. Agreement of the Rust solver with the independent Python reference solver (56 reference runs, compared with 81 Rust runs). & COMPUTED: 62 of the 63 checks of the final cross-check are true, among them the Delta-SCF energy of the hardest run after erratum E4.13. One check failed, \texttt{canonica\allowbreak{}l\_eigenv\allowbreak{}alues}: in the 301-point Rust run of the smeared ensemble ($N=1016$, $-\hat\lambda_2$) two deep Dirac-sea levels deviate from the reference by up to $2.195\times10^{-6}\,m$ against the tolerance $1.054\times10^{-6}\,m$, while its 601-point run agrees to $1.8\times10^{-7}\,m$; the cause is OPEN (a grid error of the 301-point run is a hypothesis to test), and the Stage-4 gate has not been run & \texttt{python-c\allowbreak{}heck-rep\allowbreak{}ort.json} (entries \texttt{\detokenize{checks}}, \texttt{measurem\allowbreak{}ents}, \texttt{comparis\allowbreak{}ons}); \texttt{handoff/\allowbreak{}specs/ST\allowbreak{}AGE4\_SPE\allowbreak{}C.md}, errata E4.12 and E4.13; Chapter 13 (Sections 13.12 and 13.15); Section 18.7 \\
L35. Whether a Kohn--Sham state can supply the source that the static field requires, $mS=-36H^2/\kappa$ and $\lambda S^2=30H^2/\kappa$. & COMPUTED: no computed Kohn--Sham state meets the conditions; every state with band or bulk levels has a positive average energy density, which would need a negative gravitational coupling $\kappa$ (the reference solver agrees, check \texttt{canonica\allowbreak{}l\_sourci\allowbreak{}ngCondit\allowbreak{}ions}). Whether any many-particle state can be the source, in the fixed field or with back-reaction, is OPEN & \texttt{rust/emt\allowbreak{}/emt-sum\allowbreak{}mary.csv}, column \texttt{sourcing\allowbreak{}\_conditi\allowbreak{}ons\_met}; \texttt{handoff/\allowbreak{}specs/ST\allowbreak{}AGE4\_SPE\allowbreak{}C.md}, erratum E4.1; Chapter 13 (Section 13.14); Section 18.6 \\
\bottomrule
\end{longtable}
\endgroup

\textbf{Pairs of universes, matter and antimatter} (Stage-5 reports in \texttt{artifact\allowbreak{}s/dirac1\allowbreak{}6complex\allowbreak{}/pair-cr\allowbreak{}eation/}: \texttt{wolfram-\allowbreak{}dirac16c\allowbreak{}omplex00\allowbreak{}-report.\allowbreak{}json} with 46 of 46 checks, \texttt{python-d\allowbreak{}irac16co\allowbreak{}mplex00-\allowbreak{}report.j\allowbreak{}son} with 49 of 49, \texttt{wolfram-\allowbreak{}pairing-\allowbreak{}report.j\allowbreak{}son} with 141 of 141 and \texttt{python-p\allowbreak{}airing-r\allowbreak{}eport.js\allowbreak{}on} with 172 of 172; the Python twin of a Wolfram check \texttt{\detokenize{C00_...}} or \texttt{\detokenize{PAIR_...}} is the check \texttt{\detokenize{S5_...}} with the same ending. The Stage-5 numerics are partial and no Stage-5 gate has been run).

\begingroup
\small
\begin{longtable}{@{}>{\raggedright\arraybackslash}p{0.267\linewidth}>{\raggedright\arraybackslash}p{0.267\linewidth}>{\raggedright\arraybackslash}p{0.267\linewidth}@{}}
\toprule
\textbf{Statement} & \textbf{Status} & \textbf{Where verified} \\
\midrule
\endfirsthead
\toprule
\textbf{Statement} & \textbf{Status} & \textbf{Where verified} \\
\midrule
\endhead
L36. dirac16complex00 has a non-trivial Lagrangian with an explicit mass term and a non-trivial coupling to the spin connection; its field equations have the form of those of dirac16complex, its energy--momentum tensor and equations of state follow in every gravitational field, its charge density and its classical energy have no fixed sign, and an exact classical solution is a source of the static primordial field ($m=-30$, $\lambda=125/6$, $S=6/5$ with $H=\kappa=1$, energy density $-21$). & PROVED in general by derivation; exact checks at test points and exact symbolic computations & \texttt{C00\_mass\allowbreak{}Term\_*}, \texttt{C00\_conn\allowbreak{}ection\_n\allowbreak{}onTrivia\allowbreak{}lCouplin\allowbreak{}g}, \texttt{\detokenize{C00_EL_*}}, \texttt{C00\_EMT\_\allowbreak{}*}, \texttt{C00\_char\allowbreak{}ge\_indef\allowbreak{}inite}, \texttt{C00\_ener\allowbreak{}gy\_unbou\allowbreak{}ndedBelo\allowbreak{}w}, \texttt{C00\_prim\allowbreak{}ordial\_s\allowbreak{}taticFie\allowbreak{}ldSource\allowbreak{}dExactly} in \texttt{wolfram-\allowbreak{}dirac16c\allowbreak{}omplex00\allowbreak{}-report.\allowbreak{}json} and their twins in \texttt{python-d\allowbreak{}irac16co\allowbreak{}mplex00-\allowbreak{}report.j\allowbreak{}son}; Chapters 6 and 7 \\
L37. Theorem T1: for both fields and in every gravitational field, a pair $\Psi$ with $(m,\lambda)$ and $\gamma^8\Psi$ with $(-m,-\lambda)$, treated as classical fields, has zero total energy--momentum and zero total charge at every point and every time (the map itself is L16). & PROVED by derivation; exact checks at test points and in the primordial field. At the quantum level the image $\gamma^8\Psi$ is the same quantum system with the anticommutator $-B$, and its cancellation is an identity $X+(-X)=0$ (an exact Fock model, cited from Stage 5 as provisional); a cancellation between two independently quantized universes is OPEN & \texttt{PAIR\_T1g\allowbreak{}eneric\_*}, \texttt{PAIR\_T1g\allowbreak{}rassmann\allowbreak{}\_*}, \texttt{PAIR\_T1j\allowbreak{}ets\_*}, \texttt{PAIR\_T1p\allowbreak{}rimordia\allowbreak{}l\_*}, \texttt{PAIR\_tot\allowbreak{}als\_fiel\allowbreak{}dLevelCh\allowbreak{}iralPair}, \texttt{PAIR\_T1k\allowbreak{}rein\_*} in \texttt{wolfram-\allowbreak{}pairing-\allowbreak{}report.j\allowbreak{}son} and their twins; \texttt{\detokenize{MA_M4_*}}; Chapter 15 (Sections 15.3 and 15.4); Chapter 16 \\
L38. Theorem T2: the reflection of a space-like frame direction, acting on the field and the frame, maps every solution with $(m,\lambda)$ to one with $(-m,\lambda)$, with the same energy--momentum and charge; across the brane of the static primordial field it gives the mirror universe of mass $-M$. & PROVED by derivation in every gravitational field; exact checks & \texttt{PAIR\_T2f\allowbreak{}rame\_*}, \texttt{PAIR\_T2z\allowbreak{}2\_*} in \texttt{wolfram-\allowbreak{}pairing-\allowbreak{}report.j\allowbreak{}son} and their twins; Chapter 15 (Sections 15.5 and 15.6) \\
L39. Theorem T3: in the Kohn--Sham model, for both statistics, the universe of mass $-M$ with the same $\lambda$ and the transformed tip condition has exactly the ground state, first excited state, thermodynamics and energy--momentum tensor of the universe of mass $+M$, and the opposite scalar density; with the untransformed tip condition the pairing fails. & PROVED (Theorem 15.13; exact checks); COMPUTED in part: 25 configurations of dirac16complex in the Rust solver, with the total energies of the two universes equal to at most $1.6\times10^{-8}$; no Rust run of dirac16complex00 and no cross-check report of the pair runs (OPEN) & \texttt{PAIR\_T3b\allowbreak{}lock\_*}, \texttt{PAIR\_T3k\allowbreak{}s\_*}, \texttt{PAIR\_T3e\allowbreak{}mt\_*} in \texttt{wolfram-\allowbreak{}pairing-\allowbreak{}report.j\allowbreak{}son} and their twins; \texttt{rust/pai\allowbreak{}rs/*/pai\allowbreak{}ring.jso\allowbreak{}n}; Chapter 15 (Sections 15.7 and 15.10) \\
L40. The Kohn--Sham approximation for dirac16complex00: within its model the statistics enters only through the sign of the exchange term, $e_x^{00}=+(\lambda/32)(n^2+S^2)$, so that $M_{\text{eff}}^{00}=m+\tfrac{17}{16}\lambda S_p$ and $v^{00}=+\tfrac\lambda{16}n_p$. & PROVED within the model (exact checks); the model itself ASSUMED (a prescription for a classical field, not a quantization); ground and first excited states COMPUTED in one configuration ($m=3H$, $L=3/H$, $N=112$, $\hat\lambda\in\{0,+\hat\lambda_1\}$, $T=0$, and the free state at $T=0.3H$) by the reference solver only, without an independent cross-check; all other configurations not computed & \texttt{PAIR\_sta\allowbreak{}t\_*} in \texttt{wolfram-\allowbreak{}pairing-\allowbreak{}report.j\allowbreak{}son} and \texttt{S5\_stat\_\allowbreak{}*} in \texttt{python-p\allowbreak{}airing-r\allowbreak{}eport.js\allowbreak{}on}; \texttt{referenc\allowbreak{}e/refere\allowbreak{}nce-pair\allowbreak{}s-summar\allowbreak{}y.json}; Chapter 14 \\
L41. "The big bang creates universes of masses $\pm M$ in pairs" (notebook cells 6, 7 and 17). & HYPOTHESIS, not derived: the totals of a cancelling pair of classical fields ($\gamma^8$, Corollary 15.4 and Proposition 16.5) are those of the empty state, so the conservation of the totals does not forbid such a creation, but in the theory as built the separate conservation of each member's charge forbids a pair whose members are charged; at the quantum level the image field $\gamma^8\Psi$ is the same quantum system with the Krein metric $-B$, and no cancellation between two independently quantized universes follows; a creation process, rate or amplitude is OPEN & \texttt{handoff/\allowbreak{}surveys/\allowbreak{}survey\_n\allowbreak{}otebook-\allowbreak{}physics.\allowbreak{}md}; Chapter 16 \\
L42. "This theory explains why the universe contains more matter than antimatter." & not established, and its central step is excluded within the theory as built. The exact U(1) symmetry conserves the charge $Q$ (row L44), so Sakharov's first condition fails. Each field has an exact symmetry that reverses $Q$ without reversing $x_4$ (row L45): an exact C for dirac16complex00 in every gravitational field; for dirac16complex no exact C when $m\ne0$, but an exact CP in flat space and wherever the reflection is an isometry. So the second condition fails. No departure from equilibrium is computed (the third condition is not addressed). The theory contains no Standard-Model baryons and predicts no baryon-to-photon ratio; every scenario is a HYPOTHESIS (row L47) & \texttt{MA\_M2\_sy\allowbreak{}mmetrySu\allowbreak{}mmaryAnd\allowbreak{}ChargeRe\allowbreak{}versal} in \texttt{wolfram-\allowbreak{}matter-a\allowbreak{}ntimatte\allowbreak{}r-report\allowbreak{}.json}; \texttt{MA\_M2\_C\_\allowbreak{}and\_CP\_s\allowbreak{}tatus} and the measurement \texttt{M6\_sakha\allowbreak{}rovScore\allowbreak{}card} in \texttt{python-m\allowbreak{}atter-an\allowbreak{}timatter\allowbreak{}-report.\allowbreak{}json} (both in \texttt{artifact\allowbreak{}s/dirac1\allowbreak{}6complex\allowbreak{}/matter-\allowbreak{}antimatt\allowbreak{}er/}); Chapter 17 (Section 17.15) \\
L43. Signature (4,4) as a description of the observed world, which has one time. & ASSUMED as the setting of the model; the relation to observed spacetime is OPEN & \texttt{README.m\allowbreak{}d}, "Scientific boundary"; Chapter 18 \\
\bottomrule
\end{longtable}
\endgroup

\textbf{Matter and antimatter in detail} (reports in \texttt{artifact\allowbreak{}s/dirac1\allowbreak{}6complex\allowbreak{}/matter-\allowbreak{}antimatt\allowbreak{}er/}: \texttt{wolfram-\allowbreak{}matter-a\allowbreak{}ntimatte\allowbreak{}r-report\allowbreak{}.json} with 44 of 44 checks and \texttt{python-m\allowbreak{}atter-an\allowbreak{}timatter\allowbreak{}-report.\allowbreak{}json} with 77 of 77; the document \texttt{provenan\allowbreak{}ce/DIRAC\allowbreak{}16COMPLE\allowbreak{}X\_MATTER\allowbreak{}\_ANTIMAT\allowbreak{}TER.md}).

\begingroup
\small
\begin{longtable}{@{}>{\raggedright\arraybackslash}p{0.267\linewidth}>{\raggedright\arraybackslash}p{0.267\linewidth}>{\raggedright\arraybackslash}p{0.267\linewidth}@{}}
\toprule
\textbf{Statement} & \textbf{Status} & \textbf{Where verified} \\
\midrule
\endfirsthead
\toprule
\textbf{Statement} & \textbf{Status} & \textbf{Where verified} \\
\midrule
\endhead
L44. Theorem M1: for both fields, every potential $U(S)$ and every gravitational field the Lagrangian is invariant under $\Psi\to e^{i\alpha}\Psi$, and the charge $Q$ of every slice $x_4=\text{const}$ is conserved when no charge flows through the boundary. & PROVED in general by derivation; exact checks at test points & \texttt{MA\_M1\_u1\allowbreak{}Invarian\allowbreak{}ceGrassm\allowbreak{}ann\_G1}, \texttt{MA\_M1\_no\allowbreak{}etherIde\allowbreak{}ntity\_gr\allowbreak{}assmann\_\allowbreak{}G1}, \texttt{MA\_M1\_no\allowbreak{}etherIde\allowbreak{}ntity\_co\allowbreak{}mmuting\_\allowbreak{}G1} in the Wolfram report; the \texttt{MA\_M1\_no\allowbreak{}etherIde\allowbreak{}ntity\_} checks for both statistics and both geometries in the Python report; Chapter 17 (Section 17.9) \\
L45. Theorem M2: among the maps built from a constant matrix, complex conjugation and reflections of frame directions, dirac16complex00 has the exact charge conjugation $\Psi\to\Psi^\ast$ in every gravitational field; dirac16complex has no exact C when $m\ne0$, but an exact CP (a charge conjugation with the reflection of an odd number of space-like directions) in flat space and in gravitational fields with the reflection isometry; each reverses $Q$ and preserves $x_4$. & PROVED (all 1024 maps of the class classified for each statistics by one transformation rule: finite exact computation plus derivation); whether the maps are implemented by operators on the positive Fock space is OPEN & \texttt{MA\_M2\_si\allowbreak{}gnPatter\allowbreak{}nClassif\allowbreak{}ication}, \texttt{MA\_M2\_na\allowbreak{}medTrans\allowbreak{}formatio\allowbreak{}ns}, \texttt{MA\_M2\_sy\allowbreak{}mmetrySu\allowbreak{}mmaryAnd\allowbreak{}ChargeRe\allowbreak{}versal} in the Wolfram report; \texttt{MA\_M2\_di\allowbreak{}screteGr\allowbreak{}oupChara\allowbreak{}cterTabl\allowbreak{}e}, \texttt{MA\_M2\_C\_\allowbreak{}and\_CP\_s\allowbreak{}tatus}, \texttt{MA\_M2\_cp\allowbreak{}ScopeInC\allowbreak{}urvedFie\allowbreak{}lds} in the Python report; Chapter 17 (Sections 17.10 and 17.11); Chapter 6 (Section 6.12) \\
L46. Theorem M3: the bilinears $\Psi^TM\Psi$ invariant under $\mathrm{Spin}_0(4,4)$ are exactly those with $M=C(\alpha P_-+\beta P_+)$, all symmetric, so dirac16complex admits no Majorana-type mass term and dirac16complex00 admits two; every such term carries the charge 2 and is absent from the Lagrangian. & PROVED (finite exact computation plus derivation) & \texttt{MA\_M3\_sp\allowbreak{}inInvari\allowbreak{}antForms}, \texttt{MA\_M3\_gr\allowbreak{}assmannS\allowbreak{}urvival}, \texttt{MA\_M3\_co\allowbreak{}mmutingS\allowbreak{}urvival}, \texttt{MA\_M3\_u1\allowbreak{}Charge} in the Wolfram report; \texttt{MA\_M3\_al\allowbreak{}lInvaria\allowbreak{}ntFormsS\allowbreak{}ymmetric}, \texttt{MA\_M3\_in\allowbreak{}variantF\allowbreak{}ormsSpan\allowbreak{}\_C\_Cgamm\allowbreak{}a8}, \texttt{MA\_M3\_ma\allowbreak{}ssTypeSu\allowbreak{}rvivalAn\allowbreak{}dCharge} in the Python report; Chapter 17 (Section 17.12) \\
L47. The conditional scenario M5: if our universe is one member of a $\gamma^8$ pair of two classical fields in the correlated configuration $\Psi_-=\gamma^8\Psi_+$ (H1), the pair was created with $Q_+\ne0$ (H2), and the charge is baryon number (H3), then the total baryon number of the pair is zero while each member has an excess. & the implication PROVED; H1, H2 and H3 are HYPOTHESES, none derived; the observed baryon-to-photon ratio is not predicted, and the scenario has no quantum counterpart with two independent universes & \texttt{MA\_M5\_im\allowbreak{}plicatio\allowbreak{}n} in the Wolfram report (Wolfram only); Chapter 17 (Section 17.14) \\
\bottomrule
\end{longtable}
\endgroup

\subsection{What we proved and what we assumed}

This chapter proves nothing new; it organizes. It fixes the conventions of the book (counting from 0, the roles of the coordinates, the metric $\eta=\mathrm{diag}(+1,+1,+1,+1,-1,-1,-1,-1)$, units with $\hbar=c=1$), all of which are ASSUMED choices. It defines the five status words, with the three ways in which a statement is PROVED (a derivation, an exact computation of finitely many cases, an exact symbolic computation), and applies them in the ledger. The status of every ledger row is taken from the committed reports and project documents named in its last column. The rows on Stages 4 and 5 describe the state of 2026-09-30: the Stage-4 cross-check has one failed check and the Stage-4 gate has not been run; the exact theory of Stage 5 is complete, but its numerics are partial and its documents and gate do not exist yet. The two items of the request that the project has not established, pair creation of universes by the big bang and the explanation of the matter--antimatter asymmetry, are recorded as a HYPOTHESIS and as not established; for the second, the exact theorems M1 and M2 show that within the theory as built the first two of Sakharov's conditions fail (rows L42, L44 and L45). Chapters 16 and 17 treat both items in full. The consequence of the pairing theorems quoted in Section 0.2, a pair with zero total energy, momentum and charge, holds for classical fields with the second member of mass $-m$ and self-coupling $-\lambda$; the conservation of the totals then does not forbid pair creation, while the separate conservation of each member's charge forbids a pair whose members are charged, and at the quantum level it is an identity between a field and its own image, not a cancellation between two independent universes.

\subsection{Exercises}

Exercise 0.1. Give each statement one of the five status words and one sentence of reason. (a) The matrix $\gamma^4$ of the notebook satisfies $(\gamma^4)^2=-I_{16}$. (b) In EXP-3 the condensate tuned to $w_0=-0.861$ has $w_a=-4.81$. (c) The field lives on an eight-dimensional space of signature (4,4). (d) At $x_4=0$ a pair of universes of masses $+M$ and $-M$ is created. (e) Which physical law fixes the function $a_4$ of the primordial field?

Exercise 0.2. The covariant constancy of the gamma matrices (ledger row L6) was checked exactly at a few test points. (a) Why do these checks alone not prove it for every gravitational field, and what does prove it? (b) The failure of the notebook's contraction (row L7) was also found only at test points. Why is row L7 nevertheless completely proved?

Exercise 0.3. With the method of Section 0.8, count the true checks of the two Stage-2 reports

\begin{Verbatim}[fontsize=\small]
artifacts/dirac16complex/primordial-field/wolfram-primordial-report.json
artifacts/dirac16complex/primordial-field/python-primordial-report.json
\end{Verbatim}

and compare with the table of Section 0.5.

Exercise 0.4. Counting from 0. (a) What is the entry $\eta_{44}$ of the metric $\eta$, and what is $\eta_{33}$? (b) Which coordinate is the fifth entry of the list $(x_0,x_1,\dots,x_7)$, and what is its role? (c) Mathematica counts from 1, so the notebook's matrix entry SAB[[a,b]] with $a,b=1,\dots,8$ is the book's $S^{a-1\,b-1}$. Which $S^{ab}$ is SAB[[3,5]]? (d) How many components does the field have, and what is the label of the last one?

Exercise 0.5. Row L41 of the ledger is a HYPOTHESIS. (a) Explain why a theorem that maps every solution with mass $+m$ to a solution with mass $-m$ (and self-coupling $-\lambda$), and shows that the pair has zero total energy, momentum and charge, cannot by itself turn L41 into a PROVED statement. (b) Name one kind of calculation that would be needed in addition.

Exercise 0.6. The Stage-4 row of the table in Section 0.5 contains results of two different kinds. Name them and give the status of each.

\subsection{Answers to the exercises}

Answer 0.1. (a) PROVED: it is one of the 64 equations $\{\gamma^a,\gamma^b\}=2\eta^{ab}I_{16}$ of row L1 (take $a=b=4$: by the definition $\{A,B\}=AB+BA$ of Section 1.7, $\{\gamma^4,\gamma^4\}=2(\gamma^4)^2$, and $\eta^{44}=\eta_{44}=-1$ (Section 1.11), so $2(\gamma^4)^2=-2I_{16}$), and a finite computation with integers is a complete proof. (b) COMPUTED: it is the result of a numerical solution (row L28), reproducible but not a proof, and it holds within the accuracy of the solver. (c) ASSUMED: it is the setting of the model, chosen by the notebook, not derived (row L43). (d) HYPOTHESIS: it is the notebook's proposal (row L41); nothing in the project derives it. (e) OPEN: no computation of the project determines $a_4$ (row L25).

Answer 0.2. (a) Row L6 is a statement about every metric, that is about infinitely many cases; checks at finitely many points can only confirm it, since an error could show up at a point that was not tested. It is proved by the general derivation of Chapter 4, which uses only the definition of the spin connection by the vielbein postulate and the definition of the covariant derivative. The checks confirm that derivation exactly and would have detected many kinds of error. (b) Row L7 says that the notebook's contraction does NOT make the gammas covariantly constant in general. To disprove "for every metric, X holds" one example of a metric in which X fails is enough, and the exact checks exhibit such examples (in G1 all 64 pairs of indices fail). A single counterexample is a complete proof of a "not".

Answer 0.3. Following Section 0.8 with the two paths gives \texttt{(126, 12\allowbreak{}6)} for the Wolfram report and \texttt{\detokenize{(16, 16)}} for the Python report: all 126 and all 16 checks are true, as the table of Section 0.5 says.

Answer 0.4. (a) The diagonal of $\eta$ is $(+1,+1,+1,+1,-1,-1,-1,-1)$ with the entries numbered $0,1,\dots,7$. Entry 4 is the fifth one, $-1$, so $\eta_{44}=-1$; entry 3 is the fourth one, $+1$, so $\eta_{33}=+1$. (b) The fifth entry is $x_4$, the time in which everything evolves. (c) SAB[[3,5]] is $S^{ab}$ with $a=3-1=2$ and $b=5-1=4$, written $S^{2\,4}$ (two indices, 2 and 4). (d) Sixteen components, $\Psi_0,\dots,\Psi_{15}$; the last one is $\Psi_{15}$.

Answer 0.5. (a) Such a theorem says, at the level of classical fields: if a universe of mass $+M$ exists as a solution, then a partner of mass $-M$ (and self-coupling $-\lambda$) exists as a solution, and the two together have the totals of the empty state, so the conservation of the totals does not forbid them to appear from it (in the theory as built the charge of each member is still conserved on its own, which forbids a pair whose members are charged, Chapter 16). That is a statement about which configurations are allowed, not about which configurations occur, when, or how often. The same situation exists in ordinary physics: when two photons of enough energy collide, the conservation of energy, momentum and charge allows them to turn into an electron and a positron, but whether this happens, and with which probability, is computed from the dynamics (quantum electrodynamics), not from the conservation laws. (b) One needs a dynamical calculation: for example a quantum amplitude, or a probability per unit time, for the transition from a state without universes to a state with the pair, or a classical solution of the coupled field and gravity equations that evolves from an initial singular state into the pair. Neither is derived in this project (Chapter 16 and Chapter 18).

Answer 0.6. The exact theory of Stage 4 (the reduction of the Kohn--Sham equation to $2\times2$ blocks, the exchange formula and the other exact statements, 125 of 125 Wolfram and 157 of 157 Python checks) is PROVED (rows L30 and L31). The numerical Kohn--Sham results are COMPUTED (row L33). Their comparison with the independent reference solver is COMPUTED as well (62 of 63 checks true), and the cause of the one failed comparison is OPEN (row L34).

\section{Mathematical toolkit from zero}

The physics of this book is written in a small number of mathematical languages: vectors and matrices, complex numbers, indices with a summation rule, derivatives of functions of several variables, and metrics, which measure lengths in spaces whose directions need not all behave alike. This chapter builds these languages from school algebra and the calculus of one variable. Every later chapter uses them. The examples are deliberately small (two or three dimensions, matrices of size $2\times2$), because every rule that holds for them holds in the same form for the eight dimensions and $16\times16$ matrices of the rest of the book.

Where a standard fact of mathematics is used without proof, the text says so, and Section 1.13 lists all such facts.

\subsection{Numbers, lists, sums and powers}

\textbf{Numbers.} We use the natural numbers $0$, $1$, $2$, $3$, and so on (we include 0), the integers, which also contain the negative whole numbers $-1$, $-2$, $-3$, and so on, the fractions (rational numbers) such as $-7/3$, and the real numbers, which include numbers like $\sqrt2$ and $\pi$ that are not fractions. The set of real numbers is written $\mathbb R$. The symbol $\in$ means "is an element of": $x\in\mathbb R$ says that $x$ is a real number. An \textbf{open interval} $(a,b)$ is the set of real numbers $x$ with $a<x<b$; for example $z\in(0,\pi/2)$ means $0<z<\pi/2$. Complex numbers are introduced in Section 1.5.

\textbf{Functions.} A function $f$ assigns to each allowed input exactly one output. A function of several variables takes several inputs, such as $f(x_0,x_4)=x_0^2\sin x_4$, which assigns to the pair $(x_0,x_4)$ the number $x_0^2\sin x_4$.

\textbf{Lists, counted from 0.} A list (also called a tuple) is an ordered collection of entries. This book counts entries from 0: the list $(x_0,x_1,\dots,x_7)$ has the eight entries $x_0$ to $x_7$, entry number $k$ is $x_k$, and the last entry has the number 7, one less than the number of entries. In the list $(5,7,2)$ entry 0 is 5, entry 1 is 7 and entry 2 is 2. Counting from 0 is a convention (Chapter 0); it is the convention of every document of the project, and it makes many formulas shorter.

\textbf{Sums and products.} The symbol $\sum$ (capital Greek sigma) abbreviates a sum and $\prod$ (capital pi) a product:

\[
\sum_{k=0}^{n-1}a_k=a_0+a_1+\dots+a_{n-1},\qquad \prod_{k=0}^{n-1}a_k=a_0\,a_1\cdots a_{n-1}.
\]

The letter $k$ is a \textbf{dummy index}: it can be renamed without changing the value, $\sum_{k=0}^{3}k^2=\sum_{j=0}^{3}j^2=0+1+4+9=14$. Two rules follow from the ordinary laws of arithmetic and are used constantly: a constant factor can be taken out of a sum, $\sum_kc\,a_k=c\sum_ka_k$, and a sum of sums can be added term by term, $\sum_k(a_k+b_k)=\sum_ka_k+\sum_kb_k$. A double sum $\sum_j\sum_k a_{jk}$ adds all entries $a_{jk}$ of a rectangular array, and it can be done in either order, because addition is commutative and associative.

\textbf{Powers and roots.} For positive real numbers $a,b$ and real exponents $p,q$:

\[
a^pa^q=a^{p+q},\qquad (a^p)^q=a^{pq},\qquad (ab)^p=a^pb^p,\qquad a^{-p}=\frac1{a^p},\qquad \sqrt a=a^{1/2}.
\]

These rules are school algebra and are used here without proof. They are exactly the rules that decide whether a simplification is correct. For example, for $0<s$ and real $a$,

\[
\frac1{\sqrt{s^{1/3}/e^{2a}}}=\Bigl(\frac{s^{1/3}}{e^{2a}}\Bigr)^{-1/2}=\frac{(s^{1/3})^{-1/2}}{(e^{2a})^{-1/2}}=\frac{s^{-1/6}}{e^{-a}}=\frac{e^{a}}{s^{1/6}} .
\]

Exercise 1.3 uses this computation to check a substitution rule of the author's notebook.

\textbf{Elementary functions.} We use the exponential function $e^x$, its inverse, the natural logarithm $\ln x$ (for $x>0$), and the trigonometric functions $\sin$, $\cos$, $\tan=\sin/\cos$ and $\cot=\cos/\sin$, with $\cos^2x+\sin^2x=1$ and the addition theorems

\[
\cos(\alpha+\beta)=\cos\alpha\cos\beta-\sin\alpha\sin\beta,\qquad \sin(\alpha+\beta)=\sin\alpha\cos\beta+\cos\alpha\sin\beta ,
\]

which are also school mathematics and are used without proof. The hyperbolic functions are defined by

\[
\cosh x=\frac{e^x+e^{-x}}2,\qquad \sinh x=\frac{e^x-e^{-x}}2 .
\]

They satisfy $\cosh^2x-\sinh^2x=1$. Proof: $\cosh^2x-\sinh^2x=\tfrac14\bigl[(e^{2x}+2+e^{-2x})-(e^{2x}-2+e^{-2x})\bigr]=\tfrac14\cdot4=1$.

\subsection{Vectors}

\textbf{Definition.} A \textbf{vector} with $n$ components is a column of $n$ numbers,

\[
v=\begin{pmatrix}v_0\\ v_1\\ \vdots\\ v_{n-1}\end{pmatrix}.
\]

The numbers $v_0,\dots,v_{n-1}$ are its \textbf{components}, counted from 0. Vectors with the same number of components are added component by component, and a vector is multiplied by a number $c$ (a \textbf{scalar}) by multiplying every component:

\[
(v+w)_k=v_k+w_k,\qquad (cv)_k=c\,v_k,\qquad k=0,\dots,n-1 .
\]

The \textbf{zero vector} $0$ has all components 0. The set of all vectors with $n$ real components is written $\mathbb R^n$; with complex components (Section 1.5) it is written $\mathbb C^n$. The fields of this book have sixteen components at every point, written $\Psi=(\Psi_0,\dots,\Psi_{15})^T$ (the $T$ is explained in Section 1.3). For the commuting field dirac16complex00 the components are complex numbers, so at each point $\Psi$ is a vector in $\mathbb C^{16}$. For dirac16complex they are complex anticommuting (Grassmann) quantities, introduced in Chapter 5. Matrices act on them by the same rule $(A\Psi)_i=\sum_jA_{ij}\Psi_j$ (Section 1.3), but a product of two components changes sign when the two factors are exchanged (Section 1.8 and Chapter 5).

\textbf{Example.} In $\mathbb R^3$, with $v=(1,2,0)^T$ and $w=(0,1,1)^T$: $v+w=(1,3,1)^T$ and $3v=(3,6,0)^T$.

\textbf{Linear combinations and bases.} A \textbf{linear combination} of vectors $u_0,\dots,u_{m-1}$ is a vector $c_0u_0+\dots+c_{m-1}u_{m-1}$ with numbers $c_k$. The vectors are \textbf{linearly independent} when the only linear combination that gives the zero vector is the one with all $c_k=0$; otherwise they are \textbf{linearly dependent}, and then one of them is a linear combination of the others (if $\sum_kc_ku_k=0$ with some $c_j\ne0$, then $u_j=-\sum_{k\ne j}(c_k/c_j)\,u_k$). The \textbf{standard basis} of $\mathbb R^n$ consists of the $n$ vectors $e_0,\dots,e_{n-1}$, where $e_k$ has the component 1 in place $k$ and 0 elsewhere. Every vector is a linear combination of them, $v=\sum_{k=0}^{n-1}v_ke_k$, because component $j$ of the right-hand side is $\sum_kv_k(e_k)_j=v_j$ (only the term $k=j$ contributes). The standard basis vectors are linearly independent: $\sum_kc_ke_k$ is the vector with components $c_k$, and it is zero only if every $c_k$ is zero.

A \textbf{basis} is a list of linearly independent vectors of which every vector is a linear combination; the expansion coefficients are then unique (if $v=\sum_kc_ku_k=\sum_kc'_ku_k$, then $\sum_k(c_k-c'_k)u_k=0$, so $c_k=c'_k$ for every $k$ by independence). All bases of $\mathbb R^n$ (or $\mathbb C^n$) have exactly $n$ vectors; this number is the \textbf{dimension}. (This last fact is standard linear algebra and is used without proof.) The dimension counts basis vectors, not vectors: $\mathbb R^2$ contains infinitely many vectors but has dimension 2.

\textbf{Example.} In $\mathbb R^2$ the vectors $(1,1)^T$ and $(1,-1)^T$ are linearly independent: $c_0(1,1)^T+c_1(1,-1)^T=(c_0+c_1,\,c_0-c_1)^T$ is zero only if $c_0+c_1=0$ and $c_0-c_1=0$, that is $c_0=c_1=0$. The vectors $(1,2)^T$ and $(2,4)^T$ are linearly dependent, since $2(1,2)^T-(2,4)^T=0$.

\subsection{Matrices}

\textbf{Definition.} A \textbf{matrix} with $m$ rows and $n$ columns (an $m\times n$ matrix) is a rectangular array of numbers $A_{ij}$, where $i=0,\dots,m-1$ numbers the row and $j=0,\dots,n-1$ the column. A vector with $n$ components is an $n\times1$ matrix. Matrices of the same shape are added entry by entry and multiplied by numbers entry by entry.

\textbf{Matrix times vector.} An $m\times n$ matrix $A$ turns a vector $v$ with $n$ components into a vector $Av$ with $m$ components:

\[
(Av)_i=\sum_{j=0}^{n-1}A_{ij}v_j .
\]

Component $i$ of $Av$ is row $i$ of $A$ multiplied entry by entry with $v$ and added. For example

\[
\begin{pmatrix}1&2\\ 3&4\end{pmatrix}\begin{pmatrix}5\\ 6\end{pmatrix}=\begin{pmatrix}1\cdot5+2\cdot6\\ 3\cdot5+4\cdot6\end{pmatrix}=\begin{pmatrix}17\\ 39\end{pmatrix}.
\]

The map $v\mapsto Av$ is \textbf{linear}: $A(v+w)=Av+Aw$ and $A(cv)=c\,Av$, as follows directly from the formula.

\textbf{Matrix product.} The product of an $m\times n$ matrix $A$ and an $n\times p$ matrix $B$ is the $m\times p$ matrix

\[
(AB)_{ik}=\sum_{j=0}^{n-1}A_{ij}B_{jk} .
\]

It is defined so that $(AB)v=A(Bv)$ for every vector $v$: component $i$ of $A(Bv)$ is $\sum_jA_{ij}\sum_kB_{jk}v_k=\sum_k\bigl(\sum_jA_{ij}B_{jk}\bigr)v_k$. Example:

\[
\begin{pmatrix}1&2\\ 0&1\end{pmatrix}\begin{pmatrix}0&1\\ 1&0\end{pmatrix}=\begin{pmatrix}1\cdot0+2\cdot1&1\cdot1+2\cdot0\\ 0\cdot0+1\cdot1&0\cdot1+1\cdot0\end{pmatrix}=\begin{pmatrix}2&1\\ 1&0\end{pmatrix}.
\]

In the other order the same two matrices give $\begin{pmatrix}0&1\\ 1&0\end{pmatrix}\begin{pmatrix}1&2\\ 0&1\end{pmatrix}=\begin{pmatrix}0&1\\ 1&2\end{pmatrix}$, a different matrix. \textbf{Matrix multiplication is not commutative}: in general $AB\ne BA$. This single fact is the reason why the order of factors matters in every formula of this book. Matrix multiplication is, however, \textbf{associative}, $(AB)C=A(BC)$: both sides have the entries $\sum_j\sum_kA_{ij}B_{jk}C_{kl}$, and it is \textbf{distributive}, $A(B+C)=AB+AC$ and $(A+B)C=AC+BC$.

\textbf{Special matrices.} A \textbf{square} matrix has as many rows as columns. The \textbf{identity matrix} $I_n$ (or simply $I$, or $1$) is the $n\times n$ matrix with 1 on the diagonal ($i=j$) and 0 elsewhere; $Iv=v$ and $IA=AI=A$. Its entries are written with the \textbf{Kronecker delta}, $\delta_{ij}=1$ if $i=j$ and $0$ otherwise. A \textbf{diagonal matrix} has zero entries off the diagonal; $\mathrm{diag}(d_0,\dots,d_{n-1})$ denotes the diagonal matrix with diagonal $d_0,\dots,d_{n-1}$. Two diagonal matrices are multiplied by multiplying their diagonals: $\mathrm{diag}(d_k)\,\mathrm{diag}(d'_k)=\mathrm{diag}(d_kd'_k)$, so diagonal matrices commute with each other.

\textbf{Transpose.} The \textbf{transpose} $A^T$ of an $m\times n$ matrix is the $n\times m$ matrix with $(A^T)_{ij}=A_{ji}$: rows become columns. The transpose of a column vector is a row vector, which is why we write $v=(v_0,\dots,v_{n-1})^T$ to save space. The transpose of a product is the product of the transposes in the reverse order,

\[
(AB)^T=B^TA^T .
\]

Proof: $((AB)^T)_{ki}=(AB)_{ik}=\sum_jA_{ij}B_{jk}=\sum_j(B^T)_{kj}(A^T)_{ji}=(B^TA^T)_{ki}$. For three factors, $(ABC)^T=C^TB^TA^T$ by applying the rule twice.

\textbf{Symmetric and antisymmetric matrices.} A square matrix is \textbf{symmetric} if $A^T=A$ ($A_{ij}=A_{ji}$) and \textbf{antisymmetric} if $A^T=-A$ ($A_{ij}=-A_{ji}$). The diagonal of an antisymmetric matrix is zero, because $A_{ii}=-A_{ii}$ forces $A_{ii}=0$. Every square matrix $M$ is, in exactly one way, the sum of a symmetric and an antisymmetric matrix:

\[
M=M_S+M_A,\qquad M_S=\tfrac12(M+M^T),\qquad M_A=\tfrac12(M-M^T).
\]

Proof: $M_S^T=\tfrac12(M^T+M)=M_S$, $M_A^T=\tfrac12(M^T-M)=-M_A$, and $M_S+M_A=M$. If also $M=S+A$ with $S$ symmetric and $A$ antisymmetric, then $M^T=S-A$, so $S=\tfrac12(M+M^T)=M_S$ and $A=M_A$: the split is unique. This split decides, in Chapter 5, which parts of the notebook's Lagrangian survive for anticommuting fields.

\textbf{Trace.} The \textbf{trace} of a square matrix is the sum of its diagonal entries, $\mathrm{tr}A=\sum_iA_{ii}$. Although $AB\ne BA$ in general, the two products have the same trace:

\[
\mathrm{tr}(AB)=\sum_i\sum_jA_{ij}B_{ji}=\sum_j\sum_iB_{ji}A_{ij}=\mathrm{tr}(BA).
\]

In the example above both products have trace 2.

\textbf{Inverse.} A square matrix $A$ is \textbf{invertible} if there is a matrix $A^{-1}$ with $AA^{-1}=A^{-1}A=I$. For a $2\times2$ matrix,

\[
A=\begin{pmatrix}a&b\\ c&d\end{pmatrix},\qquad A^{-1}=\frac1{ad-bc}\begin{pmatrix}d&-b\\ -c&a\end{pmatrix}\qquad(ad-bc\ne0),
\]

which is checked by multiplying out: $\begin{pmatrix}a&b\\ c&d\end{pmatrix}\begin{pmatrix}d&-b\\ -c&a\end{pmatrix}=\begin{pmatrix}ad-bc&0\\ 0&ad-bc\end{pmatrix}$. The inverse of a product is the product of the inverses in reverse order, $(AB)^{-1}=B^{-1}A^{-1}$, because $(AB)(B^{-1}A^{-1})=A(BB^{-1})A^{-1}=AA^{-1}=I$ (and similarly in the other order). A diagonal matrix $\mathrm{diag}(d_k)$ with all $d_k\ne0$ has the inverse $\mathrm{diag}(1/d_k)$.

\textbf{Block matrices.} A large matrix can be cut into smaller matrices, called blocks, and multiplied block by block as if the blocks were numbers, provided the order of the factors inside each block product is kept. For matrices cut into $2\times2$ blocks of matching sizes,

\[
\begin{pmatrix}A&B\\ C&D\end{pmatrix}\begin{pmatrix}E&F\\ G&H\end{pmatrix}=\begin{pmatrix}AE+BG&AF+BH\\ CE+DG&CF+DH\end{pmatrix}.
\]

Proof: an entry of the product in, say, the upper-left block is a sum over the column index of the first factor; splitting that sum into the columns belonging to the first block column and those belonging to the second gives exactly the entry of $AE$ plus the entry of $BG$. A case that occurs throughout the book is the off-diagonal block matrix

\[
M=\begin{pmatrix}0&X\\ Y&0\end{pmatrix},\qquad M^2=\begin{pmatrix}0\cdot0+XY&0\cdot X+X\cdot0\\ Y\cdot0+0\cdot Y&YX+0\cdot0\end{pmatrix}=\begin{pmatrix}XY&0\\ 0&YX\end{pmatrix}.
\]

The notebook builds its $16\times16$ gamma matrices in exactly this form, with $8\times8$ blocks (Chapters 2 and 3). For square matrices $D_0,\dots,D_{k-1}$, $\mathrm{diag}(D_0,\dots,D_{k-1})$ denotes the block matrix with these blocks on the diagonal and zero blocks elsewhere; for example $M^2=\mathrm{diag}(XY,YX)$ above, and $\mathrm{diag}(-I_8,I_8)$ is the $16\times16$ diagonal matrix with eight entries $-1$ followed by eight entries $+1$.

\textbf{Tensor (Kronecker) product.} For an $m\times m$ matrix $A$ and an $n\times n$ matrix $B$, the \textbf{tensor product} $A\otimes B$ is the $mn\times mn$ matrix made of an $m\times m$ arrangement of $n\times n$ blocks, block $(i,j)$ being $A_{ij}B$. For $2\times2$ matrices:

\[
A\otimes B=\begin{pmatrix}A_{00}B&A_{01}B\\ A_{10}B&A_{11}B\end{pmatrix},\qquad \text{for example}\qquad \begin{pmatrix}1&0\\ 0&-1\end{pmatrix}\otimes\begin{pmatrix}0&1\\ 1&0\end{pmatrix}=\begin{pmatrix}0&1&0&0\\ 1&0&0&0\\ 0&0&0&-1\\ 0&0&-1&0\end{pmatrix}.
\]

It obeys the \textbf{mixed-product rule} $(A\otimes B)(C\otimes D)=(AC)\otimes(BD)$. Proof: by the block rule, block $(i,k)$ of the product is $\sum_j(A_{ij}B)(C_{jk}D)=\bigl(\sum_jA_{ij}C_{jk}\bigr)BD=(AC)_{ik}\,BD$, which is block $(i,k)$ of $(AC)\otimes(BD)$. Four-fold tensor products of $2\times2$ matrices give $16\times16$ matrices; this is the second standard way of writing the gamma matrices of this book (the tensor-product picture of Section 2.7; Section 3.11 relates it to the notebook's matrices).

\subsection{Determinants}

\textbf{The $2\times2$ case.} The \textbf{determinant} of a $2\times2$ matrix is

\[
\det\begin{pmatrix}a&b\\ c&d\end{pmatrix}=ad-bc .
\]

It is the factor by which the matrix changes areas (up to sign): the square with corners $0$, $e_0$, $e_1$, $e_0+e_1$ is mapped to the parallelogram spanned by the two columns of the matrix, whose area is $|ad-bc|$ (a fact of plane geometry, used here without proof and not needed later in the book).

The determinant of a product is the product of the determinants. For $2\times2$ matrices this is a direct computation: with $A=\begin{pmatrix}a&b\\ c&d\end{pmatrix}$ and $B=\begin{pmatrix}e&f\\ g&h\end{pmatrix}$,

\[
\begin{aligned}
\det(AB)&=(ae+bg)(cf+dh)-(af+bh)(ce+dg)\\
&=aecf+aedh+bgcf+bgdh-afce-afdg-bhce-bhdg\\
&=ad(eh-fg)-bc(eh-fg)=\det A\,\det B ,
\end{aligned}
\]

where in the second line the terms $aecf$ and $afce$ cancel, and so do $bgdh$ and $bhdg$. Consequently a $2\times2$ matrix is invertible exactly when its determinant is not zero: if $ad-bc\ne0$, the formula of Section 1.3 gives the inverse, and if $A$ is invertible, then $\det A\,\det A^{-1}=\det I=1$, so $\det A\ne0$.

\textbf{Permutations and their signs.} A \textbf{permutation} of $(0,1,\dots,n-1)$ is a reordering $(\sigma(0),\sigma(1),\dots,\sigma(n-1))$ of these numbers. An \textbf{inversion} of $\sigma$ is a pair of places $i<j$ with $\sigma(i)>\sigma(j)$. The \textbf{sign} of $\sigma$ is $\mathrm{sgn}\,\sigma=(-1)^{\text{number of inversions}}$: $+1$ for an even and $-1$ for an odd number of inversions. For $n=3$ the six permutations are

\begingroup
\small
\begin{longtable}{@{}>{\raggedright\arraybackslash}p{0.267\linewidth}>{\raggedright\arraybackslash}p{0.267\linewidth}>{\raggedright\arraybackslash}p{0.267\linewidth}@{}}
\toprule
\textbf{permutation} & \textbf{inversions} & \textbf{sign} \\
\midrule
\endfirsthead
\toprule
\textbf{permutation} & \textbf{inversions} & \textbf{sign} \\
\midrule
\endhead
(0,1,2) & none & $+1$ \\
(1,2,0) & (1,0), (2,0) & $+1$ \\
(2,0,1) & (2,0), (2,1) & $+1$ \\
(1,0,2) & (1,0) & $-1$ \\
(0,2,1) & (2,1) & $-1$ \\
(2,1,0) & (2,1), (2,0), (1,0) & $-1$ \\
\bottomrule
\end{longtable}
\endgroup

(an inversion is listed by the two entries that stand in the wrong order). Exchanging two entries of a permutation always changes its sign; this is a standard fact that we use without proof.

\textbf{The general determinant.} For an $n\times n$ matrix,

\[
\det A=\sum_{\sigma}\mathrm{sgn}(\sigma)\,A_{0\sigma(0)}A_{1\sigma(1)}\cdots A_{n-1\,\sigma(n-1)},
\]

the sum running over all $n!$ permutations. Each term takes exactly one entry from every row and from every column. For $n=2$ the permutations $(0,1)$ and $(1,0)$ give $A_{00}A_{11}-A_{01}A_{10}$, the formula above. For $n=3$ the six permutations of the table give

\[
\det A=A_{00}A_{11}A_{22}+A_{01}A_{12}A_{20}+A_{02}A_{10}A_{21}-A_{01}A_{10}A_{22}-A_{00}A_{12}A_{21}-A_{02}A_{11}A_{20}.
\]

For a \textbf{diagonal} matrix every term except the one of the identity permutation contains an off-diagonal entry, which is zero, so

\[
\det\mathrm{diag}(d_0,\dots,d_{n-1})=d_0d_1\cdots d_{n-1}.
\]

This formula is used in Section 1.11 to compute the determinant of the metric of the primordial field.

Four further facts hold for every $n$ and are used without proof: $\det(AB)=\det A\det B$ (proved above for $n=2$); $\det A^T=\det A$; a square matrix $M$ is invertible exactly when $\det M\ne0$ (proved above for $n=2$); and a square matrix $M$ is invertible exactly when $Mv=0$ has only the solution $v=0$. (One direction of the last fact is immediate: if $M$ is invertible and $Mv=0$, then $v=M^{-1}Mv=M^{-1}0=0$. The other direction, that $M$ is invertible whenever $Mv=0$ forces $v=0$, is the part used without proof.) Together the last two facts say:

\[
Mv=0\ \text{has a solution}\ v\ne0\quad\text{exactly when}\quad \det M=0 .
\]

\textbf{Example.} For the matrix with rows $(1,2,0)$, $(0,1,1)$, $(1,3,1)$ the $3\times3$ formula gives $1\cdot1\cdot1+2\cdot1\cdot1+0-2\cdot0\cdot1-1\cdot1\cdot3-0=1+2-3=0$. The determinant vanishes because the third row is the sum of the first two, so the rows are linearly dependent.

\subsection{Complex numbers}

\textbf{Definition.} The equation $x^2=-1$ has no real solution. We introduce a new number $i$ with $i^2=-1$ and consider all numbers $z=a+ib$ with real $a$ and $b$; they are the \textbf{complex numbers}, and their set is $\mathbb C$. The real number $a=\mathrm{Re}\,z$ is the real part and $b=\mathrm{Im}\,z$ the imaginary part. Complex numbers are added and multiplied with the ordinary rules of algebra, replacing $i^2$ by $-1$ wherever it appears:

\[
(a+ib)+(c+id)=(a+c)+i(b+d),\qquad (a+ib)(c+id)=(ac-bd)+i(ad+bc).
\]

For example $(1+2i)(3-i)=3-i+6i-2i^2=5+5i$.

\textbf{Conjugate and modulus.} The \textbf{complex conjugate} of $z=a+ib$ is $z^\ast=a-ib$. The product $zz^\ast=a^2+b^2$ is real and not negative; its square root $|z|=\sqrt{a^2+b^2}$ is the \textbf{modulus} of $z$. Every $z\ne0$ has the inverse

\[
\frac1z=\frac{z^\ast}{zz^\ast}=\frac{z^\ast}{|z|^2},\qquad\text{for example}\qquad \frac1{1+i}=\frac{1-i}{2}.
\]

Conjugation respects sums and products, $(z+w)^\ast=z^\ast+w^\ast$ and $(zw)^\ast=z^\ast w^\ast$. Proof of the second rule: with $z=a+ib$ and $w=c+id$, $(zw)^\ast=(ac-bd)-i(ad+bc)$, and $z^\ast w^\ast=(a-ib)(c-id)=ac-iad-ibc+i^2bd=(ac-bd)-i(ad+bc)$. It follows that $|zw|^2=zw\,z^\ast w^\ast=|z|^2|w|^2$. A number is real exactly when $z^\ast=z$, and purely imaginary (of the form $ib$) exactly when $z^\ast=-z$.

\textbf{The complex exponential.} For real $\theta$ we define

\[
e^{i\theta}=\cos\theta+i\sin\theta .
\]

This definition has the property that justifies the notation, the law of exponents:

\[
e^{i\alpha}e^{i\beta}=(\cos\alpha\cos\beta-\sin\alpha\sin\beta)+i(\sin\alpha\cos\beta+\cos\alpha\sin\beta)=\cos(\alpha+\beta)+i\sin(\alpha+\beta)=e^{i(\alpha+\beta)},
\]

by the addition theorems of Section 1.1. Moreover $|e^{i\theta}|^2=\cos^2\theta+\sin^2\theta=1$ and $(e^{i\theta})^\ast=\cos\theta-i\sin\theta=e^{-i\theta}$.

\textbf{Derivatives of complex-valued functions.} A complex-valued function $f(\theta)=a(\theta)+i\,b(\theta)$ of a real variable, with real functions $a$ and $b$, is differentiated part by part: $f'=a'+i\,b'$. The sum rule $(f+g)'=f'+g'$ follows at once. The product rule $(fg)'=f'g+fg'$ also holds. Proof: with $g=c+id$ ($c$, $d$ real), the multiplication rule above gives $fg=(ac-bd)+i(ad+bc)$. Applying the real product rule to each of the four products and regrouping,

\[
\begin{aligned}
(fg)'&=(a'c+ac'-b'd-bd')+i(a'd+ad'+b'c+bc')\\
&=\bigl[(a'c-b'd)+i(a'd+b'c)\bigr]+\bigl[(ac'-bd')+i(ad'+bc')\bigr]=f'g+fg' ,
\end{aligned}
\]

because the first bracket is $(a'+ib')(c+id)$ and the second is $(a+ib)(c'+id')$. The same part-by-part rule is used for matrices with complex entries in Section 1.9. With it, the derivative of $e^{i\theta}=\cos\theta+i\sin\theta$ with respect to $\theta$ is

\[
\frac{d}{d\theta}e^{i\theta}=-\sin\theta+i\cos\theta=i(\cos\theta+i\sin\theta)=i\,e^{i\theta},
\]

the same rule as for the real exponential with the constant $i$ in place of a real constant. For a general complex number $z=x+iy$ we define $e^z=e^xe^{iy}$; then $e^ze^w=e^{z+w}$ for all complex $z,w$, by the two laws of exponents. Every complex number can be written in \textbf{polar form} $z=r\,e^{i\theta}$ with $r=|z|$ and a real angle $\theta$: for $z=a+ib\ne0$ the point $(a/r,b/r)$ of the plane has distance 1 from the origin, so it is $(\cos\theta,\sin\theta)$ for some angle $\theta$ (a fact of school trigonometry, used without proof), and then $z=r(\cos\theta+i\sin\theta)$. Multiplying by $e^{i\alpha}$ rotates $z$ by the angle $\alpha$ without changing its modulus, since $re^{i\theta}e^{i\alpha}=re^{i(\theta+\alpha)}$. A factor $e^{i\alpha}$ is called a \textbf{phase}.

A wave that oscillates in time, such as $e^{-i\varepsilon x_4}$ with a real $\varepsilon$, has modulus 1 at every time. If $\varepsilon$ is imaginary, $\varepsilon=i\gamma$ with real $\gamma>0$, the same expression is $e^{\gamma x_4}$, which grows without bound; this is how "imaginary frequencies" signal an instability in Chapter 8.

\textbf{Complex vectors and matrices.} Vectors and matrices may have complex entries; all rules of Sections 1.2 to 1.4 hold unchanged. The conjugate $A^\ast$ of a matrix conjugates every entry, and the \textbf{Hermitian conjugate} (or adjoint) is

\[
A^\dagger=(A^\ast)^T,\qquad (A^\dagger)_{ij}=(A_{ji})^\ast .
\]

For a column vector $v$, $v^\dagger$ is the row of conjugated components. Like the transpose, the Hermitian conjugate reverses products, $(AB)^\dagger=B^\dagger A^\dagger$: indeed $(AB)^\ast=A^\ast B^\ast$ entry by entry (by the product rule of conjugation), and the transpose reverses the order. The number

\[
v^\dagger w=\sum_kv_k^\ast w_k
\]

is the \textbf{scalar product} of two complex vectors, and $v^\dagger v=\sum_k|v_k|^2$ is real, not negative, and zero only for $v=0$. It is the squared length of $v$. Two vectors are \textbf{orthogonal} if their scalar product is zero, $v^\dagger w=0$; the order does not matter, because $(v^\dagger w)^\ast=\sum_kv_kw_k^\ast=w^\dagger v$, so $v^\dagger w=0$ exactly when $w^\dagger v=0$. A vector has \textbf{length 1} if $v^\dagger v=1$. For real vectors $v^\dagger w=\sum_kv_kw_k$, which is the dot product of Section 1.11.

A square matrix is \textbf{Hermitian} if $A^\dagger=A$, \textbf{anti-Hermitian} if $A^\dagger=-A$, and \textbf{unitary} if $U^\dagger U=I$. A unitary matrix preserves scalar products: $(Uv)^\dagger(Uw)=v^\dagger U^\dagger Uw=v^\dagger w$. Four facts are used repeatedly:

\begin{enumerate}
\item A real symmetric matrix is Hermitian, and a real antisymmetric matrix is anti-Hermitian. Proof: for a real matrix $A^\ast=A$, so $A^\dagger=A^T$.
\item If $A$ is anti-Hermitian, then $iA$ is Hermitian: $(iA)^\dagger=i^\ast A^\dagger=(-i)(-A)=iA$.
\item For a Hermitian $H$ and every vector $v$ the number $v^\dagger Hv$ is real. Proof: a $1\times1$ matrix equals its transpose, so $(v^\dagger Hv)^\ast=(v^\dagger Hv)^\dagger=v^\dagger H^\dagger(v^\dagger)^\dagger=v^\dagger Hv$.
\item For an anti-Hermitian $A$ the number $v^\dagger Av$ is purely imaginary, by the same computation with $A^\dagger=-A$.
\end{enumerate}

\textbf{Example.} $A=\begin{pmatrix}1&i\\ -i&2\end{pmatrix}$ is Hermitian: $A^T=\begin{pmatrix}1&-i\\ i&2\end{pmatrix}$, and conjugating every entry gives back $A$. The real antisymmetric matrix $\begin{pmatrix}0&-1\\ 1&0\end{pmatrix}$ is anti-Hermitian, and $i$ times it, $\begin{pmatrix}0&-i\\ i&0\end{pmatrix}$, is Hermitian. In Chapter 6 fact 1 is the reason why the kinetic term of the dirac16complex Lagrangian needs no factor $i$: the matrices $C\gamma^a$ are real and antisymmetric, hence anti-Hermitian.

\subsection{Eigenvalues and eigenvectors}

\textbf{Definition.} A number $\lambda$ is an \textbf{eigenvalue} of a square matrix $A$, with \textbf{eigenvector} $v\ne0$, if

\[
Av=\lambda v .
\]

The matrix acts on its eigenvectors like multiplication by a number. Since $Av=\lambda v$ is the same as $(\lambda I-A)v=0$, and such a $v\ne0$ exists exactly when $\det(\lambda I-A)=0$ (Section 1.4), the eigenvalues are the roots of the \textbf{characteristic polynomial}

\[
p(x)=\det(xI-A).
\]

For a $2\times2$ matrix $A=\begin{pmatrix}a&b\\ c&d\end{pmatrix}$,

\[
p(x)=(x-a)(x-d)-bc=x^2-(a+d)x+(ad-bc)=x^2-(\mathrm{tr}A)\,x+\det A .
\]

If $\lambda_0$ and $\lambda_1$ are the two roots, then $p(x)=(x-\lambda_0)(x-\lambda_1)=x^2-(\lambda_0+\lambda_1)x+\lambda_0\lambda_1$, and comparing coefficients gives $\mathrm{tr}A=\lambda_0+\lambda_1$ and $\det A=\lambda_0\lambda_1$: the trace is the sum and the determinant the product of the eigenvalues. (The same holds for every $n$, counting each eigenvalue as often as it occurs as a root; we use this without proof.)

\textbf{Examples.} For $X=\begin{pmatrix}0&1\\ 1&0\end{pmatrix}$, $p(x)=x^2-1=(x-1)(x+1)$: the eigenvalues are $+1$, with eigenvector $(1,1)^T$, and $-1$, with eigenvector $(1,-1)^T$; indeed $X(1,1)^T=(1,1)^T$ and $X(1,-1)^T=(-1,1)^T=-(1,-1)^T$. For $Y=\begin{pmatrix}0&-1\\ 1&0\end{pmatrix}$, $p(x)=x^2+1$, whose roots are $+i$ and $-i$: a real matrix can have complex eigenvalues, which is one reason why complex numbers are needed even when all matrices are real.

\textbf{Multiplicity.} A root may occur several times. The notation $p(x)=(x-1)^8(x+1)^8$ means that $+1$ and $-1$ each occur eight times; this is the characteristic polynomial of the $16\times16$ matrix $C$ of Chapter 2 (Stage-1 document, Result 3.3). For the Hermitian matrices of this book the multiplicity of an eigenvalue equals the number of linearly independent eigenvectors that belong to it (a standard theorem, used without proof).

\textbf{Squares equal to $\pm I$.} If $A^2=I$, every eigenvalue is $+1$ or $-1$: from $Av=\lambda v$ follows $v=A^2v=\lambda Av=\lambda^2v$, so $\lambda^2=1$. If $A^2=-I$, the same argument gives $\lambda^2=-1$, so $\lambda=\pm i$. The gamma matrices of Chapter 2 square to $+I$ or $-I$, so their eigenvalues are known at once.

\textbf{Projectors.} A matrix with $P^2=P$ is a \textbf{projector}. If $A^2=I$, the two matrices

\[
P_+=\tfrac12(I+A),\qquad P_-=\tfrac12(I-A)
\]

are projectors with $P_++P_-=I$ and $P_+P_-=P_-P_+=0$, and $AP_\pm=\pm P_\pm$. Proof: $P_+^2=\tfrac14(I+2A+A^2)=\tfrac14(2I+2A)=P_+$, similarly $P_-^2=P_-$; $P_+P_-=\tfrac14(I-A^2)=0$; and $AP_\pm=\tfrac12(A\pm A^2)=\tfrac12(A\pm I)=\pm P_\pm$. So $P_+$ keeps the part of a vector with eigenvalue $+1$ and removes the part with eigenvalue $-1$, and $P_-$ does the opposite. Every vector splits as $v=P_+v+P_-v$. The chirality projectors $P_\mp=\tfrac12(1\mp\gamma^8)$ of Chapter 2 are of this kind.

\textbf{Hermitian matrices have real eigenvalues.} If $H^\dagger=H$ and $Hv=\lambda v$ with $v\ne0$, then $v^\dagger Hv=\lambda\,v^\dagger v$. The left side is real (fact 3 of Section 1.5) and $v^\dagger v>0$, so $\lambda$ is real. Moreover, eigenvectors of a Hermitian matrix that belong to different eigenvalues are orthogonal: if $Hv=\lambda v$ and $Hw=\mu w$ with $\lambda\ne\mu$, then $w^\dagger Hv=\lambda\,w^\dagger v$, and also $w^\dagger Hv=(Hw)^\dagger v=\mu\,w^\dagger v$ (using $H^\dagger=H$ and $\mu^\ast=\mu$), so $(\lambda-\mu)\,w^\dagger v=0$ and $w^\dagger v=0$. A stronger standard theorem, the \textbf{spectral theorem}, says that every Hermitian $n\times n$ matrix has $n$ eigenvectors that are orthogonal to each other and of length 1; it is used without proof.

\subsection{Commutators and anticommutators}

Because matrices need not commute, two combinations of a product and its reverse are given names. The \textbf{commutator} and the \textbf{anticommutator} of two square matrices are

\[
[A,B]=AB-BA,\qquad \{A,B\}=AB+BA .
\]

$A$ and $B$ \textbf{commute} if $[A,B]=0$ and \textbf{anticommute} if $\{A,B\}=0$, that is $AB=-BA$. Directly from the definitions: $[B,A]=-[A,B]$, $\{B,A\}=\{A,B\}$, $\{A,A\}=2A^2$, $[A,A]=0$, and every product splits into the two parts,

\[
AB=\tfrac12[A,B]+\tfrac12\{A,B\}.
\]

The commutator obeys a product rule (the \textbf{Leibniz rule}), which is proved by writing out both sides:

\[
[A,BC]=[A,B]\,C+B\,[A,C],\qquad\text{since}\qquad ABC-BCA=(AB-BA)C+B(AC-CA).
\]

\textbf{A worked example that anticipates Chapter 2.} Take the three real $2\times2$ matrices

\[
X=\begin{pmatrix}0&1\\ 1&0\end{pmatrix},\qquad Z=\begin{pmatrix}1&0\\ 0&-1\end{pmatrix},\qquad Y=\begin{pmatrix}0&-1\\ 1&0\end{pmatrix}.
\]

Multiplying out,

\[
XZ=\begin{pmatrix}0&-1\\ 1&0\end{pmatrix}=Y,\qquad ZX=\begin{pmatrix}0&1\\ -1&0\end{pmatrix}=-Y,\qquad X^2=Z^2=I,\qquad Y^2=-I .
\]

Hence $\{X,Z\}=0$ and $[X,Z]=2Y$. Further, $XY=X(XZ)=X^2Z=Z$ and $YX=XZX=X(ZX)=-X(XZ)=-Z$, so $\{X,Y\}=0$; and $ZY=Z(XZ)=(ZX)Z=-XZ^2=-X$ and $YZ=XZ^2=X$, so $\{Z,Y\}=0$.

A word on notation before the relations are collected (Section 1.8 explains it in full). We now number the three matrices with a raised index. On the letters $\Gamma$ and $\gamma$ a raised index is a label that numbers the matrices of a list, not a power: $\Gamma^2$ below is the third matrix of the list, not $\Gamma\cdot\Gamma$, and its square is written $(\Gamma^2)^2=Y^2=-I$. Likewise the two raised indices of $\hat\eta^{ab}$ and $\eta^{ab}$ are labels: $\hat\eta^{ab}$ is the entry in row $a$ and column $b$ of $\hat\eta$ (Section 1.11 explains why such entries are written with raised indices; for the diagonal matrices with entries $\pm1$ used here it makes no difference). A raised number on a matrix that carries no label keeps its usual meaning, a power, as in $X^2=I$ above. With the names $\Gamma^0=X$, $\Gamma^1=Z$, $\Gamma^2=Y$ and the diagonal matrix $\hat\eta=\mathrm{diag}(+1,+1,-1)$, all nine relations fit into one formula:

\[
\{\Gamma^a,\Gamma^b\}=2\hat\eta^{ab}I\qquad(a,b=0,1,2).
\]

Matrices obeying such a relation generate a \textbf{Clifford algebra}; Chapter 2 builds the eight $16\times16$ matrices $\gamma^0,\dots,\gamma^7$ with $\{\gamma^a,\gamma^b\}=2\eta^{ab}I_{16}$ for the $\eta$ of this book. Notice also that $X$ and $Z$, which square to $+I$, are symmetric, while $Y$, which squares to $-I$, is antisymmetric. Part of this is forced: a real symmetric matrix is Hermitian, so its eigenvalues are real (Section 1.6), while a matrix with square $-I$ has only the eigenvalues $\pm i$; since every square matrix has at least one complex eigenvalue (the fundamental theorem of algebra, used without proof), a real symmetric matrix can never square to $-I$. The square alone does not force the rest of the pattern: $\begin{pmatrix}1&-2\\ 1&-1\end{pmatrix}$ has square $-I$ and is not antisymmetric, and $\begin{pmatrix}1&1\\ 0&-1\end{pmatrix}$ has square $+I$ and is not symmetric (multiply out to check). The full pattern holds for the real matrices with $M^TM=I$, called \textbf{orthogonal} matrices. $X$, $Z$ and $Y$ are orthogonal: $X^TX=X^2=I$, $Z^TZ=Z^2=I$ and $Y^TY=(-Y)Y=-Y^2=I$. For an orthogonal $M$, $M^T=M^{-1}$; and $M^2=\pm I$ says that $M(\pm M)=I$, that is $M^{-1}=\pm M$. Together, $M^T=\pm M$: an orthogonal matrix with square $+I$ is symmetric, and one with square $-I$ is antisymmetric. The notebook's gamma matrices are of this kind: each is a signed permutation matrix (every row and every column contains exactly one nonzero entry, $+1$ or $-1$), hence orthogonal (Section 2.2 proves this), which is why they are symmetric for the space-like directions 0 to 3 and antisymmetric for the time-like directions 4 to 7 (Chapter 2, Section 2.8, facts (N2) and (N3)).

\subsection{Indices and the summation convention}

\textbf{Free and dummy indices.} In the formula $(Av)_i=\sum_jA_{ij}v_j$ the index $i$ is \textbf{free}: the formula holds for each value of $i$, and it appears once in every term. The index $j$ is a \textbf{dummy}: it is summed over, and it can be renamed, $\sum_jA_{ij}v_j=\sum_kA_{ik}v_k$, as long as the new name is not already in use in the term. In a correct equation every term has the same free indices. An index formula is a compact way of writing many equations at once; for example $\{\Gamma^a,\Gamma^b\}=2\hat\eta^{ab}I$ of Section 1.7 stands for nine matrix equations, one for each pair $(a,b)$.

\textbf{Upper and lower indices.} In Sections 1.2 to 1.7 the components of plain columns of numbers were written with lower indices, $v_k$. For the vectors of the eight-dimensional space two kinds of objects occur, and they are distinguished by the position of their index. The components of a displacement or a velocity carry an upper index, $v^\mu$; the coefficients of a metric, $g_{\mu\nu}$ (Section 1.11), and the derivatives $\partial_\mu$ (Section 1.9) carry lower indices. An upper index is not an exponent: $v^2$ means component 2 of $v$, and the square of that component is written $(v^2)^2$.

\textbf{The summation convention.} In a single term, an index that appears twice, once as an upper and once as a lower index, is summed over its whole range, without writing $\sum$:

\[
v^\mu w_\mu=\sum_{\mu=0}^{7}v^\mu w_\mu,\qquad g_{\mu\nu}v^\mu w^\nu=\sum_{\mu=0}^7\sum_{\nu=0}^7g_{\mu\nu}v^\mu w^\nu,\qquad \partial_\mu V^\mu=\sum_{\mu=0}^7\frac{\partial V^\mu}{\partial x^\mu}.
\]

This is the \textbf{Einstein summation convention}. In this book the indices $\mu,\nu,\rho,\sigma,\lambda$ of the coordinates (coordinate indices, also called curved indices; Chapter 4) and the indices $a,b,c$ that are raised and lowered with the flat metric $\eta$ of Section 1.11 (frame indices, Section 4.10) run over $0,\dots,7$. An index that appears twice in the same position (both lower or both upper) is not summed automatically; where such a sum is meant, as for matrix entries $A_{ij}$, the symbol $\sum$ is written.

\textbf{Coordinates and their names.} In the index formulas of this chapter and of Chapter 4 the coordinates are written $x^\mu$ with an upper index. Other chapters (Chapter 9, which follows the notebook closely, and the line elements of later chapters, such as $dx_1^2$ in Chapters 7, 11 and 13), the notebook and the project documents write the same coordinates as $x_0,x_1,\dots,x_7$, also inside formulas such as $\partial_\mu=\partial/\partial x_\mu$; this chapter does the same in words and in examples such as $e^{-i\varepsilon x_4}$. There the lower index is only part of the name, and $\partial/\partial x_\mu$ is the same derivative as $\partial/\partial x^\mu$. The two notations denote the same numbers: $x^4$ in a formula is the coordinate named $x_4$, the time. The index of a coordinate is never lowered with a metric in this book.

\textbf{The Kronecker delta.} With one upper and one lower index, $\delta^\mu{}_\nu$ is 1 for $\mu=\nu$ and 0 otherwise. Contracting with it renames an index:

\[
\delta^\mu{}_\nu v^\nu=v^\mu,\qquad \delta^\mu{}_\nu\delta^\nu{}_\rho=\delta^\mu{}_\rho,\qquad \delta^\mu{}_\mu=8 .
\]

Proof of the first: in $\sum_\nu\delta^\mu{}_\nu v^\nu$ only the term with $\nu=\mu$ is not zero, and it equals $v^\mu$. The last formula counts the eight values of $\mu$.

\textbf{Symmetric times antisymmetric gives zero.} If $S_{\mu\nu}=S_{\nu\mu}$ and $A^{\mu\nu}=-A^{\nu\mu}$, then

\[
S_{\mu\nu}A^{\mu\nu}=0 .
\]

Proof: renaming the two dummy indices ($\mu$ becomes $\nu$ and $\nu$ becomes $\mu$) does not change the value, so $S_{\mu\nu}A^{\mu\nu}=S_{\nu\mu}A^{\nu\mu}$; using the two symmetries, the right side equals $S_{\mu\nu}(-A^{\mu\nu})$. A number equal to its own negative is 0. Example with $2\times2$ arrays: for $S=\begin{pmatrix}1&2\\ 2&3\end{pmatrix}$ and $A=\begin{pmatrix}0&1\\ -1&0\end{pmatrix}$ the full double sum is $1\cdot0+2\cdot1+2\cdot(-1)+3\cdot0=0$.

\textbf{Consequence for quadratic expressions.} For a vector $v$ with ordinary (commuting) components and a square matrix $M$, the quadratic expression $v^TMv=\sum_j\sum_kv_jM_{jk}v_k$ contains the products $v_jv_k$, which are symmetric in $j$ and $k$. By the rule just proved, the antisymmetric part of $M$ drops out:

\[
v^TMv=v^TM_Sv,\qquad\text{in particular}\qquad v^TAv=0\ \text{for antisymmetric }A .
\]

Chapter 5 meets the opposite case: for anticommuting quantities $\theta_j\theta_k=-\theta_k\theta_j$ the products are antisymmetric, and then the symmetric part drops out, $\theta^TM\theta=\theta^TM_A\theta$. This one difference decides whether the notebook's Lagrangian is trivial (Stage-1 document, §6).

\textbf{Sums over ordered pairs.} If $A_{ab}$ and $B^{ab}$ are both antisymmetric, the product $A_{ab}B^{ab}$ (no sum) is symmetric under exchanging $a$ and $b$, and it vanishes for $a=b$. Hence the sum over all ordered pairs is twice the sum over the pairs with $a<b$:

\[
\tfrac12A_{ab}B^{ab}=\sum_{a<b}A_{ab}B^{ab}.
\]

This is why the spinor connection $\Omega_\mu$ of Chapter 4 can be written either as $\tfrac12\omega_{\mu ab}S^{ab}$ with the summation convention or as $\sum_{a<b}\omega_{\mu ab}S^{ab}$ without the factor $\tfrac12$ (there $\omega_{\mu ab}$, the spin connection, and $S^{ab}$ are both antisymmetric in $a$ and $b$).

\textbf{The Levi-Civita symbol.} For $n$ indices, each taking the values $0,\dots,n-1$, the symbol $\varepsilon_{j_0j_1\dots j_{n-1}}$ is the sign of the permutation $(j_0,\dots,j_{n-1})$ when the indices are all different, and 0 when two of them are equal. For three indices: $\varepsilon_{012}=\varepsilon_{120}=\varepsilon_{201}=+1$, $\varepsilon_{102}=\varepsilon_{021}=\varepsilon_{210}=-1$, and for example $\varepsilon_{011}=0$. With it the determinant formula of Section 1.4 reads $\det A=\sum_{j_0,j_1,j_2}\varepsilon_{j_0j_1j_2}A_{0j_0}A_{1j_1}A_{2j_2}$ for $n=3$ (and similarly for every $n$). The notebook uses a four-index symbol with the 1-based indices $1,\dots,4$ of Mathematica to build its gamma matrices (Chapter 3); this is one of the labelled exceptions to counting from 0.

\subsection{Functions of several variables and partial derivatives}

\textbf{Partial derivatives.} For a function $f(x^0,\dots,x^7)$ of eight variables, the \textbf{partial derivative} with respect to $x^\mu$ is the ordinary derivative with respect to $x^\mu$ while all the other variables are held fixed. It is written

\[
\partial_\mu f=\frac{\partial f}{\partial x^\mu}.
\]

Example: for $f(x^0,x^4)=(x^0)^2\sin x^4$,

\[
\partial_0f=2x^0\sin x^4,\qquad \partial_4f=(x^0)^2\cos x^4,\qquad \partial_4\partial_0f=2x^0\cos x^4=\partial_0\partial_4f .
\]

The last equality is an instance of a general theorem (Schwarz's theorem): for a function whose second partial derivatives are continuous, the order of two partial derivatives does not matter, $\partial_\mu\partial_\nu f=\partial_\nu\partial_\mu f$. We use it without proof. All functions in this book are \textbf{smooth}, that is, they have as many continuous derivatives as are needed, unless the text says otherwise.

\textbf{Product rule.} Because a partial derivative is a derivative in one variable, the rules of one-variable calculus hold for it; in particular $\partial_\mu(fg)=(\partial_\mu f)g+f\,\partial_\mu g$. For matrices whose entries are functions, derivatives are taken entry by entry, and the product rule keeps the order of the factors:

\[
\partial_\mu(AB)=(\partial_\mu A)B+A\,\partial_\mu B .
\]

Proof: entry $(i,k)$ of the left side is $\partial_\mu\sum_jA_{ij}B_{jk}=\sum_j\bigl((\partial_\mu A_{ij})B_{jk}+A_{ij}\partial_\mu B_{jk}\bigr)$, which is entry $(i,k)$ of the right side. Applying it to $AA^{-1}=I$, whose derivative is zero, gives $(\partial_\mu A)A^{-1}+A\,\partial_\mu(A^{-1})=0$, and multiplying from the left by $A^{-1}$:

\[
\partial_\mu(A^{-1})=-A^{-1}(\partial_\mu A)A^{-1}.
\]

For a $1\times1$ matrix this is the familiar $(1/a)'=-a'/a^2$.

\textbf{Chain rule.} Let $f(u,v)$ depend on two variables that in turn depend on $t$. Then

\[
\frac{d}{dt}f(u(t),v(t))=\frac{\partial f}{\partial u}\,\frac{du}{dt}+\frac{\partial f}{\partial v}\,\frac{dv}{dt}.
\]

Derivation. Write $\Delta u=u(t+\Delta t)-u(t)$, $\Delta v=v(t+\Delta t)-v(t)$ and split the change of $f$ into two steps, first in $u$ and then in $v$:

\[
f(u+\Delta u,v+\Delta v)-f(u,v)=\bigl[f(u+\Delta u,v+\Delta v)-f(u,v+\Delta v)\bigr]+\bigl[f(u,v+\Delta v)-f(u,v)\bigr].
\]

By the mean value theorem of one-variable calculus (for a differentiable function, $g(b)-g(a)=g'(c)(b-a)$ for some $c$ between $a$ and $b$; used without proof), the first bracket equals $\partial_uf(u+\theta_1\Delta u,v+\Delta v)\,\Delta u$ and the second $\partial_vf(u,v+\theta_2\Delta v)\,\Delta v$, with numbers $\theta_1,\theta_2$ between 0 and 1. Divide by $\Delta t$ and let $\Delta t\to0$: then $\Delta u/\Delta t\to du/dt$, $\Delta v/\Delta t\to dv/dt$, and, because the partial derivatives are continuous and $\Delta u,\Delta v\to0$, the two partial derivatives tend to their values at $(u,v)$. This gives the formula. With more variables the same argument gives one term per variable: $\frac{d}{dt}f(x(t))=\partial_\mu f\,\frac{dx^\mu}{dt}$ (summation convention).

\textbf{Change of variables.} If new coordinates are functions of old ones, the chain rule converts derivatives. The project uses, for the primordial field, $z=6Hx^0$ and $t=Hx^4$ with a constant $H>0$. A function of $z$ then depends on $x^0$ through $z$, and

\[
\partial_0=\frac{dz}{dx^0}\,\frac{\partial}{\partial z}=6H\,\partial_z,\qquad \partial_4=\frac{dt}{dx^4}\,\frac{\partial}{\partial t}=H\,\partial_t .
\]

For example $\partial_0\cot z=6H\,\partial_z\cot z=-6H/\sin^2z$.

\textbf{Linear approximation.} For a small displacement $h$, a smooth function changes by

\[
f(x+h)=f(x)+\partial_\mu f(x)\,h^\mu+(\text{terms of second and higher order in }h).
\]

The linear term is the chain rule applied to the function $t\mapsto f(x+th)$ at $t=0$, whose derivative there is $\partial_\mu f(x)\,h^\mu$; that the remainder is of second order in $h$ for a smooth $f$ is Taylor's theorem, used without proof. "First order" in this book means: keep the term linear in the small quantity and drop the rest.

\textbf{The derivative of a determinant.} For an invertible matrix $g$ whose entries depend on $x$,

\[
\partial_\mu\det g=\det g\;\mathrm{tr}\bigl(g^{-1}\partial_\mu g\bigr).
\]

This is \textbf{Jacobi's formula}. We prove it in the two cases that the book needs most and use it without proof otherwise. For a diagonal matrix $g=\mathrm{diag}(g_0,\dots,g_{n-1})$, $\det g=g_0\cdots g_{n-1}$ and the product rule give $\partial_\mu\det g=\sum_k g_0\cdots(\partial_\mu g_k)\cdots g_{n-1}=\det g\sum_k\partial_\mu g_k/g_k$, and $\sum_k\partial_\mu g_k/g_k$ is the trace of $g^{-1}\partial_\mu g=\mathrm{diag}(\partial_\mu g_k/g_k)$. For a $2\times2$ matrix with entries $a,b,c,d$, $\partial(ad-bc)=d\,\partial a+a\,\partial d-c\,\partial b-b\,\partial c$, while

\[
\det g\;\mathrm{tr}(g^{-1}\partial g)=\mathrm{tr}\left[\begin{pmatrix}d&-b\\ -c&a\end{pmatrix}\begin{pmatrix}\partial a&\partial b\\ \partial c&\partial d\end{pmatrix}\right]=d\,\partial a-b\,\partial c-c\,\partial b+a\,\partial d ,
\]

the same expression.

\textbf{The derivative of the volume factor.} The determinant $\det g$ is a sum of products of entries of $g$ (Section 1.4), so it is continuous when the entries are. Near a point where $\det g\ne0$ it therefore has a fixed sign, because a continuous function that is not zero at a point keeps its sign near that point (a standard fact, used without proof). Write $\det g=\epsilon\,|\det g|$ with this fixed sign $\epsilon=\pm1$; then also $|\det g|=\epsilon\det g$, since $\epsilon^2=1$, and $\partial_\mu|\det g|=\epsilon\,\partial_\mu\det g$, since $\epsilon$ is constant near the point. By the chain rule for $\sqrt u$, whose derivative is $1/(2\sqrt u)$, and by Jacobi's formula,

\[
\begin{aligned}
\partial_\mu\sqrt{|\det g|}&=\frac{\partial_\mu|\det g|}{2\sqrt{|\det g|}}=\frac{\epsilon\,\partial_\mu\det g}{2\sqrt{|\det g|}}\\
&=\frac{\epsilon\det g\;\mathrm{tr}\bigl(g^{-1}\partial_\mu g\bigr)}{2\sqrt{|\det g|}}=\frac{|\det g|\;\mathrm{tr}\bigl(g^{-1}\partial_\mu g\bigr)}{2\sqrt{|\det g|}},
\end{aligned}
\]

and $|\det g|/\sqrt{|\det g|}=\sqrt{|\det g|}$ gives

\[
\partial_\mu\sqrt{|\det g|}=\tfrac12\sqrt{|\det g|}\;\mathrm{tr}\bigl(g^{-1}\partial_\mu g\bigr).
\]

This formula is used in Chapter 4 to relate the volume factor to the Christoffel symbols.

\subsection{Integrals in several variables, integration by parts and the delta function}

\textbf{One variable.} The definite integral $\int_a^bf(x)\,dx$ and the fundamental theorem of calculus, $\int_a^bF'(x)\,dx=F(b)-F(a)$, are school mathematics and are used without proof.

\textbf{Integration by parts.} Integrating the product rule $(fg)'=f'g+fg'$ from $a$ to $b$ gives $f(b)g(b)-f(a)g(a)=\int_a^bf'g\,dx+\int_a^bfg'\,dx$, that is

\[
\int_a^bf\,g'\,dx=\bigl[fg\bigr]_a^b-\int_a^bf'\,g\,dx,\qquad \bigl[fg\bigr]_a^b=f(b)g(b)-f(a)g(a).
\]

Example: with $f=x$ and $g'=e^x$ (so $g=e^x$), $\int_0^1xe^x\,dx=\bigl[xe^x\bigr]_0^1-\int_0^1e^x\,dx=e-(e-1)=1$. When the product $fg$ vanishes at both ends, the boundary term drops out, and the derivative moves from one factor to the other at the price of a minus sign. This is the key step in the derivation of the Euler--Lagrange equations (Chapter 5).

\textbf{Several variables.} An integral over a box $a_\mu\le x^\mu\le b_\mu$ is computed as an iterated integral, one variable after the other. For a continuous function on a box the order of the variables does not matter (Fubini's theorem, used without proof). The volume element of the eight coordinates is written $d^8x=dx^0dx^1\cdots dx^7$.

\textbf{Integrals of derivatives vanish when the boundary values vanish.} Let $V^\mu(x)$, $\mu=0,\dots,n-1$, be smooth functions on a box. Then

\[
\int_{\mathrm{box}}\partial_\mu V^\mu\,d^nx=\sum_{\mu}\int\Bigl[V^\mu\Bigr]_{x^\mu=a_\mu}^{x^\mu=b_\mu}\,\prod_{\nu\ne\mu}dx^\nu ,
\]

and in particular the integral is zero if every $V^\mu$ vanishes on the boundary of the box. Proof: the left side is a sum of $n$ integrals, one for each $\mu$. In the integral of $\partial_\mu V^\mu$ (no sum), do the integral over $x^\mu$ first, with the other variables fixed; by the fundamental theorem it equals the difference of the values of $V^\mu$ at $x^\mu=b_\mu$ and $x^\mu=a_\mu$. This is the simplest form of the divergence theorem (Gauss's theorem). It is the reason why adding a total derivative $\partial_\mu V^\mu$ to a Lagrangian does not change the field equations (Chapter 5). Example on the unit square: for $V^0=\sin(\pi x^0)\,x^1$ and $V^1=0$, $\int_0^1\!\int_0^1\partial_0V^0\,dx^0dx^1=\int_0^1x^1\,[\sin\pi-\sin0]\,dx^1=0$.

\textbf{The delta function.} In Chapter 8 the equal-time anticommutators contain the symbol $\delta(x-y)$, Dirac's delta function. It is not an ordinary function but a rule for integrals: for every continuous $f$,

\[
\int f(x)\,\delta(x-y)\,dx=f(y).
\]

One can picture it as the limit of ever narrower and taller rectangles of area 1. Let $r_\epsilon(u)=1/\epsilon$ for $|u|<\epsilon/2$ and $0$ otherwise. Then $\int f(x)\,r_\epsilon(x-y)\,dx=\frac1\epsilon\int_{y-\epsilon/2}^{y+\epsilon/2}f(x)\,dx$ is the average of $f$ over an interval of length $\epsilon$ around $y$. A continuous function takes its average value somewhere in the interval (the mean value theorem for integrals, used without proof), so the average tends to $f(y)$ as $\epsilon\to0$. In several variables the delta function is the product of one-dimensional ones; for example $\delta^7(x-y)$ in Chapter 8 is the product of seven deltas, one for each coordinate other than the time $x^4$.

\subsection{Metrics and signatures}

\textbf{The dot product.} In ordinary 3-space the dot product of two vectors is $u\cdot v=u_0v_0+u_1v_1+u_2v_2=u^TIv$, and the length of $v$ is $\sqrt{v\cdot v}$. A \textbf{metric} generalizes it: a symmetric, invertible $n\times n$ matrix $g$ defines the scalar product

\[
g(u,v)=u^Tg\,v=g_{\mu\nu}u^\mu v^\nu ,
\]

which is symmetric in $u$ and $v$ because $g$ is symmetric. The number $g(v,v)$ is the squared length of $v$. For $g=I$ it is the dot product. But a metric may have negative diagonal entries, and then $g(v,v)$ can be negative or zero for a vector $v\ne0$.

\textbf{Space-like, time-like and null.} In the convention of this book a vector is \textbf{space-like} if $g(v,v)>0$, \textbf{time-like} if $g(v,v)<0$, and \textbf{null} (or light-like) if $g(v,v)=0$ and $v\ne0$. The special relativity of the observed world has one time-like and three space-like directions; with time first and this sign convention its metric is $\mathrm{diag}(-1,+1,+1,+1)$. Many books of particle physics, and the scalar-field reference of the project, use the opposite overall sign, $\mathrm{diag}(+1,-1,-1,-1)$; the two conventions describe the same geometry, since one metric is the negative of the other (the Stage-1 document, §2.3, compares them for a scalar field). The flat metric of this book is

\[
\begin{aligned}
&\eta=\mathrm{diag}(+1,+1,+1,+1,-1,-1,-1,-1),\\
&\eta(v,v)=(v^0)^2+(v^1)^2+(v^2)^2+(v^3)^2-(v^4)^2-(v^5)^2-(v^6)^2-(v^7)^2 .
\end{aligned}
\]

The basis vector $e_0$ is space-like ($\eta(e_0,e_0)=\eta_{00}=+1$), $e_4$ is time-like ($\eta_{44}=-1$), and $e_0+e_4$ is null: $\eta(e_0+e_4,e_0+e_4)=1-1=0$.

\textbf{Signature.} The \textbf{signature} of a metric is the pair $(p,q)$ of the numbers of its positive and its negative eigenvalues. (A metric is real and symmetric, so its eigenvalues are real, and it is invertible, so none is zero.) For a diagonal metric the eigenvalues are the diagonal entries, so $\eta$ has signature (4,4). If the basis is changed, so that $g$ becomes $g'=\Lambda^Tg\Lambda$ with an invertible matrix $\Lambda$ whose columns are the new basis vectors, the eigenvalues change but their signs do not: $g'$ has the same signature. This is \textbf{Sylvester's law of inertia}, used without proof. The signature is therefore a property of the geometry, not of the basis.

\textbf{Example.} $g=\begin{pmatrix}1&2\\ 2&1\end{pmatrix}$ has $p(x)=(x-1)^2-4=(x-3)(x+1)$, eigenvalues 3 and $-1$, signature (1,1). In the basis $u=(1,1)^T$, $w=(1,-1)^T$ one finds $g(u,u)=1+2+2+1=6$, $g(w,w)=1-2-2+1=-2$ and $g(u,w)=(1,1)\,(1-2,\,2-1)^T=0$, so in this basis the metric is $\mathrm{diag}(6,-2)$: again one positive and one negative entry, as Sylvester's law says.

\textbf{Inverse metric, raising and lowering.} The inverse matrix of $g_{\mu\nu}$ is written with upper indices, $g^{\mu\nu}$, and satisfies $g^{\mu\nu}g_{\nu\rho}=\delta^\mu{}_\rho$. For $\eta$, whose diagonal entries are $\pm1$, $\eta^2=I$, so the inverse has the same entries: $\eta^{ab}=\eta_{ab}$. A metric \textbf{lowers} an index and its inverse \textbf{raises} it:

\[
v_\mu=g_{\mu\nu}v^\nu,\qquad v^\mu=g^{\mu\nu}v_\nu,\qquad g(u,v)=u^\mu v_\mu .
\]

With $\eta$, lowering keeps the components 0 to 3 and changes the sign of the components 4 to 7: $v_0=v^0$, and $v_4=\eta_{44}v^4=-v^4$.

\textbf{Curved metrics and the line element.} In curved space (Chapter 4) the metric depends on the point, $g_{\mu\nu}(x)$, and it measures small displacements $dx^\mu$ at $x$. The \textbf{line element}

\[
ds^2=g_{\mu\nu}(x)\,dx^\mu dx^\nu
\]

is the squared length of the displacement from $x$ to $x+dx$. The same geometry can be described in different coordinates. Example: the flat plane in polar coordinates, $x=r\cos\theta$ and $y=r\sin\theta$. The chain rule gives $dx=\cos\theta\,dr-r\sin\theta\,d\theta$ and $dy=\sin\theta\,dr+r\cos\theta\,d\theta$, hence

\[
\begin{aligned}
ds^2=dx^2+dy^2&=(\cos^2\theta+\sin^2\theta)\,dr^2+2(-r\cos\theta\sin\theta+r\sin\theta\cos\theta)\,dr\,d\theta\\
&\quad+r^2(\sin^2\theta+\cos^2\theta)\,d\theta^2\\
&=dr^2+r^2d\theta^2 .
\end{aligned}
\]

In the coordinates $(r,\theta)$ the metric is $\mathrm{diag}(1,r^2)$. In general, if the old coordinates $x^\mu$ are functions of new ones $x'^\alpha$, then $dx^\mu=\frac{\partial x^\mu}{\partial x'^\alpha}dx'^\alpha$, and inserting this into the line element gives the \textbf{transformation law of the metric}

\[
g'_{\alpha\beta}=\frac{\partial x^\mu}{\partial x'^\alpha}\,\frac{\partial x^\nu}{\partial x'^\beta}\,g_{\mu\nu},\qquad\text{in matrix form}\qquad g'=J^TgJ,\quad J^\mu{}_\alpha=\frac{\partial x^\mu}{\partial x'^\alpha}.
\]

\textbf{The volume element.} From $g'=J^TgJ$ and the product rule for determinants, $\det g'=(\det J)^2\det g$, so $\sqrt{|\det g'|}=|\det J|\sqrt{|\det g|}$. The change-of-variables formula for multiple integrals (used without proof) says that $d^nx=|\det J|\,d^nx'$. Together these show that $\sqrt{|\det g|}\,d^nx$ has the same value in every coordinate system: it is the \textbf{volume element} of the geometry, and it multiplies every Lagrangian of this book. The book writes $\sqrt{|g|}$ for $\sqrt{|\det g|}$. For polar coordinates $\sqrt{|g|}=r$, and the area of the disc of radius $R$ is $\int_0^R\int_0^{2\pi}r\,d\theta\,dr=2\pi\cdot\tfrac12R^2=\pi R^2$, as it must be.

\textbf{Worked example: the metric of the primordial field.} In Chapter 9 the notebook's metric is, in the coordinates $x^0,\dots,x^7$, with $z=6Hx^0\in(0,\pi/2)$, $s=\sin z$ and a real function $a_4$ of the time,

\[
g=\mathrm{diag}\bigl(\cot^2z,\ s^{1/3}e^{2a_4},\ s^{1/3}e^{2a_4},\ s^{1/3}e^{2a_4},\ -1,\ -s^{1/3}e^{-2a_4},\ -s^{1/3}e^{-2a_4},\ -s^{1/3}e^{-2a_4}\bigr).
\]

For $0<z<\pi/2$ the first four entries are positive and the last four negative, so the signature is (4,4). The determinant is the product of the diagonal (Section 1.4):

\[
\begin{aligned}
\det g&=\cot^2z\cdot\bigl(s^{1/3}e^{2a_4}\bigr)^3\cdot(-1)\cdot\bigl(-s^{1/3}e^{-2a_4}\bigr)^3\\
&=\cot^2z\cdot s\,e^{6a_4}\cdot(-1)\cdot(-1)\,s\,e^{-6a_4}=\cot^2z\,\sin^2z=\cos^2z ,
\end{aligned}
\]

using $(-x)^3=-x^3$, $(s^{1/3})^3=s$ and $e^{6a_4}e^{-6a_4}=1$. The sign is $+$ because there are four negative entries, an even number. Hence $\sqrt{|g|}=\cos z$, which does not depend on the time. The exact Stage-2 verifier confirms $\det g=+\cos^2z$ (check \texttt{P\_metric\allowbreak{}\_detG\_eq\allowbreak{}uals\_plu\allowbreak{}s\_cos2z} in \texttt{artifact\allowbreak{}s/dirac1\allowbreak{}6complex\allowbreak{}/primord\allowbreak{}ial-fiel\allowbreak{}d/wolfra\allowbreak{}m-primor\allowbreak{}dial-rep\allowbreak{}ort.json}); the Stage-2 specification had first written $-\cos^2z$, and the sign was corrected when it was computed. Jacobi's formula of Section 1.9 gives a check of the derivative. The diagonal entries contribute to $\mathrm{tr}(g^{-1}\partial_zg)=\sum_\mu\partial_z\ln|g_{\mu\mu}|$ the terms $\partial_z\ln\cot^2z=-2/(\sin z\cos z)$, three times $\partial_z\ln s^{1/3}=\tfrac13\cot z$ from 3-space, zero from the entry $-1$, and three times $\tfrac13\cot z$ from the extra times, in total

\[
\mathrm{tr}(g^{-1}\partial_zg)=-\frac2{\sin z\cos z}+2\,\frac{\cos z}{\sin z}=\frac{-2+2\cos^2z}{\sin z\cos z}=-2\tan z,
\]

so $\partial_z\sqrt{|g|}=\tfrac12\cos z\,(-2\tan z)=-\sin z$, which is indeed the derivative of $\cos z$.

\subsection{The (4,4) world of this book}

The space of this book has eight coordinates, four space-like and four time-like:

\begingroup
\small
\begin{longtable}{@{}>{\raggedright\arraybackslash}p{0.267\linewidth}>{\raggedright\arraybackslash}p{0.267\linewidth}>{\raggedright\arraybackslash}p{0.267\linewidth}@{}}
\toprule
\textbf{coordinate} & \textbf{role} & \textbf{$\eta$ entry} \\
\midrule
\endfirsthead
\toprule
\textbf{coordinate} & \textbf{role} & \textbf{$\eta$ entry} \\
\midrule
\endhead
$x_0$ & hidden space direction & $+1$ \\
$x_1,x_2,x_3$ & ordinary 3-space & $+1$ \\
$x_4$ & the time in which everything evolves & $-1$ \\
$x_5,x_6,x_7$ & the three extra times & $-1$ \\
\bottomrule
\end{longtable}
\endgroup

The roles are those of the notebook. Choosing $x_4$ as the time means: initial data are given at one value of $x_4$, and the field equations determine the field at later values of $x_4$ (Chapters 8 and 10).

The observed world has one time and three space directions, signature (3,1) in the sign convention of this book. The signature (4,4) is therefore not the signature of observed spacetime; it is the setting of the model, chosen by the notebook, and the book treats it as an assumption (ledger row L43 of Chapter 0). Two consequences are visible already. First, the Clifford algebra of an eight-dimensional space with this metric is the algebra of all real $16\times16$ matrices (Chapter 2), which is why the field has sixteen components. Second, three extra time-like directions change the relation between frequency and wave number: Chapter 8 derives that in flat space a wave of mass $m$ with momenta $k_0,\dots,k_3$ along the space-like directions and $k_5,k_6,k_7$ along the extra times has the squared frequency $E^2=m^2+k_0^2+k_1^2+k_2^2+k_3^2-k_5^2-k_6^2-k_7^2$, which becomes negative for large extra-time momenta. By Section 1.5 a negative $E^2$ means an imaginary frequency and a wave that grows exponentially in time; this instability is why the physics of the book is restricted to the sector without extra-time momenta.

\subsection{What we proved and what we assumed}

We \textbf{proved}, from the definitions: the uniqueness of the expansion coefficients in a basis; the rules for sums, products, transposes and inverses of matrices, including $(AB)^T=B^TA^T$, $\mathrm{tr}(AB)=\mathrm{tr}(BA)$ and $(AB)^{-1}=B^{-1}A^{-1}$; the unique split of a square matrix into symmetric and antisymmetric parts; the block product and the mixed-product rule of tensor products; the product rule of $2\times2$ determinants and the determinant of a diagonal matrix; the arithmetic of complex numbers, the product rule for complex-valued functions of a real variable, the law of exponents for $e^{i\theta}$ and its derivative; the four facts about Hermitian and anti-Hermitian matrices of Section 1.5; that Hermitian matrices have real eigenvalues and orthogonal eigenvectors for different eigenvalues; that $A^2=\pm I$ forces the eigenvalues $\pm1$ or $\pm i$ and gives the projectors $\tfrac12(I\pm A)$; the Clifford relations of the three $2\times2$ matrices $X$, $Z$, $Y$; that a real symmetric matrix cannot square to $-I$, and that an orthogonal matrix with square $+I$ is symmetric and one with square $-I$ antisymmetric; that a symmetric array contracted with an antisymmetric one gives zero, and its consequence for quadratic expressions; the product rule for matrix derivatives and the derivative of the inverse; the chain rule (from the mean value theorem); Jacobi's formula for diagonal and for $2\times2$ matrices, and from it the derivative of $\sqrt{|\det g|}$; integration by parts; the vanishing of integrals of derivatives with vanishing boundary values on a box; the polar-coordinate metric and the transformation law of the metric; the invariance of the volume element (from the change-of-variables formula); and $\det g=+\cos^2z$, $\sqrt{|g|}=\cos z$ and signature (4,4) for the metric of the primordial field.

We \textbf{assumed} (standard mathematics used without proof): the rules for powers, the addition theorems of trigonometry, and that every point at distance 1 from the origin of the plane is $(\cos\theta,\sin\theta)$ for some angle $\theta$; that all bases of $\mathbb R^n$ have $n$ vectors; the area formula $|ad-bc|$ for a parallelogram (not used later); for general $n$, that $\det(AB)=\det A\det B$, that $\det A^T=\det A$, that $A$ is invertible exactly when $\det A\ne0$, that a square matrix $M$ is invertible exactly when $Mv=0$ forces $v=0$ (only the direction from "forces $v=0$" to "invertible" is assumed), and that exchanging two entries of a permutation changes its sign; that the trace and the determinant are the sum and the product of all eigenvalues; the fundamental theorem of algebra; that multiplicities of eigenvalues of Hermitian matrices count independent eigenvectors, and the spectral theorem; Schwarz's theorem on mixed partial derivatives; that a continuous function which is not zero at a point keeps its sign near that point; the mean value theorems of differential and integral calculus; Taylor's theorem; Jacobi's formula for general matrices; the fundamental theorem of calculus; Fubini's theorem; the change-of-variables formula for multiple integrals; and Sylvester's law of inertia. We also used the \textbf{conventions} of the book (ASSUMED choices): counting from 0, the summation convention for one upper and one lower index, the metric $\eta=\mathrm{diag}(+1,+1,+1,+1,-1,-1,-1,-1)$, and the sign convention in which space-like vectors have positive squared length.

\subsection{Exercises}

Exercise 1.1. (a) How many entries does the list $(\Psi_0,\dots,\Psi_{15})$ have, and what is the label of its tenth entry? (b) Compute $\sum_{k=0}^{4}(2k+1)$. (c) Compute $\sum_{a=0}^{7}\eta_{aa}$ and $\prod_{a=0}^{7}\eta_{aa}$.

Exercise 1.2. Are the vectors $u_0=(1,2,0)^T$, $u_1=(0,1,1)^T$ and $u_2=(1,3,1)^T$ linearly independent? If not, give a linear combination with coefficients not all zero that vanishes, and relate your answer to the determinant example of Section 1.4.

Exercise 1.3. (A substitution rule of the notebook.) Let $0<z<\pi/2$ and let $a$ be real. (a) Simplify $1/\sqrt{\sin^{1/3}z/e^{2a}}$. (b) Cell 1058 of the notebook replaces this expression by $1/(e^{a}\sin^{1/6}z)$. Compute the ratio of the correct value to the notebook's value. For which $a$ do they agree?

Exercise 1.4. For $A=\begin{pmatrix}1&2\\ 0&1\end{pmatrix}$ and $B=\begin{pmatrix}0&1\\ 1&0\end{pmatrix}$ compute $AB$, $BA$, $[A,B]$, $\{A,B\}$, $\mathrm{tr}(AB)$, $\mathrm{tr}(BA)$, $\det A$, $\det B$ and $\det(AB)$.

Exercise 1.5. Let $X$ and $Y$ be $n\times n$ matrices and $M=\begin{pmatrix}0&X\\ Y&0\end{pmatrix}$. (a) Show that $M^T=\begin{pmatrix}0&Y^T\\ X^T&0\end{pmatrix}$, so that $M$ is symmetric exactly when $Y=X^T$ and antisymmetric exactly when $Y=-X^T$. (b) For $n=1$, $X=1$, $Y=-1$ write down $M$ and $M^2$. (c) For $n=2$ and $X=Y=\begin{pmatrix}0&1\\ 1&0\end{pmatrix}$ compute $M^2$.

Exercise 1.6. (a) Compute $(2+i)(1-3i)$. (b) Write $1/(3+4i)$ in the form $a+ib$. (c) Compute $|3+4i|$. (d) Show that $(e^{i\theta})^\ast=e^{-i\theta}$ and $e^{i\pi}=-1$. (e) Show that $|e^{-i\varepsilon t}|=1$ for real $\varepsilon$ and $t$, and compute $|e^{-i\varepsilon t}|$ for $\varepsilon=i\gamma$ with real $\gamma$.

Exercise 1.7. Let $A=\begin{pmatrix}2&1-i\\ 1+i&3\end{pmatrix}$. (a) Show that $A$ is Hermitian. (b) Compute its characteristic polynomial and eigenvalues. (c) Find an eigenvector for each eigenvalue and check that the two are orthogonal.

Exercise 1.8. Show that $Y=\begin{pmatrix}0&-1\\ 1&0\end{pmatrix}$ has the eigenvalues $\pm i$, find the eigenvectors, and show that they are orthogonal ($w^\dagger v=0$), although $Y$ is not Hermitian. Explain why with fact 2 of Section 1.5.

Exercise 1.9. (a) For $Z=\mathrm{diag}(1,-1)$ compute $P_\pm=\tfrac12(I\pm Z)$ and verify $P_\pm^2=P_\pm$, $P_+P_-=0$ and $P_++P_-=I$. (b) The chirality matrix of Chapter 2 is the $16\times16$ matrix $\gamma^8=\mathrm{diag}(-I_8,I_8)$. Compute $P_-=\tfrac12(I_{16}-\gamma^8)$ and $P_+=\tfrac12(I_{16}+\gamma^8)$, and say which components of $\Psi=(\Psi_0,\dots,\Psi_{15})^T$ each of them keeps.

Exercise 1.10. (a) Write $\delta^\mu{}_\nu\delta^\nu{}_\rho$ for $\mu=\rho=4$ as an explicit sum and evaluate it. (b) Give $\varepsilon_{210}$, $\varepsilon_{120}$ and $\varepsilon_{112}$. (c) Compute the determinant of the matrix with rows $(2,1,0)$, $(0,1,3)$, $(1,0,1)$ with the formula of Section 1.4. (d) Check $\sum_a\sum_bS_{ab}A_{ab}=0$ for $S=\begin{pmatrix}1&4\\ 4&0\end{pmatrix}$ and $A=\begin{pmatrix}0&-3\\ 3&0\end{pmatrix}$.

Exercise 1.11. For $f(x^0,x^4)=e^{2x^4}\sin x^0$ compute $\partial_0f$, $\partial_4f$, $\partial_0\partial_4f$, $\partial_4\partial_0f$ and $\partial_0\partial_0f$.

Exercise 1.12. (a) With $z=6Hx^0$, compute $\partial_0(\sin^{1/3}z)$. (b) For $f(u,v)=uv^2$ with $u=\cos t$ and $v=\sin t$, compute $df/dt$ with the chain rule and, as a check, directly from $f(t)=\cos t\sin^2t$.

Exercise 1.13. Compute $\int_0^\pi x\sin x\,dx$ by integration by parts.

Exercise 1.14. Let $g=\begin{pmatrix}0&1\\ 1&0\end{pmatrix}$ be a metric on $\mathbb R^2$. (a) What is its signature? (b) Show that $e_0$ and $e_1$ are null. (c) Find a space-like and a time-like vector. (d) Lower the index of $v=(v^0,v^1)$.

Exercise 1.15. With the metric $\eta$ of this book: (a) for $v$ with components $(1,0,0,0,1,0,0,0)$ compute $\eta(v,v)$ and the lowered components $v_\mu$; (b) for $w=(0,0,0,0,1,1,0,0)$ compute $\eta(w,w)$; (c) for $u=(1,1,1,1,1,1,0,0)$ compute $\eta(u,u)$. Classify each vector as space-like, time-like or null.

Exercise 1.16. (The warped coordinate of the primordial field.) With $z=6Hx^0\in(0,\pi/2)$ define $\zeta=\ln(\sin z)/(6H)$. (a) Show that $d\zeta/dx^0=\cot z$, so that the first term of the line element, $\cot^2z\,(dx^0)^2$, equals $d\zeta^2$. (b) Show that $\sin^{1/3}z=e^{2H\zeta}$. (c) Which values does $\zeta$ take?

Exercise 1.17. Let $g(x)=\begin{pmatrix}1&x\\ x&-1\end{pmatrix}$. (a) Compute $\det g$ and its derivative with respect to $x$. (b) Compute $g^{-1}$ and $\mathrm{tr}(g^{-1}\,dg/dx)$, and check Jacobi's formula. (c) What is the signature of $g$, for every $x$?

\subsection{Answers to the exercises}

Answer 1.1. (a) Sixteen entries, labelled 0 to 15; the tenth entry has the label 9, so it is $\Psi_9$. (b) $1+3+5+7+9=25$. (c) The diagonal is $(+1,+1,+1,+1,-1,-1,-1,-1)$, so the sum is $4-4=0$ and the product is $(+1)^4(-1)^4=+1$.

Answer 1.2. They are dependent: $u_0+u_1=(1,3,1)^T=u_2$, so $u_0+u_1-u_2=0$. These are the rows of the $3\times3$ example of Section 1.4, whose determinant is 0. The two facts agree: since $\det A^T=\det A=0$, the equation $A^Tc=0$ has a solution $c\ne0$ (Section 1.4), and the components of such a $c$ are the coefficients of a vanishing combination of the rows of $A$; here $c=(1,1,-1)^T$.

Answer 1.3. (a) By the power rules of Section 1.1, $1/\sqrt{\sin^{1/3}z/e^{2a}}=(\sin^{1/3}z)^{-1/2}(e^{2a})^{1/2}=e^{a}/\sin^{1/6}z$. (b) The ratio is $\dfrac{e^a/\sin^{1/6}z}{1/(e^a\sin^{1/6}z)}=e^{2a}$. The two agree only for $a=0$. For $a\ne0$ the notebook's rule is wrong, which is row L18 of the ledger; Chapter 9 follows its consequences for the notebook's equations.

Answer 1.4. $AB=\begin{pmatrix}2&1\\ 1&0\end{pmatrix}$ and $BA=\begin{pmatrix}0&1\\ 1&2\end{pmatrix}$ (Section 1.3), so $[A,B]=\begin{pmatrix}2&0\\ 0&-2\end{pmatrix}$ and $\{A,B\}=\begin{pmatrix}2&2\\ 2&2\end{pmatrix}$. Both traces are 2. $\det A=1\cdot1-2\cdot0=1$, $\det B=0-1=-1$, and $\det(AB)=2\cdot0-1\cdot1=-1=\det A\det B$.

Answer 1.5. (a) Transposing a block matrix transposes the arrangement of the blocks and each block: block $(i,j)$ of $M^T$ is the transpose of block $(j,i)$ of $M$. This gives $M^T=\begin{pmatrix}0&Y^T\\ X^T&0\end{pmatrix}$. Comparing with $M$ and $-M$ block by block: $M^T=M$ exactly when $Y^T=X$ (equivalently $Y=X^T$), and $M^T=-M$ exactly when $Y=-X^T$. (b) $M=\begin{pmatrix}0&1\\ -1&0\end{pmatrix}$ and $M^2=\mathrm{diag}(XY,YX)=\mathrm{diag}(-1,-1)=-I$; $M$ is antisymmetric (it is $-Y$ in the notation of Section 1.7). (c) $XY=YX=X^2=I_2$, so $M^2=I_4$.

Answer 1.6. (a) $(2+i)(1-3i)=2-6i+i-3i^2=5-5i$. (b) $\dfrac1{3+4i}=\dfrac{3-4i}{(3+4i)(3-4i)}=\dfrac{3-4i}{25}=\dfrac3{25}-\dfrac4{25}i$. (c) $\sqrt{9+16}=5$. (d) $(e^{i\theta})^\ast=\cos\theta-i\sin\theta=\cos(-\theta)+i\sin(-\theta)=e^{-i\theta}$, since $\cos$ is even and $\sin$ is odd; $e^{i\pi}=\cos\pi+i\sin\pi=-1$. (e) For real $\varepsilon t$ the modulus of $e^{-i\varepsilon t}$ is 1 by Section 1.5. For $\varepsilon=i\gamma$, $e^{-i\varepsilon t}=e^{-i\cdot i\gamma t}=e^{\gamma t}$, whose modulus $e^{\gamma t}$ grows without bound for $\gamma>0$.

Answer 1.7. (a) $A^T=\begin{pmatrix}2&1+i\\ 1-i&3\end{pmatrix}$, and conjugating every entry gives $A$ back, so $A^\dagger=A$. (b) $p(x)=(x-2)(x-3)-(1-i)(1+i)=x^2-5x+6-2=x^2-5x+4=(x-1)(x-4)$, so the eigenvalues are 1 and 4, both real as Section 1.6 requires. (c) For 4: $(A-4I)v=0$ gives $-2v_0+(1-i)v_1=0$; with $v_1=2$ we get $v=(1-i,\,2)^T$ (check the second row: $(1+i)(1-i)-2=0$). For 1: $(A-I)w=0$ gives $w_0+(1-i)w_1=0$; with $w_1=1$ we get $w=(-1+i,\,1)^T$ (second row: $(1+i)(-1+i)+2=-2+2=0$). Then $w^\dagger v=(-1-i)(1-i)+1\cdot2=(-1+i-i+i^2)+2=0$.

Answer 1.8. $p(x)=x^2+1$, roots $\pm i$. For $+i$: $Yv=iv$ reads $(-v_1,v_0)=(iv_0,iv_1)$, so $v_1=-iv_0$ and $v=(1,-i)^T$. For $-i$: $w=(1,i)^T$. Then $w^\dagger v=1\cdot1+(i)^\ast(-i)=1+(-i)(-i)=1+i^2=0$. The reason: $Y$ is real and antisymmetric, hence anti-Hermitian (fact 1), so $iY$ is Hermitian (fact 2); $iY$ has the same eigenvectors with the real eigenvalues $i\cdot(\pm i)=\mp1$, and eigenvectors of a Hermitian matrix with different eigenvalues are orthogonal (Section 1.6).

Answer 1.9. (a) $P_+=\mathrm{diag}(1,0)$ and $P_-=\mathrm{diag}(0,1)$. Then $P_\pm^2=P_\pm$ (the entries 0 and 1 are their own squares), $P_+P_-=\mathrm{diag}(0,0)=0$ and $P_++P_-=I$. (b) $P_-=\tfrac12\mathrm{diag}(I_8+I_8,\,I_8-I_8)=\mathrm{diag}(I_8,0)$ keeps $\Psi_0,\dots,\Psi_7$ and removes $\Psi_8,\dots,\Psi_{15}$; $P_+=\mathrm{diag}(0,I_8)$ keeps $\Psi_8,\dots,\Psi_{15}$. So $\Psi_0,\dots,\Psi_7$ are the components with $\gamma^8=-1$ (chirality $-1$) and $\Psi_8,\dots,\Psi_{15}$ those with chirality $+1$, as in Stage 1 (Stage-1 document, Result 3.6).

Answer 1.10. (a) $\sum_{\nu=0}^7\delta^4{}_\nu\delta^\nu{}_4$; only $\nu=4$ contributes, giving $\delta^4{}_4\delta^4{}_4=1=\delta^4{}_4$. (b) $\varepsilon_{210}=-1$ (three inversions), $\varepsilon_{120}=+1$ (two inversions), $\varepsilon_{112}=0$ (a repeated index). (c) With the six terms of Section 1.4: $A_{00}A_{11}A_{22}=2\cdot1\cdot1=2$, $A_{01}A_{12}A_{20}=1\cdot3\cdot1=3$, and the other four terms each contain a zero entry, so $\det A=5$. (d) $1\cdot0+4\cdot(-3)+4\cdot3+0\cdot0=0$.

Answer 1.11. $\partial_0f=e^{2x^4}\cos x^0$, $\partial_4f=2e^{2x^4}\sin x^0$, $\partial_0\partial_4f=\partial_4\partial_0f=2e^{2x^4}\cos x^0$ (as Schwarz's theorem says), and $\partial_0\partial_0f=-e^{2x^4}\sin x^0$.

Answer 1.12. (a) $\partial_0(\sin^{1/3}z)=6H\cdot\tfrac13\sin^{-2/3}z\,\cos z=2H\cos z\,\sin^{-2/3}z$. (b) Chain rule: $\partial_uf\,u'+\partial_vf\,v'=v^2(-\sin t)+2uv\cos t=-\sin^3t+2\sin t\cos^2t$. Directly: $\frac d{dt}(\cos t\sin^2t)=-\sin t\sin^2t+\cos t\cdot2\sin t\cos t$, the same.

Answer 1.13. With $f=x$ and $g'=\sin x$, $g=-\cos x$: $\int_0^\pi x\sin x\,dx=\bigl[-x\cos x\bigr]_0^\pi+\int_0^\pi\cos x\,dx=\pi+\bigl[\sin x\bigr]_0^\pi=\pi$.

Answer 1.14. (a) $g$ is the matrix $X$ of Section 1.6, with eigenvalues $+1$ and $-1$: signature (1,1). (b) $g(e_0,e_0)=g_{00}=0$ and $g(e_1,e_1)=g_{11}=0$, and neither vector is zero. (c) $g(v,v)=2v^0v^1$, so $(1,1)^T$ is space-like ($g=2$) and $(1,-1)^T$ is time-like ($g=-2$). (d) $v_0=g_{0\nu}v^\nu=v^1$ and $v_1=g_{1\nu}v^\nu=v^0$: lowering exchanges the two components.

Answer 1.15. (a) $\eta(v,v)=1-1=0$: $v$ is null; $v_\mu=(1,0,0,0,-1,0,0,0)$. (b) $\eta(w,w)=-1-1=-2$: time-like. (c) $\eta(u,u)=1+1+1+1-1-1=2$: space-like.

Answer 1.16. (a) By the chain rule, $\dfrac{d\zeta}{dx^0}=\dfrac1{6H}\,\dfrac{\cos z}{\sin z}\,\dfrac{dz}{dx^0}=\dfrac1{6H}\cot z\cdot6H=\cot z$, so $d\zeta=\cot z\,dx^0$ and $d\zeta^2=\cot^2z\,(dx^0)^2$. (b) $e^{6H\zeta}=\sin z$, so $\sin^{1/3}z=(e^{6H\zeta})^{1/3}=e^{2H\zeta}$. (c) For $0<z<\pi/2$, $0<\sin z<1$, so $\ln\sin z<0$ and $\zeta$ takes all values in $(-\infty,0)$. This is the warped form of the metric used in Chapter 9 (Stage-2 document, §4.4).

Answer 1.17. (a) $\det g=-1-x^2$ and $d(\det g)/dx=-2x$. (b) By the $2\times2$ formula, $g^{-1}=\dfrac1{-1-x^2}\begin{pmatrix}-1&-x\\ -x&1\end{pmatrix}=\dfrac1{1+x^2}\begin{pmatrix}1&x\\ x&-1\end{pmatrix}$. With $dg/dx=\begin{pmatrix}0&1\\ 1&0\end{pmatrix}$, $g^{-1}\,dg/dx=\dfrac1{1+x^2}\begin{pmatrix}x&1\\ -1&x\end{pmatrix}$, whose trace is $2x/(1+x^2)$. Jacobi's formula: $\det g\cdot\mathrm{tr}(g^{-1}dg/dx)=(-1-x^2)\cdot2x/(1+x^2)=-2x$, which is the derivative computed in (a). (c) The determinant, the product of the two eigenvalues, is negative for every $x$, so one eigenvalue is positive and one negative: the signature is (1,1) for every $x$.

\section{Clifford algebras and spinors from zero}

\subsection{What this chapter does}

The field of this book, dirac16complex, has 16 components $\Psi_0,\dots,\Psi_{15}$ at every point of an 8-dimensional spacetime, and its field equation (Chapter 7) contains eight real 16 by 16 matrices $\gamma^0,\dots,\gamma^7$, the \textbf{gamma matrices}. Neither the number 16 nor these matrices are arbitrary. This chapter explains, from nothing, where they come from and what they do.

The story has five steps.

\begin{enumerate}
\item The \textbf{square-root problem.} The length of a vector is the square root of a sum of squares. Paul Dirac discovered that a sum of squares can be written as the square of a \emph{linear} expression, provided the coefficients are matrices that anticommute. The Pauli matrices (three directions) and the Dirac matrices (four directions) are the first examples (Section 2.3 and Section 2.4).
\item The \textbf{Clifford algebra} $\mathrm{Cl}(p,q)$ is the general rule behind these examples, for any number of directions with any signs (Section 2.5).
\item \textbf{How big the matrices must be.} In 8 directions the anticommuting matrices must be at least 16 by 16, and for signature (4,4) real 16 by 16 matrices exist. We construct them in two ways: by a tensor-product recipe (Section 2.7) and by the author's notebook, whose matrices are called $\mathrm{T16}^A$ (Section 2.8). This is why a spinor in 4+4 dimensions has 16 components.
\item The \textbf{special matrices} of the theory: the charge matrix $C=\sigma_{16}$, the chirality $\gamma^8$, the spin generators $S^{ab}$ and the matrix $B$ (Section 2.9 to Section 2.12).
\item \textbf{Symmetry groups.} After defining groups and representations from zero (Section 2.13) we build the groups Pin(4,4) and Spin(4,4), which act on spinors, and we prove the \textbf{representation theorem}: under Pin(4,4) the 16 complex components form one indivisible (irreducible) block, while under Spin(4,4) they split into two blocks of 8 that are irreducible and genuinely different from each other (Section 2.14 and Section 2.15).
\end{enumerate}

Every statement of this chapter is proved in the text, step by step, with two kinds of exceptions, each labelled where it occurs. First, a few standard results of mathematics are quoted without proof: three basic facts of linear algebra (Section 2.2), the addition theorems of trigonometry (quoted in Chapter 1 and used in Section 2.11), the power series of the exponential, sine and cosine functions (calculus of one variable) and the convergence of the exponential series for matrices (Section 2.11), and the Cartan--Dieudonné theorem on reflections (Section 2.14). Second, a few properties of the notebook's particular matrices (for example that its eight gammas satisfy the Clifford relation) are taken in this chapter from the exact machine checks of the repository; Chapter 3 then proves them by hand. Where the repository has a machine check of a statement, the check is named; Section 2.16 lists the checks with their report files.

The sources of this chapter are these files of the repository:

\begin{itemize}
\item the Stage-1 document \texttt{provenan\allowbreak{}ce/DIRAC\allowbreak{}16COMPLE\allowbreak{}X\_ARBITR\allowbreak{}ARY\_FIEL\allowbreak{}D.md}, its §2 (notation), §3 (the gamma matrices) and §4 (Pin(4,4) and Spin(4,4));
\item the contract \texttt{handoff/\allowbreak{}specs/CO\allowbreak{}NTRACT.m\allowbreak{}d}, its §1, §2 and the errata in §11;
\item the exact Python module \texttt{scripts/\allowbreak{}d16c\_exa\allowbreak{}ct.py} and the checker \texttt{scripts/\allowbreak{}check\_di\allowbreak{}rac16com\allowbreak{}plex\_alg\allowbreak{}ebra.py};
\item the Wolfram package \texttt{wolfram/\allowbreak{}Dirac16C\allowbreak{}omplexAl\allowbreak{}gebra.wl};
\item the reports \texttt{python-a\allowbreak{}lgebra-r\allowbreak{}eport.js\allowbreak{}on} and \texttt{wolfram-\allowbreak{}algebra-\allowbreak{}report.j\allowbreak{}son} and the exact data file \texttt{algebra-\allowbreak{}fixture.\allowbreak{}json}, all in the folder \texttt{artifact\allowbreak{}s/dirac1\allowbreak{}6complex\allowbreak{}/arbitra\allowbreak{}ry-field\allowbreak{}/}, and, for one check, the report \texttt{python-g\allowbreak{}eometry-\allowbreak{}report.j\allowbreak{}son} of the exact Python geometry checker \texttt{scripts/\allowbreak{}check\_di\allowbreak{}rac16com\allowbreak{}plex\_geo\allowbreak{}metry.py}, in the same folder. Every report named in this chapter lies in that folder, and a check name such as \texttt{ALG\_clif\allowbreak{}ford} refers to the check of that name in both algebra reports, unless the text names a single report; Section 2.16 lists them.
\end{itemize}

We use the conventions of the whole book (Chapter 1): everything is counted from 0; the eight coordinates are $x_0,\dots,x_7$; $x_4$ is the time; the flat metric is $\eta=\mathrm{diag}(+1,+1,+1,+1,-1,-1,-1,-1)$, so directions 0 to 3 are space-like and directions 4 to 7 are time-like. We write $\eta_{ab}$ for the entry in row $a$ and column $b$ of $\eta$ and $\eta^{ab}$ for the entry of its inverse; because $\eta$ is diagonal with entries $\pm1$, $\eta^{ab}=\eta_{ab}$ as numbers. So $\eta_{aa}=+1$ for $a\le3$, $\eta_{aa}=-1$ for $a\ge4$, and $\eta_{ab}=0$ for $a\ne b$.

\subsection{Matrix tools used in this chapter}

This section collects every matrix fact the chapter uses. Chapter 1 introduces vectors, matrices and complex numbers; here we add what is special to this chapter, with proofs.

\textbf{Entries and products.} An $n$ by $n$ matrix $M$ has entries $M_{ij}$, where $i$ is the row and $j$ the column, both counted from 0 to $n-1$. The product of two matrices is $(AB)_{ij}=\sum_{k=0}^{n-1}A_{ik}B_{kj}$, and a matrix acts on a column $u$ by $(Mu)_i=\sum_jM_{ij}u_j$. Matrix multiplication is associative, $(AB)D=A(BD)$, but in general not commutative, $AB\ne BA$. $I_n$ (or simply $I$, or $1$) is the identity matrix. The \textbf{transpose} $M^T$ has entries $(M^T)_{ij}=M_{ji}$, and $(AB)^T=B^TA^T$. The \textbf{conjugate transpose} is $M^\dagger=(M^\ast)^T$, where ${}^\ast$ is complex conjugation of every entry, and $(AB)^\dagger=B^\dagger A^\dagger$. The \textbf{inverse} $M^{-1}$, if it exists, satisfies $MM^{-1}=M^{-1}M=I$. The \textbf{elementary matrix} $E_{ij}$ has a 1 in row $i$, column $j$, and 0 everywhere else.

\textbf{Linear combinations, span, independence, dimension.} A \textbf{linear combination} of matrices $M_0,\dots,M_{r-1}$ is $c_0M_0+\dots+c_{r-1}M_{r-1}$ with numbers $c_k$ (real or complex). The set of all linear combinations is the \textbf{span}. The matrices are \textbf{linearly independent} if the only combination that gives the zero matrix is the one with all $c_k=0$. A set of matrices (or of columns) that contains every linear combination of its elements is a \textbf{subspace}; the span of any list is a subspace. A \textbf{basis} of a subspace is an independent list that spans it, and the number of elements of a basis is its \textbf{dimension}. The set $\mathrm{Mat}_n(\mathbb R)$ of all real $n$ by $n$ matrices has the basis $E_{ij}$ ($i,j=0,\dots,n-1$), because $M=\sum_{ij}M_{ij}E_{ij}$ and this is zero only if every $M_{ij}=0$; so its dimension is $n^2$. The same holds for complex matrices, $\mathrm{Mat}_n(\mathbb C)$. The \textbf{rank} of a list of matrices (or of the rows of a matrix) is the largest number of independent elements in it, which is the dimension of its span. The \textbf{image} of a $d$ by $d$ matrix $M$ is the set of all columns $Mu$; it is a subspace, because $cMu+c'Mu'=M(cu+c'u')$.

\textbf{Three facts of linear algebra used without proof.} They belong to elementary linear algebra, and we quote them. (Chapter 1, in Section 1.2 and Section 1.4, already quotes (iii), the fact that all bases have the same number of vectors, and the fact that a square matrix is invertible exactly when $Mu=0$ has only the solution $u=0$, the last part of (ii); the two statements of (i) and the count $m-r$ of (ii) are quoted here.) The facts are: (i) in a subspace of dimension $m$ there are never more than $m$ independent elements, and $m$ independent elements always form a basis; (ii) the solutions $u$ of $r$ independent homogeneous linear equations in $m$ unknowns form a subspace of dimension $m-r$ (so one nonzero equation in 8 unknowns has a 7-dimensional solution space), and a square matrix $M$ whose equation $Mu=0$ has only the solution $u=0$ is invertible; (iii) $\det(AB)=\det A\,\det B$, so that the determinant does not change under a change of basis, $\det(S^{-1}MS)=\det M$, and the determinant of a diagonal matrix is the product of its diagonal entries.

\textbf{Commutator and anticommutator.} For two matrices,

\[
[A,B]:=AB-BA,\qquad\{A,B\}:=AB+BA .
\]

$A$ and $B$ \textbf{commute} if $[A,B]=0$ and \textbf{anticommute} if $\{A,B\}=0$, that is, if $AB=-BA$. Example:

\[
A=\begin{pmatrix}0&1\\1&0\end{pmatrix},\quad B=\begin{pmatrix}1&0\\0&-1\end{pmatrix}:\qquad AB=\begin{pmatrix}0&-1\\1&0\end{pmatrix},\quad BA=\begin{pmatrix}0&1\\-1&0\end{pmatrix},
\]

so $\{A,B\}=0$ and $[A,B]=2AB$. These two matrices anticommute.

\textbf{Trace.} The trace is the sum of the diagonal entries, $\operatorname{tr}M=\sum_iM_{ii}$. It is linear, $\operatorname{tr}(cA+dB)=c\operatorname{tr}A+d\operatorname{tr}B$, and it has the \textbf{cyclic property}

\[
\operatorname{tr}(AB)=\operatorname{tr}(BA).
\]

\emph{Proof.} $\operatorname{tr}(AB)=\sum_i\sum_kA_{ik}B_{ki}=\sum_k\sum_iB_{ki}A_{ik}=\operatorname{tr}(BA)$; only the order of two finite sums was exchanged. $\square$ A consequence: for an invertible $S$, $\operatorname{tr}(S^{-1}AS)=\operatorname{tr}(ASS^{-1})=\operatorname{tr}A$. The trace does not change when the basis is changed.

\textbf{Block matrices.} A 16 by 16 matrix can be cut into four 8 by 8 blocks, and block matrices multiply like 2 by 2 matrices whose entries are matrices, keeping the order of the factors:

\[
\begin{pmatrix}A&B\\ C&D\end{pmatrix}\begin{pmatrix}E&F\\ G&H\end{pmatrix}=\begin{pmatrix}AE+BG&AF+BH\\ CE+DG&CF+DH\end{pmatrix}.
\]

(Write out the sum $\sum_k$ in the product and split it into $k<8$ and $k\ge8$.) The transpose of a block matrix transposes the pattern and every block: $\begin{pmatrix}A&B\\ C&D\end{pmatrix}^T=\begin{pmatrix}A^T&C^T\\ B^T&D^T\end{pmatrix}$.

\textbf{Kronecker (tensor) product.} For an $n$ by $n$ matrix $A$ and an $m$ by $m$ matrix $B$, the Kronecker product $A\otimes B$ is the $nm$ by $nm$ matrix made of $n\times n$ blocks, the block in position $(i,j)$ being $A_{ij}B$. In indices,

\[
(A\otimes B)_{im+k,\;jm+l}=A_{ij}B_{kl}\qquad(i,j=0,\dots,n-1;\ k,l=0,\dots,m-1).
\]

Example with $n=m=2$:

\[
\begin{pmatrix}0&1\\1&0\end{pmatrix}\otimes\begin{pmatrix}1&0\\0&-1\end{pmatrix}=\begin{pmatrix}0&0&1&0\\0&0&0&-1\\1&0&0&0\\0&-1&0&0\end{pmatrix}.
\]

\textbf{Mixed-product rule.} $(A\otimes B)(D\otimes F)=(AD)\otimes(BF)$.

\emph{Proof.} The entry in row $im+k$ and column $jm+l$ of the left side is a sum over all columns $rm+s$ of the first factor:

\[
\sum_{r=0}^{n-1}\sum_{s=0}^{m-1}(A\otimes B)_{im+k,\,rm+s}(D\otimes F)_{rm+s,\,jm+l}=\sum_{r,s}A_{ir}B_{ks}D_{rj}F_{sl}=\Bigl(\sum_rA_{ir}D_{rj}\Bigr)\Bigl(\sum_sB_{ks}F_{sl}\Bigr),
\]

which is $(AD)_{ij}(BF)_{kl}$, the entry of $(AD)\otimes(BF)$. $\square$ The same rule holds for products of more factors, $A\otimes B\otimes D$ with three or four slots, applied slot by slot; and $(A\otimes B)^T=A^T\otimes B^T$ (exchange $i\leftrightarrow j$ and $k\leftrightarrow l$ in the index formula).

\textbf{Special kinds of matrices.} A real matrix $M$ is \textbf{symmetric} if $M^T=M$, \textbf{antisymmetric} if $M^T=-M$, and \textbf{orthogonal} if $M^TM=I$, that is, $M^{-1}=M^T$. A \textbf{signed permutation} matrix has exactly one nonzero entry in every row and every column, and that entry is $+1$ or $-1$; acting on a column it reorders the components and flips some signs. Every signed permutation matrix is orthogonal: $(M^TM)_{ij}=\sum_kM_{ki}M_{kj}$, and for $i\ne j$ no row $k$ has nonzero entries in both columns $i$ and $j$, while for $i=j$ exactly one term is $(\pm1)^2=1$. A complex matrix is \textbf{Hermitian} if $M^\dagger=M$, \textbf{anti-Hermitian} if $M^\dagger=-M$, and \textbf{unitary} if $M^\dagger M=I$.

\textbf{Matrices whose square is $\pm1$, and projectors.} A matrix $\Pi$ with $\Pi^2=\Pi$ is a \textbf{projector}. Let $M$ be a $d$ by $d$ matrix with $M^2=I$. Then

\[
P_\pm:=\tfrac12(I\pm M)\quad\text{satisfy}\quad P_\pm^2=P_\pm,\qquad P_+P_-=P_-P_+=0,\qquad P_++P_-=I,\qquad M=P_+-P_- .
\]

\emph{Proof.} $P_\pm^2=\tfrac14(I\pm2M+M^2)=\tfrac14(2I\pm2M)=P_\pm$ and $P_+P_-=\tfrac14(I-M^2)=0$; the last two follow by adding and subtracting. $\square$ So every column $u$ splits as $u=P_+u+P_-u$, with $M(P_\pm u)=\pm P_\pm u$ (because $MP_\pm=\tfrac12(M\pm I)=\pm P_\pm$). The images of $P_+$ and $P_-$ are the \textbf{eigenspaces} of $M$ for the \textbf{eigenvalues} $+1$ and $-1$: every column of the image of $P_\pm$ satisfies $Mu=\pm u$, and conversely $Mu=\pm u$ gives $P_\pm u=\tfrac12(u+u)=u$, so $u$ lies in the image of $P_\pm$. Every element $a=P_+w$ of the image of $P_+$ satisfies $P_+a=P_+^2w=a$, and in the same way $P_-b=b$ for every $b$ in the image of $P_-$.

Choose a basis of each image. The two bases together are a basis of the whole space. They span it, because $u=P_+u+P_-u$ and each of the two terms is a combination of the basis of its image. They are independent: a vanishing combination of the joint list has the form $a+b=0$, with $a$ a combination of the first basis (so $a=P_+a$) and $b$ a combination of the second (so $b=P_-b$). Applying $P_+$ gives $0=P_+a+P_+P_-b=a$ (since $P_+P_-=0$), then $b=-a=0$, and the independence of each basis makes all coefficients zero. In this joint basis $M$ multiplies the first basis vectors by $+1$ and the others by $-1$, so it becomes $\mathrm{diag}(+1,\dots,+1,-1,\dots,-1)$, and $P_+=\mathrm{diag}(1,\dots,1,0,\dots,0)$, whose trace is the number of $+1$'s. Since the trace does not depend on the basis,

\[
\dim(\text{eigenspace }+1)=\operatorname{tr}P_+=\tfrac12(d+\operatorname{tr}M),\qquad\dim(\text{eigenspace }-1)=\tfrac12(d-\operatorname{tr}M).
\]

In particular a \textbf{traceless} matrix ($\operatorname{tr}M=0$) with $M^2=I$ has the eigenvalue $+1$ exactly $d/2$ times and $-1$ exactly $d/2$ times, and its determinant is $(+1)^{d/2}(-1)^{d/2}$. If instead $M^2=-I$, the same argument, joint basis included, works with complex coefficients and $P_\pm=\tfrac12(I\mp iM)$, which satisfy the same projector relations and $MP_\pm=\pm iP_\pm$; it gives eigenvalues $+i$ and $-i$; for a traceless $M$ each occurs $d/2$ times, and $\det M=i^{d/2}(-i)^{d/2}=(i\cdot(-i))^{d/2}=1$. For a real symmetric or complex Hermitian matrix with $M^2=I$, the pair (number of $+1$ eigenvalues, number of $-1$ eigenvalues) is called its \textbf{signature}.

\subsection{The square-root problem and the Pauli matrices}

\textbf{Pythagoras.} A vector $(p,q)$ in the plane has squared length $p^2+q^2$. Can the square root of $p^2+q^2$ be written as a \emph{linear} expression $\alpha p+\beta q$? Squaring,

\[
(\alpha p+\beta q)^2=\alpha^2p^2+(\alpha\beta+\beta\alpha)\,pq+\beta^2q^2 .
\]

We need $\alpha^2=1$, $\beta^2=1$ and $\alpha\beta+\beta\alpha=0$. With ordinary numbers this is impossible: $\alpha\beta+\beta\alpha=2\alpha\beta=0$ forces $\alpha=0$ or $\beta=0$. With matrices it is possible, because matrices can anticommute. The two matrices of the example in Section 2.2,

\[
\sigma_x=\begin{pmatrix}0&1\\1&0\end{pmatrix},\qquad\sigma_z=\begin{pmatrix}1&0\\0&-1\end{pmatrix},
\]

satisfy $\sigma_x^2=\sigma_z^2=I$ and $\sigma_x\sigma_z+\sigma_z\sigma_x=0$. Hence $(p\sigma_x+q\sigma_z)^2=(p^2+q^2)I$.

\textbf{Worked example.} For $p=3$, $q=4$:

\[
3\sigma_x+4\sigma_z=\begin{pmatrix}4&3\\3&-4\end{pmatrix},\qquad\begin{pmatrix}4&3\\3&-4\end{pmatrix}^2=\begin{pmatrix}16+9&12-12\\12-12&9+16\end{pmatrix}=25\,I .
\]

The matrix $3\sigma_x+4\sigma_z$ is a "square root" of $25=3^2+4^2$.

\textbf{Minus signs.} Spacetime needs minus signs: in this book $\eta(v,v)=v_0^2+v_1^2+v_2^2+v_3^2-v_4^2-\dots-v_7^2$. A real matrix can square to $-I$:

\[
N=\begin{pmatrix}0&1\\-1&0\end{pmatrix},\qquad N^2=\begin{pmatrix}-1&0\\0&-1\end{pmatrix}=-I,
\]

and $\sigma_xN+N\sigma_x=\begin{pmatrix}-1&0\\0&1\end{pmatrix}+\begin{pmatrix}1&0\\0&-1\end{pmatrix}=0$. So $(p\sigma_x+qN)^2=(p^2-q^2)I$: a square root of a difference of squares, with real entries. This little matrix $N$ is the seed of everything real in this book.

\textbf{The general requirement.} Suppose we want matrices $\gamma^0,\dots,\gamma^{n-1}$ such that for all numbers $p_0,\dots,p_{n-1}$

\[
\Bigl(\sum_ap_a\gamma^a\Bigr)^2=\Bigl(\sum_a\eta^{aa}p_a^2\Bigr)I
\]

for a given diagonal metric $\eta$. Expanding the square, $\sum_{a,b}p_ap_b\gamma^a\gamma^b$, and pairing the terms $(a,b)$ and $(b,a)$ (they have the same coefficient $p_ap_b$),

\[
\Bigl(\sum_ap_a\gamma^a\Bigr)^2=\tfrac12\sum_{a,b}p_ap_b\bigl(\gamma^a\gamma^b+\gamma^b\gamma^a\bigr).
\]

This equals $\sum_{a,b}\eta^{ab}p_ap_bI$ for all $p$ exactly when

\[
\gamma^a\gamma^b+\gamma^b\gamma^a=2\eta^{ab}I\qquad\text{for all }a,b .
\]

Indeed, if this relation holds, the right-hand side of the pairing formula is $\tfrac12\sum_{a,b}p_ap_b\,2\eta^{ab}I=\sum_{a,b}\eta^{ab}p_ap_bI$. Conversely, suppose that the square equals $\sum_a\eta^{aa}p_a^2I$ for all $p$. Take $p=e_a$, the list with $p_a=1$ and all other entries 0 (the standard basis of Chapter 1, Section 1.2): the square is $(\gamma^a)^2$, so $(\gamma^a)^2=\eta^{aa}I$, which is the relation for $b=a$. Then take $p=e_a+e_b$ with $a\ne b$: the square is $(\gamma^a+\gamma^b)^2=(\gamma^a)^2+(\gamma^b)^2+\gamma^a\gamma^b+\gamma^b\gamma^a$, and it must equal $(\eta^{aa}+\eta^{bb})I$; so $\gamma^a\gamma^b+\gamma^b\gamma^a=0=2\eta^{ab}I$, since $\eta^{ab}=0$ for $a\ne b$.

This is the \textbf{Clifford relation}. It says: each $\gamma^a$ squares to $\eta^{aa}I$ ($+I$ for a space-like direction, $-I$ for a time-like one), and two different gammas anticommute.

\textbf{The Pauli matrices.} For three space directions (metric $+1,+1,+1$) we need a third 2 by 2 matrix anticommuting with $\sigma_x$ and $\sigma_z$ and squaring to $+I$. There is no real one (Exercise 2.2), but there is a complex one. The three \textbf{Pauli matrices} are

\[
\sigma_x=\begin{pmatrix}0&1\\1&0\end{pmatrix},\qquad\sigma_y=\begin{pmatrix}0&-i\\ i&0\end{pmatrix},\qquad\sigma_z=\begin{pmatrix}1&0\\0&-1\end{pmatrix}.
\]

They are named by letters, not numbers, to keep them apart from the gammas. Many books, and the two-dimensional examples of Chapters 4 and 6, number them instead: $\sigma_1=\sigma_x$, $\sigma_2=\sigma_y$, $\sigma_3=\sigma_z$; the numbered names always mean exactly these three matrices. (The real matrices $X$, $Z$ and $Y$ of Section 1.7 are $X=\sigma_x$, $Z=\sigma_z$ and $Y=-i\sigma_y=-N$; from Section 2.5 on, $\sigma_x$ and $\sigma_z$ are also called $P$ and $G$.) The Pauli matrices are Hermitian and traceless, each squares to $I$, and any two anticommute. For instance

\[
\sigma_x\sigma_y=\begin{pmatrix}0&1\\1&0\end{pmatrix}\begin{pmatrix}0&-i\\ i&0\end{pmatrix}=\begin{pmatrix}i&0\\0&-i\end{pmatrix}=i\sigma_z,\qquad\sigma_y\sigma_x=\begin{pmatrix}-i&0\\0&i\end{pmatrix}=-i\sigma_z,
\]

so $\{\sigma_x,\sigma_y\}=0$, and in the same way $\sigma_y\sigma_z=i\sigma_x$ and $\sigma_z\sigma_x=i\sigma_y$. The product of all three is $\sigma_x\sigma_y\sigma_z=i\sigma_z\sigma_z=iI$. Wolfgang Pauli introduced these matrices (1927) to describe the spin of the electron; they return in Chapter 3, where they build the quaternions.

A warning about names: the notebook uses the letter $\sigma$ for two other matrices, the 8 by 8 matrix $\sigma=\begin{pmatrix}0&I_4\\ I_4&0\end{pmatrix}$ and the 16 by 16 charge matrix $\sigma_{16}$ (Section 2.9). The subscripts $x,y,z$ always mean Pauli matrices.

\subsection{Dirac matrices in four dimensions}

\textbf{Why a first-order equation.} In special relativity (units with the speed of light equal to 1) a particle of mass $m$, energy $E$ and momentum $(k_1,k_2,k_3)$ satisfies $E^2=k_1^2+k_2^2+k_3^2+m^2$. In quantum mechanics a plane wave $e^{i(k_1x_1+k_2x_2+k_3x_3-Et)}$ turns $\partial/\partial x_j$ into $ik_j$ and $\partial/\partial t$ into $-iE$, so the energy relation becomes a second-order wave equation. Dirac (1928) wanted a \emph{first-order} equation whose square is that wave equation; he needed the square root of $E^2-k_1^2-k_2^2-k_3^2$, four directions with one sign opposite to the other three. By Section 2.3 this needs four matrices satisfying the Clifford relation.

\textbf{Four anticommuting matrices.} Two by two matrices are too small: Section 2.6 proves that four matrices satisfying the Clifford relation must be at least 4 by 4. Dirac's choice, in the particle-physics convention $\eta=\mathrm{diag}(+1,-1,-1,-1)$ with time first, is built from 2 by 2 blocks and the Pauli matrices:

\[
\gamma_{\mathrm D}^0=\begin{pmatrix}I_2&0\\0&-I_2\end{pmatrix},\qquad\gamma_{\mathrm D}^j=\begin{pmatrix}0&\sigma_j\\-\sigma_j&0\end{pmatrix}\quad(j=x,y,z).
\]

\emph{Check.} By the block rule, $(\gamma_{\mathrm D}^0)^2=\mathrm{diag}(I_2,I_2)=I_4$, and

\[
(\gamma_{\mathrm D}^j)^2=\begin{pmatrix}-\sigma_j^2&0\\0&-\sigma_j^2\end{pmatrix}=-I_4,\qquad\gamma_{\mathrm D}^0\gamma_{\mathrm D}^j=\begin{pmatrix}0&\sigma_j\\ \sigma_j&0\end{pmatrix}=-\gamma_{\mathrm D}^j\gamma_{\mathrm D}^0,
\]

\[
\gamma_{\mathrm D}^j\gamma_{\mathrm D}^k+\gamma_{\mathrm D}^k\gamma_{\mathrm D}^j=\begin{pmatrix}-(\sigma_j\sigma_k+\sigma_k\sigma_j)&0\\0&-(\sigma_j\sigma_k+\sigma_k\sigma_j)\end{pmatrix}=0\quad(j\ne k).
\]

So $\{\gamma_{\mathrm D}^\mu,\gamma_{\mathrm D}^\nu\}=2\,\mathrm{diag}(+1,-1,-1,-1)^{\mu\nu}I_4$. The subscript D (for Dirac) keeps these matrices apart from the 16 by 16 gammas of this book.

\textbf{Changing the sign convention.} This book counts space-like directions with $+1$ and time-like ones with $-1$. If matrices $\gamma^a$ satisfy the Clifford relation with a metric $\eta$, then the matrices $i\gamma^a$ satisfy it with $-\eta$, because $\{i\gamma^a,i\gamma^b\}=i^2\{\gamma^a,\gamma^b\}=-2\eta^{ab}I$. So $i\gamma_{\mathrm D}^\mu$ are Dirac matrices for the metric $(-1,+1,+1,+1)$ with time first. Neither set is real: $\gamma_{\mathrm D}^y$ contains the imaginary $\sigma_y$, and the factor $i$ makes $i\gamma_{\mathrm D}^0$, $i\gamma_{\mathrm D}^x$ and $i\gamma_{\mathrm D}^z$ imaginary (only $i\gamma_{\mathrm D}^y$ is real, because $i\sigma_y=N$). Real 4 by 4 matrices for $(-1,+1,+1,+1)$ nevertheless exist, for example $N\otimes I_2$, $\sigma_x\otimes I_2$, $\sigma_z\otimes\sigma_x$, $\sigma_z\otimes\sigma_z$ with the $N$ of Section 2.3 (check with the mixed-product rule of Section 2.2, using that any two of $N$, $\sigma_x$ and $\sigma_z$ anticommute, Section 2.3 and Section 2.5). For signature (4,4), Section 2.7 constructs real 16 by 16 matrices.

\textbf{What the gammas do in a field equation.} Take any matrices with the Clifford relation for the metric of this book, $\eta=\mathrm{diag}(+1,+1,+1,+1,-1,-1,-1,-1)$, and consider the flat-space equation $\sum_\mu\gamma^\mu\partial_\mu\Psi=m\Psi$ for a column $\Psi$ of functions ($\partial_\mu=\partial/\partial x^\mu$, with the coordinates written with an upper index in index formulas, as Section 1.8 explains; in words they keep the names $x_0,\dots,x_7$). Applying the operator $\sum_\nu\gamma^\nu\partial_\nu$ once more and using $\partial_\mu\partial_\nu=\partial_\nu\partial_\mu$ and the pairing argument of Section 2.3,

\[
m^2\Psi=\sum_{\mu,\nu}\gamma^\mu\gamma^\nu\partial_\mu\partial_\nu\Psi=\sum_\mu\eta^{\mu\mu}\partial_\mu^2\Psi .
\]

For a plane wave $\Psi=u\,e^{i(k_0x_0+\dots+k_7x_7)}$ with a constant column $u$, $\partial_\mu^2\Psi=-k_\mu^2\Psi$, so

\[
k_4^2=m^2+k_0^2+k_1^2+k_2^2+k_3^2-k_5^2-k_6^2-k_7^2 .
\]

This is the energy--momentum relation of the 8-dimensional theory, with $k_4$ playing the role of the energy. (It is the relation found in Stage 1, §10.6, for the flat-space modes; Chapter 8 discusses what the minus signs of the extra times do.) The gamma matrices are exactly what makes a first-order equation square to the correct second-order one.

\subsection{Clifford algebras in general}

\textbf{Generators and relations.} Fix a number of directions $n$ and a diagonal metric $\eta$ with $p$ entries $+1$ followed by $q$ entries $-1$ ($p+q=n$). The \textbf{Clifford algebra} $\mathrm{Cl}(p,q)$ is built from $n$ symbols $\gamma^0,\dots,\gamma^{n-1}$, the \textbf{generators}, which can be multiplied (associatively) and added, subject only to the Clifford relation

\[
\gamma^a\gamma^b+\gamma^b\gamma^a=2\eta^{ab}\,1\qquad(a,b=0,\dots,n-1).
\]

In this book the generators are always concrete matrices, so the reader may simply read "$\mathrm{Cl}(p,q)$" as "the set of all real linear combinations of products of the $n$ gamma matrices". For this book $n=8$, $p=q=4$, and the algebra is $\mathrm{Cl}(4,4)$.

\textbf{Sets.} A set is written by listing its elements in braces, as in $\{0,\dots,n-1\}$. $A\subseteq S$ ("$A$ is a \textbf{subset} of $S$") means that every element of $A$ is an element of $S$. $\varnothing$ is the \textbf{empty set}, the set with no elements; it is a subset of every set. $|A|$ is the number of elements of $A$, and $\{a_1<a_2<\dots<a_k\}$ is the set of the $k$ numbers $a_1,\dots,a_k$, listed in increasing order.

\textbf{Binomial coefficients.} The \textbf{binomial coefficient} $\binom nk$ (read "$n$ choose $k$") is the number of subsets with $k$ elements of a set with $n$ elements. With the \textbf{factorial} $n!:=1\cdot2\cdots n$ and $0!:=1$ it equals

\[
\binom nk=\frac{n!}{k!\,(n-k)!}\qquad(k=0,\dots,n):
\]

$k$ different elements can be picked one after the other in $n(n-1)\cdots(n-k+1)=n!/(n-k)!$ ways, and each subset of $k$ elements is obtained in this way exactly $k!$ times, once for each order of its elements. For example $\binom82=\frac{8\cdot7}{2}=28$. The \textbf{binomial theorem} states $(x+y)^n=\sum_{k=0}^n\binom nkx^ky^{n-k}$. \emph{Proof.} Multiplying out the $n$ factors of $(x+y)(x+y)\cdots(x+y)$ gives one term for each way of choosing $x$ or $y$ in every factor. The terms in which $x$ is chosen in exactly $k$ factors are all equal to $x^ky^{n-k}$, and there is one of them for each $k$-element subset of the $n$ factors, that is, $\binom nk$ of them. $\square$

\textbf{Monomials.} For a set $A=\{a_1<a_2<\dots<a_k\}\subseteq\{0,\dots,n-1\}$ write

\[
\gamma_A:=\gamma^{a_1}\gamma^{a_2}\cdots\gamma^{a_k},\qquad\gamma_\varnothing:=1 .
\]

The number $k=|A|$ is the \textbf{degree} (or grade) of the monomial; $\gamma_A$ is \textbf{even} if $k$ is even and \textbf{odd} if $k$ is odd. There are $2^n$ subsets of $\{0,\dots,n-1\}$ (each element is in or out), so there are $2^n$ monomials: $\binom n0=1$ of degree 0, $\binom n1=n$ of degree 1, and so on. For $n=8$ there are $2^8=256$ monomials, 128 of them even and 128 odd (Exercise 2.4).

\textbf{Three rules for monomials.} Everything in this chapter rests on three rules, each a direct consequence of the Clifford relation.

(R1) \emph{Every product of generators is $\pm$ a monomial.} To bring a product into increasing order, exchange neighbours: two different neighbours anticommute (a factor $-1$), and two equal neighbours combine to $\gamma^a\gamma^a=\eta^{aa}=\pm1$. In particular $\gamma_A\gamma_B=\pm\gamma_{A\triangle B}$, where $A\triangle B$ (the \textbf{symmetric difference}) is the set of indices that lie in exactly one of $A$ and $B$: the indices in both meet a partner and turn into $\pm1$.

(R2) \emph{The square of a monomial of degree $k$ is}

\[
\gamma_A\gamma_A=(-1)^{k(k-1)/2}\prod_{a\in A}\eta^{aa}\;1 .
\]

\emph{Proof.} In $\gamma^{a_1}\cdots\gamma^{a_k}\gamma^{a_1}\cdots\gamma^{a_k}$ move the second $\gamma^{a_1}$ to the left past the $k-1$ factors $\gamma^{a_k},\dots,\gamma^{a_2}$, all different from it: a sign $(-1)^{k-1}$. It then stands next to the first $\gamma^{a_1}$, and the two give $\eta^{a_1a_1}$. What remains is the square of the monomial $\gamma^{a_2}\cdots\gamma^{a_k}$ of degree $k-1$. Repeating, the total sign is $(-1)^{(k-1)+(k-2)+\dots+1+0}=(-1)^{k(k-1)/2}$. $\square$ So every monomial squares to $+1$ or $-1$, and its inverse is $\gamma_A^{-1}=\pm\gamma_A$ with the same sign.

(R3) \emph{Commuting a generator through a monomial.} If $b\notin A$, then $\gamma^b\gamma_A=(-1)^k\gamma_A\gamma^b$; if $b\in A$, then $\gamma^b\gamma_A=(-1)^{k-1}\gamma_A\gamma^b$.

\emph{Proof.} Moving $\gamma^b$ from the left end to the right end passes each of the $k$ factors once. Passing a factor with a different index gives $-1$; passing the factor $\gamma^b$ itself (if $b\in A$) gives $+1$. $\square$

\textbf{Worked example: $\mathrm{Cl}(1,1)$ is all real 2 by 2 matrices.} Take $n=2$, $\eta=\mathrm{diag}(+1,-1)$, and the generators $\gamma^0=P:=\sigma_x$ and $\gamma^1=N$ of Section 2.3. The four monomials are

\[
1=\begin{pmatrix}1&0\\0&1\end{pmatrix},\quad P=\begin{pmatrix}0&1\\1&0\end{pmatrix},\quad N=\begin{pmatrix}0&1\\-1&0\end{pmatrix},\quad PN=\begin{pmatrix}-1&0\\0&1\end{pmatrix}=-G,
\]

with $G:=\sigma_z=\mathrm{diag}(1,-1)$. Rule (R2) predicts $(PN)^2=(-1)^{1}\eta^{00}\eta^{11}=(-1)(+1)(-1)=+1$, and indeed $G^2=I$. Every real 2 by 2 matrix is a combination of them:

\[
\begin{pmatrix}\alpha&\beta\\ \gamma&\delta\end{pmatrix}=\tfrac{\alpha+\delta}2\,1+\tfrac{\beta+\gamma}2\,P+\tfrac{\beta-\gamma}2\,N+\tfrac{\delta-\alpha}2\,PN .
\]

(Check the four entries.) The four monomials are also independent: $c_0\,1+c_1P+c_2N+c_3PN=\begin{pmatrix}c_0-c_3&c_1+c_2\\ c_1-c_2&c_0+c_3\end{pmatrix}$ vanishes only for $c_0=c_1=c_2=c_3=0$. They are therefore a basis of $\mathrm{Mat}_2(\mathbb R)$, whose dimension is $4=2^2$, and the real linear combinations of products of $P$ and $N$ are exactly the real 2 by 2 matrices: $\mathrm{Cl}(1,1)\cong\mathrm{Mat}_2(\mathbb R)$. The symbol $\cong$ ("is isomorphic to") says that two algebras (sets in which one can add, multiply by numbers and multiply) are the same up to renaming: there is a one-to-one and onto map between them that respects sums, multiplication by numbers and products. Here, with concrete matrices as generators, the map is simply the identity. The rest of the chapter proves the same statement one size up, $\mathrm{Cl}(4,4)\cong\mathrm{Mat}_{16}(\mathbb R)$, where $256=16^2$.

\textbf{The three matrices $P$, $N$, $G$.} We record the facts about them used below. They are real, $P$ and $G$ are symmetric, $N$ is antisymmetric, $P^2=G^2=I$, $N^2=-I$, and any two of them anticommute:

\[
PN=-G=-NP,\qquad PG=-N=-GP,\qquad NG=-P=-GN .
\]

(Each product is a 2 by 2 multiplication; for example $PG=\begin{pmatrix}0&1\\1&0\end{pmatrix}\begin{pmatrix}1&0\\0&-1\end{pmatrix}=\begin{pmatrix}0&-1\\1&0\end{pmatrix}=-N$.)

\subsection{How big must the matrices be? The trace lemma}

Suppose $\gamma^0,\dots,\gamma^{n-1}$ are $d$ by $d$ matrices (real or complex) that satisfy the Clifford relation for some signature $(p,q)$, and suppose $n=p+q$ is \textbf{even}. This section proves that the $2^n$ monomials are linearly independent. Since at most $d^2$ matrices of size $d$ can be independent, this forces $d^2\ge2^n$.

\textbf{Theorem 2.1 (trace lemma).} If $n$ is even, every monomial $\gamma_A$ with $A\ne\varnothing$ has trace zero.

\emph{Proof.} Let $k=|A|\ge1$. We find a generator $\gamma^b$ that anticommutes with $\gamma_A$.

\begin{itemize}
\item If $k$ is odd, then $k\le n-1$ (because $n$ is even), so some index $b$ is not in $A$. By (R3), $\gamma^b\gamma_A=(-1)^k\gamma_A\gamma^b=-\gamma_A\gamma^b$.
\item If $k$ is even, pick any $b\in A$. By (R3), $\gamma^b\gamma_A=(-1)^{k-1}\gamma_A\gamma^b=-\gamma_A\gamma^b$.
\end{itemize}

In both cases $\gamma^b$ is invertible ($(\gamma^b)^{-1}=\eta^{bb}\gamma^b$, since $(\gamma^b)^2=\eta^{bb}$) and $\gamma_A=-(\gamma^b)^{-1}\gamma_A\gamma^b$. Taking the trace and using the cyclic property,

\[
\operatorname{tr}\gamma_A=-\operatorname{tr}\bigl((\gamma^b)^{-1}\gamma_A\gamma^b\bigr)=-\operatorname{tr}\gamma_A,
\]

so $\operatorname{tr}\gamma_A=0$. $\square$

(For odd $n$ the lemma fails for the top monomial: the Pauli product $\sigma_x\sigma_y\sigma_z=iI$ has trace $2i$.)

\textbf{Theorem 2.2 (independence of the monomials).} If $n$ is even, the $2^n$ monomials are linearly independent. Consequently $d^2\ge2^n$, that is, $d\ge2^{n/2}$.

\emph{Proof.} Suppose $\sum_Bc_B\gamma_B=0$ with numbers $c_B$. Fix a set $A$, multiply on the left by $\gamma_A^{-1}=\pm\gamma_A$ and take the trace:

\[
0=\sum_Bc_B\operatorname{tr}\bigl(\gamma_A^{-1}\gamma_B\bigr).
\]

By (R1), $\gamma_A^{-1}\gamma_B=\pm\gamma_{A\triangle B}$, and $A\triangle B=\varnothing$ exactly when $B=A$. For $B\ne A$ the trace vanishes by Theorem 2.1; for $B=A$ the product is $\gamma_A^{-1}\gamma_A=1$, with trace $d$. So $0=d\,c_A$, and $c_A=0$. Since $A$ was arbitrary, all coefficients vanish. $\square$

\textbf{Consequences for this book.}

\begin{itemize}
\item \emph{Dirac:} $n=4$ gives $d\ge4$. Dirac's 4 by 4 matrices are the smallest possible.
\item \emph{4+4 dimensions:} $n=8$ gives $d\ge16$. \textbf{A spinor in eight dimensions has at least 16 components}, and dirac16complex has exactly the minimum.
\item \emph{The full matrix algebra.} If $n=8$ and $d=16$, the 256 monomials are 256 independent elements of $\mathrm{Mat}_{16}(\mathbb R)$ (or of $\mathrm{Mat}_{16}(\mathbb C)$), whose dimension is $16^2=256$. Hence they are a \textbf{basis}: every 16 by 16 matrix is a unique linear combination of monomials. For real gammas this says $\mathrm{Cl}(4,4)\cong\mathrm{Mat}_{16}(\mathbb R)$.
\item \emph{The even and the odd monomials} are independent too (they are part of the 256), so the even part of the algebra has dimension 128 and the odd part has dimension 128, and a matrix that is a combination of even monomials and also a combination of odd monomials is zero.
\end{itemize}

\textbf{Corollary 2.3 (matrices that commute with all gammas).} Let $\gamma^0,\dots,\gamma^7$ be 16 by 16 matrices with the Clifford relation of signature (4,4). If a matrix $X$ (real or complex) satisfies $X\gamma^a=\gamma^aX$ for all $a$, then $X=c\,I$ for a number $c$.

\emph{Proof.} $X$ commutes with every product of gammas, hence with every monomial, hence (by the basis property) with every matrix, in particular with every elementary matrix $E_{ij}$. Compare the entries in position $(k,j)$ of $XE_{ij}$ and $E_{ij}X$: $(XE_{ij})_{kj}=X_{ki}$ and $(E_{ij}X)_{kj}=\delta_{ki}X_{jj}$, where $\delta_{ki}$ is 1 for $k=i$ and 0 otherwise. So $X_{ki}=0$ for $k\ne i$, and $X_{ii}=X_{jj}$ for all $i,j$: $X$ is a multiple of $I$. $\square$

In the language of Section 2.13, the \textbf{commutant} of the gammas consists of the scalar matrices. The repository computes this commutant directly (Section 2.15).

\subsection{A 16 by 16 construction: the tensor-product picture}

We now show that real 16 by 16 gammas for signature (4,4) exist, by an explicit recipe that uses only the matrices $P$, $N$ and $G$ of Section 2.5. It is the "Clifford picture" of the reference implementation dirac-main (Stage 1, §3.3). For $k=1,2,3,4$ (a 1-based slot label, one of the labelled exceptions to counting from 0) define

\[
\hat\gamma_k^+:=\underbrace{G\otimes\dots\otimes G}_{k-1}\otimes P\otimes\underbrace{I_2\otimes\dots\otimes I_2}_{4-k},\qquad\hat\gamma_k^-:=\underbrace{G\otimes\dots\otimes G}_{k-1}\otimes N\otimes\underbrace{I_2\otimes\dots\otimes I_2}_{4-k},
\]

each a Kronecker product of four 2 by 2 matrices, hence 16 by 16. Order them as frame matrices:

\[
(\hat\gamma^0,\hat\gamma^1,\hat\gamma^2,\hat\gamma^3,\hat\gamma^4,\hat\gamma^5,\hat\gamma^6,\hat\gamma^7):=(\hat\gamma_1^+,\hat\gamma_2^+,\hat\gamma_3^+,\hat\gamma_4^+,\hat\gamma_1^-,\hat\gamma_2^-,\hat\gamma_3^-,\hat\gamma_4^-).
\]

For example $\hat\gamma^0=P\otimes I_2\otimes I_2\otimes I_2$, $\hat\gamma^1=G\otimes P\otimes I_2\otimes I_2$ and $\hat\gamma^4=N\otimes I_2\otimes I_2\otimes I_2$.

\textbf{Theorem 2.4.} The $\hat\gamma^a$ are real 16 by 16 signed permutation matrices that satisfy $\{\hat\gamma^a,\hat\gamma^b\}=2\eta^{ab}I_{16}$ with $\eta=\mathrm{diag}(+1,+1,+1,+1,-1,-1,-1,-1)$.

\emph{Proof.} A Kronecker product of signed permutation matrices is a signed permutation matrix: each row of $A\otimes B$ combines one row of $A$ with one row of $B$ and so contains exactly one product of nonzero entries, and likewise for columns. By the mixed-product rule, products of these matrices are computed slot by slot.

\emph{Squares.} In $(\hat\gamma_k^\pm)^2$ every slot is squared: $G^2=I$, $I^2=I$, $P^2=I$ and $N^2=-I$. So $(\hat\gamma_k^+)^2=+I_{16}$ and $(\hat\gamma_k^-)^2=-I_{16}$, as required for $a\le3$ and $a\ge4$.

\emph{Same slot, $\hat\gamma_k^+$ and $\hat\gamma_k^-$.} The two products $\hat\gamma_k^+\hat\gamma_k^-$ and $\hat\gamma_k^-\hat\gamma_k^+$ agree in every slot except slot $k$, which is $PN$ in the first and $NP=-PN$ in the second. So they anticommute.

\emph{Different slots $k<l$.} Take one matrix from slot $k$ (with $X=P$ or $N$ in slot $k$) and one from slot $l$ (with $Y=P$ or $N$ in slot $l$). Compare the two orders of multiplication slot by slot: in the slots before $k$ both matrices have $G$, and $GG=GG$; in slot $k$ the first has $X$ and the second $G$, and $XG=-GX$; in the slots strictly between $k$ and $l$ the first has $I_2$ and the second $G$, which commute; in slot $l$ the first has $I_2$ and the second $Y$, which commute; in the slots after $l$ both have $I_2$. Exactly one slot contributes a minus sign, so the two matrices anticommute. $\square$

\textbf{The chirality of this picture.} Define $\hat\gamma^8:=\hat\gamma^0\hat\gamma^1\cdots\hat\gamma^7$. Move $\hat\gamma^4=\hat\gamma_1^-$ to the right of $\hat\gamma^0=\hat\gamma_1^+$ (past three different factors, sign $(-1)^3$), then $\hat\gamma_2^-$ to the right of $\hat\gamma_2^+$ (two factors, $(-1)^2$), then $\hat\gamma_3^-$ to the right of $\hat\gamma_3^+$ (one factor, $-1$). The total sign is $(-1)^6=+1$, and

\[
\hat\gamma^8=(\hat\gamma_1^+\hat\gamma_1^-)(\hat\gamma_2^+\hat\gamma_2^-)(\hat\gamma_3^+\hat\gamma_3^-)(\hat\gamma_4^+\hat\gamma_4^-).
\]

Slot by slot, $\hat\gamma_k^+\hat\gamma_k^-$ has $G^2=I$ before slot $k$, $PN=-G$ in slot $k$ and $I_2$ after it. The product of the four pairs is therefore

\[
\hat\gamma^8=(-G)\otimes(-G)\otimes(-G)\otimes(-G)=G\otimes G\otimes G\otimes G .
\]

This is a diagonal matrix. Number the 16 basis columns by $j=8b_0+4b_1+2b_2+b_3$ with bits $b_0,\dots,b_3\in\{0,1\}$, one bit per slot (the first slot is the most significant, as in the index formula of the Kronecker product). Since $G=\mathrm{diag}(+1,-1)$, the diagonal entry of $\hat\gamma^8$ in position $j$ is $(-1)^{b_0+b_1+b_2+b_3}$: it is $-1$ exactly for the eight numbers $j\in\{1,2,4,7,8,11,13,14\}$, which have an odd number of 1-bits, and $+1$ for $j\in\{0,3,5,6,9,10,12,15\}$. These two lists reappear in Chapter 3.

\textbf{The charge matrix of this picture} is $C_{\mathrm{dm}}:=\hat\gamma^0\hat\gamma^1\hat\gamma^2\hat\gamma^3$. Slot by slot, with $G^2=I$,

\[
C_{\mathrm{dm}}=(PGGG)\otimes(PGG)\otimes(PG)\otimes P=(PG)\otimes P\otimes(PG)\otimes P=N\otimes P\otimes N\otimes P,
\]

using $PG=-N$ twice. The subscript dm stands for dirac-main. The repository checks that these tensor matrices satisfy the Clifford relation (check \texttt{ALG\_clif\allowbreak{}fordPict\allowbreak{}ureInter\allowbreak{}twiner}); the comparison with dirac-main's committed generators is made by the Wolfram verifier (measurement \texttt{ALG\_clif\allowbreak{}fordPict\allowbreak{}ureInter\allowbreak{}twiner.m\allowbreak{}atchesDi\allowbreak{}racMainC\allowbreak{}l44SeedG\allowbreak{}enerator\allowbreak{}s} of \texttt{wolfram-\allowbreak{}algebra-\allowbreak{}report.j\allowbreak{}son}).

\subsection{The notebook's gamma matrices \texorpdfstring{$\mathrm{T16}^A$}{T16\textasciicircum{}A}}

The author's notebook uses a different set of real 16 by 16 matrices, called $(\mathrm{T16}^A)[a]$ there. In this book they are written $\gamma^a$ without a hat, and they are the \textbf{primary gamma matrices} of the whole project (contract, §1). The notebook builds them in \textbf{block form} from eight real 8 by 8 matrices $\tau_a$ and eight partners $\bar\tau_a$:

\[
\gamma^a=\begin{pmatrix}0&\bar\tau_a\\ \tau_a&0\end{pmatrix}\qquad(a=0,\dots,7).
\]

Chapter 3 constructs the $\tau_a$ from 4 by 4 building blocks and explains their origin in the split octonions. For now we only need the finished matrices. Each $\gamma^a$ is a signed permutation matrix, so it can be written down completely in a small table: for each row $i$ the table gives the column $j$ of the single nonzero entry and its sign. The entry $-13$ in row 0 of the column $\gamma^4$ means $(\gamma^4)_{0,13}=-1$, that is, $(\gamma^4u)_0=-u_{13}$ for every column $u$; the entry $+0$ means $+u_0$, and $-0$ means $-u_0$.

\begingroup
\small
\begin{longtable}{@{}>{\raggedright\arraybackslash}p{0.160\linewidth}>{\raggedright\arraybackslash}p{0.160\linewidth}>{\raggedright\arraybackslash}p{0.160\linewidth}>{\raggedright\arraybackslash}p{0.160\linewidth}>{\raggedright\arraybackslash}p{0.160\linewidth}@{}}
\toprule
\textbf{$i$} & \textbf{$\gamma^0$} & \textbf{$\gamma^1$} & \textbf{$\gamma^2$} & \textbf{$\gamma^3$} \\
\midrule
\endfirsthead
\toprule
\textbf{$i$} & \textbf{$\gamma^0$} & \textbf{$\gamma^1$} & \textbf{$\gamma^2$} & \textbf{$\gamma^3$} \\
\midrule
\endhead
0 & $+8$ & $-15$ & $+14$ & $-13$ \\
1 & $+9$ & $-14$ & $-15$ & $+12$ \\
2 & $+10$ & $+13$ & $-12$ & $-15$ \\
3 & $+11$ & $+12$ & $+13$ & $+14$ \\
4 & $+12$ & $-11$ & $+10$ & $-9$ \\
5 & $+13$ & $-10$ & $-11$ & $+8$ \\
6 & $+14$ & $+9$ & $-8$ & $-11$ \\
7 & $+15$ & $+8$ & $+9$ & $+10$ \\
8 & $+0$ & $+7$ & $-6$ & $+5$ \\
9 & $+1$ & $+6$ & $+7$ & $-4$ \\
10 & $+2$ & $-5$ & $+4$ & $+7$ \\
11 & $+3$ & $-4$ & $-5$ & $-6$ \\
12 & $+4$ & $+3$ & $-2$ & $+1$ \\
13 & $+5$ & $+2$ & $+3$ & $-0$ \\
14 & $+6$ & $-1$ & $+0$ & $+3$ \\
15 & $+7$ & $-0$ & $-1$ & $-2$ \\
\bottomrule
\end{longtable}
\endgroup

\begingroup
\small
\begin{longtable}{@{}>{\raggedright\arraybackslash}p{0.160\linewidth}>{\raggedright\arraybackslash}p{0.160\linewidth}>{\raggedright\arraybackslash}p{0.160\linewidth}>{\raggedright\arraybackslash}p{0.160\linewidth}>{\raggedright\arraybackslash}p{0.160\linewidth}@{}}
\toprule
\textbf{$i$} & \textbf{$\gamma^4$} & \textbf{$\gamma^5$} & \textbf{$\gamma^6$} & \textbf{$\gamma^7$} \\
\midrule
\endfirsthead
\toprule
\textbf{$i$} & \textbf{$\gamma^4$} & \textbf{$\gamma^5$} & \textbf{$\gamma^6$} & \textbf{$\gamma^7$} \\
\midrule
\endhead
0 & $-13$ & $+14$ & $+15$ & $+8$ \\
1 & $+12$ & $+15$ & $-14$ & $+9$ \\
2 & $+15$ & $-12$ & $+13$ & $+10$ \\
3 & $-14$ & $-13$ & $-12$ & $+11$ \\
4 & $+9$ & $-10$ & $-11$ & $-12$ \\
5 & $-8$ & $-11$ & $+10$ & $-13$ \\
6 & $-11$ & $+8$ & $-9$ & $-14$ \\
7 & $+10$ & $+9$ & $+8$ & $-15$ \\
8 & $+5$ & $-6$ & $-7$ & $-0$ \\
9 & $-4$ & $-7$ & $+6$ & $-1$ \\
10 & $-7$ & $+4$ & $-5$ & $-2$ \\
11 & $+6$ & $+5$ & $+4$ & $-3$ \\
12 & $-1$ & $+2$ & $+3$ & $+4$ \\
13 & $+0$ & $+3$ & $-2$ & $+5$ \\
14 & $+3$ & $-0$ & $+1$ & $+6$ \\
15 & $-2$ & $-1$ & $-0$ & $+7$ \\
\bottomrule
\end{longtable}
\endgroup

The tables were generated from the function \texttt{notebook\allowbreak{}\_gammas} of \texttt{scripts/\allowbreak{}d16c\_exa\allowbreak{}ct.py}; the same matrices are stored exactly under the key \texttt{\detokenize{gamma}} of the data file \texttt{algebra-\allowbreak{}fixture.\allowbreak{}json}. The Wolfram verifier also rebuilds them from the 11 verbatim definition cells of the notebook and finds them identical (within check \texttt{ALG\_expr\allowbreak{}ession1}). Every row of every table points into the other half: rows 0 to 7 point to columns 8 to 15 and rows 8 to 15 to columns 0 to 7. That is the block form: the upper-left and lower-right blocks are zero.

\textbf{Three checks by hand.}

\begin{itemize}
\item $(\gamma^0)^2=+1$ in row 0: $(\gamma^0u)_0=u_8$ and row 8 of $\gamma^0$ says $(\gamma^0v)_8=v_0$, so $(\gamma^0\gamma^0u)_0=(\gamma^0u)_8=u_0$.
\item $(\gamma^4)^2=-1$ in row 0: $(\gamma^4u)_0=-u_{13}$ and row 13 of $\gamma^4$ says $(\gamma^4v)_{13}=+v_0$, so $(\gamma^4\gamma^4u)_0=-(\gamma^4u)_{13}=-u_0$.
\item $\gamma^0\gamma^4=-\gamma^4\gamma^0$ in row 0: $(\gamma^0\gamma^4u)_0=(\gamma^4u)_8=+u_5$ (row 8 of $\gamma^4$), while $(\gamma^4\gamma^0u)_0=-(\gamma^0u)_{13}=-u_5$ (row 13 of $\gamma^0$). The two differ by a sign.
\end{itemize}

\textbf{Facts about the notebook's gammas.}

(N1) \emph{They satisfy the Clifford relation.} $\{\gamma^a,\gamma^b\}=2\eta^{ab}I_{16}$ for all 64 ordered pairs $(a,b)$, with squares $+1,+1,+1,+1,-1,-1,-1,-1$. In this chapter we take this from the exact machine check, which tests all 64 pairs (Stage 1, Result 3.1; check \texttt{ALG\_clif\allowbreak{}ford}). Chapter 3 proves it by hand, reducing it to a few 4 by 4 multiplications. The reader can also verify any single row with the tables, as above.

(N2) \emph{They are orthogonal.} Each $\gamma^a$ is a signed permutation matrix, hence orthogonal (Section 2.2): $(\gamma^a)^T=(\gamma^a)^{-1}$.

(N3) \emph{$\gamma^a$ is symmetric for $a\le3$ and antisymmetric for $a\ge4$.} \emph{Proof.} By (N1), $(\gamma^a)^2=\eta^{aa}$, so $(\gamma^a)^{-1}=\eta^{aa}\gamma^a$. By (N2), $(\gamma^a)^T=(\gamma^a)^{-1}=\eta^{aa}\gamma^a$. $\square$ (Stage 1, Result 3.1; check \texttt{ALG\_gamm\allowbreak{}aTranspo\allowbreak{}seSymmet\allowbreak{}ry}.) The tensor matrices of Section 2.7 have the same pattern, since $P,G,I_2$ are symmetric and $N$ is antisymmetric.

(N4) \emph{Their 256 monomials are a basis of $\mathrm{Mat}_{16}(\mathbb R)$.} This follows from Theorem 2.2. The computer confirms it independently: the 256 monomials, written as rows of 256 numbers, have rank 256, and the 128 even ones have rank 128 (Stage 1, Result 3.2; check \texttt{ALG\_fait\allowbreak{}hful}). The notebook itself tabulates the 256 ordered products of its gammas (its cell 594, according to the survey of the notebook \texttt{handoff/\allowbreak{}surveys/\allowbreak{}survey\_n\allowbreak{}otebook-\allowbreak{}physics.\allowbreak{}md}; we call it the notebook survey below).

\subsection{The charge matrix \texorpdfstring{$C=\sigma_{16}$}{C=σ\_16} and expression [1]}

\textbf{Definition.} The \textbf{charge matrix} is the product of the four space-like gammas,

\[
C:=\gamma^0\gamma^1\gamma^2\gamma^3 .
\]

In the notebook it is called σ16; the notebook itself identifies σ16 with the product of its first four gammas (cells 612 to 615, per the notebook survey). It is used to build the \textbf{Dirac adjoint} $\bar\Psi=\Psi^\dagger C$ of Chapter 6. We derive its properties from the Clifford relation and fact (N3) only, so they hold for the notebook's gammas and for the tensor picture alike.

\textbf{(C1) How $C$ commutes with the gammas:} $C\gamma^a=-\eta^{aa}\gamma^aC$. \emph{Proof.} $C$ is the monomial with $A=\{0,1,2,3\}$, of degree 4. By (R3): for $a\le3$ ($a\in A$), $\gamma^aC=(-1)^3C\gamma^a=-C\gamma^a$; for $a\ge4$ ($a\notin A$), $\gamma^aC=(-1)^4C\gamma^a=C\gamma^a$. Both cases are $C\gamma^a=-\eta^{aa}\gamma^aC$. $\square$

\textbf{(C2) $C$ is symmetric.} \emph{Proof.} $C^T=(\gamma^3)^T(\gamma^2)^T(\gamma^1)^T(\gamma^0)^T=\gamma^3\gamma^2\gamma^1\gamma^0$ by (N3). To restore the order, move $\gamma^0$ to the front (past three different factors), then $\gamma^1$ to the second place (two factors), then $\gamma^2$ to the third place (one factor): the sign is $(-1)^{3+2+1}=+1$. So $C^T=C$. $\square$

\textbf{(C3) $C^2=1$.} \emph{Proof.} $C^2=CC^T=\gamma^0\gamma^1\gamma^2\gamma^3\gamma^3\gamma^2\gamma^1\gamma^0$. The two $\gamma^3$ in the middle give $(\gamma^3)^2=1$, then the two $\gamma^2$, and so on: $C^2=1$. $\square$

\textbf{(C4) Expression [1].} The task statement of the project contains the following expression, quoted verbatim in WolframScript in Stage 1, §3.2:

\begin{Verbatim}[fontsize=\small]
\[Sigma]16.(T16^A)[#]==-Transpose[\[Sigma]16.(T16^A)[#]]&/@Range[0,7]
\end{Verbatim}

It asks whether $(C\gamma^a)^T=-C\gamma^a$ for $a=0,\dots,7$. \emph{Proof that it holds.} $(C\gamma^a)^T=(\gamma^a)^TC^T=\eta^{aa}\gamma^aC$ by (N3) and (C2). By (C1), $\gamma^aC=-\eta^{aa}C\gamma^a$. Hence $(C\gamma^a)^T=\eta^{aa}(-\eta^{aa})C\gamma^a=-C\gamma^a$, because $(\eta^{aa})^2=1$. $\square$ The Wolfram verifier evaluates the expression literally and obtains \{True, True, True, True, True, True, True, True\} (Stage 1, Result 3.4; check \texttt{ALG\_expr\allowbreak{}ession1}, measurement \texttt{ALG\_expr\allowbreak{}ession1.\allowbreak{}listWolf\allowbreak{}ramForm}). Expression [1] is the reason why the notebook's Lagrangian carries no dynamics for a real Grassmann field (Chapter 5).

\textbf{(C5) The signature of $C$ is (8,8).} \emph{Proof.} $C$ is a monomial of degree 4, so $\operatorname{tr}C=0$ by Theorem 2.1. With $C^2=1$ the projector argument of Section 2.2 gives eigenvalue $+1$ eight times and $-1$ eight times; $C$ is real symmetric, so its signature is (8,8). $\square$ (Stage 1, Result 3.3; check \texttt{ALG\_char\allowbreak{}geMatrix}; the Wolfram report also records the characteristic polynomial $(x-1)^8(x+1)^8$.)

\textbf{(C6) The notebook's $C$.} For the notebook's gammas the product is the block matrix

\[
C=\begin{pmatrix}-\sigma&0\\0&\sigma\end{pmatrix},\qquad\sigma=\begin{pmatrix}0&I_4\\ I_4&0\end{pmatrix},
\]

a signed permutation matrix (part of check \texttt{ALG\_char\allowbreak{}geMatrix}; Chapter 3 derives it by hand). Here it is as a table, together with the matrix $B$ of Section 2.12:

\begingroup
\small
\begin{longtable}{@{}>{\raggedright\arraybackslash}p{0.267\linewidth}>{\raggedright\arraybackslash}p{0.267\linewidth}>{\raggedright\arraybackslash}p{0.267\linewidth}@{}}
\toprule
\textbf{$i$} & \textbf{$C$} & \textbf{$B/i$} \\
\midrule
\endfirsthead
\toprule
\textbf{$i$} & \textbf{$C$} & \textbf{$B/i$} \\
\midrule
\endhead
0 & $-4$ & $+9$ \\
1 & $-5$ & $-8$ \\
2 & $-6$ & $-11$ \\
3 & $-7$ & $+10$ \\
4 & $-0$ & $-13$ \\
5 & $-1$ & $+12$ \\
6 & $-2$ & $+15$ \\
7 & $-3$ & $-14$ \\
8 & $+12$ & $+1$ \\
9 & $+13$ & $-0$ \\
10 & $+14$ & $-3$ \\
11 & $+15$ & $+2$ \\
12 & $+8$ & $-5$ \\
13 & $+9$ & $+4$ \\
14 & $+10$ & $+7$ \\
15 & $+11$ & $-6$ \\
\bottomrule
\end{longtable}
\endgroup

\textbf{Worked example: $C$ defines an indefinite form.} For a real column $u$ the number $u^TCu=\sum_{ij}u_iC_{ij}u_j$ can have either sign. With $u=e_0+e_4$ (components 0 and 4 equal to 1, all others 0) only the entries $C_{04}=C_{40}=-1$ contribute, and $u^TCu=-2$. With $u=e_8+e_{12}$ only $C_{8,12}=C_{12,8}=+1$ contribute, and $u^TCu=+2$. This is the matrix form of signature (8,8); it is why energies built with $C$ in Chapter 7 have no fixed sign.

\subsection{The chirality \texorpdfstring{$\gamma^8$}{γ\textasciicircum{}8} and the two halves}

\textbf{Definition.} The \textbf{chirality} (or volume element) is the product of all eight gammas in increasing order,

\[
\gamma^8:=\gamma^0\gamma^1\gamma^2\gamma^3\gamma^4\gamma^5\gamma^6\gamma^7 .
\]

The label 8 is a name, not a ninth direction. In the notebook it is T16A[8] (cells 620 and 621, per the notebook survey).

\textbf{(X1) $(\gamma^8)^2=1$.} By (R2) with $k=8$: $(-1)^{8\cdot7/2}\prod_{a=0}^7\eta^{aa}=(-1)^{28}(+1)^4(-1)^4=+1$.

\textbf{(X2) $\gamma^8$ anticommutes with every $\gamma^a$.} By (R3) with $k=8$ and $a\in A$: $\gamma^a\gamma^8=(-1)^7\gamma^8\gamma^a=-\gamma^8\gamma^a$.

\textbf{(X3) $\gamma^8$ commutes with every even monomial and anticommutes with every odd one.} Passing $\gamma^8$ through a monomial of degree $k$ passes $k$ gammas, each giving $-1$ by (X2): the sign is $(-1)^k$. In particular $\gamma^8$ commutes with $C$ (degree 4) and with the spin generators $S^{ab}$ of Section 2.11 (degree 2).

\textbf{(X4) Eight $+1$'s and eight $-1$'s.} $\gamma^8$ is a monomial, so $\operatorname{tr}\gamma^8=0$ (Theorem 2.1); with (X1) and Section 2.2 it has the eigenvalue $+1$ eight times and $-1$ eight times.

\textbf{(X5) The notebook's chirality is diagonal:}

\[
\gamma^8=\begin{pmatrix}-I_8&0\\0&I_8\end{pmatrix}.
\]

(Stage 1, Result 3.6; check \texttt{ALG\_chir\allowbreak{}ality}; Chapter 3 derives it by hand.) In the tensor picture the chirality is $G\otimes G\otimes G\otimes G$ (Section 2.7), also diagonal but with a different pattern of signs; the two are related by the change of basis of Chapter 3.

\textbf{Chiral projectors and the two halves.} By Section 2.2 the matrices

\[
P_-:=\tfrac12\bigl(1-\gamma^8\bigr),\qquad P_+:=\tfrac12\bigl(1+\gamma^8\bigr)
\]

are projectors with $P_-+P_+=1$ and $P_-P_+=0$. For the notebook's gammas, (X5) gives $P_-=\mathrm{diag}(I_8,0)$ and $P_+=\mathrm{diag}(0,I_8)$. So the 16 components of a spinor split into two halves of 8:

\begin{itemize}
\item the upper components $\Psi_0,\dots,\Psi_7$ (the notebook's Ψ16upper) have chirality $-1$: $\gamma^8\Psi=-\Psi$ if only they are nonzero;
\item the lower components $\Psi_8,\dots,\Psi_{15}$ (Ψ16lower) have chirality $+1$.
\end{itemize}

\textbf{(X6) Odd matrices exchange the halves, even matrices keep them.} From (X2), $\gamma^a(1-\gamma^8)=(1+\gamma^8)\gamma^a$, that is, $\gamma^aP_-=P_+\gamma^a$: a gamma maps the chirality $-1$ half into the chirality $+1$ half and vice versa. Even matrices commute with $\gamma^8$ (X3) and therefore map each half into itself.

This explains the block form of the notebook's gammas. Write any 16 by 16 matrix in blocks, $M=\begin{pmatrix}A&B\\ D&E\end{pmatrix}$. With $\gamma^8=\mathrm{diag}(-I_8,I_8)$,

\[
\gamma^8M=\begin{pmatrix}-A&-B\\ D&E\end{pmatrix},\qquad M\gamma^8=\begin{pmatrix}-A&B\\ -D&E\end{pmatrix}.
\]

$M$ anticommutes with $\gamma^8$ exactly when $A=E=0$ (block off-diagonal, like every $\gamma^a$), and commutes with it exactly when $B=D=0$ (block diagonal, like $C$ and every even monomial).

\subsection{The spin generators \texorpdfstring{$S^{ab}$}{S\textasciicircum{}ab} and the rotation of spinors}

\textbf{Definition.} For $a,b=0,\dots,7$ the \textbf{spin generators} are

\[
S^{ab}:=\tfrac14\bigl[\gamma^a,\gamma^b\bigr]=\tfrac14\bigl(\gamma^a\gamma^b-\gamma^b\gamma^a\bigr).
\]

In the notebook they are SAB (Section title of cell 430, per the notebook survey), with 1-based list indices: SAB[[a+1,b+1]] is $S^{ab}$.

\textbf{(S1) Simple facts.} $S^{ba}=-S^{ab}$ and $S^{aa}=0$. For $a\ne b$ the two gammas anticommute, so $[\gamma^a,\gamma^b]=2\gamma^a\gamma^b$ and $S^{ab}=\tfrac12\gamma^a\gamma^b$.

\textbf{(S2) There are 28 independent ones.} The $S^{ab}$ with $a<b$ are one half times the $\binom82=28$ different monomials of degree 2, which are independent by Theorem 2.2. Of the 28, twelve have both indices of the same kind (6 pairs inside $\{0,1,2,3\}$ and 6 inside $\{4,5,6,7\}$), and sixteen mix a space-like with a time-like index.

\textbf{(S3) The spin generators rotate the gammas:}

\[
\bigl[S^{ab},\gamma^c\bigr]=\eta^{bc}\gamma^a-\eta^{ac}\gamma^b .
\]

\emph{Proof.} For $a=b$ both sides vanish. For $a\ne b$ use the Clifford relation twice, first as $\gamma^b\gamma^c=2\eta^{bc}-\gamma^c\gamma^b$ and then as $\gamma^a\gamma^c=2\eta^{ac}-\gamma^c\gamma^a$:

\[
\gamma^a\gamma^b\gamma^c=2\eta^{bc}\gamma^a-\gamma^a\gamma^c\gamma^b=2\eta^{bc}\gamma^a-2\eta^{ac}\gamma^b+\gamma^c\gamma^a\gamma^b .
\]

Hence $[\gamma^a\gamma^b,\gamma^c]=\gamma^a\gamma^b\gamma^c-\gamma^c\gamma^a\gamma^b=2\eta^{bc}\gamma^a-2\eta^{ac}\gamma^b$, and dividing by 2 gives the claim. $\square$ (Stage 1, Result 3.5. The exact Python geometry checker \texttt{scripts/\allowbreak{}check\_di\allowbreak{}rac16com\allowbreak{}plex\_geo\allowbreak{}metry.py} tests the relation for all 512 triples $(a,b,c)$ as check \texttt{ALG\_SabG\allowbreak{}ammaComm\allowbreak{}utator}, recorded in \texttt{python-g\allowbreak{}eometry-\allowbreak{}report.j\allowbreak{}son}; the Wolfram verifier records it, as true, in the measurement \texttt{suppleme\allowbreak{}ntarySpi\allowbreak{}nVectorR\allowbreak{}elationH\allowbreak{}olds} of the check \texttt{ALG\_spin\allowbreak{}Transpos\allowbreak{}ePropert\allowbreak{}ies} in \texttt{wolfram-\allowbreak{}algebra-\allowbreak{}report.j\allowbreak{}son}, a measurement that is not part of the pass condition of that check.)

For a \textbf{Clifford vector} $v:=\sum_cv_c\gamma^c$ with numbers $v_c$ this reads $[S^{ab},v]=\eta^{bb}v_b\gamma^a-\eta^{aa}v_a\gamma^b$. For example $[S^{01},v]=v_1\gamma^0-v_0\gamma^1$: the components $(v_0,v_1)$ are turned a little in their plane and nothing else changes. The spin generators are the "infinitesimal rotations" of spinors.

\textbf{(S4) Finite transformations: the exponential.} For a square matrix $M$ define

\[
\exp(M):=\sum_{k=0}^\infty\frac{M^k}{k!}=1+M+\tfrac12M^2+\tfrac16M^3+\dots
\]

(the series converges for every square matrix, like the series of $e^x$; we quote this from analysis). If a matrix $J$ satisfies $J^2=-1$, then $J^{2m}=(-1)^m$ and $J^{2m+1}=(-1)^mJ$; splitting the series into even and odd powers and using the power series of cosine and sine (quoted from calculus),

\[
\exp(\varphi J)=\sum_m\frac{(-1)^m\varphi^{2m}}{(2m)!}+J\sum_m\frac{(-1)^m\varphi^{2m+1}}{(2m+1)!}=\cos\varphi+J\sin\varphi .
\]

If instead $J^2=+1$, the same steps with the series of $\cosh\varphi=\tfrac12(e^\varphi+e^{-\varphi})$ and $\sinh\varphi=\tfrac12(e^\varphi-e^{-\varphi})$ give $\exp(\varphi J)=\cosh\varphi+J\sinh\varphi$.

For $a\ne b$ put $J:=\gamma^a\gamma^b=2S^{ab}$. By (R2), $J^2=-\eta^{aa}\eta^{bb}$. So $\theta S^{ab}=\tfrac\theta2J$ and

\[
\exp(\theta S^{ab})=\begin{cases}\cos\tfrac\theta2+\gamma^a\gamma^b\sin\tfrac\theta2&\text{if }\eta^{aa}\eta^{bb}=+1\ \text{(a rotation)},\\[2pt] \cosh\tfrac\theta2+\gamma^a\gamma^b\sinh\tfrac\theta2&\text{if }\eta^{aa}\eta^{bb}=-1\ \text{(a boost)}.\end{cases}
\]

In both cases the inverse is obtained by $\theta\to-\theta$: $(c+sJ)(c-sJ)=c^2-s^2J^2$, which is $\cos^2+\sin^2=1$ for a rotation and $\cosh^2-\sinh^2=1$ for a boost.

\textbf{Worked example: a rotation by $2\pi$ gives $-1$.} Let $R(\theta)=\exp(\theta S^{01})=c+sJ$ with $J=\gamma^0\gamma^1$, $c=\cos\tfrac\theta2$, $s=\sin\tfrac\theta2$. By (R3), $J$ anticommutes with $\gamma^0$ and $\gamma^1$ and commutes with $\gamma^2,\dots,\gamma^7$. Therefore $R\gamma^0=(c+sJ)\gamma^0=\gamma^0(c-sJ)=\gamma^0R^{-1}$, and

\[
R\gamma^0R^{-1}=\gamma^0R^{-2}=\gamma^0(c-sJ)^2=\gamma^0\bigl(c^2-s^2-2csJ\bigr)=\gamma^0\bigl(\cos\theta-\sin\theta\,\gamma^0\gamma^1\bigr)=\cos\theta\,\gamma^0-\sin\theta\,\gamma^1 .
\]

Here $(c-sJ)^2=c^2-2csJ+s^2J^2=c^2-s^2-2csJ$ because $J^2=-1$, and $c^2-s^2=\cos\theta$ and $2cs=\sin\theta$ by the addition theorems of Chapter 1, Section 1.1, with $\alpha=\beta=\theta/2$. In the same way $R\gamma^1R^{-1}=\gamma^1(\cos\theta-\sin\theta\,\gamma^0\gamma^1)=\cos\theta\,\gamma^1+\sin\theta\,\gamma^0$ (because $\gamma^1\gamma^0\gamma^1=-\gamma^0\gamma^1\gamma^1=-\gamma^0$), and $R\gamma^cR^{-1}=\gamma^c$ for $c\ge2$. So conjugation by $R(\theta)$ rotates the vector directions 0 and 1 by the angle $\theta$. But $R$ itself contains only half angles. At $\theta=2\pi$ the vectors are back where they started ($\cos2\pi=1$, $\sin2\pi=0$), while

\[
R(2\pi)=\cos\pi+\gamma^0\gamma^1\sin\pi=-1 .
\]

A full turn multiplies every spinor by $-1$; only a turn by $4\pi$ brings it back. This is the defining property of spinors, and it is why the spinor groups of Section 2.14 are "double covers" of the vector groups.

\textbf{Worked example: a boost.} Let $R=\exp(\theta S^{04})=c+sJ$ with $J=\gamma^0\gamma^4$, $c=\cosh\tfrac\theta2$, $s=\sinh\tfrac\theta2$, and $R^{-1}=c-sJ$. By (R3), $J$ anticommutes with $\gamma^0$ and $\gamma^4$, so, as for the rotation, $R\gamma^0R^{-1}=\gamma^0(c-sJ)^2$ and $R\gamma^4R^{-1}=\gamma^4(c-sJ)^2$. Now $J^2=+1$, so

\[
(c-sJ)^2=c^2-2csJ+s^2J^2=c^2+s^2-2csJ=\cosh\theta-\sinh\theta\,J ,
\]

because, from the definitions of Chapter 1, Section 1.1, $\cosh^2x+\sinh^2x=\tfrac14\bigl[(e^{2x}+2+e^{-2x})+(e^{2x}-2+e^{-2x})\bigr]=\tfrac12(e^{2x}+e^{-2x})=\cosh2x$ and $2\sinh x\cosh x=\tfrac12(e^x-e^{-x})(e^x+e^{-x})=\tfrac12(e^{2x}-e^{-2x})=\sinh2x$, here with $x=\theta/2$. Using $(\gamma^0)^2=1$ and $\gamma^4\gamma^0\gamma^4=-\gamma^0\gamma^4\gamma^4=\gamma^0$,

\[
R\gamma^0R^{-1}=\cosh\theta\,\gamma^0-\sinh\theta\,\gamma^4,\qquad R\gamma^4R^{-1}=\cosh\theta\,\gamma^4-\sinh\theta\,\gamma^0 .
\]

The pair of directions (0,4) is mixed by a hyperbolic rotation, which preserves $v_0^2-v_4^2$ because $\cosh^2\theta-\sinh^2\theta=1$. This is a Lorentz boost between the space direction $x_0$ and the time $x_4$.

\textbf{(S5) A transpose property.} $(CS^{ab})^T=-CS^{ab}$ for all $a,b$. \emph{Proof.} For $a=b$ both sides vanish. For $a\ne b$, by (N3) and (C2),

\[
(CS^{ab})^T=\tfrac12(\gamma^b)^T(\gamma^a)^TC^T=\tfrac12\eta^{aa}\eta^{bb}\gamma^b\gamma^aC=-\tfrac12\eta^{aa}\eta^{bb}\gamma^a\gamma^bC .
\]

From (C1), $\gamma^aC=-\eta^{aa}C\gamma^a$ (multiply (C1) by $-\eta^{aa}$). Using it twice, $\gamma^a\gamma^bC=\eta^{aa}\eta^{bb}C\gamma^a\gamma^b$. So $(CS^{ab})^T=-\tfrac12(\eta^{aa}\eta^{bb})^2C\gamma^a\gamma^b=-CS^{ab}$. $\square$ Two related statements, that $C\{\gamma^c,S^{ab}\}$ is symmetric and $C[\gamma^c,S^{ab}]$ antisymmetric, are Exercise 2.8. The Wolfram verifier checks the first property for all 64 pairs and the other two for all 512 triples; the Python checker counts 476 verified identities of these kinds for $a<b$ (Stage 1, Result 3.5; check \texttt{ALG\_spin\allowbreak{}Transpos\allowbreak{}ePropert\allowbreak{}ies}).

\textbf{Consequence: exponentials preserve $C$.} Let $R=\exp(\theta S^{ab})=c+sJ$ as in (S4). From (S5) and $C^T=C$, $J^TC=(CJ)^T=-CJ$. Hence $R^TC=cC+sJ^TC=C(c-sJ)=CR^{-1}$, and

\[
R^TCR=C .
\]

This is the matrix fact behind the invariance of the Lagrangian under rotations and boosts (Chapter 6); Section 2.14 generalizes it to every element of the spinor groups.

\subsection{The matrix \texorpdfstring{$B=-iC\gamma^4$}{B=-iCγ\textasciicircum{}4}}

\textbf{Definition.} $B:=-iC\gamma^4=-i\gamma^0\gamma^1\gamma^2\gamma^3\gamma^4$. It is $-i$ times a monomial of degree 5.

\textbf{(B1) $B$ is Hermitian.} $B^\dagger=(-i)^\ast(C\gamma^4)^T=i(\gamma^4)^TC^T=i(-\gamma^4)C=-i\gamma^4C$, using (N3) ($\gamma^4$ is antisymmetric) and (C2). By (C1) with $a=4$, $\gamma^4C=C\gamma^4$, so $B^\dagger=-iC\gamma^4=B$.

\textbf{(B2) $B^2=1$.} $B^2=(-i)^2C\gamma^4C\gamma^4=-CC\gamma^4\gamma^4=-(1)(-1)=1$, again using $\gamma^4C=C\gamma^4$.

\textbf{(B3) $B$ commutes with $C$, and $BC=-i\gamma^4$.} $BC=-iC\gamma^4C=-iC^2\gamma^4=-i\gamma^4$ and $CB=-iC^2\gamma^4=-i\gamma^4$.

\textbf{(B4) Signature (8,8).} $\operatorname{tr}B=-i\operatorname{tr}(\gamma^0\gamma^1\gamma^2\gamma^3\gamma^4)=0$ by Theorem 2.1; with (B2) and Section 2.2, $B$ has the eigenvalue $+1$ eight times and $-1$ eight times.

(Stage 1, Result 10.2; check \texttt{ALG\_char\allowbreak{}geFormB}, which records the multiplicities 8 and 8 and the identity $BC=-i\gamma^4$.)

For the notebook's gammas, $B$ is $i$ times the signed permutation in the table of Section 2.9; for example $(Bu)_0=i\,u_9$ and $(Bu)_9=-i\,u_0$, so that $B_{09}=i$ and $B_{90}=-i=(B_{09})^\ast$, as a Hermitian matrix requires.

\textbf{Where $B$ is used.} In the quantum theory of Chapter 8, $B$ is the matrix of the equal-time anticommutator of the field in the simplest (Gaussian normal) coordinates, $\{\Psi_a(x),\Psi^\dagger_b(y)\}=B_{ab}\,\delta^7(x-y)/\sqrt{|g|}$ (Stage 1, §10.3). Its signature (8,8) is the algebraic root of the indefinite ("Krein") inner product found there.

\subsection{Groups and representations from zero}

\textbf{Groups.} A \textbf{group} is a set $G$ with a product $gh$ such that: the product is associative, $(gh)k=g(hk)$; there is an identity element $e$ with $eg=ge=g$; and every $g$ has an inverse $g^{-1}$ with $gg^{-1}=g^{-1}g=e$. Examples: the nonzero real numbers under multiplication; the two numbers $\{+1,-1\}$ under multiplication; the invertible $d$ by $d$ matrices under matrix multiplication; the rotations of the plane. A \textbf{subgroup} is a subset that is itself a group with the same product. The subgroup \textbf{generated} by some elements is the set of all finite products of them and their inverses.

\textbf{Homomorphisms.} A map $\varphi$ from a group $G$ to a group $H$ is a \textbf{homomorphism} if $\varphi(gh)=\varphi(g)\varphi(h)$ for all $g,h$. Its \textbf{kernel} is the set of $g$ with $\varphi(g)=e$. Example: the determinant is a homomorphism from invertible matrices to nonzero numbers, $\det(AB)=\det A\det B$; its kernel is the set of matrices of determinant 1.

\textbf{Representations.} A \textbf{representation} of a group $G$ on $V=\mathbb R^d$ or $\mathbb C^d$ is a homomorphism $\rho$ from $G$ to the invertible $d$ by $d$ matrices: every group element $g$ becomes a matrix $\rho(g)$ that acts on columns, and $\rho(gh)=\rho(g)\rho(h)$. We then say that $G$ \textbf{acts} on $V$. A subspace $W\subseteq V$ (Section 2.2: a subset closed under addition and multiplication by numbers) is \textbf{invariant} if $\rho(g)w\in W$ for every $g$ and every $w\in W$. The subspaces $\{0\}$ and $V$ are always invariant. The representation is \textbf{irreducible} if $V\ne\{0\}$ and there is no other invariant subspace: $V$ cannot be cut into smaller pieces that the group respects.

\textbf{Intertwiners, equivalence, commutant.} Let $\rho_1$ act on $V_1$ and $\rho_2$ on $V_2$. A matrix $T$ from $V_1$ to $V_2$ is an \textbf{intertwiner} if $T\rho_1(g)=\rho_2(g)T$ for all $g$. The two representations are \textbf{equivalent} if an \emph{invertible} intertwiner exists; then $\rho_2(g)=T\rho_1(g)T^{-1}$, and the two representations are the same up to a change of basis. The intertwiners from a representation to itself form its \textbf{commutant}: the matrices that commute with every $\rho(g)$.

\textbf{Example: the field of numbers matters.} The rotations of the plane act on $\mathbb R^2$ by $R(\varphi)=\begin{pmatrix}\cos\varphi&-\sin\varphi\\ \sin\varphi&\cos\varphi\end{pmatrix}$. Over the real numbers this is irreducible: an invariant line would be a direction that every rotation maps to itself, and a rotation by $90^{\circ}$ has no such direction. Over the complex numbers it is reducible: the column $w=(1,-i)^T$ satisfies

\[
R(\varphi)w=\begin{pmatrix}\cos\varphi+i\sin\varphi\\ \sin\varphi-i\cos\varphi\end{pmatrix}=e^{i\varphi}\begin{pmatrix}1\\-i\end{pmatrix},
\]

so the complex line through $w$ is invariant. This is why Stage 1 says precisely "irreducible \emph{complex} representation".

\textbf{Lemma 2.5.} For an intertwiner $T$ from $\rho_1$ to $\rho_2$, the \textbf{kernel} $\{v:Tv=0\}$ is invariant under $\rho_1$ and the \textbf{image} $\{Tv\}$ is invariant under $\rho_2$. \emph{Proof.} If $Tv=0$ then $T\rho_1(g)v=\rho_2(g)Tv=0$. And $\rho_2(g)(Tv)=T(\rho_1(g)v)$ lies in the image. $\square$

\textbf{Theorem 2.6 (Schur's lemma, first part).} An intertwiner between two irreducible representations is either zero or invertible. \emph{Proof.} By Lemma 2.5 the kernel is $\{0\}$ or all of $V_1$, and the image is $\{0\}$ or all of $V_2$. If $T\ne0$, the kernel is not all of $V_1$, so it is $\{0\}$ ($T$ is one-to-one), and the image is not $\{0\}$, so it is $V_2$ ($T$ is onto). A one-to-one and onto linear map is invertible. $\square$

\textbf{Theorem 2.7 (span criterion).} Let $\rho$ be a representation of $G$ on $\mathbb C^d$ whose matrices $\rho(g)$ span, by complex linear combinations, all of $\mathrm{Mat}_d(\mathbb C)$. Then $\rho$ is irreducible, and its commutant consists of the multiples of $I$.

\emph{Proof.} Let $W\ne\{0\}$ be invariant and pick $w\ne0$ in $W$. Since $W$ is invariant under every $\rho(g)$ and closed under linear combinations, it is invariant under every matrix. For any column $v$, the matrix $M=v\,w^\dagger/(w^\dagger w)$ satisfies $Mw=v$; hence $v\in W$, and $W=\mathbb C^d$. For the commutant: a matrix commuting with every $\rho(g)$ commutes with every linear combination, hence with every matrix, and the elementary-matrix argument of Corollary 2.3 shows that it is a multiple of $I$. $\square$

The second statement is the form of \textbf{Schur's lemma} that Stage 1 uses: here the commutant consists of the scalar matrices only. Section 2.15 also uses the commutant in the opposite direction: when a representation is a sum of two pieces, the size of its commutant tells whether the two pieces are equivalent.

\subsection{Reflections, O(4,4), Pin(4,4) and Spin(4,4)}

\textbf{Vectors inside the algebra.} To a vector $v=(v_0,\dots,v_7)$ of $\mathbb R^8$ we attach the matrix

\[
\gamma(v):=\sum_{a=0}^7v_a\gamma^a,
\]

and we write $\eta(u,v):=\sum_a\eta^{aa}u_av_a$ for the metric product of two vectors. The pairing argument of Section 2.3 gives

\[
\gamma(u)\gamma(v)+\gamma(v)\gamma(u)=2\eta(u,v)\,1,\qquad\gamma(v)^2=\eta(v,v)\,1 .
\]

A vector $u$ with $n(u):=\eta(u,u)=+1$ or $-1$ is a \textbf{unit vector}, and $n(u)$ is its \textbf{norm}. For a unit vector $\gamma(u)^2=n(u)$, so $\gamma(u)$ is invertible with $\gamma(u)^{-1}=n(u)\gamma(u)$. (A vector with $\eta(u,u)=0$, such as $u=e_0+e_4$, is called \textbf{null}; then $\gamma(u)^2=0$ and $\gamma(u)$ has no inverse, which is why null vectors are excluded below.)

\textbf{The groups O(4,4) and SO(4,4).} $\mathrm O(4,4)$ is the set of real 8 by 8 matrices $\Lambda$ that preserve the metric, $\eta(\Lambda u,\Lambda v)=\eta(u,v)$ for all $u,v$; in matrix form $\Lambda^T\eta\Lambda=\eta$. It is a group: products and inverses of such matrices again preserve $\eta$. Taking determinants, $(\det\Lambda)^2\det\eta=\det\eta$, and $\det\eta=(+1)^4(-1)^4=1$, so $\det\Lambda=\pm1$. The elements with $\det\Lambda=+1$ form the subgroup $\mathrm{SO}(4,4)$. These are the "Lorentz transformations" of the 4+4 dimensional flat spacetime, including those that reverse orientation.

\textbf{Reflections.} For a unit vector $u$ define the \textbf{reflection in the hyperplane orthogonal to $u$}:

\[
R_u(v):=v-2\,\frac{\eta(u,v)}{\eta(u,u)}\,u .
\]

It sends $u$ to $u-2u=-u$ and leaves every $w$ with $\eta(u,w)=0$ unchanged. The vectors $w$ with $\eta(u,w)=0$ form a 7-dimensional subspace (the solutions of one nonzero linear equation), and $u$ does not lie in it because $\eta(u,u)\ne0$. In a basis made of $u$ and seven vectors of that subspace, $R_u=\mathrm{diag}(-1,1,1,1,1,1,1,1)$; hence $R_u^2=1$ and $\det R_u=-1$. $R_u$ preserves $\eta$ (Exercise 2.11), so $R_u\in\mathrm O(4,4)$ but $R_u\notin\mathrm{SO}(4,4)$.

\textbf{The key identity.} For a unit vector $u$ and any vector $v$,

\[
-\gamma(u)\,\gamma(v)\,\gamma(u)^{-1}=\gamma\bigl(R_u(v)\bigr).
\]

\emph{Proof.} From the anticommutator, $\gamma(u)\gamma(v)=2\eta(u,v)-\gamma(v)\gamma(u)$. Multiply on the right by $\gamma(u)^{-1}=n(u)\gamma(u)$:

\[
\gamma(u)\gamma(v)\gamma(u)^{-1}=2\eta(u,v)\,n(u)\,\gamma(u)-\gamma(v)=-\gamma\Bigl(v-2\frac{\eta(u,v)}{\eta(u,u)}u\Bigr),
\]

because $n(u)=\eta(u,u)=1/\eta(u,u)$ when $\eta(u,u)=\pm1$. $\square$ So conjugation by a unit vector is a reflection, up to a sign.

\textbf{Definition of Pin(4,4) and Spin(4,4).} $\mathrm{Pin}(4,4)$ is the set of all products

\[
g=\gamma(u_1)\gamma(u_2)\cdots\gamma(u_k)
\]

of finitely many unit vectors ($k=0,1,2,\dots$; the empty product is 1). It is a group: the product of two such products is again one, and $g^{-1}=\gamma(u_k)^{-1}\cdots\gamma(u_1)^{-1}$ with $\gamma(u)^{-1}=\gamma(n(u)u)$, again a unit vector. $\mathrm{Spin}(4,4)$ is the subgroup of products with an \textbf{even} number $k$ of factors. Every element acts on spinors by matrix multiplication, $\Psi\mapsto g\Psi$ (Stage 1, §4.2). Examples: each $\gamma^a$ is a unit vector ($n=\eta^{aa}$), so $\gamma^a\in\mathrm{Pin}(4,4)$; the numbers $1=\gamma^0\gamma^0$ and $-1=\gamma^4\gamma^4$ lie in $\mathrm{Spin}(4,4)$; every monomial is $\pm$ a product of gammas and so lies in $\mathrm{Pin}(4,4)$, in $\mathrm{Spin}(4,4)$ if it is even; and $\exp(\theta S^{ab})$ lies in $\mathrm{Spin}(4,4)$, because (Section 2.11) $\cos\tfrac\theta2+\sin\tfrac\theta2\gamma^a\gamma^b=\gamma^a\bigl(\cos\tfrac\theta2\,\eta^{aa}\gamma^a+\sin\tfrac\theta2\,\gamma^b\bigr)$ is a product of two unit vectors of norm $\eta^{aa}$ (and similarly with cosh and sinh for a boost, where $\cosh^2-\sinh^2=1$ gives the norm).

\textbf{The parity is well defined.} A product of $k$ vectors is, by (R1), a combination of monomials whose degrees all have the parity of $k$ (the degree of $\gamma_A\gamma_B=\pm\gamma_{A\triangle B}$ has the parity of $|A|+|B|$). If an element could be written both with an even and with an odd number of factors, it would be both a combination of even monomials and a combination of odd monomials, hence zero (Section 2.6), which is impossible for an invertible matrix. So every $g$ has a definite \textbf{parity} $|g|\in\{0,1\}$ (0 for even, 1 for odd), and we put $\alpha(g):=(-1)^{|g|}g$.

\textbf{Theorem 2.8 (spinor norm).} For $g=\gamma(u_1)\cdots\gamma(u_k)\in\mathrm{Pin}(4,4)$,

\[
g^TCg=(-1)^k\,N(g)\,C,\qquad N(g):=n(u_1)\,n(u_2)\cdots n(u_k)=\pm1 .
\]

\emph{Proof.} First, $C\gamma(u)^TC^{-1}=-\gamma(u)$ for every vector: by (N3) and (C1), $C(\gamma^a)^TC^{-1}=\eta^{aa}C\gamma^aC^{-1}=\eta^{aa}(-\eta^{aa}\gamma^a)=-\gamma^a$, and the rest is linearity. Hence $Cg^TC^{-1}=(C\gamma(u_k)^TC^{-1})\cdots(C\gamma(u_1)^TC^{-1})=(-1)^k\gamma(u_k)\cdots\gamma(u_1)$, and, since $C^{-1}=C$, $g^TC=(-1)^kC\,\gamma(u_k)\cdots\gamma(u_1)$. Multiplying by $g$ on the right, the product $\gamma(u_k)\cdots\gamma(u_1)\gamma(u_1)\cdots\gamma(u_k)$ collapses from the middle, $\gamma(u_1)^2=n(u_1)$, then $\gamma(u_2)^2=n(u_2)$, and so on. $\square$

The number $N(g)$ is the \textbf{spinor norm}. The formula shows that it depends only on $g$, not on how $g$ is written as a product, and that $N(gh)=N(g)N(h)$. Examples: $N(\gamma^0)=+1$, $N(\gamma^4)=-1$, $N(\gamma^0\gamma^4)=-1$, and $N(\exp(\theta S^{ab}))=+1$ (it is a product of two unit vectors of equal norm). For a spinor, $\Psi\mapsto g\Psi$ multiplies the bilinear $\Psi^\dagger C\Psi$ by $(-1)^kN(g)$, because $(g\Psi)^\dagger C(g\Psi)=\Psi^\dagger g^TCg\Psi$ for a real $g$; for a single unit vector the factor is $-n(u)$. This is the "Pin character" of Stage 1 (§7.7; check \texttt{ALG\_pinL\allowbreak{}iftChara\allowbreak{}cter}), which Chapter 6 uses.

\textbf{Theorem 2.9 (how spinor transformations move vectors).} For every $g\in\mathrm{Pin}(4,4)$ there are matrices $\Lambda_{\mathrm t}(g)$ and $\Lambda_{\mathrm u}(g)$ in $\mathrm O(4,4)$ with

\[
\alpha(g)\,\gamma(v)\,g^{-1}=\gamma\bigl(\Lambda_{\mathrm t}(g)v\bigr),\qquad g\,\gamma(v)\,g^{-1}=\gamma\bigl(\Lambda_{\mathrm u}(g)v\bigr)\qquad\text{for all }v .
\]

Both are homomorphisms, $\Lambda_{\mathrm t}(\gamma(u))=R_u$, $\Lambda_{\mathrm u}(\gamma(u))=-R_u$, $\Lambda_{\mathrm u}(g)=(-1)^{|g|}\Lambda_{\mathrm t}(g)$, and $\det\Lambda_{\mathrm t}(g)=\det\Lambda_{\mathrm u}(g)=(-1)^{|g|}$.

\emph{Proof.} For one unit vector, $\alpha(\gamma(u))=-\gamma(u)$ and the key identity gives $\Lambda_{\mathrm t}(\gamma(u))=R_u$; without the sign, $\Lambda_{\mathrm u}(\gamma(u))=-R_u$. For a product $gh$, parities add, so $\alpha(gh)=\alpha(g)\alpha(h)$ and

\[
\alpha(gh)\gamma(v)(gh)^{-1}=\alpha(g)\bigl(\alpha(h)\gamma(v)h^{-1}\bigr)g^{-1}=\alpha(g)\,\gamma\bigl(\Lambda_{\mathrm t}(h)v\bigr)\,g^{-1}=\gamma\bigl(\Lambda_{\mathrm t}(g)\Lambda_{\mathrm t}(h)v\bigr).
\]

So, step by step along the factors, a matrix $\Lambda_{\mathrm t}(g)$ with the defining property exists for every product, and it is a product of reflections (hence in $\mathrm O(4,4)$). It depends only on $g$, not on the chosen product: $\alpha(g)$ depends only on $g$ (the parity is well defined), and if two matrices $\Lambda$ and $\Lambda'$ both satisfy $\alpha(g)\gamma(v)g^{-1}=\gamma(\Lambda v)=\gamma(\Lambda'v)$ for all $v$, then $\gamma(\Lambda v-\Lambda'v)=0$; since $\gamma(w)=\sum_aw_a\gamma^a=0$ forces $w=0$ (the $\gamma^a$ are independent, Theorem 2.2), $\Lambda v=\Lambda'v$ for every $v$, that is, $\Lambda=\Lambda'$. Therefore the last display proves $\Lambda_{\mathrm t}(gh)=\Lambda_{\mathrm t}(g)\Lambda_{\mathrm t}(h)$. The untwisted version differs by the sign $(-1)^{|g|}$ and is a homomorphism for the same reason; $-1_8$ is in $\mathrm O(4,4)$. Determinants: $\det R_u=-1$ and $\det(-R_u)=(-1)^8\det R_u=-1$, so a product of $k$ factors has determinant $(-1)^k$ under both maps. $\square$

The map $\Lambda_{\mathrm t}$ with the factor $\alpha$ is the \textbf{twisted adjoint action} and $\Lambda_{\mathrm u}$ the \textbf{untwisted} one (Stage 1, §4.2; contract §2). They agree on $\mathrm{Spin}(4,4)$. The computer confirms, for several unit vectors, that the untwisted matrix lies exactly in $\mathrm O(4,4)$ (check \texttt{ALG\_pinL\allowbreak{}iftChara\allowbreak{}cter}).

\textbf{Theorem 2.10 (the kernel is $\{\pm1\}$).} $\Lambda_{\mathrm t}(g)=1_8$ exactly when $g=\pm1$, and the same holds for $\Lambda_{\mathrm u}$.

\emph{Proof.} Both $\pm1$ lie in $\mathrm{Spin}(4,4)$ and act trivially. Conversely, let $\Lambda_{\mathrm t}(g)=1_8$, that is, $\alpha(g)\gamma^a=\gamma^ag$ for all $a$.

\begin{itemize}
\item If $g$ is even, $\alpha(g)=g$, so $g$ commutes with every $\gamma^a$. By Corollary 2.3, $g=c\cdot1$ with a real number $c$. Theorem 2.8 with even $k$ gives $c^2C=N(g)C$, so $c^2=N(g)=\pm1$; a real $c$ has $c^2=1$, hence $c=\pm1$.
\item If $g$ is odd, $\alpha(g)=-g$, so $g$ anticommutes with every $\gamma^a$. Then $g\gamma^8$ commutes with every $\gamma^a$: $g\gamma^8\gamma^a=-g\gamma^a\gamma^8=\gamma^ag\gamma^8$, by (X2). By Corollary 2.3, $g\gamma^8=c\cdot1$, so $g=c\gamma^8$ (since $(\gamma^8)^2=1$). But $\gamma^8$ is an even monomial and $g$ is odd; a nonzero matrix cannot be both, so $c=0$ and $g=0$, impossible.
\end{itemize}

For $\Lambda_{\mathrm u}$ the even case is identical, and in the odd case $g$ would commute with every $\gamma^a$, so $g=c\cdot1$ would be even: impossible. $\square$

\textbf{Onto, and the double cover (quoted theorem).} The \textbf{Cartan--Dieudonné theorem} of linear algebra states that every element of $\mathrm O(4,4)$ is a product of at most 8 reflections $R_u$ in non-null vectors, which may be normalized to unit vectors. We quote it without proof; it is not machine-checked in the repository. With it, $\Lambda_{\mathrm t}$ maps $\mathrm{Pin}(4,4)$ \textbf{onto} $\mathrm O(4,4)$. Together with Theorem 2.10: if $\Lambda_{\mathrm t}(g)=\Lambda_{\mathrm t}(h)$, then $\Lambda_{\mathrm t}(gh^{-1})=1_8$, so $g=\pm h$. Every element of $\mathrm O(4,4)$ is the image of exactly two elements $\pm g$: \textbf{Pin(4,4) is a double cover of O(4,4)}. Since determinants are $(-1)^{|g|}$, the even elements are exactly those mapped into $\mathrm{SO}(4,4)$, and \textbf{Spin(4,4) is a double cover of SO(4,4)}, the determinant-1 transformations, as the repository's \texttt{README.m\allowbreak{}d} puts it. The untwisted map is onto $\mathrm O(4,4)$ as well: on even elements it agrees with $\Lambda_{\mathrm t}$, and on odd elements it is $-\Lambda_{\mathrm t}$; multiplication by $-1_8$, whose determinant is $+1$ in 8 dimensions, maps the determinant $-1$ part of $\mathrm O(4,4)$ onto itself. (Stage 1 records the stronger fact that $-1_8$ lies in the part of $\mathrm{SO}(4,4)$ connected to $1_8$: it is the product of rotations by $\pi$ in the planes (0,1), (2,3), (4,5) and (6,7).) The rotation by $2\pi$ of Section 2.11 shows the double cover at work: $R(2\pi)=-1$ and $R(0)=1$ have the same vector image $1_8$.

\textbf{"Determinant 1" refers to the vector image.} As 16 by 16 matrices, all elements of $\mathrm{Pin}(4,4)$ have determinant $+1$. \emph{Proof.} For a unit vector $u$, $\gamma(u)$ is traceless (a combination of traceless monomials). If $n(u)=+1$, $\gamma(u)^2=1$ and Section 2.2 gives eight eigenvalues $+1$ and eight $-1$, so $\det\gamma(u)=(+1)^8(-1)^8=1$. If $n(u)=-1$, $\gamma(u)^2=-1$ and Section 2.2 gives $\det\gamma(u)=1$ as well. A product of determinant-1 matrices has determinant 1. $\square$ So the spinor determinant does not distinguish $\mathrm{Spin}(4,4)$ from $\mathrm{Pin}(4,4)$; "determinant 1" always means $\det\Lambda(g)=1$ (Stage 1, §4.2).

\textbf{The part connected to 1.} The products of exponentials $\exp(\theta S^{ab})$ form a subgroup of $\mathrm{Spin}(4,4)$, called $\mathrm{Spin}_0(4,4)$; Stage 1 calls it the identity component (the elements that can be reached from 1 by a continuous path, a topological description that we do not need). Every element of it has spinor norm $N=+1$, because $N$ is multiplicative and $N(\exp(\theta S^{ab}))=+1$. Since $N(\gamma^0\gamma^4)=-1$, the element $\gamma^0\gamma^4$ of $\mathrm{Spin}(4,4)$ is \emph{not} a product of exponentials: $\mathrm{Spin}(4,4)$ contains more than its part connected to 1. By Theorem 2.8, elements of $\mathrm{Spin}(4,4)$ with $N=-1$ reverse the sign of $\Psi^\dagger C\Psi$; Chapter 6 uses this to decide which transformations are symmetries of the Lagrangian.

\subsection{The representation theorem}

We can now prove the central theorem of this chapter (Stage 1, §4.3 and §4.4; contract §2). In this section "the gammas" are the notebook's, so that $\gamma^8=\mathrm{diag}(-I_8,I_8)$ by (X5); the statements hold for any 16 by 16 gammas of signature (4,4) after the change of basis that diagonalizes $\gamma^8$.

\textbf{Theorem 2.11 (irreducibility under Pin(4,4)).} The representation $g\mapsto g$ of $\mathrm{Pin}(4,4)$ on $\mathbb C^{16}$ is irreducible, and its commutant consists of the multiples of $I_{16}$.

\emph{Proof.} Every monomial is a product of gammas, and each $\gamma^a$ is a unit vector; so every monomial lies in $\mathrm{Pin}(4,4)$, up to a sign that can be absorbed because $-1\in\mathrm{Pin}(4,4)$. Hence the span of $\mathrm{Pin}(4,4)$ contains the 256 monomials, which span all real 16 by 16 matrices (Section 2.6), and with complex coefficients all of $\mathrm{Mat}_{16}(\mathbb C)$. Theorem 2.7 gives both claims. $\square$

\textbf{Theorem 2.12 (two inequivalent halves under Spin(4,4)).} Let $\mathbb C^8_-$ and $\mathbb C^8_+$ be the images of $P_-$ and $P_+$: the columns whose lower, respectively upper, eight components vanish. Then

\begin{enumerate}
\item both halves are invariant under $\mathrm{Spin}(4,4)$;
\item the span of $\mathrm{Spin}(4,4)$ is the set of all block-diagonal matrices $\mathrm{diag}(X,Y)$ with arbitrary $X,Y\in\mathrm{Mat}_8(\mathbb C)$;
\item each half is an irreducible representation of $\mathrm{Spin}(4,4)$ whose commutant consists of the multiples of $I_8$;
\item the two halves are \textbf{inequivalent}: the only intertwiner between them is zero;
\item the commutant of $\mathrm{Spin}(4,4)$ on $\mathbb C^{16}$ is two-dimensional, spanned by $P_-$ and $P_+$.
\end{enumerate}

\emph{Proof.} (1) Every element of $\mathrm{Spin}(4,4)$ is a combination of even monomials, which commute with $\gamma^8$ (X3) and are therefore block diagonal (Section 2.10); a block-diagonal matrix maps each half into itself.

(2) Every even monomial is $\pm$ an element of $\mathrm{Spin}(4,4)$, and every element of $\mathrm{Spin}(4,4)$ is a combination of even monomials, so the two spans agree. The 128 even monomials are independent (Section 2.6) and block diagonal. The real block-diagonal matrices form a set of dimension $64+64=128$, so the even monomials are a basis of it; with complex coefficients they span all complex block-diagonal matrices.

(3) By (2), the upper-left blocks $X$ of the matrices in the span run through all of $\mathrm{Mat}_8(\mathbb C)$ (take $Y=0$), and so do the lower-right blocks. Theorem 2.7, applied to each half, gives irreducibility and scalar commutants.

(4) An intertwiner is an 8 by 8 matrix $T$ with $Tg_-=g_+T$ for every $g=\mathrm{diag}(g_-,g_+)$ in $\mathrm{Spin}(4,4)$. By linearity $TX=YT$ for every $\mathrm{diag}(X,Y)$ in the span, and by (2) we may choose $X=I_8$ and $Y=0$: then $T=0$. So there is no invertible intertwiner, and the halves are inequivalent. The same argument works from $\mathbb C^8_+$ to $\mathbb C^8_-$.

(5) Write a matrix that commutes with the span in blocks, $Z=\begin{pmatrix}Z_{11}&Z_{12}\\ Z_{21}&Z_{22}\end{pmatrix}$. Commuting with $\mathrm{diag}(X,Y)$ means $Z_{11}X=XZ_{11}$, $Z_{22}Y=YZ_{22}$, $Z_{12}Y=XZ_{12}$ and $Z_{21}X=YZ_{21}$ for all $X,Y$. The first two force $Z_{11}=a\,I_8$ and $Z_{22}=b\,I_8$ (the elementary-matrix argument); the last two, with $X=I_8$, $Y=0$ and with $X=0$, $Y=I_8$, force $Z_{12}=Z_{21}=0$. So $Z=aP_-+bP_+$. $\square$

Had the two halves been equivalent through an invertible $T$, the matrix $\begin{pmatrix}0&0\\ T&0\end{pmatrix}$ would also commute with $\mathrm{Spin}(4,4)$ and the commutant would be four-dimensional (a copy of $\mathrm{Mat}_2(\mathbb C)$). The computed dimension 2 therefore also proves inequivalence; this is the form of the argument in Stage 1 (§4.4).

\textbf{Remarks.}

\begin{itemize}
\item The odd elements of $\mathrm{Pin}(4,4)$, such as the $\gamma^a$ themselves, exchange the two halves (X6). That is why $\mathrm{Pin}(4,4)$ sees one block of 16 while $\mathrm{Spin}(4,4)$ sees two blocks of 8.
\item The theorem holds unchanged for $\mathrm{Spin}_0(4,4)$: its span also contains every $\gamma^a\gamma^b$ ($a\ne b$), since for a rotation $\exp(\pi S^{ab})=\gamma^a\gamma^b$, and for a boost $\exp(\theta S^{ab})-\exp(-\theta S^{ab})=2\sinh\tfrac\theta2\,\gamma^a\gamma^b$; products of these give every even monomial.
\item The notebook's matrices are real, so the same splitting holds for real spinors, $\mathbb R^{16}=\Delta_-\oplus\Delta_+$; dirac16complex uses the complex version $\mathbb C\otimes(\Delta_-\oplus\Delta_+)$ (Stage 1, §4.5), with components that are, in addition, Grassmann-odd (Chapter 5 and Chapter 6).
\item The vector space $\mathbb R^8$ and the two halves $\Delta_-$ and $\Delta_+$ are all 8-dimensional. The notebook remarks on these three 8-dimensional representations of Spin(4,4) (cell 669, per the notebook survey); the symmetry that permutes them is called \textbf{triality}. It is a standard fact that this book neither proves nor uses; Chapter 3 shows the concrete identification of all three spaces with the split octonions.
\end{itemize}

\textbf{What the computer computed.} The Python checker writes the commutant condition $\gamma^aX-X\gamma^a=0$ ($a=0,\dots,7$) for an unknown 16 by 16 matrix $X$ as $8\cdot256=2048$ linear equations in the 256 unknowns $X_{ij}$, with integer coefficients, and computes the dimension of their solution space by exact Gaussian elimination with fractions: by fact (ii) of Section 2.2 it is 256 minus the rank of the 2048 equations (functions \texttt{intertwi\allowbreak{}ner\_equa\allowbreak{}tions}, \texttt{commutan\allowbreak{}t\_dimens\allowbreak{}ion} and \texttt{\detokenize{rank}} of \texttt{scripts/\allowbreak{}d16c\_exa\allowbreak{}ct.py}); the Wolfram verifier does the same with its own code. The solution space has dimension exactly 1 (the scalar matrices), which is Theorem 2.11 in computed form (Stage 1, Result 4.1; check \texttt{ALG\_pinI\allowbreak{}rreducib\allowbreak{}leComple\allowbreak{}x}). Gaussian elimination only adds, subtracts, multiplies and divides the integer coefficients, so it runs identically whether complex or only rational solutions are allowed; the dimension over the complex numbers is the same. For the 28 spin generators (whose products give all even monomials, so they have the same commutant as $\mathrm{Spin}(4,4)$), the computed commutant has dimension 2 and is spanned by $P_-$ and $P_+$; every $S^{ab}$ and every even monomial is block diagonal; the commutants of the two blocks have dimensions (1, 1); the even monomials restricted to each block have rank (64, 64); and the intertwiner spaces between the blocks have dimensions (0, 0) in both directions (Stage 1, Result 4.2; check \texttt{ALG\_spin\allowbreak{}Decompos\allowbreak{}ition}). Every one of these numbers is predicted by Theorem 2.12.

\subsection{What the computer checked}

Each statement below is proved by hand in this chapter (or, where marked, taken from the check and proved by hand in Chapter 3) and is also checked with exact integer or rational arithmetic by two independent programs: the Wolfram verifier \texttt{scripts/\allowbreak{}verify\_d\allowbreak{}irac16co\allowbreak{}mplex\_al\allowbreak{}gebra.wl\allowbreak{}s} (with the package \texttt{wolfram/\allowbreak{}Dirac16C\allowbreak{}omplexAl\allowbreak{}gebra.wl}) and the Python checker \texttt{scripts/\allowbreak{}check\_di\allowbreak{}rac16com\allowbreak{}plex\_alg\allowbreak{}ebra.py} (with \texttt{scripts/\allowbreak{}d16c\_exa\allowbreak{}ct.py}) (one row, on $[S^{ab},\gamma^c]$, is the exception: the Wolfram verifier records it, as true, in a measurement outside the pass condition of \texttt{ALG\_spin\allowbreak{}Transpos\allowbreak{}ePropert\allowbreak{}ies}, and the exact Python geometry checker \texttt{scripts/\allowbreak{}check\_di\allowbreak{}rac16com\allowbreak{}plex\_geo\allowbreak{}metry.py} checks it as \texttt{ALG\_SabG\allowbreak{}ammaComm\allowbreak{}utator}, recorded in \texttt{python-g\allowbreak{}eometry-\allowbreak{}report.j\allowbreak{}son}). Every other check of the table is recorded as true in both algebra reports, \texttt{wolfram-\allowbreak{}algebra-\allowbreak{}report.j\allowbreak{}son} (21 of 21 checks true) and \texttt{python-a\allowbreak{}lgebra-r\allowbreak{}eport.js\allowbreak{}on} (21 of 21 checks true); \texttt{ALG\_SabG\allowbreak{}ammaComm\allowbreak{}utator} is recorded as true in \texttt{python-g\allowbreak{}eometry-\allowbreak{}report.j\allowbreak{}son} (52 of 52 checks true). The three reports lie in the folder \texttt{artifact\allowbreak{}s/dirac1\allowbreak{}6complex\allowbreak{}/arbitra\allowbreak{}ry-field\allowbreak{}/} (Stage 1, §11).

\begingroup
\small
\begin{longtable}{@{}>{\raggedright\arraybackslash}p{0.400\linewidth}>{\raggedright\arraybackslash}p{0.400\linewidth}@{}}
\toprule
\textbf{Statement (where in this chapter; where in Stage 1)} & \textbf{Check} \\
\midrule
\endfirsthead
\toprule
\textbf{Statement (where in this chapter; where in Stage 1)} & \textbf{Check} \\
\midrule
\endhead
Clifford relation of the notebook's gammas, all 64 pairs (2.8, proved by hand in Chapter 3; Result 3.1) & \texttt{ALG\_clif\allowbreak{}ford} \\
$\gamma^a$ symmetric for $a\le3$ and antisymmetric for $a\ge4$ (2.8; Result 3.1) & \texttt{ALG\_gamm\allowbreak{}aTranspo\allowbreak{}seSymmet\allowbreak{}ry} \\
the 256 monomials have rank 256, the 128 even ones rank 128 (2.6 and 2.8; Result 3.2) & \texttt{ALG\_fait\allowbreak{}hful} \\
$C=\gamma^0\gamma^1\gamma^2\gamma^3=\mathrm{diag}(-\sigma,\sigma)$ is symmetric, $C^2=1$, signature (8,8) (2.9, the form $\mathrm{diag}(-\sigma,\sigma)$ proved by hand in Chapter 3; Result 3.3) & \texttt{ALG\_char\allowbreak{}geMatrix} \\
expression [1], $(C\gamma^a)^T=-C\gamma^a$: eight times True (2.9; Result 3.4) & \texttt{ALG\_expr\allowbreak{}ession1} \\
$(CS^{ab})^T=-CS^{ab}$ and the triple identities of Exercise 2.8 (2.11; Result 3.5) & \texttt{ALG\_spin\allowbreak{}Transpos\allowbreak{}ePropert\allowbreak{}ies} \\
$[S^{ab},\gamma^c]=\eta^{bc}\gamma^a-\eta^{ac}\gamma^b$, all 512 triples (2.11; Result 3.5) & \texttt{ALG\_SabG\allowbreak{}ammaComm\allowbreak{}utator} (Python geometry report) \\
$\gamma^8=\mathrm{diag}(-I_8,I_8)$ and its commutation rules (2.10, the form $\mathrm{diag}(-I_8,I_8)$ proved by hand in Chapter 3; Result 3.6) & \texttt{ALG\_chir\allowbreak{}ality} \\
the tensor-picture gammas satisfy the Clifford relation (2.7; Result 3.7) & \texttt{ALG\_clif\allowbreak{}fordPict\allowbreak{}ureInter\allowbreak{}twiner} \\
$B$ is Hermitian, $B^2=1$, eight eigenvalues $+1$ and eight $-1$, $[C,B]=0$, $BC=-i\gamma^4$ (2.12; Result 10.2) & \texttt{ALG\_char\allowbreak{}geFormB} \\
$\Psi^\dagger C\Psi\to-n(u)\,\Psi^\dagger C\Psi$ for a unit vector; the untwisted image lies in O(4,4) (2.14; §4.2 and §7.7) & \texttt{ALG\_pinL\allowbreak{}iftChara\allowbreak{}cter} \\
the commutant of the eight gammas has dimension 1 (2.15; Result 4.1) & \texttt{ALG\_pinI\allowbreak{}rreducib\allowbreak{}leComple\allowbreak{}x} \\
the commutant of the $S^{ab}$ has dimension 2; block commutants (1, 1); even ranks (64, 64); cross intertwiners (0, 0) (2.15; Result 4.2) & \texttt{ALG\_spin\allowbreak{}Decompos\allowbreak{}ition} \\
exactly 13 of the 28 $S^{ab}$ commute with $B$ (Exercise 2.10; Result 10.6) & \texttt{QNT\_unit\allowbreak{}aryAndKr\allowbreak{}einSubgr\allowbreak{}oups} \\
\bottomrule
\end{longtable}
\endgroup

The two programs also check each other. The Wolfram verifier reads the exact data file \texttt{algebra-\allowbreak{}fixture.\allowbreak{}json} written by Python and finds all 75 stored matrices equal to its own (check \texttt{ALG\_fixt\allowbreak{}ureAgree\allowbreak{}ment}), and the Python checker reads the Wolfram report and finds all 47 shared measurements equal (check \texttt{ALG\_wolf\allowbreak{}ramAgree\allowbreak{}ment}, Python report only; Stage 1, §3.4 and §11.6). To rerun the Python algebra checker without touching the committed reports, run from the repository root, in Git Bash

\begin{Verbatim}[fontsize=\small]
export PYTHONUTF8=1
python scripts/check_dirac16complex_algebra.py --wolfram-report= \
    --output build/textbook/python-algebra-report.json
\end{Verbatim}

or in PowerShell (where the line continuation is a backtick)

\begin{Verbatim}[fontsize=\small]
$env:PYTHONUTF8 = "1"
python scripts/check_dirac16complex_algebra.py --wolfram-report= `
    --output build/textbook/python-algebra-report.json
\end{Verbatim}

It prints one line per check (for example \texttt{check\_AL\allowbreak{}G\_cliffo\allowbreak{}rd=true}) and one line per measurement (about a hundred lines), then \texttt{check\_co\allowbreak{}unt=20} and \texttt{failed\_c\allowbreak{}heck\_cou\allowbreak{}nt=0}, and finally \texttt{report=b\allowbreak{}uild/tex\allowbreak{}tbook/py\allowbreak{}thon-alg\allowbreak{}ebra-rep\allowbreak{}ort.json}. The 20 checks are the 21 checks of the committed report except the comparison with the Wolfram report, which the empty \texttt{--wolfra\allowbreak{}m-report\allowbreak{}=} switches off. The complete Stage-1 gate, which also runs the Wolfram verifier, is \texttt{bash scr\allowbreak{}ipts/ver\allowbreak{}ify\_stag\allowbreak{}e1\_arbit\allowbreak{}rary\_fie\allowbreak{}ld.sh} (Git Bash) or \texttt{pwsh -No\allowbreak{}Profile \allowbreak{}-File sc\allowbreak{}ripts/ve\allowbreak{}rify\_sta\allowbreak{}ge1\_arbi\allowbreak{}trary\_fi\allowbreak{}eld.ps1} (PowerShell); it ends with the line \texttt{stage1\_a\allowbreak{}rbitrary\allowbreak{}\_field\_v\allowbreak{}erificat\allowbreak{}ion=OK} (Stage 1, §12; Chapter 19). Without the separately published dirac-main folder the comparisons with dirac-main record "not-run" and the gate runs in its public-clone mode, described in Stage 1, §12.

\subsection{What we proved and what we assumed}

\textbf{Proved in this chapter, step by step.} The matrix tools of Section 2.2: the cyclic property of the trace, the mixed-product rule, and, for a matrix with square $\pm1$, the splitting of the space into the two eigenspaces, whose dimensions are fixed by the trace. The formula for the binomial coefficients and the binomial theorem (Section 2.5). A linear expression $\sum_ap_a\gamma^a$ squares to the quadratic form $\sum_a\eta^{aa}p_a^2$ exactly when the matrices satisfy the Clifford relation, in both directions (Section 2.3). The Pauli and Dirac matrices satisfy it (Section 2.3 and Section 2.4). The three monomial rules (R1) to (R3) follow from the relation. For an even number $n$ of generators every monomial except 1 is traceless, the $2^n$ monomials are independent, and the matrices must be at least $2^{n/2}$ by $2^{n/2}$; for 4+4 dimensions this is 16, and 16 by 16 gammas span all 16 by 16 matrices, so only multiples of $I$ commute with all of them (Theorem 2.1, Theorem 2.2, Corollary 2.3). Real 16 by 16 gammas of signature (4,4) exist (Theorem 2.4). Given the Clifford relation of the notebook's gammas: they are symmetric for $a\le3$ and antisymmetric for $a\ge4$; $C$ is symmetric with $C^2=1$ and signature (8,8); expression [1] holds; $\gamma^8$ squares to 1, anticommutes with every $\gamma^a$ and has eight eigenvalues of each sign; the spin generators rotate the gammas, their exponentials are rotations and boosts, a rotation by $2\pi$ equals $-1$ on spinors, and $R^TCR=C$; $B$ is Hermitian with $B^2=1$, signature (8,8) and $BC=-i\gamma^4$. Schur's lemma (first part) and the span criterion (Theorem 2.6, Theorem 2.7). Conjugation by a unit vector is a reflection up to sign; Pin(4,4) and Spin(4,4) are groups with a well-defined parity; the spinor-norm formula (Theorem 2.8); both vector actions are well defined (they do not depend on how $g$ is written as a product) and are homomorphisms into O(4,4) with determinant $(-1)^{|g|}$ (Theorem 2.9) and kernel $\{\pm1\}$ (Theorem 2.10); every element of Pin(4,4) has spinor determinant 1; the products of exponentials all have spinor norm $+1$. Finally the \textbf{representation theorem}: $\mathbb C^{16}$ is an irreducible complex representation of Pin(4,4) with scalar commutant (Theorem 2.11), and under Spin(4,4) (and under its part connected to 1) it is the sum of two irreducible, inequivalent 8-dimensional halves, the images of $P_-$ and $P_+$, with a two-dimensional commutant (Theorem 2.12).

\textbf{Assumed or quoted.} Three facts of elementary linear algebra (dimension counting, solution spaces of linear equations, and the product rule of determinants; Section 2.2) are used without proof. The addition theorems of trigonometry are quoted from Chapter 1 (the double-angle formulas of cosh and sinh are derived from their definitions in Section 2.11). The power series of the exponential, sine, cosine, cosh and sinh functions, and the convergence of the exponential series of a matrix, are quoted from calculus and analysis. The Cartan--Dieudonné theorem is quoted; it is needed only for the word "onto", that is, for the statements that Pin(4,4) and Spin(4,4) are double covers of all of O(4,4) and SO(4,4). Three properties of the notebook's particular matrices, their Clifford relation, $C=\mathrm{diag}(-\sigma,\sigma)$ and $\gamma^8=\mathrm{diag}(-I_8,I_8)$, are taken in this chapter from the exact machine checks and proved by hand in Chapter 3. The conventions (counting from 0, the metric $\eta$ with $x_4$ as time, the notebook's basis as the primary basis, the names of the matrices) are choices, not results. Nothing in this chapter is a physical claim: it is algebra. Triality and the topology of the groups (how many connected pieces Spin(4,4) has) are only described, not proved.

\subsection{Exercises}

\textbf{Exercise 2.1.} Compute $\sigma_y\sigma_z$, $\sigma_z\sigma_y$ and $\sigma_z\sigma_x$, and confirm $\sigma_y\sigma_z=i\sigma_x$ and $\sigma_z\sigma_x=i\sigma_y$. What is $(\sigma_x\sigma_y\sigma_z)^2$?

\textbf{Exercise 2.2.} Show that no real 2 by 2 matrix $M$ anticommutes with both $\sigma_x$ and $\sigma_z$ and satisfies $M^2=I$. Find all complex matrices that do.

\textbf{Exercise 2.3.} Compute $(3\sigma_x+4N)^2$. Of which quadratic form in $(p,q)$ is $p\sigma_x+qN$ a square root?

\textbf{Exercise 2.4.} Show that exactly half of the $2^n$ monomials are even ($n\ge1$). For $n=8$, list the number of monomials of each degree $k=0,\dots,8$.

\textbf{Exercise 2.5.} Use (R2) to count the monomials of $\mathrm{Cl}(4,4)$ that square to $+1$. Then show that for the notebook's gammas a monomial is a symmetric matrix exactly when it squares to $+1$, and deduce from a dimension count, without the first part, that there must be exactly $136=16\cdot17/2$ of them.

\textbf{Exercise 2.6.} Compute $(P\otimes G)(N\otimes P)$ with the mixed-product rule, and check the result by multiplying the two 4 by 4 matrices in block form.

\textbf{Exercise 2.7.} In the tensor picture write $\hat\gamma^1\hat\gamma^5$ as one Kronecker product and verify that it squares to $+1$, as (R2) predicts.

\textbf{Exercise 2.8.} Prove that $C\{\gamma^c,S^{ab}\}$ is symmetric and $C[\gamma^c,S^{ab}]$ is antisymmetric, for all $a,b,c$.

\textbf{Exercise 2.9.} Compute the commutator $[S^{01},S^{12}]$.

\textbf{Exercise 2.10.} Which of the 28 spin generators $S^{ab}$ ($a<b$) commute with $B$? Count them.

\textbf{Exercise 2.11.} Show directly from the definition that the reflection $R_u$ preserves the metric, $\eta(R_uv,R_uw)=\eta(v,w)$, and that $R_u(R_uv)=v$.

\textbf{Exercise 2.12.} For the null vector $u=e_0+e_4$ show that $\gamma(u)^2=0$ and conclude that $\gamma(u)$ has no inverse.

\textbf{Exercise 2.13.} For $g=\gamma^0\gamma^1\gamma^4$ find the parity and the spinor norm, and verify $g^TCg=(-1)^3N(g)C$ by a direct computation. What does $\Psi\mapsto g\Psi$ do to $\Psi^\dagger C\Psi$?

\textbf{Exercise 2.14.} Show that $\exp(4\pi S^{01})=1$ and $\exp(\pi S^{01})=\gamma^0\gamma^1$. Which vector transformation does $\exp(\pi S^{01})$ produce?

\textbf{Exercise 2.15.} Using the tables of Section 2.8, compute $\gamma^0e_3$ and $\gamma^4e_3$, where $e_3$ is the column with a 1 in position 3 and 0 elsewhere. Check that both lie in the chirality $+1$ half, and compute $P_-(e_3+e_{11})$.

\subsection{Answers to the exercises}

\textbf{Answer 2.1.} $\sigma_y\sigma_z=\begin{pmatrix}0&-i\\ i&0\end{pmatrix}\begin{pmatrix}1&0\\0&-1\end{pmatrix}=\begin{pmatrix}0&i\\ i&0\end{pmatrix}=i\sigma_x$, and $\sigma_z\sigma_y=\begin{pmatrix}0&-i\\ -i&0\end{pmatrix}=-i\sigma_x$. Next, $\sigma_z\sigma_x=\begin{pmatrix}1&0\\0&-1\end{pmatrix}\begin{pmatrix}0&1\\1&0\end{pmatrix}=\begin{pmatrix}0&1\\-1&0\end{pmatrix}$, and $i\sigma_y=\begin{pmatrix}0&-i^2\\ i^2&0\end{pmatrix}=\begin{pmatrix}0&1\\-1&0\end{pmatrix}$, the same. Since $\sigma_x\sigma_y\sigma_z=iI$, its square is $i^2I=-I$.

\textbf{Answer 2.2.} Write $M=\begin{pmatrix}\alpha&\beta\\ \gamma&\delta\end{pmatrix}$. Then $\sigma_zM+M\sigma_z=\begin{pmatrix}2\alpha&0\\0&-2\delta\end{pmatrix}$, so $\alpha=\delta=0$. With $M=\begin{pmatrix}0&\beta\\ \gamma&0\end{pmatrix}$, $\sigma_xM+M\sigma_x=\begin{pmatrix}\beta+\gamma&0\\0&\beta+\gamma\end{pmatrix}$, so $\gamma=-\beta$ and $M=\beta N$. Then $M^2=\beta^2N^2=-\beta^2I$, which equals $I$ only for $\beta^2=-1$. No real $\beta$ does this; the complex solutions are $\beta=\pm i$, that is, $M=\pm iN=\mp\sigma_y$.

\textbf{Answer 2.3.} $3\sigma_x+4N=\begin{pmatrix}0&7\\-1&0\end{pmatrix}$ and its square is $\begin{pmatrix}-7&0\\0&-7\end{pmatrix}=-7I=(3^2-4^2)I$. In general $(p\sigma_x+qN)^2=(p^2-q^2)I$, a square root of $p^2-q^2$ (Section 2.3).

\textbf{Answer 2.4.} By the binomial theorem, $0=(1-1)^n=\sum_k\binom nk(-1)^k$, so the number of monomials of even degree equals the number of odd degree; they add up to $2^n$, so each is $2^{n-1}$. For $n=8$ the numbers $\binom8k$ are 1, 8, 28, 56, 70, 56, 28, 8, 1 for $k=0,\dots,8$; the even degrees give $1+28+70+28+1=128$ and the odd ones $8+56+56+8=128$.

\textbf{Answer 2.5.} By (R2) the sign of $\gamma_A^2$ is $(-1)^{k(k-1)/2}$ times $(-1)^t$, where $t$ is the number of time-like indices (indices $\ge4$) in $A$ and $s=k-t$ the number of space-like ones. Counting subsets with $s$ space-like and $t$ time-like indices, $\binom4s\binom4t$ of them, the monomials with square $+1$ are, by degree: $k=0$: 1; $k=1$: $t=0$, 4; $k=2$ (sign $-(-1)^t$): $t=1$, $4\cdot4=16$; $k=3$ (sign $-(-1)^t$): $t=1$ or 3, $6\cdot4+1\cdot4=28$; $k=4$ (sign $+(-1)^t$): $t=0,2,4$, $1+36+1=38$; $k=5$ (sign $+(-1)^t$): $t=2$ or 4, $4\cdot6+4\cdot1=28$; $k=6$ (sign $-(-1)^t$): $t=3$, $4\cdot4=16$; $k=7$ (sign $-(-1)^t$): $t=3$, 4; $k=8$: 1. The total is $1+4+16+28+38+28+16+4+1=136$. For the second part: every monomial is a product of orthogonal matrices (N2), hence orthogonal, so $\gamma_A^T=\gamma_A^{-1}=\pm\gamma_A$ with the sign of $\gamma_A^2$. So a monomial is symmetric exactly when its square is $+1$, and antisymmetric otherwise. The symmetric monomials are independent elements of the space of symmetric 16 by 16 matrices, whose dimension is $16\cdot17/2=136$ (the free entries on and above the diagonal), so there are at most 136 of them; likewise at most $16\cdot15/2=120$ antisymmetric ones. Since there are $256=136+120$ monomials in all, both bounds are reached. The notebook counts the same 136 and 120 (its cells 597 to 599 and 608 to 609, according to the notebook survey).

\textbf{Answer 2.6.} $(P\otimes G)(N\otimes P)=(PN)\otimes(GP)=(-G)\otimes N$, because $PN=-G$ and $GP=N$. In block form $P\otimes G=\begin{pmatrix}0&G\\ G&0\end{pmatrix}$ and $N\otimes P=\begin{pmatrix}0&P\\-P&0\end{pmatrix}$, whose product is $\begin{pmatrix}-GP&0\\0&GP\end{pmatrix}=\begin{pmatrix}-N&0\\0&N\end{pmatrix}$; and $(-G)\otimes N$ has the blocks $-1\cdot N$ and $+1\cdot N$ on the diagonal, the same matrix.

\textbf{Answer 2.7.} $\hat\gamma^1=G\otimes P\otimes I_2\otimes I_2$ and $\hat\gamma^5=\hat\gamma_2^-=G\otimes N\otimes I_2\otimes I_2$, so $\hat\gamma^1\hat\gamma^5=G^2\otimes PN\otimes I_2\otimes I_2=I_2\otimes(-G)\otimes I_2\otimes I_2$. Its square is $I_2\otimes G^2\otimes I_2\otimes I_2=I_{16}$. (R2) with $k=2$ predicts $(-1)^1\eta^{11}\eta^{55}=(-1)(+1)(-1)=+1$.

\textbf{Answer 2.8.} Put $A=C\gamma^c$ and $\Sigma=CS^{ab}$; both are antisymmetric, by expression [1] and by (S5). Since $C^2=1$, $C\gamma^cS^{ab}=C\gamma^cCCS^{ab}=AC\Sigma$, and $(AC\Sigma)^T=\Sigma^TC^TA^T=(-\Sigma)C(-A)=\Sigma CA=CS^{ab}CC\gamma^c=CS^{ab}\gamma^c$. So $(C\gamma^cS^{ab})^T=CS^{ab}\gamma^c$, and transposing $C\{\gamma^c,S^{ab}\}=C\gamma^cS^{ab}+CS^{ab}\gamma^c$ gives $CS^{ab}\gamma^c+C\gamma^cS^{ab}$, the same matrix (symmetric), while transposing $C[\gamma^c,S^{ab}]=C\gamma^cS^{ab}-CS^{ab}\gamma^c$ gives $CS^{ab}\gamma^c-C\gamma^cS^{ab}$, its negative (antisymmetric).

\textbf{Answer 2.9.} $S^{01}S^{12}=\tfrac14\gamma^0\gamma^1\gamma^1\gamma^2=\tfrac14\gamma^0\gamma^2$. For the other order, $\gamma^1\gamma^2\gamma^0\gamma^1$: move the last $\gamma^1$ to the front past $\gamma^0$ and $\gamma^2$ (two signs), which gives $\gamma^1\gamma^1\gamma^2\gamma^0=\gamma^2\gamma^0=-\gamma^0\gamma^2$; so $S^{12}S^{01}=-\tfrac14\gamma^0\gamma^2$. Hence $[S^{01},S^{12}]=\tfrac12\gamma^0\gamma^2=S^{02}$. Two infinitesimal rotations in the planes (0,1) and (1,2) fail to commute by an infinitesimal rotation in the plane (0,2), exactly as for rotations of ordinary space.

\textbf{Answer 2.10.} $B$ is $-i$ times the monomial $\gamma_A$ with $A=\{0,1,2,3,4\}$ of degree 5. By (R3), $\gamma^c\gamma_A=(-1)^4\gamma_A\gamma^c=\gamma_A\gamma^c$ if $c\in A$ and $\gamma^c\gamma_A=(-1)^5\gamma_A\gamma^c=-\gamma_A\gamma^c$ if $c\notin A$. So $S^{ab}=\tfrac12\gamma^a\gamma^b$ commutes with $B$ when $a,b$ are both in $A$ or both outside $A$ (the two signs cancel), and anticommutes when exactly one is in $A$. The commuting ones are the $\binom52=10$ pairs inside $\{0,1,2,3,4\}$ and the $\binom32=3$ pairs inside $\{5,6,7\}$: 13 in all. This is erratum E1 of the contract: 13, not 9 (Stage 1, Result 10.6); the four boosts $S^{a4}$ with $a\le3$ are among the 13.

\textbf{Answer 2.11.} Write $c=\eta(u,u)=\pm1$. Expanding with the bilinearity of $\eta$,

\[
\eta(R_uv,R_uw)=\eta(v,w)-\frac{2\eta(u,w)\eta(v,u)}{c}-\frac{2\eta(u,v)\eta(u,w)}{c}+\frac{4\eta(u,v)\eta(u,w)\eta(u,u)}{c^2}=\eta(v,w),
\]

because $\eta(u,u)/c^2=1/c$ and the last three terms add up to $(-4+4)\eta(u,v)\eta(u,w)/c=0$. Next, $\eta(u,R_uv)=\eta(u,v)-2\eta(u,v)=-\eta(u,v)$, so $R_u(R_uv)=R_uv+2\frac{\eta(u,v)}cu=v$.

\textbf{Answer 2.12.} $\gamma(u)^2=(\gamma^0+\gamma^4)^2=(\gamma^0)^2+(\gamma^4)^2+\{\gamma^0,\gamma^4\}=1-1+0=0$. If $\gamma(u)$ had an inverse, then $\gamma(u)=\gamma(u)^{-1}\gamma(u)^2=0$; but $\gamma(u)$ is a nonzero combination of two independent monomials. So it has no inverse.

\textbf{Answer 2.13.} Three factors: odd parity. $N(g)=n(e_0)n(e_1)n(e_4)=(+1)(+1)(-1)=-1$, so the formula predicts $g^TCg=(-1)^3(-1)C=C$. Directly: $g^T=(\gamma^4)^T(\gamma^1)^T(\gamma^0)^T=-\gamma^4\gamma^1\gamma^0$ by (N3). From (C1), $C\gamma^0=-\gamma^0C$, $C\gamma^1=-\gamma^1C$ and $C\gamma^4=\gamma^4C$. So $\gamma^0C\gamma^0=-\gamma^0\gamma^0C=-C$, then $\gamma^1(-C)\gamma^1=+\gamma^1\gamma^1C=C$, then $-\gamma^4C\gamma^4=-\gamma^4\gamma^4C=C$. Hence $g^TCg=-\gamma^4\gamma^1(\gamma^0C\gamma^0)\gamma^1\gamma^4=C$. Under $\Psi\mapsto g\Psi$ the bilinear $\Psi^\dagger C\Psi$ becomes $\Psi^\dagger g^TCg\Psi=\Psi^\dagger C\Psi$: it is unchanged.

\textbf{Answer 2.14.} By (S4), $\exp(\theta S^{01})=\cos\tfrac\theta2+\gamma^0\gamma^1\sin\tfrac\theta2$. For $\theta=4\pi$: $\cos2\pi+\gamma^0\gamma^1\sin2\pi=1$. For $\theta=\pi$: $\cos\tfrac\pi2+\gamma^0\gamma^1\sin\tfrac\pi2=\gamma^0\gamma^1$. By the worked example of Section 2.11 it maps $\gamma^0\to\cos\pi\,\gamma^0-\sin\pi\,\gamma^1=-\gamma^0$ and $\gamma^1\to-\gamma^1$ and leaves the other six directions alone: the rotation by $\pi$ in the plane (0,1), $(v_0,v_1)\mapsto(-v_0,-v_1)$.

\textbf{Answer 2.15.} $\gamma^0e_3$ is column 3 of $\gamma^0$. The table says that row 11 of $\gamma^0$ has its entry $+1$ in column 3 (entry $+3$), and no other row points to column 3; so $\gamma^0e_3=e_{11}$. In the same way row 14 of $\gamma^4$ has the entry $+3$, so $\gamma^4e_3=e_{14}$. Both $e_{11}$ and $e_{14}$ have their nonzero component among positions 8 to 15, the chirality $+1$ half, while $e_3$ lies in the $-1$ half, as (X6) requires. Finally $P_-=\mathrm{diag}(I_8,0)$ keeps the upper eight components and deletes the lower ones, so $P_-(e_3+e_{11})=e_3$.

\section{Split octonions and the notebook's construction of the gamma matrices}

\subsection{What this chapter does}

Chapter 2 used the notebook's sixteen by sixteen gamma matrices $\gamma^a=(\mathrm{T16}^A)[a]$ as finished objects and took three of their properties from the exact machine checks: the Clifford relation, the form $C=\mathrm{diag}(-\sigma,\sigma)$ of the charge matrix, and the form $\gamma^8=\mathrm{diag}(-I_8,I_8)$ of the chirality. This chapter opens the box. The notebook does not write its gammas down entry by entry; it builds them from eight 8 by 8 matrices $\tau_0,\dots,\tau_7$, which in turn are built from six 4 by 4 matrices. The notebook's own section titles connect this construction with the \textbf{split octonions}, an 8-dimensional number system whose natural metric is exactly the metric $\eta$ of signature (4,4) used in this book (the notebook's sections "For Spin(4,4); τ tau;T16;OCTAD:Nash", "split octonions; evalues, evecs of σ" and "split octonion multiplication constants", cells 308, 837 and 891 according to the notebook survey in \texttt{handoff/\allowbreak{}surveys/}).

The chapter has four goals.

\begin{enumerate}
\item Build the 4 by 4 blocks and the $\tau_a$ explicitly, recognize the blocks as multiplications of \textbf{quaternions}, and prove by hand the three properties that Chapter 2 borrowed from the computer (Section 3.3 to Section 3.6).
\item Introduce the split octonions from zero, in the vector-matrix form of Max Zorn, and prove their basic identities (Section 3.7 and Section 3.8).
\item Show that the octonions give a third set of gamma matrices, and that the notebook's $\tau_a$ are exactly octonion multiplications written in a \textbf{light-cone basis} (Section 3.9 and Section 3.12).
\item Prove that any two sets of 16 by 16 gamma matrices of signature (4,4) are related by an \textbf{intertwiner} that exists and is unique up to a factor, and present the two intertwiners of the project, $K_{\mathrm{clifford}}$ and $K_{\mathrm{octonion}}$, with the directions fixed by the errata E4 and E5 of the contract (Section 3.10 to Section 3.13).
\end{enumerate}

The sources of this chapter are these files of the repository:

\begin{itemize}
\item the Stage-1 document \texttt{provenan\allowbreak{}ce/DIRAC\allowbreak{}16COMPLE\allowbreak{}X\_ARBITR\allowbreak{}ARY\_FIEL\allowbreak{}D.md}, its §3 (Results 3.1 to 3.8);
\item the contract \texttt{handoff/\allowbreak{}specs/CO\allowbreak{}NTRACT.m\allowbreak{}d}, its §1 and the errata E4 and E5 in §11;
\item the exact Python module \texttt{scripts/\allowbreak{}d16c\_exa\allowbreak{}ct.py} (the constructions of the blocks, the $\tau_a$, the gammas of all three pictures, the Zorn product and the intertwiners) and the Wolfram package \texttt{wolfram/\allowbreak{}Dirac16C\allowbreak{}omplexAl\allowbreak{}gebra.wl};
\item the reports \texttt{python-a\allowbreak{}lgebra-r\allowbreak{}eport.js\allowbreak{}on} and \texttt{wolfram-\allowbreak{}algebra-\allowbreak{}report.j\allowbreak{}son} in \texttt{artifact\allowbreak{}s/dirac1\allowbreak{}6complex\allowbreak{}/arbitra\allowbreak{}ry-field\allowbreak{}/}. As in Chapter 2, a check name refers to the check of that name in both algebra reports.
\end{itemize}

\textbf{Labelled exceptions to counting from 0.} The notebook's 4 by 4 building blocks use Mathematica's 1-based indices $p,q\in\{1,2,3,4\}$ and $h\in\{1,2,3\}$ (Stage 1, §2.1). We keep them in Section 3.3 and Section 3.4, because the formulas are the notebook's. The three components of the ordinary three-dimensional vectors of Section 3.2 and Section 3.8 are also labelled 1, 2, 3, as usual. Everything built from these (the $\tau_a$, the octonion basis and the gammas) is counted from 0.

\subsection{Tools: permutation signs, and vectors in three dimensions}

\textbf{Permutations and their sign.} A permutation of a list of distinct numbers is a reordering of it. An \textbf{inversion} is a pair of positions in which the larger number stands first. The \textbf{sign} of the permutation is $+1$ if the number of inversions is even and $-1$ if it is odd. Examples (1-based): $(1,2,3,4)$ has no inversion, sign $+1$; $(2,1,3,4)$ has one (the pair 2,1), sign $-1$; $(2,3,1,4)$ has two (2,1 and 3,1), sign $+1$; $(1,4,3,2)$ has three (4,3; 4,2; 3,2), sign $-1$. Exchanging two entries always changes the sign (Exercise 3.1). Mathematica's function Signature returns this sign for a list of distinct entries and 0 for a list with a repeated entry. The notebook uses

\[
\epsilon_{hpq4}:=\mathrm{Signature}[\{h,p,q,4\}],
\]

which is $\pm1$ when $h,p,q,4$ are all different and 0 otherwise. The \textbf{Kronecker delta} is $\delta_{pq}=1$ for $p=q$ and $0$ otherwise.

\textbf{Vectors in three dimensions.} For $u=(u_1,u_2,u_3)$ and $r=(r_1,r_2,r_3)$ (three components, labelled 1 to 3 in the usual way) the \textbf{dot product} and the \textbf{cross product} are

\[
u\cdot r=u_1r_1+u_2r_2+u_3r_3,\qquad u\times r=\bigl(u_2r_3-u_3r_2,\ u_3r_1-u_1r_3,\ u_1r_2-u_2r_1\bigr).
\]

The dot product is a number and symmetric, $u\cdot r=r\cdot u$; the cross product is a vector and antisymmetric, $u\times r=-r\times u$. We need four identities, for all vectors $a,b,c,d$:

\[
\begin{aligned}
&\text{(V1)}\quad a\times a=0,\qquad a\cdot(a\times b)=0,\\
&\text{(V2)}\quad a\cdot(b\times c)=(a\times b)\cdot c,\\
&\text{(V3)}\quad a\times(b\times c)=b\,(a\cdot c)-c\,(a\cdot b),\\
&\text{(V4)}\quad (a\times b)\cdot(c\times d)=(a\cdot c)(b\cdot d)-(a\cdot d)(b\cdot c).
\end{aligned}
\]

\emph{Proofs.} (V1): $a\times a=-a\times a$ by antisymmetry, so it is zero; and $a\cdot(a\times b)=a_1(a_2b_3-a_3b_2)+a_2(a_3b_1-a_1b_3)+a_3(a_1b_2-a_2b_1)$, in which the six terms cancel in pairs. (V2): both sides expand to the same six terms,

\[
a_1b_2c_3-a_1b_3c_2+a_2b_3c_1-a_2b_1c_3+a_3b_1c_2-a_3b_2c_1 .
\]

(V3): the first component of the left side is $a_2(b\times c)_3-a_3(b\times c)_2=a_2(b_1c_2-b_2c_1)-a_3(b_3c_1-b_1c_3)$, and the first component of the right side is $b_1(a_1c_1+a_2c_2+a_3c_3)-c_1(a_1b_1+a_2b_2+a_3b_3)$; in the right side the terms $a_1b_1c_1$ cancel and what remains is $a_2b_1c_2+a_3b_1c_3-a_2b_2c_1-a_3b_3c_1$, the same four terms. The other two components follow by renaming $1\to2\to3\to1$, which does not change either side's form. (V4): (V2) with $c$ replaced by $c\times d$ gives $(a\times b)\cdot(c\times d)=a\cdot\bigl(b\times(c\times d)\bigr)$; by (V3), with $a,b,c$ renamed $b,c,d$, $b\times(c\times d)=c\,(b\cdot d)-d\,(b\cdot c)$, and taking the dot product with $a$ gives $a\cdot\bigl(c\,(b\cdot d)-d\,(b\cdot c)\bigr)=(a\cdot c)(b\cdot d)-(a\cdot d)(b\cdot c)$. $\square$

\subsection{Quaternions}

\textbf{Definition.} A \textbf{quaternion} is an expression $q=q_4+q_1\mathbf i+q_2\mathbf j+q_3\mathbf k$ with four real numbers $q_1,\dots,q_4$ (the 1-based labels match the notebook's blocks, with $q_4$ the "real part"). Quaternions are added componentwise and multiplied with the rules of William Rowan Hamilton (1843),

\[
\mathbf i^2=\mathbf j^2=\mathbf k^2=-1,\qquad\mathbf{ij}=\mathbf k=-\mathbf{ji},\qquad\mathbf{jk}=\mathbf i=-\mathbf{kj},\qquad\mathbf{ki}=\mathbf j=-\mathbf{ik},
\]

extended by the distributive law, with real numbers commuting with everything. The quaternion units $\mathbf i,\mathbf j,\mathbf k$ are written in bold to keep them apart from the complex number $i$. For $h\ne k$ in $\{1,2,3\}$ the rules say $u_hu_k=\sum_l\epsilon_{hkl}u_l$, where $u_1=\mathbf i$, $u_2=\mathbf j$, $u_3=\mathbf k$ and $\epsilon_{hkl}$ is the sign of $(h,k,l)$ (and 0 for a repeated index).

\textbf{Quaternions are associative.} The map

\[
1\mapsto I_2,\qquad\mathbf i\mapsto-i\sigma_x,\qquad\mathbf j\mapsto-i\sigma_y,\qquad\mathbf k\mapsto-i\sigma_z
\]

to complex 2 by 2 matrices respects the rules: $(-i\sigma_x)^2=-\sigma_x^2=-I_2$, and $(-i\sigma_x)(-i\sigma_y)=-\sigma_x\sigma_y=-i\sigma_z$, the image of $\mathbf k$; the other rules follow in the same way from the Pauli products of Section 2.3. The four matrices $I_2,-i\sigma_x,-i\sigma_y,-i\sigma_z$ are independent over the real numbers (their real combination $\begin{pmatrix}q_4-iq_3&-q_2-iq_1\\ q_2-iq_1&q_4+iq_3\end{pmatrix}$ vanishes only if all $q_p=0$), so a quaternion is determined by its matrix, and the product of quaternions is the product of their matrices. Matrix multiplication is associative, hence so is quaternion multiplication: $(xy)z=x(yz)$. Quaternions are, however, not commutative ($\mathbf{ij}\ne\mathbf{ji}$).

\textbf{Left and right multiplication.} Identify the quaternion $q$ with the column $(q_1,q_2,q_3,q_4)^T$. For a fixed quaternion $x$, the maps $q\mapsto xq$ and $q\mapsto qx$ are linear, so they are 4 by 4 matrices, $L_x$ and $R_x$. Associativity gives three rules:

\[
L_xL_y=L_{xy},\qquad R_xR_y=R_{yx},\qquad L_xR_y=R_yL_x,
\]

because $x(yq)=(xy)q$, $(qy)x=q(yx)$ and $x(qy)=(xq)y$.

\textbf{Worked example.} Multiply a general $q$ on the right by $\mathbf i$:

\[
q\,\mathbf i=q_1\mathbf i^2+q_2\mathbf{ji}+q_3\mathbf{ki}+q_4\mathbf i=-q_1-q_2\mathbf k+q_3\mathbf j+q_4\mathbf i .
\]

The components in the order $(\mathbf i,\mathbf j,\mathbf k,1)$ are $(q_4,\ q_3,\ -q_2,\ -q_1)$. In the same way one finds all six products with a unit:

\begingroup
\small
\begin{longtable}{@{}>{\raggedright\arraybackslash}p{0.160\linewidth}>{\raggedright\arraybackslash}p{0.160\linewidth}>{\raggedright\arraybackslash}p{0.160\linewidth}>{\raggedright\arraybackslash}p{0.160\linewidth}>{\raggedright\arraybackslash}p{0.160\linewidth}@{}}
\toprule
\textbf{product} & \textbf{component $\mathbf i$} & \textbf{component $\mathbf j$} & \textbf{component $\mathbf k$} & \textbf{component 1} \\
\midrule
\endfirsthead
\toprule
\textbf{product} & \textbf{component $\mathbf i$} & \textbf{component $\mathbf j$} & \textbf{component $\mathbf k$} & \textbf{component 1} \\
\midrule
\endhead
$q\,\mathbf i$ & $q_4$ & $q_3$ & $-q_2$ & $-q_1$ \\
$q\,\mathbf j$ & $-q_3$ & $q_4$ & $q_1$ & $-q_2$ \\
$q\,\mathbf k$ & $q_2$ & $-q_1$ & $q_4$ & $-q_3$ \\
$\mathbf i\,q$ & $q_4$ & $-q_3$ & $q_2$ & $-q_1$ \\
$\mathbf j\,q$ & $q_3$ & $q_4$ & $-q_1$ & $-q_2$ \\
$\mathbf k\,q$ & $-q_2$ & $q_1$ & $q_4$ & $-q_3$ \\
\bottomrule
\end{longtable}
\endgroup

\subsection{The notebook's 4 by 4 building blocks}

\textbf{Definition} (notebook; Stage 1, §3.1; contract §1). For $h\in\{1,2,3\}$ and $p,q\in\{1,2,3,4\}$ (1-based),

\[
Q_a[h]_{pq}=\epsilon_{hpq4},\qquad Q_b[h]_{pq}=\delta_{p4}\delta_{qh}-\delta_{ph}\delta_{q4},\qquad s_4[h]=Q_a[h]-Q_b[h],\qquad t_4[h]=Q_a[h]+Q_b[h].
\]

$Q_a[h]$ has nonzero entries only in the rows and columns different from $h$ and 4, where it is the sign of $(h,p,q,4)$; $Q_b[h]$ has the entry $+1$ in position $(4,h)$ and $-1$ in position $(h,4)$. For $h=1$, for instance, $Q_a[1]_{23}=\epsilon_{1234}=+1$ and $Q_a[1]_{32}=\epsilon_{1324}=-1$. The six resulting matrices are

\[
s_4[1]=\begin{pmatrix}0&0&0&1\\0&0&1&0\\0&-1&0&0\\-1&0&0&0\end{pmatrix},\quad s_4[2]=\begin{pmatrix}0&0&-1&0\\0&0&0&1\\1&0&0&0\\0&-1&0&0\end{pmatrix},\quad s_4[3]=\begin{pmatrix}0&1&0&0\\-1&0&0&0\\0&0&0&1\\0&0&-1&0\end{pmatrix},
\]

\[
t_4[1]=\begin{pmatrix}0&0&0&-1\\0&0&1&0\\0&-1&0&0\\1&0&0&0\end{pmatrix},\quad t_4[2]=\begin{pmatrix}0&0&-1&0\\0&0&0&-1\\1&0&0&0\\0&1&0&0\end{pmatrix},\quad t_4[3]=\begin{pmatrix}0&1&0&0\\-1&0&0&0\\0&0&0&-1\\0&0&1&0\end{pmatrix}.
\]

(They are produced by the functions \texttt{notebook\allowbreak{}\_s4} and \texttt{notebook\allowbreak{}\_t4} of \texttt{scripts/\allowbreak{}d16c\_exa\allowbreak{}ct.py}.) The contract calls the $s_4[h]$ "self-dual" and the $t_4[h]$ "anti-self-dual"; Exercise 3.5 explains the words. We write $s_h:=s_4[h]$ and $t_h:=t_4[h]$.

\textbf{Theorem 3.1 (the blocks are quaternion multiplications).} With the identification of Section 3.3,

\[
s_1=R_{\mathbf i},\quad s_2=R_{\mathbf j},\quad s_3=R_{\mathbf k},\qquad t_1=-L_{\mathbf i},\quad t_2=-L_{\mathbf j},\quad t_3=-L_{\mathbf k}.
\]

\emph{Proof.} The $p$-th row of $R_{\mathbf i}$ lists how component $p$ of $q\,\mathbf i$ is made from $q_1,\dots,q_4$. By the table of Section 3.3, $(q\,\mathbf i)=(q_4,q_3,-q_2,-q_1)$, so the rows of $R_{\mathbf i}$ are $(0,0,0,1)$, $(0,0,1,0)$, $(0,-1,0,0)$, $(-1,0,0,0)$: this is $s_1$. Likewise $\mathbf i\,q=(q_4,-q_3,q_2,-q_1)$ gives $L_{\mathbf i}$ with rows $(0,0,0,1)$, $(0,0,-1,0)$, $(0,1,0,0)$, $(-1,0,0,0)$, whose negative is $t_1$. The four remaining cases are read off the table in the same way. $\square$

\textbf{Corollary 3.2 (the algebra of the blocks).} For $h,k\in\{1,2,3\}$:

\begin{enumerate}
\item $s_h^2=t_h^2=-I_4$;
\item $s_hs_k=-\sum_l\epsilon_{hkl}s_l$ and $t_ht_k=-\sum_l\epsilon_{hkl}t_l$ for $h\ne k$; in particular $s_1s_2=-s_3$ and $t_1t_2=-t_3$;
\item $s_ht_k=t_ks_h$ for all $h,k$;
\item $s_1s_2s_3=I_4$ and $t_3t_2t_1=-I_4$;
\item all six matrices are antisymmetric signed permutation matrices, hence orthogonal.
\end{enumerate}

\emph{Proof.} (1) $s_h^2=R_{u_h}R_{u_h}=R_{u_h^2}=R_{-1}=-I_4$, and $t_h^2=L_{u_h}L_{u_h}=L_{u_h^2}=-I_4$. (2) $s_hs_k=R_{u_h}R_{u_k}=R_{u_ku_h}$ and $u_ku_h=-u_hu_k=-\sum_l\epsilon_{hkl}u_l$; and $t_ht_k=(-L_{u_h})(-L_{u_k})=L_{u_hu_k}=\sum_l\epsilon_{hkl}L_{u_l}=-\sum_l\epsilon_{hkl}t_l$. (3) $s_ht_k=-R_{u_h}L_{u_k}=-L_{u_k}R_{u_h}=t_ks_h$, by associativity. (4) $s_1s_2s_3=(-s_3)s_3=-s_3^2=I_4$, and $t_3t_2=-\epsilon_{321}t_1=t_1$, so $t_3t_2t_1=t_1^2=-I_4$. (5) Read off the matrices above. $\square$

As a check of (2) by direct multiplication:

\[
s_1s_2=\begin{pmatrix}0&0&0&1\\0&0&1&0\\0&-1&0&0\\-1&0&0&0\end{pmatrix}\begin{pmatrix}0&0&-1&0\\0&0&0&1\\1&0&0&0\\0&-1&0&0\end{pmatrix}=\begin{pmatrix}0&-1&0&0\\1&0&0&0\\0&0&0&-1\\0&0&1&0\end{pmatrix}=-s_3 .
\]

(Row 1 of $s_1$ picks row 4 of $s_2$, which is $(0,-1,0,0)$; row 2 picks row 3, $(1,0,0,0)$; row 3 picks minus row 2, $(0,0,0,-1)$; row 4 picks minus row 1, $(0,0,1,0)$.)

\subsection{The eight \texorpdfstring{$\tau$}{τ} matrices}

From here on everything is counted from 0 again. The notebook assembles 8 by 8 matrices from the 4 by 4 blocks (Stage 1, §3.1):

\[
\tau_0=I_8,\qquad\tau_h=\begin{pmatrix}0&s_h\\ s_h&0\end{pmatrix},\qquad\tau_{7-h}=\begin{pmatrix}0&t_h\\-t_h&0\end{pmatrix}\quad(h=1,2,3),\qquad\tau_7=\tau_1\tau_2\tau_3\tau_4\tau_5\tau_6 ,
\]

together with

\[
\sigma=\begin{pmatrix}0&I_4\\ I_4&0\end{pmatrix},\qquad\bar\tau_0=I_8,\qquad\bar\tau_a=\sigma\,\tau_a^T\,\sigma\quad(a=1,\dots,7).
\]

So $\tau_1,\tau_2,\tau_3$ are made from $s_1,s_2,s_3$, while $\tau_6$, $\tau_5$, $\tau_4$ are made from $t_1$, $t_2$, $t_3$ (in this order, because of the index $7-h$). In the notebook $\bar\tau_a$ is OverBar[τ][a]. As signed permutation matrices (read as in Section 2.8: the entry $-6$ in row 0 of $\tau_2$ means $(\tau_2u)_0=-u_6$):

\begingroup
\small
\begin{longtable}{@{}>{\raggedright\arraybackslash}p{0.100\linewidth}>{\raggedright\arraybackslash}p{0.100\linewidth}>{\raggedright\arraybackslash}p{0.100\linewidth}>{\raggedright\arraybackslash}p{0.100\linewidth}>{\raggedright\arraybackslash}p{0.100\linewidth}>{\raggedright\arraybackslash}p{0.100\linewidth}>{\raggedright\arraybackslash}p{0.100\linewidth}>{\raggedright\arraybackslash}p{0.100\linewidth}@{}}
\toprule
\textbf{$i$} & \textbf{$\tau_1$} & \textbf{$\tau_2$} & \textbf{$\tau_3$} & \textbf{$\tau_4$} & \textbf{$\tau_5$} & \textbf{$\tau_6$} & \textbf{$\tau_7$} \\
\midrule
\endfirsthead
\toprule
\textbf{$i$} & \textbf{$\tau_1$} & \textbf{$\tau_2$} & \textbf{$\tau_3$} & \textbf{$\tau_4$} & \textbf{$\tau_5$} & \textbf{$\tau_6$} & \textbf{$\tau_7$} \\
\midrule
\endhead
0 & $+7$ & $-6$ & $+5$ & $+5$ & $-6$ & $-7$ & $-0$ \\
1 & $+6$ & $+7$ & $-4$ & $-4$ & $-7$ & $+6$ & $-1$ \\
2 & $-5$ & $+4$ & $+7$ & $-7$ & $+4$ & $-5$ & $-2$ \\
3 & $-4$ & $-5$ & $-6$ & $+6$ & $+5$ & $+4$ & $-3$ \\
4 & $+3$ & $-2$ & $+1$ & $-1$ & $+2$ & $+3$ & $+4$ \\
5 & $+2$ & $+3$ & $-0$ & $+0$ & $+3$ & $-2$ & $+5$ \\
6 & $-1$ & $+0$ & $+3$ & $+3$ & $-0$ & $+1$ & $+6$ \\
7 & $-0$ & $-1$ & $-2$ & $-2$ & $-1$ & $-0$ & $+7$ \\
\bottomrule
\end{longtable}
\endgroup

(Generated by \texttt{notebook\allowbreak{}\_tau} of \texttt{scripts/\allowbreak{}d16c\_exa\allowbreak{}ct.py}; stored under the key \texttt{\detokenize{tau}} of \texttt{algebra-\allowbreak{}fixture.\allowbreak{}json}. The blocks $s_h$ and $t_h$ can be read off from them.) For example, rows 0 to 3 of $\tau_1$ are the upper-right block $s_1$ shifted by four columns: row 0 of $s_1$ is $(0,0,0,1)$ in 1-based columns, that is, column 3 in 0-based counting, and in $\tau_1$ it becomes column $3+4=7$.

\textbf{Theorem 3.3 (properties of the $\tau$'s).}

\begin{enumerate}
\item $\tau_1,\tau_2,\tau_3$ are antisymmetric, and $\tau_4,\dots,\tau_7$ are symmetric.
\item $\bar\tau_a=-\tau_a$ for $a=1,\dots,7$ (and $\bar\tau_0=\tau_0=I_8$).
\item For $a,b\in\{1,\dots,7\}$: $\tau_a\tau_b+\tau_b\tau_a=-2\eta_{ab}I_8$. In words: $\tau_1^2=\tau_2^2=\tau_3^2=-I_8$, $\tau_4^2=\dots=\tau_7^2=+I_8$, and any two different ones anticommute.
\item $\tau_1\tau_2\tau_3=\sigma$ and $\tau_7=\mathrm{diag}(-I_4,I_4)$.
\item $\tau_1\tau_2\tau_3\tau_4\tau_5\tau_6\tau_7=I_8$.
\end{enumerate}

\emph{Proof.} We use the block rules of Section 2.2 and Corollary 3.2; $s$ stands for any $s_h$ and $t$ for any $t_k$. Statement (4) is used in (1) to (3) for $\tau_7$; its proof below uses only Corollary 3.2, so there is no circle.

(1) With $s^T=-s$: $\begin{pmatrix}0&s\\ s&0\end{pmatrix}^T=\begin{pmatrix}0&s^T\\ s^T&0\end{pmatrix}=-\begin{pmatrix}0&s\\ s&0\end{pmatrix}$. With $t^T=-t$: $\begin{pmatrix}0&t\\-t&0\end{pmatrix}^T=\begin{pmatrix}0&-t^T\\ t^T&0\end{pmatrix}=\begin{pmatrix}0&t\\-t&0\end{pmatrix}$. And $\tau_7$ is diagonal by (4).

(2) Conjugation by $\sigma$ exchanges the two halves: $\sigma\begin{pmatrix}A&B\\ D&E\end{pmatrix}\sigma=\begin{pmatrix}E&D\\ B&A\end{pmatrix}$. For $a=1,2,3$, $\tau_a^T=-\tau_a$ and $\sigma\tau_a\sigma=\tau_a$ (both off-diagonal blocks are $s$), so $\bar\tau_a=\sigma\tau_a^T\sigma=-\tau_a$. For $a=4,5,6$, $\tau_a^T=\tau_a$ and $\sigma\begin{pmatrix}0&t\\-t&0\end{pmatrix}\sigma=\begin{pmatrix}0&-t\\ t&0\end{pmatrix}=-\tau_a$, so again $\bar\tau_a=-\tau_a$. For $a=7$, $\tau_7=\mathrm{diag}(-I_4,I_4)$ by (4), and $\sigma\tau_7\sigma=\mathrm{diag}(I_4,-I_4)=-\tau_7$.

(3) \emph{Squares.} $\tau_h^2=\mathrm{diag}(s_h^2,s_h^2)=-I_8$; for $a=7-h$, $\tau_a^2=\begin{pmatrix}0&t_h\\-t_h&0\end{pmatrix}^2=\mathrm{diag}(-t_h^2,-t_h^2)=+I_8$; and $\tau_7^2=I_8$ by (4). These are $-\eta_{aa}I_8$. \emph{Two $s$-type matrices}, $h\ne k$: $\tau_h\tau_k=\mathrm{diag}(s_hs_k,s_hs_k)$, and $s_hs_k=-s_ks_h$ by Corollary 3.2 (2), so they anticommute. \emph{Two $t$-type matrices}, built from $t$ and $t'$: $\tau\tau'=\begin{pmatrix}0&t\\-t&0\end{pmatrix}\begin{pmatrix}0&t'\\-t'&0\end{pmatrix}=\mathrm{diag}(-tt',-tt')$, and $tt'=-t't$, so they anticommute. \emph{An $s$-type and a $t$-type matrix:}

\[
\begin{pmatrix}0&s\\ s&0\end{pmatrix}\begin{pmatrix}0&t\\-t&0\end{pmatrix}=\begin{pmatrix}-st&0\\0&st\end{pmatrix},\qquad\begin{pmatrix}0&t\\-t&0\end{pmatrix}\begin{pmatrix}0&s\\ s&0\end{pmatrix}=\begin{pmatrix}ts&0\\0&-ts\end{pmatrix}.
\]

The sum is $\mathrm{diag}(ts-st,\ st-ts)=0$ because $s$ and $t$ commute (Corollary 3.2 (3)). \emph{With $\tau_7=\mathrm{diag}(-I_4,I_4)$:} $\begin{pmatrix}0&B\\ D&0\end{pmatrix}\tau_7=\begin{pmatrix}0&B\\-D&0\end{pmatrix}$ and $\tau_7\begin{pmatrix}0&B\\ D&0\end{pmatrix}=\begin{pmatrix}0&-B\\ D&0\end{pmatrix}$, whose sum is 0; every $\tau_a$ with $1\le a\le6$ has this off-diagonal form.

(4) $\tau_1\tau_2=\mathrm{diag}(s_1s_2,s_1s_2)$ and then $\tau_1\tau_2\tau_3=\begin{pmatrix}0&s_1s_2s_3\\ s_1s_2s_3&0\end{pmatrix}=\sigma$, since $s_1s_2s_3=I_4$. Next, $\tau_4\tau_5=\mathrm{diag}(-t_3t_2,-t_3t_2)=\mathrm{diag}(-t_1,-t_1)$, because $t_3t_2=t_1$ (proof of Corollary 3.2 (4)), and

\[
\tau_4\tau_5\tau_6=\begin{pmatrix}-t_1&0\\0&-t_1\end{pmatrix}\begin{pmatrix}0&t_1\\-t_1&0\end{pmatrix}=\begin{pmatrix}0&-t_1^2\\ t_1^2&0\end{pmatrix}=\begin{pmatrix}0&I_4\\-I_4&0\end{pmatrix}.
\]

Therefore $\tau_7=(\tau_1\tau_2\tau_3)(\tau_4\tau_5\tau_6)=\begin{pmatrix}0&I_4\\ I_4&0\end{pmatrix}\begin{pmatrix}0&I_4\\-I_4&0\end{pmatrix}=\begin{pmatrix}-I_4&0\\0&I_4\end{pmatrix}$.

(5) $\tau_1\cdots\tau_6\tau_7=\tau_7\tau_7=I_8$ by (3). $\square$

Statement (3) says that $\tau_1,\dots,\tau_7$ are seven anticommuting 8 by 8 matrices, three squaring to $-1$ and four to $+1$: the notebook has built a Clifford algebra with seven generators inside the 8 by 8 matrices, and the step to 16 by 16 adds the eighth direction.

\subsection{The notebook's gammas assembled: the proofs promised in Chapter 2}

The notebook's gammas are $\gamma^a=\begin{pmatrix}0&\bar\tau_a\\ \tau_a&0\end{pmatrix}$ (Section 2.8). By the block rule,

\[
\gamma^a\gamma^b=\begin{pmatrix}0&\bar\tau_a\\ \tau_a&0\end{pmatrix}\begin{pmatrix}0&\bar\tau_b\\ \tau_b&0\end{pmatrix}=\begin{pmatrix}\bar\tau_a\tau_b&0\\0&\tau_a\bar\tau_b\end{pmatrix}.
\]

\textbf{Theorem 3.4 (Stage 1, Result 3.1, by hand).} $\gamma^a\gamma^b+\gamma^b\gamma^a=2\eta^{ab}I_{16}$ for all $a,b=0,\dots,7$.

\emph{Proof.} We need $\bar\tau_a\tau_b+\bar\tau_b\tau_a=2\eta_{ab}I_8$ and $\tau_a\bar\tau_b+\tau_b\bar\tau_a=2\eta_{ab}I_8$.

\begin{itemize}
\item $a=b=0$: $\bar\tau_0=\tau_0=I_8$, and both expressions are $2I_8=2\eta_{00}I_8$.
\item $a=0$, $b\ge1$: with $\bar\tau_b=-\tau_b$, $\bar\tau_0\tau_b+\bar\tau_b\tau_0=\tau_b-\tau_b=0$, and $\tau_0\bar\tau_b+\tau_b\bar\tau_0=-\tau_b+\tau_b=0$; and $\eta_{0b}=0$.
\item $a,b\ge1$: both expressions equal $-(\tau_a\tau_b+\tau_b\tau_a)$, which is $2\eta_{ab}I_8$ by Theorem 3.3 (3). $\square$
\end{itemize}

The symmetry pattern also follows directly: for $a\ge1$, $(\gamma^a)^T=\begin{pmatrix}0&\tau_a^T\\ \bar\tau_a^T&0\end{pmatrix}$, and with $\bar\tau_a=-\tau_a$ and Theorem 3.3 (1) this is $+\gamma^a$ for $a=1,2,3$ and $-\gamma^a$ for $a=4,\dots,7$; $\gamma^0=\begin{pmatrix}0&I_8\\ I_8&0\end{pmatrix}$ is symmetric.

\textbf{Theorem 3.5 (the charge matrix and the chirality, by hand).} $C=\gamma^0\gamma^1\gamma^2\gamma^3=\mathrm{diag}(-\sigma,\sigma)$ and $\gamma^8=\gamma^0\cdots\gamma^7=\mathrm{diag}(-I_8,I_8)$.

\emph{Proof.} From the product formula with $\bar\tau_0=\tau_0=I_8$ and $\bar\tau_a=-\tau_a$ ($a\ge1$):

\[
\begin{aligned}
&\gamma^0\gamma^1=\mathrm{diag}(\tau_1,-\tau_1),&&\gamma^2\gamma^3=\mathrm{diag}(-\tau_2\tau_3,-\tau_2\tau_3),\\
&\gamma^4\gamma^5=\mathrm{diag}(-\tau_4\tau_5,-\tau_4\tau_5),&&\gamma^6\gamma^7=\mathrm{diag}(-\tau_6\tau_7,-\tau_6\tau_7).
\end{aligned}
\]

Multiplying the first two, $C=\mathrm{diag}(-\tau_1\tau_2\tau_3,\ \tau_1\tau_2\tau_3)=\mathrm{diag}(-\sigma,\sigma)$ by Theorem 3.3 (4). Multiplying the last two, $\gamma^4\gamma^5\gamma^6\gamma^7=\mathrm{diag}(\tau_4\tau_5\tau_6\tau_7,\ \tau_4\tau_5\tau_6\tau_7)$. Hence

\[
\gamma^8=C\,\gamma^4\gamma^5\gamma^6\gamma^7=\mathrm{diag}\bigl(-\tau_1\tau_2\cdots\tau_7,\ \tau_1\tau_2\cdots\tau_7\bigr)=\mathrm{diag}(-I_8,I_8)
\]

by Theorem 3.3 (5). $\square$

With Theorem 3.4 and Theorem 3.5 every statement of Chapter 2 about the notebook's gammas now rests on hand computations with six 4 by 4 matrices; the machine checks (\texttt{ALG\_clif\allowbreak{}ford}, \texttt{ALG\_char\allowbreak{}geMatrix}, \texttt{ALG\_chir\allowbreak{}ality}) confirm them. The tables of Section 2.8 can also be read off from the blocks: rows 0 to 7 of $\gamma^a$ are the rows of $\bar\tau_a$ with the column numbers shifted by 8, and rows 8 to 15 are the rows of $\tau_a$. For example, row 0 of $\tau_1$ is $+7$, so row 0 of $\bar\tau_1=-\tau_1$ is $-7$ and row 0 of $\gamma^1$ is $-(7+8)=-15$, as in the table.

\subsection{Algebras, composition algebras, and why octonions}

\textbf{Algebras.} A (real) \textbf{algebra} is the set $\mathbb R^n$ of vectors with $n$ real components (Section 1.2) together with a product $xy$, which assigns to any two vectors $x,y$ of $\mathbb R^n$ a third one and is \textbf{bilinear}, that is, linear in each factor: $(\alpha x+\beta y)z=\alpha\,xz+\beta\,yz$ and $x(\alpha y+\beta z)=\alpha\,xy+\beta\,xz$ for all vectors $x,y,z$ and numbers $\alpha,\beta$. The number $n$ is the \textbf{dimension} of the algebra. It has a \textbf{unit} $e$ if $ex=xe=x$ for all $x$; it is \textbf{associative} if $(xy)z=x(yz)$ and \textbf{commutative} if $xy=yx$. Examples: the real numbers ($n=1$); the complex numbers, as pairs with $(a,b)(c,d)=(ac-bd,\ ad+bc)$ ($n=2$); the quaternions with the columns of Section 3.3 (associative, not commutative, $n=4$); the $d$ by $d$ matrices, each read as the list of its $d^2$ entries (associative, not commutative, $n=d^2$); and $\mathbb R^3$ with the cross product, which is bilinear but neither associative nor unital. For the cross product with the unit vectors $\hat e_1,\hat e_2,\hat e_3$: $(\hat e_1\times\hat e_1)\times\hat e_2=0$, but $\hat e_1\times(\hat e_1\times\hat e_2)=\hat e_1\times\hat e_3=-\hat e_2$.

\textbf{Composition algebras.} An algebra with a unit is a \textbf{composition algebra} if it carries a \textbf{quadratic norm} $N(x)=B(x,x)$, where $B(x,y)=x^Tg\,y$ for a metric $g$ in the sense of Section 1.11 (a symmetric and invertible $n$ by $n$ matrix), such that

\[
N(xy)=N(x)\,N(y)\qquad\text{for all }x,y .
\]

Because $g$ is invertible, no nonzero $x$ has $B(x,y)=0$ for all $y$: taking for $y$ the standard basis vectors $e_0,\dots,e_{n-1}$ of Section 1.2 would give that every component of the row $x^Tg$ is zero, and multiplying $x^Tg=0$ on the right by $g^{-1}$ would give $x^T=0$. (Books on algebra call this property the \textbf{non-degeneracy} of $B$.) The real numbers with $N(x)=x^2$ ($g=(1)$) and the complex numbers with $N(z)=|z|^2$ ($g=I_2$) are examples. So are the quaternions with $N(q)=q_1^2+q_2^2+q_3^2+q_4^2$ ($g=I_4$): the determinant of the 2 by 2 matrix of $q$ (Section 3.3) is

\[
\det\begin{pmatrix}q_4-iq_3&-q_2-iq_1\\ q_2-iq_1&q_4+iq_3\end{pmatrix}=(q_4^2+q_3^2)-(-q_2-iq_1)(q_2-iq_1)=q_4^2+q_3^2+q_2^2+q_1^2,
\]

and $\det(XY)=\det X\det Y$ gives $N(xy)=N(x)N(y)$.

\textbf{Hurwitz's theorem (quoted).} A theorem of Adolf Hurwitz (1898), which we quote without proof, says that composition algebras over the real numbers exist only in dimensions 1, 2, 4 and 8. With a positive norm they are the real numbers, the complex numbers, the quaternions and the \textbf{octonions}; the octonions, of dimension 8, are not associative. In dimensions 2, 4 and 8 there is also a \textbf{split} version, whose norm is indefinite: in dimension 4 it is the algebra of real 2 by 2 matrices with $N=\det$, of signature (2,2) (we met it as $\mathrm{Cl}(1,1)$ in Section 2.5), and in dimension 8 it is the algebra of \textbf{split octonions}, whose norm has signature (4,4).

That last sentence is the reason why octonions appear in a 4+4 dimensional theory. The set $\mathbb R^8$ with the metric $\eta=\mathrm{diag}(+1,+1,+1,+1,-1,-1,-1,-1)$ of Section 1.11 can be given a product that turns it into the split octonions, with $N(x)=\eta(x,x)$. Nothing in the rest of the book depends on Hurwitz's theorem: we construct the split octonions explicitly and prove every property we use.

\subsection{The split octonions in Zorn's vector-matrix form}

\textbf{Zorn's vector matrices.} Max Zorn (1933) wrote a split octonion like a 2 by 2 matrix with two numbers on the diagonal and two three-dimensional vectors off the diagonal,

\[
x=\begin{pmatrix}a&u\\ v&b\end{pmatrix},\qquad a,b\in\mathbb R,\quad u,v\in\mathbb R^3 .
\]

The dictionary with the eight coordinates $x=(x_0,\dots,x_7)$ used by dirac-main and by the repository (Stage 1, §3.3; function \texttt{\detokenize{to_zorn}} of \texttt{scripts/\allowbreak{}d16c\_exa\allowbreak{}ct.py}) is

\[
a=x_0+x_4,\qquad b=x_0-x_4,\qquad u_i=x_i+x_{i+4},\qquad v_i=-x_i+x_{i+4}\qquad(i=1,2,3),
\]

and back: $x_0=\tfrac12(a+b)$, $x_4=\tfrac12(a-b)$, $x_i=\tfrac12(u_i-v_i)$, $x_{i+4}=\tfrac12(u_i+v_i)$. (Inside a vector matrix the diagonal letters, $a,b$ here and $c,d$ in the second factor below, denote numbers; as a subscript, as in $e_a$ below, the letter $a$ is an index from 0 to 7.) The product of two vector matrices is

\[
\begin{pmatrix}a&u\\ v&b\end{pmatrix}\begin{pmatrix}c&r\\ w&d\end{pmatrix}=\begin{pmatrix}ac+u\cdot w&ar+du-v\times w\\ cv+bw+u\times r&v\cdot r+bd\end{pmatrix},
\]

written compactly $(a,u,v,b)(c,r,w,d)=(ac+u\cdot w,\ ar+du-v\times w,\ cv+bw+u\times r,\ v\cdot r+bd)$. The formula resembles a product of 2 by 2 matrices: without the two cross-product terms it would be exactly the ordinary matrix product, with the dot product wherever two vectors are multiplied. The two cross-product terms are what makes the norm multiplicative, $N(xy)=N(x)N(y)$ for the norm $N(x)=ab-u\cdot v$ defined below (Theorem 3.6 (3)), and the product alternative (Exercise 3.8). Without them the product would be neither: in the proof of Theorem 3.6 (3) the term $(v\times w)\cdot(u\times r)$, which comes from the cross products, would be missing, and $N(xy)-N(x)N(y)$ would equal $(u\cdot w)(v\cdot r)-(u\cdot v)(r\cdot w)$, which is 1 for $u=w=\hat e_1$, $v=r=\hat e_2$. The cross products are, however, not the cause of the non-associativity: the product without them is not associative either (for $x=(0,\hat e_1,0,0)$, $y=(0,0,\hat e_1,0)$, $z=(0,\hat e_2,0,0)$ it gives $xy=(1,0,0,0)$, $(xy)z=(0,\hat e_2,0,0)$, but $yz=0$ and $x(yz)=0$), and with them the product is still not associative (see the end of this section). The product is bilinear, because each entry is a sum of products of one entry of each factor.

\textbf{The basis in Zorn form.} With $\hat e_1=(1,0,0)$, $\hat e_2=(0,1,0)$, $\hat e_3=(0,0,1)$ the eight basis vectors are

\[
\begin{aligned}
&e_0=(1,0,0,1),&&e_i=(0,\hat e_i,-\hat e_i,0),\\
&e_4=(1,0,0,-1),&&e_{i+4}=(0,\hat e_i,\hat e_i,0)\qquad(i=1,2,3).
\end{aligned}
\]

\textbf{The unit.} $e_0$ is the unit: $(1,0,0,1)(c,r,w,d)=(c,\ r,\ w,\ d)$, since every term with a zero vector or zero number drops out, and in the same way $(a,u,v,b)(1,0,0,1)=(a,u,v,b)$.

\textbf{Worked example: $e_1e_2=-e_3$.} Here $e_1=(0,\hat e_1,-\hat e_1,0)$ and $e_2=(0,\hat e_2,-\hat e_2,0)$, so $a=b=c=d=0$, $u=\hat e_1$, $v=-\hat e_1$, $r=\hat e_2$, $w=-\hat e_2$. The four entries of the product are

\[
u\cdot w=-\hat e_1\cdot\hat e_2=0,\qquad-v\times w=-\hat e_1\times\hat e_2=-\hat e_3,\qquad u\times r=\hat e_1\times\hat e_2=\hat e_3,\qquad v\cdot r=-\hat e_1\cdot\hat e_2=0 .
\]

So $e_1e_2=(0,-\hat e_3,\hat e_3,0)$. Translating back, $x_0=x_4=0$, $x_3=\tfrac12(u_3-v_3)=\tfrac12(-1-1)=-1$ and $x_7=\tfrac12(u_3+v_3)=0$: the product is $-e_3$.

\textbf{The multiplication table.} Repeating this for all 64 pairs gives the table below; the entry in row $a$ and column $j$ is $e_ae_j$ (computed with \texttt{octonion\allowbreak{}\_product} of \texttt{scripts/\allowbreak{}d16c\_exa\allowbreak{}ct.py}).

\begin{Verbatim}[fontsize=\small]
e_a e_j      (row: e_a, column: e_j)

         e0   e1   e2   e3   e4   e5   e6   e7
  e0     e0   e1   e2   e3   e4   e5   e6   e7
  e1     e1  -e0  -e3   e2  -e5   e4   e7  -e6
  e2     e2   e3  -e0  -e1  -e6  -e7   e4   e5
  e3     e3  -e2   e1  -e0  -e7   e6  -e5   e4
  e4     e4   e5   e6   e7   e0   e1   e2   e3
  e5     e5  -e4   e7  -e6  -e1   e0  -e3   e2
  e6     e6  -e7  -e4   e5  -e2   e3   e0  -e1
  e7     e7   e6  -e5  -e4  -e3  -e2   e1   e0
\end{Verbatim}

The Wolfram verifier records six of these entries as worked examples ($e_1e_1=-e_0$, $e_4e_4=e_0$, $e_1e_2=-e_3$, $e_2e_1=e_3$, $e_1e_4=-e_5$, $e_4e_1=e_5$), counts 64 nonzero structure constants, and finds the whole table equal to the multiplication tensor published by dirac-main in \texttt{split-oc\allowbreak{}tonion.j\allowbreak{}son} (Stage 1, Result 3.8; check \texttt{ALG\_octo\allowbreak{}nionPict\allowbreak{}ureInter\allowbreak{}twiner}). Note that $e_ae_a=-e_0$ for $a=1,2,3$ but $e_ae_a=+e_0$ for $a=4,\dots,7$: the squares follow the signs of $\eta$ with a minus sign, $e_ae_a=-\eta_{aa}e_0$ for $a\ge1$.

\textbf{Conjugation and norm.} The \textbf{conjugate} of $x$ is $\bar x:=(x_0,-x_1,-x_2,\dots,-x_7)$. In Zorn form it exchanges $a$ and $b$ and negates both vectors: $\bar x=(b,-u,-v,a)$ (with $x_4\to-x_4$ the numbers $x_0\pm x_4$ trade places, and $u_i=x_i+x_{i+4}$, $v_i=-x_i+x_{i+4}$ both change sign). The \textbf{norm} is the metric of this book:

\[
N(x):=\eta(x,x)=x_0^2+x_1^2+x_2^2+x_3^2-x_4^2-x_5^2-x_6^2-x_7^2=ab-u\cdot v .
\]

\emph{Proof of the last equality.} $ab=(x_0+x_4)(x_0-x_4)=x_0^2-x_4^2$, and $u_iv_i=(x_{i+4}+x_i)(x_{i+4}-x_i)=x_{i+4}^2-x_i^2$, so $-u\cdot v=\sum_{i=1}^3(x_i^2-x_{i+4}^2)$. $\square$ The norm is the "determinant" $ab-u\cdot v$ of the vector matrix, and $N(\bar x)=ba-(-u)\cdot(-v)=N(x)$.

\textbf{Theorem 3.6 (the basic identities).} For all split octonions $x,y$:

\begin{enumerate}
\item $\bar x(xy)=N(x)\,y$ and $x(\bar xy)=N(x)\,y$;
\item $\bar xx=x\bar x=N(x)\,e_0$;
\item $N(xy)=N(x)\,N(y)$ (the split octonions are a composition algebra);
\item $\overline{xy}=\bar y\,\bar x$ (conjugation reverses products).
\end{enumerate}

\emph{Proof.} Write $x=(a,u,v,b)$, $y=(c,r,w,d)$ and $xy=(A,R,W,D)$ with

\[
A=ac+u\cdot w,\qquad R=ar+du-v\times w,\qquad W=cv+bw+u\times r,\qquad D=v\cdot r+bd .
\]

(1) By the product formula with $\bar x=(b,-u,-v,a)$ on the left, $\bar x(xy)=\bigl(bA-u\cdot W,\ bR-Du+v\times W,\ -Av+aW-u\times R,\ -v\cdot R+aD\bigr)$. We expand each entry, using (V1) and (V3) of Section 3.2.

\[
\begin{aligned}
bA-u\cdot W&=abc+b\,u\cdot w-c\,u\cdot v-b\,u\cdot w-u\cdot(u\times r)=(ab-u\cdot v)\,c,\\
bR-Du+v\times W&=abr+bd\,u-b\,v\times w-(v\cdot r)u-bd\,u+c\,v\times v+b\,v\times w+v\times(u\times r)\\
&=ab\,r-(v\cdot r)\,u+\bigl(u\,(v\cdot r)-r\,(v\cdot u)\bigr)=(ab-u\cdot v)\,r,\\
-Av+aW-u\times R&=-ac\,v-(u\cdot w)v+ac\,v+ab\,w+a\,u\times r-a\,u\times r-d\,u\times u+u\times(v\times w)\\
&=-(u\cdot w)\,v+ab\,w+\bigl(v\,(u\cdot w)-w\,(u\cdot v)\bigr)=(ab-u\cdot v)\,w,\\
-v\cdot R+aD&=-a\,v\cdot r-d\,v\cdot u+v\cdot(v\times w)+a\,v\cdot r+abd=(ab-u\cdot v)\,d .
\end{aligned}
\]

So $\bar x(xy)=N(x)(c,r,w,d)=N(x)y$. Applying this to $\bar x$ in place of $x$, with $\bar{\bar x}=x$ and $N(\bar x)=N(x)$, gives $x(\bar xy)=N(x)y$.

(2) Put $y=e_0$ in (1): $\bar x(xe_0)=\bar xx=N(x)e_0$ and $x(\bar xe_0)=x\bar x=N(x)e_0$.

(3) $N(xy)=AD-R\cdot W$. First,

\[
AD=ac\,(v\cdot r)+abcd+(u\cdot w)(v\cdot r)+bd\,(u\cdot w).
\]

Next, multiplying out $R\cdot W$ term by term and dropping the terms that vanish by (V1) ($r\cdot(u\times r)=0$, $u\cdot(u\times r)=0$, $(v\times w)\cdot v=0$, $(v\times w)\cdot w=0$),

\[
R\cdot W=ac\,(r\cdot v)+ab\,(r\cdot w)+cd\,(u\cdot v)+bd\,(u\cdot w)-(v\times w)\cdot(u\times r).
\]

By (V4), $(v\times w)\cdot(u\times r)=(v\cdot u)(w\cdot r)-(v\cdot r)(w\cdot u)$. Subtracting,

\[
N(xy)=abcd-ab\,(r\cdot w)-cd\,(u\cdot v)+(u\cdot v)(r\cdot w)=(ab-u\cdot v)(cd-r\cdot w)=N(x)N(y).
\]

(4) By the product formula, $\bar y\,\bar x=(d,-r,-w,c)(b,-u,-v,a)=\bigl(db+r\cdot v,\ -du-ar-w\times v,\ -bw-cv+r\times u,\ w\cdot u+ca\bigr)$. The conjugate of $xy=(A,R,W,D)$ is $(D,-R,-W,A)$, and indeed $D=v\cdot r+bd$, $-R=-ar-du+v\times w$, $-W=-cv-bw-u\times r$ and $A=ac+u\cdot w$ agree entry by entry with $\bar y\bar x$ (use $-w\times v=v\times w$ and $r\times u=-u\times r$). $\square$

The Wolfram verifier checks the unit law, the composition law (3), the reversal (4) and the identification of $N$ with $\eta$ (check \texttt{ALG\_octo\allowbreak{}nionPict\allowbreak{}ureInter\allowbreak{}twiner}), and the Python checker tests the unit law, $x\bar x=N(x)e_0$, (3), (4) and the alternative laws of Exercise 3.8 on three rational sample octonions (all within the same check).

\textbf{The split octonions are not associative.} From the table: $e_1e_2=-e_3$ and $e_3e_4=-e_7$, so $(e_1e_2)e_4=-e_3e_4=e_7$; and $e_2e_4=-e_6$ and $e_1e_6=e_7$, so $e_1(e_2e_4)=-e_1e_6=-e_7$. The \textbf{associator} is

\[
(e_1e_2)e_4-e_1(e_2e_4)=2e_7\ne0 ,
\]

as recorded by the Wolfram verifier (Stage 1, Result 3.8). The algebra is nevertheless \textbf{alternative}: $x(xy)=(xx)y$ and $(yx)x=y(xx)$ for all $x,y$ (Exercise 3.8).

\subsection{Octonion gamma matrices}

\textbf{Left multiplication matrices.} For a split octonion $x$, the map $y\mapsto xy$ is linear; its 8 by 8 matrix is $L_x$, whose column $j$ is the coordinate column of $xe_j$. (From here on $L_x$ denotes the 8 by 8 octonion multiplication; the 4 by 4 quaternion matrices $L_x$ and $R_x$ of Section 3.3 are not used again, apart from one comparison in Section 3.13.) Because the product is bilinear, $L_x$ depends linearly on $x$. For the basis vectors, $L_{e_a}$ is read off row $a$ of the multiplication table: $(L_{e_a})_{ij}$ is the coefficient of $e_i$ in $e_ae_j$. As signed permutations ($L_{e_0}=I_8$; the same matrices are the lower-left blocks of the octonion gammas stored under the key \texttt{gammaOct\allowbreak{}onion} of \texttt{algebra-\allowbreak{}fixture.\allowbreak{}json}):

\begingroup
\small
\begin{longtable}{@{}>{\raggedright\arraybackslash}p{0.100\linewidth}>{\raggedright\arraybackslash}p{0.100\linewidth}>{\raggedright\arraybackslash}p{0.100\linewidth}>{\raggedright\arraybackslash}p{0.100\linewidth}>{\raggedright\arraybackslash}p{0.100\linewidth}>{\raggedright\arraybackslash}p{0.100\linewidth}>{\raggedright\arraybackslash}p{0.100\linewidth}>{\raggedright\arraybackslash}p{0.100\linewidth}@{}}
\toprule
\textbf{$i$} & \textbf{$L_{e_1}$} & \textbf{$L_{e_2}$} & \textbf{$L_{e_3}$} & \textbf{$L_{e_4}$} & \textbf{$L_{e_5}$} & \textbf{$L_{e_6}$} & \textbf{$L_{e_7}$} \\
\midrule
\endfirsthead
\toprule
\textbf{$i$} & \textbf{$L_{e_1}$} & \textbf{$L_{e_2}$} & \textbf{$L_{e_3}$} & \textbf{$L_{e_4}$} & \textbf{$L_{e_5}$} & \textbf{$L_{e_6}$} & \textbf{$L_{e_7}$} \\
\midrule
\endhead
0 & $-1$ & $-2$ & $-3$ & $+4$ & $+5$ & $+6$ & $+7$ \\
1 & $+0$ & $-3$ & $+2$ & $+5$ & $-4$ & $-7$ & $+6$ \\
2 & $+3$ & $+0$ & $-1$ & $+6$ & $+7$ & $-4$ & $-5$ \\
3 & $-2$ & $+1$ & $+0$ & $+7$ & $-6$ & $+5$ & $-4$ \\
4 & $+5$ & $+6$ & $+7$ & $+0$ & $-1$ & $-2$ & $-3$ \\
5 & $-4$ & $+7$ & $-6$ & $+1$ & $+0$ & $+3$ & $-2$ \\
6 & $-7$ & $-4$ & $+5$ & $+2$ & $-3$ & $+0$ & $+1$ \\
7 & $+6$ & $-5$ & $-4$ & $+3$ & $+2$ & $-1$ & $+0$ \\
\bottomrule
\end{longtable}
\endgroup

For example, row 0 of $L_{e_1}$ is $-1$: the coefficient of $e_0$ in $e_1e_1$ is $-1$. Notice that $L_{e_4}=\sigma$, the matrix $\begin{pmatrix}0&I_4\\ I_4&0\end{pmatrix}$ of the notebook, and that $L_{e_1},L_{e_2},L_{e_3}$ keep the two halves $\{e_0,\dots,e_3\}$ and $\{e_4,\dots,e_7\}$ separate while $L_{e_4},\dots,L_{e_7}$ exchange them. In matrix language, Theorem 3.6 (1) says

\[
L_{\bar x}L_x=L_xL_{\bar x}=N(x)\,I_8 .
\]

\textbf{The octonion gammas} (Stage 1, §3.3; function \texttt{octonion\allowbreak{}\_gammas}) are the 16 by 16 matrices

\[
\Gamma(x):=\begin{pmatrix}0&L_{\bar x}\\ L_x&0\end{pmatrix},\qquad\Gamma^a:=\Gamma(e_a)\quad(a=0,\dots,7).
\]

\textbf{Theorem 3.7.} $\Gamma(x)\Gamma(y)+\Gamma(y)\Gamma(x)=2\eta(x,y)\,I_{16}$ for all $x,y$. In particular $\{\Gamma^a,\Gamma^b\}=2\eta^{ab}I_{16}$.

\emph{Proof.} By the block rule and the matrix form of Theorem 3.6 (1), $\Gamma(x)^2=\mathrm{diag}(L_{\bar x}L_x,\ L_xL_{\bar x})=N(x)I_{16}$. Conjugation is linear, so $\Gamma$ depends linearly on $x$, and applying this to $x+y$,

\[
\Gamma(x)^2+\Gamma(y)^2+\Gamma(x)\Gamma(y)+\Gamma(y)\Gamma(x)=N(x+y)\,I_{16}=\bigl(N(x)+N(y)+2\eta(x,y)\bigr)I_{16}.
\]

Subtracting $\Gamma(x)^2=N(x)I_{16}$ and $\Gamma(y)^2=N(y)I_{16}$ leaves the claim. For $x=e_a$, $y=e_b$, $\eta(e_a,e_b)=\eta_{ab}=\eta^{ab}$. $\square$

So the split octonions give a \emph{third} set of real 16 by 16 gammas of signature (4,4), next to the tensor picture of Section 2.7 and the notebook's. Their chirality $\Gamma^0\cdots\Gamma^7$ is again $\mathrm{diag}(-I_8,I_8)$ (Section 3.12). In this picture a spinor is a pair of split octonions $(y,z)$, the upper and the lower half, and a vector $x$ acts by $\Gamma(x)(y,z)=(\bar xz,\ xy)$: the vectors and both halves of the spinor space are three copies of the same 8-dimensional algebra, a concrete face of the triality mentioned in Section 2.15.

\textbf{Lemma 3.8 (transposes of the multiplication matrices).} For every $x$, $L_x^T\eta L_x=N(x)\,\eta$; and if $N(x)\ne0$, then $L_x^T\eta=\eta L_{\bar x}$, that is, $L_x^T=\eta L_{\bar x}\eta$.

\emph{Proof.} Replace $y$ by $y+z$ in $N(xy)=N(x)N(y)$. On the left, $N(xy+xz)=N(xy)+N(xz)+2\eta(xy,xz)$; on the right, $N(x)\bigl(N(y)+N(z)+2\eta(y,z)\bigr)$. The first two terms agree on both sides by Theorem 3.6 (3), so $\eta(xy,xz)=N(x)\eta(y,z)$, which in matrix form is $(L_xy)^T\eta(L_xz)=N(x)\,y^T\eta z$ for all columns $y,z$, that is, $L_x^T\eta L_x=N(x)\eta$. Multiply on the right by $L_{\bar x}$ and use $L_xL_{\bar x}=N(x)I_8$: $N(x)L_x^T\eta=N(x)\eta L_{\bar x}$; divide by $N(x)\ne0$. Finally $\eta^{-1}=\eta$. $\square$

For a basis vector $e_a$ with $a\ge1$ we have $\bar e_a=-e_a$ and $N(e_a)=\eta_{aa}\ne0$, so $L_{e_a}^T=-\eta L_{e_a}\eta$. For $a=1,2,3$, $L_{e_a}$ keeps the two halves separate and therefore commutes with $\eta=\mathrm{diag}(I_4,-I_4)$: $L_{e_a}$ is \textbf{antisymmetric}. For $a=4,\dots,7$ it exchanges the halves, so $\eta L_{e_a}\eta=-L_{e_a}$: $L_{e_a}$ is \textbf{symmetric}. It follows (Exercise 3.10) that $\Gamma^a$ is symmetric for $a\le3$ and antisymmetric for $a\ge4$, the same pattern as the notebook's gammas and the tensor ones. The repository checks that the $\Gamma^a$ satisfy the Clifford relation and that they equal the octonion Clifford generators published by dirac-main in \texttt{triality\allowbreak{}44.json} (check \texttt{ALG\_octo\allowbreak{}nionPict\allowbreak{}ureInter\allowbreak{}twiner}).

\subsection{Intertwiners: why every 16 by 16 picture is the same}

We now have three sets of real 16 by 16 gamma matrices of signature (4,4): the notebook's $\gamma^a$, the tensor matrices $\hat\gamma^a$ of Section 2.7 and the octonion matrices $\Gamma^a$ of Section 3.9. They look completely different. This section proves that they are the same object written in different bases.

\textbf{Definition.} Let $\gamma^a$ and $\gamma'^a$ ($a=0,\dots,7$) be two sets of gamma matrices. An \textbf{intertwiner} from the first set to the second is a matrix $K$ with

\[
\gamma'^a\,K=K\,\gamma^a\qquad\text{for all }a .
\]

If $K$ is invertible, then $\gamma'^a=K\gamma^aK^{-1}$: the second set is the first one expressed in another basis. A spinor $\Psi$ of the first picture corresponds to the spinor $K\Psi$ of the second, because $\gamma'^a(K\Psi)=K(\gamma^a\Psi)$: applying a gamma and then translating is the same as translating and then applying the gamma. (This is the intertwiner of Section 2.13 for the representations of $\mathrm{Pin}(4,4)$ given by the two sets.)

\textbf{Theorem 3.9 (existence and uniqueness of the intertwiner).} Let $\gamma^a$ and $\gamma'^a$ be two sets of real 16 by 16 matrices, both satisfying the Clifford relation of signature (4,4). Then:

\begin{enumerate}
\item there is a real intertwiner $K\ne0$;
\item every nonzero intertwiner, real or complex, is invertible;
\item any two intertwiners are proportional: the intertwiners form a one-dimensional space.
\end{enumerate}

\emph{Proof.} Write $\gamma_A$ and $\gamma'_A$ for the monomials of the two sets. The sign $s(b,A)=\pm1$ in the rule $\gamma^b\gamma_A=s(b,A)\,\gamma_{A\triangle\{b\}}$ comes only from reordering with the Clifford relation (rule (R1) of Section 2.5), so the same sign appears for the second set: $\gamma'^b\gamma'_A=s(b,A)\,\gamma'_{A\triangle\{b\}}$.

(1) For any real 16 by 16 matrix $X$ define the \textbf{average}

\[
K_X:=\sum_{A}\gamma'_A\,X\,\gamma_A^{-1},
\]

the sum running over all 256 subsets $A$. Fix $b$, write $A'=A\triangle\{b\}$ and $s=s(b,A)$. Inverting $\gamma^b\gamma_A=s\gamma_{A'}$ gives $\gamma_A^{-1}(\gamma^b)^{-1}=s\gamma_{A'}^{-1}$ (as $s^{-1}=s$), that is, $\gamma_A^{-1}=s\,\gamma_{A'}^{-1}\gamma^b$. Therefore

\[
\gamma'^bK_X=\sum_A\bigl(\gamma'^b\gamma'_A\bigr)X\gamma_A^{-1}=\sum_A\bigl(s\gamma'_{A'}\bigr)X\bigl(s\gamma_{A'}^{-1}\gamma^b\bigr)=\Bigl(\sum_{A'}\gamma'_{A'}X\gamma_{A'}^{-1}\Bigr)\gamma^b=K_X\gamma^b,
\]

because $s^2=1$ and $A\mapsto A'$ runs through all subsets exactly once. So every $K_X$ is an intertwiner. Not all of them vanish. Suppose they did, and take $X=E_{jl}$. The entry of $K_X$ in position $(i,m)$ is $\sum_A(\gamma'_A)_{ij}(\gamma_A^{-1})_{lm}$. For fixed $l,m$ put $c_A:=(\gamma_A^{-1})_{lm}$; then the matrix $\sum_Ac_A\gamma'_A$ has every entry $(i,j)$ equal to zero, so all $c_A=0$ by the independence of the monomials (Theorem 2.2). As this holds for all $l,m$, every $\gamma_A^{-1}$ would be the zero matrix, which is absurd. So some $K_X$ is a nonzero real intertwiner.

(2) Let $K\ne0$ be an intertwiner. If $Kv=0$, then $K\gamma^av=\gamma'^aKv=0$: the kernel of $K$ is invariant under every $\gamma^a$, hence under every complex combination of monomials, which is every matrix (Section 2.6). By the argument in the proof of Theorem 2.7, a nonzero subspace invariant under all matrices is all of $\mathbb C^{16}$. The kernel is not all of $\mathbb C^{16}$ because $K\ne0$, so it is $\{0\}$: $K$ is one-to-one, and a one-to-one square matrix is invertible.

(3) Let $K_1$ be an invertible intertwiner and $K_2$ any intertwiner. From $\gamma'^aK_1=K_1\gamma^a$ we get $K_1^{-1}\gamma'^a=\gamma^aK_1^{-1}$, so $Z:=K_1^{-1}K_2$ satisfies $\gamma^aZ=K_1^{-1}\gamma'^aK_2=K_1^{-1}K_2\gamma^a=Z\gamma^a$. By Corollary 2.3, $Z=c\,I$, and $K_2=cK_1$. $\square$

\textbf{Not every $X$ works.} Pair each set $A$ with its complement $A^c$. By (R1), $\gamma_{A^c}=t_A\,\gamma_A\gamma^8$ with a sign $t_A=\pm1$ that again comes only from the relations, so that also $\gamma'_{A^c}=t_A\,\gamma'_A\gamma'^8$ (where $\gamma'^8=\gamma'^0\cdots\gamma'^7$). Then $\gamma_{A^c}^{-1}=t_A\gamma^8\gamma_A^{-1}$, and the term of $A^c$ in the average is $\gamma'_A\,(\gamma'^8X\gamma^8)\,\gamma_A^{-1}$. Summing over all $A$ gives $K_X=K_{\gamma'^8X\gamma^8}$, hence $K_X=\tfrac12K_{X+\gamma'^8X\gamma^8}$. For the notebook gammas ($\gamma^8=\mathrm{diag}(-I_8,I_8)$) and the tensor gammas ($\gamma'^8=G\otimes G\otimes G\otimes G$, whose entry in position $(0,0)$ is $+1$) and $X=E_{0,0}$, both chiralities are diagonal and $\gamma'^8E_{0,0}\gamma^8=(+1)(-1)E_{0,0}=-E_{0,0}$, so $K_{E_{0,0}}=0$. This is why the proof only claims that \emph{some} $X$ gives a nonzero average: the notebook basis spinor $e_0$ has chirality $-1$, the tensor basis spinor $e_0$ has chirality $+1$, and an intertwiner never maps one chirality to the other.

\textbf{Consequences.} An invertible intertwiner maps products to products, $K\gamma^{a_1}\cdots\gamma^{a_k}K^{-1}=(K\gamma^{a_1}K^{-1})\cdots(K\gamma^{a_k}K^{-1})=\gamma'^{a_1}\cdots\gamma'^{a_k}$ (insert $K^{-1}K=1$ between the factors). So the charge matrix, the chirality and the spin generators of one picture are mapped to those of the other, and everything that does not depend on the basis (traces, eigenvalue counts, signatures, dimensions of commutants) is the same in every picture. In the language of Section 2.13: any two 16-dimensional representations of $\mathrm{Pin}(4,4)$ built from Clifford matrices are equivalent, the equivalence is unique up to a factor, and $\mathrm{Cl}(4,4)$ has, up to a change of basis, exactly one 16-dimensional module. Stage 1 states it as "one Clifford module written in three bases" (§3.3).

\textbf{How the computer finds $K$.} As for the commutant (Section 2.15), the equations $\gamma'^aK-K\gamma^a=0$ are $8\cdot256$ linear equations with integer coefficients for the 256 entries of $K$, solved exactly (function \texttt{primitiv\allowbreak{}e\_intert\allowbreak{}winer} of \texttt{scripts/\allowbreak{}d16c\_exa\allowbreak{}ct.py}). The solution space has dimension 1, as Theorem 3.9 predicts. Its basis vector is scaled to integers without a common factor, with the sign chosen so that the first nonzero entry in \textbf{row-major order} (row 0 read from left to right, then row 1, and so on) is positive. This \textbf{primitive integer normalization} fixes $K$ completely.

\subsection{The tensor picture and \texorpdfstring{$K_{\mathrm{clifford}}$}{K\_clifford}}

\textbf{Direction (erratum E4 of the contract).} $K_{\mathrm{clifford}}$ is the primitive integer solution of

\[
\hat\gamma^a\,K=K\,\gamma^a\qquad(a=0,\dots,7):
\]

it maps the notebook picture to the tensor picture of dirac-main.

\textbf{Stage 1, Result 3.7.} The solution space has dimension exactly 1; $K_{\mathrm{clifford}}$ has rank 16 (it is invertible); $K_{\mathrm{clifford}}\,C\,K_{\mathrm{clifford}}^{-1}=C_{\mathrm{dm}}$; $K_{\mathrm{clifford}}\,\gamma^8K_{\mathrm{clifford}}^{-1}=+G\otimes G\otimes G\otimes G$; the chirality $-1$ block $\Psi_0,\dots,\Psi_7$ of the notebook is mapped onto the dirac-main basis indices $\{1,2,4,7,8,11,13,14\}$ and the $+1$ block onto $\{0,3,5,6,9,10,12,15\}$; and $K_{\mathrm{clifford}}$ is not a signed permutation: its entries are $0,\pm1$ with four nonzero entries in every row (check \texttt{ALG\_clif\allowbreak{}fordPict\allowbreak{}ureInter\allowbreak{}twiner}, whose true/false verdict covers the dimension, the rank, the intertwining and $KCK^{-1}=C_{\mathrm{dm}}$, and in the Wolfram verifier also $K\gamma^8K^{-1}=\pm G\otimes G\otimes G\otimes G$; the sign $+1$, the images of the two blocks and the signed-permutation test are measurements that the check records, see Section 3.14; the Python checker also finds the same matrix in the Wolfram report, check \texttt{ALG\_wolf\allowbreak{}ramAgree\allowbreak{}ment}).

\begin{Verbatim}[fontsize=\small]
K_clifford =
[ 0  0  0  0  0  0  0  0  1 -1  0  0  1  1  0  0 ]
[-1  1  0  0  1  1  0  0  0  0  0  0  0  0  0  0 ]
[ 0  0 -1  1  0  0 -1 -1  0  0  0  0  0  0  0  0 ]
[ 0  0  0  0  0  0  0  0  0  0  1 -1  0  0 -1 -1 ]
[ 0  0  1  1  0  0 -1  1  0  0  0  0  0  0  0  0 ]
[ 0  0  0  0  0  0  0  0  0  0 -1 -1  0  0 -1  1 ]
[ 0  0  0  0  0  0  0  0 -1 -1  0  0  1 -1  0  0 ]
[ 1  1  0  0  1 -1  0  0  0  0  0  0  0  0  0  0 ]
[ 1 -1  0  0  1  1  0  0  0  0  0  0  0  0  0  0 ]
[ 0  0  0  0  0  0  0  0 -1  1  0  0  1  1  0  0 ]
[ 0  0  0  0  0  0  0  0  0  0 -1  1  0  0 -1 -1 ]
[ 0  0  1 -1  0  0 -1 -1  0  0  0  0  0  0  0  0 ]
[ 0  0  0  0  0  0  0  0  0  0  1  1  0  0 -1  1 ]
[ 0  0 -1 -1  0  0 -1  1  0  0  0  0  0  0  0  0 ]
[-1 -1  0  0  1 -1  0  0  0  0  0  0  0  0  0  0 ]
[ 0  0  0  0  0  0  0  0  1  1  0  0  1 -1  0  0 ]
\end{Verbatim}

(Copied from the measurement \texttt{K\_cliffo\allowbreak{}rd} of \texttt{python-a\allowbreak{}lgebra-r\allowbreak{}eport.js\allowbreak{}on}; the Wolfram report and \texttt{algebra-\allowbreak{}fixture.\allowbreak{}json} hold the same matrix.)

\textbf{Why the listed facts hold.} Rank 16 and uniqueness are Theorem 3.9. The charge matrix: $KCK^{-1}=K\gamma^0\gamma^1\gamma^2\gamma^3K^{-1}=\hat\gamma^0\hat\gamma^1\hat\gamma^2\hat\gamma^3=C_{\mathrm{dm}}$ by the consequence in Section 3.10. The chirality: in the same way $K\gamma^8K^{-1}=\hat\gamma^8$, which is $G\otimes G\otimes G\otimes G$ by Section 2.7. The blocks: if $\gamma^8\Psi=-\Psi$, then $\hat\gamma^8(K\Psi)=K\gamma^8\Psi=-K\Psi$, so $K\Psi$ lies in the $-1$ eigenspace of $G\otimes G\otimes G\otimes G$, which is spanned by the basis columns with an odd number of 1-bits, the list $\{1,2,4,7,8,11,13,14\}$ of Section 2.7. These eight standard basis columns are independent, so this eigenspace has dimension 8. Since $K$ is invertible, the eight columns $Ke_0,\dots,Ke_7$ (the images of the notebook basis spinors of the $-1$ block) are independent: if $\sum_kc_kKe_k=0$, then $K\bigl(\sum_kc_ke_k\bigr)=0$, so $\sum_kc_ke_k=0$ and all $c_k=0$. They lie in the 8-dimensional $-1$ eigenspace, so by fact (i) of Section 2.2 they form a basis of it and span the whole eigenspace: the block is mapped \emph{onto} it, not only into it. In the same way the $+1$ block $e_8,\dots,e_{15}$ is mapped onto the $+1$ eigenspace, spanned by the columns $\{0,3,5,6,9,10,12,15\}$ with an even number of 1-bits.

\textbf{Worked example.} Column 0 of $K_{\mathrm{clifford}}$ is the image of the notebook basis spinor $e_0$ (chirality $-1$). Reading down column 0 of the matrix, its nonzero entries are $-1$ in row 1, $+1$ in row 7, $+1$ in row 8 and $-1$ in row 14. All four row numbers belong to the list $\{1,2,4,7,8,11,13,14\}$, as they must.

Two further remarks. First, $K_{\mathrm{clifford}}^TK_{\mathrm{clifford}}=4I_{16}$, so $\tfrac12K_{\mathrm{clifford}}$ is orthogonal (Exercise 3.11 proves this with Schur's lemma). Second, the check that the tensor matrices equal the committed generators of dirac-main (file \texttt{cl44-see\allowbreak{}d.json}), which only the Wolfram verifier makes, needs the separately published dirac-main folder; in a public clone without it this sub-check records "not-run" (Stage 1, §12).

\subsection{The octonion picture, \texorpdfstring{$K_{\mathrm{octonion}}$}{K\_octonion} and the light-cone basis}

\textbf{Direction (erratum E4).} $K_{\mathrm{octonion}}$ is the primitive integer solution of

\[
\gamma^a\,K=K\,\Gamma^a\qquad(a=0,\dots,7):
\]

it maps the octonion picture to the notebook picture. (Erratum E4 records the directions of both intertwiners explicitly, because the two are easily confused: $K_{\mathrm{clifford}}$ starts from the notebook picture, $K_{\mathrm{octonion}}$ ends in it.)

\textbf{Stage 1, Result 3.8.} The solution space has dimension exactly 1; $K_{\mathrm{octonion}}$ has rank 16 and is block diagonal with two equal blocks, $K_{\mathrm{octonion}}=\mathrm{diag}(Q,Q)$, where

\begin{Verbatim}[fontsize=\small]
Q =
[ 1  0  0  0  0  0  0 -1 ]
[ 0  0  0 -1 -1  0  0  0 ]
[ 0  0  1  0  0  1  0  0 ]
[ 0 -1  0  0  0  0  1  0 ]
[ 1  0  0  0  0  0  0  1 ]
[ 0  0  0 -1  1  0  0  0 ]
[ 0  0  1  0  0 -1  0  0 ]
[ 0 -1  0  0  0  0 -1  0 ]
\end{Verbatim}

(This $Q$ is unrelated to the 4 by 4 blocks $Q_a[h]$, $Q_b[h]$ of Section 3.4.) Moreover $Q^TQ=2I_8$ and $Q^T\sigma Q=2\eta$. (Check \texttt{ALG\_octo\allowbreak{}nionPict\allowbreak{}ureInter\allowbreak{}twiner} certifies the dimension, the rank and the intertwining; the matrix $Q$ and the two products are measurements that the check records: the entries \texttt{\detokenize{Q}}, \texttt{Qtranspo\allowbreak{}seQ} (recorded as the text \texttt{\detokenize{2 I8}}) and \texttt{Qtranspo\allowbreak{}seSigmaQ\allowbreak{}OverEta} (recorded as the number 2) of the measurement \texttt{ALG\_octo\allowbreak{}nionPict\allowbreak{}ureInter\allowbreak{}twiner.n\allowbreak{}otebookB\allowbreak{}asisChan\allowbreak{}ge} of \texttt{wolfram-\allowbreak{}algebra-\allowbreak{}report.j\allowbreak{}son}. In the Python report, \texttt{K\_octoni\allowbreak{}on} is $\mathrm{diag}(Q,Q)$, with the measurements \texttt{ALG\_octo\allowbreak{}nionKBlo\allowbreak{}ckDiagon\allowbreak{}al} and \texttt{ALG\_octo\allowbreak{}nionKEqu\allowbreak{}alBlocks} true and \texttt{ALG\_octo\allowbreak{}nionKGra\allowbreak{}mFactor} $=2$, which is $Q^TQ=2I_8$; see Section 3.14.)

\textbf{What the block form means.} With $K=\mathrm{diag}(Q,Q)$ the block rule gives

\[
\gamma^aK=\begin{pmatrix}0&\bar\tau_aQ\\ \tau_aQ&0\end{pmatrix},\qquad K\Gamma^a=\begin{pmatrix}0&QL_{\bar e_a}\\ QL_{e_a}&0\end{pmatrix}.
\]

So the intertwining property says exactly

\[
\tau_a=Q\,L_{e_a}\,Q^{-1},\qquad\bar\tau_a=Q\,L_{\bar e_a}\,Q^{-1}\qquad(a=0,\dots,7).
\]

\textbf{The notebook's $\tau_a$ are octonion multiplications in another basis.} Each $\tau_a$ is the left multiplication by the octonion unit $e_a$, written in the basis given by the columns of $Q^{-1}=\tfrac12Q^T$ (the rows of $Q$ divided by 2), that is, in the coordinates $y=Qx$ of the light-cone paragraph below. (In general, if the columns of an invertible matrix $S$ are taken as a new basis, a vector with new coordinates $y$ is $x=Sy$, and a matrix $M$ acting on $x$ acts on $y$ as $S^{-1}MS$, because $Mx=MSy=S\,(S^{-1}MSy)$; here $S=Q^{-1}$ and $S^{-1}L_{e_a}S=QL_{e_a}Q^{-1}=\tau_a$. The inverse $Q^{-1}=\tfrac12Q^T$ follows from $Q^TQ=2I_8$.) For example, the first new basis vector is $f_0=\tfrac12(e_0-e_7)$ (row 0 of $Q$ divided by 2), and by the multiplication table $e_4f_0=\tfrac12(e_4-e_3)$, which is $f_5$ (row 5 of $Q$ divided by 2), in agreement with the entry $+1$ in row 5, column 0 of $\tau_4$ (Section 3.5: row 5 of $\tau_4$ is $+0$). For $a\ge1$ the second formula follows from the first, since $\bar\tau_a=-\tau_a$ (Theorem 3.3) and $L_{\bar e_a}=L_{-e_a}=-L_{e_a}$; for $a=0$ both sides of both formulas are $I_8$.

\textbf{Worked check ($a=4$, column 0).} We check column 0 of the matrix equation $\tau_4Q=QL_{e_4}$, which is $\tau_4=QL_{e_4}Q^{-1}$ multiplied on the right by $Q$. Here $L_{e_4}=\sigma$ (Section 3.9), and $Q\sigma$ is $Q$ with its two halves of columns exchanged, so column 0 of $Q\sigma$ is column 4 of $Q$, namely $(0,-1,0,0,0,1,0,0)^T$. On the other side, column 0 of $Q$ is $u=(1,0,0,0,1,0,0,0)^T$, and with the table of Section 3.5, $\tau_4u=(u_5,-u_4,-u_7,u_6,-u_1,u_0,u_3,-u_2)=(0,-1,0,0,0,1,0,0)$. The two agree.

\textbf{Erratum E5.} The notebook's $\tau_a$ are \emph{not} the octonion left multiplications themselves, nor these multiplications with the basis reordered and signs flipped. The equality $\tau_a=L_{e_a}$ holds only for $a=0$ (measurement \texttt{ALG\_tauE\allowbreak{}qualsLef\allowbreak{}tMultipl\allowbreak{}icationP\allowbreak{}erIndex} $=[$true, false, $\dots$, false$]$ in \texttt{python-a\allowbreak{}lgebra-r\allowbreak{}eport.js\allowbreak{}on}). The equation $\tau_a=P^{-1}L_{e_a}P$ for all $a$ has, up to a factor, exactly one solution $P$, namely $P\propto Q^{-1}=\tfrac12Q^T$ (computed dimension 1, measurement \texttt{ALG\_tauL\allowbreak{}eftMulti\allowbreak{}plicatio\allowbreak{}nIntertw\allowbreak{}inerDime\allowbreak{}nsion}). Since $Q^T$ has two nonzero entries in every row, no multiple of it is a signed permutation (measurement \texttt{ALG\_tauE\allowbreak{}qualsLef\allowbreak{}tMultipl\allowbreak{}icationU\allowbreak{}pToSigne\allowbreak{}dPermuta\allowbreak{}tion}: false).

\textbf{The light-cone basis.} Let $x=(x_0,\dots,x_7)$ be the coordinates of a split octonion and $y=Qx$ its coordinates in the notebook's basis. Reading off the rows of $Q$,

\[
\begin{aligned}
&y_0=x_0-x_7,\quad y_4=x_0+x_7,\qquad y_1=-x_3-x_4,\quad y_5=-x_3+x_4,\\
&y_2=x_2+x_5,\quad y_6=x_2-x_5,\qquad y_3=-x_1+x_6,\quad y_7=-x_1-x_6 .
\end{aligned}
\]

Each pair of new coordinates is a sum and a difference of one space-like and one time-like old coordinate, and

\[
y_0y_4+y_1y_5+y_2y_6+y_3y_7=(x_0^2-x_7^2)+(x_3^2-x_4^2)+(x_2^2-x_5^2)+(x_1^2-x_6^2)=N(x).
\]

Coordinates of this kind, like $x_0\pm x_7$, are called \textbf{light-cone} (or null) coordinates: a vector with only $y_0\ne0$ has $N=0$. In matrix form this is $Q^T\sigma Q=2\eta$, because $y^T\sigma y=2(y_0y_4+y_1y_5+y_2y_6+y_3y_7)$. So the notebook's 8-component halves are split octonions written in a null basis, in which the octonion norm $\eta$ appears as $\tfrac12\sigma$; and $Q^TQ=2I_8$ says that $Q/\sqrt2$ is orthogonal: it combines four pairs of coordinates into their sums and differences divided by $\sqrt2$ (the contract calls it a 45-degree light-cone rotation).

\textbf{The chirality of the octonion picture.} From $\Gamma^a=K^{-1}\gamma^aK$ with $K=\mathrm{diag}(Q,Q)$,

\[
\Gamma^0\Gamma^1\cdots\Gamma^7=K^{-1}\gamma^8K=\mathrm{diag}(Q^{-1},Q^{-1})\,\mathrm{diag}(-I_8,I_8)\,\mathrm{diag}(Q,Q)=\mathrm{diag}(-I_8,I_8),
\]

the same as in the notebook picture (measurement \texttt{ALG\_octo\allowbreak{}nionPict\allowbreak{}ureInter\allowbreak{}twiner.o\allowbreak{}ctonionP\allowbreak{}ictureVo\allowbreak{}lumeElem\allowbreak{}entDiago\allowbreak{}nal} in the Wolfram report).

\textbf{Composing the two intertwiners.} The product $K_{\mathrm{clifford}}K_{\mathrm{octonion}}$ goes from the octonion picture to the tensor picture: $\hat\gamma^a(K_{\mathrm{clifford}}K_{\mathrm{octonion}})=K_{\mathrm{clifford}}\gamma^aK_{\mathrm{octonion}}=(K_{\mathrm{clifford}}K_{\mathrm{octonion}})\Gamma^a$. Made primitive, it equals the "canonical Clifford intertwiner" published by dirac-main in \texttt{triality\allowbreak{}44.json} (erratum E4; Stage 1, Result 3.8; this comparison, like the other dirac-main comparisons, needs the dirac-main folder).

\subsection{What a change of picture does, and what it does not}

\begin{itemize}
\item \textbf{It changes coordinates only.} By Theorem 3.9 the three pictures are one Clifford module in three bases. Every statement of Chapter 2 (the trace lemma, the charge matrix, the chirality, Pin(4,4), Spin(4,4), the representation theorem) holds in each picture with the corresponding matrices. The project uses the notebook's picture as its primary basis (a convention, Stage 1, §1.2).
\item \textbf{It does not make the octonions associative.} Matrix multiplication is associative, and the Clifford relation involves only two vectors at a time; the non-associativity of the split octonions is invisible there. It appears when multiplication matrices are multiplied: $L_{e_1}L_{e_2}$ (first multiply by $e_2$, then by $e_1$, since $L_{e_1}L_{e_2}y=e_1(e_2y)$) is not $L_{e_1e_2}=-L_{e_3}$ (which gives $(e_1e_2)y$), so, unlike the 4 by 4 quaternion matrices, which obey $L_xL_y=L_{xy}$ by associativity (Section 3.3), the 8 by 8 octonion multiplication matrices do not turn products of octonions into products of matrices; the two differ exactly on $e_4,\dots,e_7$ (Exercise 3.9), which is the associator $(e_1e_2)e_4-e_1(e_2e_4)=2e_7$ of Section 3.8 seen as matrices. Stage 1 says it the same way: the intertwiners "do not make the octonion product associative".
\item \textbf{The notebook's Clifford algebra is correct.} Every algebra check of Stage 1 about the notebook's gammas, $\sigma_{16}$ and expression [1] is true. The errors of the notebook that the project found lie elsewhere, for example in the contraction of the spin connection (Chapter 4) and in its field equations (Chapter 9).
\end{itemize}

\subsection{What the computer checked}

All checks below are true in both algebra reports, \texttt{wolfram-\allowbreak{}algebra-\allowbreak{}report.j\allowbreak{}son} and \texttt{python-a\allowbreak{}lgebra-r\allowbreak{}eport.js\allowbreak{}on} (folder \texttt{artifact\allowbreak{}s/dirac1\allowbreak{}6complex\allowbreak{}/arbitra\allowbreak{}ry-field\allowbreak{}/}), unless marked. Some items in a row are not part of the check's true/false verdict but are measurements that the check records: the block form of $K_{\mathrm{octonion}}$, $Q^TQ$, $Q^T\sigma Q$, the $\tau_a$ comparisons, the sign of $K\gamma^8K^{-1}$ and the images of the two chirality blocks under $K_{\mathrm{clifford}}$, and the signed-permutation tests. Their committed values, which agree with the table, are the measurements \texttt{ALG\_clif\allowbreak{}fordPict\allowbreak{}ureInter\allowbreak{}twiner.*} and \texttt{ALG\_octo\allowbreak{}nionPict\allowbreak{}ureInter\allowbreak{}twiner.*} of the Wolfram report and \texttt{ALG\_clif\allowbreak{}fordK*}, \texttt{ALG\_octo\allowbreak{}nionK*} and \texttt{\detokenize{ALG_tau*}} of the Python report. $Q^T\sigma Q$, the block images, the associator $2e_7$ and the 64 structure constants are computed by the Wolfram verifier only. The rerun commands are those of Section 2.16.

\begingroup
\small
\begin{longtable}{@{}>{\raggedright\arraybackslash}p{0.400\linewidth}>{\raggedright\arraybackslash}p{0.400\linewidth}@{}}
\toprule
\textbf{Statement (where in this chapter; where in Stage 1)} & \textbf{Check} \\
\midrule
\endfirsthead
\toprule
\textbf{Statement (where in this chapter; where in Stage 1)} & \textbf{Check} \\
\midrule
\endhead
the notebook's $\eta$, $\sigma$, $\tau_a$, $\bar\tau_a$, $\gamma^a$ and $\sigma_{16}$, rebuilt from the 11 verbatim definition cells, equal the package's matrices; expression [1] is eight times True (3.5, 3.6; Result 3.4) & \texttt{ALG\_expr\allowbreak{}ession1} \\
the Wolfram verifier finds all 75 matrices of \texttt{algebra-\allowbreak{}fixture.\allowbreak{}json} (among them $\tau_a$, $\bar\tau_a$, the gammas of all three pictures and both intertwiners) equal to its own (Stage 1, §3.4) & \texttt{ALG\_fixt\allowbreak{}ureAgree\allowbreak{}ment} \\
Clifford relation of the notebook's gammas (3.6, Theorem 3.4; Result 3.1) & \texttt{ALG\_clif\allowbreak{}ford} \\
$C=\mathrm{diag}(-\sigma,\sigma)$ (3.6, Theorem 3.5; Result 3.3) & \texttt{ALG\_char\allowbreak{}geMatrix} \\
$\gamma^8=\mathrm{diag}(-I_8,I_8)$ (3.6, Theorem 3.5; Result 3.6) & \texttt{ALG\_chir\allowbreak{}ality} \\
tensor gammas satisfy the Clifford relation; $K_{\mathrm{clifford}}$: dimension 1, rank 16, $KCK^{-1}=C_{\mathrm{dm}}$, $K\gamma^8K^{-1}=G\otimes G\otimes G\otimes G$, the images of the two blocks, not a signed permutation (3.11; Result 3.7) & \texttt{ALG\_clif\allowbreak{}fordPict\allowbreak{}ureInter\allowbreak{}twiner} \\
Zorn product: unit, composition law, reversal of products under conjugation, $N=\eta$, associator $2e_7$, 64 structure constants, equal to dirac-main's table; octonion gammas satisfy the Clifford relation; $K_{\mathrm{octonion}}$: dimension 1, rank 16, $\mathrm{diag}(Q,Q)$, $Q^TQ=2I_8$, $Q^T\sigma Q=2\eta$, $\tau_a=QL_{e_a}Q^{-1}$; $\tau_a=L_{e_a}$ only for $a=0$; no signed permutation (3.8, 3.9, 3.12; Result 3.8) & \texttt{ALG\_octo\allowbreak{}nionPict\allowbreak{}ureInter\allowbreak{}twiner} \\
the Python checker reads the Wolfram report and finds the same $K_{\mathrm{clifford}}$, $K_{\mathrm{octonion}}$ and chirality, and 47 of 47 shared measurements equal (Python report only; Stage 1, §3.4) & \texttt{ALG\_wolf\allowbreak{}ramAgree\allowbreak{}ment} \\
\bottomrule
\end{longtable}
\endgroup

The comparisons with the three files of dirac-main (\texttt{cl44-see\allowbreak{}d.json}, \texttt{split-oc\allowbreak{}tonion.j\allowbreak{}son}, \texttt{triality\allowbreak{}44.json}) are sub-checks of the sixth and seventh rows, \texttt{ALG\_clif\allowbreak{}fordPict\allowbreak{}ureInter\allowbreak{}twiner} and \texttt{ALG\_octo\allowbreak{}nionPict\allowbreak{}ureInter\allowbreak{}twiner}; they need the separately published dirac-main folder and record "not-run" without it (Stage 1, §12). The \texttt{cl44-see\allowbreak{}d.json} comparison is made by the Wolfram verifier only. The Python checker records its comparisons with the other two files as measurements, all true in the committed report: \texttt{ALG\_zorn\allowbreak{}ProductE\allowbreak{}qualsDir\allowbreak{}acMainSp\allowbreak{}litOcton\allowbreak{}ionJson} (the multiplication table of Section 3.8), \texttt{ALG\_octo\allowbreak{}nionGamm\allowbreak{}asEqualD\allowbreak{}iracMain\allowbreak{}Triality\allowbreak{}44Json} (the octonion gammas of Section 3.9) and \texttt{ALG\_Kcli\allowbreak{}ffordKoc\allowbreak{}tonionEq\allowbreak{}ualsDira\allowbreak{}cMainCan\allowbreak{}onicalIn\allowbreak{}tertwine\allowbreak{}r} (the composed intertwiner of Section 3.12).

\subsection{What we proved and what we assumed}

\textbf{Proved in this chapter.} The sign rules of permutations and the four vector identities (V1) to (V4). Quaternions are associative (through the Pauli matrices), and the notebook's six 4 by 4 blocks are right and (minus) left multiplications by the quaternion units (Theorem 3.1), which gives all their products (Corollary 3.2). The seven matrices $\tau_1,\dots,\tau_7$ anticommute with squares $-\eta_{aa}$, $\bar\tau_a=-\tau_a$, $\tau_7=\mathrm{diag}(-I_4,I_4)$ and $\tau_1\cdots\tau_7=I_8$ (Theorem 3.3). From this, by hand: the Clifford relation of the notebook's gammas and their symmetry pattern (Theorem 3.4), and $C=\mathrm{diag}(-\sigma,\sigma)$ and $\gamma^8=\mathrm{diag}(-I_8,I_8)$ (Theorem 3.5); these are the three facts that Chapter 2 had borrowed from the machine checks. In Zorn's model of the split octonions: the unit, $N=\eta$, the identities $\bar x(xy)=x(\bar xy)=N(x)y$, the composition law $N(xy)=N(x)N(y)$, the reversal of products under conjugation (Theorem 3.6), the non-associativity and (Exercise 3.8) the alternative laws. The octonion gammas satisfy the Clifford relation (Theorem 3.7) and have the same symmetry pattern (Lemma 3.8, Exercise 3.10). For any two sets of real 16 by 16 gammas of signature (4,4) an intertwiner exists, every nonzero one is invertible, and it is unique up to a factor (Theorem 3.9); hence the images of $C$ and $\gamma^8$ and, with a dimension count, of the two chirality blocks under $K_{\mathrm{clifford}}$. The block form $\mathrm{diag}(Q,Q)$ of $K_{\mathrm{octonion}}$ is equivalent to $\tau_a=QL_{e_a}Q^{-1}$, that is, $\tau_a$ is $L_{e_a}$ written in the basis of the columns of $Q^{-1}=\tfrac12Q^T$; the light-cone formula for the norm and the chirality of the octonion picture follow.

\textbf{Computed and taken from the reports.} The two particular matrices $K_{\mathrm{clifford}}$ and $Q$ are the exact solutions of linear equations computed by the verifiers; that they satisfy their equations for all eight $a$ is a machine result (we verified one column of one equation by hand). The same holds for the statement of erratum E5 that the equation $\tau_a=P^{-1}L_{e_a}P$ has a one-dimensional solution space, and for the comparisons with the published files of dirac-main.

\textbf{Quoted, and conventions.} Hurwitz's theorem is quoted as motivation only; triality is only described. The primary basis (the notebook's), the primitive integer normalization of the intertwiners and the Zorn coordinates of dirac-main are conventions. Nothing in this chapter is a physical statement.

\subsection{Exercises}

\textbf{Exercise 3.1.} Show that exchanging two entries of a list of distinct numbers changes the sign of the permutation. Compute Signature[\{2,1,3,4\}], Signature[\{3,1,2,4\}], Signature[\{1,2,2,4\}], $\epsilon_{1324}$ and $\epsilon_{2314}$.

\textbf{Exercise 3.2.} With the Pauli-matrix images of Section 3.3, verify $\mathbf{ki}=\mathbf j$ and $\mathbf{ik}=-\mathbf j$.

\textbf{Exercise 3.3.} Multiply $s_2s_3$ and $t_2t_3$ directly and compare with Corollary 3.2 (2).

\textbf{Exercise 3.4.} Verify $s_1t_1=t_1s_1$ by direct multiplication.

\textbf{Exercise 3.5.} For an antisymmetric 4 by 4 matrix $A$ (1-based indices) define the \textbf{dual} $({\ast}A)_{pq}:=\tfrac12\sum_{r,s}\epsilon_{pqrs}A_{rs}$ with $\epsilon_{1234}=+1$. Show that ${\ast}s_1=s_1$ (self-dual) and ${\ast}t_1=-t_1$ (anti-self-dual).

\textbf{Exercise 3.6.} Using the table of Section 3.5, check in row 0 that $\tau_1\tau_4=-\tau_4\tau_1$.

\textbf{Exercise 3.7.} Compute $e_4e_1$ and $e_1e_4$ in Zorn form.

\textbf{Exercise 3.8.} Prove the alternative laws $x(xy)=(xx)y$ and $(yx)x=y(xx)$ for all split octonions $x,y$. (Hint: $x+\bar x=2x_0e_0$.)

\textbf{Exercise 3.9.} Apply $L_{e_1}L_{e_2}$ and $L_{e_1e_2}$ to $e_0$ and to $e_4$, and compare.

\textbf{Exercise 3.10.} Show that $\Gamma^a$ is symmetric for $a\le3$ and antisymmetric for $a\ge4$.

\textbf{Exercise 3.11.} Prove that $K_{\mathrm{clifford}}^TK_{\mathrm{clifford}}=4I_{16}$. (Hint: both the notebook's and the tensor gammas satisfy $(\gamma^a)^T=\eta^{aa}\gamma^a$.)

\textbf{Exercise 3.12.} Derive $\bar\tau_a=QL_{\bar e_a}Q^{-1}$ from $\tau_a=QL_{e_a}Q^{-1}$, $Q^TQ=2I_8$, $Q^T\sigma Q=2\eta$ and Lemma 3.8, without using $\bar\tau_a=-\tau_a$.

\textbf{Exercise 3.13.} Compute $y=Qx$ and $y_0y_4+y_1y_5+y_2y_6+y_3y_7$ for $x=e_0$, $x=e_0+e_7$ and $x=e_3+e_4$, and compare with $N(x)$.

\textbf{Exercise 3.14.} Show that, in the notebook's basis, $B=-iC\gamma^4=-i\begin{pmatrix}0&\sigma\tau_4\\ \sigma\tau_4&0\end{pmatrix}$, and use it to compute $(Bu)_0$.

\textbf{Exercise 3.15.} The notebook basis spinor $e_8$ has chirality $+1$. Explain why its image $K_{\mathrm{clifford}}e_8$ is a combination of the tensor-basis columns numbered 0, 3, 5, 6, 9, 10, 12 and 15, and check this with the displayed matrix.

\subsection{Answers to the exercises}

\textbf{Answer 3.1.} Exchanging two \emph{neighbouring} entries changes the relative order of exactly one pair (those two entries) and of no other pair, so the number of inversions changes by one and the sign flips. Exchanging the entries in positions $p<r$ can be done by $2(r-p)-1$ neighbour exchanges (move the first entry right by $r-p$ steps, then the other one left by $r-p-1$ steps), an odd number, so the sign flips again. Values: $(2,1,3,4)$ has one inversion, sign $-1$; $(3,1,2,4)$ has two, $+1$; $\{1,2,2,4\}$ has a repeated entry, 0; $\epsilon_{1324}$: $(1,3,2,4)$ has one inversion, $-1$; $\epsilon_{2314}$: $(2,3,1,4)$ has two, $+1$.

\textbf{Answer 3.2.} $(-i\sigma_z)(-i\sigma_x)=-\sigma_z\sigma_x=-i\sigma_y$, the image of $\mathbf j$ (using $\sigma_z\sigma_x=i\sigma_y$ from Exercise 2.1). And $(-i\sigma_x)(-i\sigma_z)=-\sigma_x\sigma_z=+\sigma_z\sigma_x=i\sigma_y$, the image of $-\mathbf j$.

\textbf{Answer 3.3.} Row $p$ of $s_2s_3$ is the combination of the rows of $s_3$ given by row $p$ of $s_2$. With the rows of $s_3$ being $(0,1,0,0)$, $(-1,0,0,0)$, $(0,0,0,1)$, $(0,0,-1,0)$ and the rows of $s_2$ picking $-$(row 3), row 4, row 1, $-$(row 2):

\[
s_2s_3=\begin{pmatrix}0&0&0&-1\\0&0&-1&0\\0&1&0&0\\1&0&0&0\end{pmatrix}=-s_1 .
\]

In the same way, the rows of $t_2$ pick $-$(row 3), $-$(row 4), row 1, row 2 of $t_3$, whose rows are $(0,1,0,0)$, $(-1,0,0,0)$, $(0,0,0,-1)$, $(0,0,1,0)$, so $t_2t_3=\begin{pmatrix}0&0&0&1\\0&0&-1&0\\0&1&0&0\\-1&0&0&0\end{pmatrix}=-t_1$. Both agree with $X_2X_3=-\epsilon_{231}X_1=-X_1$.

\textbf{Answer 3.4.} The rows of $s_1$ pick (row 4), (row 3), $-$(row 2), $-$(row 1) of the second factor; the rows of $t_1$ pick $-$(row 4), (row 3), $-$(row 2), (row 1). With the rows of $t_1$, $(0,0,0,-1)$, $(0,0,1,0)$, $(0,-1,0,0)$, $(1,0,0,0)$, the product $s_1t_1$ has the rows $(1,0,0,0)$, $(0,-1,0,0)$, $(0,0,-1,0)$, $(0,0,0,1)$. With the rows of $s_1$, $(0,0,0,1)$, $(0,0,1,0)$, $(0,-1,0,0)$, $(-1,0,0,0)$, the product $t_1s_1$ has the rows $(1,0,0,0)$, $(0,-1,0,0)$, $(0,0,-1,0)$, $(0,0,0,1)$ as well: both equal $\mathrm{diag}(1,-1,-1,1)$.

\textbf{Answer 3.5.} Because $A$ and $\epsilon$ are both antisymmetric in the pair $(r,s)$, the sum over ordered pairs is twice the sum over $r<s$, so $({\ast}A)_{pq}=\sum_{r<s}\epsilon_{pqrs}A_{rs}$. This gives $({\ast}A)_{12}=A_{34}$, $({\ast}A)_{13}=-A_{24}$, $({\ast}A)_{14}=A_{23}$, $({\ast}A)_{23}=A_{14}$, $({\ast}A)_{24}=-A_{13}$, $({\ast}A)_{34}=A_{12}$ (for example $\epsilon_{1423}=+1$, since $(1,4,2,3)$ has two inversions). The only nonzero entries above the diagonal of $s_1$ are $A_{14}=1$ and $A_{23}=1$, so ${\ast}s_1$ has $({\ast}s_1)_{14}=A_{23}=1$, $({\ast}s_1)_{23}=A_{14}=1$ and zeros elsewhere: ${\ast}s_1=s_1$. For $t_1$, $A_{14}=-1$ and $A_{23}=1$, so $({\ast}t_1)_{14}=1=-A_{14}$ and $({\ast}t_1)_{23}=-1=-A_{23}$: ${\ast}t_1=-t_1$.

\textbf{Answer 3.6.} $(\tau_1\tau_4u)_0=(\tau_4u)_7=-u_2$ (row 0 of $\tau_1$ is $+7$, row 7 of $\tau_4$ is $-2$), and $(\tau_4\tau_1u)_0=(\tau_1u)_5=+u_2$ (row 0 of $\tau_4$ is $+5$, row 5 of $\tau_1$ is $+2$). They are opposite.

\textbf{Answer 3.7.} $e_4=(1,0,0,-1)$ and $e_1=(0,\hat e_1,-\hat e_1,0)$. For $e_4e_1$: $a=1$, $u=v=0$, $b=-1$; $c=d=0$, $r=\hat e_1$, $w=-\hat e_1$. The product is $(0,\ \hat e_1,\ (-1)(-\hat e_1),\ 0)=(0,\hat e_1,\hat e_1,0)=e_5$. For $e_1e_4$: $a=b=0$, $u=\hat e_1$, $v=-\hat e_1$; $c=1$, $r=w=0$, $d=-1$. The product is $(0,\ -\hat e_1,\ -\hat e_1,\ 0)=-e_5$. Both agree with the table and with the Wolfram worked examples.

\textbf{Answer 3.8.} From the definition of the conjugate, $\bar x=2x_0e_0-x$. By Theorem 3.6 (1), $N(x)y=\bar x(xy)=2x_0\,xy-x(xy)$, so $x(xy)=2x_0\,xy-N(x)y$. On the other hand $xx=x(2x_0e_0-\bar x)=2x_0x-N(x)e_0$ by Theorem 3.6 (2), so $(xx)y=2x_0\,xy-N(x)y$. The two agree. For the second law apply the first to $\bar x,\bar y$, $\bar x(\bar x\bar y)=(\bar x\bar x)\bar y$, and conjugate both sides with Theorem 3.6 (4): the left side becomes $\overline{\bar x\bar y}\,x=(yx)x$ and the right side $y\,\overline{\bar x\bar x}=y(xx)$.

\textbf{Answer 3.9.} On $e_0$: $e_1(e_2e_0)=e_1e_2=-e_3$ and $(e_1e_2)e_0=-e_3$: equal. On $e_4$: $e_1(e_2e_4)=e_1(-e_6)=-e_7$ but $(e_1e_2)e_4=-e_3e_4=e_7$: opposite. So $L_{e_1}L_{e_2}\ne L_{e_1e_2}$. Going on with the table: on $e_5$, $e_1(e_2e_5)=-e_1e_7=e_6$ but $-e_3e_5=-e_6$; on $e_6$, $e_1(e_2e_6)=e_1e_4=-e_5$ but $-e_3e_6=e_5$; on $e_7$, $e_1(e_2e_7)=e_1e_5=e_4$ but $-e_3e_7=-e_4$. On $e_0,\dots,e_3$ the two agree, because these four elements multiply among themselves like the quaternions ($-e_1,-e_2,-e_3$ play the roles of $\mathbf i,\mathbf j,\mathbf k$), which are associative. So the two matrices agree on $e_0,\dots,e_3$ and differ by a sign on each of $e_4,\dots,e_7$.

\textbf{Answer 3.10.} $(\Gamma^a)^T=\begin{pmatrix}0&L_{e_a}^T\\ L_{\bar e_a}^T&0\end{pmatrix}$. For $a=0$ this is $\Gamma^0=\begin{pmatrix}0&I_8\\ I_8&0\end{pmatrix}$, symmetric. For $a\ge1$, $L_{\bar e_a}=-L_{e_a}$. If $a\le3$, $L_{e_a}^T=-L_{e_a}=L_{\bar e_a}$ and $L_{\bar e_a}^T=-L_{e_a}^T=L_{e_a}$, so $(\Gamma^a)^T=\Gamma^a$. If $a\ge4$, $L_{e_a}^T=L_{e_a}=-L_{\bar e_a}$ and $L_{\bar e_a}^T=-L_{e_a}$, so $(\Gamma^a)^T=-\Gamma^a$.

\textbf{Answer 3.11.} Transpose $\hat\gamma^aK=K\gamma^a$: $K^T(\hat\gamma^a)^T=(\gamma^a)^TK^T$, that is, $\eta^{aa}K^T\hat\gamma^a=\eta^{aa}\gamma^aK^T$, so $K^T\hat\gamma^a=\gamma^aK^T$. Then $K^TK\gamma^a=K^T\hat\gamma^aK=\gamma^aK^TK$: $K^TK$ commutes with every $\gamma^a$, so $K^TK=c\,I_{16}$ by Corollary 2.3. The entry $(0,0)$ of $K^TK$ is the sum of the squares of column 0 of $K$, which has four entries $\pm1$ (worked example of Section 3.11): $c=4$. So $\tfrac12K_{\mathrm{clifford}}$ is orthogonal. The same argument for $K_{\mathrm{octonion}}$ gives $Q^TQ=2I_8$ (column 0 of $Q$ has two entries $\pm1$).

\textbf{Answer 3.12.} From $Q^TQ=2I_8$: $Q^{-1}=\tfrac12Q^T$ and $(Q^T)^{-1}=\tfrac12Q$. From $Q^T\sigma Q=2\eta$: $\sigma=2(Q^T)^{-1}\eta Q^{-1}=\tfrac12Q\eta Q^T$. Next $\tau_a^T=(Q^{-1})^TL_{e_a}^TQ^T=\tfrac12QL_{e_a}^TQ^T$. Hence

\[
\bar\tau_a=\sigma\tau_a^T\sigma=\tfrac18\,Q\eta(Q^TQ)L_{e_a}^T(Q^TQ)\eta Q^T=\tfrac12\,Q\,\bigl(\eta L_{e_a}^T\eta\bigr)Q^T=Q\bigl(\eta L_{e_a}^T\eta\bigr)Q^{-1}.
\]

By Lemma 3.8 (with $N(e_a)=\eta_{aa}\ne0$ and $\eta^{-1}=\eta$), $\eta L_{e_a}^T\eta=L_{\bar e_a}$. So $\bar\tau_a=QL_{\bar e_a}Q^{-1}$.

\textbf{Answer 3.13.} For $x=e_0$: $y_0=y_4=1$ and all other $y_i=0$, so the sum is $1=N(e_0)$. For $x=e_0+e_7$: $y_0=1-1=0$, $y_4=1+1=2$, all other $y_i=0$; the sum is $0\cdot2=0=1-1=N(x)$. For $x=e_3+e_4$: $y_1=-1-1=-2$, $y_5=-1+1=0$, all other $y_i=0$; the sum is $0=1-1=N(x)$. The last two are null vectors, and in the light-cone coordinates only one member of each pair is nonzero.

\textbf{Answer 3.14.} $C=\mathrm{diag}(-\sigma,\sigma)$ and $\gamma^4=\begin{pmatrix}0&\bar\tau_4\\ \tau_4&0\end{pmatrix}=\begin{pmatrix}0&-\tau_4\\ \tau_4&0\end{pmatrix}$. By the block rule $C\gamma^4=\begin{pmatrix}0&(-\sigma)(-\tau_4)\\ \sigma\tau_4&0\end{pmatrix}=\begin{pmatrix}0&\sigma\tau_4\\ \sigma\tau_4&0\end{pmatrix}$, and $B=-iC\gamma^4$. Row 0 of $\sigma\tau_4$ is row 4 of $\tau_4$ ($\sigma$ exchanges the halves of the rows), which is $-1$: $(\sigma\tau_4w)_0=-w_1$. The upper half of $Bu$ is $-i\,\sigma\tau_4$ applied to the lower half $(u_8,\dots,u_{15})$ of $u$, so $(Bu)_0=-i\cdot(-u_{8+1})=i\,u_9$, as in the table of Section 2.9.

\textbf{Answer 3.15.} The notebook basis spinor $e_8$ lies in the lower half, so $\gamma^8e_8=+e_8$. As in Section 3.11, $\hat\gamma^8(Ke_8)=K\gamma^8e_8=+Ke_8$: $Ke_8$ lies in the $+1$ eigenspace of $G\otimes G\otimes G\otimes G$, which is spanned by the tensor-basis columns with an even number of 1-bits, $\{0,3,5,6,9,10,12,15\}$. In the displayed matrix, column 8 has the nonzero entries $+1$ in row 0, $-1$ in row 6, $-1$ in row 9 and $+1$ in row 15, and $0,6,9,15$ are all in the list.

\section{Curved space from zero}

In Einstein's theory, gravity is not a force that acts inside space and time. It is the shape of space and time itself. A heavy body bends the geometry around it, and every other body moves along the straightest paths that the bent geometry allows. This chapter builds, from nothing, the geometry that the dirac16complex project uses. We start with coordinates on a space of eight dimensions and the metric, which measures lengths and time intervals. Then come the Christoffel symbols, which say how to differentiate when the coordinates are curved, the geodesics (the straightest paths), the curvature, and the field equations of gravity: the Einstein equations and their higher-order relatives, the Einstein-Lovelock equations. A spinor field such as dirac16complex needs more structure than a metric. It needs, at every point, a set of eight reference directions that are perpendicular to each other and of unit length. This set is the vielbein. It also needs a rule, the spin connection, that says how these directions turn from point to point. The vielbein postulate fixes that rule, and with it we build the covariant derivative of a spinor. At the end of the chapter we can say exactly why the contraction that the author's notebook uses for the spin connection is wrong, and we check the claim with small numbers.

\subsection{What this chapter does, and what it needs}

\textbf{What it delivers.} The chapter derives, step by step:

\begin{itemize}
\item the transformation rules of vectors, covectors and tensors, and the metric with its signature and volume element;
\item the covariant derivative, and the unique Christoffel symbols of a metric (the Levi-Civita connection);
\item the geodesic equation, solved completely in one worked example;
\item the Riemann curvature tensor, its symmetries, the Ricci tensor, the scalar curvature, the Bianchi identities and the Einstein tensor;
\item the Einstein equations and the Einstein-Lovelock equations, with an honest account of what the project computes with them;
\item the vielbein, the spin connection, the vielbein postulate and its unique solution;
\item the covariant derivative of a spinor, its behaviour under a change of frame, the curvature of the spinor connection, and the square of the Dirac operator;
\item the error in the notebook's contraction of the spin connection.
\end{itemize}

Every general formula is tried on small examples (the flat plane in polar coordinates, the sphere, a two-dimensional spinor), and then on the primordial field of the author's notebook, whose Christoffel symbols and spin connection we derive completely.

\textbf{What it needs from Chapter 1.} Columns of numbers (vectors), matrices, the matrix product, the transpose, the inverse and the determinant; indices counted from 0; the Kronecker delta $\delta^\mu{}_\nu$, which is 1 if $\mu=\nu$ and 0 otherwise; partial derivatives and the chain rule in several variables. We repeat the chain rule because it is the engine of this chapter. If a function $f$ depends on variables $y^0,\dots,y^{n-1}$, and each $y^\nu$ depends on variables $x^0,\dots,x^{n-1}$, then

\[
\frac{\partial f}{\partial x^\mu}=\sum_{\nu=0}^{n-1}\frac{\partial f}{\partial y^\nu}\,\frac{\partial y^\nu}{\partial x^\mu}.
\]

For smooth functions (functions whose partial derivatives of every order exist and are continuous) mixed partial derivatives do not depend on the order: $\partial_\mu\partial_\nu f=\partial_\nu\partial_\mu f$, where $\partial_\mu$ is short for $\partial/\partial x^\mu$.

\textbf{The summation convention.} When an index appears twice in a product, once as an upper index and once as a lower index, it is summed over all its values. So $V^\mu W_\mu$ means $V^0W_0+V^1W_1+\dots+V^{n-1}W_{n-1}$, and $g_{\mu\nu}V^\nu$ means $\sum_\nu g_{\mu\nu}V^\nu$. An index that appears only once is free: an equation with a free index $\mu$ is $n$ equations, one for each value of $\mu$. Where we do not want a sum over a repeated index we write "(no sum)".

\textbf{Names of indices.} $n$ is the dimension of the space: $n=8$ for the project, $n=2$ in the small examples. Greek letters $\alpha,\beta,\lambda,\mu,\nu,\rho,\sigma$ are coordinate indices and run from 0 to $n-1$. Latin letters $a,b,c,d,e,f$ are frame indices (Section 4.10) and also run from 0 to $n-1$. We never use $\gamma$ as an index, because $\gamma$ names the gamma matrices.

\textbf{Coordinates.} The notebook and the Stage documents name the eight coordinates $x_0,\dots,x_7$. In formulas with indices we write them $x^\mu$, with the index up, because their small changes $dx^\mu$ are the components of a vector (Section 4.3). The names $x_4$ and $x^4$ denote the same coordinate, the evolution time; we never lower the index of a coordinate. The roles are those of the notebook: $x^0$ is a hidden space direction, $x^1,x^2,x^3$ are ordinary 3-space, $x^4$ is the evolution time, and $x^5,x^6,x^7$ are three extra time directions. The flat metric is

\[
\eta=\mathrm{diag}(+1,+1,+1,+1,-1,-1,-1,-1),
\]

so directions 0 to 3 are space-like and 4 to 7 are time-like: the signature is (4,4).

\textbf{What it needs from Chapter 2.} The eight real 16 by 16 gamma matrices $\gamma^0,\dots,\gamma^7$ of the notebook, with the Clifford relation $\gamma^a\gamma^b+\gamma^b\gamma^a=2\eta^{ab}\,1$; the 28 spin generators $S^{ab}=\tfrac14[\gamma^a,\gamma^b]$, where $[A,B]=AB-BA$ is the commutator (and $\{A,B\}=AB+BA$ the anticommutator); the charge matrix $C=\gamma^0\gamma^1\gamma^2\gamma^3$; and three facts proved there and verified by the Stage-1 algebra reports:

\begin{itemize}
\item (F1) the only complex 16 by 16 matrices that commute with all eight $\gamma^a$ are the multiples of the identity (Stage 1, Result 4.1; check \texttt{ALG\_pinI\allowbreak{}rreducib\allowbreak{}leComple\allowbreak{}x} in \texttt{wolfram-\allowbreak{}algebra-\allowbreak{}report.j\allowbreak{}son} and \texttt{python-a\allowbreak{}lgebra-r\allowbreak{}eport.js\allowbreak{}on});
\item (F2) the 256 ordered products $\gamma^{a_1}\cdots\gamma^{a_k}$ with $a_1<\dots<a_k$ (including the identity, $k=0$) are linearly independent; in particular the eight $\gamma^a$ are linearly independent, and so are the 28 matrices $S^{ab}=\tfrac12\gamma^a\gamma^b$ with $a<b$ (Stage 1, Result 3.2; check \texttt{ALG\_fait\allowbreak{}hful});
\item (F3) $(S^{ab})^TC=-CS^{ab}$ for all $a,b$ (Stage 1, Result 3.5; check \texttt{ALG\_spin\allowbreak{}Transpos\allowbreak{}ePropert\allowbreak{}ies}).
\end{itemize}

For the two-dimensional spinor examples we use the Pauli matrices $\sigma_x,\sigma_y,\sigma_z$ of Section 2.3 under their numbered names, defined there, $\sigma_1=\sigma_x$, $\sigma_2=\sigma_y$, $\sigma_3=\sigma_z$:

\[
\sigma_1=\begin{pmatrix}0&1\\ 1&0\end{pmatrix},\qquad \sigma_2=\begin{pmatrix}0&-i\\ i&0\end{pmatrix},\qquad \sigma_3=\begin{pmatrix}1&0\\ 0&-1\end{pmatrix},
\]

which satisfy $\sigma_k^2=1$, $\sigma_1\sigma_2=i\sigma_3$, $\sigma_2\sigma_3=i\sigma_1$, $\sigma_3\sigma_1=i\sigma_2$, and anticommute in pairs. Multiplying out the matrices checks each of these in a line.

\subsection{Manifolds and coordinates}

\textbf{The idea.} The surface of the Earth is curved, but a small piece of it looks flat, and on such a piece we can label every point by two numbers, for example latitude and longitude. A manifold is a space with this property: near each of its points it can be labelled by $n$ real numbers, and the labels of overlapping pieces are related by smooth formulas.

\textbf{Definition (chart, manifold).} A set $U$ of $\mathbb R^n$ is open if with each of its points it also contains a small ball around that point. A chart on a set $M$ is a part $U$ of $M$ together with a one-to-one map that assigns to each point $p$ of $U$ its $n$ coordinates $(x^0(p),\dots,x^{n-1}(p))$, such that the image is an open set of $\mathbb R^n$. A manifold of dimension $n$ is a set covered by charts such that, wherever two charts overlap, the coordinates of one chart are smooth functions of the coordinates of the other, and conversely. Such a pair of formulas $x'^\mu=x'^\mu(x)$, $x^\mu=x^\mu(x')$ is a change of coordinates.

\textbf{The Jacobian matrix.} For a change of coordinates the $n$ by $n$ matrices $\partial x'^\mu/\partial x^\nu$ and $\partial x^\mu/\partial x'^\nu$ are inverse to each other. To see this, apply the chain rule to the identity $x'^\mu(x(x'))=x'^\mu$ and differentiate with respect to $x'^\rho$:

\[
\frac{\partial x'^\mu}{\partial x^\nu}\,\frac{\partial x^\nu}{\partial x'^\rho}=\frac{\partial x'^\mu}{\partial x'^\rho}=\delta^\mu{}_\rho .
\]

The same argument with the primes exchanged gives $(\partial x^\mu/\partial x'^\nu)(\partial x'^\nu/\partial x^\rho)=\delta^\mu{}_\rho$.

\textbf{Example 1: the plane.} The Cartesian chart $(x,y)$ covers the whole plane. The polar chart $(r,\varphi)$ with $r>0$ and $0<\varphi<2\pi$ covers the plane without the half-line $y=0$, $x\ge0$; the change of coordinates is $x=r\cos\varphi$, $y=r\sin\varphi$. Its Jacobian matrix and determinant are

\[
\frac{\partial(x,y)}{\partial(r,\varphi)}=\begin{pmatrix}\cos\varphi&-r\sin\varphi\\ \sin\varphi&r\cos\varphi\end{pmatrix},\qquad \det=r\cos^2\varphi+r\sin^2\varphi=r>0 .
\]

At $r=0$ the determinant vanishes and the polar chart fails, although nothing is wrong with the plane there. This is our first example of a coordinate singularity: a failure of the labels, not of the space.

\textbf{Example 2: the sphere.} On the sphere of radius $a$ the chart $(\theta,\varphi)$ with $0<\theta<\pi$ and $0<\varphi<2\pi$ labels the point with Cartesian position $(a\sin\theta\cos\varphi,\ a\sin\theta\sin\varphi,\ a\cos\theta)$. It misses the two poles and one half-circle between them; a second chart covers those.

\textbf{Example 3: the project.} The project works on an 8-dimensional manifold $M$ with coordinates $x^0,\dots,x^7$. For the primordial field of the notebook (Chapter 9) the chart is described by $z=6Hx^0$ and $t=Hx^4$, where $H>0$ is the notebook's single inverse length, and $z$ runs over the interval $(0,\pi/2)$. At $z=\pi/2$ this chart ends: the metric component $g_{00}=\cot^2z$ of Section 4.4 vanishes there. In the coordinate $\zeta=\ln(\sin z)/(6H)$ the same metric (Section 4.4) is perfectly regular at $\zeta=0$, which corresponds to $z=\pi/2$. (The Stage-2 document, §4.4, verifies this form of the metric with the check \texttt{P\_zeta\_w\allowbreak{}arpedMet\allowbreak{}ric} in \texttt{wolfram-\allowbreak{}primordi\allowbreak{}al-repor\allowbreak{}t.json} and records its regularity at $\zeta=0$ as an observation read off from it.) As in the plane, a formula that breaks down may be a failure of the chart, not of the space.

\subsection{Vectors, covectors and tensors}

\textbf{Tangent vectors.} A curve is a smooth map from an interval of a parameter $\lambda$ into the manifold, written in a chart as $x^\mu(\lambda)$. Its velocity at a point has the components $u^\mu=dx^\mu/d\lambda$. In another chart the same curve is $x'^\mu(\lambda)=x'^\mu(x(\lambda))$, and by the chain rule

\[
u'^\mu=\frac{dx'^\mu}{d\lambda}=\frac{\partial x'^\mu}{\partial x^\nu}\,\frac{dx^\nu}{d\lambda}=\frac{\partial x'^\mu}{\partial x^\nu}\,u^\nu .
\]

\textbf{Definition (vector).} A vector at a point $p$ is an assignment of $n$ numbers $V^\mu$ to every chart around $p$ such that the numbers of two charts are related by

\[
V'^\mu=\frac{\partial x'^\mu}{\partial x^\nu}\,V^\nu .
\]

The index of a vector is written up. A vector field assigns a vector to every point.

\textbf{Covectors.} The gradient of a function $f$ has the components $W_\mu=\partial_\mu f$. By the chain rule, $\partial f/\partial x'^\mu=(\partial x^\nu/\partial x'^\mu)\,\partial f/\partial x^\nu$. A covector at $p$ is an assignment of $n$ numbers $W_\mu$ to every chart such that

\[
W'_\mu=\frac{\partial x^\nu}{\partial x'^\mu}\,W_\nu .
\]

The index of a covector is written down. Note that the Jacobian matrix is the inverse one: $\partial x/\partial x'$ instead of $\partial x'/\partial x$.

\textbf{The contraction of an upper with a lower index is the same in every chart.} For a vector $V$ and a covector $W$,

\[
V'^\mu W'_\mu=\frac{\partial x'^\mu}{\partial x^\alpha}\,\frac{\partial x^\beta}{\partial x'^\mu}\,V^\alpha W_\beta=\delta^\beta{}_\alpha V^\alpha W_\beta=V^\alpha W_\alpha ,
\]

where we used the Jacobian identity of Section 4.2. In particular the rate of change of $f$ along a curve, $df/d\lambda=u^\mu\partial_\mu f$, does not depend on the chart, as it must not.

\textbf{Tensors.} A tensor with $p$ upper and $q$ lower indices has components that transform with one factor $\partial x'/\partial x$ for each upper index and one factor $\partial x/\partial x'$ for each lower index; for example

\[
T'^\mu{}_\nu=\frac{\partial x'^\mu}{\partial x^\alpha}\,\frac{\partial x^\beta}{\partial x'^\nu}\,T^\alpha{}_\beta .
\]

Five rules follow directly from the definition. (i) Sums of tensors of the same type are tensors. (ii) Products of tensors are tensors (the factors simply multiply). (iii) Contracting one upper index with one lower index of a tensor gives a tensor; the proof is the computation just made for $V^\mu W_\mu$. (iv) The Kronecker delta is a tensor with the same components in every chart: $(\partial x'^\mu/\partial x^\alpha)(\partial x^\beta/\partial x'^\nu)\delta^\alpha{}_\beta=\delta^\mu{}_\nu$ by the Jacobian identity. (v) If a tensor equation $A=B$ holds in one chart, it holds in every chart, because both sides transform by the same linear rule; in particular a tensor that vanishes at a point in one chart vanishes there in every chart. Rule (v) is used again and again: to prove a tensor identity at a point it is enough to prove it in one well-chosen chart.

\textbf{A warning about index positions.} Summing an upper index with another upper index does not give a chart-independent result. The smallest example has one dimension: under the change $x'=2x$ two vectors $V$ and $W$ become $V'=2V$ and $W'=2W$, so the "sum" $V'W'=4VW$ depends on the chart. Section 4.14 shows that the notebook makes exactly this kind of error with the frame indices of the spin connection.

\textbf{Worked example: a constant vector in polar coordinates.} In the plane, take the vector field $V$ that points along the $x$-axis with unit length everywhere: Cartesian components $(V^x,V^y)=(1,0)$. With $r=\sqrt{x^2+y^2}$ and $\varphi$ the polar angle, $\partial r/\partial x=x/r=\cos\varphi$ and $\partial\varphi/\partial x=-y/r^2=-\sin\varphi/r$. The polar components are therefore

\[
V^r=\frac{\partial r}{\partial x}\cdot1+\frac{\partial r}{\partial y}\cdot0=\cos\varphi,\qquad V^\varphi=\frac{\partial\varphi}{\partial x}\cdot1+\frac{\partial\varphi}{\partial y}\cdot0=-\frac{\sin\varphi}{r}.
\]

At the point $(x,y)=(0,1)$, where $r=1$ and $\varphi=\pi/2$, this gives $V^r=0$ and $V^\varphi=-1$: moving along $+x$ from that point decreases the polar angle, as a sketch shows. The field is constant, yet its polar components change from point to point. The plain derivative $\partial_\mu V^\nu$ is therefore not a good measure of how a vector field changes. Section 4.5 repairs this.

\subsection{The metric}

\textbf{Lengths.} In the flat plane with Cartesian coordinates, Pythagoras gives the squared length of a small step: $ds^2=dx^2+dy^2$. In general coordinates, and in a curved space, the squared length of a small step $dx^\mu$ is a quadratic expression

\[
ds^2=g_{\mu\nu}(x)\,dx^\mu dx^\nu ,
\]

with coefficients $g_{\mu\nu}=g_{\nu\mu}$ that may depend on the point. The symmetric matrix $(g_{\mu\nu})$ is the metric. We require it to be invertible at every point, and we write $g^{\mu\nu}$ for the entries of its inverse: $g^{\mu\lambda}g_{\lambda\nu}=\delta^\mu{}_\nu$.

\textbf{The metric is a tensor with two lower indices.} The squared length of a step is a property of the step, not of the chart. With $dx^\alpha=(\partial x^\alpha/\partial x'^\mu)\,dx'^\mu$,

\[
g_{\alpha\beta}\,dx^\alpha dx^\beta=g_{\alpha\beta}\frac{\partial x^\alpha}{\partial x'^\mu}\frac{\partial x^\beta}{\partial x'^\nu}\,dx'^\mu dx'^\nu ,
\qquad\text{hence}\qquad
g'_{\mu\nu}=\frac{\partial x^\alpha}{\partial x'^\mu}\frac{\partial x^\beta}{\partial x'^\nu}\,g_{\alpha\beta}.
\]

(The coefficients of a quadratic form are fixed once we require them to be symmetric, so the two sides can be compared term by term.) This is the transformation law of a tensor with two lower indices. The inverse $g^{\mu\nu}$ is then a tensor with two upper indices: the inverse of the transformed matrix is the transformed inverse, as one checks by multiplying.

\textbf{Worked example: the plane in polar coordinates.} From $x=r\cos\varphi$, $y=r\sin\varphi$: $dx=\cos\varphi\,dr-r\sin\varphi\,d\varphi$ and $dy=\sin\varphi\,dr+r\cos\varphi\,d\varphi$. Squaring and adding, the cross terms $\mp2r\sin\varphi\cos\varphi\,dr\,d\varphi$ cancel and

\[
ds^2=dr^2+r^2d\varphi^2,\qquad g=\begin{pmatrix}1&0\\ 0&r^2\end{pmatrix},\qquad g^{-1}=\begin{pmatrix}1&0\\ 0&r^{-2}\end{pmatrix}.
\]

\textbf{Worked example: the sphere.} The point $(X,Y,Z)=(a\sin\theta\cos\varphi,\ a\sin\theta\sin\varphi,\ a\cos\theta)$ moves by $dX=a(\cos\theta\cos\varphi\,d\theta-\sin\theta\sin\varphi\,d\varphi)$, $dY=a(\cos\theta\sin\varphi\,d\theta+\sin\theta\cos\varphi\,d\varphi)$, $dZ=-a\sin\theta\,d\theta$. In $dX^2+dY^2$ the cross terms cancel, and $\cos^2\theta\,d\theta^2+\sin^2\theta\,d\theta^2=d\theta^2$ gives

\[
ds^2=a^2\bigl(d\theta^2+\sin^2\theta\,d\varphi^2\bigr).
\]

\textbf{Raising and lowering indices.} The metric turns a vector into a covector, $V_\mu:=g_{\mu\nu}V^\nu$, and the inverse metric turns a covector into a vector, $W^\mu:=g^{\mu\nu}W_\nu$. Doing both returns the original, because $g^{\mu\lambda}g_{\lambda\nu}=\delta^\mu{}_\nu$. The same rule raises or lowers any index of any tensor. The squared length of a vector is $g(V,V)=g_{\mu\nu}V^\mu V^\nu=V^\mu V_\mu$.

\textbf{Signature.} At a point, the symmetric matrix $(g_{\mu\nu})$ has real eigenvalues. The numbers of positive and of negative eigenvalues form the signature. A theorem of linear algebra, Sylvester's law of inertia, says that the signature does not depend on the chart; we quote it without proof. The flat space of the project has, in Cartesian coordinates, $g_{\mu\nu}=\eta_{\mu\nu}$: signature (4,4). A vector is called space-like if $g(V,V)>0$, time-like if $g(V,V)<0$ and null if $g(V,V)=0$ but $V\ne0$. In signature (4,4) there are four mutually perpendicular time-like directions, and null vectors are easy to write down: in flat space the vector with components $V^0=V^4=1$ and all others 0 has $g(V,V)=1-1=0$. The project chooses $x^4$ as the time in which everything evolves; the other three time directions are the extra times. Chapter 8 shows what this choice costs in the quantum theory.

\textbf{Worked example: the primordial field.} The notebook's metric MatrixMetric44 (Stage-2 document, §4) is diagonal:

\[
\begin{aligned}
ds^2={}&\cot^2z\,(dx^0)^2+s^{1/3}e^{2a_4}\bigl((dx^1)^2+(dx^2)^2+(dx^3)^2\bigr)-(dx^4)^2\\
&-s^{1/3}e^{-2a_4}\bigl((dx^5)^2+(dx^6)^2+(dx^7)^2\bigr),\qquad z=6Hx^0,\quad t=Hx^4,\quad s=\sin z,
\end{aligned}
\]

where $a_4(t)$ is an arbitrary smooth real function and a prime means $d/dt$, so that $\partial_4a_4=Ha_4'$. For $z$ in $(0,\pi/2)$ the diagonal entries have the signs $(+,+,+,+,-,-,-,-)$, so the signature is (4,4) (check \texttt{P\_metric\allowbreak{}\_signatu\allowbreak{}re44} in \texttt{wolfram-\allowbreak{}primordi\allowbreak{}al-repor\allowbreak{}t.json}). Three-space expands with the scale factor $s^{1/6}e^{a_4}$ when $a_4$ grows, and the three extra times contract with the scale factor $s^{1/6}e^{-a_4}$. Chapter 9 discusses this field in detail; here it serves as the running example.

\textbf{Worked example: a change of coordinates in the primordial field.} Put $\zeta=\ln(\sin z)/(6H)$, which runs over $(-\infty,0)$ when $z$ runs over $(0,\pi/2)$. Then $d\zeta=\frac{1}{6H}\,\frac{\cos z}{\sin z}\,dz=\cot z\,dx^0$, because $dz=6H\,dx^0$. So $\cot^2z\,(dx^0)^2=d\zeta^2$. Also $\sin z=e^{6H\zeta}$, so $s^{1/3}=e^{2H\zeta}$, and the metric becomes the warped form

\[
ds^2=d\zeta^2-(dx^4)^2+e^{2H\zeta}\Bigl[e^{2a_4}\bigl((dx^1)^2+(dx^2)^2+(dx^3)^2\bigr)-e^{-2a_4}\bigl((dx^5)^2+(dx^6)^2+(dx^7)^2\bigr)\Bigr]
\]

(Stage-2 document, §4.4; check \texttt{P\_zeta\_w\allowbreak{}arpedMet\allowbreak{}ric}). All six transverse directions share the warp factor $e^{2H\zeta}$.

\textbf{The volume element.} Take determinants in the transformation law of the metric: $\det g'=(\det J)^2\det g$ with $J=(\partial x^\alpha/\partial x'^\mu)$. Hence $\sqrt{\lvert\det g'\rvert}=\lvert\det J\rvert\sqrt{\lvert\det g\rvert}$. The change-of-variables rule of integral calculus says $d^nx=\lvert\det J\rvert\,d^nx'$ for the coordinate volume elements. Together,

\[
\sqrt{\lvert g'\rvert}\;d^nx'=\sqrt{\lvert g\rvert}\;d^nx,\qquad \lvert g\rvert:=\lvert\det(g_{\mu\nu})\rvert :
\]

$\sqrt{\lvert g\rvert}\,d^nx$ is the invariant volume element. In polar coordinates $\sqrt{\lvert g\rvert}=r$, giving the familiar area element $r\,dr\,d\varphi$. On the sphere $\sqrt{\lvert g\rvert}=a^2\sin\theta$, and the total area is $\int_0^\pi\int_0^{2\pi}a^2\sin\theta\,d\varphi\,d\theta=4\pi a^2$.

\textbf{Worked example: the determinant of the primordial field.} The determinant of a diagonal matrix is the product of its diagonal entries:

\[
\det g=\cot^2z\cdot\bigl(s^{1/3}e^{2a_4}\bigr)^3\cdot(-1)\cdot\bigl(-s^{1/3}e^{-2a_4}\bigr)^3=\cot^2z\cdot s^2\cdot(-1)(-1)^3=\cos^2z .
\]

The sign is $+$ because there are four negative entries. So $\sqrt{\lvert g\rvert}=\cos z$ on $(0,\pi/2)$ (checks \texttt{P\_metric\allowbreak{}\_detG\_eq\allowbreak{}uals\_plu\allowbreak{}s\_cos2z} and \texttt{P\_metric\allowbreak{}\_sqrtAbs\allowbreak{}DetG\_cos\allowbreak{}z}). The Stage-2 specification had written $\det g=-\cos^2z$; the exact computation corrected it. Since $\sqrt{\lvert g\rvert}$ does not depend on $x^4$, the growth of 3-space is exactly compensated by the shrinking of the extra times. To say this precisely, take a region with fixed coordinate ranges in the seven directions other than $x^4$ (a comoving region: its points keep their coordinates while $x^4$ runs, like the freely falling observers of Section 4.6). Because $g_{44}=-1$ and $g_{4\mu}=0$ for $\mu\ne4$, expanding $\det g$ along the row of $x^4$, whose only nonzero entry is $g_{44}=-1$, shows that $\det g$ is $-1$ times the determinant of the remaining 7 by 7 block of the metric (the block that measures lengths within a slice $x^4=$ constant). So the absolute value of that determinant is $\lvert g\rvert=\cos^2z$, and the invariant volume element of the block is $\sqrt{\cos^2z}\,d^7x=\cos z\,d^7x$. The 7-volume of the region is its integral, $\int\cos z\,d^7x$ over the seven coordinates other than $x^4$, and it does not change with $x^4$.

\subsection{The covariant derivative and the Christoffel symbols}

\textbf{The problem.} Differentiate the transformation law of a vector field, $V'^\mu=(\partial x'^\mu/\partial x^\alpha)V^\alpha$, with respect to $x'^\nu$, using the chain rule $\partial/\partial x'^\nu=(\partial x^\beta/\partial x'^\nu)\,\partial/\partial x^\beta$:

\[
\partial'_\nu V'^\mu=\frac{\partial x'^\mu}{\partial x^\alpha}\frac{\partial x^\beta}{\partial x'^\nu}\,\partial_\beta V^\alpha+\frac{\partial x^\beta}{\partial x'^\nu}\,\frac{\partial^2x'^\mu}{\partial x^\beta\partial x^\alpha}\,V^\alpha .
\]

The first term is the transformation law of a tensor; the second term, with second derivatives of the coordinate change, spoils it. So $\partial_\nu V^\mu$ is not a tensor, as the polar example of Section 4.3 showed.

\textbf{Definition (covariant derivative).} We add a correction that is linear in $V$:

\[
\nabla_\nu V^\mu:=\partial_\nu V^\mu+\Gamma^\mu{}_{\nu\lambda}V^\lambda .
\]

The $n^3$ numbers $\Gamma^\mu{}_{\nu\lambda}$ at each point are the connection coefficients. They are chosen so that $\nabla_\nu V^\mu$ is a tensor. Requiring $\nabla'_\nu V'^\mu=(\partial x'^\mu/\partial x^\alpha)(\partial x^\beta/\partial x'^\nu)\nabla_\beta V^\alpha$ for every $V$ and inserting the formula above gives the rule by which the $\Gamma$ must change between charts:

\[
\Gamma'^\mu{}_{\nu\lambda}=\frac{\partial x'^\mu}{\partial x^\alpha}\frac{\partial x^\beta}{\partial x'^\nu}\frac{\partial x^\sigma}{\partial x'^\lambda}\,\Gamma^\alpha{}_{\beta\sigma}-\frac{\partial x^\beta}{\partial x'^\nu}\frac{\partial x^\sigma}{\partial x'^\lambda}\,\frac{\partial^2x'^\mu}{\partial x^\beta\partial x^\sigma}.
\]

(To obtain it, write $\Gamma'^\mu{}_{\nu\lambda}V'^\lambda$ with $V'^\lambda=(\partial x'^\lambda/\partial x^\sigma)V^\sigma$, collect the coefficient of $V^\sigma$, and multiply by $\partial x^\sigma/\partial x'^\lambda$.) The last term is the inhomogeneous part; it is what cancels the unwanted term of $\partial'_\nu V'^\mu$. Because of it, the $\Gamma$ are not the components of a tensor. The inhomogeneous part is symmetric in $\nu$ and $\lambda$, because second partial derivatives do not depend on the order.

\textbf{Covectors and general tensors.} For a scalar function $f$ (a tensor without indices) we set $\nabla_\nu f=\partial_\nu f$. We require the product rule $\nabla_\nu(V^\mu W_\mu)=(\nabla_\nu V^\mu)W_\mu+V^\mu\nabla_\nu W_\mu$. The left side is $\partial_\nu(V^\mu W_\mu)=(\partial_\nu V^\mu)W_\mu+V^\mu\partial_\nu W_\mu$. Subtracting $(\nabla_\nu V^\mu)W_\mu=(\partial_\nu V^\mu)W_\mu+\Gamma^\mu{}_{\nu\lambda}V^\lambda W_\mu$ leaves, for every $V$,

\[
\nabla_\nu W_\mu=\partial_\nu W_\mu-\Gamma^\lambda{}_{\nu\mu}W_\lambda .
\]

A tensor with several indices gets one term $+\Gamma$ for each upper index and one term $-\Gamma$ for each lower index; for the metric, for example,

\[
\nabla_\rho g_{\mu\nu}=\partial_\rho g_{\mu\nu}-\Gamma^\lambda{}_{\rho\mu}g_{\lambda\nu}-\Gamma^\lambda{}_{\rho\nu}g_{\mu\lambda}.
\]

With this rule the product rule holds for all products and contractions.

\textbf{The two conditions of Levi-Civita.} A metric singles out one connection by two natural conditions.

\begin{enumerate}
\item No torsion: $\Gamma^\mu{}_{\nu\lambda}=\Gamma^\mu{}_{\lambda\nu}$. (The difference $\Gamma^\mu{}_{\nu\lambda}-\Gamma^\mu{}_{\lambda\nu}$ is a tensor, because the inhomogeneous part is symmetric, so this condition does not depend on the chart.)
\item Metric compatibility: $\nabla_\rho g_{\mu\nu}=0$. Lengths and angles are preserved when vectors are carried along.
\end{enumerate}

\textbf{Theorem (Christoffel formula).} Exactly one connection satisfies both conditions, namely

\[
\Gamma^\rho{}_{\mu\nu}=\tfrac12\,g^{\rho\sigma}\bigl(\partial_\mu g_{\nu\sigma}+\partial_\nu g_{\mu\sigma}-\partial_\sigma g_{\mu\nu}\bigr).
\]

These are the Christoffel symbols (of the Levi-Civita connection); the project uses only this connection (Stage 1, §5.1).

\emph{Proof.} Lower the upper index, $\Gamma_{\sigma\mu\nu}:=g_{\sigma\rho}\Gamma^\rho{}_{\mu\nu}$; condition 1 says $\Gamma_{\sigma\mu\nu}=\Gamma_{\sigma\nu\mu}$. Condition 2 reads $\partial_\rho g_{\mu\nu}=\Gamma_{\nu\rho\mu}+\Gamma_{\mu\rho\nu}$. Write it three times with the indices permuted cyclically:

\[
\partial_\mu g_{\nu\rho}=\Gamma_{\rho\mu\nu}+\Gamma_{\nu\mu\rho},\qquad
\partial_\nu g_{\mu\rho}=\Gamma_{\rho\nu\mu}+\Gamma_{\mu\nu\rho},\qquad
\partial_\rho g_{\mu\nu}=\Gamma_{\nu\rho\mu}+\Gamma_{\mu\rho\nu}.
\]

Add the first two and subtract the third. By the symmetry in the last two indices, $\Gamma_{\nu\mu\rho}=\Gamma_{\nu\rho\mu}$ and $\Gamma_{\mu\nu\rho}=\Gamma_{\mu\rho\nu}$ cancel against the third line, and $\Gamma_{\rho\nu\mu}=\Gamma_{\rho\mu\nu}$, so

\[
\partial_\mu g_{\nu\rho}+\partial_\nu g_{\mu\rho}-\partial_\rho g_{\mu\nu}=2\,\Gamma_{\rho\mu\nu}.
\]

Raising the index $\rho$ with $g^{\sigma\rho}$ gives the formula. So any connection with the two properties is this one (uniqueness). Conversely, the formula is symmetric in $\mu,\nu$, and adding $\Gamma_{\nu\rho\mu}+\Gamma_{\mu\rho\nu}$ computed from it gives back $\partial_\rho g_{\mu\nu}$, so it has both properties in the chart where it was computed. It remains to see that the formulas computed in two different charts describe the same $\nabla$. Take the connection computed in the chart $x$, and transform its coefficients to the chart $x'$ with the transformation rule above. The result defines the same operator $\nabla$ in the new chart, so it is still free of torsion and compatible with the metric (both conditions are tensor equations, rule (v) of Section 4.3). By uniqueness in the chart $x'$ it equals the Christoffel formula computed from $g'$. $\square$

\textbf{Lemma (derivative of a determinant).} For an invertible matrix $A(x)$ with entries depending on $x$, $\partial_\mu\ln\lvert\det A\rvert=\mathrm{tr}\bigl(A^{-1}\partial_\mu A\bigr)$.

\emph{Proof.} Write $O(\epsilon^2)$ for terms that are at most a constant times $\epsilon^2$ for small $\epsilon$. For a small number $\epsilon$ and a matrix $Y$, $\det(1+\epsilon Y)$ is a sum over permutations of products of entries of $1+\epsilon Y$. The identity permutation gives $\prod_i(1+\epsilon Y_{ii})=1+\epsilon\,\mathrm{tr}\,Y+O(\epsilon^2)$. Every other permutation moves at least two indices and therefore picks at least two off-diagonal entries, each of size $\epsilon$; it contributes $O(\epsilon^2)$. So $\det(1+\epsilon Y)=1+\epsilon\,\mathrm{tr}\,Y+O(\epsilon^2)$. Let $A_\epsilon$ be the value of $A$ at the point where the coordinate $x^\mu$ is increased by $\epsilon$. Then $A_\epsilon=A+\epsilon\,\partial_\mu A+O(\epsilon^2)=A\bigl(1+\epsilon A^{-1}\partial_\mu A+O(\epsilon^2)\bigr)$, and since the determinant of a product is the product of the determinants, $\det A_\epsilon=\det A\,\bigl(1+\epsilon\,\mathrm{tr}(A^{-1}\partial_\mu A)+O(\epsilon^2)\bigr)$. Hence $\ln\lvert\det A_\epsilon\rvert-\ln\lvert\det A\rvert=\epsilon\,\mathrm{tr}(A^{-1}\partial_\mu A)+O(\epsilon^2)$; divide by $\epsilon$ and let $\epsilon\to0$. $\square$

\textbf{Consequence: the contracted Christoffel symbol and the divergence.} Contract $\rho$ with $\mu$ in the Christoffel formula:

\[
\Gamma^\rho{}_{\rho\nu}=\tfrac12g^{\rho\sigma}\bigl(\partial_\rho g_{\nu\sigma}+\partial_\nu g_{\rho\sigma}-\partial_\sigma g_{\rho\nu}\bigr)=\tfrac12g^{\rho\sigma}\partial_\nu g_{\rho\sigma}=\partial_\nu\ln\sqrt{\lvert g\rvert},
\]

because $g^{\rho\sigma}\partial_\rho g_{\nu\sigma}$ and $g^{\rho\sigma}\partial_\sigma g_{\rho\nu}$ are equal (rename $\rho\leftrightarrow\sigma$ and use the symmetry of $g$), and $\tfrac12g^{\rho\sigma}\partial_\nu g_{\rho\sigma}=\tfrac12\mathrm{tr}(g^{-1}\partial_\nu g)=\tfrac12\partial_\nu\ln\lvert g\rvert$ by the lemma. Hence the divergence of a vector field is

\[
\nabla_\mu V^\mu=\partial_\mu V^\mu+\Gamma^\mu{}_{\mu\lambda}V^\lambda=\frac{1}{\sqrt{\lvert g\rvert}}\,\partial_\mu\bigl(\sqrt{\lvert g\rvert}\,V^\mu\bigr).
\]

\textbf{Diagonal metrics.} All metrics of the project's examples are diagonal: $g_{\mu\nu}=0$ for $\mu\ne\nu$. We write the diagonal entries as $g_{\mu\mu}=\eta_{\mu\mu}h_\mu^2$ (no sum), with positive functions $h_\mu$ (the scale factors) and the signs $\eta_{\mu\mu}=\pm1$. Then $g^{\mu\mu}=\eta_{\mu\mu}/h_\mu^2$. Insert this into the Christoffel formula; only terms with $\sigma=\rho$ survive, and there are four cases (in each, the repeated index $\mu$ is not summed):

\[
\begin{aligned}
&\text{(a)}\quad \Gamma^\mu{}_{\mu\mu}=\tfrac12g^{\mu\mu}\partial_\mu g_{\mu\mu}=\partial_\mu\ln h_\mu ,\\
&\text{(b)}\quad \Gamma^\mu{}_{\mu\nu}=\Gamma^\mu{}_{\nu\mu}=\tfrac12g^{\mu\mu}\partial_\nu g_{\mu\mu}=\partial_\nu\ln h_\mu\qquad(\nu\ne\mu),\\
&\text{(c)}\quad \Gamma^\mu{}_{\nu\nu}=-\tfrac12g^{\mu\mu}\partial_\mu g_{\nu\nu}=-\eta_{\mu\mu}\eta_{\nu\nu}\,\frac{h_\nu}{h_\mu^2}\,\partial_\mu h_\nu\qquad(\nu\ne\mu),\\
&\text{(d)}\quad \Gamma^\mu{}_{\nu\lambda}=0\qquad(\mu,\nu,\lambda\ \text{all different}).
\end{aligned}
\]

For (b): in $\tfrac12g^{\mu\mu}(\partial_\mu g_{\nu\mu}+\partial_\nu g_{\mu\mu}-\partial_\mu g_{\mu\nu})$ the first and last terms vanish because $g_{\nu\mu}=0$ for $\nu\ne\mu$. For (c): $\tfrac12g^{\mu\mu}(2\partial_\nu g_{\nu\mu}-\partial_\mu g_{\nu\nu})$, and $g_{\nu\mu}=0$. For (d): every metric entry that appears has two different indices.

\textbf{Worked example: polar coordinates.} With $h_r=1$, $h_\varphi=r$ and both signs $+1$: case (b) gives $\Gamma^\varphi{}_{r\varphi}=\Gamma^\varphi{}_{\varphi r}=\partial_r\ln r=1/r$; case (c) gives $\Gamma^r{}_{\varphi\varphi}=-(r/1)\,\partial_r r=-r$; all others vanish. Now return to the constant field $V^r=\cos\varphi$, $V^\varphi=-\sin\varphi/r$ of Section 4.3:

\[
\begin{aligned}
\nabla_rV^r&=\partial_r\cos\varphi=0,\\
\nabla_\varphi V^r&=\partial_\varphi\cos\varphi+\Gamma^r{}_{\varphi\varphi}V^\varphi=-\sin\varphi+(-r)\Bigl(-\frac{\sin\varphi}{r}\Bigr)=0,\\
\nabla_rV^\varphi&=\partial_r\Bigl(-\frac{\sin\varphi}{r}\Bigr)+\Gamma^\varphi{}_{r\varphi}V^\varphi=\frac{\sin\varphi}{r^2}+\frac1r\Bigl(-\frac{\sin\varphi}{r}\Bigr)=0,\\
\nabla_\varphi V^\varphi&=\partial_\varphi\Bigl(-\frac{\sin\varphi}{r}\Bigr)+\Gamma^\varphi{}_{\varphi r}V^r=-\frac{\cos\varphi}{r}+\frac1r\cos\varphi=0 .
\end{aligned}
\]

The covariant derivative of the constant field vanishes, as it should: the Christoffel terms exactly remove the change of the polar components that comes from the turning of the coordinate lines.

\textbf{Worked example: the 37 Christoffel symbols of the primordial field.} The scale factors are $h=(\cot z,\ s^{1/6}e^{a_4}\ (\times3),\ 1,\ s^{1/6}e^{-a_4}\ (\times3))$ with the signs of $\eta$, and they depend on $x^0$ and $x^4$ only. With $\partial_0=6H\,d/dz$ and $\partial_4=H\,d/dt$ the logarithmic derivatives are

\[
\begin{aligned}
&\partial_0\ln\cot z=-\frac{6H}{\sin z\cos z},\qquad \partial_0\ln h_k=H\cot z,\\
&\partial_4\ln h_i=Ha_4',\qquad \partial_4\ln h_j=-Ha_4',
\end{aligned}
\]

for $k\in\{1,2,3,5,6,7\}$, $i\in\{1,2,3\}$ and $j\in\{5,6,7\}$. (For the first: $\frac{d}{dz}\ln\cot z=-\frac{1}{\sin^2z}\cdot\frac{\sin z}{\cos z}$. For the second: $\ln h_k=\tfrac16\ln\sin z\pm a_4$, and $\frac{d}{dz}\tfrac16\ln\sin z=\tfrac16\cot z$.) Now apply the four cases. Case (a) gives only $\Gamma^0{}_{00}$, since $h_0$ is the only scale factor that depends on its own coordinate. Case (b) gives $\Gamma^k{}_{k0}=\Gamma^k{}_{0k}=H\cot z$ and $\Gamma^k{}_{k4}=\Gamma^k{}_{4k}=\pm Ha_4'$. Case (c) with $\mu=0$ gives $\Gamma^0{}_{kk}=-\eta_{kk}\,h_k^2\,\tan^2z\,\partial_0\ln h_k=-\eta_{kk}H\tan z\,h_k^2$, and with $\mu=4$ it gives $\Gamma^4{}_{kk}=\eta_{kk}h_k^2\,\partial_4\ln h_k$. Nothing depends on $x^1,x^2,x^3,x^5,x^6,x^7$, so case (c) with those values of $\mu$ gives zero. The complete list:

\begingroup
\small
\begin{longtable}{@{}>{\raggedright\arraybackslash}p{0.267\linewidth}>{\raggedright\arraybackslash}p{0.267\linewidth}>{\raggedright\arraybackslash}p{0.267\linewidth}@{}}
\toprule
\textbf{Symbol} & \textbf{Value} & \textbf{Number} \\
\midrule
\endfirsthead
\toprule
\textbf{Symbol} & \textbf{Value} & \textbf{Number} \\
\midrule
\endhead
$\Gamma^0{}_{00}$ & $-6H/(\sin z\cos z)$ & 1 \\
$\Gamma^k{}_{0k}=\Gamma^k{}_{k0}$, $k\in\{1,2,3,5,6,7\}$ & $H\cot z$ & 12 \\
$\Gamma^i{}_{4i}=\Gamma^i{}_{i4}$, $i\in\{1,2,3\}$ & $Ha_4'$ & 6 \\
$\Gamma^j{}_{4j}=\Gamma^j{}_{j4}$, $j\in\{5,6,7\}$ & $-Ha_4'$ & 6 \\
$\Gamma^0{}_{ii}$ & $-H\tan z\,s^{1/3}e^{2a_4}$ & 3 \\
$\Gamma^0{}_{jj}$ & $H\tan z\,s^{1/3}e^{-2a_4}$ & 3 \\
$\Gamma^4{}_{ii}$ & $Ha_4'\,s^{1/3}e^{2a_4}$ & 3 \\
$\Gamma^4{}_{jj}$ & $Ha_4'\,s^{1/3}e^{-2a_4}$ & 3 \\
\bottomrule
\end{longtable}
\endgroup

In total $1+12+6+6+3+3+3+3=37$ of the $8^3=512$ symbols are nonzero. This is exactly the count and the closed forms that the Stage-2 verifiers found (measurement \texttt{christof\allowbreak{}felNonze\allowbreak{}roCount} = 37 and checks \texttt{P\_christ\allowbreak{}offel\_co\allowbreak{}unt} and \texttt{P\_christ\allowbreak{}offel\_cl\allowbreak{}osedForm\allowbreak{}s512} in \texttt{wolfram-\allowbreak{}primordi\allowbreak{}al-repor\allowbreak{}t.json}; \texttt{nonzeroO\allowbreak{}rderedCo\allowbreak{}unt} = 37 and check \texttt{P\_christ\allowbreak{}offel} in \texttt{python-p\allowbreak{}rimordia\allowbreak{}l-report\allowbreak{}.json}; the table of the Stage-2 document, §5, lists the same values with $\sec z/\sin z$ written for $1/(\sin z\cos z)$). As a check of the contracted formula $\Gamma^\rho{}_{\rho\nu}=\partial_\nu\ln\sqrt{\lvert g\rvert}$ with $\sqrt{\lvert g\rvert}=\cos z$:

\[
\sum_\rho\Gamma^\rho{}_{\rho0}=-\frac{6H}{\sin z\cos z}+6H\frac{\cos z}{\sin z}=\frac{6H(\cos^2z-1)}{\sin z\cos z}=-6H\tan z=\partial_0\ln\cos z,
\]

and $\sum_\rho\Gamma^\rho{}_{\rho4}=3Ha_4'-3Ha_4'=0=\partial_4\ln\cos z$.

\subsection{Geodesics, with one worked example}

\textbf{Derivative along a curve.} Let $x^\mu(\lambda)$ be a curve with velocity $u^\mu=dx^\mu/d\lambda$, and let $V^\mu(\lambda)$ be a vector given at each point of the curve. If $V$ is the restriction of a vector field, the chain rule gives $dV^\mu/d\lambda=u^\alpha\partial_\alpha V^\mu$, and so

\[
u^\alpha\nabla_\alpha V^\mu=\frac{dV^\mu}{d\lambda}+\Gamma^\mu{}_{\alpha\beta}\,u^\alpha V^\beta=:\frac{DV^\mu}{d\lambda}.
\]

The right-hand side uses only the values of $V$ on the curve, so it defines the covariant derivative along the curve for any $V(\lambda)$. $V$ is parallel transported along the curve if $DV^\mu/d\lambda=0$.

\textbf{Definition (geodesic).} A geodesic is a curve whose velocity is parallel transported along itself:

\[
\frac{d^2x^\mu}{d\lambda^2}+\Gamma^\mu{}_{\alpha\beta}\,\frac{dx^\alpha}{d\lambda}\frac{dx^\beta}{d\lambda}=0 .
\]

In flat space with Cartesian coordinates all $\Gamma$ vanish, and the equation says $d^2x^\mu/d\lambda^2=0$: straight lines traversed at constant speed. A geodesic is the curved-space version of a straight line. (Geodesics can also be characterized as curves of stationary length; that needs the calculus of variations of Chapter 5, and we do not use it.) In general relativity, a body on which no force other than gravity acts moves on a time-like geodesic.

\textbf{The squared speed is constant along a geodesic.} Differentiate $g_{\mu\nu}u^\mu u^\nu$ along the curve and use the geodesic equation:

\[
\frac{d}{d\lambda}\bigl(g_{\mu\nu}u^\mu u^\nu\bigr)=u^\rho u^\mu u^\nu\,\partial_\rho g_{\mu\nu}+2g_{\mu\nu}u^\mu\frac{du^\nu}{d\lambda}=u^\rho u^\mu u^\nu\,\partial_\rho g_{\mu\nu}-2\Gamma_{\mu\alpha\beta}\,u^\mu u^\alpha u^\beta ,
\]

with $\Gamma_{\mu\alpha\beta}=g_{\mu\nu}\Gamma^\nu{}_{\alpha\beta}$. By the proof of the Christoffel formula, $2\Gamma_{\mu\alpha\beta}=\partial_\alpha g_{\beta\mu}+\partial_\beta g_{\alpha\mu}-\partial_\mu g_{\alpha\beta}$. Contracted with the symmetric product $u^\mu u^\alpha u^\beta$, each of the three terms equals $u^\rho u^\mu u^\nu\partial_\rho g_{\mu\nu}$ after renaming the summed indices, with the signs $+,+,-$. So $2\Gamma_{\mu\alpha\beta}u^\mu u^\alpha u^\beta=u^\rho u^\mu u^\nu\partial_\rho g_{\mu\nu}$, and the derivative is zero. A geodesic that starts time-like stays time-like, and so on.

\textbf{Worked example: the straight lines of the plane in polar coordinates.} With the polar Christoffel symbols ($\Gamma^r{}_{\varphi\varphi}=-r$, $\Gamma^\varphi{}_{r\varphi}=\Gamma^\varphi{}_{\varphi r}=1/r$) the geodesic equations are, with a prime meaning $d/d\lambda$ in this example,

\[
r''-r\,\varphi'^2=0,\qquad \varphi''+\frac2r\,r'\varphi'=0 .
\]

The factor 2 in the second equation comes from the two equal terms $\Gamma^\varphi{}_{r\varphi}r'\varphi'$ and $\Gamma^\varphi{}_{\varphi r}\varphi'r'$. We solve them completely.

\emph{Step 1: a conserved quantity.} $\frac{d}{d\lambda}(r^2\varphi')=2rr'\varphi'+r^2\varphi''=r^2\bigl(\varphi''+\tfrac2rr'\varphi'\bigr)=0$. So $r^2\varphi'=L$ is constant (in mechanics, the angular momentum).

\emph{Step 2: the speed.} By the result just proved, $g(u,u)=r'^2+r^2\varphi'^2$ is constant; choose the parameter so that it equals 1 (unit speed, $\lambda$ is then the length along the curve). With $\varphi'=L/r^2$ this reads $r'^2=1-L^2/r^2$.

\emph{Step 3: solve for $r$.} Take $L>0$. (If $L=0$, then $\varphi'=0$, so $\varphi$ is constant, and Step 2 gives $r'^2=1$, in agreement with the first equation, which becomes $r''=0$. So $r=\lambda-\lambda_0$ for $\lambda>\lambda_0$, or $r=\lambda_0-\lambda$ for $\lambda<\lambda_0$: in both cases $r=\lvert\lambda-\lambda_0\rvert$, a radial half-line, the part of a straight line through the origin that the polar chart sees, since the chart excludes $r=0$. If $L<0$, replace $\varphi$ by $-\varphi$ (mod $2\pi$). This leaves the two geodesic equations unchanged, because each term of the second contains exactly one factor $\varphi'$ or $\varphi''$ and the first contains $\varphi'^2$, and it flips the sign of $L$. So the solutions with $L<0$ are the mirror images of those found below: Steps 3 to 5 hold with $\lvert L\rvert$ in place of $L$ and $\varphi$ replaced by $-\varphi$, which gives the straight lines $r\cos(\varphi-\varphi_0)=\lvert L\rvert$ traversed in the opposite sense.) Since $r'^2\ge0$, $r\ge L$. Where $r'>0$, $\frac{d}{d\lambda}\sqrt{r^2-L^2}=\frac{rr'}{\sqrt{r^2-L^2}}=\frac{r}{\sqrt{r^2-L^2}}\sqrt{1-\frac{L^2}{r^2}}=1$. So $\sqrt{r^2-L^2}=\lambda-\lambda_0$ for $\lambda>\lambda_0$, and with the other sign of $r'$ for $\lambda<\lambda_0$: in both cases

\[
r(\lambda)^2=L^2+(\lambda-\lambda_0)^2 .
\]

\emph{Step 4: solve for $\varphi$.} $\varphi'=L/r^2=L/(L^2+(\lambda-\lambda_0)^2)$, whose antiderivative is $\arctan((\lambda-\lambda_0)/L)$. So $\varphi(\lambda)=\varphi_0+\arctan((\lambda-\lambda_0)/L)$.

\emph{Step 5: recognize the curve.} From step 4, $\cos(\varphi-\varphi_0)=L/\sqrt{L^2+(\lambda-\lambda_0)^2}=L/r$, so

\[
r\cos(\varphi-\varphi_0)=L ,
\]

which is the straight line at distance $L$ from the origin, perpendicular to the direction $\varphi_0$. The geodesics of the plane are its straight lines, now derived in coordinates in which they look nothing like straight lines.

\emph{A check with small numbers.} Take $L=1$, $\varphi_0=0$, $\lambda_0=0$: $r=\sqrt{1+\lambda^2}$ and $\varphi=\arctan\lambda$, so $x=r\cos\varphi=1$ and $y=r\sin\varphi=\lambda$: the vertical line $x=1$, traversed at unit speed. At $\lambda=1$: $r=\sqrt2$, $\varphi=\pi/4$, $r'=\lambda/\sqrt{1+\lambda^2}=1/\sqrt2$, $\varphi'=1/(1+\lambda^2)=1/2$, $r''=(1+\lambda^2)^{-3/2}=1/(2\sqrt2)$ and $\varphi''=-2\lambda/(1+\lambda^2)^2=-1/2$. Then $r''-r\varphi'^2=\frac{1}{2\sqrt2}-\frac{\sqrt2}{4}=0$, $\varphi''+\frac2rr'\varphi'=-\frac12+\frac{2}{\sqrt2}\cdot\frac{1}{\sqrt2}\cdot\frac12=0$, and $r'^2+r^2\varphi'^2=\frac12+2\cdot\frac14=1$.

\textbf{Geodesics in the primordial field.} Exercise 4.3 shows that the curves on which only $x^4$ changes, $x^\mu(\lambda)=(c^0,c^1,c^2,c^3,\lambda,c^5,c^6,c^7)$ with constants $c$, are time-like geodesics for every function $a_4$. The proper time along a time-like curve is $\tau=\int\sqrt{-g_{\mu\nu}u^\mu u^\nu}\,d\lambda$, the time shown by a clock carried along it (the time-like analogue of the length of Exercise 4.1(b)). For these curves $u^\mu=\delta^\mu{}_4$ and $g_{44}=-1$, so $d\tau/d\lambda=1$, and the proper time counted from $x^4=0$ is $\tau=\lambda=x^4$. Observers at fixed $x^0,\dots,x^3,x^5,\dots,x^7$ are therefore freely falling, and $x^4$ is their proper time. A chart with this property ($g_{44}=-1$ and $g_{4\mu}=0$ for $\mu\ne4$) is called Gaussian normal; the project uses it whenever it speaks of the energy density $\rho=T_{44}$ (Chapter 7).

\subsection{Curvature}

\textbf{The idea.} On a flat sheet of paper, a vector carried around a closed loop by parallel transport comes back unchanged. On a sphere it does not: carry a vector from the north pole down to the equator, a quarter of the way along the equator, and back up to the pole, keeping it parallel all the time, and it returns turned by a right angle. Curvature measures this failure. For a very small loop it is measured by the commutator of two covariant derivatives: in flat Cartesian coordinates $\nabla_\mu\nabla_\nu=\partial_\mu\partial_\nu$, and partial derivatives commute; in a curved space covariant derivatives do not.

\textbf{Theorem (the Riemann tensor).} For every vector field $V$,

\[
\nabla_\mu\nabla_\nu V^\rho-\nabla_\nu\nabla_\mu V^\rho=R^\rho{}_{\sigma\mu\nu}V^\sigma ,\qquad
R^\rho{}_{\sigma\mu\nu}=\partial_\mu\Gamma^\rho{}_{\nu\sigma}-\partial_\nu\Gamma^\rho{}_{\mu\sigma}+\Gamma^\rho{}_{\mu\lambda}\Gamma^\lambda{}_{\nu\sigma}-\Gamma^\rho{}_{\nu\lambda}\Gamma^\lambda{}_{\mu\sigma}.
\]

This is the convention of the whole project (Stage 1, §5.5; measurement \texttt{conventi\allowbreak{}on.curva\allowbreak{}ture} of \texttt{wolfram-\allowbreak{}geometry\allowbreak{}-report.\allowbreak{}json}).

\emph{Proof.} $\nabla_\nu V^\rho$ has one lower index $\nu$ and one upper index $\rho$, so by the rule of Section 4.5

\[
\nabla_\mu\nabla_\nu V^\rho=\partial_\mu\bigl(\partial_\nu V^\rho+\Gamma^\rho{}_{\nu\sigma}V^\sigma\bigr)+\Gamma^\rho{}_{\mu\lambda}\bigl(\partial_\nu V^\lambda+\Gamma^\lambda{}_{\nu\sigma}V^\sigma\bigr)-\Gamma^\lambda{}_{\mu\nu}\nabla_\lambda V^\rho .
\]

Expand and subtract the same expression with $\mu$ and $\nu$ exchanged. Three kinds of terms drop out because they are symmetric in $\mu,\nu$: $\partial_\mu\partial_\nu V^\rho$ (partial derivatives commute); $\Gamma^\rho{}_{\nu\sigma}\partial_\mu V^\sigma+\Gamma^\rho{}_{\mu\sigma}\partial_\nu V^\sigma$ (after renaming $\lambda$ to $\sigma$); and $\Gamma^\lambda{}_{\mu\nu}\nabla_\lambda V^\rho$ (no torsion). What remains is $(\partial_\mu\Gamma^\rho{}_{\nu\sigma}-\partial_\nu\Gamma^\rho{}_{\mu\sigma}+\Gamma^\rho{}_{\mu\lambda}\Gamma^\lambda{}_{\nu\sigma}-\Gamma^\rho{}_{\nu\lambda}\Gamma^\lambda{}_{\mu\sigma})V^\sigma$. $\square$

The left-hand side is a tensor for every $V$ and depends linearly on $V$ with no derivatives of $V$; therefore $R^\rho{}_{\sigma\mu\nu}$ transforms as a tensor with one upper and three lower indices (write the transformation law of both sides and compare the coefficients of $V^\sigma$).

\textbf{Other tensors.} For a function $f$, $\nabla_\mu\nabla_\nu f-\nabla_\nu\nabla_\mu f=-(\Gamma^\lambda{}_{\mu\nu}-\Gamma^\lambda{}_{\nu\mu})\partial_\lambda f=0$. For a covector, apply the commutator to the function $V^\rho W_\rho$. By the product rule, $\nabla_\mu\nabla_\nu(V^\rho W_\rho)$ has four terms; the two mixed ones, $(\nabla_\mu V^\rho)(\nabla_\nu W_\rho)+(\nabla_\nu V^\rho)(\nabla_\mu W_\rho)$, are symmetric in $\mu,\nu$ and cancel in the commutator. So $0=(R^\rho{}_{\sigma\mu\nu}V^\sigma)W_\rho+V^\rho[\nabla_\mu,\nabla_\nu]W_\rho$ for all $V$, that is

\[
[\nabla_\mu,\nabla_\nu]W_\sigma=-R^\rho{}_{\sigma\mu\nu}W_\rho .
\]

Every tensor is a sum of products of vectors and covectors (expand it in a basis), so for a general tensor the commutator gives one term $+R$ for each upper index and one term $-R$ for each lower index.

\textbf{Symmetries.} Lower the first index, $R_{\rho\sigma\mu\nu}:=g_{\rho\lambda}R^\lambda{}_{\sigma\mu\nu}$.

\begin{itemize}
\item (S1) $R_{\rho\sigma\mu\nu}=-R_{\rho\sigma\nu\mu}$. This is visible in the definition.
\item (S2) $R_{\rho\sigma\mu\nu}=-R_{\sigma\rho\mu\nu}$. Proof: since $\nabla g=0$, $0=[\nabla_\mu,\nabla_\nu]g_{\rho\sigma}=-R^\lambda{}_{\rho\mu\nu}g_{\lambda\sigma}-R^\lambda{}_{\sigma\mu\nu}g_{\rho\lambda}=-R_{\sigma\rho\mu\nu}-R_{\rho\sigma\mu\nu}$.
\item (S3) First Bianchi identity: $R^\rho{}_{\sigma\mu\nu}+R^\rho{}_{\mu\nu\sigma}+R^\rho{}_{\nu\sigma\mu}=0$. Proof: write the three terms with the formula. The six derivative terms cancel in pairs, for example $\partial_\mu\Gamma^\rho{}_{\nu\sigma}$ from the first term against $-\partial_\mu\Gamma^\rho{}_{\sigma\nu}$ from the third, because $\Gamma$ is symmetric in its lower indices; the six products of two $\Gamma$ cancel in pairs in the same way.
\item (S4) Pair symmetry: $R_{\alpha\nu\beta\mu}=R_{\beta\mu\alpha\nu}$. Proof: write (S3) with the first index lowered four times,
\end{itemize}

\[
\begin{aligned}
&R_{\alpha\beta\mu\nu}+R_{\alpha\mu\nu\beta}+R_{\alpha\nu\beta\mu}=0, &\qquad &R_{\beta\mu\nu\alpha}+R_{\beta\nu\alpha\mu}+R_{\beta\alpha\mu\nu}=0,\\
&R_{\mu\nu\alpha\beta}+R_{\mu\alpha\beta\nu}+R_{\mu\beta\nu\alpha}=0, &\qquad &R_{\nu\alpha\beta\mu}+R_{\nu\beta\mu\alpha}+R_{\nu\mu\alpha\beta}=0 .
\end{aligned}
\]

Add the first two: by (S2) $R_{\alpha\beta\mu\nu}+R_{\beta\alpha\mu\nu}=0$, so $R_{\alpha\mu\nu\beta}+R_{\alpha\nu\beta\mu}+R_{\beta\mu\nu\alpha}+R_{\beta\nu\alpha\mu}=0$. Add the last two: $R_{\mu\nu\alpha\beta}+R_{\nu\mu\alpha\beta}=0$, so $R_{\mu\alpha\beta\nu}+R_{\mu\beta\nu\alpha}+R_{\nu\alpha\beta\mu}+R_{\nu\beta\mu\alpha}=0$. With (S1) and (S2), $R_{\mu\alpha\beta\nu}=R_{\alpha\mu\nu\beta}$, $R_{\mu\beta\nu\alpha}=-R_{\beta\mu\nu\alpha}$, $R_{\nu\alpha\beta\mu}=-R_{\alpha\nu\beta\mu}$ and $R_{\nu\beta\mu\alpha}=R_{\beta\nu\alpha\mu}$, so the second sum reads $R_{\alpha\mu\nu\beta}-R_{\beta\mu\nu\alpha}-R_{\alpha\nu\beta\mu}+R_{\beta\nu\alpha\mu}=0$. Subtracting it from the first sum leaves $2R_{\alpha\nu\beta\mu}+2R_{\beta\mu\nu\alpha}=0$, and $-R_{\beta\mu\nu\alpha}=R_{\beta\mu\alpha\nu}$ by (S1).

\textbf{Ricci tensor and scalar curvature.} The Ricci tensor and the scalar curvature are the contractions

\[
R_{\sigma\nu}:=R^\rho{}_{\sigma\rho\nu},\qquad R:=g^{\sigma\nu}R_{\sigma\nu}.
\]

The Ricci tensor is symmetric: $R_{\sigma\nu}=g^{\rho\alpha}R_{\alpha\sigma\rho\nu}=g^{\rho\alpha}R_{\rho\nu\alpha\sigma}=R^\alpha{}_{\nu\alpha\sigma}=R_{\nu\sigma}$, by (S4) in the middle step.

\textbf{Worked example: the sphere.} With $h_\theta=a$ and $h_\varphi=a\sin\theta$ the diagonal formulas give $\Gamma^\theta{}_{\varphi\varphi}=-\frac{a\sin\theta}{a^2}\,a\cos\theta=-\sin\theta\cos\theta$ and $\Gamma^\varphi{}_{\theta\varphi}=\Gamma^\varphi{}_{\varphi\theta}=\partial_\theta\ln(a\sin\theta)=\cot\theta$; all others vanish. Then

\[
\begin{aligned}
R^\theta{}_{\varphi\theta\varphi}&=\partial_\theta\Gamma^\theta{}_{\varphi\varphi}-\partial_\varphi\Gamma^\theta{}_{\theta\varphi}+\Gamma^\theta{}_{\theta\lambda}\Gamma^\lambda{}_{\varphi\varphi}-\Gamma^\theta{}_{\varphi\lambda}\Gamma^\lambda{}_{\theta\varphi}\\
&=(\sin^2\theta-\cos^2\theta)-0+0-(-\sin\theta\cos\theta)\cot\theta=\sin^2\theta ,\\
R^\varphi{}_{\theta\varphi\theta}&=\partial_\varphi\Gamma^\varphi{}_{\theta\theta}-\partial_\theta\Gamma^\varphi{}_{\varphi\theta}+\Gamma^\varphi{}_{\varphi\lambda}\Gamma^\lambda{}_{\theta\theta}-\Gamma^\varphi{}_{\theta\lambda}\Gamma^\lambda{}_{\varphi\theta}\\
&=0+\frac{1}{\sin^2\theta}+0-\cot^2\theta=1 .
\end{aligned}
\]

So $R_{\varphi\varphi}=R^\theta{}_{\varphi\theta\varphi}=\sin^2\theta$, $R_{\theta\theta}=R^\varphi{}_{\theta\varphi\theta}=1$, $R_{\theta\varphi}=0$, and

\[
R=g^{\theta\theta}R_{\theta\theta}+g^{\varphi\varphi}R_{\varphi\varphi}=\frac1{a^2}+\frac{\sin^2\theta}{a^2\sin^2\theta}=\frac{2}{a^2}>0 .
\]

A small sphere is strongly curved, a large one weakly. For the plane in polar coordinates the same computation gives $R^r{}_{\varphi r\varphi}=\partial_r(-r)-\Gamma^r{}_{\varphi\varphi}\Gamma^\varphi{}_{r\varphi}=-1-(-r)\frac1r=0$: the plane is flat, whatever coordinates we use.

\textbf{Flat spaces.} If there is a chart in which $g_{\mu\nu}$ is constant, all $\Gamma$ vanish in it, so the Riemann tensor vanishes there and hence (rule (v)) in every chart. The converse is also true: if the Riemann tensor vanishes on a region, then near each point there is a chart in which $g_{\mu\nu}=\eta_{\mu\nu}$. We quote this standard theorem without proof; the book uses it only in Section 4.13.

\textbf{Worked example: $R_{44}$ of the primordial field.} From the table of Section 4.5, no Christoffel symbol has the lower index pair $(4,4)$, so $\Gamma^\lambda{}_{44}=0$ for every $\lambda$, and $\sum_\rho\Gamma^\rho{}_{\rho4}=0$. Hence

\[
R_{44}=R^\rho{}_{4\rho4}=\partial_\rho\Gamma^\rho{}_{44}-\partial_4\Gamma^\rho{}_{\rho4}+\Gamma^\rho{}_{\rho\lambda}\Gamma^\lambda{}_{44}-\Gamma^\rho{}_{4\lambda}\Gamma^\lambda{}_{\rho4}=-\sum_{\rho,\lambda}\Gamma^\rho{}_{4\lambda}\Gamma^\lambda{}_{\rho4}.
\]

The only nonzero $\Gamma^\rho{}_{4\lambda}$ are $\Gamma^k{}_{4k}=\pm Ha_4'$ for $k\in\{1,2,3,5,6,7\}$, so the double sum is $\sum_k(\Gamma^k{}_{4k})^2=6H^2a_4'^2$ and

\[
R_{44}=-6H^2a_4'^2 ,
\]

in agreement with the Stage-2 verifier (check \texttt{P\_einste\allowbreak{}in\_R44}; Stage-2 document, §15.1). The same method, applied to all components, gives the scalar curvature

\[
R=6H^2\bigl(a_4'^2-7\bigr)
\]

(check \texttt{P\_einste\allowbreak{}in\_ricci\allowbreak{}Scalar} in \texttt{wolfram-\allowbreak{}primordi\allowbreak{}al-repor\allowbreak{}t.json}; measurement \texttt{P\_einste\allowbreak{}in.ricci\allowbreak{}Scalar} in \texttt{python-p\allowbreak{}rimordia\allowbreak{}l-report\allowbreak{}.json}; it also equals the notebook's own stored output of cell 583, check \texttt{P\_einste\allowbreak{}in\_noteb\allowbreak{}ookCell5\allowbreak{}83}). Section 4.15 gives a short program that reproduces it. \emph{Numbers.} At the first exact test point of Stage 1, $H=2/3$ and $a_4'=3/7$, so $R=6\cdot\frac49\cdot\bigl(\frac{9}{49}-7\bigr)=\frac83\cdot\frac{-334}{49}=-\frac{2672}{147}$, the value recorded as \texttt{G2.p1.sc\allowbreak{}alarCurv\allowbreak{}ature} in \texttt{wolfram-\allowbreak{}geometry\allowbreak{}-report.\allowbreak{}json}. The Riemann tensor of the primordial field vanishes nowhere: where $a_4'\ne0$ we have $R_{44}\ne0$, and where $a_4'=0$ we have $R=-42H^2\ne0$. The primordial field is curved everywhere.

\textbf{The second Bianchi identity.} For every metric,

\[
\nabla_\lambda R^\rho{}_{\sigma\mu\nu}+\nabla_\mu R^\rho{}_{\sigma\nu\lambda}+\nabla_\nu R^\rho{}_{\sigma\lambda\mu}=0 .
\]

\emph{Proof.} First we show that around any point $p$ there is a chart in which all $\Gamma^\rho{}_{\mu\nu}(p)=0$. Start from any chart with $x(p)=0$ and define new coordinates $x'^\rho=x^\rho+\tfrac12\Gamma^\rho{}_{\alpha\beta}(p)\,x^\alpha x^\beta$. At $p$, $\partial x'^\rho/\partial x^\alpha=\delta^\rho{}_\alpha$ (so near $p$ this is an allowed change of coordinates, by the inverse function theorem of calculus, which we quote) and $\partial^2x'^\rho/\partial x^\alpha\partial x^\beta=\Gamma^\rho{}_{\alpha\beta}(p)$, because $\Gamma$ is symmetric. The transformation rule of Section 4.5 then gives $\Gamma'^\rho{}_{\mu\nu}(p)=\Gamma^\rho{}_{\mu\nu}(p)-\Gamma^\rho{}_{\mu\nu}(p)=0$. In this chart, at $p$, a covariant derivative is a plain derivative, and the derivative of a product of two $\Gamma$ vanishes at $p$ because each term contains an undifferentiated $\Gamma(p)=0$. So at $p$

\[
\nabla_\lambda R^\rho{}_{\sigma\mu\nu}=\partial_\lambda\partial_\mu\Gamma^\rho{}_{\nu\sigma}-\partial_\lambda\partial_\nu\Gamma^\rho{}_{\mu\sigma}.
\]

Add the three cyclic permutations of $(\lambda,\mu,\nu)$: the six terms are $\partial_\lambda\partial_\mu\Gamma^\rho{}_{\nu\sigma}-\partial_\lambda\partial_\nu\Gamma^\rho{}_{\mu\sigma}+\partial_\mu\partial_\nu\Gamma^\rho{}_{\lambda\sigma}-\partial_\mu\partial_\lambda\Gamma^\rho{}_{\nu\sigma}+\partial_\nu\partial_\lambda\Gamma^\rho{}_{\mu\sigma}-\partial_\nu\partial_\mu\Gamma^\rho{}_{\lambda\sigma}$, and they cancel in pairs because partial derivatives commute. The left-hand side of the identity is a tensor that vanishes at $p$ in one chart, so it vanishes at $p$ in every chart (rule (v)), and $p$ was arbitrary. $\square$

\subsection{The Einstein equations}

\textbf{The contracted Bianchi identity.} First a small fact: the inverse metric is also covariantly constant. Differentiate $g^{\mu\lambda}g_{\lambda\nu}=\delta^\mu{}_\nu$ with the product rule; since $\nabla\delta=0$ and $\nabla g=0$, we get $(\nabla_\rho g^{\mu\lambda})g_{\lambda\nu}=0$, and multiplying by $g^{\nu\sigma}$ gives $\nabla_\rho g^{\mu\sigma}=0$. So $g$ and $g^{-1}$ can be moved through $\nabla$ freely. Now contract the second Bianchi identity over $\rho$ and $\mu$:

\[
\nabla_\lambda R^\rho{}_{\sigma\rho\nu}+\nabla_\rho R^\rho{}_{\sigma\nu\lambda}+\nabla_\nu R^\rho{}_{\sigma\lambda\rho}=0 .
\]

The first term is $\nabla_\lambda R_{\sigma\nu}$; by (S1) the third is $-\nabla_\nu R_{\sigma\lambda}$. Multiply by $g^{\sigma\nu}$. The first term becomes $\nabla_\lambda R$ and the third $-\nabla_\nu R^\nu{}_\lambda$. In the middle term, by (S2),

\[
g^{\sigma\nu}R^\rho{}_{\sigma\nu\lambda}=g^{\rho\alpha}g^{\sigma\nu}R_{\alpha\sigma\nu\lambda}=-g^{\rho\alpha}g^{\sigma\nu}R_{\sigma\alpha\nu\lambda}=-g^{\rho\alpha}R^\nu{}_{\alpha\nu\lambda}=-g^{\rho\alpha}R_{\alpha\lambda}=-R^\rho{}_\lambda .
\]

So $\nabla_\lambda R-2\nabla_\rho R^\rho{}_\lambda=0$. Since $\nabla_\rho(\delta^\rho{}_\lambda R)=\nabla_\lambda R$, this is

\[
\nabla_\rho G^\rho{}_\lambda=0,\qquad G_{\mu\nu}:=R_{\mu\nu}-\tfrac12\,g_{\mu\nu}R .
\]

$G_{\mu\nu}$ is the Einstein tensor. It is symmetric, it is built from the metric and its first and second derivatives, and its divergence vanishes identically, for every metric.

\textbf{Why the Einstein equations have this form.} Matter is described by its energy-momentum tensor $T_{\mu\nu}$, a symmetric tensor whose components are the densities and currents of energy and momentum; Chapter 7 derives it for the fields of this book. Conservation of energy and momentum is the statement $\nabla_\mu T^\mu{}_\nu=0$. Einstein's idea is that matter curves the geometry, with a proportionality: a tensor built from the metric equals a constant times $T$. For this to be consistent with conservation, the geometric side must have zero divergence for every metric. $G$ has this property, and so does $g$ itself (because $\nabla g=0$). The Einstein equations are

\[
G_{\mu\nu}+\Lambda\,g_{\mu\nu}=\kappa\,T_{\mu\nu},
\]

with a constant $\kappa$ and a cosmological constant $\Lambda$. In four dimensions, with $\kappa=8\pi G_{\mathrm N}/c^4$ where $G_{\mathrm N}$ is Newton's constant, these equations reproduce Newton's law of gravitation for slow bodies in weak fields (a standard result that we quote). In eight dimensions $\kappa$ is a free constant; the project's worked examples set $\kappa=1$ (Stage-2 document, §15.5, with $H=\kappa=1$ in the primordial field; the Stage-3 experiments EXP-1, with $H=\kappa=1$ and the mass of the field equal to 1, in the primordial field, and EXP-2, with $\kappa=1$, written $\kappa_8$ there, and the mass equal to 1; \texttt{provenan\allowbreak{}ce/DIRAC\allowbreak{}16COMPLE\allowbreak{}X\_DARK\_S\allowbreak{}ECTOR\_NU\allowbreak{}MERICS.m\allowbreak{}d}, §3.1). The project writes the equations without $\Lambda$ and in the mixed form

\[
G^\mu{}_\nu=\kappa\,T^\mu{}_\nu
\]

(Stage-2 document, §15.2). For an observer moving along $x^4$ in a Gaussian normal chart, the energy density is $\rho=T_{44}$ (Chapter 7), and because $g^{44}=-1$ the mixed component is $T^4{}_4=g^{44}T_{44}=-\rho$. The transverse pressures are $p_{(\mu)}=T^\mu{}_\mu$ (no sum, $\mu\ne4$).

\textbf{The trace-reversed form.} Take the trace of $G_{\mu\nu}=\kappa T_{\mu\nu}$ with $g^{\mu\nu}$. Since $g^{\mu\nu}g_{\mu\nu}=\delta^\mu{}_\mu=n$, we get $R-\tfrac n2R=\kappa T$ with $T:=T^\mu{}_\mu$, so $R=-2\kappa T/(n-2)$. Inserting this back,

\[
R_{\mu\nu}=\kappa\Bigl(T_{\mu\nu}-\frac{1}{n-2}\,T\,g_{\mu\nu}\Bigr)\qquad\bigl(\text{for }n=8:\ R_{\mu\nu}=\kappa\bigl(T_{\mu\nu}-\tfrac16Tg_{\mu\nu}\bigr)\bigr).
\]

The factor $\tfrac16=\tfrac1{n-2}$ is the one that appears in the evolution equations of the Stage-3 experiment EXP-2. Without matter ($T_{\mu\nu}=0$) the equations say $R_{\mu\nu}=0$: a vacuum solution has vanishing Ricci tensor.

\textbf{Worked example: what the primordial field would need as a source.} The primordial field is a prescribed metric, not a solution of the vacuum equations: $R_{44}=-6H^2a_4'^2$ is not zero unless $a_4$ is constant, and we will see that $G^4{}_4$ is never zero. Its mixed Einstein tensor is diagonal, with (Stage-2 document, §15.1; checks \texttt{P\_einste\allowbreak{}in\_Gmixe\allowbreak{}dClosedF\allowbreak{}orms} and \texttt{P\_einste\allowbreak{}in\_offDi\allowbreak{}agonalZe\allowbreak{}ro} in \texttt{wolfram-\allowbreak{}primordi\allowbreak{}al-repor\allowbreak{}t.json}, measurement \texttt{P\_einste\allowbreak{}in.einst\allowbreak{}einMixed\allowbreak{}Diagonal} in \texttt{python-p\allowbreak{}rimordia\allowbreak{}l-report\allowbreak{}.json})

\[
\begin{aligned}
G^0{}_0&=-3H^2\bigl(a_4'^2-5\bigr), &\qquad G^i{}_i&=H^2\bigl(15-3a_4'^2+a_4''\bigr)\quad(i=1,2,3),\\
G^4{}_4&=3H^2\bigl(7+a_4'^2\bigr), &\qquad G^j{}_j&=H^2\bigl(15-3a_4'^2-a_4''\bigr)\quad(j=5,6,7).
\end{aligned}
\]

The component $G^4{}_4$ follows from our two results: $G^4{}_4=g^{44}R_{44}-\tfrac12R=6H^2a_4'^2-3H^2(a_4'^2-7)=3H^2(7+a_4'^2)$. The Python verifier also checks the contracted Bianchi identity $\nabla_\mu G^\mu{}_\nu=0$ for all eight values of $\nu$ (measurement \texttt{P\_einste\allowbreak{}in.contr\allowbreak{}actedBia\allowbreak{}nchi}). If the field obeyed the 8-dimensional Einstein equations, its source would need the energy density

\[
\rho_{\mathrm{req}}=-T^4{}_4=-\frac{G^4{}_4}{\kappa}=-\frac{3H^2\bigl(7+a_4'^2\bigr)}{\kappa}\le-\frac{21H^2}{\kappa}<0
\]

for every $a_4$. This negative energy density is the Stage-2 result checked as \texttt{P\_einste\allowbreak{}in\_rhoRe\allowbreak{}quiredNe\allowbreak{}gative}.

For the notebook's choice $a_4=t$ ($a_4'=1$, $a_4''=0$) the formulas give $R=-36H^2$, $G^\mu{}_\nu=\mathrm{diag}(12,12,12,12,24,12,12,12)\,H^2$, $\rho_{\mathrm{req}}=-24H^2/\kappa$ and the equal pressures $p=12H^2/\kappa$ in all seven transverse directions, so $w=p/\rho=-\tfrac12$ (\texttt{python-p\allowbreak{}rimordia\allowbreak{}l-report\allowbreak{}.json}, measurement \texttt{P\_a4line\allowbreak{}ar}, row "a4 = t"). A cosmological constant cannot supply this: $\Lambda$ contributes equally to all diagonal components, but $G^0{}_0=G^4{}_4$ would require $-3H^2(a_4'^2-5)=3H^2(7+a_4'^2)$, that is $a_4'^2=-1$, which no real $a_4$ satisfies. Chapter 9 discusses which states of dirac16complex can supply such a source.

\subsection{The Einstein-Lovelock equations}

\textbf{The question.} Is $G_{\mu\nu}$ the only possible left-hand side? Lovelock answered this in 1971 (D. Lovelock, J. Math. Phys. 12, 498 (1971)). The notebook's title and its cell 14 refer to his result, in the form of equation (4.38) of the book by Lovelock and Rund, \emph{Tensors, Differential Forms, and Variational Principles} (\texttt{handoff/\allowbreak{}surveys/\allowbreak{}survey\_n\allowbreak{}otebook-\allowbreak{}physics.\allowbreak{}md}, item 1). We define the objects and prove what can be proved at this level; the rest is quoted.

\textbf{Generalized Kronecker delta.} For $2p$ indices,

\[
\delta^{\mu_1\cdots\mu_p}{}_{\nu_1\cdots\nu_p}:=\det\begin{pmatrix}\delta^{\mu_1}{}_{\nu_1}&\cdots&\delta^{\mu_1}{}_{\nu_p}\\ \vdots&&\vdots\\ \delta^{\mu_p}{}_{\nu_1}&\cdots&\delta^{\mu_p}{}_{\nu_p}\end{pmatrix}.
\]

It equals $+1$ if the upper indices are distinct and are an even rearrangement of the lower ones, $-1$ for an odd rearrangement, and 0 otherwise. Exchanging two upper indices exchanges two rows and flips the sign. Hence it vanishes whenever two upper indices are equal, and therefore it vanishes identically when $p>n$: among more than $n$ indices, each running over $n$ values, two must be equal. For $p=2$, $\delta^{\alpha\beta}{}_{\mu\nu}=\delta^\alpha{}_\mu\delta^\beta{}_\nu-\delta^\alpha{}_\nu\delta^\beta{}_\mu$.

\textbf{Lovelock scalars and tensors.} Write $R^{\alpha\beta}{}_{\mu\nu}:=g^{\beta\lambda}R^\alpha{}_{\lambda\mu\nu}$; it is antisymmetric in $\alpha,\beta$ by (S2) and in $\mu,\nu$ by (S1). For $k=0,1,2,\dots$ define

\[
\begin{aligned}
L_k&=\frac{1}{2^k}\,\delta^{\mu_1\nu_1\cdots\mu_k\nu_k}{}_{\alpha_1\beta_1\cdots\alpha_k\beta_k}\,R^{\alpha_1\beta_1}{}_{\mu_1\nu_1}\cdots R^{\alpha_k\beta_k}{}_{\mu_k\nu_k},\\
E^{(k)\mu}{}_\nu&=-\frac{1}{2^{k+1}}\,\delta^{\mu\,\mu_1\nu_1\cdots\mu_k\nu_k}{}_{\nu\,\alpha_1\beta_1\cdots\alpha_k\beta_k}\,R^{\alpha_1\beta_1}{}_{\mu_1\nu_1}\cdots R^{\alpha_k\beta_k}{}_{\mu_k\nu_k}.
\end{aligned}
\]

Normalizations differ between books; any other choice is absorbed into the coefficients below. For $k=0$: $L_0=1$ and $E^{(0)\mu}{}_\nu=-\tfrac12\delta^\mu{}_\nu$.

\textbf{Order 1 is Einstein.} For $k=1$, $L_1=\tfrac12\delta^{\mu\nu}{}_{\alpha\beta}R^{\alpha\beta}{}_{\mu\nu}=\tfrac12\bigl(R^{\mu\nu}{}_{\mu\nu}-R^{\nu\mu}{}_{\mu\nu}\bigr)=R^{\mu\nu}{}_{\mu\nu}=R$. For the tensor, expand the 3 by 3 determinant along its first row:

\[
\delta^{\mu\alpha\beta}{}_{\nu\sigma\lambda}=\delta^\mu{}_\nu\,\delta^{\alpha\beta}{}_{\sigma\lambda}-\delta^\mu{}_\sigma\,\delta^{\alpha\beta}{}_{\nu\lambda}+\delta^\mu{}_\lambda\,\delta^{\alpha\beta}{}_{\nu\sigma}.
\]

Contract with $R^{\sigma\lambda}{}_{\alpha\beta}$. The first term gives $\delta^\mu{}_\nu\cdot2R$. The second gives $-\delta^{\alpha\beta}{}_{\nu\lambda}R^{\mu\lambda}{}_{\alpha\beta}=-2R^{\mu\lambda}{}_{\nu\lambda}=-2R^\mu{}_\nu$; here $R^{\mu\lambda}{}_{\nu\lambda}=g^{\mu\alpha}g^{\lambda\beta}R_{\alpha\beta\nu\lambda}=g^{\mu\alpha}g^{\lambda\beta}R_{\beta\alpha\lambda\nu}=g^{\mu\alpha}R_{\alpha\nu}$ by (S1) and (S2). The third gives $\delta^{\alpha\beta}{}_{\nu\sigma}R^{\sigma\mu}{}_{\alpha\beta}=2R^{\sigma\mu}{}_{\nu\sigma}=-2R^{\mu\sigma}{}_{\nu\sigma}=-2R^\mu{}_\nu$. So

\[
E^{(1)\mu}{}_\nu=-\tfrac14\bigl(2R\,\delta^\mu{}_\nu-4R^\mu{}_\nu\bigr)=R^\mu{}_\nu-\tfrac12R\,\delta^\mu{}_\nu=G^\mu{}_\nu .
\]

\textbf{Lovelock's theorem (quoted, not proved here).} Each $E^{(k)}$ is symmetric (with both indices lowered) and has identically vanishing divergence, and, up to a constant factor, it is what one obtains from $\int\sqrt{\lvert g\rvert}\,L_k\,d^nx$ by varying the metric. Conversely, every symmetric tensor with identically vanishing divergence that is built from the metric and its first and second derivatives is a combination of the $E^{(k)}$. The Einstein-Lovelock field equations are therefore

\[
\sum_{k\ge0}\alpha_k\,E^{(k)\mu}{}_\nu=\kappa\,T^\mu{}_\nu ,
\]

with constants $\alpha_k$; the choice $\alpha_0=-2\Lambda$, $\alpha_1=1$ and all others zero gives back $G^\mu{}_\nu+\Lambda\delta^\mu{}_\nu=\kappa T^\mu{}_\nu$. For $k=1$ we have proved the divergence property ourselves (the contracted Bianchi identity).

\textbf{Which orders can exist (the vanishing is proved).} $E^{(k)}$ contains the generalized delta with $p=2k+1$ upper indices, so it vanishes identically when $2k+1>n$. In $n=4$ only $k=0$ and $k=1$ can contribute, so by Lovelock's theorem (quoted) Einstein's equations with a cosmological constant are the only choice in four dimensions. In $n=8$ only $k=0,1,2,3$ can contribute; $E^{(0)\mu}{}_\nu=-\tfrac12\delta^\mu{}_\nu$ and $E^{(1)}=G$ are visibly nonzero (Section 4.8 gives a metric with $G^4{}_4\ne0$), and that $E^{(2)}$ and $E^{(3)}$ do not vanish identically is quoted. The order-2 scalar is the Gauss-Bonnet combination $L_2=R^2-4R_{\mu\nu}R^{\mu\nu}+R_{\mu\nu\rho\sigma}R^{\mu\nu\rho\sigma}$ (a standard expansion of the definition that we quote; this book does not use it). Each order has a different power of length: $R$ has the dimension $1/\text{length}^2$, so $L_k$ has $1/\text{length}^{2k}$.

\textbf{The notebook and the project.} The notebook's cell 14 states "m-1 = 8/2 - 1 = 3", keeps the orders 1, 2 and 3, and writes the vacuum equations as $0=-\Lambda+H^{-2}w_1\,\mathrm{Lovelock1}+H^{-4}w_2\,\mathrm{Lovelock2}+H^{-6}w_3\,\mathrm{Lovelock3}$ with pure numbers $w_1,w_2,w_3,\Lambda$. The inverse length $H$ is introduced exactly to make the orders, which have different dimensions, comparable (\texttt{survey\_n\allowbreak{}otebook-\allowbreak{}physics.\allowbreak{}md}, item 1). This agrees with the counting above: orders higher than 3 vanish identically in eight dimensions. But no code of the notebook defines Lovelock2 or Lovelock3, and its curvature code stops at the Einstein tensor (cells 583 and 584, which the Stage-2 verifier reproduces exactly: checks \texttt{P\_einste\allowbreak{}in\_noteb\allowbreak{}ookCell5\allowbreak{}83} and \texttt{P\_einste\allowbreak{}in\_noteb\allowbreak{}ookCell5\allowbreak{}84}). The project does not compute the orders 2 and 3 either (Stage-2 document, §1, non-claim 4). Every gravitational statement in this book therefore uses the 8-dimensional Einstein equations only. Whether the Lovelock terms change the picture is an open problem (Chapter 18).

\subsection{The vielbein}

\textbf{Why a spinor needs more than a metric.} Chapter 2 built spinors from gamma matrices with the constant relation $\gamma^a\gamma^b+\gamma^b\gamma^a=2\eta^{ab}$. In a curved space the metric $g_{\mu\nu}(x)$ changes from point to point and is not $\eta$. The way out is to describe every point with its own set of $n$ reference directions that are perpendicular to each other and of unit length. In such a frame the metric looks like $\eta$, and the constant gamma matrices of Chapter 2 can be used unchanged. The spinor components are always measured with respect to such a frame.

\textbf{Definition (vielbein).} At every point choose $n$ vectors $E_0,\dots,E_{n-1}$ with

\[
g(E_a,E_b)=\eta_{ab}\qquad(a,b=0,\dots,n-1).
\]

Their components are written $E_a{}^\mu=:e_a{}^\mu$, the inverse vielbein. The matrix $(e_a{}^\mu)$ is invertible (orthonormal vectors are linearly independent), and its inverse is the vielbein $e_\mu{}^a$ ("viel Beine", German for "many legs"; in four dimensions it is called a vierbein or tetrad, and the notebook calls its eight-legged version an octad, cell 281). As in the project we write $e_\mu{}^a$ with row index $\mu$ and column index $a$ (Stage 1, §5.1). By definition of the inverse,

\[
e_\mu{}^a\,e_a{}^\nu=\delta_\mu{}^\nu,\qquad e_a{}^\mu\,e_\mu{}^b=\delta_a{}^b .
\]

\textbf{The metric from the vielbein.} The orthonormality condition reads $g_{\mu\nu}e_a{}^\mu e_b{}^\nu=\eta_{ab}$. Multiply by $e_\rho{}^a e_\sigma{}^b$ and sum over $a$ and $b$; with $e_a{}^\mu e_\rho{}^a=\delta^\mu{}_\rho$ this gives

\[
g_{\rho\sigma}=e_\rho{}^a\,\eta_{ab}\,e_\sigma{}^b,\qquad g^{\mu\nu}=e_a{}^\mu\,\eta^{ab}\,e_b{}^\nu .
\]

For the second formula, multiply the candidate $e_c{}^\nu\eta^{cd}e_d{}^\rho$ by $g_{\mu\nu}=e_\mu{}^a\eta_{ab}e_\nu{}^b$: using $e_\nu{}^be_c{}^\nu=\delta^b{}_c$ and $\eta_{ab}\eta^{bd}=\delta_a{}^d$ the product is $e_\mu{}^de_d{}^\rho=\delta_\mu{}^\rho$. In matrix language $g=e\,\eta\,e^T$, so $\det g=(\det e)^2\det\eta$. For signature (4,4), $\det\eta=(-1)^4=+1$, hence $\det g=(\det e)^2>0$ and $\sqrt{\lvert g\rvert}=\lvert\det e\rvert$. This is why the primordial field has $\det g=+\cos^2z$.

\textbf{Frame components.} A vector $V$ has frame components $V^a:=e_\mu{}^aV^\mu$, and $V^\mu=e_a{}^\mu V^a$. Then $g(V,W)=g_{\mu\nu}e_a{}^\mu e_b{}^\nu V^aW^b=\eta_{ab}V^aW^b$: in frame components, lengths are computed with $\eta$. Frame indices are raised and lowered with $\eta$: $V_a=\eta_{ab}V^b$. Frame components do not change when the coordinates change, because $e_\mu{}^a$ is a covector for each fixed $a$; they change when the frame is turned.

\textbf{Existence and freedom.} Near every point of a manifold with a metric such frames exist (orthonormalize any basis step by step; we quote this standard fact, and for all metrics of the book we write the frame down explicitly). They are not unique. If $\Lambda(x)$ is, at each point, a matrix with $\Lambda^T\eta\Lambda=\eta$ (an element of the group O(4,4) of Chapter 2), then $e'_\mu{}^a=\Lambda^a{}_b\,e_\mu{}^b$ is another vielbein for the same metric:

\[
e'_\mu{}^a\eta_{ab}e'_\nu{}^b=e_\mu{}^c\bigl(\Lambda^a{}_c\,\eta_{ab}\,\Lambda^b{}_d\bigr)e_\nu{}^d=e_\mu{}^c\,(\Lambda^T\eta\Lambda)_{cd}\,e_\nu{}^d=e_\mu{}^c\eta_{cd}e_\nu{}^d=g_{\mu\nu}.
\]

Such a point-dependent turning of the frame is a local frame rotation (local Lorentz transformation). Physics must not depend on the choice of frame; for spinors this requirement is local spin invariance (Section 4.12).

\textbf{Curved gamma matrices.} With the frame gammas $\gamma^a$ of Chapter 2 define

\[
\gamma^\mu:=e_a{}^\mu\gamma^a,\qquad \gamma_\mu:=g_{\mu\nu}\gamma^\nu=e_\mu{}^a\eta_{ab}\gamma^b .
\]

They obey the curved Clifford relation: $\gamma^\mu\gamma^\nu+\gamma^\nu\gamma^\mu=e_a{}^\mu e_b{}^\nu(\gamma^a\gamma^b+\gamma^b\gamma^a)=2e_a{}^\mu\eta^{ab}e_b{}^\nu=2g^{\mu\nu}$. As in the Stage documents, a gamma matrix with a numeric upper index, $\gamma^0,\dots,\gamma^7$, is always the constant frame matrix; a curved gamma with a numeric label is written $\gamma^{x_0},\dots,\gamma^{x_7}$.

\textbf{Worked example: the polar plane.} The unit radial vector $E_0$ has polar components $(1,0)$ and the unit angular vector $E_1$ has $(0,1/r)$, because $g(E_1,E_1)=r^2(1/r)^2=1$. So $e_a{}^\mu=\mathrm{diag}(1,1/r)$ and $e_\mu{}^a=\mathrm{diag}(1,r)$, with frame indices $a=0$ (radial) and $a=1$ (angular). The plane is Euclidean, $\eta=\mathrm{diag}(1,1)$, and a two-dimensional Clifford algebra is given by $\gamma^0=\sigma_1$, $\gamma^1=\sigma_2$. The curved gammas are $\gamma^r=\sigma_1$ and $\gamma^\varphi=\sigma_2/r$; for example $\gamma^\varphi\gamma^\varphi+\gamma^\varphi\gamma^\varphi=2\sigma_2^2/r^2=2/r^2=2g^{\varphi\varphi}$.

\textbf{Worked example: the primordial field.} A diagonal metric $g_{\mu\mu}=\eta_{\mu\mu}h_\mu^2$ has the diagonal vielbein $e_\mu{}^a=h_\mu\,\delta_\mu{}^a$ (no sum), since then $e\,\eta\,e^T=\mathrm{diag}(\eta_{\mu\mu}h_\mu^2)$. For the primordial field (Stage-2 document, §4.2; check \texttt{P\_metric\allowbreak{}\_vielbei\allowbreak{}nProduct})

\[
e_\mu{}^a=\mathrm{diag}\bigl(\cot z,\ s^{1/6}e^{a_4}\ (\times3),\ 1,\ s^{1/6}e^{-a_4}\ (\times3)\bigr),
\]

and the curved gammas are

\[
\begin{aligned}
&\gamma^{x_0}=\tan z\,\gamma^0,\qquad \gamma^{x_i}=s^{-1/6}e^{-a_4}\gamma^i\ \ (i=1,2,3),\\
&\gamma^{x_4}=\gamma^4,\qquad \gamma^{x_j}=s^{-1/6}e^{a_4}\gamma^j\ \ (j=5,6,7).
\end{aligned}
\]

In the notebook these are the matrices \texttt{(T16\textasciicircum{}α)} of cell 475, built from its diagonal octad (cells 279, 283 and 302).

\subsection{The spin connection and the vielbein postulate}

\textbf{The problem.} We want to differentiate the frame components $V^a=e_\mu{}^aV^\mu$ of a vector field. The frame turns from point to point, so $\partial_\mu V^a$ mixes the change of $V$ with the turning of the frame, just as $\partial_\mu V^\nu$ mixed the change of $V$ with the turning of the coordinate lines. The cure is the same: a correction, linear in $V$,

\[
\nabla_\mu V^a:=\partial_\mu V^a+\omega_\mu{}^a{}_b\,V^b ,
\]

with coefficients $\omega_\mu{}^a{}_b(x)$, the spin connection. We fix them by one natural requirement: the frame components of the covariant derivative must be the covariant derivative of the frame components,

\[
e_\nu{}^a\,\nabla_\mu V^\nu=\nabla_\mu V^a\qquad\text{for every vector field }V .
\]

\textbf{Derivation of the vielbein postulate.} Write out both sides. The left is $e_\nu{}^a(\partial_\mu V^\nu+\Gamma^\nu{}_{\mu\lambda}V^\lambda)$. The right is $\partial_\mu(e_\lambda{}^aV^\lambda)+\omega_\mu{}^a{}_be_\lambda{}^bV^\lambda=(\partial_\mu e_\lambda{}^a)V^\lambda+e_\lambda{}^a\partial_\mu V^\lambda+\omega_\mu{}^a{}_be_\lambda{}^bV^\lambda$. The terms with $\partial_\mu V$ agree, and what remains is

\[
\bigl(\partial_\mu e_\lambda{}^a-\Gamma^\nu{}_{\mu\lambda}e_\nu{}^a+\omega_\mu{}^a{}_b\,e_\lambda{}^b\bigr)V^\lambda=0 .
\]

Since this holds for every $V$, the bracket vanishes. Renaming the indices, this is the vielbein postulate (Stage 1, §5.2):

\[
\nabla_\mu e_\nu{}^a:=\partial_\mu e_\nu{}^a-\Gamma^\rho{}_{\mu\nu}\,e_\rho{}^a+\omega_\mu{}^a{}_b\,e_\nu{}^b=0 .
\]

In words: the total covariant derivative of the vielbein, with the Christoffel symbols acting on its curved index and the spin connection acting on its frame index, vanishes.

\textbf{Theorem (the canonical spin connection).} The vielbein postulate has exactly one solution,

\[
\omega_\mu{}^a{}_b=e_b{}^\nu\bigl(\Gamma^\rho{}_{\mu\nu}\,e_\rho{}^a-\partial_\mu e_\nu{}^a\bigr),
\]

and the lowered connection $\omega_{\mu ab}:=\eta_{ac}\,\omega_\mu{}^c{}_b$ is antisymmetric: $\omega_{\mu ab}=-\omega_{\mu ba}$.

\emph{Proof of the formula.} Multiply the postulate by $e_b{}^\nu$ and sum over $\nu$. The last term becomes $\omega_\mu{}^a{}_c\,e_\nu{}^ce_b{}^\nu=\omega_\mu{}^a{}_c\,\delta^c{}_b=\omega_\mu{}^a{}_b$, and solving for it gives the formula. Conversely, inserting the formula back into the postulate and using $e_\nu{}^be_b{}^\sigma=\delta_\nu{}^\sigma$ makes it an identity, so the formula is a solution, and the derivation shows it is the only one.

\emph{Proof of the antisymmetry.} With $e_{\rho a}:=\eta_{ac}e_\rho{}^c=g_{\rho\sigma}e_a{}^\sigma$ (the second form follows from $g_{\rho\sigma}=e_\rho{}^c\eta_{cd}e_\sigma{}^d$) and $\Gamma_{\sigma\mu\nu}=g_{\sigma\rho}\Gamma^\rho{}_{\mu\nu}$,

\[
\omega_{\mu ab}=e_a{}^\sigma e_b{}^\nu\,\Gamma_{\sigma\mu\nu}-\eta_{ac}\,e_b{}^\nu\,\partial_\mu e_\nu{}^c .
\]

Add the same expression with $a$ and $b$ exchanged. In the first terms, $\Gamma_{\sigma\mu\nu}+\Gamma_{\nu\mu\sigma}=\partial_\mu g_{\nu\sigma}$, as one sees by adding the two Christoffel formulas; so they give $e_a{}^\sigma e_b{}^\nu\,\partial_\mu g_{\sigma\nu}$. Differentiating $g_{\sigma\nu}=e_\sigma{}^c\eta_{cd}e_\nu{}^d$ with the product rule,

\[
e_a{}^\sigma e_b{}^\nu\,\partial_\mu g_{\sigma\nu}=\eta_{cb}\,e_a{}^\sigma\partial_\mu e_\sigma{}^c+\eta_{ad}\,e_b{}^\nu\partial_\mu e_\nu{}^d ,
\]

which is exactly the sum of the two second terms with the opposite sign. Hence $\omega_{\mu ab}+\omega_{\mu ba}=0$. $\square$

For each $\mu$ the antisymmetric $\omega_{\mu ab}$ has $n(n-1)/2$ independent components, 28 for $n=8$: one for each of the 28 generators $S^{ab}$ ($a<b$) of Chapter 2. So for each $\mu$ the matrix $A=(\omega_\mu{}^a{}_b)$ satisfies $A^T\eta=-\eta A$, because $\eta A=(\omega_{\mu ab})$ is antisymmetric and $\eta$ is symmetric. This is the condition for $1+\epsilon A$ to lie in O(4,4) to first order in a small number $\epsilon$: $(1+\epsilon A)^T\eta(1+\epsilon A)=\eta+\epsilon(A^T\eta+\eta A)+\epsilon^2A^T\eta A=\eta+O(\epsilon^2)$. So $\omega_\mu$ is an infinitesimal frame rotation (the set of such matrices is called so(4,4)).

\textbf{The inverse vielbein postulate.} Differentiate $e_a{}^\nu e_\nu{}^b=\delta_a{}^b$ and replace $\partial_\mu e_\nu{}^b$ by $\Gamma^\rho{}_{\mu\nu}e_\rho{}^b-\omega_\mu{}^b{}_ce_\nu{}^c$ from the postulate: $(\partial_\mu e_a{}^\nu)e_\nu{}^b+e_a{}^\nu\Gamma^\rho{}_{\mu\nu}e_\rho{}^b-\omega_\mu{}^b{}_a=0$. Multiplying by $e_b{}^\sigma$ gives

\[
\partial_\mu e_a{}^\sigma+\Gamma^\sigma{}_{\mu\nu}\,e_a{}^\nu-\omega_\mu{}^b{}_a\,e_b{}^\sigma=0 .
\]

\textbf{Behaviour under a frame rotation.} Write the postulate as a statement about the column $e_\nu=(e_\nu{}^a)_a$ and the matrices $\omega_\mu=(\omega_\mu{}^a{}_b)$: $\partial_\mu e_\nu-\Gamma^\rho{}_{\mu\nu}e_\rho+\omega_\mu e_\nu=0$. For the rotated frame $e'_\nu=\Lambda e_\nu$ (same metric, so the same $\Gamma$),

\[
\partial_\mu(\Lambda e_\nu)-\Gamma^\rho{}_{\mu\nu}\Lambda e_\rho+\omega'_\mu\Lambda e_\nu=(\partial_\mu\Lambda)e_\nu-\Lambda\omega_\mu e_\nu+\omega'_\mu\Lambda e_\nu .
\]

This vanishes for all $\nu$, and the columns $e_\nu$ span all columns, so

\[
\omega'_\mu=\Lambda\,\omega_\mu\,\Lambda^{-1}-(\partial_\mu\Lambda)\,\Lambda^{-1}.
\]

The second term shows that $\omega$ is not a tensor in its frame indices: it contains the rate at which the frame is turned.

\textbf{Worked example: the polar plane.} With $e_\mu{}^a=\mathrm{diag}(1,r)$, $e_a{}^\mu=\mathrm{diag}(1,1/r)$ and the Christoffel symbols of Section 4.5, the formula gives, for $\mu=\varphi$,

\[
\omega_\varphi{}^0{}_1=e_1{}^\varphi\bigl(\Gamma^r{}_{\varphi\varphi}e_r{}^0-\partial_\varphi e_\varphi{}^0\bigr)=\frac1r(-r)(1)=-1,\qquad
\omega_\varphi{}^1{}_0=e_0{}^r\bigl(\Gamma^\varphi{}_{\varphi r}e_\varphi{}^1-\partial_\varphi e_r{}^1\bigr)=1\cdot\frac1r\cdot r=1,
\]

and $\omega_\varphi{}^0{}_0=\omega_\varphi{}^1{}_1=0$. For $\mu=r$: $\omega_r{}^1{}_1=e_1{}^\varphi(\Gamma^\varphi{}_{r\varphi}e_\varphi{}^1-\partial_re_\varphi{}^1)=\frac1r(\frac1r\cdot r-1)=0$, and the other three vanish as well. Since $\eta=1$ here, $\omega_{\varphi01}=-1=-\omega_{\varphi10}$: antisymmetric, as the theorem says. The meaning is simple. The radial and angular unit vectors at the angle $\varphi$ are the Cartesian unit vectors turned by the angle $\varphi$; moving in $\varphi$ turns the frame at the rate $d\varphi/d\varphi=1$, and the spin connection records this rate.

\textbf{Diagonal vielbeins.} For $e_\mu{}^a=h_\mu\delta_\mu{}^a$ (no sum) the formula gives $\omega_\mu{}^a{}_b=\frac{1}{h_b}\bigl(\Gamma^a{}_{\mu b}h_a-\partial_\mu(h_b\delta_b{}^a)\bigr)$. For $a=b$ this is $\Gamma^a{}_{\mu a}-\partial_\mu\ln h_a=0$ by cases (a) and (b) of Section 4.5. For $a\ne b$ it is $\frac{h_a}{h_b}\Gamma^a{}_{\mu b}$, which by cases (b), (c) and (d) is nonzero only for $\mu=a$ or $\mu=b$. Lowering with $\eta$:

\[
\omega_{\mu\mu b}=\eta_{\mu\mu}\,\frac{\partial_bh_\mu}{h_b},\qquad \omega_{\mu b\mu}=-\omega_{\mu\mu b}\qquad(b\ne\mu,\ \text{no sum}),
\]

and all other components vanish. (For $\mu=a$: $\omega_a{}^a{}_b=\frac{h_a}{h_b}\partial_b\ln h_a$. For $\mu=b$: $\omega_b{}^a{}_b=\frac{h_a}{h_b}\bigl(-\eta_{aa}\eta_{bb}\frac{h_b}{h_a^2}\partial_ah_b\bigr)=-\eta_{aa}\eta_{bb}\frac{\partial_ah_b}{h_a}$, and $\eta_{aa}$ times this is $-\eta_{bb}\partial_ah_b/h_a$, the antisymmetric partner.) Here $\partial_b$ is the derivative with respect to the coordinate whose number equals the frame index $b$.

\textbf{Worked example: the 24 components of the primordial field.} Write $\alpha_i=H\,s^{1/6}e^{a_4}$ for $i=1,2,3$ and $\alpha_j=H\,s^{1/6}e^{-a_4}$ for $j=5,6,7$. The scale factors $h_k$ ($k\ne0,4$) depend on $x^0$ and $x^4$, with $\partial_0h_k=H\cot z\,h_k$, $\partial_4h_i=Ha_4'h_i$ and $\partial_4h_j=-Ha_4'h_j$; $h_0$ depends on $x^0$ only and $h_4=1$. So the nonzero components have $\mu=k\in\{1,2,3,5,6,7\}$ and $b\in\{0,4\}$:

\begingroup
\small
\begin{longtable}{@{}>{\raggedright\arraybackslash}p{0.267\linewidth}>{\raggedright\arraybackslash}p{0.267\linewidth}>{\raggedright\arraybackslash}p{0.267\linewidth}@{}}
\toprule
\textbf{Component} & \textbf{Formula} & \textbf{Value} \\
\midrule
\endfirsthead
\toprule
\textbf{Component} & \textbf{Formula} & \textbf{Value} \\
\midrule
\endhead
$\omega_{i\,i0}=-\omega_{i\,0i}$ & $\eta_{ii}\,\partial_0h_i/h_0$ & $H\cot z\,h_i\tan z=\alpha_i$ \\
$\omega_{i\,i4}=-\omega_{i\,4i}$ & $\eta_{ii}\,\partial_4h_i/h_4$ & $Ha_4'h_i=\alpha_ia_4'$ \\
$\omega_{j\,j0}=-\omega_{j\,0j}$ & $\eta_{jj}\,\partial_0h_j/h_0$ & $-H\cot z\,h_j\tan z=-\alpha_j$ \\
$\omega_{j\,j4}=-\omega_{j\,4j}$ & $\eta_{jj}\,\partial_4h_j/h_4$ & $-(-Ha_4'h_j)=\alpha_ja_4'$ \\
\bottomrule
\end{longtable}
\endgroup

That is $6\times2\times2=24$ nonzero components, and $\omega_{0ab}=\omega_{4ab}=0$. In the notation of the Stage-2 document (§6): $\omega_{i\,0i}=-H s^{1/6}e^{a_4}$, $\omega_{i\,i4}=Hs^{1/6}e^{a_4}a_4'$, $\omega_{j\,0j}=Hs^{1/6}e^{-a_4}$, $\omega_{j\,j4}=Hs^{1/6}e^{-a_4}a_4'$. The Stage-2 verifiers find the same: the vielbein postulate holds in all 512 components, exactly 24 components are nonzero, with these closed forms (checks \texttt{P\_spinco\allowbreak{}nn\_vielb\allowbreak{}einPostu\allowbreak{}late512}, \texttt{P\_spinco\allowbreak{}nn\_count\allowbreak{}24}, \texttt{P\_spinco\allowbreak{}nn\_antis\allowbreak{}ymmetry} and \texttt{P\_spinco\allowbreak{}nn\_close\allowbreak{}dForms} in \texttt{wolfram-\allowbreak{}primordi\allowbreak{}al-repor\allowbreak{}t.json}; \texttt{nonzeroC\allowbreak{}ount} = 24 and check \texttt{P\_spinco\allowbreak{}nn} in \texttt{python-p\allowbreak{}rimordia\allowbreak{}l-report\allowbreak{}.json}).

\textbf{The mixed components and the space-time pairs.} Since $\omega_{\mu ab}=\eta_{aa}\omega_\mu{}^a{}_b$ (no sum, $\eta_{aa}=\pm1$),

\[
\omega_\mu{}^b{}_a=\eta_{bb}\,\omega_{\mu ba}=-\eta_{bb}\,\omega_{\mu ab}=-\eta_{aa}\eta_{bb}\,\omega_\mu{}^a{}_b\qquad(\text{no sum}).
\]

For a pair of two space-like or two time-like frame directions, $\eta_{aa}\eta_{bb}=+1$ and the mixed components are antisymmetric in $(a,b)$, like the lowered ones. For a space-time pair, one space-like and one time-like direction (a boost pair), $\eta_{aa}\eta_{bb}=-1$ and the mixed components are symmetric: $\omega_\mu{}^b{}_a=+\omega_\mu{}^a{}_b$. In the primordial field the 12 components of the pairs $(i,4)$ for $\mu=i$ and $(0,j)$ for $\mu=j$ are of this kind; for example $\omega_1{}^1{}_4=\omega_1{}^4{}_1=Hs^{1/6}e^{a_4}a_4'$ (Stage-2 document, §6; measurement \texttt{G2.p1.no\allowbreak{}nzeroOme\allowbreak{}gaMixedS\allowbreak{}ymmetric\allowbreak{}Part} = 12 in \texttt{wolfram-\allowbreak{}geometry\allowbreak{}-report.\allowbreak{}json}). The notebook's own stored mixed connection (the output of its cell 501) equals the correct $\omega_\mu{}^a{}_b$ entry by entry (check \texttt{P\_spinco\allowbreak{}nn\_noteb\allowbreak{}ookCell5\allowbreak{}01OmegaM\allowbreak{}uIJEqual\allowbreak{}sMixedOm\allowbreak{}ega}): the notebook computes $\omega$ correctly. Its error, which Section 4.14 explains, is in how it contracts $\omega$ with $S^{ab}$.

\textbf{An arbitrary gravitational field.} The general statements of this section are proved above for every vielbein. Stage 1 also tested them with exact rational arithmetic in a generic non-diagonal vielbein, the test geometry G1, $e_\mu{}^a=\delta_\mu{}^a+P_\mu{}^a(x)$ with polynomial entries $P$, at three rational points (Stage 1, §5.7). There all 512 components of the vielbein postulate vanish together with their 4096 first derivatives, $\omega_{\mu ab}=-\omega_{\mu ba}$, and 448 of the 512 components $\omega_{\mu ab}$ are nonzero at each point; in the primordial field G2 the count is 24 (Stage 1, Result 5.1; checks \texttt{GEO\_viel\allowbreak{}beinPost\allowbreak{}ulate\_G1}, \texttt{GEO\_viel\allowbreak{}beinPost\allowbreak{}ulate\_G2}, \texttt{GEO\_omeg\allowbreak{}aAntisym\allowbreak{}metry\_G1} and \texttt{GEO\_omeg\allowbreak{}aAntisym\allowbreak{}metry\_G2} in \texttt{wolfram-\allowbreak{}geometry\allowbreak{}-report.\allowbreak{}json} and \texttt{python-g\allowbreak{}eometry-\allowbreak{}report.j\allowbreak{}son}). The Python checker computes $\omega$ in two independent ways, from the Christoffel symbols as above and from the derivatives of the vielbein alone (the anholonomy route), and the two agree.

\subsection{The covariant derivative of a spinor}

\textbf{The requirement.} A spinor field $\Psi(x)$ has 16 components measured in the local frame. When the frame is turned by a local rotation, the components change by a spin transformation $R(x)$ (Chapter 2): $\Psi\to R\Psi$. We look for a derivative of the form

\[
D_\mu\Psi=\partial_\mu\Psi+\Omega_\mu\Psi ,
\]

with 16 by 16 matrices $\Omega_\mu(x)$, the spinor connection. It must be compatible with the frame connection in the following precise sense. For two spinor fields $\Psi$ and $\Phi$ (here ordinary columns of complex functions), the 8 bilinears $v^a=\bar\Psi\gamma^a\Phi$ transform like the frame components of a vector under frame rotations that can be reached continuously from no rotation at all (Chapter 2). The derivative of the spinors must induce the frame derivative of this vector:

\[
\partial_\mu v^a+\omega_\mu{}^a{}_b\,v^b=(D_\mu\bar\Psi)\gamma^a\Phi+\bar\Psi\gamma^aD_\mu\Phi ,\qquad D_\mu\bar\Psi:=\partial_\mu\bar\Psi-\bar\Psi\,\Omega_\mu ,
\]

where $\bar\Psi=\Psi^\dagger C$. (The rule for $D_\mu\bar\Psi$ is the one that makes $\bar\Psi\Phi$ behave like a scalar: $(D_\mu\bar\Psi)\Phi+\bar\Psi D_\mu\Phi=\partial_\mu(\bar\Psi\Phi)$. We check below that it agrees with conjugating $D_\mu\Psi$.) Expanding the right-hand side gives $\partial_\mu v^a+\bar\Psi[\gamma^a,\Omega_\mu]\Phi$, so the requirement is, for all $\Psi$ and $\Phi$, $\bar\Psi[\gamma^a,\Omega_\mu]\Phi=\omega_\mu{}^a{}_b\bar\Psi\gamma^b\Phi$, that is

\[
[\Omega_\mu,\gamma^a]=-\omega_\mu{}^a{}_b\,\gamma^b .
\]

\textbf{A commutator identity.} From the Clifford relation, for all $a,c,d$:

\[
[S^{cd},\gamma^a]=\eta^{da}\gamma^c-\eta^{ca}\gamma^d .
\]

\emph{Proof.} For $c=d$ both sides vanish. For $c\ne d$, $S^{cd}=\tfrac14(\gamma^c\gamma^d-\gamma^d\gamma^c)=\tfrac12\gamma^c\gamma^d$, because $\gamma^d\gamma^c=-\gamma^c\gamma^d$. Moving $\gamma^a$ through step by step with $\gamma^d\gamma^a=2\eta^{da}-\gamma^a\gamma^d$ and $\gamma^c\gamma^a=2\eta^{ca}-\gamma^a\gamma^c$:

\[
\gamma^c\gamma^d\gamma^a-\gamma^a\gamma^c\gamma^d=\gamma^c(2\eta^{da}-\gamma^a\gamma^d)-\gamma^a\gamma^c\gamma^d=2\eta^{da}\gamma^c-(2\eta^{ca}-\gamma^a\gamma^c)\gamma^d-\gamma^a\gamma^c\gamma^d=2\eta^{da}\gamma^c-2\eta^{ca}\gamma^d .
\]

Divide by 2. $\square$ (This is Stage 1, Result 3.5, checked in all 512 cases: check \texttt{ALG\_SabG\allowbreak{}ammaComm\allowbreak{}utator} in \texttt{python-g\allowbreak{}eometry-\allowbreak{}report.j\allowbreak{}son}.)

\textbf{Theorem (the spinor connection).} The matrices

\[
\Omega_\mu=\tfrac12\,\omega_{\mu ab}\,S^{ab}=\tfrac18\,\omega_{\mu ab}\,[\gamma^a,\gamma^b]
\]

(summed over all ordered pairs $(a,b)$; summed over $a<b$ only, the coefficient of $[\gamma^a,\gamma^b]$ is $\tfrac14$) satisfy $[\Omega_\mu,\gamma^a]=-\omega_\mu{}^a{}_b\gamma^b$. Every other solution differs from this one by a multiple of the identity matrix, and this one is the only traceless solution.

\emph{Proof.} By the identity,

\[
[\Omega_\mu,\gamma^a]=\tfrac12\,\omega_{\mu cd}\bigl(\eta^{da}\gamma^c-\eta^{ca}\gamma^d\bigr)=\tfrac12\,\omega_{\mu c}{}^{a}\gamma^c-\tfrac12\,\omega_\mu{}^a{}_d\,\gamma^d ,
\]

where $\omega_{\mu c}{}^a:=\omega_{\mu cd}\eta^{da}$. By the antisymmetry, $\omega_{\mu c}{}^a=-\eta^{ad}\omega_{\mu dc}=-\omega_\mu{}^a{}_c$, so both terms equal $-\tfrac12\omega_\mu{}^a{}_c\gamma^c$ and the sum is $-\omega_\mu{}^a{}_c\gamma^c$. If $\Omega'_\mu$ is another solution, $X=\Omega'_\mu-\Omega_\mu$ commutes with all eight $\gamma^a$, so by fact (F1) of Section 4.1 it is a multiple of the identity. Every $S^{ab}$ is traceless (it is a multiple of a commutator, and $\mathrm{tr}(AB)=\mathrm{tr}(BA)$), so $\Omega_\mu$ is traceless, and a traceless multiple of the identity is zero. $\square$

A multiple of the identity, $\Omega'_\mu=\Omega_\mu+iA_\mu\,1$, would be an additional electromagnetic-type field, not part of gravity; the canonical choice of the project has none (Stage 1, §5.3).

\textbf{The derivative of the Dirac adjoint.} The matrices $S^{ab}$ are real, and $(S^{ab})^TC=-CS^{ab}$ by fact (F3). Hence $\Omega_\mu$ is real with $\Omega_\mu^TC=-C\Omega_\mu$, and the conjugate transpose of $D_\mu\Psi$, multiplied by $C$, is

\[
(D_\mu\Psi)^\dagger C=\partial_\mu\Psi^\dagger C+\Psi^\dagger\Omega_\mu^TC=\partial_\mu\bar\Psi-\bar\Psi\,\Omega_\mu=D_\mu\bar\Psi ,
\]

consistent with the rule used above.

\textbf{Theorem (local spin covariance).} Let the frame be rotated, $e'_\mu{}^a=\Lambda^a{}_b(x)e_\mu{}^b$, and let $R(x)$ be a spin transformation that covers $\Lambda$ in the sense $R^{-1}\gamma^aR=\Lambda^a{}_b\gamma^b$, with $\det R=1$. (This is the covering of Section 2.14 written for $R^{-1}$. Theorem 2.9 there states $g\,\gamma(v)\,g^{-1}=\gamma(\Lambda_{\mathrm u}(g)v)$ with $\gamma(v)=\sum_av_a\gamma^a$; for $g=R^{-1}$ and $v$ the $a$-th unit vector it says $R^{-1}\gamma^aR=\sum_b\Lambda_{\mathrm u}(R^{-1})_{ba}\gamma^b$. Hence $\Lambda^a{}_b=\Lambda_{\mathrm u}(R^{-1})_{ba}=(\eta\,\Lambda_{\mathrm u}(R)\,\eta)_{ab}$, because $\Lambda_{\mathrm u}$ is a homomorphism, so $\Lambda_{\mathrm u}(R^{-1})=\Lambda_{\mathrm u}(R)^{-1}$, and $\Lambda^{-1}=\eta\Lambda^T\eta$ for every $\Lambda$ in O(4,4). For $R$ in Spin(4,4), $\Lambda_{\mathrm u}=\Lambda_{\mathrm t}$. For a rotation in a space-space or time-time plane this $\Lambda$ is the vector image of $R$; for a boost it is its inverse. That every element of Pin(4,4) has determinant 1 as a 16 by 16 matrix is proved in Chapter 2, Section 2.14, and derived in Stage 1, §4.2.) Then the spinor connection of the rotated frame is

\[
\Omega'_\mu=R\,\Omega_\mu R^{-1}-(\partial_\mu R)R^{-1},\qquad\text{and}\qquad D'_\mu(R\Psi)=R\,D_\mu\Psi .
\]

\emph{Proof.} By the uniqueness part of the previous theorem it suffices to show that the right-hand side, call it $X_\mu$, is traceless and satisfies $[X_\mu,\gamma^a]=-\omega'_\mu{}^a{}_b\gamma^b$ with $\omega'_\mu=\Lambda\omega_\mu\Lambda^{-1}-(\partial_\mu\Lambda)\Lambda^{-1}$ (Section 4.11). First, $[R\Omega_\mu R^{-1},\gamma^a]=R\,[\Omega_\mu,R^{-1}\gamma^aR]\,R^{-1}=\Lambda^a{}_bR[\Omega_\mu,\gamma^b]R^{-1}=-\Lambda^a{}_b\,\omega_\mu{}^b{}_c\,R\gamma^cR^{-1}$, and $R\gamma^cR^{-1}=(\Lambda^{-1})^c{}_d\gamma^d$ (invert the covering relation), so this term is $-(\Lambda\omega_\mu\Lambda^{-1})^a{}_d\gamma^d$. Second, differentiate $R^{-1}\gamma^aR=\Lambda^a{}_b\gamma^b$, using $\partial_\mu(R^{-1})=-R^{-1}(\partial_\mu R)R^{-1}$, and multiply by $R$ on the left and $R^{-1}$ on the right:

\[
-(\partial_\mu R)R^{-1}\gamma^a+\gamma^a(\partial_\mu R)R^{-1}=(\partial_\mu\Lambda^a{}_b)\,R\gamma^bR^{-1}=\bigl((\partial_\mu\Lambda)\Lambda^{-1}\bigr)^a{}_d\,\gamma^d .
\]

The left side is $-[(\partial_\mu R)R^{-1},\gamma^a]$. Adding the two parts, $[X_\mu,\gamma^a]=-(\Lambda\omega_\mu\Lambda^{-1}-(\partial_\mu\Lambda)\Lambda^{-1})^a{}_d\gamma^d$, as required. Traces: $\mathrm{tr}(R\Omega_\mu R^{-1})=\mathrm{tr}\,\Omega_\mu=0$, and $\mathrm{tr}((\partial_\mu R)R^{-1})=\mathrm{tr}(R^{-1}\partial_\mu R)=(\partial_\mu\det R)/\det R=0$, because $\det R=1$ is constant (the middle step is the determinant lemma of Section 4.5, whose proof works unchanged for complex matrices). Finally $D'_\mu(R\Psi)=(\partial_\mu R)\Psi+R\partial_\mu\Psi+R\Omega_\mu\Psi-(\partial_\mu R)\Psi=R\,D_\mu\Psi$. $\square$

Stage 1 tested this with finite, point-dependent products of two reflections built from exact rational vector fields: recomputing $\Omega$ from the rotated frame gives exactly $R\Omega R^{-1}-(\partial R)R^{-1}$ and $\gamma'^\mu=R\gamma^\mu R^{-1}$ (Stage 1, Result 5.4; checks \texttt{LAG\_loca\allowbreak{}lSpinInv\allowbreak{}ariance\_\allowbreak{}G1} and \texttt{LAG\_loca\allowbreak{}lSpinInv\allowbreak{}ariance\_\allowbreak{}G2} in \texttt{wolfram-\allowbreak{}geometry\allowbreak{}-report.\allowbreak{}json}).

\textbf{Theorem (the gammas are covariantly constant).} With the canonical connection,

\[
D_\mu\gamma^\nu:=\partial_\mu\gamma^\nu+\Gamma^\nu{}_{\mu\lambda}\gamma^\lambda+[\Omega_\mu,\gamma^\nu]=0\qquad\text{for all }\mu,\nu ,
\]

and this statement is equivalent to the vielbein postulate.

\emph{Proof.} Insert $\gamma^\nu=e_a{}^\nu\gamma^a$ and use the defining property: $e_a{}^\nu[\Omega_\mu,\gamma^a]=-e_a{}^\nu\,\omega_\mu{}^a{}_c\,\gamma^c=-\omega_\mu{}^b{}_a\,e_b{}^\nu\,\gamma^a$ (in the last step the summed indices were renamed, $a\to b$ and $c\to a$). So

\[
D_\mu\gamma^\nu=\bigl(\partial_\mu e_a{}^\nu+\Gamma^\nu{}_{\mu\lambda}e_a{}^\lambda-\omega_\mu{}^b{}_a\,e_b{}^\nu\bigr)\gamma^a ,
\]

and the bracket is the inverse vielbein postulate of Section 4.11, which vanishes. Conversely, if $D_\mu\gamma^\nu=0$, the bracket vanishes because the eight $\gamma^a$ are linearly independent (fact (F2)), which gives back the inverse postulate and hence the postulate. $\square$

This is Stage 1, Result 5.2, verified in all 64 pairs $(\mu,\nu)$ in G1 and G2 (checks \texttt{GEO\_gamm\allowbreak{}aCovaria\allowbreak{}ntConsta\allowbreak{}ncy\_G1} and \texttt{GEO\_gamm\allowbreak{}aCovaria\allowbreak{}ntConsta\allowbreak{}ncy\_G2}), and in the primordial field by the Stage-2 check \texttt{P\_gammaC\allowbreak{}onst\_Dmu\allowbreak{}GammaNuZ\allowbreak{}ero64}.

\textbf{The divergence identity.} Contract $\mu$ with $\nu$ in $D_\mu\gamma^\nu=0$ and use $\Gamma^\mu{}_{\mu\lambda}=\partial_\lambda\ln\sqrt{\lvert g\rvert}$ (Section 4.5):

\[
\partial_\mu\gamma^\mu+\bigl(\partial_\lambda\ln\sqrt{\lvert g\rvert}\bigr)\gamma^\lambda=-[\Omega_\mu,\gamma^\mu]
\qquad\Longleftrightarrow\qquad
\partial_\mu\bigl(\sqrt{\lvert g\rvert}\,\gamma^\mu\bigr)=\sqrt{\lvert g\rvert}\,[\gamma^\mu,\Omega_\mu].
\]

(Stage 1, Result 5.3; checks \texttt{GEO\_dive\allowbreak{}rgenceId\allowbreak{}entity\_G\allowbreak{}1} and \texttt{GEO\_dive\allowbreak{}rgenceId\allowbreak{}entity\_G\allowbreak{}2}.) This identity is used in Chapter 5, to show that the notebook's Lagrangian carries no dynamics for a Grassmann field, and in Chapter 7, to derive the field equations.

\textbf{Diagonal vielbeins: the contraction $\gamma^\mu\Omega_\mu$.} For $e_\mu{}^a=h_\mu\delta_\mu{}^a$ the field equations of Chapter 7 need only the combination $\gamma^\mu\Omega_\mu$ (summed over $\mu$). We show

\[
\gamma^\mu\Omega_\mu=\frac12\sum_{b=0}^{n-1}\frac{1}{h_b}\,\partial_b\ln\Bigl(\prod_{c\ne b}h_c\Bigr)\gamma^b ,\qquad \{\gamma^\mu,\Omega_\mu\}=0 .
\]

\emph{Proof.} Here $\gamma^\mu=\gamma^\mu_{\text{frame}}/h_\mu$, where $\gamma^\mu_{\text{frame}}$ is the frame matrix with the same number, and $S^{ab}=\tfrac12\gamma^a\gamma^b$ for $a\ne b$. By Section 4.11 only $\omega_{\mu\mu b}=-\omega_{\mu b\mu}$ with $b\ne\mu$ are nonzero, so $\Omega_\mu=\omega_{\mu\mu b}S^{\mu b}$ summed over $b\ne\mu$ (the two orders of each pair give equal terms, which cancels the $\tfrac12$). Then

\[
\gamma^\mu\Omega_\mu=\sum_\mu\sum_{b\ne\mu}\frac{1}{h_\mu}\,\eta_{\mu\mu}\frac{\partial_bh_\mu}{h_b}\cdot\tfrac12\gamma^\mu\gamma^\mu\gamma^b=\frac12\sum_b\frac1{h_b}\sum_{\mu\ne b}\frac{\partial_bh_\mu}{h_\mu}\,\gamma^b ,
\]

because $\gamma^\mu\gamma^\mu=\eta^{\mu\mu}$ for a frame matrix and $\eta_{\mu\mu}\eta^{\mu\mu}=1$. The inner sum is $\partial_b\ln\prod_{c\ne b}h_c$. For the anticommutator: $\gamma^\mu S^{\mu b}=\tfrac12\eta^{\mu\mu}\gamma^b$ and $S^{\mu b}\gamma^\mu=\tfrac12\gamma^\mu\gamma^b\gamma^\mu=-\tfrac12\eta^{\mu\mu}\gamma^b$, so every term of $\{\gamma^\mu,\Omega_\mu\}$ cancels. $\square$

This is the formula of Stage 1, §8.5 (checks \texttt{GEO\_diag\allowbreak{}onalSlas\allowbreak{}hFormula\allowbreak{}\_G2} and \texttt{GEO\_diag\allowbreak{}onalSlas\allowbreak{}hFormula\allowbreak{}\_G3}; \texttt{GEO\_anti\allowbreak{}commutat\allowbreak{}orGammaO\allowbreak{}megaVani\allowbreak{}shesDiag\allowbreak{}onal\_G2} for the anticommutator).

\textbf{Worked example: the polar plane.} Only $\omega_{\varphi01}=-1=-\omega_{\varphi10}$ is nonzero, so $\Omega_r=0$ and $\Omega_\varphi=\omega_{\varphi01}S^{01}=-\tfrac12\gamma^0\gamma^1=-\tfrac12\sigma_1\sigma_2=-\tfrac i2\sigma_3$. Directly,

\[
\gamma^\mu\Omega_\mu=\frac{\sigma_2}{r}\Bigl(-\frac i2\sigma_3\Bigr)=-\frac{i}{2r}\,\sigma_2\sigma_3=-\frac{i}{2r}\,i\sigma_1=\frac{1}{2r}\,\sigma_1 ,
\]

in agreement with the formula: only $b=0$ contributes, $\frac12\cdot\frac11\cdot\partial_r\ln r\,\gamma^0=\frac1{2r}\sigma_1$. The Dirac operator of the plane in polar coordinates is therefore

\[
\gamma^\mu D_\mu=\sigma_1\Bigl(\partial_r+\frac{1}{2r}\Bigr)+\frac{\sigma_2}{r}\,\partial_\varphi .
\]

The term $1/(2r)$ is exactly what the turning polar frame requires; Exercise 4.8 shows where it comes from.

\textbf{Worked example: the primordial field.} With the 24 components of Section 4.11 (only the pairs $(i,0)$, $(i,4)$ for $\mu=i$ and $(j,0)$, $(j,4)$ for $\mu=j$), each $\Omega_\mu$ has two terms:

\[
\begin{aligned}
\Omega_0&=\Omega_4=0,\\
\Omega_i&=\omega_{i\,i0}S^{i0}+\omega_{i\,i4}S^{i4}=-\alpha_i\bigl(S^{0i}+a_4'S^{4i}\bigr)\qquad(i=1,2,3),\\
\Omega_j&=\omega_{j\,j0}S^{j0}+\omega_{j\,j4}S^{j4}=\alpha_j\bigl(S^{0j}-a_4'S^{4j}\bigr)\qquad(j=5,6,7),
\end{aligned}
\]

using $S^{ba}=-S^{ab}$ (Stage-2 document, §7.1; check \texttt{P\_Omega\_\allowbreak{}closedFo\allowbreak{}rms}). Now contract with $\gamma^{x_i}=\gamma^i/h_i$ and $\gamma^{x_j}=\gamma^j/h_j$; note $\alpha_i/h_i=\alpha_j/h_j=H$. With $S^{0i}=\tfrac12\gamma^0\gamma^i$ and $\gamma^i\gamma^0\gamma^i=-\gamma^0(\gamma^i)^2=-\gamma^0$, $\gamma^i\gamma^4\gamma^i=-\gamma^4$ (since $(\gamma^i)^2=+1$), and, since $(\gamma^j)^2=-1$, $\gamma^j\gamma^0\gamma^j=\gamma^0$, $\gamma^j\gamma^4\gamma^j=\gamma^4$:

\[
\begin{aligned}
\gamma^{x_i}\Omega_i&=-\tfrac H2\bigl(\gamma^i\gamma^0\gamma^i+a_4'\gamma^i\gamma^4\gamma^i\bigr)=\tfrac H2\bigl(\gamma^0+a_4'\gamma^4\bigr),\\
\gamma^{x_j}\Omega_j&=\tfrac H2\bigl(\gamma^j\gamma^0\gamma^j-a_4'\gamma^j\gamma^4\gamma^j\bigr)=\tfrac H2\bigl(\gamma^0-a_4'\gamma^4\bigr).
\end{aligned}
\]

Summing over the three $i$ and the three $j$,

\[
\gamma^\mu\Omega_\mu=\tfrac{3H}2\bigl(\gamma^0+a_4'\gamma^4\bigr)+\tfrac{3H}2\bigl(\gamma^0-a_4'\gamma^4\bigr)=3H\gamma^0 .
\]

The free function $a_4$ cancels: the boosts of 3-space and the rotations of the extra times contribute opposite amounts. The general formula agrees: only $b=0$ contributes, because $\prod_{c\ne0}h_c=s^{1/2}e^{3a_4}\cdot1\cdot s^{1/2}e^{-3a_4}=s$ gives $\frac12\tan z\,\partial_0\ln s=\frac12\tan z\cdot6H\cot z=3H$, while $\prod_{c\ne4}h_c=\cot z\cdot s=\cos z$ does not depend on $x^4$. In words, $a_4$ cancels because the 7-volume of a comoving region is constant (Section 4.4; Stage-2 document, §8.1; checks \texttt{P\_Omega\_\allowbreak{}gammaSla\allowbreak{}sh3Hgamm\allowbreak{}a0} and \texttt{P\_Omega\_\allowbreak{}contract\allowbreak{}Diagonal\allowbreak{}Formula}; Stage-1 check \texttt{GEO\_prim\allowbreak{}ordialIn\allowbreak{}variants\allowbreak{}\_G2}). With $\{\gamma^\mu,\Omega_\mu\}=0$ we get $[\gamma^\mu,\Omega_\mu]=6H\gamma^0$, and the divergence identity can be checked by hand: $\partial_\mu(\sqrt{\lvert g\rvert}\gamma^\mu)=\partial_0(\cos z\tan z\,\gamma^0)=\partial_0(\sin z)\gamma^0=6H\cos z\,\gamma^0=\sqrt{\lvert g\rvert}\cdot6H\gamma^0$ (check \texttt{P\_gammaC\allowbreak{}onst\_div\allowbreak{}ergenceI\allowbreak{}dentity}).

\subsection{The curvature of the spinor connection}

\textbf{Definition.} The curvature of the spinor connection is the commutator of two spinor covariant derivatives. For a spinor field $\Psi$ (no curved index),

\[
D_\mu(D_\nu\Psi)-D_\nu(D_\mu\Psi)=F_{\mu\nu}\Psi,\qquad F_{\mu\nu}:=\partial_\mu\Omega_\nu-\partial_\nu\Omega_\mu+[\Omega_\mu,\Omega_\nu],
\]

where here $D_\mu$ acts on $D_\nu\Psi$ as on a spinor, $D_\mu(D_\nu\Psi)=\partial_\mu(D_\nu\Psi)+\Omega_\mu D_\nu\Psi$. (Expand: the terms $\partial_\mu\partial_\nu\Psi$ and $\Omega_\nu\partial_\mu\Psi+\Omega_\mu\partial_\nu\Psi$ are symmetric in $\mu,\nu$ and cancel; what remains is $(\partial_\mu\Omega_\nu)\Psi+\Omega_\mu\Omega_\nu\Psi$ minus the same with $\mu\leftrightarrow\nu$.)

\textbf{Lemma ($\Omega$ respects commutators).} For two antisymmetric arrays $A_{ab}$ and $B_{ab}$ put $\Omega(A)=\tfrac12A_{ab}S^{ab}$, and regard $A^a{}_b=\eta^{ac}A_{cb}$ as a matrix. Then

\[
[\Omega(A),\Omega(B)]=\Omega\bigl([A,B]\bigr),
\]

where $[A,B]$ is the matrix commutator, lowered with $\eta$.

\emph{Proof.} First, $[A,B]$ is again antisymmetric after lowering. In matrix form, "$\eta A$ is antisymmetric" means $A^T\eta=-\eta A$. Then $(AB)^T\eta=B^TA^T\eta=-B^T\eta A=\eta BA$, so $[A,B]^T\eta=\eta BA-\eta AB=-\eta[A,B]$. Second, by the theorem of Section 4.12, $[\Omega(A),\gamma^e]=-A^e{}_f\gamma^f$ for any such $A$. The Jacobi identity $[[X,Y],Z]=[X,[Y,Z]]-[Y,[X,Z]]$, which one checks by expanding all commutators, gives

\[
[[\Omega(A),\Omega(B)],\gamma^e]=[\Omega(A),-B^e{}_f\gamma^f]-[\Omega(B),-A^e{}_f\gamma^f]=B^e{}_fA^f{}_g\gamma^g-A^e{}_fB^f{}_g\gamma^g=-[A,B]^e{}_g\gamma^g .
\]

So $[\Omega(A),\Omega(B)]$ and $\Omega([A,B])$ have the same commutator with every $\gamma^e$. Their difference commutes with all $\gamma^e$ and is traceless (a commutator has zero trace, and so has every $S^{ab}$), hence it is zero by fact (F1). $\square$

\textbf{Lemma (frame curvature).} Define the curvature of the frame connection,

\[
R(\omega)^a{}_{b\mu\nu}:=\partial_\mu\omega_\nu{}^a{}_b-\partial_\nu\omega_\mu{}^a{}_b+\omega_\mu{}^a{}_c\,\omega_\nu{}^c{}_b-\omega_\nu{}^a{}_c\,\omega_\mu{}^c{}_b .
\]

Then $R(\omega)^a{}_{b\mu\nu}=e_\rho{}^a\,R^\rho{}_{\sigma\mu\nu}\,e_b{}^\sigma$: it is the Riemann tensor written with two frame indices.

\emph{Proof.} For a vector field with frame components $V^a$, let the total covariant derivative act with $\omega$ on frame indices and $\Gamma$ on curved indices: $\nabla_\nu V^a=\partial_\nu V^a+\omega_\nu{}^a{}_bV^b$ and $\nabla_\mu(\nabla_\nu V^a)=\partial_\mu(\nabla_\nu V^a)+\omega_\mu{}^a{}_c\nabla_\nu V^c-\Gamma^\lambda{}_{\mu\nu}\nabla_\lambda V^a$. Exactly as in the proof of the Riemann formula (Section 4.7), the commutator is $[\nabla_\mu,\nabla_\nu]V^a=R(\omega)^a{}_{b\mu\nu}V^b$: the $\Gamma$ term drops out by the symmetry of $\Gamma$, the second derivatives and the mixed first-derivative terms drop out by symmetry. On the other hand, frame components of covariant derivatives are covariant derivatives of frame components, and this stays true one step further. For $V$ itself it is the requirement that defined $\omega$ in Section 4.11: $\nabla_\nu V^a=e_\rho{}^a\nabla_\nu V^\rho$. The object $T^\rho{}_\nu:=\nabla_\nu V^\rho$ carries one extra lower curved index, which the total covariant derivative treats with $\Gamma$, both in $T^\rho{}_\nu$ and in its frame components $e_\rho{}^aT^\rho{}_\nu=\nabla_\nu V^a$. By the product rule,

\[
\begin{aligned}
\nabla_\mu\bigl(e_\rho{}^aT^\rho{}_\nu\bigr)&=\partial_\mu\bigl(e_\rho{}^aT^\rho{}_\nu\bigr)+\omega_\mu{}^a{}_b\,e_\rho{}^bT^\rho{}_\nu-\Gamma^\lambda{}_{\mu\nu}\,e_\rho{}^aT^\rho{}_\lambda\\
&=\bigl(\partial_\mu e_\rho{}^a+\omega_\mu{}^a{}_b\,e_\rho{}^b\bigr)T^\rho{}_\nu+e_\rho{}^a\partial_\mu T^\rho{}_\nu-\Gamma^\lambda{}_{\mu\nu}\,e_\rho{}^aT^\rho{}_\lambda .
\end{aligned}
\]

By the vielbein postulate $\partial_\mu e_\rho{}^a+\omega_\mu{}^a{}_b\,e_\rho{}^b=\Gamma^\lambda{}_{\mu\rho}\,e_\lambda{}^a$, so, after renaming summed indices, this equals

\[
e_\lambda{}^a\bigl(\partial_\mu T^\lambda{}_\nu+\Gamma^\lambda{}_{\mu\rho}T^\rho{}_\nu-\Gamma^\rho{}_{\mu\nu}T^\lambda{}_\rho\bigr)=e_\lambda{}^a\,\nabla_\mu T^\lambda{}_\nu .
\]

So $\nabla_\mu(\nabla_\nu V^a)=e_\lambda{}^a\nabla_\mu\nabla_\nu V^\lambda$, and hence $[\nabla_\mu,\nabla_\nu]V^a=e_\rho{}^a[\nabla_\mu,\nabla_\nu]V^\rho=e_\rho{}^aR^\rho{}_{\sigma\mu\nu}V^\sigma=e_\rho{}^aR^\rho{}_{\sigma\mu\nu}e_b{}^\sigma V^b$. Both results hold for every $V$. $\square$

\textbf{Theorem (spinor curvature).} $F_{\mu\nu}=\tfrac12\,R_{ab\mu\nu}\,S^{ab}$ with $R_{ab\mu\nu}:=\eta_{ac}\,e_\rho{}^c\,R^\rho{}_{\sigma\mu\nu}\,e_b{}^\sigma$.

\emph{Proof.} $\Omega_\mu=\Omega(\omega_\mu)$ is linear in $\omega$, so $\partial_\mu\Omega_\nu-\partial_\nu\Omega_\mu=\Omega(\partial_\mu\omega_\nu-\partial_\nu\omega_\mu)$, and by the first lemma $[\Omega_\mu,\Omega_\nu]=\Omega([\omega_\mu,\omega_\nu])$. So $F_{\mu\nu}=\Omega(R(\omega)_{\mu\nu})$, and the second lemma identifies $R(\omega)$. $\square$

The sign and the placement of $\eta$ matter. Stage 1 verified the form with $+\tfrac12$ and the lowered first index at all six Wolfram points and at the Python points; the forms with $-\tfrac12$ and without $\eta$ (mixed indices) fail there (Stage 1, Result 5.5; checks \texttt{GEO\_curv\allowbreak{}ature\_G1} and \texttt{GEO\_curv\allowbreak{}ature\_G2}; measurements \texttt{G1.p1.cu\allowbreak{}rvatureC\allowbreak{}andidate\allowbreak{}.plusHal\allowbreak{}fLowered} = true, \texttt{minusHal\allowbreak{}fLowered} = false and \texttt{plusHalf\allowbreak{}MixedNoE\allowbreak{}ta} = false in \texttt{wolfram-\allowbreak{}geometry\allowbreak{}-report.\allowbreak{}json}).

\textbf{Theorem (the spin connection cannot be removed in a curved field).} Near a point, the spinor connection is pure gauge, $\Omega_\mu=-(\partial_\mu R)R^{-1}$ for some spin transformation $R(x)$, if and only if the Riemann tensor vanishes there.

\emph{Proof.} If $\Omega_\mu=-(\partial_\mu R)R^{-1}$, then, with $\partial_\mu(R^{-1})=-R^{-1}(\partial_\mu R)R^{-1}$,

\[
\partial_\mu\Omega_\nu-\partial_\nu\Omega_\mu=(\partial_\nu R)R^{-1}(\partial_\mu R)R^{-1}-(\partial_\mu R)R^{-1}(\partial_\nu R)R^{-1}=-[\Omega_\mu,\Omega_\nu],
\]

because the second derivatives $\partial_\mu\partial_\nu R$ cancel. So $F_{\mu\nu}=0$. By the theorem above $R_{ab\mu\nu}S^{ab}=2\sum_{a<b}R_{ab\mu\nu}S^{ab}=0$; the 28 matrices $S^{ab}$ with $a<b$ are linearly independent (fact (F2)), so every $R_{ab\mu\nu}=0$, and since the vielbein is invertible the Riemann tensor vanishes. Conversely, if the Riemann tensor vanishes near the point, the flatness theorem quoted in Section 4.7 gives a chart with $g_{\mu\nu}=\eta_{\mu\nu}$; in it the frame $e_\mu{}^a=\delta_\mu{}^a$ has $\Gamma=0$, $\omega=0$ and $\Omega=0$. Any other frame is a rotation of this one, and by the covariance theorem of Section 4.12 its spinor connection is $0-(\partial_\mu R)R^{-1}$. (That a smooth frame rotation $\Lambda(x)$ has a smooth covering $R(x)$ near a point is a standard fact that we quote.) $\square$

Since the primordial field is curved everywhere (Section 4.7), its spin connection cannot be removed by any choice of frame (Stage-2 document, §7.3).

\textbf{The square of the Dirac operator.} The Dirac operator is $\gamma^\mu D_\mu$. For an object with one curved index and one spinor index, such as $D_\nu\Psi$, the full covariant derivative is $\nabla_\mu(D_\nu\Psi)=\partial_\mu(D_\nu\Psi)+\Omega_\mu D_\nu\Psi-\Gamma^\lambda{}_{\mu\nu}D_\lambda\Psi$.

\textbf{Theorem (Lichnerowicz formula).}

\[
(\gamma^\mu D_\mu)^2\Psi=g^{\mu\nu}\bigl(D_\mu D_\nu\Psi-\Gamma^\lambda{}_{\mu\nu}D_\lambda\Psi\bigr)-\tfrac14\,R\,\Psi .
\]

\emph{Proof.} Step 1. $(\gamma^\mu D_\mu)(\gamma^\nu D_\nu\Psi)=\gamma^\mu\bigl[\partial_\mu(\gamma^\nu D_\nu\Psi)+\Omega_\mu\gamma^\nu D_\nu\Psi\bigr]$. Replace $\partial_\mu\gamma^\nu$ by $-\Gamma^\nu{}_{\mu\lambda}\gamma^\lambda-[\Omega_\mu,\gamma^\nu]$ (covariant constancy of the gammas). The two terms $\mp\Omega_\mu\gamma^\nu D_\nu\Psi$ cancel, and

\[
(\gamma^\mu D_\mu)^2\Psi=\gamma^\mu\gamma^\nu\bigl(\partial_\mu D_\nu\Psi+\Omega_\mu D_\nu\Psi\bigr)-\gamma^\mu\Gamma^\nu{}_{\mu\lambda}\gamma^\lambda D_\nu\Psi=\gamma^\mu\gamma^\nu\,\nabla_\mu(D_\nu\Psi),
\]

after renaming $\lambda\leftrightarrow\nu$ in the last term. Step 2. Split $\gamma^\mu\gamma^\nu=g^{\mu\nu}+\tfrac12[\gamma^\mu,\gamma^\nu]$ (the curved Clifford relation). The first part, $g^{\mu\nu}\nabla_\mu(D_\nu\Psi)$, is the first term of the formula. The second part is $\tfrac12[\gamma^\mu,\gamma^\nu]\nabla_\mu(D_\nu\Psi)$. Since $[\gamma^\mu,\gamma^\nu]$ is antisymmetric in $\mu,\nu$, only the antisymmetric part of $\nabla_\mu(D_\nu\Psi)$ contributes, and that is $\tfrac12\bigl(\nabla_\mu(D_\nu\Psi)-\nabla_\nu(D_\mu\Psi)\bigr)=\tfrac12F_{\mu\nu}\Psi$ (the $\Gamma$ terms cancel because $\Gamma$ is symmetric). So the second part is $\tfrac14[\gamma^\mu,\gamma^\nu]F_{\mu\nu}\Psi=\tfrac12\gamma^\mu\gamma^\nu F_{\mu\nu}\Psi$, where the last step uses $\gamma^\nu\gamma^\mu F_{\mu\nu}=-\gamma^\nu\gamma^\mu F_{\nu\mu}=-\gamma^\mu\gamma^\nu F_{\mu\nu}$ (rename the summed indices). Step 3. In frame components, $\gamma^\mu\gamma^\nu F_{\mu\nu}=\gamma^c\gamma^dF_{cd}$ with $F_{cd}=\tfrac12R_{abcd}S^{ab}=\tfrac14R_{abcd}\gamma^a\gamma^b$ (the terms with $a=b$ vanish, and $S^{ab}=\tfrac12\gamma^a\gamma^b$ otherwise). The symmetries (S1) to (S4) hold equally for the frame components $R_{abcd}$, which are obtained from $R_{\rho\sigma\mu\nu}$ by multiplying with inverse vielbeins. By pair symmetry (S4), $R_{abcd}\gamma^c\gamma^d\gamma^a\gamma^b=R_{cdab}\gamma^c\gamma^d\gamma^a\gamma^b$, which after renaming the summed indices is $R_{abcd}\gamma^a\gamma^b\gamma^c\gamma^d$. Step 4. For three frame gammas,

\[
\gamma^b\gamma^c\gamma^d=\gamma^{[bcd]}+\eta^{bc}\gamma^d-\eta^{bd}\gamma^c+\eta^{cd}\gamma^b ,
\]

where $\gamma^{[bcd]}$ is totally antisymmetric in $b,c,d$ (it equals $\gamma^b\gamma^c\gamma^d$ when the three are different and is zero otherwise). To check it, try the five cases: all different; $b=c\ne d$; $b=d\ne c$; $c=d\ne b$; all equal. For example for $b=d\ne c$ the left side is $\gamma^b\gamma^c\gamma^b=-\eta^{bb}\gamma^c$ and the right side is $-\eta^{bb}\gamma^c$. Contract with $R_{abcd}$. The antisymmetric part gives zero by the first Bianchi identity (S3) in the last three indices, which for frame components reads $R_{abcd}+R_{acdb}+R_{adbc}=0$. Because $\gamma^{[bcd]}$ is unchanged by a cyclic permutation of $(b,c,d)$ (a cyclic permutation of three indices is an even rearrangement), renaming the summed indices gives $R_{abcd}\gamma^{[bcd]}=R_{acdb}\gamma^{[bcd]}=R_{adbc}\gamma^{[bcd]}$; for example, in $R_{acdb}\gamma^{[bcd]}$ rename $c\to b$, $d\to c$, $b\to d$ to get $R_{abcd}\gamma^{[dbc]}=R_{abcd}\gamma^{[bcd]}$. Hence $3R_{abcd}\gamma^{[bcd]}=(R_{abcd}+R_{acdb}+R_{adbc})\gamma^{[bcd]}=0$. The three other terms give $\eta^{bc}R_{abcd}\gamma^d=-R_{ad}\gamma^d$, $-\eta^{bd}R_{abcd}\gamma^c=-R_{ac}\gamma^c$ and $\eta^{cd}R_{abcd}\gamma^b=0$, where we used (S1), (S2) and the frame Ricci tensor $R_{bd}=\eta^{ac}R_{abcd}$ with its symmetry. Hence $R_{abcd}\gamma^b\gamma^c\gamma^d=-2R_{ad}\gamma^d$ and

\[
R_{abcd}\gamma^a\gamma^b\gamma^c\gamma^d=-2R_{ad}\gamma^a\gamma^d=-2R_{ad}\,\eta^{ad}=-2R ,
\]

because $R_{ad}$ is symmetric and $\tfrac12(\gamma^a\gamma^d+\gamma^d\gamma^a)=\eta^{ad}$. Step 5. The second part of Step 2 is $\tfrac12\cdot\tfrac14\cdot(-2R)\Psi=-\tfrac14R\Psi$. $\square$

The constant $-\tfrac14$ is exactly the value that Stage 1 measured at all its test points: \texttt{lichnero\allowbreak{}wiczC} = -1/4 at the three G1 points and the three Wolfram G2 points, with $+\tfrac14$ failing (Stage 1, Result 8.2; checks \texttt{GEO\_lich\allowbreak{}nerowicz\allowbreak{}\_G1}, \texttt{GEO\_lich\allowbreak{}nerowicz\allowbreak{}\_G2} and \texttt{GEO\_lich\allowbreak{}nerowicz\allowbreak{}Constant\allowbreak{}SameG1G2}). In the primordial field, with $R=6H^2(a_4'^2-7)$, the curvature term is $-\tfrac R4\Psi=\tfrac32H^2(7-a_4'^2)\Psi$, which is $9H^2\Psi$ for $a_4=t$ (Stage-2 document, §7.3). Because of this term, curvature enters the second-order form of the field equations explicitly (Chapter 7).

\subsection{Why the notebook's contraction is wrong}

\textbf{What the notebook does.} The notebook's Lagrangian Lg[] (cell 1064, In[1034]) is quoted verbatim in the Stage-1 document (§6.1; line breaks only for the page width):

\begin{Verbatim}[fontsize=\small]
Lg[]:=Sqrt[detgg] *( Transpose[\[CapitalPsi]16].\[Sigma]16.
Sum[FullSimplify[((T16^\[Alpha])[\[Alpha]1-1]/.sg),constraintVars].
(D[ \[CapitalPsi]16,X[[\[Alpha]1]]]+(Q1/2)*Sum[\[Omega]mat[[\[Alpha]1,a,b]]*SAB[[a,
b]].\[CapitalPsi]16,{a,1,8},{b,1,8}]),{\[Alpha]1,1,Length[X]}]+
(H*M)*Transpose[\[CapitalPsi]16].\[Sigma]16.\[CapitalPsi]16)//Simplify[#,
constraintVars]&
\end{Verbatim}

Mathematica lists count from 1, so \texttt{ωmat[[μ+\allowbreak{}1,a+1,b+\allowbreak{}1]]} is our $\omega_\mu{}^a{}_b$, the mixed spin connection with the first frame index up, and \texttt{SAB[[a+1\allowbreak{},b+1]]} is $S^{ab}$, with both indices up. \texttt{\detokenize{Q1}} (written $Q_1$ below) is the notebook's book-keeping switch for the spin connection; the notebook keeps it as a symbol (cell 1063; \texttt{handoff/\allowbreak{}surveys/\allowbreak{}survey\_n\allowbreak{}otebook-\allowbreak{}physics.\allowbreak{}md}, item 4), and setting it to 1, as Stage 1 does, switches the connection on. So, with $Q_1=1$, the notebook's spinor connection is

\[
\Omega^{\mathrm{nb}}_\mu=\tfrac12\,\omega_\mu{}^a{}_b\,S^{ab}\qquad\text{instead of}\qquad \Omega_\mu=\tfrac12\,\omega_{\mu ab}\,S^{ab}.
\]

The factor $\tfrac12$ is right. What is wrong is the index placement: the contraction pairs the upper index $a$ of $\omega$ with the upper index $a$ of $S^{ab}$. One $\eta$ is missing. Section 4.3 warned that summing two upper indices does not give a frame-independent result; here is what it does.

\textbf{What the missing $\eta$ deletes.} Since $\omega_\mu{}^a{}_b=\eta^{aa}\omega_{\mu ab}$ (no sum, $\eta^{aa}=\pm1$),

\[
\Omega^{\mathrm{nb}}_\mu=\tfrac12\sum_{a,b}\eta^{aa}\,\omega_{\mu ab}\,S^{ab}.
\]

Sort the pairs $(a,b)$ into three kinds.

\begin{itemize}
\item Both space-like ($a,b\in\{0,1,2,3\}$): $\eta^{aa}=+1$, the term is the same as in $\Omega_\mu$.
\item Both time-like ($a,b\in\{4,5,6,7\}$): $\eta^{aa}=-1$, the term has the wrong sign.
\item One space-like and one time-like (a boost pair): the two orders give $\eta^{aa}\omega_{\mu ab}S^{ab}+\eta^{bb}\omega_{\mu ba}S^{ba}=(\eta^{aa}+\eta^{bb})\,\omega_{\mu ab}S^{ab}=0$, because $\omega_{\mu ba}S^{ba}=\omega_{\mu ab}S^{ab}$ and $\eta^{aa}+\eta^{bb}=0$.
\end{itemize}

So, writing $\Omega_\mu=\Omega^{\mathrm{ss}}_\mu+\Omega^{\mathrm{tt}}_\mu+\Omega^{\mathrm{st}}_\mu$ for the space-space, time-time and boost parts,

\[
\Omega^{\mathrm{nb}}_\mu=\Omega^{\mathrm{ss}}_\mu-\Omega^{\mathrm{tt}}_\mu,\qquad \Omega^{\mathrm{nb}}_\mu-\Omega_\mu=-\Omega^{\mathrm{st}}_\mu-2\,\Omega^{\mathrm{tt}}_\mu .
\]

Of the 28 pairs $a<b$ in signature (4,4), 6 are space-space, 6 are time-time and $4\times4=16$ are boost pairs: for each $\mu$, the notebook's contraction deletes 16 of the 28 independent components of the connection and reverses the sign of 6 more. (This is the index argument of Stage 1, §5.6, and of \texttt{handoff/\allowbreak{}specs/CO\allowbreak{}NTRACT.m\allowbreak{}d}, §3, item 4.) In a positive-definite space, such as the polar plane or the sphere, all $\eta^{aa}=+1$ and the two contractions coincide. The mistake is therefore invisible in every textbook example with a positive metric; it appears only with an indefinite metric.

\textbf{Consequence: the gammas are no longer covariantly constant.} With $\Omega^{\mathrm{nb}}$ in place of $\Omega$ the covariant derivative of the gammas is, by the theorem of Section 4.12,

\[
D^{\mathrm{nb}}_\mu\gamma^\nu=\partial_\mu\gamma^\nu+\Gamma^\nu{}_{\mu\lambda}\gamma^\lambda+[\Omega^{\mathrm{nb}}_\mu,\gamma^\nu]=[\Omega^{\mathrm{nb}}_\mu-\Omega_\mu,\gamma^\nu].
\]

The difference $X_\mu=\Omega^{\mathrm{nb}}_\mu-\Omega_\mu$ is a combination of the $S^{ab}$, hence traceless. If it commuted with all $\gamma^\nu$, it would be a multiple of the identity by fact (F1), hence zero. So $D^{\mathrm{nb}}_\mu\gamma^\nu=0$ for all $\nu$ holds only if the boost and time-time parts of $\omega_\mu$ vanish, which is not the case in a generic field. Everything that the book derives from $D_\mu\gamma^\nu=0$, or from the defining property $[\Omega_\mu,\gamma^a]=-\omega_\mu{}^a{}_b\gamma^b$, then fails: the divergence identity, the Lichnerowicz formula, local spin covariance, and the derivation of the field equations in Chapter 7.

\textbf{Worked example: the primordial field.} From Section 4.12, $\Omega_i=-\alpha_i(S^{0i}+a_4'S^{4i})$ and $\Omega_j=\alpha_j(S^{0j}-a_4'S^{4j})$. The pair $(0,i)$ is space-space, $(4,i)$ is a boost, $(0,j)$ is a boost and $(4,j)$ is time-time. So

\[
\Omega^{\mathrm{nb}}_i=-\alpha_i\,S^{0i},\qquad \Omega^{\mathrm{nb}}_j=+\alpha_j\,a_4'\,S^{4j}.
\]

Repeat the computation of Section 4.12 with these: $\gamma^{x_i}\Omega^{\mathrm{nb}}_i=-\tfrac H2\gamma^i\gamma^0\gamma^i=\tfrac H2\gamma^0$ and $\gamma^{x_j}\Omega^{\mathrm{nb}}_j=\tfrac{Ha_4'}{2}\gamma^j\gamma^4\gamma^j=\tfrac{Ha_4'}2\gamma^4$. Summing,

\[
\gamma^\mu\Omega^{\mathrm{nb}}_\mu=\tfrac{3H}{2}\bigl(\gamma^0+a_4'\gamma^4\bigr)\qquad\text{instead of the correct}\qquad \gamma^\mu\Omega_\mu=3H\gamma^0 .
\]

Half of the $\gamma^0$ term is lost, and the $a_4'$ terms no longer cancel. \emph{Numbers.} At the three exact G2 test points of Stage 1, $(H,a_4')=(2/3,3/7)$, $(1/5,-4/9)$ and $(5/7,6/5)$, this gives $\gamma^0+\tfrac37\gamma^4$, $\tfrac3{10}\gamma^0-\tfrac2{15}\gamma^4$ and $\tfrac{15}{14}\gamma^0+\tfrac97\gamma^4$, exactly the values recorded as \texttt{G2.p1.no\allowbreak{}tebookSl\allowbreak{}ash}, \texttt{G2.p2.no\allowbreak{}tebookSl\allowbreak{}ash} and \texttt{G2.p3.no\allowbreak{}tebookSl\allowbreak{}ash} in \texttt{wolfram-\allowbreak{}geometry\allowbreak{}-report.\allowbreak{}json}; the correct values are $2\gamma^0$, $\tfrac35\gamma^0$ and $\tfrac{15}7\gamma^0$.

\textbf{Worked example: one component of $D_\mu\gamma^\nu$.} Take $\mu=1$, $\nu=4$. Since $\gamma^{x_4}=\gamma^4$ is constant, $\partial_1\gamma^{x_4}=0$. The only nonzero $\Gamma^4{}_{1\lambda}$ is $\Gamma^4{}_{11}=Ha_4's^{1/3}e^{2a_4}$, so $\Gamma^4{}_{1\lambda}\gamma^{x_\lambda}=Ha_4's^{1/3}e^{2a_4}\cdot s^{-1/6}e^{-a_4}\gamma^1=\alpha_1a_4'\gamma^1$. For the commutator we need $[S^{01},\gamma^4]=0$ ($\gamma^4$ anticommutes with both $\gamma^0$ and $\gamma^1$, so it commutes with their product) and $[S^{41},\gamma^4]=\eta^{14}\gamma^4-\eta^{44}\gamma^1=\gamma^1$ (the identity of Section 4.12). With the correct connection,

\[
D_1\gamma^{x_4}=\alpha_1a_4'\gamma^1-\alpha_1\bigl([S^{01},\gamma^4]+a_4'[S^{41},\gamma^4]\bigr)=\alpha_1a_4'\gamma^1-\alpha_1a_4'\gamma^1=0 .
\]

With the notebook's connection the boost term is missing, and

\[
D^{\mathrm{nb}}_1\gamma^{x_4}=\alpha_1a_4'\gamma^1-\alpha_1[S^{01},\gamma^4]=H\,s^{1/6}e^{a_4}a_4'\,\gamma^1\ne0 ,
\]

which is the entry "$(i,4)$" of the table in the Stage-2 document, §8.3 (check \texttt{P\_gammaC\allowbreak{}onst\_not\allowbreak{}ebookCon\allowbreak{}traction\allowbreak{}ClosedFo\allowbreak{}rms}).

\textbf{What the verifiers measured.} In the primordial field the notebook's contraction violates $D_\mu\gamma^\nu=0$ in 15 of the 64 pairs $(\mu,\nu)$, with 288 nonzero matrix entries; in the generic test geometry G1 it violates it in all 64 pairs, with 4096 nonzero entries. The largest violation is $2/3$, $1/5$ and $9/5$ at the three G2 points and approximately 1.46822, 1.52979 and 1.68937 at the three G1 points. The symmetric part of $\omega_\mu{}^a{}_b$ that the contraction deletes has 12 nonzero entries in G2 and 256 at each G1 point (in G1 every boost pair is populated: 8 values of $\mu$, 16 pairs, 2 orders). Replacing $\Omega$ by $\Omega^{\mathrm{nb}}$ makes all eleven connection-sensitive G2 checks fail, and it breaks local spin invariance already at first order. (Stage 1, Result 5.6; checks \texttt{GEO\_note\allowbreak{}bookCont\allowbreak{}ractionF\allowbreak{}ails\_G1}, \texttt{GEO\_note\allowbreak{}bookCont\allowbreak{}ractionF\allowbreak{}ails\_G2} and \texttt{NEG\_note\allowbreak{}bookConn\allowbreak{}ectionDe\allowbreak{}tected\_G\allowbreak{}2}; Stage-2 checks \texttt{P\_gammaC\allowbreak{}onst\_not\allowbreak{}ebookCon\allowbreak{}traction\allowbreak{}Fails} and \texttt{P\_gammaC\allowbreak{}onst\_div\allowbreak{}ergenceI\allowbreak{}dentityF\allowbreak{}ailsNote\allowbreak{}bookCont\allowbreak{}raction}; measurements \texttt{notebook\allowbreak{}DGammaMa\allowbreak{}xAbs}, \texttt{notebook\allowbreak{}DGammaNo\allowbreak{}nzeroEnt\allowbreak{}ries} and \texttt{nonzeroO\allowbreak{}megaMixe\allowbreak{}dSymmetr\allowbreak{}icPart} in \texttt{wolfram-\allowbreak{}geometry\allowbreak{}-report.\allowbreak{}json}.)

\textbf{The repair.} Lower the index first, $\omega_{\mu ab}=\eta_{ac}\omega_\mu{}^c{}_b$, and contract that with $S^{ab}$. In WolframScript, with the notebook's own names (Stage-1 document, §7.2):

\begin{Verbatim}[fontsize=\small]
(* lowered spin connection omega_{mu ab} = eta_{ac} omega_mu^c_b *)
\[Omega]low = Table[Sum[\[Eta]4488[[a, c]] \[Omega]mat[[\[Alpha]1, c, b]], {c, 1, 8}],
   {\[Alpha]1, 1, 8}, {a, 1, 8}, {b, 1, 8}];
\[CapitalOmega]16[\[Alpha]1_] := (1/2) Sum[\[Omega]low[[\[Alpha]1, a, b]] SAB[[a, b]],
   {a, 1, 8}, {b, 1, 8}];
\end{Verbatim}

The notebook has two further problems, separate from this one. First, for an anticommuting (Grassmann) field Lg[] carries no dynamics at all, whichever contraction is used (Chapter 5). Second, the extra term $q$ in the field equations that the notebook stores (from its evaluation form La[] with commuting fields, cell 1066) is not caused by the contraction error: with curved gammas that obey the Clifford relation, every $Q_1$ term drops out of those equations, for the notebook's contraction and for the correct one alike (Chapter 9). Chapter 9 shows that the stored $q$ term is reproduced exactly by extra-time curved gammas that violate the Clifford relation; it enters through the spin-connection term, multiplied by the switch $Q_1$. Chapter 9 attributes these non-Clifford gammas to a wrong substitution rule of cell 1058; that attribution is a reconstruction, labelled a HYPOTHESIS in ledger row L20 of Chapter 0, because the notebook does not store the gammas it used and a literal re-execution of cell 1058 gives no $q$ term.

\subsection{Where the computations live in the repository}

\textbf{The objects and their checks.} The formulas of Sections 4.4 to 4.14 that the project uses in eight dimensions are verified by exact computer algebra, most of them twice and independently: in Wolfram Language and in Python. (The two-dimensional examples and the Lovelock tensors of orders 2 and 3 are not part of any verifier.) The reports are JSON files; each check is a named entry that is either true or false, and each measurement a named value.

In the list below, "Stage 1" names checks of the arbitrary-field geometry reports and "Stage 2" checks of \texttt{wolfram-\allowbreak{}primordi\allowbreak{}al-repor\allowbreak{}t.json}; the numbers in parentheses are the sections of this chapter.

\begin{Verbatim}[fontsize=\small]
metric from the vielbein, det g (4.4, 4.10)
    Stage 1: GEO_frameNondegenerate_G1, GEO_sqrtgSquaredEqualsDetg_G1
    Stage 2: P_metric_vielbeinProduct, P_metric_detG_equals_plus_cos2z
Christoffel symbols (4.5)
    Stage 2: P_christoffel_count, P_christoffel_closedForms512
Riemann tensor, Ricci tensor, R, Einstein tensor (4.7, 4.8)
    Stage 1: GEO_curvature_G1, GEO_curvature_G2
    Stage 2: P_einstein_ricciScalar, P_einstein_GmixedClosedForms, P_einstein_R44
vielbein postulate, spin connection (4.11)
    Stage 1: GEO_vielbeinPostulate_G1, GEO_vielbeinPostulate_G2,
             GEO_omegaAntisymmetry_G1, GEO_omegaAntisymmetry_G2
    Stage 2: P_spinconn_vielbeinPostulate512, P_spinconn_count24,
             P_spinconn_closedForms
spinor connection, gamma^mu Omega_mu (4.12)
    Stage 1: GEO_diagonalSlashFormula_G2, GEO_primordialInvariants_G2
    Stage 2: P_Omega_closedForms, P_Omega_gammaSlash3Hgamma0
covariant constancy of the gammas, divergence identity (4.12)
    Stage 1: GEO_gammaCovariantConstancy_G1, GEO_gammaCovariantConstancy_G2,
             GEO_divergenceIdentity_G1, GEO_divergenceIdentity_G2
    Stage 2: P_gammaConst_DmuGammaNuZero64, P_gammaConst_divergenceIdentity
local spin covariance (4.12)
    Stage 1: LAG_localSpinInvariance_G1, LAG_localSpinInvariance_G2
spinor curvature, Lichnerowicz formula (4.13)
    Stage 1: GEO_curvature_G1, GEO_curvature_G2, GEO_lichnerowicz_G1,
             GEO_lichnerowicz_G2, GEO_lichnerowiczConstantSameG1G2
the notebook's contraction (4.14)
    Stage 1: GEO_notebookContractionFails_G1, GEO_notebookContractionFails_G2,
             NEG_notebookConnectionDetected_G2
    Stage 2: P_gammaConst_notebookContractionFails,
             P_gammaConst_notebookContractionClosedForms
\end{Verbatim}

The reports, with the number of true checks (Stage-1 document, §11; Stage-2 document, abstract), and the programs that write them:

\begin{Verbatim}[fontsize=\small]
artifacts/dirac16complex/arbitrary-field/
    wolfram-geometry-report.json     43 of 43 true
    python-geometry-report.json      52 of 52 true
artifacts/dirac16complex/primordial-field/
    wolfram-primordial-report.json   126 of 126 true
    python-primordial-report.json    16 of 16 true

Stage 1, Wolfram:  wolfram/Dirac16ComplexGeometry.wl
                   scripts/verify_dirac16complex_geometry.wls
Stage 1, Python:   scripts/d16c_geometry_sympy.py
                   scripts/check_dirac16complex_geometry.py
Stage 2, Wolfram:  wolfram/Dirac16ComplexPrimordial.wl
                   scripts/verify_dirac16complex_primordial.wls
Stage 2, Python:   scripts/check_dirac16complex_primordial.py
\end{Verbatim}

Chapter 19 gives the commands that rerun them and the lines they print.

\textbf{How the Python checker computes the geometry.} The class Geometry in \texttt{scripts/\allowbreak{}d16c\_geo\allowbreak{}metry\_sy\allowbreak{}mpy.py} follows Sections 4.4 to 4.12 line by line. Its arrays are exact "jets": the value and the derivatives of every quantity at one point. The function jein(spec, A, B) is a sum over repeated indices, written in the letter notation of numpy's einsum (for example "rs,smn->rmn" means $\sum_s A_{rs}B_{smn}$). An excerpt of lines 705-735, without the indentation, without three comment lines and without the fourteen lines 719-732 of the anholonomy route mentioned below, and with one comment shortened:

\begin{Verbatim}[fontsize=\small]
self.einv = jinv(e, dom)                                    # [a, mu] = e_a^mu
self.g = jein("mb,nb->mn", jein("ma,ab->mb", e, ETA), e)   # g = e eta e^T
self.ginv = jinv(self.g, dom)
dg = self.g.grad()                                           # [l, m, n] = d_l g_mn
t1 = dg.map(lambda a: np.einsum("mns->smn", a))              # d_m g_ns
t2 = dg.map(lambda a: np.einsum("nms->smn", a))              # d_n g_ms
gam1 = (t1 + t2 - dg).scale(half)                            # Gamma_{s m n}
self.Gamma = jein("rs,smn->rmn", self.ginv, gam1)           # Gamma^r_{mn}
self.de = e.grad()                                           # [m, n, a] = d_m e_n^a
inner = jein("rmn,ra->mna", self.Gamma, e) - self.de
self.omega_mixed = jein("bn,mna->mab", self.einv, inner)    # omega_mu^a_b
self.omega_low = jein("ac,mcb->mab", ETA, self.omega_mixed)  # omega_{mu a b}
self.Omega = jein("mab,abij->mij", self.omega_low, gd.S).scale(half)
self.Omega_nb = jein("mab,abij->mij", self.omega_mixed, gd.S).scale(half)
\end{Verbatim}

Read it against the formulas: gam1 is $\Gamma_{\sigma\mu\nu}=\tfrac12(\partial_\mu g_{\nu\sigma}+\partial_\nu g_{\mu\sigma}-\partial_\sigma g_{\mu\nu})$, self.Gamma raises its first index, inner is $\Gamma^\rho{}_{\mu\nu}e_\rho{}^a-\partial_\mu e_\nu{}^a$, and self.omega\_mixed multiplies by $e_b{}^\nu$: the solution of the vielbein postulate. The last two lines are the correct contraction, with the lowered $\omega_{\mu ab}$, and the notebook's contraction, with the mixed $\omega_\mu{}^a{}_b$; gd.S holds the 64 matrices $S^{ab}$. The same class also computes $\omega$ by a second, independent route from the derivatives of the vielbein alone (the anholonomy route) and compares the two.

\textbf{Try it yourself.} The following short program (not part of the repository; it needs Python with sympy, see Chapter 19) computes the Christoffel symbols and the scalar curvature of the primordial metric from the formulas of Sections 4.5 and 4.7. Save it as \texttt{curvatur\allowbreak{}e.py} and run \texttt{python c\allowbreak{}urvature\allowbreak{}.py}; it takes a few seconds.

\begin{Verbatim}[fontsize=\small]
import sympy as sp

H = sp.symbols("H", positive=True)
x = sp.symbols("x0:8")
a4 = sp.Function("a4")(x[4])        # a4 as a function of x4
z = 6 * H * x[0]
s = sp.sin(z)
third = sp.Rational(1, 3)
g = sp.diag(sp.cot(z)**2, *[s**third * sp.exp(2 * a4)] * 3, -1,
            *[-s**third * sp.exp(-2 * a4)] * 3)
n = 8
gi = g.inv()


def christoffel(r, m, b):           # Gamma^r_{m b}
    total = 0
    for q in range(n):
        total += gi[r, q] * (sp.diff(g[b, q], x[m]) + sp.diff(g[m, q], x[b])
                             - sp.diff(g[m, b], x[q]))
    return sp.simplify(total / 2)


Gam = [[[christoffel(r, m, b) for b in range(n)] for m in range(n)]
       for r in range(n)]
count = sum(1 for r in range(n) for m in range(n) for b in range(n)
            if Gam[r][m][b] != 0)
print("nonzero Christoffel symbols:", count)


def riemann(r, q, m, b):            # R^r_{q m b}
    val = sp.diff(Gam[r][b][q], x[m]) - sp.diff(Gam[r][m][q], x[b])
    val += sum(Gam[r][m][l] * Gam[l][b][q] - Gam[r][b][l] * Gam[l][m][q]
               for l in range(n))
    return val


ricci = [sum(riemann(r, q, r, q) for r in range(n)) for q in range(n)]
R = sum(gi[q, q] * ricci[q] for q in range(n))     # g is diagonal
A1 = sp.Symbol("A1")                # A1 stands for d a4 / d t
zs = sp.Symbol("z", positive=True)
R = R.subs(sp.Derivative(a4, x[4]), H * A1).subs(x[0], zs / (6 * H))
R = sp.simplify(sp.expand_trig(sp.simplify(R)))
print("R =", sp.factor(R))
\end{Verbatim}

It prints

\begin{Verbatim}[fontsize=\small]
nonzero Christoffel symbols: 37
R = 6*H**2*(A1**2 - 7)
\end{Verbatim}

that is, 37 nonzero symbols and $R=6H^2(a_4'^2-7)$ (the sum for R uses only the diagonal of $g^{-1}$ because the metric is diagonal). These are the numbers of Sections 4.5 and 4.7; the authoritative values are the ones in the committed reports cited there.

\subsection{What we proved and what we assumed}

\textbf{Proved in this chapter} (complete derivations): the transformation rules of vectors, covectors, tensors and the metric, and the invariance of $\sqrt{\lvert g\rvert}\,d^nx$; $\det g=+\cos^2z$ and the constant 7-volume of a comoving region of the primordial field; the transformation rule of the connection coefficients; existence and uniqueness of the Levi-Civita connection (the Christoffel formula); the derivative of a determinant and $\Gamma^\rho{}_{\rho\nu}=\partial_\nu\ln\sqrt{\lvert g\rvert}$; the Christoffel symbols of diagonal metrics and the 37 symbols of the primordial field; the constancy of $g(u,u)$ along geodesics and the complete solution of the geodesic equation in polar coordinates (all values of the constant $L$); the Riemann formula, the symmetries (S1) to (S4), the symmetry of the Ricci tensor, the curvature of the sphere, $R_{44}=-6H^2a_4'^2$ and $G^4{}_4=3H^2(7+a_4'^2)$ in the primordial field (using the quoted $R$); both Bianchi identities and $\nabla_\mu G^\mu{}_\nu=0$; the trace-reversed Einstein equations; $E^{(1)}=G$ and the vanishing of $E^{(k)}$ for $2k+1>n$; the vielbein relations and $\det g=(\det e)^2$ in signature (4,4); the vielbein postulate as a consequence of a natural requirement, its unique solution and the antisymmetry of $\omega_{\mu ab}$; the transformation of $\omega$ under frame rotations; the 24 components of $\omega$ in the primordial field; the spinor connection $\Omega_\mu=\tfrac12\omega_{\mu ab}S^{ab}$ as the only traceless solution of its defining property; local spin covariance; $D_\mu\gamma^\nu=0$ and its equivalence with the vielbein postulate; the divergence identity; the formula for $\gamma^\mu\Omega_\mu$ with a diagonal vielbein and $\gamma^\mu\Omega_\mu=3H\gamma^0$ in the primordial field; $F_{\mu\nu}=\tfrac12R_{ab\mu\nu}S^{ab}$; the Lichnerowicz formula with the constant $-\tfrac14$; and the analysis of the notebook's contraction (it deletes the 16 boost components, reverses the 6 time-time components, and gives $\gamma^\mu\Omega^{\mathrm{nb}}_\mu=\tfrac{3H}2(\gamma^0+a_4'\gamma^4)$ and $D^{\mathrm{nb}}_1\gamma^{x_4}\ne0$ in the primordial field).

\textbf{Computed by the verifiers and quoted here:} the full scalar curvature $R=6H^2(a_4'^2-7)$ and the other components of the Einstein tensor of the primordial field, and all numbers at the test points (counts of nonzero components, curvature values, sizes of the violations), each with its report and check.

\textbf{Assumed or quoted without proof:} that spacetime is a smooth manifold with a smooth metric of signature (4,4) and admits a spin structure (Stage 1, §4.5, assumes the latter); the standard theorems we named (Sylvester's law of inertia, the change-of-variables rule for integrals, the inverse function theorem, the flatness theorem, the local existence of orthonormal frames and of smooth spin coverings of frame rotations, Newton's limit of Einstein's equations, the Gauss-Bonnet expansion, and Lovelock's theorem, together with the fact that $E^{(2)}$ and $E^{(3)}$ do not vanish identically in eight dimensions); the facts (F1) to (F3) about the gamma matrices and the determinant 1 of the elements of Pin(4,4) as 16 by 16 matrices, which Chapter 2 proves; and two choices of the project: the connection is the torsion-free Levi-Civita connection, and the law of gravity is the 8-dimensional Einstein equation $G^\mu{}_\nu=\kappa T^\mu{}_\nu$ without the Lovelock terms of order 2 and 3, which neither the notebook nor the project computes. The Stage-1 test-point checks do not replace the general proofs: they test them exactly at three points of one generic vielbein, in the primordial family and in homogeneous frames, not symbolically for every vielbein (Stage-1 document, §13, limitation 1).

\subsection{Exercises}

\textbf{Exercise 4.1.} \emph{The polar plane.} (a) Write down $g^{\mu\nu}$ for $ds^2=dr^2+r^2d\varphi^2$. (b) The length of a curve is $\int\sqrt{g_{\mu\nu}u^\mu u^\nu}\,d\lambda$. Compute the length of the circle $r=2$, $\varphi=\lambda$, $0\le\lambda\le2\pi$. (c) Lower the index of the constant field $V^r=\cos\varphi$, $V^\varphi=-\sin\varphi/r$ and check that $g(V,V)=1$.

\textbf{Exercise 4.2.} \emph{Geodesics of the sphere.} (a) With the Christoffel symbols of the sphere (Section 4.7), write the two geodesic equations. (b) Show that the equator, $\theta=\pi/2$, $\varphi=\lambda/a$, is a geodesic, and compute $g(u,u)$. (c) Show that the circle of latitude $\theta=\pi/3$, traversed as $\varphi=\lambda/(a\sin(\pi/3))$, is not a geodesic.

\textbf{Exercise 4.3.} \emph{Freely falling observers in the primordial field.} Show that for every constant $c^0,\dots,c^3,c^5,\dots,c^7$ and every function $a_4$, the curve $x^\mu(\lambda)=(c^0,c^1,c^2,c^3,\lambda,c^5,c^6,c^7)$ is a geodesic of the primordial field, and compute $g(u,u)$.

\textbf{Exercise 4.4.} \emph{A space of constant negative curvature.} Take the two-dimensional metric $ds^2=dy^2+e^{2Hy}dw^2$ with coordinates $(y,w)$ and a constant $H>0$. Compute its Christoffel symbols, $R^y{}_{wyw}$, the Ricci tensor and the scalar curvature.

\textbf{Exercise 4.5.} \emph{The static primordial field.} Put $a_4$ equal to a constant. Using the formulas of Sections 4.7 and 4.8, compute $R$, $G^\mu{}_\nu$, the required energy density $\rho_{\mathrm{req}}$, the required pressures and $w=p/\rho$.

\textbf{Exercise 4.6.} \emph{Einstein's tensor in two dimensions.} (a) Show that $G_{\mu\nu}=0$ for the sphere. (b) Explain why $G_{\mu\nu}=0$ for every two-dimensional metric.

\textbf{Exercise 4.7.} \emph{Lovelock orders.} For $n=4,5,6,8,10$, list the orders $k$ that can contribute, that is, the orders for which $E^{(k)}$ is not identically zero. Compare with the notebook's rule "m-1 = n/2 - 1" for the highest order in even dimension $n$.

\textbf{Exercise 4.8.} \emph{A spinor in the polar frame.} In the plane, use $\gamma^0=\sigma_1$, $\gamma^1=\sigma_2$ in every orthonormal frame. In the Cartesian frame $\Omega=0$. The polar frame is the Cartesian frame turned by the angle $\varphi$. Let $R(\varphi)=\exp(i\varphi\sigma_3/2)$. (a) Show that $R=\mathrm{diag}(e^{i\varphi/2},e^{-i\varphi/2})$ and $R\,(\cos\varphi\,\sigma_1+\sin\varphi\,\sigma_2)\,R^{-1}=\sigma_1$. (The left side contains the curved gamma $\gamma^r$ built with the Cartesian frame.) (b) Show that $-(\partial_\varphi R)R^{-1}$ is the spinor connection $\Omega_\varphi$ of Section 4.12. (c) Compute $R(2\pi)$ and interpret it. (d) Where does the term $1/(2r)$ in the polar Dirac operator come from?

\textbf{Exercise 4.9.} \emph{Homogeneous frames.} For $ds^2=-dt^2+\sum_{j\ne4}\eta_{jj}h_j(t)^2(dx^j)^2$ with $t=x^4$, show that $\gamma^\mu\Omega_\mu=\tfrac12\Theta\gamma^4$ with $\Theta=\sum_{j\ne4}\dot h_j/h_j$, where a dot is $d/dt$.

\textbf{Exercise 4.10.} \emph{Counting what the notebook deletes.} (a) In signature (4,4), how many of the 28 pairs $a<b$ are space-space, time-time and boost pairs? (b) The same question in four dimensions with three space directions and one time direction. (c) For a generic vielbein in eight dimensions, how many nonzero entries does the symmetric part of $\omega_\mu{}^a{}_b$ have, and how many nonzero components does $\omega_{\mu ab}$ have? Compare with the G1 measurements of Section 4.14.

\textbf{Exercise 4.11.} \emph{The notebook's contraction at a test point.} At the first G2 test point of Stage 1, $H=2/3$ and $a_4'=3/7$. Compute $\gamma^\mu\Omega_\mu$, $\gamma^\mu\Omega^{\mathrm{nb}}_\mu$, and $D_5\gamma^{x_5}$ with both connections.

\textbf{Exercise 4.12.} \emph{A consistency check with a trace.} (a) Show that in $n=8$ dimensions $G^\mu{}_\mu=-3R$. (b) Check this for the primordial field with $a_4=t$, using the values of Section 4.8.

\subsection{Answers to the exercises}

\textbf{Answer 4.1.} (a) The inverse of a diagonal matrix is diagonal with the reciprocal entries: $g^{rr}=1$, $g^{\varphi\varphi}=1/r^2$, $g^{r\varphi}=0$. (b) Along the circle $u^r=0$ and $u^\varphi=1$, so $\sqrt{g_{\varphi\varphi}}=r=2$ and the length is $\int_0^{2\pi}2\,d\lambda=4\pi$. (c) $V_r=g_{rr}V^r=\cos\varphi$ and $V_\varphi=g_{\varphi\varphi}V^\varphi=r^2(-\sin\varphi/r)=-r\sin\varphi$. Then $g(V,V)=V^rV_r+V^\varphi V_\varphi=\cos^2\varphi+(-\sin\varphi/r)(-r\sin\varphi)=\cos^2\varphi+\sin^2\varphi=1$.

\textbf{Answer 4.2.} (a) With $\Gamma^\theta{}_{\varphi\varphi}=-\sin\theta\cos\theta$ and $\Gamma^\varphi{}_{\theta\varphi}=\Gamma^\varphi{}_{\varphi\theta}=\cot\theta$ (primes are $d/d\lambda$): $\theta''-\sin\theta\cos\theta\,\varphi'^2=0$ and $\varphi''+2\cot\theta\,\theta'\varphi'=0$. (b) On the equator $\theta'=\theta''=0$, $\varphi'=1/a$ and $\varphi''=0$. The first equation holds because $\cos(\pi/2)=0$; the second because $\theta'=0$. $g(u,u)=a^2(\theta'^2+\sin^2\theta\,\varphi'^2)=a^2\cdot\frac1{a^2}=1$. (c) Now $\theta'=\theta''=0$ and $\varphi'=\frac{1}{a\sin(\pi/3)}=\frac{2}{\sqrt3\,a}$, so $\varphi'^2=\frac{4}{3a^2}$, while $\sin\theta\cos\theta=\frac{\sqrt3}2\cdot\frac12=\frac{\sqrt3}4$. The first equation gives $0-\frac{\sqrt3}{4}\cdot\frac{4}{3a^2}=-\frac{1}{\sqrt3\,a^2}\ne0$. Circles of latitude other than the equator are not geodesics; the geodesics of the sphere are its great circles, which is why long-distance flights do not follow lines of constant latitude.

\textbf{Answer 4.3.} The velocity is $u^\mu=\delta^\mu{}_4$ (only $u^4=1$), so $du^\mu/d\lambda=0$ and the geodesic equation reduces to $\Gamma^\mu{}_{44}=0$ for every $\mu$. By the diagonal formulas of Section 4.5: $\Gamma^4{}_{44}=\partial_4\ln h_4=0$ because $h_4=1$, and for $\mu\ne4$, $\Gamma^\mu{}_{44}=-\eta_{\mu\mu}\eta_{44}\frac{h_4}{h_\mu^2}\partial_\mu h_4=0$, again because $h_4=1$ is constant. (The table of the 37 symbols contains no symbol with the lower pair $(4,4)$.) So the curve is a geodesic for every $a_4$, and $g(u,u)=g_{44}=-1$: it is time-like, and by the definition of proper time in Section 4.6, $\tau=\int\sqrt{-g_{44}}\,d\lambda=\lambda=x^4$ (counted from $x^4=0$) is its proper time.

\textbf{Answer 4.4.} Here $h_y=1$ and $h_w=e^{Hy}$, both signs $+1$. By the diagonal formulas: $\Gamma^w{}_{yw}=\Gamma^w{}_{wy}=\partial_y\ln e^{Hy}=H$ and $\Gamma^y{}_{ww}=-\frac{h_w}{h_y^2}\partial_yh_w=-He^{2Hy}$; all others vanish (in particular $\Gamma^\lambda{}_{yy}=0$). Then

\[
R^y{}_{wyw}=\partial_y\Gamma^y{}_{ww}-\partial_w\Gamma^y{}_{yw}+\Gamma^y{}_{y\lambda}\Gamma^\lambda{}_{ww}-\Gamma^y{}_{w\lambda}\Gamma^\lambda{}_{yw}=-2H^2e^{2Hy}-0+0-(-He^{2Hy})H=-H^2e^{2Hy}.
\]

So $R_{ww}=R^y{}_{wyw}=-H^2e^{2Hy}$. For $R_{yy}=R^w{}_{ywy}=\partial_w\Gamma^w{}_{yy}-\partial_y\Gamma^w{}_{wy}+\Gamma^w{}_{w\lambda}\Gamma^\lambda{}_{yy}-\Gamma^w{}_{y\lambda}\Gamma^\lambda{}_{wy}=0-0+0-H\cdot H=-H^2$. $R_{yw}=0$. Finally $R=g^{yy}R_{yy}+g^{ww}R_{ww}=-H^2+e^{-2Hy}(-H^2e^{2Hy})=-2H^2$: constant and negative, a space of constant negative curvature (the hyperbolic plane). The warped form of the primordial field (Section 4.4) has the same kind of warp factor, $e^{2H\zeta}$, in each of its six transverse directions.

\textbf{Answer 4.5.} With $a_4'=a_4''=0$: $R=6H^2(0-7)=-42H^2$; $G^0{}_0=-3H^2(0-5)=15H^2$, $G^i{}_i=15H^2$, $G^4{}_4=21H^2$, $G^j{}_j=15H^2$, so $G^\mu{}_\nu=\mathrm{diag}(15,15,15,15,21,15,15,15)\,H^2$. The required source has $\rho_{\mathrm{req}}=-G^4{}_4/\kappa=-21H^2/\kappa$ and the equal pressures $p=15H^2/\kappa$ in all seven transverse directions, so $w=-15/21=-5/7$. These are the values of the Stage-4 exact theory for the static field (checks \texttt{KS\_geome\allowbreak{}try\_ricc\allowbreak{}iScalarM\allowbreak{}inus42H2}, \texttt{KS\_geome\allowbreak{}try\_eins\allowbreak{}teinMixe\allowbreak{}dDiag} and \texttt{KS\_geome\allowbreak{}try\_requ\allowbreak{}iredSour\allowbreak{}ce} in \texttt{artifact\allowbreak{}s/dirac1\allowbreak{}6complex\allowbreak{}/kohn-sh\allowbreak{}am/wolfr\allowbreak{}am-kohn-\allowbreak{}sham-rep\allowbreak{}ort.json}; \texttt{handoff/\allowbreak{}specs/ST\allowbreak{}AGE4\_SPE\allowbreak{}C.md}, erratum E4.3, for the sign of $p$), and $w=-5/7$ is the value of the Stage-2 formula $w=(c^2-5)/(c^2+7)$ at $c=0$ (Stage-2 document, §15.4).

\textbf{Answer 4.6.} (a) $G_{\theta\theta}=R_{\theta\theta}-\tfrac12g_{\theta\theta}R=1-\tfrac12a^2\cdot\tfrac2{a^2}=0$, $G_{\varphi\varphi}=\sin^2\theta-\tfrac12a^2\sin^2\theta\cdot\tfrac{2}{a^2}=0$ and $G_{\theta\varphi}=0$. (b) By Section 4.9, $G^\mu{}_\nu=E^{(1)\mu}{}_\nu$ is a contraction of the generalized delta with $2k+1=3$ upper indices. In two dimensions each index takes only two values, so two of the three upper indices are always equal and the delta vanishes. Hence $G_{\mu\nu}=0$ for every two-dimensional metric, and Einstein's equations in two dimensions would force $T_{\mu\nu}=0$.

\textbf{Answer 4.7.} $E^{(k)}$ vanishes identically when $2k+1>n$ (that the lower orders do not vanish identically is quoted, Section 4.9). So: $n=4$: $k=0,1$; $n=5$: $k=0,1,2$; $n=6$: $k=0,1,2$; $n=8$: $k=0,1,2,3$; $n=10$: $k=0,\dots,4$. For even $n$ the highest order is $n/2-1$: 1, 2, 3 and 4 for $n=4,6,8,10$, in agreement with the notebook's rule ($m=n/2$).

\textbf{Answer 4.8.} (a) $\sigma_3^2=1$, so the exponential series splits into even and odd powers: $\exp(i\varphi\sigma_3/2)=\cos(\varphi/2)+i\sigma_3\sin(\varphi/2)=\mathrm{diag}(e^{i\varphi/2},e^{-i\varphi/2})$. Since $\sigma_3$ anticommutes with $\sigma_1$ and $\sigma_2$, moving $\sigma_1$ through $R$ reverses the sign of the exponent: $R\sigma_1=\sigma_1R^{-1}$. Hence $R\sigma_1R^{-1}=\sigma_1R^{-2}=\sigma_1(\cos\varphi-i\sigma_3\sin\varphi)=\cos\varphi\,\sigma_1-\sin\varphi\,\sigma_2$, using $\sigma_1\sigma_3=-i\sigma_2$. In the same way $R\sigma_2R^{-1}=\sigma_2(\cos\varphi-i\sigma_3\sin\varphi)=\cos\varphi\,\sigma_2+\sin\varphi\,\sigma_1$, using $\sigma_2\sigma_3=i\sigma_1$. Therefore $R(\cos\varphi\,\sigma_1+\sin\varphi\,\sigma_2)R^{-1}=(\cos^2\varphi+\sin^2\varphi)\sigma_1+(-\cos\varphi\sin\varphi+\sin\varphi\cos\varphi)\sigma_2=\sigma_1$: the curved gamma $\gamma^r$ of the Cartesian frame, transformed with $R$, is the curved gamma $\gamma^r=\sigma_1$ of the polar frame, as the covariance theorem of Section 4.12 requires. (b) $\partial_\varphi R=\tfrac i2\sigma_3R$, so $-(\partial_\varphi R)R^{-1}=-\tfrac i2\sigma_3$, which is $\Omega_\varphi$ of Section 4.12. By the covariance theorem the polar connection is $R\cdot0\cdot R^{-1}-(\partial_\varphi R)R^{-1}$, and the two computations agree. (c) $R(2\pi)=\mathrm{diag}(e^{i\pi},e^{-i\pi})=-1$. After one full turn around the origin the polar frame is back where it started, but the spin transformation that connects it with the Cartesian frame has become $-1$. So the polar-frame components $R\Psi$ of a smooth spinor field change sign after a full turn. This is the sign of Chapter 2: a rotation by $2\pi$ acts on spinors as $-1$. (d) From $\gamma^\varphi\Omega_\varphi=\frac{\sigma_2}{r}\bigl(-\frac i2\sigma_3\bigr)=\frac{1}{2r}\sigma_1$: it is the spinor connection of the turning polar frame. By covariance, $\gamma^\mu D_\mu(R\Psi)=R\,(\sigma_1\partial_x+\sigma_2\partial_y)\Psi$, so the term $1/(2r)$ is exactly what is needed for the polar form to describe the same flat-space Dirac operator.

\textbf{Answer 4.9.} Use the diagonal formula of Section 4.12. The scale factors depend on $t=x^4$ only, so only $b=4$ contributes; $h_4=1$, and $\prod_{c\ne4}h_c=V$ is the 7-volume of a comoving region with unit coordinate ranges (Section 4.4: here the 7 by 7 block is diagonal with the entries $\eta_{jj}h_j^2$, so the square root of the absolute value of its determinant is $V$). So $\gamma^\mu\Omega_\mu=\tfrac12\,\partial_t\ln V\,\gamma^4=\tfrac12\sum_{j\ne4}\frac{\dot h_j}{h_j}\,\gamma^4=\tfrac12\Theta\gamma^4$. This is the one line of geometry that the Stage-3 experiments EXP-2 to EXP-5 need (student guide, §6.6; check \texttt{GEO\_diag\allowbreak{}onalSlas\allowbreak{}hFormula\allowbreak{}\_G3} in \texttt{python-g\allowbreak{}eometry-\allowbreak{}report.j\allowbreak{}son}).

\textbf{Answer 4.10.} (a) Space-space: choose 2 of the 4 space-like directions, $\binom42=6$; time-time: likewise 6; boost pairs: $4\times4=16$; total $6+6+16=28$. (b) With three space directions and one time direction: space-space $\binom32=3$, time-time 0, boosts $3\times1=3$. The same mistake in ordinary four-dimensional spacetime would delete all three boost components of the connection. (c) The symmetric part of $\omega_\mu{}^a{}_b$ is the boost part: for each of the 8 values of $\mu$, 16 pairs in 2 orders, $8\times16\times2=256$ entries. $\omega_{\mu ab}$ has, for each $\mu$, $8\times7=56$ ordered pairs with $a\ne b$, in total $8\times56=448$. Both numbers are exactly the measurements \texttt{nonzeroO\allowbreak{}megaMixe\allowbreak{}dSymmetr\allowbreak{}icPart} = 256 and \texttt{nonzeroO\allowbreak{}megaLowe\allowbreak{}r} = 448 at each G1 point in \texttt{wolfram-\allowbreak{}geometry\allowbreak{}-report.\allowbreak{}json}: the generic vielbein G1 populates every component.

\textbf{Answer 4.11.} Correct: $\gamma^\mu\Omega_\mu=3H\gamma^0=2\gamma^0$. Notebook: $\tfrac{3H}{2}(\gamma^0+a_4'\gamma^4)=\gamma^0+\tfrac37\gamma^4$, the value \texttt{G2.p1.no\allowbreak{}tebookSl\allowbreak{}ash} of \texttt{wolfram-\allowbreak{}geometry\allowbreak{}-report.\allowbreak{}json}. For $D_5\gamma^{x_5}$: $\gamma^{x_5}=\gamma^5/h_5$ does not depend on $x^5$, so $\partial_5\gamma^{x_5}=0$. The nonzero $\Gamma^5{}_{5\lambda}$ are $\Gamma^5{}_{50}=H\cot z$ and $\Gamma^5{}_{54}=-Ha_4'$, so $\Gamma^5{}_{5\lambda}\gamma^{x_\lambda}=H\cot z\tan z\,\gamma^0-Ha_4'\gamma^4=H\gamma^0-Ha_4'\gamma^4$. With the identity of Section 4.12, $[S^{05},\gamma^5]=\eta^{55}\gamma^0-\eta^{05}\gamma^5=-\gamma^0$ and $[S^{45},\gamma^5]=\eta^{55}\gamma^4-\eta^{45}\gamma^5=-\gamma^4$. With $\Omega_5=\alpha_5(S^{05}-a_4'S^{45})$ and $\alpha_5/h_5=H$: $[\Omega_5,\gamma^{x_5}]=H(-\gamma^0+a_4'\gamma^4)$, and $D_5\gamma^{x_5}=0$. With $\Omega^{\mathrm{nb}}_5=\alpha_5a_4'S^{45}$: $[\Omega^{\mathrm{nb}}_5,\gamma^{x_5}]=-Ha_4'\gamma^4$, and $D^{\mathrm{nb}}_5\gamma^{x_5}=H\gamma^0-2Ha_4'\gamma^4=\tfrac23\gamma^0-\tfrac47\gamma^4$. This is the entry "$(j,j)$" of the table of the Stage-2 document, §8.3.

\textbf{Answer 4.12.} (a) $G^\mu{}_\mu=R^\mu{}_\mu-\tfrac12\delta^\mu{}_\mu R=R-\tfrac82R=-3R$. (b) With $a_4=t$, $G^\mu{}_\mu=(7\cdot12+24)H^2=108H^2$ and $-3R=-3(-36H^2)=108H^2$. The two values of Section 4.8 are consistent.

\section{Classical field theory, Grassmann numbers, and why the notebook's Lg[] is empty}

\subsection{What this chapter does}

Physics describes how things change. For a single particle the thing that changes is its position; for a field, such as the electromagnetic field or the dirac16complex field of this book, it is a list of numbers attached to every point of spacetime. In both cases the most economical way to write down the laws of change is to write down one function, the \textbf{Lagrangian}, and to demand that a certain integral of it, the \textbf{action}, does not change to first order when the motion is changed slightly. This one demand produces the equations of motion, called the \textbf{Euler--Lagrange equations}. The same Lagrangian also tells us which quantities are conserved (energy, momentum, charge), through a theorem of Emmy Noether, and it gives the \textbf{energy--momentum tensor}, the object that tells gravity where the energy is.

This chapter builds that machinery from zero. It then introduces \textbf{Grassmann numbers}, the anticommuting numbers ($\theta_1\theta_2=-\theta_2\theta_1$) that are needed to describe fermions such as electrons, and it ends with one of the first results of the dirac16complex project: the Lagrangian $\mathrm{Lg}[\,]$ that the author wrote in the notebook (cell 1064) contains no dynamics at all when its real field (the notebook's Ψ16) is taken to be anticommuting (Grassmann), as this book assumes a fermion field must be (Section 5.16). With the correct (canonical) spin connection its Euler--Lagrange equations read $0=0$. With the notebook's own contraction they are the algebraic equations $X\Psi=0$, which force $\Psi=0$ wherever the matrix $X$ is invertible, as it is at all six points where the repository computed it (Sections 5.13 and 5.14). Either way there is no wave equation. This is why the project had to replace $\mathrm{Lg}[\,]$ by a new Lagrangian (Chapter 6).

We use the conventions of the whole book (Chapter 1): everything is counted from 0; there are eight coordinates $x=(x_0,x_1,\dots,x_7)$; $x_4$ is the time; the flat metric is

\[
\eta=\mathrm{diag}(+1,+1,+1,+1,-1,-1,-1,-1),
\]

so directions 0 to 3 are space-like and directions 4 to 7 are time-like, and $x_5,x_6,x_7$ are called the three extra times. In words we name the coordinates $x_0,\dots,x_7$ as the notebook does; in index formulas we write $x^\mu$ and $\partial_\mu=\partial/\partial x^\mu$ (Section 1.8). The names $x_4$ and $x^4$ denote the same coordinate; no metric is used to lower the index of a coordinate. A repeated index, one up and one down, is summed from 0 to 7 (the sum convention of Chapter 1).

The sources of this chapter are the Stage-1 document \texttt{provenan\allowbreak{}ce/DIRAC\allowbreak{}16COMPLE\allowbreak{}X\_ARBITR\allowbreak{}ARY\_FIEL\allowbreak{}D.md} (its §6 ``Why the notebook Lagrangian Lg[] is empty for a Grassmann field'' and §3 on the matrices), the contract \texttt{handoff/\allowbreak{}specs/CO\allowbreak{}NTRACT.m\allowbreak{}d} (its §3), the exact Grassmann-algebra code \texttt{scripts/\allowbreak{}grassman\allowbreak{}n\_algebr\allowbreak{}a.py} and \texttt{scripts/\allowbreak{}demo\_gra\allowbreak{}ssmann\_l\allowbreak{}agrangia\allowbreak{}ns.py}, and the report \texttt{artifact\allowbreak{}s/dirac1\allowbreak{}6complex\allowbreak{}/arbitra\allowbreak{}ry-field\allowbreak{}/grassma\allowbreak{}nn-demo-\allowbreak{}report.j\allowbreak{}son}.

\subsection{The action of a particle and the Euler--Lagrange equation}

\textbf{Setting.} A particle moves on a line. Its position at time $t$ is a number $q(t)$, and its velocity is $\dot q(t)=dq/dt$. A \textbf{Lagrangian} is a function $L(q,\dot q)$ of the position and the velocity; for a particle of mass $m$ in a potential $V(q)$ it is

\[
L(q,\dot q)=\tfrac12m\dot q^2-V(q),
\]

the kinetic energy minus the potential energy. A \textbf{path} is a function $q(t)$ on a time interval $t_A\le t\le t_B$. The \textbf{action} of a path is the number

\[
S[q]=\int_{t_A}^{t_B}L\bigl(q(t),\dot q(t)\bigr)\,dt .
\]

The square brackets say that $S$ depends on the whole function $q$, not on one number. A rule that turns a function into a number is called a \textbf{functional}.

\textbf{The principle of stationary action.} Among all paths with fixed end points $q(t_A)=q_A$ and $q(t_B)=q_B$, the path that the particle actually follows is one for which $S$ does not change to first order when the path is changed a little. To make this precise, pick any smooth function $\xi(t)$ with $\xi(t_A)=\xi(t_B)=0$ (a \textbf{variation}) and a small number $\epsilon$, and compare $q$ with $q+\epsilon\xi$. The path $q$ is \textbf{stationary} if

\[
\frac{d}{d\epsilon}S[q+\epsilon\xi]\Big|_{\epsilon=0}=0\qquad\text{for every such }\xi .
\]

\textbf{Derivation of the Euler--Lagrange equation.} Differentiate under the integral sign with the chain rule; the Lagrangian depends on $\epsilon$ through its first argument $q+\epsilon\xi$ and its second argument $\dot q+\epsilon\dot\xi$:

\[
\frac{d}{d\epsilon}S[q+\epsilon\xi]\Big|_{\epsilon=0}=\int_{t_A}^{t_B}\Bigl(\frac{\partial L}{\partial q}\,\xi+\frac{\partial L}{\partial\dot q}\,\dot\xi\Bigr)dt .
\]

The second term contains $\dot\xi$. Integrate it by parts, which is the product rule $\frac{d}{dt}(fg)=\dot fg+f\dot g$ integrated from $t_A$ to $t_B$:

\[
\int_{t_A}^{t_B}\frac{\partial L}{\partial\dot q}\,\dot\xi\,dt=\Bigl[\frac{\partial L}{\partial\dot q}\,\xi\Bigr]_{t_A}^{t_B}-\int_{t_A}^{t_B}\frac{d}{dt}\Bigl(\frac{\partial L}{\partial\dot q}\Bigr)\xi\,dt .
\]

The bracket vanishes because $\xi(t_A)=\xi(t_B)=0$. Hence

\[
\frac{d}{d\epsilon}S[q+\epsilon\xi]\Big|_{\epsilon=0}=\int_{t_A}^{t_B}E(t)\,\xi(t)\,dt,\qquad E:=\frac{\partial L}{\partial q}-\frac{d}{dt}\frac{\partial L}{\partial\dot q}.
\]

\textbf{The fundamental lemma.} If a continuous function $E(t)$ satisfies $\int E\xi\,dt=0$ for every smooth $\xi$ that vanishes at the end points, then $E=0$ everywhere. \emph{Proof.} Suppose $E(t_0)>0$ at some interior time $t_0$. By continuity $E>0$ on a small interval around $t_0$. Choose $\xi$ positive inside that interval and zero outside (a smooth ``bump''). Then $\int E\xi\,dt>0$, a contradiction. The case $E(t_0)<0$ is the same with the sign reversed. $\square$

So a stationary path satisfies the \textbf{Euler--Lagrange equation}

\[
\frac{\partial L}{\partial q}-\frac{d}{dt}\frac{\partial L}{\partial\dot q}=0 .
\]

With several coordinates $q_0,\dots,q_{n-1}$ one varies each separately and obtains one equation for each $q_i$.

\textbf{Worked example (the harmonic oscillator).} Take $m=1$ and $V=\tfrac12\omega^2q^2$. Then $\partial L/\partial q=-\omega^2q$ and $\partial L/\partial\dot q=\dot q$, so the Euler--Lagrange equation is $-\omega^2q-\ddot q=0$, that is $\ddot q=-\omega^2q$. Its solutions are $q(t)=A\cos\omega t+B\sin\omega t$. For $\omega=2$, $A=1$, $B=0$: $q(t)=\cos2t$, $\ddot q=-4\cos2t=-\omega^2q$, as required.

\textbf{Energy.} For a Lagrangian that does not depend explicitly on $t$ the quantity $H:=\dot q\,\partial L/\partial\dot q-L$ is constant along every solution. \emph{Proof.} $\frac{dH}{dt}=\ddot q\frac{\partial L}{\partial\dot q}+\dot q\frac{d}{dt}\frac{\partial L}{\partial\dot q}-\frac{\partial L}{\partial q}\dot q-\frac{\partial L}{\partial\dot q}\ddot q=\dot q\Bigl(\frac{d}{dt}\frac{\partial L}{\partial\dot q}-\frac{\partial L}{\partial q}\Bigr)=0$. $\square$ For the oscillator $H=\tfrac12\dot q^2+\tfrac12\omega^2q^2$, the kinetic plus the potential energy. $H$ written as a function of $q$ and the \textbf{momentum} $p:=\partial L/\partial\dot q$ is the \textbf{Hamiltonian}; it returns in Chapter 8.

\subsection{Fields, Lagrangian densities and the field equations}

\textbf{A field} is a rule that assigns numbers to every point $x=(x_0,\dots,x_7)$ of spacetime. A \textbf{scalar field} $\phi(x)$ assigns one real number; the dirac16complex field assigns 16 numbers $\Psi_0(x),\dots,\Psi_{15}(x)$ (of a new kind, Section 5.9). We write a general field as a list $\phi_A(x)$, $A=0,1,\dots$, of components.

\textbf{Lagrangian density and action.} For a field the Lagrangian is replaced by a \textbf{Lagrangian density} $\mathcal L$, a function of the field components $\phi_A$, their first derivatives $\partial_\mu\phi_A$, and possibly the point $x$ itself. The action of a field configuration in a region $R$ of spacetime is the eight-fold integral

\[
S[\phi]=\int_R\mathcal L\bigl(\phi_A(x),\partial_\mu\phi_A(x),x\bigr)\,d^8x,\qquad d^8x=dx_0\,dx_1\cdots dx_7 .
\]

Time is now just one of the eight coordinates, and the region $R$ plays the role of the time interval $[t_A,t_B]$.

\textbf{The divergence theorem in the form we need.} Let $V^\mu(x)$, $\mu=0,\dots,7$, be eight smooth functions that vanish near the boundary of a box $R=[a_0,b_0]\times\dots\times[a_7,b_7]$. Then

\[
\int_R\partial_\mu V^\mu\,d^8x=0 .
\]

\emph{Proof.} The sum $\partial_\mu V^\mu$ has eight terms. In the term $\partial_0V^0$ do the $x_0$ integral first: by the fundamental theorem of calculus it equals $V^0(b_0,\dots)-V^0(a_0,\dots)=0$. The same holds for each of the other seven terms, doing the corresponding integral first. $\square$ (For a general region the integral equals a flux through the boundary, which vanishes when $V$ vanishes near the boundary; the box is enough for everything in this book.) An expression of the form $\partial_\mu V^\mu$ is called a \textbf{total divergence} or \textbf{total derivative}.

\textbf{Derivation of the field equations.} Vary every component, $\phi_A\to\phi_A+\epsilon\xi_A$, where the functions $\xi_A(x)$ vanish near the boundary of $R$. Exactly as in Section 5.2,

\[
\frac{d}{d\epsilon}S\Big|_{\epsilon=0}=\int_R\Bigl(\frac{\partial\mathcal L}{\partial\phi_A}\xi_A+\frac{\partial\mathcal L}{\partial(\partial_\mu\phi_A)}\partial_\mu\xi_A\Bigr)d^8x .
\]

Here $\partial\mathcal L/\partial(\partial_\mu\phi_A)$ means: treat the eight numbers $\partial_0\phi_A,\dots,\partial_7\phi_A$ as independent variables of the function $\mathcal L$ and differentiate with respect to one of them. The product rule gives $\frac{\partial\mathcal L}{\partial(\partial_\mu\phi_A)}\partial_\mu\xi_A=\partial_\mu\bigl(\frac{\partial\mathcal L}{\partial(\partial_\mu\phi_A)}\xi_A\bigr)-\partial_\mu\bigl(\frac{\partial\mathcal L}{\partial(\partial_\mu\phi_A)}\bigr)\xi_A$. The first piece is a total divergence of a function that vanishes near the boundary, so its integral is zero. Therefore

\[
\frac{d}{d\epsilon}S\Big|_{\epsilon=0}=\int_R E^A\xi_A\,d^8x,\qquad E^A:=\frac{\partial\mathcal L}{\partial\phi_A}-\partial_\mu\frac{\partial\mathcal L}{\partial(\partial_\mu\phi_A)} .
\]

The fundamental lemma, applied with a bump in one component $\xi_A$ at a time, gives the \textbf{Euler--Lagrange field equations}

\[
E^A=\frac{\partial\mathcal L}{\partial\phi_A}-\partial_\mu\frac{\partial\mathcal L}{\partial(\partial_\mu\phi_A)}=0\qquad\text{for every component }A .
\]

In $\partial_\mu(\dots)$ the derivative is the \textbf{total} derivative with respect to $x^\mu$: it acts on the explicit $x$-dependence of $\mathcal L$ and, through the chain rule, on the fields inside it. The expression $E^A$ is called the \textbf{Euler--Lagrange expression} of the component $A$.

\subsection{Total derivatives change nothing}

\textbf{Proposition 5.1.} If $\mathcal L=\partial_\mu V^\mu$, where each $V^\mu$ is a function of the fields $\phi_A$ and of $x$ (but not of the derivatives of the fields), then every Euler--Lagrange expression vanishes identically: $E^A=0$ for every field configuration, whether or not it satisfies any equation.

\emph{Proof by the action.} Take any variation $\xi_A$ that vanishes near the boundary. The Lagrangian evaluated on $\phi+\epsilon\xi$ is the total divergence of the functions $x\mapsto V^\mu(\phi(x)+\epsilon\xi(x),x)$, so

\[
S[\phi+\epsilon\xi]-S[\phi]=\int_R\partial_\mu\Bigl(V^\mu(\phi+\epsilon\xi,x)-V^\mu(\phi,x)\Bigr)d^8x=0\qquad\text{for every }\epsilon,
\]

by the divergence theorem, because the bracket vanishes near the boundary (there $\xi=0$). Hence $\frac{d}{d\epsilon}S[\phi+\epsilon\xi]=0$, that is $\int E^A\xi_A\,d^8x=0$ for every $\xi$, and the fundamental lemma gives $E^A=0$ for every $\phi$. $\square$

\emph{Proof by direct computation.} By the chain rule $\partial_\mu V^\mu=\partial^{\mathrm{expl}}_\mu V^\mu+\frac{\partial V^\mu}{\partial\phi_B}\partial_\mu\phi_B$, where $\partial^{\mathrm{expl}}_\mu$ differentiates only the explicit $x$-dependence. Then $\frac{\partial\mathcal L}{\partial(\partial_\nu\phi_A)}=\frac{\partial V^\nu}{\partial\phi_A}$ and $\frac{\partial\mathcal L}{\partial\phi_A}=\partial^{\mathrm{expl}}_\mu\frac{\partial V^\mu}{\partial\phi_A}+\frac{\partial^2V^\mu}{\partial\phi_A\partial\phi_B}\partial_\mu\phi_B=\partial_\mu\Bigl(\frac{\partial V^\mu}{\partial\phi_A}\Bigr)$, the total derivative of the function $\partial V^\mu/\partial\phi_A$. Hence $E^A=\partial_\mu(\partial V^\mu/\partial\phi_A)-\partial_\nu(\partial V^\nu/\partial\phi_A)=0$. $\square$

The consequence is used again and again: \textbf{adding a total divergence to a Lagrangian density does not change its field equations}, and a Lagrangian density that is nothing but a total divergence has the empty field equations $0=0$. It describes nothing.

\textbf{Worked example.} One field $\phi$ in one dimension $x$, $V=\phi^2$, so $\mathcal L=\frac{d}{dx}(\phi^2)=2\phi\phi'$. Then $\partial\mathcal L/\partial\phi=2\phi'$ and $\partial\mathcal L/\partial\phi'=2\phi$, so $E=2\phi'-\frac{d}{dx}(2\phi)=0$ for every function $\phi$.

\subsection{Worked example: a scalar field in 4+4 dimensions}

Take one real scalar field $\phi$ in flat space with the metric $\eta$ and

\[
\mathcal L=-\tfrac12\eta^{\mu\nu}\partial_\mu\phi\,\partial_\nu\phi-V(\phi),\qquad V(\phi)=\tfrac12m^2\phi^2 .
\]

Written out, $\eta^{\mu\nu}\partial_\mu\phi\,\partial_\nu\phi=\sum_{i=0}^{3}(\partial_i\phi)^2-\sum_{j=4}^{7}(\partial_j\phi)^2$. For a field that depends only on the time $x_4$ this Lagrangian is $\tfrac12\dot\phi^2-V$ (a dot is $\partial_4$), the kinetic minus the potential energy, exactly as for the particle. This is the sign convention of the whole project (\texttt{handoff/\allowbreak{}specs/CO\allowbreak{}NTRACT.m\allowbreak{}d}, its §0); it is the same function as the $\tfrac12\partial_\mu\phi\,\partial^\mu\phi-V$ of the scalar-field reference that the task supplied (a private input that is not in the repository; the Stage-1 document quotes its formulas in its §2.3), which uses the signature $(+,-,-,-)$ with time first.

\textbf{Field equation.} $\partial\mathcal L/\partial\phi=-m^2\phi$ and $\partial\mathcal L/\partial(\partial_\mu\phi)=-\eta^{\mu\nu}\partial_\nu\phi$ (the factor $\tfrac12$ is cancelled because $\partial_\mu\phi$ occurs twice in the product). The Euler--Lagrange equation $-m^2\phi+\partial_\mu(\eta^{\mu\nu}\partial_\nu\phi)=0$ is

\[
\eta^{\mu\nu}\partial_\mu\partial_\nu\phi=m^2\phi,\qquad\text{i.e.}\qquad \sum_{i=0}^{3}\partial_i^2\phi-\partial_4^2\phi-\sum_{j=5}^{7}\partial_j^2\phi=m^2\phi .
\]

With one time (signature (7,1)) this would be the ordinary Klein--Gordon wave equation. With four time-like directions it is called \textbf{ultrahyperbolic}.

\textbf{Plane waves.} Try $\phi=\cos\bigl(\sum_{j\ne4}k_jx^j-\omega x^4\bigr)$ with constant real $k_j$ and $\omega$. Each $\partial_j^2$ produces $-k_j^2\phi$ and $\partial_4^2$ produces $-\omega^2\phi$, so the field equation becomes $-\sum_{i=0}^{3}k_i^2+\omega^2+\sum_{j=5}^{7}k_j^2=m^2$, that is

\[
\omega^2=m^2+k_0^2+k_1^2+k_2^2+k_3^2-k_5^2-k_6^2-k_7^2 .
\]

A wave number $k_5,k_6,k_7$ along an extra time \textbf{lowers} $\omega^2$. For $m=1$, $k_1=1$ and $k_5=1$ we get $\omega^2=1$: an ordinary oscillation. For $m=1$, $k_1=1$ and $k_5=2$ we get $\omega^2=1+1-4=-2$: no real $\omega$ exists. The solution of the equation with this $k$ is then $\cos(k\cdot x)\,e^{\pm\sqrt2\,x_4}$, which grows exponentially in time. The same phenomenon, for the dirac16complex field, is the ``unstable extra-time modes'' of Chapter 8.

\subsection{First-order Lagrangians and complex fields}

The Lagrangians so far were quadratic in the velocities. The Lagrangian of a fermion is instead \textbf{linear} in the time derivative. This section shows, with two small examples, that such Lagrangians are perfectly good and what is special about them.

\textbf{Example A (two real coordinates).} Let $q_0(t)$, $q_1(t)$ be two real coordinates and

\[
L=\tfrac12\bigl(q_0\dot q_1-q_1\dot q_0\bigr)-\tfrac12\omega\bigl(q_0^2+q_1^2\bigr).
\]

For $q_0$: $\partial L/\partial q_0=\tfrac12\dot q_1-\omega q_0$ and $\partial L/\partial\dot q_0=-\tfrac12q_1$, whose time derivative is $-\tfrac12\dot q_1$. The Euler--Lagrange equation is $\tfrac12\dot q_1-\omega q_0+\tfrac12\dot q_1=0$, that is $\dot q_1=\omega q_0$. For $q_1$: $\partial L/\partial q_1=-\tfrac12\dot q_0-\omega q_1$ and $\partial L/\partial\dot q_1=\tfrac12q_0$, so $-\tfrac12\dot q_0-\omega q_1-\tfrac12\dot q_0=0$, that is $\dot q_0=-\omega q_1$. The solution with $q_0(0)=1$, $q_1(0)=0$ is $q_0=\cos\omega t$, $q_1=\sin\omega t$: a rotation. Three features matter later. The equations are of \textbf{first order} in time, so the initial positions alone fix the motion (no initial velocities are needed). The \textbf{momenta} $p_0=\partial L/\partial\dot q_0=-\tfrac12q_1$ and $p_1=\partial L/\partial\dot q_1=\tfrac12q_0$ are not new variables: they are fixed functions of the coordinates. Such relations are called \textbf{constraints}; they are the reason why the quantization of the dirac16complex field in Chapter 8 needs care. And the part $\tfrac12(q_0\dot q_1-q_1\dot q_0)$ is \textbf{not} a total derivative, although it is antisymmetric in $(q_0,q_1)$: its Euler--Lagrange expressions are $\dot q_1$ and $-\dot q_0$, not zero. Keep this in mind; for anticommuting variables exactly the same expression will turn out to be a total derivative (Section 5.11).

\textbf{Example B (one complex coordinate).} A \textbf{complex number} is $z=x+iy$ with real $x$, $y$ and $i^2=-1$; its complex conjugate is $z^\ast=x-iy$, and $z^\ast z=x^2+y^2\ge0$ (Chapter 1). Combine the two real coordinates into one complex coordinate $\psi=(q_0+iq_1)/\sqrt2$. Then $\psi^\ast\psi=\tfrac12(q_0^2+q_1^2)$, and multiplying out gives $\psi^\ast\dot\psi-\dot\psi^\ast\psi=i(q_0\dot q_1-q_1\dot q_0)$. The real Lagrangian

\[
L=\tfrac i2\bigl(\psi^\ast\dot\psi-\dot\psi^\ast\psi\bigr)-\omega\,\psi^\ast\psi=\tfrac12\bigl(q_1\dot q_0-q_0\dot q_1\bigr)-\tfrac12\omega\bigl(q_0^2+q_1^2\bigr)
\]

is Example A with the sense of rotation reversed. \textbf{One may vary $\psi$ and $\psi^\ast$ as if they were independent.} The reason is that $(q_0,q_1)\mapsto(\psi,\psi^\ast)$ is a linear change of variables with constant coefficients that can be inverted, $q_0=(\psi+\psi^\ast)/\sqrt2$ and $q_1=(\psi-\psi^\ast)/(i\sqrt2)$. With the derivatives $\partial/\partial\psi=\tfrac1{\sqrt2}(\partial/\partial q_0-i\,\partial/\partial q_1)$ and $\partial/\partial\psi^\ast=\tfrac1{\sqrt2}(\partial/\partial q_0+i\,\partial/\partial q_1)$ one checks $\partial\psi/\partial\psi=1$ and $\partial\psi^\ast/\partial\psi=0$, and the Euler--Lagrange expressions combine in the same way, $E_{\psi^\ast}=\tfrac1{\sqrt2}(E_{q_0}+iE_{q_1})$ and $E_\psi=\tfrac1{\sqrt2}(E_{q_0}-iE_{q_1})$. So $E_\psi=E_{\psi^\ast}=0$ holds exactly when $E_{q_0}=E_{q_1}=0$.

Varying $\psi^\ast$: $\partial L/\partial\psi^\ast=\tfrac i2\dot\psi-\omega\psi$ and $\partial L/\partial\dot\psi^\ast=-\tfrac i2\psi$, so the Euler--Lagrange equation is $\tfrac i2\dot\psi-\omega\psi+\tfrac i2\dot\psi=0$, that is

\[
i\dot\psi=\omega\psi,\qquad\psi(t)=e^{-i\omega t}\psi(0).
\]

This has the form of the Schrödinger equation $i\,df/dt=Hf$ of quantum mechanics (Chapter 8), with the number $\omega$ in the role of the Hamiltonian $H$. Varying $\psi$ gives the complex conjugate equation $-i\dot\psi^\ast=\omega\psi^\ast$. The dirac16complex Lagrangian of Chapter 6 has exactly this structure, with $\psi$ replaced by a column of 16 anticommuting components and $\omega$ by a $16\times16$ matrix operator (Chapter 8).

\subsection{Symmetries and Noether's theorem}

\textbf{Definition.} A \textbf{continuous symmetry} of a Lagrangian density is a family of changes of the fields, $\phi_A\to\phi_A+\epsilon\,\Delta\phi_A+O(\epsilon^2)$ (with $O(\epsilon^2)$ as in Section 4.5), under which $\mathcal L$ changes, to first order in $\epsilon$, only by a total divergence: $\mathcal L\to\mathcal L+\epsilon\,\partial_\mu K^\mu$ for every field configuration (not only for solutions). By Section 5.4 such a change does not alter the field equations.

\textbf{Theorem 5.2 (Noether).} For a continuous symmetry the \textbf{current}

\[
j^\mu=\frac{\partial\mathcal L}{\partial(\partial_\mu\phi_A)}\,\Delta\phi_A-K^\mu
\]

satisfies $\partial_\mu j^\mu=0$ for every solution of the field equations.

\emph{Proof.} To first order in $\epsilon$ the change of $\mathcal L$ is $\epsilon\bigl(\frac{\partial\mathcal L}{\partial\phi_A}\Delta\phi_A+\frac{\partial\mathcal L}{\partial(\partial_\mu\phi_A)}\partial_\mu\Delta\phi_A\bigr)$. On a solution the field equation replaces $\partial\mathcal L/\partial\phi_A$ by $\partial_\mu\bigl(\partial\mathcal L/\partial(\partial_\mu\phi_A)\bigr)$, and the product rule turns the bracket into $\partial_\mu\bigl(\frac{\partial\mathcal L}{\partial(\partial_\mu\phi_A)}\Delta\phi_A\bigr)$. By assumption the change equals $\epsilon\,\partial_\mu K^\mu$. Subtracting gives $\partial_\mu j^\mu=0$. $\square$

\textbf{Conserved charge.} Integrate $j^4$ over a slice $x_4=\text{const}$, whose seven coordinates are $x_0,x_1,x_2,x_3,x_5,x_6,x_7$: $Q(x_4)=\int j^4\,d^7x$. Then $dQ/dx_4=\int\partial_4j^4\,d^7x=-\sum_{j\ne4}\int\partial_jj^j\,d^7x=0$ when the fields vanish far away, by the divergence theorem of Section 5.3 on the slice. So $Q$ does not change in time. (The overall sign of a current is a convention.)

\textbf{Example 1 (phase symmetry).} In Example B of Section 5.6 the change $\psi\to e^{i\alpha}\psi$, $\psi^\ast\to e^{-i\alpha}\psi^\ast$ leaves $L$ unchanged ($K=0$). To first order $\Delta\psi=i\psi$ and $\Delta\psi^\ast=-i\psi^\ast$. The ``current'' has only a time component, $j=\frac{\partial L}{\partial\dot\psi}\,i\psi+\frac{\partial L}{\partial\dot\psi^\ast}(-i\psi^\ast)=\tfrac i2\psi^\ast\,i\psi-\tfrac i2\psi\,(-i\psi^\ast)=-\psi^\ast\psi$. Its conservation says $\psi^\ast\psi$ is constant, which the solution $e^{-i\omega t}\psi(0)$ confirms. For the dirac16complex field (Chapter 6; its Dirac adjoint $\bar\Psi=\Psi^\dagger C$ is introduced in Section 5.14), Theorem 5.2 applied to the same symmetry, $\Delta\Psi=i\Psi$ and $\Delta\bar\Psi=-i\bar\Psi$, gives the current $+i\bar\Psi\gamma^\mu\Psi$. (In flat space, for the Lagrangian $\mathcal L$ of Section 5.14 with commuting components, $\partial\mathcal L/\partial(\partial_\mu\Psi)=\tfrac12\bar\Psi\gamma^\mu$, $\partial\mathcal L/\partial(\partial_\mu\bar\Psi)=-\tfrac12\gamma^\mu\Psi$ and $K^\mu=0$, so $j^\mu=\tfrac12\bar\Psi\gamma^\mu(i\Psi)+(-i\bar\Psi)(-\tfrac12\gamma^\mu\Psi)=i\bar\Psi\gamma^\mu\Psi$; with Grassmann components, left derivatives and the variation written on the left, as in Section 5.10, the same current comes out.) The project reverses this overall sign, which is allowed because the overall sign of a current is a convention, and uses the conserved current $J^\mu=-i\bar\Psi\gamma^\mu\Psi$. Its time component in flat space is $J^4=-i\Psi^\dagger C\gamma^4\Psi=\Psi^\dagger B\Psi$, with the matrix $B=-iC\gamma^4$ of Chapter 2, and with this sign the particles of Chapter 8 have charge $+1$ (Chapters 7 and 8). For the oscillator the same choice of sign would give $+\psi^\ast\psi$ instead of $-\psi^\ast\psi$.

\textbf{Example 2 (translations and the canonical energy--momentum tensor).} If $\mathcal L$ does not depend explicitly on $x$, shifting the field, $\phi_A(x)\to\phi_A(x+\epsilon e_\nu)$ with $e_\nu$ the unit step in direction $\nu$, is a symmetry: $\Delta\phi_A=\partial_\nu\phi_A$, and $\mathcal L$ changes by $\epsilon\,\partial_\nu\mathcal L=\epsilon\,\partial_\mu(\delta^\mu{}_\nu\mathcal L)$, so $K^\mu=\delta^\mu{}_\nu\mathcal L$ ($\delta^\mu{}_\nu$ is 1 for $\mu=\nu$ and 0 otherwise). Noether's theorem gives eight conserved currents, one for each $\nu$, collected in the \textbf{canonical energy--momentum tensor}

\[
\Theta^\mu{}_\nu=\frac{\partial\mathcal L}{\partial(\partial_\mu\phi_A)}\,\partial_\nu\phi_A-\delta^\mu{}_\nu\,\mathcal L,\qquad \partial_\mu\Theta^\mu{}_\nu=0\ \text{on solutions}.
\]

The component $\Theta^4{}_4=\frac{\partial\mathcal L}{\partial(\partial_4\phi_A)}\partial_4\phi_A-\mathcal L$ is the field version of the energy $H=\dot q\,\partial L/\partial\dot q-L$ of Section 5.2: it is the \textbf{energy density}, and $\int\Theta^4{}_4\,d^7x$ is the conserved energy.

\textbf{The scalar field in 4+4 dimensions.} For the $\mathcal L$ of Section 5.5, $\partial\mathcal L/\partial(\partial_\mu\phi)=-\eta^{\mu\lambda}\partial_\lambda\phi$, so $\Theta^\mu{}_\nu=-\eta^{\mu\lambda}\partial_\lambda\phi\,\partial_\nu\phi-\delta^\mu{}_\nu\mathcal L$. With $\eta^{44}=-1$ the energy density is

\[
\rho=\Theta^4{}_4=\tfrac12(\partial_4\phi)^2+\tfrac12\sum_{i=0}^{3}(\partial_i\phi)^2-\tfrac12\sum_{j=5}^{7}(\partial_j\phi)^2+V(\phi).
\]

(Derivation: $\Theta^4{}_4=(\partial_4\phi)^2-\mathcal L$ and $-\mathcal L=\tfrac12\sum_{i\le3}(\partial_i\phi)^2-\tfrac12(\partial_4\phi)^2-\tfrac12\sum_{j\ge5}(\partial_j\phi)^2+V$.) The gradients along the extra times enter with a \textbf{minus} sign, and this makes the energy unbounded below. To see it, take $m=1$ and, at the moment $x_4=0$, a field with $\partial_4\phi=0$ and

\[
\phi=\epsilon\,f\,\sin(kx_5),
\]

where $\epsilon$ and $k>0$ are numbers and $f$ is a fixed smooth bump on the slice: a function of the seven coordinates $x_0,x_1,x_2,x_3,x_5,x_6,x_7$ that vanishes outside a bounded region and is not zero everywhere, so that $I:=\int f^2\,d^7x>0$. This field vanishes far away, so its energy $E=\int\rho\,d^7x$ is a finite number. We show that $E$ becomes as negative as we like when $k$ is made large. Write $s=\sin(kx_5)$ and $c=\cos(kx_5)$.

\emph{The terms without $\partial_5\phi$.} With $\partial_4\phi=0$ and $m=1$, $\rho=\tfrac12\sum_{i=0}^{3}(\partial_i\phi)^2-\tfrac12\sum_{j=5}^{7}(\partial_j\phi)^2+\tfrac12\phi^2$. For $i\le3$ and for $j=6,7$ the derivative does not act on $s$, so $\partial_i\phi=\epsilon s\,\partial_if$ and $\partial_j\phi=\epsilon s\,\partial_jf$. The terms $-\tfrac12(\partial_6\phi)^2$ and $-\tfrac12(\partial_7\phi)^2$ are $\le0$, and since $s^2\le1$ the others satisfy $\tfrac12\sum_{i\le3}(\partial_i\phi)^2+\tfrac12\phi^2\le\tfrac12\epsilon^2\bigl(\sum_{i\le3}(\partial_if)^2+f^2\bigr)$. So these terms contribute at most the number $K:=\tfrac12\epsilon^2\int\bigl(\sum_{i\le3}(\partial_if)^2+f^2\bigr)d^7x$, which does not depend on $k$.

\emph{The term $-\tfrac12(\partial_5\phi)^2$.} Here $\partial_5\phi=\epsilon(s\,\partial_5f+kcf)$, so, using $2f\,\partial_5f=\partial_5(f^2)$,

\[
-\tfrac12(\partial_5\phi)^2=-\tfrac12\epsilon^2\bigl(s^2(\partial_5f)^2+k\,sc\,\partial_5(f^2)+k^2c^2f^2\bigr).
\]

The first term is $\le0$. For the other two use $sc=\tfrac12\sin(2kx_5)$ and $c^2=\tfrac12\bigl(1+\cos(2kx_5)\bigr)$, and write $J_k:=\int\cos(2kx_5)\,f^2\,d^7x$. Integrating by parts in $x_5$ (the boundary terms vanish because $f$ does), $\int\sin(2kx_5)\,\partial_5(f^2)\,d^7x=-2k\,J_k$. Hence the integral of the last two terms is $-\tfrac14\epsilon^2k\,(-2kJ_k)-\tfrac14\epsilon^2k^2(I+J_k)=-\tfrac14\epsilon^2k^2I+\tfrac14\epsilon^2k^2J_k$.

\emph{The oscillating remainder stays bounded.} Two more integrations by parts in $x_5$ give

\[
J_k=-\frac1{2k}\int\sin(2kx_5)\,\partial_5(f^2)\,d^7x=-\frac1{4k^2}\int\cos(2kx_5)\,\partial_5^2(f^2)\,d^7x ,
\]

so, because $|\cos|\le1$, the term $\tfrac14\epsilon^2k^2J_k$ is at most $L:=\tfrac1{16}\epsilon^2\int|\partial_5^2(f^2)|\,d^7x$, again a number that does not depend on $k$. Adding up,

\[
E=\int\rho\,d^7x\le K+L-\tfrac14\epsilon^2k^2I ,
\]

which is negative for large $k$ and tends to $-\infty$ as $k\to\infty$. The energy of a scalar field in 4+4 dimensions is not bounded below, even for fields that vanish far away; this is the classical face of the extra-time instability of Chapter 8. (For the plane wave $\epsilon\sin(kx_5)$ alone, without the bump, $\rho=\tfrac12\epsilon^2\bigl(\sin^2(kx_5)-k^2\cos^2(kx_5)\bigr)$, whose average over $x_5$ is $\tfrac14\epsilon^2(1-k^2)$; but that field does not vanish far away, and the integral of its energy density over the whole slice does not exist.) This observation is derived here; it is not a check of the repository.

\subsection{The energy--momentum tensor}

\textbf{Flat space.} For the scalar field define

\[
T_{\mu\nu}:=-\eta_{\mu\rho}\,\Theta^\rho{}_\nu=\partial_\mu\phi\,\partial_\nu\phi+\eta_{\mu\nu}\,\mathcal L .
\]

It is symmetric in $\mu,\nu$, it is conserved, and its sign is chosen so that $T_{44}=\Theta^4{}_4=\rho$ is the energy density (the project fixes the sign of every energy--momentum tensor this way; Stage-1 document, §1.2).

\textbf{Curved space.} In a gravitational field (Chapter 4) the flat metric $\eta_{\mu\nu}$ is replaced by a metric $g_{\mu\nu}(x)$ with inverse $g^{\mu\nu}$, and the volume element $d^8x$ by $\sqrt{|g|}\,d^8x$, where $g=\det(g_{\mu\nu})$. The scalar action becomes

\[
S=\int\sqrt{|g|}\,\Bigl(-\tfrac12g^{\mu\nu}\partial_\mu\phi\,\partial_\nu\phi-V(\phi)\Bigr)d^8x=\int\sqrt{|g|}\,\mathcal L_\phi\,d^8x .
\]

The energy--momentum tensor is \textbf{defined} by how the action responds to a change of the metric:

\[
T_{\mu\nu}:=-\frac{2}{\sqrt{|g|}}\,\frac{\delta S}{\delta g^{\mu\nu}},
\]

where $\delta S/\delta g^{\mu\nu}$ is the function that multiplies $\delta g^{\mu\nu}(x)$ in the first-order change of $S$ when $g^{\mu\nu}$ is changed by a small symmetric $\delta g^{\mu\nu}(x)$ that vanishes near the boundary. It is the object that appears on the right-hand side of Einstein's equations (Chapter 4 and Chapter 9).

\textbf{Jacobi's formula.} We need the change of $\sqrt{|g|}$. First, for a square matrix $A$ and a small number $\epsilon$, $\det(1+\epsilon A)=1+\epsilon\,\mathrm{tr}A+O(\epsilon^2)$. \emph{Proof.} The determinant is a sum over permutations of products of one entry from each row and each column. The identity permutation gives $\prod_i(1+\epsilon A_{ii})=1+\epsilon\sum_iA_{ii}+O(\epsilon^2)$. Every other permutation moves at least two indices, so its product contains at least two off-diagonal entries $\epsilon A_{ij}$ and is $O(\epsilon^2)$. $\square$ Hence $\det(g+\delta g)=\det g\,\det(1+g^{-1}\delta g)=\det g\,(1+g^{\mu\nu}\delta g_{\mu\nu})$ to first order, that is $\delta\ln|\det g|=g^{\mu\nu}\delta g_{\mu\nu}$. Second, varying $g^{\mu\lambda}g_{\lambda\nu}=\delta^\mu{}_\nu$ gives $\delta g^{\mu\lambda}\,g_{\lambda\nu}+g^{\mu\lambda}\,\delta g_{\lambda\nu}=0$; setting $\nu=\mu$ and summing gives $g^{\mu\nu}\delta g_{\mu\nu}=-g_{\mu\nu}\delta g^{\mu\nu}$. Together,

\[
\delta\sqrt{|g|}=\tfrac12\sqrt{|g|}\,g^{\mu\nu}\delta g_{\mu\nu}=-\tfrac12\sqrt{|g|}\,g_{\mu\nu}\,\delta g^{\mu\nu}.
\]

\textbf{The scalar energy--momentum tensor.} Now vary $S$. The explicit $g^{\mu\nu}$ in $\mathcal L_\phi$ and the factor $\sqrt{|g|}$ both change, and to first order

\[
\delta S=\int\Bigl(-\tfrac12\sqrt{|g|}\,\partial_\mu\phi\,\partial_\nu\phi-\tfrac12\sqrt{|g|}\,g_{\mu\nu}\,\mathcal L_\phi\Bigr)\delta g^{\mu\nu}\,d^8x .
\]

Multiplying the bracket by $-2/\sqrt{|g|}$,

\[
T_{\mu\nu}=\partial_\mu\phi\,\partial_\nu\phi+g_{\mu\nu}\,\mathcal L_\phi ,
\]

which in flat space is the tensor above: the Noether route and the metric route agree for a scalar field. For a spinor field the canonical tensor $\Theta$ is not symmetric, and the metric route has to be carried out through the vielbein of Chapter 4; the result for dirac16complex is derived in Chapter 7.

\textbf{A derivative of $\sqrt{|g|}$ that we need in Section 5.13.} Jacobi's formula with $\delta$ replaced by $\partial_\lambda$ gives $\partial_\lambda\ln|\det g|=g^{\rho\sigma}\partial_\lambda g_{\rho\sigma}$. The Christoffel symbols of Chapter 4, $\Gamma^\rho{}_{\mu\nu}=\tfrac12g^{\rho\sigma}(\partial_\mu g_{\nu\sigma}+\partial_\nu g_{\mu\sigma}-\partial_\sigma g_{\mu\nu})$, contracted over $\rho=\mu$, give $\Gamma^\rho{}_{\rho\lambda}=\tfrac12g^{\rho\sigma}(\partial_\rho g_{\lambda\sigma}+\partial_\lambda g_{\rho\sigma}-\partial_\sigma g_{\rho\lambda})$. The first and the third term cancel (rename $\rho\leftrightarrow\sigma$ in the third and use the symmetry of $g$). Therefore

\[
\Gamma^\rho{}_{\rho\lambda}=\tfrac12g^{\rho\sigma}\partial_\lambda g_{\rho\sigma}=\partial_\lambda\ln\sqrt{|g|},\qquad\text{i.e.}\qquad \partial_\lambda\sqrt{|g|}=\sqrt{|g|}\,\Gamma^\rho{}_{\rho\lambda}.
\]

\textbf{Energy density, pressure and equation of state of a homogeneous field.} In flat space let $\phi$ depend only on $x_4$. Then $\rho=T_{44}=\dot\phi^2+\eta_{44}\mathcal L=\dot\phi^2-(\tfrac12\dot\phi^2-V)=\tfrac12\dot\phi^2+V$. For a direction $i\ne4$, $T_{ii}=\eta_{ii}\mathcal L$ and the \textbf{pressure} along $i$ is $p_{(i)}:=T^i{}_i=\eta^{ii}T_{ii}=\mathcal L=\tfrac12\dot\phi^2-V$ (no sum over $i$). The \textbf{equation-of-state parameter} is

\[
w=\frac p\rho=\frac{\tfrac12\dot\phi^2-V}{\tfrac12\dot\phi^2+V}.
\]

These are the standard formulas $\rho_\phi=\tfrac12\dot\phi^2+V$ and $P_\phi=\tfrac12\dot\phi^2-V$ of a homogeneous scalar field (the same as in the scalar-field reference of Section 5.5, which the Stage-1 document quotes in its §2.3). With small numbers: $\dot\phi=2$ and $V=1$ give $\rho=3$, $p=1$, $w=1/3$; $\dot\phi=0$ gives $w=-1$, the behaviour of a cosmological constant; $V=0$ gives $w=+1$. Chapter 7 derives the corresponding formulas for dirac16complex, and Chapter 11 uses them in the dark-sector experiments.

\subsection{Grassmann numbers from zero}

\textbf{Why a new kind of number.} Electrons, quarks and, by construction, the quanta of the dirac16complex field are \textbf{fermions}: two of them can never occupy the same state (the Pauli principle). In quantum theory this is expressed by operators that \textbf{anticommute}, $ab=-ba$ (Chapter 8). The classical field whose quantization produces such operators must itself take values that anticommute. Ordinary numbers cannot do that, so one introduces new symbols with exactly this property. They are called \textbf{Grassmann numbers} after Hermann Grassmann. They are not results of measurements; they are bookkeeping symbols with precise algebraic rules, and every statement about them in this book is a statement about those rules.

\textbf{Definition.} Choose $n$ symbols $\theta_0,\theta_1,\dots,\theta_{n-1}$, the \textbf{generators}. The \textbf{Grassmann algebra} $\Lambda_n$ consists of all finite sums of complex numbers times products of generators, added and multiplied with the usual distributive rules, with ordinary numbers commuting with everything, and with the single new rule

\[
\theta_i\theta_j=-\theta_j\theta_i\qquad\text{for all }i,j .
\]

Taking $i=j$ gives $\theta_i\theta_i=-\theta_i\theta_i$, hence $\theta_i^2=0$: a generator squares to zero.

\textbf{A basis.} Any product of generators can be reordered into increasing index order; each exchange of two neighbours gives a factor $-1$, so the reordering gives the sign of the permutation (the number $(-1)^s$, $s$ the number of neighbour exchanges needed). A product in which a generator occurs twice is zero. So every element is a unique combination of the $2^n$ \textbf{monomials}

\[
1,\qquad \theta_{i_1}\theta_{i_2}\cdots\theta_{i_k}\quad(i_1<i_2<\dots<i_k),
\]

and $\Lambda_n$ is a vector space of dimension $2^n$. This is exactly how the repository stores an element on a computer: \texttt{scripts/\allowbreak{}grassman\allowbreak{}n\_algebr\allowbreak{}a.py} keeps a dictionary from sorted tuples of generator indices to exact coefficients, and every product sorts the merged tuple and records the sign of the sorting permutation (its function \texttt{\detokenize{_merge}}).

\textbf{Worked example ($n=2$).} Every element is $F=a+b\theta_0+c\theta_1+d\theta_0\theta_1$ with complex $a,b,c,d$. For a second element $F'=a'+b'\theta_0+c'\theta_1+d'\theta_0\theta_1$,

\[
FF'=aa'+(ab'+ba')\theta_0+(ac'+ca')\theta_1+(ad'+da'+bc'-cb')\theta_0\theta_1 .
\]

The term $bc'-cb'$ comes from $b\theta_0\,c'\theta_1=bc'\,\theta_0\theta_1$ and $c\theta_1\,b'\theta_0=cb'\,\theta_1\theta_0=-cb'\,\theta_0\theta_1$; all products with a repeated generator vanish. In particular $(b\theta_0+c\theta_1)^2=bc\,\theta_0\theta_1+cb\,\theta_1\theta_0=0$.

\textbf{Even and odd.} A monomial of $k$ generators has \textbf{degree} $k$; it is \textbf{even} if $k$ is even and \textbf{odd} if $k$ is odd, and an element is even (odd) if all its monomials are. Moving a monomial of degree $k$ past one of degree $l$ takes $kl$ exchanges, so

\[
XY=(-1)^{kl}\,YX\qquad\text{for monomials of degrees }k,l .
\]

Hence an even element commutes with every element, and two odd elements anticommute. An odd element $X$ satisfies $X^2=-X^2$, so $X^2=0$.

\textbf{Worked example: powers of a sum of commuting pairs.} Let $P_0,\dots,P_{15}$ be 16 products of two generators each, with no generator shared between two of them, for example $P_a=\theta_{2a}\theta_{2a+1}$ with 32 generators. Each $P_a$ is even, so they all commute with each other, and $P_a^2=0$. For $S=P_0+P_1+\dots+P_{15}$ the multinomial expansion therefore keeps only products of \textbf{distinct} $P$'s, each appearing $k!$ times in $S^k$ (once for every order):

\[
S^k=k!\sum_{a_1<a_2<\dots<a_k}P_{a_1}P_{a_2}\cdots P_{a_k}.
\]

So $S^k$ has $\binom{16}{k}$ monomials with coefficients $\pm k!$, $S^{16}=16!\,P_0P_1\cdots P_{15}$ is a single monomial, and $S^{17}=0$. With two pairs only: $S=P_0+P_1$, $S^2=P_0P_1+P_1P_0=2P_0P_1$ and $S^3=0$. This is exactly the structure of the scalar density $S=\bar\Psi\Psi$ of dirac16complex at one point (Chapter 6): there the pairs are $\Psi_a^\ast\Psi_b$, and the repository computed all powers in an explicit Grassmann algebra with 32 generators, finding 16, 120, 560, 1820, 4368, 8008, 11440, 12870, 11440, 8008, 4368, 1820, 560, 120, 16 and 1 monomials for $k=1,\dots,16$, coefficients of absolute value $k!$, and $S^{17}=0$ (check \texttt{GR\_quart\allowbreak{}icTermPo\allowbreak{}lynomial} in \texttt{artifact\allowbreak{}s/dirac1\allowbreak{}6complex\allowbreak{}/arbitra\allowbreak{}ry-field\allowbreak{}/grassma\allowbreak{}nn-demo-\allowbreak{}report.j\allowbreak{}son}; Stage-1 document, Result 7.2). A consequence: every function of $S$ is a polynomial of degree at most 16 (a power series in $S$ stops after the term $S^{16}$). The project requires $U(0)=U'(0)=0$ (no constant and no second mass term) and chooses the simplest such polynomial, $U(S)=\tfrac\lambda2S^2$; this is a convention, not a consequence (Stage-1 document, §1.2 and §7.6; Chapter 6).

\textbf{Complex conjugation.} A \textbf{conjugation} on the algebra is a rule $F\mapsto F^\ast$ with $(F+G)^\ast=F^\ast+G^\ast$, $(cF)^\ast=c^\ast F^\ast$ for numbers $c$, $(F^\ast)^\ast=F$, and the \textbf{order-reversing} product rule

\[
(FG)^\ast=G^\ast F^\ast .
\]

A generator with $\theta^\ast=\theta$ is called \textbf{real}. A \textbf{complex} Grassmann variable is described by two generators $\theta$ and $\theta^\ast$ that are exchanged by the conjugation. Two consequences are worth seeing once. The product of two real generators is imaginary: $(\theta_0\theta_1)^\ast=\theta_1^\ast\theta_0^\ast=\theta_1\theta_0=-\theta_0\theta_1$, so $i\theta_0\theta_1$ is real. The product $\theta^\ast\theta$ of a complex generator with its conjugate is real: $(\theta^\ast\theta)^\ast=\theta^\ast\theta^{\ast\ast}=\theta^\ast\theta$. More generally, for a column $\Psi=(\Psi_0,\dots,\Psi_{n-1})^T$ of complex odd generators, $\Psi^\dagger=(\Psi^\ast)^T$ and a numerical matrix $M$,

\[
\bigl(\Psi^\dagger M\Psi\bigr)^\ast=\sum_{a,b}M_{ab}^\ast\,\Psi_b^\ast\Psi_a=\Psi^\dagger M^\dagger\Psi,
\]

so a \textbf{bilinear} $\Psi^\dagger M\Psi$ is real, that is, equal to its own conjugate (this book also calls such an element \textbf{Hermitian}, as the Stage-1 document does in its §7.4), exactly when the matrix $M$ is Hermitian, $M^\dagger=M$, just as for ordinary numbers. These rules, including the order reversal, were checked in an explicit Grassmann algebra (check \texttt{GR\_conju\allowbreak{}gationRu\allowbreak{}les}, \texttt{grassman\allowbreak{}n-demo-r\allowbreak{}eport.js\allowbreak{}on}).

\subsection{Calculus with Grassmann numbers}

\textbf{Grassmann-valued fields.} A Grassmann field assigns to every point $x$ odd elements $\Psi_a(x)$; its derivatives $\partial_\mu\Psi_a(x)$ are odd as well. At one point, the values $\Psi_a$, the first derivatives $\partial_\mu\Psi_a$ and the second derivatives $\partial_\mu\partial_\nu\Psi_a$ ($\mu\le\nu$) can be treated as independent generators, called \textbf{jets}. For 16 real components this is $16\times(1+8+36)=720$ generators, and for a complex field with $\Psi$ and $\Psi^\ast$ it is 1440; these are the generator counts of the repository's Grassmann computations (measurement \texttt{GR\_gener\allowbreak{}atorCoun\allowbreak{}ts} in \texttt{grassman\allowbreak{}n-demo-r\allowbreak{}eport.js\allowbreak{}on}; the class \texttt{\detokenize{JetSpace}} in \texttt{scripts/\allowbreak{}grassman\allowbreak{}n\_algebr\allowbreak{}a.py}). The total derivative $\partial_\mu$ acts on jets by raising the derivative order, $\partial_\mu(\partial_\nu\Psi_a)=\partial_\mu\partial_\nu\Psi_a$, and on products by the ordinary product rule, with no extra signs, because $\partial_\mu$ itself is even (it does not change the degree).

\textbf{Left and right derivatives.} Let $F$ be an element and $\theta_k$ a generator. The \textbf{left derivative} $\partial_LF/\partial\theta_k$ is computed monomial by monomial: move $\theta_k$ to the far left, collecting a factor $-1$ for every generator it passes, and then delete it; monomials without $\theta_k$ give zero. The \textbf{right derivative} $\partial_RF/\partial\theta_k$ moves $\theta_k$ to the far right instead. Examples: for $F=\theta_0\theta_1$, $\partial_LF/\partial\theta_0=\theta_1$, $\partial_LF/\partial\theta_1=-\theta_0$ (write $F=-\theta_1\theta_0$), $\partial_RF/\partial\theta_1=\theta_0$ and $\partial_RF/\partial\theta_0=-\theta_1$.

\textbf{Proposition 5.3.} For an even element $F$, $\partial_RF/\partial\theta_k=-\partial_LF/\partial\theta_k$.

\emph{Proof.} In a monomial of even degree $n$ let $\theta_k$ stand at position $p$ (counting from 0). Moving it to the left passes $p$ generators, moving it to the right passes $n-1-p$. The two signs differ by $(-1)^{n-1-2p}=(-1)^{n-1}=-1$. $\square$

\textbf{The first-order change.} For an odd ``small'' change $\delta\theta_k$ of one generator,

\[
F(\theta_k+\delta\theta_k)=F(\theta_k)+\delta\theta_k\,\frac{\partial_LF}{\partial\theta_k}\qquad\text{exactly, when }(\delta\theta_k)^2=0 .
\]

\emph{Proof.} Write $F=G+\theta_kH$ with $G$, $H$ free of $\theta_k$ (move $\theta_k$ to the left in every monomial that contains it). Then $H=\partial_LF/\partial\theta_k$ and $F(\theta_k+\delta\theta_k)=G+\theta_kH+\delta\theta_kH$. $\square$ This is why the left derivative is the natural one when the variation is written on the left.

\textbf{Euler--Lagrange equations for Grassmann fields.} Let $\mathcal L$ be an even element built from odd fields $\Psi_a$ and their derivatives (a Lagrangian must be even: every term contains the odd fields in pairs, so that it commutes with everything like an ordinary function). Vary $\Psi_a\to\Psi_a+\varepsilon\,\xi(x)$ for one component $a$, where $\varepsilon$ is a new odd generator (so $\varepsilon^2=0$) and $\xi$ an ordinary bump function. By the first-order formula, and because $\xi$ is an ordinary function,

\[
\delta\mathcal L=\varepsilon\,\xi\,\frac{\partial_L\mathcal L}{\partial\Psi_a}+\varepsilon\,(\partial_\mu\xi)\,\frac{\partial_L\mathcal L}{\partial(\partial_\mu\Psi_a)} .
\]

Integrating by parts exactly as in Section 5.3 gives $\delta S=\varepsilon\int\xi\,E_a\,d^8x$ with

\[
E_a=\frac{\partial_L\mathcal L}{\partial\Psi_a}-\partial_\mu\frac{\partial_L\mathcal L}{\partial(\partial_\mu\Psi_a)} .
\]

Since $\varepsilon$ does not occur in $E_a$, $\varepsilon\int\xi E_a\,d^8x=0$ for all $\xi$ forces $E_a=0$: these are the field equations. With right derivatives, and the variation written on the right, one obtains the expression $E^R_a$; by Proposition 5.3, applied to the even $\mathcal L$, $E^R_a=-E_a$, so the two conventions give the same equations with opposite overall signs. (The Stage-1 document uses left derivatives for $\Psi^\dagger$ and right derivatives for $\Psi$; its §8.1 records the sign.) For a complex Grassmann field one varies $\Psi_a$ and $\Psi_a^\ast$ independently, by the same linear-substitution argument as in Section 5.6.

\textbf{Worked example (a complex Grassmann oscillator).} Let $\theta(t)$ be a complex odd variable and

\[
L=\tfrac i2\bigl(\theta^\ast\dot\theta-\dot\theta^\ast\theta\bigr)-\omega\,\theta^\ast\theta .
\]

This $L$ is real by the conjugation rules of Section 5.9. Varying $\theta^\ast$ (left derivative): $\partial_LL/\partial\theta^\ast=\tfrac i2\dot\theta-\omega\theta$, because $\theta^\ast$ already stands on the left; and $\partial_LL/\partial\dot\theta^\ast=-\tfrac i2\theta$. So $E=\tfrac i2\dot\theta-\omega\theta+\tfrac i2\dot\theta=i\dot\theta-\omega\theta$, and the field equation $i\dot\theta=\omega\theta$ has the same form as for the ordinary complex oscillator of Section 5.6. For \textbf{complex} Grassmann fields written with every conjugate factor to the left, the anticommutation does not change the form of the equations. For \textbf{real} Grassmann fields it changes everything, as the next section shows.

\subsection{Bilinear forms in Grassmann variables: two lemmas}

Let $\Psi=(\Psi_0,\dots,\Psi_{n-1})^T$ be a column of \textbf{real} odd fields and $M$ an $n\times n$ matrix of ordinary numbers (or of ordinary functions of $x$). Write $\Psi^TM\Psi=\sum_{a,b}\Psi_aM_{ab}\Psi_b$, and split $M$ into its symmetric and antisymmetric parts, $M_S=\tfrac12(M+M^T)$ and $M_A=\tfrac12(M-M^T)$.

\textbf{Lemma 5.4.} $\Psi^TM\Psi=\Psi^TM_A\Psi$. In particular $\Psi^TM\Psi=0$ for every symmetric $M$.

\emph{Proof.} Exchange the two odd factors and rename the summation indices: $\Psi^TM\Psi=\sum_{a,b}\Psi_aM_{ab}\Psi_b=-\sum_{a,b}\Psi_bM_{ab}\Psi_a=-\sum_{a,b}\Psi_a(M^T)_{ab}\Psi_b=-\Psi^TM^T\Psi$. Hence $\Psi^TM_S\Psi=\tfrac12(\Psi^TM\Psi+\Psi^TM^T\Psi)=0$. $\square$

\emph{Worked example.} For $n=2$ and $M=\begin{pmatrix}1&2\\3&4\end{pmatrix}$: $\Psi^TM\Psi=\Psi_0\Psi_0+2\Psi_0\Psi_1+3\Psi_1\Psi_0+4\Psi_1\Psi_1=(2-3)\Psi_0\Psi_1=-\Psi_0\Psi_1$. The antisymmetric part is $M_A=\begin{pmatrix}0&-1/2\\1/2&0\end{pmatrix}$, and $\Psi^TM_A\Psi=-\tfrac12\Psi_0\Psi_1+\tfrac12\Psi_1\Psi_0=-\Psi_0\Psi_1$, the same. For commuting variables the opposite holds: $q^TMq=q^TM_Sq$, since $q^TM_Aq=0$.

\textbf{Lemma 5.5.} If $A(x)$ is an antisymmetric matrix function, $A^T=-A$, then

\[
\Psi^TA\,\partial_\mu\Psi=\tfrac12\,\partial_\mu\bigl(\Psi^TA\Psi\bigr)-\tfrac12\,\Psi^T(\partial_\mu A)\Psi .
\]

\emph{Proof.} The product rule gives $\partial_\mu(\Psi^TA\Psi)=(\partial_\mu\Psi)^TA\Psi+\Psi^T(\partial_\mu A)\Psi+\Psi^TA\,\partial_\mu\Psi$. In the first term exchange the two odd factors: $(\partial_\mu\Psi)^TA\Psi=\sum_{a,b}\partial_\mu\Psi_aA_{ab}\Psi_b=-\sum_{a,b}\Psi_bA_{ab}\partial_\mu\Psi_a=-\Psi^TA^T\partial_\mu\Psi=\Psi^TA\,\partial_\mu\Psi$. So $\partial_\mu(\Psi^TA\Psi)=2\Psi^TA\,\partial_\mu\Psi+\Psi^T(\partial_\mu A)\Psi$, which is the claim. $\square$

So for real Grassmann fields a first-order kinetic term with an antisymmetric matrix is \textbf{half a total derivative plus a term without derivatives}. For a constant $A$ it is a total derivative and has no field equations at all (Proposition 5.1; the proof by the action works unchanged for Grassmann fields, with $\varepsilon\xi$ as the variation).

\emph{Worked example.} One coordinate $x$, two real odd fields $\theta_0(x)$, $\theta_1(x)$, a prime for $d/dx$, and $A=\begin{pmatrix}0&1\\-1&0\end{pmatrix}$. Then

\[
\theta^TA\theta'=\theta_0\theta_1'-\theta_1\theta_0'=\theta_0\theta_1'+\theta_0'\theta_1=(\theta_0\theta_1)'=\tfrac12\bigl(\theta^TA\theta\bigr)',
\]

because $\theta^TA\theta=\theta_0\theta_1-\theta_1\theta_0=2\theta_0\theta_1$. Directly: for $\theta_0$, $\partial_L/\partial\theta_0$ of $\theta_0\theta_1'-\theta_1\theta_0'$ is $\theta_1'$, and $\partial_L/\partial\theta_0'$ is $+\theta_1$ (write $-\theta_1\theta_0'=\theta_0'\theta_1$), so $E_0=\theta_1'-(\theta_1)'=0$; similarly $E_1=-\theta_0'-(-\theta_0)'=0$. Compare Example A of Section 5.6: the same expression with commuting $q_0,q_1$ gave the non-trivial expressions $\dot q_1$ and $-\dot q_0$. Now take instead the symmetric identity matrix $I_2$: $\theta^TI_2\theta=\theta_0\theta_0+\theta_1\theta_1=0$, but $\theta^TI_2\theta'=\theta_0\theta_0'+\theta_1\theta_1'$ is not a derivative, and its Euler--Lagrange expression for $\theta_0$ is $\theta_0'-(-\theta_0)'=2\theta_0'\ne0$. The repository found the same contrast for the $16\times16$ matrix $C$ of Chapter 2: with $C$ in place of an antisymmetric matrix the Euler--Lagrange expression is $2(C\partial_0\Psi)_0$ (check \texttt{GR\_kinet\allowbreak{}icSymmet\allowbreak{}ricMatri\allowbreak{}xContras\allowbreak{}t}).

The two lemmas side by side:

\begingroup
\small
\begin{longtable}{@{}>{\raggedright\arraybackslash}p{0.267\linewidth}>{\raggedright\arraybackslash}p{0.267\linewidth}>{\raggedright\arraybackslash}p{0.267\linewidth}@{}}
\toprule
\textbf{property} & \textbf{commuting real fields} & \textbf{real Grassmann fields} \\
\midrule
\endfirsthead
\toprule
\textbf{property} & \textbf{commuting real fields} & \textbf{real Grassmann fields} \\
\midrule
\endhead
part of $M$ that survives in $\Psi^TM\Psi$ & symmetric $M_S$ & antisymmetric $M_A$ \\
mass-type term $\Psi^TC\Psi$ with symmetric $C$ & survives & vanishes identically \\
kinetic term $\Psi^TA\,\partial\Psi$ with antisymmetric $A$ & genuine first-order dynamics & total derivative plus a derivative-free term \\
\bottomrule
\end{longtable}
\endgroup

\subsection{The author's Lagrangian Lg[] and the matrices it uses}

\textbf{The matrices.} Chapter 2 constructs the eight real $16\times16$ integer matrices $\gamma^0,\dots,\gamma^7$ of the notebook (its T16\textasciicircum{}A) and proves their properties. We need the following definitions and facts, all facts verified exactly by the repository (Stage-1 document, Results 3.1, 3.3, 3.4 and 3.5; checks \texttt{ALG\_clif\allowbreak{}ford}, \texttt{ALG\_gamm\allowbreak{}aTranspo\allowbreak{}seSymmet\allowbreak{}ry}, \texttt{ALG\_char\allowbreak{}geMatrix}, \texttt{ALG\_expr\allowbreak{}ession1} and \texttt{ALG\_spin\allowbreak{}Transpos\allowbreak{}ePropert\allowbreak{}ies} in \texttt{wolfram-\allowbreak{}algebra-\allowbreak{}report.j\allowbreak{}son} and \texttt{python-a\allowbreak{}lgebra-r\allowbreak{}eport.js\allowbreak{}on}; all Stage-1 reports are in \texttt{artifact\allowbreak{}s/dirac1\allowbreak{}6complex\allowbreak{}/arbitra\allowbreak{}ry-field\allowbreak{}/}):

\begin{itemize}
\item the Clifford relation $\gamma^a\gamma^b+\gamma^b\gamma^a=2\eta^{ab}I_{16}$, so different gammas anticommute, $(\gamma^a)^2=+1$ for $a\le3$ and $(\gamma^a)^2=-1$ for $a\ge4$;
\item $\gamma^a$ is symmetric for $a\le3$ and antisymmetric for $a\ge4$;
\item the charge matrix $C:=\gamma^0\gamma^1\gamma^2\gamma^3$ (the notebook's σ16) is symmetric and $C^2=I_{16}$;
\item the spin matrices $S^{ab}=\tfrac14(\gamma^a\gamma^b-\gamma^b\gamma^a)$ generate the rotations of spinors (Chapter 2);
\item expression [1] of the task: $(C\gamma^a)^T=-C\gamma^a$ for every $a$.
\end{itemize}

The properties of $C$ and expression [1] follow from the first two facts, and we prove them because the whole theorem rests on them. \emph{$C$ is symmetric:} $C^T=(\gamma^3)^T(\gamma^2)^T(\gamma^1)^T(\gamma^0)^T=\gamma^3\gamma^2\gamma^1\gamma^0$, and reversing four mutually anticommuting factors takes $3+2+1=6$ exchanges, so $C^T=(+1)C$. \emph{$C^2=1$:} $C\cdot C=C\cdot C^T=\gamma^0\gamma^1\gamma^2\gamma^3\gamma^3\gamma^2\gamma^1\gamma^0=1$, collapsing the squares $(\gamma^a)^2=1$ from the middle. \emph{Expression [1]:} for $a\le3$, $(C\gamma^a)^T=(\gamma^a)^TC^T=\gamma^aC$, and $\gamma^a$ anticommutes with the three factors of $C$ other than itself and commutes with itself, so $\gamma^aC=(-1)^3C\gamma^a=-C\gamma^a$. For $a\ge4$, $(C\gamma^a)^T=-\gamma^aC$, and $\gamma^a$ anticommutes with all four factors, so $\gamma^aC=C\gamma^a$ and again $(C\gamma^a)^T=-C\gamma^a$. $\square$

\textbf{Two consequences.} Multiplying $(\gamma^a)^TC=-C\gamma^a$ on the right by $C$ gives $(\gamma^a)^T=-C\gamma^aC$. Then $(S^{ab})^T=\tfrac14[(\gamma^b)^T,(\gamma^a)^T]=\tfrac14C[\gamma^b,\gamma^a]C=-CS^{ab}C$, hence

\[
(CS^{ab})^T=(S^{ab})^TC=-CS^{ab},\qquad (C\gamma^cS^{ab})^T=(S^{ab})^T(\gamma^c)^TC=(-CS^{ab}C)(-C\gamma^cC)C=CS^{ab}\gamma^c .
\]

So the matrix $C\gamma^cS^{ab}$ splits into a symmetric part $\tfrac12C\{\gamma^c,S^{ab}\}$ and an antisymmetric part $\tfrac12C[\gamma^c,S^{ab}]$, where $\{X,Y\}=XY+YX$ is the \textbf{anticommutator} and $[X,Y]=XY-YX$ the \textbf{commutator}. This is Lemma 6.3 of the Stage-1 document (checks \texttt{ALG\_CAnt\allowbreak{}icommuta\allowbreak{}torGamma\allowbreak{}SSymmetr\allowbreak{}ic} and \texttt{ALG\_CCom\allowbreak{}mutatorG\allowbreak{}ammaSAnt\allowbreak{}isymmetr\allowbreak{}ic}, \texttt{python-g\allowbreak{}eometry-\allowbreak{}report.j\allowbreak{}son}).

\textbf{Worked example with the actual matrices.} The Student Guide (\texttt{provenan\allowbreak{}ce/DIRAC\allowbreak{}16COMPLE\allowbreak{}X\_STUDEN\allowbreak{}T\_GUIDE.\allowbreak{}md}, its §6.3 and §6.4) lists every gamma and $C$ as signed permutations: row 0 of $C$ reads ``$-4$'', meaning $(Cu)_0=-u_4$, and row 4 reads ``$-0$'', meaning $(Cu)_4=-u_0$. So $C_{0,4}=C_{4,0}=-1$, and in $\Psi^TC\Psi$ these two entries give $-\Psi_0\Psi_4-\Psi_4\Psi_0=0$; all 16 entries of $C$ cancel in pairs in the same way, and $\Psi^TC\Psi=0$. For $C\gamma^0$: row 4 of $\gamma^0$ reads ``$+12$'', so $(C\gamma^0u)_0=-(\gamma^0u)_4=-u_{12}$, while row 12 of $C$ reads ``$+8$'' and row 8 of $\gamma^0$ reads ``$+0$'', so $(C\gamma^0u)_{12}=+(\gamma^0u)_8=+u_0$. The entries $-1$ at $(0,12)$ and $+1$ at $(12,0)$ give $-\Psi_0\Psi_{12}+\Psi_{12}\Psi_0=-2\Psi_0\Psi_{12}$. Doing this for all rows,

\[
\Psi^TC\gamma^0\Psi=-2\bigl(\Psi_0\Psi_{12}+\Psi_1\Psi_{13}+\Psi_2\Psi_{14}+\Psi_3\Psi_{15}+\Psi_4\Psi_8+\Psi_5\Psi_9+\Psi_6\Psi_{10}+\Psi_7\Psi_{11}\bigr).
\]

The repository's Grassmann algebra finds exactly this pattern: $\Psi^TC\Psi$ has 0 monomials, and each of the eight $\Psi^TC\gamma^a\Psi$ has 8 monomials with coefficients $\pm2$ (the three measurements of \texttt{grassman\allowbreak{}n-demo-r\allowbreak{}eport.js\allowbreak{}on} whose names begin with \texttt{GR\_massT\allowbreak{}ermVanis\allowbreak{}hesReal\_}).

\textbf{Geometry.} From Chapter 4 we need: the vielbein $e_\mu{}^a(x)$ and its inverse $e_a{}^\mu$, the metric $g_{\mu\nu}=e_\mu{}^a\eta_{ab}e_\nu{}^b$, the curved gammas $\gamma^\mu=e_a{}^\mu\gamma^a$ (ordinary functions times constant matrices), the canonical spin connection $\omega_{\mu ab}=-\omega_{\mu ba}$ fixed by the vielbein postulate, the spinor connection $\Omega_\mu=\tfrac12\omega_{\mu ab}S^{ab}$, and the covariant constancy of the gammas,

\[
D_\mu\gamma^\nu:=\partial_\mu\gamma^\nu+\Gamma^\nu{}_{\mu\lambda}\gamma^\lambda+[\Omega_\mu,\gamma^\nu]=0,
\]

which is equivalent to the vielbein postulate (Stage-1 document, Result 5.2; checks \texttt{GEO\_gamm\allowbreak{}aCovaria\allowbreak{}ntConsta\allowbreak{}ncy\_G1} and \texttt{\detokenize{_G2}}).

\textbf{The divergence identity.} Set $\nu=\mu$ in $D_\mu\gamma^\nu=0$ and sum: $\partial_\mu\gamma^\mu+\Gamma^\mu{}_{\mu\lambda}\gamma^\lambda+[\Omega_\mu,\gamma^\mu]=0$. Multiply by $\sqrt{|g|}$ and use $\partial_\lambda\sqrt{|g|}=\sqrt{|g|}\,\Gamma^\mu{}_{\mu\lambda}$ from Section 5.8: the first two terms become $\partial_\mu(\sqrt{|g|}\gamma^\mu)$. Hence

\[
\partial_\mu\bigl(\sqrt{|g|}\,\gamma^\mu\bigr)=\sqrt{|g|}\,[\gamma^\mu,\Omega_\mu] .
\]

This is Result 5.3 of the Stage-1 document (checks \texttt{GEO\_dive\allowbreak{}rgenceId\allowbreak{}entity\_G\allowbreak{}1} and \texttt{GEO\_dive\allowbreak{}rgenceId\allowbreak{}entity\_G\allowbreak{}2} in both geometry reports).

\textbf{The notebook's Lagrangian.} Cell 1064 of the notebook defines $\mathrm{Lg}[\,]$. Verbatim, as transcribed in the Stage-1 document (§6.1; the line breaks are only for the page width):

\begin{Verbatim}[fontsize=\small]
Lg[]:=Sqrt[detgg] *( Transpose[\[CapitalPsi]16].\[Sigma]16.
Sum[FullSimplify[((T16^\[Alpha])[\[Alpha]1-1]/.sg),constraintVars].
(D[ \[CapitalPsi]16,X[[\[Alpha]1]]]+(Q1/2)*Sum[\[Omega]mat[[\[Alpha]1,a,b]]*SAB[[a,
b]].\[CapitalPsi]16,{a,1,8},{b,1,8}]),{\[Alpha]1,1,Length[X]}]+
(H*M)*Transpose[\[CapitalPsi]16].\[Sigma]16.\[CapitalPsi]16)//Simplify[#,
constraintVars]&
\end{Verbatim}

\textbf{Reading the code.} In the code, \texttt{\textbackslash{}[Capita\allowbreak{}lPsi]}, \texttt{\textbackslash{}[Sigma]}, \texttt{\textbackslash{}[Alpha]} and \texttt{\textbackslash{}[Omega]} are Mathematica's names for the Greek letters Ψ, σ, α and ω. Mathematica lists count from 1, so the summation index α1 $=1,\dots,8$ stands for the coordinate index $\mu=\alpha1-1=0,\dots,7$, and the indices a, b $=1,\dots,8$ of \texttt{ωmat} and \texttt{\detokenize{SAB}} stand for the frame indices $0,\dots,7$. The pieces of the cell mean (Stage-1 document, §2.2 and §6.1):

\begin{itemize}
\item \texttt{Sqrt[det\allowbreak{}gg]} is $\sqrt{|g|}$. Here \texttt{\detokenize{detgg}} is \texttt{Det[g448\allowbreak{}8 /. sg]} (cell 275), the determinant of the notebook's metric \texttt{\detokenize{g4488}} (cell 147). It is positive, because in signature (4,4) $\det g=(\det e)^2\det\eta=(\det e)^2$ (Section 4.10).
\item \texttt{\detokenize{sg}} is the notebook's list of rules that turns the metric symbols into the diagonal functions of $(x_0,x_4)$ of its field (cell 269), and \texttt{\detokenize{/.sg}} applies these rules.
\item \texttt{Transpos\allowbreak{}e[Ψ16]} is the row $\Psi^T$, \texttt{σ16} is $C$, and the dot \texttt{\detokenize{.}} is the matrix product.
\item \texttt{(T16\textasciicircum{}α)[\allowbreak{}α1-1]} is the curved gamma $\gamma^\mu$ with $\mu=\alpha1-1$ (cell 475, built from the notebook's diagonal vielbein; Section 4.10).
\item \texttt{\detokenize{X}} is the list of the eight coordinates (cell 61), so \texttt{D[Ψ16, X\allowbreak{}[[α1]]]} is $\partial_\mu\Psi$, and the outer \texttt{Sum[...,\allowbreak{} \{α1, 1,\allowbreak{} Length[\allowbreak{}X]\}]} is the sum over the eight values of $\mu$.
\item \texttt{ωmat[[α1\allowbreak{},a,b]]} is the mixed spin connection $\omega_\mu{}^{a-1}{}_{b-1}$ (cells 529 to 532) and \texttt{SAB[[a,b\allowbreak{}]]} is $S^{a-1\,b-1}$, so \texttt{(Q1/2)*S\allowbreak{}um[...]} is Q1 times $\tfrac12\,\omega_\mu{}^a{}_b\,S^{ab}\Psi$, with the indices counted from 0.
\item \texttt{\detokenize{Q1}} is the notebook's book-keeping switch for the spin connection (1 switches it on), and \texttt{\detokenize{H}} and \texttt{\detokenize{M}} are the notebook's inverse length and mass.
\item \texttt{FullSimp\allowbreak{}lify[...\allowbreak{}, constr\allowbreak{}aintVars\allowbreak{}]} and \texttt{Simplify\allowbreak{}[\#, cons\allowbreak{}traintVa\allowbreak{}rs]\&} only simplify the expression, with the positivity assumptions listed in \texttt{constrai\allowbreak{}ntVars} (cell 73); they do not change its value.
\end{itemize}

In the notation of this book, with the switch Q1 = 1 and the mass $m:=-HM$ (so that $+HM\,\Psi^TC\Psi=-m\,\Psi^TC\Psi$),

\[
\mathrm{Lg}=\sqrt{|g|}\,\Bigl[\Psi^TC\gamma^\mu\bigl(\partial_\mu\Psi+\Omega^{\mathrm{nb}}_\mu\Psi\bigr)-m\,\Psi^TC\Psi\Bigr],\qquad \Omega^{\mathrm{nb}}_\mu=\tfrac12\,\omega_\mu{}^a{}_b\,S^{ab}.
\]

Here $\Psi$ is the notebook's Ψ16, a \textbf{real} column of 16 components f16[k]. The notebook contracts the mixed connection $\omega_\mu{}^a{}_b$ with $S^{ab}$; one factor of $\eta$ is missing, and Chapter 4 shows that this $\Omega^{\mathrm{nb}}$ is not the correct spinor connection. The notebook treats the f16[k] as ordinary commuting functions; a search of the notebook finds no Grassmann or anticommuting variable anywhere (\texttt{handoff/\allowbreak{}surveys/\allowbreak{}survey\_n\allowbreak{}otebook-\allowbreak{}physics.\allowbreak{}md}, its item 6). A fermion field, however, must have anticommuting components. We therefore ask what $\mathrm{Lg}[\,]$ gives for a real Grassmann field, first with the correct connection $\Omega_\mu$ in place of $\Omega^{\mathrm{nb}}_\mu$, then with the notebook's own.

\subsection{Theorem: Lg[] is a pure divergence for a real Grassmann field}

\textbf{Theorem 5.6.} Let $\Psi$ be a real 16-component Grassmann-odd field in an arbitrary gravitational field, and let the connection in $\mathrm{Lg}$ be the canonical $\Omega_\mu$. Then

\[
\mathrm{Lg}=\partial_\mu V^\mu,\qquad V^\mu=\tfrac12\sqrt{|g|}\,\Psi^TC\gamma^\mu\Psi,
\]

and all 16 Euler--Lagrange expressions of $\mathrm{Lg}$ vanish identically. The field equations of $\mathrm{Lg}[\,]$ are $0=0$.

\emph{Proof.} We treat the three terms of $\mathrm{Lg}$ one by one.

\emph{Step 1 (mass term).} $C$ is symmetric, so $\Psi^TC\Psi=0$ by Lemma 5.4. The mass term vanishes identically, whatever $m$ is.

\emph{Step 2 (derivative term).} Put $A^\mu:=\sqrt{|g|}\,C\gamma^\mu=\sqrt{|g|}\,e_a{}^\mu\,C\gamma^a$. It is a sum of the antisymmetric matrices $C\gamma^a$ (expression [1]) with ordinary functions as coefficients, hence antisymmetric. Lemma 5.5 gives

\[
\sqrt{|g|}\,\Psi^TC\gamma^\mu\partial_\mu\Psi=\tfrac12\,\partial_\mu\bigl(\Psi^TA^\mu\Psi\bigr)-\tfrac12\,\Psi^TC\,\partial_\mu\bigl(\sqrt{|g|}\gamma^\mu\bigr)\Psi .
\]

\emph{Step 3 (connection term).} $\Omega_\mu=\tfrac12\omega_{\mu ab}S^{ab}$ and $\gamma^\mu=e_c{}^\mu\gamma^c$, so $C\gamma^\mu\Omega_\mu$ is a combination, with ordinary-function coefficients, of the matrices $C\gamma^cS^{ab}$. By Section 5.12 its symmetric part is $\tfrac12C\{\gamma^\mu,\Omega_\mu\}$ and its antisymmetric part is $\tfrac12C[\gamma^\mu,\Omega_\mu]$. By Lemma 5.4 only the antisymmetric part survives:

\[
\sqrt{|g|}\,\Psi^TC\gamma^\mu\Omega_\mu\Psi=\tfrac12\sqrt{|g|}\,\Psi^TC[\gamma^\mu,\Omega_\mu]\Psi=\tfrac12\,\Psi^TC\,\partial_\mu\bigl(\sqrt{|g|}\gamma^\mu\bigr)\Psi,
\]

where the last step is the divergence identity of Section 5.12.

\emph{Step 4 (cancellation).} The result of Step 3 is exactly minus the second term of Step 2. Adding Steps 1 to 3,

\[
\mathrm{Lg}=\tfrac12\,\partial_\mu\bigl(\Psi^TA^\mu\Psi\bigr)=\partial_\mu V^\mu .
\]

\emph{Step 5 (no field equations).} $V^\mu$ depends on the field $\Psi$ and on $x$ (through $\sqrt{|g|}$ and $e_a{}^\mu$) but not on derivatives of $\Psi$. By Proposition 5.1, in its Grassmann form (Section 5.11), every Euler--Lagrange expression of $\partial_\mu V^\mu$ vanishes identically. $\square$

The proof uses only the algebra of Section 5.12, the two lemmas, and the divergence identity, which holds in \textbf{every} gravitational field with the canonical connection. So the theorem holds in every gravitational field, not only at the points where the repository tested it (Section 5.15).

\textbf{What the theorem means.} $\mathrm{Lg}[\,]$ is the Lagrangian of nothing, for the kind of field that a fermion is. Its mass term is zero, its derivative term is a total derivative, and its spin-connection term exactly cancels the only part of the derivative term that is not a total derivative. This is Theorem 6.1 of the Stage-1 document, the first of the three findings about the notebook that the README lists.

\subsection{The notebook's own contraction, commuting fields, and the way out}

\textbf{With the notebook's contraction.} Repeat the proof with $\Omega^{\mathrm{nb}}_\mu$ in place of $\Omega_\mu$. Steps 1 and 2 do not involve the connection. $\Omega^{\mathrm{nb}}_\mu$ is also a combination of the $S^{ab}$, so Step 3 still gives $\tfrac12\sqrt{|g|}\,\Psi^TC[\gamma^\mu,\Omega^{\mathrm{nb}}_\mu]\Psi$; but the divergence identity holds for $\Omega_\mu$, not for $\Omega^{\mathrm{nb}}_\mu$, so the cancellation of Step 4 is incomplete. What remains is

\[
\mathrm{Lg}=\partial_\mu V^\mu+\tfrac12\,\Psi^TX\Psi,\qquad X=\sqrt{|g|}\;C\sum_{\mu}\bigl[\gamma^\mu,\Omega^{\mathrm{nb}}_\mu-\Omega_\mu\bigr],
\]

with an antisymmetric matrix $X$ (a sum of matrices $C[\gamma^c,S^{ab}]$). The left derivative of $\tfrac12\sum_{a,b}X_{ab}\Psi_a\Psi_b$ with respect to $\Psi_c$ is $\tfrac12\sum_bX_{cb}\Psi_b-\tfrac12\sum_aX_{ac}\Psi_a=(X\Psi)_c$, because $X^T=-X$; no derivative of $\Psi$ appears. So the Euler--Lagrange equations are the \textbf{algebraic} equations

\[
X\Psi=0 .
\]

At every point where the $16\times16$ matrix $X$ is invertible (rank 16) they force $\Psi=0$. The repository computed $X$ exactly at six points, three points of a generic non-diagonal test vielbein called G1 and three points of the notebook's primordial field called G2 (Chapter 9): $X$ has rank 16 at all six, with largest entries of absolute value about 9.927, 12.36 and 10.72 at the G1 points and $3\sqrt{455}/32\approx1.99976$, about 0.5977 and about 2.530 at the G2 points (measurements \texttt{G1.p1.no\allowbreak{}tebookLg\allowbreak{}Residual\allowbreak{}XRank}, \texttt{G1.p1.no\allowbreak{}tebookLg\allowbreak{}Residual\allowbreak{}XMaxAbsD\allowbreak{}ecimal} and their analogues for the other five points in \texttt{artifact\allowbreak{}s/dirac1\allowbreak{}6complex\allowbreak{}/arbitra\allowbreak{}ry-field\allowbreak{}/wolfram\allowbreak{}-geometr\allowbreak{}y-report\allowbreak{}.json}). Either way, $\mathrm{Lg}[\,]$ gives no wave equation for a Grassmann field. That the rank is 16 is a computed fact at these six points, not a theorem for every field.

\textbf{With commuting fields (the notebook's actual use).} For ordinary commuting components the roles of the symmetric and antisymmetric parts are exchanged (Lemma 5.4 and the remark after it). Then the mass term $\Psi^TC\Psi$ survives, the derivative term $\Psi^TC\gamma^\mu\partial_\mu\Psi$ is a genuine first-order kinetic term (as in Example A of Section 5.6), and from the connection term only the symmetric part $\tfrac12C\{\gamma^\mu,\Omega_\mu\}$ survives. That part is small in a precise sense. Write $\omega_{cab}:=e_c{}^\mu\,\omega_{\mu ab}$ for the components of the connection whose first (coordinate) index has also been turned into a frame index with the inverse vielbein. Since $\gamma^\mu=e_c{}^\mu\gamma^c$, we get $\gamma^\mu\Omega_\mu=\tfrac12\,\omega_{cab}\,\gamma^cS^{ab}$ and $\{\gamma^\mu,\Omega_\mu\}=\tfrac12\,\omega_{cab}\,\{\gamma^c,S^{ab}\}$ (summed over $c,a,b$). For distinct $c,a,b$ one has $\gamma^cS^{ab}=\tfrac12\gamma^c\gamma^a\gamma^b=S^{ab}\gamma^c$ (moving $\gamma^c$ through two anticommuting factors), so $\{\gamma^c,S^{ab}\}=\gamma^c\gamma^a\gamma^b$; for $c=a\ne b$ one has $\gamma^aS^{ab}=\tfrac12\eta^{aa}\gamma^b=-S^{ab}\gamma^a$, so $\{\gamma^a,S^{ab}\}=0$, and the same holds for $c=b\ne a$ because $S^{ab}=-S^{ba}$; finally $S^{aa}=0$. Hence $\{\gamma^\mu,\Omega_\mu\}$ contains only the components $\omega_{cab}$ with three different indices, and, because $\gamma^c\gamma^a\gamma^b$ changes sign when two of its three different indices are exchanged, only their totally antisymmetric combination. For a diagonal vielbein $e_\mu{}^a=h_\mu\,\delta_\mu{}^a$ (no sum), such as the notebook's field, these components vanish. Because the vielbein is diagonal, frame and coordinate indices can be identified: $e_c{}^\mu$ is $1/h_c$ for $\mu=c$ and 0 otherwise, so $\omega_{cab}=\omega_{\mu ab}/h_c$ with $\mu=c$. Section 4.11 showed that for a diagonal vielbein $\omega_{\mu ab}$ is nonzero only when $\mu=a$ or $\mu=b$; so $\omega_{cab}=0$ whenever $c,a,b$ are all different. The notebook's contraction uses the mixed components $\omega_\mu{}^a{}_b=\eta_{aa}\,\omega_{\mu ab}$ (no sum) in place of $\omega_{\mu ab}$; they vanish in the same cases, so the argument holds for both contractions (Stage-1 document, §6.5). The repository checks the result, $\sum_\mu\{\gamma^\mu,\Omega_\mu\}=0$ for both contractions, in the notebook's field G2 (check \texttt{GEO\_anti\allowbreak{}commutat\allowbreak{}orGammaO\allowbreak{}megaVani\allowbreak{}shesDiag\allowbreak{}onal\_G2} in \texttt{python-g\allowbreak{}eometry-\allowbreak{}report.j\allowbreak{}son}). So for commuting fields in the notebook's diagonal field the spin connection drops out of $\mathrm{Lg}[\,]$ altogether, for either contraction, and the gravitational term of the field equations comes from $\partial_\mu(\sqrt{|g|}\gamma^\mu)$. The notebook's stored equations contain an additional term, which Chapter 9 reconstructs from a substitution rule in cell 1058; that this rule produced the term is a hypothesis (ledger row L20), not a proved fact. The commuting 16-component field, taken seriously as a classical field, is the second field of this book, dirac16complex00 (Chapter 6).

\textbf{The way out: a complex Grassmann field.} For a \textbf{complex} Grassmann field the Dirac adjoint $\bar\Psi:=\Psi^\dagger C$ pairs the conjugate components $\Psi_a^\ast$ with the components $\Psi_b$. These are independent generators, so Lemma 5.4 does not apply: $\bar\Psi\Psi=\sum_{a,b}\Psi_a^\ast C_{ab}\Psi_b$ has 16 nonzero monomials, one for each nonzero entry of $C$ (Stage-1 document, Results 7.1 and 7.2; checks \texttt{GR\_scala\allowbreak{}rBilinea\allowbreak{}rHermiti\allowbreak{}an} and \texttt{GR\_quart\allowbreak{}icTermPo\allowbreak{}lynomial}), whereas the real $\Psi^TC\Psi$ has none. Chapter 6 builds the Lagrangian of dirac16complex from $\bar\Psi$:

\[
\mathcal L=\sqrt{|g|}\,\Bigl[\tfrac12\bigl(\bar\Psi\gamma^\mu D_\mu\Psi-(D_\mu\bar\Psi)\gamma^\mu\Psi\bigr)-m\,\bar\Psi\Psi-U(\bar\Psi\Psi)\Bigr],\qquad D_\mu\Psi=\partial_\mu\Psi+\Omega_\mu\Psi .
\]

Its field equation $\gamma^\mu D_\mu\Psi=(m+U'(\bar\Psi\Psi))\Psi$ is derived in Chapter 7. In the repository's Grassmann algebra at the G1 point p1 the Euler--Lagrange expression of this $\mathcal L$ for component 0 contains 32 terms with derivatives of the field and 8 spin-connection terms, and $\gamma^\mu\Omega_\mu$ has 128 nonzero entries there (measurements \texttt{GR\_compl\allowbreak{}exLagran\allowbreak{}gianNonT\allowbreak{}rivial\_E\allowbreak{}L0\_deriv\allowbreak{}ativeJet\allowbreak{}Terms\_G1\allowbreak{}\_p1}, \texttt{\_EL0\_spi\allowbreak{}nConnect\allowbreak{}ionTerms\allowbreak{}\_G1\_p1} and \texttt{\_gammaMu\allowbreak{}OmegaMu\_\allowbreak{}nonzeroE\allowbreak{}ntries\_G\allowbreak{}1\_p1}). Unlike $\mathrm{Lg}[\,]$, it is a genuine field theory.

\subsection{How the repository checks all of this}

The statements of Sections 5.12 to 5.14 were checked by exact programs (Stage-1 document, §3.4, §6.4, §6.5 and §11). The matrix identities were checked by the Wolfram and the Python algebra checkers; the split of $C\gamma^cS^{ab}$ was also checked by the Python geometry checker. The covariant constancy of the gammas and the divergence identity were checked by both geometry verifiers. Theorem 5.6 and the form $X\Psi$ of the notebook-contraction equations were checked in genuine Grassmann algebras by two independent programs, the Python demonstration and the Wolfram geometry verifier. Some items come from one program only. The rank and the largest entries of $X$ come from the Wolfram geometry verifier. The monomial counts of $\Psi^TC\gamma^a\Psi$ and of the complex Lagrangian come from the Python demonstration. The vanishing of $\{\gamma^\mu,\Omega_\mu\}$ in the diagonal field G2 comes from the Python geometry checker. The remark about the notebook's stored equations is derived, not machine-checked. The main checks:

\begingroup
\small
\begin{longtable}{@{}>{\raggedright\arraybackslash}p{0.400\linewidth}>{\raggedright\arraybackslash}p{0.400\linewidth}@{}}
\toprule
\textbf{statement} & \textbf{check (report)} \\
\midrule
\endfirsthead
\toprule
\textbf{statement} & \textbf{check (report)} \\
\midrule
\endhead
$\Psi^TC\Psi=0$; $\Psi^TC\gamma^a\Psi$ has 8 monomials & \texttt{GR\_massT\allowbreak{}ermVanis\allowbreak{}hesReal} (grassmann-demo) \\
only the antisymmetric part of a random integer matrix survives & \texttt{GR\_bilin\allowbreak{}earOnlyA\allowbreak{}ntisymme\allowbreak{}tricPart\allowbreak{}Survives} (grassmann-demo) \\
Lemma 5.5 for each constant $C\gamma^a$ and a random antisymmetric matrix & \texttt{GR\_kinet\allowbreak{}icTotalD\allowbreak{}erivativ\allowbreak{}eReal} (grassmann-demo) \\
the symmetric-matrix contrast & \texttt{GR\_kinet\allowbreak{}icSymmet\allowbreak{}ricMatri\allowbreak{}xContras\allowbreak{}t} (grassmann-demo) \\
$\Psi^TC\Psi=0$, $\Psi^TC\gamma^0\Psi\ne0$, derivatives of $\tfrac\lambda2S^2$ & \texttt{ALG\_gras\allowbreak{}smannLem\allowbreak{}mas} (wolfram-geometry) \\
the split of $C\gamma^cS^{ab}$ of Section 5.12 & \texttt{ALG\_spin\allowbreak{}Transpos\allowbreak{}ePropert\allowbreak{}ies} (both algebra reports) \\
Theorem 5.6 at G1 p1 and in the symbolic G2 & \texttt{GR\_noteb\allowbreak{}ookLgELT\allowbreak{}rivial}, \texttt{GR\_noteb\allowbreak{}ookLgPur\allowbreak{}eDiverge\allowbreak{}nce} (grassmann-demo) \\
Theorem 5.6 and $X\Psi$ at all six G1 and G2 points & \texttt{LAG\_note\allowbreak{}bookLgGr\allowbreak{}assmannT\allowbreak{}rivial\_G\allowbreak{}1}, \texttt{\detokenize{_G2}} (wolfram-geometry) \\
the covariant constancy of the gammas & \texttt{GEO\_gamm\allowbreak{}aCovaria\allowbreak{}ntConsta\allowbreak{}ncy\_G1}, \texttt{\detokenize{_G2}} (wolfram-geometry, python-geometry) \\
the divergence identity & \texttt{GEO\_dive\allowbreak{}rgenceId\allowbreak{}entity\_G\allowbreak{}1}, \texttt{\detokenize{_G2}} (wolfram-geometry, python-geometry) \\
$\sum_\mu\{\gamma^\mu,\Omega_\mu\}=0$ in G2, for both contractions & \texttt{GEO\_anti\allowbreak{}commutat\allowbreak{}orGammaO\allowbreak{}megaVani\allowbreak{}shesDiag\allowbreak{}onal\_G2} (python-geometry) \\
\bottomrule
\end{longtable}
\endgroup

The Python demonstration works with 720 generators for the real field and 1440 for the complex one. With the canonical connection 0 of the 16 Euler--Lagrange components of $\mathrm{Lg}$ are nonzero, and $\mathrm{Lg}$ (544 monomials at G1 p1, 136 in G2) equals $\partial_\mu V^\mu$ exactly; with the notebook contraction all 16 components are nonzero, contain no derivative, and equal $X\Psi$ (measurements \texttt{GR\_noteb\allowbreak{}ookLgEL\_\allowbreak{}*} and \texttt{GR\_noteb\allowbreak{}ookLg\_la\allowbreak{}grangian\allowbreak{}Terms\_ca\allowbreak{}nonical\_\allowbreak{}G1\_p1}, \texttt{\detokenize{_G2}} in \texttt{grassman\allowbreak{}n-demo-r\allowbreak{}eport.js\allowbreak{}on}; the Wolfram report counts the same 544 and 136 monomials, \texttt{G1.p1.no\allowbreak{}tebookLg\allowbreak{}Monomial\allowbreak{}s} and \texttt{G2.p1.no\allowbreak{}tebookLg\allowbreak{}Monomial\allowbreak{}s}, and the same counts at the other four points). All 16 checks of \texttt{grassman\allowbreak{}n-demo-r\allowbreak{}eport.js\allowbreak{}on} are true. To rerun the Python demonstration from the repository root (it rewrites that report, so run it in a scratch clone if you want to keep the committed file untouched), in PowerShell

\begin{Verbatim}[fontsize=\small]
$env:PYTHONUTF8 = "1"
python scripts/demo_grassmann_lagrangians.py
\end{Verbatim}

and in Git Bash, macOS or Linux

\begin{Verbatim}[fontsize=\small]
export PYTHONUTF8=1
python scripts/demo_grassmann_lagrangians.py
\end{Verbatim}

The last two lines it writes to standard output are \texttt{check\_co\allowbreak{}unt=16} and \texttt{failed\_c\allowbreak{}heck\_cou\allowbreak{}nt=0}; a timing line ending in \texttt{\detokenize{done}} follows on standard error. The run takes under two minutes on a desktop computer and needs the Python packages numpy and sympy (Chapter 19). The whole Stage-1 gate, which runs this step and all others, is described in Chapter 19.

\subsection{What we proved and what we assumed}

\textbf{What we proved.} With complete arguments in this chapter: the Euler--Lagrange equations of a particle and of a field follow from the principle of stationary action (Sections 5.2 and 5.3); a Lagrangian that is a total divergence has identically vanishing Euler--Lagrange expressions (Proposition 5.1), for ordinary and for Grassmann fields; Noether's theorem (Theorem 5.2) and the canonical energy--momentum tensor; the metric energy--momentum tensor of a scalar field and its agreement with the Noether one; $\rho=\tfrac12\dot\phi^2+V$, $p=\tfrac12\dot\phi^2-V$ for a homogeneous scalar field; the classical energy of a scalar field in 4+4 dimensions is not bounded below, even for fields that vanish far away (Section 5.7); the rules of Grassmann algebra, of conjugation and of left and right derivatives (Proposition 5.3); the two bilinear lemmas (Lemmas 5.4 and 5.5); the matrix identities $C^T=C$, $C^2=1$, $(C\gamma^a)^T=-C\gamma^a$ and the symmetric/antisymmetric split of $C\gamma^cS^{ab}$ from the Clifford relations and the transposition pattern of the gammas; the divergence identity from the covariant constancy of the gammas; and Theorem 5.6: in every gravitational field, with the canonical spin connection, the notebook's $\mathrm{Lg}[\,]$ is the total divergence $\partial_\mu(\tfrac12\sqrt{|g|}\,\Psi^TC\gamma^\mu\Psi)$ for a real Grassmann field, with the field equations $0=0$. With the notebook's own contraction the field equations are the algebraic $X\Psi=0$. For commuting fields in a diagonal vielbein the spin connection drops out of $\mathrm{Lg}[\,]$ for either contraction (Section 5.14).

\textbf{What we assumed or took from elsewhere.} The Clifford relations, the symmetry pattern of the gammas and $C=\gamma^0\gamma^1\gamma^2\gamma^3$ are properties of the notebook's matrices, proved in Chapter 2 and verified exactly by the repository. The canonical spin connection and the covariant constancy $D_\mu\gamma^\nu=0$ come from Chapter 4 (they are equivalent to the vielbein postulate; the repository verifies them exactly at test points). The reading of cell 1064 as the formula of Section 5.12, and the meaning of the notebook's names in it, are taken from the Stage-1 document (§2.2 and §6.1). Two choices of the project are conventions, not results: the potential $U(S)=\tfrac\lambda2S^2$ named in Sections 5.9 and 5.14 (Stage-1 document, §1.2), and the overall sign of the current $J^\mu=-i\bar\Psi\gamma^\mu\Psi$, which is opposite to the sign that Theorem 5.2 gives (Section 5.7). That a fermion field must have anticommuting (Grassmann) components is the physical premise of the whole construction, the classical counterpart of the anticommutators of Chapter 8; no spin--statistics theorem is proved for signature (4,4) in this book. That the matrix $X$ of the notebook contraction has rank 16 is a computed fact at six test points, not a theorem for every field. Fields are assumed smooth and to vanish near the boundary of the region of integration whenever we integrate by parts.

\subsection{Exercises}

\textbf{Exercise 5.1.} Show that $L=\tfrac12\dot q^2-\tfrac12\omega^2q^2+\frac{d}{dt}\bigl(q^3\bigr)$ has the same Euler--Lagrange equation as the harmonic oscillator.

\textbf{Exercise 5.2.} For the scalar field of Section 5.5 with $m=2$, compute $\omega^2$ for the wave numbers $k_0=1$, $k_7=1$ (all others 0) and for $k_0=1$, $k_7=3$. Which one grows in time, and how fast?

\textbf{Exercise 5.3.} A homogeneous scalar field has $\dot\phi=1$ and $V=3/2$ at some moment. Compute $\rho$, $p$ and $w$. Repeat for $\dot\phi=2$, $V=0$.

\textbf{Exercise 5.4.} For the complex oscillator $L=\tfrac i2(\psi^\ast\dot\psi-\dot\psi^\ast\psi)-\omega\psi^\ast\psi$ (commuting $\psi$) compute the energy $H=\dot\psi\,\partial L/\partial\dot\psi+\dot\psi^\ast\,\partial L/\partial\dot\psi^\ast-L$.

\textbf{Exercise 5.5.} In the Grassmann algebra with generators $\theta_0,\theta_1,\theta_2$ compute $(\theta_0+\theta_1\theta_2)^2$ and $(\theta_0+\theta_1)(\theta_0-\theta_1)$.

\textbf{Exercise 5.6.} For $F=\theta_0\theta_1\theta_2+\theta_1\theta_2$ compute the left and the right derivatives with respect to $\theta_1$ and with respect to $\theta_2$. Which part of $F$ changes sign between left and right, and why?

\textbf{Exercise 5.7.} For a real Grassmann column $\Psi=(\Psi_0,\Psi_1,\Psi_2)^T$ and $M=\begin{pmatrix}0&1&2\\3&0&4\\5&6&0\end{pmatrix}$, compute $\Psi^TM\Psi$ and check Lemma 5.4.

\textbf{Exercise 5.8.} Show by an example that Lemma 5.5 is false for commuting fields: take $A=\begin{pmatrix}0&1\\-1&0\end{pmatrix}$ and $q(x)=(\cos x,\sin x)^T$.

\textbf{Exercise 5.9.} Using the tables of the Student Guide (row 4 of $\gamma^4$ reads ``$+9$'', row 13 of $\gamma^4$ reads ``$+0$'', row 0 of $C$ reads ``$-4$'', row 9 of $C$ reads ``$+13$''), find the entries $(0,9)$ and $(9,0)$ of $C\gamma^4$, confirm the antisymmetry of expression [1] there, and write the corresponding term of $\Psi^TC\gamma^4\Psi$.

\textbf{Exercise 5.10.} Let $\theta$ be one complex Grassmann variable. Show that $\theta^\ast\theta\ne0$ although $\theta\theta=0$, and that $(\theta^\ast\theta)^2=0$. For two complex variables and $S=\theta_0^\ast\theta_0-\theta_1^\ast\theta_1$, compute $S^2$ and $S^3$.

\textbf{Exercise 5.11.} For two real Grassmann fields and the antisymmetric $X=\begin{pmatrix}0&x\\-x&0\end{pmatrix}$ with an ordinary number $x\ne0$, compute $\tfrac12\Psi^TX\Psi$, its Euler--Lagrange expressions, and the solutions of the resulting equations.

\textbf{Exercise 5.12.} Check $\det(1+\epsilon A)=1+\epsilon\,\mathrm{tr}A+O(\epsilon^2)$ for $A=\begin{pmatrix}1&2\\3&4\end{pmatrix}$.

\subsection{Answers to the exercises}

\textbf{Answer 5.1.} The extra term is $F=3q^2\dot q$. Its contribution to the Euler--Lagrange expression is $\partial F/\partial q-\frac{d}{dt}\partial F/\partial\dot q=6q\dot q-\frac{d}{dt}(3q^2)=6q\dot q-6q\dot q=0$, so the equation is still $\ddot q=-\omega^2q$ (Proposition 5.1 in one dimension).

\textbf{Answer 5.2.} $\omega^2=m^2+k_0^2-k_7^2$. For $k_7=1$: $\omega^2=4+1-1=4$, $\omega=\pm2$, an oscillation. For $k_7=3$: $\omega^2=4+1-9=-4$; the time dependence is $e^{\pm2x_4}$, and the solution with $e^{+2x_4}$ grows by the factor $e^2\approx7.39$ per unit of $x_4$.

\textbf{Answer 5.3.} $\rho=\tfrac12\cdot1+\tfrac32=2$, $p=\tfrac12-\tfrac32=-1$, $w=-1/2$. For $\dot\phi=2$, $V=0$: $\rho=p=2$ and $w=1$.

\textbf{Answer 5.4.} $\partial L/\partial\dot\psi=\tfrac i2\psi^\ast$ and $\partial L/\partial\dot\psi^\ast=-\tfrac i2\psi$, so $H=\tfrac i2\psi^\ast\dot\psi-\tfrac i2\dot\psi^\ast\psi-\tfrac i2(\psi^\ast\dot\psi-\dot\psi^\ast\psi)+\omega\psi^\ast\psi=\omega\,\psi^\ast\psi$. The time derivatives drop out: for a first-order Lagrangian they carry no energy, a fact that returns for dirac16complex in Chapter 7.

\textbf{Answer 5.5.} $(\theta_0+\theta_1\theta_2)^2=\theta_0\theta_0+\theta_0\theta_1\theta_2+\theta_1\theta_2\theta_0+\theta_1\theta_2\theta_1\theta_2=0+\theta_0\theta_1\theta_2+\theta_0\theta_1\theta_2+0=2\theta_0\theta_1\theta_2$ (moving $\theta_0$ to the front of $\theta_1\theta_2\theta_0$ takes two exchanges). So an element with an odd and an even part need not square to zero. $(\theta_0+\theta_1)(\theta_0-\theta_1)=\theta_0\theta_0-\theta_0\theta_1+\theta_1\theta_0-\theta_1\theta_1=-2\theta_0\theta_1$.

\textbf{Answer 5.6.} With respect to $\theta_1$: left $-\theta_0\theta_2+\theta_2$, right $-\theta_0\theta_2-\theta_2$. With respect to $\theta_2$: left $\theta_0\theta_1-\theta_1$ (in $\theta_0\theta_1\theta_2$ the generator $\theta_2$ passes two others, in $\theta_1\theta_2$ one), right $\theta_0\theta_1+\theta_1$. The even part $\theta_1\theta_2$ changes sign between left and right, the odd part $\theta_0\theta_1\theta_2$ does not: by the proof of Proposition 5.3 the two signs differ by $(-1)^{n-1}$, which is $-1$ for even degree $n$ and $+1$ for odd degree.

\textbf{Answer 5.7.} Collect each pair $a<b$: the coefficient of $\Psi_a\Psi_b$ is $M_{ab}-M_{ba}$. So $\Psi^TM\Psi=(1-3)\Psi_0\Psi_1+(2-5)\Psi_0\Psi_2+(4-6)\Psi_1\Psi_2=-2\Psi_0\Psi_1-3\Psi_0\Psi_2-2\Psi_1\Psi_2$; the diagonal of $M$ is zero anyway. With $M_A=\tfrac12(M-M^T)$ the coefficients are $2(M_A)_{ab}=M_{ab}-M_{ba}$, the same numbers.

\textbf{Answer 5.8.} $q^TAq'=q_0q_1'-q_1q_0'=\cos x\cos x-\sin x\,(-\sin x)=1$, while $q^TAq=q_0q_1-q_1q_0=0$ for commuting numbers. Lemma 5.5 would claim $1=\tfrac12\cdot0'-0=0$. Its proof needs the exchange $\partial_\mu\Psi_a\,\Psi_b=-\Psi_b\,\partial_\mu\Psi_a$, which holds only for anticommuting fields.

\textbf{Answer 5.9.} $(C\gamma^4u)_0=-(\gamma^4u)_4=-u_9$, so the entry $(0,9)$ is $-1$. $(C\gamma^4u)_9=+(\gamma^4u)_{13}=+u_0$, so the entry $(9,0)$ is $+1$. The two entries are opposite, as $(C\gamma^4)^T=-C\gamma^4$ requires. In $\Psi^TC\gamma^4\Psi$ they give $-\Psi_0\Psi_9+\Psi_9\Psi_0=-2\Psi_0\Psi_9$.

\textbf{Answer 5.10.} $\theta$ and $\theta^\ast$ are different generators, so $\theta^\ast\theta$ is a nonzero monomial, while $\theta\theta=0$ because a generator squares to zero. $(\theta^\ast\theta)^2=\theta^\ast\theta\theta^\ast\theta$ contains $\theta$ twice and vanishes. With $P_0=\theta_0^\ast\theta_0$ and $P_1=-\theta_1^\ast\theta_1$ (even, commuting, $P_a^2=0$): $S^2=2P_0P_1=-2\,\theta_0^\ast\theta_0\theta_1^\ast\theta_1$ and $S^3=0$. With 16 components the same reasoning gives $S^{17}=0$ (Section 5.9).

\textbf{Answer 5.11.} $\tfrac12\Psi^TX\Psi=\tfrac12(x\Psi_0\Psi_1-x\Psi_1\Psi_0)=x\Psi_0\Psi_1$. There are no derivatives, so $E_0=\partial_L(x\Psi_0\Psi_1)/\partial\Psi_0=x\Psi_1=(X\Psi)_0$ and $E_1=\partial_L(-x\Psi_1\Psi_0)/\partial\Psi_1=-x\Psi_0=(X\Psi)_1$. The equations $x\Psi_1=0$ and $x\Psi_0=0$ with $x\ne0$ give $\Psi=0$: no field survives.

\textbf{Answer 5.12.} $\det\begin{pmatrix}1+\epsilon&2\epsilon\\3\epsilon&1+4\epsilon\end{pmatrix}=(1+\epsilon)(1+4\epsilon)-6\epsilon^2=1+5\epsilon-2\epsilon^2$, and $\mathrm{tr}A=5$.

\section{The two fields and their Lagrangians}

\subsection{What this chapter does}

Chapter 5 ended with a negative result: the Lagrangian $\mathrm{Lg}[\,]$ of the author's notebook describes nothing when its sixteen components are real anticommuting (Grassmann) numbers, as the components of a fermion field must be. This chapter writes down the Lagrangians that the project uses instead, and it does so for two fields that share the same sixteen components and the same formula but differ in one respect, the kind of numbers their components are.

\begin{itemize}
\item \textbf{dirac16complex} has complex Grassmann-odd components. It is the classical form of a field that is meant to be quantized with anticommutators, a "second-quantized" field (Section 6.2 explains the words); Chapter 8 carries out that quantization.
\item \textbf{dirac16complex00} has ordinary commuting complex components. It is a classical field, the sixteen-component analogue in 4+4 dimensions of the four-component wave function that Dirac wrote down in 1928.
\end{itemize}

For both fields the chapter derives, step by step: why the Lagrangian contains the matrix $C$ (Section 6.3); the Lagrangian itself (Section 6.4); why it is real (Section 6.5); the explicit mass term, which is linear in the mass and quadratic in the components, and why it is neither zero nor trivial (Section 6.6); the interaction $U(S)$ (Section 6.7); the coupling to gravity through the vielbein and the spin connection (Section 6.8); the continuous symmetries (Section 6.9); why the notebook's $\mathrm{Lg}[\,]$ is empty for real Grassmann fields but a genuine Lagrangian for real commuting ones, and how it sits inside the Lagrangian of dirac16complex00 (Section 6.10); and how the Lagrangians behave under the chirality map $\Psi\to\gamma^8\Psi$, under reflections and under complex conjugation (Sections 6.11 and 6.12). The field equations, the energy--momentum tensor and the equations of state that follow from these Lagrangians are the subject of Chapter 7.

\textbf{Sources.} The Stage-1 document \texttt{provenan\allowbreak{}ce/DIRAC\allowbreak{}16COMPLE\allowbreak{}X\_ARBITR\allowbreak{}ARY\_FIEL\allowbreak{}D.md} (its §4, §6 and §7) with its reports in \texttt{artifact\allowbreak{}s/dirac1\allowbreak{}6complex\allowbreak{}/arbitra\allowbreak{}ry-field\allowbreak{}/} (153 of 153 checks true, counted in \texttt{stage1-s\allowbreak{}ummary.j\allowbreak{}son}); the Stage-5 specification \texttt{handoff/\allowbreak{}specs/ST\allowbreak{}AGE5\_SPE\allowbreak{}C.md} (its §1 and §5); the exact Stage-5 reports on dirac16complex00 in \texttt{artifact\allowbreak{}s/dirac1\allowbreak{}6complex\allowbreak{}/pair-cr\allowbreak{}eation/}, namely \texttt{wolfram-\allowbreak{}dirac16c\allowbreak{}omplex00\allowbreak{}-report.\allowbreak{}json} (46 of 46 checks true; written by \texttt{scripts/\allowbreak{}verify\_d\allowbreak{}irac16co\allowbreak{}mplex00.\allowbreak{}wls} with the package \texttt{wolfram/\allowbreak{}Dirac16C\allowbreak{}omplex00\allowbreak{}.wl}) and \texttt{python-d\allowbreak{}irac16co\allowbreak{}mplex00-\allowbreak{}report.j\allowbreak{}son} (49 of 49 checks true; written by the independent sympy checker \texttt{scripts/\allowbreak{}check\_di\allowbreak{}rac16com\allowbreak{}plex00.p\allowbreak{}y}), whose exact formulas are collected in \texttt{dirac16c\allowbreak{}omplex00\allowbreak{}-theory.\allowbreak{}json} in the same folder; the Stage-5 pairing reports \texttt{wolfram-\allowbreak{}pairing-\allowbreak{}report.j\allowbreak{}son} (141 of 141) and \texttt{python-p\allowbreak{}airing-r\allowbreak{}eport.js\allowbreak{}on} (172 of 172) with \texttt{pairing-\allowbreak{}theory.j\allowbreak{}son} in the same folder; and the matter--antimatter document \texttt{provenan\allowbreak{}ce/DIRAC\allowbreak{}16COMPLE\allowbreak{}X\_MATTER\allowbreak{}\_ANTIMAT\allowbreak{}TER.md} (its §5) with its reports in \texttt{artifact\allowbreak{}s/dirac1\allowbreak{}6complex\allowbreak{}/matter-\allowbreak{}antimatt\allowbreak{}er/} (44 of 44 and 75 of 75 checks true). Check names that begin with \texttt{\detokenize{C00_}} belong to the Wolfram dirac16complex00 report, and names that begin with \texttt{\detokenize{PAIR_}} to the Wolfram pairing report \texttt{wolfram-\allowbreak{}pairing-\allowbreak{}report.j\allowbreak{}son}. The two independent Python checkers give all their checks the prefix \texttt{\detokenize{S5_}}: the Python twin of a \texttt{\detokenize{C00_}} check is the \texttt{\detokenize{S5_}} check with the same ending in \texttt{python-d\allowbreak{}irac16co\allowbreak{}mplex00-\allowbreak{}report.j\allowbreak{}son}, and the Python twin of a \texttt{\detokenize{PAIR_}} check is the \texttt{\detokenize{S5_}} check with the same ending in \texttt{python-p\allowbreak{}airing-r\allowbreak{}eport.js\allowbreak{}on}. An \texttt{\detokenize{S5_}} name that this chapter cites on its own is a check of \texttt{python-d\allowbreak{}irac16co\allowbreak{}mplex00-\allowbreak{}report.j\allowbreak{}son}. Names that begin with \texttt{\detokenize{MA_}} belong to the matter--antimatter reports; the other names are Stage-1 checks. The prose documents of Stage 5 (planned as \texttt{provenan\allowbreak{}ce/DIRAC\allowbreak{}16COMPLE\allowbreak{}X00\_FIEL\allowbreak{}D\_THEORY\allowbreak{}.md} and \texttt{provenan\allowbreak{}ce/DIRAC\allowbreak{}16COMPLE\allowbreak{}X\_PAIR\_C\allowbreak{}REATION.\allowbreak{}md}) had not been written when this chapter was written, so the chapter cites the committed reports directly. Every general statement of this chapter is derived here. Three kinds of statement are only quoted from the reports, and each is marked as such where it occurs: numbers measured at test points (numbers of terms and of nonzero entries); the curved-space form of the reflections of Section 6.12, which Chapter 15 proves; and the completeness of the classification of the discrete maps in Section 6.12, which Chapter 17 treats.

\textbf{Conventions} (Chapter 0 and Chapter 1): counting from 0; coordinates $x_0,\dots,x_7$ with $x_4$ the time; $\eta=\mathrm{diag}(+1,+1,+1,+1,-1,-1,-1,-1)$; Greek indices $\mu,\nu,\dots$ for coordinates, Latin $a,b,c,\dots$ for frame directions, both from 0 to 7; a repeated index, one up and one down, is summed. The gamma matrices $\gamma^0,\dots,\gamma^7$ are the notebook's (Section 2.8), $C=\gamma^0\gamma^1\gamma^2\gamma^3$, $\gamma^8=\gamma^0\gamma^1\cdots\gamma^7$ and $B=-iC\gamma^4$.

\subsection{Two fields with the same sixteen components}

\textbf{The spinor space.} Chapter 2 showed that sixteen complex numbers form the smallest space on which the eight gamma matrices of signature (4,4) can act, and that the group Pin(4,4) acts on it by $\Psi\to g\Psi$. Under Pin(4,4) the space $\mathbb C^{16}$ is irreducible; under its even part Spin(4,4) it splits into the two inequivalent halves $\Psi_0,\dots,\Psi_7$ (chirality $-1$) and $\Psi_8,\dots,\Psi_{15}$ (chirality $+1$) (Theorems 2.11 and 2.12). A \textbf{spinor field} attaches to every point $x$ a column $\Psi(x)=(\Psi_0(x),\dots,\Psi_{15}(x))^T$ whose components are measured in the local frame of the vielbein (Chapter 4); when the frame is turned by a local frame rotation, the column changes by the covering spin transformation, $\Psi\to R(x)\Psi$ (Section 4.12).

\textbf{Dirac's wave function of 1928.} Dirac's equation for the electron is an equation for a column of four complex functions $\psi(x)$ of the four coordinates of ordinary spacetime. Dirac read $\psi$ as a \textbf{wave function} in the sense of quantum mechanics: the number $\psi^\dagger\psi\ge0$ is the probability density of finding the electron at $x$, and it obeys a conservation law. The four components are ordinary complex numbers. This is called \textbf{first quantization}: a single particle is described by a wave function. Later physics found that this reading fails (the equation has solutions of negative energy, and particles can be created and destroyed), and the electron is described instead by a \textbf{quantum field}, whose components are operators that create and annihilate electrons and positrons and that anticommute. Turning the classical field into such operators is called \textbf{second quantization} (Chapter 8 does it for dirac16complex).

\textbf{dirac16complex.} Its components are complex Grassmann-odd numbers (Section 5.9):

\[
\Psi_a\Psi_b=-\Psi_b\Psi_a,\qquad \Psi_a\Psi_b^\ast=-\Psi_b^\ast\Psi_a,\qquad (\theta_1\theta_2)^\ast=\theta_2^\ast\theta_1^\ast .
\]

A classical Grassmann field is the classical counterpart of a field of anticommuting operators: its Lagrangian is the starting point of the quantization of Chapter 8, where the components become operators with the anticommutator $\{\Psi_a,\Psi_b^\dagger\}=B_{ab}\,\delta^7/\sqrt{|g|}$ at equal times (in the frames of Section 8.6, where $\gamma^{x_4}=\gamma^4$). This is what "second quantized" means for dirac16complex. It was introduced in Stage 1.

\textbf{dirac16complex00.} Its sixteen components are ordinary commuting complex numbers ("c-numbers") that transform together as one Pin(4,4) spinor: $\Psi\to g\Psi$ under Pin(4,4) and $\Psi\to R(x)\Psi$ under local frame rotations, exactly as for dirac16complex. It is the classical analogue, in 4+4 dimensions, of Dirac's four-component wave function. It was introduced in Stage 5 (\texttt{handoff/\allowbreak{}specs/ST\allowbreak{}AGE5\_SPE\allowbreak{}C.md}, §0, items 2 and 3). The two fields carry the same representation, because the representation is a property of the matrices, not of the numbers they act on: the matrices that commute with all eight gammas are the multiples of the identity (Theorem 2.11); those that commute with all 28 spin generators $S^{ab}$ form the two-dimensional space spanned by the projectors $P_-$ and $P_+$ (Theorem 2.12 (5)); and the invariant bilinear forms form the two-dimensional space spanned by $CP_-$ and $CP_+$ (checks \texttt{C00\_alge\allowbreak{}bra\_pinM\allowbreak{}oduleAnd\allowbreak{}SpinInva\allowbreak{}riance} and \texttt{S5\_algeb\allowbreak{}ra\_pinMo\allowbreak{}duleAndS\allowbreak{}pinInvar\allowbreak{}iance}).

The last statement has a short proof. A matrix $G$ defines an \textbf{invariant bilinear form} $\Psi^TG\Phi$ if the form does not change, to first order, under the infinitesimal spin rotations $\Psi\to(1+\epsilon S^{ab})\Psi$, $\Phi\to(1+\epsilon S^{ab})\Phi$, that is if $(S^{ab})^TG+GS^{ab}=0$ for all $a,b$ (the definition of Section 8.8 and of the checks). Property (S5) of Section 2.11 says $(CS^{ab})^T=-CS^{ab}$; since $C$ is symmetric this is $(S^{ab})^TC=-CS^{ab}$, and multiplying on the right by $C$ (with $C^2=1$) gives $(S^{ab})^T=-CS^{ab}C$. So the condition reads $-CS^{ab}CG+GS^{ab}=0$. Multiply it by $C$ from the left: $-S^{ab}(CG)+(CG)S^{ab}=0$. The matrix $Z:=CG$ therefore commutes with every $S^{ab}$, hence with every $\gamma^a\gamma^b=2S^{ab}$ ($a\ne b$), hence with every even monomial (a product of such pairs), hence with every combination of even monomials, which is the whole span of Spin(4,4) (Theorem 2.12 (2)). By Theorem 2.12 (5), $Z=r_-P_-+r_+P_+$ with two numbers $r_\mp$, that is $G=CZ=r_-CP_-+r_+CP_+$. Conversely every such $G$ satisfies the condition, because $P_\mp$ commute with every $S^{ab}$ (even elements commute with $\gamma^8$). The two matrices $CP_-$ and $CP_+$ are independent: if $x\,CP_-+y\,CP_+=0$, multiply by $C$ from the left and then by $P_-$ from the right ($P_-^2=P_-$, $P_+P_-=0$) to get $x\,P_-=0$, so $x=0$, and in the same way $y=0$. $\square$ The scalar density uses $G=C=CP_-+CP_+$.

\textbf{What differs, and what does not.} The Lagrangian, the field equations, the energy--momentum tensor and the current are given by the same formulas for both fields. The statistics of the components enters in five places, listed in \texttt{dirac16c\allowbreak{}omplex00\allowbreak{}-theory.\allowbreak{}json} (key \texttt{fields.d\allowbreak{}ifferenc\allowbreak{}e}) and treated in this book as follows:

\begin{enumerate}
\item which potentials $U(S)$ are possible (Section 6.7);
\item the real restriction and the notebook's $\mathrm{Lg}[\,]$ (Section 6.10);
\item the meaning of the energy--momentum tensor: for dirac16complex an operator of the quantized theory, \textbf{normal-ordered} (a prescription of Section 8.10 that removes the constant energy of the filled "Dirac sea" of negative-energy states), for dirac16complex00 an ordinary number field (Chapter 7);
\item positivity: the charge and the classical energy of dirac16complex00 have no fixed sign (Chapter 7), while for the quantized dirac16complex Chapter 8 derives a Hamiltonian that is not negative in the \textbf{good sector}, the fields that do not depend on the three extra times $x_5,x_6,x_7$ (Section 8.10);
\item the sign of the exchange term in a density-functional treatment (Chapter 14).
\end{enumerate}

\textbf{What dirac16complex00 is not.} It is not a quantum field. In ordinary 3+1 dimensional physics the spin--statistics theorem states that fields of half-integer spin must be quantized with anticommutators; a quantum theory of commuting spinors would contradict it. This book proves no spin--statistics theorem for signature (4,4) (Section 5.16), and the project does not quantize dirac16complex00: it treats it only as a classical field (\texttt{handoff/\allowbreak{}specs/ST\allowbreak{}AGE5\_SPE\allowbreak{}C.md}, §4; \texttt{pairing-\allowbreak{}theory.j\allowbreak{}son}, key \texttt{statisti\allowbreak{}cs}). Where Chapter 14 uses dirac16complex00 in a density-functional model, that model is labelled there as a prescription.

\textbf{A small example of the difference.} Take two components $\Psi_0,\Psi_1$. For commuting components $\Psi_0\Psi_1-\Psi_1\Psi_0=0$ and $\Psi_0\Psi_0=\Psi_0^2$ can be any complex number. For Grassmann components $\Psi_0\Psi_1-\Psi_1\Psi_0=2\Psi_0\Psi_1$ and $\Psi_0\Psi_0=0$. Every construction below keeps all conjugate factors $\Psi^\ast$ (or $\Psi^\dagger$) to the left of all factors $\Psi$, so that a formula written for one kind of component can be read, without reordering, for the other; where a reordering is unavoidable (Sections 6.10 and 6.12) the sign it produces is written out.

\subsection{The building blocks, and why the adjoint contains \texorpdfstring{$C$}{C}}

We collect the objects of Chapters 2 and 4 that the Lagrangian is built from.

\begin{itemize}
\item (B1) The gamma matrices satisfy $\gamma^a\gamma^b+\gamma^b\gamma^a=2\eta^{ab}$; $\gamma^a$ is symmetric for $a\le3$ and antisymmetric for $a\ge4$ (Section 2.8).
\item (B2) $C=\gamma^0\gamma^1\gamma^2\gamma^3$ is real and symmetric, $C^2=1$, and it has eight eigenvalues $+1$ and eight $-1$. Every $C\gamma^a$ is real and antisymmetric (expression [1] of Section 2.9), and so is every $CS^{ab}$, where $S^{ab}=\tfrac14[\gamma^a,\gamma^b]$ (Section 2.11). Equivalently $(\gamma^a)^TC=-C\gamma^a$ and $(S^{ab})^TC=-CS^{ab}$.
\item (B3) The \textbf{Dirac adjoint} of a column $\Psi$ is the row $\bar\Psi:=\Psi^\dagger C$, with the components $\bar\Psi_b=\sum_a\Psi_a^\ast C_{ab}$.
\item (B4) The vielbein $e_\mu{}^a(x)$, its inverse $e_a{}^\mu$, the metric $g_{\mu\nu}=e_\mu{}^a\eta_{ab}e_\nu{}^b$, the volume factor $\sqrt{|g|}=|\det e|$, and the curved gammas $\gamma^\mu=e_a{}^\mu\gamma^a$ with $\gamma^\mu\gamma^\nu+\gamma^\nu\gamma^\mu=2g^{\mu\nu}$ (Section 4.10). Since $e_a{}^\mu$ is real, every $\gamma^\mu$ is a real matrix and every $C\gamma^\mu$ is real and antisymmetric.
\item (B5) The canonical spin connection $\omega_{\mu ab}=-\omega_{\mu ba}$ of Section 4.11, the spinor connection $\Omega_\mu=\tfrac12\omega_{\mu ab}S^{ab}$ (a real matrix), and the covariant derivatives $D_\mu\Psi=\partial_\mu\Psi+\Omega_\mu\Psi$ and $D_\mu\bar\Psi=\partial_\mu\bar\Psi-\bar\Psi\,\Omega_\mu$, which satisfy $D_\mu\bar\Psi=(D_\mu\Psi)^\dagger C$ (Section 4.12).
\item (B6) The divergence identity $\partial_\mu\bigl(\sqrt{|g|}\,\gamma^\mu\bigr)=\sqrt{|g|}\,[\gamma^\mu,\Omega_\mu]$ (Sections 4.12 and 5.12).
\end{itemize}

\textbf{Why not $\Psi^\dagger\Psi$?} A Lagrangian must not depend on the choice of frame, so it must be built from combinations of the field that do not change under frame rotations. The simplest candidate, $\Psi^\dagger\Psi=\sum_a|\Psi_a|^2$, does not qualify. Under a spin transformation $R$ (a real matrix, Section 2.11) the column becomes $R\Psi$ and the row $\Psi^\dagger$ becomes $\Psi^\dagger R^T$, so $\Psi^\dagger\Psi\to\Psi^\dagger R^TR\,\Psi$. For a rotation $R^TR=1$, but not for a boost.

\emph{Worked example.} Take the boost $R=\cosh\tfrac\theta2+\sinh\tfrac\theta2\,J$ with $J=\gamma^0\gamma^4$ (Section 2.11). By (B1), $J^T=(\gamma^4)^T(\gamma^0)^T=-\gamma^4\gamma^0=\gamma^0\gamma^4=J$, so $R^T=R$, and $J^2=-(\gamma^0)^2(\gamma^4)^2=+1$. Hence $R^TR=R^2=\cosh^2\tfrac\theta2+\sinh^2\tfrac\theta2+2\cosh\tfrac\theta2\sinh\tfrac\theta2\,J=\cosh\theta+\sinh\theta\,J$. For $\Psi=e_0+e_4$ ($e_n$ is the column with 1 in place $n$) the tables of Section 2.8 give $\gamma^4e_0=e_{13}$ and $\gamma^0e_{13}=e_5$, so $Je_0=e_5$, and $\gamma^4e_4=-e_9$, $\gamma^0e_9=e_1$, so $Je_4=-e_1$. Then $\Psi^TJ\Psi=(e_0+e_4)^T(e_5-e_1)=0$ and $(R\Psi)^\dagger(R\Psi)=2\cosh\theta$, which is not the original value 2 unless $\theta=0$. The Dirac adjoint cures this.

\textbf{$\bar\Psi\Psi$ is invariant.} Section 2.11 proved $R^TCR=C$ for every exponential $R=\exp(\theta S^{ab})$ (and Section 2.14 for all products of them). Therefore

\[
\bar\Psi\Psi=\Psi^\dagger C\Psi\ \to\ \Psi^\dagger R^TCR\,\Psi=\Psi^\dagger C\Psi .
\]

In the example, $\Psi=e_0+e_4$ has $\bar\Psi\Psi=\Psi^TC\Psi=C_{04}+C_{40}=-2$ (row 0 of $C$ reads $-4$ and row 4 reads $-0$, Section 2.9), and so has $R\Psi$ for every $\theta$.

\textbf{$\bar\Psi\gamma^a\Phi$ is a vector.} From $R^TCR=C$ we get $R^TC=CR^{-1}$. If $R$ covers the frame rotation $\Lambda$, that is $R^{-1}\gamma^aR=\Lambda^a{}_b\gamma^b$ (Section 4.12), then

\[
(R\Psi)^\dagger C\gamma^a(R\Phi)=\Psi^\dagger R^TC\gamma^aR\,\Phi=\Psi^\dagger CR^{-1}\gamma^aR\,\Phi=\Lambda^a{}_b\,\bar\Psi\gamma^b\Phi ,
\]

so the eight numbers $\bar\Psi\gamma^a\Phi$ change like the frame components of a vector. Contracting the frame index with $e_a{}^\mu$ and then with a covariant derivative $D_\mu$ gives combinations such as $\bar\Psi\gamma^\mu D_\mu\Psi$ that do not change at all (Section 6.9 completes the argument). These are the building blocks of the Lagrangian.

\subsection{The Lagrangian of both fields}

\textbf{Definition.} For both fields, in an arbitrary gravitational field,

\[
\mathcal L=\sqrt{|g|}\,\Bigl[\tfrac12\bigl(\bar\Psi\gamma^\mu D_\mu\Psi-(D_\mu\bar\Psi)\gamma^\mu\Psi\bigr)-m\,\bar\Psi\Psi-U(\bar\Psi\Psi)\Bigr].
\]

This is the Lagrangian of Stage 1 for dirac16complex (Stage-1 document, §7.1) and equation (L1) of \texttt{handoff/\allowbreak{}specs/ST\allowbreak{}AGE5\_SPE\allowbreak{}C.md}, §1, for both fields. The action of a region $R$ of spacetime is $\int_R\mathcal L\,d^8x$. We use the names

\[
\begin{aligned}
&K:=\tfrac12\bigl(\bar\Psi\gamma^\mu D_\mu\Psi-(D_\mu\bar\Psi)\gamma^\mu\Psi\bigr)\quad\text{(kinetic term)},\qquad S:=\bar\Psi\Psi\quad\text{(scalar density)},\\
&\mathcal L_s:=\mathcal L/\sqrt{|g|}=K-mS-U(S),\qquad M_{\mathrm{eff}}:=m+U'(S)\quad\text{(effective mass)},
\end{aligned}
\]

and we write $\mathcal L_{m,U}$, or $\mathcal L_{m,\lambda}$ for the default potential $U(S)=\tfrac\lambda2S^2$, when the parameters matter.

\textbf{Reading the formula term by term.}

\begin{itemize}
\item $\sqrt{|g|}$ turns the coordinate volume $d^8x$ into proper volume (Section 4.4), so the action does not depend on the coordinates.
\item The \textbf{kinetic term} $K$ contains one derivative of the field, like the first-order Lagrangians of Sections 5.6 and 5.10. It is \textbf{symmetrized}: the derivative acts once on $\Psi$ and once on $\bar\Psi$, with opposite signs. Section 6.5 shows that this makes the Lagrangian real without any factor $i$, and Chapter 7 shows that its time part plays the role of the $\tfrac i2(\psi^\ast\dot\psi-\dot\psi^\ast\psi)$ of the oscillator of Section 5.6.
\item The \textbf{mass term} $-mS=-m\bar\Psi\Psi$ is linear in the mass $m$ and quadratic in the components (Section 6.6).
\item The \textbf{interaction} $-U(S)$ is a function of the scalar density (Section 6.7).
\item Gravity enters through $\sqrt{|g|}$, through the curved gammas $\gamma^\mu=e_a{}^\mu\gamma^a$ and through the spin connection in $D_\mu$ (Section 6.8).
\end{itemize}

\textbf{Relation to the notebook.} Stage 1 compared the formula with the notebook's $\mathrm{Lg}[\,]$ of Section 5.12 term by term (Stage-1 document, §7.3):

\begingroup
\small
\begin{longtable}{@{}>{\raggedright\arraybackslash}p{0.400\linewidth}>{\raggedright\arraybackslash}p{0.400\linewidth}@{}}
\toprule
\textbf{notebook $\mathrm{Lg}[\,]$} & \textbf{this Lagrangian} \\
\midrule
\endfirsthead
\toprule
\textbf{notebook $\mathrm{Lg}[\,]$} & \textbf{this Lagrangian} \\
\midrule
\endhead
Transpose[Ψ16], a real column & ConjugateTranspose: $\Psi^\dagger$ of a complex column \\
Transpose[Ψ16].σ16 & $\bar\Psi=\Psi^\dagger C$ with the same $C=\sigma_{16}$ \\
unsymmetrized $\Psi^TC\gamma^\mu D_\mu\Psi$ & symmetrized kinetic term $K$ \\
mixed connection $\omega_\mu{}^a{}_b$ & lowered connection $\omega_{\mu ab}=\eta_{ac}\omega_\mu{}^c{}_b$ \\
Sqrt[detgg] & $\sqrt{\lvert g\rvert}$ \\
$+(HM)\,\Psi^TC\Psi$ & $-m\bar\Psi\Psi$ with $m=-HM$ (same sign) \\
no self-interaction & $-U(\bar\Psi\Psi)$ \\
\bottomrule
\end{longtable}
\endgroup

\textbf{Worked example.} In flat space ($e_\mu{}^a=\delta_\mu{}^a$, $\Omega_\mu=0$) take the constant commuting field $\Psi=e_0+e_4$. All derivatives vanish, so $K=0$; $S=-2$ (Section 6.3); and for $U=\tfrac\lambda2S^2$ the Lagrangian is $\mathcal L=-m(-2)-\tfrac\lambda2\cdot4=2m-2\lambda$.

\subsection{Reality: the Lagrangian is real for both kinds of components}

A Lagrangian must be real, for commuting components an ordinary real number, and for Grassmann components an element that is unchanged by the conjugation of the Grassmann algebra (the book calls such an element Hermitian, Section 5.9). Otherwise the action is a complex number, the equations obtained by varying $\Psi$ and by varying $\Psi^\ast$ are in general not the conjugates of each other, and there are more independent equations than field components.

\textbf{Lemma 6.1 (conjugating a bilinear).} For two columns $\Psi$ and $\Phi$ (for example $\Phi=D_\mu\Psi$), both commuting or both Grassmann, and a matrix $M$ of ordinary numbers,

\[
\bigl(\Psi^\dagger M\Phi\bigr)^\ast=\Phi^\dagger M^\dagger\Psi .
\]

\emph{Proof.} Commuting components: $\bigl(\sum_{a,b}\Psi_a^\ast M_{ab}\Phi_b\bigr)^\ast=\sum_{a,b}\Psi_aM_{ab}^\ast\Phi_b^\ast=\sum_{a,b}\Phi_b^\ast(M^\dagger)_{ba}\Psi_a$, because ordinary numbers may be written in any order. Grassmann components: the conjugation reverses the order of a product, $(\Psi_a^\ast\Phi_b)^\ast=\Phi_b^\ast\Psi_a$ (Section 5.9), so $\bigl(\Psi_a^\ast M_{ab}\Phi_b\bigr)^\ast=M_{ab}^\ast\,\Phi_b^\ast\Psi_a$, and the sum is the same expression. $\square$

The two computations end in the same formula for opposite reasons: for commuting numbers nothing is reordered, and for Grassmann numbers the reversal of the order is built into the conjugation. This is why the reality proof below holds for both fields at once (\texttt{dirac16c\allowbreak{}omplex00\allowbreak{}-theory.\allowbreak{}json}, key \texttt{lagrangi\allowbreak{}an.reali\allowbreak{}ty}).

\textbf{Theorem 6.2 (reality).} For both kinds of components the Lagrangian $\mathcal L$ is real (Hermitian): $\mathcal L^\ast=\mathcal L$.

\emph{Proof.} Put $X:=\bar\Psi\gamma^\mu D_\mu\Psi=\Psi^\dagger\,(C\gamma^\mu)\,(D_\mu\Psi)$, summed over $\mu$. By (B4), $C\gamma^\mu$ is real and antisymmetric, so $(C\gamma^\mu)^\dagger=(C\gamma^\mu)^T=-C\gamma^\mu$. Lemma 6.1 and (B5) give

\[
X^\ast=(D_\mu\Psi)^\dagger(C\gamma^\mu)^\dagger\Psi=-(D_\mu\Psi)^\dagger C\gamma^\mu\Psi=-(D_\mu\bar\Psi)\gamma^\mu\Psi .
\]

Hence $K=\tfrac12\bigl(X-(D_\mu\bar\Psi)\gamma^\mu\Psi\bigr)=\tfrac12(X+X^\ast)$ is real. Next, $S^\ast=(\Psi^\dagger C\Psi)^\ast=\Psi^\dagger C^\dagger\Psi=S$ because $C$ is real and symmetric. For commuting components $U(S)$ is a real function of the real number $S$. For Grassmann components $U$ is a polynomial with real coefficients (Section 6.7); since $S$ is even, $(S^k)^\ast=(S^\ast)^k=S^k$, so $U(S)^\ast=U(S)$. Finally $m$ and $\sqrt{|g|}$ are real. $\square$

\textbf{Why the kinetic term is symmetrized.} The unsymmetrized $X$ alone is not real. Its imaginary part is a total derivative: Section 7.5 proves the identity $\sqrt{|g|}\,\bigl(X+(D_\mu\bar\Psi)\gamma^\mu\Psi\bigr)=\partial_\mu\bigl(\sqrt{|g|}\,\bar\Psi\gamma^\mu\Psi\bigr)$, which says $\sqrt{|g|}(X-X^\ast)=\partial_\mu(\sqrt{|g|}\,\bar\Psi\gamma^\mu\Psi)$, so

\[
\sqrt{|g|}\,X=\sqrt{|g|}\,K+\tfrac12\,\partial_\mu\bigl(\sqrt{|g|}\,\bar\Psi\gamma^\mu\Psi\bigr).
\]

By Section 5.4 the two versions give the same field equations, but only the symmetrized one gives a real action. The repository checks both statements: the symmetrized Lagrangian is real and the unsymmetrized kinetic term is not (checks \texttt{LAG\_herm\allowbreak{}iticity\_\allowbreak{}G1}, \texttt{LAG\_herm\allowbreak{}iticity\_\allowbreak{}G2}, \texttt{GR\_lagra\allowbreak{}ngianHer\allowbreak{}mitian}, \texttt{GR\_unsym\allowbreak{}metrized\allowbreak{}KineticN\allowbreak{}otHermit\allowbreak{}ian} of Stage 1; \texttt{C00\_lagr\allowbreak{}angian\_r\allowbreak{}ealCommu\allowbreak{}ting}, which writes $\Psi=X_r+iY_r$ with real commuting symbols and finds $\mathrm{Im}\,\mathcal L=0$ exactly at a point of the general test geometry G1 and at three points of the primordial field; \texttt{S5\_lagra\allowbreak{}ngian\_re\allowbreak{}alCommut\allowbreak{}ing}, which also records \texttt{unsymmet\allowbreak{}rizedKin\allowbreak{}eticNotR\allowbreak{}eal}).

\textbf{The same statement for the coefficients.} Collect the terms of $\mathcal L$ by the number of derivatives: $\mathcal L=\Psi^\dagger K_0\Psi+\Psi^\dagger K^\mu\partial_\mu\Psi+\partial_\mu\Psi^\dagger L^\mu\Psi-\sqrt{|g|}\,(mS+U)$ with

\[
K^\mu=\tfrac12\sqrt{|g|}\,C\gamma^\mu,\qquad L^\mu=-\tfrac12\sqrt{|g|}\,C\gamma^\mu=(K^\mu)^\dagger,\qquad K_0=\tfrac12\sqrt{|g|}\,C\{\gamma^\mu,\Omega_\mu\} .
\]

($K_0$ collects $\tfrac12\bar\Psi\gamma^\mu\Omega_\mu\Psi$ from the first half of $K$ and $+\tfrac12\bar\Psi\Omega_\mu\gamma^\mu\Psi$ from the second.) $K_0$ is real and symmetric because $C\{\gamma^c,S^{ab}\}$ is symmetric (Section 5.12). The conditions $K_0^\dagger=K_0$ and $L^\mu=(K^\mu)^\dagger$ are exactly what makes $\mathcal L$ real (Stage-1 document, Result 7.1; checks \texttt{C00\_lagr\allowbreak{}angian\_c\allowbreak{}oefficie\allowbreak{}ntHermit\allowbreak{}icity} and \texttt{S5\_lagra\allowbreak{}ngian\_co\allowbreak{}efficien\allowbreak{}tHermiti\allowbreak{}city}, which read the coefficients off the symbolic Lagrangian).

\textbf{Worked example.} Flat space, commuting field depending on $x_1$ only: $\Psi=f(x_1)\,e_0+g(x_1)\,e_{11}$. Only $\gamma^1\partial_1$ contributes to $X$, and we need two entries of $C\gamma^1$. From the tables (Sections 2.8 and 2.9): $(C\gamma^1u)_0=-(\gamma^1u)_4=-(-u_{11})=u_{11}$ and $(C\gamma^1u)_{11}=+(\gamma^1u)_{15}=-u_0$. So $(C\gamma^1)_{0,11}=+1$ and $(C\gamma^1)_{11,0}=-1$, antisymmetric as it must be. Then $X=f^\ast g'-g^\ast f'$ (a prime is $d/dx_1$). Take $f=1$ and $g=e^{ikx_1}$ with a real $k$. Then

\[
X=ik\,e^{ikx_1},\qquad K=\tfrac12(X+X^\ast)=\tfrac12\bigl(ik\,e^{ikx_1}-ik\,e^{-ikx_1}\bigr)=-k\sin(kx_1),
\]

a real number, while $X$ itself is complex. And $X-X^\ast=2ik\cos(kx_1)$, which is the derivative of $\bar\Psi\gamma^1\Psi=f^\ast g-g^\ast f=2i\sin(kx_1)$, as the total-derivative identity says.

\subsection{The mass term}

\textbf{The term.} The mass term of both Lagrangians is

\[
\mathcal L_m=-m\sqrt{|g|}\;\bar\Psi\Psi=-m\sqrt{|g|}\sum_{a,b=0}^{15}\Psi_a^\ast\,C_{ab}\,\Psi_b .
\]

It is \textbf{linear in the mass} $m$ and \textbf{quadratic in the components}: every term contains exactly one factor $\Psi_a^\ast$ and one factor $\Psi_b$ (a bilinear). The Stage-5 specification requires such a term for both fields and requires it to be non-zero and non-trivial (\texttt{handoff/\allowbreak{}specs/ST\allowbreak{}AGE5\_SPE\allowbreak{}C.md}, §1). Both properties are proved now.

\textbf{Non-zero for commuting components.} $C$ is a signed permutation matrix (Section 2.9): $(Cu)_i=-u_{i+4}$ for $i=0,\dots,3$, $(Cu)_i=-u_{i-4}$ for $i=4,\dots,7$, $(Cu)_i=+u_{i+4}$ for $i=8,\dots,11$ and $(Cu)_i=+u_{i-4}$ for $i=12,\dots,15$. Hence

\[
S=\bar\Psi\Psi=-\sum_{i=0}^{3}\bigl(\Psi_i^\ast\Psi_{i+4}+\Psi_{i+4}^\ast\Psi_i\bigr)+\sum_{i=8}^{11}\bigl(\Psi_i^\ast\Psi_{i+4}+\Psi_{i+4}^\ast\Psi_i\bigr).
\]

Three exact values: $S=-2$ for $\Psi=e_0+e_4$, $S=+2$ for $\Psi=e_8+e_{12}$, and $S=0$ for $\Psi=e_0+i\,e_4$ (there $\Psi_0^\ast\Psi_4+\Psi_4^\ast\Psi_0=i-i=0$). The scalar density takes both signs because $C$ has eight eigenvalues $+1$ and eight $-1$ (checks \texttt{C00\_mass\allowbreak{}Term\_com\allowbreak{}mutingEx\allowbreak{}plicitSp\allowbreak{}inors} and \texttt{S5\_massT\allowbreak{}erm\_comm\allowbreak{}utingExp\allowbreak{}licitSpi\allowbreak{}nors}, which record the three values and the characteristic polynomial $(x-1)^8(x+1)^8$ of $C$).

\textbf{Non-zero for Grassmann components.} For Grassmann components the formula above is an element of the Grassmann algebra generated by the 32 odd elements $\Psi_a,\Psi_a^\ast$ at a point. Its 16 terms $\Psi_a^\ast C_{ab}\Psi_b$ are 16 different monomials, one for each nonzero entry of $C$, each with coefficient $\pm1$, and none cancels, because $\Psi_a^\ast$ and $\Psi_b$ are different generators. So $S\ne0$. (Contrast the real Grassmann field of Section 5.13, for which $\Psi^TC\Psi=0$.) $S^2$ has 120 monomials and $S^3$ has 560 (checks \texttt{C00\_mass\allowbreak{}Term\_gra\allowbreak{}ssmannNo\allowbreak{}nzero} and \texttt{S5\_massT\allowbreak{}erm\_gras\allowbreak{}smannNon\allowbreak{}zero}, which also records all the counts of Section 5.9 up to $S^{17}=0$).

\textbf{Non-trivial: the mass changes the field equation.} A term of the Lagrangian is trivial if it does not change the field equations, for example a total derivative (Section 5.4). The mass term is not trivial, because the mass appears in the relation between frequency and wave number. We show it in flat space with $U=0$, where Chapter 7 derives the field equation $\gamma^a\partial_a\Psi=m\Psi$. Try a plane wave

\[
\Psi(x)=u\,\exp\Bigl(i\sum_{a=0}^7k_ax_a\Bigr)
\]

with a constant column $u$ and real wave numbers $k_a$. Each $\partial_a$ produces $ik_a$, so the equation becomes $\bigl(i\gamma(k)-m\bigr)u=0$ with $\gamma(k):=\sum_ak_a\gamma^a$, the notation of Section 2.14. By the Clifford relation $\gamma(k)^2=\sum_a\eta^{aa}k_a^2$ (Section 2.14; the cross terms cancel in pairs), so

\[
\bigl(i\gamma(k)-m\bigr)\bigl(i\gamma(k)+m\bigr)=-\gamma(k)^2-m^2=\Bigl(k_4^2+k_5^2+k_6^2+k_7^2-k_0^2-k_1^2-k_2^2-k_3^2-m^2\Bigr)\cdot1 .
\]

If the bracket is not zero, $i\gamma(k)-m$ has an inverse and the only solution is $u=0$. If the bracket is zero, nonzero solutions exist. Writing $\omega:=|k_4|$ for the frequency in the time $x_4$, the condition is

\[
\omega^2=m^2+k_0^2+k_1^2+k_2^2+k_3^2-k_5^2-k_6^2-k_7^2 ,
\]

the same relation as for the scalar field of Section 5.5: the mass enters as $m^2$. On the mass shell, for $m\ne0$, the solutions form a space of dimension exactly 8. \emph{Proof.} Write $N_\pm=i\gamma(k)\pm m$ and $\ker N$ for the set of columns $u$ with $Nu=0$ (a subspace). Both $N_\pm$ are built from the single matrix $\gamma(k)$, so they commute, and on shell the product computed above is zero: $N_-N_+=N_+N_-=0$. Also $N_+-N_-=2m\cdot1$. Step 1: every column $u$ is a sum of a column of $\ker N_-$ and a column of $\ker N_+$. Indeed $u=\tfrac1{2m}(N_+-N_-)u=p+q$ with $p=\tfrac1{2m}N_+u$ and $q=-\tfrac1{2m}N_-u$, and $N_-p=\tfrac1{2m}N_-N_+u=0$, $N_+q=-\tfrac1{2m}N_+N_-u=0$. Step 2: the two subspaces meet only in 0: if $N_-u=N_+u=0$ then $2mu=(N_+-N_-)u=0$, so $u=0$. Step 3: take a basis $p_1,\dots,p_r$ of $\ker N_-$ and a basis $q_1,\dots,q_s$ of $\ker N_+$. By Step 1 the $r+s$ columns together span all of $\mathbb C^{16}$. They are also independent: if $\sum a_ip_i+\sum b_jq_j=0$, the column $\sum a_ip_i=-\sum b_jq_j$ lies in both subspaces, so it is 0 by Step 2, and then all $a_i$ and all $b_j$ vanish because each set is a basis. So the $r+s$ columns form a basis of $\mathbb C^{16}$, and since every basis of $\mathbb C^{16}$ has 16 elements (all bases of a space have the same number of elements, Section 2.2), $r+s=16$. Step 4: $\gamma^8N_-\gamma^8=-i\gamma(k)-m=-N_+$ (because $\gamma^8$ anticommutes with every $\gamma^a$ and $(\gamma^8)^2=1$). So if $N_-u=0$, then $N_+(\gamma^8u)=-\gamma^8N_-\gamma^8\gamma^8u=-\gamma^8N_-u=0$: the matrix $\gamma^8$ maps $\ker N_-$ into $\ker N_+$, and in the same way $\ker N_+$ into $\ker N_-$. The columns $\gamma^8p_1,\dots,\gamma^8p_r$ are independent (if $\sum c_i\gamma^8p_i=0$, multiplying by $\gamma^8$ gives $\sum c_ip_i=0$, so all $c_i=0$), and they lie in the $s$-dimensional $\ker N_+$, so $r\le s$ by fact (i) of Section 2.2; in the same way $s\le r$. Hence $r=s=8$. $\square$ The repository checks the product identity and the two cases $m=1$, $k=(k_0,\dots,k_7)=(1,1,2,3,4,0,0,0)$, on shell because $16-(1+1+4+9)=1=m^2$, with an 8-dimensional solution space, and $k_4=3$ instead of 4, off shell, with no solution (checks \texttt{C00\_mass\allowbreak{}Term\_dis\allowbreak{}persionF\allowbreak{}lat} and \texttt{S5\_massT\allowbreak{}erm\_disp\allowbreak{}ersionFl\allowbreak{}at}). The plane-wave analysis concerns the linear field equation, which is the same for both fields (Chapter 7). In a curved field the mass enters the second-order form of the field equation in the same way, as $m^2$ (Section 7.4).

\subsection{The interaction \texorpdfstring{$U(S)$}{U(S)}}

\textbf{Grassmann components: only polynomials.} At one point, $S=\bar\Psi\Psi$ is an even element of the Grassmann algebra of the 32 odd generators $\Psi_a,\Psi_a^\ast$. Every monomial of $S^k$ contains $2k$ distinct generators, so $S^{17}=0$ (Section 5.9: $S^k$ has $\binom{16}{k}$ monomials for $k\le16$). A function of $S$ can therefore only mean its Taylor polynomial, which stops at degree 16:

\[
U(S)=u_2S^2+u_3S^3+\dots+u_{16}S^{16}\qquad(\text{real }u_k).
\]

The project requires $U(0)=U'(0)=0$ (Stage-1 document, §7.6): a constant term would be a cosmological constant, and a linear term would be a second mass term. The default is the four-fermion contact interaction

\[
U(S)=\tfrac\lambda2S^2,\qquad U'(S)=\lambda S,\qquad M_{\mathrm{eff}}=m+\lambda S .
\]

Chapter 11 calls $\lambda<0$ attractive and $\lambda>0$ repulsive.

\textbf{Commuting components: any smooth function.} For dirac16complex00, $S$ is an ordinary real number at each point, so $U$ may be any smooth real function of one variable; power series such as $e^{S}-1-S$ are allowed. The default is the same $U=\tfrac\lambda2S^2$, so that the two theories differ \textbf{only} in the statistics (\texttt{handoff/\allowbreak{}specs/ST\allowbreak{}AGE5\_SPE\allowbreak{}C.md}, §1). The Stage-5 verifiers check the field equations and the energy--momentum tensor with an abstract smooth $U$, whose values $U(S_0)$, $U'(S_0)$, $U''(S_0)$ at the point are independent symbols (checks \texttt{C00\_EL\_c\allowbreak{}ommuting\allowbreak{}GeneralS\allowbreak{}moothU}, \texttt{C00\_EMT\_\allowbreak{}generalS\allowbreak{}moothU} and their \texttt{\detokenize{S5_}} twins).

\textbf{Worked example.} Two complex Grassmann components with $S=\theta_0^\ast\theta_0-\theta_1^\ast\theta_1$ (Exercise 5.10). The two pieces $P_0=\theta_0^\ast\theta_0$ and $P_1=-\theta_1^\ast\theta_1$ are even, commute, and square to zero. So $S^2=2P_0P_1=-2\,\theta_0^\ast\theta_0\theta_1^\ast\theta_1$ and $S^3=0$. The "exponential" $e^S=1+S+\tfrac12S^2$ is an exact finite sum, and $U=\tfrac\lambda2S^2=-\lambda\,\theta_0^\ast\theta_0\theta_1^\ast\theta_1$. For commuting numbers the same symbols would give an infinite series for $e^S$.

\subsection{The coupling to gravity: vielbein and spin connection}

Gravity enters the Lagrangian in three places: the volume factor $\sqrt{|g|}$, the curved gammas $\gamma^\mu=e_a{}^\mu\gamma^a$, and the spinor connection $\Omega_\mu$ inside $D_\mu$. The first two carry the vielbein itself, the third its derivatives. This section shows precisely how much of the spin connection the Lagrangian actually sees, and why the coupling is nevertheless not trivial in a curved field.

\textbf{Lemma 6.3.} For $a\ne b$ and any $c$,

\[
\{\gamma^c,S^{ab}\}=\begin{cases}\gamma^c\gamma^a\gamma^b&\text{if }c\ne a\text{ and }c\ne b,\\ 0&\text{if }c=a\text{ or }c=b.\end{cases}
\]

\emph{Proof.} For $a\ne b$, $S^{ab}=\tfrac12\gamma^a\gamma^b$. If $c$ differs from $a$ and $b$, moving $\gamma^c$ past $\gamma^a\gamma^b$ costs two signs, so $\gamma^cS^{ab}=S^{ab}\gamma^c=\tfrac12\gamma^c\gamma^a\gamma^b$ and the anticommutator is $\gamma^c\gamma^a\gamma^b$. If $c=a$: $\gamma^aS^{ab}=\tfrac12(\gamma^a)^2\gamma^b=\tfrac12\eta^{aa}\gamma^b$ and $S^{ab}\gamma^a=\tfrac12\gamma^a\gamma^b\gamma^a=-\tfrac12(\gamma^a)^2\gamma^b=-\tfrac12\eta^{aa}\gamma^b$; they cancel. The case $c=b$ is the same with the roles exchanged. $\square$

\textbf{Theorem 6.4 (what the Lagrangian sees of the connection).} Write $\omega_{cab}:=e_c{}^\mu\,\omega_{\mu ab}$ for the frame components of the spin connection. The part of the kinetic term that contains $\Omega_\mu$ is

\[
\tfrac12\bigl(\bar\Psi\gamma^\mu\Omega_\mu\Psi+\bar\Psi\Omega_\mu\gamma^\mu\Psi\bigr)=\tfrac12\,\bar\Psi\{\gamma^\mu,\Omega_\mu\}\Psi=\tfrac14\sum_{c,a,b\ \text{distinct}}\omega_{cab}\,\bar\Psi\gamma^c\gamma^a\gamma^b\Psi ,
\]

and only the \textbf{totally antisymmetric part} of $\omega_{cab}$ contributes. In a field with a diagonal vielbein this part is zero, so there the Lagrangian contains no spin connection at all.

\emph{Proof.} The first equality: $K=\tfrac12(\bar\Psi\gamma^\mu(\partial_\mu+\Omega_\mu)\Psi-(\partial_\mu\bar\Psi-\bar\Psi\Omega_\mu)\gamma^\mu\Psi)$, and the terms with $\Omega_\mu$ are $\tfrac12\bar\Psi\gamma^\mu\Omega_\mu\Psi+\tfrac12\bar\Psi\Omega_\mu\gamma^\mu\Psi$. The second: $\gamma^\mu=e_c{}^\mu\gamma^c$ and $\Omega_\mu=\tfrac12\omega_{\mu ab}S^{ab}$ give $\{\gamma^\mu,\Omega_\mu\}=\tfrac12\omega_{cab}\{\gamma^c,S^{ab}\}$, and Lemma 6.3 keeps only the terms with $c,a,b$ distinct. For three distinct indices $\gamma^c\gamma^a\gamma^b$ changes sign when any two of them are exchanged (the gammas anticommute), so in the sum $\omega_{cab}$ may be replaced by its totally antisymmetric part $\tfrac16(\omega_{cab}+\omega_{abc}+\omega_{bca}-\omega_{acb}-\omega_{bac}-\omega_{cba})$. For a diagonal vielbein $e_\mu{}^a=h_\mu\delta_\mu{}^a$, Section 4.11 showed that $\omega_{\mu ab}\ne0$ only when $\mu=a$ or $\mu=b$. The inverse vielbein is diagonal too, $e_c{}^\mu=\delta_c{}^\mu/h_c$, so $\omega_{cab}$ equals $1/h_c$ times the component $\omega_{\mu ab}$ with the coordinate index $\mu$ equal to the number $c$, and that component vanishes unless $c=a$ or $c=b$. So $\omega_{cab}=0$ whenever $c,a,b$ are distinct. $\square$

(Checks: \texttt{C00\_alge\allowbreak{}bra\_anti\allowbreak{}commutat\allowbreak{}orTotall\allowbreak{}yAntisym\allowbreak{}metric} and \texttt{S5\_algeb\allowbreak{}ra\_antic\allowbreak{}ommutato\allowbreak{}rTotally\allowbreak{}Antisymm\allowbreak{}etric} for the lemma; in \texttt{C00\_conn\allowbreak{}ection\_n\allowbreak{}onTrivia\allowbreak{}lCouplin\allowbreak{}g} the measurement \texttt{onlyTota\allowbreak{}llyAntis\allowbreak{}ymmetric\allowbreak{}Omega} is true at every test point, \texttt{anticomm\allowbreak{}utatorTe\allowbreak{}rmNonzer\allowbreak{}o} is true at the three points of the non-diagonal geometry G1 and false at the three points of the diagonal primordial field G2. Stage 1 recorded the diagonal statement as \texttt{GEO\_anti\allowbreak{}commutat\allowbreak{}orGammaO\allowbreak{}megaVani\allowbreak{}shesDiag\allowbreak{}onal\_G2}.)

\textbf{The connection enters the field equation.} Theorem 6.4 does not mean that gravity decouples in a diagonal field. When the field equation is derived (Section 7.2), the derivative of $\sqrt{|g|}\,\gamma^\mu$ appears through an integration by parts, and the divergence identity (B6) turns it into the commutator $[\gamma^\mu,\Omega_\mu]$. The field equation is $\gamma^\mu D_\mu\Psi=M_{\mathrm{eff}}\Psi$, which contains the matrix

\[
\gamma^\mu\Omega_\mu .
\]

For a diagonal vielbein $\{\gamma^\mu,\Omega_\mu\}=0$ (Section 4.12), so $[\gamma^\mu,\Omega_\mu]=2\gamma^\mu\Omega_\mu$, and (B6) gives

\[
\gamma^\mu\Omega_\mu=\frac{1}{2\sqrt{|g|}}\,\partial_\mu\bigl(\sqrt{|g|}\,\gamma^\mu\bigr).
\]

This shows how the coupling through $\sqrt{|g|}$ and $\gamma^\mu$ becomes a connection term. In the notebook's primordial field this matrix is $3H\gamma^0$ (Section 4.12), and at the points of the general test geometry G1 it has 128 nonzero entries. The spin-connection term of the field equation is nonzero in all 16 components at every G1 and G2 point (checks \texttt{C00\_conn\allowbreak{}ection\_n\allowbreak{}onTrivia\allowbreak{}lCouplin\allowbreak{}g}, \texttt{C00\_conn\allowbreak{}ection\_O\allowbreak{}megaTerm\allowbreak{}InFieldE\allowbreak{}quation}, \texttt{S5\_conne\allowbreak{}ction\_no\allowbreak{}nTrivial\allowbreak{}Coupling}, \texttt{S5\_conne\allowbreak{}ction\_Om\allowbreak{}egaTermI\allowbreak{}nFieldEq\allowbreak{}uation}). It cannot be removed by any choice of frame unless the Riemann tensor vanishes (Section 4.13), and the curvature appears explicitly in the second-order form of the field equation through the term $-\tfrac14R\Psi$ of the Lichnerowicz formula (Section 4.13), with the constant $-\tfrac14$ measured again for the commuting field at seven test points (G1 p1 to p3, G2 p1 to p3 and a point of a further generic non-diagonal integer frame, G4; measurement \texttt{lichnero\allowbreak{}wiczC} in \texttt{C00\_conn\allowbreak{}ection\_n\allowbreak{}onTrivia\allowbreak{}lCouplin\allowbreak{}g}). The operator $\gamma^\mu D_\mu$ does not know the statistics of the components, so all of this holds for both fields.

\subsection{Continuous symmetries of the Lagrangian}

\textbf{Local frame rotations.} Let the frame be turned, $e'_\mu{}^a=\Lambda^a{}_b(x)e_\mu{}^b$, and let the field change by a covering spin transformation $R(x)$ of spinor norm $+1$, in particular any element of $\mathrm{Spin}_0(4,4)$ (Sections 2.14 and 4.12). Then $\mathcal L$ does not change. \emph{Proof.} Chapter 4 showed $\gamma'^\mu=R\gamma^\mu R^{-1}$ and $D'_\mu(R\Psi)=R\,D_\mu\Psi$ (Section 4.12). From $R^TCR=C$ we get $(R\Psi)^\dagger C=\Psi^\dagger R^TC=\bar\Psi R^{-1}$ (the matrix $R$ is real). Hence $\bar\Psi'\gamma'^\mu D'_\mu\Psi'=\bar\Psi R^{-1}R\gamma^\mu R^{-1}R\,D_\mu\Psi=\bar\Psi\gamma^\mu D_\mu\Psi$, the same for the second half of $K$, and $S'=S$. The metric, and so $\sqrt{|g|}$, is unchanged. $\square$ (Stage-1 document, §7.5; checks \texttt{LAG\_loca\allowbreak{}lSpinInv\allowbreak{}ariance\_\allowbreak{}G1} and \texttt{LAG\_loca\allowbreak{}lSpinInv\allowbreak{}ariance\_\allowbreak{}G2}.) The proof used $R^TCR=C$. By Theorem 2.8 a product $g$ of $k$ unit vectors has $g^TCg=(-1)^kN(g)\,C$, so the factor that decides whether $S$ is kept is $(-1)^kN(g)$, not the spinor norm $N(g)$ alone. Elements with $(-1)^kN(g)=-1$ reverse the sign of $S$: for example a single space-like unit vector $\gamma^b$ ($b\le3$, $k=1$, $N=+1$) and the even element $\gamma^0\gamma^4$ ($k=2$, $N=-1$), whereas a single time-like $\gamma^b$ ($b\ge4$, $k=1$, $N=-1$) keeps $S$. The single unit vectors $\gamma^b$, used as reflections, are treated in Section 6.12.

\textbf{Changes of coordinates.} Every index in $\mathcal L_s$ is contracted, so $\mathcal L_s$ is a scalar, and $\sqrt{|g|}\,d^8x$ is invariant (Section 4.4): the action does not depend on the coordinates.

\textbf{The global phase.} For a constant real $\alpha$, the change $\Psi\to e^{i\alpha}\Psi$ sends $\Psi^\dagger\to e^{-i\alpha}\Psi^\dagger$. Every term of $\mathcal L$ contains as many factors from $\Psi^\dagger$ as from $\Psi$, so each is multiplied by $e^{-i\alpha}e^{i\alpha}=1$. For Grassmann components the substitution $\Psi_a\to e^{i\alpha}\Psi_a$, $\Psi_a^\ast\to e^{-i\alpha}\Psi_a^\ast$ preserves the order of every product and all the algebra rules, and the same counting applies to every polynomial $U(S)$. So the phase is a symmetry of both fields for every potential (the matter--antimatter document, §4, Theorem M1; checks \texttt{MA\_M1\_u1\allowbreak{}Invarian\allowbreak{}ceCommut\allowbreak{}ingGener\allowbreak{}icU\_G1} and \texttt{MA\_M1\_u1\allowbreak{}Invarian\allowbreak{}ceGrassm\allowbreak{}ann\_G1}). Chapter 7 derives the conserved current $J^\mu=-i\bar\Psi\gamma^\mu\Psi$ that belongs to it.

\subsection{The real restriction and the notebook's \texorpdfstring{$\mathrm{Lg}[\,]$}{Lg[ ]}}

The notebook's Lagrangian of Section 5.12 uses a \textbf{real} column $\Psi$ and the transpose instead of the Dirac adjoint. In the notation of this book, with the canonical connection,

\[
\mathrm{Lg}=\sqrt{|g|}\,\Bigl[\Psi^TC\gamma^\mu D_\mu\Psi+(HM)\,\Psi^TC\Psi\Bigr].
\]

\textbf{Real Grassmann components: nothing.} Theorem 5.6 proved that for a real Grassmann column $\mathrm{Lg}$ is the total divergence $\partial_\mu\bigl(\tfrac12\sqrt{|g|}\,\Psi^TC\gamma^\mu\Psi\bigr)$ in every gravitational field; its mass term vanishes identically because $C$ is symmetric (Lemma 5.4), and all its Euler--Lagrange expressions vanish. Stage 5 re-ran this in flat space and at the point G1 p1 of the general test geometry: the Wolfram check with Stage 1's Grassmann algebra (as its report states), the Python check with the exact Grassmann module \texttt{scripts/\allowbreak{}grassman\allowbreak{}n\_algebr\allowbreak{}a.py} of the Python checkers, a separate program (checks \texttt{C00\_real\allowbreak{}Restrict\allowbreak{}ion\_gras\allowbreak{}smannFla\allowbreak{}tTrivial}, \texttt{C00\_real\allowbreak{}Restrict\allowbreak{}ion\_gras\allowbreak{}smannCur\allowbreak{}vedTrivi\allowbreak{}al} and their \texttt{\detokenize{S5_}} twins; $\mathrm{Lg}$ has 544 monomials at G1 p1, a number measured there, and is nevertheless a pure divergence).

\textbf{Real commuting components: a genuine Lagrangian.} For ordinary numbers Lemma 5.4 is reversed: in $q^TMq$ only the symmetric part of $M$ survives. So now the mass term survives, and the kinetic term is not a total derivative.

\textbf{Theorem 6.5.} Let $X(x)$ be a real column of 16 commuting functions. The Euler--Lagrange expression of $\mathrm{Lg}[X]$ with the canonical connection is

\[
E_{\mathrm{Lg}}=2\sqrt{|g|}\;C\,\bigl(\gamma^\mu D_\mu X+HM\,X\bigr),
\]

so the field equation of $\mathrm{Lg}$ is the Dirac equation $\gamma^\mu D_\mu X=mX$ with $m=-HM$.

\emph{Proof.} Write $\mathrm{Lg}=\sqrt{|g|}\bigl[X^TC\gamma^\mu\partial_\mu X+X^TC\gamma^\mu\Omega_\mu X+HM\,X^TCX\bigr]$ and vary the component $X_c$ (Section 5.3). The derivative of $\mathrm{Lg}$ with respect to $X_c$ gets one contribution from each factor $X$ in a product:

\[
\frac{\partial\mathrm{Lg}}{\partial X_c}=\sqrt{|g|}\,\Bigl[(C\gamma^\mu\partial_\mu X)_c+(C\gamma^\mu\Omega_\mu X)_c+\bigl((C\gamma^\mu\Omega_\mu)^TX\bigr)_c+2HM\,(CX)_c\Bigr],
\]

where the factor 2 comes from the two factors $X$ of $X^TCX$ and the symmetry of $C$. The derivative with respect to $\partial_\mu X_c$ is $\sqrt{|g|}\,\bigl((C\gamma^\mu)^TX\bigr)_c=-\sqrt{|g|}\,(C\gamma^\mu X)_c$, by expression [1]. So the Euler--Lagrange expression is

\[
E_{\mathrm{Lg}}=\sqrt{|g|}\Bigl[C\gamma^\mu\partial_\mu X+C\gamma^\mu\Omega_\mu X+(C\gamma^\mu\Omega_\mu)^TX+2HM\,CX\Bigr]+\partial_\mu\bigl(\sqrt{|g|}\,C\gamma^\mu X\bigr).
\]

Two facts finish the computation. First, $(C\gamma^cS^{ab})^T=CS^{ab}\gamma^c$ (Section 5.12), so $(C\gamma^\mu\Omega_\mu)^T=C\Omega_\mu\gamma^\mu$. Second, by (B6), $\partial_\mu(\sqrt{|g|}\,C\gamma^\mu X)=\sqrt{|g|}\,C\bigl([\gamma^\mu,\Omega_\mu]X+\gamma^\mu\partial_\mu X\bigr)$. Inserting both,

\[
\begin{aligned}
E_{\mathrm{Lg}}&=\sqrt{|g|}\,C\Bigl[2\gamma^\mu\partial_\mu X+\gamma^\mu\Omega_\mu X+\Omega_\mu\gamma^\mu X+\gamma^\mu\Omega_\mu X-\Omega_\mu\gamma^\mu X+2HM\,X\Bigr]\\
&=2\sqrt{|g|}\,C\bigl(\gamma^\mu D_\mu X+HM\,X\bigr). \qquad\square
\end{aligned}
\]

Neither part is trivial. The mass part $2\sqrt{|g|}\,HM\,CX$ is nonzero for every $X\ne0$ because $C$ is invertible, and the coefficient of $\partial_4X$ is $2\sqrt{|g|}\,C\gamma^{x_4}$, which is invertible whenever $g^{44}\ne0$ (Section 8.5 proves $(C\gamma^{x_4})(\gamma^{x_4}C)=g^{44}$), so the equation determines the time derivative. The matrices $\gamma^\mu$ and $\Omega_\mu$ are real, so real solutions exist. (Checks \texttt{C00\_real\allowbreak{}Restrict\allowbreak{}ion\_comm\allowbreak{}utingNon\allowbreak{}Trivial} and \texttt{S5\_realR\allowbreak{}estricti\allowbreak{}on\_commu\allowbreak{}tingNonT\allowbreak{}rivial}: at three G1 points and in the primordial field the Euler--Lagrange expression is exactly $2\sqrt{|g|}\,C(\gamma^\mu D_\mu X+HM\,X)$, and neither the mass part nor the kinetic part is a total divergence.)

\textbf{With the notebook's own contraction.} If $\Omega_\mu$ in $\mathrm{Lg}$ is replaced by the notebook's $\Omega^{\mathrm{nb}}_\mu$ (Section 4.14), the step $(C\gamma^\mu\Omega^{\mathrm{nb}}_\mu)^T=C\Omega^{\mathrm{nb}}_\mu\gamma^\mu$ still holds ($\Omega^{\mathrm{nb}}_\mu$ is also a combination of the $S^{ab}$), but the divergence identity still contains the correct $\Omega_\mu$. Repeating the computation leaves the extra term

\[
E_{\mathrm{Lg}}^{\mathrm{nb}}=2\sqrt{|g|}\,C\bigl(\gamma^\mu D_\mu X+HM\,X\bigr)+\sqrt{|g|}\,C\,\{\gamma^\mu,\Omega^{\mathrm{nb}}_\mu-\Omega_\mu\}\,X .
\]

For a diagonal vielbein both anticommutators vanish (Section 5.14), so in the notebook's primordial field the wrong contraction drops out of the equations for commuting fields; in the non-diagonal G1 the extra term is nonzero (checks \texttt{C00\_real\allowbreak{}Restrict\allowbreak{}ion\_comm\allowbreak{}utingNot\allowbreak{}ebookCon\allowbreak{}traction} and \texttt{S5\_realR\allowbreak{}estricti\allowbreak{}on\_commu\allowbreak{}tingNote\allowbreak{}bookCont\allowbreak{}raction}).

\textbf{How $\mathrm{Lg}$ sits inside the Lagrangian of dirac16complex00.} For commuting components write $\Psi=X+iY$ with real columns $X,Y$, and put $K_R[Z]:=Z^TC\gamma^\mu D_\mu Z$ and $S_Z:=Z^TCZ$ for a real column $Z$.

\textbf{Theorem 6.6.} For commuting components, $S=S_X+S_Y$, $K=K_R[X]+K_R[Y]$, and therefore

\[
\mathcal L[X+iY]=\sqrt{|g|}\,\Bigl[K_R[X]+K_R[Y]-m\bigl(S_X+S_Y\bigr)-U\bigl(S_X+S_Y\bigr)\Bigr].
\]

For $U=0$ this is $\mathrm{Lg}[X]+\mathrm{Lg}[Y]$ with $HM=-m$.

\emph{Proof.} $S=(X^T-iY^T)C(X+iY)=X^TCX+Y^TCY+i(X^TCY-Y^TCX)$, and $X^TCY=Y^TCX$ for commuting numbers and symmetric $C$. For the kinetic term, $\bar\Psi\gamma^\mu D_\mu\Psi=(X^T-iY^T)C\gamma^\mu(D_\mu X+iD_\mu Y)$ has the real part $X^TC\gamma^\mu D_\mu X+Y^TC\gamma^\mu D_\mu Y$ and the imaginary part $X^TC\gamma^\mu D_\mu Y-Y^TC\gamma^\mu D_\mu X$. In the same way $(D_\mu\bar\Psi)\gamma^\mu\Psi$ has the real part $(D_\mu X)^TC\gamma^\mu X+(D_\mu Y)^TC\gamma^\mu Y$ and the imaginary part $(D_\mu X)^TC\gamma^\mu Y-(D_\mu Y)^TC\gamma^\mu X$. A number equals its own transpose, and $(C\gamma^\mu)^T=-C\gamma^\mu$, so $(D_\mu Z)^TC\gamma^\mu W=-W^TC\gamma^\mu D_\mu Z$ for any real $Z,W$. Therefore the real part of $(D_\mu\bar\Psi)\gamma^\mu\Psi$ is $-K_R[X]-K_R[Y]$, and its imaginary part is $-Y^TC\gamma^\mu D_\mu X+X^TC\gamma^\mu D_\mu Y$, equal to the imaginary part of the first product. In $K=\tfrac12(\text{first}-\text{second})$ the imaginary parts cancel and the real parts add: $K=K_R[X]+K_R[Y]$. $\square$

So the Lagrangian of dirac16complex00 is the \textbf{complexification} of the notebook's real Lagrangian: two copies of $\mathrm{Lg}$, one for the real part and one for the imaginary part of the field, coupled only through $U$. Setting $Y=0$ is consistent: the Euler--Lagrange expression for $Y$ is $2\sqrt{|g|}\,C\bigl(\gamma^\mu D_\mu Y-(m+U'(S))Y\bigr)$ (Theorem 6.5 with $U$ added), which vanishes at $Y=0$. With $Y=0$ the field equation of $X$ is $\gamma^\mu D_\mu X=\bigl(m+U'(S_X)\bigr)X$. For $U=0$ this is the notebook's equation $\gamma^\mu D_\mu X=mX$ with $m=-HM$ (Theorem 6.5); with a potential the mass $m$ is replaced by $m+U'(S_X)$, a term that the notebook's $\mathrm{Lg}[\,]$ does not contain. In this precise sense, \textbf{for $U=0$ the real restriction of dirac16complex00 is the notebook's $\mathrm{Lg}[\,]$ with the canonical connection}, with $U\ne0$ it is $\mathrm{Lg}[\,]$ plus the self-interaction $-\sqrt{|g|}\,U(S_X)$, and $\mathrm{Lg}[\,]$ is a genuine Lagrangian for commuting fields (checks \texttt{C00\_real\allowbreak{}Restrict\allowbreak{}ion\_rela\allowbreak{}tionToL1}, whose recorded $X$ equation is $2\sqrt{|g|}\,C\bigl(\gamma^\mu D_\mu X-(m+\lambda S_X)X\bigr)$ for the default potential, \texttt{C00\_real\allowbreak{}Restrict\allowbreak{}ion\_comm\allowbreak{}utingFla\allowbreak{}tDecompo\allowbreak{}sition} and their \texttt{\detokenize{S5_}} twins).

\textbf{For Grassmann components the opposite happens.} Write $\Psi=\rho+i\iota$ with real Grassmann columns $\rho,\iota$. Then $\rho^TC\rho=\iota^TC\iota=0$ (Lemma 5.4), and $\iota^TC\rho=\sum_{a,b}\iota_aC_{ab}\rho_b=-\sum_{a,b}\rho_bC_{ab}\iota_a=-\rho^TC\iota$, so

\[
S=\rho^TC\rho+\iota^TC\iota+i\bigl(\rho^TC\iota-\iota^TC\rho\bigr)=2i\,\rho^TC\iota :
\]

the mass term of the complex Grassmann field is a pure cross term between its real and imaginary parts, and each real part alone carries neither a mass term nor (Theorem 5.6) any dynamics (checks \texttt{C00\_real\allowbreak{}Restrict\allowbreak{}ion\_gras\allowbreak{}smannCom\allowbreak{}plexDeco\allowbreak{}mpositio\allowbreak{}n} and \texttt{S5\_realR\allowbreak{}estricti\allowbreak{}on\_grass\allowbreak{}mannComp\allowbreak{}lexDecom\allowbreak{}position}). This is why dirac16complex must be complex.

\subsection{The chirality map \texorpdfstring{$\gamma^8$}{γ\textasciicircum{}8} of the Lagrangians}

The chirality $\gamma^8=\gamma^0\gamma^1\cdots\gamma^7=\mathrm{diag}(-I_8,I_8)$ of Section 2.10 multiplies the upper eight components by $-1$ and leaves the lower eight unchanged. Mapping $\Psi$ to $\gamma^8\Psi$ turns the Lagrangian with mass $m$ into minus the Lagrangian with mass $-m$. This is the Lagrangian form of the pairing of the masses $\pm m$ that Chapter 15 develops into the pairing theorems.

\textbf{Five facts about $\gamma^8$} (all proved in Section 2.10 or immediate from them): (G1) $\gamma^8$ is real and diagonal, hence Hermitian; (G2) $(\gamma^8)^2=1$; (G3) $\gamma^8\gamma^a=-\gamma^a\gamma^8$ for every $a$; (G4) $\gamma^8$ commutes with $C$ and with every $S^{ab}$, both even products of gammas; (G5) consequently $\gamma^8$ commutes with every $\Omega_\mu$, in every gravitational field, and anticommutes with every $\gamma^\mu=e_a{}^\mu\gamma^a$. (Checks \texttt{PAIR\_alg\allowbreak{}ebra\_gam\allowbreak{}ma8Prope\allowbreak{}rties} and \texttt{C00\_alge\allowbreak{}bra\_lead\allowbreak{}Facts}.)

\textbf{Theorem 6.7 (the $\gamma^8$ map).} For every vielbein, every mass $m$, every potential $U$ and every field configuration $\Psi$ (not necessarily a solution), for both kinds of components,

\[
\mathcal L_{m,U}[\gamma^8\Psi]=-\mathcal L_{-m,-U}[\Psi],\qquad\text{in particular}\qquad \mathcal L_{m,\lambda}[\gamma^8\Psi]=-\mathcal L_{-m,-\lambda}[\Psi].
\]

\emph{Proof.} Put $\Psi_-:=\gamma^8\Psi$.

\begin{enumerate}
\item The adjoint: $\bar\Psi_-=(\gamma^8\Psi)^\dagger C=\Psi^\dagger\gamma^8C=\Psi^\dagger C\gamma^8=\bar\Psi\gamma^8$, by (G1) and (G4).
\item The derivatives: by (G5), $D_\mu(\gamma^8\Psi)=\gamma^8\partial_\mu\Psi+\Omega_\mu\gamma^8\Psi=\gamma^8D_\mu\Psi$, and in the same way $D_\mu\bar\Psi_-=(D_\mu\bar\Psi)\gamma^8$.
\item The scalar density: $S[\Psi_-]=\bar\Psi\gamma^8\gamma^8\Psi=S[\Psi]$ by (G2).
\item The kinetic term: $\bar\Psi_-\gamma^\mu D_\mu\Psi_-=\bar\Psi\gamma^8\gamma^\mu\gamma^8D_\mu\Psi=-\bar\Psi\gamma^\mu D_\mu\Psi$ by (G5) and (G2), and likewise $(D_\mu\bar\Psi_-)\gamma^\mu\Psi_-=-(D_\mu\bar\Psi)\gamma^\mu\Psi$. So $K[\Psi_-]=-K[\Psi]$.
\item Therefore $\mathcal L_{m,U}[\Psi_-]=\sqrt{|g|}\bigl[-K-mS-U(S)\bigr]=-\sqrt{|g|}\bigl[K-(-m)S-(-U)(S)\bigr]=-\mathcal L_{-m,-U}[\Psi]$.
\end{enumerate}

For Grassmann components nothing was reordered: in the notebook's basis the map is the substitution $\Psi_a\to-\Psi_a$, $\Psi_a^\ast\to-\Psi_a^\ast$ for $a\le7$ (and nothing for $a\ge8$), which respects every rule of the Grassmann algebra. $\square$

In words: the mass term and the interaction do not change, the kinetic term changes sign, and so the image is the theory with the opposite sign of the whole Lagrangian and the opposite mass. Because the potential does not change sign under the map, it must be given the opposite sign on the right-hand side: for the default potential the map relates $(m,\lambda)$ to $(-m,-\lambda)$, \textbf{not} to $(-m,\lambda)$. Indeed

\[
\mathcal L_{m,\lambda}[\gamma^8\Psi]+\mathcal L_{-m,\lambda}[\Psi]=\sqrt{|g|}\bigl[-K-mS-\tfrac\lambda2S^2+K+mS-\tfrac\lambda2S^2\bigr]=-\lambda\sqrt{|g|}\,S^2\ne0
\]

for $\lambda\ne0$. The simpler statement "$\mathcal L_m\to-\mathcal L_{-m}$" of the original project contract holds only for $U=0$; this is erratum E2 of \texttt{handoff/\allowbreak{}specs/CO\allowbreak{}NTRACT.m\allowbreak{}d} (§11). (Checks: \texttt{ALG\_gamm\allowbreak{}a8Map} of Stage 1; \texttt{PAIR\_T1g\allowbreak{}eneric\_l\allowbreak{}agrangia\allowbreak{}n} as an identity in completely generic symbols, \texttt{PAIR\_T1g\allowbreak{}rassmann\allowbreak{}\_lagrang\allowbreak{}ian} in a Grassmann algebra of 288 odd generators, \texttt{PAIR\_T1j\allowbreak{}ets\_lagr\allowbreak{}angian} with exact jets of the general vielbein G1, \texttt{PAIR\_T1p\allowbreak{}rimordia\allowbreak{}l\_lagran\allowbreak{}gian} for arbitrary fields in the primordial field; \texttt{PAIR\_T1g\allowbreak{}eneric\_n\allowbreak{}aiveFixe\allowbreak{}dLambdaF\allowbreak{}ails} for the fixed-$\lambda$ failure.)

\textbf{The same map, seen as reversing the frame.} Change the vielbein to $-e_\mu{}^a$ and keep $\Psi$. The metric $g=e\eta e^T$ and $\sqrt{|g|}=|\det e|$ do not change ($\det(-e)=(-1)^8\det e$); the inverse vielbein and the curved gammas change sign; the Christoffel symbols do not change, and neither does the spin connection, because in $\omega_\mu{}^a{}_b=e_b{}^\nu(\Gamma^\rho{}_{\mu\nu}e_\rho{}^a-\partial_\mu e_\nu{}^a)$ (Section 4.11) both factors change sign. So the kinetic term changes sign and $S$ does not: $\mathcal L_{m,\lambda}[-e,\Psi]=-\mathcal L_{-m,-\lambda}[e,\Psi]$, the same as Theorem 6.7. Together the two changes cancel: $\mathcal L_{m,\lambda}[-e,\gamma^8\Psi]=\mathcal L_{m,\lambda}[e,\Psi]$. This is no accident. The frame change $-1_8$ is a rotation by $\pi$ in each of the four planes $(0,1)$, $(2,3)$, $(4,5)$, $(6,7)$, and its spin cover is $\exp(\pi S^{01})\exp(\pi S^{23})\exp(\pi S^{45})\exp(\pi S^{67})=\gamma^0\gamma^1\,\gamma^2\gamma^3\,\gamma^4\gamma^5\,\gamma^6\gamma^7=\gamma^8$ (use $\exp(\pi S^{ab})=\gamma^a\gamma^b$ for a rotation, Exercise 2.14, and note that $(4,5)$ and $(6,7)$ are rotations because $\eta^{44}\eta^{55}=\eta^{66}\eta^{77}=+1$). So $(e,\Psi)\to(-e,\gamma^8\Psi)$ is a local frame rotation of Section 6.9, under which $\mathcal L$ is invariant, and the $\gamma^8$ map of Theorem 6.7 is the same as reversing all eight frame directions at a fixed field (checks \texttt{PAIR\_alg\allowbreak{}ebra\_gam\allowbreak{}ma8InIde\allowbreak{}ntityCom\allowbreak{}ponent}, \texttt{PAIR\_T1j\allowbreak{}ets\_viel\allowbreak{}beinSign\allowbreak{}FlipGeom\allowbreak{}etry}, \texttt{PAIR\_T1j\allowbreak{}ets\_viel\allowbreak{}beinSign\allowbreak{}FlipIsT1}, \texttt{PAIR\_T1j\allowbreak{}ets\_gamm\allowbreak{}a8WithFr\allowbreak{}ameSignI\allowbreak{}sSymmetr\allowbreak{}y}).

\textbf{Worked example.} Flat space, the constant commuting field $\Psi=e_0+e_4+i\,e_{12}$. Its scalar density is $S=-2$ (the pair $e_0,e_4$ gives $-2$ as in Section 6.3, and $e_{12}$ pairs only with $e_8$, which is absent). The current $J^a=-i\bar\Psi\gamma^a\Psi$ of Chapter 7 has two nonzero components. For example $(C\gamma^0)_{0,12}=-1$ and $(C\gamma^0)_{12,0}=+1$ (Section 5.12), so $\bar\Psi\gamma^0\Psi=\Psi_0^\ast(-1)\Psi_{12}+\Psi_{12}^\ast(+1)\Psi_0=-i-i=-2i$ and $J^0=-2$; in the same way $J^7=+2$. The image $\gamma^8\Psi=-e_0-e_4+i\,e_{12}$ has $S=-2$ again and $J^0=+2$, $J^7=-2$: the scalar is unchanged and the vector bilinears change sign. The field is constant, so $K=0$, and for $U=\tfrac\lambda2S^2$ both $\mathcal L_{m,\lambda}[\gamma^8\Psi]$ and $-\mathcal L_{-m,-\lambda}[\Psi]=-\bigl(2(-m)-2(-\lambda)\bigr)$ equal $2m-2\lambda$.

\textbf{What the map does and does not say.} Chapter 7 derives the consequences for the field equations (solutions with $(m,\lambda)$ go to solutions with $(-m,-\lambda)$), for the energy--momentum tensor ($T_{\mu\nu}\to-T_{\mu\nu}$) and for the current ($J^\mu\to-J^\mu$). Chapter 15 turns these into the pairing theorems. The Stage-1 document states their nature in one sentence, which this book repeats: the pairing of $\pm m$ is "a structural property of the equations, not a claim that universes of masses $\pm M$ are created in pairs" (Stage-1 document, §1.3, item 6). Chapter 16 explains exactly what is and what is not established.

\subsection{Two more discrete maps: reflections and complex conjugation}

Two other kinds of discrete map relate the Lagrangians with different signs of the mass, and one of them shows the first place where the statistics decides whether a map is a symmetry. We derive the reflections in flat space (their curved-space form, in which the frame is reflected together with the field, is recorded in the reports named below and proved in Chapter 15) and the complex conjugations in every gravitational field.

\textbf{Reflections.} For a frame direction $b$ let $R_b$ reverse the coordinate $x_b$, and write $r_a=-1$ for $a=b$ and $r_a=+1$ otherwise. Define

\[
\Psi'(x)=\gamma^b\,\Psi(R_bx).
\]

This is the action of the unit vector $\gamma^b$ of Pin(4,4) (Section 2.14), in the lift that Stage 5 calls twisted.

\textbf{Theorem 6.8.} In flat space, for both kinds of components, with $n_b=\eta^{bb}$: $K\to n_b\,K$ and $S\to-n_b\,S$ (evaluated at the reflected point). Consequently

\[
\mathcal L_{m,\lambda}[\Psi']=\mathcal L_{-m,\lambda}[\Psi]\quad(b\le3,\ \text{space-like}),\qquad \mathcal L_{m,\lambda}[\Psi']=-\mathcal L_{-m,-\lambda}[\Psi]\quad(b\ge4,\ \text{time-like}).
\]

\emph{Proof.} $\Psi'^\dagger=\Psi^\dagger(\gamma^b)^T$ and, by the chain rule, $\partial_a\Psi'(x)=r_a\gamma^b(\partial_a\Psi)(R_bx)$. From (B2), $(\gamma^b)^TC=-C\gamma^b$. Hence $(\gamma^b)^TC\gamma^b=-C(\gamma^b)^2=-n_bC$, which gives $S\to-n_bS$. For the kinetic matrices, $(\gamma^b)^TC\gamma^a\gamma^b=-C\gamma^b\gamma^a\gamma^b$. For $a\ne b$, $\gamma^b\gamma^a\gamma^b=-\gamma^a(\gamma^b)^2=-n_b\gamma^a$, so the product is $n_bC\gamma^a$, and $r_a=+1$. For $a=b$, $\gamma^b\gamma^b\gamma^b=n_b\gamma^b$, so the product is $-n_bC\gamma^b$, and $r_b=-1$. In both cases $r_a(\gamma^b)^TC\gamma^a\gamma^b=n_b\,C\gamma^a$, and both halves of $K$ are multiplied by $n_b$. For $n_b=+1$: $\mathcal L\to K+mS-\tfrac\lambda2S^2=\mathcal L_{-m,\lambda}$. For $n_b=-1$: $\mathcal L\to-K-mS-\tfrac\lambda2S^2=-\mathcal L_{-m,-\lambda}$. $\square$

A reflection of a space-like direction therefore relates $m$ to $-m$ \textbf{with the same} $\lambda$ and the same sign of the Lagrangian. In a curved field the same statement holds when the frame is reflected together with the field (Stage 5 calls this the mirror map T2; Chapter 15 proves it). The other lift, $\Psi'=\Gamma\Psi(R_bx)$ with $\Gamma$ the product of the seven gammas other than $\gamma^b$ (which is $\pm\gamma^8\gamma^b$), has $K\to-n_bK$, $S\to-n_bS$: it gives $-\mathcal L_{m,-\lambda}$ for a space-like $b$ and an exact symmetry for a time-like $b$ (erratum E3 of the contract). The factor $-n_b$ of $S$ is the Pin character of Section 2.14. (Checks \texttt{ALG\_pinL\allowbreak{}iftChara\allowbreak{}cter}, \texttt{MA\_M2\_na\allowbreak{}medRefle\allowbreak{}ctionsFl\allowbreak{}atJets}, \texttt{PAIR\_T2f\allowbreak{}rame\_twi\allowbreak{}stedSpac\allowbreak{}elikeMap\allowbreak{}sToMinus\allowbreak{}MSameLam\allowbreak{}bda}, \texttt{PAIR\_T2f\allowbreak{}rame\_twi\allowbreak{}stedTime\allowbreak{}like}, \texttt{PAIR\_T2f\allowbreak{}rame\_unt\allowbreak{}wistedSp\allowbreak{}acelikeC\allowbreak{}ontractE\allowbreak{}3}.)

\textbf{Complex conjugation.} Now let $\Psi'=\Psi^\ast$ (the map called $C_0$) or $\Psi'=\gamma^8\Psi^\ast$ (called $C_8$). Here a reordering cannot be avoided, and it produces the \textbf{statistics sign} $s$, with $s=+1$ for commuting and $s=-1$ for Grassmann components.

\textbf{Lemma 6.9 (statistics sign).} For columns $\Psi,\Phi$ of the same kind and any numerical matrix $Y$, $\Psi^TY\Phi^\ast=s\,\Phi^\dagger Y^T\Psi$.

\emph{Proof.} $\sum_{a,b}\Psi_aY_{ab}\Phi_b^\ast=s\sum_{a,b}\Phi_b^\ast Y_{ab}\Psi_a$, because exchanging two odd factors costs a sign and exchanging two commuting ones does not. $\square$

\textbf{Theorem 6.10.} In every gravitational field: $C_0$ sends $K\to s\,K$ and $S\to s\,S$; $C_8$ sends $K\to-s\,K$ and $S\to s\,S$. Hence

\begingroup
\small
\begin{longtable}{@{}>{\raggedright\arraybackslash}p{0.267\linewidth}>{\raggedright\arraybackslash}p{0.267\linewidth}>{\raggedright\arraybackslash}p{0.267\linewidth}@{}}
\toprule
\textbf{map} & \textbf{commuting components ($s=+1$)} & \textbf{Grassmann components ($s=-1$)} \\
\midrule
\endfirsthead
\toprule
\textbf{map} & \textbf{commuting components ($s=+1$)} & \textbf{Grassmann components ($s=-1$)} \\
\midrule
\endhead
$C_0$: $\Psi\to\Psi^\ast$ & $\mathcal L_{m,\lambda}$ (exact symmetry) & $-\mathcal L_{m,-\lambda}$ \\
$C_8$: $\Psi\to\gamma^8\Psi^\ast$ & $-\mathcal L_{-m,-\lambda}$ & $\mathcal L_{-m,\lambda}$ \\
\bottomrule
\end{longtable}
\endgroup

\emph{Proof for $C_0$.} The conjugate of $\Psi'=\Psi^\ast$ is $\Psi$, so $\Psi'^\dagger=\Psi^T$. The gammas and $\Omega_\mu$ are real, so $D_\mu\Psi^\ast=(D_\mu\Psi)^\ast$. By Lemma 6.9 with $Y=C$ (symmetric), $S'=\Psi^TC\Psi^\ast=s\,\Psi^\dagger C\Psi=sS$. With $Y=C\gamma^\mu$ (antisymmetric), $\Psi^TC\gamma^\mu D_\mu\Psi^\ast=s\,(D_\mu\Psi)^\dagger(C\gamma^\mu)^T\Psi=-s\,(D_\mu\bar\Psi)\gamma^\mu\Psi$ and $-(D_\mu\Psi)^TC\gamma^\mu\Psi^\ast=-s\,\Psi^\dagger(C\gamma^\mu)^TD_\mu\Psi=s\,\bar\Psi\gamma^\mu D_\mu\Psi$. So $K'=\tfrac12\bigl(s\bar\Psi\gamma^\mu D_\mu\Psi-s(D_\mu\bar\Psi)\gamma^\mu\Psi\bigr)=sK$. For $s=+1$ nothing changes; for $s=-1$, $\mathcal L\to\sqrt{|g|}[-K+mS-\tfrac\lambda2S^2]=-\mathcal L_{m,-\lambda}$. \emph{For $C_8$}, compose with Theorem 6.7 ($K\to-K$, $S\to S$). $\square$

The charge density $J^4=\Psi^\dagger B\Psi$ of Chapter 7 transforms under $C_0$ into $\Psi^TB\Psi^\ast=s\Psi^\dagger B^T\Psi=-s\,J^4$ (since $B^T=-B$). So for the commuting field $C_0$ is an exact symmetry that \textbf{reverses the charge}: it is the charge conjugation of dirac16complex00. For the Grassmann field no constant charge conjugation is a symmetry when $m\ne0$: $C_0$ reverses the sign of the Lagrangian and of $\lambda$, and $C_8$ reverses the mass; and Lemma A of Theorem M2 (quoted here, not proved) shows that every constant charge conjugation compatible with the Dirac operator is, up to a constant factor, $C_0$ or $C_8$. But $C_8$ combined with the reflection of one space-like direction (Theorem 6.8, which also sends $m\to-m$ with the same $\lambda$) is an exact symmetry of dirac16complex, of the kind called CP, and it reverses the charge. These classifications are Theorem M2 of the matter--antimatter document (its §5.3 to §5.5, verified for all 1024 maps of this kind for each statistics; checks \texttt{MA\_M2\_la\allowbreak{}grangian\allowbreak{}FlatComm\allowbreak{}uting\_al\allowbreak{}l}, \texttt{MA\_M2\_la\allowbreak{}grangian\allowbreak{}FlatGras\allowbreak{}smann\_al\allowbreak{}l}, \texttt{MA\_M2\_in\allowbreak{}ternalMa\allowbreak{}psCurved\allowbreak{}Jets}, \texttt{MA\_M2\_C\_\allowbreak{}and\_CP\_s\allowbreak{}tatus}). Chapter 17 explains what they mean for matter and antimatter.

\subsection{How the repository checks all of this}

Every statement of Sections 6.2 to 6.12 that the repository verifies is listed with its checks. "Both" means the Wolfram check (\texttt{\detokenize{C00_...}}) and its independent Python twin (\texttt{\detokenize{S5_...}} with the same ending, in \texttt{python-d\allowbreak{}irac16co\allowbreak{}mplex00-\allowbreak{}report.j\allowbreak{}son}). Each \texttt{\detokenize{PAIR_}} check of \texttt{wolfram-\allowbreak{}pairing-\allowbreak{}report.j\allowbreak{}son} has a Python twin as well: the \texttt{\detokenize{S5_}} check with the same ending in \texttt{python-p\allowbreak{}airing-r\allowbreak{}eport.js\allowbreak{}on} (for example \texttt{S5\_T1gen\allowbreak{}eric\_lag\allowbreak{}rangian}).

\begingroup
\small
\begin{longtable}{@{}>{\raggedright\arraybackslash}p{0.400\linewidth}>{\raggedright\arraybackslash}p{0.400\linewidth}@{}}
\toprule
\textbf{statement (section)} & \textbf{checks} \\
\midrule
\endfirsthead
\toprule
\textbf{statement (section)} & \textbf{checks} \\
\midrule
\endhead
same Pin(4,4) module for both fields; invariant forms $CP_\mp$; $\bar\Psi\Psi$ invariant (6.2, 6.3) & \texttt{C00\_alge\allowbreak{}bra\_pinM\allowbreak{}oduleAnd\allowbreak{}SpinInva\allowbreak{}riance} (both) \\
$\gamma^8$, $C$, $B$ facts (6.3, 6.11) & \texttt{C00\_alge\allowbreak{}bra\_lead\allowbreak{}Facts} (both), \texttt{PAIR\_alg\allowbreak{}ebra\_gam\allowbreak{}ma8Prope\allowbreak{}rties} \\
reality for both statistics; unsymmetrized term not real (6.5) & \texttt{C00\_lagr\allowbreak{}angian\_r\allowbreak{}ealCommu\allowbreak{}ting}, \texttt{C00\_lagr\allowbreak{}angian\_c\allowbreak{}oefficie\allowbreak{}ntHermit\allowbreak{}icity}, \texttt{C00\_lagr\allowbreak{}angian\_g\allowbreak{}rassmann\allowbreak{}FlatHerm\allowbreak{}itianAnd\allowbreak{}EL} (both); \texttt{LAG\_herm\allowbreak{}iticity\_\allowbreak{}G1}, \texttt{\detokenize{_G2}}; \texttt{GR\_lagra\allowbreak{}ngianHer\allowbreak{}mitian} \\
mass term non-zero; dispersion $\omega^2=m^2+\dots$ (6.6) & \texttt{C00\_mass\allowbreak{}Term\_com\allowbreak{}mutingEx\allowbreak{}plicitSp\allowbreak{}inors}, \texttt{C00\_mass\allowbreak{}Term\_gra\allowbreak{}ssmannNo\allowbreak{}nzero}, \texttt{C00\_mass\allowbreak{}Term\_dis\allowbreak{}persionF\allowbreak{}lat} (both) \\
any smooth $U$ for the commuting field (6.7) & \texttt{C00\_EL\_c\allowbreak{}ommuting\allowbreak{}GeneralS\allowbreak{}moothU}, \texttt{C00\_EMT\_\allowbreak{}generalS\allowbreak{}moothU} (both) \\
only $\omega_{[cab]}$ in the Lagrangian; connection in the field equation (6.8) & \texttt{C00\_alge\allowbreak{}bra\_anti\allowbreak{}commutat\allowbreak{}orTotall\allowbreak{}yAntisym\allowbreak{}metric}, \texttt{C00\_conn\allowbreak{}ection\_n\allowbreak{}onTrivia\allowbreak{}lCouplin\allowbreak{}g}, \texttt{C00\_conn\allowbreak{}ection\_O\allowbreak{}megaTerm\allowbreak{}InFieldE\allowbreak{}quation} (both) \\
local spin invariance (6.9) & \texttt{LAG\_loca\allowbreak{}lSpinInv\allowbreak{}ariance\_\allowbreak{}G1}, \texttt{LAG\_loca\allowbreak{}lSpinInv\allowbreak{}ariance\_\allowbreak{}G2} \\
real restriction: trivial (Grassmann), non-trivial (commuting) (6.10) & \texttt{C00\_real\allowbreak{}Restrict\allowbreak{}ion\_gras\allowbreak{}smannFla\allowbreak{}tTrivial}, \texttt{C00\_real\allowbreak{}Restrict\allowbreak{}ion\_gras\allowbreak{}smannCur\allowbreak{}vedTrivi\allowbreak{}al}, \texttt{C00\_real\allowbreak{}Restrict\allowbreak{}ion\_comm\allowbreak{}utingNon\allowbreak{}Trivial}, \texttt{C00\_real\allowbreak{}Restrict\allowbreak{}ion\_comm\allowbreak{}utingNot\allowbreak{}ebookCon\allowbreak{}traction} (both) \\
$\mathcal L[X+iY]=\mathrm{Lg}[X]+\mathrm{Lg}[Y]$; Grassmann cross term (6.10) & \texttt{C00\_real\allowbreak{}Restrict\allowbreak{}ion\_rela\allowbreak{}tionToL1}, \texttt{C00\_real\allowbreak{}Restrict\allowbreak{}ion\_comm\allowbreak{}utingFla\allowbreak{}tDecompo\allowbreak{}sition}, \texttt{C00\_real\allowbreak{}Restrict\allowbreak{}ion\_gras\allowbreak{}smannCom\allowbreak{}plexDeco\allowbreak{}mpositio\allowbreak{}n} (both) \\
$\mathcal L_{m,\lambda}[\gamma^8\Psi]=-\mathcal L_{-m,-\lambda}[\Psi]$ (6.11) & \texttt{ALG\_gamm\allowbreak{}a8Map}; \texttt{PAIR\_T1g\allowbreak{}eneric\_l\allowbreak{}agrangia\allowbreak{}n}, \texttt{PAIR\_T1g\allowbreak{}rassmann\allowbreak{}\_lagrang\allowbreak{}ian}, \texttt{PAIR\_T1j\allowbreak{}ets\_lagr\allowbreak{}angian}, \texttt{PAIR\_T1p\allowbreak{}rimordia\allowbreak{}l\_lagran\allowbreak{}gian} \\
reflections and conjugations (6.12) & \texttt{ALG\_pinL\allowbreak{}iftChara\allowbreak{}cter}; \texttt{PAIR\_T2f\allowbreak{}rame\_*}; \texttt{\detokenize{MA_M2_*}} \\
\bottomrule
\end{longtable}
\endgroup

The reports live in three folders, for Stage 1, for Stage 5 and for the matter--antimatter analysis:

\begin{Verbatim}[fontsize=\small]
artifacts/dirac16complex/arbitrary-field/
artifacts/dirac16complex/pair-creation/
artifacts/dirac16complex/matter-antimatter/
\end{Verbatim}

Every check named here is true in the committed reports. To count the checks of the two dirac16complex00 reports, use the method of Section 0.8 with the paths

\begin{Verbatim}[fontsize=\small]
artifacts/dirac16complex/pair-creation/wolfram-dirac16complex00-report.json
artifacts/dirac16complex/pair-creation/python-dirac16complex00-report.json
\end{Verbatim}

which gives \texttt{\detokenize{(46, 46)}} and \texttt{\detokenize{(49, 49)}}. The Python checker also reads the Wolfram report and compares the exact values that both compute (check \texttt{S5\_agree\allowbreak{}sWithWol\allowbreak{}fram}).

\subsection{What we proved and what we assumed}

\textbf{What we proved.} With complete derivations in this chapter: the invariant bilinear forms of the spin generators are exactly the combinations of $CP_-$ and $CP_+$ (Section 6.2, from Theorem 2.12); $\bar\Psi\Psi$ is invariant under spin rotations while $\Psi^\dagger\Psi$ is not (worked example with a boost), and $\bar\Psi\gamma^a\Phi$ is a vector; the conjugate of a bilinear is $\Phi^\dagger M^\dagger\Psi$ for commuting and for Grassmann components (Lemma 6.1), and the Lagrangian is real for both (Theorem 6.2), the symmetrization differing from the plain kinetic term by an imaginary total derivative; the mass term is linear in $m$, quadratic in the components and non-zero for both kinds of components, and it is non-trivial because it enters the dispersion relation as $m^2$, with an 8-dimensional space of plane-wave solutions on the mass shell for $m\ne0$; for Grassmann components the interaction can only be a polynomial of degree at most 16, while for commuting components it can be any smooth function; only the totally antisymmetric part of the spin connection appears in the Lagrangian, and none of it for a diagonal vielbein (Lemma 6.3, Theorem 6.4), while the field equation contains $\gamma^\mu\Omega_\mu=\tfrac1{2\sqrt{|g|}}\partial_\mu(\sqrt{|g|}\gamma^\mu)$ in diagonal frames; the Lagrangian is invariant under local frame rotations of spinor norm $+1$, changes of coordinates and the global phase, while Pin elements with $(-1)^kN(g)=-1$ reverse the sign of $S$; for real commuting fields the notebook's $\mathrm{Lg}[\,]$ gives the Dirac equation with $m=-HM$ (Theorem 6.5), and the Lagrangian of dirac16complex00 is two copies of it coupled through $U$ (Theorem 6.6), so that its real restriction is $\mathrm{Lg}[\,]$ for $U=0$ and $\mathrm{Lg}[\,]$ plus the self-interaction otherwise, whereas for Grassmann fields the mass term is a pure cross term between real and imaginary parts; the chirality map gives $\mathcal L_{m,U}[\gamma^8\Psi]=-\mathcal L_{-m,-U}[\Psi]$ for both statistics and fails at fixed $\lambda\ne0$ (Theorem 6.7), and it is equivalent to reversing all eight frame vectors; reflections of a space-like direction give $\mathcal L_{-m,\lambda}$ (Theorem 6.8); complex conjugation is an exact, charge-reversing symmetry of the commuting field but not of the Grassmann field, for which $C_8$ combined with a space-like reflection is (Theorem 6.10).

\textbf{What we assumed or took from elsewhere.} The algebraic properties of the gamma matrices, $C$, $\gamma^8$ and the spin generators, including Theorems 2.8, 2.11 and 2.12 and the facts of linear algebra quoted in Section 2.2 (Chapter 2), and the geometry of the vielbein, the canonical spin connection, the covariant derivatives and the divergence identity (Chapter 4). The form of the Lagrangian is a choice of the project: the symmetrized kinetic term, the Dirac adjoint with $C$, the default potential $\tfrac\lambda2S^2$ and the conditions $U(0)=U'(0)=0$ are conventions, not derivations (Stage-1 document, §1.2). That fermion components must be anticommuting is the physical premise of dirac16complex; that dirac16complex00 is only a classical field is a decision of the project, and no spin--statistics theorem is proved for signature (4,4). The reflections were derived here in flat space only; their curved-space form, with the frame reflected together with the field, is taken from the Stage-5 and matter--antimatter reports and proved in Chapter 15; the completeness of the classification of the constant discrete maps (Theorem M2, which examines all 1024 maps of the kind of Section 6.12 for each statistics) is quoted from the matter--antimatter reports and treated in Chapter 17. Numbers measured at test points, such as the 544 monomials of $\mathrm{Lg}$ at G1 p1 or the 128 nonzero entries of $\gamma^\mu\Omega_\mu$ at the G1 points, are quoted from the reports, not derived. The machine checks at test points confirm the general derivations; they do not replace them.

\subsection{Exercises}

\textbf{Exercise 6.1.} Using the table of $C$ in Section 2.9, compute $S=\bar\Psi\Psi$ for the commuting columns $\Psi=e_0+2e_4+e_8$, $\Psi=e_0+(1+i)e_4$ and $\Psi=e_0+(1-i)e_4$.

\textbf{Exercise 6.2.} Show that $\Psi^\dagger B\Psi$ is real for every commuting $\Psi$, where $B=-iC\gamma^4$, using Lemma 6.1.

\textbf{Exercise 6.3.} For the boost $R=\cosh\tfrac\theta2+\sinh\tfrac\theta2\,\gamma^0\gamma^4$ verify $R^TCR=C$ directly from $\gamma^0C=-C\gamma^0$ and $\gamma^4C=C\gamma^4$.

\textbf{Exercise 6.4.} In flat space with $m=2$, find the frequency $\omega$ of the plane waves with (a) $k_1=1$, $k_5=1$; (b) $k_0=3$, $k_6=5$; (c) $k_1=4$, $k_7=5$. In which case does no real frequency exist?

\textbf{Exercise 6.5.} For the field $\Psi=f(x_1)e_0+g(x_1)e_{11}$ of Section 6.5 take $f=\cos(kx_1)$ and $g=i\sin(kx_1)$ with real $k$. Compute $X=\bar\Psi\gamma^1\partial_1\Psi$, the kinetic term $K$ and $\bar\Psi\gamma^1\Psi$.

\textbf{Exercise 6.6.} For Grassmann components show that $U(S)=e^{\mu S}-1-\mu S$ is a polynomial, and write it out for the two-component example of Section 6.7.

\textbf{Exercise 6.7.} Show from Lemma 6.3 that $\{\gamma^c,S^{ab}\}$ is totally antisymmetric in $(c,a,b)$ when the three indices are distinct, and compute $\{\gamma^0,S^{12}\}$ and $\{\gamma^1,S^{12}\}$.

\textbf{Exercise 6.8.} Take the polar-plane vielbein of Section 4.10 ($e_\mu{}^a=\mathrm{diag}(1,r)$, $\sqrt{|g|}=r$, $\gamma^r=\sigma_1$, $\gamma^\varphi=\sigma_2/r$). Compute $\frac1{2\sqrt{|g|}}\partial_\mu(\sqrt{|g|}\gamma^\mu)$ and compare with the value $\gamma^\mu\Omega_\mu=\frac1{2r}\sigma_1$ found in Section 4.12.

\textbf{Exercise 6.9.} Show that the Euler--Lagrange expressions of the flat-space Lagrangian $\mathcal L=\tfrac12\bigl(q^TA\dot q-\dot q^TAq\bigr)$ of two real commuting functions $q(t)$ with $A=\begin{pmatrix}0&1\\-1&0\end{pmatrix}$ are those of $q^TA\dot q$, and compute them. (This is the one-dimensional shadow of Theorem 6.5: for real commuting fields the symmetrization is automatic.)

\textbf{Exercise 6.10.} Show $\mathcal L_{m,\lambda}[\gamma^8\Psi]+\mathcal L_{-m,\lambda}[\Psi]=-\lambda\sqrt{|g|}\,S^2$ for $U=\tfrac\lambda2S^2$, and find for which potentials $U$ the map sends $\mathcal L_{m,U}$ to $-\mathcal L_{-m,U}$.

\textbf{Exercise 6.11.} Check $\exp(\pi S^{45})=\gamma^4\gamma^5$ and conclude that the product of the four rotations by $\pi$ of Section 6.11 is $\gamma^8$.

\textbf{Exercise 6.12.} Use Theorem 6.8 to find the image of $\mathcal L_{m,\lambda}$ under the reflection $\Psi'(x)=\gamma^b\Psi(R_bx)$ for $b=1$ and for $b=5$, and then under the composition of the reflections for $b=1$ and $b=2$.

\textbf{Exercise 6.13.} For Grassmann components show that $C_8$ followed by the reflection with $b=0$ is an exact symmetry, using Theorems 6.8 and 6.10.

\textbf{Exercise 6.14.} Let $\Psi=\rho+i\iota$ with real Grassmann columns $\rho,\iota$, and let $A$ be a constant real antisymmetric matrix (for example $C\gamma^a$). Show that $\Psi^\dagger A\,\partial_a\Psi-(\partial_a\Psi^\dagger)A\Psi=2i\bigl(\rho^TA\,\partial_a\iota-\iota^TA\,\partial_a\rho\bigr)$, and write the right-hand side as a total derivative plus a multiple of $\rho^TA\,\partial_a\iota$.

\subsection{Answers to the exercises}

\textbf{Answer 6.1.} With $(Cu)_0=-u_4$, $(Cu)_4=-u_0$, $(Cu)_8=u_{12}$, $(Cu)_{12}=u_8$: for $\Psi=e_0+2e_4+e_8$, $S=1\cdot(-2)+2\cdot(-1)+1\cdot0=-4$ (the component 8 pairs with 12, which is zero). For $\Psi=e_0+(1+i)e_4$, $S=\Psi_0^\ast(-\Psi_4)+\Psi_4^\ast(-\Psi_0)=-(1+i)-(1-i)=-2$. For $\Psi=e_0+(1-i)e_4$, $S=-(1-i)-(1+i)=-2$: the same, as the symmetry $C_0$ of Section 6.12 requires for a commuting field.

\textbf{Answer 6.2.} By Section 2.12, $B$ is Hermitian. Lemma 6.1 gives $(\Psi^\dagger B\Psi)^\ast=\Psi^\dagger B^\dagger\Psi=\Psi^\dagger B\Psi$, so the number is real.

\textbf{Answer 6.3.} $R^T=R$ (Section 6.3). Write $c=\cosh\tfrac\theta2$, $s=\sinh\tfrac\theta2$ and $J=\gamma^0\gamma^4$. Then $JC=\gamma^0\gamma^4C=\gamma^0C\gamma^4=-C\gamma^0\gamma^4=-CJ$, so $RC=C(c-sJ)$ and $R^TCR=RCR=C(c-sJ)(c+sJ)=C(c^2-s^2J^2)=C(c^2-s^2)=C$, because $J^2=1$ and $\cosh^2-\sinh^2=1$.

\textbf{Answer 6.4.} $\omega^2=m^2+k_0^2+\dots+k_3^2-k_5^2-k_6^2-k_7^2$. (a) $4+1-1=4$, $\omega=2$. (b) $4+9-25=-12$: no real frequency; the mode grows like $e^{\sqrt{12}\,x_4}$ (Chapter 8). (c) $4+16-25=-5$: no real frequency either. Only (a) oscillates.

\textbf{Answer 6.5.} $X=f^\ast g'-g^\ast f'=\cos(kx_1)\cdot ik\cos(kx_1)-(-i\sin(kx_1))\cdot(-k\sin(kx_1))=ik\cos^2(kx_1)-ik\sin^2(kx_1)=ik\cos(2kx_1)$. It is purely imaginary, so $K=\tfrac12(X+X^\ast)=0$. And $\bar\Psi\gamma^1\Psi=f^\ast g-g^\ast f=i\cos\sin-(-i)\sin\cos=2i\sin(kx_1)\cos(kx_1)=i\sin(2kx_1)$, whose derivative $2ik\cos(2kx_1)$ equals $X-X^\ast$, as the identity of Section 6.5 requires.

\textbf{Answer 6.6.} $e^{\mu S}-1-\mu S=\sum_{k\ge2}\mu^kS^k/k!$, and $S^k=0$ for $k\ge17$, so the sum stops at $k=16$: a polynomial with real coefficients and $U(0)=U'(0)=0$. For two components $S^3=0$, so $U=\tfrac{\mu^2}2S^2=-\mu^2\,\theta_0^\ast\theta_0\theta_1^\ast\theta_1$.

\textbf{Answer 6.7.} For distinct $c,a,b$, $\{\gamma^c,S^{ab}\}=\gamma^c\gamma^a\gamma^b$, and exchanging two of the three anticommuting factors changes the sign, so the expression is totally antisymmetric. $\{\gamma^0,S^{12}\}=\gamma^0\gamma^1\gamma^2$. $\{\gamma^1,S^{12}\}=0$ because $1\in\{1,2\}$.

\textbf{Answer 6.8.} $\sqrt{|g|}\gamma^r=r\sigma_1$ and $\sqrt{|g|}\gamma^\varphi=\sigma_2$. Only $\partial_r(r\sigma_1)=\sigma_1$ is nonzero ($\sigma_2$ does not depend on $\varphi$). So $\frac1{2r}\sigma_1$, the value of Section 4.12.

\textbf{Answer 6.9.} For commuting numbers $\dot q^TAq=q^TA^T\dot q=-q^TA\dot q$, so $\mathcal L=q^TA\dot q=q_0\dot q_1-q_1\dot q_0$. Then $\partial\mathcal L/\partial q_0=\dot q_1$, $\partial\mathcal L/\partial\dot q_0=-q_1$, so $E_0=\dot q_1+\dot q_1=2\dot q_1$; and $E_1=-\dot q_0-\dot q_0=-2\dot q_0$. These are twice the expressions of Example A of Section 5.6 (whose Lagrangian is half of this one).

\textbf{Answer 6.10.} By Theorem 6.7, $\mathcal L_{m,\lambda}[\gamma^8\Psi]=\sqrt{|g|}(-K-mS-\tfrac\lambda2S^2)$ and $\mathcal L_{-m,\lambda}[\Psi]=\sqrt{|g|}(K+mS-\tfrac\lambda2S^2)$; the sum is $-\lambda\sqrt{|g|}S^2$. In general $\mathcal L_{m,U}[\gamma^8\Psi]=-\mathcal L_{-m,-U}[\Psi]$, which equals $-\mathcal L_{-m,U}[\Psi]$ for every field exactly when $-U=U$, that is $U=0$.

\textbf{Answer 6.11.} $S^{45}=\tfrac12\gamma^4\gamma^5$ and $(\gamma^4\gamma^5)^2=-(\gamma^4)^2(\gamma^5)^2=-(-1)(-1)=-1$, so Section 2.11 gives $\exp(\pi S^{45})=\cos\tfrac\pi2+\gamma^4\gamma^5\sin\tfrac\pi2=\gamma^4\gamma^5$. In the same way $\exp(\pi S^{01})=\gamma^0\gamma^1$, $\exp(\pi S^{23})=\gamma^2\gamma^3$ and $\exp(\pi S^{67})=\gamma^6\gamma^7$ (each square is $-1$). The product in the order $(0,1),(2,3),(4,5),(6,7)$ is $\gamma^0\gamma^1\gamma^2\gamma^3\gamma^4\gamma^5\gamma^6\gamma^7=\gamma^8$.

\textbf{Answer 6.12.} $b=1$ is space-like: $\mathcal L_{-m,\lambda}$. $b=5$ is time-like: $-\mathcal L_{-m,-\lambda}$. Two space-like reflections: the first gives $K\to K$, $S\to-S$, and the second again $K\to K$, $S\to-S$, so altogether $K\to K$, $S\to S$: the composition is an exact symmetry, $\mathcal L_{m,\lambda}$ (it is the rotation by $\pi$ in the plane $(1,2)$, up to a sign of the spinor).

\textbf{Answer 6.13.} For $s=-1$, $C_8$ gives $K\to K$, $S\to-S$ (Theorem 6.10). The reflection with $b=0$ ($n_0=+1$) gives $K\to K$, $S\to-S$ (Theorem 6.8). Together $K\to K$, $S\to S$, and $U(S)$ is unchanged: $\mathcal L_{m,\lambda}\to\mathcal L_{m,\lambda}$, an exact symmetry for every $m$ and $\lambda$.

\textbf{Answer 6.14.} The conjugate of $\rho+i\iota$ is $\rho-i\iota$ (real generators are unchanged, and $i\to-i$), so $\Psi^\dagger=\rho^T-i\iota^T$. Expand both products:

\[
\begin{aligned}
\Psi^\dagger A\,\partial_a\Psi&=\rho^TA\,\partial_a\rho+\iota^TA\,\partial_a\iota+i\bigl(\rho^TA\,\partial_a\iota-\iota^TA\,\partial_a\rho\bigr),\\
(\partial_a\Psi^\dagger)A\Psi&=(\partial_a\rho)^TA\rho+(\partial_a\iota)^TA\iota+i\bigl((\partial_a\rho)^TA\iota-(\partial_a\iota)^TA\rho\bigr).
\end{aligned}
\]

For odd columns $\phi,\chi$ and antisymmetric $A$, exchanging the two odd factors gives $\phi^TA\chi=\sum_{i,j}\phi_iA_{ij}\chi_j=-\sum_{i,j}\chi_jA_{ij}\phi_i=-\chi^TA^T\phi=\chi^TA\phi$. With it, $(\partial_a\rho)^TA\rho=\rho^TA\,\partial_a\rho$, $(\partial_a\iota)^TA\iota=\iota^TA\,\partial_a\iota$, $(\partial_a\rho)^TA\iota=\iota^TA\,\partial_a\rho$ and $(\partial_a\iota)^TA\rho=\rho^TA\,\partial_a\iota$. So in the difference the terms without $i$ cancel exactly, and the difference is $i(\rho^TA\,\partial_a\iota-\iota^TA\,\partial_a\rho)-i(\iota^TA\,\partial_a\rho-\rho^TA\,\partial_a\iota)=2i\bigl(\rho^TA\,\partial_a\iota-\iota^TA\,\partial_a\rho\bigr)$. The product rule and the same exchange give $\iota^TA\,\partial_a\rho=\partial_a(\iota^TA\rho)-(\partial_a\iota)^TA\rho=\partial_a(\iota^TA\rho)-\rho^TA\,\partial_a\iota$, so the difference is $4i\,\rho^TA\,\partial_a\iota-2i\,\partial_a(\iota^TA\rho)$: a total derivative plus a genuine cross term. As for the mass term, the dynamics of the complex Grassmann field lives entirely in the cross terms between its real and imaginary parts.

\section{Field equations, energy--momentum tensor and equations of state}

\subsection{What this chapter does}

Chapter 6 wrote down one Lagrangian for two fields, dirac16complex (complex Grassmann components) and dirac16complex00 (complex commuting components):

\[
\mathcal L=\sqrt{|g|}\,\bigl[K-mS-U(S)\bigr],\qquad K=\tfrac12\bigl(\bar\Psi\gamma^\mu D_\mu\Psi-(D_\mu\bar\Psi)\gamma^\mu\Psi\bigr),\qquad S=\bar\Psi\Psi .
\]

This chapter derives everything that the physics needs from it, in an arbitrary gravitational field and for both fields:

\begin{itemize}
\item the \textbf{field equations}, from the principle of stationary action of Chapter 5 (Sections 7.2 to 7.4);
\item the \textbf{conserved current} and the \textbf{charge}, from the phase symmetry (Section 7.5);
\item the \textbf{energy--momentum tensor}, from the response of the action to a change of the gravitational field, with its symmetry, reality, trace and conservation (Sections 7.6 to 7.8);
\item the \textbf{energy density}, the \textbf{pressures}, two definitions of \textbf{kinetic and potential energy}, and the \textbf{equations of state} (Section 7.9).
\end{itemize}

Then it specializes these results to the situations used later in the book: homogeneous backgrounds (Section 7.10), plane waves in flat space, where the commuting field shows an indefinite charge and an energy without lower bound (Section 7.11), and the notebook's primordial field, including an exact classical dirac16complex00 source of its static member (Section 7.12). Section 7.13 derives how the chirality map $\gamma^8$ of Section 6.11 acts on all of these.

\textbf{Sources.} The Stage-1 document \texttt{provenan\allowbreak{}ce/DIRAC\allowbreak{}16COMPLE\allowbreak{}X\_ARBITR\allowbreak{}ARY\_FIEL\allowbreak{}D.md} (its §8, §9 and §10.9) with its reports in \texttt{artifact\allowbreak{}s/dirac1\allowbreak{}6complex\allowbreak{}/arbitra\allowbreak{}ry-field\allowbreak{}/}; the Stage-5 specification \texttt{handoff/\allowbreak{}specs/ST\allowbreak{}AGE5\_SPE\allowbreak{}C.md} (its §2 and §3); the exact Stage-5 reports on dirac16complex00, \texttt{wolfram-\allowbreak{}dirac16c\allowbreak{}omplex00\allowbreak{}-report.\allowbreak{}json} (46 of 46 checks) and \texttt{python-d\allowbreak{}irac16co\allowbreak{}mplex00-\allowbreak{}report.j\allowbreak{}son} (49 of 49), with \texttt{dirac16c\allowbreak{}omplex00\allowbreak{}-theory.\allowbreak{}json}, all in \texttt{artifact\allowbreak{}s/dirac1\allowbreak{}6complex\allowbreak{}/pair-cr\allowbreak{}eation/}; the pairing reports in the same folder (141 of 141 and 172 of 172 checks); and Theorem M1 of the matter--antimatter document \texttt{provenan\allowbreak{}ce/DIRAC\allowbreak{}16COMPLE\allowbreak{}X\_MATTER\allowbreak{}\_ANTIMAT\allowbreak{}TER.md} (its §4). The prefixes of the check names are those of Section 6.1: \texttt{\detokenize{C00_}} for the Wolfram dirac16complex00 report, \texttt{\detokenize{PAIR_}} for the Wolfram pairing report \texttt{wolfram-\allowbreak{}pairing-\allowbreak{}report.j\allowbreak{}son}, and \texttt{\detokenize{S5_}} for both independent Python reports. The Python twin of a \texttt{\detokenize{C00_}} check is the \texttt{\detokenize{S5_}} check with the same ending in \texttt{python-d\allowbreak{}irac16co\allowbreak{}mplex00-\allowbreak{}report.j\allowbreak{}son}, and the Python twin of a \texttt{\detokenize{PAIR_}} check is the \texttt{\detokenize{S5_}} check with the same ending in \texttt{python-p\allowbreak{}airing-r\allowbreak{}eport.js\allowbreak{}on}; an \texttt{\detokenize{S5_}} name cited on its own in this chapter is a check of \texttt{python-d\allowbreak{}irac16co\allowbreak{}mplex00-\allowbreak{}report.j\allowbreak{}son}. \texttt{\detokenize{MA_}} marks the matter--antimatter reports, and the Stage-1 checks have no prefix of this kind. The Stage-5 prose documents had not been written when this chapter was written; the chapter cites the committed reports.

\textbf{Notation used throughout} (Chapter 6): $\bar\Psi=\Psi^\dagger C$; $\gamma^\mu=e_a{}^\mu\gamma^a$; $D_\mu\Psi=\partial_\mu\Psi+\Omega_\mu\Psi$ and $D_\mu\bar\Psi=\partial_\mu\bar\Psi-\bar\Psi\Omega_\mu$; $\mathcal L_s=\mathcal L/\sqrt{|g|}=K-mS-U(S)$; $M_{\mathrm{eff}}=m+U'(S)$; and the default potential $U=\tfrac\lambda2S^2$, for which $M_{\mathrm{eff}}=m+\lambda S$. The facts (B1) to (B6) of Section 6.3 are used without further comment; the most important are that $C\gamma^\mu$ is real and antisymmetric and the divergence identity $\partial_\mu(\sqrt{|g|}\,\gamma^\mu)=\sqrt{|g|}\,[\gamma^\mu,\Omega_\mu]$.

\subsection{The Euler--Lagrange equations of dirac16complex00}

For commuting components we vary $\Psi_a$ and $\Psi_a^\ast$ as if they were independent, as justified in Section 5.6. We first need the Lagrangian written so that the dependence on $\Psi^\ast$ and on its derivatives is visible. Using $D_\mu\bar\Psi=\partial_\mu\Psi^\dagger C-\Psi^\dagger C\Omega_\mu$,

\[
\frac{\mathcal L}{\sqrt{|g|}}=\tfrac12\,\Psi^\dagger C\gamma^\mu D_\mu\Psi-\tfrac12\,(\partial_\mu\Psi^\dagger)\,C\gamma^\mu\Psi+\tfrac12\,\Psi^\dagger C\Omega_\mu\gamma^\mu\Psi-m\,\Psi^\dagger C\Psi-U(\Psi^\dagger C\Psi).
\]

\textbf{Step 1 (derivative with respect to $\Psi_a^\ast$).} Every term except the second contains $\Psi^\dagger$ undifferentiated, as the left factor of a row--matrix--column product; the derivative of $\Psi^\dagger Y$ with respect to $\Psi_a^\ast$ is the entry $a$ of the column $Y$. For the potential, the chain rule and $\partial S/\partial\Psi_a^\ast=(C\Psi)_a$ give $U'(S)(C\Psi)_a$. So

\[
\frac{\partial\mathcal L}{\partial\Psi_a^\ast}=\sqrt{|g|}\,\Bigl(\tfrac12\,C\gamma^\mu D_\mu\Psi+\tfrac12\,C\Omega_\mu\gamma^\mu\Psi-\bigl(m+U'(S)\bigr)C\Psi\Bigr)_a .
\]

\textbf{Step 2 (derivative with respect to $\partial_\mu\Psi_a^\ast$).} Only the second term contains it:

\[
\frac{\partial\mathcal L}{\partial(\partial_\mu\Psi_a^\ast)}=-\tfrac12\sqrt{|g|}\,\bigl(C\gamma^\mu\Psi\bigr)_a .
\]

\textbf{Step 3 (the Euler--Lagrange expression).} By Section 5.3, $\mathcal E=\partial\mathcal L/\partial\Psi^\ast-\partial_\mu\bigl(\partial\mathcal L/\partial(\partial_\mu\Psi^\ast)\bigr)$, a column of 16 entries:

\[
\mathcal E=\sqrt{|g|}\,C\Bigl(\tfrac12\gamma^\mu D_\mu\Psi+\tfrac12\Omega_\mu\gamma^\mu\Psi-M_{\mathrm{eff}}\Psi\Bigr)+\tfrac12\,\partial_\mu\bigl(\sqrt{|g|}\,C\gamma^\mu\Psi\bigr).
\]

\textbf{Step 4 (the divergence identity).} $C$ is constant, so by the product rule and (B6)

\[
\partial_\mu\bigl(\sqrt{|g|}\,C\gamma^\mu\Psi\bigr)=C\,\partial_\mu\bigl(\sqrt{|g|}\,\gamma^\mu\bigr)\Psi+\sqrt{|g|}\,C\gamma^\mu\partial_\mu\Psi=\sqrt{|g|}\,C\bigl(\gamma^\mu\Omega_\mu\Psi-\Omega_\mu\gamma^\mu\Psi+\gamma^\mu\partial_\mu\Psi\bigr).
\]

\textbf{Step 5 (collect).} Inserting Step 4 into Step 3, the two terms $\pm\tfrac12\Omega_\mu\gamma^\mu\Psi$ cancel, and $\tfrac12\gamma^\mu\Omega_\mu\Psi+\tfrac12\gamma^\mu\partial_\mu\Psi=\tfrac12\gamma^\mu D_\mu\Psi$:

\[
\mathcal E=\sqrt{|g|}\,C\bigl(\gamma^\mu D_\mu\Psi-M_{\mathrm{eff}}\Psi\bigr).
\]

\textbf{Step 6 (the equation).} $\sqrt{|g|}\ne0$ and $C$ is invertible ($C^2=1$), so $\mathcal E=0$ is equivalent to

\[
\gamma^\mu D_\mu\Psi=\bigl(m+U'(S)\bigr)\Psi .
\]

\textbf{The equation for $\bar\Psi$.} Varying $\Psi_b$ in the same way: the terms that contain $\Psi$ undifferentiated give $\partial\mathcal L/\partial\Psi_b=\sqrt{|g|}\bigl(\tfrac12\bar\Psi\gamma^\mu\Omega_\mu-\tfrac12(D_\mu\bar\Psi)\gamma^\mu-M_{\mathrm{eff}}\bar\Psi\bigr)_b$, and $\partial\mathcal L/\partial(\partial_\mu\Psi_b)=\tfrac12\sqrt{|g|}\,(\bar\Psi\gamma^\mu)_b$. The Euler--Lagrange expression is

\[
\sqrt{|g|}\Bigl(\tfrac12\bar\Psi\gamma^\mu\Omega_\mu-\tfrac12(D_\mu\bar\Psi)\gamma^\mu-M_{\mathrm{eff}}\bar\Psi\Bigr)-\tfrac12(\partial_\mu\bar\Psi)\sqrt{|g|}\gamma^\mu-\tfrac12\bar\Psi\,\partial_\mu\bigl(\sqrt{|g|}\gamma^\mu\bigr)=-\sqrt{|g|}\Bigl((D_\mu\bar\Psi)\gamma^\mu+M_{\mathrm{eff}}\bar\Psi\Bigr),
\]

where the last step used (B6) once more. So the second field equation is

\[
(D_\mu\bar\Psi)\gamma^\mu=-\bigl(m+U'(S)\bigr)\bar\Psi .
\]

The Stage-5 verifiers derive both Euler--Lagrange expressions for commuting components in two independent ways (directly and from the momenta $\partial\mathcal L/\partial(\partial_\mu\Psi^\ast)=-\tfrac12\sqrt{|g|}\,C\gamma^\mu\Psi$), in flat space, at the three points of the general test geometry G1 and at three points of the primordial field G2, with $\lambda\ne0$ and also with an abstract smooth $U$ (checks \texttt{C00\_EL\_c\allowbreak{}ommuting\allowbreak{}Flat}, \texttt{C00\_EL\_c\allowbreak{}ommuting\allowbreak{}Curved}, \texttt{C00\_EL\_c\allowbreak{}ommuting\allowbreak{}GeneralS\allowbreak{}moothU} and their \texttt{\detokenize{S5_}} twins; the Lagrangian has 1304 terms at each G1 point).

\subsection{The same equations for dirac16complex}

For Grassmann components the derivatives must be taken with the rules of Section 5.10: the left derivative for $\Psi^\ast$ and the right derivative for $\Psi$. In every term of $\mathcal L$ the factor $\Psi^\dagger$ (or $\partial_\mu\Psi^\dagger$) already stands on the far left and the factor $\Psi$ (or $\partial_\mu\Psi$) on the far right, so moving the varied generator to the left, respectively to the right, passes no other odd factor and costs no sign. For the potential, $S$ is even and commutes with everything, so the left derivative of $S^k$ with respect to $\Psi_a^\ast$ is $kS^{k-1}(C\Psi)_a$, and hence that of $U(S)$ is $U'(S)(C\Psi)_a$, exactly as for commuting numbers. Steps 1 to 6 of Section 7.2 therefore go through word for word, and the two field equations are the same:

\[
\gamma^\mu D_\mu\Psi=\bigl(m+U'(S)\bigr)\Psi,\qquad (D_\mu\bar\Psi)\gamma^\mu=-\bigl(m+U'(S)\bigr)\bar\Psi .
\]

(If the derivative with respect to $\Psi$ is taken from the left instead, the second expression changes its overall sign, by Proposition 5.3; the equation is the same.)

Stage 1 verified these expressions in a genuine Grassmann algebra at the G1 and G2 points with $\lambda\ne0$. Stage 5 repeated the computation at the point G1 p1 in its own Grassmann algebra, where the Lagrangian has 1288 monomials, and recorded that the two statistics give identical equations. The checks are

\begin{Verbatim}[fontsize=\small]
Stage 1:  LAG_eulerLagrangePsibar_G1   LAG_eulerLagrangePsibar_G2
          LAG_eulerLagrangePsi_G1      LAG_eulerLagrangePsi_G2
          GR_complexQuarticEL          GR_complexPsiEquation
Stage 5:  S5_EL_grassmannCurved
          C00_EL_identicalFormBothStatistics   S5_EL_identicalFormBothStatistics
\end{Verbatim}

\subsection{What the field equations say}

\textbf{The two equations are one.} The second equation is the conjugate of the first. Take the conjugate transpose of $\gamma^\mu D_\mu\Psi$ and multiply by $C$ on the right: $(\gamma^\mu D_\mu\Psi)^\dagger C=(D_\mu\Psi)^\dagger(\gamma^\mu)^TC=-(D_\mu\Psi)^\dagger C\gamma^\mu=-(D_\mu\bar\Psi)\gamma^\mu$, using $(\gamma^\mu)^TC=-C\gamma^\mu$ and $(D_\mu\Psi)^\dagger C=D_\mu\bar\Psi$. And $(M_{\mathrm{eff}}\Psi)^\dagger C=M_{\mathrm{eff}}\bar\Psi$ because $M_{\mathrm{eff}}$ is real. So the conjugate of the first equation is the second. For commuting components this means that one complex equation of 16 components is the whole content.

\textbf{Written out} (Stage-1 document, §8.3): with $\Omega_\mu=\tfrac18\omega_{\mu bc}[\gamma^b,\gamma^c]$,

\[
\sum_{\mu=0}^7\sum_{a=0}^7\sum_{k=0}^{15}e_a{}^\mu(\gamma^a)_{jk}\Bigl(\partial_\mu\Psi_k+\sum_{l=0}^{15}(\Omega_\mu)_{kl}\Psi_l\Bigr)=\bigl(m+U'(S)\bigr)\Psi_j\qquad(j=0,\dots,15).
\]

\textbf{First order in time.} Suppose the frame's time direction is aligned with $x_4$, $\gamma^{x_4}=\gamma^4$ (this holds in the backgrounds that the book uses for physics: the homogeneous frames, the primordial field and its static member). Multiplying the equation by $-\gamma^4$ and using $(\gamma^4)^2=-1$ gives

\[
D_4\Psi=-\gamma^4\Bigl(M_{\mathrm{eff}}\Psi-\sum_{\mu\ne4}\gamma^\mu D_\mu\Psi\Bigr):
\]

the time derivative of $\Psi$ is fixed by $\Psi$ and its derivatives along the slice. As for the first-order oscillators of Section 5.6, the initial values of $\Psi$ alone determine the evolution; no initial velocities are needed.

\textbf{The on-shell value of $K$ and of $\mathcal L_s$.} Multiply the first equation by $\bar\Psi$ from the left: $\bar\Psi\gamma^\mu D_\mu\Psi=M_{\mathrm{eff}}S$. Multiply the second by $\Psi$ from the right: $(D_\mu\bar\Psi)\gamma^\mu\Psi=-M_{\mathrm{eff}}S$. (For Grassmann components $M_{\mathrm{eff}}=m+U'(S)$ is even and commutes with every component, so it may be moved out.) Hence on every solution

\[
K=M_{\mathrm{eff}}\,S,\qquad \mathcal L_s=K-mS-U=S\,U'(S)-U(S),
\]

and for the default potential $\mathcal L_s=\tfrac\lambda2S^2$ on shell (Stage-1 document, §9.2; checks \texttt{EMT\_onsh\allowbreak{}ellLagra\allowbreak{}ngianSUp\allowbreak{}rimeMinu\allowbreak{}sU\_G1} and \texttt{\detokenize{_G2}}).

\textbf{The mass shell in a curved field.} Let $U=0$. Apply the operator $\gamma^\nu D_\nu$ to both sides of $\gamma^\mu D_\mu\Psi=m\Psi$; since $m$ is a constant, $(\gamma^\nu D_\nu)^2\Psi=m\,\gamma^\nu D_\nu\Psi=m^2\Psi$. The Lichnerowicz formula of Section 4.13 rewrites the left side, so every solution satisfies the second-order equation

\[
g^{\mu\nu}\bigl(D_\mu D_\nu\Psi-\Gamma^\lambda{}_{\mu\nu}D_\lambda\Psi\bigr)-\tfrac14R\,\Psi=m^2\Psi .
\]

In flat space ($R=0$, $\Omega=0$, $\Gamma=0$) each component obeys $\eta^{ab}\partial_a\partial_b\Psi_k=m^2\Psi_k$, the ultrahyperbolic equation of Section 5.5, which gives the dispersion relation of Section 6.6. In a curved field the scalar curvature appears explicitly. The Stage-5 verifier confirms the second-order form on shell at seven test points with $R\ne0$ (measurement \texttt{onShellS\allowbreak{}econdOrd\allowbreak{}er} of check \texttt{C00\_conn\allowbreak{}ection\_n\allowbreak{}onTrivia\allowbreak{}lCouplin\allowbreak{}g}).

\textbf{The connection term.} The equation reads $\gamma^\mu\partial_\mu\Psi+(\gamma^\mu\Omega_\mu)\Psi=M_{\mathrm{eff}}\Psi$. Section 6.8 showed that for a diagonal vielbein $\gamma^\mu\Omega_\mu=\frac1{2\sqrt{|g|}}\partial_\mu(\sqrt{|g|}\,\gamma^\mu)$. Two cases used later:

\begin{itemize}
\item a homogeneous diagonal background, $ds^2=-dx_4^2+\sum_{i\ne4}\eta_{ii}h_i(x_4)^2dx_i^2$, with $\sqrt{|g|}=V:=\prod_{i\ne4}h_i$, the curved gammas $\gamma^{x_4}=\gamma^4$ and $\gamma^{x_i}=\gamma^i/h_i(x_4)$. Nothing depends on the coordinates $x_i$ with $i\ne4$, so in $\partial_\mu(\sqrt{|g|}\,\gamma^\mu)$ the terms with $\mu=i$ vanish, and the term with $\mu=4$ is $\partial_4(V\gamma^4)=(\partial_4V)\,\gamma^4$ because $\gamma^4$ is a constant matrix. Hence $\gamma^\mu\Omega_\mu=\frac{\partial_4V}{2V}\gamma^4=\tfrac12\Theta\,\gamma^4$ with the total expansion rate $\Theta=\partial_4\ln V=\sum_{i\ne4}H_i$, $H_i=\partial_4h_i/h_i$ (the last equality is the derivative of the logarithm of a product; Chapter 11 uses the same formula);
\item the notebook's primordial field: $\gamma^\mu\Omega_\mu=3H\gamma^0$, independent of the free function $a_4$ (Sections 4.12 and 9.13).
\end{itemize}

\textbf{Worked example (flat space, at rest).} Let $U=0$, $m>0$, and look for solutions that depend on $x_4$ only, $\Psi=e^{-i\omega x_4}u$. The equation is $\gamma^4(-i\omega)u=mu$, that is $-i\gamma^4u=(m/\omega)u$. The matrix $-i\gamma^4$ is Hermitian (Section 8.9) with square $1$, so its eigenvalues are $\pm1$: the solutions have $\omega=+m$ with $-i\gamma^4u=u$, or $\omega=-m$ with $-i\gamma^4u=-u$. The column $u_+=\tfrac12(e_0-e_4-ie_9-ie_{13})$ of Section 8.8 is of the first kind. Check with the tables of Section 2.8: $(\gamma^4u_+)_0=-(u_+)_{13}=\tfrac i2$, so $(-i\gamma^4u_+)_0=\tfrac12=(u_+)_0$, and in the same way for the components 4, 9 and 13.

\subsection{The conserved current and the charge}

\textbf{Theorem 7.1 (the current identity).} For both kinds of components, every potential and every vielbein, and for every field configuration (not only solutions), with $E:=\gamma^\mu D_\mu\Psi-M_{\mathrm{eff}}\Psi$ and $\bar E:=(D_\mu\bar\Psi)\gamma^\mu+M_{\mathrm{eff}}\bar\Psi$,

\[
\partial_\mu\bigl(\sqrt{|g|}\,\bar\Psi\gamma^\mu\Psi\bigr)=\sqrt{|g|}\,\bigl((D_\mu\bar\Psi)\gamma^\mu\Psi+\bar\Psi\gamma^\mu D_\mu\Psi\bigr)=\sqrt{|g|}\,\bigl(\bar E\,\Psi+\bar\Psi\,E\bigr).
\]

\emph{Proof.} The product rule gives $\partial_\mu(\sqrt{|g|}\,\bar\Psi\gamma^\mu\Psi)=\sqrt{|g|}(\partial_\mu\bar\Psi)\gamma^\mu\Psi+\bar\Psi\,\partial_\mu(\sqrt{|g|}\gamma^\mu)\,\Psi+\sqrt{|g|}\,\bar\Psi\gamma^\mu\partial_\mu\Psi$. By (B6) the middle term is $\sqrt{|g|}\,\bar\Psi(\gamma^\mu\Omega_\mu-\Omega_\mu\gamma^\mu)\Psi$. Collecting, $(\partial_\mu\bar\Psi-\bar\Psi\Omega_\mu)\gamma^\mu\Psi+\bar\Psi\gamma^\mu(\partial_\mu\Psi+\Omega_\mu\Psi)$, which is the first equality. For the second, $\bar E\Psi+\bar\Psi E=(D_\mu\bar\Psi)\gamma^\mu\Psi+M_{\mathrm{eff}}S+\bar\Psi\gamma^\mu D_\mu\Psi-M_{\mathrm{eff}}S$ ($M_{\mathrm{eff}}$ commutes with the components). No factor was reordered, so the proof holds for both statistics. $\square$

(Theorem M1 of the matter--antimatter document, §4.1, item 3; checks \texttt{MA\_M1\_no\allowbreak{}etherIde\allowbreak{}ntity\_co\allowbreak{}mmuting\_\allowbreak{}G1} and \texttt{MA\_M1\_no\allowbreak{}etherIde\allowbreak{}ntity\_gr\allowbreak{}assmann\_\allowbreak{}G1}, the latter with 1088 monomials at each G1 point; and, as a negative control, the identity fails with the notebook's contraction of the spin connection, because the divergence identity fails for it, check \texttt{MA\_M1\_ne\allowbreak{}gativeCo\allowbreak{}ntrolNot\allowbreak{}ebookCon\allowbreak{}nection}.) The first equality is also the identity used in Section 6.5 to show that the unsymmetrized kinetic term differs from $K$ by an imaginary total derivative.

\textbf{The Hermitian current.} The vector $\bar\Psi\gamma^\mu\Psi=\Psi^\dagger(C\gamma^\mu)\Psi$ has an anti-Hermitian matrix, so for commuting components it is purely imaginary (Lemma 6.1). The project therefore uses

\[
J^\mu:=-i\,\bar\Psi\gamma^\mu\Psi,
\]

whose matrices $-iC\gamma^\mu$ are Hermitian: $(-iC\gamma^\mu)^\dagger=i(C\gamma^\mu)^T=-iC\gamma^\mu$. $J^\mu$ is real for commuting components and Hermitian for Grassmann components (checks \texttt{QNT\_curr\allowbreak{}entHermi\allowbreak{}ticity}, \texttt{GR\_curre\allowbreak{}ntHermit\allowbreak{}ian}; \texttt{C00\_curr\allowbreak{}ent\_cons\allowbreak{}ervation\allowbreak{}AndReali\allowbreak{}ty}, which records that $\bar\Psi\gamma^\mu\Psi$ is purely imaginary for commuting fields).

\textbf{Conservation.} On every solution $E=\bar E=0$, so Theorem 7.1 gives $\partial_\mu(\sqrt{|g|}\,J^\mu)=0$, which is $\nabla_\mu J^\mu=0$ because $\nabla_\mu J^\mu=\partial_\mu J^\mu+\Gamma^\mu{}_{\mu\lambda}J^\lambda$ and $\Gamma^\mu{}_{\mu\lambda}=\partial_\lambda\ln\sqrt{|g|}$ (Section 5.8). The current is the Noether current of the phase symmetry of Section 6.9: with $\Delta\Psi=i\Psi$ and $\Delta\Psi^\ast=-i\Psi^\ast$, the formula of Theorem 5.2 and the momenta of Section 7.2 give $\tfrac12\sqrt{|g|}\bar\Psi\gamma^\mu(i\Psi)+(-i\Psi^\dagger)\bigl(-\tfrac12\sqrt{|g|}C\gamma^\mu\Psi\bigr)=i\sqrt{|g|}\,\bar\Psi\gamma^\mu\Psi=-\sqrt{|g|}\,J^\mu$: the same current up to the constant factor $-1$, whose sign is a convention (Section 5.7). The Stage-5 verifiers find $\partial_\mu(\sqrt{|g|}J^\mu)=0$ on shell and $\ne0$ off shell at all test points, also for an arbitrary smooth $U$ (checks \texttt{C00\_curr\allowbreak{}ent\_cons\allowbreak{}ervation\allowbreak{}AndReali\allowbreak{}ty}, \texttt{C00\_EMT\_\allowbreak{}generalS\allowbreak{}moothU}, and their \texttt{\detokenize{S5_}} twins).

\textbf{The charge.} For a slice $x_4=\text{const}$ define

\[
Q(x_4)=\int\sqrt{|g|}\,J^{x_4}\,d^7x .
\]

The argument is the one of Section 5.7, with $\sqrt{|g|}\,J^\mu$ in place of $j^\mu$. Differentiate under the integral sign and use the conservation law $\partial_4(\sqrt{|g|}J^{x_4})=-\sum_{j\ne4}\partial_j(\sqrt{|g|}J^{x_j})$:

\[
\frac{dQ}{dx_4}=\int\partial_4\bigl(\sqrt{|g|}\,J^{x_4}\bigr)\,d^7x=-\sum_{j\ne4}\int\partial_j\bigl(\sqrt{|g|}\,J^{x_j}\bigr)\,d^7x=0 .
\]

Each of the seven integrals on the right vanishes by the divergence theorem of Section 5.3, applied on the seven-dimensional slice (do the $x_j$ integral first, by the fundamental theorem of calculus), provided the current vanishes far away; if instead the slice is closed up periodically, the two end values in that integral are equal and cancel. So $Q$ does not change with $x_4$. In a frame with $\gamma^{x_4}=\gamma^4$, the density is

\[
J^4=-i\,\Psi^\dagger C\gamma^4\Psi=\Psi^\dagger B\Psi,\qquad B=-iC\gamma^4 ,
\]

the matrix of Section 2.12, Hermitian with eight eigenvalues $+1$ and eight $-1$.

\textbf{The charge of dirac16complex00 has no sign.} For commuting components $J^4$ is an ordinary number that can be positive, negative or zero. Two exact examples from the report: for $\Psi=i\,e_6+e_{15}$, $J^4=+2=\Psi^\dagger\Psi$; for $\Psi=-i\,e_6+e_{15}$, $J^4=-2=-\Psi^\dagger\Psi$ (check \texttt{C00\_char\allowbreak{}ge\_indef\allowbreak{}inite}). With the table of $B$ in Section 2.9, $(Bu)_6=i\,u_{15}$ and $(Bu)_{15}=-i\,u_6$, so for $\Psi=c_6e_6+c_{15}e_{15}$ one gets $J^4=c_6^\ast\,i\,c_{15}-c_{15}^\ast\,i\,c_6$; for $c_6=i$, $c_{15}=1$ this is $(-i)(i)-(1)(i)(i)=1+1=2$. This is not an accident of the choice $B$: Section 8.8 shows that every candidate charge density built from a Spin(4,4)-invariant form has as many negative as positive eigenvalues. So the reading of Dirac's 1928 theory, in which the time component of the conserved current is a probability density, \textbf{fails for dirac16complex00 in signature (4,4)} (\texttt{dirac16c\allowbreak{}omplex00\allowbreak{}-theory.\allowbreak{}json}, key \texttt{commutin\allowbreak{}gFieldSp\allowbreak{}ecifics.\allowbreak{}chargeIn\allowbreak{}definite}). For the quantized dirac16complex the same indefinite $B$ is the matrix of the canonical anticommutator (Section 8.6). Chapter 8 then restricts to the \textbf{good sector}, the fields that do not depend on the three extra times $x_5,x_6,x_7$ (Section 8.10). There the field is expanded in operators $b$ that annihilate particles and $d$ that annihilate antiparticles, with $b^\ast$, $d^\ast$ creating them, and after \textbf{normal ordering} (the prescription that removes the constant contribution of the filled sea of negative-energy states, the "Dirac sea") the charge becomes $Q=\sum(b^\ast b-d^\ast d)$, summed over all modes: each particle carries $+1$ and each antiparticle $-1$ (Section 8.12). This paragraph is a preview; Chapter 8 derives it.

\subsection{The energy--momentum tensor: what is varied}

\textbf{The definition.} Section 5.8 defined the energy--momentum tensor by the response of the action to a small change of the metric,

\[
T_{\mu\nu}:=-\frac{2}{\sqrt{|g|}}\,\frac{\delta\mathcal S}{\delta g^{\mu\nu}},
\]

with the sign chosen so that $T_{44}$ is the energy density; here $\mathcal S=\int\mathcal L\,d^8x$ is the action (we write $\mathcal S$ to keep the letter $S$ for the scalar density $\bar\Psi\Psi$). It is the source in Einstein's equations (Chapter 4). For a spinor field there is a difficulty: the action depends on the vielbein, not only on the metric, and the spinor components are measured in the frame of the vielbein. We therefore change the vielbein, in a way that changes the metric by the prescribed amount and does not turn the frame.

\textbf{Every change of the vielbein.} Any small change of the vielbein can be written as

\[
\delta e_\mu{}^c=X^c{}_d(x)\,e_\mu{}^d,\qquad X^c{}_d:=\delta e_\mu{}^c\,e_d{}^\mu ,
\]

with a matrix $X(x)$ of small numbers; we lower its first index with $\eta$, $X_{cd}:=\eta_{ce}X^e{}_d$. The metric $g_{\mu\nu}=e_\mu{}^c\eta_{cd}e_\nu{}^d$ changes by

\[
\delta g_{\mu\nu}=X^c{}_{c'}e_\mu{}^{c'}\eta_{cd}e_\nu{}^d+e_\mu{}^c\eta_{cd}X^d{}_{d'}e_\nu{}^{d'}=e_\mu{}^c\,e_\nu{}^d\,\bigl(X_{cd}+X_{dc}\bigr).
\]

Only the symmetric part of $X$ changes the metric. The antisymmetric part is an infinitesimal frame rotation: a matrix $\Lambda=1+X$ satisfies $\Lambda^T\eta\Lambda=\eta$ to first order exactly when $\eta X+(\eta X)^T=0$ (Section 4.10). For a prescribed $\delta g_{\mu\nu}$ there is exactly one \textbf{symmetric} $X$, namely $X_{cd}=\tfrac12e_c{}^\mu e_d{}^\nu\,\delta g_{\mu\nu}$. The energy--momentum tensor is computed with this symmetric change and with the spinor components held fixed. (This is the rule $\delta e_a{}^\mu=\tfrac12e_{a\nu}\,\delta g^{\mu\nu}$ of the Stage-1 document, §9.1, written for the vielbein instead of its inverse.)

\textbf{What has to be computed.} From $\delta g^{\mu\nu}=-g^{\mu\alpha}g^{\nu\beta}\delta g_{\alpha\beta}$ (vary $g^{\mu\alpha}g_{\alpha\nu}=\delta^\mu{}_\nu$ as in Section 5.8), the definition is the same as $\delta\mathcal S=\tfrac12\int\sqrt{|g|}\,T^{\mu\nu}\delta g_{\mu\nu}\,d^8x$. Inserting $\delta g_{\mu\nu}=2e_\mu{}^ce_\nu{}^dX_{cd}$ for symmetric $X$,

\[
\delta\mathcal S=\int\sqrt{|g|}\;T^{cd}\,X_{cd}\;d^8x,\qquad T^{cd}:=e_\mu{}^c\,e_\nu{}^d\,T^{\mu\nu}\ \ (\text{the frame components}).
\]

So we must compute the first-order change of $\sqrt{|g|}\,\mathcal L_s$ when the vielbein changes by a symmetric $X$ and $\Psi$ is held fixed, and read off the coefficient of $X_{cd}$.

\subsection{The energy--momentum tensor: the derivation}

We write $\partial_a:=e_a{}^\mu\partial_\mu$ and $D_a:=e_a{}^\mu D_\mu$ for derivatives along the frame directions, raise frame indices with $\eta$ ($D^a=\eta^{ab}D_b$, $\gamma_a=\eta_{ab}\gamma^b$), and put $e_{\nu a}:=\eta_{ad}\,e_\nu{}^d$.

\textbf{Step 1 (volume and inverse vielbein).} $\sqrt{|g|}=|\det e|$ (Section 4.10). Jacobi's formula of Section 5.8 gives $\delta\ln|\det e|=e_a{}^\mu\,\delta e_\mu{}^a=e_a{}^\mu X^a{}_de_\mu{}^d=X^a{}_a$, so

\[
\delta\sqrt{|g|}=\sqrt{|g|}\,X^a{}_a=\sqrt{|g|}\,\eta^{cd}X_{cd}.
\]

Varying $e_c{}^\mu e_\mu{}^b=\delta_c{}^b$ gives $\delta e_c{}^\mu=-X^d{}_c\,e_d{}^\mu$.

\textbf{Step 2 (what depends on the vielbein).} Split the kinetic term into its derivative part and its connection part (Theorem 6.4),

\[
K=K_\partial+K_\Omega,\qquad K_\partial=\tfrac12\,e_a{}^\mu\bigl(\bar\Psi\gamma^a\partial_\mu\Psi-\partial_\mu\bar\Psi\,\gamma^a\Psi\bigr),\qquad K_\Omega=\tfrac14\,\omega_{cab}\,A^{cab},
\]

where $\omega_{cab}=e_c{}^\mu\omega_{\mu ab}$ and $A^{cab}:=\bar\Psi\{\gamma^c,S^{ab}\}\Psi$. By Lemma 6.3, $A^{cab}$ vanishes unless $c,a,b$ are distinct and is then $\bar\Psi\gamma^c\gamma^a\gamma^b\Psi$: it is \textbf{totally antisymmetric} in its three indices, and it contains no vielbein. Neither do $S$ and $U(S)$. So

\[
\delta\bigl(\sqrt{|g|}\,\mathcal L_s\bigr)=\sqrt{|g|}\,\bigl(\eta^{cd}X_{cd}\,\mathcal L_s+\delta K_\partial+\delta K_\Omega\bigr).
\]

\textbf{Step 3 (the derivative part).} Only the explicit inverse vielbein changes: $\delta K_\partial=-X^d{}_a\cdot\tfrac12\bigl(\bar\Psi\gamma^a\partial_d\Psi-\partial_d\bar\Psi\,\gamma^a\Psi\bigr)$. Write $X^d{}_a=\eta^{de}X_{ea}$ and use the symmetry of $X$ to replace the coefficient of $X_{ea}$ by its symmetric part:

\[
\delta K_\partial=-\tfrac14\,X_{ea}\Bigl(\bar\Psi\gamma^a\partial^e\Psi+\bar\Psi\gamma^e\partial^a\Psi-\partial^e\bar\Psi\,\gamma^a\Psi-\partial^a\bar\Psi\,\gamma^e\Psi\Bigr).
\]

\textbf{Step 4 (the connection in terms of the vielbein).} Define

\[
F_{cab}:=e_c{}^\mu\,e_b{}^\nu\,\partial_\mu e_{\nu a},
\]

the derivative of the vielbein along the frame direction $c$, with its two indices converted to frame indices $a$ and $b$.

\textbf{Lemma 7.2.} (i) For every totally antisymmetric $A^{cab}$, $A^{cab}\omega_{cab}=-A^{cab}F_{cab}$. (ii) For all $d,c,f$,

\[
\omega_{dcf}=\tfrac12\bigl(-F_{dcf}+F_{dfc}+F_{fcd}+F_{fdc}-F_{cfd}-F_{cdf}\bigr).
\]

\emph{Proof.} Section 4.11 wrote the canonical connection as $\omega_{\mu ab}=e_a{}^\sigma e_b{}^\nu\,\Gamma_{\sigma\mu\nu}-e_b{}^\nu\,\partial_\mu e_{\nu a}$, with $\Gamma_{\sigma\mu\nu}=g_{\sigma\rho}\Gamma^\rho{}_{\mu\nu}$. Contracting with $e_c{}^\mu$,

\[
\omega_{cab}=G_{cab}-F_{cab},\qquad G_{cab}:=e_c{}^\mu e_a{}^\sigma e_b{}^\nu\,\Gamma_{\sigma\mu\nu}.
\]

(i) $\Gamma_{\sigma\mu\nu}$ is symmetric in $\mu,\nu$, so $G_{cab}$ is symmetric in $c,b$, and its contraction with $A^{cab}$, which is antisymmetric in $c,b$, vanishes. (ii) We express $G$ through $F$. From $g_{\nu\sigma}=e_\nu{}^be_{\sigma b}$ and the product rule, $\partial_\mu g_{\nu\sigma}=(\partial_\mu e_{\nu b})e_\sigma{}^b+e_\nu{}^b\,\partial_\mu e_{\sigma b}$. Contracting with $e_x{}^\mu e_y{}^\nu e_z{}^\sigma$ and using $e_z{}^\sigma e_\sigma{}^b=\delta_z{}^b$,

\[
e_x{}^\mu e_y{}^\nu e_z{}^\sigma\,\partial_\mu g_{\nu\sigma}=F_{xzy}+F_{xyz}.
\]

Now $\Gamma_{\sigma\mu\nu}=\tfrac12(\partial_\mu g_{\nu\sigma}+\partial_\nu g_{\mu\sigma}-\partial_\sigma g_{\mu\nu})$ (Section 4.5). Contract with $e_d{}^\mu e_c{}^\sigma e_f{}^\nu$ and apply the last formula to each of the three terms (derivative directions $d$, $f$, $c$ respectively):

\[
G_{dcf}=\tfrac12\bigl(F_{dcf}+F_{dfc}+F_{fcd}+F_{fdc}-F_{cfd}-F_{cdf}\bigr),
\]

and $\omega_{dcf}=G_{dcf}-F_{dcf}$ is the claim. $\square$

\textbf{Corollary 7.3.} If $A^{acf}$ is antisymmetric in $c,f$, then $\omega_{dcf}A^{acf}=-\bigl(F_{dcf}+F_{cfd}+F_{cdf}\bigr)A^{acf}$.

\emph{Proof.} Group the six terms of Lemma 7.2 (ii) as $\tfrac12\bigl[(F_{dfc}-F_{dcf})+(F_{fcd}-F_{cfd})+(F_{fdc}-F_{cdf})\bigr]$. In each pair the first term is the second with $c$ and $f$ exchanged. After contraction with $A^{acf}$ and renaming the summation indices $c\leftrightarrow f$ in the first term, $A^{afc}=-A^{acf}$ turns each pair into twice the negative of its second term. $\square$

\textbf{Step 5 (the connection part).} $K_\Omega=\tfrac14A^{cab}\omega_{cab}=-\tfrac14A^{cab}F_{cab}$ by Lemma 7.2 (i), and $A$ does not change, so $\delta K_\Omega=-\tfrac14A^{cab}\,\delta F_{cab}$. Vary the three factors of $F_{cab}$, using Step 1 and $\delta e_{\nu a}=X_{ad}\,e_\nu{}^d$:

\[
\begin{aligned}
\delta F_{cab}&=\delta e_c{}^\mu\,e_b{}^\nu\,\partial_\mu e_{\nu a}+e_c{}^\mu\,\delta e_b{}^\nu\,\partial_\mu e_{\nu a}+e_c{}^\mu e_b{}^\nu\,\partial_\mu\bigl(X_{ad}\,e_\nu{}^d\bigr)\\
&=-X^d{}_c\,F_{dab}-X^d{}_b\,F_{cad}+\partial_cX_{ab}+X_a{}^d\,F_{cdb}.
\end{aligned}
\]

(In the last term, $X_{ad}\,e_c{}^\mu e_b{}^\nu\partial_\mu e_\nu{}^d=X_{ad}\eta^{de}F_{ceb}$.) The term $\partial_cX_{ab}$, the only one with a derivative of $X$, is symmetric in $a,b$; contracted with $A^{cab}$, which is antisymmetric in $a,b$, it vanishes. In the remaining three terms we rename the summation indices so that $X$ always carries the first index of $A^{acf}$. A summation index is a dummy: it may be given any letter not used elsewhere in the term. We treat the three terms one at a time.

\begin{itemize}
\item First term. Rename $c\to a$, $a\to c$, $b\to f$ (all three at once):
\end{itemize}

\[
-A^{cab}X^d{}_c\,F_{dab}=-A^{acf}X^d{}_a\,F_{dcf} .
\]

\begin{itemize}
\item Second term. Rename $b\to a$, $a\to f$ and keep $c$; then use that $A$ does not change under the cyclic shift of its three indices, $A^{cfa}=A^{acf}$ (a cyclic shift is two exchanges, each of which costs a sign):
\end{itemize}

\[
-A^{cab}X^d{}_b\,F_{cad}=-A^{cfa}X^d{}_a\,F_{cfd}=-A^{acf}X^d{}_a\,F_{cfd} .
\]

\begin{itemize}
\item Third term. Exchange the first two indices of $A$, $A^{cab}=-A^{acb}$; rename $b\to f$; and use $X_a{}^d=\eta^{de}X_{ae}=\eta^{de}X_{ea}=X^d{}_a$ for symmetric $X$:
\end{itemize}

\[
+A^{cab}X_a{}^d\,F_{cdb}=-A^{acb}X_a{}^d\,F_{cdb}=-A^{acf}X^d{}_a\,F_{cdf} .
\]

Adding the three lines and using Corollary 7.3 (with $A^{acf}$ antisymmetric in $c,f$),

\[
A^{cab}\,\delta F_{cab}=-X^d{}_a\,A^{acf}\bigl(F_{dcf}+F_{cfd}+F_{cdf}\bigr)=X^d{}_a\,\omega_{dcf}\,A^{acf} .
\]

Hence

\[
\delta K_\Omega=-\tfrac14\,X^d{}_a\,\omega_{dcf}\,A^{acf}.
\]

Finally we write this with the spinor connection. The frame component $\Omega_d:=e_d{}^\mu\Omega_\mu=\tfrac12\omega_{dcf}S^{cf}$ gives $\bar\Psi\{\gamma^a,\Omega_d\}\Psi=\tfrac12\omega_{dcf}A^{acf}$, so $\delta K_\Omega=-\tfrac12X^d{}_a\bar\Psi\{\gamma^a,\Omega_d\}\Psi$. With $X^d{}_a=\eta^{de}X_{ea}$, $\Omega^e=\eta^{ed}\Omega_d$ and the symmetry of $X$,

\[
\delta K_\Omega=-\tfrac14\,X_{ea}\Bigl(\bar\Psi\gamma^a\Omega^e\Psi+\bar\Psi\Omega^e\gamma^a\Psi+\bar\Psi\gamma^e\Omega^a\Psi+\bar\Psi\Omega^a\gamma^e\Psi\Bigr).
\]

\textbf{Step 6 (collect).} Adding Steps 3 and 5, the connection terms complete the derivatives to covariant ones: $\bar\Psi\gamma^a(\partial^e+\Omega^e)\Psi=\bar\Psi\gamma^aD^e\Psi$ and $-(\partial^e\bar\Psi-\bar\Psi\Omega^e)\gamma^a\Psi=-(D^e\bar\Psi)\gamma^a\Psi$. So

\[
\delta\bigl(\sqrt{|g|}\,\mathcal L_s\bigr)=\sqrt{|g|}\,X_{ea}\Bigl[\eta^{ea}\mathcal L_s-\tfrac14\bigl(\bar\Psi\gamma^aD^e\Psi+\bar\Psi\gamma^eD^a\Psi-(D^e\bar\Psi)\gamma^a\Psi-(D^a\bar\Psi)\gamma^e\Psi\bigr)\Bigr].
\]

No integration by parts was needed: the derivatives of $X$ dropped out in Step 5, so this holds at every point. The bracket is symmetric in $e,a$, so by Section 7.6 it is $T^{ea}$. Converting to coordinate indices with $T_{\mu\nu}=e_{\mu c}e_{\nu d}T^{cd}$, where $e_{\mu c}\gamma^c=\gamma_\mu=g_{\mu\nu}\gamma^\nu$ and $e_{\nu d}D^d=D_\nu$:

\textbf{Theorem 7.4 (the energy--momentum tensor).} For both fields, in every gravitational field,

\[
T_{\mu\nu}=-\tfrac14\Bigl[\bar\Psi\gamma_\mu D_\nu\Psi+\bar\Psi\gamma_\nu D_\mu\Psi-(D_\mu\bar\Psi)\gamma_\nu\Psi-(D_\nu\bar\Psi)\gamma_\mu\Psi\Bigr]+g_{\mu\nu}\,\mathcal L_s .
\]

This is the formula of the Stage-1 document (§9.2) and of \texttt{handoff/\allowbreak{}specs/CO\allowbreak{}NTRACT.m\allowbreak{}d} (§7). For dirac16complex every step above is a statement about bilinears in which $\Psi^\dagger$ stays to the left of $\Psi$, so the same tensor results with Grassmann components; after quantization it is an operator, normal-ordered with respect to the Dirac sea as in Chapter 8. For dirac16complex00 it is an ordinary real symmetric tensor field.

\textbf{How the repository checks it.} Stage 1 carried out the metric variation explicitly for diagonal frames (checks \texttt{EMT\_vari\allowbreak{}ation\_G3} and \texttt{EMT\_vari\allowbreak{}ation}). Stage 5 computed the full functional derivative $\delta\mathcal S/\delta e_\mu{}^a$, including the change of the spin connection, exactly at points of two general non-diagonal frames (G1, and a further integer frame G4) and of a homogeneous frame: its symmetric part equals $\sqrt{|g|}\,T^{\mu\nu}$ of Theorem 7.4 in all 36 independent components \textbf{off shell}, its antisymmetric part (the frame-rotation part) is nonzero off shell and vanishes on shell. The derivative with respect to the derivatives of the vielbein was computed in closed form and compared with an exact linearisation in all 512 directions at each point; the Wolfram verifier also recomputed the functional derivative independently by a generic jet identity, for all 64 entries at the G4 point and for three spot entries at the point G1 p1 (at the homogeneous point it lists no spot entries) (checks \texttt{C00\_EMT\_\allowbreak{}vielbein\allowbreak{}Variatio\allowbreak{}n} and \texttt{S5\_EMT\_v\allowbreak{}ielbeinV\allowbreak{}ariation}; measurement \texttt{jetIdent\allowbreak{}itySpotE\allowbreak{}ntries} of the Wolfram report). The report notes that this lifts a limitation of Stage 1, which had not varied the off-diagonal components explicitly.

\subsection{Properties of the energy--momentum tensor}

\textbf{Symmetric.} $T_{\mu\nu}=T_{\nu\mu}$: the formula is visibly unchanged when $\mu$ and $\nu$ are exchanged.

\textbf{Real (Hermitian).} Group the bracket as $\bigl[\bar\Psi\gamma_\mu D_\nu\Psi-(D_\nu\bar\Psi)\gamma_\mu\Psi\bigr]+\bigl[\bar\Psi\gamma_\nu D_\mu\Psi-(D_\mu\bar\Psi)\gamma_\nu\Psi\bigr]$. Exactly as for $K$ in Theorem 6.2, the conjugate of $\bar\Psi\gamma_\mu D_\nu\Psi=\Psi^\dagger(C\gamma_\mu)D_\nu\Psi$ is $-(D_\nu\bar\Psi)\gamma_\mu\Psi$, because $C\gamma_\mu$ is real and antisymmetric; so each bracket is of the form $Y+Y^\ast$. And $\mathcal L_s$ is real (Theorem 6.2). (Checks \texttt{EMT\_symm\allowbreak{}etricHer\allowbreak{}mitian\_G\allowbreak{}1}, \texttt{EMT\_symm\allowbreak{}etricHer\allowbreak{}mitian\_G\allowbreak{}2}, \texttt{GR\_emtHe\allowbreak{}rmitian}, \texttt{C00\_EMT\_\allowbreak{}symmetri\allowbreak{}cAndReal}, \texttt{S5\_EMT\_s\allowbreak{}ymmetric\allowbreak{}AndReal}.)

\textbf{Trace.} Contract with $g^{\mu\nu}$: $g^{\mu\nu}\gamma_\mu=\gamma^\nu$, so the bracket becomes $2\bar\Psi\gamma^\nu D_\nu\Psi-2(D_\nu\bar\Psi)\gamma^\nu\Psi=4K$, and $g^{\mu\nu}g_{\mu\nu}=8$. Therefore, off shell and on shell,

\[
T^\mu{}_\mu=-K+8\,\mathcal L_s,\qquad\text{and on shell}\qquad T^\mu{}_\mu=-M_{\mathrm{eff}}S+8\bigl(SU'-U\bigr)=-mS+7SU'(S)-8U(S),
\]

which is $-mS+3\lambda S^2$ for the default potential. (Checks \texttt{EMT\_trac\allowbreak{}eOffShel\allowbreak{}lIdentit\allowbreak{}y\_G1}, \texttt{EMT\_trac\allowbreak{}e\_G1}, \texttt{EMT\_trac\allowbreak{}e\_G2} of Stage 1; \texttt{C00\_EMT\_\allowbreak{}traceOnS\allowbreak{}hell} and \texttt{S5\_EMT\_t\allowbreak{}raceOnSh\allowbreak{}ell}, whose recorded on-shell values include $923549/26880$ at the Wolfram point G1 p1; for a general smooth $U$ in \texttt{C00\_EMT\_\allowbreak{}generalS\allowbreak{}moothU}.)

\textbf{Theorem 7.5 (conservation).} On every solution of the field equations, for both fields, $\nabla_\mu T^{\mu\nu}=0$.

The proof uses the fact that the action does not depend on the coordinates. We give it in five steps.

\emph{Step 1 (a change of coordinates changes the fields).} Let $\xi^\mu(x)$ be a small vector field that vanishes near the boundary of the region, and introduce new coordinates $x'^\mu=x^\mu-\xi^\mu(x)$. By the transformation rule of covectors (Section 4.3), the vielbein in the new coordinates is $e'_\mu{}^a(x')=\frac{\partial x^\nu}{\partial x'^\mu}e_\nu{}^a(x)$, and the spinor components, which refer to the frame and not to the coordinates, are unchanged: $\Psi'(x')=\Psi(x)$. To first order $x=x'+\xi(x')$ and $\partial x^\nu/\partial x'^\mu=\delta^\nu{}_\mu+\partial_\mu\xi^\nu$, so, as functions of their argument,

\[
e'_\mu{}^a=e_\mu{}^a+\delta_\xi e_\mu{}^a,\quad \delta_\xi e_\mu{}^a=\xi^\lambda\partial_\lambda e_\mu{}^a+e_\lambda{}^a\,\partial_\mu\xi^\lambda;\qquad \Psi'=\Psi+\delta_\xi\Psi,\quad \delta_\xi\Psi=\xi^\lambda\partial_\lambda\Psi .
\]

The action is the same integral written in other coordinates ($\mathcal L_s$ is a scalar and $\sqrt{|g|}\,d^8x$ is invariant, Section 6.9), and the region is mapped onto itself because $\xi$ vanishes near its boundary. Renaming the integration variable, $\mathcal S[e',\Psi']=\mathcal S[e,\Psi]$, so the first-order change vanishes: $\delta_\xi\mathcal S=0$.

\emph{Step 2 (split the change of the vielbein).} Insert the vielbein postulate of Section 4.11, $\partial_\lambda e_\mu{}^c=\Gamma^\rho{}_{\lambda\mu}e_\rho{}^c-\omega_\lambda{}^c{}_d\,e_\mu{}^d$:

\[
\delta_\xi e_\mu{}^c=e_\rho{}^c\bigl(\partial_\mu\xi^\rho+\Gamma^\rho{}_{\mu\lambda}\xi^\lambda\bigr)-\xi^\lambda\omega_\lambda{}^c{}_d\,e_\mu{}^d=e_\rho{}^c\,\nabla_\mu\xi^\rho-\xi^\lambda\omega_\lambda{}^c{}_d\,e_\mu{}^d .
\]

In the notation of Section 7.6, $X_{cd}=e_c{}^\sigma e_d{}^\mu\,\nabla_\mu\xi_\sigma-\xi^\lambda\omega_{\lambda cd}$ (with $\xi_\sigma=g_{\sigma\rho}\xi^\rho$). Its symmetric part is $X^{\mathrm S}_{cd}=\tfrac12e_c{}^\sigma e_d{}^\mu(\nabla_\mu\xi_\sigma+\nabla_\sigma\xi_\mu)$, and its antisymmetric part $\varepsilon_{cd}$ is the rest (it contains the whole $\omega$ term, because $\omega_{\lambda cd}$ is antisymmetric).

\emph{Step 3 (three changes).} To first order, changes add. So the change $(\delta_\xi e,\delta_\xi\Psi)$ is the sum of: (1) the symmetric change $X^{\mathrm S}$ of the vielbein with $\Psi$ fixed; (2) the frame rotation $\varepsilon$ of the vielbein together with the spin rotation $\delta\Psi=\tfrac12\varepsilon_{cd}S^{cd}\Psi$ that covers it (by Section 4.13, $[\tfrac12\varepsilon_{cd}S^{cd},\gamma^a]=-\varepsilon^a{}_b\gamma^b$, the infinitesimal form of the covering relation); (3) the change $\delta_3\Psi=\xi^\lambda\partial_\lambda\Psi-\tfrac12\varepsilon_{cd}S^{cd}\Psi$ of the field alone. The change (1) contributes $\int\sqrt{|g|}\,T^{cd}X^{\mathrm S}_{cd}\,d^8x=\int\sqrt{|g|}\,T^{\sigma\mu}\nabla_\mu\xi_\sigma\,d^8x$ (Section 7.6, and $T$ is symmetric). The change (2) contributes nothing: it is an infinitesimal local frame rotation, under which $\mathcal L$ is invariant (Section 6.9, differentiated at zero rotation angle). The change (3) is a change of the field alone that vanishes near the boundary; by Sections 5.3 and 5.10 it changes the action by the integral of the Euler--Lagrange expressions times the change, which is zero on a solution.

\emph{Step 4 (the result of Steps 1 to 3).} On every solution, for every small $\xi$ vanishing near the boundary,

\[
0=\int\sqrt{|g|}\;T^{\sigma\mu}\,\nabla_\mu\xi_\sigma\;d^8x .
\]

\emph{Step 5 (integrate by parts).} By the product rule, $T^{\sigma\mu}\nabla_\mu\xi_\sigma=\nabla_\mu(T^{\sigma\mu}\xi_\sigma)-(\nabla_\mu T^{\sigma\mu})\,\xi_\sigma$. For any vector field $V^\mu$, $\sqrt{|g|}\,\nabla_\mu V^\mu=\partial_\mu(\sqrt{|g|}\,V^\mu)$ (the computation used for the current in Section 7.5), so the first term integrates to zero by the divergence theorem. Hence $\int\sqrt{|g|}\,(\nabla_\mu T^{\sigma\mu})\,\xi_\sigma\,d^8x=0$ for every $\xi$, and the fundamental lemma of Section 5.2 (applied to each component) gives $\nabla_\mu T^{\sigma\mu}=0$. $\square$

For Grassmann components the same proof applies: the changes (1) and (2) involve only ordinary functions multiplying bilinears in the fixed order, and the first-order change of the action under the change (3) of the odd field is again the integral of the Euler--Lagrange expressions times the change (Section 5.10). The repository verifies the conservation with exact on-shell field data at the three G1 points and at three points of the primordial field, and checks that the divergence is not zero for data that do not satisfy the field equations (checks \texttt{EMT\_cons\allowbreak{}ervation\allowbreak{}\_G1} and \texttt{EMT\_cons\allowbreak{}ervation\allowbreak{}\_G2} of Stage 1; \texttt{C00\_EMT\_\allowbreak{}conserva\allowbreak{}tionOnSh\allowbreak{}ell} and \texttt{S5\_EMT\_c\allowbreak{}onservat\allowbreak{}ionOnShe\allowbreak{}ll}, including a general smooth $U$).

\subsection{Energy density, pressures, kinetic and potential energy, equations of state}

\textbf{The observer.} From here on the time is \textbf{Gaussian normal}: $g_{44}=-1$ and $g_{4i}=0$ for $i\ne4$, and the frame's time direction is aligned with $x_4$, so $\gamma^{x_4}=\gamma^4$ and $\gamma_4=g_{44}\gamma^{x_4}=-\gamma^4$ (the homogeneous frames, the primordial field and its static member have this form; Section 4.6 defines it). The observer at rest in these coordinates has the 8-velocity $u=\partial_4$, a unit time-like vector.

\textbf{Definitions} (Stage-1 document, §9.4). The energy density, the pressure along each of the seven transverse directions, their mean, and the equation-of-state parameters are

\[
\begin{aligned}
&\rho:=T_{\mu\nu}u^\mu u^\nu=T_{44}=-T^4{}_4,\qquad p_{(i)}:=T^i{}_i\ \ (\text{no sum},\ i\ne4),\\
&\bar p=\tfrac17\sum_{i\ne4}p_{(i)},\qquad w_{(i)}=\frac{p_{(i)}}\rho,\qquad \bar w=\frac{\bar p}\rho .
\end{aligned}
\]

When all seven pressures are equal we write $p$ and $w=p/\rho$. Four of the transverse directions are space-like ($i=0,1,2,3$) and three are time-like ($i=5,6,7$). For the space-like ones, $T^i{}_i$ is the flux of momentum along $i$ through a surface of constant $x_i$, the usual pressure; for a scalar field in the homogeneous case it is $\tfrac12\dot\phi^2-V$ (Section 5.8). For the time-like transverse directions the same formula is used by definition; there is no everyday meaning of a pressure along a time.

\textbf{The energy density.} Use $\gamma_4=-\gamma^4$ and $g_{44}=-1$ in Theorem 7.4 with $\mu=\nu=4$:

\[
\rho=T_{44}=-\tfrac14\bigl[2\bar\Psi\gamma_4D_4\Psi-2(D_4\bar\Psi)\gamma_4\Psi\bigr]-\mathcal L_s=K_4-\mathcal L_s,
\]

with the time part and the transverse part of the kinetic term

\[
K_4:=\tfrac12\bigl(\bar\Psi\gamma^4D_4\Psi-(D_4\bar\Psi)\gamma^4\Psi\bigr),\qquad K_\perp:=\tfrac12\sum_{\mu\ne4}\bigl(\bar\Psi\gamma^\mu D_\mu\Psi-(D_\mu\bar\Psi)\gamma^\mu\Psi\bigr),\qquad K=K_4+K_\perp .
\]

Since $\mathcal L_s=K_4+K_\perp-mS-U$,

\[
\rho=T_{44}=-K_\perp+mS+U(S)\qquad\text{identically (off shell)}.
\]

\textbf{The time part is the oscillator of Section 5.6.} With $C\gamma^4=iB$ (from $B=-iC\gamma^4$), $\bar\Psi\gamma^4=i\Psi^\dagger B$ and $(D_4\bar\Psi)\gamma^4=i(D_4\Psi)^\dagger B$, so

\[
K_4=\tfrac i2\Bigl(\Psi^\dagger B\,D_4\Psi-(D_4\Psi)^\dagger B\,\Psi\Bigr).
\]

With $B$ replaced by $1$ this is exactly the time part $\tfrac i2(\psi^\ast\dot\psi-\dot\psi^\ast\psi)$ of the complex oscillator of Section 5.6. The indefinite matrix $B$ is where signature (4,4) enters.

\textbf{Two splits into kinetic and potential energy} (Stage-1 document, §9.5; \texttt{handoff/\allowbreak{}specs/CO\allowbreak{}NTRACT.m\allowbreak{}d}, §7).

\begin{itemize}
\item \textbf{(A) Lagrangian split:} $\mathrm{KE}_L:=\tfrac12K_4$ and $\mathrm{PE}_L:=\rho-\mathrm{KE}_L$. This imitates the scalar field, whose kinetic energy $\tfrac12\dot\phi^2$ is half of the time-derivative term $\dot\phi\,\partial\mathcal L/\partial\dot\phi$ of its Lagrangian; by definition $\rho=\mathrm{KE}_L+\mathrm{PE}_L$.
\item \textbf{(B) Hamiltonian split:} $\mathrm{KE}_H:=-K_\perp$ (the energy of gradients and of the transverse connection) and $\mathrm{PE}_H:=mS+U(S)$ (rest energy and interaction). Then $\rho=\mathrm{KE}_H+\mathrm{PE}_H$ holds identically, by the formula above. The time derivatives do not appear in split (B): for a first-order field they carry no energy, as for the oscillator of Exercise 5.4.
\end{itemize}

The Stage-5 verifiers confirm $\rho=\mathrm{KE}_H+\mathrm{PE}_H$ and $\rho=K_4-\mathcal L_s$ off shell, with arbitrary field data, at three points of the primordial field (check \texttt{C00\_EMT\_\allowbreak{}observer\allowbreak{}SplitGau\allowbreak{}ssianNor\allowbreak{}mal} and its \texttt{\detokenize{S5_}} twin; measurements \texttt{rhoEqual\allowbreak{}sKEHplus\allowbreak{}PEH\_offS\allowbreak{}hell} and \texttt{rhoEqual\allowbreak{}sK4minus\allowbreak{}Ls\_offSh\allowbreak{}ell}).

\subsection{Homogeneous backgrounds}

\textbf{The setting.} Take a homogeneous diagonal background (a Bianchi-I metric, as in Chapter 11),

\[
ds^2=-dx_4^2+\sum_{i\ne4}\eta_{ii}\,h_i(x_4)^2\,dx_i^2,
\]

with the Hubble rates $H_i=\partial_4h_i/h_i$, the 7-volume $V=\prod_{i\ne4}h_i=\sqrt{|g|}$ and $\Theta=\sum_{i\ne4}H_i=\partial_4\ln V$, and a field that depends on $x_4$ only, $\Psi=\Psi(x_4)$. The vielbein is diagonal with $h_4=1$, and by Section 4.11 the only nonzero connection components are $\omega_{i\,i4}=-\omega_{i\,4i}=\eta_{ii}\,\partial_4h_i$ (the pair $(i,4)$ for $\mu=i$). So $\Omega_4=0$ and $\Omega_i=\omega_{i\,i4}S^{i4}$ (both orders of the pair give equal terms, which cancels the factor $\tfrac12$).

\textbf{The transverse pressures.} For $i\ne4$ the field does not depend on $x_i$, so $D_i\Psi=\Omega_i\Psi$ and $D_i\bar\Psi=-\bar\Psi\Omega_i$. Then

\[
T_{ii}=-\tfrac14\bigl[2\bar\Psi\gamma_i\Omega_i\Psi+2\bar\Psi\Omega_i\gamma_i\Psi\bigr]+g_{ii}\mathcal L_s=-\tfrac12\bar\Psi\{\gamma_i,\Omega_i\}\Psi+g_{ii}\mathcal L_s .
\]

$\gamma_i$ is a multiple of $\gamma^i$, and $\{\gamma^i,S^{i4}\}=0$ by Lemma 6.3 (the index $i$ belongs to the pair). So $T^i{}_i=g^{ii}T_{ii}=\mathcal L_s$ for all seven transverse directions, off shell.

\textbf{The energy density.} $K_\perp=\tfrac12\sum_{i\ne4}\bar\Psi\{\gamma^{x_i},\Omega_i\}\Psi=0$ for the same reason, so $\rho=mS+U$, off shell.

\textbf{On shell} $\mathcal L_s=SU'-U$ and $K=K_4=M_{\mathrm{eff}}S$ (Section 7.4). Collecting:

\[
\begin{aligned}
&\rho=mS+U,\qquad p_{(i)}=p=SU'(S)-U(S)\ \ \text{for all seven }i\ne4,\qquad w=\frac{SU'-U}{mS+U},\\
&\mathrm{KE}_L=\tfrac12S\bigl(m+U'\bigr),\qquad \mathrm{PE}_L=\tfrac12\bigl(mS+2U-SU'\bigr),\qquad \mathrm{KE}_H=0,\qquad \mathrm{PE}_H=\rho,
\end{aligned}
\]

and therefore $\rho+p=2\,\mathrm{KE}_L=S\,M_{\mathrm{eff}}$ and $p=\mathrm{KE}_L-\mathrm{PE}_L$, so $w=(\mathrm{KE}_L-\mathrm{PE}_L)/(\mathrm{KE}_L+\mathrm{PE}_L)$, the same formula as for the scalar field. The pressure is the same in all seven directions although the scale factors $h_i$ differ. For the free field ($U=0$) $p=0$ and $w=0$: dust. For the default potential,

\[
\rho=mS+\tfrac\lambda2S^2,\qquad p=\tfrac\lambda2S^2,\qquad w=\frac{\lambda S}{2m+\lambda S},\qquad \mathrm{KE}_L=\tfrac12S(m+\lambda S),\qquad \mathrm{PE}_L=\tfrac12mS .
\]

\textbf{The off-diagonal components.} For $i\ne j$, both $\ne4$, the same computation gives $T_{ij}=-\tfrac14\bar\Psi\bigl(\{\gamma_i,\Omega_j\}+\{\gamma_j,\Omega_i\}\bigr)\Psi$. With $\gamma_i=\eta_{ii}h_i\gamma^i$, $\Omega_j=\eta_{jj}H_jh_j\cdot\tfrac12\gamma^j\gamma^4$ and $\{\gamma^i,\gamma^j\gamma^4\}=2\gamma^i\gamma^j\gamma^4$ for three distinct indices,

\[
T_{ij}=\tfrac14\,\eta_{ii}\eta_{jj}\,h_ih_j\,(H_i-H_j)\,\bar\Psi\gamma^i\gamma^j\gamma^4\Psi ,
\]

a three-gamma bilinear that vanishes when the two directions expand at the same rate. For $T_{4i}$: $T_{4i}=-\tfrac14\bigl[\bar\Psi\{\gamma_4,\Omega_i\}\Psi+\bar\Psi\gamma_i\partial_4\Psi-\partial_4\bar\Psi\,\gamma_i\Psi\bigr]$, and $\{\gamma^4,S^{i4}\}=0$, so $T_{4i}=-\tfrac14\bigl(\bar\Psi\gamma_i\partial_4\Psi-\partial_4\bar\Psi\,\gamma_i\Psi\bigr)$ off shell. On shell $\gamma^4\partial_4\Psi+\tfrac12\Theta\gamma^4\Psi=M_{\mathrm{eff}}\Psi$ (Section 7.4), that is $\partial_4\Psi=-M_{\mathrm{eff}}\gamma^4\Psi-\tfrac12\Theta\Psi$ and, conjugating, $\partial_4\bar\Psi=M_{\mathrm{eff}}\bar\Psi\gamma^4-\tfrac12\Theta\bar\Psi$. Inserting, the $\Theta$ terms cancel and what remains is $-M_{\mathrm{eff}}\bar\Psi\{\gamma_i,\gamma^4\}\Psi=0$: \textbf{$T_{4i}=0$ on shell.} (Stage-1 document, §9.7; checks \texttt{EMT\_homo\allowbreak{}geneousR\allowbreak{}eduction\allowbreak{}\_G3}, \texttt{EMT\_homo\allowbreak{}geneousR\allowbreak{}eduction}, \texttt{EMT\_homo\allowbreak{}geneousT\allowbreak{}imeSpace\allowbreak{}Vanishes\allowbreak{}OnShell\_\allowbreak{}G3}; for dirac16complex00 with the default and with a general smooth $U$, checks \texttt{C00\_EMT\_\allowbreak{}homogene\allowbreak{}ousEquat\allowbreak{}ionsOfSt\allowbreak{}ate} and \texttt{S5\_EMT\_h\allowbreak{}omogeneo\allowbreak{}usEquati\allowbreak{}onsOfSta\allowbreak{}te}, at three times of the frame $h_i=1+x_4^2/(i+2)$.)

\textbf{Energy conservation and dilution.} The tensor is not diagonal in general: the $T_{ij}$ above are nonzero for generic states when the rates differ. Nevertheless the component $\nu=4$ of $\nabla_\mu T^\mu{}_\nu=0$ contains only diagonal components. Write it out with the rule of Section 4.5 for the covariant derivative of a tensor (one term $+\Gamma$ for the upper index, one term $-\Gamma$ for the lower index), contracted over $\mu$:

\[
\nabla_\mu T^\mu{}_4=\partial_\mu T^\mu{}_4+\Gamma^\mu{}_{\mu\lambda}T^\lambda{}_4-\Gamma^\lambda{}_{\mu4}T^\mu{}_\lambda .
\]

Three facts reduce it. (a) On shell $T_{4i}=0$, so $T^i{}_4=g^{ii}T_{i4}=0$ and only $T^4{}_4=g^{44}T_{44}=-\rho$ remains in the first two terms. (b) Nothing depends on $x_i$, so $\partial_\mu T^\mu{}_4=\partial_4T^4{}_4=-\partial_4\rho$, and $\Gamma^\mu{}_{\mu4}=\partial_4\ln\sqrt{|g|}=\Theta$ (Section 5.8). (c) The only Christoffel symbols with a lower index 4 are $\Gamma^i{}_{i4}=\Gamma^i{}_{4i}=H_i$ (no sum; $\Gamma^4{}_{44}=\Gamma^4{}_{4i}=0$ because $g_{44}=-1$ is constant and $g_{4i}=0$), so the last term is $-\sum_{i\ne4}H_i\,T^i{}_i=-\sum_{i\ne4}H_i\,p_{(i)}$: the off-diagonal $T^i{}_j$ ($i\ne j$) would need $\Gamma^j{}_{i4}$ with $i\ne j$, which vanishes. Altogether $-\partial_4\rho-\Theta\rho-\sum_iH_ip_{(i)}=0$, that is

\[
\partial_4\rho+\sum_{i\ne4}H_i\,\bigl(\rho+p_{(i)}\bigr)=0 .
\]

With $\rho=mS+U$ and $\rho+p=S(m+U')$ this is $(m+U')\bigl(\partial_4S+\Theta S\bigr)=0$. Wherever $M_{\mathrm{eff}}\ne0$, $\partial_4(VS)=V(\partial_4S+\Theta S)=0$: \textbf{$S$ dilutes as the inverse 7-volume}, $S\,V=\text{constant}$ (Stage-1 document, §9.6, derived there; Chapter 11 derives the same law from the mode equation). For $U=0$ the energy density $mS$ then falls like $1/V$, as dust does.

\textbf{Beyond the scalar field.} Unlike $\tfrac12\dot\phi^2\ge0$, the quantity $\mathrm{KE}_L=\tfrac12SM_{\mathrm{eff}}$ has no fixed sign, because $S$ can have either sign (Section 6.6). Two consequences, both at the level of these classical formulas (Stage-1 document, §9.6): for $\rho>0$, $w<-1$ ("phantom") exactly when $\mathrm{KE}_L<0$, and $w>1$ exactly when $\mathrm{PE}_L<0$. Whether such states are physical is a separate question (Chapter 11 discusses it for dirac16complex).

\begingroup
\small
\begin{longtable}{@{}>{\raggedright\arraybackslash}p{0.267\linewidth}>{\raggedright\arraybackslash}p{0.267\linewidth}>{\raggedright\arraybackslash}p{0.267\linewidth}@{}}
\toprule
\textbf{quantity} & \textbf{homogeneous scalar field (Section 5.8)} & \textbf{homogeneous dirac16complex or dirac16complex00} \\
\midrule
\endfirsthead
\toprule
\textbf{quantity} & \textbf{homogeneous scalar field (Section 5.8)} & \textbf{homogeneous dirac16complex or dirac16complex00} \\
\midrule
\endhead
Lagrangian & $\tfrac12\dot\phi^2-V$ & $\mathcal L_s=K-mS-U$ \\
energy density & $\tfrac12\dot\phi^2+V$ & $mS+U$ \\
pressure & $\tfrac12\dot\phi^2-V$ & $SU'-U$ (all seven directions) \\
kinetic energy & $\tfrac12\dot\phi^2\ge0$ & $\mathrm{KE}_L=\tfrac12S(m+U')$, either sign \\
potential energy & $V$ & $\mathrm{PE}_L=\tfrac12(mS+2U-SU')$ \\
$\rho=\mathrm{KE}+\mathrm{PE}$, $p=\mathrm{KE}-\mathrm{PE}$ & yes & yes (split A) \\
\bottomrule
\end{longtable}
\endgroup

\textbf{Worked example with exact numbers.} The Stage-5 report records three exact on-shell states in the frame above. At $x_4=3/7$ with $m=5/9$, $\lambda=-3/8$ and $S=-83/15$ (\texttt{dirac16c\allowbreak{}omplex00\allowbreak{}-theory.\allowbreak{}json}, key \texttt{exactVal\allowbreak{}ues.homo\allowbreak{}geneousG\allowbreak{}3.G3t3}):

\[
\begin{aligned}
&mS=-\tfrac{83}{27},\qquad \tfrac\lambda2S^2=-\tfrac3{16}\cdot\tfrac{6889}{225}=-\tfrac{6889}{1200},\qquad \rho=-\tfrac{83}{27}-\tfrac{6889}{1200}=-\tfrac{95201}{10800},\qquad p=-\tfrac{6889}{1200},\\
&w=\frac{p}{\rho}=\frac{6889\cdot9}{95201}=\frac{747}{1147},\qquad \mathrm{KE}_L=\tfrac12S\bigl(m+\lambda S\bigr)=-\tfrac{78601}{10800},\qquad \mathrm{PE}_L=\tfrac12mS=-\tfrac{83}{54}.
\end{aligned}
\]

Here both the energy density and the pressure are negative, which the classical formulas allow because $S<0$ while $m>0$. The first recorded state ($x_4=1/3$, $m=3/7$, $\lambda=5/11$, $S=5617/567$) has $\rho=187781927/7072758$, $p=157753445/7072758$ and $w=28085/33431$, all positive. Stage 1 recorded a state with $w=75899/24059\approx3.155>1$ ($m=4$, $\lambda=7/6$, $S=-75899/7560$, $\mathrm{PE}_L=-75899/3780<0$; Stage-1 document, Result 9.3).

\subsection{Plane waves: the charge and the energy of dirac16complex00 have no sign}

\textbf{Theorem 7.6 (a Gordon-type identity).} Let $\Psi=u\exp\bigl(i\sum_ak_ax_a\bigr)$ be a plane-wave solution of the flat-space equation $\gamma^a\partial_a\Psi=M\Psi$ with a constant $M\ne0$ (for example the free field, $M=m$). Then

\[
J^a=\frac{k^a}{M}\,S\qquad(k^a:=\eta^{ab}k_b),\qquad\text{in particular}\qquad S=\frac{M}{\omega}\,J^4\ \ \text{with}\ \ \omega:=k^4=-k_4 .
\]

\emph{Proof.} With $\gamma(k)=\sum_ak_a\gamma^a$ the equation is $i\gamma(k)u=Mu$, so $\gamma(k)u=-iMu$. For the adjoint, $\bar\Psi$ carries the factor $\exp(-i\sum k_ax_a)$, so the second field equation gives $-i\,\bar u\gamma(k)=-M\bar u$, that is $\bar u\gamma(k)=-iM\bar u$ ($\bar u=u^\dagger C$). The Clifford relation gives $\gamma^a\gamma(k)+\gamma(k)\gamma^a=2k^a$. Sandwich it: $\bar u\gamma^a\gamma(k)u+\bar u\gamma(k)\gamma^au=-iM\,\bar u\gamma^au-iM\,\bar u\gamma^au=2k^a\,\bar uu$. Hence $\bar u\gamma^au=i(k^a/M)\,\bar uu$ and $J^a=-i\bar u\gamma^au=(k^a/M)S$. For $a=4$: $k^4=\eta^{44}k_4=-k_4=\omega$, the frequency of $e^{-i\omega x_4}$. $\square$

\textbf{The energy density of a free plane wave.} Let $U=0$. On shell $\mathcal L_s=0$ (Section 7.4), and $\partial_4\Psi=-i\omega\Psi$ gives $K_4=\tfrac i2(-i\omega-i\omega)\Psi^\dagger B\Psi=\omega\,J^4$. So

\[
\rho=K_4-\mathcal L_s=\omega\,J^4=\frac{\omega^2}{m}\,S .
\]

For a given momentum and frequency the solutions form an 8-dimensional space (Section 6.6), and on it the Hermitian form $J^4=u^\dagger Bu$ has four positive and four negative directions: the space has a basis of eight solutions, orthogonal for this form, on four of which $u^\dagger Bu>0$ and on four $u^\dagger Bu<0$. The pair of numbers (4,4) is called the \textbf{signature} of the form; Chapter 8 calls the indefinite inner product $f^\dagger Bh$ the \textbf{Krein} form, so the statement is the "Krein signature (4,4)" of check \texttt{QNT\_krei\allowbreak{}nSignatu\allowbreak{}re}. This book proves it only at rest (Section 8.8; the worked example below exhibits one rest solution of each sign); for moving waves it is taken from the check. So there are solutions of positive frequency with negative energy density.

\textbf{Worked example (at rest).} Take $m>0$ and the two positive-frequency rest solutions $\Psi_\pm=e^{-imx_4}u_\pm$ with $u_+=\tfrac12(e_0-e_4-ie_9-ie_{13})$ and $u_-=\tfrac12(e_0+e_4+ie_9-ie_{13})$ of Section 8.8. Both satisfy $-i\gamma^4u_\pm=u_\pm$, so both solve the field equation with $\omega=m$ (Section 7.4). With the table of $C$: $(Cu)_0=-u_4$, $(Cu)_4=-u_0$, $(Cu)_9=u_{13}$, $(Cu)_{13}=u_9$, so $Cu_+=u_+$ and $Cu_-=-u_-$. Hence $S=u_\pm^\dagger Cu_\pm=\pm1$, $J^4=u_\pm^\dagger Bu_\pm=u_\pm^\dagger C(-i\gamma^4)u_\pm=u_\pm^\dagger Cu_\pm=\pm1$, and

\[
\rho(\Psi_+)=+m,\qquad \rho(\Psi_-)=-m .
\]

Both are solutions of positive frequency and unit norm $u^\dagger u=1$; the second has negative energy density. Multiplying a solution by a number $c$ multiplies $\rho$ by $|c|^2$, so the energy density of solutions of the free dirac16complex00 field takes every real value: \textbf{the classical energy of dirac16complex00 is not bounded below.} The report exhibits the same with moving plane waves ($m=1$, $k=(1,1,2,3)$, $\omega=\pm4$: $\rho=\omega J^4$ with both signs of $J^4$ on each frequency), and with the interaction: for $\lambda=-1$ ($m=1$) exact homogeneous rest states with $S=1,10,100$ have $\rho=mS+\tfrac\lambda2S^2=\tfrac12,-40,-4900$, and for $\lambda=+1$ a family of exact self-consistent plane waves has $\rho=-56/75,\,-1512/845,\dots$, decreasing without bound (check \texttt{C00\_ener\allowbreak{}gy\_unbou\allowbreak{}ndedBelo\allowbreak{}w} and its \texttt{\detokenize{S5_}} twin; \texttt{dirac16c\allowbreak{}omplex00\allowbreak{}-theory.\allowbreak{}json}, key \texttt{commutin\allowbreak{}gFieldSp\allowbreak{}ecifics.\allowbreak{}energySt\allowbreak{}atement}).

\textbf{The contrast with dirac16complex.} For the Grassmann field the classical "energy" is not a number at all but an element of the Grassmann algebra; the physical energy exists only after quantization. There the same indefinite $B$ becomes the matrix of the canonical anticommutator, and in the good sector (the fields independent of $x_5,x_6,x_7$, Section 7.5) the filled Dirac sea and normal ordering give a Hamiltonian that is not negative (a preview of Section 8.10, derived in the Stage-1 document, §10.7, not a machine check). This mechanism needs the anticommutators of Fermi statistics; the classical commuting field has no counterpart of it.

\subsection{The primordial field}

\textbf{The field equation in the notebook's chart.} In the primordial field of Chapter 9, with $z=6Hx_0$, $s=\sin z$ and the diagonal vielbein $h=(\cot z,\ s^{1/6}e^{a_4}\,(\times3),\ 1,\ s^{1/6}e^{-a_4}\,(\times3))$, the curved gammas are $\gamma^{x_0}=\tan z\,\gamma^0$, $\gamma^{x_i}=s^{-1/6}e^{-a_4}\gamma^i$ ($i=1,2,3$), $\gamma^{x_4}=\gamma^4$ and $\gamma^{x_j}=s^{-1/6}e^{a_4}\gamma^j$ ($j=5,6,7$), and $\gamma^\mu\Omega_\mu=3H\gamma^0$ (Section 4.12). The field equation of Section 7.2 becomes, for both fields,

\[
\tan z\,\gamma^0\partial_0\Psi+s^{-1/6}e^{-a_4}\sum_{i=1}^3\gamma^i\partial_i\Psi+\gamma^4\partial_4\Psi+s^{-1/6}e^{a_4}\sum_{j=5}^7\gamma^j\partial_j\Psi+3H\gamma^0\Psi=M_{\mathrm{eff}}\,\Psi
\]

(the Stage-2 component equations; checks \texttt{C00\_prim\allowbreak{}ordial\_E\allowbreak{}Lcommuti\allowbreak{}ng}, \texttt{C00\_prim\allowbreak{}ordial\_g\allowbreak{}eometryM\allowbreak{}atchesSt\allowbreak{}age2} and their \texttt{\detokenize{S5_}} twins). Chapter 9 studies it in detail.

\textbf{An exact homogeneous solution for every $a_4$.} Look for a field that depends on $x_4$ only, $\Psi=u(x_4)$ (in the language of Chapter 9 it is the plane wave in the hidden coordinate with the imaginary wave number $K=-3iH$). All derivatives except $\partial_4$ vanish, and the equation becomes $\gamma^4\partial_4u+3H\gamma^0u=M_{\mathrm{eff}}u$. Multiplying by $-\gamma^4$:

\[
\partial_4u=N\,u,\qquad N:=-\gamma^4\bigl(M_{\mathrm{eff}}-3H\gamma^0\bigr)=-M_{\mathrm{eff}}\gamma^4+3H\gamma^4\gamma^0 .
\]

\textbf{$S$ is constant.} $N$ is real, so for commuting $u$, $\partial_4(u^\dagger Cu)=u^\dagger(N^TC+CN)u$. Now $N^T=M_{\mathrm{eff}}\gamma^4+3H\gamma^4\gamma^0$ (transpose with $(\gamma^4)^T=-\gamma^4$, $(\gamma^0)^T=\gamma^0$, and $\gamma^0\gamma^4=-\gamma^4\gamma^0$), and with $\gamma^4C=C\gamma^4$, $\gamma^0C=-C\gamma^0$:

\[
N^TC+CN=M_{\mathrm{eff}}\bigl(\gamma^4C-C\gamma^4\bigr)+3H\bigl(\gamma^4\gamma^0C+C\gamma^4\gamma^0\bigr)=0+3H\bigl(-C\gamma^4\gamma^0+C\gamma^4\gamma^0\bigr)=0 .
\]

So $S=u^\dagger Cu$ is constant, hence $M_{\mathrm{eff}}=m+U'(S)$ is constant, and $u(x_4)=\exp(Nx_4)\,u(0)$ solves the nonlinear equation exactly, for every potential $U$ and every function $a_4$: the free function does not appear (checks \texttt{C00\_prim\allowbreak{}ordial\_h\allowbreak{}omogeneo\allowbreak{}usState} and \texttt{S5\_primo\allowbreak{}rdial\_ho\allowbreak{}mogeneou\allowbreak{}sState}). For dirac16complex00 this is an exact classical solution; for dirac16complex the same formulas are the mean-field reading in which the bilinears are replaced by numbers.

\textbf{Its energy--momentum tensor.} For $\mu=0$: $\Omega_0=0$ (Section 4.12) and $\partial_0u=0$, so $D_0\Psi=0$ and $T^0{}_0=\mathcal L_s$. For the six directions $p\in\{1,2,3,5,6,7\}$: $T^p{}_p=\mathcal L_s$, by the argument of Section 7.10 ($\Omega_p$ is a combination of $S^{0p}$ and $S^{4p}$, and $\{\gamma^p,S^{0p}\}=\{\gamma^p,S^{4p}\}=0$; Proposition 9.2). $K_\perp=\tfrac12\bar\Psi\{\gamma^\mu,\Omega_\mu\}\Psi=0$ for the diagonal vielbein, so $\rho=mS+U$. On shell, therefore,

\[
\begin{aligned}
&\rho=mS+U,\qquad p_{(\mu)}=SU'-U\ \ (\text{all seven }\mu\ne4),\\
&\mathrm{KE}_L=\tfrac12SM_{\mathrm{eff}},\qquad \mathrm{PE}_L=\tfrac12(mS+2U-SU'),\qquad \mathrm{KE}_H=0,
\end{aligned}
\]

independent of $a_4$: the same equation of state as in Section 7.10.

\textbf{An exact dirac16complex00 source of the static field.} Now take the static member of the primordial field, $a_4=a_{4,0}$ constant (Section 9.15). Its Einstein tensor is diagonal, $G^\mu{}_\nu=\mathrm{diag}(15,15,15,15,21,15,15,15)H^2$ in the coordinate order $x_0,\dots,x_7$ (Chapter 9). So the Einstein equations $G^\mu{}_\nu=\kappa T^\mu{}_\nu$ ($\kappa$ is the gravitational coupling) say two things: every off-diagonal component of $T$ vanishes, and the diagonal components take the prescribed values. Proposition 9.4 with the slope $c=0$ states the result: the state $\Psi=u(x_4)$ satisfies all 64 components \textbf{if and only if} (1) the six bilinears $V_p:=\bar\Psi\gamma^0\gamma^p\gamma^4\Psi$ ($p=1,2,3,5,6,7$) vanish, (2) $mS=-36H^2/\kappa$ and (3) $\lambda S^2=30H^2/\kappa$. Condition (1) does \textbf{not} say that all three-gamma bilinears vanish; only these six enter. We derive all three conditions for this state.

\emph{The off-diagonal components.} For $a_4'=0$ Section 9.13 gives $\Omega_0=\Omega_4=0$ and $\Omega_p=-\eta_{pp}H\,h_p\,S^{0p}$ for $p\in\{1,2,3,5,6,7\}$, with the scale factors $h_p=s^{1/6}e^{a_{4,0}}$ for $p\le3$ and $h_p=s^{1/6}e^{-a_{4,0}}$ for $p\ge5$. The lowered curved gammas are $\gamma_\mu=g_{\mu\mu}\gamma^{x_\mu}=\eta_{\mu\mu}h_\mu\gamma^\mu$ (no sum), in particular $\gamma_0=\cot z\,\gamma^0$, $\gamma_4=-\gamma^4$ and $\gamma_p=\eta_{pp}h_p\gamma^p$. For the state on shell, $D_0\Psi=0$, $D_p\Psi=\Omega_p\Psi$ and $D_4\Psi=N\Psi$; and $D_0\bar\Psi=0$, $D_p\bar\Psi=-\bar\Psi\Omega_p$ and $D_4\bar\Psi=(N\Psi)^\dagger C=\Psi^\dagger N^TC=-\bar\Psi N$, the last step by $N^TC=-CN$ (shown above). Insert these into Theorem 7.4, where $g_{\mu\nu}=0$ for $\mu\ne\nu$:

\begin{itemize}
\item $T_{0p}=-\tfrac14\bar\Psi\{\gamma_0,\Omega_p\}\Psi=0$, because $\{\gamma^0,S^{0p}\}=0$ (Lemma 6.3: the index 0 belongs to the pair).
\item $T_{04}=-\tfrac14\bar\Psi\{\gamma_0,N\}\Psi=0$, because $\{\gamma^0,\gamma^4\}=0$ and $\{\gamma^0,\gamma^4\gamma^0\}=\gamma^0\gamma^4\gamma^0+\gamma^4\gamma^0\gamma^0=-\gamma^4+\gamma^4=0$.
\item For two different transverse directions $p\ne q$, $T_{pq}=-\tfrac14\bar\Psi\bigl(\{\gamma_p,\Omega_q\}+\{\gamma_q,\Omega_p\}\bigr)\Psi=\tfrac14H\eta_{pp}\eta_{qq}h_ph_q\,\bar\Psi\bigl(\gamma^p\gamma^0\gamma^q+\gamma^q\gamma^0\gamma^p\bigr)\Psi=0$, by Lemma 6.3 and because reversing the order of three distinct anticommuting gammas costs three signs.
\item $T_{p4}=-\tfrac14\bigl[\bar\Psi\{\gamma_p,N\}\Psi+\bar\Psi\{\gamma_4,\Omega_p\}\Psi\bigr]$. In the first term $\{\gamma^p,\gamma^4\}=0$, while $\gamma^p$ commutes with the pair $\gamma^4\gamma^0$ (it passes two gammas), so $\{\gamma_p,N\}=3H\eta_{pp}h_p\cdot2\gamma^p\gamma^4\gamma^0=6H\eta_{pp}h_p\,\gamma^0\gamma^p\gamma^4$. In the second, the two minus signs of $\gamma_4$ and $\Omega_p$ cancel and Lemma 6.3 gives $\{\gamma_4,\Omega_p\}=\eta_{pp}Hh_p\,\gamma^4\gamma^0\gamma^p=\eta_{pp}Hh_p\,\gamma^0\gamma^p\gamma^4$ (move $\gamma^4$ two places to the right). Hence
\end{itemize}

\[
T_{p4}=T_{4p}=-\tfrac74\,\eta_{pp}\,H\,h_p\,V_p\qquad(p=1,2,3,5,6,7).
\]

Since $H\ne0$ and $h_p\ne0$, the off-diagonal Einstein equations hold exactly when all six $V_p$ vanish: condition (1). (For $p\le3$ the coefficient is $-\tfrac74Hs^{1/6}e^{a_{4,0}}$, the entry of the Stage-2 table of these bilinears, \texttt{provenan\allowbreak{}ce/DIRAC\allowbreak{}16COMPLE\allowbreak{}X\_PRIMOR\allowbreak{}DIAL\_FIE\allowbreak{}LD.md} §13. For a slope $c\ne0$ nine further bilinears $\bar\Psi\gamma^i\gamma^4\gamma^j\Psi$ appear, with coefficients proportional to $c$; that is why Proposition 9.4 needs 15 bilinears in general and only these 6 here.)

\emph{The diagonal components.} On shell $T^0{}_0=T^p{}_p=\mathcal L_s=SU'-U=\tfrac\lambda2S^2$ and $T^4{}_4=g^{44}T_{44}=-\rho=-(mS+\tfrac\lambda2S^2)$. The seven equations with $\mu=\nu\ne4$ read $15H^2=\tfrac\kappa2\lambda S^2$, which is condition (3). The equation $G^4{}_4=\kappa T^4{}_4$ reads $21H^2=-\kappa mS-\tfrac\kappa2\lambda S^2=-\kappa mS-15H^2$, so $\kappa mS=-36H^2$, which is condition (2).

\emph{The example.} The Stage-5 report exhibits such a state for dirac16complex00 (\texttt{dirac16c\allowbreak{}omplex00\allowbreak{}-theory.\allowbreak{}json}, key \texttt{primordi\allowbreak{}al.stati\allowbreak{}cSourceE\allowbreak{}xample}). Units $H=\kappa=1$; parameters

\[
m=-30,\qquad \lambda=\tfrac{125}6,\qquad S=\tfrac65,\qquad M_{\mathrm{eff}}=m+\lambda S=-30+25=-5 .
\]

Conditions (2) and (3) hold: $mS=-36$ and $\lambda S^2=\tfrac{125}6\cdot\tfrac{36}{25}=30$. The matrix $N=5\gamma^4+3\gamma^4\gamma^0$ satisfies $N^2=25(\gamma^4)^2+9(\gamma^4\gamma^0)^2+15(\gamma^4\gamma^4\gamma^0+\gamma^4\gamma^0\gamma^4)=-25+9+0=-16$, so its eigenvalues are $\pm4i$; the column

\[
v=(1,\,3i,\,0,\,0,\,1,\,-3i,\,0,\,0,\,-3,\,-i,\,0,\,0,\,-3,\,i,\,0,\,0)^T
\]

satisfies $Nv=4i\,v$ and $v^\dagger Cv=32$ (Exercises 7.13 and 7.10). Condition (1) holds as well. The column $v$ is nonzero only in the eight places $\{0,1,4,5,8,9,12,13\}$; call this set $I$. Because $C$ and $\gamma^0$ are real and symmetric, $v^\dagger C\gamma^0=\bigl((C\gamma^0)^Tv\bigr)^\dagger=(\gamma^0Cv)^\dagger$, so

\[
V_p(v)=v^\dagger C\gamma^0\gamma^p\gamma^4v=a^\dagger\gamma^pb,\qquad a:=\gamma^0Cv,\quad b:=\gamma^4v .
\]

With the tables of Sections 2.8 and 2.9 both $a$ and $b$ are again nonzero only on $I$ (Exercise 7.13 lists them). For $p=1,2,5,6$ the rows $0,1,4,5,8,9,12,13$ of $\gamma^p$ all point to columns outside $I$ (for $\gamma^1$ they are $15,14,11,10,7,6,3,2$), so $\gamma^pb$ vanishes on $I$, and every term $a_k^\ast(\gamma^pb)_k$ of $a^\dagger\gamma^pb$ contains a zero factor: $V_1=V_2=V_5=V_6=0$. For $p=3$ and $p=7$ the rows of $I$ point into $I$, and the eight terms cancel in pairs (Exercise 7.13): $V_3=V_7=0$. Other three-gamma bilinears of $v$ need not vanish and do not, for example $\bar v\gamma^1\gamma^2\gamma^3v=24$; they do not enter $T$ for this state. So $\Psi=e^{i\omega x_4}u_0$ with $\omega=4=\sqrt{M_{\mathrm{eff}}^2-9H^2}$ and $u_0=\sqrt{S/32}\;v=\sqrt{3/80}\;v$ is an exact solution with $S=6/5$ (the factor $e^{i\omega x_4}$ and its conjugate cancel in every bilinear, and $V_p(u_0)=\tfrac3{80}V_p(v)=0$), and

\[
\begin{aligned}
&\rho=mS+\tfrac\lambda2S^2=-36+15=-21,\qquad p=\tfrac\lambda2S^2=15,\qquad w=-\tfrac57,\\
&\mathrm{KE}_L=\tfrac12SM_{\mathrm{eff}}=-3,\qquad \mathrm{PE}_L=\rho-\mathrm{KE}_L=-18 .
\end{aligned}
\]

These are exactly the required source $\rho_{\mathrm{req}}=-21H^2/\kappa$, $p_{\mathrm{req}}=+15H^2/\kappa$ of Section 9.15. Together with condition (1) this proves $G^\mu{}_\nu=\kappa T^\mu{}_\nu$ in all 64 components, and the verifiers confirm it exactly (checks \texttt{C00\_prim\allowbreak{}ordial\_s\allowbreak{}taticFie\allowbreak{}ldSource\allowbreak{}dExactly}, \texttt{S5\_primo\allowbreak{}rdial\_st\allowbreak{}aticFiel\allowbreak{}dSourced\allowbreak{}Exactly}, \texttt{S5\_stati\allowbreak{}c\_source\allowbreak{}ProperCh\allowbreak{}art}). In words: \textbf{within eight-dimensional Einstein gravity, the static member of the primordial field is sourced exactly by a homogeneous classical dirac16complex00 field, with negative energy density.} The statement concerns Einstein's equations without the Lovelock terms (which the project does not compute, Section 9.9) and the smooth field of the notebook's chart, which in the proper hidden coordinate $y=\ln(\sin z)/(6H)$ of Section 9.15 is the region $y<0$ (because $0<\sin z<1$ there). Section 9.16 also considers a model in which two copies of this field are glued along the surface $y=0$, one of them reflected (the "Z2 brane"; Z2 names the reflection $y\to-y$). There the slope of the warp factor jumps at $y=0$, and this kink needs an additional source concentrated on the surface, which this state does not supply. Two further limitations are part of the statement. The state has negative energy density, which the classical commuting field allows (Section 7.11) and which is not bounded below; its stability is not examined, and nothing says that such a state is realized. For dirac16complex the same numbers are only the mean-field reading of $\langle\bar\Psi\Psi\rangle=S$; Chapter 13 found that none of the computed Kohn--Sham states meets the two conditions.

\textbf{The static warped form (a preview of Chapters 13 and 14).} In the proper hidden coordinate $y$ the static field is $ds^2=dy^2-dx_4^2+e^{2Hy}\bigl[e^{2a_{4,0}}(dx_1^2+dx_2^2+dx_3^2)-e^{-2a_{4,0}}(dx_5^2+dx_6^2+dx_7^2)\bigr]$ with $\sqrt{|g|}=e^{6Hy}$ (Section 9.15). Section 13.3 derives that the stationary ansatz $\Psi=e^{-i\varepsilon x_4}e^{i(k_1x_1+k_2x_2+k_3x_3)}e^{-3Hy}\chi(y)$, with a frequency $\varepsilon$ and a 3-momentum $(k_1,k_2,k_3)$, removes the term $3H\gamma^0$ exactly and gives the reduced equation

\[
\gamma^0\chi'+i\,\xi(y)\sum_{l=1}^3k_l\gamma^l\chi-i\varepsilon\gamma^4\chi=M_{\mathrm{eff}}\,\chi,\qquad \xi(y)=e^{-Hy-a_{4,0}} .
\]

The Greek letter $\xi$ (xi) names the weight with which the momentum enters, as in Chapter 13 (it has nothing to do with the small vector field $\xi^\mu$ of Section 7.8); the letter $\kappa$ stays reserved for the gravitational coupling (the theory files and the Rust code call this weight \texttt{\detokenize{kappa}}). The field equation is linear in $\Psi$ at fixed $M_{\mathrm{eff}}$, so this reduction is the same for both statistics (checks \texttt{C00\_stat\allowbreak{}ic\_geome\allowbreak{}tryAndRe\allowbreak{}ducedEqu\allowbreak{}ation}, \texttt{S5\_stati\allowbreak{}c\_geomet\allowbreak{}ryAndRed\allowbreak{}ucedEqua\allowbreak{}tion}, which take the momentum along $x_1$). What differs is the meaning of the energy. Section 13.4 finds a basis of $\mathbb C^{16}$ in which the equation splits into eight separate pairs of components, called \textbf{blocks}; each block is labelled by three signs $(j,s_2,s_3)$, each $+1$ or $-1$, and on each block the charge matrix $B$ acts as multiplication by the number $js_2$. This number is the \textbf{Krein sign} of the block (the sign of the form $u^\dagger Bu$ of Section 7.11 on that block); four blocks have $+1$ and four $-1$. Chapter 13 describes the quantized dirac16complex in the Kohn--Sham approximation (Chapter 12), in which each quantum occupies an \textbf{orbital}, one solution $\chi$ of this equation, and it assigns to each orbital a \textbf{one-body} energy--momentum tensor, obtained from the classical formula by replacing $\Psi^\dagger$ with $u^\dagger B$ (the expectation-value rule of Section 8.12). For a classical commuting mode in one block, every component of $T_{\mu\nu}$ equals $js_2$ times this one-body value of the same orbital. So the classical energy density of dirac16complex00 is negative in four of the eight blocks, also in the static primordial field; explicit exact box modes show both signs (checks \texttt{C00\_stat\allowbreak{}ic\_class\allowbreak{}icalMode\allowbreak{}IsKreinW\allowbreak{}eightedK\allowbreak{}S}, \texttt{C00\_stat\allowbreak{}ic\_class\allowbreak{}icalEner\allowbreak{}gyKreinS\allowbreak{}igned} and their \texttt{\detokenize{S5_}} twins). Chapter 14 builds the Kohn--Sham model of dirac16complex00 on this.

\subsection{The chirality map of the equations, the current and the energy--momentum tensor}

Section 6.11 proved $\mathcal L_{m,\lambda}[\gamma^8\Psi]=-\mathcal L_{-m,-\lambda}[\Psi]$. The same five facts about $\gamma^8$ act on everything derived in this chapter.

\textbf{Theorem 7.7.} For every vielbein, every field configuration and both kinds of components, with $E_{m,\lambda}[\Psi]:=\gamma^\mu D_\mu\Psi-(m+\lambda S)\Psi$:

\[
\begin{aligned}
&E_{-m,-\lambda}[\gamma^8\Psi]=-\gamma^8E_{m,\lambda}[\Psi],\qquad J^\mu[\gamma^8\Psi]=-J^\mu[\Psi],\qquad S[\gamma^8\Psi]=S[\Psi],\\
&T_{\mu\nu}[\gamma^8\Psi;-m,-\lambda]=-T_{\mu\nu}[\Psi;m,\lambda].
\end{aligned}
\]

Hence $\Psi$ solves the field equations with $(m,\lambda)$ exactly when $\gamma^8\Psi$ solves them with $(-m,-\lambda)$.

\emph{Proof.} By (G5) of Section 6.11, $\gamma^\mu D_\mu(\gamma^8\Psi)=\gamma^\mu\gamma^8D_\mu\Psi=-\gamma^8\gamma^\mu D_\mu\Psi$; and $S$ is unchanged, so $(-m-\lambda S)\gamma^8\Psi=-\gamma^8(m+\lambda S)\Psi$. Adding, $E_{-m,-\lambda}[\gamma^8\Psi]=-\gamma^8E_{m,\lambda}[\Psi]$; since $\gamma^8$ is invertible, one vanishes exactly when the other does. The current: $\bar\Psi\gamma^8\gamma^\mu\gamma^8\Psi=-\bar\Psi\gamma^\mu\Psi$. The tensor: the bracket of Theorem 7.4 consists of bilinears with one gamma matrix and one covariant derivative, each of which changes sign exactly like $K$ in Section 6.11; and $\mathcal L_s[\gamma^8\Psi;-m,-\lambda]=-\mathcal L_s[\Psi;m,\lambda]$ by Theorem 6.7 with $(m,\lambda)$ replaced by $(-m,-\lambda)$. $\square$

The pairing reports check the theorem as an identity in completely generic symbols, in a Grassmann algebra, with exact jets of the general vielbein G1, and in the primordial field. The names below are those of \texttt{wolfram-\allowbreak{}pairing-\allowbreak{}report.j\allowbreak{}son}; each has a Python twin in \texttt{python-p\allowbreak{}airing-r\allowbreak{}eport.js\allowbreak{}on} with \texttt{\detokenize{S5_}} in place of \texttt{\detokenize{PAIR_}} (for example \texttt{S5\_T1gen\allowbreak{}eric\_fie\allowbreak{}ldEquati\allowbreak{}on}):

\begin{Verbatim}[fontsize=\small]
generic symbols:   PAIR_T1generic_fieldEquation  PAIR_T1generic_conjugateFieldEquation
                   PAIR_T1generic_current        PAIR_T1generic_emtAll36
Grassmann algebra: PAIR_T1grassmann_diracOperator  PAIR_T1grassmann_emt
                   PAIR_T1grassmann_current
jets of G1:        PAIR_T1jets_fieldEquations  PAIR_T1jets_emt
                   PAIR_T1jets_onShellImage    PAIR_T1jets_conservationBoth
primordial field:  PAIR_T1primordial_emt64
\end{Verbatim}

\textbf{Example.} For the homogeneous state of Section 7.10, $\rho=mS+\tfrac\lambda2S^2$. Its image $\gamma^8\Psi$, a solution with $(-m,-\lambda)$, has the same $S$, hence $\rho=-mS-\tfrac\lambda2S^2$: exactly the negative, as the theorem says.

\textbf{What follows, and what does not.} Adding a solution with $(m,\lambda)$ and its image with $(-m,-\lambda)$ in the same gravitational field, the two energy--momentum tensors and the two currents cancel at every point and at every time. Chapter 15 states this as the pair corollary of Theorem T1, with its hypotheses: the image must carry the coupling $-\lambda$ (or $\lambda=0$), and for the quantized dirac16complex the image field carries the opposite Krein metric $-B$. (The Krein metric of a quantized field is the matrix of its equal-time anticommutator, $\{\Psi,\Psi^\dagger\}=B\,\delta^7/\sqrt{|g|}$ for dirac16complex, Section 8.6; for $\gamma^8\Psi$ it is $\gamma^8B\gamma^8=-B$, because $\gamma^8$ anticommutes with $\gamma^4$ and commutes with $C$.) An independently quantized theory with mass $-m$, with its own metric $+B$ and positive energies, does not cancel but doubles (\texttt{pairing-\allowbreak{}theory.j\allowbreak{}son}, keys \texttt{\detokenize{T1}} and \texttt{\detokenize{T1krein}}; a preview of Chapter 15). These are statements about conservation laws and constraints. They compute no process that creates such a pair, and Chapter 16 explains exactly what they do and do not establish about the notebook's hypothesis.

\subsection{How the repository checks all of this}

\begingroup
\small
\begin{longtable}{@{}>{\raggedright\arraybackslash}p{0.400\linewidth}>{\raggedright\arraybackslash}p{0.400\linewidth}@{}}
\toprule
\textbf{statement (section)} & \textbf{checks} \\
\midrule
\endfirsthead
\toprule
\textbf{statement (section)} & \textbf{checks} \\
\midrule
\endhead
field equations, commuting components (flat, G1, G2, general $U$) (7.2) & \texttt{C00\_EL\_c\allowbreak{}ommuting\allowbreak{}Flat}, \texttt{C00\_EL\_c\allowbreak{}ommuting\allowbreak{}Curved}, \texttt{C00\_EL\_c\allowbreak{}ommuting\allowbreak{}GeneralS\allowbreak{}moothU} (both) \\
field equations, Grassmann components; same form (7.3) & \texttt{LAG\_eule\allowbreak{}rLagrang\allowbreak{}ePsibar\_\allowbreak{}G1}, \texttt{LAG\_eule\allowbreak{}rLagrang\allowbreak{}ePsi\_G1} (and G2); \texttt{S5\_EL\_gr\allowbreak{}assmannC\allowbreak{}urved}; \texttt{C00\_EL\_i\allowbreak{}dentical\allowbreak{}FormBoth\allowbreak{}Statisti\allowbreak{}cs} (both) \\
second-order form with $-\tfrac14R$ (7.4) & \texttt{GEO\_lich\allowbreak{}nerowicz\allowbreak{}\_G1}, \texttt{GEO\_lich\allowbreak{}nerowicz\allowbreak{}\_G2}; \texttt{C00\_conn\allowbreak{}ection\_n\allowbreak{}onTrivia\allowbreak{}lCouplin\allowbreak{}g} (both) \\
current identity, conservation, Hermiticity (7.5) & \texttt{MA\_M1\_no\allowbreak{}etherIde\allowbreak{}ntity\_co\allowbreak{}mmuting\_\allowbreak{}G1}, \texttt{MA\_M1\_no\allowbreak{}etherIde\allowbreak{}ntity\_gr\allowbreak{}assmann\_\allowbreak{}G1}; \texttt{QNT\_curr\allowbreak{}entHermi\allowbreak{}ticity}; \texttt{C00\_curr\allowbreak{}ent\_cons\allowbreak{}ervation\allowbreak{}AndReali\allowbreak{}ty} (both) \\
indefinite charge of dirac16complex00 (7.5) & \texttt{C00\_char\allowbreak{}ge\_indef\allowbreak{}inite} (both) \\
$T_{\mu\nu}$ from the full vielbein variation (7.7) & \texttt{EMT\_vari\allowbreak{}ation\_G3}, \texttt{EMT\_vari\allowbreak{}ation}; \texttt{C00\_EMT\_\allowbreak{}vielbein\allowbreak{}Variatio\allowbreak{}n} (both) \\
symmetry, reality, trace, conservation (7.8) & \texttt{EMT\_symm\allowbreak{}etricHer\allowbreak{}mitian\_G\allowbreak{}1}, \texttt{EMT\_trac\allowbreak{}e\_G1}, \texttt{EMT\_cons\allowbreak{}ervation\allowbreak{}\_G1} (and G2); \texttt{C00\_EMT\_\allowbreak{}symmetri\allowbreak{}cAndReal}, \texttt{C00\_EMT\_\allowbreak{}traceOnS\allowbreak{}hell}, \texttt{C00\_EMT\_\allowbreak{}conserva\allowbreak{}tionOnSh\allowbreak{}ell}, \texttt{C00\_EMT\_\allowbreak{}generalS\allowbreak{}moothU} (both) \\
$\rho=K_4-\mathcal L_s=\mathrm{KE}_H+\mathrm{PE}_H$ (7.9) & \texttt{C00\_EMT\_\allowbreak{}observer\allowbreak{}SplitGau\allowbreak{}ssianNor\allowbreak{}mal} (both) \\
homogeneous $\rho$, $p$, KE, PE, $w$, $T_{ij}$, $T_{4i}$ (7.10) & \texttt{EMT\_homo\allowbreak{}geneousR\allowbreak{}eduction\allowbreak{}\_G3}, \texttt{EMT\_homo\allowbreak{}geneousR\allowbreak{}eduction}; \texttt{C00\_EMT\_\allowbreak{}homogene\allowbreak{}ousEquat\allowbreak{}ionsOfSt\allowbreak{}ate} (both) \\
energy of dirac16complex00 unbounded below (7.11) & \texttt{C00\_ener\allowbreak{}gy\_unbou\allowbreak{}ndedBelo\allowbreak{}w} (both); \texttt{QNT\_krei\allowbreak{}nSignatu\allowbreak{}re} \\
primordial field: equation, homogeneous state, static source (7.12) & \texttt{C00\_prim\allowbreak{}ordial\_E\allowbreak{}Lcommuti\allowbreak{}ng}, \texttt{C00\_prim\allowbreak{}ordial\_h\allowbreak{}omogeneo\allowbreak{}usState}, \texttt{C00\_prim\allowbreak{}ordial\_s\allowbreak{}taticFie\allowbreak{}ldSource\allowbreak{}dExactly} (both); \texttt{S5\_stati\allowbreak{}c\_source\allowbreak{}ProperCh\allowbreak{}art} \\
static warped form and block signs (7.12) & \texttt{C00\_stat\allowbreak{}ic\_geome\allowbreak{}tryAndRe\allowbreak{}ducedEqu\allowbreak{}ation}, \texttt{C00\_stat\allowbreak{}ic\_class\allowbreak{}icalMode\allowbreak{}IsKreinW\allowbreak{}eightedK\allowbreak{}S}, \texttt{C00\_stat\allowbreak{}ic\_class\allowbreak{}icalEner\allowbreak{}gyKreinS\allowbreak{}igned} (both) \\
$\gamma^8$ map of equations, current and $T_{\mu\nu}$ (7.13) & \texttt{PAIR\_T1g\allowbreak{}eneric\_*}, \texttt{PAIR\_T1g\allowbreak{}rassmann\allowbreak{}\_*}, \texttt{PAIR\_T1j\allowbreak{}ets\_*}, \texttt{PAIR\_T1p\allowbreak{}rimordia\allowbreak{}l\_emt64} \\
\bottomrule
\end{longtable}
\endgroup

"Both" means the \texttt{\detokenize{C00_}} check of \texttt{wolfram-\allowbreak{}dirac16c\allowbreak{}omplex00\allowbreak{}-report.\allowbreak{}json} and the \texttt{\detokenize{S5_}} check with the same ending in \texttt{python-d\allowbreak{}irac16co\allowbreak{}mplex00-\allowbreak{}report.j\allowbreak{}son}. The \texttt{\detokenize{PAIR_}} checks are those of \texttt{wolfram-\allowbreak{}pairing-\allowbreak{}report.j\allowbreak{}son}; each has a twin with the same ending and the prefix \texttt{\detokenize{S5_}} in \texttt{python-p\allowbreak{}airing-r\allowbreak{}eport.js\allowbreak{}on}. The Stage-1 checks are in the reports of \texttt{artifact\allowbreak{}s/dirac1\allowbreak{}6complex\allowbreak{}/arbitra\allowbreak{}ry-field\allowbreak{}/}. Chapter 19 lists the commands that rerun the verifiers. The Stage-5 gate script that \texttt{handoff/\allowbreak{}specs/ST\allowbreak{}AGE5\_SPE\allowbreak{}C.md} (§6) plans, \texttt{scripts/\allowbreak{}verify\_s\allowbreak{}tage5\_pa\allowbreak{}ir\_creat\allowbreak{}ion.ps1} and \texttt{\detokenize{.sh}}, did not exist yet when this chapter was written, so the Stage-5 reports have not been regenerated from a fresh copy of the repository.

\subsection{What we proved and what we assumed}

\textbf{What we proved.} With complete derivations in this chapter: the Euler--Lagrange equations $\gamma^\mu D_\mu\Psi=(m+U'(S))\Psi$ and $(D_\mu\bar\Psi)\gamma^\mu=-(m+U'(S))\bar\Psi$ in every gravitational field, for commuting components (Section 7.2) and for Grassmann components (Section 7.3); that the second is the conjugate of the first, that the equations are of first order in time, that $K=M_{\mathrm{eff}}S$ and $\mathcal L_s=SU'-U$ on shell, and that for $U=0$ every solution obeys the second-order equation with the curvature term $-\tfrac14R$; the off-shell current identity (Theorem 7.1), the Hermiticity and conservation of $J^\mu=-i\bar\Psi\gamma^\mu\Psi$, its relation to the Noether current, the conservation of $Q$ (by the slice argument of Section 5.7) and $J^4=\Psi^\dagger B\Psi$, which is indefinite for dirac16complex00; the energy--momentum tensor of Theorem 7.4 from the symmetric variation of the vielbein, including the change of the spin connection (Lemma 7.2, Corollary 7.3); its symmetry, reality, trace $-mS+7SU'-8U$ on shell, and conservation (Theorem 7.5); $\rho=K_4-\mathcal L_s=\mathrm{KE}_H+\mathrm{PE}_H$ and the oscillator form of $K_4$; the homogeneous formulas $\rho=mS+U$, $p=SU'-U$ in all seven directions, the off-diagonal components, $T_{4i}=0$ on shell and, although $T_{ij}\ne0$ in general, the energy equation with only diagonal components and the dilution $SV=\text{constant}$; the Gordon-type identity (Theorem 7.6), $\rho=\omega J^4$ for free plane waves and the negative energy density of the positive-frequency rest solution $u_-$, so that the classical energy of dirac16complex00 is unbounded below; the exact homogeneous solution in the primordial field for every $a_4$ and every $U$, with its equation of state; for this state in the static field, the off-diagonal components of its energy--momentum tensor ($T_{0p}=T_{04}=T_{pq}=0$ and $T_{p4}=-\tfrac74\eta_{pp}Hh_pV_p$ with the six bilinears $V_p=\bar\Psi\gamma^0\gamma^p\gamma^4\Psi$), hence the three conditions under which it sources the static field in all 64 components (the six $V_p$ vanish, $mS=-36H^2/\kappa$, $\lambda S^2=30H^2/\kappa$: Proposition 9.4 with $c=0$), and the exact classical dirac16complex00 source that meets them ($m=-30$, $\lambda=125/6$, $S=6/5$: $Nv=4iv$, all six $V_p=0$, $\rho=-21$, $p=15$; Exercises 7.10 and 7.13); and the $\gamma^8$ map of the equations, the current and the energy--momentum tensor (Theorem 7.7).

\textbf{What we assumed or took from elsewhere.} The Lagrangian and its building blocks (Chapter 6), the geometry of Chapter 4 (the vielbein postulate, the divergence identity, the Lichnerowicz formula, the covariant constancy of the gammas), the variational calculus of Chapter 5 for ordinary and Grassmann variables, and the invariance of the action under changes of coordinates and local frame rotations (Section 6.9). The definition of the energy--momentum tensor by the metric variation with the spinor held fixed, the symmetric (non-rotating) change of the vielbein, the sign convention $\rho=T_{44}$, and the names "pressure" for the time-like transverse directions and "kinetic" and "potential" energy for the two splits are conventions of the project. The Krein signature (4,4) of the moving plane waves is taken from the check \texttt{QNT\_krei\allowbreak{}nSignatu\allowbreak{}re} (Section 8.8 proves the rest case). For the primordial field we used from Chapter 9 its vielbein, its spinor connection (Section 9.13) and the Einstein tensor of its static member, and from Section 9.15 the required source. The reduction of the static field to $2\times2$ blocks, the block Krein signs and the Kohn--Sham one-body tensor are derived in Chapters 13 and 14 and only previewed here; the statement that a classical mode's $T_{\mu\nu}$ is $js_2$ times the one-body value is quoted from the Stage-5 checks. The formulas for dirac16complex are classical statements about the Grassmann field; their meaning as operator statements needs the quantization of Chapter 8, and the positivity of the quantized energy in the good sector is derived there, not machine-checked. Where the bilinears are replaced by numbers for dirac16complex (Sections 7.10 and 7.12) this is the mean-field approximation. No statement of this chapter says that any of the classical states exists in nature.

\subsection{Exercises}

\textbf{Exercise 7.1.} In flat space, for the Lagrangian of one direction only, $\mathcal L=\tfrac12\bigl(\Psi^\dagger C\gamma^1\partial_1\Psi-\partial_1\Psi^\dagger C\gamma^1\Psi\bigr)-m\Psi^\dagger C\Psi$ (commuting components), compute $\partial\mathcal L/\partial\Psi_a^\ast$ and $\partial\mathcal L/\partial(\partial_1\Psi_a^\ast)$ and derive the Euler--Lagrange equation.

\textbf{Exercise 7.2.} Show that on shell $\mathcal L_s=\tfrac\lambda2S^2+\tfrac{2\mu}3S^3$ for the commuting field with $U=\tfrac\lambda2S^2+\tfrac\mu3S^3$, and compute the on-shell trace $T^\mu{}_\mu$.

\textbf{Exercise 7.3.} Using the table of $B$ in Section 2.9, compute $J^4=\Psi^\dagger B\Psi$ for $\Psi=e_0+e_9$, $\Psi=e_0+i\,e_9$ and $\Psi=e_0-i\,e_9$.

\textbf{Exercise 7.4.} Check the claim of Section 7.6 that $\Lambda=1+X$ satisfies $\Lambda^T\eta\Lambda=\eta$ to first order exactly when $\eta X$ is antisymmetric, and give an example of a nonzero $2\times2$ such $X$ for $\eta=\mathrm{diag}(1,-1)$.

\textbf{Exercise 7.5.} A homogeneous commuting state with $m=2$, $\lambda=-3$ and $S=1$ is given. Compute $\rho$, $p$, $w$, $\mathrm{KE}_L$ and $\mathrm{PE}_L$. Is it phantom?

\textbf{Exercise 7.6.} Verify the Stage-1 example of Section 7.10: $m=4$, $\lambda=7/6$, $S=-75899/7560$. Compute $\mathrm{PE}_L$ and $w$ and confirm $w>1$.

\textbf{Exercise 7.7.} In the homogeneous background of Section 7.10 with $U=0$, the 7-volume doubles between two times. What happens to $S$, $\rho$ and $w$?

\textbf{Exercise 7.8.} For the rest solutions with negative frequency, $\Psi=e^{+imx_4}u$ with $-i\gamma^4u=-u$, show that $S=-J^4$ and $\rho=-mJ^4=mS$. Which sign of $J^4$ gives a negative energy density now?

\textbf{Exercise 7.9.} For a moving plane wave with $m=3$, $k_1=4$ (all other $k$ zero except $k_4$) and positive frequency, find $\omega$ and the ratio $S/J^4$.

\textbf{Exercise 7.10.} For the static source of Section 7.12 verify $v^\dagger Cv=32$ with the table of $C$, and compute $\omega$ from $\omega^2=M_{\mathrm{eff}}^2-9H^2$.

\textbf{Exercise 7.11.} Show from Theorem 7.7 that for the homogeneous state the image $\gamma^8\Psi$ has the pressure $-p$ and the same $w$.

\textbf{Exercise 7.12.} In the proof of Theorem 7.5, take flat space with the vielbein $e_\mu{}^a=\delta_\mu{}^a$ and a small vector field $\xi(x)$ that vanishes near the boundary. Compute $\delta_\xi e_\mu{}^a$, the matrix $X_{cd}$, its symmetric part and its antisymmetric part, and state what the proof gives.

\textbf{Exercise 7.13.} For the column $v$ of the static source of Section 7.12 ($H=1$, $M_{\mathrm{eff}}=-5$, so $N=5\gamma^4+3\gamma^4\gamma^0=\gamma^4(5+3\gamma^0)$), use the tables of Sections 2.8 and 2.9 to (a) compute $\gamma^0v$ and verify $Nv=4i\,v$; (b) compute $b=\gamma^4v$ and $a=\gamma^0Cv$; (c) compute $\gamma^3b$ and $\gamma^7b$, and show $V_3(v)=a^\dagger\gamma^3b=0$ and $V_7(v)=a^\dagger\gamma^7b=0$.

\subsection{Answers to the exercises}

\textbf{Answer 7.1.} Write $A=C\gamma^1$ (real antisymmetric). $\partial\mathcal L/\partial\Psi_a^\ast=\tfrac12(A\partial_1\Psi)_a-m(C\Psi)_a$ and $\partial\mathcal L/\partial(\partial_1\Psi_a^\ast)=-\tfrac12(A\Psi)_a$. The Euler--Lagrange expression is $\tfrac12A\partial_1\Psi-mC\Psi+\tfrac12A\partial_1\Psi=C(\gamma^1\partial_1\Psi-m\Psi)$, so the equation is $\gamma^1\partial_1\Psi=m\Psi$: Section 7.2 with only one derivative direction and no connection.

\textbf{Answer 7.2.} $U'=\lambda S+\mu S^2$, so $SU'-U=\lambda S^2+\mu S^3-\tfrac\lambda2S^2-\tfrac\mu3S^3=\tfrac\lambda2S^2+\tfrac{2\mu}3S^3$. The trace is $-mS+7SU'-8U=-mS+7\lambda S^2+7\mu S^3-4\lambda S^2-\tfrac{8\mu}3S^3=-mS+3\lambda S^2+\tfrac{13\mu}3S^3$.

\textbf{Answer 7.3.} From the table, $(Bu)_0=i\,u_9$ and $(Bu)_9=-i\,u_0$. So $J^4=\Psi_0^\ast\,i\,\Psi_9-\Psi_9^\ast\,i\,\Psi_0$. For $e_0+e_9$: $i-i=0$. For $e_0+ie_9$: $i\cdot i-(-i)\,i=-1-1=-2$. For $e_0-ie_9$: $i(-i)-(i)(i)=1+1=2$. All three have $\Psi^\dagger\Psi=2$, yet the charge density is $0$, $-2$ and $+2$.

\textbf{Answer 7.4.} $\Lambda^T\eta\Lambda=\eta+X^T\eta+\eta X+O(X^2)$, so the condition is $X^T\eta+\eta X=0$, which says $(\eta X)^T=-\eta X$ because $\eta^T=\eta$. Example: $X=\begin{pmatrix}0&\beta\\ \beta&0\end{pmatrix}$ gives $\eta X=\begin{pmatrix}0&\beta\\-\beta&0\end{pmatrix}$, antisymmetric: an infinitesimal boost, $\Lambda=\begin{pmatrix}1&\beta\\ \beta&1\end{pmatrix}$.

\textbf{Answer 7.5.} $\rho=mS+\tfrac\lambda2S^2=2-\tfrac32=\tfrac12$, $p=\tfrac\lambda2S^2=-\tfrac32$, $w=-3$, $\mathrm{KE}_L=\tfrac12S(m+\lambda S)=\tfrac12(2-3)=-\tfrac12$, $\mathrm{PE}_L=\tfrac12mS=1$. Check: $\mathrm{KE}_L+\mathrm{PE}_L=\tfrac12=\rho$ and $\mathrm{KE}_L-\mathrm{PE}_L=-\tfrac32=p$. With $\rho>0$ and $\mathrm{KE}_L<0$ it is phantom ($w<-1$), at the level of these classical formulas.

\textbf{Answer 7.6.} $\mathrm{PE}_L=\tfrac12mS=2\cdot(-75899/7560)=-75899/3780<0$. $w=\lambda S/(2m+\lambda S)$ with $\lambda S=\tfrac76\cdot(-\tfrac{75899}{7560})=-\tfrac{531293}{45360}$ and $2m=8=\tfrac{362880}{45360}$, so $2m+\lambda S=-\tfrac{168413}{45360}$ and $w=\tfrac{531293}{168413}=\tfrac{75899}{24059}\approx3.155>1$ (divide numerator and denominator by 7). The energy density is positive, $\rho=mS+\tfrac\lambda2S^2\approx-40.16+58.80\approx18.64$, while $\mathrm{PE}_L<0$: exactly the situation in which Section 7.10 says $w>1$.

\textbf{Answer 7.7.} $S\,V$ is constant, so $S$ halves. With $U=0$, $\rho=mS$ halves as well (dust), and $p=0$, $w=0$ throughout.

\textbf{Answer 7.8.} $-i\gamma^4u=-u$ gives $J^4=u^\dagger C(-i\gamma^4)u=-u^\dagger Cu=-S$. With $\partial_4\Psi=+im\Psi$, $K_4=\tfrac i2(im+im)\Psi^\dagger B\Psi=-mJ^4$, and $\mathcal L_s=0$ on shell, so $\rho=-mJ^4=mS$. Now a positive $J^4$ (equivalently $S<0$) gives a negative energy density.

\textbf{Answer 7.9.} $\omega^2=m^2+k_1^2=9+16=25$, $\omega=5$. Theorem 7.6: $S/J^4=M/\omega=3/5$.

\textbf{Answer 7.10.} With $(Cu)_0=-u_4$, $(Cu)_1=-u_5$, $(Cu)_4=-u_0$, $(Cu)_5=-u_1$, $(Cu)_8=u_{12}$, $(Cu)_9=u_{13}$, $(Cu)_{12}=u_8$, $(Cu)_{13}=u_9$: $v^\dagger Cv=-2\,\mathrm{Re}(v_0^\ast v_4)-2\,\mathrm{Re}(v_1^\ast v_5)+2\,\mathrm{Re}(v_8^\ast v_{12})+2\,\mathrm{Re}(v_9^\ast v_{13})=-2(1)-2\,\mathrm{Re}\bigl((-3i)(-3i)\bigr)+2(9)+2\,\mathrm{Re}\bigl((i)(i)\bigr)=-2+18+18-2=32$. And $\omega^2=25-9=16$, $\omega=4$.

\textbf{Answer 7.11.} Theorem 7.7 gives $T_{\mu\nu}\to-T_{\mu\nu}$ with $(m,\lambda)\to(-m,-\lambda)$, so $\rho\to-\rho$ and $p\to-p$; directly, $p=\tfrac\lambda2S^2\to-\tfrac\lambda2S^2$ with the same $S$. The ratio $w=p/\rho$ is unchanged.

\textbf{Answer 7.12.} In flat space $\Gamma=0$ and $\omega=0$. The vielbein is constant, so $\delta_\xi e_\mu{}^a=\xi^\lambda\partial_\lambda\delta_\mu{}^a+\delta_\lambda{}^a\partial_\mu\xi^\lambda=\partial_\mu\xi^a$. Then $X^c{}_d=\delta_\xi e_\mu{}^c\,e_d{}^\mu=\partial_d\xi^c$ and $X_{cd}=\partial_d\xi_c$ with $\xi_c=\eta_{ca}\xi^a$. Its symmetric part is $\tfrac12(\partial_d\xi_c+\partial_c\xi_d)$ and its antisymmetric part $\varepsilon_{cd}=\tfrac12(\partial_d\xi_c-\partial_c\xi_d)$, an infinitesimal frame rotation that changes from point to point. The field changes by $\xi^\lambda\partial_\lambda\Psi$. The proof then gives $\int T^{\sigma\mu}\partial_\mu\xi_\sigma\,d^8x=0$ on shell for every such $\xi$, and after the integration by parts $\partial_\mu T^{\sigma\mu}=0$: the flat-space conservation law of Section 5.7, now for the symmetric tensor of Theorem 7.4.

\textbf{Answer 7.13.} The column is $v=(1,3i,0,0,1,-3i,0,0,-3,-i,0,0,-3,i,0,0)^T$; only the places $I=\{0,1,4,5,8,9,12,13\}$ are nonzero, and every column below is written by its entries in these eight places, in this order (the other eight entries are zero). (a) The table of $\gamma^0$ says $(\gamma^0u)_i=u_{i+8}$ for $i\le7$ and $u_{i-8}$ for $i\ge8$, so $\gamma^0v$ has the entries $(-3,-i,-3,i,1,3i,1,-3i)$. Then $w:=5v+3\gamma^0v$ has the entries $(-4,12i,-4,-12i,-12,4i,-12,-4i)$; for example $5\cdot1+3\cdot(-3)=-4$ and $5\cdot3i+3\cdot(-i)=12i$. The rows of $\gamma^4$ in the places of $I$ read $(\gamma^4u)_0=-u_{13}$, $(\gamma^4u)_1=u_{12}$, $(\gamma^4u)_4=u_9$, $(\gamma^4u)_5=-u_8$, $(\gamma^4u)_8=u_5$, $(\gamma^4u)_9=-u_4$, $(\gamma^4u)_{12}=-u_1$, $(\gamma^4u)_{13}=u_0$, so $Nv=\gamma^4w$ has the entries $(4i,-12,4i,12,-12i,4,-12i,-4)$. And $4iv$ has the same entries: for example $4i\cdot3i=-12$, $4i\cdot(-3i)=12$, $4i\cdot(-3)=-12i$ and $4i\cdot(-i)=4$. (b) With the same rows of $\gamma^4$, $b=\gamma^4v$ has the entries $(-v_{13},v_{12},v_9,-v_8,v_5,-v_4,-v_1,v_0)$, that is $(-i,-3,-i,3,-3i,-1,-3i,1)$. The table of $C$ says $(Cu)_i=-u_{i+4}$ for $i\le3$, $-u_{i-4}$ for $4\le i\le7$, $u_{i+4}$ for $8\le i\le11$ and $u_{i-4}$ for $i\ge12$. So $Cv$ has the entries $(-v_4,-v_5,-v_0,-v_1,v_{12},v_{13},v_8,v_9)$, that is $(-1,3i,-1,-3i,-3,i,-3,-i)$, and then $a=\gamma^0(Cv)$ has the entries $(-3,i,-3,-i,-1,3i,-1,-3i)$. Both $a$ and $b$ vanish outside $I$. (c) The rows of $\gamma^3$ in the places of $I$ read $-13,+12,-9,+8,+5,-4,+1,-0$, so $\gamma^3b$ has the entries $(-b_{13},b_{12},-b_9,b_8,b_5,-b_4,b_1,-b_0)$, that is $(-1,-3i,1,-3i,3,i,-3,i)$. The conjugates of the entries of $a$ are $(-3,-i,-3,i,-1,-3i,-1,3i)$, and the eight products $a_k^\ast(\gamma^3b)_k$ are

\[
3,\quad -3,\quad -3,\quad 3,\quad -3,\quad 3,\quad 3,\quad -3,
\]

whose sum is $V_3(v)=0$. The rows of $\gamma^7$ in the places of $I$ read $+8,+9,-12,-13,-0,-1,+4,+5$, so $\gamma^7b$ has the entries $(b_8,b_9,-b_{12},-b_{13},-b_0,-b_1,b_4,b_5)$, that is $(-3i,-1,3i,-1,i,3,-i,3)$, and the eight products $a_k^\ast(\gamma^7b)_k$ are

\[
9i,\quad i,\quad -9i,\quad -i,\quad -i,\quad -9i,\quad i,\quad 9i,
\]

whose sum is $V_7(v)=0$. Together with the support argument of Section 7.12 for $p=1,2,5,6$, all six $V_p$ vanish. (A computer check with the fixture gammas of \texttt{artifact\allowbreak{}s/dirac1\allowbreak{}6complex\allowbreak{}/arbitra\allowbreak{}ry-field\allowbreak{}/algebra\allowbreak{}-fixture\allowbreak{}.json} gives the same, and finds that the three-gamma bilinears $\bar v\gamma^1\gamma^2\gamma^3v$, $\bar v\gamma^1\gamma^6\gamma^7v$, $\bar v\gamma^2\gamma^5\gamma^7v$ and $\bar v\gamma^3\gamma^5\gamma^6v$ equal $24$, $24$, $-24$ and $24$: these are not among the $V_p$ and do not enter the Einstein equations for this state.)

\section{Canonical quantization in 4+4 dimensions}

\subsection{What this chapter does}

Chapters 6 and 7 treat dirac16complex as a classical field: a column $\Psi=(\Psi_0,\dots,\Psi_{15})^T$ of complex Grassmann-valued functions of the eight coordinates, with the Lagrangian

\[
\mathcal L=\sqrt{|g|}\,\Bigl[\tfrac12\bigl(\bar\Psi\gamma^\mu D_\mu\Psi-(D_\mu\bar\Psi)\gamma^\mu\Psi\bigr)-m\,\bar\Psi\Psi-U(\bar\Psi\Psi)\Bigr],\qquad \bar\Psi=\Psi^\dagger C,
\]

and the field equation $\gamma^\mu D_\mu\Psi=(m+U'(\bar\Psi\Psi))\Psi$. \textbf{Quantization} turns the components into operators acting on a space of states, with rules that make the classical equations come back as equations between operators. This chapter carries out the standard procedure, \textbf{canonical quantization with respect to the time} $x_4$, and finds three things that do not happen in ordinary four-dimensional physics:

\begin{enumerate}
\item the canonical anticommutator $\{\Psi,\Psi^\dagger\}$ is not a positive matrix but the matrix $B=-iC\gamma^4$, which has eight eigenvalues $+1$ and eight eigenvalues $-1$; so the natural space of states is a \textbf{Krein space}, a space with an inner product that is not positive, and no Spin(4,4)-invariant choice of the charge density avoids this in signature (4,4) (Section 8.8);
\item in the \textbf{good sector}, where nothing depends on the three extra times $x_5,x_6,x_7$, everything becomes ordinary again: there is a positive \textbf{Fock space} of particles and antiparticles and a Hamiltonian that is not negative;
\item a wave with momentum along an extra time has a non-Hermitian mode Hamiltonian; if that momentum is large enough, $k_5^2+k_6^2+k_7^2\ge m^2+k_0^2+\dots+k_3^2$, the wave \textbf{grows} exponentially (linearly at equality) instead of oscillating (Section 8.9).
\end{enumerate}

The chapter follows §10 of the Stage-1 document \texttt{provenan\allowbreak{}ce/DIRAC\allowbreak{}16COMPLE\allowbreak{}X\_ARBITR\allowbreak{}ARY\_FIEL\allowbreak{}D.md} and its verification records (§11), §17 of the Stage-2 document \texttt{provenan\allowbreak{}ce/DIRAC\allowbreak{}16COMPLE\allowbreak{}X\_PRIMOR\allowbreak{}DIAL\_FIE\allowbreak{}LD.md}, the erratum E1 of \texttt{handoff/\allowbreak{}specs/CO\allowbreak{}NTRACT.m\allowbreak{}d} (§11 there), and the experiment EXP-5 of \texttt{provenan\allowbreak{}ce/DIRAC\allowbreak{}16COMPLE\allowbreak{}X\_DARK\_S\allowbreak{}ECTOR\_NU\allowbreak{}MERICS.m\allowbreak{}d}. The quantization is \textbf{canonical and formal}: no interacting quantum theory, no regularization and no renormalization is constructed (Stage-1 document, §1.3, item 4, and §13, item 6). We use $\hbar=1$, count from 0, and take $x_4$ as the time.

\subsection{States, operators and inner products}

\textbf{States.} In quantum mechanics the state of a system is a vector $f$ in a complex vector space (Chapter 1). In the examples of this chapter the space is finite-dimensional, so a state is a column of complex numbers. An \textbf{inner product} assigns to two states the complex number $\langle f,h\rangle=f^\dagger h=\sum_if_i^\ast h_i$. It is \textbf{positive}: $\langle f,f\rangle=\sum_i|f_i|^2>0$ for every $f\ne0$. A complex vector space with a positive inner product is a \textbf{Hilbert space}. Its role is to give probabilities: for normalized states ($\langle f,f\rangle=1$) the number $|\langle f,h\rangle|^2$ is the probability of finding $h$ in the state $f$, and this needs positivity.

\textbf{Operators.} A (linear) \textbf{operator} is a rule that maps states to states; in finite dimensions it is a matrix. The \textbf{adjoint} $O^\dagger$ of an operator $O$ is defined by $\langle f,Oh\rangle=\langle O^\dagger f,h\rangle$ for all $f,h$; for the inner product above it is the conjugate transpose. An operator is \textbf{Hermitian} (self-adjoint) if $O^\dagger=O$. Then its eigenvalues are real and its \textbf{expectation values} $\langle f,Of\rangle$ are real; measurable quantities (energy, charge) are represented by Hermitian operators. An operator is \textbf{unitary} if $U^\dagger U=1$; unitary operators preserve inner products, and symmetries and time evolution are represented by them.

\textbf{Commutators.} For two operators, $[X,Y]=XY-YX$ is the \textbf{commutator} and $\{X,Y\}=XY+YX$ the \textbf{anticommutator}. We shall use the identity

\[
[XY,Z]=X\{Y,Z\}-\{X,Z\}Y,
\]

which one checks by multiplying out: $X(YZ+ZY)-(XZ+ZX)Y=XYZ-ZXY$.

\textbf{Time evolution.} The \textbf{Hamiltonian} $H$ is the Hermitian operator of the energy. States evolve by the \textbf{Schrödinger equation} $i\,df/dt=Hf$, whose solution is $f(t)=e^{-iHt}f(0)$ (the exponential of a matrix is its power series). Equivalently, one keeps the states fixed and lets the operators evolve, $O(t)=e^{iHt}Oe^{-iHt}$; differentiating gives the \textbf{Heisenberg equation}

\[
\frac{dO}{dt}=i\,[H,O].
\]

\textbf{Canonical quantization in one line.} For a particle with coordinate $q$ and momentum $p$ (Section 5.2) the classical \textbf{Poisson bracket} of two functions $F(q,p)$ and $G(q,p)$ is $\{F,G\}_{\mathrm P}=\frac{\partial F}{\partial q}\frac{\partial G}{\partial p}-\frac{\partial F}{\partial p}\frac{\partial G}{\partial q}$, so $\{q,p\}_{\mathrm P}=1$. Section 5.2 defined $p=\partial L/\partial\dot q$ and the energy $H=\dot q\,p-L$. Regard $\dot q$ as a function of $q$ and $p$ (solve $p=\partial L/\partial\dot q$ for $\dot q$), so that $H(q,p)=\dot q(q,p)\,p-L\bigl(q,\dot q(q,p)\bigr)$. By the chain rule (Chapter 1),

\[
\frac{\partial H}{\partial p}=\dot q+\Bigl(p-\frac{\partial L}{\partial\dot q}\Bigr)\frac{\partial\dot q}{\partial p}=\dot q,\qquad \frac{\partial H}{\partial q}=\Bigl(p-\frac{\partial L}{\partial\dot q}\Bigr)\frac{\partial\dot q}{\partial q}-\frac{\partial L}{\partial q}=-\frac{\partial L}{\partial q}=-\dot p,
\]

because the brackets vanish by the definition of $p$, and the last step is the Euler--Lagrange equation $\dot p=\frac{d}{dt}\frac{\partial L}{\partial\dot q}=\frac{\partial L}{\partial q}$. These are \textbf{Hamilton's equations}, $\dot q=\partial H/\partial p$ and $\dot p=-\partial H/\partial q$. For any function $F(q,p)$ the chain rule then gives

\[
\dot F=\frac{\partial F}{\partial q}\dot q+\frac{\partial F}{\partial p}\dot p=\frac{\partial F}{\partial q}\frac{\partial H}{\partial p}-\frac{\partial F}{\partial p}\frac{\partial H}{\partial q}=\{F,H\}_{\mathrm P}.
\]

(For the first-order Lagrangians of Section 8.5, $p=\partial L/\partial\dot q$ cannot be solved for $\dot q$; that is why Section 8.6 derives the rule in a different way.) The quantization rule is $[q,p]=i$: the Poisson bracket, times $i$, becomes the commutator, and $\dot F=\{F,H\}_{\mathrm P}$ becomes the Heisenberg equation. For fermions the commutator is replaced by the anticommutator. The next two sections show why this gives the Pauli principle.

\subsection{One fermion mode: anticommutators and the Pauli principle}

Let $b$ be an operator with

\[
\{b,b^\dagger\}=1,\qquad b^2=0,\qquad (b^\dagger)^2=0 .
\]

The \textbf{number operator} $N=b^\dagger b$ satisfies $N^2=b^\dagger bb^\dagger b=b^\dagger(1-b^\dagger b)b=b^\dagger b-(b^\dagger)^2b^2=N$. So every eigenvalue $n$ of $N$ satisfies $n^2=n$: $n=0$ or $n=1$. \textbf{A fermion mode is either empty or occupied once}; this is the Pauli principle, and it follows from the anticommutator alone.

\textbf{Worked example.} On the two states $|0\rangle=(1,0)^T$ (empty) and $|1\rangle=(0,1)^T$ (occupied) take

\[
b=\begin{pmatrix}0&1\\0&0\end{pmatrix},\qquad b^\dagger=\begin{pmatrix}0&0\\1&0\end{pmatrix},\qquad bb^\dagger=\begin{pmatrix}1&0\\0&0\end{pmatrix},\qquad b^\dagger b=\begin{pmatrix}0&0\\0&1\end{pmatrix}.
\]

Then $bb^\dagger+b^\dagger b=1$, $b^2=0$, $b|1\rangle=|0\rangle$, $b|0\rangle=0$, $b^\dagger|0\rangle=|1\rangle$ and $b^\dagger|1\rangle=0$: $b$ \textbf{annihilates} a fermion and $b^\dagger$ \textbf{creates} one, and a second one cannot be created. With the Hamiltonian $H=\omega\,b^\dagger b$ the two states have the energies 0 and $\omega$.

\textbf{Where it comes from.} The complex Grassmann oscillator of Section 5.10, $L=\tfrac i2(\theta^\ast\dot\theta-\dot\theta^\ast\theta)-\omega\theta^\ast\theta$, is quantized exactly like this: its quantum version has $\{\theta,\theta^\dagger\}=1$ and $H=\omega\theta^\dagger\theta$ (Section 8.6 derives the rule in general). The classical Grassmann variable is the ``shadow'' of the anticommuting operator.

\subsection{Many fermion modes: Fock space}

With $n$ modes $b_0,\dots,b_{n-1}$ the rules are

\[
\{b_i,b_j^\dagger\}=\delta_{ij},\qquad \{b_i,b_j\}=0,\qquad \{b_i^\dagger,b_j^\dagger\}=0,
\]

where $\delta_{ij}$ is 1 for $i=j$ and 0 otherwise. A \textbf{vacuum} $|0\rangle$ is a normalized state with $b_i|0\rangle=0$ for all $i$. The states $b_{i_1}^\dagger b_{i_2}^\dagger\cdots b_{i_k}^\dagger|0\rangle$ with $i_1<i_2<\dots<i_k$ are the states with the modes $i_1,\dots,i_k$ occupied once each; there are $2^n$ of them (one for each set of occupied modes), and together they span the \textbf{Fock space}.

\textbf{They are orthonormal.} Here $b_i^\dagger$ is the adjoint of $b_i$ for the positive inner product of Section 8.2, so $\langle b_i^\dagger f,h\rangle=\langle f,b_ih\rangle$. Take two such states with the sets of occupied modes $I=\{i_1,\dots,i_k\}$ and $J=\{j_1,\dots,j_l\}$. Moving every creation operator of the first state to the other side turns their inner product into $\langle0|\,b_{i_k}\cdots b_{i_1}\,b_{j_1}^\dagger\cdots b_{j_l}^\dagger|0\rangle$ (we write $\langle f|X|h\rangle$ for $\langle f,Xh\rangle$, and $\langle f|h\rangle$ for $\langle f,h\rangle$). Now move one annihilator $b_i$ to the right through the creation operators with the rules: $b_ib_j^\dagger=-b_j^\dagger b_i$ for $j\ne i$, and $b_ib_i^\dagger=1-b_i^\dagger b_i$. When $b_i$ reaches $|0\rangle$ it gives zero, so only the term with the $1$ survives:

\[
b_i\,b_{j_1}^\dagger\cdots b_{j_l}^\dagger|0\rangle=\begin{cases}\pm\,b_{j_1}^\dagger\cdots(\text{without }b_i^\dagger)\cdots b_{j_l}^\dagger|0\rangle&\text{if }i\in J,\\ 0&\text{if }i\notin J.\end{cases}
\]

Applying $b_{i_1},\dots,b_{i_k}$ in turn therefore gives zero unless every $i$ of $I$ lies in $J$, and otherwise leaves $\pm$ the state of the modes of $J$ that are not in $I$. If one of those is left, the remaining state has the form $b_j^\dagger\varphi$, and its inner product with the vacuum vanishes, because $\langle f,b_j^\dagger\varphi\rangle=\langle b_jf,\varphi\rangle$ and $b_jf=0$ for $f=|0\rangle$. So the inner product is zero unless $I=J$. For $I=J$ each $b_{i_r}$ meets its own $b_{i_r}^\dagger$ first, all signs are $+$, and the result is $\langle0|0\rangle=1$. The simplest case is $\langle0|b_ib_j^\dagger|0\rangle=\langle0|(\delta_{ij}-b_j^\dagger b_i)|0\rangle=\delta_{ij}$. Hence a state $f=\sum_Ic_I\,b_I^\dagger|0\rangle$, with $b_I^\dagger:=b_{i_1}^\dagger\cdots b_{i_k}^\dagger$ for $I=\{i_1<\dots<i_k\}$, has $\langle f,f\rangle=\sum_I|c_I|^2>0$ unless all $c_I$ vanish: the Fock space is a Hilbert space, its inner product is positive.

\textbf{Worked example ($n=2$).} The four states are $|00\rangle=|0\rangle$, $|10\rangle=b_0^\dagger|0\rangle$, $|01\rangle=b_1^\dagger|0\rangle$ and $|11\rangle=b_0^\dagger b_1^\dagger|0\rangle$. Using only the rules:

\[
b_0|11\rangle=b_0b_0^\dagger b_1^\dagger|0\rangle=(1-b_0^\dagger b_0)b_1^\dagger|0\rangle=b_1^\dagger|0\rangle+b_0^\dagger b_1^\dagger b_0|0\rangle=|01\rangle,
\]

\[
b_1|11\rangle=b_1b_0^\dagger b_1^\dagger|0\rangle=-b_0^\dagger b_1b_1^\dagger|0\rangle=-b_0^\dagger(1-b_1^\dagger b_1)|0\rangle=-|10\rangle .
\]

The minus sign in $b_1|11\rangle=-|10\rangle$ is the price of anticommutation. Check: $b_0b_1|11\rangle=-b_0|10\rangle=-|00\rangle$ and $b_1b_0|11\rangle=b_1|01\rangle=|00\rangle$, so $\{b_0,b_1\}|11\rangle=0$ as required.

\textbf{Changing the basis of modes.} If $\Psi_a=\sum_sU_{as}b_s$ with a unitary matrix $U$ ($UU^\dagger=1$), then $\{\Psi_a,\Psi_b^\dagger\}=\sum_{s,s'}U_{as}U^\ast_{bs'}\{b_s,b_{s'}^\dagger\}=\sum_sU_{as}U^\ast_{bs}=(UU^\dagger)_{ab}=\delta_{ab}$. Any orthonormal choice of modes gives the same standard anticommutator. This is used in Section 8.10.

\subsection{The Hamiltonian form of the dirac16complex Lagrangian}

\textbf{Slices.} We quantize with respect to $x_4$. The \textbf{slices} $x_4=\text{const}$ are seven-dimensional, with the coordinates $x_0,x_1,x_2,x_3,x_5,x_6,x_7$. The procedure needs $g^{44}\ne0$: the slices must not be ``characteristic'' (Stage-1 document, §10.1). We write $\gamma^{x_4}$ for the curved gamma $\gamma^\mu$ with $\mu=4$, to distinguish it from the constant frame matrix $\gamma^4$.

\textbf{Removing the time derivative of $\Psi^\dagger$.} In $\mathcal L$ the time derivatives appear in the combination $\tfrac12\sqrt{|g|}\,\bigl(\bar\Psi\gamma^{x_4}\partial_4\Psi-\partial_4\bar\Psi\,\gamma^{x_4}\Psi\bigr)$. Add the total derivative $\tfrac12\partial_4\bigl(\sqrt{|g|}\,\bar\Psi\gamma^{x_4}\Psi\bigr)$, which by the product rule equals $\tfrac12\sqrt{|g|}\,\partial_4\bar\Psi\,\gamma^{x_4}\Psi+\tfrac12\sqrt{|g|}\,\bar\Psi\gamma^{x_4}\partial_4\Psi+\tfrac12\bar\Psi\,\partial_4(\sqrt{|g|}\gamma^{x_4})\Psi$. The terms with $\partial_4\bar\Psi$ cancel, and

\[
\mathcal L'=\mathcal L+\tfrac12\partial_4\bigl(\sqrt{|g|}\,\bar\Psi\gamma^{x_4}\Psi\bigr)=\Psi^\dagger K\,\partial_4\Psi-\mathcal H,\qquad K:=\sqrt{|g|}\,C\gamma^{x_4},
\]

where $\mathcal H$ collects (with a minus sign) all terms without a time derivative of $\Psi$ or $\Psi^\dagger$. By Section 5.4, $\mathcal L'$ has the same field equations as $\mathcal L$. It has the form of the first-order Lagrangians of Section 5.6: linear in the time derivative, with a matrix $K$ in front.

\textbf{Momenta and constraints.} The momentum conjugate to $\Psi_a$ (right derivative, Section 5.10) is

\[
\Pi_a=\frac{\partial_R\mathcal L'}{\partial(\partial_4\Psi_a)}=\sqrt{|g|}\,\bigl(\Psi^\dagger C\gamma^{x_4}\bigr)_a ,
\]

and $\mathcal L'$ does not contain $\partial_4\Psi^\dagger$ at all. As in Example A of Section 5.6, the momenta are fixed functions of the fields. In Dirac's theory of constrained systems the relations $\Pi_a-\sqrt{|g|}(\Psi^\dagger C\gamma^{x_4})_a\approx0$ and $\Pi_{\Psi^\dagger,a}\approx0$ are called \textbf{primary constraints}. Their brackets with each other form, up to sign and transposition, the matrix $K$ (Stage-1 document, §10.2). The canonical momentum formula was verified in a genuine Grassmann algebra at all G1 and G2 test points (checks \texttt{QNT\_cano\allowbreak{}nicalMom\allowbreak{}entum\_G1} and \texttt{QNT\_cano\allowbreak{}nicalMom\allowbreak{}entum\_G2}, \texttt{wolfram-\allowbreak{}geometry\allowbreak{}-report.\allowbreak{}json}).

\textbf{$K$ is invertible exactly when $g^{44}\ne0$.} The curved gamma satisfies $(\gamma^{x_4})^2=\tfrac12\{\gamma^{x_4},\gamma^{x_4}\}=g^{44}$ (the curved Clifford relation of Chapter 4), and $C^2=1$. Hence

\[
(C\gamma^{x_4})(\gamma^{x_4}C)=C\,g^{44}\,C=g^{44}\,I_{16},\qquad (\gamma^{x_4}C)(C\gamma^{x_4})=g^{44}\,I_{16},\qquad (C\gamma^{x_4})^{-1}=\frac{\gamma^{x_4}C}{g^{44}} .
\]

The constraints are therefore of the kind Dirac calls \textbf{second class}: they can be solved, and the resulting bracket (the \textbf{Dirac bracket}) is built from $K^{-1}$. The identity above was verified symbolically for an arbitrary $e_a{}^4$ and at the G1 points (check \texttt{QNT\_curv\allowbreak{}edAntico\allowbreak{}mmutator\allowbreak{}Matrix}, \texttt{wolfram-\allowbreak{}algebra-\allowbreak{}report.j\allowbreak{}son}). The graded constraint computation itself is a derivation, not a separate machine check (Stage-1 document, §10.2).

\textbf{$K$ is anti-Hermitian.} $K=\sqrt{|g|}\,e_a{}^4\,C\gamma^a$ is a real combination of the real antisymmetric matrices $C\gamma^a$ (expression [1], Section 5.12), so $K^T=-K$ and, being real, $K^\dagger=-K$. This is what makes the anticommutator matrix $iK^{-1}$ of Section 8.6 Hermitian. (When $K$ depends on $x_4$, the total derivative added above also puts the derivative-free term $-\tfrac12\bar\Psi\,\partial_4(\sqrt{|g|}\gamma^{x_4})\Psi$ into $\mathcal H$; the Stage-1 document, §9.5 and §10.4, discusses it. It vanishes in flat space and in the primordial field.)

\textbf{The Hamiltonian in flat space.} In flat space ($g=\eta$, $\sqrt{|g|}=1$, $\Omega_\mu=0$, $\gamma^{x_4}=\gamma^4$) with $U=0$,

\[
\mathcal H=-\tfrac12\sum_{j\ne4}\bigl(\bar\Psi\gamma^j\partial_j\Psi-\partial_j\bar\Psi\,\gamma^j\Psi\bigr)+m\,\bar\Psi\Psi .
\]

Integrated over a slice, the term $-\partial_j\bar\Psi\,\gamma^j\Psi$ can be moved by parts, $-\int\partial_j\bar\Psi\,\gamma^j\Psi\,d^7x=+\int\bar\Psi\gamma^j\partial_j\Psi\,d^7x$ (the fields vanish far away), so the Hamiltonian is

\[
H=\int\Psi^\dagger\Bigl(mC-\sum_{j\ne4}C\gamma^j\partial_j\Bigr)\Psi\,d^7x .
\]

In a curved field $\mathcal H$ contains in addition the spin-connection terms; Chapter 7 and the Stage-1 document (§10.4) give its general form. The canonical structure, that is $K$, does not depend on the potential $U$, which contains no derivatives.

\subsection{The equal-time anticommutator}

\textbf{Result.} For two points $x$ and $y$ on the same slice,

\[
\begin{aligned}
&\bigl\{\Psi_a(x),\Psi^\dagger_b(y)\bigr\}=i\,\bigl[K^{-1}\bigr]_{ab}\,\delta^7(x-y)=i\,\frac{[\gamma^{x_4}C]_{ab}}{g^{44}}\,\frac{\delta^7(x-y)}{\sqrt{|g|}},\\
&\{\Psi_a(x),\Psi_b(y)\}=\{\Psi^\dagger_a(x),\Psi^\dagger_b(y)\}=0 .
\end{aligned}
\]

Here $\delta^7(x-y)$ is the seven-dimensional \textbf{Dirac delta}: the idealization of a very narrow bump of total integral 1, defined by $\int f(y)\,\delta^7(x-y)\,d^7y=f(x)$. On a lattice of points with cell volume $v$ it would be $\delta_{xy}/v$. The matrix $iK^{-1}$ is Hermitian, because $K^\dagger=-K$ gives $(iK^{-1})^\dagger=-i(K^\dagger)^{-1}=iK^{-1}$.

\textbf{Derivation.} We derive the rule from the requirement that the quantum Heisenberg equation reproduce the classical field equation. To keep every step visible, replace the slice by finitely many points, so that $\Psi$ is a column of $N$ odd operators, and take a Hamiltonian that is bilinear, $H=\Psi^\dagger h\Psi=\sum_{a,b}\Psi_a^\dagger h_{ab}\Psi_b$, with a numerical matrix $h$. The classical Lagrangian is $L'=\Psi^\dagger K\dot\Psi-\Psi^\dagger h\Psi$, and its Euler--Lagrange equation for $\Psi^\dagger$ (left derivative) is $K\dot\Psi=h\Psi$, that is

\[
\dot\Psi=K^{-1}h\Psi\qquad\text{(classical)}.
\]

Now suppose $\{\Psi_a,\Psi_b^\dagger\}=A_{ab}$ and $\{\Psi_a,\Psi_b\}=0$ for some matrix $A$. By the identity of Section 8.2,

\[
\begin{aligned}
&[\Psi_a^\dagger\Psi_b,\Psi_c]=\Psi_a^\dagger\{\Psi_b,\Psi_c\}-\{\Psi_a^\dagger,\Psi_c\}\Psi_b=-A_{ca}\Psi_b,\\
&[H,\Psi_c]=-\sum_{a,b}A_{ca}h_{ab}\Psi_b=-(Ah\Psi)_c .
\end{aligned}
\]

The Heisenberg equation gives $\dot\Psi=i[H,\Psi]=-iAh\Psi$. Agreement with the classical equation for every $h$ (for instance for every value of the mass) requires $-iA=K^{-1}$, that is

\[
\{\Psi_a,\Psi_b^\dagger\}=A_{ab}=i\,[K^{-1}]_{ab} .
\]

For the field the sums over points become integrals over the slice, $h$ becomes the differential operator of $H$, and $\delta_{ab}$ between points becomes $\delta^7(x-y)$; the computation is otherwise the same. Dirac's bracket method gives the same result (Stage-1 document, §10.2 and §10.3). The same computation with commutators instead of anticommutators would also reproduce the equation of motion; the choice of anticommutators is part of the definition of dirac16complex as a fermion field (Grassmann components, Section 5.9), and it is what gives the Pauli principle of Section 8.3.

\textbf{Check with the Grassmann oscillator.} For $L=\tfrac i2(\theta^\ast\dot\theta-\dot\theta^\ast\theta)-\omega\theta^\ast\theta$, adding $\tfrac i2\frac{d}{dt}(\theta^\ast\theta)$ gives $L'=i\theta^\ast\dot\theta-\omega\theta^\ast\theta$, so $K=i$ and $\{\theta,\theta^\dagger\}=i\cdot\tfrac1i=1$: the rule of Section 8.3.

\textbf{Gaussian normal gauge and the matrix $B$.} In the frames used throughout the book the time is ``Gaussian normal'': $g^{44}=-1$ and $\gamma^{x_4}=\gamma^4$ (for example in the primordial field of Chapter 9, where the vielbein component along $x_4$ is 1). Then $iK^{-1}=i(C\gamma^4)^{-1}/\sqrt{|g|}=-i\gamma^4C/\sqrt{|g|}$, and since $\gamma^4$ anticommutes with all four factors of $C=\gamma^0\gamma^1\gamma^2\gamma^3$, it commutes with $C$. Therefore

\[
\bigl\{\Psi_a(x),\Psi^\dagger_b(y)\bigr\}=B_{ab}\,\frac{\delta^7(x-y)}{\sqrt{|g|}},\qquad B:=-iC\gamma^4 .
\]

In the primordial field $\sqrt{|g|}=\cos z$, and this is the anticommutator $B\,\delta^7/\cos z$ of the Stage-2 document (§17.1).

\textbf{Properties of $B$.} Each follows from $C^2=1$, $(\gamma^4)^2=-1$, $C\gamma^4=\gamma^4C$, $C^T=C$ and $(\gamma^4)^T=-\gamma^4$:

\begin{itemize}
\item $B^\dagger=i(\gamma^4)^TC^T=-i\gamma^4C=-iC\gamma^4=B$: \textbf{Hermitian};
\item $B^2=(-i)^2C\gamma^4C\gamma^4=-C^2(\gamma^4)^2=1$;
\item $\mathrm{tr}\,B=0$: $C\gamma^4=\gamma^0\gamma^1\gamma^2\gamma^3\gamma^4$ is a product of an odd number of gammas, and such a product $X$ anticommutes with the chirality $\gamma^8$ of Chapter 2, so $\mathrm{tr}X=\mathrm{tr}(\gamma^8\gamma^8X)=\mathrm{tr}(\gamma^8X\gamma^8)=-\mathrm{tr}X$ (cyclicity of the trace and $(\gamma^8)^2=1$);
\item so the eigenvalues of $B$ are $+1$ and $-1$ (from $B^2=1$), with multiplicity 8 each (from $\mathrm{tr}B=0$);
\item $[C,B]=0$ and $BC=-iC\gamma^4C=-i\gamma^4$.
\end{itemize}

These are Result 10.2 of the Stage-1 document (check \texttt{ALG\_char\allowbreak{}geFormB}, both algebra reports; the characteristic polynomial $(x-1)^8(x+1)^8$ is recorded in \texttt{wolfram-\allowbreak{}algebra-\allowbreak{}report.j\allowbreak{}son}). $B$ is $i$ times a signed permutation; the Student Guide (§6.4) tabulates it, for example $(Bu)_0=i\,u_9$.

\textbf{Heisenberg equation in flat space.} With $H$ of Section 8.5 and $\{\Psi_a,\Psi^\dagger_b\}=B_{ab}\delta^7$ the computation above gives $\partial_4\Psi=-iB\bigl(mC-\sum_jC\gamma^j\partial_j\bigr)\Psi$. With $BC=-i\gamma^4$ this is

\[
\partial_4\Psi=-m\gamma^4\Psi+\sum_{j\ne4}\gamma^4\gamma^j\partial_j\Psi ,
\]

which is the field equation $\gamma^\mu\partial_\mu\Psi=m\Psi$ multiplied from the left by $\gamma^4$ (use $(\gamma^4)^2=-1$). The quantum theory reproduces the classical equation, as it was built to do.

\subsection{Why the anticommutator forces an indefinite inner product}

\textbf{Theorem 8.1.} There is no Hilbert space (positive inner product) on which the operators $\Psi_a$ act with $\Psi_a^\dagger$ equal to their Hilbert adjoints and with $\{\Psi_a(x),\Psi_b^\dagger(y)\}=B_{ab}\delta^7(x-y)/\sqrt{|g|}$.

\emph{Proof.} Suppose there were. Take a numerical column $v$ with $Bv=-v$ (it exists: $B$ has the eigenvalue $-1$) and an ordinary bump function $f$ on the slice, and form the smeared operator $\Psi[f,v]:=\sum_av_a^\ast\int f(x)\Psi_a(x)\,d^7x$. Its adjoint is $\Psi[f,v]^\dagger=\sum_bv_b\int f(y)\Psi_b^\dagger(y)\,d^7y$, and

\[
\{\Psi[f,v],\Psi[f,v]^\dagger\}=\sum_{a,b}v_a^\ast B_{ab}v_b\int\frac{f(x)^2}{\sqrt{|g|}}\,d^7x=-\,v^\dagger v\int\frac{f^2}{\sqrt{|g|}}\,d^7x<0 .
\]

But for any operator $O$ on a Hilbert space and any state $\varphi$, $\langle\varphi,\{O,O^\dagger\}\varphi\rangle=\langle O^\dagger\varphi,O^\dagger\varphi\rangle+\langle O\varphi,O\varphi\rangle\ge0$. A negative multiple of the identity cannot satisfy this. $\square$

\textbf{The way out: a Krein space.} A \textbf{Hermitian form} on a space of columns is a rule $[f,h]=f^\dagger Gh$ with a Hermitian matrix $G$; then $[h,f]=h^\dagger Gf=(f^\dagger G^\dagger h)^\ast=[f,h]^\ast$, so $[f,f]$ is always real. It is \textbf{non-degenerate} if no nonzero $f$ has $[f,h]=0$ for all $h$; since $f^\dagger Gh=0$ for all $h$ means $f^\dagger G=0$, that is $Gf=0$ (take the adjoint and use $G^\dagger=G$), this says that $Gf=0$ only for $f=0$. A \textbf{Krein space} is a complex vector space with a Hermitian form $[f,h]$ that is non-degenerate but \textbf{indefinite} (some vectors have $[f,f]<0$), together with a \textbf{fundamental symmetry} $J$, an operator with $J=J^\dagger=J^{-1}$, such that $[f,Jh]$ is a positive inner product. The canonical anticommutator defines exactly such a structure on the one-particle wave functions: the form is $[f,h]=f^\dagger Bh$, of signature (8,8) and non-degenerate ($Bf=0$ gives $f=B^2f=0$), and $J=B$ is a fundamental symmetry, because $B$ is Hermitian with $B^2=1$ and $[f,Bh]=f^\dagger B^2h=f^\dagger h$ is the standard positive product (Stage-1 document, §10.5).

\textbf{Hilbert adjoint and Krein adjoint.} In a representation on a positive Hilbert space, define

\[
\chi:=\Psi^\dagger B,\qquad\text{i.e.}\qquad \chi_b=\sum_a\Psi_a^\dagger B_{ab}.
\]

Then, in flat space and Gaussian normal gauge, $\{\Psi_a,\chi_b\}=\sum_c\{\Psi_a,\Psi_c^\dagger\}B_{cb}\,=(B^2)_{ab}\,\delta^7=\delta_{ab}\delta^7$, the standard positive anticommutator: $\chi$ is the \textbf{Hilbert adjoint} of $\Psi$. The canonical $\Psi^\dagger=\chi B$ is not the Hilbert adjoint; it is the adjoint with respect to the indefinite form, the \textbf{Krein adjoint}. Every bilinear $\Psi^\dagger M\Psi$ of the classical theory becomes $\chi BM\Psi$ in the positive representation; Section 8.12 draws the consequences.

\textbf{Worked example (a two-mode Krein toy).} Take two modes with $B=\mathrm{diag}(+1,-1)$: $\{\psi_0,\psi_0^\dagger\}=1$ and $\{\psi_1,\psi_1^\dagger\}=-1$. In a positive representation, $\chi_0=\psi_0^\dagger$ and $\chi_1=-\psi_1^\dagger$ are the Hilbert adjoints, with $\{\psi_1,\chi_1\}=1$. The operator $\psi_1^\dagger\psi_1=-\chi_1\psi_1$ has the eigenvalues 0 and $-1$ (Section 8.3 with $b=\psi_1$): the second mode carries negative ``norm''. The $16\times16$ matrix $B$ contains eight such negative directions.

\subsection{No Spin(4,4)-invariant positive density in signature (4,4)}

One might hope that a cleverer choice of ``charge density'' would be positive. In ordinary Dirac theory with one time direction it is: the density $\bar\psi\gamma^0\psi=\psi^\dagger\psi$ is positive. The Stage-1 document proves that in signature (4,4) no such choice built from a Spin(4,4)-invariant form exists (its Theorem 10.1). We give the argument in steps.

\textbf{Step 1 (invariant forms).} A matrix $G$ defines a \textbf{Spin(4,4)-invariant bilinear form} $\Psi^TG\Phi$ if it does not change under the infinitesimal spin rotations $\Psi\to(1+\epsilon S^{ab})\Psi$, $\Phi\to(1+\epsilon S^{ab})\Phi$, that is if $(S^{ab})^TG+GS^{ab}=0$ for all $a,b$. The charge matrix $C$ is one: $(S^{ab})^TC=-CS^{ab}C\,C=-CS^{ab}$ by Section 5.12. The chiral projectors $P_\mp=\tfrac12(1\mp\gamma^8)$ of Chapter 2 commute with every $S^{ab}$ and with $C$, so $CP_-$ and $CP_+$ are invariant as well. The repository computed the space of all invariant bilinear forms exactly: it is 2-dimensional and spanned by $CP_-$ and $CP_+$ (Result 10.3; check \texttt{ALG\_inva\allowbreak{}riantFor\allowbreak{}ms}, both algebra reports).

A density needs a \textbf{Hermitian} form $\Psi^\dagger G\Phi$ (Section 8.7), which changes under the same rotation by $\epsilon\,\Psi^\dagger\bigl((S^{ab})^\dagger G+GS^{ab}\bigr)\Phi$ to first order. Because the $S^{ab}$ are real, $(S^{ab})^\dagger=(S^{ab})^T$, so the invariance condition $(S^{ab})^\dagger G+GS^{ab}=0$ is exactly the condition above, and the invariant Hermitian forms are found among $G=r_-CP_-+r_+CP_+$ with complex numbers $r_\mp$. Which of these are Hermitian? $C$ and $\gamma^8=\mathrm{diag}(-I_8,I_8)$ (Section 2.10) are real symmetric and commute, so $P_\mp$ is Hermitian and $(CP_\mp)^\dagger=P_\mp C=CP_\mp$. Hence $G^\dagger=\bar r_-CP_-+\bar r_+CP_+$. The two matrices $CP_-$ and $CP_+$ are linearly independent: if $x\,CP_-+y\,CP_+=0$, multiplying from the right by $P_-$ (with $P_-^2=P_-$ and $P_+P_-=0$) gives $x\,CP_-=0$, hence $x=0$ because $C$ is invertible and $P_-\ne0$, and in the same way $y=0$. So $G^\dagger-G=(\bar r_--r_-)CP_-+(\bar r_+-r_+)CP_+$ vanishes exactly when $r_-$ and $r_+$ are real. So every invariant Hermitian form is $G=r_-CP_-+r_+CP_+$ with real $r_-,r_+$. Such a $G$ is \textbf{even}: it commutes with $\gamma^8$, because $C$ and $P_\mp$ do.

\textbf{Step 2 (the density is odd).} The time component of a current built with $G$ has the matrix $G\gamma^4$, or more generally the Hermitian part $X=\tfrac12\bigl(cG\gamma^4+(cG\gamma^4)^\dagger\bigr)$ of $cG\gamma^4$ for some constant $c$. The frame matrix $\gamma^4$ anticommutes with $\gamma^8$, and $G$ commutes with it, so $\gamma^8(cG\gamma^4)\gamma^8=-cG\gamma^4$; taking adjoints (with $\gamma^8$ Hermitian) gives the same for $(cG\gamma^4)^\dagger$. Hence

\[
\gamma^8X\gamma^8=-X .
\]

\textbf{Step 3 (equal numbers of signs).} If $Xv=\lambda v$ then $X(\gamma^8v)=-\gamma^8Xv=-\lambda\,\gamma^8v$. So $\gamma^8$ maps the eigenvectors of $\lambda$ onto eigenvectors of $-\lambda$, one to one, and the Hermitian matrix $X$ has as many positive as negative eigenvalues. Unless $X=0$, the density $\Psi^\dagger X\Psi$ is indefinite. The repository shows more: $X^2=|k|^2I_{16}$ with $k=\tfrac12(\bar c\,r_+-c\,r_-)$, so $X$ has signature (8,8) whenever it is not zero (\texttt{ALG\_inva\allowbreak{}riantFor\allowbreak{}ms}).

The hypothesis of invariance matters. The density $\Psi^\dagger\Psi$ is positive, but its matrix $1$ commutes with $\gamma^8$, so by Step 2 it is not the time component of a current built from an invariant form.

The reason is the one stated in the Stage-1 document: every invariant form is even in chirality while the time-like $\gamma^4$ is odd. In four dimensions with one time the invariant form is itself odd ($\gamma^0$), and the product with $\gamma^0$ is the positive identity. The indefinite (Krein) structure of dirac16complex is therefore not an accident of the choice $B$: in signature (4,4) no density built from a Spin(4,4)-invariant form avoids it.

\textbf{Positive energy does not help.} One might restrict to states of positive energy. At rest in flat space the mode Hamiltonian of Section 8.9 is $h_0=m(-i\gamma^4)$, so for $m>0$ the positive-energy space $V_+$ is the eigenspace of $-i\gamma^4$ with eigenvalue $+1$. The matrix $-i\gamma^4$ is Hermitian, squares to 1 and is traceless, so $V_+$ has dimension 8. On $V_+$ the matrix $B=C(-i\gamma^4)$ acts as $C$, which maps $V_+$ to itself (it commutes with $\gamma^4$) and squares to 1. Its trace on $V_+$ is $\mathrm{tr}(C\Lambda_+)$ with the projector $\Lambda_+=\tfrac12(1-i\gamma^4)$ onto $V_+$ (the energy projector of Section 8.9 at rest): $\mathrm{tr}(C\Lambda_+)=\tfrac12\mathrm{tr}C-\tfrac i2\mathrm{tr}(C\gamma^4)=0$, because $C$ is a product of four anticommuting gammas ($\mathrm{tr}\,\gamma^0\gamma^1\gamma^2\gamma^3=\mathrm{tr}\,\gamma^1\gamma^2\gamma^3\gamma^0=-\mathrm{tr}\,\gamma^0\gamma^1\gamma^2\gamma^3$) and $C\gamma^4$ is odd. So the Krein form $u^\dagger Bu$ has signature (4,4) even on the positive-energy states at rest. The repository finds (4,4) on both energy eigenspaces at rest and also for moving states in the good sector, for example at $m=1$, $k=(1,1,2,3)$, $E=4$ (Result 10.4; check \texttt{QNT\_krei\allowbreak{}nSignatu\allowbreak{}re}, both algebra reports).

\textbf{Worked example with explicit vectors.} Write $e_n$ for the column with 1 in position $n$ and 0 elsewhere, and use the tables of the Student Guide (§6.3, §6.4): $(\gamma^4u)_0=-u_{13}$, $(\gamma^4u)_4=u_9$, $(\gamma^4u)_9=-u_4$, $(\gamma^4u)_{13}=u_0$, and $(Cu)_0=-u_4$, $(Cu)_4=-u_0$, $(Cu)_9=u_{13}$, $(Cu)_{13}=u_9$. Consider

\[
u_+=\tfrac12\bigl(e_0-e_4-i\,e_9-i\,e_{13}\bigr),\qquad u_-=\tfrac12\bigl(e_0+e_4+i\,e_9-i\,e_{13}\bigr).
\]

For $u_+$: $(\gamma^4u_+)_0=-(-\tfrac i2)=\tfrac i2$, so $(-i\gamma^4u_+)_0=\tfrac12$; in the same way components 4, 9 and 13 reproduce $-\tfrac12$, $-\tfrac i2$, $-\tfrac i2$, so $-i\gamma^4u_+=u_+$. And $(Cu_+)_0=-(-\tfrac12)=\tfrac12$, $(Cu_+)_4=-\tfrac12$, $(Cu_+)_9=-\tfrac i2$, $(Cu_+)_{13}=-\tfrac i2$, so $Cu_+=u_+$ and $Bu_+=u_+$. For $u_-$ the same bookkeeping gives $-i\gamma^4u_-=u_-$ but $Cu_-=-u_-$, so $Bu_-=-u_-$. Both are normalized positive-energy states at rest ($u^\dagger u=1$), and $u_+^\dagger Bu_+=+1$ while $u_-^\dagger Bu_-=-1$. ($u_+$ is the standard rest state of the Stage-3 programs, called $u_0$ in the Student Guide, §6.10.)

\subsection{Flat-space modes and the mode Hamiltonian}

\textbf{Schrödinger form.} In flat space with $U=0$ the field equation is $\gamma^4\partial_4\Psi+\sum_{j\ne4}\gamma^j\partial_j\Psi=m\Psi$. Multiply from the left by $\gamma^4$ and use $(\gamma^4)^2=-1$: $-\partial_4\Psi+\sum_j\gamma^4\gamma^j\partial_j\Psi=m\gamma^4\Psi$. For a plane wave $\Psi=e^{i\sum_jk_jx_j}u(x_4)$ with real wave numbers $k_j$ ($j\in\{0,1,2,3,5,6,7\}$) each $\partial_j$ becomes $ik_j$, and multiplying by $i$ gives

\[
i\,\partial_4u=h_k\,u,\qquad h_k=-im\gamma^4-\gamma^4\sum_{j\ne4}k_j\gamma^j .
\]

$h_k$ is the \textbf{mode Hamiltonian}, a $16\times16$ complex matrix.

\textbf{Its square.} Three facts follow from the Clifford relations: $(-im\gamma^4)^2=-m^2(\gamma^4)^2=m^2$; $(\gamma^4\gamma^j)^2=\gamma^4\gamma^j\gamma^4\gamma^j=-(\gamma^4)^2(\gamma^j)^2=\eta^{jj}$ for $j\ne4$; and all cross terms cancel, because $\{\gamma^4,\gamma^4\gamma^j\}=(\gamma^4)^2\gamma^j+\gamma^4\gamma^j\gamma^4=-\gamma^j+\gamma^j=0$ and, for $i\ne j$, $\{\gamma^4\gamma^i,\gamma^4\gamma^j\}=-(\gamma^4)^2(\gamma^i\gamma^j+\gamma^j\gamma^i)=0$. Hence

\[
h_k^2=E^2\,I_{16},\qquad E^2=m^2+k_0^2+k_1^2+k_2^2+k_3^2-k_5^2-k_6^2-k_7^2 .
\]

\textbf{Hermitian or not.} $(-i\gamma^4)^\dagger=i(\gamma^4)^T=-i\gamma^4$, so the mass term is Hermitian. $(\gamma^4\gamma^j)^\dagger=(\gamma^j)^T(\gamma^4)^T=-(\gamma^j)^T\gamma^4$. For $j\le3$, $\gamma^j$ is symmetric and this is $-\gamma^j\gamma^4=\gamma^4\gamma^j$: \textbf{Hermitian}. For $j\ge5$, $\gamma^j$ is antisymmetric and this is $\gamma^j\gamma^4=-\gamma^4\gamma^j$: \textbf{anti-Hermitian}. So the anti-Hermitian part of $h_k$ is $-\gamma^4(k_5\gamma^5+k_6\gamma^6+k_7\gamma^7)$, whose square is $-(k_5^2+k_6^2+k_7^2)I_{16}$ by the same cross-term cancellation; it vanishes only when $k_5=k_6=k_7=0$. \textbf{$h_k$ is Hermitian if and only if there is no momentum along the extra times.} The same split shows how $h_k$ behaves with respect to $B\propto\gamma^0\gamma^1\gamma^2\gamma^3\gamma^4$: the Hermitian part commutes with $B$ and the anti-Hermitian part anticommutes with it (a gamma with index $\le4$ commutes with this product of five gammas, one with index $\ge5$ anticommutes), which gives $[h_k,B]=2i\sum_{j=5}^7k_jC\gamma^j$ (Stage-1 document, §10.6).

\textbf{The Krein norm is always conserved.} Write $h_k=h_H+h_A$ with the Hermitian part $h_H$ (commuting with $B$) and the anti-Hermitian part $h_A$ (anticommuting with $B$). Then $h_k^\dagger B=(h_H-h_A)B=B(h_H+h_A)=Bh_k$. For a solution of $i\dot u=h_ku$ (so $\dot u^\dagger=iu^\dagger h_k^\dagger$),

\[
\frac{d}{dx_4}\bigl(u^\dagger Bu\bigr)=i\,u^\dagger\bigl(h_k^\dagger B-Bh_k\bigr)u=0,\qquad \frac{d}{dx_4}\bigl(u^\dagger u\bigr)=i\,u^\dagger\bigl(h_k^\dagger-h_k\bigr)u=-2i\,u^\dagger h_Au .
\]

The Krein norm $u^\dagger Bu$ never changes; the Hilbert norm $u^\dagger u$ changes whenever there is extra-time momentum.

\textbf{Samples from the repository.} These identities hold for symbolic $m$ and $k$ (Wolfram) and by a structural proof for all real $m,k$ plus ten exact samples (Python) (Result 10.5; check \texttt{QNT\_flat\allowbreak{}ModeHami\allowbreak{}ltonian}, both algebra reports). Some of the samples of \texttt{artifact\allowbreak{}s/dirac1\allowbreak{}6complex\allowbreak{}/arbitra\allowbreak{}ry-field\allowbreak{}/python-\allowbreak{}algebra-\allowbreak{}report.j\allowbreak{}son}:

\begingroup
\small
\begin{longtable}{@{}>{\raggedright\arraybackslash}p{0.200\linewidth}>{\raggedright\arraybackslash}p{0.200\linewidth}>{\raggedright\arraybackslash}p{0.200\linewidth}>{\raggedright\arraybackslash}p{0.200\linewidth}@{}}
\toprule
\textbf{$m$} & \textbf{nonzero wave numbers} & \textbf{Hermitian} & \textbf{$E^2$} \\
\midrule
\endfirsthead
\toprule
\textbf{$m$} & \textbf{nonzero wave numbers} & \textbf{Hermitian} & \textbf{$E^2$} \\
\midrule
\endhead
1 & none & yes & 1 \\
3/2 & $k_0=1/2$, $k_1=-1/3$, $k_2=2$ & yes & 119/18 \\
1 & $k_0=1/2$, $k_6=-2/3$ & no & 29/36 \\
1 & $k_5=1$ & no & 0 \\
1 & $k_7=2$ & no & $-3$ \\
1/2 & $k_0,\dots,k_3=\tfrac17,\tfrac27,\tfrac37,\tfrac47$; $k_5,k_6,k_7=\tfrac57,\tfrac67,1$ & no & $-271/196$ \\
0 & $k_1=3$, $k_5=3$ & no & 0 \\
\bottomrule
\end{longtable}
\endgroup

For instance $\tfrac14+\tfrac{1+4+9+16}{49}-\tfrac{25+36+49}{49}=\tfrac14-\tfrac{80}{49}=\tfrac{49-320}{196}=-\tfrac{271}{196}$.

\textbf{The three cases.} Because $h_k^2=E^2$, the power series of the exponential collapses:

\[
e^{-ih_kx_4}=\cos(Ex_4)-i\,\frac{\sin(Ex_4)}{E}\,h_k .
\]

(The even powers give $\cos$, the odd powers give $h_k$ times $\sin/E$.) Three cases occur (Stage-1 document, §10.10).

\begin{itemize}
\item $E^2>0$: the frequencies $\pm E$ are real and the modes oscillate, even when $h_k$ is not Hermitian (the sample $E^2=29/36$). $h_k$ is then diagonalizable, with the \textbf{energy projectors} $\Lambda_\pm=\tfrac12(1\pm h_k/E)$ (not the chiral projectors $P_\mp$ of Section 8.8): one checks $\Lambda_\pm^2=\Lambda_\pm$, $\Lambda_++\Lambda_-=1$ and $h_k\Lambda_\pm=\pm E\Lambda_\pm$ from $h_k^2=E^2$.
\item $E^2=0$ and $h_k\ne0$: $h_k$ is \textbf{nilpotent}, $h_k^2=0$, and $e^{-ih_kx_4}=1-ih_kx_4$: the mode \textbf{grows linearly} (the samples $m=1$, $k_5=1$ and $m=0$, $k_1=k_5=3$).
\item $E^2<0$, that is exactly when $k_5^2+k_6^2+k_7^2>m^2+k_0^2+k_1^2+k_2^2+k_3^2$: with $E=i\varkappa$, $\varkappa=\sqrt{-E^2}$ (the letter $\varkappa$ is a variant of the Greek kappa; this book uses it for growth rates, so that it is not confused with the gravitational coupling $\kappa$ of Chapter 4), one has $\cos(i\varkappa x_4)=\cosh(\varkappa x_4)$ and $\sin(i\varkappa x_4)/(i\varkappa)=\sinh(\varkappa x_4)/\varkappa$, so $e^{-ih_kx_4}=\cosh(\varkappa x_4)-i\,h_k\sinh(\varkappa x_4)/\varkappa$: the mode \textbf{grows exponentially}, like $e^{\varkappa x_4}$. Two of the ten samples are of this kind.
\end{itemize}

This is the quantum version of the classical observation of Section 5.5, where a plane wave with enough extra-time momentum grows exponentially. It means that an equation with several time directions does not have a well-posed initial-value problem on the slices $x_4=\text{const}$, because those slices themselves contain three time-like directions (the \textbf{ultrahyperbolic} problem). An initial-value problem is \textbf{well-posed} (Hadamard's definition) if for given data on the slice a solution exists, is unique, and depends continuously on the data: a small enough change of the data changes the solution at a later time only a little. The continuity fails here. Keep $m$ fixed, let only $k_5$ be nonzero, with $|k_5|>m$, so that $\varkappa=\sqrt{k_5^2-m^2}$, and choose a column $u$ with $u^\dagger u=1$ and $-ih_ku=\varkappa u$. Such columns exist: $(-ih_k)^2=-h_k^2=\varkappa^2$, so the matrix $\Lambda_+=\tfrac12(1+h_k/E)$ of the first case, taken at $E=i\varkappa$, that is $\Lambda_+=\tfrac12(1-ih_k/\varkappa)$, satisfies $\Lambda_+^2=\Lambda_+$ and $-ih_k\Lambda_+=\varkappa\Lambda_+$, and $\Lambda_+\ne0$ because $h_k$ is traceless and therefore not $-i\varkappa$ times the identity; take $u$ proportional to a nonzero column of $\Lambda_+$. The formula of the third case gives $e^{-ih_kx_4}u=\bigl(\cosh(\varkappa x_4)+\sinh(\varkappa x_4)\bigr)u=e^{\varkappa x_4}u$. So a change of the data by $\epsilon\,e^{ik_5x_5}u$, of size $\epsilon$, changes the solution at the time $x_4>0$ by $\epsilon\,e^{\varkappa x_4}e^{ik_5x_5}u$, and this exceeds 1 in size once $|k_5|$ is large enough, however small $\epsilon$ is.

\subsection{The good sector and its Fock space}

\textbf{Definition.} The \textbf{good sector} is the sector of fields that do not depend on $x_5,x_6,x_7$, that is $k_5=k_6=k_7=0$. It is a dimensional reduction to a theory in 4+1 dimensions on the coordinates $(x_0,x_1,x_2,x_3,x_4)$. It is \textbf{not} a subspace of the one-particle space of the full theory: on the seven-dimensional slice the set $k_5=k_6=k_7=0$ has zero volume, so no normalizable state of the full theory lies in it, and a field independent of $x_5,x_6,x_7$ cannot satisfy the $\delta^7$ anticommutator. One keeps instead the zero modes in $x_5,x_6,x_7$ on a normalization volume $V_3$, for which $\{\Psi_a,\chi_b\}=\delta_{ab}\,\delta^4/V_3$ (Stage-1 document, §10.6).

\textbf{Boxes: where $\delta^4/V_3$ comes from.} Put the extra times in a box: let $x_5,x_6,x_7$ run over intervals of lengths $L_5,L_6,L_7$ with periodic ends (a field takes the same value at both ends of each interval), so that the box has the volume $V_3=L_5L_6L_7$. A function that does not depend on $x_5,x_6,x_7$ is constant on the box: it is the \textbf{zero mode}. The zero-mode part of the field is its average over the box, $\Psi^{(0)}(x)=\frac1{V_3}\int\Psi\,d^3x'$, where $x'=(x_5,x_6,x_7)$ and $x=(x_0,x_1,x_2,x_3)$ now denotes the remaining four slice coordinates. In flat space and Gaussian normal gauge $\{\Psi_a,\chi_b\}=\delta_{ab}\,\delta^7$ (Section 8.7), with $\delta^7=\delta^4(x-y)\,\delta^3(x'-y')$, so

\[
\bigl\{\Psi^{(0)}_a(x),\chi^{(0)}_b(y)\bigr\}=\frac{1}{V_3^2}\int\!\!\int\delta_{ab}\,\delta^4(x-y)\,\delta^3(x'-y')\,d^3x'\,d^3y'=\frac{\delta_{ab}\,\delta^4(x-y)}{V_3},
\]

because the $y'$ integral of $\delta^3$ is 1 and the remaining $x'$ integral is $V_3$ (here $\chi^{(0)}=\Psi^{(0)\dagger}B$ is the average of $\chi$). This is the anticommutator quoted above. The rescaled field $\tilde\Psi:=\sqrt{V_3}\,\Psi^{(0)}$ has $\{\tilde\Psi_a(x),\tilde\chi_b(y)\}=\delta_{ab}\,\delta^4(x-y)$ (the Stage-1 document absorbs $V_3$ into the field in this way, §10.7). For a field that does not depend on $x'$ the derivatives $\partial_5,\partial_6,\partial_7$ vanish and the $x'$ integral of the Hamiltonian of Section 8.5 gives a factor $V_3$, so $H=\int\tilde\Psi^\dagger\bigl(mC-\sum_{a=0}^{3}C\gamma^a\partial_a\bigr)\tilde\Psi\,d^4x$.

\textbf{Momenta.} In the same way let $x_0,\dots,x_3$ run over a periodic box of volume $V_4=L_0L_1L_2L_3$. The allowed wave numbers are $k_a=2\pi n_a/L_a$ with integers $n_a$ (so that $e^{ik_ax_a}$ has the same value at both ends), and with $k\cdot x:=\sum_{a=0}^{3}k_ax_a$ the functions $e^{ik\cdot x}/\sqrt{V_4}$ are orthonormal: $\frac1{V_4}\int e^{i(k'-k)\cdot x}\,d^4x$ is 1 for $k'=k$ and 0 otherwise, because $\int_0^Le^{2\pi inx/L}\,dx=0$ for every integer $n\ne0$. For each allowed $k$ define the column of 16 \textbf{mode operators} $\Psi_k:=\frac1{\sqrt{V_4}}\int e^{-ik\cdot x}\,\tilde\Psi(x)\,d^4x$, with Hilbert adjoint $\Psi_k^\ast=\frac1{\sqrt{V_4}}\int e^{ik\cdot x}\,\tilde\chi(x)\,d^4x$. The anticommutator of $\tilde\Psi$ gives

\[
\bigl\{\Psi_{k,a},\Psi_{k',b}^\ast\bigr\}=\frac1{V_4}\int\!\!\int e^{-ik\cdot x}e^{ik'\cdot y}\,\delta_{ab}\,\delta^4(x-y)\,d^4x\,d^4y=\delta_{ab}\,\delta_{kk'} ,
\]

where $\delta_{kk'}$ is 1 for $k=k'$ and 0 otherwise. Conversely $\tilde\Psi(x)=\sum_k\Psi_k\,e^{ik\cdot x}/\sqrt{V_4}$ (the Fourier series; that every periodic function is such a sum is a standard fact of analysis, used without proof). Inserting this into $H$, each $\partial_a$ becomes $ik_a$, and the orthonormality of the exponentials leaves one term for each $k$:

\[
H=\sum_k\Psi_k^\dagger\Bigl(mC-iC\sum_{a=0}^{3}k_a\gamma^a\Bigr)\Psi_k ,
\]

with $\Psi_k^\dagger=\Psi_k^\ast B$ (the Krein adjoint of Section 8.7). Different momenta do not talk to each other: every allowed $k$ is a copy of the same one-momentum problem, which we now solve.

\textbf{The Dirac Hamiltonian in 4+1 dimensions.} Put $\beta:=-i\gamma^4$ and $\alpha^a:=-\gamma^4\gamma^a$ for $a=0,1,2,3$. All five are Hermitian (Section 8.9), $\beta^2=1$, $(\alpha^a)^2=-(\gamma^4)^2(\gamma^a)^2=1$, $\{\alpha^a,\alpha^b\}=2\delta^{ab}$ and $\{\alpha^a,\beta\}=0$. In the good sector

\[
h_k=m\beta+\sum_{a=0}^{3}k_a\alpha^a,\qquad h_k^2=E_k^2,\qquad E_k=\sqrt{m^2+k_0^2+k_1^2+k_2^2+k_3^2},
\]

the standard Dirac Hamiltonian in 4+1 dimensions, acting on 16 components. It is Hermitian and traceless ($\mathrm{tr}\gamma^4=0$ and $\mathrm{tr}(\gamma^4\gamma^a)=0$ for $a\ne4$), so its eigenvalues are $+E_k$ and $-E_k$ with multiplicity 8 each: for every momentum $k=(k_0,k_1,k_2,k_3)$ there are 8 particle states and 8 antiparticle states.

\textbf{One momentum, step by step.} Fix $k$ and choose orthonormal eigenvectors $u_s$ ($h_ku_s=+E_ku_s$) and $v_s$ ($h_kv_s=-E_kv_s$), $s=0,\dots,7$. Together they form an orthonormal basis of $\mathbb C^{16}$ (eigenvectors of a Hermitian matrix for different eigenvalues are orthogonal, Chapter 1), so the \textbf{completeness relation} $\sum_s(u_s)_a(u_s)_b^\ast+\sum_s(v_s)_a(v_s)_b^\ast=\delta_{ab}$ holds. \emph{Derivation.} Call the 16 basis columns $e_0,\dots,e_{15}$. Every column $w$ is a combination $w=\sum_tc_te_t$; multiplying from the left by $e_s^\dagger$ and using $e_s^\dagger e_t=\delta_{st}$ gives $c_s=e_s^\dagger w$. So $w=\sum_se_s(e_s^\dagger w)=\bigl(\sum_se_se_s^\dagger\bigr)w$ for every $w$, hence $\sum_se_se_s^\dagger=I_{16}$, and the entry $(a,b)$ of this matrix equation is the completeness relation. Expand the mode operator $\Psi_k$ of this momentum (we drop the label $k$ and write $\Psi$) as

\[
\Psi=\sum_{s=0}^{7}b_s\,u_s\,e^{-iE_kx_4}+\sum_{s=0}^{7}d_s^\ast\,v_s\,e^{+iE_kx_4},\qquad \{b_s,b_{s'}^\ast\}=\{d_s,d_{s'}^\ast\}=\delta_{ss'},
\]

all other anticommutators zero, where ${}^\ast$ is the Hilbert adjoint. Then $\{\Psi_a,\Psi_b^\ast\}=\sum_s(u_s)_a(u_s)_b^\ast+\sum_s(v_s)_a(v_s)_b^\ast=\delta_{ab}$ by the completeness relation (the phases $e^{\mp iE_kx_4}$ cancel, and the computation is that of Section 8.4): this is the anticommutator $\{\Psi_{k,a},\Psi_{k,b}^\ast\}=\delta_{ab}$ found above, the positive anticommutator of Section 8.7. The \textbf{Hamiltonian} of the mode (its term in the sum over $k$ above) is $\Psi^\dagger(mC-iC\sum_ak_a\gamma^a)\Psi=\chi B(mC-iC\sum_ak_a\gamma^a)\Psi=\chi h_k\Psi$, where $\chi=\Psi^\ast$ is the Hilbert adjoint row of Section 8.7 and $\Psi^\dagger=\chi B$, because $BC=-i\gamma^4$ gives $B(mC)=m\beta$ and $B(-ik_aC\gamma^a)=-k_a\gamma^4\gamma^a=k_a\alpha^a$. Inserting the expansion and using orthonormality,

\[
H_k=\sum_sE_k\,b_s^\ast b_s-\sum_sE_k\,d_sd_s^\ast=\sum_sE_k\bigl(b_s^\ast b_s+d_s^\ast d_s\bigr)-8E_k .
\]

\textbf{Vacuum, Dirac sea, normal ordering.} The \textbf{vacuum} $|0\rangle$ is annihilated by all $b_s$ and all $d_s$. In the language of the classical modes, all negative-energy states are filled (the \textbf{Dirac sea}): $d_s^\ast$, which multiplies the negative-energy wave function $v_s$, creates an \textbf{antiparticle}, a hole in the sea, with positive energy $E_k$. The constant $-8E_k$ is the energy of the filled sea for this momentum; \textbf{normal ordering} removes it by definition.

\textbf{All momenta, and the continuum.} Every allowed $k$ has its own operators $b_s(k)$ and $d_s(k)$, and operators of different momenta anticommute, because $\{\Psi_{k,a},\Psi_{k',b}^\ast\}=0$ for $k\ne k'$. Adding the normal-ordered $H_k$ gives $H=\sum_kE_k\sum_s\bigl(b_s(k)^\ast b_s(k)+d_s(k)^\ast d_s(k)\bigr)$. The allowed $k$ form a grid with the spacing $2\pi/L_a$ in direction $a$, so each grid point owns a cell of volume $(2\pi)^4/V_4$, and for a large box a sum over the grid becomes an integral, $\sum_kf(k)\approx V_4\int f(k)\,d^4k/(2\pi)^4$. Rescale the operators, $\sqrt{V_4}\,b_s(k)\to b_s(k)$ and $\sqrt{V_4}\,d_s(k)\to d_s(k)$. Then $H=\frac1{V_4}\sum_kE_k\sum_s\bigl(b_s(k)^\ast b_s(k)+d_s(k)^\ast d_s(k)\bigr)$, which becomes the integral below, and the anticommutator of the new operators becomes $V_4\,\delta_{kk'}=(2\pi)^4\,\delta_{kk'}/\bigl((2\pi)^4/V_4\bigr)$, a grid delta divided by the cell volume, which for a large box is the idealized $(2\pi)^4\delta^4(k-k')$ (as for the lattice of points in Section 8.6). The result, for all momenta (Stage-1 document, §10.7), is

\[
H=\int\frac{d^4k}{(2\pi)^4}\,E_k\sum_{s=0}^{7}\bigl(b_s(k)^\ast b_s(k)+d_s(k)^\ast d_s(k)\bigr)\ \ge0,\qquad \{b_s(k),b_{s'}(k')^\ast\}=\delta_{ss'}(2\pi)^4\delta^4(k-k') .
\]

(The Stage-1 document writes the antiparticle operator of the mode $e^{ik\cdot x}$ as $d_s(-k)$; renaming $k\to-k$ in the sum does not change $H$, because $E_{-k}=E_k$.)

This is a positive-definite Fock space with a Hamiltonian that is not negative. It is the quantization of the reduced (4+1)-dimensional theory, not a subspace of a Fock space of the full eight-dimensional theory. The construction of this section is derived from the checked Results 10.2, 10.4 and 10.5; it is not itself a machine check (Stage-1 document, §10.7 and §13, item 5).

\subsection{Symmetries: unitary and Krein-unitary}

\textbf{Two notions.} A spin transformation $R=\exp(\theta S)$ that maps the slices $x_4=\text{const}$ to themselves preserves the canonical anticommutator if $R^\dagger BR=B$; it is then called \textbf{Krein-unitary}. (The anticommutator of the transformed field $R\Psi$ is $\sum_{c,d}R_{ac}R_{bd}^\ast\{\Psi_c,\Psi_d^\dagger\}$, that is $RBR^\dagger$ in place of $B$. The two conditions are equivalent because $B^2=1$: if $R^\dagger BR=B$, then $(BR^\dagger B)R=B(R^\dagger BR)=B^2=1$, so $R^{-1}=BR^\dagger B$, hence $R^\dagger=BR^{-1}B$ and $RBR^\dagger=RB\,BR^{-1}B=B$; conversely, if $RBR^\dagger=B$, then $R(BR^\dagger B)=B^2=1$, so again $R^\dagger=BR^{-1}B$ and $R^\dagger BR=BR^{-1}B\,BR=B$.) It is implemented by a unitary operator on the positive Fock space if in addition $R$ is unitary, $R^\dagger R=1$, and commutes with $J=B$ (Stage-1 document, §10.8). For the generator $S$ the conditions read: $S$ anti-Hermitian ($S^\dagger=-S$) for unitarity, and $S^\dagger B+BS=0$ for Krein-unitarity (the first-order terms of $R^\dagger R=1$ and $R^\dagger BR=B$ for $R=1+\theta S+\dots$).

\textbf{Names of the groups.} Take a set of the directions $0,\dots,7$ containing $p$ space-like ones (indices $\le3$) and $q$ time-like ones (indices $\ge4$). $\mathrm{Spin}(p,q)$ denotes the analogue of $\mathrm{Spin}(4,4)$ (Section 2.14) built only from those directions: the products $\gamma(u_1)\cdots\gamma(u_k)$, with $k$ even, of unit vectors $u_i$ whose components outside the set are zero. Below, Spin(4) uses the directions 0 to 3, Spin(3) the directions 5 to 7, Spin(4,3) all directions except 4, and Spin(4,1) the directions 0 to 4. Every $\exp(\theta S^{ab})$ with $a,b$ in the set lies in $\mathrm{Spin}(p,q)$, because Section 2.14 writes it as a product of two unit vectors built from $\gamma^a$ and $\gamma^b$, both of norm $\eta^{aa}$; the products of such exponentials form a group, and we say, with the Stage-1 document, that the generators $S^{ab}$ of the set \textbf{generate} $\mathrm{Spin}(p,q)$. When $p$ and $q$ are both nonzero, these products are not all of $\mathrm{Spin}(p,q)$: each exponential has the spinor norm $N=(\eta^{aa})^2=+1$ of Theorem 2.8, so every product of them has $N=+1$, while $\gamma^0\gamma^4\in\mathrm{Spin}(4,1)$ and $\gamma^0\gamma^5\in\mathrm{Spin}(4,3)$ have $N=(+1)(-1)=-1$ ($N$ is multiplicative: it is the product of the norms of the factors). ``Generate'' therefore means: generate the part of the group that the products of exponentials reach, its \textbf{part connected to 1}, as for Spin(4,4) in Section 2.14. $G_1\times G_2$ denotes the group of products $g_1g_2$, $g_1$ in $G_1$ and $g_2$ in $G_2$, when every element of $G_1$ commutes with every element of $G_2$. This holds for Spin(4) and Spin(3): a gamma with index $\le3$ anticommutes with one with index $\ge5$, so a product of an even number of the first kind commutes with a product of an even number of the second kind. (Since $-1$ lies in both groups, $g_1g_2=(-g_1)(-g_2)$ has two such descriptions; mathematicians record this by writing $(G_1\times G_2)/\{\pm1\}$, and we keep the physicists' name.) Finally, U(1) is the group of the phases $e^{i\alpha}$, $\alpha$ real, of Section 5.7.

\textbf{Counting by hand.} The 28 generators $S^{ab}=\tfrac12\gamma^a\gamma^b$ ($a<b$) are real, so $(S^{ab})^\dagger=(S^{ab})^T=\tfrac14[(\gamma^b)^T,(\gamma^a)^T]$.

\begin{itemize}
\item \emph{Hermiticity.} If $a,b\le3$ both gammas are symmetric and $(S^{ab})^T=\tfrac14[\gamma^b,\gamma^a]=-S^{ab}$; if $a,b\ge4$ both are antisymmetric and the two signs cancel, again $-S^{ab}$. If one index is $\le3$ and the other $\ge4$ one sign remains and $(S^{ab})^T=+S^{ab}$. So $6+6=12$ generators are anti-Hermitian and 16 (the ``boosts'' mixing a space-like and a time-like direction) are Hermitian.
\item \emph{Commuting with $B$.} $B=-i\gamma^0\gamma^1\gamma^2\gamma^3\gamma^4$. A gamma with index $\le4$ anticommutes with four of these five factors and commutes with itself, so it commutes with $B$; a gamma with index $5$, $6$ or $7$ anticommutes with all five, so it anticommutes with $B$. Hence $S^{ab}$ commutes with $B$ exactly when both indices lie in $\{0,\dots,4\}$ or both in $\{5,6,7\}$: $\binom52+\binom32=10+3=13$ generators; the other 15 anticommute with $B$.
\item \emph{Unitarily implemented and slice-preserving.} Anti-Hermitian and commuting with $B$: both indices in $\{0,1,2,3\}$ (6) or both in $\{5,6,7\}$ (3), in total \textbf{9}, generating $\mathrm{Spin}(4)\times\mathrm{Spin}(3)$.
\item \emph{Krein-unitary.} For an anti-Hermitian $S$ the condition $S^\dagger B+BS=0$ means $[B,S]=0$; for a Hermitian $S$ it means $\{B,S\}=0$. The first kind gives the 9 above; the second kind gives the 12 Hermitian generators with one index in $\{0,1,2,3\}$ and the other in $\{5,6,7\}$. In total \textbf{21}, exactly the generators with $a,b\ne4$, generating the Krein-unitary slice group Spin(4,3). The condition fails for the 7 generators with an index 4: the 4 Hermitian $S^{a4}$ ($a\le3$), which commute with $B$, and the 3 anti-Hermitian $S^{4b}$ ($b\ge5$), which anticommute with it.
\end{itemize}

These counts, 13, 9, 21 and 12, are Result 10.6 of the Stage-1 document (check \texttt{QNT\_unit\allowbreak{}aryAndKr\allowbreak{}einSubgr\allowbreak{}oups}, both algebra reports). They correct the first statement of the project contract, ``Spin(4)×Spin(3) (compact, commuting with B)'', which suggested that exactly 9 generators commute with $B$; the reports record that literal claim as false (erratum E1 in \texttt{handoff/\allowbreak{}specs/CO\allowbreak{}NTRACT.m\allowbreak{}d}, §11).

\textbf{Boosts in the good sector.} The four boosts $S^{a4}$, $a\le3$, commute with $B$ but are Hermitian, so as $16\times16$ matrices $\exp(\theta S^{a4})$ is neither unitary nor Krein-unitary. This is a statement about matrices, not about physics: a boost tilts the slices $x_4=\text{const}$, and in the good sector it is implemented unitarily on the Fock space exactly as in ordinary Dirac theory, whose boost generators are also Hermitian. The Stage-1 document (§10.8) derives this; the matrix statements behind it were checked numerically at sample momenta while that document was revised, not as recorded checks. The symmetries that act unitarily in the good sector include the translations in $x_0,\dots,x_4$, the Spin(4,1) generated by the 10 $S^{ab}$ with $a,b\le4$, the Spin(3) of the extra times, and the U(1) phase.

\subsection{The U(1) charge and the expectation-value rule}

\textbf{The current.} The phase symmetry $\Psi\to e^{i\alpha}\Psi$ (Section 5.7; the real number $\alpha$ has nothing to do with the matrices $\alpha^a$ of Section 8.10) gives the conserved current $J^\mu=-i\bar\Psi\gamma^\mu\Psi$ (Chapter 7). Its matrices $-iC\gamma^a$ are Hermitian: $(-iC\gamma^a)^\dagger=i(\gamma^a)^TC=i(-C\gamma^aC)C=-iC\gamma^a$ (Section 5.12). In Gaussian normal gauge $J^4=-i\Psi^\dagger C\gamma^4\Psi=\Psi^\dagger B\Psi$: the charge density is the Krein form (checks \texttt{QNT\_curr\allowbreak{}entHermi\allowbreak{}ticity}, both algebra reports, and \texttt{GR\_curre\allowbreak{}ntHermit\allowbreak{}ian}, \texttt{grassman\allowbreak{}n-demo-r\allowbreak{}eport.js\allowbreak{}on}). The conserved charge is $Q=\int\sqrt{|g|}\,J^4\,d^7x$. As a quadratic form in classical spinors it is indefinite, of signature (8,8). In the positive representation, $\Psi^\dagger B\Psi=\chi BB\Psi=\chi\Psi$, and in the good sector the mode expansion of Section 8.10 gives, after normal ordering (summed over all modes),

\[
Q=\sum\bigl(b^\ast b-d^\ast d\bigr).
\]

Particles have charge $+1$ and antiparticles charge $-1$ (derived, Stage-1 document §10.9).

\textbf{The expectation-value rule.} Every classical bilinear becomes, in the positive representation, $\Psi^\dagger M\Psi=\chi BM\Psi$ (Section 8.7). For a one-particle state $|u\rangle=b_u^\ast|0\rangle$, whose wave function $u$ is a normalized ($u^\dagger u=1$) positive-energy solution in the good sector, the normal-ordered expectation value is

\[
\bigl\langle\Psi^\dagger M\Psi\bigr\rangle=u^\dagger BM\,u .
\]

\emph{Derivation.} Choose the positive-energy basis of Section 8.10 with $u_0=u$. In $\chi_a(BM)_{ab}\Psi_b$ the particle part is $\sum_{s,s'}b_s^\ast b_{s'}\,(u_s^\dagger BMu_{s'})$, and $\langle u|b_s^\ast b_{s'}|u\rangle=1$ for $s=s'=0$ and 0 otherwise. The terms with $d\,d^\ast$ become $-d^\ast d$ after normal ordering and vanish in a state without antiparticles; the mixed terms $b^\ast d^\ast$ and $d\,b$ change the number of particles and have zero expectation value. $\square$

\textbf{Examples} (Stage-1 document, §10.11):

\begin{itemize}
\item the charge, $M=B$: $u^\dagger B^2u=u^\dagger u=1$;
\item the energy, $M=Bh$ (because $H=\chi h\Psi=\Psi^\dagger Bh\Psi$): $u^\dagger B^2hu=u^\dagger hu=E$;
\item the scalar density $S=\bar\Psi\Psi$, $M=C$: $u^\dagger BCu=u^\dagger(-i\gamma^4)u$, which equals $+1$ for a positive-energy state at rest with $m>0$. So the energy density at rest is $\langle\rho\rangle=m\langle S\rangle=m=E$, as it should be.
\end{itemize}

The two rest states $u_\pm$ of Section 8.8 both have $u^\dagger(-i\gamma^4)u=1$, but the naive density $\Psi^\dagger\Psi$ would give $u^\dagger Bu=+1$ for $u_+$ and $-1$ for $u_-$. The naive density is not the physical one. This rule is the one used by the dark-sector experiments of Chapter 11 (Student Guide, §6.10).

\textbf{Which bilinears are observables.} In the positive representation $\chi BM\Psi$ is self-adjoint exactly when $BM$ is Hermitian; for a Hermitian $M$ this holds exactly when $[M,B]=0$. The current components $J^0,\dots,J^4$ pass this test. For $J^5,J^6,J^7$ the Hermitian matrix $-iC\gamma^j$ ($j\ge5$) anticommutes with $B$ (Section 8.11: $\gamma^j$ anticommutes with $B$, and $C$ commutes with it), so $BM$ is anti-Hermitian and these components have purely imaginary expectation values: they are not observables in this Fock space (Stage-1 document, §10.9 and §10.11).

\subsection{The extra-time modes grow}

Section 8.9 showed that a wave with enough momentum along an extra time grows instead of oscillating. The fifth experiment of Stage 3, \textbf{EXP-5}, follows this growth numerically in a background in which the extra times shrink (\texttt{provenan\allowbreak{}ce/DIRAC\allowbreak{}16COMPLE\allowbreak{}X\_DARK\_S\allowbreak{}ECTOR\_NU\allowbreak{}MERICS.m\allowbreak{}d}, §10). The extra-time scale factor is $e^{-x_4}$, and the mode has comoving momentum $q$ along $x_5$ ($m=1$), so its physical momentum $Q=q\,e^{x_4}$ (not the charge $Q$ of Section 8.12) grows and

\[
h=-im\gamma^4-Q\,\gamma^4\gamma^5,\qquad E^2=m^2-q^2e^{2x_4},\qquad E^2<0\ \text{ for }\ x_4>t_\ast=\ln(m/q).
\]

For $q=0.05$ and $q=0.1$ the turning times are $t_\ast=2.99573$ and $2.30259$, and each run starts on a positive-energy eigenvector of $h$ with Krein norm $u^\dagger Bu=\pm E_0/m$, $E_0=0.998749$ and $0.994987$ (\texttt{artifact\allowbreak{}s/dirac1\allowbreak{}6complex\allowbreak{}/numeric\allowbreak{}s/exp5/s\allowbreak{}ummary.j\allowbreak{}son}).

\textbf{The WKB estimate at leading order.} The \textbf{WKB approximation} (named after Wentzel, Kramers and Brillouin) assumes that $Q$ changes slowly compared with the local growth rate. At leading order one treats $Q$ as constant. After $t_\ast$ the eigenvalues of $h$ are $\pm i\varkappa$ with $\varkappa=\sqrt{Q^2-m^2}$ (because $h^2=E^2=-\varkappa^2$); for the eigenvalue $+i\varkappa$ the equation $i\dot u=i\varkappa u$ gives $u\propto e^{\varkappa x_4}$, so the norm $u^\dagger u$ grows at the rate $2\varkappa$. Adding up these rates gives $\ln(u^\dagger u)\approx2W(x_4)+\text{const}$ with

\[
W(x_4)=\int_{t_\ast}^{x_4}\varkappa\,dx_4'=\sqrt{Q^2-m^2}-m\arccos(m/Q) .
\]

To check the closed form, differentiate with $\frac{d}{dQ}\arccos(m/Q)=m/(Q\sqrt{Q^2-m^2})$: $\frac{d}{dQ}\bigl[\sqrt{Q^2-m^2}-m\arccos(m/Q)\bigr]=\frac{Q}{\sqrt{Q^2-m^2}}-\frac{m^2}{Q\sqrt{Q^2-m^2}}=\frac{\sqrt{Q^2-m^2}}{Q}$, and $dQ/dx_4=Q$, so the derivative with respect to $x_4$ is $\varkappa$; both sides vanish at $t_\ast$, where $Q=m$.

\textbf{The first-order correction.} The next correction comes from the turning of the growing eigenvector while $Q$ changes. The problem splits into eight identical two-dimensional ones. Put $A:=-i\gamma^4$ and $M:=-i\gamma^4\gamma^5$. $A$ is Hermitian with square 1 (Section 8.9). $\gamma^4\gamma^5$ is real antisymmetric, $(\gamma^4\gamma^5)^T=(\gamma^5)^T(\gamma^4)^T=\gamma^5\gamma^4=-\gamma^4\gamma^5$, and $(\gamma^4\gamma^5)^2=-(\gamma^4)^2(\gamma^5)^2=-1$, so $M$ is Hermitian with $M^2=1$. They anticommute, $AM=-\gamma^4\gamma^4\gamma^5=\gamma^5$ and $MA=-\gamma^4\gamma^5\gamma^4=-\gamma^5$, and $h=mA-iQM$. Choose eight orthonormal columns $w_0,\dots,w_7$ with $Aw_i=w_i$ ($A$ is traceless with square 1, so its eigenvalue $+1$ occurs eight times). Then $A(Mw_i)=-MAw_i=-Mw_i$, and the 16 columns $w_i$ and $Mw_i$ are orthonormal ($M^\dagger M=M^2=1$, and eigenvectors of $A$ for different eigenvalues are orthogonal). In the plane of $w_i$ and $Mw_i$ the matrices act as $A=\mathrm{diag}(1,-1)$ and $M=\begin{pmatrix}0&1\\1&0\end{pmatrix}$. So, writing $u=\sum_i(x_iw_i+y_iMw_i)$, every pair $v=(x_i,y_i)^T$ obeys the same equation

\[
\dot v=-ih\,v=\begin{pmatrix}-im&-Q\\-Q&im\end{pmatrix}v .
\]

This symmetric matrix has the eigenvector $r=(Q,\,-\varkappa-im)^T$ with the eigenvalue $+\varkappa$ and $s=(Q,\,\varkappa-im)^T$ with the eigenvalue $-\varkappa$ (insert them and use $Q^2=\varkappa^2+m^2$), and $r^Ts=Q^2-(\varkappa^2+m^2)=0$. Write $v=c\,r+d\,s$, with numbers $c$ and $d$ that depend on the pair $i$. Multiply the equation from the left by $r^T$; because the matrix is symmetric, $r^T(-ih)=\varkappa\,r^T$, and so $\frac{d}{dx_4}(r^Tv)-\dot r^Tv=\varkappa\,r^Tv$. Here $r^Tv=c\,n$ with $n:=r^Tr=Q^2+(\varkappa+im)^2=2\varkappa(\varkappa+im)$, and $\dot r^Tr=\tfrac12\dot n$. Now comes the approximation: drop the term $d\,\dot r^Ts$, which couples the growing part to the decaying one (the book does not estimate it; the late-rate check below compares the result with the computed growth). What remains is $\dot c\,n+\tfrac12c\,\dot n=\varkappa\,c\,n$, so $\ln c=W-\tfrac12\ln n+\text{const}$. Once the decaying part has died away, $u^\dagger u\approx\sum_i|c_i|^2\,r^\dagger r$, and with $|n|=2\varkappa\,|\varkappa+im|=2\varkappa Q$ and $r^\dagger r=Q^2+\varkappa^2+m^2=2Q^2$,

\[
\ln(u^\dagger u)\approx2W+\ln\frac Q\varkappa+\text{const},\qquad \frac{d}{dx_4}\ln(u^\dagger u)\approx2\varkappa+\frac{\dot Q}{Q}-\frac{\dot\varkappa}{\varkappa}=2\varkappa-\frac{m^2}{\varkappa^2},
\]

because $\dot Q=Q$ and $\varkappa\dot\varkappa=Q\dot Q=Q^2$ give $\frac{\dot Q}{Q}-\frac{\dot\varkappa}{\varkappa}=1-\frac{Q^2}{\varkappa^2}=-\frac{m^2}{\varkappa^2}$. This is the corrected rate of the numerics document (§10.2); integrating it gives the \textbf{first-order} value $2\Delta W+\Delta\ln(Q/\varkappa)$ of an increase, where $\Delta$ is the change over an interval. Near $t_\ast$ both approximations fail ($\varkappa\to0$, and $\ln(Q/\varkappa)$ becomes infinite), and the constant depends on how the solution passes the turning point; it drops out of differences. So the repository compares increases over a \textbf{window} away from the turning point, from the first output time $x_A$ at or after $t_\ast+1.5$ ($x_A=4.4968$ for $q=0.05$ and $3.8090$ for $q=0.1$) to the end, $t_\ast+3$.

\textbf{The committed results.}

\begin{itemize}
\item The Hilbert norm $u^\dagger u$ grows to $1.258\times10^{16}$ ($q=0.05$) and $1.313\times10^{16}$ ($q=0.1$) by $x_4=t_\ast+3$. Its average growth rate rises from about 2.4 over the first unit of time after $t_\ast$ to 25.3 over the last unit (at the end the rate $2\varkappa$ itself is $2\sqrt{e^6-1}\approx40.1$, because $Q(t_\ast+3)=m\,e^3$): super-exponential growth, because $Q$ keeps growing.
\item Over the window the increase of $\ln(u^\dagger u)$, 30.9922 for $q=0.05$, agrees with the leading WKB value $2[W(t_\ast+3)-W(x_A)]=31.0241$ to $1.03\times10^{-3}$ relative, and with the first-order value 30.99986, which adds the change of $\ln(Q/\varkappa)$ over the window ($-0.0242$), to $2.47\times10^{-4}$ (for $q=0.1$: 30.9455 against 30.9770 and 30.9530). The self-checks are \texttt{growth\_m\allowbreak{}atches\_w\allowbreak{}kb\_leadi\allowbreak{}ng\_order} and \texttt{growth\_m\allowbreak{}atches\_w\allowbreak{}kb\_first\allowbreak{}\_order}, and the checker's \texttt{wkbLeadi\allowbreak{}ng} and \texttt{wkbFirst\allowbreak{}Order}; its check \texttt{wkbLateR\allowbreak{}ate} finds that the computed rate $d\ln(u^\dagger u)/dx_4$ near the end (a finite-difference derivative two output steps before the end) agrees with $2\varkappa-m^2/\varkappa^2$ to $4.6\times10^{-6}$ relative (numerics document, §10.5; \texttt{python-c\allowbreak{}heck-rep\allowbreak{}ort.json} in the same folder).
\item The Krein norm is conserved to the accuracy of 16-digit arithmetic: its drift divided by $\max(u^\dagger u,1)$ is at most $1.08\times10^{-9}$, and its absolute drift up to $t_\ast$ is at most $1.28\times10^{-9}$. At the end, where $u^\dagger u\approx10^{16}$, the computed $u^\dagger Bu$ is a difference of two numbers of order $10^{16}$, and rounding limits its absolute accuracy to about 1. The conservation law of Section 8.9 says that the exact drift is zero.
\end{itemize}

\begin{figure}[htbp]
\centering
\includegraphics[width=0.92\linewidth]{artifacts/dirac16complex/numerics/figures/exp5_growth.png}
\caption{EXP-5: (a) \texorpdfstring{$\ln(u^\dagger u)$}{(u\textasciicircum{}†u)} for both momenta and both signs of \texorpdfstring{$C$}{C} together with the WKB curve \texorpdfstring{$2W$}{2W}, shifted by a constant to agree with \texorpdfstring{$\ln(u^\dagger u)$}{(u\textasciicircum{}†u)} at the window start \texorpdfstring{$x_A\approx t_\ast+1.5$}{x\_A≈t\_∗+1.5} (dashed); (b) \texorpdfstring{$E^2$}{E\textasciicircum{}2}, which turns negative at \texorpdfstring{$t_\ast=\ln(m/q)$}{t\_∗=(m/q)} (dotted).}
\end{figure}

The Rust program passes 7 of 7 self-checks for EXP-5 and its independent checker 21 of 21 (numerics document, §10.6). The growth conserves the indefinite Krein norm, so the positive-norm and the negative-norm parts grow together.

\textbf{In the primordial field.} The Stage-2 document (§14.3) finds the same onset in the author's primordial field, where 3-space inflates as $e^{a_4}$ and the extra times deflate as $e^{-a_4}$: with frozen coefficients (that is, treating the coefficients of the equation as constants near one point, as in the leading WKB step above) a mode with momentum $k_5$ along $x_5$ has

\[
E^2=M_{\mathrm{eff}}^2+K^2+\bigl(k_1e^{-H\zeta-a_4}\bigr)^2-\bigl(k_5e^{-H\zeta+a_4}\bigr)^2 .
\]

Here the mode is, near the point considered, $\Psi=e^{-3H\zeta}e^{iK\zeta}e^{ik_1x_1+ik_5x_5}u(x_4)$ in the warped coordinates of Section 4.4 (Stage-2 document, §14.1 to §14.3): $\zeta$ is the hidden coordinate of Section 4.4, $H>0$ the notebook's inverse length (Section 4.2; not a Hamiltonian), $a_4$ the free function of the metric, $K$ the wave number along $\zeta$ (not the matrix $K$ of Section 8.5), $k_1$ and $k_5$ the wave numbers along $x_1$ and $x_5$, and $M_{\mathrm{eff}}=m+U'(S)$ the effective mass of Chapter 6, which equals $m+\lambda S$ for $U=\tfrac\lambda2S^2$; Chapter 9 treats this field in detail. $E^2$ is negative exactly when $e^{-2H\zeta}\bigl(k_5^2e^{2a_4}-k_1^2e^{-2a_4}\bigr)>M_{\mathrm{eff}}^2+K^2$. For $k_1=0$ this happens once $a_4>H\zeta+\ln\bigl(\sqrt{M_{\mathrm{eff}}^2+K^2}/|k_5|\bigr)$ (take the logarithm of $k_5^2e^{2a_4-2H\zeta}>M_{\mathrm{eff}}^2+K^2$); a nonzero $k_1$ makes the left side smaller and so delays the onset. An unboundedly growing $a_4$, such as the notebook's $a_4=t$, reaches this onset, because the $k_5$ term grows like $e^{2a_4}$ while the $k_1$ term decays like $e^{-2a_4}$.

\textbf{Consequence.} The positive Fock space exists only in the good sector. Every quantum statement and every numerical result of Chapters 11 (EXP-1 to EXP-4; EXP-5 studies the excluded modes), 13 and 14 is restricted to that sector. The restriction is \textbf{imposed, not derived}: whether some dynamical mechanism suppresses the extra-time sector is an open problem (Chapter 18).

\subsection{How the repository checks all of this}

The algebraic ingredients of the quantization are machine-checked; the constructions built on them are derivations (Stage-1 document, §13, item 5):

\begingroup
\small
\begin{longtable}{@{}>{\raggedright\arraybackslash}p{0.400\linewidth}>{\raggedright\arraybackslash}p{0.400\linewidth}@{}}
\toprule
\textbf{statement} & \textbf{check (report)} \\
\midrule
\endfirsthead
\toprule
\textbf{statement} & \textbf{check (report)} \\
\midrule
\endhead
$\Pi=\sqrt{\lvert g\rvert}\,\Psi^\dagger C\gamma^{x_4}$ in a Grassmann algebra & \texttt{QNT\_cano\allowbreak{}nicalMom\allowbreak{}entum\_G1}, \texttt{\detokenize{_G2}} (wolfram-geometry) \\
$(C\gamma^{x_4})(\gamma^{x_4}C)=g^{44}$; reduction to $B$ in the primordial field & \texttt{QNT\_curv\allowbreak{}edAntico\allowbreak{}mmutator\allowbreak{}Matrix} (wolfram-algebra) \\
$B$ Hermitian, $B^2=1$, spectrum $(8,8)$, $[C,B]=0$, $BC=-i\gamma^4$ & \texttt{ALG\_char\allowbreak{}geFormB} (both algebra reports) \\
invariant forms $CP_\pm$; every invariant density indefinite & \texttt{ALG\_inva\allowbreak{}riantFor\allowbreak{}ms} (both algebra reports) \\
Krein signature (4,4) on the energy eigenspaces & \texttt{QNT\_krei\allowbreak{}nSignatu\allowbreak{}re} (both algebra reports) \\
$h_k^2=E^2$, Hermiticity, $[h_k,B]$ & \texttt{QNT\_flat\allowbreak{}ModeHami\allowbreak{}ltonian} (both algebra reports) \\
the counts 13, 9, 21, 12 & \texttt{QNT\_unit\allowbreak{}aryAndKr\allowbreak{}einSubgr\allowbreak{}oups} (both algebra reports) \\
Hermitian current, $J^4$ matrix equals $B$ & \texttt{QNT\_curr\allowbreak{}entHermi\allowbreak{}ticity} (both), \texttt{GR\_curre\allowbreak{}ntHermit\allowbreak{}ian} (grassmann-demo) \\
extra-time growth, Krein-norm conservation & EXP-5: 7 self-checks, 21 checker checks \\
\bottomrule
\end{longtable}
\endgroup

Not machine-checked, but derived in the Stage-1 document and in this chapter: the Dirac-bracket computation, the Hamiltonian density, the Heisenberg equations, the mode expansion, the Fock space and the expectation-value rule. The algebra reports are \texttt{artifact\allowbreak{}s/dirac1\allowbreak{}6complex\allowbreak{}/arbitra\allowbreak{}ry-field\allowbreak{}/wolfram\allowbreak{}-algebra\allowbreak{}-report.\allowbreak{}json} (21 of 21 checks true) and \texttt{python-a\allowbreak{}lgebra-r\allowbreak{}eport.js\allowbreak{}on} (21 of 21), the geometry report \texttt{wolfram-\allowbreak{}geometry\allowbreak{}-report.\allowbreak{}json} (43 of 43); the Stage-1 gate that regenerates and compares all of them is described in Chapter 19.

\subsection{What we proved and what we assumed}

\textbf{What we proved.} With complete arguments in this chapter: Hamilton's equations and $\dot F=\{F,H\}_{\mathrm P}$ follow from the Lagrangian form (Section 8.2); the anticommutator of a fermion mode implies the Pauli principle, and many modes give a Fock space whose occupation-number states are orthonormal, so that its inner product is positive (Sections 8.3 and 8.4); the dirac16complex Lagrangian can be written, up to a total derivative, in the first-order form $\Psi^\dagger K\partial_4\Psi-\mathcal H$ with the anti-Hermitian $K=\sqrt{|g|}\,C\gamma^{x_4}$, which is invertible exactly when $g^{44}\ne0$; requiring that the Heisenberg equation reproduce the classical field equation gives $\{\Psi,\Psi^\dagger\}=iK^{-1}\delta^7$, which in Gaussian normal gauge is $B\,\delta^7/\sqrt{|g|}$ with $B=-iC\gamma^4$ Hermitian, $B^2=1$ and eight eigenvalues of each sign; no positive Hilbert space can carry this anticommutator with $\Psi^\dagger$ as the Hilbert adjoint (Theorem 8.1), so the one-particle space is a Krein space with fundamental symmetry $J=B$; every Spin(4,4)-invariant charge density is indefinite, given the computed list of invariant forms (the invariant Hermitian forms are then exactly the real combinations of $CP_-$ and $CP_+$); the Krein form has signature (4,4) on the positive-energy rest states; the mode Hamiltonian satisfies $h_k^2=E^2$, it is Hermitian exactly when there is no extra-time momentum, and it always conserves the Krein norm; with extra-time momentum the modes oscillate, grow linearly or grow exponentially according to the sign of $E^2$, and the exponential growth makes the initial-value problem ill-posed; in the good sector, with box normalization of the extra times and of the momenta, there is a positive Fock space with a non-negative normal-ordered Hamiltonian, particles of charge $+1$ and antiparticles of charge $-1$; the expectation-value rule $\langle\Psi^\dagger M\Psi\rangle=u^\dagger BMu$; the counts 13, 9, 21 and 12 of the symmetry generators, and the equivalence of $R^\dagger BR=B$ and $RBR^\dagger=B$; and, within the WKB approximation, the EXP-5 growth rates $2\varkappa$ (leading order) and $2\varkappa-m^2/\varkappa^2$ (first order).

\textbf{What we assumed or took from elsewhere.} Anticommutators rather than commutators are part of the definition of dirac16complex as a fermion field; no spin--statistics theorem for signature (4,4) is proved. The quantization is canonical with respect to $x_4$ and requires $g^{44}\ne0$; the matrix $B$ appears in Gaussian normal gauge. Dirac's bracket method is described, not carried out step by step (our derivation uses the Heisenberg equation on a finite set of points instead). That the invariant bilinear forms are spanned by $CP_-$ and $CP_+$ is an exact computation of the repository (\texttt{ALG\_inva\allowbreak{}riantFor\allowbreak{}ms}), not a hand proof. Normal ordering (subtracting the energy of the filled sea) is a prescription. That every periodic function is a Fourier series, and the passage from the sum over a momentum grid to the integral for a large box, are standard analysis used without proof. The WKB approximation drops the coupling of the growing to the decaying part of an EXP-5 mode without estimating it; the late-rate check of the repository compares the result with the computation. The good sector is a restriction \textbf{imposed by hand}; no mechanism that removes the extra-time modes is derived. The unitary implementation of the boosts in the good sector is derived in the Stage-1 document, not a recorded check. The quantization is formal: no interacting theory, regularization or renormalization is constructed, and the Krein space is not given a physical interpretation beyond the good sector. The EXP-5 numbers are floating-point results with the tolerances stated in the numerics document.

\subsection{Exercises}

\textbf{Exercise 8.1.} With the $2\times2$ matrices $b$ and $b^\dagger$ of Section 8.3, verify $\{b,b^\dagger\}=1$ and $(b^\dagger)^2=0$, and find the eigenvalues of $N=b^\dagger b$.

\textbf{Exercise 8.2.} In the two-mode Fock space of Section 8.4 compute $b_0^\dagger|01\rangle$, $b_1b_0^\dagger|01\rangle$ and $b_0^\dagger b_1|01\rangle$, and check $\{b_1,b_0^\dagger\}|01\rangle=0$.

\textbf{Exercise 8.3.} For the first-order Lagrangian $L'=2i\,\theta^\ast\dot\theta-\omega\,\theta^\ast\theta$ (complex Grassmann $\theta$), find the classical equation of motion and the anticommutator $\{\theta,\theta^\dagger\}$ given by the rule $iK^{-1}$ of Section 8.6, and check that the Heisenberg equation with $H=\omega\theta^\dagger\theta$ reproduces the classical equation.

\textbf{Exercise 8.4.} Show $B^\dagger=B$ and $B^2=1$ for $B=-iC\gamma^4$, using only $C^T=C$, $C^2=1$, $(\gamma^4)^T=-\gamma^4$, $(\gamma^4)^2=-1$ and $C\gamma^4=\gamma^4C$.

\textbf{Exercise 8.5.} Compute $E^2$ for (a) $m=2$, $k_1=1$, $k_6=2$; (b) $m=1$, $k_0=1$, $k_5=1$, $k_7=1$; (c) $m=1$, $k_5=3$. For each case say whether the mode oscillates, grows linearly or grows exponentially, and give the growth rate $\varkappa$ where there is one.

\textbf{Exercise 8.6.} For $m=1$ and $k_5=1$ (all other $k_j=0$) show directly that $h=-i\gamma^4-\gamma^4\gamma^5$ satisfies $h^2=0$, and write the solution $u(x_4)$ of $i\partial_4u=hu$.

\textbf{Exercise 8.7.} Verify with the tables quoted in Section 8.8 that $u_-=\tfrac12(e_0+e_4+ie_9-ie_{13})$ satisfies $-i\gamma^4u_-=u_-$ and $Cu_-=-u_-$. Compute $u_-^\dagger u_-$, $u_-^\dagger Bu_-$ and the scalar density $u_-^\dagger(-i\gamma^4)u_-$.

\textbf{Exercise 8.8.} Among the seven generators $S^{01},\dots,S^{07}$, which commute with $B$, which are anti-Hermitian, and which satisfy the Krein condition?

\textbf{Exercise 8.9.} A toy with one positive-energy level (energy $E$, operator $b$) and one negative-energy level (energy $-E$, operator $d^\ast$) has $H=E\,b^\ast b-E\,dd^\ast$. Find the energy of the vacuum before normal ordering and the normal-ordered Hamiltonian.

\textbf{Exercise 8.10.} For EXP-5 with $m=1$ and $q=0.1$ compute $t_\ast=\ln(m/q)$ and $E_0=\sqrt{m^2-q^2}$, and compare with the values quoted in Section 8.13.

\textbf{Exercise 8.11.} A two-dimensional toy: $B=\mathrm{diag}(1,-1)$ and $h=\begin{pmatrix}0&1\\-1&0\end{pmatrix}$. Show $h^\dagger B=Bh$ and $h^2=-1$, solve $i\dot u=hu$ with $u(0)=(1,0)^T$, and compute $u^\dagger u$ and $u^\dagger Bu$ as functions of time.

\subsection{Answers to the exercises}

\textbf{Answer 8.1.} $bb^\dagger=\mathrm{diag}(1,0)$ and $b^\dagger b=\mathrm{diag}(0,1)$ add to the identity; $(b^\dagger)^2=\begin{pmatrix}0&0\\1&0\end{pmatrix}\begin{pmatrix}0&0\\1&0\end{pmatrix}=0$. $N=\mathrm{diag}(0,1)$ has the eigenvalues 0 and 1.

\textbf{Answer 8.2.} $b_0^\dagger|01\rangle=b_0^\dagger b_1^\dagger|0\rangle=|11\rangle$; $b_1b_0^\dagger|01\rangle=b_1|11\rangle=-|10\rangle$ (Section 8.4); $b_0^\dagger b_1|01\rangle=b_0^\dagger|00\rangle=|10\rangle$. The sum is 0.

\textbf{Answer 8.3.} The left derivative with respect to $\theta^\ast$ gives $2i\dot\theta-\omega\theta=0$, that is $i\dot\theta=\tfrac\omega2\theta$. Here $K=2i$, so $\{\theta,\theta^\dagger\}=i/(2i)=\tfrac12$. With $A=\tfrac12$ the computation of Section 8.6 gives $[H,\theta]=-A\omega\theta=-\tfrac\omega2\theta$ and $\dot\theta=i[H,\theta]=-\tfrac{i\omega}2\theta$, the same equation.

\textbf{Answer 8.4.} $B^\dagger=\bigl(-iC\gamma^4\bigr)^\dagger=i(\gamma^4)^TC^T=-i\gamma^4C=-iC\gamma^4=B$. $B^2=-C\gamma^4C\gamma^4=-C^2(\gamma^4)^2=-(1)(-1)=1$.

\textbf{Answer 8.5.} (a) $E^2=4+1-4=1$: oscillation with frequency 1, although $h$ is not Hermitian. (b) $E^2=1+1-1-1=0$ with $h\ne0$: linear growth. (c) $E^2=1-9=-8$: exponential growth with $\varkappa=\sqrt8=2\sqrt2\approx2.83$.

\textbf{Answer 8.6.} $(-i\gamma^4)^2=-(\gamma^4)^2=1$, $(\gamma^4\gamma^5)^2=-(\gamma^4)^2(\gamma^5)^2=-(-1)(-1)=-1$, and the cross terms cancel because $\{\gamma^4,\gamma^4\gamma^5\}=0$. So $h^2=1-1=0$, and $u(x_4)=e^{-ihx_4}u(0)=(1-ihx_4)u(0)$, which grows linearly unless $hu(0)=0$.

\textbf{Answer 8.7.} Components: $u_0=\tfrac12$, $u_4=\tfrac12$, $u_9=\tfrac i2$, $u_{13}=-\tfrac i2$. $(\gamma^4u_-)_0=-u_{13}=\tfrac i2$, so $(-i\gamma^4u_-)_0=\tfrac12$; $(\gamma^4u_-)_4=u_9=\tfrac i2$, giving $\tfrac12$; $(\gamma^4u_-)_9=-u_4=-\tfrac12$, giving $\tfrac i2$; $(\gamma^4u_-)_{13}=u_0=\tfrac12$, giving $-\tfrac i2$. So $-i\gamma^4u_-=u_-$. $(Cu_-)_0=-u_4=-\tfrac12$, $(Cu_-)_4=-u_0=-\tfrac12$, $(Cu_-)_9=u_{13}=-\tfrac i2$, $(Cu_-)_{13}=u_9=\tfrac i2$: $Cu_-=-u_-$. Hence $u_-^\dagger u_-=4\cdot\tfrac14=1$, $u_-^\dagger Bu_-=u_-^\dagger C(-i\gamma^4)u_-=u_-^\dagger Cu_-=-1$, and $u_-^\dagger(-i\gamma^4)u_-=u_-^\dagger u_-=1$.

\textbf{Answer 8.8.} Commuting with $B$: $S^{01},S^{02},S^{03},S^{04}$ (both indices $\le4$); $S^{05},S^{06},S^{07}$ anticommute. Anti-Hermitian: $S^{01},S^{02},S^{03}$ (both indices $\le3$); the other four are Hermitian. Krein condition: $S^{01},S^{02},S^{03}$ (anti-Hermitian and commuting) and $S^{05},S^{06},S^{07}$ (Hermitian and anticommuting); $S^{04}$ fails it.

\textbf{Answer 8.9.} $dd^\ast=1-d^\ast d$, so $H=E\,b^\ast b+E\,d^\ast d-E$. The vacuum ($b|0\rangle=d|0\rangle=0$) has energy $-E$, the energy of the filled level. The normal-ordered Hamiltonian is $E(b^\ast b+d^\ast d)\ge0$.

\textbf{Answer 8.10.} $t_\ast=\ln10\approx2.302585$ and $E_0=\sqrt{0.99}\approx0.994987$, the values of Section 8.13 (from \texttt{artifact\allowbreak{}s/dirac1\allowbreak{}6complex\allowbreak{}/numeric\allowbreak{}s/exp5/s\allowbreak{}ummary.j\allowbreak{}son}).

\textbf{Answer 8.11.} $h^\dagger=\begin{pmatrix}0&-1\\1&0\end{pmatrix}$, $h^\dagger B=\begin{pmatrix}0&1\\1&0\end{pmatrix}=Bh$, and $h^2=-1$, so $E^2=-1$ and $\varkappa=1$. Then $u(t)=(\cosh t-ih\sinh t)u(0)=(\cosh t,\ i\sinh t)^T$. The Hilbert norm $u^\dagger u=\cosh^2t+\sinh^2t=\cosh2t$ grows exponentially, while the Krein norm $u^\dagger Bu=\cosh^2t-\sinh^2t=1$ stays constant: the pattern of EXP-5 in two dimensions.

\section{The primordial gravitational field of the notebook}

\subsection{What this chapter does}

The author's notebook opens with a hypothesis. In the notebook's words (cell 6, quoted through the survey \texttt{handoff/\allowbreak{}surveys/\allowbreak{}survey\_n\allowbreak{}otebook-\allowbreak{}physics.\allowbreak{}md}), if superluminal inflation or deflation exists in Einstein or Einstein--Lovelock gravity, then at the time $x_4=0$, before the particles of the standard model exist, "a pair of universes with MASSES ± M is created". To express this picture the notebook writes down one particular gravitational field, the metric that it calls MatrixMetric44. In this book we call it the \textbf{primordial field}.

This chapter studies that field from zero. It is a \textbf{prescribed background}: we do not solve an equation to find the metric. We take the notebook's metric as given, compute everything about its geometry exactly, and then ask two questions. First, what would gravity need as a source in order to produce exactly this metric? Second, how does the dirac16complex field of Chapters 6 and 7 behave in it? The answer to the first question is a surprise: in 8-dimensional Einstein gravity the field needs a negative energy density, for every choice of its free function. The chapter also records four errors (errata) that the project found while checking the notebook and its own specifications. We label the three that Stage 2 found (a), (b) and (c); the fourth is erratum E4.3 of the Stage-4 specification. The letters keep them apart from the numbered errata E1 to E5 of the project contract \texttt{handoff/\allowbreak{}specs/CO\allowbreak{}NTRACT.m\allowbreak{}d} and E4.1, E4.2 and so on of the Stage-4 specification, which this and other chapters cite. Finally, the chapter prepares the ground for Part IV: the static member of the family of fields, the hidden coordinate $y$, and a construction with two mirror copies of the field glued along a wall, the \textbf{Z2 brane}.

The chapter uses the tools of Chapter 4 (metric, Christoffel symbols, curvature, vielbein, spin connection) and the field equations and energy--momentum tensor of Chapter 7. Every formula that the repository verifies is derived here step by step, and the name of the machine check is given in typewriter type, for example \texttt{P\_metric\allowbreak{}\_detG\_eq\allowbreak{}uals\_plu\allowbreak{}s\_cos2z}. Checks whose names begin with \texttt{\detokenize{P_}} belong to Stage 2; they are listed in \texttt{artifact\allowbreak{}s/dirac1\allowbreak{}6complex\allowbreak{}/primord\allowbreak{}ial-fiel\allowbreak{}d/wolfra\allowbreak{}m-primor\allowbreak{}dial-rep\allowbreak{}ort.json} (WolframScript, 126 of 126 checks true). Except for the checks named in the next sentence, they are repeated by the independent Python checker \texttt{scripts/\allowbreak{}check\_di\allowbreak{}rac16com\allowbreak{}plex\_pri\allowbreak{}mordial.\allowbreak{}py} (report \texttt{python-p\allowbreak{}rimordia\allowbreak{}l-report\allowbreak{}.json} in the same folder, 16 of 16 check groups true). Among the checks cited in this chapter, WolframScript alone performs all checks of the group \texttt{P\_notebo\allowbreak{}okCompar\allowbreak{}e\_*} (the comparison with the notebook's stored cells and the algebra of the reconstruction behind it), \texttt{P\_metric\allowbreak{}\_noteboo\allowbreak{}kCell106\allowbreak{}0Det} and \texttt{P\_spinco\allowbreak{}nn\_noteb\allowbreak{}ookCell5\allowbreak{}01OmegaM\allowbreak{}uIJEqual\allowbreak{}sMixedOm\allowbreak{}ega}. The Python checker has no notebook reader. It compares with the notebook only through transcriptions typed into the checker: the stored outputs of cells 583, 584 and 1137 (Stage-2 document §18.2) and the four coupling sets of cell 1089. Its cell-1137 comparison finds the same $\pm q$ differences as Section 9.14, but it is recorded as a measurement, not a check. The reconstruction of the $q$ term in Section 9.14 therefore rests on the WolframScript verifier alone. Checks whose names begin with \texttt{\detokenize{KS_}} belong to the exact theory of Stage 4, recorded in \texttt{artifact\allowbreak{}s/dirac1\allowbreak{}6complex\allowbreak{}/kohn-sh\allowbreak{}am/wolfr\allowbreak{}am-kohn-\allowbreak{}sham-rep\allowbreak{}ort.json} (125 of 125 true) and \texttt{artifact\allowbreak{}s/dirac1\allowbreak{}6complex\allowbreak{}/kohn-sh\allowbreak{}am/pytho\allowbreak{}n-theory\allowbreak{}-report.\allowbreak{}json} (157 of 157 true). The main text sources are \texttt{provenan\allowbreak{}ce/DIRAC\allowbreak{}16COMPLE\allowbreak{}X\_PRIMOR\allowbreak{}DIAL\_FIE\allowbreak{}LD.md} (the Stage-2 document; we cite its sections as §N) and \texttt{handoff/\allowbreak{}specs/ST\allowbreak{}AGE4\_SPE\allowbreak{}C.md} (its §1 and the errata §7 to §9).

Conventions are those of the whole book. Everything is counted from 0. The eight coordinates are $x_0,\dots,x_7$. Because this chapter follows the notebook closely, it keeps the notebook's lower-index names also inside formulas, for example $\partial_\mu=\partial/\partial x_\mu$ and $dx_\mu$, instead of the upper-index coordinates $x^\mu$ of Section 1.8. Section 1.8 allows this exception, and it changes nothing: the lower index is only part of the name and is never lowered with a metric, $\partial/\partial x_\mu$ is the same derivative as $\partial/\partial x^\mu$, and a sum over two coordinate indices that are both written low is always written out with $\sum$. The coordinate $x_0$ is the \textbf{hidden space}, $x_1,x_2,x_3$ are ordinary \textbf{3-space}, $x_4$ is the \textbf{time} in which everything evolves, and $x_5,x_6,x_7$ are the three \textbf{extra times} (the notebook calls them "superluminal deflating time"). The flat (tangent) metric is $\eta=\mathrm{diag}(+1,+1,+1,+1,-1,-1,-1,-1)$. A repeated index, one up and one down, is summed from 0 to 7 unless we say "no sum". The letter $i$ always runs over $\{1,2,3\}$ and the letter $j$ over $\{5,6,7\}$.

\subsection{The metric MatrixMetric44}

\textbf{Two abbreviations.} The notebook uses one constant $H>0$, an inverse length (a number with the unit 1/length), and the two combinations

\[
z=6Hx_0,\qquad t=Hx_4 .
\]

The coordinate $x_0$ is restricted so that $0<z<\pi/2$; in this range $\sin z$, $\cos z$, $\tan z$ and $\cot z=\cos z/\sin z$ are all positive. We write $s=\sin z$. The dimensionless time $t$ is the argument of the one free function of the metric, $a_4(t)$, an arbitrary smooth real function. Its derivatives are written with primes, and \textbf{primes always mean derivatives with respect to $t$}:

\[
a_4'=\frac{da_4}{dt},\qquad a_4''=\frac{d^2a_4}{dt^2},\qquad \partial_4a_4=\frac{\partial a_4(Hx_4)}{\partial x_4}=H\,a_4' ,
\]

the last by the chain rule. In the same way $\partial_0=6H\,\partial_z$ on any function of $z$.

\textbf{The line element.} A metric is a rule that assigns a squared length $ds^2=\sum_{\mu,\nu}g_{\mu\nu}\,dx_\mu\,dx_\nu$ to a small coordinate step $(dx_0,\dots,dx_7)$ (Chapter 4). The primordial field is \textbf{diagonal}: $g_{\mu\nu}=0$ for $\mu\ne\nu$. Its line element is

\[
\begin{aligned}
ds^2={}&\cot^2z\,dx_0^2+s^{1/3}e^{2a_4}\bigl(dx_1^2+dx_2^2+dx_3^2\bigr)-dx_4^2\\
&-s^{1/3}e^{-2a_4}\bigl(dx_5^2+dx_6^2+dx_7^2\bigr),
\end{aligned}
\]

that is,

\[
g=\mathrm{diag}\bigl(\cot^2z,\ s^{1/3}e^{2a_4},\ s^{1/3}e^{2a_4},\ s^{1/3}e^{2a_4},\ -1,\ -s^{1/3}e^{-2a_4},\ -s^{1/3}e^{-2a_4},\ -s^{1/3}e^{-2a_4}\bigr),
\]

with $z=6Hx_0$, $s=\sin(6Hx_0)$ and $a_4=a_4(Hx_4)$. This is the notebook's MatrixMetric44 (Stage-2 document §4.1). The metric depends on only two of the eight coordinates, $x_0$ and $x_4$.

\textbf{Signature.} For $0<z<\pi/2$ and real $a_4$ the four entries $g_{00},g_{11},g_{22},g_{33}$ are positive and the four entries $g_{44},\dots,g_{77}$ are negative. So the metric has the same sign pattern as $\eta$: signature (4,4), four space-like and four time-like directions (check \texttt{P\_metric\allowbreak{}\_signatu\allowbreak{}re44}).

\textbf{A worked example.} Take the point where $\sin z=3/5$. Then $\cos z=4/5$ (because $\sin^2z+\cos^2z=1$), $\cot z=4/3$ and $\tan z=3/4$. Take also $a_4=0$ there. The metric at that point is

\[
g=\mathrm{diag}\bigl(\tfrac{16}{9},\ (\tfrac35)^{1/3},\ (\tfrac35)^{1/3},\ (\tfrac35)^{1/3},\ -1,\ -(\tfrac35)^{1/3},\ -(\tfrac35)^{1/3},\ -(\tfrac35)^{1/3}\bigr),\qquad (\tfrac35)^{1/3}\approx0.8434 .
\]

Four positive and four negative entries, as claimed.

\textbf{The vielbein.} A vielbein (Chapter 4) is a matrix $e_\mu{}^a$ with $g_{\mu\nu}=e_\mu{}^a\,\eta_{ab}\,e_\nu{}^b$. For a diagonal metric the simplest choice is diagonal, $e_\mu{}^a=h_\mu\,\delta_\mu^a$ (no sum), with $h_\mu^2\,\eta_{\mu\mu}=g_{\mu\mu}$. Taking positive square roots,

\[
h=\bigl(\cot z,\ s^{1/6}e^{a_4},\ s^{1/6}e^{a_4},\ s^{1/6}e^{a_4},\ 1,\ s^{1/6}e^{-a_4},\ s^{1/6}e^{-a_4},\ s^{1/6}e^{-a_4}\bigr).
\]

Check one entry of each kind: $h_0^2\eta_{00}=\cot^2z$; $h_1^2\eta_{11}=s^{1/3}e^{2a_4}$; $h_4^2\eta_{44}=-1$; $h_5^2\eta_{55}=-s^{1/3}e^{-2a_4}$. All agree with $g$ (check \texttt{P\_metric\allowbreak{}\_vielbei\allowbreak{}nProduct}). The inverse vielbein is $e_a{}^\mu=\mathrm{diag}(1/h_\mu)$, and the curved gamma matrices of Chapter 4, $\gamma^{x_\mu}=e_a{}^\mu\gamma^a$ (no sum over $\mu$), are

\[
\gamma^{x_0}=\tan z\,\gamma^0,\qquad \gamma^{x_i}=s^{-1/6}e^{-a_4}\gamma^i,\qquad \gamma^{x_4}=\gamma^4,\qquad \gamma^{x_j}=s^{-1/6}e^{a_4}\gamma^j .
\]

Here $\gamma^0,\dots,\gamma^7$ are the constant 16 by 16 frame matrices of Chapter 2, with $\{\gamma^a,\gamma^b\}=2\eta^{ab}$. Keep the last formula in mind: the factor $s^{-1/6}e^{+a_4}$ of the extra-time gammas is exactly where the notebook made the error of Section 9.14.

\subsection{Determinant, volume and erratum (a)}

\textbf{The determinant.} The determinant of a diagonal matrix is the product of its diagonal entries. Hence

\[
\det g=\cot^2z\cdot\bigl(s^{1/3}e^{2a_4}\bigr)^3\cdot(-1)\cdot\bigl(-s^{1/3}e^{-2a_4}\bigr)^3 .
\]

Work it out factor by factor. $(s^{1/3}e^{2a_4})^3=s\,e^{6a_4}$. $(-s^{1/3}e^{-2a_4})^3=-s\,e^{-6a_4}$. The two minus signs, one from $g_{44}=-1$ and one from the cube of the extra-time factor, multiply to $+1$. The exponentials cancel, $e^{6a_4}e^{-6a_4}=1$. What is left is

\[
\det g=\cot^2z\cdot s^2=\frac{\cos^2z}{\sin^2z}\,\sin^2z=+\cos^2z,\qquad \sqrt{\lvert g\rvert}=\cos z ,
\]

where $\lvert g\rvert$ is the absolute value of $\det g$ and $\cos z>0$ on the chart. The sign is $+$ for a simple reason: there are four negative diagonal entries, an even number (checks \texttt{P\_metric\allowbreak{}\_detG\_eq\allowbreak{}uals\_plu\allowbreak{}s\_cos2z} and \texttt{P\_metric\allowbreak{}\_sqrtAbs\allowbreak{}DetG\_cos\allowbreak{}z}).

In the worked example $\sin z=3/5$: $\det g=\tfrac{16}{9}\cdot\tfrac{9}{25}=\tfrac{16}{25}=\cos^2z$, whatever $a_4$ is.

\textbf{Erratum (a): the sign of the determinant.} The Stage-2 specification \texttt{handoff/\allowbreak{}specs/ST\allowbreak{}AGE2\_SPE\allowbreak{}C.md} had written $\det g=-\cos^2z$. That is wrong; the correct value is $+\cos^2z$, and the notebook itself has it right: its stored output of cell 1060 shows the same value (check \texttt{P\_metric\allowbreak{}\_noteboo\allowbreak{}kCell106\allowbreak{}0Det}). The error did not propagate, because only $\sqrt{\lvert g\rvert}=\cos z$ enters the physics, and that is the same for both signs (Stage-2 document §4.3; the report records the discrepancy under the key \texttt{specDisc\allowbreak{}repancy\_\allowbreak{}detG}).

\textbf{Volume.} The factor $\sqrt{\lvert g\rvert}$ converts coordinate volume into proper volume (Chapter 4). Because $g_{44}=-1$, the determinant of the seven remaining directions, $x_0,x_1,x_2,x_3,x_5,x_6,x_7$, equals $-\det g$ and has absolute value $\cos^2z$ as well. So the proper 7-volume of a region of fixed coordinate size is $\cos z$ times its coordinate volume. It does not depend on $x_4$. \textbf{The 7-volume is constant in time.} Remember this fact: it is the reason why the dirac16complex field is frozen in this background (Chapter 11, experiment EXP-1).

\subsection{Inflating 3-space and deflating extra times}

A coordinate interval $\Delta x_1$ along $x_1$ has the proper length $L=\sqrt{g_{11}}\,\Delta x_1=s^{1/6}e^{a_4}\,\Delta x_1$. The number that multiplies the coordinate interval is called the \textbf{scale factor} of that direction. Its logarithmic rate of change in the time $x_4$ is the \textbf{Hubble rate} of the direction:

\[
H_1=\frac{1}{L}\frac{\partial L}{\partial x_4}=\partial_4\ln\bigl(s^{1/6}e^{a_4}\bigr)=\partial_4a_4=H\,a_4' .
\]

The same holds for $x_2$ and $x_3$. For the extra times the proper length is $\sqrt{\lvert g_{55}\rvert}\,\Delta x_5=s^{1/6}e^{-a_4}\,\Delta x_5$ and the rate is $H_5=-H\,a_4'$. So whenever $a_4'>0$, ordinary 3-space expands and the three extra times shrink at the same rate. The hidden direction does not change with time ($g_{00}$ depends on $x_0$ only).

\textbf{The notebook's example.} For $a_4=t=Hx_4$ we have $a_4'=1$: 3-space grows like $e^{Hx_4}$, exponentially, at the constant rate $H$; the notebook calls this superluminal inflation. The extra times shrink like $e^{-Hx_4}$ (deflation).

\textbf{Why the 7-volume stays constant.} The six transverse scale factors multiply to

\[
\bigl(s^{1/6}e^{a_4}\bigr)^3\bigl(s^{1/6}e^{-a_4}\bigr)^3=s ,
\]

independent of $a_4$. The inflation of 3-space is exactly compensated by the deflation of the extra times. Together with $h_0=\cot z$ this gives $\prod_{\mu\ne4}h_\mu=\cot z\cdot s=\cos z=\sqrt{\lvert g\rvert}$, the result of Section 9.3.

\subsection{The warped form and the hidden coordinate}

The coefficient $\cot^2z$ of $dx_0^2$ makes the $x_0$ direction awkward. A better coordinate measures proper distance along the hidden direction. Define

\[
\zeta=\frac{\ln\sin z}{6H} .
\]

\textbf{Step 1: the range.} For $0<z<\pi/2$, $\sin z$ runs through $(0,1)$ and $\ln\sin z$ through $(-\infty,0)$. So $\zeta\in(-\infty,0)$, and $\zeta\to0$ as $z\to\pi/2$, $\zeta\to-\infty$ as $z\to0$ (check \texttt{P\_zeta\_r\allowbreak{}ange}).

\textbf{Step 2: the differential.} By the chain rule, with $dz/dx_0=6H$,

\[
\frac{d\zeta}{dx_0}=\frac{1}{6H}\,\frac{\cos z}{\sin z}\,6H=\cot z,\qquad d\zeta=\cot z\,dx_0,\qquad d\zeta^2=\cot^2z\,dx_0^2=g_{00}\,dx_0^2 .
\]

So in the coordinate $\zeta$ the hidden direction has the coefficient 1: $\zeta$ is proper distance (check \texttt{P\_zeta\_g\allowbreak{}ZetaZeta\allowbreak{}IsOne}).

\textbf{Step 3: the warp factor.} Exponentiating the definition gives $e^{6H\zeta}=\sin z=s$, hence $s^{1/3}=e^{2H\zeta}$ and $s^{1/6}=e^{H\zeta}$ (check \texttt{P\_zeta\_w\allowbreak{}arpFacto\allowbreak{}r}).

\textbf{Step 4: the result.} Substituting Steps 2 and 3 into the line element of Section 9.2,

\[
ds^2=d\zeta^2-dx_4^2+e^{2H\zeta}\Bigl[e^{2a_4}\bigl(dx_1^2+dx_2^2+dx_3^2\bigr)-e^{-2a_4}\bigl(dx_5^2+dx_6^2+dx_7^2\bigr)\Bigr]
\]

(check \texttt{P\_zeta\_w\allowbreak{}arpedMet\allowbreak{}ric}). This is called a \textbf{warped} metric: the six transverse directions carry a common factor $e^{2H\zeta}$, the \textbf{warp factor}, that depends on the position in the hidden direction. In the coordinates $(\zeta,x_1,\dots,x_7)$ the volume factor is $\sqrt{\lvert g\rvert}=e^{6H\zeta}$: the old value $\cos z$ divided by $d\zeta/dx_0=\cot z$, which gives $\sin z=e^{6H\zeta}$.

\textbf{Where the chart ends.} At $\zeta=0$ (that is $z=\pi/2$) the $x_0$ description breaks down: $g_{00}=\cot^2z=0$ and $\sqrt{\lvert g\rvert}=\cos z=0$ there. The $\zeta$ form of the line element is perfectly regular at $\zeta=0$; the breakdown belongs to the coordinate $x_0$, not to the geometry (Section 9.16 continues this). As $\zeta\to-\infty$ the warp factor tends to 0 and the six transverse directions shrink to nothing; we call this end the \textbf{tip}.

\textbf{Worked example.} At $\sin z=3/5$ and $H=1$: $\zeta=\tfrac16\ln0.6=-0.08514$, and $e^{2\zeta}=e^{-0.17028}=0.8434=(0.6)^{1/3}$, the transverse coefficient of the example of Section 9.2.

\subsection{Christoffel symbols of a diagonal metric}

The Christoffel symbols (Chapter 4) are

\[
\Gamma^\rho{}_{\mu\nu}=\tfrac12\,g^{\rho\sigma}\bigl(\partial_\mu g_{\nu\sigma}+\partial_\nu g_{\mu\sigma}-\partial_\sigma g_{\mu\nu}\bigr).
\]

For a diagonal metric the inverse is diagonal too, $g^{\rho\rho}=1/g_{\rho\rho}$, so only $\sigma=\rho$ survives in the sum:

\[
\Gamma^\rho{}_{\mu\nu}=\frac{1}{2g_{\rho\rho}}\bigl(\partial_\mu g_{\nu\rho}+\partial_\nu g_{\mu\rho}-\partial_\rho g_{\mu\nu}\bigr)\qquad(\text{no sum over }\rho).
\]

Each of the three terms is nonzero only when its two metric indices are equal. Going through the possible coincidences of $\rho,\mu,\nu$ gives four rules.

\begin{itemize}
\item (C1) $\nu=\rho$, any $\mu$: the three terms are $\partial_\mu g_{\rho\rho}$, $+\partial_\rho g_{\mu\rho}$ and $-\partial_\rho g_{\mu\rho}$. For $\mu\ne\rho$ the last two are zero; for $\mu=\nu=\rho$ they are $+\partial_\rho g_{\rho\rho}$ and $-\partial_\rho g_{\rho\rho}$, which need not be zero (for $\Gamma^0{}_{00}$ below they are not), and they cancel. In both cases only the first term remains, $\Gamma^\rho{}_{\mu\rho}=\Gamma^\rho{}_{\rho\mu}=\partial_\mu g_{\rho\rho}/(2g_{\rho\rho})=\tfrac12\partial_\mu\ln\lvert g_{\rho\rho}\rvert$.
\item (C2) $\mu=\nu\ne\rho$: only the third term survives, $\Gamma^\rho{}_{\mu\mu}=-\partial_\rho g_{\mu\mu}/(2g_{\rho\rho})$.
\item (C3) $\rho,\mu,\nu$ all different: every term has two different metric indices, so $\Gamma^\rho{}_{\mu\nu}=0$.
\item (C4) The symbols are symmetric, $\Gamma^\rho{}_{\mu\nu}=\Gamma^\rho{}_{\nu\mu}$, because the formula is symmetric in $\mu$ and $\nu$.
\end{itemize}

\textbf{Which derivatives exist.} In the primordial field $g_{00}$ depends on $x_0$ only, $g_{44}=-1$ is constant, and the six transverse entries depend on $x_0$ and $x_4$. Their logarithmic derivatives are all we need. With $\partial_0=6H\partial_z$, $\partial_z\ln s=\cot z$ and $\partial_4a_4=Ha_4'$:

\[
\partial_0\ln g_{ii}=\partial_0\ln\lvert g_{jj}\rvert=6H\cdot\tfrac13\cot z=2H\cot z,\qquad \partial_4\ln g_{ii}=2Ha_4',\qquad \partial_4\ln\lvert g_{jj}\rvert=-2Ha_4' ,
\]

and $\partial_0\ln g_{00}=2\,\partial_0\ln\cot z=2\cdot6H\cdot\dfrac{-1/\sin^2z}{\cot z}=-\dfrac{12H}{\sin z\cos z}$.

\textbf{The symbols one by one.} Rule (C1) with $\rho=\mu=0$: $\Gamma^0{}_{00}=\tfrac12\partial_0\ln g_{00}=-6H/(\sin z\cos z)=-6H\,s^{-1}\sec z$. Rule (C1) with $\rho=i$: $\Gamma^i{}_{0i}=\tfrac12\cdot2H\cot z=H\cot z$ and $\Gamma^i{}_{4i}=\tfrac12\cdot2Ha_4'=Ha_4'$. Rule (C1) with $\rho=j$: $\Gamma^j{}_{0j}=H\cot z$ and $\Gamma^j{}_{4j}=-Ha_4'$. Rule (C2) with $\rho=0$: $\Gamma^0{}_{ii}=-\partial_0g_{ii}/(2g_{00})=-2H\cot z\,g_{ii}/(2\cot^2z)=-H\tan z\,g_{ii}$, and in the same way $\Gamma^0{}_{jj}=-H\tan z\,g_{jj}$. Rule (C2) with $\rho=4$ and $g_{44}=-1$: $\Gamma^4{}_{ii}=\partial_4g_{ii}/2=Ha_4'\,g_{ii}$ and $\Gamma^4{}_{jj}=\partial_4g_{jj}/2=-Ha_4'\,g_{jj}$. Inserting $g_{ii}=s^{1/3}e^{2a_4}$ and $g_{jj}=-s^{1/3}e^{-2a_4}$:

\begingroup
\small
\begin{longtable}{@{}>{\raggedright\arraybackslash}p{0.267\linewidth}>{\raggedright\arraybackslash}p{0.267\linewidth}>{\raggedright\arraybackslash}p{0.267\linewidth}@{}}
\toprule
\textbf{Symbol} & \textbf{Closed form} & \textbf{How many} \\
\midrule
\endfirsthead
\toprule
\textbf{Symbol} & \textbf{Closed form} & \textbf{How many} \\
\midrule
\endhead
$\Gamma^0{}_{00}$ & $-6H\,s^{-1}\sec z$ & 1 \\
$\Gamma^0{}_{ii}$ & $-H\tan z\,s^{1/3}e^{2a_4}$ & 3 \\
$\Gamma^0{}_{jj}$ & $+H\tan z\,s^{1/3}e^{-2a_4}$ & 3 \\
$\Gamma^4{}_{ii}$ & $H\,a_4'\,s^{1/3}e^{2a_4}$ & 3 \\
$\Gamma^4{}_{jj}$ & $H\,a_4'\,s^{1/3}e^{-2a_4}$ & 3 \\
$\Gamma^i{}_{0i}=\Gamma^i{}_{i0}$ & $H\cot z$ & 6 \\
$\Gamma^i{}_{4i}=\Gamma^i{}_{i4}$ & $H\,a_4'$ & 6 \\
$\Gamma^j{}_{0j}=\Gamma^j{}_{j0}$ & $H\cot z$ & 6 \\
$\Gamma^j{}_{4j}=\Gamma^j{}_{j4}$ & $-H\,a_4'$ & 6 \\
\bottomrule
\end{longtable}
\endgroup

\textbf{Everything else vanishes.} Nothing depends on $x_1,\dots,x_3,x_5,\dots,x_7$, so every rule that needs such a derivative gives 0; $g_{00}$ does not depend on $x_4$, so $\Gamma^0{}_{04}=\Gamma^4{}_{00}=0$; $g_{44}$ is constant, so every symbol with $\rho=4$ from rule (C1) vanishes. The count is $1+3+3+3+3+6+6+6+6=37$ nonzero symbols out of $8^3=512$ (25 of them with $\mu\le\nu$), each with at least one index 0 or 4 (checks \texttt{P\_christ\allowbreak{}offel\_cl\allowbreak{}osedForm\allowbreak{}s512} and \texttt{P\_christ\allowbreak{}offel\_co\allowbreak{}unt}; the same table is in the Stage-2 document §5).

\textbf{Worked example.} At $\sin z=3/5$, $a_4=0$, $a_4'=1$, $H=1$: $\Gamma^0{}_{00}=-6/(\tfrac35\cdot\tfrac45)=-12.5$; $\Gamma^0{}_{11}=-\tfrac34(\tfrac35)^{1/3}=-0.6326$; $\Gamma^1{}_{01}=\tfrac43$; $\Gamma^5{}_{45}=-1$.

\subsection{Curvature of a warped metric: one lemma for everything}

Curvature needs second derivatives and products of Christoffel symbols; done by brute force it is long. The primordial field, its static member (Section 9.15) and the mirror construction (Section 9.16) all have the same structure, so we prove one lemma once and use it three times.

\textbf{Curvature convention.} We use the convention of the repository (Stage-2 document §7.3):

\[
R^\rho{}_{\sigma\mu\nu}=\partial_\mu\Gamma^\rho{}_{\nu\sigma}-\partial_\nu\Gamma^\rho{}_{\mu\sigma}+\Gamma^\rho{}_{\mu\lambda}\Gamma^\lambda{}_{\nu\sigma}-\Gamma^\rho{}_{\nu\lambda}\Gamma^\lambda{}_{\mu\sigma},\qquad R_{\sigma\nu}=R^\rho{}_{\sigma\rho\nu},\qquad R=g^{\sigma\nu}R_{\sigma\nu} .
\]

Setting $\mu=\rho$ and renaming, the Ricci tensor is

\[
R_{\mu\nu}=\partial_\rho\Gamma^\rho{}_{\mu\nu}-\partial_\nu\Gamma^\rho{}_{\rho\mu}+\Gamma^\rho{}_{\rho\lambda}\Gamma^\lambda{}_{\mu\nu}-\Gamma^\rho{}_{\nu\lambda}\Gamma^\lambda{}_{\rho\mu} .
\]

With this convention a round sphere has positive $R$.

\textbf{Lemma 9.1 (Ricci tensor of a metric warped over a flat plane).} Split the eight coordinates into two \textbf{base} coordinates $x_A$, $A\in\{0,4\}$, and six \textbf{fibre} coordinates $x_p$, $p\in F=\{1,2,3,5,6,7\}$. Let

\[
ds^2=\epsilon_0\,dx_0^2+\epsilon_4\,dx_4^2+\sum_{p\in F}\epsilon_p\,e^{2f_p(x_0,x_4)}\,dx_p^2,
\]

where every $\epsilon$ is a constant sign $\pm1$ and the six functions $f_p$ depend on the base coordinates only, and $\partial_A=\partial/\partial x_A$. (Only the form of the metric matters, not the names: below the warped coordinate $\zeta$ of Section 9.5, or the proper coordinate $y$ of Section 9.15, takes the place of $x_0$.) Write $\Theta_A=\sum_{p\in F}\partial_Af_p$. Then the only nonzero components of the Ricci tensor are

\[
R_{AB}=-\sum_{p\in F}\bigl(\partial_A\partial_Bf_p+\partial_Af_p\,\partial_Bf_p\bigr),\qquad R^p{}_p=-\sum_{A\in\{0,4\}}\epsilon_A\bigl(\partial_A^2f_p+\Theta_A\,\partial_Af_p\bigr)\quad(\text{no sum over }p),
\]

and $R_{Ap}=0$, $R_{pq}=0$ for $p\ne q$.

\textbf{Proof.} \emph{Step 1: the Christoffel symbols.} The base entries $g_{AA}=\epsilon_A$ are constant and $1/\epsilon_A=\epsilon_A$. Rules (C1) to (C4) of Section 9.6 give exactly two kinds of nonzero symbols:

\[
\Gamma^A{}_{pp}=-\frac{\partial_A g_{pp}}{2g_{AA}}=-\epsilon_A\epsilon_p\,e^{2f_p}\,\partial_Af_p,\qquad \Gamma^p{}_{Ap}=\Gamma^p{}_{pA}=\tfrac12\partial_A\ln\lvert g_{pp}\rvert=\partial_Af_p .
\]

All others vanish: $\Gamma^A{}_{BC}=0$ and $\Gamma^A{}_{Bp}=0$ because the base entries are constant; $\Gamma^p{}_{AB}=0$ because $g_{AA}$ does not depend on $x_p$; and every symbol that needs a derivative along a fibre coordinate vanishes. The contracted symbols are therefore $\Gamma^\rho{}_{\rho A}=\sum_p\partial_Af_p=\Theta_A$ and $\Gamma^\rho{}_{\rho p}=0$.

\emph{Step 2: $R_{AB}$.} Take the four terms of the Ricci formula in order. First term: $\Gamma^\rho{}_{AB}=0$ for every $\rho$, so it is 0. Second term: $-\partial_B\Gamma^\rho{}_{\rho A}=-\partial_B\Theta_A=-\sum_p\partial_A\partial_Bf_p$. Third term: contains $\Gamma^\lambda{}_{AB}=0$. Fourth term: $\Gamma^\lambda{}_{\rho A}$ is nonzero only for $\lambda=\rho=p$, where it is $\partial_Af_p$, and then $\Gamma^\rho{}_{B\lambda}=\Gamma^p{}_{Bp}=\partial_Bf_p$; so the fourth term is $-\sum_p\partial_Bf_p\,\partial_Af_p$. Adding gives the formula for $R_{AB}$.

\emph{Step 3: $R_{pp}$.} First term: $\partial_\rho\Gamma^\rho{}_{pp}=\sum_A\partial_A\bigl(-\epsilon_A\epsilon_pe^{2f_p}\partial_Af_p\bigr)=-\epsilon_pe^{2f_p}\sum_A\epsilon_A\bigl(\partial_A^2f_p+2(\partial_Af_p)^2\bigr)$, by the product rule. Second term: $-\partial_p\Gamma^\rho{}_{\rho p}=0$. Third term: $\Gamma^\lambda{}_{pp}$ is nonzero only for $\lambda=A$, so the term is $\sum_A\Theta_A\Gamma^A{}_{pp}=-\epsilon_pe^{2f_p}\sum_A\epsilon_A\Theta_A\partial_Af_p$. Fourth term: $-\Gamma^\rho{}_{p\lambda}\Gamma^\lambda{}_{\rho p}$ has two kinds of nonzero products, $(\rho,\lambda)=(A,p)$, giving $\Gamma^A{}_{pp}\Gamma^p{}_{Ap}$, and $(\rho,\lambda)=(p,A)$, giving $\Gamma^p{}_{pA}\Gamma^A{}_{pp}$. Each equals $-\epsilon_A\epsilon_pe^{2f_p}(\partial_Af_p)^2$, so the fourth term is $+2\epsilon_pe^{2f_p}\sum_A\epsilon_A(\partial_Af_p)^2$. In the sum the terms $\pm2(\partial_Af_p)^2$ cancel:

\[
R_{pp}=-\epsilon_p\,e^{2f_p}\sum_A\epsilon_A\bigl(\partial_A^2f_p+\Theta_A\,\partial_Af_p\bigr).
\]

Raising one index with $g^{pp}=\epsilon_pe^{-2f_p}$ (and $\epsilon_p^2=1$) gives $R^p{}_p$ as stated.

\emph{Step 4: the rest.} For $R_{Ap}$: the first term is $\partial_p\Gamma^p{}_{Ap}=0$ (no fibre dependence), the second is 0, the third needs $\Gamma^\rho{}_{\rho p}=0$, and the fourth needs $\Gamma^q{}_{pq}$ with $q$ a fibre index, which vanishes. For $R_{pq}$ with $p\ne q$: every term needs a symbol with two different fibre indices, such as $\Gamma^A{}_{pq}$ or $\Gamma^p{}_{qA}$, and all of those vanish. $\square$

\subsection{The Ricci scalar and the Einstein tensor of the primordial field}

\textbf{Applying the lemma.} The warped form of Section 9.5 is of the lemma's type, with the base coordinates $\zeta$ in the place of $x_0$ ($\epsilon_0=+1$) and $x_4$ ($\epsilon_4=-1$), fibre signs $\epsilon_p=+1$ for $p=1,2,3$ and $-1$ for $p=5,6,7$, and

\[
f_p=H\zeta+\sigma_p\,a_4(Hx_4),\qquad \sigma_p=+1\ (p=1,2,3),\quad \sigma_p=-1\ (p=5,6,7).
\]

The derivatives we need are

\[
\partial_\zeta f_p=H,\qquad \partial_\zeta^2f_p=0,\qquad \partial_4f_p=\sigma_pHa_4',\qquad \partial_4^2f_p=\sigma_pH^2a_4'',\qquad \partial_\zeta\partial_4f_p=0,
\]

and the two sums $\Theta_\zeta=6H$ and $\Theta_4=Ha_4'\sum_p\sigma_p=0$. The last zero, three inflating directions balanced by three deflating ones, is the constancy of the 7-volume again.

\textbf{The Ricci tensor.} From the lemma:

\[
\begin{aligned}
R_{\zeta\zeta}&=-\sum_p H^2=-6H^2,\qquad R_{44}=-\sum_p\bigl(\sigma_pH^2a_4''+H^2a_4'^2\bigr)=-6H^2a_4'^2,\\
R_{\zeta4}&=-\sum_p\bigl(0+H\cdot\sigma_pHa_4'\bigr)=0,\\
R^p{}_p&=-\bigl[(+1)(0+6H\cdot H)+(-1)(\sigma_pH^2a_4''+0)\bigr]=-6H^2+\sigma_pH^2a_4'' ,
\end{aligned}
\]

where $\sum_p\sigma_p=0$ removed the $a_4''$ term from $R_{44}$. With $g^{\zeta\zeta}=1$ and $g^{44}=-1$ the mixed components are $R^\zeta{}_\zeta=-6H^2$ and $R^4{}_4=+6H^2a_4'^2$. All off-diagonal components vanish (check \texttt{P\_einste\allowbreak{}in\_offDi\allowbreak{}agonalZe\allowbreak{}ro}), and $R_{44}=-6H^2a_4'^2$ is check \texttt{P\_einste\allowbreak{}in\_R44}.

\textbf{The Ricci scalar.} The trace of the mixed Ricci tensor is

\[
R=R^\zeta{}_\zeta+R^4{}_4+\sum_pR^p{}_p=-6H^2+6H^2a_4'^2-36H^2+H^2a_4''\sum_p\sigma_p=6H^2\bigl(a_4'^2-7\bigr)
\]

(check \texttt{P\_einste\allowbreak{}in\_ricci\allowbreak{}Scalar}). Chapter 4 derived $R_{44}$ and $G^4{}_4$ directly and quoted this value of $R$; the lemma has now derived every component.

\textbf{The Einstein tensor.} $G^\mu{}_\nu=R^\mu{}_\nu-\tfrac12\delta^\mu_\nu R$, with $-\tfrac12R=3H^2(7-a_4'^2)=21H^2-3H^2a_4'^2$. Adding this to each diagonal entry of the mixed Ricci tensor:

\[
\begin{aligned}
G^\zeta{}_\zeta&=-6H^2+21H^2-3H^2a_4'^2=-3H^2\bigl(a_4'^2-5\bigr),\\
G^i{}_i&=-6H^2+H^2a_4''+21H^2-3H^2a_4'^2=H^2\bigl(15-3a_4'^2+a_4''\bigr)\qquad(i=1,2,3),\\
G^4{}_4&=6H^2a_4'^2+21H^2-3H^2a_4'^2=3H^2\bigl(7+a_4'^2\bigr),\\
G^j{}_j&=-6H^2-H^2a_4''+21H^2-3H^2a_4'^2=H^2\bigl(15-3a_4'^2-a_4''\bigr)\qquad(j=5,6,7).
\end{aligned}
\]

\textbf{Back to the $x_0$ chart.} A tensor with one upper and one lower index transforms under a change of coordinates with one factor $\partial x'/\partial x$ and one factor $\partial x/\partial x'$ (Chapter 4). The change $x_0\to\zeta$ touches only one coordinate, so the two factors are $d\zeta/dx_0$ and $dx_0/d\zeta$, whose product is 1. Hence $G^0{}_0=G^\zeta{}_\zeta$, and every mixed component above is also the mixed component in the notebook's coordinates. These closed forms are check \texttt{P\_einste\allowbreak{}in\_Gmixe\allowbreak{}dClosedF\allowbreak{}orms}. The Ricci scalar equals the notebook's stored cell-583 output, and the covariant tensor $G_{\mu\nu}=g_{\mu\rho}G^\rho{}_\nu$ equals the stored cell-584 output entry by entry (checks \texttt{P\_einste\allowbreak{}in\_noteb\allowbreak{}ookCell5\allowbreak{}83} and \texttt{P\_einste\allowbreak{}in\_noteb\allowbreak{}ookCell5\allowbreak{}84}; Stage-2 document §15.1). The notebook computed this much correctly. Nothing here depends on $x_0$.

\textbf{A consistency test: the Bianchi identity.} Every Einstein tensor satisfies $\nabla_\mu G^\mu{}_\nu=0$ identically (Chapter 4). For a diagonal $G$ this reads $\nabla_\mu G^\mu{}_\nu=\partial_\nu G^\nu{}_\nu+\Gamma^\mu{}_{\mu\nu}G^\nu{}_\nu-\sum_\mu\Gamma^\mu{}_{\mu\nu}G^\mu{}_\mu$ (no sum over $\nu$). Test it for $\nu=4$ in the $\zeta$ chart, where $\Gamma^p{}_{p4}=\sigma_pHa_4'$, $\Gamma^\zeta{}_{\zeta4}=\Gamma^4{}_{44}=0$ and hence $\Gamma^\mu{}_{\mu4}=\Theta_4=0$:

\[
\nabla_\mu G^\mu{}_4=\partial_4\bigl[3H^2(7+a_4'^2)\bigr]-Ha_4'\sum_p\sigma_pG^p{}_p=6H^3a_4'a_4''-Ha_4'\cdot6H^2a_4''=0 ,
\]

using $\partial_4(a_4'^2)=2a_4'a_4''\,H$ and $\sum_p\sigma_pG^p{}_p=3(G^i{}_i-G^j{}_j)=6H^2a_4''$. Exercise 9.4 does $\nu=\zeta$.

\textbf{Worked example: the notebook's $a_4=t$.} Then $a_4'=1$ and $a_4''=0$, and

\[
R=-36H^2,\qquad G^\mu{}_\nu=\mathrm{diag}(12,12,12,12,24,12,12,12)\,H^2 .
\]

\subsection{Einstein and Einstein--Lovelock field equations}

\textbf{Einstein's equations in eight dimensions.} Gravity is the metric; its field equations relate the curvature of the metric to the energy and momentum of matter. In $D$ dimensions Einstein's equations read

\[
G_{\mu\nu}=\kappa\,T_{\mu\nu},\qquad G_{\mu\nu}=R_{\mu\nu}-\tfrac12g_{\mu\nu}R ,
\]

where $T_{\mu\nu}$ is the energy--momentum tensor of the matter (Chapter 7) and $\kappa>0$ is the gravitational coupling (in four dimensions $\kappa=8\pi G_N$; in eight dimensions it is a new constant, $\kappa_8$, which we simply call $\kappa$). They follow from the Einstein--Hilbert action $\tfrac1{2\kappa}\int\sqrt{\lvert g\rvert}\,R\,d^Dx$ plus the matter action, with $T_{\mu\nu}=-(2/\sqrt{\lvert g\rvert})\,\delta I_{\mathrm{matter}}/\delta g^{\mu\nu}$, the sign convention of Chapter 7; we quote this standard variational result without deriving it here. Two consequences are used below.

\begin{enumerate}
\item \textbf{Conservation.} The left side obeys the Bianchi identity $\nabla_\mu G^\mu{}_\nu=0$ (Section 9.8 tested it), so the equations can only be solved by a source with $\nabla_\mu T^\mu{}_\nu=0$.
\item \textbf{Trace-reversed form.} Taking the trace, $G^\mu{}_\mu=R-\tfrac D2R=-\tfrac{D-2}{2}R=\kappa T$ with $T=T^\mu{}_\mu$, so $R=-2\kappa T/(D-2)$ and, substituting back, $R_{\mu\nu}=\kappa\bigl(T_{\mu\nu}-g_{\mu\nu}T/(D-2)\bigr)$. For $D=8$ the factor is $\tfrac16$; it is the $\tfrac16$ in the evolution equations of Chapter 11.
\end{enumerate}

\textbf{Why Einstein--Lovelock.} Chapter 4 (Section 4.9) introduced the Einstein--Lovelock equations; we recall what this chapter needs, in the same normalization and with the same tensors $E^{(k)\mu}{}_\nu$. Einstein's tensor is not the only possibility. David Lovelock proved in 1971 (we quote the theorem, we do not prove it) that in $D$ dimensions the most general symmetric tensor $E^\mu{}_\nu$ that is built from the metric and its first and second derivatives only, and that is divergence-free for every metric, is a sum

\[
E^\mu{}_\nu=\sum_{k\ge0}c_k\,E^{(k)\mu}{}_\nu,\qquad E^{(k)\mu}{}_\nu=-\frac{1}{2^{k+1}}\,\delta^{\mu\,\alpha_1\beta_1\cdots\alpha_k\beta_k}_{\nu\,\gamma_1\delta_1\cdots\gamma_k\delta_k}\,R^{\gamma_1\delta_1}{}_{\alpha_1\beta_1}\cdots R^{\gamma_k\delta_k}{}_{\alpha_k\beta_k},
\]

with constant coefficients $c_k$. (Chapter 4 calls them $\alpha_k$; we write $c_k$ because the letters $\alpha_1,\alpha_2,\dots$ serve as indices in the formula above. A cosmological constant $\Lambda$ corresponds to $c_0=-2\Lambda$, since $c_0E^{(0)\mu}{}_\nu=\Lambda\delta^\mu_\nu$ by the order-0 formula below.) Here $R^{\gamma\delta}{}_{\alpha\beta}=g^{\delta\sigma}R^\gamma{}_{\sigma\alpha\beta}$, and the \textbf{generalized Kronecker delta} $\delta^{\mu_1\cdots\mu_n}_{\nu_1\cdots\nu_n}$ is the determinant of the $n$ by $n$ matrix whose entry in row $r$ and column $c$ is $\delta^{\mu_r}_{\nu_c}$. For $n=2$, for example, $\delta^{\mu\nu}_{\alpha\beta}=\delta^\mu_\alpha\delta^\nu_\beta-\delta^\mu_\beta\delta^\nu_\alpha$. Such a determinant is zero whenever two of its upper (or two of its lower) indices are equal, because two rows (or columns) are then equal.

\textbf{Order 0 and order 1.} For $k=0$ the formula gives $E^{(0)\mu}{}_\nu=-\tfrac12\delta^\mu_\nu$: a cosmological-constant term. For $k=1$ expand the 3 by 3 determinant along its first row:

\[
\delta^{\mu\alpha\beta}_{\nu\gamma\delta}=\delta^\mu_\nu\,\delta^{\alpha\beta}_{\gamma\delta}-\delta^\mu_\gamma\,\delta^{\alpha\beta}_{\nu\delta}+\delta^\mu_\delta\,\delta^{\alpha\beta}_{\nu\gamma}.
\]

Contract with $R^{\gamma\delta}{}_{\alpha\beta}$ and use $R^{\gamma\delta}{}_{\gamma\beta}=R^\delta{}_\beta$ (the Ricci tensor with one index raised) and the antisymmetry of the curvature in each pair. The first term gives $\delta^\mu_\nu(R^{\alpha\beta}{}_{\alpha\beta}-R^{\beta\alpha}{}_{\alpha\beta})=2R\,\delta^\mu_\nu$. The second gives $-(R^{\mu\delta}{}_{\nu\delta}-R^{\mu\delta}{}_{\delta\nu})=-2R^\mu{}_\nu$. The third gives $R^{\gamma\mu}{}_{\nu\gamma}-R^{\gamma\mu}{}_{\gamma\nu}=-2R^\mu{}_\nu$. Altogether

\[
E^{(1)\mu}{}_\nu=-\tfrac14\bigl(2R\,\delta^\mu_\nu-4R^\mu{}_\nu\bigr)=R^\mu{}_\nu-\tfrac12\delta^\mu_\nu R=G^\mu{}_\nu :
\]

the order-1 Lovelock tensor is Einstein's tensor.

\textbf{Higher orders.} For $k=2$ the tensor is quadratic in the curvature; it comes from the Lagrangian $L_2=R^2-4R_{\mu\nu}R^{\mu\nu}+R_{\mu\nu\rho\sigma}R^{\mu\nu\rho\sigma}$, called the \textbf{Gauss--Bonnet} combination. It is the value of $\tfrac14\delta^{\mu_1\nu_1\mu_2\nu_2}_{\alpha_1\beta_1\alpha_2\beta_2}R^{\alpha_1\beta_1}{}_{\mu_1\nu_1}R^{\alpha_2\beta_2}{}_{\mu_2\nu_2}$; the 24 permutations in the determinant fall into three families: the 4 that keep the two index pairs separate give $4R^2$, the 4 that exchange the pairs give $4R_{\mu\nu\rho\sigma}R^{\mu\nu\rho\sigma}$, and each of the 16 that mix the pairs gives $-R^\mu{}_\nu R^\nu{}_\mu$ (Exercise 9.6 checks one of them). For $k=3$ the tensor is cubic. The generalized delta in $E^{(k)}$ has $2k+1$ upper indices; in $D$ dimensions two of them must coincide when $2k+1>D$, and then the delta is zero. So in $D=8$ only $k=0,1,2,3$ can contribute. This is the notebook's count "$m-1=8/2-1=3$" (cell 14, quoted from the survey \texttt{handoff/\allowbreak{}surveys/\allowbreak{}survey\_n\allowbreak{}otebook-\allowbreak{}physics.\allowbreak{}md}).

\textbf{The notebook's plan and what was done.} The notebook (cell 14) writes the vacuum Einstein--Lovelock equations as $0=-\Lambda+H^{-2}w_1\,\mathrm{Lovelock1}+H^{-4}w_2\,\mathrm{Lovelock2}+H^{-6}w_3\,\mathrm{Lovelock3}$ with pure numbers $w_1,w_2,w_3,\Lambda$. The powers of $H$ are there because $E^{(k)}$ has the unit $1/\mathrm{length}^{2k}$; the notebook calls $H$ "a fundamental inverse length" that exists because the Lovelock tensors have different units. According to the survey, \textbf{no code in the notebook defines Lovelock2 or Lovelock3}; its curvature code stops at the Einstein tensor (cells 583 and 584). The project does not compute them either (Stage-2 document §1, non-claim 4, and §20). Everything below uses \textbf{8-dimensional Einstein gravity only}, $c_1=1$ and all other $c_k=0$ (Einstein gravity without a cosmological constant). What the order-2 and order-3 terms would change is an open problem (Chapter 18).

\subsection{The source the field requires}

Suppose the primordial field is produced by some matter through $G^\mu{}_\nu=\kappa T^\mu{}_\nu$. Then the matter must have $T^\mu{}_\nu=G^\mu{}_\nu/\kappa$. We read off its \textbf{energy density} and \textbf{pressures} for the observer who moves along $x_4$ (Chapter 7): $\rho=T_{44}=-T^4{}_4$ (because $g_{44}=-1$) and $p_{(\mu)}=T^\mu{}_\mu$ (no sum) for each of the seven transverse directions $\mu\ne4$. From Section 9.8:

\[
\begin{aligned}
\rho_{\mathrm{req}}&=-\frac{3H^2\bigl(7+a_4'^2\bigr)}{\kappa}, &\qquad p_{(0)}&=-\frac{3H^2\bigl(a_4'^2-5\bigr)}{\kappa},\\
p_{(1,2,3)}&=\frac{H^2\bigl(15-3a_4'^2+a_4''\bigr)}{\kappa}, &\qquad p_{(5,6,7)}&=\frac{H^2\bigl(15-3a_4'^2-a_4''\bigr)}{\kappa}
\end{aligned}
\]

(checks \texttt{P\_einste\allowbreak{}in\_requi\allowbreak{}redSourc\allowbreak{}e} and \texttt{P\_einste\allowbreak{}in\_rhoRe\allowbreak{}quiredNe\allowbreak{}gative}). Since $a_4'^2\ge0$,

\[
\rho_{\mathrm{req}}\le-\frac{21H^2}{\kappa}<0\qquad\text{for every function }a_4 .
\]

\textbf{The primordial field needs a negative energy density.} This is the main physical fact of the chapter, and it does not depend on any choice.

\textbf{Not a vacuum and not a cosmological constant.} A vacuum solution would need $G^\mu{}_\nu=0$, impossible because $G^4{}_4=3H^2(7+a_4'^2)>0$. A pure cosmological constant would need $G^\mu{}_\nu=-\Lambda\,\delta^\mu_\nu$, all eight diagonal entries equal; already $G^0{}_0=G^4{}_4$ would require $-3H^2(a_4'^2-5)=3H^2(7+a_4'^2)$, that is $-a_4'^2+5=7+a_4'^2$, that is $a_4'^2=-1$, impossible for real $a_4$ (Stage-2 document §3). So the notebook's hope, recorded in its cell 23 (survey), that the source tensor of its spinor field be $\Lambda g$ cannot be realized in Einstein gravity.

\textbf{Linear $a_4$.} If $a_4''=0$, so that $a_4=c\,t$ plus a constant, all seven pressures coincide, $p=H^2(15-3c^2)/\kappa$, and the equation of state $w=p/\rho$ is

\[
w=\frac{H^2(15-3c^2)}{-3H^2(7+c^2)}=\frac{c^2-5}{c^2+7}.
\]

For the notebook's $a_4=t$ ($c=1$): $\rho_{\mathrm{req}}=-24H^2/\kappa$, $p=12H^2/\kappa$ in all seven directions, $w=-\tfrac12$ (check \texttt{P\_a4line\allowbreak{}ar\_value\allowbreak{}s}). Note that $w<0$ here comes from $\rho<0$ with $p>0$, not from a negative pressure; the usual reading of $w$ (for example "$w<-1/3$ means acceleration") does not carry over.

\textbf{The slopes of the notebook's cell 150.} Cell 150 of the notebook lists, without derivation, the slopes $a_4'=\tfrac23(M-1)$ and $a_4'=\tfrac23(M+1)$ with $a_4''=0$ (survey). Write both as $a_4'=\tfrac23(M\mp1)$. Then $a_4'^2=\tfrac49(M\mp1)^2$ and

\[
\kappa\,\rho_{\mathrm{req}}=-3H^2\Bigl(7+\tfrac49(M\mp1)^2\Bigr)=-\tfrac13H^2\bigl(4M^2\mp8M+67\bigr)=-\tfrac13H^2\bigl(4(M\mp1)^2+63\bigr)<0
\]

for every real $M$ (check \texttt{P\_a4line\allowbreak{}ar\_cell1\allowbreak{}50Altern\allowbreak{}ativesRh\allowbreak{}oNegativ\allowbreak{}e}).

\textbf{A picture.} Experiment EXP-1 of Stage 3 (Chapter 11) used a smooth window for $a_4'(t)$, rising from 0 to a plateau $A$ and falling back, and recorded $\rho_{\mathrm{req}}$ along the way. On the plateau $A=1$ the formula gives $-24$ and on $A=2$ it gives $-3(7+4)=-33$ (units $H=\kappa=1$); outside the window $-21$. The committed summary \texttt{artifact\allowbreak{}s/dirac1\allowbreak{}6complex\allowbreak{}/numeric\allowbreak{}s/exp1/s\allowbreak{}ummary.j\allowbreak{}son} records the largest value over the whole run, $-21.000000000113$, still negative.

\begin{figure}[htbp]
\centering
\includegraphics[width=0.92\linewidth]{artifacts/dirac16complex/numerics/figures/exp1_einstein_requirement.png}
\caption{EXP-1 background and Einstein requirement, from Stage 3. (a) the window \texorpdfstring{$a_4'(t)$}{a\_4'(t)} for plateaus \texorpdfstring{$A=1$}{A=1} and \texorpdfstring{$A=2$}{A=2}; (b) the 3-space factor \texorpdfstring{$e^{a_4}$}{e\textasciicircum{}a\_4}, the extra-time factor \texorpdfstring{$e^{-a_4}$}{e\textasciicircum{}-a\_4} and the constant 7-volume; (c) the required energy density \texorpdfstring{$\rho_{\mathrm{req}}$}{ρ\_req} (between \texorpdfstring{$-21$}{-21} and \texorpdfstring{$-24$}{-24} for \texorpdfstring{$A=1$}{A=1}, between \texorpdfstring{$-21$}{-21} and \texorpdfstring{$-33$}{-33} for \texorpdfstring{$A=2$}{A=2}) beside the positive energy density of a free dirac16complex mode; (d) \texorpdfstring{$w_{\mathrm{req}}=\bar p_{\mathrm{req}}/\rho_{\mathrm{req}}$}{w\_req=p\_req/ρ\_req}, where \texorpdfstring{$\bar p_{\mathrm{req}}=\tfrac17\sum_{\mu\ne4}p_{(\mu)}=H^2(15-3a_4'^2)/\kappa$}{p\_req=17∑\_μ≠4p\_(μ)=H\textasciicircum{}2(15-3a\_4'\textasciicircum{}2)/κ} is the mean of the seven required pressures (for linear \texorpdfstring{$a_4$}{a\_4} it equals the common pressure).}
\end{figure}

\subsection{Energy conditions}

An \textbf{energy condition} is an inequality that ordinary matter is expected to satisfy. We evaluate them in the orthonormal frame of the vielbein, where the metric is $\eta$ and the components of $T$ are $T_{\hat4\hat4}=\rho$, $T_{\hat\mu\hat\mu}=p_{(\mu)}$ for a space-like $\mu\in\{0,1,2,3\}$ and $T_{\hat\mu\hat\mu}=-p_{(\mu)}$ for a time-like $\mu\in\{5,6,7\}$ (because $\eta_{\mu\mu}=-1$ there).

\begin{itemize}
\item \textbf{Weak energy condition (WEC)}: $T(u,u)\ge0$ for time-like $u$; for $u=e_4$ it is $\rho\ge0$.
\item \textbf{Null energy condition (NEC)}: $T(k,k)\ge0$ for every null vector $k$ (one with $\eta(k,k)=0$). A null vector is formed from one time-like and one space-like frame vector, for example $k=e_4+e_0$, since $\eta(k,k)=-1+1=0$. Then $T(k,k)=T_{\hat4\hat4}+T_{\hat0\hat0}=\rho+p_{(0)}$. With a time-like $e_j$ instead of $e_4$, $T(e_j+e_0,e_j+e_0)=-p_{(j)}+p_{(0)}$.
\item \textbf{Strong energy condition (SEC)}: in $D$ dimensions, $\bigl(T_{\mu\nu}-g_{\mu\nu}T/(D-2)\bigr)u^\mu u^\nu\ge0$; by the trace-reversed Einstein equations of Section 9.9 this is $R_{\mu\nu}u^\mu u^\nu\ge0$ (time-like convergence).
\item \textbf{Dominant energy condition (DEC)}: WEC plus the requirement that energy does not flow faster than light; it needs $\rho\ge0$ in particular.
\end{itemize}

Inserting the required source:

\begingroup
\small
\begin{longtable}{@{}>{\raggedright\arraybackslash}p{0.267\linewidth}>{\raggedright\arraybackslash}p{0.267\linewidth}>{\raggedright\arraybackslash}p{0.267\linewidth}@{}}
\toprule
\textbf{Condition} & \textbf{Quantity} & \textbf{Status} \\
\midrule
\endfirsthead
\toprule
\textbf{Condition} & \textbf{Quantity} & \textbf{Status} \\
\midrule
\endhead
WEC & $\rho=-\tfrac{3H^2}{\kappa}\bigl(7+a_4'^2\bigr)$ & violated for every $a_4$ \\
NEC, $k=e_4+e_0$ & $\rho+p_{(0)}=-\tfrac{6H^2}{\kappa}\bigl(1+a_4'^2\bigr)$ & violated for every $a_4$ \\
NEC, $k=e_4+e_i$ & $\rho+p_{(i)}=-\tfrac{H^2}{\kappa}\bigl(6+6a_4'^2-a_4''\bigr)$ & violated unless $a_4''\ge6(1+a_4'^2)$ \\
NEC, $k=e_j+e_0$ & $p_{(0)}-p_{(j)}=\tfrac{H^2}{\kappa}a_4''$ & holds if and only if $a_4''\ge0$ \\
NEC, $k=e_j+e_i$ & $p_{(i)}-p_{(j)}=\tfrac{2H^2}{\kappa}a_4''$ & holds if and only if $a_4''\ge0$ \\
SEC, $u=e_4$ & $R_{44}=-6H^2a_4'^2$ & violated whenever $a_4'\ne0$ \\
DEC & needs $\rho\ge0$ & violated \\
\bottomrule
\end{longtable}
\endgroup

Each line is two or three lines of algebra; for example $\rho+p_{(0)}=\tfrac{H^2}{\kappa}\bigl[-21-3a_4'^2-3a_4'^2+15\bigr]=-\tfrac{6H^2}{\kappa}(1+a_4'^2)$. The forms are check \texttt{P\_einste\allowbreak{}in\_energ\allowbreak{}yConditi\allowbreak{}onForms} (Stage-2 document §15.3). With the trace-reversed equation and $g_{44}=-1$, $T=-\rho+\sum_{\mu\ne4}p_{(\mu)}$ and $R_{44}=\kappa\bigl(\rho+T/6\bigr)$, so the strong condition can also be written $5\rho+\sum_{\mu\ne4}p_{(\mu)}=6R_{44}/\kappa=-36H^2a_4'^2/\kappa$. \textbf{Caution:} in signature (4,4) there are four time-like directions, so these are the conditions for the chosen observer and the chosen null vectors, not a complete classification.

\subsection{Which dirac16complex states can be the source}

Could the dirac16complex field itself be the negative-energy source? The Stage-2 document (§15.5) answers this completely when $a_4''\ne0$: then no state that depends on $x_0$ and $x_4$ only can be the source (Proposition 9.2 below). For linear $a_4$ it analyses two families, the real-$K$ plane waves (Proposition 9.3) and the $x_0$-independent state (Proposition 9.4). Other states of $x_0$ and $x_4$ are not analysed, for example the waves $s^{-1/2+iK/(6H)}u(x_4)$ of Proposition 9.3 with a complex $K$ whose imaginary part is $\mathrm{Im}\,K=-3H$ and whose real part is $\mathrm{Re}\,K\ne0$ (Stage-2 document §20, item 2). Throughout, the bilinears are read as ordinary numbers (the classical, mean-field reading of Chapter 7). We need the energy--momentum tensor of Chapter 7:

\[
T_{\mu\nu}=-\tfrac14\Bigl[\bar\Psi\gamma_\mu D_\nu\Psi+\bar\Psi\gamma_\nu D_\mu\Psi-(D_\mu\bar\Psi)\gamma_\nu\Psi-(D_\nu\bar\Psi)\gamma_\mu\Psi\Bigr]+g_{\mu\nu}\mathcal L_s ,
\]

where $\gamma_\mu=g_{\mu\nu}\gamma^{x_\nu}$ are the lowered curved gammas, $D_\mu\Psi=\partial_\mu\Psi+\Omega_\mu\Psi$, $D_\mu\bar\Psi=\partial_\mu\bar\Psi-\bar\Psi\Omega_\mu$ with the spinor connection $\Omega_\mu$ of Section 9.13, and $\mathcal L_s$ is the Lagrangian divided by $\sqrt{\lvert g\rvert}$; on shell $\mathcal L_s=SU'(S)-U(S)$ with $S=\bar\Psi\Psi$ and $U(S)=\tfrac\lambda2S^2$.

\textbf{Proposition 9.2 (no source when $a_4''\ne0$).} No dirac16complex state that depends on $x_0$ and $x_4$ only satisfies $G^\mu{}_\nu=\kappa T^\mu{}_\nu$ on an open set where $a_4''\ne0$.

\textbf{Proof.} Take a transverse direction $p\in\{1,2,3,5,6,7\}$. The state does not depend on $x_p$, so $D_p\Psi=\Omega_p\Psi$ and $D_p\bar\Psi=-\bar\Psi\Omega_p$, and the $pp$ component of $T$ is

\[
T_{pp}=-\tfrac14\bigl[2\bar\Psi\gamma_p\Omega_p\Psi+2\bar\Psi\Omega_p\gamma_p\Psi\bigr]+g_{pp}\mathcal L_s=-\tfrac12\bar\Psi\{\gamma_p,\Omega_p\}\Psi+g_{pp}\mathcal L_s .
\]

By Section 9.13, $\Omega_p$ is a combination of $S^{0p}$ and $S^{4p}$, and $\gamma_p$ is a multiple of $\gamma^p$. Now $S^{0p}=\tfrac12\gamma^0\gamma^p$ (because $\gamma^0$ and $\gamma^p$ anticommute), and $\{\gamma^p,\gamma^0\gamma^p\}=\gamma^p\gamma^0\gamma^p+\gamma^0\gamma^p\gamma^p=-\gamma^0(\gamma^p)^2+\gamma^0(\gamma^p)^2=0$; the same holds for $S^{4p}$. So $\{\gamma_p,\Omega_p\}=0$ and $T^p{}_p=g^{pp}T_{pp}=\mathcal L_s$ for all six transverse directions, off shell and for every $U$ (check \texttt{P\_source\allowbreak{}\_transve\allowbreak{}rsePress\allowbreak{}uresEqua\allowbreak{}lForEver\allowbreak{}yX0X4Sta\allowbreak{}te}). In particular $T^i{}_i-T^j{}_j=0$. But $G^i{}_i-G^j{}_j=2H^2a_4''$ by Section 9.8 (check \texttt{P\_source\allowbreak{}\_einstei\allowbreak{}nTransve\allowbreak{}rseDiffe\allowbreak{}renceIs2\allowbreak{}H2a4pp}). The equation $\kappa\cdot0=2H^2a_4''$ fails wherever $a_4''\ne0$. $\square$

\textbf{Proposition 9.3 (no real-$K$ plane wave).} A plane wave in the hidden coordinate, $\Psi=e^{-3H\zeta}e^{iK\zeta}u(x_4)=s^{-1/2+iK/(6H)}u(x_4)$ with real wave number $K$, is not a source for any $a_4$.

\textbf{Proof.} For real $K$ the modulus of the prefactor is $\lvert s^{-1/2+iK/(6H)}\rvert^2=s^{-1}$, so every bilinear $\bar\Psi X\Psi$ is a function of $x_4$ divided by $\sin z$. The energy density is then (Stage-2 document §15.5, check \texttt{P\_source\allowbreak{}\_realKDi\allowbreak{}agonalIs\allowbreak{}C1OverSP\allowbreak{}lusC2Ove\allowbreak{}rS2})

\[
\rho=\frac{m\,u^\dagger Cu-iK\,u^\dagger C\gamma^0u}{\sin z}+\frac{\lambda\,(u^\dagger Cu)^2}{2\sin^2z},
\]

of the form $c_1/\sin z+c_2/\sin^2z$ at each $x_4$, where $c_1$ and $c_2$ depend on $x_4$ only. The required $\rho_{\mathrm{req}}=-3H^2(7+a_4'^2)/\kappa$ of Section 9.10, on the other hand, does not depend on $z$ at all. Fix $x_4$ and suppose that $\rho=\rho_{\mathrm{req}}$ held for all $z$ in $(0,\pi/2)$, or only on some interval of $z$. Put $v=1/\sin z$. As $z$ runs through $(0,\pi/2)$, $\sin z$ runs through $(0,1)$ and $v$ through $(1,\infty)$, and an interval of $z$ gives an interval of $v$. Multiplied by $\kappa$, the equation $\rho=\rho_{\mathrm{req}}$ reads

\[
3H^2\bigl(7+a_4'^2\bigr)+\kappa c_1\,v+\kappa c_2\,v^2=0
\]

for infinitely many values of $v$. The left side is a polynomial in $v$ of degree at most 2. Such a polynomial, unless all its coefficients are zero, has at most two roots: a nonzero constant has none, a polynomial of degree 1 has one, and one of degree 2 has at most two (school algebra: the quadratic formula). So all three coefficients vanish, in particular the constant one: $3H^2(7+a_4'^2)=0$. That is impossible, because $7+a_4'^2\ge7$ (check \texttt{P\_source\allowbreak{}\_realKPl\allowbreak{}aneWaveC\allowbreak{}annotSou\allowbreak{}rce}). $\square$

\textbf{Is the plane wave a solution?} The proof uses only the form of $\rho$, which holds off shell, that is, for every $u(x_4)$ (Stage-2 document §13.2). But a source must also solve its own field equation, so we ask when the plane wave does. For this ansatz $\partial_0=\cot z\,\partial_\zeta$ (Section 9.5) gives $\tan z\,\gamma^0\partial_0\Psi=(-3H+iK)\gamma^0\Psi$, hence $\tan z\,\gamma^0\partial_0\Psi+3H\gamma^0\Psi=iK\gamma^0\Psi$. The Dirac equation of Section 9.13 for fields of $(x_0,x_4)$ becomes, after division by the common prefactor $s^{-1/2+iK/(6H)}$,

\[
\gamma^4\partial_4u=\bigl(M_{\mathrm{eff}}-iK\gamma^0\bigr)u,\qquad M_{\mathrm{eff}}=m+\lambda S .
\]

For $\lambda=0$ ($M_{\mathrm{eff}}=m$) this is an equation for $u(x_4)$ alone, and every solution $u$ of it gives an exact solution $\Psi$. For $\lambda\ne0$ the plane wave is not an exact solution, because $S=u^\dagger Cu/\sin z$ makes $M_{\mathrm{eff}}$ depend on $z$ (unless $u^\dagger Cu=0$), while the left side does not. Proposition 9.3 then holds in the pointwise mean-field reading of Stage 2, in which $M_{\mathrm{eff}}$ is treated as a given number at each point (Stage-2 document §13.1 and §20, item 3).

\textbf{Proposition 9.4 (an exact source for linear $a_4$).} Let $a_4=c\,t$ plus a constant. The state that does not depend on $x_0$, $\Psi=u(x_4)$ (the plane wave above with the imaginary value $K=-3iH$, which makes the prefactor equal to 1), solves the Dirac equation exactly for every $\lambda$ with a constant $S=u^\dagger Cu$ (check \texttt{P\_source\allowbreak{}\_x0Indep\allowbreak{}endentSt\allowbreak{}ateSolve\allowbreak{}sDiracEx\allowbreak{}actly}). Its energy--momentum tensor on shell has $T^4{}_4=-(mS+U)$ and $T^\mu{}_\mu=SU'-U$ for all seven $\mu\ne4$ (check \texttt{P\_source\allowbreak{}\_x0Indep\allowbreak{}endentDi\allowbreak{}agonalOn\allowbreak{}Shell}), and its off-diagonal components are multiples of 15 three-gamma bilinears (check \texttt{P\_source\allowbreak{}\_x0Indep\allowbreak{}endentOf\allowbreak{}fDiagona\allowbreak{}lAre15Bi\allowbreak{}linears}). It satisfies $G^\mu{}_\nu=\kappa T^\mu{}_\nu$ in all 64 components if and only if (1) those bilinears vanish (for $c=0$ only 6 of them are needed), (2) $mS=-36H^2/\kappa$ and (3) $\lambda S^2=2H^2(15-3c^2)/\kappa$.

\textbf{Derivation of (2) and (3).} With $U=\tfrac\lambda2S^2$, $U'=\lambda S$ and $SU'-U=\tfrac\lambda2S^2$. The $\zeta\zeta$ equation, $G^0{}_0=\kappa T^0{}_0$, reads $-3H^2(c^2-5)=\kappa\tfrac\lambda2S^2$, that is $H^2(15-3c^2)=\tfrac\kappa2\lambda S^2$: condition (3). The $44$ equation, $G^4{}_4=\kappa T^4{}_4$, reads $3H^2(7+c^2)=-\kappa mS-\tfrac\kappa2\lambda S^2=-\kappa mS-H^2(15-3c^2)$, so $\kappa mS=-21H^2-3H^2c^2-15H^2+3H^2c^2=-36H^2$: condition (2). The six transverse equations read $H^2(15-3c^2\pm a_4'')=\tfrac\kappa2\lambda S^2=H^2(15-3c^2)$ and hold because $a_4''=0$ (check \texttt{P\_source\allowbreak{}\_x0Indep\allowbreak{}endentSo\allowbreak{}urceCond\allowbreak{}itions}). The off-diagonal equations are condition (1), since $G$ is diagonal.

\textbf{The energy density is still negative.} For such a source $\rho=mS+\tfrac\lambda2S^2=\tfrac1\kappa\bigl[-36H^2+H^2(15-3c^2)\bigr]=-3H^2(7+c^2)/\kappa$, exactly $\rho_{\mathrm{req}}$. It is possible because $C$ has eigenvalues $+1$ and $-1$, so $S=u^\dagger Cu$ can have either sign.

\textbf{Two exact examples}, with the effective mass $M_{\mathrm{eff}}=m+\lambda S$ (units $H=\kappa=1$; check \texttt{P\_source\allowbreak{}\_x0Indep\allowbreak{}endentEx\allowbreak{}actExamp\allowbreak{}les}, Stage-2 document §15.5):

\begingroup
\small
\begin{longtable}{@{}>{\raggedright\arraybackslash}p{0.100\linewidth}>{\raggedright\arraybackslash}p{0.100\linewidth}>{\raggedright\arraybackslash}p{0.100\linewidth}>{\raggedright\arraybackslash}p{0.100\linewidth}>{\raggedright\arraybackslash}p{0.100\linewidth}>{\raggedright\arraybackslash}p{0.100\linewidth}>{\raggedright\arraybackslash}p{0.100\linewidth}>{\raggedright\arraybackslash}p{0.100\linewidth}@{}}
\toprule
\textbf{Case} & \textbf{$a_4$} & \textbf{$m$} & \textbf{$\lambda$} & \textbf{$S$} & \textbf{$M_{\mathrm{eff}}$} & \textbf{$\rho$} & \textbf{$p$} \\
\midrule
\endfirsthead
\toprule
\textbf{Case} & \textbf{$a_4$} & \textbf{$m$} & \textbf{$\lambda$} & \textbf{$S$} & \textbf{$M_{\mathrm{eff}}$} & \textbf{$\rho$} & \textbf{$p$} \\
\midrule
\endhead
A & $\sqrt5\,t$ & 5 & 0 & $-36/5$ & 5 & $-36$ & 0 \\
B & $t$ & $-15$ & $25/6$ & $12/5$ & $-5$ & $-24$ & 12 \\
\bottomrule
\end{longtable}
\endgroup

The equation of state is $w=p/\rho=0$ in Case A (dust of negative energy) and $w=-\tfrac12$ in Case B. Case B is the notebook's $a_4=t$. Check it with the conditions: $mS=-15\cdot\tfrac{12}{5}=-36$; $\lambda S^2=\tfrac{25}{6}\cdot\tfrac{144}{25}=24=2(15-3)$; $\rho=-36+12=-24$. So within Einstein gravity the dirac16complex field \textbf{can} be the source of the primordial field when $a_4$ is linear, but only as a state of negative energy density, and nothing here says that such a state is realized, stable or preferred (Stage-2 document §1, non-claim 2). For the static field ($c=0$) the conditions become $mS=-36H^2/\kappa$ and $\lambda S^2=30H^2/\kappa$ (so $\lambda>0$ and $mS<0$); whether a Kohn--Sham state can meet them is a question of Chapter 13 (\texttt{handoff/\allowbreak{}specs/ST\allowbreak{}AGE4\_SPE\allowbreak{}C.md}, erratum E4.1). The older statement of the project contract that a condensate "cannot" source this field is correct only for the real-$K$ plane waves of Proposition 9.3 and must not be repeated in general (erratum E4.1).

\subsection{The spin connection and the Dirac operator in this field}

\textbf{The spin connection of a diagonal vielbein.} Chapter 4 fixes the spin connection by the vielbein postulate:

\[
\omega_\mu{}^a{}_b=e_b{}^\nu\bigl(\Gamma^\rho{}_{\mu\nu}\,e_\rho{}^a-\partial_\mu e_\nu{}^a\bigr),\qquad \omega_{\mu ab}=\eta_{ac}\,\omega_\mu{}^c{}_b=-\omega_{\mu ba},
\]

and the spinor connection is $\Omega_\mu=\tfrac12\omega_{\mu ab}S^{ab}$ with $S^{ab}=\tfrac14[\gamma^a,\gamma^b]$, summed over all ordered pairs $(a,b)$. For the diagonal vielbein $e_\nu{}^a=h_\nu\delta_\nu^a$, with inverse $e_b{}^\nu=\delta_b^\nu/h_b$, the sums collapse (no sums below):

\[
\omega_\mu{}^a{}_b=\frac{1}{h_b}\bigl(h_a\,\Gamma^a{}_{\mu b}-\delta^a_b\,\partial_\mu h_b\bigr).
\]

\emph{Case $a=b$.} By rule (C1), $\Gamma^a{}_{\mu a}=\tfrac12\partial_\mu\ln\lvert g_{aa}\rvert=\partial_\mu\ln h_a$, so $\omega_\mu{}^a{}_a=\partial_\mu\ln h_a-\partial_\mu\ln h_a=0$.

\emph{Case $a\ne b$.} Then $\omega_\mu{}^a{}_b=(h_a/h_b)\Gamma^a{}_{\mu b}$, which by rule (C3) vanishes unless $\mu=a$ or $\mu=b$. For $\mu=b$ rule (C2) gives $\Gamma^a{}_{bb}=-\partial_ag_{bb}/(2g_{aa})=-\eta_{bb}h_b\,\partial_ah_b/(\eta_{aa}h_a^2)$, and lowering with $\eta_{aa}$:

\[
\omega_{b\,ab}=-\eta_{bb}\,\frac{\partial_ah_b}{h_a}\qquad(a\ne b,\ \text{no sums}).
\]

For $\mu=a$ the result is the antisymmetric partner of the same number, $\omega_{a\,ab}=\eta_{aa}\partial_bh_a/h_b=-\omega_{a\,ba}$. So a component is nonzero exactly when some scale factor $h_b$ depends on some other coordinate $x_a$.

\textbf{The 24 components.} In the primordial field $h_i=s^{1/6}e^{a_4}$ and $h_j=s^{1/6}e^{-a_4}$ depend on $x_0$ and $x_4$, and $h_0=\cot z$ depends on $x_0$ only (which gives nothing, since $a\ne b$ is required). With $\partial_0h_i=H\cot z\,h_i$, $\partial_4h_i=Ha_4'h_i$, $\partial_0h_j=H\cot z\,h_j$, $\partial_4h_j=-Ha_4'h_j$, $h_0=\cot z$ and $h_4=1$:

\[
\begin{aligned}
\omega_{i\,0i}&=-\omega_{i\,i0}=-H\,s^{1/6}e^{a_4}, &\qquad \omega_{i\,i4}&=-\omega_{i\,4i}=H\,s^{1/6}e^{a_4}\,a_4',\\
\omega_{j\,0j}&=-\omega_{j\,j0}=H\,s^{1/6}e^{-a_4}, &\qquad \omega_{j\,j4}&=-\omega_{j\,4j}=H\,s^{1/6}e^{-a_4}\,a_4'.
\end{aligned}
\]

For example $\omega_{1\,01}=-\eta_{11}\,\partial_0h_1/h_0=-H\cot z\,h_1/\cot z=-H\,s^{1/6}e^{a_4}$. That is $6\times2\times2=24$ nonzero components; $\omega_{0ab}=\omega_{4ab}=0$ (checks \texttt{P\_spinco\allowbreak{}nn\_close\allowbreak{}dForms}, \texttt{P\_spinco\allowbreak{}nn\_count\allowbreak{}24}, \texttt{P\_spinco\allowbreak{}nn\_antis\allowbreak{}ymmetry}; the vielbein postulate holds in all 512 components, check \texttt{P\_spinco\allowbreak{}nn\_vielb\allowbreak{}einPostu\allowbreak{}late512}; Stage-2 document §6). The mixed components $\omega_\mu{}^a{}_b$ are exactly the notebook's ωμIJ of cell 501 (check \texttt{P\_spinco\allowbreak{}nn\_noteb\allowbreak{}ookCell5\allowbreak{}01OmegaM\allowbreak{}uIJEqual\allowbreak{}sMixedOm\allowbreak{}ega}).

\textbf{The spinor connection.} In the double sum $\tfrac12\omega_{\mu ab}S^{ab}$ every antisymmetric pair appears twice, once as $(a,b)$ and once as $(b,a)$, and both terms are equal because $\omega_{\mu ab}$ and $S^{ab}$ both change sign. Hence the spinor connection is a single sum over the pairs with $a<b$,

\[
\Omega_\mu=\sum_{a<b}\omega_{\mu ab}S^{ab},
\]

and inserting the 24 components (with $S^{i4}=-S^{4i}$ and $S^{j4}=-S^{4j}$) gives

\[
\Omega_0=\Omega_4=0,\qquad \Omega_i=-H\,s^{1/6}e^{a_4}\bigl(S^{0i}+a_4'S^{4i}\bigr),\qquad \Omega_j=H\,s^{1/6}e^{-a_4}\bigl(S^{0j}-a_4'S^{4j}\bigr)
\]

(check \texttt{P\_Omega\_\allowbreak{}closedFo\allowbreak{}rms}).

\textbf{The contraction $\gamma^\mu\Omega_\mu=3H\gamma^0$.} Only the Clifford relations $\{\gamma^a,\gamma^b\}=2\eta^{ab}$ are needed. For $a\ne b$ the matrices anticommute, so $S^{ab}=\tfrac12\gamma^a\gamma^b$, and

\[
\gamma^bS^{ab}=\tfrac12\gamma^b\gamma^a\gamma^b=-\tfrac12\gamma^a(\gamma^b)^2=-\tfrac12\eta^{bb}\gamma^a .
\]

Now use the curved gammas of Section 9.2. For $i=1,2,3$ ($\eta^{ii}=+1$):

\[
\gamma^{x_i}\Omega_i=s^{-1/6}e^{-a_4}\gamma^i\cdot\bigl(-Hs^{1/6}e^{a_4}\bigr)\bigl(S^{0i}+a_4'S^{4i}\bigr)=-H\Bigl(-\tfrac12\gamma^0-\tfrac12a_4'\gamma^4\Bigr)=\tfrac H2\bigl(\gamma^0+a_4'\gamma^4\bigr).
\]

For $j=5,6,7$ ($\eta^{jj}=-1$):

\[
\gamma^{x_j}\Omega_j=s^{-1/6}e^{a_4}\gamma^j\cdot Hs^{1/6}e^{-a_4}\bigl(S^{0j}-a_4'S^{4j}\bigr)=H\Bigl(\tfrac12\gamma^0-\tfrac12a_4'\gamma^4\Bigr)=\tfrac H2\bigl(\gamma^0-a_4'\gamma^4\bigr).
\]

Summing three of each,

\[
\gamma^\mu\Omega_\mu=\tfrac{3H}2\bigl(\gamma^0+a_4'\gamma^4\bigr)+\tfrac{3H}2\bigl(\gamma^0-a_4'\gamma^4\bigr)=3H\gamma^0 .
\]

\textbf{The free function $a_4$ cancels}: the inflating and the deflating directions contribute opposite $a_4'$ terms (checks \texttt{P\_Omega\_\allowbreak{}gammaSla\allowbreak{}sh3Hgamm\allowbreak{}a0} and \texttt{P\_Omega\_\allowbreak{}slashA4I\allowbreak{}ndepende\allowbreak{}nt}). In the other order, $S^{ab}\gamma^b=\tfrac12\gamma^a(\gamma^b)^2=+\tfrac12\eta^{bb}\gamma^a$, every term flips sign, so $\Omega_\mu\gamma^\mu=-3H\gamma^0$, $\{\gamma^\mu,\Omega_\mu\}=0$ and $[\gamma^\mu,\Omega_\mu]=6H\gamma^0$ (check \texttt{P\_Omega\_\allowbreak{}OmegaGam\allowbreak{}maAndAnt\allowbreak{}icommuta\allowbreak{}tor}). The divergence identity $\partial_\mu(\sqrt{\lvert g\rvert}\,\gamma^{x_\mu})=\sqrt{\lvert g\rvert}\,[\gamma^\mu,\Omega_\mu]$, proved in Section 4.12 (and again in Section 5.12) and used in Chapter 7 to derive the field equations, can be checked directly: only $\mu=0$ contributes on the left, and

\[
\partial_\mu\bigl(\sqrt{\lvert g\rvert}\,\gamma^{x_\mu}\bigr)=\partial_0\bigl(\cos z\tan z\,\gamma^0\bigr)=\partial_0(\sin z)\,\gamma^0=6H\cos z\,\gamma^0=\sqrt{\lvert g\rvert}\,[\gamma^\mu,\Omega_\mu]
\]

(check \texttt{P\_gammaC\allowbreak{}onst\_div\allowbreak{}ergenceI\allowbreak{}dentity}). With this connection the gammas are covariantly constant, $D_\mu\gamma^\nu=0$ for all 64 pairs (check \texttt{P\_gammaC\allowbreak{}onst\_Dmu\allowbreak{}GammaNuZ\allowbreak{}ero64}).

\textbf{The Dirac equation in the primordial field.} The field equation of Chapter 7, $\gamma^\mu D_\mu\Psi=(m+\lambda S)\Psi$, becomes

\[
\tan z\,\gamma^0\partial_0\Psi+s^{-1/6}e^{-a_4}\sum_i\gamma^i\partial_i\Psi+\gamma^4\partial_4\Psi+s^{-1/6}e^{a_4}\sum_j\gamma^j\partial_j\Psi+3H\gamma^0\Psi=(m+\lambda S)\Psi .
\]

For a field that depends on $x_0$ and $x_4$ only, the $a_4$-dependent coefficients multiply vanishing derivatives, and

\[
\tan z\,\gamma^0\partial_0\Psi+\gamma^4\partial_4\Psi+3H\gamma^0\Psi=(m+\lambda S)\Psi .
\]

\textbf{No $a_4$ appears at all.} Multiplying by $-\gamma^4$ (note $(\gamma^4)^2=\eta^{44}=-1$) solves for the time derivative. Every row of every $\gamma^a$ has exactly one nonzero entry $\pm1$ (Chapter 2), so each equation links a component to exactly one other component through $\gamma^4\gamma^0$ and to one through $\gamma^4$. Following these links, the 16 equations split into four independent blocks of four components, $\{0,5,8,13\}$, $\{1,4,9,12\}$, $\{2,7,10,15\}$, $\{3,6,11,14\}$ (the notebook's coupling sets of cell 1089; checks \texttt{P\_blocks\allowbreak{}\_fourBlo\allowbreak{}cksOfFou\allowbreak{}r} and \texttt{P\_blocks\allowbreak{}\_matchNo\allowbreak{}tebookSe\allowbreak{}ts}). For $\lambda\ne0$ the blocks are coupled only through the single number $S$, a bilinear in all 16 components; the linear part is block diagonal. In the variables $z,t$ ($\partial_0=6H\partial_z$, $\partial_4=H\partial_t$) the first block is (Stage-2 document §10.5)

\[
\begin{aligned}
H\,\partial_t\Psi_0&=-\bigl(6H\tan z\,\partial_z+3H\bigr)\Psi_5+(m+\lambda S)\Psi_{13},\\
H\,\partial_t\Psi_5&=-\bigl(6H\tan z\,\partial_z+3H\bigr)\Psi_0+(m+\lambda S)\Psi_8,\\
H\,\partial_t\Psi_8&=\bigl(6H\tan z\,\partial_z+3H\bigr)\Psi_{13}-(m+\lambda S)\Psi_5,\\
H\,\partial_t\Psi_{13}&=\bigl(6H\tan z\,\partial_z+3H\bigr)\Psi_8-(m+\lambda S)\Psi_0 ,
\end{aligned}
\]

and the other three blocks have the same shape with other signs.

\textbf{The notebook's contraction in this field.} The notebook builds its connection as $\Omega^{\mathrm{NB}}_\mu=\tfrac12\omega_\mu{}^a{}_b\,S^{ab}$, with the mixed $\omega$ contracted against $S^{ab}$ with both indices up; one factor $\eta$ is missing (Chapter 4). Since $\omega_\mu{}^a{}_b=\eta^{aa}\omega_{\mu ab}$ (no sum), the notebook's connection is $\tfrac12\sum_{a,b}\eta^{aa}\omega_{\mu ab}S^{ab}$, and its effect on a pair $(a,b)$ depends on the two signs. The two ordered terms of a pair are $\tfrac12\eta^{aa}\omega_{\mu ab}S^{ab}$ and $\tfrac12\eta^{bb}\omega_{\mu ba}S^{ba}$, and since $\omega_{\mu ba}S^{ba}=(-\omega_{\mu ab})(-S^{ab})=\omega_{\mu ab}S^{ab}$ their sum is $\tfrac12(\eta^{aa}+\eta^{bb})\,\omega_{\mu ab}S^{ab}$; the correct connection has $\tfrac12(1+1)\,\omega_{\mu ab}S^{ab}$ instead. So a pair of two space-like frame indices is unchanged, a pair of two time-like indices changes sign, and a mixed pair (one space-like, one time-like: a \textbf{boost}) cancels, because then $\eta^{bb}=-\eta^{aa}$. Applied to our field, $\Omega^{\mathrm{NB}}_i=-Hs^{1/6}e^{a_4}S^{0i}$ (the boost $S^{4i}$ is lost) and $\Omega^{\mathrm{NB}}_j=+Hs^{1/6}e^{-a_4}a_4'S^{4j}$ (the boost $S^{0j}$ is lost and the time-time term flips sign). Repeating the contraction above gives

\[
\gamma^\mu\Omega^{\mathrm{NB}}_\mu=\tfrac{3H}{2}\bigl(\gamma^0+a_4'\gamma^4\bigr)\ne3H\gamma^0 ,
\]

and $D_\mu\gamma^\nu=0$ fails in 15 of the 64 pairs, with 288 nonzero matrix entries in total. For example $D_1\gamma^4=H\,s^{1/6}e^{a_4}a_4'\,\gamma^1$ instead of zero. The failure is the check \texttt{P\_gammaC\allowbreak{}onst\_not\allowbreak{}ebookCon\allowbreak{}traction\allowbreak{}Fails}, and its closed forms are the check \texttt{P\_gammaC\allowbreak{}onst\_not\allowbreak{}ebookCon\allowbreak{}traction\allowbreak{}ClosedFo\allowbreak{}rms} (Stage-2 document §8.3). For commuting fields, which the notebook actually uses, this error drops out of the $(x_0,x_4)$ equations (Stage-2 document §11.3, fact 1); the error that does survive is the next one.

\subsection{The cell-1058 rule and the q term}

\textbf{The notebook's chain.} The notebook's production Lagrangian La[] (cell 1066) is built from commuting fields f16[k]$(x_0,x_4)$, the curved gamma matrices useT16 built in cell 1058, the mixed connection of cell 501 multiplied by a switch $Q_1$ (the notebook's book-keeping factor for the spin connection), and the mass term $HM\,\Psi_{16}^T\sigma_{16}\Psi_{16}$. Its Euler--Lagrange expressions are stored in cell 1079; cell 1096 multiplies them by $1/(2H)$ and changes to the variables $z,t$; cell 1111 relabels the components block by block into yZ$_0$ to yZ$_{15}$ (block $\{0,5,8,13\}$ becomes yZ$_0$ to yZ$_3$, and so on); and cell 1137 stores the equations solved for $\partial_t$ (Stage-2 document §11.1). The dictionary to our notation is $\Psi_k=\mathrm{f16}[k]$ and $m=-HM$.

\textbf{What the stored equations say.} For the first block the stored cell-1137 equations are (Stage-2 document §11.2)

\[
\begin{aligned}
\partial_t yZ_0&=q\,yZ_0-3\,yZ_1-M\,yZ_3-6\tan z\,\partial_z yZ_1,\\
\partial_t yZ_1&=-3\,yZ_0-q\,yZ_1-M\,yZ_2-6\tan z\,\partial_z yZ_0,\\
\partial_t yZ_2&=M\,yZ_1-q\,yZ_2+3\,yZ_3+6\tan z\,\partial_z yZ_3,\\
\partial_t yZ_3&=M\,yZ_0+3\,yZ_2+q\,yZ_3+6\tan z\,\partial_z yZ_2,
\end{aligned}
\qquad q=Q_1\sinh(a_4)\,a_4'\,e^{-a_4}.
\]

Compare with the correct first block of Section 9.13 for $\lambda=0$: divide it by $H$ and put $m/H=-M$, $yZ_0=\Psi_0$, $yZ_1=\Psi_5$, $yZ_2=\Psi_8$, $yZ_3=\Psi_{13}$. The first line becomes $\partial_tyZ_0=-3\,yZ_1-M\,yZ_3-6\tan z\,\partial_zyZ_1$. \textbf{The stored equations agree with the correct ones term by term, except for the terms $\pm q\,yZ_k$.} The same holds in all 16 rows; the $q$ terms occur only in the first two blocks, yZ$_0$ to yZ$_7$ (checks \texttt{P\_notebo\allowbreak{}okCompar\allowbreak{}e\_cell11\allowbreak{}37Reprod\allowbreak{}uced16of\allowbreak{}16}, \texttt{P\_notebo\allowbreak{}okCompar\allowbreak{}e\_correc\allowbreak{}tVsStore\allowbreak{}dDifferO\allowbreak{}nlyByQ}, \texttt{P\_notebo\allowbreak{}okCompar\allowbreak{}e\_qOnlyI\allowbreak{}nYZ0to7}). Since the correct $(x_0,x_4)$ equations contain no $a_4$ at all (Section 9.13), the term $q$, and with it every $a_4$ dependence of the notebook's block equations, is spurious.

\textbf{Erratum (b): the substitution rule of cell 1058.} Cell 1058 contains the rule

\begin{Verbatim}[fontsize=\small]
1/Sqrt[Sin[6*H*x0]^(1/3)/E^(2*a4[H*x4])] -> 1/(E^a4[H*x4]*Sin[6*H*x0]^(1/6))
\end{Verbatim}

The left side is $1/\sqrt{s^{1/3}e^{-2a_4}}$. The square root is $s^{1/6}e^{-a_4}$ and its reciprocal is $s^{-1/6}e^{+a_4}$, so the correct rule is

\begin{Verbatim}[fontsize=\small]
1/Sqrt[Sin[6*H*x0]^(1/3)/E^(2*a4[H*x4])] -> E^a4[H*x4]/Sin[6*H*x0]^(1/6)
\end{Verbatim}

The rule as written puts $e^{-a_4}$ where $e^{+a_4}$ belongs. This quantity is precisely the coefficient $1/h_j=s^{-1/6}e^{a_4}$ of the extra-time curved gammas $\gamma^{x_5},\gamma^{x_6},\gamma^{x_7}$ of Section 9.2. \textbf{A numerical test} makes the error concrete: take $s=1/64$ and $a_4=\ln2$. Then $s^{1/3}=1/4$, $e^{2a_4}=4$, the left side is $1/\sqrt{1/16}=4$, the correct right side is $e^{a_4}/s^{1/6}=2/(1/2)=4$, and the notebook's right side is $1/(2\cdot\tfrac12)=1$.

\textbf{Erratum (c): how the q term arises (a reconstruction).} The notebook does not store the value of useT16, so the path from the rule to the $q$ term cannot be read off; it has to be reconstructed. Three facts are established by exact computation (Stage-2 document §11.3):

\begin{enumerate}
\item With Clifford-consistent curved gammas, every $Q_1$ term drops out of the commuting-field equations, for the notebook's contraction and for the correct one alike (check \texttt{P\_notebo\allowbreak{}okCompar\allowbreak{}e\_commut\allowbreak{}ingQ1Dro\allowbreak{}psOutCli\allowbreak{}ffordGam\allowbreak{}mas}).
\item The stored cell-1079 expressions, parsed from the committed notebook (check \texttt{P\_notebo\allowbreak{}okCompar\allowbreak{}e\_stored\allowbreak{}Ela16}), contain $Q_1$ terms in exactly the 8 rows of the first two blocks (Stage-2 document §11.3, fact 2).
\item These terms are reproduced exactly, 16 of 16 rows with zero residual, if the extra-time curved gammas carry the wrong factor $s^{-1/6}e^{-a_4}$ on their $+1$ entries and the correct factor $s^{-1/6}e^{+a_4}$ on their $-1$ entries. The same reconstruction, carried through the notebook's own chain, reproduces the stored outputs of cells 1096, 1111 and 1137 exactly (checks \texttt{P\_notebo\allowbreak{}okCompar\allowbreak{}e\_recons\allowbreak{}truction\allowbreak{}Reproduc\allowbreak{}esStored\allowbreak{}Ela}, \texttt{P\_notebo\allowbreak{}okCompar\allowbreak{}e\_eLaztC\allowbreak{}ell1096}, \texttt{P\_notebo\allowbreak{}okCompar\allowbreak{}e\_relabe\allowbreak{}lCell111\allowbreak{}1}).
\end{enumerate}

In formulas the reconstructed gamma is

\[
\gamma'^{\,x_j}=s^{-1/6}\bigl(\cosh(a_4)\,\gamma^j-\sinh(a_4)\,\lvert\gamma^j\rvert\bigr)\qquad(j=5,6,7),
\]

where $\lvert\gamma^j\rvert$ is the matrix of absolute values of the entries. Check it entry by entry: where $\gamma^j$ has $+1$, the entry is $s^{-1/6}(\cosh a_4-\sinh a_4)=s^{-1/6}e^{-a_4}$ (wrong factor); where $\gamma^j$ has $-1$, it is $s^{-1/6}(-\cosh a_4-\sinh a_4)=-s^{-1/6}e^{a_4}$ (correct factor times $-1$) (check \texttt{P\_notebo\allowbreak{}okCompar\allowbreak{}e\_recons\allowbreak{}tructedG\allowbreak{}ammaForm}).

\textbf{The reconstructed gamma violates the Clifford relation.} As a 16 by 16 matrix, $\gamma'^{\,x_j}$ is, like every such matrix, an element of $\mathrm{Cl}(4,4)$ (Section 2.6: the 256 monomials of the gammas are a basis of all 16 by 16 matrices). What fails is the curved Clifford relation $\{\gamma^{x_\mu},\gamma^{x_\nu}\}=2g^{\mu\nu}$ that the curved gammas of a metric must obey (Section 4.10). Write $\gamma^5$ as a signed permutation: row $r$ has its single entry $\epsilon_r=\pm1$ in column $\pi(r)$. Since $(\gamma^5)^2=\eta^{55}=-1$ is diagonal, $\pi(\pi(r))=r$, and since $\gamma^5$ is antisymmetric (Chapter 2), the entry in row $\pi(r)$, column $r$ is $-\epsilon_r$. The matrix $D=s^{1/6}\gamma'^{\,x_5}$ has, in row $r$, the entry $e^{-a_4}$ if $\epsilon_r=+1$ and $-e^{a_4}$ if $\epsilon_r=-1$. The diagonal entry $(D^2)_{rr}$ is the product of the entries at $(r,\pi(r))$ and $(\pi(r),r)$: for $\epsilon_r=+1$ it is $e^{-a_4}\cdot(-e^{a_4})=-1$, for $\epsilon_r=-1$ it is $(-e^{a_4})\cdot e^{-a_4}=-1$; the off-diagonal entries vanish because $\pi(\pi(r))=r$. Hence

\[
\bigl(\gamma'^{\,x_5}\bigr)^2=-s^{-1/3}\cdot1,\qquad\text{whereas the Clifford relation requires}\qquad \bigl(\gamma^{x_5}\bigr)^2=g^{55}=-s^{-1/3}e^{2a_4}.
\]

The two agree only when $a_4=0$ (check \texttt{P\_notebo\allowbreak{}okCompar\allowbreak{}e\_recons\allowbreak{}tructedG\allowbreak{}amma5Not\allowbreak{}Clifford}; the report also records that $\gamma'^{\,x_5}$ anticommutes neither with $\gamma'^{\,x_6}$ nor with $\gamma^{x_0}$).

\textbf{The size of the spurious term, step by step.} Because $\gamma'^{\,x_j}$ violates the Clifford relation $\{\gamma^{x_\mu},\gamma^{x_\nu}\}=2g^{\mu\nu}$, a product such as $\sigma_{16}\gamma'^{\,x_j}S^{4j}$ acquires a symmetric part, and commuting fields do not annihilate a symmetric part. We now derive the $q$ term from this in five steps (Stage-2 document §11.3, "Mechanism").

\emph{Step 1: where the extra-time gammas enter.} The fields depend on $x_0$ and $x_4$ only. In the kinetic part of La the gammas useT16 of $x_5,x_6,x_7$ therefore multiply the vanishing derivatives $\partial_5\Psi_{16}$, $\partial_6\Psi_{16}$, $\partial_7\Psi_{16}$, and they survive only in the $Q_1$ term of La,

\[
\cos z\;Q_1\,\Psi_{16}^T\sigma_{16}X\,\Psi_{16},\qquad X=\sum_\mu\mathrm{useT16}_\mu\,\Omega^{\mathrm{NB}}_\mu ,
\]

where $\Omega^{\mathrm{NB}}_\mu=\tfrac12\omega_\mu{}^a{}_b\,S^{ab}$ is the notebook's connection of Section 9.13 and $\cos z=\sqrt{\lvert g\rvert}$ is the volume factor (cell 1061) that multiplies all of La.

\emph{Step 2: the Euler--Lagrange derivative of a quadratic form.} For commuting components the derivative of $\Psi^TA\Psi=\sum_{r,c}\Psi_rA_{rc}\Psi_c$ with respect to $\Psi_k$ is $\sum_cA_{kc}\Psi_c+\sum_r\Psi_rA_{rk}$, which is entry $k$ of $(A+A^T)\Psi$. The $Q_1$ term contains no derivative of $\Psi$, so this is all it gives. The notebook's Euler--Lagrange expressions (cell 1075) are divided by the volume factor $\cos z$ that multiplies La. So in the normalisation of cell 1079 the $Q_1$ term contributes

\[
Q_1\bigl(\sigma_{16}X+(\sigma_{16}X)^T\bigr)\Psi_{16} .
\]

Only the \textbf{symmetric part} $A+A^T$ of $A=\sigma_{16}X$ counts; an antisymmetric $A$ (one with $A^T=-A$) contributes nothing.

\emph{Step 3: split off the Clifford-consistent part.} Write the reconstructed gammas as the correct ones plus a defect, $X=\sum_\mu\gamma^{x_\mu}\Omega^{\mathrm{NB}}_\mu+Y$, with

\[
Y=\sum_{j=5}^{7}\bigl(\gamma'^{\,x_j}-\gamma^{x_j}\bigr)\,\Omega^{\mathrm{NB}}_j .
\]

By fact 1, $\sigma_{16}\sum_\mu\gamma^{x_\mu}\Omega^{\mathrm{NB}}_\mu$ has no symmetric part. So the whole $Q_1$ content of the reconstructed (and hence of the stored) cell-1079 expressions is $Q_1\bigl(\sigma_{16}Y+(\sigma_{16}Y)^T\bigr)\Psi_{16}$. It is the difference between the stored expressions and the Clifford-consistent rebuild, which the Stage-2 document calls the \textbf{$q$ vector} (check \texttt{P\_notebo\allowbreak{}okCompar\allowbreak{}e\_qTermF\allowbreak{}romNonCl\allowbreak{}iffordEx\allowbreak{}traTimeG\allowbreak{}ammas}).

\emph{Step 4: the two factors of $Y$.} By the pair rule of Section 9.13, the boost pair $(0,j)$ cancels in $\Omega^{\mathrm{NB}}_j$ and only the time-time pair $(4,j)$ remains: $\Omega^{\mathrm{NB}}_j=\omega_j{}^4{}_j\,S^{4j}$, with the mixed component $\omega_j{}^4{}_j=\eta^{44}\omega_{j\,4j}=\omega_{j\,j4}=H\,a_4'\,s^{1/6}e^{-a_4}$. This is the $\Omega^{\mathrm{NB}}_j=+Hs^{1/6}e^{-a_4}a_4'S^{4j}$ found there. The defect $\gamma'^{\,x_j}-\gamma^{x_j}$ is zero on the $-1$ entries of $\gamma^j$, where both matrices have $-s^{-1/6}e^{a_4}$. On the $+1$ entries it is $s^{-1/6}e^{-a_4}-s^{-1/6}e^{a_4}=-s^{-1/6}(e^{a_4}-e^{-a_4})$. Let $P_j$ be the matrix that keeps the $+1$ entries of $\gamma^j$ and has 0 everywhere else. Then $\gamma'^{\,x_j}-\gamma^{x_j}=-s^{-1/6}(e^{a_4}-e^{-a_4})P_j$ and, with $e^{a_4}-e^{-a_4}=2\sinh a_4$,

\[
Y=-s^{-1/6}\bigl(e^{a_4}-e^{-a_4}\bigr)\cdot H\,a_4'\,s^{1/6}e^{-a_4}\sum_{j=5}^{7}P_jS^{4j}=-2H\sinh(a_4)\,a_4'\,e^{-a_4}\sum_{j=5}^{7}P_jS^{4j} .
\]

The powers of $s$ cancel, so nothing here depends on $z$. Multiplying by $Q_1$ and using $q=Q_1\sinh(a_4)\,a_4'\,e^{-a_4}$:

\[
Q_1\bigl(\sigma_{16}Y+(\sigma_{16}Y)^T\bigr)=2Hq\,N,\qquad N=-\sum_{j=5}^{7}\Bigl(\sigma_{16}P_jS^{4j}+\bigl(\sigma_{16}P_jS^{4j}\bigr)^T\Bigr).
\]

The size of the spurious term is thus the difference of the two factors, $s^{-1/6}(e^{a_4}-e^{-a_4})$, times the mixed connection component $\omega_j{}^4{}_j=H\,a_4'\,s^{1/6}e^{-a_4}$, times $Q_1$, which is $2Hq$.

\emph{Step 5: the numbers.} $N$ is a 16 by 16 matrix of pure numbers, built from $\sigma_{16}$, the $P_j$ and $S^{4j}=\tfrac12\gamma^4\gamma^j$. Multiplied out, it has exactly 8 nonzero entries, each $+1$ or $-1$, one in each of the rows 0, 1, 4, 5, 8, 9, 12 and 13, which are the components of the first two blocks $\{0,5,8,13\}$ and $\{1,4,9,12\}$, as fact 2 requires. As (row, column, entry) they are (0, 9, $-1$), (1, 8, $+1$), (4, 13, $-1$), (5, 12, $+1$), (8, 1, $+1$), (9, 0, $-1$), (12, 5, $+1$) and (13, 4, $-1$) (report measurement \texttt{notebook\allowbreak{}Ela\_qVec\allowbreak{}torRows\_\allowbreak{}sign\_par\allowbreak{}tner}). So each nonzero row of the $q$ vector is $\pm2Hq\,\Psi_k$ for one component $k$, for example $-2Hq\,\Psi_9$ in row 0 (check \texttt{P\_notebo\allowbreak{}okCompar\allowbreak{}e\_qVecto\allowbreak{}rIs2HqAt\allowbreak{}Cell1079}). The factor $1/(2H)$ of cell 1096 turns $2Hq$ into $q$, and solving for the time derivatives turns the eight rows into the terms $\pm q\,yZ_k$ of the stored equations of the first two blocks (checks \texttt{P\_notebo\allowbreak{}okCompar\allowbreak{}e\_eLaztC\allowbreak{}ell1096} and \texttt{P\_notebo\allowbreak{}okCompar\allowbreak{}e\_cell11\allowbreak{}37Reprod\allowbreak{}uced16of\allowbreak{}16}). The identity $\sinh(a_4)e^{-a_4}=\tfrac12(1-e^{-2a_4})$ gives the equivalent form $q=\tfrac12Q_1a_4'(1-e^{-2a_4})$; for $a_4=t$, $q=\tfrac12Q_1(1-e^{-2t})$.

\textbf{What is established and what is reconstructed.} Established by checks of the WolframScript verifier alone (Section 9.1; the Python checker only measures the same $\pm q$ differences in a transcription of cell 1137): the stored outputs of cells 1079, 1096, 1111 and 1137 are reproduced exactly, and the correct equations differ from the stored ones only by the $q$ terms. Reconstructed: that the rule of cell 1058 produced the reconstructed useT16 by acting on the $+1$ entries only. When cell 1058 is re-executed literally in the verifier's kernel (Mathematica 15.0.1), the wrong factor $e^{-a_4}$ is applied to all entries of the three extra-time gammas; the resulting equations then contain no $q$ term and differ from the stored ones by exactly the $q$ terms (check \texttt{P\_notebo\allowbreak{}okCompar\allowbreak{}e\_litera\allowbreak{}lRebuild\allowbreak{}Residual\allowbreak{}IsExactl\allowbreak{}yQTerms}). The notebook's metadata indicate that the stored chain ran in one older kernel session (cell 1096 stores the date 2026-01-30), in which the simplified form of the $-1$ entries, and hence whether the rule matched them, may have differed. That explanation is a reconstruction as well (Stage-2 document §11.3 and §20, item 7).

\textbf{Consequences.} Later steps of the notebook that remove or use the $q$ term, a rescaling of the yZ components and an equation for $a_4$ built from the same coefficient (cells 100, 1157 and 1158 in the survey), have no counterpart in the correct equations. This last remark is a reading of the notebook's text, not a machine check (Stage-2 document §11.4). The correct statement is simple: \textbf{for fields of $(x_0,x_4)$ the dynamics does not see $a_4$ at all.}

\subsection{The static member of the family}

\textbf{Why a static field.} Part IV of this book looks for the ground state and the first excited state of many dirac16complex quanta (Chapter 13). A ground state is a state that does not change in time, and that needs a background that does not change in time either. The primordial family has exactly one such member: $a_4$ constant, $a_4=a_{4,0}$, so $a_4'=a_4''=0$ (\texttt{handoff/\allowbreak{}specs/ST\allowbreak{}AGE4\_SPE\allowbreak{}C.md} §1).

\textbf{The metric.} Stage 4 calls the hidden coordinate $y$; it is the $\zeta$ of Section 9.5, $y=\ln(\sin z)/(6H)$, now including the end point $y=0$ ($z=\pi/2$):

\[
ds^2=dy^2-dx_4^2+e^{2Hy}\Bigl[e^{2a_{4,0}}\bigl(dx_1^2+dx_2^2+dx_3^2\bigr)-e^{-2a_{4,0}}\bigl(dx_5^2+dx_6^2+dx_7^2\bigr)\Bigr],\qquad y\le0 .
\]

The warp factor is $W(y)=e^{Hy}$ and the volume factor is $\sqrt{\lvert g\rvert}=W^6=e^{6Hy}$ (checks \texttt{KS\_geome\allowbreak{}try\_sqrt\allowbreak{}DetG\_W6}, \texttt{KS\_geome\allowbreak{}try\_sign\allowbreak{}ature44} and \texttt{KS\_geome\allowbreak{}try\_note\allowbreak{}bookChar\allowbreak{}t}).

\textbf{The constant $a_{4,0}$ can be absorbed.} New coordinates $x_i'=e^{a_{4,0}}x_i$ and $x_j'=e^{-a_{4,0}}x_j$ remove $a_{4,0}$ from the line element. It returns only through quantities measured in the old coordinates; in the Kohn--Sham problem of Chapter 13 it rescales every 3-space momentum, $k\to k\,e^{-a_{4,0}}$ (check \texttt{KS\_reduc\allowbreak{}tion\_a4I\allowbreak{}sMomentu\allowbreak{}mRescali\allowbreak{}ng}).

\textbf{Curvature.} Lemma 9.1 with $f_p=Hy+\sigma_pa_{4,0}$ gives $\partial_yf_p=H$, all other derivatives zero, $\Theta_y=6H$, $\Theta_4=0$. So $R^y{}_y=-6H^2$, $R^4{}_4=0$, $R^p{}_p=-6H^2$ for the six fibre directions, and

\[
R=-42H^2,\qquad G^\mu{}_\nu=\mathrm{diag}(15,15,15,15,21,15,15,15)\,H^2
\]

in the order $(y,x_1,x_2,x_3,x_4,x_5,x_6,x_7)$. This is also the case $a_4'=a_4''=0$ of Section 9.8. The Stage-4 verifier checks the scalar in \texttt{KS\_geome\allowbreak{}try\_ricc\allowbreak{}iScalarM\allowbreak{}inus42H2}, the Einstein tensor in \texttt{KS\_geome\allowbreak{}try\_eins\allowbreak{}teinMixe\allowbreak{}dDiag} and the Ricci tensor in \texttt{KS\_geome\allowbreak{}try\_ricc\allowbreak{}iMixed}. Only 18 Christoffel symbols are nonzero, namely $\Gamma^y{}_{pp}=-H\,g_{pp}$ and $\Gamma^p{}_{yp}=\Gamma^p{}_{py}=H$ (check \texttt{KS\_geome\allowbreak{}try\_chri\allowbreak{}stoffelC\allowbreak{}ount18}). The exact theory of Stage 4 found more. The seven directions $(y,x_1,x_2,x_3,x_5,x_6,x_7)$ form a space of constant curvature $-H^2$, which means $R_{rsmn}=-H^2(g_{rm}g_{sn}-g_{rn}g_{sm})$ for those directions, and the time line $x_4$ is flat (check \texttt{KS\_geome\allowbreak{}try\_cons\allowbreak{}tantCurv\allowbreak{}atureSev\allowbreak{}enSpace}). For such a space in $n=7$ dimensions the Ricci tensor is $-(n-1)H^2g$ and $R=-n(n-1)H^2=-42H^2$, in agreement with the lemma, and the square of the curvature is $R_{\mu\nu\rho\sigma}R^{\mu\nu\rho\sigma}=2n(n-1)H^4=84H^4$ (check \texttt{KS\_geome\allowbreak{}try\_kret\allowbreak{}schmannC\allowbreak{}onstant}). Every curvature invariant is a constant: \textbf{nothing singular happens anywhere in $y$}, in particular not at $y=0$.

\textbf{The required source.} From $G^\mu{}_\nu=\kappa T^\mu{}_\nu$:

\[
\rho_{\mathrm{req}}=-\frac{21H^2}{\kappa}<0,\qquad p_{\mathrm{req}}=+\frac{15H^2}{\kappa}\quad\text{in all seven transverse directions},\qquad w_{\mathrm{req}}=-\frac57
\]

(checks \texttt{KS\_geome\allowbreak{}try\_requ\allowbreak{}iredSour\allowbreak{}ce} and \texttt{KS\_geome\allowbreak{}try\_rhoR\allowbreak{}equiredN\allowbreak{}egative}). \textbf{Stage-4 erratum E4.3 (the sign of $p_{\mathrm{req}}$).} The Stage-4 specification §1 had written $p_{\mathrm{req}}=-15H^2/\kappa$; the exact mixed components give $+15H^2/\kappa$ (erratum E4.3 of \texttt{handoff/\allowbreak{}specs/ST\allowbreak{}AGE4\_SPE\allowbreak{}C.md}; the Wolfram report records it under \texttt{specDisc\allowbreak{}repancy\_\allowbreak{}pReqSign}). By Proposition 9.4 with $c=0$, an $x_0$-independent dirac16complex state is an exact source of the static field if and only if its three-gamma bilinears vanish, $mS=-36H^2/\kappa$ and $\lambda S^2=30H^2/\kappa$.

\textbf{Extrinsic curvature of the slices.} Each surface $y=\mathrm{const}$ is a 7-dimensional slice. Its \textbf{unit normal} is $n=\partial_y$ (a unit vector because $g_{yy}=1$). Its \textbf{extrinsic curvature} measures how the slice bends inside the 8-dimensional space: $K_{ab}=\nabla_an_b$ for directions $a,b$ along the slice. The lower components of the normal are $n_c=\delta_c^y$, so $K_{ab}=\partial_an_b-\Gamma^c{}_{ab}n_c=-\Gamma^y{}_{ab}$, and by rule (C2) with $g_{yy}=1$, $\Gamma^y{}_{ab}=-\tfrac12\partial_yg_{ab}$. Hence

\[
K_{ab}=\tfrac12\partial_yg_{ab},\qquad K^a{}_b=\tfrac12g^{ac}\partial_yg_{cb} .
\]

For the six warped directions $g_{pp}\propto e^{2Hy}$ and $K^p{}_p=H$; for $x_4$, $K^4{}_4=0$; the trace is $K=6H$ (check \texttt{KS\_geome\allowbreak{}try\_extr\allowbreak{}insicCur\allowbreak{}vature}).

\subsection{Beyond y = 0: the smooth extension and the Z2 brane}

The notebook's chart ends at $y=0$, where its coordinate $x_0$ degenerates. The geometry itself does not end there. Stage 4 documents two ways to continue it (\texttt{handoff/\allowbreak{}specs/ST\allowbreak{}AGE4\_SPE\allowbreak{}C.md} §1, where they are labelled (E1) and (E2); these labels name extensions, not errata, and we call them Extension 1 and Extension 2).

\textbf{Extension 1: the smooth extension.} Keep $W=e^{Hy}$ for all real $y$. By Section 9.15 all curvature invariants are constant, so this is a regular homogeneous space with no special point at $y=0$. The tip $y\to-\infty$, where $W^6\to0$ and the transverse space pinches off, lies at infinite proper distance, because $y$ itself measures proper distance.

\textbf{Extension 2: the Z2 mirror.} Take two copies of the notebook's patch $y\le0$ and glue them along $y=0$, the second copy reflected. Equivalently, use the warp factor

\[
W(y)=e^{-H\lvert y\rvert}:\qquad W=e^{Hy}\ \text{for}\ y\le0,\qquad W=e^{-Hy}\ \text{for}\ y\ge0 .
\]

The name Z2 refers to the group with two elements, the identity and the reflection $y\to-y$, under which this geometry is symmetric. This is the construction that Stage 4 adopts as a model of the notebook's picture of a pair of universes; it is a \textbf{choice of model, not a result derived from any equation}. The price is a kink: $W$ is continuous at $y=0$, but its slope jumps from $+H$ to $-H$.

\textbf{A tool: the delta function.} The derivative of $\lvert y\rvert$ is the sign function, $\mathrm{sgn}(y)=-1$ for $y<0$ and $+1$ for $y>0$. The derivative of the sign function is zero everywhere except at $y=0$, where it jumps by 2. Physicists write it as $2\delta(y)$, where the \textbf{delta function} $\delta(y)$ is the limit of ever narrower spikes of total area 1. A concrete way to see this is to smooth the corner: $\lvert y\rvert\approx\sqrt{y^2+\varepsilon^2}$ for small $\varepsilon>0$. Its first derivative $y/\sqrt{y^2+\varepsilon^2}$ tends to $\mathrm{sgn}(y)$, and its second derivative $\varepsilon^2/(y^2+\varepsilon^2)^{3/2}$ is a spike of height $1/\varepsilon$ and total area 2 (Exercise 9.11). So for $f(y)=-H\lvert y\rvert$:

\[
f'(y)=-H\,\mathrm{sgn}(y),\qquad f'(y)^2=H^2\ (y\ne0),\qquad f''(y)=-2H\,\delta(y).
\]

\textbf{The Einstein tensor with a kink.} Use Lemma 9.1 once more, with the base coordinates $(y,x_4)$ and $f_p=f(y)+\sigma_pa_{4,0}$, so that $\partial_yf_p=f'$, $\partial_y^2f_p=f''$, $\Theta_y=6f'$ and all $x_4$ derivatives vanish. The lemma was proved for smooth functions; its result contains $f''$ only linearly and $f'$ only through the bounded function $f'^2$, so it extends to the kinked $f$ when $f''$ is read as the delta function (the usual justification of thin-wall sources):

\[
R^y{}_y=-6\bigl(f''+f'^2\bigr),\qquad R^4{}_4=0,\qquad R^p{}_p=-\bigl(f''+6f'^2\bigr),\qquad R=-12f''-42f'^2 ,
\]

and therefore

\[
G^y{}_y=15f'^2,\qquad G^4{}_4=6f''+21f'^2,\qquad G^p{}_p=5f''+15f'^2 .
\]

For $f=Hy$ this returns the static values $15H^2$, $21H^2$, $15H^2$. For the mirror, $f=-H\lvert y\rvert$:

\[
G^y{}_y=15H^2,\qquad G^4{}_4=21H^2-12H\,\delta(y),\qquad G^p{}_p=15H^2-10H\,\delta(y).
\]

\textbf{The brane.} Einstein's equations $G^\mu{}_\nu=\kappa T^\mu{}_\nu$ now require, besides the bulk source $\rho_{\mathrm{req}}=-21H^2/\kappa$ of Section 9.15 on both sides, a source concentrated on the surface $y=0$: $T^\mu{}_\nu\supset S^\mu{}_\nu\,\delta(y)$ with

\[
S^y{}_y=0,\qquad S^4{}_4=-\frac{12H}{\kappa},\qquad S^p{}_p=-\frac{10H}{\kappa}\quad(p=1,2,3,5,6,7).
\]

A surface that carries energy and momentum of its own is called a \textbf{brane} (from "membrane"). Its surface energy density and surface pressure are

\[
\rho_{\mathrm{brane}}=-S^4{}_4=+\frac{12H}{\kappa}>0,\qquad p_{\mathrm{brane}}=S^p{}_p=-\frac{10H}{\kappa}
\]

(checks \texttt{KS\_geome\allowbreak{}try\_isra\allowbreak{}elStress} and \texttt{KS\_geome\allowbreak{}try\_bran\allowbreak{}eEnergyP\allowbreak{}ositive}). The pressure is negative, like a stretched membrane, but it is not a pure tension, for which $\rho=-p$ would hold. The component $S^y{}_y=0$ says that the concentrated stress lies entirely along the brane, as the stress of a thin wall must. The metric induced on the brane, in the order $(x_1,x_2,x_3,x_4,x_5,x_6,x_7)$,

\[
h=\mathrm{diag}\bigl(e^{2a_{4,0}},e^{2a_{4,0}},e^{2a_{4,0}},-1,-e^{-2a_{4,0}},-e^{-2a_{4,0}},-e^{-2a_{4,0}}\bigr),
\]

is constant, so the brane itself is flat (check \texttt{KS\_geome\allowbreak{}try\_indu\allowbreak{}cedMetri\allowbreak{}cFlat}).

\textbf{The same result from the Israel junction conditions.} The general rule for a thin wall in Einstein gravity is the \textbf{Israel junction condition}. Let $[X]=X(0^+)-X(0^-)$ denote the jump of a quantity across the wall, let the normal be $n=+\partial_y$, pointing from $y<0$ to $y>0$, and let $K_{ab}=\tfrac12\partial_yg_{ab}$ be the extrinsic curvature on each side (Section 9.15). Then

\[
[K_{ab}]-h_{ab}[K]=-\kappa\,S_{ab},
\]

with $h_{ab}$ the induced metric. On the side $y<0$, $K^p{}_p=+H$; on the side $y>0$, where $g_{pp}\propto e^{-2Hy}$, $K^p{}_p=-H$; along $x_4$, $K^4{}_4=0$ on both sides. So $[K^p{}_p]=-2H$, $[K^4{}_4]=0$ and $[K]=-12H$, and with one index raised

\[
S^p{}_p=-\frac1\kappa\bigl(-2H+12H\bigr)=-\frac{10H}{\kappa},\qquad S^4{}_4=-\frac1\kappa\bigl(0+12H\bigr)=-\frac{12H}{\kappa},
\]

the same as before (check \texttt{KS\_geome\allowbreak{}try\_isra\allowbreak{}elJump}). The sign convention is recorded in the report (\texttt{geometry\allowbreak{}\_israelS\allowbreak{}ignConve\allowbreak{}ntion} in \texttt{wolfram-\allowbreak{}kohn-sha\allowbreak{}m-report\allowbreak{}.json}); with it, the brane of the Randall--Sundrum model, a well-known five-dimensional model with warp factor $e^{-k\lvert y\rvert}$, gets a positive tension $6k/\kappa$ (Exercise 9.9; check \texttt{KS\_geome\allowbreak{}try\_isra\allowbreak{}elConven\allowbreak{}tionRand\allowbreak{}allSundr\allowbreak{}um} of the Python theory checker).

\textbf{What the mirror construction does and does not say.} The brane stress is what the geometry requires; nothing in this chapter derives it from matter. The bulk on each side still requires the negative energy density $-21H^2/\kappa$. A structural remark connects the construction to the notebook's hypothesis: the chirality map $\Psi\to\gamma^8\Psi$ sends the Lagrangian with mass $m$ and potential $U$ to minus the Lagrangian with $-m$ and $-U$ (erratum E2 of \texttt{handoff/\allowbreak{}specs/CO\allowbreak{}NTRACT.m\allowbreak{}d}): a field $\gamma^8\Psi$ has the opposite mass and the opposite coupling, $(m,\lambda)\to(-m,-\lambda)$. The reflection across the brane is a different map: it relates a solution with the mass function $M(y)$ to one with $-M(-y)$ and the same coupling $\lambda$ (\texttt{artifact\allowbreak{}s/dirac1\allowbreak{}6complex\allowbreak{}/pair-cr\allowbreak{}eation/p\allowbreak{}airing-t\allowbreak{}heory.js\allowbreak{}on}, entries T1 and T2z2). Both are structural echoes of the notebook's $\pm M$ pair, and they are two different pairings (Chapter 15). In the reduced Kohn--Sham equation of Chapter 13 the reflection $\Psi(-y)=\pm\gamma^0\Psi(y)$ is a symmetry only if the mass function is odd, $M(-y)=-M(y)$, that is, if the mirror universe carries the opposite mass; for an even mass function the symmetry is a different matrix, $i\gamma^0\gamma^8$ (erratum E4.6 of \texttt{handoff/\allowbreak{}specs/ST\allowbreak{}AGE4\_SPE\allowbreak{}C.md}). These are statements about the structure of the equations. They do not describe a creation process at $x_4=0$, and they do not show that universes are created in pairs; Chapter 15 proves the exact pairing theorems and Chapter 16 states precisely what they establish.

\subsection{What we proved and what we assumed}

\textbf{Proved in this chapter} (each step derived above and verified by the named checks of the Stage-2 and Stage-4 verifiers): the notebook's MatrixMetric44 has signature (4,4), $\det g=+\cos^2z$ and a constant 7-volume; in the hidden proper coordinate it is a warped metric; it has exactly 37 nonzero Christoffel symbols; by Lemma 9.1 its Ricci scalar is $6H^2(a_4'^2-7)$ and its Einstein tensor is diagonal with the closed forms of Section 9.8, which agree with the notebook's stored outputs; in 8-dimensional Einstein gravity it requires $\rho_{\mathrm{req}}=-3H^2(7+a_4'^2)/\kappa<0$ for every $a_4$, so it is neither a vacuum nor a cosmological-constant solution, and the weak and dominant energy conditions and the null condition along $e_4+e_0$ fail for every $a_4$; no dirac16complex state of $(x_0,x_4)$ can be its source when $a_4''\ne0$, no real-$K$ plane wave can be its source at all, and for linear $a_4$ an $x_0$-independent state is an exact source, necessarily of negative energy density; the spin connection has 24 nonzero components and $\gamma^\mu\Omega_\mu=3H\gamma^0$, so the equations of fields of $(x_0,x_4)$ contain no $a_4$; the notebook's cell-1058 rule is mathematically wrong, its stored block equations differ from the correct ones exactly by the $q$ terms, and the reconstructed curved gamma violates the Clifford relation, which through the symmetric part of $\sigma_{16}Y$ produces exactly the $q$ terms; the static member has $R=-42H^2$, $G=\mathrm{diag}(15,15,15,15,21,15,15,15)H^2$, $\rho_{\mathrm{req}}=-21H^2/\kappa$ and $p_{\mathrm{req}}=+15H^2/\kappa$; and the Z2 mirror geometry requires a flat brane with $\rho_{\mathrm{brane}}=12H/\kappa$ and $p_{\mathrm{brane}}=-10H/\kappa$. Four errata were recorded: (a) the sign of $\det g$ in the Stage-2 specification, (b) the rule of cell 1058, (c) the $q$ term, and Stage-4 erratum E4.3, the sign of $p_{\mathrm{req}}$ in the Stage-4 specification.

\textbf{Assumed, quoted or open.} The metric is a prescribed background; nothing here derives it or its function $a_4$ from dynamics, and the slopes of cell 150 are tabulated, not derived. Only 8-dimensional Einstein gravity is used: Lovelock's theorem and the Einstein--Hilbert variational principle are quoted, and the Lovelock terms of order 2 and 3, which the notebook names, are computed neither by the notebook nor by the project. The energy--momentum tensor and field equations are taken from Chapter 7, and the source analysis reads the bilinears as ordinary numbers (mean field); the stability and quantum status of the negative-energy source are not examined. For linear $a_4$ the source analysis covers only two families, the real-$K$ plane waves and the $x_0$-independent state; other states of $x_0$ and $x_4$, for example the plane waves with $\mathrm{Im}\,K=-3H$ and $\mathrm{Re}\,K\ne0$, are open. For $\lambda\ne0$ the real-$K$ plane wave is not an exact solution of the Dirac equation, and Proposition 9.3 then holds in the pointwise mean-field reading. The comparisons with the notebook, in particular the reconstruction of the $q$ term, rest on the WolframScript verifier alone; the Python checker compares only with transcriptions of a few stored outputs. The attribution of the $q$ term to the rule of cell 1058 acting on the $+1$ entries only is a reconstruction. The Z2 mirror is a modelling choice of Stage 4; the brane stress it requires is derived here as a requirement of the geometry, not from any matter; and the pairings of $\pm M$ under $\gamma^8$ and under the reflection across the brane are structural and are two different maps. The general Israel junction condition is quoted (its content was derived here for our metric from the Einstein tensor with a kink), and the Randall--Sundrum comparison is quoted from the report.

\subsection{Exercises}

\textbf{Exercise 9.1.} At the point where $\sin z=3/5$, compute $\det g$ for an arbitrary value of $a_4$ and check that $\sqrt{\lvert g\rvert}=\cos z$.

\textbf{Exercise 9.2.} Let $H=1$. At which $\sin z$ is the hidden proper coordinate $\zeta=-1$? Check that the warp factor $e^{2H\zeta}$ equals $s^{1/3}$ there.

\textbf{Exercise 9.3.} Derive $\Gamma^j{}_{4j}=-Ha_4'$ and $\Gamma^4{}_{jj}=Ha_4'\,s^{1/3}e^{-2a_4}$ from the rules (C1) and (C2) of Section 9.6.

\textbf{Exercise 9.4.} Verify the Bianchi identity $\nabla_\mu G^\mu{}_\zeta=0$ for the Einstein tensor of Section 9.8 in the $\zeta$ chart.

\textbf{Exercise 9.5.} Take $a_4=2t$. Compute $R$, the mixed Einstein tensor, $\rho_{\mathrm{req}}$, the seven pressures and $w$, and compare $\rho_{\mathrm{req}}$ and $w$ with the plateau $A=2$ of the figure in Section 9.10.

\textbf{Exercise 9.6.} In the Gauss--Bonnet sum $\delta^{\mu_1\nu_1\mu_2\nu_2}_{\alpha_1\beta_1\alpha_2\beta_2}R^{\alpha_1\beta_1}{}_{\mu_1\nu_1}R^{\alpha_2\beta_2}{}_{\mu_2\nu_2}$, the determinant is a sum over permutations. Show that the permutation that exchanges the second and third lower positions contributes $-R^\mu{}_\nu R^\nu{}_\mu$.

\textbf{Exercise 9.7.} For the slopes $a_4'=\tfrac23(M\mp1)$ of cell 150, compute $R$ and show that $\rho_{\mathrm{req}}<0$ for every real $M$.

\textbf{Exercise 9.8.} Derive the NEC entries $\rho+p_{(i)}$ and $p_{(i)}-p_{(j)}$ of the table in Section 9.11.

\textbf{Exercise 9.9.} The Randall--Sundrum model has five coordinates $(y,x_0,x_1,x_2,x_3)$ with $x_0$ the time and $ds^2=dy^2+e^{-2k\lvert y\rvert}\bigl(-dx_0^2+dx_1^2+dx_2^2+dx_3^2\bigr)$. Use the Israel condition of Section 9.16 to find the brane stress $S^\mu{}_\nu$, its energy density and its pressure.

\textbf{Exercise 9.10.} Verify Case B of Proposition 9.4: with $H=\kappa=1$, $a_4=t$, $m=-15$, $\lambda=25/6$ and $S=12/5$, check conditions (2) and (3), and compute $M_{\mathrm{eff}}$, $\rho$, $p$ and $w$.

\textbf{Exercise 9.11.} Show that $\int_{-\infty}^{\infty}\varepsilon^2(y^2+\varepsilon^2)^{-3/2}\,dy=2$ for every $\varepsilon>0$, and that the integrand at $y=0$ equals $1/\varepsilon$.

\textbf{Exercise 9.12.} Test the rule of cell 1058 at $s=1/729$ and $a_4=\ln3$. Then take the 2 by 2 toy matrix $\gamma=\begin{pmatrix}0&1\\-1&0\end{pmatrix}$ (antisymmetric, $\gamma^2=-1$), form $\gamma'=\cosh(a)\,\gamma-\sinh(a)\,\lvert\gamma\rvert$, and compare $\gamma'^2$ with $(e^{a}\gamma)^2$.

\subsection{Answers to the exercises}

\textbf{Answer 9.1.} The determinant is the product of the diagonal entries, $\det g=\cot^2z\cdot s\,e^{6a_4}\cdot(-1)\cdot(-s\,e^{-6a_4})=\cot^2z\,s^2$; the exponentials cancel for every $a_4$. With $\sin z=3/5$: $\cos z=4/5$, $\cot^2z=16/9$, $s^2=9/25$, so $\det g=16/25=(4/5)^2=\cos^2z$ and $\sqrt{\lvert g\rvert}=4/5=\cos z$.

\textbf{Answer 9.2.} $\zeta=\ln(\sin z)/6=-1$ gives $\ln\sin z=-6$, so $\sin z=e^{-6}\approx0.00248$, a point close to the tip. There $e^{2H\zeta}=e^{-2}\approx0.1353$ and $s^{1/3}=(e^{-6})^{1/3}=e^{-2}$: the same number.

\textbf{Answer 9.3.} Rule (C1) with $\rho=j$ and $\mu=4$: $\Gamma^j{}_{4j}=\tfrac12\partial_4\ln\lvert g_{jj}\rvert=\tfrac12\partial_4\bigl(\tfrac13\ln s-2a_4\bigr)=\tfrac12(-2Ha_4')=-Ha_4'$, because $s$ does not depend on $x_4$ and $\partial_4a_4=Ha_4'$. Rule (C2) with $\rho=4$ and $\mu=j$: $\Gamma^4{}_{jj}=-\partial_4g_{jj}/(2g_{44})=\tfrac12\partial_4g_{jj}$. With $g_{jj}=-s^{1/3}e^{-2a_4}$, $\partial_4g_{jj}=-s^{1/3}e^{-2a_4}\cdot(-2Ha_4')=2Ha_4'\,s^{1/3}e^{-2a_4}$, and half of it is $Ha_4'\,s^{1/3}e^{-2a_4}$.

\textbf{Answer 9.4.} For diagonal $G$, $\nabla_\mu G^\mu{}_\zeta=\partial_\zeta G^\zeta{}_\zeta+\Gamma^\mu{}_{\mu\zeta}G^\zeta{}_\zeta-\sum_\mu\Gamma^\mu{}_{\mu\zeta}G^\mu{}_\mu$. In the $\zeta$ chart $G$ does not depend on $\zeta$, $\Gamma^\zeta{}_{\zeta\zeta}=\Gamma^4{}_{4\zeta}=0$ and $\Gamma^p{}_{p\zeta}=\partial_\zeta f_p=H$, so $\Gamma^\mu{}_{\mu\zeta}=6H$. Then

\[
\begin{aligned}
\nabla_\mu G^\mu{}_\zeta&=6H\bigl(-3H^2(a_4'^2-5)\bigr)-H\bigl[3H^2(15-3a_4'^2+a_4'')+3H^2(15-3a_4'^2-a_4'')\bigr]\\
&=(90-18a_4'^2)H^3-(90-18a_4'^2)H^3=0 .
\end{aligned}
\]

\textbf{Answer 9.5.} With $a_4'=2$ and $a_4''=0$: $R=6H^2(4-7)=-18H^2$; $G^0{}_0=-3H^2(4-5)=3H^2$; $G^i{}_i=H^2(15-12)=3H^2$; $G^4{}_4=3H^2(7+4)=33H^2$; $G^j{}_j=3H^2$. So $G^\mu{}_\nu=\mathrm{diag}(3,3,3,3,33,3,3,3)H^2$, $\rho_{\mathrm{req}}=-33H^2/\kappa$, all seven pressures $3H^2/\kappa$, and $w=3/(-33)=-1/11\approx-0.091$, in agreement with $w=(c^2-5)/(c^2+7)$ at $c=2$. On the plateau $A=2$ of the figure (units $H=\kappa=1$) panel (c) shows $\rho_{\mathrm{req}}=-33$ and panel (d) shows $w_{\mathrm{req}}\approx-0.09$.

\textbf{Answer 9.6.} Write the delta as the determinant of the matrix with entries $\delta^{u_r}_{l_c}$, upper list $u=(\mu_1,\nu_1,\mu_2,\nu_2)$ and lower list $l=(\alpha_1,\beta_1,\alpha_2,\beta_2)$. A determinant is the sum over permutations $\pi$ of $\mathrm{sgn}(\pi)\prod_r\delta^{u_r}_{l_{\pi(r)}}$. The permutation that exchanges positions 2 and 3 has sign $-1$ and sets $\alpha_1=\mu_1$, $\alpha_2=\nu_1$, $\beta_1=\mu_2$, $\beta_2=\nu_2$. Its term is $-R^{\mu_1\mu_2}{}_{\mu_1\nu_1}\,R^{\nu_1\nu_2}{}_{\mu_2\nu_2}$. In the first factor the first upper index is contracted with the first lower index, which gives the Ricci tensor $R^{\mu_2}{}_{\nu_1}$. In the second factor, reversing both pairs (two sign changes) gives $R^{\nu_2\nu_1}{}_{\nu_2\mu_2}=R^{\nu_1}{}_{\mu_2}$. The term is $-R^{\mu_2}{}_{\nu_1}R^{\nu_1}{}_{\mu_2}=-R^\mu{}_\nu R^\nu{}_\mu$.

\textbf{Answer 9.7.} $a_4'^2=\tfrac49(M\mp1)^2$ and $a_4''=0$. Then $R=6H^2\bigl(\tfrac49(M\mp1)^2-7\bigr)=\tfrac23H^2\bigl(4(M\mp1)^2-63\bigr)=\tfrac23H^2(4M^2\mp8M-59)$, and $\kappa\rho_{\mathrm{req}}=-3H^2\bigl(7+\tfrac49(M\mp1)^2\bigr)=-\tfrac13H^2\bigl(4(M\mp1)^2+63\bigr)$. A square is never negative, so the bracket is at least 63 and $\rho_{\mathrm{req}}<0$ for every real $M$. The common pressure is $\kappa p=H^2\bigl(15-\tfrac43(M\mp1)^2\bigr)=-\tfrac13H^2(4M^2\mp8M-41)$.

\textbf{Answer 9.8.} $\kappa(\rho+p_{(i)})=H^2\bigl[-3(7+a_4'^2)+15-3a_4'^2+a_4''\bigr]=H^2\bigl[-6-6a_4'^2+a_4''\bigr]$, which is negative unless $a_4''\ge6(1+a_4'^2)$. For the null vector $e_j+e_i$ the frame components are $T_{\hat j\hat j}=-p_{(j)}$ and $T_{\hat i\hat i}=p_{(i)}$, so $T(k,k)=p_{(i)}-p_{(j)}$, and $\kappa(p_{(i)}-p_{(j)})=H^2\bigl[(15-3a_4'^2+a_4'')-(15-3a_4'^2-a_4'')\bigr]=2H^2a_4''$.

\textbf{Answer 9.9.} Here all four directions along the brane are warped. For $y<0$, $g_{ab}=e^{2ky}\eta_{ab}$ and $K_{ab}=\tfrac12\partial_yg_{ab}=k\,g_{ab}$, so $K^a{}_b=k\,\delta^a_b$; for $y>0$, $K^a{}_b=-k\,\delta^a_b$. The jumps are $[K^a{}_b]=-2k\,\delta^a_b$ and $[K]=-8k$. The Israel condition gives $S^a{}_b=-\tfrac1\kappa\bigl(-2k+8k\bigr)\delta^a_b=-\tfrac{6k}{\kappa}\delta^a_b$. The energy density is $\rho=-S^0{}_0=6k/\kappa>0$ and every pressure is $p=S^i{}_i=-6k/\kappa=-\rho$: a pure positive tension $6k/\kappa$, as stated in Section 9.16.

\textbf{Answer 9.10.} Condition (2): $mS=-15\cdot\tfrac{12}{5}=-36=-36H^2/\kappa$. Condition (3): $\lambda S^2=\tfrac{25}{6}\cdot\tfrac{144}{25}=24$ and $2H^2(15-3c^2)/\kappa=2\cdot12=24$ for $c=1$. Then $M_{\mathrm{eff}}=-15+\tfrac{25}{6}\cdot\tfrac{12}{5}=-15+10=-5$, $\rho=mS+\tfrac\lambda2S^2=-36+12=-24$, $p=\tfrac\lambda2S^2=12$ and $w=12/(-24)=-\tfrac12$: exactly the required source of the notebook's $a_4=t$ in Section 9.10.

\textbf{Answer 9.11.} By the quotient rule, $\dfrac{d}{dy}\dfrac{y}{\sqrt{y^2+\varepsilon^2}}=\dfrac{\sqrt{y^2+\varepsilon^2}-y^2/\sqrt{y^2+\varepsilon^2}}{y^2+\varepsilon^2}=\dfrac{\varepsilon^2}{(y^2+\varepsilon^2)^{3/2}}$. So the integral is $\bigl[y/\sqrt{y^2+\varepsilon^2}\bigr]_{-\infty}^{\infty}=1-(-1)=2$, independent of $\varepsilon$. At $y=0$ the integrand is $\varepsilon^2/\varepsilon^3=1/\varepsilon$. As $\varepsilon\to0$ the spike becomes infinitely high and narrow with area 2: it tends to $2\delta(y)$, the second derivative of $\lvert y\rvert$.

\textbf{Answer 9.12.} $s=3^{-6}$ gives $s^{1/6}=1/3$ and $s^{1/3}=1/9$; $a_4=\ln3$ gives $e^{a_4}=3$ and $e^{2a_4}=9$. The left side is $1/\sqrt{(1/9)/9}=1/\sqrt{1/81}=9$; the correct right side is $e^{a_4}/s^{1/6}=3\cdot3=9$; the notebook's right side is $1/(3\cdot\tfrac13)=1$. For the toy matrix, $\lvert\gamma\rvert=\begin{pmatrix}0&1\\1&0\end{pmatrix}$ and

\[
\gamma'=\begin{pmatrix}0&\cosh a-\sinh a\\-\cosh a-\sinh a&0\end{pmatrix}=\begin{pmatrix}0&e^{-a}\\-e^{a}&0\end{pmatrix},\qquad \gamma'^2=\begin{pmatrix}-1&0\\0&-1\end{pmatrix},
\]

whereas $(e^{a}\gamma)^2=e^{2a}\gamma^2=-e^{2a}\cdot1$. They agree only for $a=0$. This is the 2 by 2 version of $(\gamma'^{\,x_5})^2=-s^{-1/3}$ against the required $-s^{-1/3}e^{2a_4}$ in Section 9.14.

\section{Solving differential equations on a computer from zero}

\subsection{What this chapter does}

Most equations of this book cannot be solved with pencil and paper once the background changes in time or the particles interact. The Stage-3 experiments of Chapter 11 follow a 16-component complex spinor through expanding and deflating universes; the Kohn--Sham calculation of Chapter 13 has to find the stationary states of many quanta in the static primordial field. In both cases the computer solves \textbf{ordinary differential equations}, equations for functions of one variable, with a program called \textbf{CVODE}. This chapter explains from zero what that program does, why each of its settings was chosen, and, most importantly, why we may trust its answers: every numerical result of Stage 3 is checked by independent programs against exact solutions, conservation laws and each other; the Stage-4 results and the status of their independent cross-check are described in Chapter 13.

The sources of this chapter are these files of the repository.

\begingroup
\small
\begin{longtable}{@{}>{\raggedright\arraybackslash}p{0.400\linewidth}>{\raggedright\arraybackslash}p{0.400\linewidth}@{}}
\toprule
\textbf{Source} & \textbf{What this chapter takes from it} \\
\midrule
\endfirsthead
\toprule
\textbf{Source} & \textbf{What this chapter takes from it} \\
\midrule
\endhead
\texttt{provenan\allowbreak{}ce/DIRAC\allowbreak{}16COMPLE\allowbreak{}X\_DARK\_S\allowbreak{}ECTOR\_NU\allowbreak{}MERICS.m\allowbreak{}d} & the Stage-3 document: §5 (method and verification strategy) and the solver and verification sections of each experiment, cited as §N \\
\texttt{provenan\allowbreak{}ce/DIRAC\allowbreak{}16COMPLE\allowbreak{}X\_STUDEN\allowbreak{}T\_GUIDE.\allowbreak{}md} & §4.5 (engines), §5.2 (FMA), §6.9 (real state layout), §6.16 (CVODE), §7 and §8 (program and files), §13 (exercises) \\
\texttt{provenan\allowbreak{}ce/DIRAC\allowbreak{}16COMPLE\allowbreak{}X\_PRIMOR\allowbreak{}DIAL\_FIE\allowbreak{}LD.md} & the Stage-2 document: §14.1 (the exact mode solution), §13.5 and §15.5 (negative controls), cited as Stage-2 §N \\
\texttt{handoff/\allowbreak{}specs/NU\allowbreak{}MERICS\_C\allowbreak{}ONTRACT.\allowbreak{}md} & the numerical programme and its errata \\
\texttt{studies/\allowbreak{}dirac16c\allowbreak{}omplex\_c\allowbreak{}osmology\allowbreak{}/src/dri\allowbreak{}ver.rs} & the CVODE driver \\
\texttt{studies/\allowbreak{}dirac16c\allowbreak{}omplex\_c\allowbreak{}osmology\allowbreak{}/src/out\allowbreak{}put.rs} and \texttt{src/exp1\allowbreak{}.rs} to \texttt{src/exp5\allowbreak{}.rs} & the output format and the solver settings of the five experiments \\
\texttt{scripts/\allowbreak{}check\_di\allowbreak{}rac16com\allowbreak{}plex\_exp\allowbreak{}1.py} to \texttt{...\_exp5\allowbreak{}.py} & the independent checkers \\
\texttt{artifact\allowbreak{}s/dirac1\allowbreak{}6complex\allowbreak{}/numeric\allowbreak{}s/expN/s\allowbreak{}ummary.j\allowbreak{}son} & settings, solver statistics and self-checks of experiment N = 1, ..., 5 \\
\texttt{artifact\allowbreak{}s/dirac1\allowbreak{}6complex\allowbreak{}/numeric\allowbreak{}s/expN/p\allowbreak{}ython-ch\allowbreak{}eck-repo\allowbreak{}rt.json} & the results of the independent checker of experiment N \\
\texttt{artifact\allowbreak{}s/dirac1\allowbreak{}6complex\allowbreak{}/numeric\allowbreak{}s/mathem\allowbreak{}atica-re\allowbreak{}port.jso\allowbreak{}n} & the checks and negative controls of the Mathematica notebook \\
\texttt{scripts/\allowbreak{}build\_di\allowbreak{}rac16com\allowbreak{}plex\_mat\allowbreak{}hematica\allowbreak{}\_noteboo\allowbreak{}k.wls} & the finite-difference test of the Mathematica notebook (function \texttt{\detokenize{fdCheck}}) \\
\texttt{studies/\allowbreak{}dirac16c\allowbreak{}omplex\_k\allowbreak{}ohn\_sham\allowbreak{}/src/sho\allowbreak{}oting.rs} and \texttt{handoff/\allowbreak{}specs/ST\allowbreak{}AGE4\_SPE\allowbreak{}C.md} & the level search of the Stage-4 shooting solver and its plan (§5) \\
\bottomrule
\end{longtable}
\endgroup

Every number about the project's runs in this chapter is copied from those files. The small examples that teach the methods (in Sections 10.2 to 10.8, 10.11 and 10.12) are ordinary arithmetic that the reader can repeat with a pocket calculator.

\subsection{Ordinary differential equations}

\textbf{Definitions.} Let $t$ be a real variable (a time, or a position along a line) and $y(t)$ an unknown function. An \textbf{ordinary differential equation} (ODE) is an equation that relates $y$ and some of its derivatives $y'=dy/dt$, $y''=d^2y/dt^2$, and so on, at the same $t$. Its \textbf{order} is the highest derivative that occurs. A \textbf{system} of ODEs has several unknown functions $y_0(t),\dots,y_{n-1}(t)$, which we collect into one column, the \textbf{state vector} $Y(t)=(y_0,\dots,y_{n-1})^T$. Every system we meet can be written in the \textbf{first-order form}

\[
Y'(t)=f\bigl(t,Y(t)\bigr),
\]

where $f$ is a given rule, the \textbf{right-hand side}, that produces a column of $n$ numbers from $t$ and $Y$. Together with a \textbf{start value} $Y(t_0)=Y_0$ it is an \textbf{initial-value problem}.

\textbf{Example 1 (decay).} $y'=-y$ with $y(0)=1$ has the solution $y(t)=e^{-t}$: indeed $\tfrac{d}{dt}e^{-t}=-e^{-t}$ and $e^0=1$.

\textbf{Example 2 (oscillation, and how to lower the order).} The oscillator $x''=-\omega^2x$ is second order. Introduce the velocity $v=x'$ as a second unknown. Then $x'=v$ and $v'=x''=-\omega^2x$, a first-order system with $Y=(x,v)^T$ and $f(t,Y)=(v,\,-\omega^2x)^T$. Every equation of higher order can be lowered in this way.

\textbf{Existence and uniqueness.} If $f$ is smooth (it has continuous derivatives with respect to $Y$), then the initial-value problem has exactly one solution, at least for some time around $t_0$ (the theorem of Picard and Lindelöf, quoted here without proof). A numerical method tries to approximate that one solution.

\textbf{The project's equation.} The mode equation of the Stage-3 experiments (Chapter 11; Stage-3 document §4.5) is

\[
i\,\dot u=h(t)\,u,\qquad u(t)\in\mathbb C^{16},
\]

where the dot is $d/dt$, $t=x_4$, and $h(t)$ is a 16 by 16 complex matrix, the mode Hamiltonian, for example $h=-iM_{\mathrm{eff}}\gamma^4-K\gamma^4\gamma^0$ in experiment EXP-1. CVODE works with real numbers only, so the complex equation is split into real and imaginary parts. Write $u=x+iy$ and $h=h_r+ih_i$ with real columns $x,y$ and real matrices $h_r,h_i$. Then $i(\dot x+i\dot y)=(h_r+ih_i)(x+iy)$. The left side is $i\dot x-\dot y$; the right side is $(h_rx-h_iy)+i(h_ix+h_ry)$. Comparing real parts and imaginary parts:

\[
\dot x=h_ry+h_ix,\qquad \dot y=h_iy-h_rx .
\]

These are 32 real first-order equations. Where further quantities are integrated together with the spinor, the state is longer: 38 reals in EXP-2 (the logarithms of the three scale factors of the metric and their three rates of change besides $u$), 35 in EXP-3 (the time, a distance and the logarithm $\ln\sigma$ of the scaled density $\sigma$ besides $u$) and 33 in the smooth-transition runs of EXP-4 (the logarithm $\ln a$ of the scale factor besides $u$) (Stage-3 document §7.2, §8.2 and §9.2; these experiments are Chapter 11). The program stores the spinor part as 32 consecutive numbers, first the 16 real parts, then the 16 imaginary parts (the columns \texttt{\detokenize{u_re_0}} to \texttt{\detokenize{u_re_15}} and \texttt{\detokenize{u_im_0}} to \texttt{\detokenize{u_im_15}} of every output file that stores the spinor, for example \texttt{exp1/run\allowbreak{}\_A1\_K0\_p\allowbreak{}os\_Bp.cs\allowbreak{}v}; student guide §6.9). \textbf{A worked component.} For zero momentum the matrix is $h=-iM\gamma^4$, so $h_r=0$ and $h_i=-M\gamma^4$. Row 0 of $\gamma^4$ has its single entry $-1$ in column 13 (Chapter 2), so $(\gamma^4x)_0=-x_{13}$ and the first equation reads $\dot x_0=(h_ix)_0=-M(-x_{13})=M\,x_{13}$: one multiplication.

\textbf{An exact solution to test against.} When $h$ is a constant matrix with $h^2=E^2\cdot1$ for a number $E>0$ (true in EXP-1, where $E=\sqrt{M_{\mathrm{eff}}^2+K^2}$; Stage-2 document §14.1), the solution is known exactly. Formally $u(t)=e^{-iht}u_0$, where the exponential of a matrix is defined by its power series $e^{A}=\sum_{n\ge0}A^n/n!$. Split the series into even and odd powers and use $h^{2k}=E^{2k}\cdot1$ and $h^{2k+1}=E^{2k}h$:

\[
e^{-iht}=\sum_{k\ge0}\frac{(-1)^k(Et)^{2k}}{(2k)!}\cdot1-i\,\frac hE\sum_{k\ge0}\frac{(-1)^k(Et)^{2k+1}}{(2k+1)!}=\cos(Et)\cdot1-i\sin(Et)\,\frac hE .
\]

\textbf{Worked example.} The 2 by 2 matrix $h=\begin{pmatrix}3&4\\4&-3\end{pmatrix}$ has $h^2=25\cdot1$, so $E=5$. For $u_0=(1,0)^T$:

\[
u(t)=\Bigl(\cos5t-\tfrac{3i}{5}\sin5t,\ -\tfrac{4i}{5}\sin5t\Bigr)^T,\qquad u^\dagger u=\cos^25t+\tfrac{9}{25}\sin^25t+\tfrac{16}{25}\sin^25t=1 .
\]

The length $u^\dagger u$ stays exactly 1 for all $t$. This is general: for a Hermitian $h$ ($h^\dagger=h$) the quantity $u^\dagger u$ is conserved, because $\tfrac{d}{dt}(u^\dagger u)=\dot u^\dagger u+u^\dagger\dot u=(iu^\dagger h)u+u^\dagger(-ihu)=0$. Exact solutions and conserved quantities are the two strongest tests of a numerical solution (Section 10.11).

\subsection{Euler's method and the idea of error}

\textbf{The method.} Choose a \textbf{step size} $h>0$ (here $h$ is a small number, not the Hamiltonian) and the times $t_n=t_0+nh$. Taylor's theorem says $y(t+h)=y(t)+h\,y'(t)+\tfrac12h^2y''(\xi)$ for some $\xi$ between $t$ and $t+h$. Dropping the last term and using $y'=f(t,y)$ gives \textbf{Euler's method}:

\[
Y_{n+1}=Y_n+h\,f(t_n,Y_n),
\]

where $Y_n$ is the computed approximation of $Y(t_n)$.

\textbf{Local and global error.} One step makes the error $\tfrac12h^2y''$, the \textbf{local error}, proportional to $h^2$. To reach a fixed end time $T$ one needs $N=(T-t_0)/h$ steps, and the local errors add up to a \textbf{global error} of order $N\cdot h^2\propto h$. A method whose global error is proportional to $h^p$ is said to have \textbf{order} $p$. Euler's method has order 1.

\textbf{Worked example.} Solve $y'=-y$, $y(0)=1$ up to $t=1$. Each Euler step multiplies by $1-h$, so $Y_N=(1-h)^N$ with $N=1/h$. The exact value is $e^{-1}=0.367879$.

\begingroup
\small
\begin{longtable}{@{}>{\raggedright\arraybackslash}p{0.160\linewidth}>{\raggedright\arraybackslash}p{0.160\linewidth}>{\raggedright\arraybackslash}p{0.160\linewidth}>{\raggedright\arraybackslash}p{0.160\linewidth}>{\raggedright\arraybackslash}p{0.160\linewidth}@{}}
\toprule
\textbf{Step $h$} & \textbf{Steps $N$} & \textbf{Euler value $(1-h)^N$} & \textbf{Error} & \textbf{Error ratio} \\
\midrule
\endfirsthead
\toprule
\textbf{Step $h$} & \textbf{Steps $N$} & \textbf{Euler value $(1-h)^N$} & \textbf{Error} & \textbf{Error ratio} \\
\midrule
\endhead
0.5 & 2 & 0.250000 & 0.117879 & --- \\
0.25 & 4 & 0.316406 & 0.051473 & 2.29 \\
0.125 & 8 & 0.343609 & 0.024271 & 2.12 \\
0.0625 & 16 & 0.356074 & 0.011805 & 2.06 \\
\bottomrule
\end{longtable}
\endgroup

Halving the step halves the error (the ratio tends to $2^1=2$): order 1 confirmed. The method \textbf{converges}, the error tends to 0 as $h\to0$, but slowly: the error is close to $0.19\,h$, so for 6 correct digits (an error below $5\times10^{-7}$) Euler would need a step of about $2.6\times10^{-6}$, that is about 400000 steps. Better methods (Section 10.6) reach the same accuracy with far fewer steps.

\subsection{Stability: decay and oscillation}

\textbf{The test equation.} Much can be learned from $y'=\lambda y$ with a constant, possibly complex, number $\lambda$. Its exact solution $y=e^{\lambda t}y_0$ decays if the real part $\mathrm{Re}\,\lambda<0$ and oscillates with constant size if $\lambda=i\omega$ is purely imaginary. A one-step method applied to it multiplies $y$ by a fixed \textbf{amplification factor} $R(z)$ per step, with $z=h\lambda$:

\begin{itemize}
\item \textbf{explicit Euler} $Y_{n+1}=Y_n+h\lambda Y_n$ gives $R(z)=1+z$;
\item \textbf{implicit (backward) Euler} $Y_{n+1}=Y_n+h\lambda Y_{n+1}$ gives, solving for $Y_{n+1}$, $R(z)=1/(1-z)$;
\item \textbf{the trapezoidal rule} $Y_{n+1}=Y_n+\tfrac h2\lambda(Y_n+Y_{n+1})$ gives $R(z)=(1+z/2)/(1-z/2)$.
\end{itemize}

A method is \textbf{stable} for a given $z$ if $\lvert R(z)\rvert\le1$, so that errors do not grow from step to step. The set of such $z$ is the method's \textbf{stability region}.

\textbf{Oscillation.} For $\lambda=i\omega$ the exact factor per step is $\lvert e^{i\omega h}\rvert=1$. Explicit Euler has $\lvert1+i\omega h\rvert=\sqrt{1+\omega^2h^2}>1$: the numerical oscillation grows. With $\omega=1$ and $h=0.1$, after 100 steps (up to $t=10$) the squared size $\lvert y\rvert^2$ has grown by $(1.01)^{100}=2.70$. Backward Euler has $\lvert R\rvert=1/\sqrt{1+\omega^2h^2}<1$: the oscillation is damped, $\lvert y\rvert^2$ falls by $(1.01)^{-100}=0.370$. The trapezoidal rule has $\lvert R\rvert=\lvert1+i\omega h/2\rvert/\lvert1-i\omega h/2\rvert=1$ exactly: it keeps the size of an oscillation (Exercise 10.3). The lesson for the mode equation $i\dot u=hu$, whose right-hand side $-ihu$ has purely imaginary eigenvalues when $h$ is Hermitian, is that a method can gain or lose norm purely by its own error. Monitoring the conserved $u^\dagger u$ therefore measures the accuracy of the integration directly (Section 10.11).

\textbf{Decay.} For $\lambda=-1$ the trapezoidal rule gives $y(1)=\bigl(\tfrac{1-h/2}{1+h/2}\bigr)^{1/h}$: $0.360000$ for $h=0.5$ (error $0.0079$) and $0.365950$ for $h=0.25$ (error $0.0019$). The error falls by a factor 4.1 when $h$ is halved: the trapezoidal rule has order 2.

\subsection{Stiffness}

\textbf{Definition.} A problem is \textbf{stiff} when it contains components that decay much faster than the solution we are interested in changes. An explicit method is then forced, by stability and not by accuracy, to take steps as small as the fastest decay time, even long after the fast components have died away.

\textbf{Example.} Consider $y'=-1000\,(y-\cos t)-\sin t$. Its solution with $y(0)=y_0$ is $y(t)=\cos t+(y_0-1)\,e^{-1000t}$ (check: $y'=-\sin t-1000(y_0-1)e^{-1000t}$, and $-1000(y-\cos t)=-1000(y_0-1)e^{-1000t}$). After a time of about $0.005$ the second term is gone and $y=\cos t$ changes slowly. But an error $\delta$ added to $y$ obeys $\delta'=-1000\,\delta$, the test equation with $\lambda=-1000$. Explicit Euler multiplies it by $1-1000h$ per step, and it stays bounded only if $\lvert1-1000h\rvert\le1$, that is $h\le0.002$. With $h=0.01$, a step that would describe $\cos t$ perfectly well, the factor is $-9$ and any rounding error grows ninefold per step. Backward Euler multiplies the error by $1/(1+1000h)=1/11$ at $h=0.01$: it is damped for every step size. The price is that each step of an implicit method requires solving an equation for $Y_{n+1}$ (Section 10.7).

\textbf{Are the project's equations stiff?} In the good sector (no momentum along the extra times) the mode equations are oscillations: $h$ is Hermitian and the right-hand side $-ihu$ has purely imaginary eigenvalues. In the extra-time experiment EXP-5 the matrix $h$ satisfies $h^2=E^2(t)\cdot1$ with $E^2(t)=m^2-Q(t)^2$, where $m$ is the mass and $Q(t)$ a momentum along an extra time that grows as the extra times deflate, so that $E^2$ turns negative at a time $t_\ast$ (Chapter 11). If $hv=\lambda v$ for a vector $v\ne0$, then $\lambda^2v=h^2v=E^2v$, so every eigenvalue of $h$ satisfies $\lambda^2=E^2$. Before $t_\ast$ ($E^2>0$) the eigenvalues of $-ih$ are therefore imaginary, $\mp iE$; after it ($E^2<0$) they are real, $\pm\varkappa$ with $\varkappa=\sqrt{Q^2-m^2}$. The runs have $m=1$ and end three time units after $t_\ast$, where $Q=me^3$ (Stage-3 document §10.2 and §10.3), so $\varkappa$ grows to at most $\sqrt{e^6-1}=20.06$ (computed here). So after $t_\ast$ one component grows and one decays, both at the same moderate rate $\varkappa$. In none of these problems does a component decay much faster than the solution changes, so none of them is \textbf{stiff} (Stage-3 document §5.1 and §10.2; header of \texttt{src/driv\allowbreak{}er.rs}). This is why four of the five experiments use a method designed for non-stiff problems (Adams--Moulton, Section 10.6), and why a test in EXP-4 found it both more accurate and cheaper than the stiff method (Section 10.8).

\subsection{Multistep methods: Adams and BDF}

A \textbf{linear multistep method} uses several past values to make one step. CVODE offers two families (student guide §6.16).

\textbf{Adams methods} integrate the equation, $Y(t_{n+1})=Y(t_n)+\int_{t_n}^{t_{n+1}}f\,dt$, and replace $f$ under the integral by the polynomial that interpolates its known values. \emph{Adams--Bashforth of order 2.} Through the two known values $f_{n-1}$ and $f_n$ (with $f_n=f(t_n,Y_n)$) passes the straight line $p(t)=f_n+(t-t_n)(f_n-f_{n-1})/h$. Its integral over $[t_n,t_n+h]$ is $hf_n+\tfrac{h^2}{2}\cdot\tfrac{f_n-f_{n-1}}{h}$, so

\[
Y_{n+1}=Y_n+\frac h2\bigl(3f_n-f_{n-1}\bigr).
\]

This is \textbf{explicit}: everything on the right is known. \emph{Adams--Moulton} methods include the unknown value $f_{n+1}$ among the interpolation points, which makes them \textbf{implicit} and more accurate. The simplest one that uses $f_n$ and $f_{n+1}$ is the trapezoidal rule of Section 10.4, $Y_{n+1}=Y_n+\tfrac h2(f_n+f_{n+1})$. CVODE's Adams--Moulton family has the orders 1 to 12.

\textbf{Backward differentiation formulas (BDF)} interpolate $Y$ instead of $f$: they pass a polynomial through the unknown $Y_{n+1}$ and the known $Y_n,Y_{n-1},\dots$, differentiate it at $t_{n+1}$, and set the derivative equal to $f(t_{n+1},Y_{n+1})$. With one past value this is backward Euler. \emph{BDF of order 2.} Let $s=(t-t_{n+1})/h$. The quadratic

\[
p(s)=Y_{n+1}+s\,(Y_{n+1}-Y_n)+\frac{s(s+1)}{2}\,(Y_{n+1}-2Y_n+Y_{n-1})
\]

takes the values $Y_{n+1}$, $Y_n$ and $Y_{n-1}$ at $s=0,-1,-2$ (put them in: at $s=-1$ the last term vanishes and $p=Y_n$; at $s=-2$ it is $Y_{n+1}-2(Y_{n+1}-Y_n)+(Y_{n+1}-2Y_n+Y_{n-1})=Y_{n-1}$). Its time derivative is $\tfrac1h\,dp/ds$, and at $s=0$

\[
\frac1h\Bigl[(Y_{n+1}-Y_n)+\tfrac12(Y_{n+1}-2Y_n+Y_{n-1})\Bigr]=\frac{\tfrac32Y_{n+1}-2Y_n+\tfrac12Y_{n-1}}{h}=f(t_{n+1},Y_{n+1}).
\]

CVODE's BDF family has the orders 1 to 5. BDF methods are very stable for decaying components, which makes them the standard choice for stiff problems, and like backward Euler they slightly damp oscillations.

\textbf{Variable step, variable order.} CVODE does not use a fixed $h$: after every step it estimates the error (Section 10.8) and chooses the next step size and the order within its family. The number of steps it takes is reported in every summary file of the project.

\textbf{One implicit equation per step.} Both implicit families lead to an equation of the form

\[
Y_{n+1}=\gamma\,f(t_{n+1},Y_{n+1})+a_n ,
\]

with a number $\gamma$ proportional to $h$ and a known vector $a_n$ built from past values. For the trapezoidal rule $\gamma=h/2$ and $a_n=Y_n+\tfrac h2f_n$; for BDF2, multiplying the BDF2 formula above, $(\tfrac32Y_{n+1}-2Y_n+\tfrac12Y_{n-1})/h=f(t_{n+1},Y_{n+1})$, by $\tfrac23h$ and solving for $Y_{n+1}$ gives $Y_{n+1}=\tfrac23h\,f(t_{n+1},Y_{n+1})+\tfrac43Y_n-\tfrac13Y_{n-1}$, so $\gamma=\tfrac23h$ and $a_n=\tfrac13(4Y_n-Y_{n-1})$.

\subsection{Solving the implicit equation: fixed point or Newton}

\textbf{Fixed-point iteration.} Start from a guess $Y^{(0)}$ (CVODE uses an explicit prediction) and repeat $Y^{(m+1)}=\gamma f(t_{n+1},Y^{(m)})+a_n$. The error of the iterate is multiplied at each iteration by about $\gamma J$, where $J=\partial f/\partial Y$ is the \textbf{Jacobian matrix} of the right-hand side (its entry in row $r$ and column $c$ is $\partial f_r/\partial Y_c$). \emph{Derivation.} Let $Y^\ast$ be the exact solution of the implicit equation, $Y^\ast=\gamma f(t_{n+1},Y^\ast)+a_n$, and let $e^{(m)}=Y^{(m)}-Y^\ast$ be the error of the $m$-th iterate. Subtracting the equation for $Y^\ast$ from the iteration gives $e^{(m+1)}=\gamma\bigl[f(t_{n+1},Y^\ast+e^{(m)})-f(t_{n+1},Y^\ast)\bigr]$. The linear approximation of Section 1.9, applied to each component $f_r$ as a function of the components $Y_0,Y_1,\dots$ of $Y$, gives for a small column $d$ the expansion $f_r(t,Y^\ast+d)=f_r(t,Y^\ast)+\sum_c(\partial f_r/\partial Y_c)\,d_c+\dots$, that is $f(t,Y^\ast+d)=f(t,Y^\ast)+J\,d+\dots$, where the dots are terms of second order in $d$. With $d=e^{(m)}$ and those terms dropped, $e^{(m+1)}\approx\gamma J\,e^{(m)}$. The iteration therefore converges when $\gamma$ times the size of $J$ is below 1. For the test equation $y'=\lambda y$ the Jacobian is the number $\lambda$, there are no second-order terms, and the factor is exactly $\gamma\lambda$. \emph{Worked example.} Backward Euler for $y'=-y$ with $h=0.1$ from $y_n=1$: $\gamma=h$, $a_n=1$, and the iteration is $Y\leftarrow1-0.1\,Y$. From $Y^{(0)}=1$ it gives $0.9$, $0.91$, $0.909$, $0.9091$, $0.90909$, approaching the exact solution $1/1.1=0.909090\ldots$; each iteration gains one digit, because $\gamma\lambda=-0.1$ (the errors $+0.0909$, $-0.0091$, $+0.0009$, ... shrink tenfold and alternate in sign). For the stiff $\lambda=-1000$ the same factor would be 100 and the iteration would diverge. Fixed-point iteration needs no Jacobian and no linear algebra, and it is the classical partner of Adams methods for non-stiff problems.

\textbf{Newton's method.} For one unknown $x$ and an equation $g(x)=0$, \textbf{Newton's method} replaces $g$ near the current guess $x^{(m)}$ by its linear approximation $g(x^{(m)})+g'(x^{(m)})\,(x-x^{(m)})$, the tangent line of the graph, and takes the zero of that line as the next guess:

\[
x^{(m+1)}=x^{(m)}-\frac{g(x^{(m)})}{g'(x^{(m)})}.
\]

\emph{Worked example.} For $g(x)=x^2-2$, with $g'(x)=2x$, from $x^{(0)}=1$: $x^{(1)}=1-(-1)/2=1.5$, $x^{(2)}=1.5-0.25/3=1.416667$ and $x^{(3)}=1.416667-0.006944/2.833333=1.414216$, against $\sqrt2=1.414214$. The errors $0.086$, $0.0025$ and $2.1\times10^{-6}$ are each roughly the square of the one before (times a fixed factor), so the number of correct digits roughly doubles at every iteration.

For the implicit equation of a step, write it as $G(Y)=Y-\gamma f(t_{n+1},Y)-a_n=0$. Near the current iterate the linear approximation of Section 1.9 gives $G(Y^{(m)}+d)=G(Y^{(m)})+(1-\gamma J)\,d+\dots$ for a small column $d$, because the matrix of derivatives of $Y$ with respect to $Y$ is the unit matrix $1$ and that of $\gamma f$ is $\gamma J$ (with $J$ taken at $Y^{(m)}$). Newton's method sets the linear part to zero, $(1-\gamma J)\,d=-G(Y^{(m)})$, and takes $Y^{(m+1)}=Y^{(m)}+d$:

\[
Y^{(m+1)}=Y^{(m)}-\bigl(1-\gamma J\bigr)^{-1}G\bigl(Y^{(m)}\bigr).
\]

For a linear right-hand side, $f(t,Y)=A(t)\,Y+b(t)$ with a matrix $A$ and a column $b$, the Jacobian is $J=A$, and $G(Y^{(m)}+d)=G(Y^{(m)})+(1-\gamma A)\,d$ holds exactly, with no dropped terms. The step $d$ therefore makes $G(Y^{(m+1)})=0$: one iteration is exact, whatever the stiffness (provided $1-\gamma A$ is invertible). The price is the Jacobian and the solution of a linear system with the matrix $1-\gamma J$. CVODE can estimate the Jacobian by \textbf{difference quotients}: column $c$ is approximately $\bigl[f(t,Y+\delta e_c)-f(t,Y)\bigr]/\delta$ for a small $\delta$, where $e_c$ is the $c$-th unit column, so one estimate costs one extra evaluation of $f$ per state component. For a linear $f$ this quotient is $\bigl[A(Y+\delta e_c)+b-AY-b\bigr]/\delta=Ae_c$, the column $c$ of $A$, so the estimate is exact up to rounding. The linear system is then solved by a \textbf{dense direct solver} (Gaussian elimination on the full matrix).

\textbf{The two configurations of the project.} The driver \texttt{studies/\allowbreak{}dirac16c\allowbreak{}omplex\_c\allowbreak{}osmology\allowbreak{}/src/dri\allowbreak{}ver.rs} offers exactly two, and every summary file records which one was used:

\begin{itemize}
\item BDF with Newton iteration and the dense direct linear solver, using the difference-quotient Jacobian, called \texttt{Method::\allowbreak{}Bdf} in the code: used for EXP-1 (according to the driver's header it is the house style of the related projects dirac-main and planet\_Mercury);
\item Adams--Moulton with fixed-point iteration and no linear solver, called \texttt{Method::\allowbreak{}Adams} in the code: used for EXP-2 to EXP-5.
\end{itemize}

The EXP-4 method test of Section 10.8 shows the cost of the Jacobian: its BDF runs spent 2464 and 164384 right-hand-side evaluations on difference-quotient Jacobians, which are $77\times32$ and $5137\times32$, one evaluation per component of the 32-dimensional state (this factorization is our reading of the recorded counts).

\subsection{Error control: rtol, atol and max\_step}

\textbf{The local error test.} After each step CVODE estimates the local error $e_c$ of every component $c$ of the state (by comparing the implicit solution with its explicit prediction; we quote this from the CVODE design and do not derive it). It then forms the \textbf{weighted root-mean-square norm}

\[
\lVert e\rVert=\sqrt{\frac1n\sum_{c=0}^{n-1}\Bigl(\frac{e_c}{\mathrm{rtol}\cdot\lvert Y_c\rvert+\mathrm{atol}}\Bigr)^2}
\]

and accepts the step if $\lVert e\rVert\le1$; otherwise it repeats the step with a smaller $h$. Here \textbf{rtol} is the \textbf{relative tolerance} (the accepted error as a fraction of the size of each component) and \textbf{atol} the \textbf{absolute tolerance} (the accepted error for components that are near zero, where $\mathrm{rtol}\cdot\lvert Y_c\rvert$ would be useless). A third setting, \textbf{max\_step}, caps the step size.

\textbf{Worked example.} Let $\mathrm{rtol}=10^{-6}$, $\mathrm{atol}=10^{-8}$, $Y=(2,0)$ and estimated errors $e=(10^{-6},5\times10^{-9})$. The weights are $10^{-6}\cdot2+10^{-8}=2.01\times10^{-6}$ and $0+10^{-8}=10^{-8}$; the ratios are $0.4975$ and $0.5$; the norm is $\sqrt{(0.4975^2+0.5^2)/2}=0.4988\le1$, so the step is accepted. Without atol the second weight would be zero and no step could ever be accepted for a component that passes through zero.

\textbf{Tolerances bound the local error, not the global error.} The per-step errors accumulate over thousands of steps. The figure shows it for EXP-1: the drifts of the conserved energy density and of the two norms grow steadily with time, while staying below $5\times10^{-9}$ (the committed values in \texttt{artifact\allowbreak{}s/dirac1\allowbreak{}6complex\allowbreak{}/numeric\allowbreak{}s/exp1/s\allowbreak{}ummary.j\allowbreak{}son}: largest drift of $\rho$ $4.19\times10^{-9}$, of the norms $2.70\times10^{-9}$). At the tolerance $\mathrm{rtol}=10^{-12}$ the largest distance from the exact solution over the 26 runs is $6.25\times10^{-9}$ (measurement \texttt{maxExact\allowbreak{}Error}), several thousand times rtol. Each experiment therefore checks itself against exact solutions or conserved quantities rather than trusting the tolerance.

\begin{figure}[htbp]
\centering
\includegraphics[width=0.92\linewidth]{artifacts/dirac16complex/numerics/figures/exp1_frozen_observables.png}
\caption{EXP-1 of Stage 3, where every quantity should stay exactly constant. Panels (a) to (d) show the energy density, the mean pressure, the Lagrangian split and the equation of state of several runs, flat as expected (the dashed mixed state oscillates for a physical reason explained in Chapter 11). Panels (e) and (f) show the numerical drifts of \texorpdfstring{$\rho$}{ρ} and of the Hilbert and Krein norms for all 26 runs: they grow with time as the local errors accumulate and stay below about \texorpdfstring{$5\times10^{-9}$}{5×10\textasciicircum{}-9}.}
\end{figure}

\textbf{The settings of the five experiments} (from the \texttt{toleranc\allowbreak{}es}, \texttt{\detokenize{solver}} and \texttt{solverTo\allowbreak{}tals} entries of the five \texttt{summary.\allowbreak{}json} files):

\begingroup
\small
\begin{longtable}{@{}>{\raggedright\arraybackslash}p{0.114\linewidth}>{\raggedright\arraybackslash}p{0.114\linewidth}>{\raggedright\arraybackslash}p{0.114\linewidth}>{\raggedright\arraybackslash}p{0.114\linewidth}>{\raggedright\arraybackslash}p{0.114\linewidth}>{\raggedright\arraybackslash}p{0.114\linewidth}>{\raggedright\arraybackslash}p{0.114\linewidth}@{}}
\toprule
\textbf{Run} & \textbf{Method} & \textbf{rtol} & \textbf{atol} & \textbf{max\_step} & \textbf{Steps} & \textbf{RHS evaluations} \\
\midrule
\endfirsthead
\toprule
\textbf{Run} & \textbf{Method} & \textbf{rtol} & \textbf{atol} & \textbf{max\_step} & \textbf{Steps} & \textbf{RHS evaluations} \\
\midrule
\endhead
EXP-1 & BDF + Newton + dense & $10^{-12}$ & $10^{-14}$ & 0.02 & 33866 & 34950 \\
EXP-2 & Adams + fixed point & $10^{-12}$ & $10^{-15}$ & 0.02 & 1779784 & 1781396 \\
EXP-3 & Adams + fixed point & $10^{-11}$ & $10^{-13}$ & 0.01 & 6074 & 10115 \\
EXP-4 & Adams + fixed point & $10^{-13}$ & $10^{-14}$ & 2.0 & 339071692 & 511425232 \\
EXP-5 & Adams + fixed point & $10^{-10}$ & $10^{-12}$ & 0.02 & 2548 & 4800 \\
\bottomrule
\end{longtable}
\endgroup

\textbf{Why these values: four measured stories.}

\begin{enumerate}
\item \textbf{EXP-1: tighten the tolerance, never relax the limit.} The checks of EXP-1 require the conserved quantities to drift by less than $10^{-8}$ (the \texttt{\detokenize{limits}} entry of its summary). At $\mathrm{rtol}=10^{-10}$ the quantities that should stay constant drifted by up to $1.2\times10^{-7}$ ($\rho$), $5.4\times10^{-8}$ (the Hilbert and Krein norms) and $1.4\times10^{-8}$ (the mass $M_{\mathrm{eff}}$ of the interacting run), and the pressure check failed as well, all against the limit $10^{-8}$, while the distance from the exact solution, $7.8\times10^{-8}$, still passed its limit $10^{-7}$ (student guide Exercise 13.1, which reruns EXP-1 with \texttt{--rtol 1\allowbreak{}e-10} and lets you watch the checks fail; the Stage-3 document §6.4 quotes drifts of about $4\times10^{-8}$, which this rerun does not reproduce). So the tolerance was tightened to $10^{-12}$ instead of loosening the limit.
\item \textbf{EXP-2: why a step cap.} Without max\_step the Adams method took steps of about 0.045 and made an amplitude error of about $5\times10^{-13}$ per step on the oscillating spinor. Over times of order $10^4$ this drifted the conserved scalar density by $1.7\times10^{-7}$ and pushed the residual of the Einstein constraint to $3\times10^{-8}$. With the cap 0.02 the relative constraint residual stays below $1.26\times10^{-10}$ (Stage-3 document §7.4; measurement \texttt{maxConst\allowbreak{}raintRel\allowbreak{}ative} in \texttt{exp2/sum\allowbreak{}mary.jso\allowbreak{}n}).
\item \textbf{EXP-4: error follows rtol.} Each thermal mode turns through about $10^5$ radians. The largest deviation of the conserved $u^\dagger u$ from 1 was $2.4\times10^{-4}$ at $\mathrm{rtol}=10^{-10}$ and $2.4\times10^{-7}$ at $10^{-13}$: a thousand times smaller tolerance, a thousand times smaller error (Stage-3 document §9.4).
\item \textbf{EXP-4: Adams against BDF, measured.} The program chooses its method at run time by a recorded test on two momentum nodes, comparing, over the window $1\le a\le2$ of the scale factor, with a reference run of the Adams method at $\mathrm{rtol}/10$ and $\mathrm{atol}/10$ (the \texttt{methodSe\allowbreak{}lection} entry of \texttt{exp4/sum\allowbreak{}mary.jso\allowbreak{}n}):
\end{enumerate}

\begingroup
\small
\begin{longtable}{@{}>{\raggedright\arraybackslash}p{0.133\linewidth}>{\raggedright\arraybackslash}p{0.133\linewidth}>{\raggedright\arraybackslash}p{0.133\linewidth}>{\raggedright\arraybackslash}p{0.133\linewidth}>{\raggedright\arraybackslash}p{0.133\linewidth}>{\raggedright\arraybackslash}p{0.133\linewidth}@{}}
\toprule
\textbf{Node, momentum} & \textbf{Method} & \textbf{Steps} & \textbf{RHS evaluations} & \textbf{Jacobian RHS evaluations} & \textbf{Largest error} \\
\midrule
\endfirsthead
\toprule
\textbf{Node, momentum} & \textbf{Method} & \textbf{Steps} & \textbf{RHS evaluations} & \textbf{Jacobian RHS evaluations} & \textbf{Largest error} \\
\midrule
\endhead
2, $k=0.9525$ & Adams & 1007 & 1370 & 0 & $1.25\times10^{-12}$ \\
2, $k=0.9525$ & BDF & 4573 & 4730 & 2464 & $2.27\times10^{-10}$ \\
47, $k=119.93$ & Adams & 68001 & 88979 & 0 & $1.71\times10^{-9}$ \\
47, $k=119.93$ & BDF & 308166 & 313247 & 164384 & $1.84\times10^{-8}$ \\
\bottomrule
\end{longtable}
\endgroup

On this oscillatory, non-stiff problem the Adams method is 11 to 180 times more accurate against that Adams reference, and it needs about 3.5 times fewer right-hand-side evaluations, or about 5.3 times fewer when the Jacobian evaluations of BDF are counted as well (ratios computed here from the table). Since the reference is itself an Adams run, a measure that needs no reference is a useful second opinion: the largest drift of the conserved $u^\dagger u$ (key \texttt{maxNormD\allowbreak{}rift} of the same entry) is $2.6\times10^{-12}$ for Adams against $4.5\times10^{-10}$ for BDF at node 2, and $9.9\times10^{-10}$ against $2.9\times10^{-8}$ at node 47, so it too is 29 to 173 times smaller for Adams (ratios computed here). This is what Sections 10.4 to 10.6 lead one to expect.

\subsection{CVODE and the project's driver}

\textbf{What CVODE is.} CVODE is the solver for initial-value problems of SUNDIALS, a suite of numerical software from Lawrence Livermore National Laboratory. It implements the variable-step, variable-order Adams--Moulton and BDF methods of Section 10.6, the two nonlinear iterations of Section 10.7 and the error test of Section 10.8. The project uses a pure-Rust translation of SUNDIALS 7.8.0 (\texttt{sundials\allowbreak{}\_rs}) from the rustSolveIt repositories of the same author. The setup scripts \texttt{scripts/\allowbreak{}setup\_so\allowbreak{}lver.ps1} and \texttt{scripts/\allowbreak{}setup\_so\allowbreak{}lver.sh} clone it at a pinned commit into the folder \texttt{vendor/r\allowbreak{}ustSolve\allowbreak{}It}, which Git ignores; the engine is licensed BSD-3-Clause and is fetched at build time, not redistributed (Stage-3 document §5.1):

\begin{Verbatim}[fontsize=\small]
https://github.com/once-ere/rustSolveIt_Win11_SUNDIALS_7_8_0
  a8fdff459adfe181573d7924b18bffbdf378fdb3
https://github.com/once-ere/rustSolveIt_macos-silicon_SUNDIALS_7_8_0
  5360157f4f6160978f66400566c31b2ae25dd44d
https://github.com/once-ere/rustSolveIt_linux_SUNDIALS_7_8_0
  6f58e02e53717a51375bd4bc5918edc57088d922
\end{Verbatim}

\textbf{The driver.} All five experiments call one function of \texttt{studies/\allowbreak{}dirac16c\allowbreak{}omplex\_c\allowbreak{}osmology\allowbreak{}/src/dri\allowbreak{}ver.rs}, \texttt{integrat\allowbreak{}e(y0, t0\allowbreak{}, target\allowbreak{}s, rhs, \allowbreak{}cfg)}. It integrates $Y'=f(t,Y)$ from $t_0$ through a list of output times and returns the state at $t_0$ and at every output time. In outline (the real code checks every return flag of SUNDIALS and frees every object in the order the C library requires):

\begin{Verbatim}[fontsize=\small]
integrate(y0, t0, targets, rhs, cfg):
    check: y0 finite, targets strictly monotone, rtol > 0, atol > 0, max_step >= 0
    create CVODE with lmm = BDF or ADAMS
    CVodeInit(rhs, t0, y0);  CVodeSStolerances(rtol, atol)
    if BDF:   dense matrix + dense linear solver (Newton, difference-quotient Jacobian)
    if Adams: fixed-point nonlinear solver, no linear solver
    CVodeSetMaxNumSteps(1000000, or 50000000 in EXP-4)
    if max_step > 0: CVodeSetMaxStep(max_step)
    if a stop time is given: CVodeSetStopTime(stop time)
    for each target: CVode(target, CV_NORMAL) -> state at target; stop on a negative flag
    record steps, rhs evaluations, Jacobian rhs evaluations, error-test failures
\end{Verbatim}

The maximum number of steps applies to each output interval; EXP-4 raises it to 50000000 for all its integrations (\texttt{MAX\_NUM\_\allowbreak{}STEPS} in \texttt{src/exp4\allowbreak{}.rs}; \texttt{print-co\allowbreak{}nfig} shows it). Three details matter for correctness. In the mode \texttt{CV\_NORMA\allowbreak{}L} CVODE may step beyond an intermediate output time and return the state there by \textbf{interpolation} with its own polynomial, so the output grid does not force small steps. Every integration of the project also sets a \textbf{stop time} equal to its last output time (EXP-2 and EXP-3 separately for their forward and backward runs, EXP-4 also at the junction $t=0$ of its sudden-transition pair modes, where the background changes its law and the integration is restarted), which keeps CVODE from stepping past the end. And the right-hand side is an ordinary Rust function; if it returns a non-finite number, CVODE is told that the failure is recoverable and retries with a smaller step, while an explicit error aborts the run with a named message.

\textbf{The program.} The Rust crate \texttt{studies/\allowbreak{}dirac16c\allowbreak{}omplex\_c\allowbreak{}osmology} has the subcommands \texttt{print-co\allowbreak{}nfig}, \texttt{\detokenize{exp1}} to \texttt{\detokenize{exp5}} and \texttt{\detokenize{all}}, and the options \texttt{--output\allowbreak{} DIR}, \texttt{\detokenize{--rtol X}}, \texttt{\detokenize{--atol X}} and \texttt{--refine\allowbreak{}d}; the last divides rtol and atol by 10 and max\_step by 2 (Stage-3 document §5.2). Following the convention of the rustSolveIt examples, every self-check prints a line \texttt{PASS - n\allowbreak{}ame: det\allowbreak{}ail} or \texttt{FAIL - n\allowbreak{}ame: det\allowbreak{}ail}, and the last line of the output is \texttt{\detokenize{SUCCESS}} (exit code 0) or \texttt{\detokenize{FAILURE}} (exit code 1). The five experiments run 69 self-checks, all passing. How to build and run it on Windows, macOS and Linux is Chapter 19 and the student guide.

\subsection{Numbers in a computer: rounding and reproducibility}

\textbf{Floating-point numbers.} A computer stores a real number as a \textbf{double}: a sign, 53 binary digits and an exponent (the IEEE 754 standard). Only finitely many numbers can be stored, so most results are rounded to the nearest double. The relative spacing of doubles is $2^{-52}\approx2.2\times10^{-16}$ (the \textbf{machine epsilon}), and the gap between a double $x$ and the next one is called $\mathrm{ulp}(x)$, the unit in the last place: $\mathrm{ulp}(1)=2.2\times10^{-16}$ and $\mathrm{ulp}(10^4)=2^{-39}\approx1.8\times10^{-12}$. The decimal number $10^{-10}$ has no exact double; the nearest one, written with 18 significant digits, is $1.00000000000000004\times10^{-10}$, which is why the program's \texttt{print-co\allowbreak{}nfig} shows the default tolerance of EXP-5 as \texttt{1.000000\allowbreak{}00000000\allowbreak{}004e-10} (student guide §7.1).

\textbf{How the files are written.} Every floating-point number in every output file is written by the C format \texttt{\%.17e}, one digit before the decimal point and 17 after it: 18 significant digits (\texttt{studies/\allowbreak{}dirac16c\allowbreak{}omplex\_c\allowbreak{}osmology\allowbreak{}/src/out\allowbreak{}put.rs}). Integers such as step counts are written as plain integers (for example \texttt{"steps":\allowbreak{} 769} in \texttt{exp1/sum\allowbreak{}mary.jso\allowbreak{}n}), and a JSON entry without a finite value is written as \texttt{\detokenize{null}}. The Stage-3 document says "17 significant digits"; the format has 17 digits after the point and hence 18 significant digits, as the student guide §8.1 states. Seventeen significant digits already identify every double uniquely, so no information is lost and two files can be compared byte by byte. The files use line feeds only and contain no dates, timings or folder names.

\textbf{Rounding can set a floor that no tolerance removes.} In EXP-2 the runs last until $t\approx10^4$; the three runs together take 1779784 steps, several hundred thousand each. CVODE advances its internal time by $t_{n+1}=t_n+h_n$, and near $t=10^4$ each addition is rounded to a multiple of about $1.8\times10^{-12}$. The spinor's phase rotates at the rate $M_{\mathrm{eff}}$, so the phase error can accumulate up to about $\max\lvert M_{\mathrm{eff}}\rvert\cdot N_{\mathrm{steps}}\cdot\mathrm{ulp}(t_{\mathrm{end}})/2$, recorded as \texttt{maxSpino\allowbreak{}rPhaseRo\allowbreak{}undingBo\allowbreak{}und} $=1.32\times10^{-6}$ in \texttt{exp2/sum\allowbreak{}mary.jso\allowbreak{}n}. The measured phase error is $2.24\times10^{-7}$ (\texttt{maxSpino\allowbreak{}rPhaseEr\allowbreak{}ror}), below the bound. The decisive test: in the refined run the phase error does \textbf{not} shrink ($2.71\times10^{-7}$), whereas the phase-invariant shape error of the spinor shrinks from $2.0\times10^{-10}$ to $5.2\times10^{-11}$ (Stage-3 document §7.6). A truncation error would have shrunk; a rounding floor does not. Only phase-invariant quantities enter the physical results.

\textbf{Determinism and byte identity.} Run twice with the same inputs, the program must write byte-identical files. That requires more than correct mathematics:

\begin{itemize}
\item a deterministic mathematical library: the engine brings its own implementation of the elementary functions (\texttt{sundials\allowbreak{}\_libm}) instead of using the operating system's;
\item the same machine instructions: the file \texttt{.cargo/c\allowbreak{}onfig.to\allowbreak{}ml} compiles with the processor feature FMA (fused multiply-add, $a\cdot b+c$ with a single rounding), which the engine's library and the pinned results assume (student guide §5.2);
\item a fixed order of operations: EXP-4 distributes its modes over up to 8 threads, but every mode is integrated independently and the results are collected in a fixed order, so the outcome does not depend on the number of threads (header of \texttt{parallel\allowbreak{}\_map} in \texttt{src/exp4\allowbreak{}.rs});
\item output without timings, dates or absolute paths.
\end{itemize}

Each independent checker reruns the program into a fresh folder and requires every output file to be byte-identical to the committed one (check \texttt{repeatBy\allowbreak{}teIdenti\allowbreak{}ty} in every \texttt{python-c\allowbreak{}heck-rep\allowbreak{}ort.json}; all five are true).

\textbf{Byte identity depends on the engine.} The three rustSolveIt engines are not identical: the Windows 11 engine's mathematical library differs from the one of the macOS and Linux engines, and some functions (the logarithm for some arguments) differ in the last bit. With the Windows 11 engine every committed file is reproduced byte for byte, on Windows 11 and on Ubuntu 24.04. With the Linux engine all 69 self-checks and all 167 checker checks still pass, but 14 committed files differ in their last digits (student guide §4.5; erratum in \texttt{handoff/\allowbreak{}specs/NU\allowbreak{}MERICS\_C\allowbreak{}ONTRACT.\allowbreak{}md}). Byte identity is a test of reproducibility with a fixed engine; agreement to the tolerance is the test of correctness.

\subsection{Independent checkers: how to trust a numerical result}

A number printed by a program is not yet a result. Stage 3 accepts a numerical result only after an \textbf{independent} program has confirmed it, and the checkers follow four rules (Stage-3 document §5.3): they are separate code, in Python with numpy and the standard library only; they rebuild the gamma matrices from the exact algebra fixture \texttt{artifact\allowbreak{}s/dirac1\allowbreak{}6complex\allowbreak{}/arbitra\allowbreak{}ry-field\allowbreak{}/algebra\allowbreak{}-fixture\allowbreak{}.json} instead of trusting the program's copy; they recompute every physical quantity from the raw state columns of the output files, never from the program's derived columns; and they test the program against mathematics that does not involve the program. Six kinds of test are used.

\textbf{1. Exact solutions.} Where an exact solution exists it is computed independently. For EXP-1 the checker computes $u(t)=e^{-iht}u_0$ from the eigenvalues and eigenvectors of $h$ (numpy's \texttt{\detokenize{eigh}}) and finds the largest distance $6.25\times10^{-9}$ from the CVODE solution over all 26 runs, against the limit $10^{-7}$ (\texttt{exactMax\allowbreak{}Error} in \texttt{exp1/pyt\allowbreak{}hon-chec\allowbreak{}k-report\allowbreak{}.json}).

\textbf{2. Conservation laws.} In the good sector the Hilbert norm $u^\dagger u$ is conserved (Section 10.2), and in every sector the Krein norm $u^\dagger Bu$ is conserved, because the mode Hamiltonian satisfies $h^\dagger B=Bh$ (Chapter 8) and therefore $\tfrac{d}{dt}(u^\dagger Bu)=(iu^\dagger h^\dagger)Bu+u^\dagger B(-ihu)=iu^\dagger(h^\dagger B-Bh)u=0$. In EXP-1 both drift by at most $2.70\times10^{-9}$ (\texttt{normMaxD\allowbreak{}rift}). EXP-5 is the hard case: there $h$ is not Hermitian, $u^\dagger u$ grows to about $10^{16}$, and the Krein norm must nevertheless stay constant; its drift, divided by $\max(u^\dagger u,1)$, stays below $1.08\times10^{-9}$ (\texttt{kreinMax\allowbreak{}Normaliz\allowbreak{}edDrift} in \texttt{exp5/pyt\allowbreak{}hon-chec\allowbreak{}k-report\allowbreak{}.json}).

\begin{figure}[htbp]
\centering
\includegraphics[width=0.92\linewidth]{artifacts/dirac16complex/numerics/figures/exp5_krein.png}
\caption{EXP-5 of Stage 3: (a) the drift of the Krein norm, divided by \texorpdfstring{$\max(u^\dagger u,1)$}{(u\textasciicircum{}†u,1)}, for the four runs, below the limit \texorpdfstring{$10^{-8}$}{10\textasciicircum{}-8} (dashed) throughout; (b) the Krein norm itself, drawn while \texorpdfstring{$u^\dagger u<10^8$}{u\textasciicircum{}†u<10\textasciicircum{}8}, stays at its initial value \texorpdfstring{$\pm E_0/m$}{±E\_0/m}. Later, when \texorpdfstring{$u^\dagger u$}{u\textasciicircum{}†u} approaches \texorpdfstring{$10^{16}$}{10\textasciicircum{}16}, the computed \texorpdfstring{$u^\dagger Bu$}{u\textasciicircum{}†Bu} is a difference of numbers of that size and loses its digits to rounding (absolute drift up to 1.005, \texttt{kreinDri\allowbreak{}ft.maxAb\allowbreak{}s} in \texttt{exp5/sum\allowbreak{}mary.jso\allowbreak{}n}), which is why panel (a) divides the drift by \texorpdfstring{$\max(u^\dagger u,1)$}{(u\textasciicircum{}†u,1)}.}
\end{figure}

\textbf{3. Finite-difference residuals.} The checker inserts the stored solution back into the equation. The derivative is estimated from five neighbouring samples with spacing $\Delta$:

\[
u'(t)\approx\frac{u(t-2\Delta)-8u(t-\Delta)+8u(t+\Delta)-u(t+2\Delta)}{12\Delta}.
\]

\emph{Derivation.} By Taylor's theorem $u(t\pm\Delta)=u\pm\Delta u'+\tfrac{\Delta^2}{2}u''\pm\tfrac{\Delta^3}{6}u'''+\tfrac{\Delta^4}{24}u''''\pm\tfrac{\Delta^5}{120}u^{(5)}+\dots$, so $u(t+\Delta)-u(t-\Delta)=2\Delta u'+\tfrac{\Delta^3}{3}u'''+\tfrac{\Delta^5}{60}u^{(5)}+\dots$ and, with $2\Delta$ in place of $\Delta$, $u(t+2\Delta)-u(t-2\Delta)=4\Delta u'+\tfrac{8\Delta^3}{3}u'''+\tfrac{32\Delta^5}{60}u^{(5)}+\dots$. Eight times the first minus the second removes the $u'''$ terms:

\[
8\bigl[u(t+\Delta)-u(t-\Delta)\bigr]-\bigl[u(t+2\Delta)-u(t-2\Delta)\bigr]=12\Delta\,u'-\tfrac{24}{60}\Delta^5u^{(5)}+\dots,
\]

and dividing by $12\Delta$ gives $u'-\tfrac{\Delta^4}{30}u^{(5)}$. The formula is exact up to $\tfrac{\Delta^4}{30}u^{(5)}$. For an oscillation with frequency at most $E$, $\lvert u^{(5)}\rvert\le E^5\lvert u\rvert$. \emph{Derivation.} For a constant $h$ each derivative of a solution multiplies it by $-ih$ ($\dot u=-ihu$, $\ddot u=-ih\dot u=(-ih)^2u$, and so on), so $u^{(5)}=(-ih)^5u=-i\,h^5u$. In EXP-1, $h^2=E^2\cdot1$ gives $h^5=E^4h$, and for the Hermitian $h$ we have $\lvert hu\rvert^2=u^\dagger h^\dagger hu=u^\dagger h^2u=E^2\lvert u\rvert^2$; hence $\lvert u^{(5)}\rvert=E^4\lvert hu\rvert=E^5\lvert u\rvert$ exactly. For a general constant Hermitian $h$ whose eigenvalues $\lambda_j$ all have modulus at most $E$, expand $u=\sum_jc_jv_j$ in orthonormal eigenvectors $v_j$ (the spectral theorem of Section 1.6); then $h^5u=\sum_jc_j\lambda_j^5v_j$ and $\lvert h^5u\rvert^2=\sum_j\lvert c_j\rvert^2\lambda_j^{10}\le E^{10}\sum_j\lvert c_j\rvert^2=E^{10}\lvert u\rvert^2$. So the residual $\lVert\text{formula}-(-ihu)\rVert$ must stay below about $\tfrac{\Delta^4}{30}E^5$ for a spinor with $\lvert u\rvert=1$, as in EXP-1. The EXP-1 checker allows $1.5\,\tfrac{\Delta^4}{30}E^5+10^{-6}$ and finds at most 0.63 of it (\texttt{fdResidu\allowbreak{}alWorstR\allowbreak{}atioToBo\allowbreak{}und}); EXP-5 finds at most 0.71.

\textbf{4. Refined-tolerance convergence.} The checker reruns the program with \texttt{--refine\allowbreak{}d} (rtol and atol divided by 10, max\_step by 2) and requires the errors to shrink. In EXP-5 the relative error against an independent reference falls from $3.34\times10^{-8}$ to $2.28\times10^{-9}$ (\texttt{canonica\allowbreak{}lMaxRela\allowbreak{}tiveErro\allowbreak{}r} and \texttt{refinedM\allowbreak{}axRelati\allowbreak{}veError} in \texttt{exp5/pyt\allowbreak{}hon-chec\allowbreak{}k-report\allowbreak{}.json}). In EXP-1 every run's distance from the exact solution strictly decreases. Where a floor was identified instead, such as the rounding floor of Section 10.10, the check requires the error to stay within that floor rather than to shrink (Stage-3 document §5.3).

\textbf{5. Independent integrators.} Where no exact solution exists, the checkers integrate selected modes again with methods that share no code with CVODE. The classical \textbf{fourth-order Runge--Kutta method} (RK4) makes a step with four evaluations of $f$,

\[
k_1=f(t_n,Y_n),\quad k_2=f\bigl(t_n+\tfrac h2,Y_n+\tfrac h2k_1\bigr),\quad k_3=f\bigl(t_n+\tfrac h2,Y_n+\tfrac h2k_2\bigr),\quad k_4=f(t_n+h,Y_n+hk_3),
\]

\[
Y_{n+1}=Y_n+\tfrac h6\bigl(k_1+2k_2+2k_3+k_4\bigr).
\]

For $y'=-y$ one RK4 step multiplies by $1-h+\tfrac{h^2}{2}-\tfrac{h^3}{6}+\tfrac{h^4}{24}$, which is $0.606771$ for $h=0.5$; two steps give $y(1)\approx0.368171$, an error of only $2.9\times10^{-4}$, against $0.118$ for Euler with the same step (Section 10.3). The error of such a reference is itself estimated by \textbf{Richardson's rule}: for a method of order $p$, halving the step reduces the error by about $2^p$, so the difference $Y_{h/2}-Y_h$ is about $(2^p-1)$ times the error of $Y_{h/2}$. The EXP-4 and EXP-5 checkers use RK4 with a Richardson estimate, and the EXP-4 checker also uses a fourth-order \textbf{Magnus method}, which advances $i\dot u=h(t)u$ by the exponential of $-i$ times a Hermitian matrix, a unitary matrix, and therefore keeps $u^\dagger u$ exactly (Stage-3 document §9.6 and §10.6). A Mathematica notebook re-integrates runs of all five experiments with Mathematica's NDSolve (for EXP-4 only a subset: four of its 48 thermal modes, two of them only up to the scale factor $a=10$ instead of 100, and ten pair-creation modes; Stage-3 document §9.6), including an order-8 Runge--Kutta method and a 32-digit arithmetic run (Stage-3 document §11.2); the figure shows its agreement with CVODE.

\begin{figure}[htbp]
\centering
\includegraphics[width=0.92\linewidth]{artifacts/dirac16complex/numerics/figures/mathematica/cross_check_deviations.png}
\caption{Agreement of Mathematica's NDSolve with the Rust CVODE solutions, as digits of agreement (\texorpdfstring{$-\log_{10}$}{-\_10} of the largest deviation) for each experiment; "aligned" means that a constant phase was removed first. From the Stage-3 Mathematica notebook.}
\end{figure}

\textbf{6. Negative controls.} A test that cannot fail proves nothing. The checks are therefore also run with deliberately wrong physics and must fail. The Mathematica notebook has its own finite-difference test, built like test 3 but with nine-point (eighth-order) derivatives: row by row (one row per output time, leaving out the rows that the output grid samples too coarsely) it compares the derivative of the stored solution with the right-hand side $f$, divides the largest difference of a component by the largest size of a component of $f$ in that row (the \textbf{relative residual}), and uses the difference between the eighth-order and a sixth-order derivative as the row's \textbf{truncation estimate}. A row passes only if its residual is at most twice its truncation estimate plus an allowance, proportional to the Rust tolerance, for the noise of the stored samples (function \texttt{\detokenize{fdCheck}} of \texttt{scripts/\allowbreak{}build\_di\allowbreak{}rac16com\allowbreak{}plex\_mat\allowbreak{}hematica\allowbreak{}\_noteboo\allowbreak{}k.wls}). With the correct sign of the mass the residual is at most $1.0\times10^{-10}$ for EXP-1 and $7.1\times10^{-5}$ for EXP-3, against truncation estimates of at most $1.4\times10^{-8}$ and $5.4\times10^{-4}$, and every row passes. With the sign of the mass flipped (in one run of each experiment) the largest residual becomes 1.19 times the size of the right-hand side for EXP-1 and 2.0 times for EXP-3. For EXP-3 the value 2 is no accident: there the right-hand side of the spinor equation is proportional to the mass, so flipping its sign turns $f$ into $-f$, and the finite-difference derivative of the stored solution, which is close to $f$, differs from $-f$ by about $2f$. Both controls therefore miss their pass bound by orders of magnitude, as they must (keys \texttt{maxRelat\allowbreak{}iveResid\allowbreak{}ual}, \texttt{maxRelat\allowbreak{}iveTrunc\allowbreak{}ationEst\allowbreak{}imate} and \texttt{negative\allowbreak{}ControlF\allowbreak{}lippedMa\allowbreak{}ssMaxRel\allowbreak{}ativeRes\allowbreak{}idual} of \texttt{exp1.fin\allowbreak{}iteDiffe\allowbreak{}rence} and \texttt{exp3.fin\allowbreak{}iteDiffe\allowbreak{}rence} in \texttt{artifact\allowbreak{}s/dirac1\allowbreak{}6complex\allowbreak{}/numeric\allowbreak{}s/mathem\allowbreak{}atica-re\allowbreak{}port.jso\allowbreak{}n}; Stage-3 document §6.6 and §8.6). The exact Stage-2 verifier contains controls of the same kind: the conservation law $\nabla_\mu T^\mu{}_\nu=0$ must fail off shell, and the exact source of Chapter 9 must fail for a wrong slope of $a_4$ (Stage-2 document §13.5 and §15.5).

\textbf{The count.} Over the five experiments: 69 of 69 self-checks of the Rust program, 167 of 167 checks of the five independent checkers, 10 of 10 checks of the EXP-3 analysis, all 49 checks of the Mathematica notebook and all 71 assertions of the Jupyter notebook are true (Stage-3 document §13). Chapter 11 describes what these experiments show.

\subsection{Boundary-value problems and shooting}

\textbf{Two-point problems.} Some problems do not give all conditions at one end. A Kohn--Sham orbital of Chapter 13 must satisfy one condition at the tip cutoff $y=-L$ and one at the brane $y=0$, and it exists only for special values of its energy $\varepsilon$: an \textbf{eigenvalue problem} for an ODE. The \textbf{shooting method} turns it into initial-value problems: guess $\varepsilon$, integrate from one end with the condition there, measure how badly the condition at the other end fails (the \textbf{mismatch}), and adjust $\varepsilon$ until the mismatch is zero. The Stage-4 Rust solver \texttt{studies/\allowbreak{}dirac16c\allowbreak{}omplex\_k\allowbreak{}ohn\_sham} shoots with CVODE from $-L$ to 0. Instead of the raw mismatch it uses the \textbf{Prüfer angle}: it writes the two real components $a$ and $b$ of the orbital as $a=r\cos\theta$ and $b=-r\sin\theta$, integrates the equation for the angle $\theta(y)$ from $\theta(-L)=0$, and takes the end value $\Theta(\varepsilon)=\theta(0;\varepsilon)$, which increases with $\varepsilon$. Each level is the solution of $\Theta(\varepsilon)=n\pi$ (or $\tfrac\pi2+n\pi$, depending on the parity of the orbital) for an integer $n$ (Section 13.10). The solver first \textbf{brackets} each target by doubling the step in $\varepsilon$ until $\Theta(\varepsilon)$ minus the target changes sign. It then refines it by the \textbf{regula falsi}: the secant rule applied to the two ends of the bracket, keeping, as bisection does, the part in which the sign changes, so that the root always stays bracketed. In the plain regula falsi one end can stay fixed for many steps, and the bracket then shrinks only slowly from the other side. The solver therefore uses the \textbf{Illinois} variant: when the same end is replaced twice in a row, the function value kept at the other end is halved, which pulls the next point towards the end that has stayed fixed, so that this end is soon replaced as well (header and \texttt{find\_lev\allowbreak{}el} of \texttt{src/shoo\allowbreak{}ting.rs}; the plan \texttt{handoff/\allowbreak{}specs/ST\allowbreak{}AGE4\_SPE\allowbreak{}C.md} §5 had named bisection and the secant rule, the two ingredients shown in the worked example below; Exercise 10.13 compares the two versions of the regula falsi). The solver's results, and the status of its independent cross-check, are the subject of Chapter 13.

\textbf{Worked example.} Find the values $E$ for which $u''=-Eu$ has a solution with $u(0)=0$ and $u(1)=0$. As a first-order system, $u'=v$ and $v'=-Eu$; start with $u(0)=0$, $v(0)=1$. For $E>0$ the solution is $u(x)=\sin(\sqrt Ex)/\sqrt E$, so the mismatch is $F(E)=u(1)=\sin\sqrt E/\sqrt E$ (a computer would get it from CVODE; here we know it exactly). Since $F(5)=0.3518>0$ and $F(15)=-0.1725<0$, a zero lies between. \textbf{Bisection} halves the bracket and keeps the half on which $F$ changes sign:

\begingroup
\small
\begin{longtable}{@{}>{\raggedright\arraybackslash}p{0.200\linewidth}>{\raggedright\arraybackslash}p{0.200\linewidth}>{\raggedright\arraybackslash}p{0.200\linewidth}>{\raggedright\arraybackslash}p{0.200\linewidth}@{}}
\toprule
\textbf{Step} & \textbf{Midpoint $E$} & \textbf{$F(E)$} & \textbf{New bracket} \\
\midrule
\endfirsthead
\toprule
\textbf{Step} & \textbf{Midpoint $E$} & \textbf{$F(E)$} & \textbf{New bracket} \\
\midrule
\endhead
1 & 10 & $-0.00654$ & [5, 10] \\
2 & 7.5 & $+0.14320$ & [7.5, 10] \\
3 & 8.75 & $+0.06170$ & [8.75, 10] \\
4 & 9.375 & $+0.02601$ & [9.375, 10] \\
5 & 9.6875 & $+0.00935$ & [9.6875, 10] \\
6 & 9.84375 & $+0.00131$ & [9.84375, 10] \\
\bottomrule
\end{longtable}
\endgroup

Each step gains one binary digit. The \textbf{secant rule} instead draws the straight line through the last two points and takes its zero, $E_{\mathrm{new}}=E_1-F(E_1)\,(E_1-E_0)/(F(E_1)-F(E_0))$. From $E_0=5$ and $E_1=15$ it gives 11.7108, 8.8043, 10.0237, 9.88156, 9.869463 and 9.8696045: after six steps seven correct digits of the exact answer $\pi^2=9.8696044$, because $u(x)=\sin(\pi x)/\pi$ vanishes at $x=1$.

\subsection{What we proved and what we assumed}

\textbf{Derived in this chapter:} how a complex 16-component mode equation becomes 32 real first-order equations; the exact solution $u(t)=\bigl(\cos Et-i\sin Et\,h/E\bigr)u_0$ when $h^2=E^2\cdot1$, and the conservation of $u^\dagger u$ for Hermitian $h$ (and of $u^\dagger Bu$ when $h^\dagger B=Bh$); the order of Euler's method; the amplification factors of explicit Euler, backward Euler and the trapezoidal rule and what they do to decaying and oscillating solutions; why stiffness forces implicit methods, and that the eigenvalues of $-ih$ are $\mp iE$ or $\pm\varkappa$ when $h^2=E^2\cdot1$ with $E^2$ positive or negative; the Adams--Bashforth, trapezoidal and BDF2 formulas and the implicit equation of each step; the convergence factor $\gamma J$ of fixed-point iteration; Newton's method from the linear approximation, its exactness for linear problems and the exactness of the difference-quotient Jacobian for a linear right-hand side; the weighted error norm; the five-point derivative formula with its error $\tfrac{\Delta^4}{30}u^{(5)}$ and the bound $\lvert u^{(5)}\rvert\le E^5\lvert u\rvert$ for a constant Hermitian $h$ whose eigenvalues have modulus at most $E$ (with equality when $h^2=E^2\cdot1$); the RK4 step and Richardson's rule; bisection, the secant rule and the regula falsi with its Illinois variant for shooting. \textbf{Computed by the project and copied from its committed reports:} the methods, tolerances and step counts of the five Stage-3 experiments; the drifts, exact-solution errors, finite-difference ratios and refined-run errors quoted above; the failing checks of EXP-1 at $\mathrm{rtol}=10^{-10}$ (student guide Exercise 13.1); the recorded Adams-against-BDF test, with its errors against an Adams reference and its norm drifts; the residuals, truncation estimates and flipped-mass negative controls of the Mathematica finite-difference test; the byte identity of repeat runs; and the totals of 69 self-checks, 167 checker checks, 49 Mathematica checks and 71 notebook assertions, all true.

\textbf{Assumed or quoted, not proved here:} the existence and uniqueness theorem of Picard and Lindelöf; Taylor's theorem and the spectral theorem of Chapter 1; how fast Newton's method and the Illinois regula falsi converge in general, which the worked examples and Exercise 10.13 only illustrate; the internal design of CVODE (its error estimate and its choice of step and order), which we describe but do not derive, and the correctness of its Rust translation, which is tested only through the checks of Section 10.11; the fact that 17 significant digits identify a double uniquely; the independence of the EXP-4 results from the number of threads, which rests on the code's design and on the repeat-run check; and the engine differences of Section 10.10, which were measured and recorded in the student guide. The Stage-4 shooting solver is only introduced here: the Prüfer angle, its increase with $\varepsilon$ and the level condition $\Theta(\varepsilon)=n\pi$ or $\tfrac\pi2+n\pi$ are derived in Section 13.10, and its accuracy and cross-check status belong to Chapter 13.

\subsection{Exercises}

\textbf{Exercise 10.1.} Write the damped oscillator $x''+2x'+5x=0$ as a first-order system $Y'=f(t,Y)$, and check that $x=e^{-t}\cos2t$ solves it.

\textbf{Exercise 10.2.} Apply Euler's method to $y'=-2y$, $y(0)=1$, up to $t=1$ with $h=0.25$ and with $h=0.125$. Compare with $e^{-2}$ and determine the order from the two errors.

\textbf{Exercise 10.3.} Show that the trapezoidal rule applied to $y'=i\omega y$ keeps $\lvert y\rvert$ exactly constant, for every real $\omega$ and every step $h$.

\textbf{Exercise 10.4.} For $y'=-50y$, find the largest step for which explicit Euler is stable, and the factor by which backward Euler multiplies $y$ per step when $h=0.1$.

\textbf{Exercise 10.5.} Derive the BDF2 formula a second way: find numbers $a,b,c$ such that $(a\,y_{n+1}+b\,y_n+c\,y_{n-1})/h$ equals $y'(t_{n+1})$ exactly for $y=1$, $y=t$ and $y=t^2$. (Put $t_{n+1}=0$, $t_n=-h$, $t_{n-1}=-2h$.)

\textbf{Exercise 10.6.} With $\mathrm{rtol}=10^{-6}$ and $\mathrm{atol}=10^{-9}$, a step has the state $Y=(1,\,10^{-3},\,0)$ and the error estimate $e=(5\times10^{-7},\,2\times10^{-9},\,10^{-9})$. Is the step accepted?

\textbf{Exercise 10.7.} Apply the five-point derivative formula of Section 10.11 to $u(t)=t^5$ at $t=0$ with $\Delta=0.1$, and compare the result with the predicted error $-\tfrac{\Delta^4}{30}u^{(5)}$.

\textbf{Exercise 10.8.} For $h=\begin{pmatrix}3&4\\4&-3\end{pmatrix}$ and $u_0=(1,0)^T$, evaluate the exact solution of $i\dot u=hu$ at $t=\pi/10$ and check its norm. Then make one explicit Euler step of size $0.01$ from $u_0$ and compute the new $u^\dagger u$.

\textbf{Exercise 10.9.} Carry out the first secant step of Section 10.12 from $E_0=5$ and $E_1=15$ with the values $F(5)=0.35184$ and $F(15)=-0.17245$.

\textbf{Exercise 10.10.} Euler's method for $y'=-y$ gave $0.316406$ with $h=0.25$ and $0.343609$ with $h=0.125$ (Section 10.3). Use Richardson's rule with $p=1$ to estimate the error of the second value and to improve it. Compare with $e^{-1}=0.367879$.

\textbf{Exercise 10.11.} Near $t=10^4$ the spacing of doubles is $2^{-39}$. Estimate how large the rounding error of CVODE's internal time can become after $6\times10^5$ additions $t_{n+1}=t_n+h_n$, and the resulting phase error of a spinor that rotates at the rate $M_{\mathrm{eff}}=2$.

\textbf{Exercise 10.12.} In EXP-5 the refined run (rtol and atol divided by 10, max\_step by 2) reduced the relative error from $3.34\times10^{-8}$ to $2.28\times10^{-9}$. What does the ratio suggest, and what ratio would a rounding floor give?

\textbf{Exercise 10.13.} Apply the regula falsi of Section 10.12 to the shooting example $F(E)=\sin\sqrt E/\sqrt E$ with the starting bracket $[5,15]$ and the values $F(5)=0.35184$, $F(15)=-0.17245$. (a) Compute the first two points with the plain regula falsi, using $F(11.7108)=-0.08090$, and say which end of the bracket is replaced each time. (b) In the Illinois variant the second step replaces the same end as the first. Halve the value kept at the other end and compute the third point, using $F(10.4562)=-0.02842$. Which end does it replace?

\subsection{Answers to the exercises}

\textbf{Answer 10.1.} With $v=x'$: $x'=v$ and $v'=x''=-5x-2v$, so $Y=(x,v)^T$ and $f(t,Y)=(v,\,-5x-2v)^T$. For $x=e^{-t}\cos2t$: $x'=-e^{-t}\cos2t-2e^{-t}\sin2t$ and $x''=e^{-t}\cos2t+2e^{-t}\sin2t+2e^{-t}\sin2t-4e^{-t}\cos2t=-3e^{-t}\cos2t+4e^{-t}\sin2t$. Then $x''+2x'+5x=e^{-t}\bigl[(-3-2+5)\cos2t+(4-4)\sin2t\bigr]=0$.

\textbf{Answer 10.2.} Each step multiplies by $1-2h$. With $h=0.25$: $(0.5)^4=0.0625$; with $h=0.125$: $(0.75)^8=0.100113$. The exact value is $e^{-2}=0.135335$, so the errors are $0.072835$ and $0.035222$. Their ratio is 2.07, close to $2^1$: order 1.

\textbf{Answer 10.3.} The factor is $R=(1+i\omega h/2)/(1-i\omega h/2)$. Numerator and denominator are complex conjugates of each other, so they have the same modulus $\sqrt{1+\omega^2h^2/4}$, and $\lvert R\rvert=1$. Hence $\lvert y_{n+1}\rvert=\lvert y_n\rvert$ at every step.

\textbf{Answer 10.4.} Explicit Euler multiplies by $1-50h$; stability needs $\lvert1-50h\rvert\le1$, that is $0\le h\le0.04$. Backward Euler multiplies by $1/(1+50h)=1/(1+5)=1/6$ at $h=0.1$, stable and decaying.

\textbf{Answer 10.5.} For $y=1$ the derivative is 0: $a+b+c=0$. For $y=t$ the derivative is 1, and the values are $0,-h,-2h$: $(0\cdot a-hb-2hc)/h=1$, that is $-b-2c=1$. For $y=t^2$ the derivative at $t=0$ is 0, and the values are $0,h^2,4h^2$: $(h^2b+4h^2c)/h=0$, that is $b=-4c$. Then $4c-2c=1$ gives $c=\tfrac12$, $b=-2$ and $a=\tfrac32$: the formula $(\tfrac32y_{n+1}-2y_n+\tfrac12y_{n-1})/h$ of Section 10.6.

\textbf{Answer 10.6.} The weights are $10^{-6}\cdot1+10^{-9}=1.001\times10^{-6}$, $10^{-6}\cdot10^{-3}+10^{-9}=2\times10^{-9}$ and $0+10^{-9}=10^{-9}$. The ratios are $0.4995$, $1.0$ and $1.0$. The norm is $\sqrt{(0.2495+1+1)/3}=\sqrt{0.7498}=0.866\le1$: the step is accepted, although two components sit exactly at their limit.

\textbf{Answer 10.7.} $u(\pm0.1)=\pm10^{-5}$ and $u(\pm0.2)=\pm3.2\times10^{-4}$. The formula gives $\bigl[-3.2\times10^{-4}-8(-10^{-5})+8(10^{-5})-3.2\times10^{-4}\bigr]/1.2=-4.8\times10^{-4}/1.2=-4\times10^{-4}$, whereas $u'(0)=0$. The predicted error is $-\tfrac{(0.1)^4}{30}\cdot120=-4\times10^{-4}$: exact agreement, because $t^5$ has no derivatives beyond the fifth.

\textbf{Answer 10.8.} At $t=\pi/10$, $5t=\pi/2$, $\cos5t=0$ and $\sin5t=1$, so $u=(-\tfrac{3i}5,-\tfrac{4i}5)^T$ and $u^\dagger u=\tfrac9{25}+\tfrac{16}{25}=1$. One Euler step of $\dot u=-ihu$ gives $u_1=(1-0.01\,ih)u_0$, and $u_1^\dagger u_1=u_0^\dagger(1+0.01\,ih)(1-0.01\,ih)u_0=u_0^\dagger(1+10^{-4}h^2)u_0=1+10^{-4}\cdot25=1.0025$: the norm has grown by $(\Delta t)^2E^2$ in one step, the explicit-Euler growth of Section 10.4.

\textbf{Answer 10.9.} $E_2=15-(-0.17245)\cdot\dfrac{15-5}{-0.17245-0.35184}=15-\dfrac{1.7245}{0.52429}=15-3.2892=11.7108$, the first secant value quoted in Section 10.12.

\textbf{Answer 10.10.} The difference is $0.343609-0.316406=0.027203$. With $p=1$, $2^p-1=1$, so the estimated error of the second value is $0.0272$ (the true error is $0.0243$), and the improved value is $0.343609+0.027203=0.370812$, whose error $0.0029$ is about eight times smaller.

\textbf{Answer 10.11.} Each rounded addition is off by at most half the spacing, $2^{-40}\approx9.1\times10^{-13}$. After $6\times10^5$ additions the time can be off by up to about $5.5\times10^{-7}$ (if all roundings go the same way), and a phase that grows at the rate 2 can then be off by about $1.1\times10^{-6}$. This is the size of the bound $1.32\times10^{-6}$ recorded for EXP-2 (Section 10.10), which is computed with the exact step counts and masses of the runs; the measured phase error, $2.24\times10^{-7}$, is below it.

\textbf{Answer 10.12.} The ratio is $3.34\times10^{-8}/2.28\times10^{-9}=14.6$ for a tenfold smaller tolerance: the error decreases at least in proportion to the tolerance, as a truncation error should, so the canonical run's error is dominated by truncation. A rounding floor would give a ratio close to 1, as the EXP-2 phase error did ($2.24\times10^{-7}$ against $2.71\times10^{-7}$).

\textbf{Answer 10.13.} (a) The point between the ends $E_{\mathrm{lo}}$ and $E_{\mathrm{hi}}$ of a bracket is the zero of the straight line through the two ends, $E=\bigl(E_{\mathrm{lo}}F(E_{\mathrm{hi}})-E_{\mathrm{hi}}F(E_{\mathrm{lo}})\bigr)/\bigl(F(E_{\mathrm{hi}})-F(E_{\mathrm{lo}})\bigr)$, which is the secant formula of Section 10.12 written for the two ends. The first point is therefore the first secant value, $11.7108$ (Answer 10.9). Since $F(11.7108)=-0.08090$ has the sign of $F(15)$, the sign change lies in $[5,11.7108]$: the right end is replaced. The second point is $\bigl(5\cdot(-0.08090)-11.7108\cdot0.35184\bigr)/(-0.08090-0.35184)=(-0.40450-4.12033)/(-0.43274)=10.4562$. Again $F(10.4562)=-0.02842<0$, so the right end is replaced again and the left end $5$ stays. The plain regula falsi goes on in this way: carrying all digits, its next points are $10.0485$, $9.9234$, $9.8857$, $9.8744$ and $9.8710$, all to the right of the root and all with the left end $5$, so that after seven points the error is still $1.4\times10^{-3}$. (b) The second step replaced the right end for the second time in a row, so the value kept at the left end is halved, $0.35184/2=0.17592$. The third point is $\bigl(5\cdot(-0.02842)-10.4562\cdot0.17592\bigr)/(-0.02842-0.17592)=(-0.14210-1.83945)/(-0.20434)=9.697$ ($9.69746$ with all digits). There $F$ is about $+0.009>0$, the sign of $F$ at the left end, so this time the \textbf{left} end is replaced and the bracket becomes $[9.697,10.4562]$. Carrying all digits, the next points are $9.87743$, $9.86971$, $9.86951$ and $9.8696044$: after seven points the value agrees with $\pi^2=9.8696044$ in all the digits shown, like the secant rule of Section 10.12, but with the root kept inside a bracket at every step.

\section{The five dark-sector experiments}

\subsection{What this chapter answers}

Astronomers infer two things about the universe that ordinary matter does not explain. First, galaxies and clusters of galaxies move and bend light as if they contained much more gravitating mass than the stars and gas that can be seen; the missing part is called \textbf{dark matter}. It clusters like ordinary matter and has almost no pressure. Second, the light of distant exploding stars (supernovae) shows that the expansion of the universe is speeding up; whatever causes this is called \textbf{dark energy}. It must have a large negative pressure. This chapter asks whether the field dirac16complex of the earlier chapters can play either role, and it answers with numbers, not words.

Stage 3 of the project posed the question in these words: ``in detail, investigate, determine and discuss, using numerical solutions and graphs, the physical behaviour and changing nature of pressure and energy density in this framework; is there any connection to dark matter and/or dark energy?'' To answer it, five \textbf{experiments} (numbered EXP-1 to EXP-5) were run. Each one is a system of ordinary differential equations that a computer program integrates, as described in Chapter 10.

\begingroup
\small
\begin{longtable}{@{}>{\raggedright\arraybackslash}p{0.267\linewidth}>{\raggedright\arraybackslash}p{0.267\linewidth}>{\raggedright\arraybackslash}p{0.267\linewidth}@{}}
\toprule
\textbf{Experiment} & \textbf{Background (the gravitational field)} & \textbf{Question} \\
\midrule
\endfirsthead
\toprule
\textbf{Experiment} & \textbf{Background (the gravitational field)} & \textbf{Question} \\
\midrule
\endhead
EXP-1 & the primordial pair-creation field of the author's notebook & what do energy density and pressure do there, and what source would the field need? \\
EXP-2 & a homogeneous 8-dimensional universe solved together with Einstein's equations & does a self-gravitating condensate make the universe isotropic, and what would a 3-space observer infer? \\
EXP-3 & a 4-dimensional late universe with frozen extra dimensions & can the condensate be the dark energy of the quoted supernova fits? \\
EXP-4 & an expanding 3-space: a radiation era, and inflation followed by radiation & do the quanta behave like dark matter, and how many does the expansion create? \\
EXP-5 & deflating extra times & why must everything be restricted to modes without momentum along the extra times? \\
\bottomrule
\end{longtable}
\endgroup

The answer, which Sections 11.13 and 11.14 argue in full:

\begin{itemize}
\item \textbf{Dark matter: a qualified yes.} In the sector without momentum along the extra times, and with the extra dimensions held static by assumption, a gas of dirac16complex quanta turns from radiation-like into dust-like, and the expansion of the universe creates such quanta for every mass $m>0$. That is a dark-matter-like equation of state and a production mechanism. It is not a dark-matter model: the abundance, the darkness, the stabilisation of the extra dimensions and the growth of structure were neither computed nor derived.
\item \textbf{Dark energy: no, within everything computed.} Negative pressure arises only from an attractive self-interaction. Tuned to today's quoted value $w_0=-0.861$ it evolves about eight times too fast, and a short distance into the past (redshift $z=0.026$) it leaves the regime in which its approximation is valid.
\end{itemize}

Every number in this chapter is copied from a committed output file of the repository. The files are the summary.json and python-check-report.json of each experiment under \texttt{artifact\allowbreak{}s/dirac1\allowbreak{}6complex\allowbreak{}/numeric\allowbreak{}s/exp1} to \texttt{\detokenize{exp5}}, the fits file \texttt{artifact\allowbreak{}s/dirac1\allowbreak{}6complex\allowbreak{}/numeric\allowbreak{}s/exp3/f\allowbreak{}its.json}, and the collected \texttt{artifact\allowbreak{}s/dirac1\allowbreak{}6complex\allowbreak{}/numeric\allowbreak{}s/numeri\allowbreak{}cs-summa\allowbreak{}ry.json}. The scientific document of Stage 3 is \texttt{provenan\allowbreak{}ce/DIRAC\allowbreak{}16COMPLE\allowbreak{}X\_DARK\_S\allowbreak{}ECTOR\_NU\allowbreak{}MERICS.m\allowbreak{}d} (below: the Stage-3 document), and its companion \texttt{provenan\allowbreak{}ce/DIRAC\allowbreak{}16COMPLE\allowbreak{}X\_STUDEN\allowbreak{}T\_GUIDE.\allowbreak{}md} shows how to rerun everything. When we write ``Stage-3 document §4.6'' we mean its numbered part 4.6; ``Section 11.4'' always means a section of this book. Where a statement below is derived by hand rather than checked by a program, the text says so.

\subsection{The words and numbers of cosmology}

This section defines the handful of cosmological notions that the experiments use. Nothing here is specific to dirac16complex.

\textbf{Scale factor.} A universe that looks the same at every point (homogeneous) and in every direction (isotropic) can only change by a common stretching of all distances. The stretching factor $a(t)$ is the \textbf{scale factor}: two distant galaxies at rest in the expanding space are a distance proportional to $a(t)$ apart. By convention $a=1$ today.

\textbf{Hubble rate.} The relative growth rate $H=\dot a/a$ (a dot means $d/dt$) is the \textbf{Hubble rate}; $H_0$ is its value today. If $H$ is constant, $a=e^{Ht}$ grows exponentially; this is called de Sitter expansion, and a phase of it in the early universe is called inflation.

\textbf{Redshift.} Light that left a galaxy when the scale factor was $a$ arrives today with its wavelength stretched by the factor $1/a$. Astronomers call $z$ with $1+z=1/a$ the \textbf{redshift}. For example $a=0.5$ gives $z=1$, and $z=0.026$ gives $a=1/1.026=0.975$: a very recent epoch.

\textbf{Comoving quantities.} Coordinates that expand with the universe, in which galaxies at rest in the expanding space keep fixed coordinate values, are called \textbf{comoving}, and so are the quantities measured in them. A comoving distance or volume stays fixed as $a$ grows; the physical one is $a$ (or $a^3$) times it. A comoving momentum $k$ corresponds to the physical momentum $K=k/a$, and a comoving number density $na^3$ (the number per unit comoving volume) stays fixed as long as no particles are created or destroyed. A comoving temperature describes the occupations of the states as a function of the comoving momentum $k$; at $a=1$, where $K=k$, it is the physical temperature.

\textbf{Units and the Hubble length.} A parsec (pc) is $3.086\times10^{16}$ m and 1 Mpc $=10^6$ pc, so $H_0=67.4$ km/s/Mpc means that galaxies 1 Mpc apart recede from each other at 67.4 km/s. The \textbf{Hubble length} $c/H$ is the distance light travels in one expansion time $1/H$. A wave is \textbf{inside the horizon} when its physical wavelength is much shorter than $c/H$ ($K\gg H$ in units with $c=1$), and it is stretched beyond the Hubble length when the expansion makes its wavelength longer than $c/H$.

\textbf{A fluid and its equation of state.} On large scales matter is described as a fluid with an \textbf{energy density} $\rho$ (energy per volume) and a \textbf{pressure} $p$ (the flux of momentum through a surface). Their ratio

\[
w=\frac{p}{\rho}
\]

is the \textbf{equation-of-state parameter}. Three standard cases: dust (slow particles, including ordinary and dark matter) has $w=0$; radiation (particles moving at the speed of light) has $w=1/3$; a cosmological constant has $w=-1$.

\textbf{How a fluid dilutes.} In an expanding 3-space, energy conservation reads $\dot\rho=-3H(\rho+p)$. (The energy $\rho a^3$ in a unit comoving volume, whose physical volume is $a^3$, changes by the work of the pressure, $d(\rho a^3)=-p\,d(a^3)$, that is $\dot\rho a^3+3\rho a^2\dot a=-3pa^2\dot a$; divide by $a^3$ and use $\dot a/a=H$.) For constant $w$ divide by $\rho$: $d\ln\rho/dt=-3(1+w)\,d\ln a/dt$, so

\[
\rho\propto a^{-3(1+w)} .
\]

Dust dilutes as $a^{-3}$ (a fixed number of particles in a volume that grows as $a^3$), radiation as $a^{-4}$ (the same, and each photon's energy also falls as $1/a$), and a cosmological constant not at all. Worked example: for $w=-0.861$ the exponent is $-3(0.139)=-0.417$, so at $a=0.5$ the density was $2^{0.417}=1.335$ times today's value.

\textbf{Friedmann equations.} Einstein's equations (Chapter 4) for such a universe with 3 space dimensions reduce to

\[
H^2=\frac{\kappa_4}{3}\,\rho_{\mathrm{tot}},\qquad \frac{\ddot a}{a}=-\frac{\kappa_4}{6}\,(\rho_{\mathrm{tot}}+3p_{\mathrm{tot}}),
\]

with $\kappa_4=8\pi G$. Section 11.7 shows that these are the $D=4$ case of the equations that EXP-2 solves in $D=8$. The second equation says that the expansion \textbf{accelerates} ($\ddot a>0$) only if $\rho+3p<0$; for a single component with $\rho>0$ that means $w<-1/3$. This is why dark energy needs negative pressure.

\textbf{Density parameters and $E(a)$.} The critical density is $\rho_{\mathrm{crit}}=3H_0^2/\kappa_4$, and $\Omega_i=\rho_i(a=1)/\rho_{\mathrm{crit}}$ is the density parameter of component $i$. Dividing the first Friedmann equation by $H_0^2$ gives $E^2:=H^2/H_0^2=\sum_i\rho_i(a)/\rho_{\mathrm{crit}}$; for radiation, matter and a component $\psi$ this is $E^2=\Omega_ra^{-4}+\Omega_ma^{-3}+\rho_\psi(a)/\rho_{\mathrm{crit}}$, and $E(1)=1$ requires $\sum_i\Omega_i=1$ (a spatially flat universe).

\textbf{Deceleration parameter.} $q_{\mathrm{dec}}=-a\ddot a/\dot a^2$. From the two Friedmann equations, $q_{\mathrm{dec}}=(\rho_{\mathrm{tot}}+3p_{\mathrm{tot}})/(2\rho_{\mathrm{tot}})$: positive when the expansion slows down, negative when it speeds up.

\textbf{Phantom.} A fluid with $\rho>0$ and $w<-1$ (equivalently $\rho+p<0$) is called \textbf{phantom}. The condition $\rho+p\ge0$ is the null energy condition. A canonical scalar field, the usual model of dark energy, has $\rho+p=\dot\phi^2\ge0$ and can never be phantom.

\textbf{The CPL form and the quoted supernova numbers.} A slowly changing $w$ is often summarised by the Chevallier-Polarski-Linder (CPL) form

\[
w(a)=w_0+w_a\,(1-a),\qquad \frac{dw}{da}=-w_a .
\]

The reference document of Stage 3 is an e-mail about the equation of state (a PDF in the author's working folder, a private input that is not committed). It quotes, for the ``Unite'' supernova compilation, the CPL fit $w_0=-0.861$, $w_a=-0.60$, and a constant-$w$ fit $w=-0.764$, ``roughly two standard deviations'' from $-1$ (a standard deviation, or one sigma, measures the statistical uncertainty of a fitted number). It gives no error bars on $(w_0,w_a)$ and cites no primary publication, so the project did not check these values against one; this chapter uses them only as quoted (Stage-3 document §3.3). Since $a$ grows with time, a field whose $w$ rises from $-1$ (called thawing) has $dw/da>0$, that is $w_a<0$. The e-mail's table writes the opposite signs; this book follows the formula: \textbf{thawing means $w_a<0$, freezing means $w_a>0$.}

\textbf{Distances.} Light from redshift $z$ has travelled the comoving distance $D_C(z)=\int_0^z c\,dz'/H(z')$. In a flat universe the luminosity distance is $d_L=(1+z)D_C$. (Light moves on $ds^2=0$; in a flat universe with the metric $ds^2=-c^2dt^2+a^2\,d\chi^2$ along its path, with $\chi$ the comoving distance, this gives $c\,dt=a\,d\chi$. With $dt=da/(aH)$, $D_C=\int c\,dt/a=\int_a^1c\,da'/(a'^2H)=\int_0^zc\,dz'/H$, since $z=1/a-1$ gives $dz=-da/a^2$. Today, with $a=1$, the light emitted by a source of luminosity $L$ (emitted power) is spread over a sphere of area $4\pi D_C^2$, and the flux $F$ (received power per area) is lowered by one factor $1+z$ from the redshift of each photon's energy and by another from the slower arrival rate of the photons: $F=L/\bigl(4\pi D_C^2(1+z)^2\bigr)$. Defining $d_L$ by $F=L/(4\pi d_L^2)$ gives $d_L=(1+z)D_C$.) Astronomers quote the \textbf{distance modulus} $\mathrm{DM}=5\log_{10}(d_L/10\,\mathrm{pc})$. DM is measured in magnitudes (mag), the logarithmic brightness unit of astronomy. A constant shift of DM is not measured independently (it depends on $H_0$ and on the true brightness of the supernovae), so fits ``profile'' it: they subtract the best constant offset.

\textbf{The standard model of cosmology}, called $\Lambda$CDM, contains radiation, ordinary matter, cold dark matter (CDM) and a cosmological constant $\Lambda$ ($w=-1$); it is the reference against which a varying $w$ is compared.

\textbf{What each role needs.} A dark-matter candidate must have $w\approx0$ at late times, the right abundance, (almost) no interaction with light, a long lifetime, and small enough random velocities to let structures form. A dark-energy candidate must have $\rho>0$ and $w$ near $-1$ today, stable small fluctuations, and a history in which matter dominated earlier. Sections 11.13 and 11.14 go through these lists.

\subsection{The field in a homogeneous background: the mode equation}

We now recall the field and put it into the simplest gravitational fields. Chapters 2, 4, 6 and 7 derive everything we quote here.

\textbf{The field.} dirac16complex is a column $\Psi=(\Psi_0,\dots,\Psi_{15})$ of 16 complex Grassmann-odd components on an 8-dimensional spacetime with coordinates $x_0,\dots,x_7$. The flat metric is $\eta=\mathrm{diag}(+1,+1,+1,+1,-1,-1,-1,-1)$: directions 0 to 3 are space-like, directions 4 to 7 time-like. $x_4$ is the time $t$; $x_0$ is a hidden space direction, $x_1,x_2,x_3$ are ordinary 3-space, and $x_5,x_6,x_7$ are three extra times. We write $T=\{0,1,2,3,5,6,7\}$ for the seven directions other than $x_4$. The gamma matrices $\gamma^0,\dots,\gamma^7$ are real 16 by 16 matrices with $\gamma^a\gamma^b+\gamma^b\gamma^a=2\eta^{ab}$; $\gamma^0,\dots,\gamma^3$ are symmetric and $\gamma^4,\dots,\gamma^7$ antisymmetric (Chapter 2). The charge matrix is $C=\gamma^0\gamma^1\gamma^2\gamma^3$, the adjoint is $\bar\Psi=\Psi^\dagger C$, and $S=\bar\Psi\Psi$ is the scalar density.

\textbf{The Lagrangian and the field equation} (Chapters 6 and 7):

\[
\mathcal L=\sqrt{|g|}\Bigl[\tfrac12\bigl(\bar\Psi\gamma^\mu D_\mu\Psi-(D_\mu\bar\Psi)\gamma^\mu\Psi\bigr)-m\,S-U(S)\Bigr],\qquad U(S)=\tfrac{\lambda}{2}S^2,
\]

\[
\gamma^\mu D_\mu\Psi=M_{\mathrm{eff}}\,\Psi,\qquad M_{\mathrm{eff}}=m+U'(S)=m+\lambda S .
\]

$m$ is the mass and $\lambda$ the strength of a four-fermion contact interaction: $\lambda<0$ is attractive, $\lambda>0$ repulsive. Because $S$ is Grassmann-even of degree 2 and the algebra at a point has 32 odd generators, $S^{17}=0$, so every function $U(S)$ is a polynomial of degree at most 16 (Section 5.9 and Chapter 6). The model requires $U(0)=U'(0)=0$ (Chapter 6; Stage-1 document §7.6): no second mass term and no constant term, so \textbf{this Lagrangian contains no cosmological constant, by choice.} The Grassmann algebra itself would allow a constant term $-\sqrt{|g|}\,U(0)$, which would act as a cosmological constant. $M_{\mathrm{eff}}$ is the effective mass: the interaction shifts the mass in proportion to $S$.

\textbf{The backgrounds.} EXP-2 to EXP-5 use a homogeneous, diagonal metric with lapse 1:

\[
ds^2=-dt^2+\sum_{i\in T}\epsilon_i\,h_i(t)^2\,dx_i^2,\qquad \epsilon_i=+1\ (i\le3),\qquad \epsilon_i=-1\ (i\ge5).
\]

EXP-1 uses the warped primordial field of Section 11.6, which has the same form along $x_1,\dots,x_3,x_5,\dots,x_7$ but whose scale factors also depend on the hidden coordinate $\zeta$; Section 11.6 reduces it to the same mode equation.

$h_i(t)$ is the scale factor of direction $i$. Its Hubble rate is $H_i=\dot h_i/h_i$, the total expansion rate is $\Theta=\sum_{i\in T}H_i$, and the 7-volume is $V=\prod_{i\in T}h_i$. Taking the logarithm, $\ln V=\sum_i\ln h_i$, and differentiating gives $\dot V/V=\Theta$. The experiments group the directions: $b=h_0$ (hidden space), $a=h_1=h_2=h_3$ (3-space) and $c=h_5=h_6=h_7$ (extra times), so that $\Theta=H_b+3H_a+3H_c$ and $V=b\,a^3c^3$.

\textbf{Vielbein and connection.} The vielbein (Chapter 4) is $e^4=dt$ and $e^i=h_i\,dx_i$, so that $ds^2=\eta_{ab}e^ae^b$. The gamma matrices with a curved index are $\gamma^{x_4}=\gamma^4$ and $\gamma^{x_i}=\gamma^i/h_i$. For a diagonal vielbein $e^b=h_b\,dx_b$ Chapter 4 (and Stage-1 document §8.5) gives

\[
\gamma^\mu\Omega_\mu=\frac12\sum_{b=0}^{7}\frac{1}{h_b}\,\partial_b\ln\Bigl(\prod_{c\ne b}h_c\Bigr)\,\gamma^b .
\]

Here everything depends on $t=x_4$ only, so only $b=4$ contributes, with $h_4=1$ and $\prod_{c\ne4}h_c=V$:

\[
\gamma^\mu\Omega_\mu=\tfrac12\,\frac{\dot V}{V}\,\gamma^4=\tfrac12\,\Theta\,\gamma^4 .
\]

\textbf{The mode ansatz.} Take a plane wave with coordinate momentum $k_j$ along the directions $j\in T$,

\[
\Psi=V^{-1/2}\,e^{i\sum_jk_jx_j}\,u(t),
\]

where $u(t)$ is an ordinary column of 16 complex numbers (the \textbf{mode amplitude}). The Grassmann nature of the field enters only through the counting of states and through the expectation-value rule of Section 11.4. Insert the ansatz into $\gamma^\mu D_\mu\Psi=\gamma^{x_4}\partial_4\Psi+\sum_j\gamma^{x_j}\partial_j\Psi+\gamma^\mu\Omega_\mu\Psi$, one term at a time, and drop the common factor $V^{-1/2}e^{ik\cdot x}$:

\begin{enumerate}
\item time derivative: $\partial_4(V^{-1/2}u)=V^{-1/2}\bigl(\dot u-\tfrac12\Theta u\bigr)$, which gives $\gamma^4\dot u-\tfrac12\Theta\gamma^4u$;
\item space derivatives: $\partial_j\Psi=ik_j\Psi$ and $\gamma^{x_j}=\gamma^j/h_j$ give $i\sum_jK_j\gamma^ju$, where $K_j:=k_j/h_j$ is the \textbf{physical momentum} along $j$;
\item connection: $+\tfrac12\Theta\gamma^4u$.
\end{enumerate}

The two $\Theta$ terms cancel exactly; that is why the factor $V^{-1/2}$ was put in front. What remains is $\gamma^4\dot u+i\sum_jK_j\gamma^ju=M_{\mathrm{eff}}u$. Multiply from the left by $\gamma^4$ and use $(\gamma^4)^2=-1$: $-\dot u+i\sum_jK_j\gamma^4\gamma^ju=M_{\mathrm{eff}}\gamma^4u$. Multiply by $-i$ and rearrange:

\[
i\,\dot u=h(t)\,u,\qquad h(t)=-iM_{\mathrm{eff}}\gamma^4-\sum_{j\in T}K_j\,\gamma^4\gamma^j .
\]

This looks like a Schrödinger equation with a 16 by 16 \textbf{mode Hamiltonian} $h(t)$.

\textbf{The square of $h$.} Write $A=-iM_{\mathrm{eff}}\gamma^4$ and $P_j=-K_j\gamma^4\gamma^j$. Three small computations with the Clifford relation, using $\gamma^j\gamma^4=-\gamma^4\gamma^j$ for $j\ne4$ and hence $\gamma^4\gamma^j\gamma^4=-\gamma^4\gamma^4\gamma^j=\gamma^j$:

\begin{itemize}
\item $A^2=-M_{\mathrm{eff}}^2(\gamma^4)^2=M_{\mathrm{eff}}^2$;
\item $AP_j+P_jA=iM_{\mathrm{eff}}K_j\bigl(\gamma^4\gamma^4\gamma^j+\gamma^4\gamma^j\gamma^4\bigr)=iM_{\mathrm{eff}}K_j(-\gamma^j+\gamma^j)=0$;
\item $P_jP_l+P_lP_j=K_jK_l\bigl(\gamma^4\gamma^j\gamma^4\gamma^l+\gamma^4\gamma^l\gamma^4\gamma^j\bigr)=K_jK_l(\gamma^j\gamma^l+\gamma^l\gamma^j)=2\eta^{jl}K_jK_l$.
\end{itemize}

Expanding $h^2=(A+\sum_jP_j)^2$ and collecting terms:

\[
h^2=E^2\cdot1,\qquad E^2=M_{\mathrm{eff}}^2+\sum_{j=0}^{3}K_j^2-\sum_{j=5}^{7}K_j^2 .
\]

So the only possible eigenvalues of $h$ are $+E$ and $-E$. Because $h$ has zero trace (the diagonal of the antisymmetric $\gamma^4$ is zero, and $\mathrm{tr}(\gamma^4\gamma^j)=-\mathrm{tr}(\gamma^j\gamma^4)=-\mathrm{tr}(\gamma^4\gamma^j)$ forces $\mathrm{tr}(\gamma^4\gamma^j)=0$), each occurs eight times when $E\ne0$.

\textbf{When is $h$ Hermitian?} A real symmetric matrix is Hermitian, a real antisymmetric one satisfies $M^\dagger=-M$. Since $\gamma^4$ is real antisymmetric, $(-iM_{\mathrm{eff}}\gamma^4)^\dagger=iM_{\mathrm{eff}}(-\gamma^4)=-iM_{\mathrm{eff}}\gamma^4$: Hermitian. For $j\le3$, $(\gamma^4\gamma^j)^\dagger=(\gamma^j)^T(\gamma^4)^T=\gamma^j(-\gamma^4)=\gamma^4\gamma^j$: Hermitian. For $j\ge5$, $(\gamma^4\gamma^j)^\dagger=(-\gamma^j)(-\gamma^4)=\gamma^j\gamma^4=-\gamma^4\gamma^j$: anti-Hermitian. Hence

\[
h^\dagger=h\quad\Longleftrightarrow\quad K_5=K_6=K_7=0 .
\]

The modes without momentum along the extra times form the \textbf{good sector}. There $E^2>0$, and the \textbf{Hilbert norm} $u^\dagger u$ is conserved: $\frac{d}{dt}(u^\dagger u)=\dot u^\dagger u+u^\dagger\dot u=(-ihu)^\dagger u+u^\dagger(-ihu)=i\,u^\dagger(h^\dagger-h)\,u=0$.

\textbf{The Krein norm.} The matrix $B=-iC\gamma^4$ is Hermitian with $B^2=1$ and has eigenvalues $+1$ and $-1$, eight each (Chapters 2 and 8). $C$ contains $\gamma^0,\dots,\gamma^3$ once each; moving $\gamma^a$ with $a\le3$ through $C$ meets three gammas that anticommute with it, so $C\gamma^a=-\gamma^aC$ for $a\le3$, while for $a\ge4$ all four gammas of $C$ anticommute with $\gamma^a$, the four signs cancel, and $C\gamma^a=\gamma^aC$. It follows that $\gamma^4$ and $\gamma^j$ ($j\le3$) commute with $B$ and $\gamma^j$ ($j\ge5$) anticommutes with it; for example, for $j\ge5$, $\gamma^jB=-i\gamma^jC\gamma^4=-iC\gamma^j\gamma^4=iC\gamma^4\gamma^j=-B\gamma^j$. Now check $h^\dagger B=Bh$ term by term. The term $-iM_{\mathrm{eff}}\gamma^4$ is Hermitian and commutes with $B$. For $j\le3$ the term $-K_j\gamma^4\gamma^j$ is Hermitian and commutes with $B$. For $j\ge5$ it is anti-Hermitian and anticommutes with $B$, so $(\gamma^4\gamma^j)^\dagger B=-\gamma^4\gamma^jB=B\gamma^4\gamma^j$. Hence $h^\dagger B=Bh$ for every momentum. Then $\frac{d}{dt}(u^\dagger Bu)=i\,u^\dagger(h^\dagger B-Bh)\,u=0$: the \textbf{Krein norm} $u^\dagger Bu$, which can be positive, negative or zero, is conserved always, in and out of the good sector.

\textbf{What the computer integrates.} A computer program integrates real numbers. Write $u=x+iy$ and $h=h_r+ih_i$ with real columns $x,y$ and real matrices $h_r,h_i$. The real and imaginary parts of $i(\dot x+i\dot y)=(h_r+ih_i)(x+iy)$ give

\[
\dot x=h_r\,y+h_i\,x,\qquad \dot y=h_i\,y-h_r\,x .
\]

The state is 32 real numbers: $\mathrm{Re}\,u_0,\dots,\mathrm{Re}\,u_{15}$ and then $\mathrm{Im}\,u_0,\dots,\mathrm{Im}\,u_{15}$, the columns \texttt{\detokenize{u_re_0}} to \texttt{\detokenize{u_re_15}} and \texttt{\detokenize{u_im_0}} to \texttt{\detokenize{u_im_15}} of the output files. For $h=-iM\gamma^4-\sum_jK_j\gamma^4\gamma^j$ the pieces are $h_i=-M\gamma^4$ and $h_r=-\sum_jK_j\gamma^4\gamma^j$.

\subsection{Energy density, pressure and the expectation-value rule}

\textbf{The energy-momentum tensor.} Chapter 7 derives the energy-momentum tensor of the field:

\[
T_{\mu\nu}=-\tfrac14\Bigl[\bar\Psi\gamma_\mu D_\nu\Psi+\bar\Psi\gamma_\nu D_\mu\Psi-(D_\mu\bar\Psi)\gamma_\nu\Psi-(D_\nu\bar\Psi)\gamma_\mu\Psi\Bigr]+g_{\mu\nu}\,\mathcal L_s,\qquad \mathcal L_s=K-mS-U(S),
\]

where $K=\tfrac12\bigl(\bar\Psi\gamma^\mu D_\mu\Psi-(D_\mu\bar\Psi)\gamma^\mu\Psi\bigr)$ is the kinetic term. An observer at rest in the coordinates (whose 8-velocity is $\partial_4$, a unit time-like vector because $g_{44}=-1$) measures the energy density $\rho=T_{44}$ and, in each of the seven directions $i\in T$, the pressure $p_{(i)}=T^i{}_i$ (no sum over $i$). Four of these directions are space-like and three are time-like. When the pressures differ one uses the mean $\bar p=\tfrac17\sum_{i\in T}p_{(i)}$ and $\bar w=\bar p/\rho$.

\textbf{The homogeneous condensate.} Contracting the field equation with $\bar\Psi$, and its adjoint with $\Psi$, gives on shell $K=M_{\mathrm{eff}}S$, hence $\mathcal L_s=M_{\mathrm{eff}}S-mS-U=SU'(S)-U(S)$. For a state that depends on $t$ only (a homogeneous condensate) in the diagonal backgrounds of Section 11.3, Stage 1 found that the kinetic bracket contributes nothing to the seven transverse diagonal components, so $T^i{}_i=\mathcal L_s$ for every $i\in T$. This is stated in Stage-1 document §9.7 and checked there in exact arithmetic at three test times of a homogeneous diagonal frame (checks \texttt{EMT\_homo\allowbreak{}geneousR\allowbreak{}eduction\allowbreak{}\_G3} of the Wolfram geometry verifier and \texttt{EMT\_homo\allowbreak{}geneousR\allowbreak{}eduction} of the Python geometry checker, reports under \texttt{artifact\allowbreak{}s/dirac1\allowbreak{}6complex\allowbreak{}/arbitra\allowbreak{}ry-field\allowbreak{}/}); we take it from there. The trace of $T$ supplies the time-time component: $g^{\mu\nu}\gamma_\mu=\gamma^\nu$ turns the bracket into $-K$ and $g^{\mu\nu}g_{\mu\nu}=8$, so $T^\mu{}_\mu=-K+8\mathcal L_s$. Subtracting the seven transverse components, $T^4{}_4=-K+\mathcal L_s$, and with $T_{44}=g_{44}T^4{}_4=-T^4{}_4$,

\[
\rho=K-\mathcal L_s=M_{\mathrm{eff}}S-(SU'-U)=mS+U,\qquad p=\mathcal L_s=SU'-U .
\]

For $U=\tfrac\lambda2S^2$:

\[
\rho=mS+\tfrac\lambda2S^2,\qquad p=\tfrac\lambda2S^2,\qquad w=\frac{\lambda S}{2m+\lambda S} .
\]

So a free condensate ($\lambda=0$) is dust, a repulsive one ($\lambda>0$) has positive pressure, and an attractive one ($\lambda<0$) has negative pressure. The pressure is the same in all seven directions although the scale factors differ.

\textbf{Kinetic and potential energy.} A canonical scalar field has $\rho=\tfrac12\dot\phi^2+V$ and $p=\tfrac12\dot\phi^2-V$. Stage 1 defines two exact analogues for dirac16complex (Chapter 7). Split the kinetic term into its time part and the rest, $K=K_4+K_\perp$, with $K_4=\tfrac12\bigl(\bar\Psi\gamma^{x_4}D_4\Psi-(D_4\bar\Psi)\gamma^{x_4}\Psi\bigr)$.

\begin{itemize}
\item \textbf{(A) Lagrangian split:} $\mathrm{KE}_L=\tfrac12K_4$ and $\mathrm{PE}_L=\rho-\mathrm{KE}_L$.
\item \textbf{(B) Hamiltonian split:} $\mathrm{KE}_H=-K_\perp$ (momentum or gradient energy) and $\mathrm{PE}_H=mS+U$ (rest mass and interaction). $\rho=\mathrm{KE}_H+\mathrm{PE}_H$ holds identically.
\end{itemize}

For the condensate $K_\perp=0$: it has no momentum, and the connection terms cancel in the antisymmetrised combination; equivalently, $\mathrm{KE}_H=\rho-\mathrm{PE}_H=0$ because $\rho=mS+U$ above. So $K_4=K=M_{\mathrm{eff}}S$ and

\[
\mathrm{KE}_L=\tfrac12S\,M_{\mathrm{eff}},\qquad \mathrm{PE}_L=mS+\tfrac\lambda2S^2-\tfrac12(m+\lambda S)S=\tfrac12mS,\qquad \mathrm{KE}_H=0,\qquad \mathrm{PE}_H=\rho .
\]

Then $\mathrm{KE}_L+\mathrm{PE}_L=\rho$ and $\mathrm{KE}_L-\mathrm{PE}_L=\tfrac12\lambda S^2=p$, exactly as for the scalar field, and

\[
\rho+p=2\,\mathrm{KE}_L=S\,M_{\mathrm{eff}} .
\]

\textbf{Phantom criterion.} For $\rho>0$, $w<-1$ holds exactly when $\rho+p<0$, that is when $\mathrm{KE}_L<0$, that is when $S\,M_{\mathrm{eff}}<0$. Unlike $\tfrac12\dot\phi^2$, $\mathrm{KE}_L$ has no fixed sign. Worked example with $m=1$ and $S>0$: for $\lambda S=-1.5$ we get $M_{\mathrm{eff}}=-0.5$, $\rho=S(1-0.75)=0.25\,S>0$, $p=-0.75\,S$ and $w=-1.5/0.5=-3$, a phantom value; for $\lambda S=-0.5$ we get $M_{\mathrm{eff}}=0.5$ and $w=-0.5/1.5=-1/3$, not phantom. Whether the phantom case describes a physical state is the question of the mean-field limits below.

\textbf{Per-mode quantities and the expectation-value rule.} Chapter 8 quantises the field with respect to $t=x_4$. The anticommutator is $\{\Psi,\Psi^\dagger\}=B\,\delta^7/\sqrt{|g|}$, and because $B$ has eigenvalues of both signs the state space carries an indefinite inner product (a Krein space). In the good sector there is nevertheless an ordinary Fock space: particles and antiparticles of energy $\sqrt{m^2+K^2}$ above a filled \textbf{Dirac sea} of negative-energy states. Energies are measured relative to the filled sea (\textbf{normal ordering}). In this positive-norm quantisation a one-particle state built on a normalised positive-energy mode $u$ ($u^\dagger u=1$) has the \textbf{expectation-value rule}

\[
\langle\Psi^\dagger M\Psi\rangle=u^\dagger BM\,u
\]

for any 16 by 16 matrix $M$ (Stage-3 document §4.6). Since $C$ commutes with $\gamma^4$ and $C^2=1$, $BC=-iC\gamma^4C=-i\gamma^4$. Three quantities per mode follow:

\[
s(u)=u^\dagger BC\,u=u^\dagger(-i\gamma^4)\,u,\qquad \varepsilon(u)=u^\dagger h\,u,\qquad p_j(u)=-K_j\,u^\dagger\gamma^4\gamma^ju\quad(\text{no sum}),
\]

the scalar density, the energy and the pressure along $j$ of one quantum.

\textbf{Why $\varepsilon$ and $p_j$ are the energy and the pressure.} The formula for $T_{\mu\nu}$ above, with $\gamma_4=g_{44}\gamma^{x_4}=-\gamma^4$ and $g^{jj}\gamma_j=\gamma^{x_j}$, gives

\[
T_{44}=K_4-\mathcal L_s,\qquad T^j{}_j=-\tfrac12\bigl(\bar\Psi\gamma^{x_j}D_j\Psi-(D_j\bar\Psi)\gamma^{x_j}\Psi\bigr)+\mathcal L_s\quad(\text{no sum}),
\]

and for $\lambda=0$ the on-shell $\mathcal L_s=SU'-U$ vanishes. Take $\Psi=V^{-1/2}e^{ik\cdot x}u$ in a good-sector background of Section 11.3. The real factor $V^{-1/2}$ drops out of these antisymmetric combinations, because its derivative gives $-\tfrac12\Theta$ times the same bilinear in both terms. The connection terms drop out too. By Section 4.12, $\Omega_\mu$ is the sum over $b\ne\mu$ of $\omega_{\mu\mu b}S^{\mu b}$ with $S^{\mu b}=\tfrac12\gamma^\mu\gamma^b$ and $\omega_{\mu\mu b}$ proportional to $\partial_bh_\mu$. Here $h_4=1$, so $\Omega_4=0$, and $\Omega_j$ is a multiple of $\gamma^j\gamma^4$. With $D_j\bar\Psi=\partial_j\bar\Psi-\bar\Psi\Omega_j$ (Chapter 6) the connection enters $T^j{}_j$ only as $\bar\Psi\{\gamma^{x_j},\Omega_j\}\Psi$, which vanishes because $\gamma^j$ anticommutes with $\gamma^j\gamma^4$. What remains is evaluated with $\dot u=-ihu$, $\partial_j\Psi=ik_j\Psi$ and the rule, using $BC\gamma^4=(-i\gamma^4)\gamma^4=i$.

\begin{itemize}
\item \textbf{Time part.} $\Psi^\dagger C\gamma^4\partial_4\Psi\mapsto u^\dagger(BC\gamma^4)(-ih)u=u^\dagger(i)(-ih)u=u^\dagger hu$. Next, $(\partial_4\Psi)^\dagger=i\Psi^\dagger h$ ($h$ is Hermitian), and $B$ commutes with $h$ in the good sector (from $h^\dagger B=Bh$ and $h^\dagger=h$), so $-(\partial_4\Psi^\dagger)C\gamma^4\Psi\mapsto-i\,u^\dagger BhC\gamma^4u=-i\,u^\dagger h(BC\gamma^4)u=u^\dagger hu$. Hence $K_4\mapsto\tfrac12(u^\dagger hu+u^\dagger hu)=\varepsilon(u)$.
\item \textbf{Space part.} $\partial_j\Psi=ik_j\Psi$ and $\partial_j\bar\Psi=-ik_j\bar\Psi$ turn the bracket into $2ik_j\bar\Psi\gamma^{x_j}\Psi=2iK_j\Psi^\dagger C\gamma^j\Psi$ (with $\gamma^{x_j}=\gamma^j/h_j$ and $K_j=k_j/h_j$). So $T^j{}_j=-iK_j\Psi^\dagger C\gamma^j\Psi\mapsto-iK_j\,u^\dagger(BC)\gamma^ju=-iK_j\,u^\dagger(-i\gamma^4)\gamma^ju=-K_j\,u^\dagger\gamma^4\gamma^ju=p_j(u)$.
\end{itemize}

For $\lambda\ne0$ the same steps, with $\mathcal L_s=SU'-U=\tfrac\lambda2S^2$, give the recorded $\rho=S_0\varepsilon-\tfrac\lambda2S^2$ and $p_j=S_0\,p_j(u)+\tfrac\lambda2S^2$ of Section 11.6 (there the factor $e^{-3H\zeta}$ plays the role of $V^{-1/2}$, and $\Omega_j$ also contains a multiple of $\gamma^j\gamma^0$, which anticommutes with $\gamma^j$ as well).

Reading off the definition of $h$, the three quantities satisfy $\varepsilon=M_{\mathrm{eff}}\,s+\sum_jp_j$. On a normalised good-sector eigenmode $hu=Eu$ they take simple values. The anticommutators computed in Section 11.3 give $\{h,-i\gamma^4\}=2M_{\mathrm{eff}}$ (only the $A$ term survives, since $(-i\gamma^4)^2=1$ and $-i\gamma^4$ anticommutes with every $\gamma^4\gamma^j$). Sandwich this between $u^\dagger$ and $u$ and use $hu=Eu$ and $u^\dagger h=Eu^\dagger$ ($h$ Hermitian): $2E\,s=2M_{\mathrm{eff}}$, so $s=M_{\mathrm{eff}}/E$. In the same way $\{h,-\gamma^4\gamma^j\}=2K_j$ for $j\le3$ gives $p_j=K_j^2/E$. So a quantum at rest has $s=1$ and $\varepsilon=M_{\mathrm{eff}}$, and a moving one has a pressure equal to its momentum flux.

\textbf{Worked example: a quantum at rest.} For $K=0$, $h=-iM\gamma^4$, and a positive-energy eigenvector satisfies $-i\gamma^4u=u$. The program's standard choice is

\[
u_0=\tfrac12\bigl(e_0-e_4-i\,e_9-i\,e_{13}\bigr),
\]

where $e_n$ is the column with 1 in place $n$ (counting from 0). Only four rows of $\gamma^4$ matter: $(\gamma^4v)_0=-v_{13}$, $(\gamma^4v)_4=+v_9$, $(\gamma^4v)_9=-v_4$ and $(\gamma^4v)_{13}=+v_0$ (the tables of Chapter 2 and of the student guide, §6.3). Then $(\gamma^4u_0)_0=-(-i/2)=i/2$, so $(-i\gamma^4u_0)_0=1/2=(u_0)_0$; similarly $(-i\gamma^4u_0)_4=-i(-i/2)=-1/2$, $(-i\gamma^4u_0)_9=-i(1/2)=-i/2$ and $(-i\gamma^4u_0)_{13}=-i/2$, which are the components of $u_0$. Hence $hu_0=Mu_0$ and $s(u_0)=u_0^\dagger u_0=4\cdot\tfrac14=1$. The same four rows of $B$, $(Bv)_0=iv_9$, $(Bv)_4=-iv_{13}$, $(Bv)_9=-iv_0$ and $(Bv)_{13}=iv_4$, give $Bu_0=u_0$: $u_0$ has Krein sign $+1$.

\textbf{The dilution law.} At $K=0$, $h=M_{\mathrm{eff}}(t)\,(-i\gamma^4)$ is a multiple of $-i\gamma^4$ at every time. Then $\frac{d}{dt}s(u)=i\,u^\dagger\bigl[h,(-i\gamma^4)\bigr]u=0$ (the same computation as for the Hilbert norm, with $h$ Hermitian), for any $M_{\mathrm{eff}}(t)$. The bilinear of $\Psi=V^{-1/2}u$ carries the factor $1/V$, so

\[
S=S_0\,\frac{V_0}{V}\,s(u),\qquad S\,V=\text{constant}.
\]

\textbf{This law drives almost everything that follows.} The mass term $mS\propto V^{-1}$ dilutes like dust in the 7-volume, while the interaction term $\tfrac\lambda2S^2\propto V^{-2}$ dominates at small volume and dies away at large volume. (Stage 1 derives the same law from energy conservation, $\partial_4\rho+\sum_iH_i(\rho+p_{(i)})=0$ with $\rho+p=SM_{\mathrm{eff}}$, wherever $M_{\mathrm{eff}}\ne0$.)

\textbf{The mean-field picture and its three limits.} In EXP-1 to EXP-3 one mode stands for a macroscopic density of quanta, with $S$ in $M_{\mathrm{eff}}=m+\lambda S$ replaced by its average value. This is a \textbf{mean-field} (Hartree) approximation: each quantum moves in the average field of all the others (Chapter 12 introduces the Hartree idea in general). The density is written $S=S_0\,s(u)$ (times $V_0/V$). Three limits decide what the results mean (Stage-3 document §4.6).

\begin{enumerate}
\item \textbf{The occupied level must stay a positive-energy level.} At $K=0$ the rest eigenvector only changes its phase, $u(t)=e^{-i\int M_{\mathrm{eff}}dt}u_0$, so $s(u)=1$ stays fixed while $M_{\mathrm{eff}}$ may change sign. When $M_{\mathrm{eff}}<0$, $\varepsilon=u^\dagger hu=M_{\mathrm{eff}}<0$: the single occupied mode has become a negative-energy state of the instantaneous $h$. The expectation-value rule is stated for a quantum above the sea, so it no longer describes the state. Every row with $M_{\mathrm{eff}}<0$ in EXP-2 and EXP-3, which includes every phantom row and every row with $\rho<0$, is therefore an \textbf{artefact of the mean field with fixed normal ordering}, not an established property of the field.
\item \textbf{The Pauli principle.} For a Grassmann field each state is occupied at most once, and there are 8 positive-energy states per momentum ($h$ has the eigenvalue $+E$ eight times). A finite density of quanta cannot all sit in the $K=0$ mode. A Pauli-consistent homogeneous state is a filled Fermi sea up to a Fermi momentum $k_F$, with a degeneracy pressure that the one-mode picture leaves out; it is negligible only when $k_F\ll m$. This pressure was not computed. For the free condensate of EXP-3 an order-of-magnitude estimate (fits.json, \texttt{meanFiel\allowbreak{}dScales}, with today's Hubble rate $H_0=67.4$ km/s/Mpc) gives $k_F=0.57$ meV at $m=1$ eV (an electron-volt, eV, is the energy unit of particle physics, and 1 meV $=10^{-3}$ eV), and $k_F/m\le0.1$ for $m\ge21$ meV. The interaction is also treated only at Hartree level; the exchange (Fock) term, the correction that the antisymmetry of fermion states adds to the average-field energy (Section 12.6), is for a contact interaction of the same order and is not included.
\item \textbf{The strength of the tuned coupling.} When the attractive models of EXP-3 are tuned to the quoted $w_0$, the coupling in physical units comes out as $\lambda m^2=-0.497\,(m/2.25\ \mathrm{meV})^4$ (fits.json, \texttt{meanFiel\allowbreak{}dScales}; here $\lambda$ is the coupling of the 3-space density $S$ of EXP-3, with the static extra dimensions absorbed, so that $\lambda m^2$ is dimensionless; in the 8-dimensional Lagrangian $\lambda$ has mass dimension $-6$, Section 13.8). Where the Pauli caveat is negligible this is ultra-strong: at $m=24$ meV (where $k_F/m=0.1$ for that model) $\lambda m^2=-6.8\times10^{3}$. Conversely $|\lambda|m^2\le2$ needs $m\le3.2$ meV, where $k_F/m\ge1.5$. So no mass makes both the one-mode picture and a weak Hartree interaction self-consistent for those models. These are estimates from closed forms, not results of the differential equations.
\end{enumerate}

\textbf{States of definite particle number are never phantom} (derived in Stage-3 document §4.4, not machine-checked). Take a state in which each positive-energy mode is either occupied or empty, with occupation numbers $n$, at Hartree level. Each occupied mode contributes its energy $E$ to $\rho$ and its momentum flux, $K^2/(dE)$ after averaging over $d$ directions, to $p$. The interaction adds $-\tfrac\lambda2S^2$ to $\rho$ (the Hartree double-counting correction) and $+\tfrac\lambda2S^2$ to $p$, which cancel in the sum:

\[
\rho+p=\sum n\,\Bigl(E+\frac{K^2}{d\,E}\Bigr)\ \ge\ 0 .
\]

Superpositions of the vacuum and pair states (such as the states created in EXP-4b) are allowed states of the same theory whose pressure contains a term linear in the pair amplitude; they can have $\rho+p<0$ with $\rho>0$ for a while, a known effect in quantum field theory. That is not the mechanism of the phantom rows of EXP-2 and EXP-3, which lie where the occupied rest level has become a negative-energy level.

\subsection{How the five experiments were computed and checked}

Every differential equation below is integrated by CVODE, the solver of the pure-Rust SUNDIALS 7.8.0 engine of the rustSolveIt repositories (Chapter 10 explains how such a solver works). The setup scripts \texttt{scripts/\allowbreak{}setup\_so\allowbreak{}lver.ps1} and \texttt{scripts/\allowbreak{}setup\_so\allowbreak{}lver.sh} clone the engine at a pinned commit into the git-ignored folder \texttt{vendor/r\allowbreak{}ustSolve\allowbreak{}It}; the committed outputs were produced with the Windows 11 engine (commit a8fdff45), which reproduces them byte for byte on Windows and Linux. To get byte identity, install that engine explicitly (\texttt{bash scr\allowbreak{}ipts/set\allowbreak{}up\_solve\allowbreak{}r.sh win\allowbreak{}11}; \texttt{scripts/\allowbreak{}setup\_so\allowbreak{}lver.ps1} defaults to it). Without the argument \texttt{setup\_so\allowbreak{}lver.sh} installs the platform's own engine. With the Linux engine all checks pass, but 14 committed files differ in their last digits and the Jupyter notebook stops at its byte-identity assertion (measured; student guide §4.5). The macOS engine is identical to the Linux one and is expected to behave the same way, but it was not run. Two CVODE configurations are used: BDF with Newton iteration and a dense linear solver (EXP-1), and Adams-Moulton with fixed-point iteration (EXP-2 to EXP-5), the classical choice for non-stiff oscillations. The program is the Rust crate \texttt{studies/\allowbreak{}dirac16c\allowbreak{}omplex\_c\allowbreak{}osmology}; its gamma matrices, $C$, $\gamma^8$ and $B$ are generated from the exact Stage-1 fixture \texttt{artifact\allowbreak{}s/dirac1\allowbreak{}6complex\allowbreak{}/arbitra\allowbreak{}ry-field\allowbreak{}/algebra\allowbreak{}-fixture\allowbreak{}.json}, and every output records the fixture's sha256.

\begingroup
\small
\begin{longtable}{@{}>{\raggedright\arraybackslash}p{0.133\linewidth}>{\raggedright\arraybackslash}p{0.133\linewidth}>{\raggedright\arraybackslash}p{0.133\linewidth}>{\raggedright\arraybackslash}p{0.133\linewidth}>{\raggedright\arraybackslash}p{0.133\linewidth}>{\raggedright\arraybackslash}p{0.133\linewidth}@{}}
\toprule
\textbf{Experiment} & \textbf{CVODE method} & \textbf{rtol, atol} & \textbf{max step} & \textbf{steps} & \textbf{right-hand sides} \\
\midrule
\endfirsthead
\toprule
\textbf{Experiment} & \textbf{CVODE method} & \textbf{rtol, atol} & \textbf{max step} & \textbf{steps} & \textbf{right-hand sides} \\
\midrule
\endhead
EXP-1 & BDF + Newton & $10^{-12}$, $10^{-14}$ & 0.02 & 33866 & 34950 \\
EXP-2 & Adams + fixed point & $10^{-12}$, $10^{-15}$ & 0.02 & 1779784 & 1781396 \\
EXP-3 & Adams + fixed point & $10^{-11}$, $10^{-13}$ & 0.01 in $N$ & 6074 & 10115 \\
EXP-4 & Adams + fixed point & $10^{-13}$, $10^{-14}$ & 2.0 & 339071692 & 511425232 \\
EXP-5 & Adams + fixed point & $10^{-10}$, $10^{-12}$ & 0.02 & 2548 & 4800 \\
\bottomrule
\end{longtable}
\endgroup

Three layers of checking stand behind every number:

\begin{enumerate}
\item \textbf{Self-checks of the program.} Each experiment tests its own results against exact solutions, conserved quantities and a priori limits and prints PASS or FAIL; the last line is SUCCESS or FAILURE. 69 of 69 self-checks pass (10, 15, 14, 23 and 7 for EXP-1 to EXP-5).
\item \textbf{Independent checkers.} \texttt{scripts/\allowbreak{}check\_di\allowbreak{}rac16com\allowbreak{}plex\_exp\allowbreak{}1.py} to \texttt{check\_di\allowbreak{}rac16com\allowbreak{}plex\_exp\allowbreak{}5.py} use only numpy and the standard library. Each rebuilds the gamma matrices from the fixture and recomputes every physical quantity from the raw spinor columns of the output files, never from the program's derived columns. Each also reruns the program into a fresh folder and requires every file to be byte-identical (\textbf{repeat-run byte identity}), and reruns it with ten times tighter tolerances and requires the errors to shrink or to stay at an identified floor (\textbf{refined convergence}). 167 of 167 checks pass (25, 35, 32, 54 and 21). The EXP-3 fit analysis \texttt{scripts/\allowbreak{}analyze\_\allowbreak{}dirac16c\allowbreak{}omplex\_e\allowbreak{}xp3.py} passes its own 10 of 10 checks.
\item \textbf{Two notebooks.} The Jupyter notebook \texttt{notebook\allowbreak{}s/dirac1\allowbreak{}6complex\allowbreak{}\_dark\_se\allowbreak{}ctor.ipy\allowbreak{}nb} reruns everything, recomputes the physics and draws the 17 figures of \texttt{artifact\allowbreak{}s/dirac1\allowbreak{}6complex\allowbreak{}/numeric\allowbreak{}s/figure\allowbreak{}s/} (71 of 71 assertions, \texttt{notebook\allowbreak{}-report.\allowbreak{}json}). The Mathematica notebook \texttt{notebook\allowbreak{}s/Dirac1\allowbreak{}6Complex\allowbreak{}DarkSect\allowbreak{}or.nb} re-derives every reduced equation symbolically and re-integrates the solutions with NDSolve (49 of 49 checks, \texttt{mathemat\allowbreak{}ica-repo\allowbreak{}rt.json}). None of the three rustSolveIt repositories contains a Mathematica notebook, so this one is new.
\end{enumerate}

The collected verdict in \texttt{numerics\allowbreak{}-summary\allowbreak{}.json} is SUCCESS. Chapter 19 lists the commands that reproduce all of this.

\subsection{EXP-1: the frozen field in the primordial background}

\textbf{Purpose.} The primordial field of the author's notebook (Chapter 9) inflates 3-space while the three extra times deflate, and the 7-volume stays constant. EXP-1 asks what dirac16complex does in this field, and what source 8-dimensional Einstein gravity would need to produce the field.

\textbf{The background.} With $H=1$ and $t=Hx_4$, and with the proper hidden coordinate $\zeta=\ln(\sin z)/(6H)$, $z=6Hx_0$ (Chapter 9), the metric is

\[
ds^2=d\zeta^2-dt^2+e^{2H\zeta}\Bigl[e^{2a_4(t)}\bigl(dx_1^2+dx_2^2+dx_3^2\bigr)-e^{-2a_4(t)}\bigl(dx_5^2+dx_6^2+dx_7^2\bigr)\Bigr].
\]

3-space grows as $e^{a_4}$ and the extra times shrink as $e^{-a_4}$, so the product of the seven scale factors does not depend on $t$: the 7-volume is constant. The experiment uses a \textbf{primordial window}, a smooth switch-on and switch-off of the expansion,

\[
a_4'(t)=\tfrac A4\Bigl(1+\tanh\tfrac{t-t_1}{\Delta}\Bigr)\Bigl(1-\tanh\tfrac{t-t_2}{\Delta}\Bigr),\qquad A\in\{1,2\},\quad t_1=2,\quad t_2=7,\quad \Delta=0.5 .
\]

Between $t_1$ and $t_2$ both brackets are close to 2, so $a_4'\approx A$ and 3-space grows by about $e^{A(t_2-t_1)}=e^{5A}$ in total. The program uses the exact integral of this formula, $a_4(t)=A\Delta\,[\ell(2(t-t_1)/\Delta)-\ell(2(t-t_2)/\Delta)]/\bigl(2(1-e^{-2(t_2-t_1)/\Delta})\bigr)$ with $\ell(x)=\ln(1+e^x)$ (Stage-3 document §6.2).

\textbf{The reduction to a mode equation.} Here the vielbein depends on $\zeta$ and $t$: $h_0=1$ for the $\zeta$ direction, $h_{1,2,3}=e^{H\zeta+a_4}$ and $h_{5,6,7}=e^{H\zeta-a_4}$. Apply the connection formula of Section 11.3 with $b=0$ and $b=4$. For $b=0$: $\prod_{c\ne0}h_c=e^{3(H\zeta+a_4)}e^{3(H\zeta-a_4)}=e^{6H\zeta}$, and $\partial_\zeta\ln e^{6H\zeta}=6H$. For $b=4$: $\prod_{c\ne4}h_c=e^{6H\zeta}$ does not depend on $t$. Hence

\[
\gamma^\mu\Omega_\mu=3H\,\gamma^0 ,
\]

and \textbf{$a_4$ has dropped out}. Take a wave in the hidden direction and no momentum elsewhere, $\Psi=e^{-3H\zeta}e^{iK\zeta}u(t)$. The $\zeta$-derivative gives $\gamma^0(-3H+iK)\Psi$; the $-3H$ cancels the connection term $3H\gamma^0\Psi$ (the factor $e^{-3H\zeta}$ plays the role of $V^{-1/2}$), and what remains is

\[
\gamma^4\dot u=(M_{\mathrm{eff}}-iK\gamma^0)\,u\quad\Longleftrightarrow\quad i\dot u=h\,u,\qquad h=-iM_{\mathrm{eff}}\gamma^4-K\gamma^4\gamma^0 .
\]

This is the equation of Section 11.3 with the single momentum $K_0=K$. Stage 2 verified the reduction exactly (Wolfram check \texttt{P\_modes\_\allowbreak{}exactRed\allowbreak{}uction}) for $\lambda=0$, and the Mathematica notebook derives it again symbolically. For $\lambda\ne0$ it holds only pointwise: the scalar density of this mode is $S=S_0\,s(u)/\sin z$, which depends on $x_0$, so $M_{\mathrm{eff}}=m+\lambda S$ is not the same at every $x_0$, and the one run with $\lambda\ne0$ treats $S_0$ as the local density factor at one fixed $x_0$ (a mean-field approximation). For the same reason the $K=0$ state is not a homogeneous condensate: its density falls off as $1/\sin z$ along the hidden direction.

\textbf{What is recorded.} With the density factor $S_0$ and $S=S_0\,s(u)$ (Stage-3 document §6.2):

\[
\rho=S_0\,u^\dagger hu-\tfrac\lambda2S^2,\qquad p_j=S_0\,p_j(u)+\tfrac\lambda2S^2,\qquad \bar p=\tfrac17\sum_{j\in T}p_j,\qquad w=\frac{\bar p}{\rho},
\]

and $\mathrm{KE}_L=\tfrac12S_0\,u^\dagger hu$, $\mathrm{PE}_L=\rho-\mathrm{KE}_L$, $\mathrm{KE}_H=S_0\sum_jp_j(u)$, $\mathrm{PE}_H=mS+\tfrac\lambda2S^2$. The term $-\tfrac\lambda2S^2$ in $\rho$ corrects the double counting of the interaction in $u^\dagger hu$, which contains $M_{\mathrm{eff}}s$.

\textbf{The exact solution.} For constant $M_{\mathrm{eff}}$ the matrix $h$ is constant and $h^2=E^2$ with $E=\sqrt{M_{\mathrm{eff}}^2+K^2}$. The solution of $i\dot u=hu$ is $u(t)=e^{-iht}u_0$. In the power series of the exponential, even powers of $h$ are powers of $E^2$ and odd powers are powers of $E^2$ times $h$, so

\[
u(t)=\Bigl(\cos(Et)-i\,\sin(Et)\,\frac hE\Bigr)u_0 .
\]

Since this matrix is unitary ($h$ is Hermitian) and commutes with $h$, $\varepsilon=u^\dagger hu$ is constant for every initial spinor. An eigenvector of $h$ only changes its phase, so all its bilinears are constant. A mixture $u=\alpha u_++\beta u_-$ of a $+E$ and a $-E$ eigenvector has bilinears with cross terms $\alpha^\ast\beta\,e^{2iEt}u_+^\dagger Mu_-$ and their conjugates, which oscillate at the frequency $2E$. This interference of the positive- and negative-energy parts is called Zitterbewegung (German for trembling motion).

\textbf{The source the geometry would need.} Chapter 9 (and the Stage-2 document) gives the Einstein tensor of the primordial field (primes are $d/dt$):

\[
G^4{}_4=3H^2(7+a_4'^2),\quad G^0{}_0=-3H^2(a_4'^2-5),\quad G^i{}_i=H^2(15-3a_4'^2+a_4''),\quad G^j{}_j=H^2(15-3a_4'^2-a_4''),
\]

for $i=1,2,3$ and $j=5,6,7$. With $G^\mu{}_\nu=\kappa T^\mu{}_\nu$ and $T^4{}_4=-\rho$, the required energy density and pressures are

\[
\rho_{\mathrm{req}}=-\frac{3H^2(7+a_4'^2)}{\kappa},\qquad p_{\mathrm{req},0}=-\frac{3H^2(a_4'^2-5)}{\kappa},\qquad p_{\mathrm{req},i}=\frac{G^i{}_i}{\kappa},\qquad p_{\mathrm{req},j}=\frac{G^j{}_j}{\kappa}.
\]

$\rho_{\mathrm{req}}$ is negative for every $a_4$. Average the seven pressures (with $\kappa=H=1$): $\bar p_{\mathrm{req}}=\tfrac17[-3(a_4'^2-5)+3(15-3a_4'^2+a_4'')+3(15-3a_4'^2-a_4'')]=\tfrac17(105-21a_4'^2)=15-3a_4'^2$; the $a_4''$ terms cancel. Hence

\[
w_{\mathrm{req}}=\frac{\bar p_{\mathrm{req}}}{\rho_{\mathrm{req}}}=-\frac{5-a_4'^2}{7+a_4'^2}.
\]

Worked values: $a_4'=0$ gives $\rho_{\mathrm{req}}=-21$ and $w_{\mathrm{req}}=-5/7=-0.714$; $a_4'=1$ gives $-24$ and $-4/8=-0.5$; $a_4'=2$ gives $-33$ and $-1/11=-0.091$.

\textbf{Initial data.} $H=m=\kappa=1$, $S_0=1$, $t\in[0,10]$ with 201 samples, hidden momenta $K\in\{0,0.5,2\}$ and both profiles $A=1,2$. For each $(A,K)$ four initial spinors: a joint eigenvector of $h$ and $B$ with energy $+E$ and Krein sign $+1$ (run name \texttt{\detokenize{pos_Bp}}), energy $+E$ and Krein sign $-1$ (\texttt{\detokenize{pos_Bm}}), energy $-E$ and Krein sign $+1$ (\texttt{\detokenize{neg_Bp}}), and a normalised generic mixture of both energy signs (\texttt{\detokenize{mix}}). Only the positive-energy eigenstates are one-particle states; the \texttt{\detokenize{neg_Bp}} runs are numerical controls (in the quantum theory those levels are filled by the sea). One more run per profile has $\lambda=0.5$ and $K=0.5$ and starts on the positive-energy eigenvector of the self-consistent mass $M_\ast=1+\lambda S_0M_\ast/\sqrt{M_\ast^2+K^2}$, solved by Newton's method: $M_\ast=1.47348266406420$. In total 26 runs. At rtol $10^{-10}$ the drifts of the conserved quantities exceeded the a priori limit $10^{-8}$, so the tolerance was tightened to $10^{-12}$ rather than the limit relaxed (Stage-3 document §6.4).

\begin{figure}[htbp]
\centering
\includegraphics[width=0.92\linewidth]{artifacts/dirac16complex/numerics/figures/exp1_frozen_observables.png}
\caption{EXP-1 observables for profile \texorpdfstring{$A=1$}{A=1} (profile \texorpdfstring{$A=2$}{A=2} is bit-identical): (a) energy density, (b) mean transverse pressure, (c) Lagrangian split, (d) \texorpdfstring{$w=\bar p/\rho$}{w=p/ρ}, (e) and (f) drifts of \texorpdfstring{$\rho$}{ρ} and of the Hilbert and Krein norms over all 26 runs.}
\end{figure}

\textbf{Results} (exp1/summary.json and exp1/python-check-report.json).

\begin{itemize}
\item \textbf{Frozen observables.} $\rho$ is constant for every initial spinor (largest drift $4.19\times10^{-9}$), because $h$ is constant and Hermitian. The pressures and $S$ are constant for the eigenstates (largest drift $2.64\times10^{-9}$). For the mixtures with $K\ne0$ the hidden-space pressure $p_0$ and $S$ oscillate at $2E$: the range of $p_0$ is 0.803 for $K=2$ and 0.341 for $K=0.5$, and the checker's exact solution reproduces the range to $3.7\times10^{-10}$.
\item \textbf{Eigenmode laws.} $K=0$: $\rho=1$, $p=0$, $w=0$, $\mathrm{KE}_L=\mathrm{PE}_L=0.5$ (dust). $K=0.5$: $\rho=E=1.1180$, $p_0=K^2/E=0.2236$, $w=K^2/(7E^2)=0.02857$. $K=2$: $\rho=2.2361$, $p_0=1.7889$, $w=4/35=0.11429$. Every free eigenmode has $\mathrm{KE}_L=\mathrm{PE}_L$ whatever $K$ is; the actual pressure is the hidden-space momentum flux.
\item \textbf{Repulsive interaction.} The $\lambda=0.5$ run (pointwise mean field) has $\rho=1.3318$, $\bar p=0.2471$, $w=0.1856$, $\mathrm{KE}_L=0.7780$ and $\mathrm{PE}_L=0.5538$; the six pressures $p_1,\dots,p_7$ equal $\tfrac\lambda2S^2=0.2242$, and $M_{\mathrm{eff}}$ stays at $M_\ast$ to $1.3\times10^{-9}$.
\item \textbf{Profile independence.} $a_4$ enters $h$ only through $k_j/h_j$ with $k_j=0$, so the $A=1$ and $A=2$ spinors are bit-identical (difference 0.0).
\item \textbf{Mixed states are coherent pair states.} A mixture $u=\alpha u_++\beta u_-$ is not a state of one quantum. It does describe a state of the positive-norm Fock space: the Dirac sea with its level $u_-$ replaced by $u$, a superposition of the vacuum and one particle-antiparticle pair. Its normal-ordered expectation values are the output columns minus a constant, so the $2E$ oscillation of $p_0$ and $S$ is physical in that state: an interference of its vacuum and pair parts. The mixtures take negative values of $p_0$, $\bar p$ and $S$ at some times, and carry tensor bilinears that the eigenstates do not have.
\item \textbf{Negative-energy controls.} The \texttt{\detokenize{neg_Bp}} runs have $\rho=-E$ in this reading; physically an antiparticle is a hole in such a level and carries $+E$.
\item \textbf{The required source has negative energy.} $\rho_{\mathrm{req}}$ lies in $[-24.0,-21.0]$ for $A=1$ and in $[-33.0,-21.0]$ for $A=2$; its largest value is $-21.000000000113$. $w_{\mathrm{req}}$ lies in $[-0.714,-0.500]$ and $[-0.714,-0.091]$. It is negative because $\rho_{\mathrm{req}}<0$ while the mean required pressure is positive, not because the pressure is negative. $a_4$ reaches 5.0 ($A=1$) and 10.0 ($A=2$), so 3-space grows by $e^5$ and $e^{10}$, while $V/V_0-1$ stays at $6.7\times10^{-16}$.
\end{itemize}

\begin{figure}[htbp]
\centering
\includegraphics[width=0.92\linewidth]{artifacts/dirac16complex/numerics/figures/exp1_einstein_requirement.png}
\caption{EXP-1 background and Einstein requirement: (a) the window \texorpdfstring{$a_4'(t)$}{a\_4'(t)}, (b) the 3-space factor \texorpdfstring{$e^{a_4}$}{e\textasciicircum{}a\_4}, the extra-time factor \texorpdfstring{$e^{-a_4}$}{e\textasciicircum{}-a\_4} and the constant 7-volume, (c) \texorpdfstring{$\rho_{\mathrm{req}}$}{ρ\_req} for both profiles beside the energy density of the free \texorpdfstring{$K=0$}{K=0} field, (d) \texorpdfstring{$w_{\mathrm{req}}$}{w\_req}.}
\end{figure}

\textbf{Can dirac16complex supply that source?} Not with the positive-energy states integrated here: they have $\rho>0$, while $\rho_{\mathrm{req}}\le-21$. Stage 2 settles the general question for states that depend on $x_0$ and $x_4$ only (Stage-2 document §15.5, Chapter 9). No such state can source the field on a time interval where $a_4''\ne0$, because such states have $T^i{}_i=T^j{}_j$ while $G^i{}_i-G^j{}_j=2H^2a_4''$. For a linear $a_4=ct$ ($a_4''=0$) there are exact sources, and their energy density is $\rho=-3H^2(7+c^2)/\kappa<0$; they need $mS<0$, which is possible because the indefinite $C$ lets $S$ take either sign. The window of EXP-1 has $a_4''\ne0$ except at one instant: writing $\tau_1=\tanh((t-t_1)/\Delta)$, $\tau_2=\tanh((t-t_2)/\Delta)$ and using $d\tanh(x)/dx=(1-\tanh x)(1+\tanh x)$,

\[
a_4''=\frac{A}{4\Delta}(1+\tau_1)(1-\tau_2)\bigl[(1-\tau_1)-(1+\tau_2)\bigr]=-\frac{A}{4\Delta}(1+\tau_1)(1-\tau_2)(\tau_1+\tau_2),
\]

which vanishes only where $\tau_1=-\tau_2$, that is at $t=(t_1+t_2)/2=4.5$ (derived here; the brackets $1+\tau_1$ and $1-\tau_2$ are always positive). Every time interval contains an open piece with $a_4''\ne0$, so no state of this kind can source the EXP-1 background on any time interval.

\textbf{Verification.} The program passes 10 of 10 self-checks and the checker 25 of 25; among them:

\begin{Verbatim}[fontsize=\small]
exp1/summary.json (Rust)     exact_solution_all_runs, rho_frozen_all_runs,
                             a4_profile_independence, einstein_source_negative_energy,
                             seven_volume_constant
exp1/python-check-report.json
                             exactSolution, eigenmodeLaws, einsteinSourceNegative,
                             mixedStateOscillationMatchesExact, repeatByteIdentity,
                             refinedConvergence
\end{Verbatim}

The largest distance from the exact propagator is $6.25\times10^{-9}$ (limit $10^{-7}$); a five-point finite-difference residual of $i\dot u=hu$ is at most 0.63 of its truncation bound. The Mathematica notebook re-integrates all 26 runs with NDSolve (largest state deviation from the Rust result $2.92\times10^{-9}$).

\begin{figure}[htbp]
\centering
\includegraphics[width=0.92\linewidth]{artifacts/dirac16complex/numerics/figures/mathematica/exp1_mixed_state_rho_p0.png}
\caption{Mathematica cross-check of the mixed state with \texorpdfstring{$K=2$}{K=2}: \texorpdfstring{$\rho$}{ρ} frozen and \texorpdfstring{$p_0$}{p\_0} oscillating at \texorpdfstring{$2E$}{2E}; dots are the Rust samples, lines the independent NDSolve solution.}
\end{figure}

\textbf{What EXP-1 shows.} In the primordial field the 7-volume is constant, so the dilution law $SV=$ constant freezes the field: energy density and pressure do not change at all while 3-space inflates by up to $e^{10}$. The frozen equation of state of the positive-energy eigenstates is dust ($K=0$) or slightly stiff ($K\ne0$, or $\lambda>0$); none of them has negative pressure. \textbf{What it does not show.} It does not produce the primordial field: that geometry needs a source of negative energy density, which the positive-energy states cannot be. The frozen density matters for dark energy only as a dilution effect seen from 3-space (Section 11.14).

\subsection{EXP-2: an eight-dimensional universe that feels the condensate}

\textbf{Purpose.} EXP-2 lets the geometry respond to the matter. It couples the condensate to 8-dimensional Einstein gravity, with the hidden space, 3-space and the extra times free to evolve, and asks how $\rho$, $p$ and the anisotropy evolve, whether the extra times keep deflating, and what a 3-space observer would infer.

\textbf{The metric and Einstein's equations.} With $\kappa=\kappa_8=1$ and $m=1$,

\[
ds^2=-dt^2+b^2dx_0^2+a^2\bigl(dx_1^2+dx_2^2+dx_3^2\bigr)-c^2\bigl(dx_5^2+dx_6^2+dx_7^2\bigr),\qquad V=b\,a^3c^3 .
\]

\textbf{The curvature of a diagonal metric that depends on $t$ only.} We compute it for the general metric of Section 11.3 with the diagonal formulas of Section 4.5, with $h_4=1$, $\eta_{44}=-1$ and $g_{ii}=\epsilon_ih_i^2$ for $i\in T$. Every $h_i$ depends on $t=x_4$ only, so case (a) gives nothing ($\partial_4\ln h_4=0$ and $\partial_i\ln h_i=0$), case (b) gives only $\partial_4\ln h_i$, case (c) gives only symbols with the upper index 4, and case (d) gives zero. The only nonzero Christoffel symbols are therefore

\[
\Gamma^4{}_{ii}=-\eta_{44}\,\epsilon_i\,\frac{h_i}{h_4^2}\,\partial_4h_i=\epsilon_i\,h_i\dot h_i\quad\text{(case (c))},\qquad \Gamma^i{}_{4i}=\Gamma^i{}_{i4}=\partial_4\ln h_i=H_i\quad\text{(case (b))},
\]

for $i\in T$ (no sum). In particular $\Gamma^\lambda{}_{44}=0$ for every $\lambda$, and $\Gamma^\rho{}_{\rho4}=\sum_{i\in T}H_i=\Theta$. The Ricci tensor is $R_{\sigma\nu}=R^\rho{}_{\sigma\rho\nu}$ with the Riemann formula of Section 4.7:

\[
R_{\sigma\nu}=\partial_\rho\Gamma^\rho{}_{\nu\sigma}-\partial_\nu\Gamma^\rho{}_{\rho\sigma}+\Gamma^\rho{}_{\rho\lambda}\Gamma^\lambda{}_{\nu\sigma}-\Gamma^\rho{}_{\nu\lambda}\Gamma^\lambda{}_{\rho\sigma},
\]

and nothing depends on the $x_i$, so every derivative along an $x_i$ vanishes. For $R_{44}$ the first and third terms contain $\Gamma^\lambda{}_{44}=0$, the second is $-\partial_4\Theta=-\sum_i\dot H_i$, and in the fourth, $-\sum_{\rho,\lambda}\Gamma^\rho{}_{4\lambda}\Gamma^\lambda{}_{\rho4}$, only $\rho=\lambda=i$ contributes, giving $-\sum_iH_i^2$. So, as in the primordial example of Section 4.7,

\[
R_{44}=-\partial_4\Theta-\sum_iH_i^2=-\sum_{i\in T}\bigl(\dot H_i+H_i^2\bigr).
\]

For $R_{ii}$ (no sum) the first term is $\partial_4\Gamma^4{}_{ii}$, the second is a derivative along $x_i$ and vanishes, the third is $\Gamma^\rho{}_{\rho4}\Gamma^4{}_{ii}=\Theta\,\Gamma^4{}_{ii}$, and in the fourth only $(\rho,\lambda)=(4,i)$ and $(\rho,\lambda)=(i,4)$ contribute, each with $\Gamma^4{}_{ii}H_i$. With $h_i\dot h_i=h_i^2H_i$ and $\tfrac{d}{dt}(h_i\dot h_i)=\dot h_i^2+h_i\ddot h_i=h_i^2(\dot H_i+2H_i^2)$ (because $\dot H_i=\ddot h_i/h_i-H_i^2$):

\[
R_{ii}=\partial_4\Gamma^4{}_{ii}+\Theta\,\Gamma^4{}_{ii}-2\,\Gamma^4{}_{ii}H_i=\epsilon_i\Bigl(\tfrac{d}{dt}(h_i\dot h_i)+\Theta h_i\dot h_i-2h_i\dot h_iH_i\Bigr)=\epsilon_ih_i^2\bigl(\dot H_i+H_i\Theta\bigr).
\]

For $\mu\ne\nu$, $R_{\mu\nu}=0$, because every term then contains a derivative along some $x_i$ or a vanishing Christoffel symbol. Raising an index with $g^{44}=-1$ and $g^{ii}=\epsilon_i/h_i^2$ (the $\epsilon_i$ cancel, since $\epsilon_i^2=1$),

\[
R^4{}_4=\sum_i\bigl(\dot H_i+H_i^2\bigr),\qquad R^i{}_i=\dot H_i+H_i\Theta\quad(\text{no sum}),\qquad R=R^4{}_4+\sum_iR^i{}_i=2\sum_i\dot H_i+\sum_iH_i^2+\Theta^2 .
\]

Since $\Theta^2=\sum_iH_i^2+2\sum_{i<j}H_iH_j$,

\[
G^4{}_4=R^4{}_4-\tfrac12R=\tfrac12\Bigl(\sum_iH_i^2-\Theta^2\Bigr)=-\sum_{i<j}H_iH_j ,
\]

the sum over the 21 unordered pairs of the seven directions. The EXP-2 checker recomputes the Einstein tensor from the metric independently (check \texttt{einstein\allowbreak{}TensorFr\allowbreak{}omMetric}), and the Mathematica notebook derives it symbolically.

Einstein's equation $G_{\mu\nu}=\kappa T_{\mu\nu}$ in $D$ dimensions can be rewritten with its trace: $g^{\mu\nu}$ applied to $R_{\mu\nu}-\tfrac12g_{\mu\nu}R=\kappa T_{\mu\nu}$ gives $R(1-D/2)=\kappa T$, so $R_{\mu\nu}=\kappa\bigl(T_{\mu\nu}-g_{\mu\nu}T/(D-2)\bigr)$. With the isotropic pressure of the condensate, $T=T^\mu{}_\mu=-\rho+(D-1)p$ and $T^i{}_i-T/(D-2)=(\rho-p)/(D-2)$. Hence, for $D=8$:

\[
\sum_{i<j}H_iH_j=\kappa\rho,\qquad \dot H_i=-H_i\Theta+\frac{\kappa}{6}(\rho-p),\qquad \dot h_i=H_ih_i .
\]

The first is a \textbf{constraint} (it contains no second derivatives and must hold at every time); the others are \textbf{evolution equations}. Counting the pairs by groups (hidden-3-space: 3 pairs; hidden-extra-time: 3; within 3-space: 3; within the extra times: 3; 3-space-extra-time: 9) gives

\[
\sum_{i<j}H_iH_j=3H_bH_a+3H_bH_c+3H_a^2+3H_c^2+9H_aH_c .
\]

\textbf{The Friedmann equations as a check.} The same equations with $D=4$ and three equal directions give $3H^2=\kappa\rho$ and $\dot H=-3H^2+\tfrac\kappa2(\rho-p)$. Then $\ddot a/a=\dot H+H^2=-\tfrac23\kappa\rho+\tfrac\kappa2(\rho-p)=-\tfrac\kappa6(\rho+3p)$: exactly the Friedmann equations of Section 11.2.

\textbf{Matter.} The condensate is the $K=0$ mode: $h=-iM_{\mathrm{eff}}\gamma^4$, $S=S_0(V_0/V)s(u)$ with $V_0=1$, $\rho=mS+\tfrac\lambda2S^2$ and $p=\tfrac\lambda2S^2$. Note that $\rho-p=mS$ does not depend on $\lambda$. CVODE integrates 38 real numbers: $\ln b$, $\ln a$, $\ln c$, $H_b$, $H_a$, $H_c$ and the 32 components of $u$.

\textbf{Is a diagonal metric consistent?} Only the diagonal Einstein equations are integrated, so the off-diagonal stresses must vanish. For a homogeneous state Stage 1 gives $T_{ij}=\tfrac14(H_i-H_j)\,\bar\Psi\gamma_i\gamma_j\gamma^{x_4}\Psi$ for $i\ne j$. They vanish if the tensor bilinears $u^\dagger BC\gamma^i\gamma^j\gamma^4u$ vanish in the 15 planes $(i,j)$ between directions of different groups. Since $u(t)$ is $u(0)$ times a phase, zero stays zero. This is a real condition on $u(0)$: a generic vector of the same joint eigenspace ($h=+M_{\mathrm{eff}}$, $B=+1$) has nonzero bilinears in the planes $(0,1)$, $(0,2)$ and $(0,3)$. The chosen $u(0)=u_0$ of Section 11.4 has zero bilinears in all 21 planes (check \texttt{offDiago\allowbreak{}nalStres\allowbreak{}sVanishe\allowbreak{}s}, value 0.0 in every row).

\textbf{The exact solution} (derived in the study and re-derived by the checker, check \texttt{closedFo\allowbreak{}rmVerifi\allowbreak{}ed}). Multiply the evolution equation by $V$ and use $\dot V=\Theta V$:

\[
\frac{d}{dt}(H_iV)=\dot H_iV+H_i\Theta V=\frac{\kappa}{6}(\rho-p)V=\frac{\kappa}{6}\,mSV=\frac{\kappa mS_0}{6}=:\beta ,
\]

a constant, because $SV=S_0$ ($s(u)=1$ for the rest eigenvector). So $H_iV=H_{i0}+\beta t$. Summing over the seven directions, $\Theta V=\dot V$ obeys $\ddot V=7\beta$, and with $V(0)=1$, $\dot V(0)=\Theta_0$:

\[
V(t)=1+\Theta_0t+\alpha t^2,\qquad \alpha=\frac{7\beta}{2},\qquad H_i(t)=\frac{H_{i0}+\beta t}{V(t)},\qquad u(t)=e^{-i\int_0^tM_{\mathrm{eff}}\,dt'}u_0 .
\]

The program uses this closed form only for its self-checks and to place the output times; CVODE integrates the full system of 38 equations.

\textbf{The singularity and Kasner exponents.} Backward in time $V$ reaches 0 at the root of $\alpha t^2+\Theta_0t+1=0$ nearest to 0, $t_s=-2/(\Theta_0+\sqrt{\mathcal D})$ with the discriminant $\mathcal D:=\Theta_0^2-4\alpha$ (the root formula with the square root moved into the denominator; we write $\mathcal D$ because $D$ is the number of dimensions). Near it $V\approx\dot V(t_s)(t-t_s)$ with $\dot V(t_s)=\Theta_0+2\alpha t_s=\sqrt{\mathcal D}$, so $H_i\approx p^{\mathrm K}_i/(t-t_s)$ and $h_i\propto(t-t_s)^{p^{\mathrm K}_i}$, with the \textbf{Kasner exponents} $p^{\mathrm K}_i$ (the superscript K distinguishes them from the pressures $p_{(i)}$ and $p_j$ of this chapter)

\[
p^{\mathrm K}_i=\frac{H_{i0}+\beta t_s}{\sqrt{\mathcal D}} .
\]

Their sum is $\sum_ip^{\mathrm K}_i=(\Theta_0+7\beta t_s)/\sqrt{\mathcal D}=(\Theta_0+2\alpha t_s)/\sqrt{\mathcal D}=1$. For the sum of squares use $\sum_i(p^{\mathrm K}_i)^2=\bigl(\sum_ip^{\mathrm K}_i\bigr)^2-2\sum_{i<j}p^{\mathrm K}_ip^{\mathrm K}_j$. By the formula for the exponents, $\sum_{i<j}p^{\mathrm K}_ip^{\mathrm K}_j$ is $1/\mathcal D$ times $\sum_{i<j}(H_{i0}+\beta t_s)(H_{j0}+\beta t_s)$, and expanding (with the constraint above at $t=0$, $\sum_{i<j}H_{i0}H_{j0}=\kappa\rho_0$), $\sum_{i<j}(H_{i0}+\beta t_s)(H_{j0}+\beta t_s)=\kappa\rho_0+6\beta\Theta_0t_s+21\beta^2t_s^2$ (each direction lies in 6 of the 21 pairs). Multiplying $\alpha t_s^2+\Theta_0t_s+1=0$ by $6\beta$ and using $6\alpha\beta=21\beta^2$ shows $6\beta\Theta_0t_s+21\beta^2t_s^2=-6\beta=-\kappa mS_0$. Define the dimensionless coupling $x_0:=\lambda S_0/(2m)$, the interaction energy relative to the rest energy (the same letter as the hidden coordinate; the reports call it x0, and as a coupling it always appears as a number). Then $\rho_0=mS_0+\tfrac\lambda2S_0^2=mS_0(1+x_0)$, the expanded sum is $\kappa mS_0(1+x_0)-\kappa mS_0=\kappa mS_0x_0$, and

\[
\sum_i(p^{\mathrm K}_i)^2=1-\frac{2\kappa m\,x_0S_0}{\mathcal D} .
\]

For the free condensate ($x_0=0$) this is the vacuum Kasner relation $\sum_ip^{\mathrm K}_i=\sum_i(p^{\mathrm K}_i)^2=1$: near the singularity the dust term is negligible. The interaction term $\tfrac\lambda2S^2$ is a stiff ($w=1$) component there and changes the exponents. \textbf{Late times:} $V\approx\alpha t^2$ and $H_i\approx\beta t/(\alpha t^2)=(2/7)/t$ in every direction: isotropic 8-dimensional dust.

\textbf{What a 3-space observer infers.} An observer who knows only 3-space and assumes $\rho\propto a^{-3(1+w_{\mathrm{eff}})}$ infers $d\ln\rho/d\ln a=-3(1+w_{\mathrm{eff}})$. The true dilution follows from energy conservation, $\dot\rho=-\Theta(\rho+p)$, so $d\ln\rho/d\ln a=-\Theta(1+w)/H_a$, and

\[
w_{\mathrm{eff}}=-1+\frac{\Theta}{3H_a}(1+w) .
\]

This is a dilution measure, not the deceleration of 3-space. If $\Theta=3H_a$ (static extra dimensions) then $w_{\mathrm{eff}}=w$; at late times in EXP-2, $\Theta\to7H_a$ and dust ($w=0$) appears as $w_{\mathrm{eff}}\to-1+7/3=4/3$.

\textbf{An energy bound} (derived here from the constraint, and monitored by the program as the quantity bound). Expanding the square, $\Theta^2=H_b^2+9H_a^2+9H_c^2+6H_bH_a+6H_bH_c+18H_aH_c$, while twice the constraint is $2\kappa\rho=6H_bH_a+6H_bH_c+6H_a^2+6H_c^2+18H_aH_c$. Subtracting,

\[
\Theta^2-3H_a^2-2\kappa\rho=H_b^2+3H_c^2\ \ge\ 0 .
\]

So wherever $\rho\ge0$ (and $H_a,\Theta>0$), $\Theta\ge\sqrt3\,H_a$, and for $1+w\ge0$ this gives $w_{\mathrm{eff}}\ge-1+(1+w)/\sqrt3$; for dust $-1+1/\sqrt3=-0.42265$. A 3-space observer can see dust dilute like a cosmological constant ($w_{\mathrm{eff}}$ near $-1$) only if the 8-dimensional energy density is negative.

\textbf{Initial data.} $\ln b=\ln a=\ln c=0$, $H_b=0$, $H_a=1$ and $H_c=-0.2$: 3-space expands and the extra times start out deflating. $u(0)=u_0$, a positive-energy quantum at rest with $s=1$. Then $\Theta_0=0+3-0.6=2.4$, and the left side of the constraint at $t=0$ is $C_0=3(1)+3(0.04)+9(1)(-0.2)=1.32$. The constraint must hold initially: $\kappa\rho(0)=mS_0(1+x_0)=C_0$, so

\[
S_0=\frac{1.32}{1+x_0},\qquad \lambda=\frac{2m\,x_0}{S_0},\qquad M_{\mathrm{eff}}(0)=m+\lambda S_0=m(1+2x_0),\qquad \beta=\frac{S_0}{6},\qquad \alpha=\frac{7S_0}{12}.
\]

Three runs, from exp2/summary.json:

\begingroup
\small
\begin{longtable}{@{}>{\raggedright\arraybackslash}p{0.114\linewidth}>{\raggedright\arraybackslash}p{0.114\linewidth}>{\raggedright\arraybackslash}p{0.114\linewidth}>{\raggedright\arraybackslash}p{0.114\linewidth}>{\raggedright\arraybackslash}p{0.114\linewidth}>{\raggedright\arraybackslash}p{0.114\linewidth}>{\raggedright\arraybackslash}p{0.114\linewidth}@{}}
\toprule
\textbf{Run} & \textbf{$S_0$} & \textbf{$\lambda$} & \textbf{$M_{\mathrm{eff}}(0)$} & \textbf{$\beta$} & \textbf{$\alpha$} & \textbf{$t_s$} \\
\midrule
\endfirsthead
\toprule
\textbf{Run} & \textbf{$S_0$} & \textbf{$\lambda$} & \textbf{$M_{\mathrm{eff}}(0)$} & \textbf{$\beta$} & \textbf{$\alpha$} & \textbf{$t_s$} \\
\midrule
\endhead
$x_0=0$ (dust) & 1.32 & 0 & 1 & 0.22 & 0.77 & $-0.4954$ \\
$x_0=-0.4$ (attractive) & 2.2 & $-0.3636$ & 0.2 & 0.3667 & 1.2833 & $-0.6266$ \\
$x_0=+0.5$ (repulsive) & 0.88 & 1.1364 & 2 & 0.1467 & 0.5133 & $-0.4624$ \\
\bottomrule
\end{longtable}
\endgroup

The Kasner exponents of the formula above, for the same three runs (exp2/summary.json):

\begingroup
\small
\begin{longtable}{@{}>{\raggedright\arraybackslash}p{0.160\linewidth}>{\raggedright\arraybackslash}p{0.160\linewidth}>{\raggedright\arraybackslash}p{0.160\linewidth}>{\raggedright\arraybackslash}p{0.160\linewidth}>{\raggedright\arraybackslash}p{0.160\linewidth}@{}}
\toprule
\textbf{Run} & \textbf{$p^{\mathrm K}_b$} & \textbf{$p^{\mathrm K}_a$} & \textbf{$p^{\mathrm K}_c$} & \textbf{$\sum_i(p^{\mathrm K}_i)^2$} \\
\midrule
\endfirsthead
\toprule
\textbf{Run} & \textbf{$p^{\mathrm K}_b$} & \textbf{$p^{\mathrm K}_a$} & \textbf{$p^{\mathrm K}_c$} & \textbf{$\sum_i(p^{\mathrm K}_i)^2$} \\
\midrule
\endhead
$x_0=0$ & $-0.066576$ & 0.544271 & $-0.188746$ & 1 \\
$x_0=-0.4$ & $-0.290250$ & 0.972978 & $-0.542895$ & 3.80851 \\
$x_0=+0.5$ & $-0.035225$ & 0.484182 & $-0.139107$ & 0.76259 \\
\bottomrule
\end{longtable}
\endgroup

The output times are uniform in $\ln V$ with step 0.05 on $[-7,18.5]$ (511 rows). Backward the run stops at $V/V_0=e^{-7}$, about $10^{-3}$ in $t$ above $t_s$; forward it ends at $V/V_0=e^{18.5}$, at $t=11855.5$, 9183.5 and 14519.6. The step cap 0.02 is needed: without it Adams takes steps of about 0.045 and its small amplitude error per step on the oscillating spinor drifts $s(u)$ by $1.7\times10^{-7}$ over $t\sim10^4$ (Stage-3 document §7.4).

\textbf{Where the attractive run leaves the mean field.} With $V_0=1$ and $s=1$, $S=S_0/V$, so $M_{\mathrm{eff}}=m(1+2x_0/V)$ and $\rho=mS(1+x_0/V)$. For $x_0=-0.4$: $M_{\mathrm{eff}}<0$ for $V<0.8$, and $\rho<0$ for $V<0.4$.

\begin{figure}[htbp]
\centering
\includegraphics[width=0.92\linewidth]{artifacts/dirac16complex/numerics/figures/exp2_hubble.png}
\caption{EXP-2 Hubble rates of the hidden space, 3-space and the extra times against \texorpdfstring{$t-t_s$}{t-t\_s} (top), and \texorpdfstring{$H_i(t-t_s)$}{H\_i(t-t\_s)} with the Kasner exponents dotted and the 8-dimensional dust value \texorpdfstring{$2/7$}{2/7} dashed (bottom), for the three runs.}
\end{figure}

\textbf{Results} (exp2/summary.json).

\begin{itemize}
\item \textbf{Isotropisation.} From $H_iV=H_{i0}+\beta t$, the anisotropy $(H_i-H_j)V=H_{i0}-H_{j0}$ is constant (measured drift $6.2\times10^{-10}$), so $H_i-H_j$ decays as $1/V$. At the forward end the relative anisotropy $7\max|H_i-H_j|/\Theta$ is $4.6\times10^{-4}$, $3.6\times10^{-4}$ and $5.6\times10^{-4}$; $H_it$ approaches $2/7=0.2857$, and $w_{\mathrm{eff}}=1.33275$, 1.33288 and 1.33261 approaches $4/3$.
\item \textbf{The extra times turn from deflation to expansion} where $H_c=0$, that is at $t=0.2/\beta=0.9091$, 0.5455 and 1.3636 (measured on the output grid: 0.9095, 0.5457, 1.3639). Nothing in this dynamics keeps them deflating.
\item \textbf{Kasner exponents.} Extrapolating the CVODE solution to $V=0$ reproduces the closed-form exponents of the second table above to $2.9\times10^{-7}$.
\item \textbf{Equation of state.} The dust run keeps $w=0$. The repulsive run is stiff near the singularity ($w\to1$) and dust at late times. The attractive run is dust at late times; backward it becomes phantom ($w<-1$ with $\rho>0$) for $0.4<V/V_0<0.8$ (14 output rows, measured volumes 0.407 to 0.779) and has negative energy density for $V/V_0<0.4$ (122 rows), down to $\rho=-1.06\times10^{6}$ at $V=e^{-7}$. Among the rows with $\rho>0$, the phantom rows are exactly the rows with $\mathrm{KE}_L<0$ (Rust check \texttt{phantom\_\allowbreak{}iff\_nega\allowbreak{}tive\_kin\allowbreak{}etic\_ene\allowbreak{}rgy}); all 136 rows with $M_{\mathrm{eff}}<0$, including the 122 with $\rho<0$ (where $w>1$, because $\rho-p=mS>0$ gives $p<\rho<0$), have $\mathrm{KE}_L=\tfrac12SM_{\mathrm{eff}}<0$. \textbf{These backward rows are outside the domain of the mean field}: in all $14+122$ rows with $M_{\mathrm{eff}}<0$ the output has $s(u)=1$ and mode energy $u^\dagger hu=M_{\mathrm{eff}}<0$, so the occupied mode is a negative-energy level (Section 11.4, limit 1).
\item \textbf{The energy bound.} For dust the smallest $w_{\mathrm{eff}}$ is $-0.38732>-0.42265$, and the smallest $\Theta/(3H_a)$ over rows with $\rho>0$ is 0.61268. In the attractive run the 14 phantom rows have $w_{\mathrm{eff}}<-1+(1+w)/\sqrt3$ (the bound assumes $1+w\ge0$), and $\Theta<\sqrt3H_a$ occurs where $\rho<0$: over all rows $\Theta/(3H_a)$ falls to 0.344, while $\Theta^2-3H_a^2-2\kappa\rho$ itself stays non-negative. Both happen in the artefact region.
\end{itemize}

\begin{figure}[htbp]
\centering
\includegraphics[width=0.92\linewidth]{artifacts/dirac16complex/numerics/figures/exp2_volume_density.png}
\caption{EXP-2 (a) the 7-volume against \texorpdfstring{$t-t_s$}{t-t\_s} with the exact quadratic, (b) energy density and pressure against \texorpdfstring{$V/V_0$}{V/V\_0}, (c) conservation of \texorpdfstring{$SV$}{SV}.}
\end{figure}

\begin{figure}[htbp]
\centering
\includegraphics[width=0.92\linewidth]{artifacts/dirac16complex/numerics/figures/exp2_eos_constraint.png}
\caption{EXP-2 (a) \texorpdfstring{$w=p/\rho$}{w=p/ρ} and (b) the dilution-inferred \texorpdfstring{$w_{\mathrm{eff}}$}{w\_eff} against \texorpdfstring{$\ln V$}{V}, with the lines \texorpdfstring{$4/3$}{4/3}, \texorpdfstring{$-1+1/\sqrt3$}{-1+1/√3} and \texorpdfstring{$-1$}{-1}, and (c) the relative constraint residual with the limit \texorpdfstring{$10^{-9}$}{10\textasciicircum{}-9}.}
\end{figure}

\textbf{Verification.} 15 of 15 self-checks and 35 of 35 checker checks; among them:

\begin{Verbatim}[fontsize=\small]
exp2/summary.json (Rust)     constraint_preserved, S_times_V_constant, exact_solution,
                             backward_kasner_exponents, extra_times_turn_to_expansion,
                             phantom_iff_negative_kinetic_energy
exp2/python-check-report.json
                             constraintPreserved, einsteinTensorFromMetric,
                             diracEquationFd, kasnerExponents, phantomStructure,
                             boundNonnegative
\end{Verbatim}

The relative constraint residual is at most $1.26\times10^{-10}$ (limit $10^{-9}$) and $|SV/(S_0V_0)-1|$ at most $4.03\times10^{-10}$. One subtlety teaches something about numerics: the spinor's phase error of $2.24\times10^{-7}$ at $t\sim10^4$ does not shrink in the refined run ($2.71\times10^{-7}$), while its shape error shrinks from $2.0\times10^{-10}$ to $5.2\times10^{-11}$. The checker identified the cause: the rounding of CVODE's internal time $t_{n+1}=t_n+h$, accumulated over the many capped steps (1779784 in the three runs), not a truncation error; it stays below the predicted rounding bound $1.32\times10^{-6}$ (checks \texttt{spinorPh\allowbreak{}aseWithi\allowbreak{}nTimeRou\allowbreak{}nding} and \texttt{phaseDri\allowbreak{}ftIsTime\allowbreak{}Rounding}). The Mathematica notebook agrees with the Rust solution to $3.16\times10^{-8}$ in $\ln h_i$.

\textbf{What EXP-2 shows.} In a self-consistent 8-dimensional cosmology the condensate's energy density falls as $1/V$ and its pressure as $1/V^2$, so every run ends as isotropic 8-dimensional dust whatever $\lambda$ is; $\lambda$ matters only near the singularity. Two consequences matter for the dark sector. First, the extra times do not stay deflating: 8-dimensional gravity with this matter drives all seven directions to the same expansion, and a 3-space observer then sees the dust dilute like $a^{-7}$ ($w_{\mathrm{eff}}\to4/3$), not like $a^{-3}$. Dust-like behaviour in 3-space therefore requires extra dimensions that are held static by something; EXP-3 and EXP-4 assume this, and nothing computed here provides it. Second, with $\rho\ge0$, $w\ge-1$ and $H_a,\Theta>0$, a 3-space observer can never infer $w_{\mathrm{eff}}<-1+(1+w)/\sqrt3$. \textbf{What it does not show.} The phantom stage and the negative energy density of the attractive run are artefacts of the one-mode approximation, not properties of the field.

\subsection{EXP-3: the condensate as late-time dark energy}

\textbf{Purpose and assumptions.} EXP-3 asks whether an attractive condensate can be the dark energy of the quoted supernova fits. It \textbf{assumes} that the hidden space and the extra times are static ($b$, $c$ constant), so that the 4-dimensional Friedmann equation of Section 11.2 applies with the condensate as one more component. This is an assumption of the experiment, not a consequence of the 8-dimensional equations; EXP-2 showed that those do not keep the extra dimensions static.

\textbf{The model.} With $b$, $c$ constant, $V\propto a^3$ and $S\propto a^{-3}$. Units: $H_0=c=1$, densities in units of $3H_0^2/\kappa_4$. The independent variable is the number of e-folds $N=\ln a$ ($N=0$ today). Write $\sigma=S/S_0$ (equal to $a^{-3}$ in exact arithmetic) and $x_0=\lambda S_0/(2m)$ as in EXP-2. Then $\rho_\psi=mS+\tfrac\lambda2S^2=mS_0\,\sigma(1+x_0\sigma)$. Requiring $\rho_\psi(a=1)=\Omega_\psi$ fixes $mS_0=\Omega_\psi/(1+x_0)=:A$, and

\[
\rho_\psi=A\,\sigma(1+x_0\sigma),\qquad p_\psi=A\,x_0\sigma^2,\qquad w=\frac{x_0\sigma}{1+x_0\sigma}=\frac{x_0}{a^3+x_0},
\]

\[
E^2=\frac{H^2}{H_0^2}=\Omega_ra^{-4}+\Omega_ma^{-3}+\rho_\psi,\qquad \mathrm{KE}_L=\frac A2\,\sigma(1+2x_0\sigma),\qquad \mathrm{PE}_L=\frac A2\,\sigma .
\]

(The splits follow from $\mathrm{KE}_L=\tfrac12SM_{\mathrm{eff}}$ with $M_{\mathrm{eff}}=m(1+2x_0\sigma)$ and $\mathrm{PE}_L=\tfrac12mS$.) \textbf{A check of energy conservation:} $d\rho_\psi/dN=A(-3\sigma-6x_0\sigma^2)=-3A\sigma(1+2x_0\sigma)=-3(\rho_\psi+p_\psi)$, the 4-dimensional continuity equation, as it must be for $\sigma=a^{-3}$.

\textbf{What the computer integrates.} The condensate is not assumed; it is obtained from the evolving spinor, $\sigma=a^{-3}s(u)/s(u_0)$. Since $dN=H\,dt=H_0E\,dt$, every time derivative becomes $d/dN=(1/(H_0E))\,d/dt$. The state has 35 real components: the time $H_0t$, the comoving distance $D_C$ (in units of $c/H_0$), the 32 components of $u$, and $\ln\sigma$:

\[
\begin{aligned}
&\frac{d(H_0t)}{dN}=\frac1E,\qquad \frac{dD_C}{dN}=-\frac{1}{aE},\qquad \frac{du}{dN}=-i\,\frac{M_{\mathrm{eff}}}{H_0}\,\frac{(-i\gamma^4)\,u}{E},\\
&\frac{d\ln\sigma}{dN}=-3+\frac{2\,\mathrm{Re}\bigl(u^\dagger(-i\gamma^4)\,du/dN\bigr)}{s(u)} .
\end{aligned}
\]

The distance equation follows from $dD_C=dz/E$ and $z=e^{-N}-1$, so $dz=-e^{-N}dN=-dN/a$. The last equation is the derivative of $\ln\sigma=-3N+\ln s(u)-\ln s(u_0)$, with $ds/dN=2\,\mathrm{Re}(u^\dagger(-i\gamma^4)\,du/dN)$ because $-i\gamma^4$ is Hermitian. The spinor frequency is $M_{\mathrm{eff}}/H_0=\mu(1+2x_0\sigma)/(1+2x_0)$, where $\mu=M_{\mathrm{eff}}(a=1)/H_0\in\{3,7\}$ is a \textbf{reduced frequency}: real fermions have $m/H_0>10^{30}$, which no integrator could follow, and two small values show that $\rho$, $p$ and $w$ do not depend on the spinor's phase. The distance modulus is $\mathrm{DM}(z)=5\log_{10}((1+z)D_C)$ plus a constant.

\textbf{Closed-form diagnostics} (derived here from the formulas above; the program and the checker test the zero of $\rho_\psi$, the crossing of $w=-1$, the bounce, the zeros of $q_{\mathrm{dec}}$ and the tangent parameters against the CVODE solutions, checks \texttt{rhoZeroL\allowbreak{}ocation}, \texttt{phantomC\allowbreak{}rossing}, \texttt{bounceAn\allowbreak{}dStop}, \texttt{decelera\allowbreak{}tionRoot\allowbreak{}s} and \texttt{tangentC\allowbreak{}PL}). $\rho_\psi=0$ where $a^3=-x_0$, at $a=|x_0|^{1/3}$. $w=-1$ where $x_0/(a^3+x_0)=-1$, that is $a^3=-2x_0$, at $a=(2|x_0|)^{1/3}$; there $M_{\mathrm{eff}}=m(1+2x_0a^{-3})=0$. A \textbf{bounce} ($E^2=0$, the expansion rate reaches zero): multiplying $E^2=0$ by $a^6$ gives $(\Omega_m+A)a^3+\Omega_ra^2+Ax_0=0$. The deceleration parameter follows from the Friedmann equations with radiation ($\rho+3p=2\rho_r$) and matter ($\rho+3p=\rho_m$):

\[
q_{\mathrm{dec}}=\frac{2\Omega_ra^{-4}+\Omega_ma^{-3}+\rho_\psi+3p_\psi}{2E^2}.
\]

The adiabatic sound speed is $c_s^2=dp_\psi/d\rho_\psi=(dp_\psi/d\sigma)/(d\rho_\psi/d\sigma)=2x_0\sigma/(1+2x_0\sigma)$. The \textbf{tangent CPL parameters} at $a=1$: $w_0=w(1)=x_0/(1+x_0)$, and from $dw/da=-3a^2x_0/(a^3+x_0)^2$,

\[
w_a=-\frac{dw}{da}\Big|_{a=1}=\frac{3x_0}{(1+x_0)^2} .
\]

\textbf{Worked example: the model tuned to $w_0=-0.861$.} Inverting $w_0=x_0/(1+x_0)$ gives $x_0=w_0/(1-w_0)=-0.861/1.861=-0.462654$. Then $1+x_0=0.537346$, $w_a=3(-0.462654)/0.288741=-4.807$, $c_s^2=-0.925308/0.074692=-12.39$, $\rho_\psi=0$ at $a=0.462654^{1/3}=0.7734$ ($z=0.293$), $w=-1$ at $a=0.925308^{1/3}=0.9745$ ($z=0.026$), and $A=0.69491/0.537346=1.2932$. Neglecting the tiny $\Omega_r$, the bounce is at $a^3=A|x_0|/(\Omega_m+A)=0.59832/1.59823=0.3744$, $a=0.721$ ($z=0.388$). Today $p_\psi=Ax_0=-0.59832$, so $q_{\mathrm{dec}}=(0.00018+0.305+0.69491-1.79495)/2=-0.397$. All of these agree with the committed values below.

\textbf{Initial data.} $\Omega_r=0.00009$, $\Omega_m=0.305$, $\Omega_\psi=0.69491$ (flat); these are inputs chosen by the numerical programme, since the e-mail gives no $\Omega_m$. Five couplings: $x_0=-0.462654$, $-0.433107$, $-0.3$, $-0.2$ and 0. The first reproduces $w_0=-0.861$, the second $w_0=-0.764$, and $x_0=0$ is dust. With $\mu\in\{3,7\}$ that makes 10 runs. The initial spinor is the positive-energy rest eigenvector with $B=+1$; as in EXP-2 its tensor bilinears vanish in the 12 planes between 3-space and the static directions (checker: 0.0 in every row). Each run goes from $N=0$ backward (towards $a=1/3.5$) and forward (to $a=2$). For every $x_0<0$ the backward run meets the bounce before $a=1/3.5$; smaller $a$ does not exist in these models, so the run stops just after it ($N\ge N_b+0.01$), and a check confirms that the stop sits exactly there. Only the dust model reaches $a=1/3.5$.

\textbf{Results} (exp3/summary.json and fits.json). Tangent parameters and today's sound speed:

\begingroup
\small
\begin{longtable}{@{}>{\raggedright\arraybackslash}p{0.200\linewidth}>{\raggedright\arraybackslash}p{0.200\linewidth}>{\raggedright\arraybackslash}p{0.200\linewidth}>{\raggedright\arraybackslash}p{0.200\linewidth}@{}}
\toprule
\textbf{$x_0$} & \textbf{$w_0$} & \textbf{$w_a$ (tangent)} & \textbf{$c_s^2$ today} \\
\midrule
\endfirsthead
\toprule
\textbf{$x_0$} & \textbf{$w_0$} & \textbf{$w_a$ (tangent)} & \textbf{$c_s^2$ today} \\
\midrule
\endhead
$-0.462654$ & $-0.861$ & $-4.807$ & $-12.39$ \\
$-0.433107$ & $-0.764$ & $-4.043$ & $-6.47$ \\
$-0.3$ & $-0.4286$ & $-1.837$ & $-1.50$ \\
$-0.2$ & $-0.25$ & $-0.9375$ & $-0.667$ \\
0 & 0 & 0 & 0 \\
\bottomrule
\end{longtable}
\endgroup

The characteristic epochs, as redshift $z$ and scale factor $a$:

\begingroup
\small
\begin{longtable}{@{}>{\raggedright\arraybackslash}p{0.200\linewidth}>{\raggedright\arraybackslash}p{0.200\linewidth}>{\raggedright\arraybackslash}p{0.200\linewidth}>{\raggedright\arraybackslash}p{0.200\linewidth}@{}}
\toprule
\textbf{$x_0$} & \textbf{$\rho_\psi=0$} & \textbf{$w=-1$} & \textbf{bounce $E^2=0$} \\
\midrule
\endfirsthead
\toprule
\textbf{$x_0$} & \textbf{$\rho_\psi=0$} & \textbf{$w=-1$} & \textbf{bounce $E^2=0$} \\
\midrule
\endhead
$-0.462654$ & $z=0.293$, $a=0.773$ & $z=0.0262$, $a=0.974$ & $z=0.388$, $a=0.721$ \\
$-0.433107$ & $z=0.322$, $a=0.757$ & $z=0.0490$, $a=0.953$ & $z=0.423$, $a=0.703$ \\
$-0.3$ & $z=0.494$, $a=0.669$ & $z=0.186$, $a=0.843$ & $z=0.633$, $a=0.612$ \\
$-0.2$ & $z=0.710$, $a=0.585$ & $z=0.357$, $a=0.737$ & $z=0.890$, $a=0.529$ \\
0 & none & none & none \\
\bottomrule
\end{longtable}
\endgroup

Deceleration and the bare masses $m/H_0=\mu/(1+2x_0)$ of the two runs per coupling:

\begingroup
\small
\begin{longtable}{@{}>{\raggedright\arraybackslash}p{0.200\linewidth}>{\raggedright\arraybackslash}p{0.200\linewidth}>{\raggedright\arraybackslash}p{0.200\linewidth}>{\raggedright\arraybackslash}p{0.200\linewidth}@{}}
\toprule
\textbf{$x_0$} & \textbf{$q_{\mathrm{dec}}$ today} & \textbf{$q_{\mathrm{dec}}=0$ at} & \textbf{$m/H_0$ ($\mu=3$, 7)} \\
\midrule
\endfirsthead
\toprule
\textbf{$x_0$} & \textbf{$q_{\mathrm{dec}}$ today} & \textbf{$q_{\mathrm{dec}}=0$ at} & \textbf{$m/H_0$ ($\mu=3$, 7)} \\
\midrule
\endhead
$-0.462654$ & $-0.397$ & $z=-0.126$, $a=1.144$ & 40.2, 93.7 \\
$-0.433107$ & $-0.296$ & $z=-0.103$, $a=1.115$ & 22.4, 52.3 \\
$-0.3$ & 0.0533 & $z=0.0290$, $a=0.972$ & 7.5, 17.5 \\
$-0.2$ & 0.239 & $z=0.191$, $a=0.840$ & 5.0, 11.7 \\
0 & 0.500 & none & 3, 7 \\
\bottomrule
\end{longtable}
\endgroup

\begin{figure}[htbp]
\centering
\includegraphics[width=0.92\linewidth]{artifacts/dirac16complex/numerics/figures/exp3_w_of_a.png}
\caption{EXP-3 \texorpdfstring{$w(a)=p_\psi/\rho_\psi$}{w(a)=p\_ψ/ρ\_ψ} of the condensate for the five couplings, with the quoted CPL line \texorpdfstring{$w_0=-0.861$}{w\_0=-0.861}, \texorpdfstring{$w_a=-0.60$}{w\_a=-0.60} and the quoted constant \texorpdfstring{$w=-0.764$}{w=-0.764}. The shaded band (\texorpdfstring{$\pm0.118$}{±0.118}) is only an approximate one-sigma band inferred from the phrase ``roughly two standard deviations from \texorpdfstring{$-1$}{-1}''; no band is drawn for the CPL line because the e-mail gives no error bars. The curves have a pole where \texorpdfstring{$\rho_\psi=0$}{ρ\_ψ=0} and end at the bounce.}
\end{figure}

\begin{figure}[htbp]
\centering
\includegraphics[width=0.92\linewidth]{artifacts/dirac16complex/numerics/figures/exp3_rho_psi.png}
\caption{EXP-3 (a) the condensate energy density \texorpdfstring{$\rho_\psi(a)$}{ρ\_ψ(a)}, which changes sign at \texorpdfstring{$a=\lvert x_0\rvert^{1/3}$}{a=x\_0\textasciicircum{}1/3} (circles), and (b) \texorpdfstring{$E^2=H^2/H_0^2$}{E\textasciicircum{}2=H\textasciicircum{}2/H\_0\textasciicircum{}2}, which reaches zero at the bounce for every attractive model.}
\end{figure}

\begin{figure}[htbp]
\centering
\includegraphics[width=0.92\linewidth]{artifacts/dirac16complex/numerics/figures/exp3_ke_pe.png}
\caption{EXP-3 Lagrangian split: \texorpdfstring{$\mathrm{KE}_L$}{KE\_L} (solid) turns negative, and \texorpdfstring{$w$}{w} drops below \texorpdfstring{$-1$}{-1}, for \texorpdfstring{$a<(2\lvert x_0\rvert)^{1/3}$}{a<(2x\_0)\textasciicircum{}1/3}, where the mean-field mode is a negative-energy level; \texorpdfstring{$\mathrm{PE}_L$}{PE\_L} (dashed) stays positive.}
\end{figure}

\begin{itemize}
\item \textbf{The changing equation of state.} For $x_0<0$, $w$ rises with time ($w_a<0$, the thawing sign of $w_a$, but not a thawing field in the usual sense, which starts at $w\approx-1$ and stays at or above it). In the equations $w$ comes up from $-\infty$ at the pole where $\rho_\psi=0$, crosses $-1$ where $M_{\mathrm{eff}}=0$, passes $w_0$ today, and tends to 0 in the future, because $p_\psi\propto a^{-6}$ dies away faster than $\rho_\psi$. Energy density and pressure change for one reason only: $\sigma=a^{-3}$.
\item \textbf{Validity of the mean field.} Before the crossing ($a<(2|x_0|)^{1/3}$), $M_{\mathrm{eff}}<0$ while $s(u)=1$: the occupied mode is a negative-energy level (Section 11.4, limit 1). In the committed output of every attractive run all phantom rows and all rows with $\rho_\psi<0$ lie there, for example all rows with $a\le0.973$ for $x_0=-0.462654$. \textbf{The phantom epoch, the negative energy density and the bounce are therefore artefacts of the mean field}; the model tuned to $w_0=-0.861$ leaves the valid regime at $z=0.026$.
\item \textbf{Acceleration.} The two models tuned to the quoted values accelerate today ($q_{\mathrm{dec}}=-0.397$ and $-0.296$) and all the way back to the bounce, and begin to decelerate again at $a=1.144$ and 1.115 when the condensate becomes dust-like. Neither has a matter-dominated era in these equations, and before $z=0.026$ and $z=0.049$ neither is inside the validity of the mean field. The weaker couplings decelerate today and accelerate only in the past, near the bounce.
\item \textbf{Phase independence.} $\mu=3$ and $\mu=7$ give the same $\rho$, $p$ and $w$ to a few times $10^{-9}$ (check \texttt{rho\_p\_w\_\allowbreak{}mu\_indep\allowbreak{}endent}, checker \texttt{muIndepe\allowbreak{}ndence}): only $s(u)=1$ matters, not the phase.
\item \textbf{Sound speed.} $c_s^2$ is negative today for every $x_0<0$. In a fluid description a negative $c_s^2$ means that small density fluctuations grow instead of oscillating (a gradient instability); the fluctuations themselves were not computed.
\end{itemize}

\subsection{EXP-3 continued: the quoted supernova fits and the deflating-extra-times variant}

\textbf{The fits the programme asked for do not exist.} The numerical programme asked for (1) a least-squares CPL fit of $w(a)$ on $a\in[1/3.26,1]$, and (2) CPL and constant-$w$ fits of $\mathrm{DM}(z)$ on $z\in[0.01,2.26]$ with the offset profiled; the analysis \texttt{scripts/\allowbreak{}analyze\_\allowbreak{}dirac16c\allowbreak{}omplex\_e\allowbreak{}xp3.py} implements its own Nelder-Mead minimiser in numpy (a method that finds the minimum of a function without derivatives, by moving and shrinking a small cluster of trial points). A fit of $w(a)$ over a range that contains the pole $a=|x_0|^{1/3}$ diverges; the pole lies inside $[1/3.26,1]$ exactly when $|x_0|>(1/3.26)^3=0.028863$. A fit of distances up to $z=2.26$ needs the model to exist up to that redshift; for $x_0<-0.041024$ the bounce comes first. So for all four attractive couplings the requested fits are undefined (fits.json records the reason for each). The supplementary fits (fits.json, \texttt{wFitRest\allowbreak{}ricted} and \texttt{muFitRes\allowbreak{}tricted}) use, for $w(a)$, the range $a\ge\max\bigl(1/3.26,(4|x_0|/3)^{1/3}\bigr)$, which stops where $w=-3$ (there $a^3+x_0=-x_0/3$) and so avoids the pole, and, for the distances, $z\in[0.01,\min(2.26,z_b)]$, which stops at the bounce redshift $z_b$ (for $x_0=0$ both are the requested ranges). They only show that the CPL form does not describe these models:

\begingroup
\small
\begin{longtable}{@{}>{\raggedright\arraybackslash}p{0.200\linewidth}>{\raggedright\arraybackslash}p{0.200\linewidth}>{\raggedright\arraybackslash}p{0.200\linewidth}>{\raggedright\arraybackslash}p{0.200\linewidth}@{}}
\toprule
\textbf{$x_0$} & \textbf{$w(a)$ fit, $(w_0,w_a)$} & \textbf{DM fit, $(w_0,w_a)$} & \textbf{DM fit, constant $w$} \\
\midrule
\endfirsthead
\toprule
\textbf{$x_0$} & \textbf{$w(a)$ fit, $(w_0,w_a)$} & \textbf{DM fit, $(w_0,w_a)$} & \textbf{DM fit, constant $w$} \\
\midrule
\endhead
$-0.462654$ & $(-0.605,-12.83)$ & $(2.106,-48.18)$ & $-2.764$ \\
$-0.433107$ & $(-0.485,-11.68)$ & $(1.846,-39.84)$ & $-2.478$ \\
$-0.3$ & $(-0.074,-7.583)$ & $(1.164,-18.11)$ & $-1.524$ \\
$-0.2$ & $(0.130,-5.267)$ & $(0.871,-9.943)$ & $-1.012$ \\
0 & $(0,0)$ & $(0,0)$ & 0 \\
\bottomrule
\end{longtable}
\endgroup

A scan over 491 couplings $x_0\in[-0.49,0]$ in steps of 0.001 (\texttt{exp3/fit\allowbreak{}s\_scan.c\allowbreak{}sv}) finds $x_0=-0.4626545$ for $w_0=-0.861$, with tangent $w_a=-4.807$, and $x_0=-0.4331066$ for $w_0=-0.764$.

\textbf{Reading the numbers.} Matching today's quoted $w_0=-0.861$ forces $w_a=-4.81$: the condensate evolves about eight times faster than the quoted fit ($w_a=-0.60$). Its effective mass vanishes and $w=-1$ at $z=0.026$ (the quoted CPL line crosses $-1$ at $a=0.768$), and beyond that redshift the mean field is outside its valid regime. The model therefore describes nothing beyond $z=0.026$ and cannot be confronted with supernovae up to $z=2.26$, let alone with the cosmic microwave background. The quoted constant $w=-0.764$ is reproduced only as today's value ($x_0=-0.433107$, with $w_a=-4.04$), and a constant-$w$ fit to that model's own distances gives $-2.48$.

\textbf{How $-0.764$ relates to the CPL numbers depends on $\Omega_m$.} This is a fact about the quoted numbers, not about dirac16complex. With $\Omega_m=0.305$ fixed and noise-free, equally weighted distances, the best constant $w$ for the distances of the quoted CPL model is $-1.028$ (offset profiled), $-0.975$ (offset fixed at 0) or $-0.979$ (a grid uniform in $\ln z$). With $\Omega_m$ free it is $-0.9115$ (best $\Omega_m=0.2777$), $-0.8788$ or $-0.8721$: about $-0.9$, not $-0.764$. Freeing $\Omega_m$ closes 44\%, 46\% and 50\% of the gap. How the quoted constant-$w$ fit treated $\Omega_m$ is not stated in the e-mail, so the remaining difference cannot be judged without the supernova likelihood, its errors and its $\Omega_m$.

\begin{figure}[htbp]
\centering
\includegraphics[width=0.92\linewidth]{artifacts/dirac16complex/numerics/figures/exp3_distance_modulus.png}
\caption{EXP-3 distance modulus minus that of the quoted CPL model (\texorpdfstring{$\Omega_m=0.305$}{Ω\_m=0.305}) for the five couplings, a constant \texorpdfstring{$w=-0.764$}{w=-0.764} and \texorpdfstring{$\Lambda$}{Λ}CDM: (a) same \texorpdfstring{$H_0$}{H\_0}, curves ending at the bounce; (b) with each curve's mean offset removed. The dash-dot \texorpdfstring{$w=-0.764$}{w=-0.764} curve uses the assumed \texorpdfstring{$\Omega_m=0.305$}{Ω\_m=0.305} and differs from the CPL curve by up to 0.084 mag with the offset profiled; the dashed one uses \texorpdfstring{$\Omega_m=0.240$}{Ω\_m=0.240}, which best fits the CPL distances, and differs by at most 0.019 mag.}
\end{figure}

\textbf{The deflating-extra-times variant.} The numerical programme also examined one way to give the condensate's $p/\rho$ the quoted tangent: let the extra times keep deflating. Take $c\propto a^{-\gamma}$ with $b$ static. Then $V\propto a^3c^3\propto a^{3(1-\gamma)}$ and $\sigma\propto a^{-n}$ with $n=3(1-\gamma)<3$: the condensate dilutes more slowly than $a^{-3}$. Now $w(a)=x_0/(a^n+x_0)$, so $w_0=x_0/(1+x_0)$ as before, and differentiating, $w_a=-dw/da|_1=n\,x_0/(1+x_0)^2$. Requiring $(w_0,w_a)=(-0.861,-0.60)$ gives $x_0=-0.4626545$, then $n=-0.60(1+x_0)^2/x_0=0.374457$, and $\gamma=1-n/3=0.875181$. By construction the tangent of $p/\rho$ equals the quoted values (analysis check \texttt{gammaVar\allowbreak{}iantRepr\allowbreak{}oducesUn\allowbreak{}iteTange\allowbreak{}nt}). It fails in five ways (fits.json, \texttt{gammaVar\allowbreak{}iant}; Stage-3 document §12.5):

\begin{enumerate}
\item \textbf{Like for like it does not match.} The quoted numbers are fit parameters over a data range, not a tangent. The least-squares CPL fit of the variant's own $p/\rho$ over $a\in[1/3.26,1]$ gives $(w_0,w_a)=(-0.640,-1.98)$ (rms residual 0.154), far from $(-0.861,-0.60)$.
\item \textbf{4-dimensional energy conservation fails.} $d\rho_\psi/dN=-nA\sigma(1+2x_0\sigma)$ while $\rho_\psi+p_\psi=A\sigma(1+2x_0\sigma)$, so $d\rho_\psi/dN+3(\rho_\psi+p_\psi)=(3-n)A\sigma(1+2x_0\sigma)\ne0$: energy flows to or from the extra dimensions. Today its value is 0.254 in units of $3H_0^2/\kappa_4$. Distances then measure $w_{\mathrm{eff}}=-1-\tfrac13\,d\ln\rho_\psi/d\ln a$, whose tangent is $(-0.9827,-0.0749)$, not the quoted pair.
\item \textbf{Newton's constant varies.} In a higher-dimensional theory the 4-dimensional $G_N$ is the higher-dimensional one divided by the volume of the extra dimensions, $G_N=G_8/V_{\mathrm{extra}}$ with $V_{\mathrm{extra}}=bc^3\propto a^{-3\gamma}$. So $G_N\propto a^{3\gamma}$ and $\dot G_N/G_N=3\gamma H_0=2.63\,H_0$, about $1.8\times10^{-10}$ per year for $H_0=67.4$ km/s/Mpc, some 1800 times the order-of-magnitude bound of $10^{-13}$ per year from lunar laser ranging (timing laser pulses reflected from the Moon) recorded in fits.json; and $G_N(z=1)/G_N(0)=2^{-3\gamma}=0.162$. The constant-$G_N$ Friedmann equation used in EXP-3 is itself inconsistent with this.
\item \textbf{It needs negative 8-dimensional energy.} Insert $H_b=0$ and $H_c=-\gamma H_a$ into the EXP-2 constraint: $\kappa\rho_8=3H_a^2+3\gamma^2H_a^2-9\gamma H_a^2=3H_a^2(1-3\gamma+\gamma^2)$. For $\gamma=0.875$ this is $-2.579\,H_a^2<0$. The quadratic $1-3\gamma+\gamma^2$ is negative for every $\gamma$ between its roots $(3-\sqrt5)/2=0.382$ and $(3+\sqrt5)/2=2.618$, so any sustained deflation in that range needs negative total energy density; and EXP-2 showed that the extra times turn to expansion by themselves.
\item \textbf{Its own history.} Its $\rho_\psi$ is negative beyond $z=6.83$, and its $p/\rho$ crosses $-1$ at $z=0.230$.
\end{enumerate}

Its distances are close to those of the quoted CPL model (largest difference 0.021 mag, 0.011 mag with the offset profiled), but so are those of $\Lambda$CDM (0.024 and 0.013 mag), so distances at this level do not distinguish it from a cosmological constant; its own distance fits give CPL $(-0.9707,-0.1611)$ and constant $w=-1.0159$.

\begin{figure}[htbp]
\centering
\includegraphics[width=0.92\linewidth]{artifacts/dirac16complex/numerics/figures/exp3_w0wa_plane.png}
\caption{EXP-3 \texorpdfstring{$(w_0,w_a)$}{(w\_0,w\_a)} plane: (a) the tangent CPL parameters of the condensate for \texorpdfstring{$x_0\in[-0.49,0]$}{x\_0∈[-0.49,0]} with the five canonical couplings, the quoted point, \texorpdfstring{$\Lambda$}{Λ}CDM, the line \texorpdfstring{$w_0+w_a=-1$}{w\_0+w\_a=-1} and the distance-inferred point of the deflation variant; (b) the supplementary restricted fits of the scan, far from the quoted point for every \texorpdfstring{$x_0<0$}{x\_0<0}.}
\end{figure}

\textbf{Verification of EXP-3.} 14 of 14 self-checks (for example \texttt{sigma\_fr\allowbreak{}om\_spino\allowbreak{}r\_equals\allowbreak{}\_a\_minus\allowbreak{}\_3}, \texttt{rho\_p\_w\_\allowbreak{}match\_cl\allowbreak{}osed\_for\allowbreak{}m}, \texttt{phantom\_\allowbreak{}crossing\allowbreak{}\_at\_cube\allowbreak{}\_root\_2a\allowbreak{}bs\_x0}, \texttt{hubble\_p\allowbreak{}ositive\_\allowbreak{}backward\allowbreak{}\_stop\_at\allowbreak{}\_predict\allowbreak{}ed\_bounc\allowbreak{}e}), 32 of 32 checker checks and 10 of 10 analysis checks. $\sigma$ computed from the spinor agrees with $a^{-3}$ to $2.50\times10^{-9}$ (refined run: $2.10\times10^{-10}$); the time and distance agree with independent quadratures to about $10^{-9}$; the Nelder-Mead minimiser recovers synthetic CPL and constant-$w$ curves to about $10^{-12}$. The Mathematica notebook derives the Friedmann constraint, the continuity equation, $w$ and $q_{\mathrm{dec}}$ symbolically and agrees with the Rust solution to $2.12\times10^{-8}$ in the spinor.

\textbf{What EXP-3 shows.} The attractive condensate is the only object in this framework whose own pressure is negative. Its negative pressure $\tfrac\lambda2S^2\propto a^{-6}$ matters only at high density, that is in the past, and there it wins over the positive $mS$: going back in time the effective mass passes through zero. Today's $w_0$ can be tuned to any value in $(-1,0)$ by the single parameter $x_0$, but then $w_a=3x_0/(1+x_0)^2$ is fixed, and the evolution is fast. \textbf{What it does not show.} The crossing of $w=-1$ is not an established phantom crossing: it is the instant where the effective mass vanishes, and beyond it the single mean-field mode occupies a negative-energy level. Everything computed there (phantom $w$, $\mathrm{KE}_L<0$, $\rho_\psi<0$, the bounce) is an artefact of the approximation. No fluctuations, no fit to data and no likelihood were computed.

\subsection{EXP-4a: a thermal gas of quanta}

\textbf{Purpose and assumptions.} EXP-4 looks at quanta instead of a condensate. It assumes a 3-space universe of the Friedmann type with the hidden space and the extra times static ($b=c=1$), and asks how the equation of state of a gas of quanta changes as the universe expands (EXP-4a), and how many quanta the end of inflation creates (EXP-4b and EXP-4c, Section 11.11). Only momenta in 3-space are used: the hidden-space momentum $k_0$ is set to zero. The hidden direction is space-like and in the good sector, so this is an \textbf{assumption}, not a consequence: it holds for a small compact hidden dimension whose lowest excitation energy is much larger than the temperatures and Hubble rates used, or it is imposed by hand. With $k_0$ allowed, the relativistic 3-space pressure would be $\rho/4$ instead of $\rho/3$.

\textbf{The mode equation.} A mode with comoving momentum $k$ along $x_1$ has physical momentum $K=k/a$, which redshifts as $\dot K=-KH$, and

\[
i\dot u=h(t)\,u,\qquad h=-im\gamma^4-K\gamma^4\gamma^1,\qquad h^2=E^2,\qquad E=\sqrt{m^2+K^2}.
\]

\textbf{Eight copies of a two-level system.} Let $Z=-i\gamma^4$ and $X=-\gamma^4\gamma^1$, so that $h=mZ+KX$. Then $Z^2=-(\gamma^4)^2=1$, $X^2=\gamma^4\gamma^1\gamma^4\gamma^1=\gamma^1\gamma^1=1$ (using $\gamma^4\gamma^1\gamma^4=\gamma^1$), and $ZX+XZ=i(\gamma^4\gamma^4\gamma^1+\gamma^4\gamma^1\gamma^4)=i(-\gamma^1+\gamma^1)=0$. Two anticommuting matrices that square to 1 obey the same rules as the Pauli matrices $\sigma_z$ and $\sigma_x$ (Chapter 2). We now show that the 16 components split into 8 two-component blocks on each of which $Z$ and $X$ act as $\sigma_z$ and $\sigma_x$. Both matrices are Hermitian (Section 11.3: $-i\gamma^4$ and $\gamma^4\gamma^j$ for $j\le3$). Since $Z$ is Hermitian with $Z^2=1$, its eigenvalues are $\pm1$, and its two eigenspaces are orthogonal. If $Zv=v$ then $Z(Xv)=-XZv=-Xv$, and $X^2=1$, so $X$ maps the $+1$ eigenspace one-to-one onto the $-1$ eigenspace (and back), and each has dimension 8. Take an orthonormal basis $v_1,\dots,v_8$ of the $+1$ eigenspace. The vectors $Xv_1,\dots,Xv_8$ are orthonormal, because $(Xv_n)^\dagger(Xv_l)=v_n^\dagger X^\dagger Xv_l=v_n^\dagger X^2v_l=v_n^\dagger v_l$, and together the 16 vectors form an orthonormal basis. On each pair $(v_n,Xv_n)$, $Z$ multiplies the first vector by $+1$ and the second by $-1$, and $X$ sends the first to the second and the second back to the first ($X(Xv_n)=X^2v_n=v_n$): in this pair of basis vectors $Z=\sigma_z$ and $X=\sigma_x$. So in this basis $h=(m\sigma_z+K\sigma_x)\otimes I_8$; since $Z$ and $X$ do not depend on $t$, the same basis serves at every time (the study also verified this numerically). So there are \textbf{8 spin states per momentum}, each a two-level system with energies $\pm E$. The relation $Ch^\ast C=-h$ maps particles to antiparticles.

\textbf{Per mode.} $\varepsilon=u^\dagger hu$, $p_1=-K\,u^\dagger\gamma^4\gamma^1u$ and $s=u^\dagger(-i\gamma^4)u$, with $\varepsilon=ms+p_1$ (Section 11.4).

\textbf{Counting states and the thermal sums.} In a box, the number of momentum states per unit volume in $d^3K$ is $d^3K/(2\pi)^3$. (In units with $\hbar=1$, as in Chapter 8, momentum and wave number coincide. In a periodic box of side $L_{\mathrm{box}}$ a plane wave $e^{iK\cdot x}$ must repeat after $L_{\mathrm{box}}$ in each direction, so the allowed momenta are $K=2\pi n/L_{\mathrm{box}}$ with integer components $n$: one state per volume $(2\pi/L_{\mathrm{box}})^3$ of momentum space, that is $L_{\mathrm{box}}^3\,d^3K/(2\pi)^3$ states in $d^3K$, and $d^3K/(2\pi)^3$ per unit volume of space.) For an isotropic distribution $d^3K=4\pi K^2dK$, which gives $K^2dK/(2\pi^2)$ per state label. With $K=k/a$ this is $k^2dk/(2\pi^2a^3)$. There are 16 states per momentum (8 particle and 8 antiparticle states), and each is occupied with the Fermi-Dirac probability $f=1/(e^{E/T}+1)\le1$: the Pauli principle. (At temperature $T$, in units with Boltzmann's constant 1 and with equal numbers of particles and antiparticles, so that the chemical potential is zero, a fermion state of energy $E$ is occupied with probability $1/(e^{E/T}+1)$, never more than 1; Section 12.14 derives this Fermi-Dirac distribution.) The pressure of an isotropic gas is the average momentum flux: a particle with momentum $K$ and speed $v=K/E$ carries $K_xv_x$ across a plane perpendicular to $x$, and averaging over directions gives $Kv/3=K^2/(3E)$. The program integrates a mode moving along $x_1$ and divides its $p_1$ by 3. The gas is treated as \textbf{collisionless} (an assumption of EXP-4: the quanta are free, and the $\lambda$ interaction between them is not included): after $a=1$ nothing changes the occupation of a comoving mode, so $f$ keeps its initial value $f(k)=1/(e^{E_i/T_i}+1)$ (here $E_i=\sqrt{m^2+k^2}$ at $a=1$), while each mode's energy and pressure evolve. The integrals over $k$ are done with \textbf{Gauss-Legendre quadrature}: 48 nodes $k_n$ and weights $w_n$ on $[0,k_{\max}]$, chosen so that $\sum_nw_nF(k_n)$ equals $\int_0^{k_{\max}}F\,dk$ exactly for every polynomial $F$ of degree up to 95. With $W_n=\tfrac{16}{2\pi^2}w_nk_n^2f_n$:

\[
\rho=\frac{1}{a^3}\sum_nW_n\,\varepsilon_n,\qquad p=\frac{1}{3a^3}\sum_nW_n\,p_{1,n},\qquad \mathrm{KE}_H=3p,\qquad \mathrm{PE}_H=\rho-3p,
\]

compared with the kinetic-theory integrals

\[
\rho_{\mathrm{kin}}=\frac{16}{2\pi^2a^3}\int_0^{k_{\max}}k^2f\,E\,dk,\qquad p_{\mathrm{kin}}=\frac{16}{2\pi^2a^3}\int_0^{k_{\max}}k^2f\,\frac{K^2}{3E}\,dk .
\]

\textbf{The two limits.} When $K\gg m$, $E\approx K$ and each mode has $p/\rho=K^2/(3E\cdot E)\approx1/3$: radiation, and $\rho\propto a^{-3}\cdot a^{-1}=a^{-4}$ because each energy redshifts as $1/a$. When $K\ll m$, $E\approx m$ and $p/\rho\approx K^2/(3m^2)\to0$: dust, and $\rho\propto a^{-3}$.

\textbf{Parameters.} $m=1$; radiation era $a(t)=(t/t_i)^{1/2}$ with $t_i=1/(2H_i)=10$, so $H_i=0.05$; comoving temperature $T_i=10$ at $a=1$ (relativistic); $a$ from 1 to 100 ($t$ from 10 to $10^5$), 61 samples uniform in $\ln a$; $k_{\max}=12T_i=120$. Each mode starts from the positive-energy $B=+1$ eigenvector of $h(t_i)$. Extras: a second spin state for four nodes, a negative-energy mode for one node, and the first-order adiabatic vacuum (Section 11.11) for two nodes. Each thermal mode turns through about $10^5$ radians, which is why the tolerance is $10^{-13}$: at $10^{-10}$ the Hilbert norm drifted by $2.4\times10^{-4}$.

\begin{figure}[htbp]
\centering
\includegraphics[width=0.92\linewidth]{artifacts/dirac16complex/numerics/figures/exp4_thermal_w.png}
\caption{EXP-4 thermal gas: (a) \texorpdfstring{$w=p/\rho$}{w=p/ρ} of the CVODE mode sum and of kinetic theory against \texorpdfstring{$a$}{a}, from \texorpdfstring{$1/3$}{1/3} towards 0; (b) relative deviations of \texorpdfstring{$\rho$}{ρ} (about \texorpdfstring{$10^{-8}$}{10\textasciicircum{}-8}) and of \texorpdfstring{$p$}{p} (up to \texorpdfstring{$1.9\times10^{-5}$}{1.9×10\textasciicircum{}-5}, the sudden-start wave) from kinetic theory.}
\end{figure}

\begin{figure}[htbp]
\centering
\includegraphics[width=0.92\linewidth]{artifacts/dirac16complex/numerics/figures/exp4_thermal_scaling.png}
\caption{EXP-4 thermal gas scaling: (a) \texorpdfstring{$\rho a^4$}{ρa\textasciicircum{}4}, constant while the gas is relativistic, and (b) \texorpdfstring{$\rho a^3$}{ρa\textasciicircum{}3}, constant once it is non-relativistic.}
\end{figure}

\begin{figure}[htbp]
\centering
\includegraphics[width=0.92\linewidth]{artifacts/dirac16complex/numerics/figures/exp4_thermal_split.png}
\caption{EXP-4 kinetic and potential energy of the gas: \texorpdfstring{$\mathrm{KE}_H/\rho=3w$}{KE\_H/ρ=3w} (momentum energy) falls from 1 and \texorpdfstring{$\mathrm{PE}_H/\rho$}{PE\_H/ρ} (rest mass) rises towards 1, while the Lagrangian split stays at \texorpdfstring{$\mathrm{KE}_L/\rho=\mathrm{PE}_L/\rho=1/2$}{KE\_L/ρ=PE\_L/ρ=1/2}.}
\end{figure}

\textbf{Results} (exp4/summary.json, key thermal).

\begin{itemize}
\item \textbf{From radiation to dust.} $w=0.3328543348$ at $a=1$ (kinetic theory: the same to ten digits; slightly below $1/3$ because $m=T_i/10$ is not zero) and $w=0.0359224$ at $a=100$ (kinetic theory 0.0359223), both on the grid $k\le k_{\max}$. $\rho a^4$ is constant early and $\rho a^3$ late. The momentum energy $\mathrm{KE}_H=3p$ gives way to rest-mass energy: $\mathrm{KE}_H/\rho=3w$ falls from about 1 to about 0.11. For every mode of a free gas $\mathrm{KE}_L=\mathrm{PE}_L=\rho/2$, so the scalar-field relation $p=\mathrm{KE}-\mathrm{PE}$ of split (A) holds only for condensates.
\item \textbf{Agreement with kinetic theory.} $\rho$ agrees to $1.98\times10^{-8}$. $p$ deviates by up to $1.87\times10^{-5}$, independently of the tolerance. The cause is physical, not numerical: starting a mode on the instantaneous eigenvector leaves a small admixture $\beta$ of the other energy level, which enters $p_1$ at first order, $p_1=(1-2|\beta|^2)K^2/E+2\,\mathrm{Re}(\alpha^\ast\beta)\,mK/E$, while $\varepsilon$ has no such term. The deviation stays inside the independently predicted envelope (ratio 0.28), and the runs that start in the adiabatic vacuum reduce the single-mode pressure deviation from 1.78 to 0.12 relative.
\item \textbf{No spurious particle production.} The occupation-weighted $|\beta|^2$ is $3.7\times10^{-8}$, and the adiabatic-vacuum runs end at $5.6\times10^{-8}$. The largest single-mode value, $6.37\times10^{-5}$ at $k=0.953$, is 0.779 times $4|c|^2$ with $|c|^2=(mkH_i/(4E_i^3))^2$: the sudden-start wave, not particle production (Exercise 11.11 derives the size of $|c|^2$).
\item \textbf{Exact symmetries.} Spin independence and particle-antiparticle symmetry hold exactly (deviation 0.0).
\item \textbf{Momentum truncation.} The cut $k\le12T_i$ misses $2.4\times10^{-3}$ of $\rho$ at $a=1$ and, at $a=100$, $9.0\times10^{-4}$ of $\rho$ and $5.3\times10^{-3}$ of $p$. Without the cut the gas has $w=0.0361$ at $a=100$ instead of 0.0359. The truncation is reported, not corrected.
\end{itemize}

\subsection{EXP-4b and EXP-4c: quanta created by the end of inflation}

\textbf{Pair creation in plain words.} In an expanding universe the Hamiltonian $h(t)$ of every mode changes with time. If it changes slowly, a state that starts in the positive-energy level stays in the positive-energy level of the instantaneous $h$ (this is the adiabatic theorem of quantum mechanics). If it changes quickly, part of the state ends in the other level. Apply this to the Dirac sea: every filled negative-energy state ends with a small weight on the positive-energy level. A particle has appeared, and the hole it leaves behind is an antiparticle. This is \textbf{gravitational pair creation}; nothing but the expansion is needed.

\textbf{Instantaneous levels.} Since $h^2=E^2$, the matrices $P_\pm=\tfrac12(1\pm h/E)$ satisfy $P_\pm^2=\tfrac14(1\pm2h/E+h^2/E^2)=P_\pm$: they project onto the positive- and negative-energy levels. For a mode that starts in its positive-energy level (in the adiabatic vacuum defined below), the weight it ends with on the negative-energy level is

\[
|\beta_k|^2=\frac{|P_-u|^2}{|u|^2}.
\]

In each two-level block the evolution over the whole run is a 2 by 2 unitary matrix $U$ (because $h$ is Hermitian). Its first column and its first row are both unit vectors, so $|U_{++}|^2+|U_{-+}|^2=1=|U_{++}|^2+|U_{+-}|^2$ and hence $|U_{-+}|=|U_{+-}|$: the weight that a filled negative-energy state ends with on the positive-energy level equals $|\beta_k|^2$. With 8 filled sea states per $k$, the expansion leaves $8|\beta_k|^2$ particles and $8|\beta_k|^2$ antiparticles per momentum, and the comoving number density of the created quanta is

\[
na^3=\frac{16}{2\pi^2}\int k^2\,|\beta_k|^2\,dk .
\]

\textbf{The adiabatic expansion} (derived here in outline, following the Stage-3 study). In one block write $h=E(\cos\theta\,\sigma_z+\sin\theta\,\sigma_x)$ with $\tan\theta=K/m$. Its eigenvectors are $u_+=(\cos\tfrac\theta2,\sin\tfrac\theta2)$ with energy $+E$ and $u_-=(-\sin\tfrac\theta2,\cos\tfrac\theta2)$ with energy $-E$, and they rotate as $\dot u_+=\tfrac{\dot\theta}2u_-$, $\dot u_-=-\tfrac{\dot\theta}2u_+$. Insert $u=\alpha(t)e^{-i\varphi}u_++\beta(t)e^{i\varphi}u_-$ with $\varphi=\int E\,dt$ into $i\dot u=hu$. The terms with $E$ cancel, and comparing the coefficients of $u_+$ and $u_-$ gives

\[
\dot\alpha=\frac{\dot\theta}{2}\,\beta\,e^{2i\varphi},\qquad \dot\beta=-\frac{\dot\theta}{2}\,\alpha\,e^{-2i\varphi}.
\]

From $\theta=\arctan(K/m)$ and $\dot K=-KH$, $\dot\theta=m\dot K/(m^2+K^2)=-mKH/E^2$. With $\alpha\approx1$, $\beta(t)=-\int^t\tfrac{\dot\theta}2e^{-2i\varphi}dt'$. Integrate by parts using $e^{-2i\varphi}=\frac{d}{dt}(e^{-2i\varphi})/(-2iE)$:

\[
\beta(t)=\frac{\dot\theta}{4iE}\,e^{-2i\varphi}+\int^t e^{-2i\varphi}\,\frac{d}{dt'}\Bigl(\frac{i\dot\theta}{4E}\Bigr)dt' .
\]

The first term, of size $|\dot\theta|/(4E)=mKH/(4E^3)$, is present at every instant and disappears when the expansion stops; it is the \textbf{adiabatic dressing} of an unexcited mode, not created particles. The program therefore measures $|\beta_k|^2$ in the \textbf{first-order adiabatic basis}, which removes this term. The remaining integral is small as long as $\dot\theta/E$ changes smoothly. Integrating by parts once more shows what happens at a kink: if $\frac{d}{dt}(\dot\theta/E)$ jumps at $t=0$, the boundary terms at $t=0^-$ and $t=0^+$ no longer cancel and leave a created amplitude of size $|\Delta\ddot\theta|/(8E^2)$, where $\Delta\ddot\theta$ is the jump of $\ddot\theta$ ($E$ and $\dot\theta$ are continuous).

\textbf{The sudden end of inflation (EXP-4b).} The background is de Sitter, $a=e^{H_{\mathrm{inf}}t}$ for $t<0$, glued at $t=0$ to radiation, $a=(1+2H_{\mathrm{inf}}t)^{1/2}$. Both give $a(0)=1$ and $H(0)=H_{\mathrm{inf}}$ (for radiation $H=H_{\mathrm{inf}}/(1+2H_{\mathrm{inf}}t)$), so $a$ and $H$ are continuous; but $\dot H$ jumps from 0 to $-2H_{\mathrm{inf}}^2$ (a $C^1$ junction). Differentiating $\dot\theta=-mKH/E^2$, the only discontinuous piece of $\ddot\theta$ is $-mK\dot H/E^2$, so $\Delta\ddot\theta=2mKH_{\mathrm{inf}}^2/E^2$, and at $a=1$ ($K=k$) the created weight is

\[
|\beta_k|^2\approx\Bigl(\frac{mk\,H_{\mathrm{inf}}^2}{4E^4}\Bigr)^2 ,\qquad E=\sqrt{m^2+k^2},
\]

the \textbf{kink formula}. It is the leading term of an expansion that is valid for $E\gg H_{\mathrm{inf}}$, so it describes the large-$k$ tail. Integrated over all $k$ it gives $na^3=H_{\mathrm{inf}}^4/(64\pi m)$ (Exercise 11.13): the number that the jump of $\dot H$ alone would create.

\textbf{Parameters.} $H_{\mathrm{inf}}=1$, masses $m/H_{\mathrm{inf}}\in\{0,0.1,0.5,1,2\}$, 64 Gauss-Legendre nodes in $\ln k$ on $[10^{-3},40]$. Each mode starts in the first-order adiabatic vacuum when it is deep inside the horizon, at $k/a=200H_{\mathrm{inf}}$, and the run ends when $H=10^{-4}m$, that is at $a_{\mathrm{end}}^2=10^4H_{\mathrm{inf}}/m$ (for $m=0$ the end of $m=0.1$ is used). The integration is restarted at $t=0$. The analytic kink tail beyond $k=40$ is added separately as a correction.

\textbf{The same transition made smooth (EXP-4c).} To measure how much of the yield is due to the sharp junction, every mode is integrated again in a background in which the transition takes about one Hubble time:

\[
\frac1H=\frac{1}{H_{\mathrm{inf}}}+\tau\ln\bigl(1+e^{2t/\tau}\bigr),\qquad \epsilon:=-\frac{\dot H}{H^2}=\frac{d}{dt}\Bigl(\frac1H\Bigr)=1+\tanh\frac t\tau,\qquad \tau=\frac{1}{H_{\mathrm{inf}}},
\]

which is de Sitter ($\epsilon=0$) long before $t=0$ and radiation ($\epsilon=2$) long after. Here $\ln a$ has no closed form and is integrated by CVODE as a 33rd state component. Because $H_{\mathrm{smooth}}\le H_{\mathrm{sudden}}$ at every time, the smooth universe expands less, and its late-time $a$ is smaller by the constant factor 0.656. An abundance depends on the number density at a given $H$, so the comparison ``at the same final $H$'' multiplies the smooth $na^3$ by $0.656^{-3}$.

\begin{figure}[htbp]
\centering
\includegraphics[width=0.92\linewidth]{artifacts/dirac16complex/numerics/figures/exp4_pair_spectra.png}
\caption{EXP-4 pair creation by the de Sitter to radiation transition: \texorpdfstring{$\lvert\beta_k\rvert^2$}{β\_k\textasciicircum{}2} at the end (first-order adiabatic basis) for \texorpdfstring{$m/H_{\mathrm{inf}}=0.1$}{m/H\_inf=0.1}, 0.5, 1 and 2, with the kink tail (dashed) and the spectra of the smooth transition (dotted); for \texorpdfstring{$m=0$}{m=0}, \texorpdfstring{$\lvert\beta_k\rvert^2\le7.7\times10^{-25}$}{β\_k\textasciicircum{}2≤7.7×10\textasciicircum{}-25} (no production).}
\end{figure}

\textbf{Results} (exp4/summary.json, key pair). The sudden transition, on the grid $k\le40$, final spectrum in the first-order adiabatic basis:

\begingroup
\small
\begin{longtable}{@{}>{\raggedright\arraybackslash}p{0.267\linewidth}>{\raggedright\arraybackslash}p{0.267\linewidth}>{\raggedright\arraybackslash}p{0.267\linewidth}@{}}
\toprule
\textbf{$m/H_{\mathrm{inf}}$} & \textbf{$na^3$ in units of $H_{\mathrm{inf}}^3$} & \textbf{largest final $\lvert\beta_k\rvert^2$} \\
\midrule
\endfirsthead
\toprule
\textbf{$m/H_{\mathrm{inf}}$} & \textbf{$na^3$ in units of $H_{\mathrm{inf}}^3$} & \textbf{largest final $\lvert\beta_k\rvert^2$} \\
\midrule
\endhead
0 & $1.8\times10^{-24}$ & $7.7\times10^{-25}$ \\
0.1 & $1.454\times10^{-3}$ & 0.347 \\
0.5 & $4.990\times10^{-3}$ & 0.0510 \\
1 & $4.412\times10^{-3}$ & 0.00588 \\
2 & $2.421\times10^{-3}$ & $3.54\times10^{-4}$ \\
\bottomrule
\end{longtable}
\endgroup

The smooth transition, compared at the same final $H$, and the kink formula integrated over all $k$:

\begingroup
\small
\begin{longtable}{@{}>{\raggedright\arraybackslash}p{0.200\linewidth}>{\raggedright\arraybackslash}p{0.200\linewidth}>{\raggedright\arraybackslash}p{0.200\linewidth}>{\raggedright\arraybackslash}p{0.200\linewidth}@{}}
\toprule
\textbf{$m/H_{\mathrm{inf}}$} & \textbf{$na^3$ smooth} & \textbf{sudden / smooth} & \textbf{$H_{\mathrm{inf}}^4/(64\pi m)$} \\
\midrule
\endfirsthead
\toprule
\textbf{$m/H_{\mathrm{inf}}$} & \textbf{$na^3$ smooth} & \textbf{sudden / smooth} & \textbf{$H_{\mathrm{inf}}^4/(64\pi m)$} \\
\midrule
\endhead
0.1 & $1.074\times10^{-3}$ & 1.35 & $4.97\times10^{-2}$ \\
0.5 & $1.536\times10^{-3}$ & 3.25 & $9.95\times10^{-3}$ \\
1 & $3.327\times10^{-4}$ & 13.3 & $4.97\times10^{-3}$ \\
2 & $3.043\times10^{-6}$ & 796 & $2.49\times10^{-3}$ \\
\bottomrule
\end{longtable}
\endgroup

The equation of state of the created gas, from its final (frozen) spectrum; ``$w$ at $a=1$'' evaluates that spectrum at $a=1$, a hypothetical value, because the occupation of the long-wavelength modes is not yet final there:

\begingroup
\small
\begin{longtable}{@{}>{\raggedright\arraybackslash}p{0.200\linewidth}>{\raggedright\arraybackslash}p{0.200\linewidth}>{\raggedright\arraybackslash}p{0.200\linewidth}>{\raggedright\arraybackslash}p{0.200\linewidth}@{}}
\toprule
\textbf{$m/H_{\mathrm{inf}}$} & \textbf{$w$ at $a=1$} & \textbf{$w$ at the end} & \textbf{$\rho a^3$ at the end} \\
\midrule
\endfirsthead
\toprule
\textbf{$m/H_{\mathrm{inf}}$} & \textbf{$w$ at $a=1$} & \textbf{$w$ at the end} & \textbf{$\rho a^3$ at the end} \\
\midrule
\endhead
0 & $1/3$ & $1/3$ & $1.3\times10^{-26}$ \\
0.1 & 0.3081 & $1.20\times10^{-4}$ & $1.454\times10^{-4}$ \\
0.5 & 0.2539 & $1.14\times10^{-4}$ & $2.495\times10^{-3}$ \\
1 & 0.2386 & $1.62\times10^{-4}$ & $4.413\times10^{-3}$ \\
2 & 0.2395 & $3.02\times10^{-4}$ & $4.845\times10^{-3}$ \\
\bottomrule
\end{longtable}
\endgroup

\begin{figure}[htbp]
\centering
\includegraphics[width=0.92\linewidth]{artifacts/dirac16complex/numerics/figures/exp4_pair_density.png}
\caption{EXP-4 (a) the created comoving number density \texorpdfstring{$na^3$}{na\textasciicircum{}3} (adiabatic basis solid, instantaneous basis dashed), and (b) the equation of state of the created gas with its final spectrum, from nearly \texorpdfstring{$1/3$}{1/3} when evaluated at \texorpdfstring{$a=1$}{a=1} to dust.}
\end{figure}

\begin{itemize}
\item \textbf{Production for $m>0$ only.} For $m=0$, $h=K(t)\,\sigma_x\otimes I_8$: its eigenvectors, those of $\sigma_x$, do not depend on $K$, so $\theta=\pi/2$ is constant, $\dot\theta=0$, and no state ever leaves its level. This is the conformal invariance of the massless equation (it does not change when all lengths are rescaled by a time-dependent factor, and an expanding universe of this type is flat space rescaled in this way), and the measured $|\beta_k|^2$ ($8.5\times10^{-25}$ at most over the run) is roundoff. For $m>0$ the sudden transition gives a few $10^{-3}H_{\mathrm{inf}}^3$, and $|\beta_k|^2\le1$ everywhere, as the Pauli principle requires (check \texttt{pairPaul\allowbreak{}i}).
\item \textbf{What sets the yield.} The large-$k$ tail follows the kink formula to 4.2\% (limit 35\%; checks \texttt{pair\_tai\allowbreak{}l\_matche\allowbreak{}s\_kink\_t\allowbreak{}heory} and \texttt{pairTail\allowbreak{}MatchesK\allowbreak{}inkTheor\allowbreak{}y}). For $m\gtrsim H_{\mathrm{inf}}$ the whole sudden yield is the response to the jump of $\dot H$: the kink formula integrated over all $k$ is 1.13 and 1.03 times the computed $na^3$ for $m=1$ and 2, and the smooth transition gives 13.3 and 796 times fewer quanta at the same final $H$. For $m=0.5$ the sharpness still matters (factor 3.25). Only the light field, $m=0.1$, is robust: its production comes from small $k$, from modes that were stretched beyond the Hubble length during inflation, and the smooth transition changes its largest $|\beta_k|^2$ only from 0.347 to 0.346 and its $na^3$ by a factor 1.35. So the ranking of the sudden yields (largest near $m=0.5H_{\mathrm{inf}}$) and their size for $m\gtrsim0.5H_{\mathrm{inf}}$ are properties of the chosen junction; the sudden values are the yields of an idealised instantaneous end of inflation.
\item \textbf{The created gas becomes dust.} Its $w$ is between 0.24 and 0.31 when its final spectrum is evaluated at $a=1$ and about $10^{-4}$ at the end, and $\rho a^3$ approaches $m\,na^3$ (for $m=1$: $4.413\times10^{-3}$ against $na^3=4.412\times10^{-3}$).
\item \textbf{Basis and truncation.} The literal instantaneous-basis final $|\beta_k|^2$ includes the adiabatic dressing at the end of the run and gives almost the same numbers ($na^3=1.455\times10^{-3}$, $4.991\times10^{-3}$, $4.414\times10^{-3}$ and $2.422\times10^{-3}$). The analytic tail beyond $k=40$ adds $1.05\times10^{-6}$ to $na^3$ for $m=2$ (a fraction $4.3\times10^{-4}$).
\end{itemize}

\textbf{Verification of EXP-4.} 23 of 23 self-checks and 54 of 54 checker checks. The checker recomputes every mode sum from the raw spinors, re-derives kinetic theory with its own quadrature, and re-integrates every one of the 320 sudden and 320 smooth pair modes with an independent fourth-order Magnus method (a step-by-step integrator for linear equations whose steps are exactly unitary) for the two-level reduction: the final $|\beta_k|^2$ agree with the Rust values to $10^{-6}$ relative (checks \texttt{pairSudd\allowbreak{}enSpectr\allowbreak{}umMagnus\allowbreak{}Referenc\allowbreak{}e} and \texttt{pairSmoo\allowbreak{}thTransi\allowbreak{}tionRefe\allowbreak{}rence}). Unitarity holds to $2.4\times10^{-7}$ (thermal) and $4.1\times10^{-8}$ (pairs). The thermal modes carry a phase error of at most $6.5\times10^{-6}$ rad at $a=100$ that the refined run does not reduce (the same time-rounding effect as in EXP-2); only phase-independent quantities ($\varepsilon$, $p_1$, $s$, $|\beta|^2$) enter the results. The Mathematica notebook re-integrates four thermal nodes and ten pair modes with NDSolve (the two hardest thermal nodes only up to $a=10$, so for them the cross-check is partial).

\textbf{What EXP-4 shows.} In the good sector, with the extra dimensions held static and no hidden-space momentum, dirac16complex quanta are an ordinary gas of massive fermions with 16 states per momentum. Their pressure is momentum flux and redshifts away, so the equation of state goes continuously from radiation to dust. The expansion itself creates such quanta for every $m>0$. \textbf{What it does not show.} No relic density in physical units: that needs physical values of $m$ and $H_{\mathrm{inf}}$ and the history after inflation, none of which the framework fixes; and for $m\gtrsim0.5H_{\mathrm{inf}}$ the yield depends strongly on how inflation ends. Whether such a relic would be hot, warm or cold dark matter depends on its free streaming, which was not computed.

\subsection{EXP-5: momentum along an extra time}

\textbf{Purpose.} Every result above uses the good sector. EXP-5 shows why: a mode with momentum along an extra time becomes unstable once the extra times have deflated enough.

\textbf{Equations.} The extra times deflate as $c=h_5=h_6=h_7=e^{-Ht}$ ($H=1$), and a mode has momentum $q$ along $x_5$ and none elsewhere ($m=1$). Its physical momentum is $Q(t)=q/h_5=q\,e^{Ht}$, and by Section 11.3 (with $\eta^{55}=-1$)

\[
h(t)=-im\gamma^4-Q(t)\,\gamma^4\gamma^5,\qquad h^2=E^2(t)=m^2-q^2e^{2Ht},\qquad E^2<0\ \text{for}\ t>t_\ast=\frac{\ln(m/q)}{H}.
\]

$h$ is not Hermitian (Section 11.3), so the Hilbert norm $u^\dagger u$ is not conserved; but $h^\dagger B=Bh$, so the Krein norm $u^\dagger Bu$ is. After $t_\ast$ the eigenvalues of $h$ are $\pm i\varkappa$ with $\varkappa=\sqrt{Q^2-m^2}$, and $i\dot u=\pm i\varkappa u$ gives $u\propto e^{\pm\int\varkappa\,dt}$: one part grows. At leading order in the WKB approximation (the approximation that the coefficients change slowly compared with the local rate), $d\ln(u^\dagger u)/dt=2\varkappa$; at the next order $2\varkappa-m^2/\varkappa^2$. The growth exponent $W(t)=\int_{t_\ast}^t\varkappa\,dt'$ has a closed form: with $dQ=QH\,dt$,

\[
W(t)=\frac1H\int_m^{Q(t)}\frac{\sqrt{Q'^2-m^2}}{Q'}\,dQ'=\frac1H\Bigl[\sqrt{Q^2-m^2}-m\arccos\frac mQ\Bigr],
\]

as one checks by differentiating (with $\frac{d}{dQ}\arccos(m/Q)=m/(Q\sqrt{Q^2-m^2})$):

\[
\frac{d}{dQ}\Bigl[\sqrt{Q^2-m^2}-m\arccos\frac mQ\Bigr]=\frac{Q}{\sqrt{Q^2-m^2}}-\frac{m^2}{Q\sqrt{Q^2-m^2}}=\frac{\sqrt{Q^2-m^2}}{Q}.
\]

\textbf{Initial data.} $q\in\{0.05,0.1\}$, so $t_\ast=\ln20=2.99573$ and $\ln10=2.30259$; $t$ runs from 0 to $t_\ast+3$ with 601 samples. The initial spinor is a positive-energy eigenvector of $h(0)$, with $E_0=\sqrt{m^2-q^2}=0.998749$ and 0.994987, that is also an eigenvector of $C$ with eigenvalue $\pm1$ (possible because $C$ commutes with $\gamma^4$ and $\gamma^5$, hence with $h(t)$). Its Krein norm is not zero: from $hu=E_0u$ and $u^\dagger u=1$, $E_0=u^\dagger hu=m\,s(u)-q\,u^\dagger\gamma^4\gamma^5u$; the first term is real, and the second is imaginary because $\gamma^4\gamma^5$ is real antisymmetric, so $s(u)=E_0/m$. With $Cu=\pm u$, $Bu=-i\gamma^4Cu=\pm(-i\gamma^4)u$ and $u^\dagger Bu=\pm E_0/m$. Four runs. The WKB comparison uses the window from the first output time at or after $t_\ast+1.5$ to $t_\ast+3$, away from the turning point.

\begin{figure}[htbp]
\centering
\includegraphics[width=0.92\linewidth]{artifacts/dirac16complex/numerics/figures/exp5_growth.png}
\caption{EXP-5 (a) \texorpdfstring{$\ln(u^\dagger u)$}{(u\textasciicircum{}†u)} for both momenta and both \texorpdfstring{$C$}{C} signs with the WKB curve \texorpdfstring{$2W(t)$}{2W(t)}, and (b) \texorpdfstring{$E^2(t)$}{E\textasciicircum{}2(t)}, which turns negative at \texorpdfstring{$t_\ast=\ln(m/q)/H$}{t\_∗=(m/q)/H} (dotted).}
\end{figure}

\begin{figure}[htbp]
\centering
\includegraphics[width=0.92\linewidth]{artifacts/dirac16complex/numerics/figures/exp5_krein.png}
\caption{EXP-5 (a) the Krein-norm drift divided by \texorpdfstring{$\max(u^\dagger u,1)$}{(u\textasciicircum{}†u,1)}, below \texorpdfstring{$10^{-8}$}{10\textasciicircum{}-8} throughout, and (b) the Krein norm, drawn while \texorpdfstring{$u^\dagger u<10^8$}{u\textasciicircum{}†u<10\textasciicircum{}8}, which stays at its initial value \texorpdfstring{$\pm E_0/m$}{±E\_0/m}. Later, when \texorpdfstring{$u^\dagger u$}{u\textasciicircum{}†u} approaches \texorpdfstring{$10^{16}$}{10\textasciicircum{}16}, the computed \texorpdfstring{$u^\dagger Bu$}{u\textasciicircum{}†Bu} is a difference of numbers of that size and loses its digits to rounding (absolute drift up to 1.005, \texttt{kreinDri\allowbreak{}ft.maxAb\allowbreak{}s} in \texttt{exp5/sum\allowbreak{}mary.jso\allowbreak{}n}), which is why panel (a) divides the drift by \texorpdfstring{$\max(u^\dagger u,1)$}{(u\textasciicircum{}†u,1)}.}
\end{figure}

\textbf{Results} (exp5/summary.json).

\begin{itemize}
\item $u^\dagger u$ reaches $1.258\times10^{16}$ ($q=0.05$) and $1.313\times10^{16}$ ($q=0.1$). The average growth rate of $\ln(u^\dagger u)$ is 2.37 ($q=0.05$) and 2.40 ($q=0.1$) over the first unit of time after $t_\ast$, and 25.3 over the last unit: ratios 10.7 and 10.6. The growth is faster than exponential.
\item Over the WKB window, $\Delta\ln(u^\dagger u)=30.9922$ against the leading WKB value 31.0241 (relative deviation $1.03\times10^{-3}$) and the first-order value 30.99986 ($2.47\times10^{-4}$) for $q=0.05$; 30.9455 against 30.9770 and 30.9530 for $q=0.1$. The late rate matches $2\varkappa-m^2/\varkappa^2$ to $4.6\times10^{-6}$.
\item The Krein norm is conserved: its drift divided by $\max(u^\dagger u,1)$ is at most $1.08\times10^{-9}$, and its absolute drift up to $t_\ast$ at most $1.28\times10^{-9}$. At the end the absolute drift is about 1 because $u^\dagger Bu$ is then a difference of two numbers of order $10^{16}$, the limit of 16-digit arithmetic.
\end{itemize}

\textbf{Verification.} 7 of 7 self-checks (for example \texttt{energy\_s\allowbreak{}quared\_c\allowbreak{}hanges\_s\allowbreak{}ign\_at\_t\allowbreak{}star}, \texttt{hilbert\_\allowbreak{}norm\_sup\allowbreak{}erexpone\allowbreak{}ntial\_gr\allowbreak{}owth}, \texttt{growth\_m\allowbreak{}atches\_w\allowbreak{}kb\_leadi\allowbreak{}ng\_order}) and 21 of 21 checker checks (for example \texttt{kreinCon\allowbreak{}servedNo\allowbreak{}rmalized}, \texttt{wkbLeadi\allowbreak{}ng}, \texttt{wkbFirst\allowbreak{}Order}, \texttt{wkbLateR\allowbreak{}ate}); an independent Runge-Kutta reference agrees to $3.34\times10^{-8}$ relative. The Mathematica notebook verifies $h^2=(m^2-Q^2)\cdot1$, $h^\dagger B=Bh$ and that $h$ is not Hermitian, and a 32-digit run agrees with machine precision to $1.8\times10^{-14}$.

\textbf{What EXP-5 shows.} A mode with extra-time momentum stops oscillating and grows once the extra-time scale factor has shrunk enough. Because the indefinite Krein norm is conserved, the growth is a growing positive-norm part together with a growing negative-norm part: the ill-posedness of wave equations with more than one time (called ultrahyperbolic). No positive Fock space exists for these modes. Any cosmology built on this field must exclude them, and every result of EXP-1 to EXP-4 is restricted to the good sector. \textbf{What it does not show.} The restriction is imposed by hand, not derived; whether some dynamical mechanism suppresses the extra-time sector is an open problem (Chapter 18).

\subsection{What the experiments show about dark matter, and what they do not}

\textbf{How pressure and energy density behave.} Three mechanisms control everything the five experiments found.

\begin{itemize}
\item \textbf{Dilution.} $SV$ is constant, so the mass term $mS$ behaves like dust in the 7-volume and the interaction term $\tfrac\lambda2S^2$ like a stiff fluid in the 7-volume.
\item \textbf{Redshift.} The pressure of quanta is momentum flux, and momenta fall as $1/a$, so the pressure of a gas falls faster than its energy density.
\item \textbf{Geometry.} What a 3-space observer infers depends on how the 7-volume grows compared with the 3-volume, $w_{\mathrm{eff}}=-1+\Theta(1+w)/(3H_a)$. In EXP-1 $\Theta=0$; in EXP-2 $\Theta\to7H_a$; in EXP-3 and EXP-4 $\Theta=3H_a$ by assumption.
\end{itemize}

Case by case:

\begin{itemize}
\item \textbf{Free condensate, any background (one-mode mean field):} $\rho=mS\propto1/V$ and $p=0$, dust, up to the Fermi-degeneracy pressure of a Pauli-consistent state, which was not computed.
\item \textbf{Repulsive condensate:} $w\to1$ at small $V$ and $w\to0$ at large $V$ (EXP-1 pointwise, EXP-2).
\item \textbf{Attractive condensate:} $p<0$, and $w\to0$ at large $V$; where $M_{\mathrm{eff}}<0$ the equations give phantom $w$ and then $\rho<0$, but these are mean-field artefacts (EXP-2, EXP-3).
\item \textbf{Primordial field:} $V$ constant, so $\rho$ and $p$ are frozen: dust or slightly stiff for the positive-energy eigenstates (EXP-1).
\item \textbf{Thermal gas:} $\rho\propto a^{-4}$ early and $a^{-3}$ late; $w$ falls from 0.3329 to 0.0359 (EXP-4a).
\item \textbf{Created pairs:} dust ($w\approx10^{-4}$) by the end of each run, with $\rho a^3\to m\,na^3$ (EXP-4b).
\item \textbf{Extra-time modes:} no positive norm and faster-than-exponential growth; $\rho$ and $p$ are not defined in a positive Fock space (EXP-5).
\end{itemize}

\textbf{What behaves like dark matter.}

\begin{enumerate}
\item A Pauli-consistent gas of quanta ($f\le1$, 8 + 8 states per momentum, $k_0=0$) goes from radiation to dust: $w=0.3329$ at $T=10m$ and $w=0.0359$ after a hundredfold expansion on the momentum grid (0.0361 without the cut), in agreement with kinetic theory; the momentum-energy fraction $\mathrm{KE}_H/\rho=3w$ falls from about 1 to about 0.11 (EXP-4a). This result does not depend on the one-mode picture.
\item The end of inflation creates quanta for every $m>0$ and none for $m=0$: with the sudden junction $na^3=1.454\times10^{-3}$, $4.990\times10^{-3}$, $4.412\times10^{-3}$ and $2.421\times10^{-3}\,H_{\mathrm{inf}}^3$ for $m/H_{\mathrm{inf}}=0.1$, 0.5, 1 and 2, and a transition lasting one Hubble time gives 1.35, 3.25, 13.3 and 796 times fewer at the same final $H$. The created gas is dust by the end of each run (EXP-4b, EXP-4c).
\item The free condensate is dust in every background in the one-mode picture (EXP-1 with $K=0$, EXP-2 and EXP-3 with $x_0=0$); an order-of-magnitude estimate says its Fermi-degeneracy pressure is negligible for the free EXP-3 condensate if $m\ge21$ meV.
\item Every condensate, repulsive or attractive, becomes dust at large volume, because the interaction energy dilutes faster (EXP-2, EXP-3).
\end{enumerate}

\textbf{What is not established.}

\begin{enumerate}
\item \textbf{Abundance.} No relic density was computed. The yields are per comoving volume in units of $H_{\mathrm{inf}}^3$; turning them into the observed dark-matter density needs physical values of $m$ and $H_{\mathrm{inf}}$ and the history after inflation, and for $m\gtrsim0.5H_{\mathrm{inf}}$ the yield itself depends strongly on how inflation ends.
\item \textbf{Darkness.} The framework has no coupling to the Standard Model of particle physics, so the quanta are dark by construction; no bounds on interactions, on annihilation through the $\lambda$ term or on decays were derived.
\item \textbf{Static extra dimensions.} In self-consistent 8-dimensional gravity all seven directions end up expanding equally, and a 3-space observer then sees dust dilute like $a^{-7}$ ($w_{\mathrm{eff}}\to4/3$, EXP-2). Dust-like dilution in 3-space, which dark matter needs, requires static extra dimensions; EXP-3 and EXP-4 assume them, nothing derives them.
\item \textbf{The extra-time sector} must be excluded by hand (EXP-5).
\item \textbf{Structure formation}, the free streaming of the quanta and their clustering were not computed.
\item \textbf{The hidden-space momentum} $k_0$ was set to zero in EXP-4; with it the relativistic pressure would be $\rho/4$.
\end{enumerate}

\textbf{The answer on dark matter: a qualified yes.} In the good sector, with stabilised extra dimensions and no hidden-space momentum, dirac16complex quanta are pressureless matter at late times, and the expansion of the universe produces them gravitationally. That is a dark-matter-like equation of state and a production mechanism. It is not a dark-matter model, and nothing computed here predicts the observed amount of dark matter.

\subsection{What the experiments show about dark energy, and what they do not}

\textbf{Which mechanisms give negative pressure or $w<-1/3$.}

\begin{enumerate}
\item \textbf{The attractive four-fermion interaction} ($\lambda<0$) is the only mechanism with genuinely negative pressure, $p=\tfrac\lambda2S^2<0$. At low density $w=x_0\sigma/(1+x_0\sigma)$ lies between $-1$ and 0. At higher density $M_{\mathrm{eff}}$ passes through zero, and beyond that point the equations give $w<-1$ with $\rho>0$, then $\rho<0$, and in EXP-3 a bounce. \textbf{This phantom crossing is a mean-field artefact}: in all those rows the single occupied mode is a negative-energy level, outside the domain of the expectation-value rule, whereas a state with definite occupation numbers of positive-energy quanta always has $\rho+p\ge0$ (Section 11.4).
\item \textbf{The frozen density of the primordial field.} In EXP-1 $\Theta=0$, so a 3-space observer would infer $w_{\mathrm{eff}}=-1+0\cdot(1+w)=-1$, the dilution of a cosmological constant (derived here from the formula of Section 11.7, not a separate machine check). But the actual pressure is zero or positive, and the geometry is not self-consistent: 8-dimensional Einstein gravity would need $\rho_{\mathrm{req}}\le-21$.
\item \textbf{Deflating extra times} (the EXP-3 variant) give $p/\rho$ the quoted tangent by construction, but its fit over the data range is $(-0.640,-1.98)$, and it violates 4-dimensional energy conservation, the constancy of Newton's constant and the 8-dimensional energy bound (Section 11.9).
\end{enumerate}

The Lagrangian has no cosmological-constant term ($U(0)=0$ is a choice of the model, Chapter 6 and Stage-1 document §7.6; adding a cosmological constant would be a new assumption, not a property of the field), and a condensate has $w=-1$ only at the instant where $M_{\mathrm{eff}}=0$.

\textbf{Where the condensate fails as dark energy} (Stage-3 document §12.5):

\begin{enumerate}
\item \textbf{Too fast.} $w_a=3x_0/(1+x_0)^2$ is fixed by $w_0$; at the quoted $w_0=-0.861$ it is $-4.81$ instead of $-0.60$.
\item \textbf{It leaves its domain.} Going back in time $M_{\mathrm{eff}}$ vanishes at $z=0.026$; beyond it the phantom epoch, the negative energy density beyond $z=0.293$ and the bounce at $z=0.388$ are artefacts, and the model gives no matter era and no early universe.
\item \textbf{Negative sound speed squared} today for every attractive model ($-12.39$, $-6.47$, $-1.50$, $-0.667$): in a fluid description, growing fluctuations (not computed).
\item \textbf{The 8-dimensional energy bound} forbids a 3-space observer from inferring $w_{\mathrm{eff}}<-1+(1+w)/\sqrt3$ when $\rho\ge0$, $w\ge-1$ and $H_a,\Theta>0$; the deflation variant needs negative 8-dimensional energy density.
\item \textbf{Newton's constant} would vary about 1800 times faster than the lunar-laser-ranging bound allows in the deflation variant.
\item \textbf{4-dimensional energy conservation} fails in the deflation variant.
\item \textbf{The coupling is outside the approximation.} For the tuned models the contact coupling is so strong wherever the Pauli caveat is negligible that no mass makes the Hartree description self-consistent (Section 11.4, limit 3).
\item \textbf{Only the good sector}, imposed by hand (EXP-5).
\item \textbf{Homogeneous mean field only}: no quantum corrections, exchange terms, stability of the condensate or inhomogeneities.
\item \textbf{No data likelihood}: the distance comparisons use noise-free curves with equal weights and the assumed $\Omega_m=0.305$ (except the $\Omega_m$-free projections of Section 11.9).
\end{enumerate}

\textbf{The answer on dark energy: no, within everything computed.} The only negative pressure comes from an attractive four-fermion interaction. Tuned to the quoted $w_0=-0.861$ it evolves eight times too fast, has a negative sound speed squared, and its effective mass vanishes at $z=0.026$, where the single mean-field mode stops being a positive-energy level. The phantom crossing, the negative energy density at earlier times and the bounce that the equations then give are artefacts of the approximation, not properties of the field. The frozen density of the primordial field mimics a cosmological constant only in its dilution, in a geometry that needs negative energy. The framework reproduces neither the quoted $(w_0,w_a)$ nor the quoted $w=-0.764$ consistently. This ``no'' refers to the mechanisms available in this Lagrangian and these backgrounds; a different potential, inhomogeneous states, other couplings or a mechanism that stabilises the extra dimensions were not examined.

\subsection{What we proved and what we assumed}

\textbf{Proved (derived exactly in this chapter or in the cited stage documents).} The mode equation $i\dot u=hu$ in every homogeneous diagonal background, with $h^2=E^2$, $E^2=M_{\mathrm{eff}}^2+\sum_{j\le3}K_j^2-\sum_{j\ge5}K_j^2$; $h$ is Hermitian exactly in the good sector, where the Hilbert norm is conserved, and the Krein norm is conserved always ($h^\dagger B=Bh$). The per-mode energy $\varepsilon(u)=u^\dagger hu$ and pressure $p_j(u)$ as the images of $T_{44}$ and $T^j{}_j$ under the expectation-value rule (for $\lambda=0$; for $\lambda\ne0$ with the $\tfrac\lambda2S^2$ terms of Section 11.6). The dilution law $SV=$ constant at zero momentum. The curvature of a diagonal metric that depends on $t$ only ($R_{44}$, $R_{ii}$, $R^i{}_i=\dot H_i+H_i\Theta$, $G^4{}_4=-\sum_{i<j}H_iH_j$), hence the EXP-2 equations and the Friedmann equations. The split of $h=mZ+KX$ into eight two-level blocks. For a homogeneous condensate, given Stage 1's verified $T^i{}_i=\mathcal L_s$: $\rho=mS+\tfrac\lambda2S^2$, $p=\tfrac\lambda2S^2$, $\rho+p=2\,\mathrm{KE}_L=SM_{\mathrm{eff}}$ and the phantom criterion. The exact EXP-2 solution $V=1+\Theta_0t+\alpha t^2$, its Kasner exponents with $\sum_ip^{\mathrm K}_i=1$ and $\sum_i(p^{\mathrm K}_i)^2=1-2\kappa mx_0S_0/\mathcal D$, the energy bound $\Theta^2-3H_a^2-2\kappa\rho=H_b^2+3H_c^2$, and the dilution-inferred $w_{\mathrm{eff}}$. The EXP-1 requirement $\rho_{\mathrm{req}}=-3H^2(7+a_4'^2)/\kappa<0$ and $w_{\mathrm{req}}=-(5-a_4'^2)/(7+a_4'^2)$. The closed forms of EXP-3 (zeros, crossing, bounce cubic, tangent CPL parameters, sound speed) and of the deflation variant. No production for $m=0$; the equality of particle and antiparticle weights by unitarity; the WKB exponent of EXP-5. Derived in outline, with its formula confirmed numerically only in the large-$k$ tail: the kink formula of EXP-4b. Derived at Hartree level but not machine-checked: $\rho+p\ge0$ for states with definite occupation numbers.

\textbf{Computed (numbers from the committed reports, verified by 69 self-checks, 167 independent checks, 10 analysis checks and two notebooks).} All the values quoted in Sections 11.6 to 11.12.

\textbf{Assumed.} $U(0)=U'(0)=0$ (no cosmological-constant term and no second mass term), a choice of the model. Homogeneous backgrounds in EXP-2 to EXP-5; in EXP-1 the prescribed warped primordial field, at one fixed $x_0$ for the $\lambda\ne0$ run; the 3-space expansion is prescribed in EXP-3 and EXP-4; only EXP-2 solves Einstein's equations, and without Einstein-Lovelock terms. Static hidden space and extra times in EXP-3 and EXP-4 (not derived; EXP-2 contradicts it for 8-dimensional Einstein gravity). The good sector only (imposed). The one-mode mean field at Hartree level, without the Pauli principle for the condensate, without exchange and with fixed normal ordering; its results are meaningful only while $M_{\mathrm{eff}}>0$. $k_0=0$ in EXP-4. A $C^1$ junction as the end of inflation (plus one smooth comparison). The density parameters $\Omega_r=0.00009$, $\Omega_m=0.305$, $\Omega_\psi=0.69491$. The supernova numbers $(w_0,w_a)=(-0.861,-0.60)$ and $w=-0.764$ as quoted in a private e-mail without a primary reference, without error bars and without a likelihood. No fluctuations, no structure formation and no fit to real data were computed. \textbf{Not claimed:} that dirac16complex is dark matter or dark energy; that the (4,4) signature is the observed spacetime.

\subsection{Exercises}

Units as in the experiments ($m=1$ unless stated). Answers are in Section 11.17.

\textbf{Exercise 11.1.} \emph{Energies of a mode.} For $h=-iM\gamma^4-K_0\gamma^4\gamma^0$ with $M=3$ and $K_0=4$, what are the eigenvalues of $h$ and how many times does each occur? Repeat for $h=-iM\gamma^4-K_5\gamma^4\gamma^5$ with $M=3$, $K_5=4$. Which of the two matrices is Hermitian, and what happens to $u^\dagger u$ in the second case?

\textbf{Exercise 11.2.} \emph{Conservation laws.} Show that a solution of $i\dot u=hu$ keeps $u^\dagger u$ constant when $h^\dagger=h$, and keeps $u^\dagger Bu$ constant when $h^\dagger B=Bh$.

\textbf{Exercise 11.3.} \emph{Condensate equations of state.} With $m=1$ and $S=1$, compute $M_{\mathrm{eff}}$, $\rho$, $p$, $w$, $\mathrm{KE}_L$ and $\mathrm{PE}_L$ for $\lambda=0.5$, $-0.5$, $-1.5$ and $-2.5$. Which cases are phantom, which have $\rho<0$, and which can the one-mode mean field describe?

\textbf{Exercise 11.4.} \emph{An EXP-1 eigenmode.} For $K=0.5$, $\lambda=0$, $S_0=1$, predict $E$, $\rho$, $p_0$, $\bar p$, $w$, $\mathrm{KE}_L$, $\mathrm{PE}_L$ and $s$ of the positive-energy eigenmode.

\textbf{Exercise 11.5.} \emph{The EXP-1 interaction run.} For $\lambda=0.5$, $K=0.5$, $S_0=1$, check that $M_\ast=1.47348266406420$ solves $M_\ast=1+\lambda S_0M_\ast/\sqrt{M_\ast^2+K^2}$, and compute $\rho$, the pressures $p_j$ ($j\ne0$) and $p_0$, $\bar p$ and $w$.

\textbf{Exercise 11.6.} \emph{The primordial source on the plateau.} Where $a_4'=1$ and $a_4''=0$ (with $H=\kappa=1$), compute the seven required pressures, $\rho_{\mathrm{req}}$ and $w_{\mathrm{req}}$.

\textbf{Exercise 11.7.} \emph{EXP-2 by hand.} For $x_0=-0.4$ compute $S_0$, $\lambda$, $\beta$, $\alpha$, the discriminant $\mathcal D$, $t_s$, the Kasner exponents $(p^{\mathrm K}_b,p^{\mathrm K}_a,p^{\mathrm K}_c)$, $\sum_ip^{\mathrm K}_i$ and $\sum_i(p^{\mathrm K}_i)^2$ (both ways), the time at which the extra times stop deflating, and the volumes at which $M_{\mathrm{eff}}=0$ and $\rho=0$.

\textbf{Exercise 11.8.} \emph{A different start.} In EXP-2 with $x_0=0$, change the initial extra-time rate to $H_c=-0.3$. Predict $C_0$, $\Theta_0$, $S_0$, $\beta$, $\alpha$ and the time at which $H_c$ changes sign. (The student guide, its exercise 13.4, shows how to run this modified case.)

\textbf{Exercise 11.9.} \emph{EXP-3 with $x_0=-0.25$.} Predict $w_0$, $w_a$, the redshifts where $\rho_\psi=0$ and where $w=-1$, today's $c_s^2$, today's $q_{\mathrm{dec}}$, and the bounce. Compare with the row $x_0=-0.25$ of \texttt{artifact\allowbreak{}s/dirac1\allowbreak{}6complex\allowbreak{}/numeric\allowbreak{}s/exp3/f\allowbreak{}its\_scan\allowbreak{}.csv}.

\textbf{Exercise 11.10.} \emph{Closer to $-1$.} Which $x_0$ gives $w_0=-0.9$, and what is its tangent $w_a$? What does this say about approaching a cosmological constant with this model?

\textbf{Exercise 11.11.} \emph{The sudden-start wave.} The admixture left by starting a thermal mode on the instantaneous eigenvector has $|c|=mkH/(4E^3)$ with $E=\sqrt{m^2+k^2}$. Find the $k$ at which $|c|$ is largest and the largest value of $|c|^2$ for $H/m=0.05$.

\textbf{Exercise 11.12.} \emph{Rotating eigenvectors.} For $h=E(\cos\theta\,\sigma_z+\sin\theta\,\sigma_x)$, verify that $u_+=(\cos\tfrac\theta2,\sin\tfrac\theta2)$ and $u_-=(-\sin\tfrac\theta2,\cos\tfrac\theta2)$ are eigenvectors with eigenvalues $\pm E$, and that $\dot u_+=\tfrac{\dot\theta}2u_-$. Why does this give no pair creation for $m=0$?

\textbf{Exercise 11.13.} \emph{The kink formula integrated.} Show that $na^3=\frac{16}{2\pi^2}\int_0^\infty k^2\bigl(\frac{mkH^2}{4E^4}\bigr)^2dk=\frac{H^4}{64\pi m}$, using the substitution $k=m\tan\theta$. Compare with the computed yields for $m=0.1$, 1 and 2 ($H=1$).

\textbf{Exercise 11.14.} \emph{The EXP-5 growth.} For $q=0.05$, compute $t_\ast$, $W(t_\ast+2)$ and $W(t_\ast+3)$, the WKB average growth rate of $\ln(u^\dagger u)$ over $[t_\ast+2,t_\ast+3]$, and $e^{2W(t_\ast+3)}$. Compare with the committed results. Why is the WKB rate over $[t_\ast,t_\ast+1]$ less reliable?

\textbf{Exercise 11.15.} \emph{The initial Krein norm in EXP-5.} For $q=0.1$ compute $E_0$ and the Krein norm $u^\dagger Bu$ of the two initial spinors.

\textbf{Exercise 11.16.} \emph{The deflation variant.} From $(w_0,w_a)=(-0.861,-0.60)$ find $x_0$, $n$ and $\gamma$; then $\kappa\rho_8/H_a^2$, $\dot G_N/(G_NH_0)$ and $G_N(z=1)/G_N(0)$.

\textbf{Exercise 11.17.} \emph{What a 3-space observer sees.} Compute $w_{\mathrm{eff}}$ for dust when $\Theta=0$ (EXP-1), $\Theta=3H_a$ (EXP-3 and EXP-4), $\Theta=\sqrt3H_a$, and $\Theta=7H_a$ (late EXP-2).

\textbf{Exercise 11.18.} \emph{The artefact.} Explain in three sentences why the phantom rows of the attractive EXP-2 run are not a prediction of the theory.

\textbf{Exercise 11.19.} \emph{From candidate to model.} List what would have to be computed or derived to turn the dark-matter-like result of EXP-4 into a dark-matter model.

\textbf{Exercise 11.20.} \emph{The dark-energy verdict.} Give two independent reasons, each sufficient by itself, why the tuned condensate of EXP-3 is not viable dark energy.

\subsection{Answers to the exercises}

\textbf{Answer 11.1.} From Section 11.3, $h^2=(M^2+K_0^2)\cdot1=25$, so the eigenvalues are $+5$ and $-5$, eight times each (the Stage-2 verifier evaluates exactly this case). In the second case $E^2=M^2-K_5^2=9-16=-7$: the eigenvalues are $\pm i\sqrt7=\pm2.6458\,i$, eight each. The first matrix is Hermitian; the second is not ($\gamma^4\gamma^5$ is anti-Hermitian). In the second case the part of $u$ along the eigenvalue $+i\sqrt7$ grows like $e^{\sqrt7t}$, so $u^\dagger u$ grows like $e^{2\sqrt7t}$, while $u^\dagger Bu$ stays constant.

\textbf{Answer 11.2.} $\frac{d}{dt}(u^\dagger Xu)=\dot u^\dagger Xu+u^\dagger X\dot u$ with $\dot u=-ihu$ and $\dot u^\dagger=iu^\dagger h^\dagger$ gives $i\,u^\dagger(h^\dagger X-Xh)u$. For $X=1$ this vanishes when $h^\dagger=h$; for $X=B$ when $h^\dagger B=Bh$.

\textbf{Answer 11.3.} With $M_{\mathrm{eff}}=1+\lambda$, $\rho=1+\lambda/2$, $p=\lambda/2$, $\mathrm{KE}_L=M_{\mathrm{eff}}/2$, $\mathrm{PE}_L=1/2$: $\lambda=0.5$: $M_{\mathrm{eff}}=1.5$, $\rho=1.25$, $p=0.25$, $w=0.2$, $\mathrm{KE}_L=0.75$. $\lambda=-0.5$: $0.5$, $0.75$, $-0.25$, $w=-1/3$, $\mathrm{KE}_L=0.25$. $\lambda=-1.5$: $-0.5$, $0.25$, $-0.75$, $w=-3$ (phantom), $\mathrm{KE}_L=-0.25$. $\lambda=-2.5$: $-1.5$, $\rho=-0.25<0$, $p=-1.25$, $w=5$, $\mathrm{KE}_L=-0.75$. $\mathrm{PE}_L=0.5$ in all four. Only the first two have $M_{\mathrm{eff}}>0$; in the other two the occupied rest level has negative energy, so the mean field does not describe a physical state there.

\textbf{Answer 11.4.} $E=\sqrt{1.25}=1.118034=\rho$; $p_0=K^2/E=0.223607$; $\bar p=p_0/7=0.031944$; $w=K^2/(7E^2)=1/35=0.028571$; $\mathrm{KE}_L=\mathrm{PE}_L=E/2=0.559017$; $s=M/E=0.894427$. These are the values of run A1\_K0p5\_pos\_Bp in exp1/summary.json (key \texttt{\detokenize{initial}}), collected under \texttt{freeEige\allowbreak{}nmodes} in numerics-summary.json.

\textbf{Answer 11.5.} $E_\ast=\sqrt{M_\ast^2+0.25}=1.556005$ and $S=s=M_\ast/E_\ast=0.946965$, so $1+0.5\cdot0.946965=1.473483=M_\ast$. Then $\tfrac\lambda2S^2=0.224186$, $\rho=E_\ast-0.224186=1.331819$, $p_j=0.224186$ for the six directions $j\ne0$, $p_0=K^2/E_\ast+0.224186=0.384854$, $\bar p=(0.384854+6\cdot0.224186)/7=0.247138$, $w=0.18556$, $\mathrm{KE}_L=E_\ast/2=0.778002$ and $\mathrm{PE}_L=0.553817$: the committed 1.3318, 0.2471, 0.1856, 0.7780 and 0.5538.

\textbf{Answer 11.6.} $p_{\mathrm{req},0}=-3(1-5)=12$, $p_{\mathrm{req},i}=15-3=12$, $p_{\mathrm{req},j}=12$: isotropic $p=12$; $\rho_{\mathrm{req}}=-3(7+1)=-24$; $w_{\mathrm{req}}=-1/2$. The pressure is positive; $w$ is negative only because $\rho$ is.

\textbf{Answer 11.7.} $S_0=1.32/0.6=2.2$, $\lambda=-0.8/2.2=-0.363636$, $\beta=2.2/6=0.366667$, $\alpha=1.283333$, $\mathcal D=5.76-5.133333=0.626667$, $\sqrt{\mathcal D}=0.791623$, $t_s=-2/3.191623=-0.626640$. $p^{\mathrm K}_b=(0-0.229768)/0.791623=-0.290250$, $p^{\mathrm K}_a=(1-0.229768)/0.791623=0.972978$, $p^{\mathrm K}_c=(-0.2-0.229768)/0.791623=-0.542895$. $\sum_ip^{\mathrm K}_i=-0.290250+3(0.972978)+3(-0.542895)=1.000000$; $\sum_i(p^{\mathrm K}_i)^2=3.808511$, and $1-2(-0.4)(2.2)/0.626667=3.808511$. $H_c=0$ at $t=0.2/\beta=0.545455$. $M_{\mathrm{eff}}=1+2x_0/V=0$ at $V=0.8$; $\rho\propto1+x_0/V=0$ at $V=0.4$. All agree with Section 11.7.

\textbf{Answer 11.8.} $C_0=3+3(0.09)+9(-0.3)=0.57$, $\Theta_0=3-0.9=2.1$, $S_0=0.57$, $\beta=0.095$, $\alpha=7S_0/12=0.3325$, and $H_c=0$ at $t=0.3/0.095=3.158$.

\textbf{Answer 11.9.} $w_0=-0.25/0.75=-1/3$; $w_a=3(-0.25)/0.5625=-4/3$; $\rho_\psi=0$ at $a=0.25^{1/3}=0.629961$, $z=0.587401$; $w=-1$ at $a=0.5^{1/3}=0.793701$, $z=0.259921$; $c_s^2=2x_0/(1+2x_0)=-0.5/0.5=-1$. Since $1+3w_0=0$, the condensate neither accelerates nor decelerates today and $q_{\mathrm{dec}}=(2\Omega_r+\Omega_m)/2=0.15259$. With $A=0.69491/0.75=0.926547$ the bounce cubic $1.231547\,a^3+0.00009\,a^2-0.231637=0$ gives $a_b=0.5729$, $z_b=0.7454$. The committed row gives $-0.3333$, $-1.3333$, $z=0.587401$, $z=0.259921$, $c_s^2=-1$, $q_0=0.15259$ and $z_b=0.745419$. Beyond $z=0.26$ the run would be outside the one-mode picture.

\textbf{Answer 11.10.} $x_0=w_0/(1-w_0)=-0.9/1.9=-0.473684$ and $w_a=3x_0/(1+x_0)^2=-5.13$; the scan row $x_0=-0.474$ gives $w_0=-0.90114$ and $w_a=-5.1396$. The closer $w_0$ is to $-1$, the faster the condensate evolves: the opposite of a slowly varying dark energy near a cosmological constant.

\textbf{Answer 11.11.} Maximise $k/E^3=k(m^2+k^2)^{-3/2}$: the derivative is $(m^2+k^2)^{-5/2}(m^2+k^2-3k^2)$, zero at $k=m/\sqrt2$, where $E^2=3m^2/2$. There $|c|=\frac{H}{4\sqrt2(3/2)^{3/2}m}=\frac{H}{10.39\,m}$, so for $H/m=0.05$ the largest $|c|^2$ is $(0.05/10.39)^2=2.3\times10^{-5}$ (the value recorded in the EXP-4 findings of exp4/summary.json). At $k=0.9525$ the same formula gives $4|c|^2=8.17\times10^{-5}$, and the measured $6.37\times10^{-5}$ is 0.779 of it.

\textbf{Answer 11.12.} $hu_+=E(\cos\theta\cos\tfrac\theta2+\sin\theta\sin\tfrac\theta2,\ \sin\theta\cos\tfrac\theta2-\cos\theta\sin\tfrac\theta2)=E(\cos\tfrac\theta2,\sin\tfrac\theta2)$ by the difference formulas for cosine and sine; similarly $hu_-=-Eu_-$. Differentiating, $\dot u_+=\tfrac{\dot\theta}2(-\sin\tfrac\theta2,\cos\tfrac\theta2)=\tfrac{\dot\theta}2u_-$. For $m=0$, $h=K\sigma_x$, so $\cos\theta=m/E=0$ and $\sin\theta=K/E=1$: $\theta=\pi/2$ for every $K>0$. Then $\dot\theta=0$, the eigenvectors never rotate, and $\dot\beta=0$: no pairs.

\textbf{Answer 11.13.} With $E^2=m^2+k^2$: $na^3=\frac{16}{2\pi^2}\frac{m^2H^4}{16}\int_0^\infty\frac{k^4}{(m^2+k^2)^4}dk$. Put $k=m\tan\theta$, $dk=m\,d\theta/\cos^2\theta$, $m^2+k^2=m^2/\cos^2\theta$: the integral becomes $m^{-3}\int_0^{\pi/2}\sin^4\theta\cos^2\theta\,d\theta=m^{-3}\bigl(\int\sin^4-\int\sin^6\bigr)=m^{-3}\bigl(\tfrac{3\pi}{16}-\tfrac{5\pi}{32}\bigr)=\frac{\pi}{32m^3}$. So $na^3=\frac{m^2H^4}{2\pi^2}\frac{\pi}{32m^3}=\frac{H^4}{64\pi m}$. For $m=1$: $4.97\times10^{-3}$, 1.13 times the computed $4.412\times10^{-3}$; for $m=2$: $2.49\times10^{-3}$, 1.03 times $2.421\times10^{-3}$; for $m=0.1$: $4.97\times10^{-2}$, 34 times the computed $1.454\times10^{-3}$. The formula assumes $E\gg H$, which fails for the small momenta that dominate the light field's production.

\textbf{Answer 11.14.} $t_\ast=\ln20=2.995732$. Since $Q=qe^{t}=e^{t-t_\ast}$: $W(t_\ast+2)=\sqrt{e^4-1}-\arccos(e^{-2})=5.88603$ and $W(t_\ast+3)=\sqrt{e^6-1}-\arccos(e^{-3})=18.53964$. The average rate is $2(18.53964-5.88603)=25.307$, against the committed 25.306. $e^{2W(t_\ast+3)}=e^{37.079}=1.27\times10^{16}$, the same order as the committed final $u^\dagger u=1.258\times10^{16}$; the leading WKB estimate ignores the amplitude factor and the region near $t_\ast$, so only the order of magnitude should be compared. Over $[t_\ast,t_\ast+1]$ WKB gives $2W(t_\ast+1)=2.667$ against the measured 2.37: near the turning point $\varkappa$ changes as fast as the solution itself, which violates the slow-change assumption of WKB.

\textbf{Answer 11.15.} $E_0=\sqrt{1-0.01}=0.994987$; $u^\dagger Bu=\pm E_0/m=\pm0.994987$ (the summary's \texttt{kreinNor\allowbreak{}m0} for runs \texttt{\detokenize{q0p1_Cp}} and \texttt{\detokenize{q0p1_Cm}}).

\textbf{Answer 11.16.} $x_0=-0.861/1.861=-0.462654$; $n=w_a(1+x_0)^2/x_0=0.374457$; $\gamma=1-n/3=0.875181$; $\kappa\rho_8/H_a^2=3(1-3\gamma+\gamma^2)=-2.5788$; $\dot G_N/(G_NH_0)=3\gamma=2.6255$; $G_N(z=1)/G_N(0)=(1+z)^{-3\gamma}=2^{-2.6255}=0.162$.

\textbf{Answer 11.17.} $w_{\mathrm{eff}}=-1+\Theta/(3H_a)$ for dust: $-1$ ($\Theta=0$), 0 ($\Theta=3H_a$), $-1+1/\sqrt3=-0.42265$ ($\Theta=\sqrt3H_a$), $-1+7/3=4/3$ ($\Theta=7H_a$).

\textbf{Answer 11.18.} The phantom rows all have $M_{\mathrm{eff}}<0$, while the occupied rest mode keeps $s(u)=1$ because it only changes its phase. Its energy $u^\dagger hu=M_{\mathrm{eff}}$ is then negative, so the one quantum of the mean-field picture sits in a negative-energy level, outside the domain of the expectation-value rule. A state with definite occupation numbers of positive-energy quanta has $\rho+p\ge0$ and cannot be phantom.

\textbf{Answer 11.19.} Physical values of $m$ and $H_{\mathrm{inf}}$ and the post-inflationary history, giving a relic density to compare with the observed one; a derived mechanism that keeps the hidden space and the extra times static; the couplings to ordinary matter and the resulting bounds (darkness, stability, annihilation through $\lambda$); the treatment of the hidden-space momentum; a justification for excluding the extra-time sector; and the growth of fluctuations and the free streaming that decide whether the relic is cold, warm or hot.

\textbf{Answer 11.20.} For example: (a) matching $w_0=-0.861$ forces $w_a=-4.81$, eight times the quoted $-0.60$, so the model cannot follow the quoted evolution; (b) the model leaves the validity of its own approximation at $z=0.026$ and gives no matter era; (c) its $c_s^2=-12.39$ today signals growing fluctuations in a fluid description; (d) no mass makes the tuned Hartree description self-consistent.

\subsection{Repository files used in this chapter}

The chapter's numbers and figures come from these committed files; Chapter 19 gives the commands that regenerate and check them.

\begin{Verbatim}[fontsize=\small]
provenance/DIRAC16COMPLEX_DARK_SECTOR_NUMERICS.md      the Stage-3 document
provenance/DIRAC16COMPLEX_STUDENT_GUIDE.md             running and checking, step by step
handoff/specs/NUMERICS_CONTRACT.md                     the numerical programme and errata
studies/dirac16complex_cosmology/src/exp1.rs ... exp5.rs
                                                       the five experiments
artifacts/dirac16complex/numerics/exp1/ ... exp5/      CSV, summary.json, checker reports
artifacts/dirac16complex/numerics/exp3/fits.json       EXP-3 fits and estimates
artifacts/dirac16complex/numerics/exp3/fits_scan.csv   the 491-coupling scan
artifacts/dirac16complex/numerics/numerics-summary.json
                                                       all checks, verdict SUCCESS
artifacts/dirac16complex/numerics/figures/             the 17 notebook figures
artifacts/dirac16complex/numerics/figures/mathematica/ the Mathematica figures
scripts/check_dirac16complex_exp1.py ... exp5.py       independent checkers
scripts/analyze_dirac16complex_exp3.py                 fits (numpy Nelder-Mead)
notebooks/dirac16complex_dark_sector.ipynb             Jupyter notebook (71 assertions)
notebooks/Dirac16ComplexDarkSector.nb                  Mathematica notebook (49 checks)
provenance/DIRAC16COMPLEX_PRIMORDIAL_FIELD.md          Stage 2: sources of the field
provenance/DIRAC16COMPLEX_ARBITRARY_FIELD.md           Stage 1: T, splits, quantization
\end{Verbatim}

\section{Many-body quantum mechanics and density functional theory from zero}

\subsection{What this chapter is for}

Chapter 13 computes the lowest-energy state (the \textbf{ground state}) and the first excited state of a finite number $N$ of dirac16complex quanta that sit in the primordial gravitational field of the notebook and interact with each other. The tool is \textbf{density functional theory} (DFT) in the form invented by Kohn and Sham. This chapter builds that tool from zero, for ordinary non-relativistic particles first, because every idea is easier to see there. Nothing in this chapter is specific to dirac16complex; Chapter 13 transfers each step to the 16-component field.

Why is a special tool needed at all? A single quantum particle in three-dimensional space is described by a complex function $\psi(\mathbf r)$ of three variables. If we store it on a grid with 10 points per coordinate we need $10^3$ complex numbers. Two particles need a function of 6 variables, $10^6$ numbers; ten particles need $10^{30}$ numbers, more than any computer can hold. This growth, by a factor $10^3$ for every added particle, is called the \textbf{exponential wall} of the many-body problem. DFT climbs over the wall by working with the \textbf{density} $n(\mathbf r)$, the expected number of particles per unit volume at the point $\mathbf r$: a function of three variables, whatever $N$ is. The price is that one ingredient, the \textbf{exchange--correlation energy}, is not known exactly and must be approximated. We will see exactly where the approximation enters and what it is.

\textbf{Plan.} Sections 12.2 to 12.8 are quantum mechanics of many identical fermions: wave functions, the Pauli principle, Slater determinants, creation and annihilation operators, and the Hartree--Fock approximation, with a small example that can be solved exactly (Section 12.8). Sections 12.9 to 12.13 are density functional theory proper: the Hohenberg--Kohn theorems, functional derivatives, the Kohn--Sham equations, the local density approximation and the numerical self-consistency loop. Section 12.14 extends everything to a finite temperature (Mermin's functional), and Section 12.15 treats excited states: the Kohn--Sham gap, particle--hole excitations and the Delta-SCF method. Section 12.16 lists what was proved and what was assumed; exercises with full answers close the chapter.

\textbf{Units and counting.} In Sections 12.2 to 12.15 we use \textbf{atomic units}: Planck's constant divided by $2\pi$, the electron mass and the electron charge are set to 1 ($\hbar=m_e=e=1$), so the kinetic energy operator of one particle is $-\tfrac12\nabla^2$ and the Coulomb repulsion of two electrons at distance $r$ is $1/r$. As everywhere in this book we count from 0: the $N$ particles are numbered $0,1,\dots,N-1$, and so are the orbitals.

\subsection{One quantum particle: states, operators and the variational principle}

\textbf{States.} The state of one particle is a complex-valued function $\psi(\mathbf r)$, the \textbf{wave function}. Its squared modulus $|\psi(\mathbf r)|^2=\psi^*(\mathbf r)\psi(\mathbf r)$ is the probability per unit volume of finding the particle at $\mathbf r$, so a physical state is \textbf{normalized}: $\int|\psi|^2\,d^3r=1$. Two states are compared with the \textbf{inner product}

\[
\langle\varphi|\psi\rangle=\int\varphi^*(\mathbf r)\,\psi(\mathbf r)\,d^3r ,
\]

which is linear in $\psi$, antilinear in $\varphi$ (a factor $c$ in $\varphi$ comes out as $c^*$), and satisfies $\langle\psi|\varphi\rangle=\langle\varphi|\psi\rangle^*$ and $\langle\psi|\psi\rangle\ge0$. Two states with $\langle\varphi|\psi\rangle=0$ are \textbf{orthogonal}. It often helps to replace space by a finite set of points: then a state is a column of $M$ complex numbers $\psi=(\psi_0,\dots,\psi_{M-1})^T$ and $\langle\varphi|\psi\rangle=\sum_x\varphi_x^*\psi_x=\varphi^\dagger\psi$ (the dagger means complex conjugate and transpose, Chapter 1). Everything below holds in both pictures.

\textbf{Operators.} An observable quantity is represented by a linear \textbf{operator} $A$, a rule that turns a state into a state and satisfies $A(c_0\psi+c_1\varphi)=c_0A\psi+c_1A\varphi$. On a finite grid an operator is an $M\times M$ matrix. The operator is \textbf{Hermitian} (self-adjoint) if $\langle\varphi|A\psi\rangle=\langle A\varphi|\psi\rangle$ for all states; for a matrix this means $A^\dagger=A$. The \textbf{expectation value} of $A$ in the normalized state $\psi$ is $\langle A\rangle=\langle\psi|A\psi\rangle$, the average of many measurements. The energy operator of one particle in an external potential $v(\mathbf r)$ is the \textbf{Hamiltonian}

\[
h=-\tfrac12\nabla^2+v(\mathbf r),\qquad \nabla^2=\partial_x^2+\partial_y^2+\partial_z^2 .
\]

A state with $h\psi=\varepsilon\psi$ is a \textbf{stationary state} (an eigenstate) with energy $\varepsilon$.

\textbf{Two facts about Hermitian operators.} (i) The eigenvalues are real: if $A\psi=a\psi$ with $\psi\ne0$, then $a\langle\psi|\psi\rangle=\langle\psi|A\psi\rangle=\langle A\psi|\psi\rangle=a^*\langle\psi|\psi\rangle$, so $a=a^*$. (ii) Eigenvectors with different eigenvalues are orthogonal: if $A\psi=a\psi$ and $A\varphi=b\varphi$ with $a\ne b$, then $b\langle\psi|\varphi\rangle=\langle\psi|A\varphi\rangle=\langle A\psi|\varphi\rangle=a\langle\psi|\varphi\rangle$, so $\langle\psi|\varphi\rangle=0$. For a Hermitian matrix one can always choose an orthonormal basis of eigenvectors (the spectral theorem of linear algebra, Chapter 1).

\textbf{Outer products and functions of an operator.} We also write $|\psi\rangle$ for the state $\psi$ (a \textbf{ket}). For two states $\varphi,\chi$ the symbol $|\varphi\rangle\langle\chi|$ is the operator $\psi\mapsto\varphi\,\langle\chi|\psi\rangle$; on a grid it is the matrix $\varphi\chi^\dagger$ (a column times a row), and its trace is $\mathrm{Tr}\,|\varphi\rangle\langle\chi|=\chi^\dagger\varphi=\langle\chi|\varphi\rangle$ (we write $\mathrm{Tr}$ for the trace $\mathrm{tr}$ of Chapter 1). A Hermitian $A$ with orthonormal eigenvectors $|i\rangle$ and eigenvalues $a_i$ is $A=\sum_ia_i|i\rangle\langle i|$: both sides give $a_j|j\rangle$ on every eigenvector $|j\rangle$, and the eigenvectors form a basis. A function of $A$ is defined in the same way, $f(A)=\sum_if(a_i)|i\rangle\langle i|$. For $f=\exp$ this agrees with the power series of Chapters 8 and 10, because $A^k=\sum_ia_i^k|i\rangle\langle i|$. For a Hermitian $B$ with eigenvalues $b_i>0$ the logarithm is $\ln B=\sum_i\ln b_i\,|i\rangle\langle i|$; if some $b_i=0$, only the product $B\ln B=\sum_ib_i\ln b_i\,|i\rangle\langle i|$ is used, with the convention $0\ln0=0$ (the limit of $b\ln b$ as $b\to0$), so that $\mathrm{Tr}(B\ln B)=\sum_ib_i\ln b_i$. The same definitions apply to operators on many-particle states (Section 12.14).

\textbf{Worked example.} Two grid points (``two sites'') with the hopping matrix

\[
h=\begin{pmatrix}0&-t\\-t&0\end{pmatrix},\qquad t>0 .
\]

The characteristic polynomial is $\det(h-\varepsilon)=\varepsilon^2-t^2$, so $\varepsilon=\pm t$. For $\varepsilon=-t$ the equation $(h+t)\psi=0$ gives $\psi_0=\psi_1$, normalized $\psi_b=(1,1)^T/\sqrt2$ (the \textbf{bonding} orbital); for $\varepsilon=+t$ it gives $\psi_a=(1,-1)^T/\sqrt2$ (antibonding). Check: $\psi_b^\dagger\psi_a=(1-1)/2=0$.

\textbf{The variational principle.} Let $\psi_0,\psi_1,\dots$ be an orthonormal basis of eigenvectors of a Hermitian $H$ with eigenvalues $E_0\le E_1\le\dots$. Every normalized state is $\psi=\sum_nc_n\psi_n$ with $\sum_n|c_n|^2=1$, and

\[
\langle\psi|H\psi\rangle=\sum_n|c_n|^2E_n\ \ge\ E_0\sum_n|c_n|^2=E_0 ,
\]

with equality exactly when only the $c_n$ with $E_n=E_0$ are nonzero. So \textbf{the ground-state energy is the minimum of the energy expectation value over all normalized states}, and the minimizers are the ground states. Every method of this chapter is an application of this one inequality.

\textbf{Internal states.} Real particles carry an internal label besides their position. For an electron it is the spin $\sigma\in\{\uparrow,\downarrow\}$, two values. We collect position and label into one symbol $x=(\mathbf r,\sigma)$ and write $\int dx=\sum_\sigma\int d^3r$. A particle with $g$ internal states has a wave function $\psi(\mathbf r,\sigma)$ with $\sigma=0,\dots,g-1$. For the quanta of dirac16complex in the sector studied in Chapter 13 the analogous number is $g=8$: eight states of positive energy for each momentum (Chapter 8).

\subsection{Many identical fermions and the Pauli principle}

\textbf{The many-particle wave function.} $N$ particles are described by one complex function of all their coordinates, $\Psi(x_0,x_1,\dots,x_{N-1})$. The number $|\Psi(x_0,\dots,x_{N-1})|^2$ is the probability density of finding particle 0 at $x_0$, particle 1 at $x_1$, and so on, and $\int|\Psi|^2\,dx_0\cdots dx_{N-1}=1$.

\textbf{Identical particles.} All electrons are identical: no experiment can tell ``electron 0'' from ``electron 1''. Hence $|\Psi|^2$ must be unchanged when two arguments are swapped. Nature realizes this in two ways. For \textbf{bosons} $\Psi$ itself is unchanged; for \textbf{fermions} it changes sign:

\[
\Psi(\dots,x_i,\dots,x_j,\dots)=-\Psi(\dots,x_j,\dots,x_i,\dots)\qquad(\text{fermions}).
\]

Electrons, protons and neutrons are fermions, and so are the quanta of dirac16complex: its components are anticommuting Grassmann fields (Chapters 5 and 8), and anticommuting fields describe fermions. Which particles are fermions is an input from experiment and from the spin--statistics connection of relativistic quantum field theory; in this book it is an \textbf{assumption about the field}, built into dirac16complex by choosing Grassmann components.

\textbf{The Pauli principle.} Put $x_i=x_j=x$ in the antisymmetry rule: $\Psi(\dots,x,\dots,x,\dots)=-\Psi(\dots,x,\dots,x,\dots)$, so $\Psi=0$ there. Two identical fermions are never found with the same position and the same internal label.

\textbf{The many-body Hamiltonian.} For $N$ particles in an external potential $v$ with a pair interaction $w$,

\[
\hat H=\sum_{i=0}^{N-1}\Bigl[-\tfrac12\nabla_i^2+v(\mathbf r_i)\Bigr]+\sum_{0\le i<j\le N-1}w(\mathbf r_i,\mathbf r_j),
\]

where $\nabla_i$ acts on the coordinates of particle $i$. For electrons $w=1/|\mathbf r_i-\mathbf r_j|$. We write $\hat H=\hat T+\hat V+\hat W$ for the three parts (kinetic, external, interaction). Only $\hat V$ distinguishes one system from another: atoms, molecules and solids with $N$ electrons all share the same $\hat T+\hat W$. This observation is the seed of DFT.

\textbf{The density.} The density is the expected number of particles per unit volume at $\mathbf r$:

\[
n(\mathbf r)=N\sum_\sigma\int|\Psi(\mathbf r\sigma,x_1,\dots,x_{N-1})|^2\,dx_1\cdots dx_{N-1}.
\]

The factor $N$ appears because any of the $N$ identical particles may be the one at $\mathbf r$, and antisymmetry makes all $N$ choices give the same integral. Integrating over $\mathbf r$ gives $\int n\,d^3r=N$. The expectation value of the external energy needs only the density: $\langle\Psi|\hat V\Psi\rangle=\int v(\mathbf r)\,n(\mathbf r)\,d^3r$, because each term $v(\mathbf r_i)$ contributes the same integral, $N$ times.

\subsection{Slater determinants}

\textbf{Two particles.} Take two orthonormal one-particle states (\textbf{orbitals}) $\varphi_a$ and $\varphi_b$. The product $\varphi_a(x_0)\varphi_b(x_1)$ is not antisymmetric, but the combination

\[
\Phi(x_0,x_1)=\frac1{\sqrt2}\bigl[\varphi_a(x_0)\varphi_b(x_1)-\varphi_b(x_0)\varphi_a(x_1)\bigr]
=\frac1{\sqrt2}\det\begin{pmatrix}\varphi_a(x_0)&\varphi_b(x_0)\\ \varphi_a(x_1)&\varphi_b(x_1)\end{pmatrix}
\]

is. Its norm is $\tfrac12[1+1-2\,\mathrm{Re}(\langle\varphi_a|\varphi_b\rangle\langle\varphi_b|\varphi_a\rangle)]=1$ for orthonormal orbitals, and it vanishes identically if $\varphi_a=\varphi_b$: two fermions cannot occupy the same orbital. This is the form of the Pauli principle that chemists and solid-state physicists use.

\textbf{$N$ particles.} For $N$ orthonormal orbitals $\varphi_0,\dots,\varphi_{N-1}$ the \textbf{Slater determinant} is

\[
\Phi(x_0,\dots,x_{N-1})=\frac1{\sqrt{N!}}\det\bigl[\varphi_a(x_i)\bigr]_{i,a=0}^{N-1}
=\frac1{\sqrt{N!}}\sum_{P}\mathrm{sgn}(P)\prod_{i=0}^{N-1}\varphi_{P(i)}(x_i),
\]

where the sum runs over the $N!$ permutations $P$ of $\{0,\dots,N-1\}$ and $\mathrm{sgn}(P)=\pm1$ is the sign of the permutation. Three properties follow from the rules for determinants. (i) Swapping two arguments $x_i,x_j$ swaps two rows, which changes the sign: $\Phi$ is antisymmetric. (ii) Two equal orbitals make two columns equal, and the determinant vanishes. (iii) $\Phi$ is normalized. For (iii) write $|\Phi|^2$ as a double sum over permutations $P,P'$ and integrate: the integral of $\prod_i\varphi^*_{P'(i)}(x_i)\varphi_{P(i)}(x_i)$ is $\prod_i\langle\varphi_{P'(i)}|\varphi_{P(i)}\rangle$, which is 1 if $P'=P$ and 0 otherwise (at least one factor is an inner product of two different orthonormal orbitals). So $\int|\Phi|^2=\frac1{N!}\sum_P1=1$.

\textbf{The density of a Slater determinant} is the sum of the orbital densities,

\[
n(x)=\sum_{a=0}^{N-1}|\varphi_a(x)|^2 .
\]

(Insert the permutation sum into the definition of $n$ and integrate over $x_1,\dots,x_{N-1}$: by orthonormality only $P'=P$ survives, and the remaining factor for the variable $x_0$ is $|\varphi_{P(0)}(x_0)|^2$; each orbital $a$ appears as $P(0)$ in $(N-1)!$ permutations, and $N\cdot(N-1)!/N!=1$.)

\textbf{Worked example.} Let ``space'' consist of three points $x\in\{0,1,2\}$ (no spin) and take the orthonormal orbitals $\varphi_0=(1,0,0)^T$ and $\varphi_1=(0,1,1)^T/\sqrt2$. The two-particle Slater determinant has the values

\begingroup
\small
\begin{longtable}{@{}>{\raggedright\arraybackslash}p{0.133\linewidth}>{\raggedright\arraybackslash}p{0.133\linewidth}>{\raggedright\arraybackslash}p{0.133\linewidth}>{\raggedright\arraybackslash}p{0.133\linewidth}>{\raggedright\arraybackslash}p{0.133\linewidth}>{\raggedright\arraybackslash}p{0.133\linewidth}@{}}
\toprule
\textbf{$(x_0,x_1)$} & \textbf{$(0,1)$} & \textbf{$(0,2)$} & \textbf{$(1,2)$} & \textbf{$(1,0)$} & \textbf{$(2,0)$} \\
\midrule
\endfirsthead
\toprule
\textbf{$(x_0,x_1)$} & \textbf{$(0,1)$} & \textbf{$(0,2)$} & \textbf{$(1,2)$} & \textbf{$(1,0)$} & \textbf{$(2,0)$} \\
\midrule
\endhead
$\Phi(x_0,x_1)$ & $1/2$ & $1/2$ & $0$ & $-1/2$ & $-1/2$ \\
\bottomrule
\end{longtable}
\endgroup

and $\Phi(x,x)=0$ on the diagonal. For example $\Phi(0,1)=\tfrac1{\sqrt2}[\varphi_0(0)\varphi_1(1)-\varphi_1(0)\varphi_0(1)]=\tfrac1{\sqrt2}\cdot\tfrac1{\sqrt2}=\tfrac12$, and $\Phi(1,2)=\tfrac1{\sqrt2}[0\cdot\tfrac1{\sqrt2}-\tfrac1{\sqrt2}\cdot0]=0$. The sum of $|\Phi|^2$ over all nine ordered pairs is $4\cdot\tfrac14=1$, and the density $n(x)=N\sum_{x_1}|\Phi(x,x_1)|^2$ is $n=(2\cdot\tfrac12,\ 2\cdot\tfrac14,\ 2\cdot\tfrac14)=(1,\tfrac12,\tfrac12)$, which equals $|\varphi_0|^2+|\varphi_1|^2$ as it must, with total 2.

A general antisymmetric $N$-particle state is not one Slater determinant but a linear combination of many (all determinants built from a complete orbital basis form a basis of the antisymmetric states). A single determinant is the simplest possible many-fermion state, the state of $N$ \textbf{independent} fermions. The whole of Kohn--Sham theory rests on the fact that a single determinant can nevertheless carry the exact density.

\subsection{Creation and annihilation operators}

Slater determinants are cumbersome to write. A shorter bookkeeping, called \textbf{second quantization}, labels a determinant only by which orbitals are occupied. It is the same formalism that Chapter 8 uses for the quantized dirac16complex field, here in its simplest form.

\textbf{Occupation numbers.} Fix an orthonormal basis of orbitals $\varphi_0,\varphi_1,\dots,\varphi_{M-1}$ (finite, to keep things simple). A Slater determinant built from some of them is fixed, up to sign, by the list of \textbf{occupation numbers} $n_p\in\{0,1\}$ ($p=0,\dots,M-1$): $n_p=1$ if $\varphi_p$ is used. We write it $|n_0n_1\cdots n_{M-1}\rangle$ and fix the sign by listing the occupied orbitals in increasing order of $p$ as the columns of the determinant. The state with no particle at all is the \textbf{vacuum} $|0\rangle=|00\cdots0\rangle$.

\textbf{The operators.} The \textbf{creation operator} $a_p^\dagger$ adds a particle in orbital $p$ and the \textbf{annihilation operator} $a_p$ removes one:

\[
\begin{aligned}
a_p^\dagger|\cdots n_p\cdots\rangle&=(-1)^{\nu_p}\,(1-n_p)\,|\cdots n_p{+}1\cdots\rangle,\\
a_p|\cdots n_p\cdots\rangle&=(-1)^{\nu_p}\,n_p\,|\cdots n_p{-}1\cdots\rangle,\qquad \nu_p=\sum_{q<p}n_q .
\end{aligned}
\]

The factor $(1-n_p)$ makes $a_p^\dagger$ give zero on an occupied orbital (Pauli), the factor $n_p$ makes $a_p$ give zero on an empty one, and the sign $(-1)^{\nu_p}$ is the sign of the determinant when the new column $p$ is moved past the $\nu_p$ occupied columns in front of it. From these definitions follow the \textbf{anticommutation relations}

\[
\{a_p,a_q^\dagger\}=a_pa_q^\dagger+a_q^\dagger a_p=\delta_{pq},\qquad \{a_p,a_q\}=\{a_p^\dagger,a_q^\dagger\}=0 .
\]

Proof for $p=q$: on a state with $n_p=0$, $a_pa_p^\dagger$ gives the state back (the two signs are equal and multiply to $+1$) and $a_p^\dagger a_p$ gives 0; with $n_p=1$ it is the other way round; in both cases the sum is 1. For $p<q$: $a_p$ or $a_p^\dagger$ changes $n_p$, which changes the sign factor $(-1)^{\nu_q}$ of the other operator, so the two orders differ exactly by a sign and the anticommutator vanishes. The same argument gives $\{a_p,a_q\}=0$; in particular $a_p^\dagger a_p^\dagger=0$, the Pauli principle again. The relations $\{a_p,a_q\}=\{a_p^\dagger,a_q^\dagger\}=0$ (in particular $(a_p^\dagger)^2=0$) are the rule of Grassmann numbers (Chapter 5); the remaining relation $\{a_p,a_q^\dagger\}=\delta_{pq}$ has no Grassmann analogue: it is what canonical quantization of an anticommuting field produces (Chapter 8). This is why an anticommuting field describes fermions.

The \textbf{number operator} $\hat n_p=a_p^\dagger a_p$ gives $n_p$ on $|\cdots n_p\cdots\rangle$, and $\hat N=\sum_p\hat n_p$ counts all particles. A determinant with the occupied set $O=\{p_0<p_1<\dots<p_{N-1}\}$ is $|\Phi\rangle=a_{p_0}^\dagger a_{p_1}^\dagger\cdots a_{p_{N-1}}^\dagger|0\rangle$.

\textbf{Operators of the many-body theory.} A one-body operator $\sum_iA(i)$, one copy of the one-particle operator $A$ for each particle, becomes

\[
\hat A=\sum_{p,q}A_{pq}\,a_p^\dagger a_q,\qquad A_{pq}=\langle\varphi_p|A\varphi_q\rangle .
\]

Reason: applied to a determinant, $\sum_iA(i)$ replaces one occupied orbital $\varphi_q$ at a time by $A\varphi_q=\sum_p\varphi_pA_{pq}$ and adds the results; ``remove $q$, put $p$, with weight $A_{pq}$'' is exactly what $A_{pq}a_p^\dagger a_q$ does, including the sign. The pair interaction $\sum_{i<j}w(i,j)$ is represented by

\[
\begin{aligned}
&\hat W=\tfrac12\sum_{p,q,r,s}w_{pqrs}\,a_p^\dagger a_q^\dagger a_sa_r,\\
&w_{pqrs}=\int\!\!\int\varphi_p^*(x)\,\varphi_q^*(x')\,w(x,x')\,\varphi_r(x)\,\varphi_s(x')\,dx\,dx' :
\end{aligned}
\]

the operator removes an occupied pair $(r,s)$ and puts $(p,q)$ in its place; the factor $\tfrac12$ compensates for counting each unordered pair twice, and the order $a_sa_r$ (with $s$ and $r$ reversed) makes the direct term come with a plus sign: $a_sa_r$ removes $r$ first and then $s$, and $a_p^\dagger a_q^\dagger$ puts back $q$ first and then $p$, the mirror order, so for $p=r$, $q=s$ the two sign factors cancel. This is a sketch; what we actually use is that $\hat W$ gives the correct energy in every determinant, which Section 12.6 derives directly from the permutation sum and checks against the expectation values below.

\textbf{Expectation values in a determinant.} Let $|\Phi\rangle$ be the determinant with occupied set $O$ in the orbital basis. Then

\[
\begin{aligned}
&\langle\Phi|a_p^\dagger a_q|\Phi\rangle=\delta_{pq}\,[p\in O],\\
&\langle\Phi|a_p^\dagger a_q^\dagger a_sa_r|\Phi\rangle=[p\in O]\,[q\in O]\,\bigl(\delta_{pr}\delta_{qs}-\delta_{ps}\delta_{qr}\bigr),
\end{aligned}
\]

where $[p\in O]$ is 1 if $p$ is occupied and 0 otherwise. The first formula holds because $a_q$ removes $q$ and $a_p^\dagger$ must put back the same orbital to return to $|\Phi\rangle$ (states with different occupations are orthogonal), and then $a_p^\dagger a_p=\hat n_p$. For the second, the pair removed and the pair added must coincide, $\{p,q\}=\{r,s\}$ with $p\ne q$. If $p=r$ and $q=s$, then $a_p^\dagger a_q^\dagger a_qa_p=a_p^\dagger\hat n_qa_p=\hat n_q\hat n_p$ (for $q\ne p$ the operator $\hat n_q$ commutes with $a_p$ and $a_p^\dagger$), which gives $[p\in O][q\in O]$. If $p=s$ and $q=r$, one anticommutation, $a_pa_q=-a_qa_p$, gives $-\hat n_p\hat n_q$. Both formulas are summarized by the \textbf{one-body density matrix} $\rho=\sum_{a\in O}|\varphi_a\rangle\langle\varphi_a|$, whose matrix elements are $\rho_{qp}=\langle\Phi|a_p^\dagger a_q|\Phi\rangle$:

\[
\langle a_p^\dagger a_q\rangle=\rho_{qp},\qquad
\langle a_p^\dagger a_q^\dagger a_sa_r\rangle=\rho_{rp}\,\rho_{sq}-\rho_{sp}\,\rho_{rq}.
\]

Both sides of these identities change in the same way under a unitary change of the orbital basis, so they hold in every orthonormal basis, not only in the one in which $\Phi$ is a single determinant. This is \textbf{Wick's theorem} for a Slater determinant: every expectation value is a sum of products of the density matrix. The same statement holds for a thermal ensemble of non-interacting fermions, with $\rho=\sum_af_a|\varphi_a\rangle\langle\varphi_a|$ and occupation probabilities $0\le f_a\le1$: in such an ensemble the occupations of different orbitals are independent (Section 12.14 derives this), so $\langle\hat n_p\hat n_q\rangle=f_pf_q$ for $p\ne q$, and the same two-line proof applies. In real space $\rho(x,x')=\sum_af_a\varphi_a(x)\varphi_a^*(x')$, and its diagonal $\rho(x,x)$ is the density $n(x)$.

\subsection{The energy of a Slater determinant: direct and exchange terms}

\textbf{The energy from the permutation sum.} We compute the energy of a determinant directly from its definition, as we computed its density in Section 12.4. For the pair interaction,

\[
\begin{aligned}
&\Bigl\langle\Phi\Bigm|\sum_{i<j}w(x_i,x_j)\Bigm|\Phi\Bigr\rangle\\
&\quad=\frac1{N!}\sum_{P,P'}\mathrm{sgn}(P)\,\mathrm{sgn}(P')\sum_{i<j}\int\prod_k\varphi^*_{P'(k)}(x_k)\;w(x_i,x_j)\prod_k\varphi_{P(k)}(x_k)\;dx_0\cdots dx_{N-1}.
\end{aligned}
\]

In the term $(i,j)$ the integral over each $x_k$ with $k\ne i,j$ is $\langle\varphi_{P'(k)}|\varphi_{P(k)}\rangle$, which vanishes unless $P'(k)=P(k)$. Then $P'$ and $P$ agree on every number except $i$ and $j$, so either $P'=P$, or $P'=P\circ(ij)$, the permutation with $P'(i)=P(j)$ and $P'(j)=P(i)$, whose sign is $\mathrm{sgn}(P')=-\mathrm{sgn}(P)$. With $a=P(i)$ and $b=P(j)$ the first case gives $w_{abab}$ and the second gives $-w_{baab}$; renaming $x\leftrightarrow x'$ in the integral and using $w(x,x')=w(x',x)$ shows $w_{baab}=w_{abba}$. Each ordered pair $(a,b)$ of different occupied orbitals arises from the $N(N-1)/2$ pairs $i<j$, each with the $(N-2)!$ permutations that have $P(i)=a$ and $P(j)=b$, and $\frac1{N!}\cdot\frac{N(N-1)}2\cdot(N-2)!=\tfrac12$. Terms with $a=b$ would contribute $w_{aaaa}-w_{aaaa}=0$, so they may be included. The one-body part is simpler: in the term $i$ the integrals over all $x_k$ with $k\ne i$ force $P'=P$, and each orbital $a$ appears as $P(i)$ in $(N-1)!$ permutations for each of the $N$ values of $i$, with $N\cdot(N-1)!/N!=1$. Together,

\[
\langle\Phi|\hat H|\Phi\rangle=\sum_{a\in O}\langle\varphi_a|h\varphi_a\rangle
+\tfrac12\sum_{a,b\in O}\bigl(w_{abab}-w_{abba}\bigr),\qquad h=-\tfrac12\nabla^2+v .
\]

\textbf{The same result from second quantization.} Wick's theorem of Section 12.5 gives it in two lines: $\langle\Phi|\hat A|\Phi\rangle=\sum_{p,q}A_{pq}\,\delta_{pq}[p\in O]=\sum_{a\in O}A_{aa}$, and $\langle\Phi|\hat W|\Phi\rangle=\tfrac12\sum_{p,q,r,s}w_{pqrs}[p\in O][q\in O](\delta_{pr}\delta_{qs}-\delta_{ps}\delta_{qr})=\tfrac12\sum_{a,b\in O}(w_{abab}-w_{abba})$. This is the check announced in Section 12.5: in every determinant the operators $\hat A$ and $\hat W$ give the energies of $\sum_iA(i)$ and $\sum_{i<j}w(i,j)$.

The first sum is the kinetic plus external energy of independent particles. The second sum has two parts. Writing out the integrals and using $\sum_a\varphi_a(x)\varphi_a^*(x')=\rho(x,x')$,

\[
\begin{aligned}
E_H&=\tfrac12\sum_{a,b}w_{abab}=\tfrac12\int\!\!\int n(x)\,w(x,x')\,n(x')\,dx\,dx',\\
E_x&=-\tfrac12\sum_{a,b}w_{abba}=-\tfrac12\int\!\!\int|\rho(x,x')|^2\,w(x,x')\,dx\,dx' .
\end{aligned}
\]

$E_H$ is the \textbf{Hartree energy}: the classical electrostatic energy of the charge cloud $n$ with itself. $E_x$ is the \textbf{exchange energy} (Fock term). It has no classical analogue; it comes from the antisymmetry of $\Phi$ and lowers the energy for a repulsive $w$. Three remarks.

\begin{enumerate}
\item \textbf{No self-interaction.} In $E_H$ the terms $a=b$ describe each particle repelling its own charge cloud, which is unphysical. In $E_x$ the terms $a=b$ are exactly the same numbers with a minus sign, so they cancel.
\item \textbf{Exchange acts only between equal internal labels.} If every orbital has a definite spin, then $\rho(\mathbf r\sigma,\mathbf r'\sigma')=0$ for $\sigma\ne\sigma'$, and $E_x$ couples only particles with the same spin (for a spin-independent $w$).
\item \textbf{Double counting.} Adding the orbital energies of Section 12.7 counts every pair twice, which is why $E_H$ and $E_x$ will be subtracted once from $\sum_a\varepsilon_a$.
\end{enumerate}

\textbf{Contact interaction.} Let the particles have $g$ internal states and interact only when they are at the same point, independently of the label: $w(x,x')=g_c\,\delta(\mathbf r-\mathbf r')$ with a strength $g_c$. Then, with $n_\sigma(\mathbf r)=\rho(\mathbf r\sigma,\mathbf r\sigma)$ and orbitals of definite label,

\[
E_H=\frac{g_c}2\int n(\mathbf r)^2\,d^3r,\qquad E_x=-\frac{g_c}2\int\sum_{\sigma}n_\sigma(\mathbf r)^2\,d^3r .
\]

Two facts follow. First, the exchange energy of a contact interaction is \textbf{exactly local}: it is an integral over one point of a function of the densities at that point, although the general exchange integral involves two points. Second, if all $g$ labels are equally occupied, $n_\sigma=n/g$, then $\sum_\sigma n_\sigma^2=n^2/g$ and

\[
E_x=-\frac1g\,E_H\qquad(\text{contact interaction, equal occupation of the }g\text{ labels}).
\]

For electrons ($g=2$) this says $E_H+E_x=g_c\int n_\uparrow n_\downarrow\,d^3r$: two electrons of the same spin never meet at one point (Pauli), so only opposite spins feel a contact force, and exchange removes precisely the unphysical same-spin part of $E_H$. Chapter 13 finds the ratio $1/8$ for a filled shell of dirac16complex quanta, the case $g=8$ with a matrix-valued vertex.

\subsection{The Hartree and Hartree--Fock equations}

The \textbf{Hartree--Fock approximation} takes the best single determinant: it minimizes $\langle\Phi|\hat H|\Phi\rangle$ over all choices of $N$ orthonormal orbitals. By the variational principle (Section 12.2) the result $E_{HF}$ is an upper bound on the true ground-state energy $E_0$.

\textbf{Minimizing under a side condition: Lagrange multipliers.} The orbitals are not free: they must stay orthonormal. The standard tool for such problems, used again in Sections 12.10, 12.11 and 12.14, is the method of \textbf{Lagrange multipliers}. Take a real function $f(y_0,\dots,y_{m-1})$ of $m$ real variables and one side condition (a \textbf{constraint}) $c(y)=c_0$, and let $y^\ast$ be a minimum of $f$ among the points that satisfy it. To first order (Section 1.9) a small step $dy$ changes $f$ by $df=\sum_i(\partial f/\partial y_i)\,dy_i$ and $c$ by $dc=\sum_i(\partial c/\partial y_i)\,dy_i$. An allowed small step keeps $c=c_0$, so it satisfies $\sum_i(\partial c/\partial y_i)\,dy_i=0$ (this describes the allowed steps correctly where the gradient of $c$, the list of the $\partial c/\partial y_i$, is not zero; with several constraints, where their gradients are linearly independent; that such first-order steps can be completed to steps that keep the constraint exactly is the implicit function theorem of analysis, used without proof). At the minimum $df=0$ for every allowed step, since otherwise the step or the opposite step would lower $f$. Now define

\[
\lambda=\frac{\sum_i(\partial f/\partial y_i)(\partial c/\partial y_i)}{\sum_i(\partial c/\partial y_i)^2},\qquad r_i=\frac{\partial f}{\partial y_i}-\lambda\,\frac{\partial c}{\partial y_i} .
\]

Then $\sum_ir_i\,\partial c/\partial y_i=0$, so $dy_i=\epsilon r_i$ is an allowed step, and $0=df=\epsilon\sum_i(\partial f/\partial y_i)\,r_i=\epsilon\sum_i(r_i+\lambda\,\partial c/\partial y_i)\,r_i=\epsilon\sum_ir_i^2$ forces every $r_i=0$. Hence

\[
\frac{\partial f}{\partial y_i}=\lambda\,\frac{\partial c}{\partial y_i}\quad(\text{all }i),\qquad c(y)=c_0 .
\]

These are exactly the conditions that the \textbf{Lagrangian} $\mathcal L(y,\lambda)=f(y)-\lambda\,\bigl(c(y)-c_0\bigr)$ be stationary with respect to every $y_i$ and to $\lambda$: the constrained problem has become an unconstrained one with one more unknown, the \textbf{multiplier} $\lambda$. With several constraints $c_k(y)=c_{k,0}$ one takes one multiplier for each, $\mathcal L=f-\sum_k\lambda_k(c_k-c_{k,0})$, and the same argument, with $r$ the part of the gradient of $f$ that is orthogonal to all the gradients of the $c_k$, gives $\partial f/\partial y_i=\sum_k\lambda_k\,\partial c_k/\partial y_i$. A complex variable counts as two real ones (next paragraph), and a function of the values on a grid becomes a functional in the limit of a fine grid (Section 12.10).

\emph{Example.} Minimize $f=y_0+y_1$ on the circle $y_0^2+y_1^2=1$. The conditions are $1=2\lambda y_0$ and $1=2\lambda y_1$, so $y_0=y_1=1/(2\lambda)$, and the constraint gives $2/(4\lambda^2)=1$, $\lambda=\pm1/\sqrt2$. The two stationary points $y_0=y_1=\pm1/\sqrt2$ are the maximum $f=\sqrt2$ and the minimum $f=-\sqrt2$ (with $\lambda=-1/\sqrt2$): the method finds all stationary points, and comparing their values picks the minimum.

\emph{What the multiplier means.} Let $f^\ast(c_0)$ be the minimum for the constraint value $c_0$, attained at $y^\ast(c_0)$ with the multiplier $\lambda(c_0)$, and suppose that $y^\ast$ depends differentiably on $c_0$. By the chain rule and the conditions above,

\[
\frac{df^\ast}{dc_0}=\sum_i\frac{\partial f}{\partial y_i}\,\frac{dy^\ast_i}{dc_0}=\lambda\sum_i\frac{\partial c}{\partial y_i}\,\frac{dy^\ast_i}{dc_0}=\lambda\,\frac{d}{dc_0}\,c\bigl(y^\ast(c_0)\bigr)=\lambda ,
\]

because $c(y^\ast(c_0))=c_0$ for every $c_0$. The multiplier is the rate at which the constrained minimum changes with the constraint value. In the example the minimum on the circle $y_0^2+y_1^2=c_0$ is $f^\ast=-\sqrt{2c_0}$, whose derivative at $c_0=1$ is $-1/\sqrt2=\lambda$.

\textbf{How to vary a complex function.} A complex function $\varphi=u+iv$ has two real parts. Instead of varying $u$ and $v$ we may vary $\varphi$ and $\varphi^*$ as if they were independent, because $u=(\varphi+\varphi^*)/2$ and $v=(\varphi-\varphi^*)/(2i)$ are recovered from them; setting the derivative with respect to $\varphi^*$ to zero is equivalent to setting both real derivatives to zero (for a real-valued function of $\varphi$, the derivative with respect to $\varphi$ is the complex conjugate of the one with respect to $\varphi^*$). The precise definition of the derivative of a functional with respect to a function is given in Section 12.10; here we only need the rule ``differentiate the integrand''.

\textbf{The Lagrangian.} The orbitals must stay orthonormal, so we add Lagrange multipliers $\Lambda_{ba}$ for the constraints $\langle\varphi_a|\varphi_b\rangle=\delta_{ab}$:

\[
\mathcal L=\sum_a\langle\varphi_a|h\varphi_a\rangle+\tfrac12\sum_{a,b}\bigl(w_{abab}-w_{abba}\bigr)-\sum_{a,b}\Lambda_{ba}\bigl(\langle\varphi_a|\varphi_b\rangle-\delta_{ab}\bigr).
\]

Differentiate with respect to $\varphi_a^*(x)$. The orbital $a$ appears in the first and in the second slot of $w_{abab}$ and $w_{abba}$; because $w(x,x')=w(x',x)$ the two contributions are equal, which cancels the $\tfrac12$. The result is

\[
\hat F\varphi_a=\sum_b\Lambda_{ba}\varphi_b,\qquad
(\hat F\varphi)(x)=h\varphi(x)+v_H(x)\,\varphi(x)-\int\rho(x,x')\,w(x,x')\,\varphi(x')\,dx',
\]

with the \textbf{Hartree potential} $v_H(x)=\int w(x,x')\,n(x')\,dx'$. $\hat F$ is the \textbf{Fock operator}. It depends on the occupied orbitals only through $\rho$, and $\rho$ does not change if the occupied orbitals are mixed among themselves by a unitary matrix $Q$ ($\varphi'_a=\sum_bQ_{ba}\varphi_b$), while the determinant only picks up the phase $\det Q$. The matrix $\Lambda_{ba}=\langle\varphi_b|\hat F\varphi_a\rangle$ is Hermitian, so we can choose $Q$ to diagonalize it. This gives the \textbf{canonical Hartree--Fock equations}

\[
\hat F\varphi_a=\varepsilon_a\varphi_a,\qquad a=0,\dots,N-1 .
\]

They look like one-particle Schrödinger equations, but $\hat F$ contains the unknown orbitals: the equations are \textbf{nonlinear} and are solved by iteration, starting from a guess and repeating until the orbitals no longer change (\textbf{self-consistency}, Section 12.13). For the ground state one usually occupies the $N$ eigenstates of $\hat F$ with the lowest $\varepsilon_a$ (the \textbf{aufbau} rule, from the German for ``building up''). The older \textbf{Hartree approximation} keeps only $v_H$ and drops the exchange integral; it therefore keeps the self-interaction of remark 1 of Section 12.6.

\textbf{Total energy and double counting.} Summing the orbital energies, $\sum_a\varepsilon_a=\sum_a\langle\varphi_a|\hat F\varphi_a\rangle=\sum_a\langle\varphi_a|h\varphi_a\rangle+2E_H+2E_x$, so

\[
E_{HF}=\sum_{a\in O}\varepsilon_a-E_H-E_x .
\]

\textbf{Koopmans' theorem.} Remove the particle in orbital $a$ and keep the other orbitals frozen. The energy lost is $\langle\varphi_a|h\varphi_a\rangle$ plus all pair terms that contain $a$, $\sum_b(w_{abab}-w_{abba})$ (the ordered pairs $(a,b)$ and $(b,a)$ each carry the factor $\tfrac12$, and $w_{baba}=w_{abab}$, $w_{baab}=w_{abba}$). That is exactly $\langle\varphi_a|\hat F\varphi_a\rangle=\varepsilon_a$. So $-\varepsilon_a$ is the energy needed to remove a particle from orbital $a$ when the others do not relax. This gives the Hartree--Fock orbital energies a physical meaning; the Kohn--Sham analogue (Janak's theorem) appears in Section 12.15.

\subsection{A worked example: two sites, two electrons}

The smallest system in which interaction matters is a two-site model: two sites L and R, one orbital on each, hopping $t$ between them (the matrix of Section 12.2), and a repulsion $U>0$ when two electrons sit on the same site. We put in two electrons with opposite spins.

\textbf{Exact solution.} With opposite spins the two-electron states split into spin-antisymmetric states $(\uparrow\downarrow-\downarrow\uparrow)/\sqrt2$ with a symmetric spatial part, and spin-symmetric states $(\uparrow\downarrow+\downarrow\uparrow)/\sqrt2$ with an antisymmetric spatial part (in both cases the product is antisymmetric, as Section 12.3 requires). The Hamiltonian does not act on the spins and treats the two electrons alike, so it maps symmetric spatial states to symmetric ones and antisymmetric to antisymmetric ones, and the two kinds can be treated separately. The only antisymmetric spatial state of two sites is $(\mathrm{LR}-\mathrm{RL})/\sqrt2$, where LR means electron 0 on L and electron 1 on R. The hopping maps LR and RL to the same state, $(h_0+h_1)\,\mathrm{LR}=-t(\mathrm{RR}+\mathrm{LL})=(h_0+h_1)\,\mathrm{RL}$, so it gives 0 on $(\mathrm{LR}-\mathrm{RL})/\sqrt2$, and the repulsion gives 0 too (the electrons are on different sites): its energy is 0. The ground state is found among the symmetric spatial states, which are spanned by LL (both on L), RR, and $S=(\mathrm{LR}+\mathrm{RL})/\sqrt2$. The hopping of either electron turns LL into LR or RL, so $(h_0+h_1)\,\mathrm{LL}=-t(\mathrm{RL}+\mathrm{LR})=-\sqrt2\,t\,S$, and likewise for RR, while $(h_0+h_1)\,S=-\sqrt2\,t(\mathrm{LL}+\mathrm{RR})$. The repulsion gives $U$ on LL and RR and 0 on $S$. In the basis (LL, RR, $S$)

\[
\hat H=\begin{pmatrix}U&0&-\sqrt2\,t\\0&U&-\sqrt2\,t\\-\sqrt2\,t&-\sqrt2\,t&0\end{pmatrix}.
\]

The combination $(\mathrm{LL}-\mathrm{RR})/\sqrt2$ decouples with energy $U$; the combination $D=(\mathrm{LL}+\mathrm{RR})/\sqrt2$ couples to $S$ with the matrix element $-2t$, which leaves the $2\times2$ problem with the matrix $\begin{pmatrix}U&-2t\\-2t&0\end{pmatrix}$ and the eigenvalues $\tfrac12\bigl(U\pm\sqrt{U^2+16t^2}\bigr)$. The ground-state energy is

\[
E_0=\tfrac12\Bigl(U-\sqrt{U^2+16t^2}\Bigr),
\]

which is below 0 (because $\sqrt{U^2+16t^2}>U$), so it is the ground state.

\textbf{Hartree--Fock.} A determinant for two electrons of opposite spins is built from one spin-up orbital and one spin-down orbital. The simplest one, called \textbf{restricted Hartree--Fock}, puts both electrons into the bonding orbital $\varphi_b=(\mathrm L+\mathrm R)/\sqrt2$ of Section 12.2, with opposite spins. Its spatial part is $\varphi_b(0)\varphi_b(1)=(\mathrm{LL}+\mathrm{LR}+\mathrm{RL}+\mathrm{RR})/2=(D+S)/\sqrt2$. With the $2\times2$ matrix above, the energy is $\tfrac12U+\tfrac12\cdot0+2\cdot\tfrac1{\sqrt2}\cdot\tfrac1{\sqrt2}\cdot(-2t)$, that is

\[
E_{\text{RHF}}=-2t+\tfrac12U .
\]

\textbf{Is it the best determinant?} Take any spin-up orbital with moduli $(\cos\alpha,\sin\alpha)$ on (L, R) and any spin-down orbital with moduli $(\cos\beta,\sin\beta)$, $0\le\alpha,\beta\le\pi/2$. Exchange acts only between equal spins (Section 12.6), so for this contact repulsion the Hartree and exchange energies add up to $U(n_{\mathrm L\uparrow}n_{\mathrm L\downarrow}+n_{\mathrm R\uparrow}n_{\mathrm R\downarrow})$, the rule $E_H+E_x=g_c\int n_\uparrow n_\downarrow$ of Section 12.6 with the integral replaced by the sum over the two sites; here $n_{\mathrm L\uparrow}=\cos^2\alpha$ is the spin-up density on L, and so on. The energy is this term plus the hopping energies of the two orbitals. The hopping energy of an orbital with the values $(\varphi_\mathrm L,\varphi_\mathrm R)$ on the two sites is $\varphi^\dagger h\varphi=-2t\,\mathrm{Re}(\varphi_\mathrm L^*\varphi_\mathrm R)\ge-2t\,|\varphi_\mathrm L|\,|\varphi_\mathrm R|$, with equality when $\varphi_\mathrm L^*\varphi_\mathrm R$ is real and nonnegative; the repulsion does not depend on the phases, so we choose the phases that make the hopping term most negative, $-2t\cos\alpha\sin\alpha=-t\sin2\alpha$ for spin up. With $s_\alpha=\sin2\alpha$, $c_\alpha=\cos2\alpha$ (the same for $\beta$), and $\cos^2\alpha=\tfrac12(1+c_\alpha)$, $\sin^2\alpha=\tfrac12(1-c_\alpha)$,

\[
E=-t\,(s_\alpha+s_\beta)+\tfrac U2\bigl(1+c_\alpha c_\beta\bigr).
\]

Since $c_\alpha c_\beta\ge-\tfrac12(c_\alpha^2+c_\beta^2)$ (this is $(c_\alpha+c_\beta)^2\ge0$) and $1-c^2=s^2$,

\[
E\ \ge\ \sum_{x=\alpha,\beta}\Bigl(-t\,s_x+\tfrac U4\,s_x^2\Bigr).
\]

Each term, a parabola in $s_x\in[0,1]$, is smallest at $s_x=\min(1,2t/U)$, and equality holds throughout for $\beta=\pi/2-\alpha$ (then $c_\beta=-c_\alpha$ and $s_\beta=s_\alpha$). Hence $E_{HF}=-2t+\tfrac12U$ for $U\le2t$, attained by the restricted state ($\alpha=\beta=\pi/4$). For $U>2t$, $E_{HF}=-2t^2/U$, attained at $\sin2\alpha=2t/U$ by the \textbf{unrestricted} determinant with spin-up orbital $(\cos\alpha,\sin\alpha)$ and spin-down orbital $(\sin\alpha,\cos\alpha)$. That determinant breaks the left--right symmetry of each spin but keeps $n=(1,1)$. For $t=1$, $U=2$ the two coincide: $U=2t$ is exactly the threshold.

\textbf{Numbers.} For $t=1$ and $U=2$: $E_0=\tfrac12(2-\sqrt{20})=-1.236068$ and $E_{HF}=-1$. The difference $E_c=E_0-E_{HF}=-0.236068$ is the \textbf{correlation energy}: the energy that no single determinant can capture. Its physical meaning is visible in the probability of double occupancy (both electrons on one site). In the restricted state it is $|\langle D|\Phi\rangle|^2=\tfrac12$, independent of $U$; for $U\le2t$, so also here at $U=2t$, this is the Hartree--Fock state. In the exact ground state it is $0.276393$ (the squared $D$-component of the lowest eigenvector). The exact electrons \textbf{avoid each other} and keep the left--right symmetry of each spin; a single determinant cannot do both. Its double occupancy is $n_{\mathrm L\uparrow}n_{\mathrm L\downarrow}+n_{\mathrm R\uparrow}n_{\mathrm R\downarrow}=\tfrac12(1+c_\alpha c_\beta)$, which equals $\tfrac12$ whenever each spin is shared equally ($c_\alpha=c_\beta=0$); for $U>2t$ the unrestricted determinant lowers it, but only by breaking that symmetry (Exercise 12.4). Yet the density is the same in the exact and in the Hartree--Fock state, one electron on each site, $n_\mathrm L=n_\mathrm R=1$, by the left--right symmetry. A theory that works with the density alone must therefore contain, somewhere, the information that turns $n=(1,1)$ into $-1.236068$ rather than $-1$. In density functional theory that information is the exchange--correlation functional.

\subsection{The density decides everything: the Hohenberg--Kohn theorems}

We now leave wave functions behind. Fix the particle number $N$ and the interaction $\hat W$, and regard the external potential $v$ as the only thing that distinguishes one system from another (Section 12.3).

\textbf{Theorem 12.1 (Hohenberg--Kohn, 1964).} Let $v$ and $v'$ be two external potentials whose Hamiltonians $\hat H=\hat T+\hat W+\hat V$ and $\hat H'=\hat T+\hat W+\hat V'$ have non-degenerate ground states $\Psi$ and $\Psi'$. If $v-v'$ is not a constant, then the ground-state densities $n$ and $n'$ are different. In other words, the ground-state density determines the potential up to a constant, hence the Hamiltonian, hence every property of the system.

\textbf{Proof.} First, $\Psi\ne\Psi'$ (as states, that is, not merely up to a phase). If they were equal, subtracting the two Schrödinger equations would give $(\hat V-\hat V')\Psi=(E_0-E_0')\Psi$, that is $\sum_i[v(\mathbf r_i)-v'(\mathbf r_i)]=E_0-E_0'$ wherever $\Psi\ne0$, which forces $v-v'$ to be constant. (This step assumes that $\Psi$ does not vanish on a region of positive volume, which holds for the potentials of physics.) Now suppose, to reach a contradiction, that $n=n'$. Because $\Psi$ is the unique ground state of $\hat H$ and $\Psi'\ne\Psi$, the variational principle holds with strict inequality:

\[
E_0<\langle\Psi'|\hat H\Psi'\rangle=\langle\Psi'|\hat H'\Psi'\rangle+\langle\Psi'|(\hat V-\hat V')\Psi'\rangle=E_0'+\int\bigl(v-v'\bigr)\,n'\,d^3r .
\]

Exchanging the roles of the two systems gives $E_0'<E_0+\int(v'-v)\,n\,d^3r$. Add the two inequalities and use $n=n'$: the integrals cancel and we get $E_0+E_0'<E_0+E_0'$, which is false. Hence $n\ne n'$. $\square$

\textbf{The universal functional (Levy--Lieb constrained search).} For every density $n$ that comes from some antisymmetric $N$-particle wave function (an \textbf{$N$-representable} density: nonnegative, integrating to $N$, and not too irregular) define

\[
F[n]=\min_{\Psi\to n}\ \langle\Psi|\hat T+\hat W|\Psi\rangle ,
\]

the smallest kinetic-plus-interaction energy among all wave functions that have the density $n$ (``$\Psi\to n$''). $F$ is \textbf{universal}: it does not depend on $v$, only on $N$ and on the interaction. (We assume that the minimum is attained; the careful mathematical theory replaces ``min'' by ``infimum'', which changes nothing below.)

\textbf{Theorem 12.2 (variational principle for the density).} For every external potential $v$ with ground-state energy $E_0$,

\[
E_0=\min_n\Bigl(F[n]+\int v\,n\,d^3r\Bigr),
\]

and the minimum is attained at the ground-state density.

\textbf{Proof.} The variational principle over wave functions reads $E_0=\min_\Psi\langle\Psi|\hat T+\hat W+\hat V|\Psi\rangle$. Organize the minimization in two steps: first over all $\Psi$ with a given density $n$, then over $n$. In the first step $\langle\Psi|\hat V|\Psi\rangle=\int v\,n$ is the same for all these $\Psi$ (Section 12.3), so the inner minimum is $F[n]+\int v\,n$. The outer minimum over $n$ then equals the minimum over all $\Psi$, which is $E_0$. $\square$

These two theorems say that the energy is a functional of the density alone, $E_v[n]=F[n]+\int vn$, and that minimizing it over densities gives the exact ground-state energy and density. They do not say what $F[n]$ is. For the two-site model of Section 12.8 the theorem guarantees that the number $-1.236068$ is a property of the density $(1,1)$ together with $v$; the task of DFT is to find good approximations for $F$.

\subsection{Functionals and functional derivatives from zero}

A \textbf{functional} is a rule that assigns a number to a whole function: $F[n]$ takes the function $n(\mathbf r)$ and returns one number. Examples: $N[n]=\int n\,d^3r$, $A[n]=\int n^2\,d^3r$, and the Hartree energy $E_H[n]=\tfrac12\int\!\int n(\mathbf r)\,w(\mathbf r,\mathbf r')\,n(\mathbf r')\,d^3r\,d^3r'$.

\textbf{Definition.} Change $n$ by a small amount $\epsilon\,\eta(\mathbf r)$, where $\eta$ is any function and $\epsilon$ a small number. If

\[
F[n+\epsilon\eta]=F[n]+\epsilon\int\frac{\delta F}{\delta n(\mathbf r)}\,\eta(\mathbf r)\,d^3r+O(\epsilon^2)
\]

for every $\eta$, the function $\delta F/\delta n(\mathbf r)$ is the \textbf{functional derivative} of $F$ at $n$. It is the continuum version of the gradient: on a grid with points $\mathbf r_k$ and volume $\Delta V$ per point, $F$ is an ordinary function of the numbers $n_k$ and $\delta F/\delta n(\mathbf r_k)=(\partial F/\partial n_k)/\Delta V$.

\textbf{Examples.} (i) $A[n]=\int n^2$: $A[n+\epsilon\eta]=\int(n^2+2\epsilon n\eta+\epsilon^2\eta^2)$, so $\delta A/\delta n=2n$. (ii) For $G[n]=\int f(n(\mathbf r))\,d^3r$ with any differentiable function $f$, Taylor's formula $f(n+\epsilon\eta)=f(n)+\epsilon f'(n)\eta+O(\epsilon^2)$ gives $\delta G/\delta n=f'(n)$; for instance $\delta\int n^{4/3}/\delta n=\tfrac43n^{1/3}$ and $\delta\int n^{5/3}/\delta n=\tfrac53n^{2/3}$. (iii) For $E_H$ the first-order change is $\tfrac12\int\!\int(\eta w n'+n w\eta')$, and because $w$ is symmetric both terms are equal: $\delta E_H/\delta n(\mathbf r)=\int w(\mathbf r,\mathbf r')\,n(\mathbf r')\,d^3r'=v_H(\mathbf r)$, the Hartree potential of Section 12.7. (iv) $\delta\int v\,n/\delta n=v$.

\textbf{Worked example on a grid.} Take three points with $\Delta V=1$ and $A=\sum_kn_k^2$. At $n=(1,\tfrac12,\tfrac12)$ (the density of Section 12.4) the gradient is $(2,1,1)$, which is $2n$ as example (i) says. Moving the density by $\epsilon\eta$ with $\eta=(1,-1,0)$ (a particle-conserving change) changes $A$ by $\epsilon(2-1)+O(\epsilon^2)=\epsilon+O(\epsilon^2)$; directly, $A(1+\epsilon,\tfrac12-\epsilon,\tfrac12)-A(1,\tfrac12,\tfrac12)=\epsilon+2\epsilon^2$.

\textbf{Minimizing with a constraint.} To minimize $E[n]$ over densities with $\int n=N$, introduce a Lagrange multiplier $\mu$ (Section 12.7), that is, make $\mathcal L=E[n]-\mu\bigl(\int n-N\bigr)$ stationary, and require $\delta E/\delta n(\mathbf r)=\mu$ at every point. (On a grid the constraint is $c=\sum_kn_k\,\Delta V=N$, with $\partial c/\partial n_k=\Delta V$, so the condition $\partial E/\partial n_k=\mu\,\Delta V$ of Section 12.7 is $\delta E/\delta n(\mathbf r_k)=\mu$ by the definition above.) The multiplier $\mu$ is the \textbf{chemical potential}: the change of the minimal energy per added particle, $\mu=dE_0/dN$. This is the meaning of a multiplier derived in Section 12.7, $df^\ast/dc_0=\lambda$, with $f=E$, $c=\int n$ and $c_0=N$, and it holds where the minimum depends differentiably on $N$.

\subsection{The Kohn--Sham equations}

The unknown functional $F[n]$ contains the kinetic energy, and the kinetic energy is the hardest part to approximate by an explicit formula in $n$ (Section 12.12 shows the crude attempt). Kohn and Sham (1965) avoided the problem with a trick: compute the kinetic energy \textbf{exactly}, but for a fictitious system of non-interacting particles that has the same density.

\textbf{The non-interacting reference system.} For non-interacting particles ($\hat W=0$) the constrained search of Section 12.9 defines

\[
T_s[n]=\min_{\Phi\to n}\langle\Phi|\hat T|\Phi\rangle ,
\]

the smallest kinetic energy of a non-interacting state with density $n$; its minimizer is (for the densities we consider) a single Slater determinant of orbitals $\varphi_a$ with $n=\sum_a|\varphi_a|^2$. The Kohn--Sham scheme \textbf{assumes} that the density of interest is \textbf{non-interacting $v$-representable}: it is the ground-state density of some non-interacting system in some potential $v_s$. This is an assumption, not a theorem; it holds in all practical applications we are aware of, and Chapter 13 simply adopts it.

\textbf{Splitting the functional.} Write the exact universal functional as

\[
F[n]=T_s[n]+E_H[n]+E_{xc}[n],
\]

which \textbf{defines} the \textbf{exchange--correlation energy} $E_{xc}=F-T_s-E_H$. It collects everything that the easy pieces miss: the exchange energy of Section 12.6, the correlation energy of Section 12.8, and the difference between the true kinetic energy and $T_s$. Nothing has been approximated yet.

\textbf{The equations.} Minimize $E[n]=T_s[n]+\int v\,n+E_H[n]+E_{xc}[n]$ by varying the orbitals of the determinant, with Lagrange multipliers $\varepsilon_a$ for their normalization. Since $n=\sum_a|\varphi_a|^2$, the chain rule gives $\delta n(\mathbf r')/\delta\varphi_a^*(\mathbf r)=\varphi_a(\mathbf r)\,\delta(\mathbf r-\mathbf r')$, so every density functional $G[n]$ contributes $(\delta G/\delta n)\,\varphi_a$. The kinetic term contributes $-\tfrac12\nabla^2\varphi_a$. Setting the derivative to zero gives the \textbf{Kohn--Sham equations}

\[
\Bigl[-\tfrac12\nabla^2+v_s(\mathbf r)\Bigr]\varphi_a=\varepsilon_a\varphi_a,\qquad
v_s=v+v_H+v_{xc},\qquad v_{xc}(\mathbf r)=\frac{\delta E_{xc}}{\delta n(\mathbf r)},
\]

\[
n(\mathbf r)=\sum_{a\in O}|\varphi_a(\mathbf r)|^2 ,
\]

with the $N$ lowest orbitals occupied. (As in Section 12.7 the multipliers can be made diagonal because the equations depend on the occupied orbitals only through $n$.) The \textbf{Kohn--Sham potential} $v_s$ is local: a multiplication by a function, simpler than the Fock operator. If $E_{xc}$ were known exactly, these one-particle equations would give the exact ground-state density and energy of the interacting system.

\textbf{Total energy.} Multiplying the equations by $\varphi_a^*$, integrating and summing, $\sum_a\varepsilon_a=T_s+\int v_s\,n$. Hence

\[
E_0=\sum_{a\in O}\varepsilon_a-E_H[n]+E_{xc}[n]-\int v_{xc}\,n\,d^3r ,
\]

the Kohn--Sham analogue of the double-counting formula of Section 12.7.

\textbf{Self-consistency.} $v_s$ depends on $n$, and $n$ on the orbitals, which depend on $v_s$. The equations are solved by the loop

\begin{enumerate}
\item guess a density $n$;
\item build $v_s=v+v_H[n]+v_{xc}[n]$;
\item solve the one-particle equations and occupy the $N$ lowest orbitals;
\item form the new density $n_{\text{out}}=\sum_a|\varphi_a|^2$;
\item if $n_{\text{out}}$ equals $n$ within a tolerance, stop; otherwise mix $n$ and $n_{\text{out}}$ (Section 12.13) and go to step 2.
\end{enumerate}

\textbf{What the orbitals mean.} The Kohn--Sham orbitals and energies belong to the fictitious non-interacting system. Only the density and the total energy are guaranteed to be physical. The orbital energies are nevertheless useful: Section 12.15 shows that each is the derivative of the total energy with respect to the occupation of its own orbital, so the highest occupied one, $\varepsilon_{\text{HOMO}}$ (HOMO: highest occupied molecular orbital, a name taken over from chemistry), is the rate at which the energy changes when a small fraction of a particle is removed from the system.

\subsection{The uniform gas and the local density approximation}

To use the Kohn--Sham equations we need an approximation for $E_{xc}[n]$. The oldest and simplest one borrows it from the only many-body system whose properties are known accurately: the \textbf{uniform gas}, infinitely many particles spread with constant density over all space.

\textbf{Spherical coordinates.} The integrals of this section are done in spherical coordinates: $x=r\sin\theta\cos\phi$, $y=r\sin\theta\sin\phi$, $z=r\cos\theta$ with $r\ge0$, $0\le\theta\le\pi$ and $0\le\phi<2\pi$. The flat metric in these coordinates follows from the polar-coordinate result of Chapter 1 used twice. With $\varrho=r\sin\theta$, the pair $(x,y)=(\varrho\cos\phi,\varrho\sin\phi)$ gives $dx^2+dy^2=d\varrho^2+\varrho^2d\phi^2$, and the pair $(z,\varrho)=(r\cos\theta,r\sin\theta)$ gives $dz^2+d\varrho^2=dr^2+r^2d\theta^2$. Hence

\[
dx^2+dy^2+dz^2=dr^2+r^2d\theta^2+r^2\sin^2\theta\,d\phi^2 ,
\]

the metric is $\mathrm{diag}(1,r^2,r^2\sin^2\theta)$, $\sqrt{|g|}=r^2\sin\theta$ (Chapter 1: $\sqrt{|g|}$ is the volume factor) and $d^3r=r^2\sin\theta\,dr\,d\theta\,d\phi$. With $u=\cos\theta$ ($du=-\sin\theta\,d\theta$, and $\theta$ from 0 to $\pi$ means $u$ from 1 to $-1$):

\[
\int d^3r\,f=\int_0^\infty r^2\,dr\int_{-1}^{1}du\int_0^{2\pi}d\phi\ f .
\]

For $f$ depending only on $r$ this is $4\pi\int_0^\infty r^2f\,dr$. For $f$ depending only on $r$ and $u$ it is $2\pi\int_0^\infty r^2\,dr\int_{-1}^1du\,f$. The same formulas hold for integrals over a momentum $\mathbf k$, with $k=|\mathbf k|$ in place of $r$, and the axis $\theta=0$ may be chosen in any direction.

\textbf{Plane waves in a box.} Put the particles in a cube of side $\ell$ and volume $V=\ell^3$ and require the wave functions to repeat themselves from one face to the opposite one (\textbf{periodic boundary conditions}, a torus). The orbitals are plane waves $\varphi_{\mathbf k}(\mathbf r)=e^{i\mathbf k\cdot\mathbf r}/\sqrt V$, eigenfunctions of $-\tfrac12\nabla^2$ with energy $\tfrac12k^2$, and periodicity allows only $\mathbf k=(2\pi/\ell)(n_1,n_2,n_3)$ with integers $n_1,n_2,n_3$. Each allowed $\mathbf k$ occupies a cube of volume $(2\pi/\ell)^3=(2\pi)^3/V$ in $\mathbf k$-space, so for a large box a sum over allowed momenta becomes an integral:

\[
\sum_{\mathbf k}\ \longrightarrow\ V\int\frac{d^3k}{(2\pi)^3}.
\]

The same torus, with the spacing $\Delta k=2\pi/\ell$, is used for the three space directions $x_1,x_2,x_3$ in Chapter 13.

\textbf{The Fermi sphere.} The non-interacting ground state fills the lowest energies: all $\mathbf k$ with $|\mathbf k|<k_F$ (the \textbf{Fermi momentum}), each with all $g$ internal states. Counting states,

\[
N=g\,V\,\frac{\tfrac43\pi k_F^3}{(2\pi)^3}\quad\Longrightarrow\quad n=\frac{g\,k_F^3}{6\pi^2}.
\]

The kinetic energy per volume is $t=g\int_{k<k_F}\frac{d^3k}{(2\pi)^3}\frac{k^2}2=\frac{g}{2\pi^2}\cdot\frac{k_F^5}{10}$ (spherical coordinates, as above: $d^3k\to4\pi k^2dk$ for an integrand that depends only on $k$). For electrons ($g=2$): $n=k_F^3/(3\pi^2)$, $t=k_F^5/(10\pi^2)$, and eliminating $k_F=(3\pi^2n)^{1/3}$,

\[
t(n)=C_F\,n^{5/3},\qquad C_F=\tfrac3{10}\bigl(3\pi^2\bigr)^{2/3}=2.871234 .
\]

The \textbf{Thomas--Fermi} model (1927) approximated the kinetic energy of any system by $\int t(n(\mathbf r))\,d^3r$. It is too crude (it predicts, for example, that molecules do not bind); the Kohn--Sham scheme keeps the kinetic energy exact through $T_s$ and approximates only $E_{xc}$.

\textbf{Exchange energy of the uniform electron gas.} We compute $E_x$ of Section 12.6 for the filled Fermi sphere with $w=1/|\mathbf r-\mathbf r'|$, $g=2$, in five steps.

\emph{Step 1: the density matrix.} For each spin, $\rho_\sigma(\mathbf r,\mathbf r')=\sum_{k<k_F}\varphi_{\mathbf k}(\mathbf r)\varphi^*_{\mathbf k}(\mathbf r')=\int_{k<k_F}\frac{d^3k}{(2\pi)^3}e^{i\mathbf k\cdot(\mathbf r-\mathbf r')}$, a function of $\mathbf R=\mathbf r-\mathbf r'$ only.

\emph{Step 2: one integral from translation invariance.} Since the integrand depends only on $\mathbf R$, $\int\!\int d^3r\,d^3r'\to V\int d^3R$, and the two spins give equal contributions:

\[
\frac{E_x}V=-\tfrac12\cdot2\int d^3R\,\frac{|\rho_\sigma(\mathbf R)|^2}{R}
=-\int_{k<k_F}\!\!\frac{d^3k}{(2\pi)^3}\int_{k'<k_F}\!\!\frac{d^3k'}{(2\pi)^3}\int d^3R\,\frac{e^{i(\mathbf k-\mathbf k')\cdot\mathbf R}}{R}.
\]

\emph{Step 3: the Coulomb integral.} For $\mathbf q\ne0$ compute $\int d^3R\,e^{i\mathbf q\cdot\mathbf R}e^{-\kappa R}/R$ with a small number $\kappa>0$ that makes it converge (a \textbf{regulator}; the letter $\mu$ is kept for the chemical potential), using spherical coordinates around $\mathbf q$ ($u$ is the cosine of the angle between $\mathbf R$ and $\mathbf q$):

\[
\begin{aligned}
\int d^3R\,\frac{e^{i\mathbf q\cdot\mathbf R}e^{-\kappa R}}R
&=2\pi\int_0^\infty R\,e^{-\kappa R}\,dR\int_{-1}^1e^{iqRu}\,du
=\frac{2\pi}{iq}\int_0^\infty\bigl(e^{(iq-\kappa)R}-e^{(-iq-\kappa)R}\bigr)\,dR\\
&=\frac{2\pi}{iq}\Bigl(\frac1{\kappa-iq}-\frac1{\kappa+iq}\Bigr)=\frac{4\pi}{q^2+\kappa^2}.
\end{aligned}
\]

Letting $\kappa\to0$ gives $4\pi/q^2$. Hence

\[
\frac{E_x}V=-\frac1{(2\pi)^6}\int_{k<k_F}d^3k\ I(k),\qquad I(k)=\int_{k'<k_F}d^3k'\,\frac{4\pi}{|\mathbf k-\mathbf k'|^2}.
\]

\emph{Step 4: the inner integral.} Spherical coordinates around $\mathbf k$ give $\int_{-1}^1\frac{du}{k^2+k'^2-2kk'u}=\frac1{kk'}\ln\Bigl|\frac{k+k'}{k-k'}\Bigr|$, so $I(k)=\frac{8\pi^2}k\int_0^{k_F}k'\ln\bigl|\frac{k+k'}{k-k'}\bigr|\,dk'$. The last integral equals $kk_F+\tfrac12(k_F^2-k^2)\ln\bigl|\frac{k+k_F}{k-k_F}\bigr|$ (both sides vanish at $k_F=0$, and their derivatives with respect to $k_F$ agree, as one checks with $\frac{d}{dk_F}\ln\bigl|\frac{k+k_F}{k-k_F}\bigr|=\frac{2k}{k^2-k_F^2}$). With $x=k/k_F$,

\[
I(k)=16\pi^2k_F\,L(x),\qquad L(x)=\frac12+\frac{1-x^2}{4x}\ln\Bigl|\frac{1+x}{1-x}\Bigr| .
\]

\emph{Step 5: the outer integral.} $\int_{k<k_F}d^3k\,I(k)=16\pi^2k_F\cdot4\pi k_F^3\int_0^1x^2L(x)\,dx$. Expand $\ln\frac{1+x}{1-x}=2\sum_{m\ge0}\frac{x^{2m+1}}{2m+1}$ for $0\le x<1$; then $\int_0^1x(1-x^2)\ln\frac{1+x}{1-x}\,dx=2\sum_m\frac1{2m+1}\bigl(\frac1{2m+3}-\frac1{2m+5}\bigr)$. The partial fractions $\frac1{(2m+1)(2m+3)}=\frac12\bigl(\frac1{2m+1}-\frac1{2m+3}\bigr)$ and $\frac1{(2m+1)(2m+5)}=\frac14\bigl(\frac1{2m+1}-\frac1{2m+5}\bigr)$ make both sums telescope, to $\tfrac12$ and $\tfrac14(1+\tfrac13)=\tfrac13$, so the integral is $2(\tfrac12-\tfrac13)=\tfrac13$ and $\int_0^1x^2L\,dx=\tfrac16+\tfrac14\cdot\tfrac13=\tfrac14$. Therefore $E_x/V=-64\pi^3k_F^4/(4\cdot64\pi^6)=-k_F^4/(4\pi^3)$, and with $k_F=(3\pi^2n)^{1/3}$:

\[
e_x(n)=\frac{E_x}V=-\frac34\Bigl(\frac3\pi\Bigr)^{1/3}n^{4/3}=-0.738559\,n^{4/3},\qquad \frac{E_x}N=-\frac{3k_F}{4\pi}.
\]

This is \textbf{Dirac's exchange energy} (1930).

\textbf{The local density approximation (LDA).} A real density varies in space. The LDA treats each small volume $d^3r$ as if it were a piece of uniform gas with the local density $n(\mathbf r)$:

\[
E_{xc}^{\text{LDA}}[n]=\int e_{xc}\bigl(n(\mathbf r)\bigr)\,d^3r,\qquad v_{xc}^{\text{LDA}}(\mathbf r)=\frac{de_{xc}}{dn}\Big|_{n(\mathbf r)}
\]

(example (ii) of Section 12.10). Its exchange part is Dirac's formula, $v_x=-(3/\pi)^{1/3}n^{1/3}$. The correlation part of the uniform electron gas is not known in closed form; it is taken from numerical (quantum Monte Carlo) simulations of the uniform gas, which we quote as a fact of the literature and do not use in this book. The LDA is exact for a uniform density and is an approximation for every other one; its accuracy must be judged case by case.

\textbf{The contact interaction is special.} For the contact interaction of Section 12.6 the exchange energy of every determinant (or non-interacting ensemble) is exactly local:

\[
E_x=-\frac{g_c}2\int\sum_{\sigma,\sigma'}\bigl|\rho(\mathbf r\sigma,\mathbf r\sigma')\bigr|^2\,d^3r ,
\]

which for orbitals of definite label is $-\tfrac{g_c}2\int\sum_\sigma n_\sigma^2\,d^3r$ (put $w=g_c\,\delta(\mathbf r-\mathbf r')$ into the exchange integral of Section 12.6). The uniform-gas formula $e_x=-\tfrac{g_c}2\sum_\sigma n_\sigma^2$ (the plane waves of the uniform gas have definite labels), evaluated with all the local label densities $n_\sigma$, is then not an approximation for the exchange part; evaluated with the total density alone, as $-\tfrac{g_c}{2g}n^2$, it is one, unless all $n_\sigma$ are equal. An ``exchange-only LDA'' of this kind is the Hartree--Fock energy (over determinants of orbitals of definite label) written as a density functional. Chapter 13 is close to this situation but not in it. There the exact Fock term is local but depends on the $2\times2$ density matrices of eight blocks, while the solver uses the uniform-gas closed form in only two local densities, a number density and a scalar density, of a relativistic gas of particles and antiparticles. That closed form equals the exact term only for equally occupied blocks without momentum current, so the choice is the approximation of the exchange part there (Section 13.8). There is also no correlation term (Section 13.8).

\subsection{Solving the Kohn--Sham equations: iteration and mixing}

The loop of Section 12.11 defines a map from an input density to an output density, $n_{\text{out}}=G[n_{\text{in}}]$, and the self-consistent density is a \textbf{fixed point}, $G[n]=n$. The naive iteration $n_{\text{in}}\leftarrow n_{\text{out}}$ often fails; this section shows why with a model small enough to follow by hand.

\textbf{Model.} Two sites L and R with site energies $-\Delta/2$ and $+\Delta/2$, hopping $t$ and on-site repulsion $U$, and two electrons with opposite spins in the lowest orbital. In the mean field each electron feels $U$ times the density of the \textbf{other} spin on the same site, $Un_\mathrm L/2$ on L and $Un_\mathrm R/2$ on R (the contact rule of Section 12.6: exchange removes the same-spin half). The one-electron Hamiltonian and the output density are

\[
h[n]=\begin{pmatrix}-\tfrac\Delta2+\tfrac U2n_\mathrm L&-t\\-t&\tfrac\Delta2+\tfrac U2n_\mathrm R\end{pmatrix},\qquad n_\mathrm R=2-n_\mathrm L,\qquad n_\mathrm L^{\text{out}}=2\,|c_\mathrm L|^2,
\]

where $(c_\mathrm L,c_\mathrm R)$ is the normalized lowest eigenvector of $h[n]$. For a real symmetric matrix $\begin{pmatrix}a&-t\\-t&b\end{pmatrix}$ the lowest eigenvector has $|c_\mathrm L|^2=\tfrac12\bigl(1+(b-a)/\sqrt{(b-a)^2+4t^2}\bigr)$. Here $b-a=\Delta+U(1-n_\mathrm L)$. With $x=n_\mathrm L-1$ (the excess on L) the whole loop becomes one function of one number, the map $G$ of the beginning of this section in reduced form:

\[
x_{\text{out}}=G(x)=\frac{\Delta-Ux}{\sqrt{(\Delta-Ux)^2+4t^2}} .
\]

\textbf{Numbers.} Take $\Delta=2$, $t=1$, $U=4$ and the start $n_\mathrm L=2$ (both electrons on the low site). Plain iteration ($n_\mathrm L\leftarrow n_\mathrm L^{\text{out}}$) gives

\begingroup
\small
\begin{longtable}{@{}>{\raggedright\arraybackslash}p{0.114\linewidth}>{\raggedright\arraybackslash}p{0.114\linewidth}>{\raggedright\arraybackslash}p{0.114\linewidth}>{\raggedright\arraybackslash}p{0.114\linewidth}>{\raggedright\arraybackslash}p{0.114\linewidth}>{\raggedright\arraybackslash}p{0.114\linewidth}>{\raggedright\arraybackslash}p{0.114\linewidth}@{}}
\toprule
\textbf{step} & \textbf{0} & \textbf{1} & \textbf{2} & \textbf{3} & \textbf{4} & \textbf{5} \\
\midrule
\endfirsthead
\toprule
\textbf{step} & \textbf{0} & \textbf{1} & \textbf{2} & \textbf{3} & \textbf{4} & \textbf{5} \\
\midrule
\endhead
$n_\mathrm L$ in & 2.000000 & 0.292893 & 1.923880 & 0.353344 & 1.916644 & 0.359836 \\
$n_\mathrm L$ out & 0.292893 & 1.923880 & 0.353344 & 1.916644 & 0.359836 & 1.915809 \\
\bottomrule
\end{longtable}
\endgroup

and continues to jump between about $0.3607$ and $1.9157$ forever: when the electrons sit on L, the repulsion pushes them to R, and vice versa. This is called \textbf{charge sloshing}. \textbf{Linear mixing} feeds back only a fraction $\beta$ of the change, $n_{\text{in}}\leftarrow(1-\beta)\,n_{\text{in}}+\beta\,n_{\text{out}}$. With $\beta=\tfrac12$:

\begingroup
\small
\begin{longtable}{@{}>{\raggedright\arraybackslash}p{0.114\linewidth}>{\raggedright\arraybackslash}p{0.114\linewidth}>{\raggedright\arraybackslash}p{0.114\linewidth}>{\raggedright\arraybackslash}p{0.114\linewidth}>{\raggedright\arraybackslash}p{0.114\linewidth}>{\raggedright\arraybackslash}p{0.114\linewidth}>{\raggedright\arraybackslash}p{0.114\linewidth}@{}}
\toprule
\textbf{step} & \textbf{0} & \textbf{1} & \textbf{2} & \textbf{3} & \textbf{4} & \textbf{5} \\
\midrule
\endfirsthead
\toprule
\textbf{step} & \textbf{0} & \textbf{1} & \textbf{2} & \textbf{3} & \textbf{4} & \textbf{5} \\
\midrule
\endhead
$n_\mathrm L$ in & 2.000000 & 1.146447 & 1.361898 & 1.314066 & 1.331307 & 1.325494 \\
$n_\mathrm L$ out & 0.292893 & 1.577350 & 1.266234 & 1.348548 & 1.319682 & 1.329519 \\
\bottomrule
\end{longtable}
\endgroup

The iteration converges to the self-consistent value $n_\mathrm L=1.326993$; the input and output agree to $10^{-6}$ after 13 steps.

\textbf{Why.} Near the fixed point $x_\ast=0.326993$ write $x=x_\ast+e$. To first order $G(x)=x_\ast+G'(x_\ast)\,e$, so the mixed iteration multiplies the error by $1-\beta\,(1-G'(x_\ast))$ at each step. Differentiating, $G'(x)=-4Ut^2/\bigl((\Delta-Ux)^2+4t^2\bigr)^{3/2}$, which gives $G'(x_\ast)=-1.687961$. The iteration converges exactly when $|1-\beta(1-G')|<1$, that is

\[
0<\beta<\frac2{1-G'(x_\ast)}=0.744058 .
\]

Plain iteration ($\beta=1$) multiplies the error by $-1.688$: it grows and alternates in sign, which is the sloshing. For $\beta=\tfrac12$ the factor is $-0.344$, and the error shrinks by about a factor 3 per step, as the table shows. The choice $\beta=1/(1-G')=0.372$ would give the factor 0.

\textbf{Anderson (Pulay) mixing.} For densities with many components the best $\beta$ differs from direction to direction. \textbf{Anderson mixing} keeps the last few input densities $n^{(i)}$ and their residuals $R^{(i)}=G[n^{(i)}]-n^{(i)}$, finds the numbers $c_i$ with $\sum_ic_i=1$ that make $\bigl\|\sum_ic_iR^{(i)}\bigr\|$ smallest, and takes as the next input $\sum_ic_i\bigl(n^{(i)}+\beta R^{(i)}\bigr)$. It estimates the derivative of $G$ from the history, a cheap substitute for Newton's method. The Rust solver of Chapter 13 uses it, with the stopping rule that the largest change of the densities, divided by their largest value, is below $10^{-10}$.

\textbf{Level crossings and smearing.} With integer (aufbau) occupations the output density jumps whenever the order of two levels near the highest occupied one changes during the iteration. If the levels cross back and forth, the loop never settles, whatever $\beta$ is. The standard remedy replaces the integer occupations by smooth Fermi--Dirac occupations of a small width (\textbf{smearing}). The converged object is then an \textbf{ensemble}: a weighted mixture of determinants with fractional occupations near the Fermi level, not a single determinant. One run of Chapter 13 needs this remedy, and its results are reported as those of the ensemble.

\subsection{Finite temperature: Mermin's functional}

So far the system was in its ground state. At a temperature $T>0$ it is in a statistical mixture of states. We set Boltzmann's constant to 1, so a temperature is an energy.

\textbf{Ensembles.} A mixture in which the state $\Psi_k$ occurs with probability $w_k\ge0$ ($\sum_kw_k=1$) is described by the \textbf{density operator} $\hat\rho=\sum_kw_k|\Psi_k\rangle\langle\Psi_k|$ (outer products, Section 12.2), a Hermitian operator with nonnegative eigenvalues and trace 1: $\langle\varphi|\hat\rho\varphi\rangle=\sum_kw_k|\langle\Psi_k|\varphi\rangle|^2\ge0$ for every state $\varphi$, and $\mathrm{Tr}\,\hat\rho=\sum_kw_k\langle\Psi_k|\Psi_k\rangle=1$ for normalized $\Psi_k$. Its eigenvalues $p_i$ therefore lie between 0 and 1 and add up to 1. The expectation value of an observable is $\mathrm{Tr}(\hat\rho\hat A)=\sum_kw_k\langle\Psi_k|\hat A\Psi_k\rangle$, and the \textbf{entropy} is $S=-\mathrm{Tr}(\hat\rho\ln\hat\rho)=-\sum_ip_i\ln p_i\ge0$, where the $p_i$ are the eigenvalues of $\hat\rho$ (the logarithm of an operator is taken in its eigenbasis, Section 12.2, and every term is $\ge0$ because $\ln p_i\le0$). If the $\Psi_k$ are orthonormal, the $p_i$ are the $w_k$ and $S=-\sum_kw_k\ln w_k$. When the particle number may vary (it is fixed on average by a chemical potential $\mu$), the equilibrium state at temperature $T$ minimizes the \textbf{grand potential}

\[
\Omega[\hat\rho]=\mathrm{Tr}\bigl[\hat\rho\,(\hat H-\mu\hat N)\bigr]+T\,\mathrm{Tr}\bigl[\hat\rho\ln\hat\rho\bigr].
\]

\textbf{Theorem 12.3 (Gibbs principle).} In a finite-dimensional state space, $\Omega$ has exactly one minimizer, the Gibbs state $\hat\rho_0=e^{-(\hat H-\mu\hat N)/T}/Z$, $Z=\mathrm{Tr}\,e^{-(\hat H-\mu\hat N)/T}$, and $\Omega[\hat\rho_0]=-T\ln Z$.

\textbf{Proof.} Since $\ln\hat\rho_0=-(\hat H-\mu\hat N)/T-\ln Z$, we have $\hat H-\mu\hat N=-T\ln\hat\rho_0-T\ln Z$, and therefore $\Omega[\hat\rho]=-T\ln Z+T\,\mathrm{Tr}\bigl[\hat\rho(\ln\hat\rho-\ln\hat\rho_0)\bigr]$. It remains to show \textbf{Klein's inequality} $D=\mathrm{Tr}[\hat\rho(\ln\hat\rho-\ln\hat\sigma)]\ge0$ for density operators $\hat\rho$ and $\hat\sigma$ with $\hat\sigma$ having positive eigenvalues. Write $\hat\rho=\sum_ip_i|i\rangle\langle i|$ and $\hat\sigma=\sum_jq_j|j\rangle\langle j|$ with orthonormal eigenbases. Then, using $\sum_j|\langle i|j\rangle|^2=1$ for every $i$ and $\sum_i|\langle i|j\rangle|^2=1$ for every $j$ (completeness of the two bases),

\[
\begin{aligned}
D&=\sum_ip_i\ln p_i-\sum_{i,j}p_i\,|\langle i|j\rangle|^2\ln q_j
=\sum_{i,j}|\langle i|j\rangle|^2\,p_i\bigl(\ln p_i-\ln q_j\bigr)\\
&\ge\sum_{i,j}|\langle i|j\rangle|^2\bigl(p_i-q_j\bigr)=1-1=0 .
\end{aligned}
\]

The inequality used, $a\ln a-a\ln b\ge a-b$ for $a\ge0$ and $b>0$, is $\ln y\le y-1$ with $y=b/a$, multiplied by $-a$ (for $a=0$ it reads $0\ge-b$). Moreover $\ln y\le y-1$ holds with equality only at $y=1$: the function $y-1-\ln y$ has the derivative $1-1/y$, negative for $y<1$ and positive for $y>1$, so its smallest value is the value 0 at $y=1$.

\emph{Equality.} Subtracting the sum $\sum_{i,j}|\langle i|j\rangle|^2(p_i-q_j)=0$ writes $D=\sum_{i,j}|\langle i|j\rangle|^2\,d(p_i,q_j)$ with the nonnegative numbers $d(a,b)=a\ln a-a\ln b-a+b$. So $D=0$ forces $d(p_i,q_j)=0$ whenever $\langle i|j\rangle\ne0$. For $a>0$, $d(a,b)=a\,(y-1-\ln y)$ with $y=b/a$ vanishes only for $b=a$; for $a=0$, $d(0,b)=b>0$ never vanishes. Hence $p_i=q_j$ whenever $\langle i|j\rangle\ne0$, and then for every $j$

\[
\hat\rho\,|j\rangle=\sum_ip_i\,|i\rangle\langle i|j\rangle=q_j\sum_i|i\rangle\langle i|j\rangle=q_j\,|j\rangle=\hat\sigma\,|j\rangle ,
\]

so $\hat\rho=\hat\sigma$. With $\hat\sigma=\hat\rho_0$, whose eigenvalues are the positive numbers $e^{-\epsilon_k/T}/Z$, $\epsilon_k$ running over the eigenvalues of $\hat H-\mu\hat N$ (a function of an operator, Section 12.2), this gives $\Omega[\hat\rho]=-T\ln Z+T\,D>\Omega[\hat\rho_0]$ for every $\hat\rho\ne\hat\rho_0$: the Gibbs state is the only minimizer. $\square$

\textbf{Mermin's theorem (1965).} Replace the variational principle of Section 12.9 by the Gibbs principle: the proofs of Theorems 12.1 and 12.2 then go through word for word (the first step becomes: if $\hat\rho_0=\hat\rho_0'$ then $\ln\hat\rho_0=\ln\hat\rho_0'$, so $\hat V-\hat V'$ is a multiple of the identity; on the vacuum it is 0, so the multiple is 0, and on one particle $v=v'$; at fixed $\mu$ even a constant difference of the potentials changes the Gibbs state), with the strict inequality $\Omega[\hat\rho_0']>\Omega[\hat\rho_0]$ for two different Gibbs states taking the place of $\langle\Psi'|\hat H\Psi'\rangle>E_0$. At fixed $T$ and $\mu$ the equilibrium density determines the external potential, and there is a universal functional

\[
F_T[n]=\min_{\hat\rho\to n}\mathrm{Tr}\bigl[\hat\rho\,(\hat T+\hat W+T\ln\hat\rho)\bigr],\qquad
\Omega=\min_n\Bigl(F_T[n]+\int(v-\mu)\,n\,d^3r\Bigr).
\]

(Here $\hat T$ is the kinetic-energy operator and $T$ the temperature.)

\textbf{The Kohn--Sham form at finite temperature.} The non-interacting reference system is now an ensemble: orbitals $\varphi_a$ with \textbf{occupation probabilities} $f_a\in[0,1]$, density $n=\sum_af_a|\varphi_a|^2$. Its entropy follows from the independence of the occupations, which we now derive. For non-interacting fermions with orbital energies $\varepsilon_a$, $\hat H-\mu\hat N=\sum_a(\varepsilon_a-\mu)\,\hat n_a$ in the basis of these orbitals, and the $\hat n_a$ commute with one another. Every determinant $|n_0n_1\cdots\rangle$ is therefore an eigenstate of $\hat H-\mu\hat N$ with the eigenvalue $\sum_a(\varepsilon_a-\mu)n_a$, and in the Gibbs state it has the weight

\[
\frac{e^{-\sum_a(\varepsilon_a-\mu)n_a/T}}Z=\prod_a\frac{e^{-(\varepsilon_a-\mu)n_a/T}}{1+e^{-(\varepsilon_a-\mu)/T}},\qquad Z=\prod_a\bigl(1+e^{-(\varepsilon_a-\mu)/T}\bigr),
\]

because the sum of $e^{-\sum_a(\varepsilon_a-\mu)n_a/T}$ over all choices $n_a\in\{0,1\}$ factorizes into one factor per orbital. The probability of the configuration $\{n_a\}$ is thus a product $\prod_ap_a(n_a)$ with $p_a(1)=e^{-(\varepsilon_a-\mu)/T}/(1+e^{-(\varepsilon_a-\mu)/T})=f_a$ and $p_a(0)=1-f_a$: each orbital is an independent two-state system, and $f_a=1/(e^{(\varepsilon_a-\mu)/T}+1)$ is the Fermi--Dirac function met again below. For a product of independent probabilities the entropies add, since the logarithm of a product is the sum of the logarithms and the probabilities of each factor add up to 1. The reference ensemble with arbitrary occupations $f_a$ is taken of the same product form, so its entropy is

\[
S_s=-\sum_a\bigl[f_a\ln f_a+(1-f_a)\ln(1-f_a)\bigr].
\]

\textbf{Mermin's Kohn--Sham functional} is the free energy

\[
F=\sum_af_a\langle\varphi_a|-\tfrac12\nabla^2|\varphi_a\rangle-T\,S_s+\int v\,n+E_H[n]+F_{xc}[n],
\]

with a temperature-dependent exchange--correlation free energy $F_{xc}$. Making $F$ stationary with respect to the orbitals gives the Kohn--Sham equations of Section 12.11 unchanged. Making it stationary with respect to the occupations under the constraint $\sum_af_a=N$ (multiplier $\mu$): the energy terms contribute $\partial/\partial f_a=\langle\varphi_a|-\tfrac12\nabla^2|\varphi_a\rangle+\int(v+v_H+v_{xc})|\varphi_a|^2=\varepsilon_a$, and $\partial/\partial f\,[f\ln f+(1-f)\ln(1-f)]=\ln\frac f{1-f}$. So

\[
\varepsilon_a+T\ln\frac{f_a}{1-f_a}-\mu=0\quad\Longleftrightarrow\quad f_a=\frac1{e^{(\varepsilon_a-\mu)/T}+1},
\]

the \textbf{Fermi--Dirac distribution}. The sum $\sum_af_a$ increases strictly with $\mu$ (each $f_a$ does), so exactly one $\mu$ gives $\sum_af_a=N$; a computer finds it by \textbf{bisection} (halving an interval that brackets it). As $T\to0$, $f_a\to1$ for $\varepsilon_a<\mu$ and $f_a\to0$ for $\varepsilon_a>\mu$: the aufbau rule.

\textbf{Thermodynamics.} With $E=F+TS_s$ (the energy), the free energy $F(T)$ is the minimum of the functional at temperature $T$. Because $F$ is stationary with respect to orbitals and occupations and the constraints do not depend on $T$, only the explicit $T$ in $-TS_s$ contributes to $dF/dT$ (the \textbf{envelope theorem}: if $F(T)=G(T,x_\ast(T))$ with $\partial_xG=0$ at $x_\ast$, then $dF/dT=\partial_TG+\partial_xG\cdot dx_\ast/dT=\partial_TG$). For a functional in which $T$ appears only there, as in Chapter 13,

\[
\frac{dF}{dT}=-S_s,\qquad C_V=\frac{dE}{dT}=\frac{dF}{dT}+S_s+T\frac{dS_s}{dT}=T\frac{dS_s}{dT},
\]

where $C_V$ is the \textbf{heat capacity} at fixed particle number.

\textbf{Worked example.} Two levels $\varepsilon_0=0$ and $\varepsilon_1=1$, one particle, $T=\tfrac12$, no interaction. By symmetry $\mu=\tfrac12$ (then $f_0+f_1=\frac1{e^{-1}+1}+\frac1{e^{1}+1}=1$). The numbers are $f_0=0.731059$, $f_1=0.268941$, $E=f_1\varepsilon_1=0.268941$, $S_s=1.164406$ (each level contributes $0.582203$), $F=E-TS_s=-0.313262$. A numerical derivative of $F(T)$ at $T=\tfrac12$ gives $-1.164406=-S_s$, and $C_V=0.393224$.

\subsection{Excited states: the Kohn--Sham gap, particle--hole pairs and Delta-SCF}

The ground-state theory says nothing directly about excited states. Three quantities are used in practice, and Chapter 13 reports all three.

\textbf{Particle--hole excitations and the Kohn--Sham gap.} In a non-interacting system, moving one particle from an occupied orbital $i$ (it leaves a \textbf{hole}) to an empty orbital $a$ (a \textbf{particle}) costs exactly $\varepsilon_a-\varepsilon_i$, because the levels do not depend on the occupations. The list of these differences is the \textbf{particle--hole spectrum}, and its smallest member,

\[
\Delta_{KS}=\varepsilon_{\text{LUMO}}-\varepsilon_{\text{HOMO}},
\]

is the \textbf{Kohn--Sham gap} (LUMO: the lowest unoccupied orbital). In the interacting system the Kohn--Sham levels move when the occupations change, so $\Delta_{KS}$ is only the first estimate of the lowest excitation energy.

\textbf{Janak's theorem.} Let the energy $E(\{f\})$ of Section 12.14 be evaluated with orbitals that are self-consistent for the given occupations. Then

\[
\frac{\partial E}{\partial f_a}=\varepsilon_a .
\]

\textbf{Proof.} $E$ depends on $f_a$ explicitly and through the orbitals. The explicit derivative is $\varepsilon_a$ (Section 12.14). The implicit part is $\sum_bf_b\bigl(\langle d\varphi_b|h_s\varphi_b\rangle+\langle h_s\varphi_b|d\varphi_b\rangle\bigr)=\sum_bf_b\varepsilon_b\,d\langle\varphi_b|\varphi_b\rangle=0$, because the Kohn--Sham equations $h_s\varphi_b=\varepsilon_b\varphi_b$ hold and the orbitals stay normalized. $\square$

\textbf{Delta-SCF.} The \textbf{Delta-SCF} method computes an excited state as a second self-consistent solution in which the occupations are held fixed by hand: one particle is taken out of the HOMO and put into the LUMO, and the Kohn--Sham loop is run again with these occupations. The excitation energy is the difference of the two total energies, $\Delta_{\text{SCF}}=E_1-E_0$. Janak's theorem connects it with the gap. Move the particle gradually, with $f_{\text{HOMO}}=1-\tau$ and $f_{\text{LUMO}}=\tau$; then $dE/d\tau=\varepsilon_{\text{LUMO}}(\tau)-\varepsilon_{\text{HOMO}}(\tau)$, and integrating from $\tau=0$ to $\tau=1$,

\[
\Delta_{\text{SCF}}=\int_0^1\bigl[\varepsilon_{\text{LUMO}}(\tau)-\varepsilon_{\text{HOMO}}(\tau)\bigr]\,d\tau .
\]

At $\tau=0$ the integrand is $\Delta_{KS}$; the difference $\Delta_{\text{SCF}}-\Delta_{KS}$ measures how much the levels move when the particle is transferred (the \textbf{orbital relaxation}). The midpoint value $\varepsilon_{\text{LUMO}}(\tfrac12)-\varepsilon_{\text{HOMO}}(\tfrac12)$ (Slater's \textbf{transition state}) is a good estimate of the integral when the integrand is nearly linear in $\tau$. That the Delta-SCF state approximates a true excited state is an assumption of the method; it is well founded for the lowest state of a given symmetry and an approximation in general.

\textbf{Toy example.} Take the model energy $E(f_H,f_L)=\varepsilon_H^0f_H+\varepsilon_L^0f_L+\tfrac U2\bigl(f_H^2+f_L^2\bigr)$ with $\varepsilon_H^0=0$, $\varepsilon_L^0=1$ and $U=0.2$ (the last term plays the role of a self-interaction of each orbital). Janak's theorem gives the levels $\varepsilon_H=Uf_H$ and $\varepsilon_L=1+Uf_L$. In the ground state $(f_H,f_L)=(1,0)$: $E_0=0.1$, $\varepsilon_H=0.2$, $\varepsilon_L=1$, so $\Delta_{KS}=0.8$. In the Delta-SCF state $(0,1)$: $E_1=1.1$, so $\Delta_{\text{SCF}}=1.0$. The integral formula reproduces it: $\int_0^1\bigl[1+U\tau-U(1-\tau)\bigr]d\tau=1$. Here the relaxation raises the excitation energy by $U=0.2$ above the Kohn--Sham gap, and the transition state is exact because the integrand is linear.

\textbf{Two further notions.} At a finite temperature the occupations are fractional, and the particle--hole spectrum is read with a threshold: a level counts as occupied (HOMO side) if $f\ge\tfrac12$ and as empty if $f<\tfrac12$ (Chapter 13 uses this rule). The \textbf{fundamental gap} $E(N+1)+E(N-1)-2E(N)$, which involves adding and removing a particle, is a different quantity again; even with the exact functional it differs from $\Delta_{KS}$ in general. We quote this as a known property of exact DFT and do not need it: Chapter 13 works at fixed $N$ and reports $\Delta_{KS}$, the particle--hole list and $\Delta_{\text{SCF}}$.

\subsection{What we proved and what we assumed}

We proved, from the definitions: that Hermitian operators have real eigenvalues and orthogonal eigenvectors, and the variational principle; that Slater determinants are antisymmetric, normalized and obey the Pauli principle, with density $\sum_a|\varphi_a|^2$; the anticommutation relations of creation and annihilation operators; Wick's theorem for a determinant and for a non-interacting thermal ensemble; the energy of a determinant from the permutation sum, with its direct and exchange terms, and its agreement with the second-quantized operators; the cancellation of the self-interaction, and the exact locality of exchange for a contact interaction with the ratio $E_x=-E_H/g$ for equally occupied labels; the method of Lagrange multipliers and the meaning of a multiplier as the derivative of the constrained minimum; the Hartree--Fock equations, the double-counting formula and Koopmans' theorem; the exact energy of the two-site model and its Hartree--Fock energy (restricted for $U\le2t$, unrestricted for $U>2t$); the Hohenberg--Kohn theorem, the constrained-search variational principle for the density, the Kohn--Sham equations and their total-energy formula; the volume element of spherical coordinates, and the kinetic and exchange energies of the uniform gas (Thomas--Fermi constant and Dirac exchange); the convergence condition of linear mixing in the two-site model; the Gibbs principle via Klein's inequality, including the uniqueness of the minimizer, and Mermin's theorem (by the same argument); the independence of the orbital occupations in the non-interacting Gibbs state, the Fermi--Dirac occupations of the Mermin--Kohn--Sham functional and the relations $dF/dT=-S_s$, $C_V=T\,dS_s/dT$; Janak's theorem and the integral formula for Delta-SCF.

We assumed: that the particles are fermions (for dirac16complex this is the choice of Grassmann components, Chapters 5 and 8); a non-degenerate ground state and a wave function that does not vanish on a region of positive volume (Theorem 12.1); that the constrained minima exist and that the densities of interest are non-interacting $v$-representable (the Kohn--Sham scheme); a finite-dimensional state space in the proof of the Gibbs principle; the implicit function theorem (in the method of Lagrange multipliers) and, for $\mu=dE_0/dN$, a minimum that depends differentiably on $N$; and that a Delta-SCF state approximates a true excited state. We quoted without derivation, and do not use later: the correlation energy of the uniform electron gas from quantum Monte Carlo, and the statement that the exact Kohn--Sham gap differs from the fundamental gap. The \textbf{approximation} that makes DFT practical is the choice of $E_{xc}$; in the local density approximation it is the uniform-gas value at the local density. For a contact interaction with scalar labels (a vertex that does not depend on the label, as in Section 12.6) the exchange part of that choice is exact for determinants and non-interacting ensembles when it is evaluated with all label densities, and what remains approximate is the omission of correlation. Chapter 13 applies to dirac16complex, with Mermin's finite-temperature functional, an exchange-only local scheme built on the uniform-gas closed form in two densities, which approximates the exact local Fock term, with no correlation term (Section 13.8); the Rust implementation is \texttt{studies/\allowbreak{}dirac16c\allowbreak{}omplex\_k\allowbreak{}ohn\_sham} (self-consistency and mixing in \texttt{src/scf.\allowbreak{}rs}, the exchange functional in \texttt{src/exch\allowbreak{}ange.rs}).

\subsection{Exercises}

\textbf{Exercise 12.1.} In the three-point example of Section 12.4, compute $\Phi(2,1)$ and $\Phi(1,1)$, and show directly from the definition that the two-particle determinant vanishes identically if $\varphi_a=\varphi_b$.

\textbf{Exercise 12.2.} Using only the anticommutation relations, show that $\hat n_p^2=\hat n_p$ (so $\hat n_p$ has the eigenvalues 0 and 1) and that $a_p^\dagger a_p^\dagger=0$.

\textbf{Exercise 12.3.} For the same example, write the one-body density matrix $\rho(x,x')$ as a $3\times3$ matrix, and check $\rho^2=\rho$ and $\mathrm{Tr}\,\rho=2$.

\textbf{Exercise 12.4.} Solve the two-site model of Section 12.8 for $t=1$, $U=4$: exact energy, Hartree--Fock energy, correlation energy, and the exact probability of double occupancy.

\textbf{Exercise 12.5.} (a) Compute $\delta/\delta n$ of $\int n^{5/3}\,d^3r$ and of the LDA exchange energy $\int e_x(n)\,d^3r$. (b) Minimize the Thomas--Fermi energy $C_F\int n^{5/3}+\int vn+E_H[n]$ at fixed $\int n=N$ and write the resulting equation for $n$.

\textbf{Exercise 12.6.} For a uniform gas with $g$ internal states in three dimensions, show that the kinetic energy per particle is $\tfrac3{10}k_F^2$, independent of $g$.

\textbf{Exercise 12.7.} Repeat the mixing analysis of Section 12.13 for $\Delta=2$, $t=1$, $U=2$. Does plain iteration ($\beta=1$) converge?

\textbf{Exercise 12.8.} For the two-level example of Section 12.14 find the limits of $f_0$, $E$ and $S_s$ as $T\to0$ and as $T\to\infty$.

\textbf{Exercise 12.9.} For non-interacting levels with a fixed spectrum and fixed $N$, show that $C_V=\frac1{T^2}\Bigl[\sum_aw_a(\varepsilon_a-\mu)^2-\bigl(\sum_aw_a(\varepsilon_a-\mu)\bigr)^2/\sum_aw_a\Bigr]$ with $w_a=f_a(1-f_a)$, and conclude $C_V\ge0$.

\textbf{Exercise 12.10.} For the model energy $E=\sum_a\varepsilon_a^0f_a+\tfrac U2\sum_af_a^2$ of Section 12.15 with arbitrary $\varepsilon_H^0<\varepsilon_L^0$, show that $\Delta_{\text{SCF}}-\Delta_{KS}=U$.

\textbf{Exercise 12.11.} Electrons with a contact interaction of strength $g_c=1$ have the spin densities $n_\uparrow=0.3$, $n_\downarrow=0.1$ at a point. Compute the Hartree and exchange energy densities and check that their sum is $n_\uparrow n_\downarrow$.

\subsection{Answers to the exercises}

\textbf{Answer 12.1.} $\Phi(2,1)=\tfrac1{\sqrt2}[\varphi_0(2)\varphi_1(1)-\varphi_1(2)\varphi_0(1)]=\tfrac1{\sqrt2}[0\cdot\tfrac1{\sqrt2}-\tfrac1{\sqrt2}\cdot0]=0$, and $\Phi(1,1)=\tfrac1{\sqrt2}[\varphi_0(1)\varphi_1(1)-\varphi_1(1)\varphi_0(1)]=0$. If $\varphi_a=\varphi_b=\varphi$, then $\Phi(x_0,x_1)=\tfrac1{\sqrt2}[\varphi(x_0)\varphi(x_1)-\varphi(x_0)\varphi(x_1)]=0$ for all arguments.

\textbf{Answer 12.2.} $\hat n_p^2=a_p^\dagger a_pa_p^\dagger a_p=a_p^\dagger(1-a_p^\dagger a_p)a_p=\hat n_p-a_p^\dagger a_p^\dagger a_pa_p$. From $\{a_p^\dagger,a_p^\dagger\}=0$ we get $2a_p^\dagger a_p^\dagger=0$, so the last term vanishes and $\hat n_p^2=\hat n_p$. An eigenvalue $\nu$ of $\hat n_p$ satisfies $\nu^2=\nu$, so $\nu\in\{0,1\}$.

\textbf{Answer 12.3.} $\rho=\varphi_0\varphi_0^\dagger+\varphi_1\varphi_1^\dagger=\begin{pmatrix}1&0&0\\0&\tfrac12&\tfrac12\\0&\tfrac12&\tfrac12\end{pmatrix}$. Squaring, the lower $2\times2$ block gives $\begin{pmatrix}\tfrac12&\tfrac12\\ \tfrac12&\tfrac12\end{pmatrix}^2=\begin{pmatrix}\tfrac12&\tfrac12\\ \tfrac12&\tfrac12\end{pmatrix}$ and the corner gives $1$, so $\rho^2=\rho$; the trace is $1+\tfrac12+\tfrac12=2$, the particle number. The diagonal $(1,\tfrac12,\tfrac12)$ is the density found in Section 12.4.

\textbf{Answer 12.4.} $E_0=\tfrac12(4-\sqrt{32})=2-2\sqrt2=-0.828427$. The restricted determinant gives $-2+2=0$. Since $U=4>2t$, the best determinant is the unrestricted one of Section 12.8 with $\sin2\alpha=2t/U=\tfrac12$: $E_{HF}=-2t^2/U=-0.5$ (check: $-t\cdot2\cdot\tfrac12+\tfrac U2\bigl(1-\cos^22\alpha\bigr)=-1+2\cdot\tfrac14=-0.5$), so $E_c=E_0-E_{HF}=-0.328427$. The lowest eigenvector of $\begin{pmatrix}4&-2\\-2&0\end{pmatrix}$ has components $(D,S)\propto(1,1+\sqrt2)$ (from $(4-E_0)D=2S$), so the exact double occupancy is $D^2=1/\bigl(1+(1+\sqrt2)^2\bigr)=1/(4+2\sqrt2)=(2-\sqrt2)/4=0.146447$, much smaller than the restricted value $\tfrac12$: the stronger the repulsion, the more the electrons avoid each other. The unrestricted value is $n_{\mathrm L\uparrow}n_{\mathrm L\downarrow}+n_{\mathrm R\uparrow}n_{\mathrm R\downarrow}=\tfrac12\sin^22\alpha=0.125$, below the exact value: by breaking the left--right symmetry of each spin, the determinant over-separates the electrons.

\textbf{Answer 12.5.} (a) $\tfrac53n^{2/3}$, and $de_x/dn=-\tfrac43\cdot\tfrac34(3/\pi)^{1/3}n^{1/3}=-(3/\pi)^{1/3}n^{1/3}$. (b) With a multiplier $\mu$: $\tfrac53C_Fn^{2/3}+v+v_H=\mu$ wherever $n>0$, that is $n=\bigl[\tfrac3{5C_F}(\mu-v-v_H)\bigr]^{3/2}$, an equation for $n$ because $v_H$ depends on $n$ (the Thomas--Fermi equation).

\textbf{Answer 12.6.} $t/n=\dfrac{g\,k_F^5/(20\pi^2)}{g\,k_F^3/(6\pi^2)}=\dfrac{6}{20}k_F^2=\dfrac3{10}k_F^2$; the factor $g$ cancels. (It is $\tfrac35$ of the largest kinetic energy $\tfrac12k_F^2$, the average of $k^2/2$ over a full sphere.)

\textbf{Answer 12.7.} The fixed point of $G(x)=(2-2x)/\sqrt{(2-2x)^2+4}$ is $x_\ast=0.468990$ ($n_\mathrm L=1.468990$), where $G'(x_\ast)=-4Ut^2/\bigl((\Delta-Ux_\ast)^2+4t^2\bigr)^{3/2}=-0.688942$. Linear mixing converges for $0<\beta<2/(1-G')=1.184174$. Plain iteration, $\beta=1$, multiplies the error by $-0.689$ per step and converges (slowly, alternating). A weaker repulsion makes the feedback weaker.

\textbf{Answer 12.8.} As $T\to0$: $f_0\to1$, $E\to0$, $S_s\to0$ (a pure state). As $T\to\infty$: $f_0,f_1\to\tfrac12$, $E\to\tfrac12$, and each level contributes $-2\cdot\tfrac12\ln\tfrac12=\ln2$, so $S_s\to2\ln2=\ln4$: the four configurations (both empty, one or the other occupied, both occupied) become equally likely, and their average particle number is still 1.

\textbf{Answer 12.9.} With a fixed spectrum $E=\sum_af_a\varepsilon_a$ and $\sum_af_a=N$, so $\sum_a\partial_Tf_a=0$ and $C_V=\sum_a\varepsilon_a\,\partial_Tf_a=\sum_a(\varepsilon_a-\mu)\,\partial_Tf_a$. Differentiating the Fermi function, $\partial_Tf_a=w_a\bigl[(\varepsilon_a-\mu)/T^2+\mu'/T\bigr]$ with $\mu'=d\mu/dT$ and $w_a=f_a(1-f_a)$. The condition $\sum_a\partial_Tf_a=0$ gives $\mu'=-\sum_aw_a(\varepsilon_a-\mu)/\bigl(T\sum_aw_a\bigr)$. Inserting, $C_V$ equals the stated expression. The bracket is $\sum_aw_a$ times the variance of $\varepsilon_a-\mu$ with the weights $w_a/\sum_bw_b$, which is never negative.

\textbf{Answer 12.10.} Janak gives $\varepsilon_a=\varepsilon_a^0+Uf_a$. Ground state $(1,0)$: $\Delta_{KS}=\varepsilon_L^0-(\varepsilon_H^0+U)$. Energies: $E_0=\varepsilon_H^0+\tfrac U2$ and $E_1=\varepsilon_L^0+\tfrac U2$, so $\Delta_{\text{SCF}}=\varepsilon_L^0-\varepsilon_H^0$ and $\Delta_{\text{SCF}}-\Delta_{KS}=U$.

\textbf{Answer 12.11.} $n=0.4$: $e_H=\tfrac12n^2=0.08$; $e_x=-\tfrac12(n_\uparrow^2+n_\downarrow^2)=-\tfrac12(0.09+0.01)=-0.05$; the sum $0.03$ equals $n_\uparrow n_\downarrow=0.3\cdot0.1$.

\section{The Kohn--Sham approximation for dirac16complex in the primordial field}

\subsection{What this chapter computes, and the status of its numbers}

This chapter applies the density functional theory of Chapter 12 to dirac16complex. The system is a finite number $N$ of dirac16complex quanta (particles above the Dirac sea, Chapter 8) held in the \textbf{static} member of the notebook's primordial gravitational field (Chapter 9), extended by the notebook's mirror construction, a \textbf{Z2 brane}. The quanta interact through the contact potential $U(S)=\tfrac\lambda2S^2$ of the Lagrangian (Chapter 6). We derive, step by step: the reduction of the 16-component field equation to eight copies of a $2\times2$ first-order system in one coordinate; the boundary conditions; exact solutions at zero momentum; the Hartree and the exact exchange energy of the contact interaction for the uniform gas, which gives a local density approximation without any fitted number; Mermin's finite-temperature Kohn--Sham functional and its equations; and the energy--momentum tensor of the resulting state. Then we report what the computer found: ground states, first excited states (Kohn--Sham gap, particle--hole list and Delta-SCF), thermodynamics, and the comparison with the source that Einstein gravity would need to produce the field.

\textbf{Sources.} The binding specification is \texttt{handoff/\allowbreak{}specs/ST\allowbreak{}AGE4\_SPE\allowbreak{}C.md}, including its errata E4.1 to E4.13 (its §7 to §9), which override its earlier text. The exact theory is \texttt{wolfram/\allowbreak{}Dirac16C\allowbreak{}omplexKo\allowbreak{}hnSham.w\allowbreak{}l} with the verifier \texttt{scripts/\allowbreak{}verify\_d\allowbreak{}irac16co\allowbreak{}mplex\_ko\allowbreak{}hn\_sham.\allowbreak{}wls}, whose report \texttt{artifact\allowbreak{}s/dirac1\allowbreak{}6complex\allowbreak{}/kohn-sh\allowbreak{}am/wolfr\allowbreak{}am-kohn-\allowbreak{}sham-rep\allowbreak{}ort.json} has 125 of 125 checks true, and the independent sympy checker \texttt{scripts/\allowbreak{}check\_di\allowbreak{}rac16com\allowbreak{}plex\_koh\allowbreak{}n\_sham\_t\allowbreak{}heory.py}, whose report \texttt{artifact\allowbreak{}s/dirac1\allowbreak{}6complex\allowbreak{}/kohn-sh\allowbreak{}am/pytho\allowbreak{}n-theory\allowbreak{}-report.\allowbreak{}json} has 157 of 157 checks true. The exact formulas are collected in \texttt{artifact\allowbreak{}s/dirac1\allowbreak{}6complex\allowbreak{}/kohn-sh\allowbreak{}am/kohn-\allowbreak{}sham-the\allowbreak{}ory.json}. Check names in typewriter type beginning with \texttt{\detokenize{KS_}} are checks of these two reports. The numbers come from the Rust solver \texttt{studies/\allowbreak{}dirac16c\allowbreak{}omplex\_k\allowbreak{}ohn\_sham} (its \texttt{README.m\allowbreak{}d} describes the implementation), whose canonical outputs are under \texttt{artifact\allowbreak{}s/dirac1\allowbreak{}6complex\allowbreak{}/kohn-sh\allowbreak{}am/rust/}.

\textbf{Status of the numbers (read this first).} Every number below is \textbf{as computed by the Rust solver} and verified by the solver's own checks: all 397 checks of the five subcommands pass (33 in \texttt{\detokenize{spectrum}}, 137 in \texttt{\detokenize{scf}}, 65 in \texttt{\detokenize{excited}}, 102 in \texttt{\detokenize{thermo}}, 60 in \texttt{\detokenize{emt}}; the \texttt{summary.\allowbreak{}json} file of each subfolder), a repeated run is byte-identical in all 330 files compared, and a run with the tolerances divided by 10 and half the largest step agrees with the canonical one to $5.96\times10^{-8}$ relative in the energies $E$ and $F$ (65 runs) and to $2.8\times10^{-8}$ in 364829 eigenvalues (\texttt{rust/det\allowbreak{}erminism\allowbreak{}-report.\allowbreak{}json}). The numbers were then compared, run by run, with those of an independent Python reference solver (\texttt{scripts/\allowbreak{}ks\_refer\allowbreak{}ence\_sol\allowbreak{}ver.py}, 56 runs under \texttt{artifact\allowbreak{}s/dirac1\allowbreak{}6complex\allowbreak{}/kohn-sh\allowbreak{}am/refer\allowbreak{}ence/}) by the checker \texttt{scripts/\allowbreak{}check\_di\allowbreak{}rac16com\allowbreak{}plex\_koh\allowbreak{}n\_sham.p\allowbreak{}y}. Its final report, \texttt{artifact\allowbreak{}s/dirac1\allowbreak{}6complex\allowbreak{}/kohn-sh\allowbreak{}am/pytho\allowbreak{}n-check-\allowbreak{}report.j\allowbreak{}son}, has \textbf{63 checks, of which 62 are true}. The one failed check, \texttt{canonica\allowbreak{}l\_eigenv\allowbreak{}alues}, concerns two deep levels of the Dirac sea in the 301-point run of a single state, the smeared ensemble at $N=1016$ and $-\hat\lambda_2$ (Sections 13.11 and 13.15); the Delta-SCF energy of the same state, which an earlier comparison had left open, now agrees (erratum E4.13; Section 13.12). The Stage-4 gate (\texttt{scripts/\allowbreak{}verify\_s\allowbreak{}tage4\_ko\allowbreak{}hn\_sham.\allowbreak{}ps1} and \texttt{\detokenize{.sh}}), which reruns the programs of the stage in a fresh copy of the repository (the reference solver in a reduced form), has not been run. Section 13.15 states the cross-check in full.

\textbf{What is not done here.} The metric is a fixed background: the quanta do not act back on it. The Kohn--Sham state is a static mean-field state with the exact exchange of the uniform gas and \textbf{no correlation energy} (Section 13.8 explains this choice). No claim is made that such a state existed in any universe.

\textbf{Units and conventions.} $H=1$ (the notebook's inverse length); the mass is $m=1$ or $m=3$ in these units, and energies, momenta and temperatures are quoted in units of $m$. The coupling is given as the dimensionless $\hat\lambda=\lambda m^6$. Coordinates are $x_0,\dots,x_7$ with $x_4$ the time, $\eta=\mathrm{diag}(+1,+1,+1,+1,-1,-1,-1,-1)$, the gammas are the notebook's $\gamma^a$ (Chapter 2), $C=\gamma^0\gamma^1\gamma^2\gamma^3$, $\bar\Psi=\Psi^\dagger C$, and $B=-iC\gamma^4$ (Chapter 8). The letter $l$ is used as an index for the six warped directions, $l\in\{1,2,3,5,6,7\}$ (3-space when $l\le3$); the letter $j$, which Chapters 0, 4 and 9 use for the extra times, is here reserved for the block type of Section 13.4. The Pauli matrices are $\sigma_x,\sigma_y,\sigma_z$ of Chapter 2 (the subscript $y$ of $\sigma_y$ names the matrix and has nothing to do with the coordinate $y$ of Section 13.2).

\subsection{The static primordial field, the hidden coordinate and the brane}

\textbf{The static field.} The notebook's primordial metric (Chapter 9) contains a free function $a_4(t)$ of the time $t=Hx_4$. A stationary ground state needs a metric that does not change in time, so Stage 4 takes the \textbf{static member} $a_4=a_{4,0}$, a constant ($a_4'=a_4''=0$).

\textbf{The hidden coordinate $y$.} Stage 4 and this chapter call the proper hidden coordinate $y$; it is the coordinate $\zeta$ of Chapters 1, 4, 9 and 11. With $z=6Hx_0\in(0,\pi/2)$ it is defined by

\[
y=\frac{\ln\sin z}{6H}\in(-\infty,0],\qquad \sin z=e^{6Hy},\qquad dy=\cot z\,dx_0 .
\]

(Differentiate: $dy/dx_0=\frac1{6H}\cdot\frac{\cos z}{\sin z}\cdot6H=\cot z$.) Then $g_{00}\,dx_0^2=\cot^2z\,dx_0^2=dy^2$ and $(\sin z)^{1/3}=e^{2Hy}$, and the line element of Chapter 9 becomes the \textbf{warped form}

\[
ds^2=dy^2-dx_4^2+e^{2Hy}\Bigl[e^{2a_{4,0}}\bigl(dx_1^2+dx_2^2+dx_3^2\bigr)-e^{-2a_{4,0}}\bigl(dx_5^2+dx_6^2+dx_7^2\bigr)\Bigr].
\]

The factor $W(y)=e^{Hy}$ is the \textbf{warp}. The square root of the determinant is the product of the eight scale factors, $\sqrt{|g|}=(e^{Hy+a_{4,0}})^3\cdot1\cdot1\cdot(e^{Hy-a_{4,0}})^3=W^6=e^{6Hy}$ (\texttt{KS\_geome\allowbreak{}try\_sqrt\allowbreak{}DetG\_W6}): the proper 7-volume per unit coordinate volume. It tends to 0 as $y\to-\infty$ ($z\to0$), where the transverse space pinches off; we call this end the \textbf{tip}. The notebook's coordinate patch ends at $y=0$ ($z=\pi/2$), where $\cot z=0$ and the $x_0$ chart degenerates, although the $y$ form of the metric is regular there.

\textbf{Curvature and the required source.} Setting $a_4'=a_4''=0$ in the curvature of Chapter 9, $R=6H^2(a_4'^2-7)$, $G^0{}_0=-3H^2(a_4'^2-5)$, $G^1{}_1=G^2{}_2=G^3{}_3=H^2(15-3a_4'^2+a_4'')$, $G^4{}_4=3H^2(7+a_4'^2)$ and $G^5{}_5=G^6{}_6=G^7{}_7=H^2(15-3a_4'^2-a_4'')$, gives

\[
R=-42H^2,\qquad G^\mu{}_\nu=\mathrm{diag}(15,15,15,15,21,15,15,15)\,H^2
\]

in the order $(y,x_1,x_2,x_3,x_4,x_5,x_6,x_7)$ (\texttt{KS\_geome\allowbreak{}try\_ricc\allowbreak{}iScalarM\allowbreak{}inus42H2}, \texttt{KS\_geome\allowbreak{}try\_eins\allowbreak{}teinMixe\allowbreak{}dDiag}). There is a short way to see this. The seven directions other than $x_4$ form a space of constant curvature $-H^2$ (\texttt{KS\_geome\allowbreak{}try\_cons\allowbreak{}tantCurv\allowbreak{}atureSev\allowbreak{}enSpace}), for which $R_{\mu\nu}=-(7-1)H^2g_{\mu\nu}=-6H^2g_{\mu\nu}$, while the time $x_4$ is flat, $R^4{}_4=0$. So $R=7\cdot(-6H^2)=-42H^2$, and $G^\mu{}_\nu=R^\mu{}_\nu-\tfrac12\delta^\mu{}_\nu R$ is $-6H^2+21H^2=15H^2$ on the seven directions and $0+21H^2=21H^2$ on $x_4$. If 8-dimensional Einstein gravity $G^\mu{}_\nu=\kappa T^\mu{}_\nu$ produced this field, the source would have to be (erratum E4.3; \texttt{KS\_geome\allowbreak{}try\_requ\allowbreak{}iredSour\allowbreak{}ce}, \texttt{KS\_geome\allowbreak{}try\_rhoR\allowbreak{}equiredN\allowbreak{}egative})

\[
\rho_{\text{req}}=-T^4{}_4=-\frac{21H^2}\kappa<0,\qquad p_{\text{req}}=T^\mu{}_\mu=+\frac{15H^2}\kappa\quad(\text{all seven }\mu\ne4),\qquad w_{\text{req}}=-\frac57 .
\]

A \textbf{negative} energy density is required. Section 13.14 compares this with the Kohn--Sham state.

\textbf{Extrinsic curvature.} A surface $y=\text{const}$ with the unit normal $\partial_y$ has the extrinsic curvature $K_{\mu\nu}=\tfrac12\partial_yg_{\mu\nu}$. For the six warped directions $g_{\mu\mu}=\pm e^{2Hy\pm2a_{4,0}}$, so $K_{\mu\mu}=Hg_{\mu\mu}$ (no sum) and $K^\mu{}_\nu=H\,\delta^\mu{}_\nu$ there, while $K^4{}_4=0$; the trace is $K=6H$ (\texttt{KS\_geome\allowbreak{}try\_extr\allowbreak{}insicCur\allowbreak{}vature}).

\textbf{The Z2 brane.} The notebook's hypothesis speaks of a pair of universes. Stage 4 realizes a geometric version of it (extension E2 of the specification): two copies of the patch $y\le0$ glued at $y=0$, with the warp $W=e^{-H|y|}$ on both sides. The surface $y=0$ is then a \textbf{brane}, a wall at which the extrinsic curvature jumps. With the normal $n=+\partial_y$ pointing from $y<0$ to $y>0$, $K^\mu{}_\mu(0^-)=+H$ and $K^\mu{}_\mu(0^+)=-H$ (no sum) on the six warped directions, so the jump $[X]=X(0^+)-X(0^-)$ is $[K^\mu{}_\mu]=-2H$ there, $0$ on $x_4$, and $[K]=6\cdot(-2H)=-12H$. The \textbf{Israel junction condition}, with the convention $[K_{\mu\nu}]-h_{\mu\nu}[K]=-\kappa S_{\mu\nu}$ ($h_{\mu\nu}$ the induced metric of the brane, $S_{\mu\nu}$ the stress carried by the brane, $\mu,\nu\ne y$), gives

\[
\begin{aligned}
&S^\mu{}_\nu=-\frac1\kappa\bigl([K^\mu{}_\nu]-\delta^\mu{}_\nu[K]\bigr):\\
&S^\mu{}_\mu=-\frac{-2H+12H}\kappa=-\frac{10H}\kappa\ \ (\text{six warped directions}),\qquad S^4{}_4=-\frac{0+12H}\kappa=-\frac{12H}\kappa ,
\end{aligned}
\]

so the brane carries the energy density $\rho_{\text{brane}}=-S^4{}_4=+12H/\kappa>0$ and the pressures $-10H/\kappa$ (\texttt{KS\_geome\allowbreak{}try\_isra\allowbreak{}elJump}, \texttt{KS\_geome\allowbreak{}try\_isra\allowbreak{}elStress}, \texttt{KS\_geome\allowbreak{}try\_bran\allowbreak{}eEnergyP\allowbreak{}ositive}; the same convention gives the Randall--Sundrum brane a positive tension, \texttt{KS\_geome\allowbreak{}try\_isra\allowbreak{}elConven\allowbreak{}tionRand\allowbreak{}allSundr\allowbreak{}um}). This is a statement about what the glued geometry would need; the chapter does not derive that such a brane exists. On the $y>0$ side the field equation below has the same form with $y\to-y$, so the Kohn--Sham problem is solved on $y\in[-L,0]$ with boundary conditions at the brane (Section 13.5).

\subsection{The stationary ansatz and the reduced equation}

\textbf{The field equation.} In the mean field the quanta obey the Dirac equation of Chapter 7 with an effective mass: $\gamma^\mu D_\mu\Psi=M_{\text{eff}}\Psi$, where $M_{\text{eff}}=m+\lambda S$ at the Hartree level (Section 13.8 adds the exchange). We work in the \textbf{good sector}: $\Psi$ does not depend on the extra times $x_5,x_6,x_7$. Modes with momentum along the extra times have a non-Hermitian single-particle operator and grow without bound (Chapter 8), so they are excluded.

\textbf{The pieces of the Dirac operator.} The vielbein is diagonal, with the entries (in the order $y,x_1,\dots,x_7$)

\[
h=\bigl(1,\ e^{Hy+a_{4,0}},\ e^{Hy+a_{4,0}},\ e^{Hy+a_{4,0}},\ 1,\ e^{Hy-a_{4,0}},\ e^{Hy-a_{4,0}},\ e^{Hy-a_{4,0}}\bigr),
\]

so the curved gammas are $\gamma^y=\gamma^0$, $\gamma^{x_l}=e^{-Hy-a_{4,0}}\gamma^l$ for $l=1,2,3$, $\gamma^{x_4}=\gamma^4$, and $\gamma^{x_l}=e^{-Hy+a_{4,0}}\gamma^l$ for $l=5,6,7$. The spin connection enters only through the contraction $\gamma^\mu\Omega_\mu$ (Chapters 4 and 9). For a diagonal vielbein that depends on one coordinate only, the formula of Chapter 4, $\gamma^\mu\Omega_\mu=\tfrac12\sum_b\frac1{h_b}\,\partial_b\ln\bigl(\prod_{c\ne b}h_c\bigr)\gamma^b$, has a single term, $b=y$, and $\prod_{c\ne y}h_c=e^{6Hy}$:

\[
\gamma^\mu\Omega_\mu=\tfrac12\,\partial_y(6Hy)\,\gamma^0=3H\gamma^0
\]

(\texttt{KS\_reduc\allowbreak{}tion\_gam\allowbreak{}maSlashO\allowbreak{}mega3Hga\allowbreak{}mma0}).

\textbf{The ansatz.} Stationary in time with frequency $\varepsilon$, a plane wave with 3-momentum $\mathbf k=(k_1,k_2,k_3)$ along 3-space, and a factor that will cancel the spin connection:

\[
\Psi=e^{-i\varepsilon x_4}\,e^{i\mathbf k\cdot\mathbf x}\,W(y)^{-3}\,\chi(y),\qquad W^{-3}=e^{-3Hy},
\]

with $\chi(y)$ a column of 16 complex functions of $y$ alone. Now compute each term. $\partial_y\bigl(e^{-3Hy}\chi\bigr)=e^{-3Hy}(\chi'-3H\chi)$, so $\gamma^0\partial_y\Psi+3H\gamma^0\Psi$ equals the common factor $e^{-i\varepsilon x_4}e^{i\mathbf k\cdot\mathbf x}e^{-3Hy}$ times $\gamma^0\chi'$: the $-3H$ from the derivative cancels the $+3H$ of the spin connection. The time derivative gives $\gamma^4\partial_4\to-i\varepsilon\gamma^4$ and the 3-space derivatives give $\gamma^{x_l}\partial_l\to ik_le^{-Hy-a_{4,0}}\gamma^l$. Dividing by the common factor,

\[
\gamma^0\chi'+i\xi(y)\sum_{l=1}^3k_l\gamma^l\chi-i\varepsilon\gamma^4\chi=M_{\text{eff}}(y)\,\chi,\qquad \xi(y)=e^{-Hy-a_{4,0}}
\]

(\texttt{KS\_reduc\allowbreak{}tion\_ans\allowbreak{}atzRemov\allowbreak{}es3H}, \texttt{KS\_reduc\allowbreak{}tion\_red\allowbreak{}ucedEqua\allowbreak{}tionODEF\allowbreak{}orm}; without the factor $W^{-3}$ the term $3H\gamma^0$ would survive, \texttt{KS\_reduc\allowbreak{}tion\_wit\allowbreak{}houtW3th\allowbreak{}e3HTermS\allowbreak{}urvives}). We write $\xi$ (the Greek letter xi) for this momentum weight, $\xi(y)=e^{-a_{4,0}}/W(y)$, because $\kappa$ is the gravitational coupling of Section 13.2; the theory file and the Rust code call the weight \texttt{\detokenize{kappa}}. Multiplying by $\gamma^0$, whose square is $+1$ because $\eta^{00}=+1$, gives an ordinary differential equation for $\chi$:

\[
\chi'=\gamma^0\Bigl[M_{\text{eff}}\chi-i\xi\sum_{l=1}^3k_l\gamma^l\chi+i\varepsilon\gamma^4\chi\Bigr].
\]

\textbf{Three properties.} (i) \textbf{Flat measure.} $\sqrt{|g|}\,\Psi^\dagger\Psi=e^{6Hy}e^{-6Hy}\chi^\dagger\chi=\chi^\dagger\chi$, so an orbital is normalized by $\int_{-L}^0\chi^\dagger\chi\,dy=1$ with the ordinary measure $dy$ (\texttt{KS\_reduc\allowbreak{}tion\_fla\allowbreak{}tMeasure}). (ii) \textbf{Confinement.} The momentum enters with the weight $\xi(y)=e^{-Hy-a_{4,0}}$, which grows toward the tip: a given momentum costs more energy far from the brane, which pushes the quanta toward $y=0$. (iii) \textbf{The constant $a_{4,0}$ is a rescaling.} It appears only in $\xi$, as $k\,e^{-a_{4,0}}$ (\texttt{KS\_reduc\allowbreak{}tion\_a4I\allowbreak{}sMomentu\allowbreak{}mRescali\allowbreak{}ng}); the canonical runs use $a_{4,0}=0$, and Section 13.15 checks the equivalence numerically.

\textbf{Choosing the direction of $\mathbf k$.} Rotations of $(x_1,x_2,x_3)$ are symmetries, so the spectrum depends only on $k=|\mathbf k|$ (\texttt{KS\_reduc\allowbreak{}tion\_rot\allowbreak{}ationalS\allowbreak{}ymmetry}), and we take $\mathbf k=(k,0,0)$. Then only $\gamma^0$, $\gamma^1$ and $\gamma^4$ appear:

\[
\chi'=\bigl[M_{\text{eff}}\,A_0-i\xi k\,A_1+i\varepsilon A_4\bigr]\chi,\qquad A_0=\gamma^0,\quad A_1=\gamma^0\gamma^1,\quad A_4=\gamma^0\gamma^4 .
\]

On the mirror side $y>0$ of the Z2 geometry the same equation holds with $\xi(y)=e^{Hy-a_{4,0}}$: $\xi$ is an even function of $y$ (\texttt{KS\_reduc\allowbreak{}tion\_mir\allowbreak{}rorPatch\allowbreak{}SameForm}).

\subsection{Eight \texorpdfstring{$2\times2$}{2×2} blocks}

The reduced equation still couples 16 components. This section shows that a fixed change of basis splits it into eight independent systems of two components each, and computes them exactly.

\textbf{Two rules for products of gamma matrices.} Let $P=\gamma^{a_1}\gamma^{a_2}\cdots\gamma^{a_k}$ be a product of $k$ different frame gammas. Moving one gamma $\gamma^c$ from the left of $P$ to its right, it passes each factor once; it anticommutes with every factor $\gamma^{a}\ne\gamma^c$ and commutes with itself. Hence

\[
\gamma^cP=(-1)^{k-1}P\gamma^c\ \ \text{if }\gamma^c\text{ is a factor of }P,\qquad
\gamma^cP=(-1)^{k}P\gamma^c\ \ \text{if it is not}.
\]

Reversing the order of the $k$ factors needs $k(k-1)/2$ swaps, so $P^2=(-1)^{k(k-1)/2}\,(\gamma^{a_1})^2\cdots(\gamma^{a_k})^2$, where $(\gamma^a)^2=\eta^{aa}=+1$ for $a\le3$ and $-1$ for $a\ge4$.

\textbf{The algebra of the equation.} With these rules: $A_0^2=(\gamma^0)^2=1$; $A_1^2=(\gamma^0\gamma^1)^2=(-1)^1\cdot1\cdot1=-1$; $A_4^2=(\gamma^0\gamma^4)^2=(-1)\cdot1\cdot(-1)=+1$; and the three matrices anticommute pairwise (for example $A_0A_1=\gamma^1$ while $A_1A_0=\gamma^0\gamma^1\gamma^0=-\gamma^1$). Three pairwise anticommuting matrices with squares $+1,-1,+1$ generate a Clifford algebra of signature (2,1), whose dimension is 8 (\texttt{KS\_reduc\allowbreak{}tion\_alg\allowbreak{}ebraDim8}).

\textbf{Three commuting labels.} Define

\[
J=\gamma^0\gamma^1\gamma^4,\qquad K_1=\gamma^2\gamma^3,\qquad K_2=\gamma^5\gamma^6 .
\]

By the first rule, $J$ ($k=3$) commutes with its own factors $\gamma^0,\gamma^1,\gamma^4$ and anticommutes with the other five gammas; $K_1$ ($k=2$) anticommutes with $\gamma^2,\gamma^3$ and commutes with the other six; $K_2$ anticommutes with $\gamma^5,\gamma^6$ and commutes with the other six. Therefore all three commute with $A_0$, $A_1$, $A_4$ (products of $\gamma^0,\gamma^1,\gamma^4$), with $C=\gamma^0\gamma^1\gamma^2\gamma^3$ (for $J$ the two anticommuting factors $\gamma^2,\gamma^3$ give two signs, which cancel) and with $B=-iC\gamma^4$, and they commute with each other (\texttt{KS\_reduc\allowbreak{}tion\_JK1\allowbreak{}K2commut\allowbreak{}e}). By the second rule, $J^2=(-1)^3\cdot(1\cdot1\cdot(-1))=+1$, $K_1^2=(-1)\cdot1\cdot1=-1$ and $K_2^2=(-1)\cdot(-1)(-1)=-1$. So $J$ has the eigenvalues $j=\pm1$, and $K_1$, $K_2$ have the eigenvalues $i s_2$, $i s_3$ with $s_2,s_3=\pm1$.

\textbf{The projectors have rank 2.} The joint eigenspace with labels $(j,s_2,s_3)$ is the range of

\[
P(j,s_2,s_3)=\frac{1+jJ}2\cdot\frac{1-is_2K_1}2\cdot\frac{1-is_3K_2}2 .
\]

(On a vector with $K_1v=is_2v$ the middle factor gives $(1-is_2\cdot is_2)/2=1$, on one with $K_1v=-is_2v$ it gives 0; likewise for the other two factors.) Multiplying out, $P$ is $\tfrac18$ times the sum of $1$ and seven products of distinct gammas, $J$, $K_1$, $K_2$, $JK_1$, $JK_2$, $K_1K_2$, $JK_1K_2$ (with coefficients). Each of these seven anticommutes with some $\gamma^c$ by the first rule: $J$, $K_1$, $JK_2$ and $K_1K_2$ with $\gamma^2$; $K_2$ and $JK_1$ with $\gamma^5$; and $JK_1K_2$, which contains every gamma except $\gamma^7$, with $\gamma^7$. A matrix $X$ that anticommutes with an invertible $\gamma^c$ is traceless: $\mathrm{tr}\,X=\mathrm{tr}\bigl(\gamma^cX(\gamma^c)^{-1}\bigr)=-\mathrm{tr}\,X$. Hence $\mathrm{tr}\,P=\tfrac18\,\mathrm{tr}\,1=16/8=2$. To turn this trace into a dimension, note that each of the three factors of $P$ is a projector of the kind treated in Section 2.2, $\tfrac12(1+X)$ with $X=jJ$, $X=-is_2K_1$ or $X=-is_3K_2$, and each of these $X$ squares to $1$ (for example $(-is_2K_1)^2=-s_2^2K_1^2=+1$). The three factors commute, so $P^2=P$; then $X=2P-1$ squares to $4P^2-4P+1=1$, and Section 2.2 gives $\dim(\text{range of }P)=\mathrm{tr}\,P$. So each of the eight label combinations has a 2-dimensional eigenspace (\texttt{KS\_reduc\allowbreak{}tion\_pro\allowbreak{}jectorsR\allowbreak{}ank2}), and the eight spaces together fill $\mathbb C^{16}$.

\textbf{The basis in each block.} Inside one block the matrices $A_0,A_1,A_4$ still satisfy the algebra above. $A_1$ anticommutes with $A_0$ and is invertible, so it maps the $A_0=+1$ eigenvector to an $A_0=-1$ eigenvector; each eigenvalue of $A_0$ therefore occurs once in the block. Choose $v_+$ with $A_0v_+=v_+$ and set $v_-=A_1v_+$. Then $A_0v_-=A_0A_1v_+=-A_1A_0v_+=-v_-$, $A_1v_+=v_-$ and $A_1v_-=A_1^2v_+=-v_+$. In the basis $(v_+,v_-)$, with the Pauli matrices $\sigma_x,\sigma_y,\sigma_z$ of Chapter 2,

\[
A_0\to\sigma_z,\qquad A_1\to\begin{pmatrix}0&-1\\1&0\end{pmatrix}=-i\sigma_y .
\]

$A_4$ anticommutes with $\sigma_z$, so it is off-diagonal, and it is fixed by one product: $A_0A_1A_4=\gamma^0(\gamma^0\gamma^1)(\gamma^0\gamma^4)=\gamma^1\gamma^0\gamma^4=-J$, which is the number $-j$ in the block. Since $\sigma_z(-i\sigma_y)=-i\sigma_z\sigma_y=-i(-i\sigma_x)=-\sigma_x$, the condition $-\sigma_xA_4=-j$ gives $A_4\to j\sigma_x$ (\texttt{KS\_reduc\allowbreak{}tion\_blo\allowbreak{}cksA0A1A\allowbreak{}4}). The same way we get the matrices of the expectation-value rule. $C=\gamma^0\gamma^1\gamma^2\gamma^3=A_1K_1\to(-i\sigma_y)(is_2)=s_2\sigma_y$. Moving $\gamma^4$ to the right past $\gamma^2\gamma^3$ (two swaps) shows $JK_1=\gamma^0\gamma^1\gamma^2\gamma^3\gamma^4=C\gamma^4$, so $B=-iC\gamma^4=-iJK_1\to-i\,j\,(is_2)=js_2$, a number. Then

\[
C\to s_2\sigma_y,\qquad B\to js_2,\qquad BC\to j\sigma_y,\qquad \gamma^4\gamma^1\to-j\sigma_z
\]

(\texttt{KS\_reduc\allowbreak{}tion\_blo\allowbreak{}cksBC}; for the last one $\gamma^4=A_0A_4\to ij\sigma_y$ and $\gamma^1=A_0A_1\to-\sigma_x$).

\textbf{A concrete block.} The exact theory builds the basis as $v_+=8P(j,s_2,s_3)\tfrac12(1+\gamma^0)e_c$ with the first standard unit vector $e_c$ that gives a nonzero result, and $v_-=\gamma^0\gamma^1v_+$; every entry is $0$, $\pm1$ or $\pm i$, and $|v_\pm|^2=8$. For the block $(j,s_2,s_3)=(1,1,1)$ ($c=4$, counting components from 0):

\[
\begin{aligned}
v_+&=(0,0,0,0,\,1,i,-1,i,\,0,0,0,0,\,1,i,-1,i)^T,\\
v_-&=(i,-1,-i,-1,\,0,0,0,0,\,-i,1,i,1,\,0,0,0,0)^T .
\end{aligned}
\]

With the gammas of the fixture \texttt{artifact\allowbreak{}s/dirac1\allowbreak{}6complex\allowbreak{}/arbitra\allowbreak{}ry-field\allowbreak{}/algebra\allowbreak{}-fixture\allowbreak{}.json} one checks $\gamma^0v_+=v_+$, $Jv_+=v_+$, $K_1v_+=iv_+$, $K_2v_+=iv_+$ and $Bv_+=v_+$ (the Krein sign $js_2=1$ of this block), and that the $2\times2$ matrices in the normalized basis $(v_+,v_-)/\sqrt8$ are exactly the ones above with $j=s_2=1$. All eight blocks are listed in \texttt{kohn-sha\allowbreak{}m-theory\allowbreak{}.json} (keys \texttt{reductio\allowbreak{}n}, \texttt{blockDia\allowbreak{}gonalisa\allowbreak{}tion}); the $16\times16$ matrix with the columns $v_\pm/\sqrt8$ is unitary and makes $A_0,A_1,A_4,B,C$ block diagonal (\texttt{KS\_reduc\allowbreak{}tion\_bas\allowbreak{}isUnitar\allowbreak{}y}, \texttt{KS\_reduc\allowbreak{}tion\_fiv\allowbreak{}eMatrice\allowbreak{}sBlockDi\allowbreak{}agonal}, \texttt{KS\_reduc\allowbreak{}tion\_rec\allowbreak{}onstruct\allowbreak{}FromBloc\allowbreak{}ks}).

\begin{figure}[htbp]
\centering
\includegraphics[width=0.92\linewidth]{artifacts/dirac16complex/kohn-sham/figures/block_structure.png}
\caption{Moduli of the matrix entries of the \texorpdfstring{$y$}{y}-current matrix \texorpdfstring{$\gamma^0\gamma^4$}{γ\textasciicircum{}0γ\textasciicircum{}4} in the original spinor basis (left) and in the block basis (middle), and of the operator of the reduced equation in the block basis at one sample point (right): eight \texorpdfstring{$2\times2$}{2×2} blocks along the diagonal (figure \texttt{block\_st\allowbreak{}ructure.\allowbreak{}png} of \texttt{artifact\allowbreak{}s/dirac1\allowbreak{}6complex\allowbreak{}/kohn-sh\allowbreak{}am/figur\allowbreak{}es}).}
\end{figure}

\textbf{Two types of block.} For the $y$ equation only $A_0,A_1,A_4$ matter, and their block forms depend on $j$ alone: there are \textbf{two inequivalent block types}, $j=+1$ (four blocks) and $j=-1$ (four blocks) (\texttt{KS\_reduc\allowbreak{}tion\_blo\allowbreak{}ckTypes}, erratum E4.4). Including $B$ and $C$ there are four types $(j,s_2)$ of two blocks each. The chirality $\gamma^8$ anticommutes with $J$ (a product of three gammas) and commutes with $K_1,K_2$, so it maps the block $(j,s_2,s_3)$ onto $(-j,s_2,s_3)$ (\texttt{KS\_reduc\allowbreak{}tion\_gam\allowbreak{}ma8Swaps\allowbreak{}J}); this is the seed of the pairing of Section 13.14 and Chapter 15.

\textbf{Convention of the Rust code.} The Rust module \texttt{studies/\allowbreak{}dirac16c\allowbreak{}omplex\_k\allowbreak{}ohn\_sham\allowbreak{}/src/blo\allowbreak{}cks.rs} labels its blocks by the eigenvalue $s$ of $A_0A_1A_4$ (the header of \texttt{blocks.r\allowbreak{}s} calls this product itself $J$), which is $-J$ in the notation of this chapter; hence $s=-j$ (Exercise 13.2). Its block forms, printed in the header of \texttt{blocks.r\allowbreak{}s}, read $A_4\to-s\sigma_x$ and $BC\to-s\sigma_y$, which are the forms $j\sigma_x$ and $j\sigma_y$ above with $s=-j$; and the detail of the check \texttt{theory\_z\allowbreak{}ero\_mode\allowbreak{}\_splitti\allowbreak{}ng} in \texttt{rust/spe\allowbreak{}ctrum/th\allowbreak{}eory-agr\allowbreak{}eement.j\allowbreak{}son} identifies the theory's block type $j=+1$ with the Rust label $s=-1$. The two descriptions are identical; this book uses $j$ throughout and converts the Rust labels.

\subsection{The block equation, its Hamiltonian, the current and the boundary conditions}

\textbf{The block equation.} In a block of type $j$, with $\chi=(\chi_1,\chi_2)^T$ the two components in the basis $(v_+,v_-)$ and with the exchange potential $v(y)$ of Section 13.8 included ($\varepsilon\to\varepsilon-v$),

\[
\chi'=N\chi,\qquad N=M_{\text{eff}}(y)\,\sigma_z-\xi(y)\,k\,\sigma_y+ij\bigl(\varepsilon-v(y)\bigr)\sigma_x
\]

(\texttt{KS\_reduc\allowbreak{}tion\_blo\allowbreak{}ckODEMat\allowbreak{}rix}), since $-i\xi k(-i\sigma_y)=-\xi k\sigma_y$. In components,

\[
\chi_1'=M_{\text{eff}}\chi_1+i\bigl(\xi k+jE\bigr)\chi_2,\qquad \chi_2'=-M_{\text{eff}}\chi_2+i\bigl(jE-\xi k\bigr)\chi_1,\qquad E=\varepsilon-v .
\]

Writing $\chi=(a,\,ib)$ with real $a,b$ makes the system real:

\[
a'=M_{\text{eff}}\,a-(jE+\xi k)\,b,\qquad b'=(jE-\xi k)\,a-M_{\text{eff}}\,b .
\]

(With the Rust label $s=-j$ this is the system $a'=Ma+(sE-\xi k)b$, $b'=-(sE+\xi k)a-Mb$ printed in \texttt{blocks.r\allowbreak{}s}.)

\textbf{The Hamiltonian form.} The same equation is an eigenvalue problem $h_j\chi=\varepsilon\chi$ with

\[
h_j=j\Bigl[-i\sigma_x\frac d{dy}+M_{\text{eff}}(y)\,\sigma_y+\xi(y)\,k\,\sigma_z\Bigr]+v(y).
\]

To see it, write $h_j\chi=\varepsilon\chi$ as $-i\sigma_x\chi'=jE\chi-M_{\text{eff}}\sigma_y\chi-\xi k\sigma_z\chi$ (using $j^2=1$) and multiply by $i\sigma_x$: with $\sigma_x\sigma_y=i\sigma_z$ and $\sigma_x\sigma_z=-i\sigma_y$ one gets $\chi'=ijE\sigma_x\chi+M_{\text{eff}}\sigma_z\chi-\xi k\sigma_y\chi$, the block equation (\texttt{KS\_reduc\allowbreak{}tion\_blo\allowbreak{}ckHamilt\allowbreak{}onianEqu\allowbreak{}ivalentT\allowbreak{}oODE}).

\textbf{Symmetries of the spectrum.} Write $h_j-v=jD_k$ with $D_k=-i\sigma_x\,d/dy+M_{\text{eff}}\sigma_y+\xi k\sigma_z$, which does not depend on $j$. Two facts follow. (i) $h_{-1}-v=-(h_{+1}-v)$: the two block types carry opposite spectra of $h-v$ (\texttt{KS\_reduc\allowbreak{}tion\_hMi\allowbreak{}nusEqual\allowbreak{}sMinusHP\allowbreak{}lus}); when $v=0$ the spectrum of $j=-1$ is the negative of that of $j=+1$. (ii) Conjugation with $\sigma_z$ gives $\sigma_zD_k\sigma_z=+i\sigma_x\,d/dy-M_{\text{eff}}\sigma_y+\xi k\sigma_z=-D_{-k}$, and the boundary conditions below are unchanged by $\chi\to\sigma_z\chi$. Hence the orbital $\sigma_z\chi$ of $(-k,-j)$ has the same energy and the same $|\chi|^2$ as the orbital $\chi$ of $(k,j)$ (\texttt{KS\_reduc\allowbreak{}tion\_sig\allowbreak{}ma3Conju\allowbreak{}gationFl\allowbreak{}ipsK}). So at a fixed momentum $\mathbf k=(k,0,0)$ every level is 4-fold (the four blocks of one type), and over a closed shell of momenta $\{\mathbf k,-\mathbf k,\dots\}$ and at $k=0$ the levels are 8-fold (erratum E4.4).

\textbf{The $y$-current.} The current of the U(1) charge along the hidden direction is $J^y=-i\bar\Psi\gamma^y\Psi=-i\Psi^\dagger C\gamma^0\Psi$. By the expectation-value rule of Section 13.7 its matrix is $B(-iC\gamma^0)=-iBC\gamma^0=-i(-i\gamma^4)\gamma^0=-\gamma^4\gamma^0=\gamma^0\gamma^4=A_4$, which is $j\sigma_x$ in the block (erratum E4.5; \texttt{KS\_bound\allowbreak{}ary\_curr\allowbreak{}entMatri\allowbreak{}x}, \texttt{KS\_bound\allowbreak{}ary\_curr\allowbreak{}entBlock\allowbreak{}Form}). The current density of an orbital is proportional to $\chi^\dagger\sigma_x\chi=2\,\mathrm{Re}(\chi_1^*\chi_2)$. It is \textbf{constant in $y$} for every solution: $\frac d{dy}(\chi^\dagger\sigma_x\chi)=\chi^\dagger(N^\dagger\sigma_x+\sigma_xN)\chi$, and with $N^\dagger=M_{\text{eff}}\sigma_z-\xi k\sigma_y-ijE\sigma_x$ (real $M_{\text{eff}}$, $\xi$, $k$, $E$) the anticommutation of the Pauli matrices gives $N^\dagger\sigma_x+\sigma_xN=M_{\text{eff}}\{\sigma_z,\sigma_x\}-\xi k\{\sigma_y,\sigma_x\}+ijE(\sigma_x\sigma_x-\sigma_x\sigma_x)=0$ (\texttt{KS\_bound\allowbreak{}ary\_curr\allowbreak{}entConse\allowbreak{}rvedAlon\allowbreak{}gY}). The plain density $\chi^\dagger\chi$, in contrast, is not constant in $y$ (\texttt{KS\_bound\allowbreak{}ary\_hilb\allowbreak{}ertNormN\allowbreak{}otConser\allowbreak{}vedAlong\allowbreak{}Y}).

\textbf{Boundary conditions at the brane.} On the Z2 geometry one asks $\Psi(-y)=\pm\gamma^0\Psi(y)$, which at $y=0$ means $(1\mp\gamma^0)\chi(0)=0$; since $\gamma^0=A_0=\sigma_z$ in every block,

\[
\text{parity }+:\ \chi_2(0)=0\ \ (b(0)=0),\qquad \text{parity }-:\ \chi_1(0)=0\ \ (a(0)=0)
\]

(\texttt{KS\_bound\allowbreak{}ary\_pari\allowbreak{}tyProjec\allowbreak{}torsBloc\allowbreak{}kForm}). Either condition kills the current at the brane, because $\chi^\dagger\sigma_x\chi=2\mathrm{Re}(\chi_1^*\chi_2)$ vanishes when one component vanishes (\texttt{KS\_bound\allowbreak{}ary\_pari\allowbreak{}tyKillsC\allowbreak{}urrent}). A caution from the exact theory (erratum E4.6): the reflection $\chi(y)\to\gamma^0\chi(-y)$ maps solutions with the mass function $M(y)$ onto solutions with $-M(-y)$ (use $\sigma_z\sigma_{x,y}\sigma_z=-\sigma_{x,y}$ and the evenness of $\xi$), so it is a \textbf{symmetry} only when the mass function is odd, which is the case of a $\pm M$ mirror pair (\texttt{KS\_bound\allowbreak{}ary\_pari\allowbreak{}tyA\_symm\allowbreak{}etryIffM\allowbreak{}assOdd}). For an even mass function (an identical mirror copy) the symmetry is instead $\chi(y)\to i\gamma^0\gamma^8\chi(-y)$, which couples the blocks $j$ and $-j$ (\texttt{KS\_bound\allowbreak{}ary\_pari\allowbreak{}tyB\_symm\allowbreak{}etryForE\allowbreak{}venMass}, \texttt{KS\_bound\allowbreak{}ary\_pari\allowbreak{}tyB\_coup\allowbreak{}lesJBloc\allowbreak{}ks}). The Kohn--Sham problem therefore uses the two parities as \textbf{boundary conditions} at $y=0$ and computes both; it does not assume a reflection symmetry.

\textbf{Boundary condition at the tip.} At the cutoff $y=-L$ the solver imposes the \textbf{bag condition} $\chi_2(-L)=0$ ($b(-L)=0$, equivalently $\gamma^0\chi(-L)=+\chi(-L)$). It is the member $\theta=0$ of the family $(1-Q(\theta))\chi(-L)=0$ with $Q(\theta)=\cos\theta\,\sigma_z+\sin\theta\,\sigma_y$ (\texttt{KS\_bound\allowbreak{}ary\_bagF\allowbreak{}amily}, \texttt{KS\_bound\allowbreak{}ary\_bagT\allowbreak{}hetaZero\allowbreak{}IsEvenPa\allowbreak{}rity}). Every member kills the current: $Q$ is Hermitian, $Q^2=1$ and $Q$ anticommutes with $\sigma_x$, so for $Q\chi=\chi$ we get $\chi^\dagger\sigma_x\chi=\chi^\dagger\sigma_xQ\chi=-\chi^\dagger Q\sigma_x\chi=-(Q\chi)^\dagger\sigma_x\chi=-\chi^\dagger\sigma_x\chi$, hence 0. The choice $\theta=0$ admits the brane zero mode of Section 13.6 for every $L$ and, for $M_{\text{eff}}>0$, no artificial mode localized at the cutoff (the header of \texttt{src/shoo\allowbreak{}ting.rs}). The tip is where the transverse volume $W^6$ goes to zero, at infinite proper distance (the curvature stays constant there, Sections 9.15 and 9.16); the cutoff $L$ truncates it, and the results are reported for $L=2,3,4$ (Section 13.15).

\textbf{Self-adjointness.} For two functions $\varphi,\chi$ on $[-L,0]$, integrating by parts,

\[
\langle\varphi|h_j\chi\rangle-\langle h_j\varphi|\chi\rangle=-ij\int_{-L}^0\bigl(\varphi^\dagger\sigma_x\chi'+\varphi'^\dagger\sigma_x\chi\bigr)dy=-ij\bigl[\varphi^\dagger\sigma_x\chi\bigr]_{-L}^{0}
\]

(the terms without derivatives cancel because $\sigma_y,\sigma_z$ and the real functions are Hermitian). If both functions satisfy the same condition at each end (one of the two components zero), then $\varphi^\dagger\sigma_x\chi=\varphi_1^*\chi_2+\varphi_2^*\chi_1=0$ there, and $h_j$ is self-adjoint (\texttt{KS\_bound\allowbreak{}ary\_self\allowbreak{}AdjointB\allowbreak{}oundaryT\allowbreak{}erm}). Consequently the levels $\varepsilon$ are real and the orbitals of different levels are orthogonal (Section 12.2).

\subsection{Exact solutions at zero momentum and the brane band}

Before any interaction, take $k=0$, a constant mass $M_{\text{eff}}=M>0$ and $v=0$. The block equation can then be solved by hand, and the results are the backbone of every numerical run.

\textbf{Eliminating one component.} From $\chi_2'=-M\chi_2+ij\varepsilon\chi_1$ we get $\chi_1=(\chi_2'+M\chi_2)/(ij\varepsilon)$ for $\varepsilon\ne0$. Inserting this into $\chi_1'=M\chi_1+ij\varepsilon\chi_2$ and multiplying by $ij\varepsilon$ gives $\chi_2''+M\chi_2'=M\chi_2'+M^2\chi_2+(ij)^2\varepsilon^2\chi_2$, that is

\[
\chi_2''=\bigl(M^2-\varepsilon^2\bigr)\chi_2 .
\]

\textbf{Parity $+$: the zero mode and the massive levels.} The conditions are $\chi_2(0)=0$ and $\chi_2(-L)=0$. (i) If $\varepsilon=0$, the equations decouple: $\chi_1'=M\chi_1$ and $\chi_2'=-M\chi_2$, and $\chi_2\equiv0$ satisfies both conditions. This is the \textbf{zero mode}

\[
\varepsilon=0,\qquad \chi=\bigl(e^{My},\,0\bigr)^T ,
\]

which grows toward $y=0$: it is \textbf{localized at the brane} (\texttt{KS\_reduc\allowbreak{}tion\_k0Z\allowbreak{}eroMode}). (ii) If $\varepsilon^2>M^2$, write $p^2=\varepsilon^2-M^2$; then $\chi_2''=-p^2\chi_2$ with $\chi_2(0)=\chi_2(-L)=0$ gives $\chi_2=\sin(py)$ with $pL=n\pi$, so

\[
\varepsilon=\pm\sqrt{M^2+(n\pi/L)^2},\qquad \chi_2=\sin\frac{n\pi y}L,\qquad \chi_1=\frac{M\sin(k_ny)+k_n\cos(k_ny)}{ij\varepsilon},\quad k_n=\frac{n\pi}L,
\]

$n=1,2,\dots$ (\texttt{KS\_reduc\allowbreak{}tion\_k0M\allowbreak{}assiveLe\allowbreak{}vels}). (iii) If $0<\varepsilon^2\le M^2$, the solutions of $\chi_2''=(M^2-\varepsilon^2)\chi_2$ that vanish at both ends are zero, and then $\chi_2\equiv0$ forces $\varepsilon\chi_1=0$: no level. So at $k=0$ the spectrum of parity $+$ is exactly $\{0\}\cup\{\pm\sqrt{M^2+(n\pi/L)^2}\}$, each level 4-fold per block type, 8-fold in total.

\textbf{Worked example.} For $M=m=1$, $L=3$: $\sqrt{1+(\pi/3)^2}=1.447972$ and $\sqrt{1+(2\pi/3)^2}=2.320881$. The Rust solver lists the first of these as $1.447972$ (the lowest positive massive $k=0$ level of parity $+$ in \texttt{rust/scf\allowbreak{}/m1\_L3\_N\allowbreak{}1016\_lam\allowbreak{}0\_T0/lev\allowbreak{}els.csv}).

\textbf{Parity $-$.} Now $\chi_1(0)=0$ and $\chi_2(-L)=0$. The second condition gives $\chi_2=\sin\bigl(p(y+L)\bigr)$ with $p^2=\varepsilon^2-M^2$, and the first, via $\chi_1=(\chi_2'+M\chi_2)/(ij\varepsilon)$, becomes $\chi_2'(0)+M\chi_2(0)=0$:

\[
p\cos(pL)+M\sin(pL)=0\quad\Longleftrightarrow\quad \tan(pL)=-\frac pM .
\]

(For $\varepsilon^2<M^2$ the analogous condition with $\sinh$ and $\cosh$, $q\cosh(qL)+M\sinh(qL)=0$, has no solution with $q>0$, and $\varepsilon=0$ is excluded too: parity $-$ has no level with $|\varepsilon|<M$.) For $M=1$, $L=3$ the smallest root lies between $\pi/6$ and $\pi/3$; bisection gives $p=0.818548$ and $\varepsilon=\sqrt{1+p^2}=1.292293$. The Rust value of this level (8 states, occupied in the free $N=1016$ state, Section 13.11) is $1.292293$; the two agree to about $2\times10^{-12}$.

\textbf{The brane band: the zero mode at small momentum.} For $k\ne0$ the zero mode moves. Its energy to first order in $k$ follows from the \textbf{Hellmann--Feynman} rule: if $h(k)\chi=\varepsilon(k)\chi$ with $\langle\chi|\chi\rangle=1$ and boundary conditions that do not depend on $k$, then $\varepsilon=\langle\chi|h\chi\rangle$ and $d\varepsilon/dk=\langle\chi|(\partial_kh)\chi\rangle$, because the terms with $d\chi/dk$ add up to $\varepsilon\,d\langle\chi|\chi\rangle/dk=0$ (self-adjointness, Section 13.5). Here $\partial_kh_j=j\xi(y)\sigma_z$ and $\chi^\dagger\sigma_z\chi=e^{2My}$ for the zero mode, so

\[
\frac{d\varepsilon}{dk}\Big|_{k=0}=j\,c,\qquad c=\frac{\int_{-L}^0e^{-Hy-a_{4,0}}e^{2My}\,dy}{\int_{-L}^0e^{2My}\,dy}
=e^{-a_{4,0}}\,\frac{2M}{2M-H}\cdot\frac{1-e^{-(2M-H)L}}{1-e^{-2ML}}
\]

(\texttt{KS\_reduc\allowbreak{}tion\_zer\allowbreak{}oModeSpl\allowbreak{}itting}). For $H=M=1$ and $a_{4,0}=0$: $c=2(1-e^{-L})/(1-e^{-2L})=2/(1+e^{-L})$, which is $1.761594$ for $L=2$, $1.905148$ for $L=3$ and $1.964028$ for $L=4$. The Rust solver measures $1.9051482525$ for $L=3$ against the exact $1.9051482536$ (\texttt{rust/spe\allowbreak{}ctrum/th\allowbreak{}eory-agr\allowbreak{}eement.j\allowbreak{}son}). So the eight zero modes split into a \textbf{brane band} $\varepsilon=+ck$ (four states, $j=+1$) and $\varepsilon=-ck$ (four states, $j=-1$) at each momentum. The band is odd in $k$ (complex conjugation maps the block equation at $(k,\varepsilon)$ to the one at $(-k,-\varepsilon)$ when $v=0$), so the next correction is of order $k^3$. At the smallest torus momentum $k=\Delta k=0.25\,m$ (Section 13.7) the first-order value is $ck=0.476287\,m$, and the exact level found by the solver is $0.430734\,m$ (\texttt{rust/scf\allowbreak{}/m1\_L3\_N\allowbreak{}8\_lam0\_T\allowbreak{}0/run.js\allowbreak{}on}, \texttt{\detokenize{epsLumo}}). The brane band behaves like a massless fermion that lives on the brane; its states are the lowest positive levels at every $k\ne0$, and together with the zero modes at $k=0$ they hold almost all the particles of the ground states of this chapter (Section 13.11 lists the few bulk levels that are occupied as well).

\begin{figure}[htbp]
\centering
\includegraphics[width=0.92\linewidth]{artifacts/dirac16complex/kohn-sham/figures/ks_spectrum.png}
\caption{Kohn--Sham levels \texorpdfstring{$\varepsilon_n(k)$}{ε\_n(k)} against the 3-momentum \texorpdfstring{$k$}{k} for parity \texorpdfstring{$+$}{+} (top) and parity \texorpdfstring{$-$}{-} (bottom), free (\texorpdfstring{$m=1$}{m=1}, \texorpdfstring{$L=3$}{L=3}, left) and self-consistent (\texorpdfstring{$N=112$}{N=112} with \texorpdfstring{$+\hat\lambda_2$}{+λ\_2}, right), for the two block types (circles and triangles). The dashed lines at the origin are the first-order brane band \texorpdfstring{$\pm ck$}{±ck}; the dotted line is the chemical potential (file ks\_spectrum.png in artifacts/dirac16complex/kohn-sham/figures).}
\end{figure}

\subsection{Densities, the torus, and particles above the sea}

\textbf{The expectation-value rule.} In the positive-norm quantization of Chapter 8, a one-particle state built on a normalized mode $u$ above the Dirac sea has $\langle\Psi^\dagger X\Psi\rangle=u^\dagger BXu$ for every matrix $X$. The number (charge) density operator is $\Psi^\dagger B\Psi$, so its expectation is $u^\dagger B^2u=u^\dagger u\ge0$; the scalar density $\bar\Psi\Psi=\Psi^\dagger C\Psi$ gives $u^\dagger BCu$; the $y$-current gives $u^\dagger A_4u$ (Section 13.5). For a whole state with occupation weights $w_n$ the density matrix is $\rho=\sum_nw_nu_nu_n^\dagger$ and $\langle\Psi^\dagger X\Psi\rangle=\mathrm{Tr}(BX\rho)$.

\textbf{Densities of one orbital.} For a block orbital $\chi$ of type $j$ with $\int_{-L}^0\chi^\dagger\chi\,dy=1$, spread over a coordinate 3-volume $\ell^3$, the four \textbf{proper densities} (per unit proper volume) are

\[
n=\frac{e^{-6Hy}}{\ell^3}\chi^\dagger\chi,\qquad s=\frac{e^{-6Hy}}{\ell^3}\,j\chi^\dagger\sigma_y\chi,\qquad t=\frac{e^{-6Hy}}{\ell^3}\,j\chi^\dagger\sigma_z\chi,\qquad c=\frac{e^{-6Hy}}{\ell^3}\,j\chi^\dagger\sigma_x\chi :
\]

the number density, the scalar density (from $BC\to j\sigma_y$), the density of the 3-momentum current (from $-\gamma^4\gamma^1\to j\sigma_z$) and the $y$-current, which vanishes for every eigenstate (it is constant and zero at the ends). The factor $e^{-6Hy}=1/\sqrt{|g|}$ converts a density per unit $y$ into a density per unit proper volume. In the real form $\chi=(a,ib)$: $\chi^\dagger\chi=a^2+b^2$, $\chi^\dagger\sigma_y\chi=2ab$, $\chi^\dagger\sigma_z\chi=a^2-b^2$.

The \textbf{coordinate densities} $n_c=e^{6Hy}n_p$ and $S_c=e^{6Hy}S_p$ count particles per unit $y$ and unit coordinate 3-volume; the total particle number is $N=\ell^3\int_{-L}^0n_c\,dy$. Because the proper volume shrinks toward the tip, a state that is concentrated near the brane in $n_c$ can still have its largest \textbf{proper} densities at the tip; the interaction, which depends on the proper densities, is strongest there (Section 13.11).

\textbf{The torus and the shells.} The 3-space directions are a coordinate torus of size $\ell$, so $\mathbf k=\Delta k\,(n_1,n_2,n_3)$ with integers $n_l$ and $\Delta k=2\pi/\ell$ (Section 12.12). The runs use $\Delta k=0.25\,m$, that is $\ell=2\pi/(0.25\,m)=25.1327/m$. Since the spectrum depends only on $|\mathbf k|$, the momenta fall into \textbf{shells} $|\mathbf k|^2=\Delta k^2\,\nu$ with the integer $\nu=n_1^2+n_2^2+n_3^2$ (called \texttt{\detokenize{n2}} in the Rust files), each with the number $g(\nu)$ of lattice vectors: $g(0)=1$, $g(1)=6$, $g(2)=12$, $g(3)=8$, $g(4)=6$, $g(5)=24$, $g(6)=24$, $g(8)=12$, $g(9)=30$, and so on. Each level of one block type at one lattice vector holds 4 states (the four blocks).

\textbf{Particles and the Dirac sea.} The quantization of Chapter 8 fills all negative-energy states (the Dirac sea) in the vacuum, and $N$ counts the quanta above it. The solver decides which levels belong to the sea by continuity from the free problem: a level is a \textbf{particle} level if its free partner (same shell, parity, block type and Prüfer index at $\lambda=0$, Section 13.10) has $\varepsilon_{\text{free}}\ge0$, and a \textbf{sea} level otherwise. The exact zero modes ($\varepsilon_{\text{free}}=0$) are counted as particle levels in both block types; this is a \textbf{convention} of the solver (the header of \texttt{src/scf.\allowbreak{}rs}). A particle level carries the weight $w=f(\varepsilon)$, the Fermi--Dirac occupation of Section 12.14, and a sea level the weight $w=-(1-f(\varepsilon))$: a hole in the sea, that is a thermal antiparticle, counts with a minus sign in the number of quanta (normal ordering). At $T=0$ the sea is full and contributes nothing.

\textbf{Closed shells.} At $T=0$ and $\lambda=0$ the $N$ lowest particle states are filled. A value of $N$ that fills complete levels is a \textbf{closed shell}; for it the Kohn--Sham gap is well defined. With the degeneracies above: $N=8$ fills the eight zero modes; the first band shell adds $4\cdot g(1)=24$ states ($N=32$), the next $4\cdot12=48$ ($N=80$), the next $4\cdot8=32$ ($N=112$), then $4\cdot6=24$ ($N=136$), and so on (the list \texttt{closedSh\allowbreak{}ells\_m1\_\allowbreak{}L3\_N\_eps\allowbreak{}Homo} in \texttt{rust/spe\allowbreak{}ctrum/su\allowbreak{}mmary.js\allowbreak{}on}, up to $N=1408$). The study uses $N=8$ (the smallest closed shell), $N=112$ (the closed shell nearest to 100, \texttt{\detokenize{nMid}} in the same file) and $N=1016$ (the one nearest to 1000, \texttt{\detokenize{nLarge}}).

\subsection{The interaction: Hartree energy and the exact exchange of the contact term}

\textbf{The interaction and its size.} The Lagrangian contains $U(S)=\tfrac\lambda2S^2$ with $S=\bar\Psi\Psi$, normal ordered with respect to the free Dirac sea. To measure its size we count \textbf{mass dimensions}. With $\hbar=c=1$ (Chapter 0) every unit is a power of a mass unit: an energy or a momentum is a mass, and a length or a time is an inverse mass (the length $\hbar/(mc)$ belongs to the mass $m$). The power is the mass dimension of the quantity: a mass has mass dimension $1$, a length $-1$, a derivative $\partial_\mu$ has $+1$ and the volume element $d^8x$ has $-8$. The action $\int\mathcal L\,d^8x$ must be a pure number, so every term of the Lagrangian density has mass dimension 8. The kinetic term $\bar\Psi\gamma^\mu\partial_\mu\Psi$ has twice the dimension of $\Psi$ plus 1, so $\Psi$ has mass dimension $7/2$ ($2\cdot\tfrac72+1=8$); then $S=\bar\Psi\Psi$ has dimension 7 (and the mass term $mS$ has $1+7=8$, as it must), $S^2$ has dimension 14, and $\lambda S^2$ has dimension 8 only if $\lambda$ has dimension $8-14=-6$. The dimensionless strength is $\hat\lambda=\lambda m^6$ (\texttt{KS\_excha\allowbreak{}nge\_coup\allowbreak{}lingDime\allowbreak{}nsion}).

\textbf{Hartree--Fock energy of the contact term.} For a quasi-free state (a Slater determinant or a non-interacting ensemble) with the density matrix $\rho$ of Section 13.7, the expectation-value rule makes the two-point function $G_{ab}=\langle\Psi_a^\dagger\Psi_b\rangle=(\rho B)_{ba}$. Now $S^2=\Psi_a^\dagger C_{ab}\Psi_b\,\Psi_c^\dagger C_{cd}\Psi_d$ (sums over repeated spinor indices), and we move $\Psi_b$ to the right of $\Psi_c^\dagger$. An ordinary reordering would give $\Psi_b\Psi_c^\dagger=-\Psi_c^\dagger\Psi_b+\{\Psi_b,\Psi_c^\dagger\}$, where the anticommutator $\{\Psi_b,\Psi_c^\dagger\}$ is not an operator but a number times the unit operator (the equal-time anticommutator of Section 8.6); this number is called the \textbf{contraction} of the pair. Inside a normal-ordered product the operators are reordered as if they anticommuted exactly, that is, without the contraction, which gives $-\Psi_a^\dagger\Psi_c^\dagger\Psi_b\Psi_d$. Wick's theorem (Section 12.5) evaluates the four-operator expectation as $\langle\Psi_a^\dagger\Psi_c^\dagger\Psi_b\Psi_d\rangle=G_{ad}G_{cb}-G_{ab}G_{cd}$. Hence

\[
\langle{:}S^2{:}\rangle=C_{ab}C_{cd}\bigl(G_{ab}G_{cd}-G_{ad}G_{cb}\bigr)=\bigl[\mathrm{Tr}(BC\rho)\bigr]^2-\mathrm{Tr}(BC\rho\,BC\rho),
\]

using $\sum_{ab}C_{ab}(\rho B)_{ba}=\mathrm{Tr}(C\rho B)=\mathrm{Tr}(BC\rho)$ and the cyclic property of the trace for the second term. The \textbf{Hartree--Fock energy density} is therefore

\[
e_{HF}=\frac\lambda2\Bigl(S^2-\mathrm{Tr}(BC\rho\,BC\rho)\Bigr),\qquad S=\mathrm{Tr}(BC\rho),
\]

the Hartree term $e_H=\tfrac\lambda2S^2$ and the exchange (Fock) term $e_x=-\tfrac\lambda2\mathrm{Tr}(BC\rho BC\rho)$ (\texttt{KS\_excha\allowbreak{}nge\_wick\allowbreak{}TheoremH\allowbreak{}F}, verified exactly on a Fock space of four modes for one-, two- and three-particle determinants). The Hartree term shifts the mass: $\partial e_H/\partial S=\lambda S$, so $M_{\text{eff}}=m+\lambda S$ as in Section 13.3.

\textbf{A filled shell: the ratio $1/8$.} Fill all eight positive-energy states at rest ($\mathbf p=0$). The rest Hamiltonian is $h_0=-im\gamma^4=m\,BC$ (Chapter 8), and $BC=-i\gamma^4$ is Hermitian, traceless and squares to 1 (since $(\gamma^4)^2=-1$). The projector onto its $+1$ eigenspace is $P_+=\tfrac12(1+BC)$, of rank 8, and $BC\,P_+=P_+$. With $\rho=P_+$:

\[
S=\mathrm{Tr}(BCP_+)=\tfrac12\mathrm{Tr}(BC)+\tfrac12\mathrm{Tr}(1)=8,\qquad \mathrm{Tr}(BCP_+BCP_+)=\mathrm{Tr}(P_+)=8=\frac{S^2}8 .
\]

So $e_x=-e_H/8$ (\texttt{KS\_excha\allowbreak{}nge\_fill\allowbreak{}edShellO\allowbreak{}neEighth}), the result $E_x=-E_H/g$ of Section 12.6 with $g=8$ internal states. (At a momentum $\mathbf p\ne0$ the same holds with $S=8m/E_p$, $E_p=\sqrt{m^2+p^2}$, \texttt{KS\_excha\allowbreak{}nge\_fill\allowbreak{}edShellS\allowbreak{}calarDen\allowbreak{}sity}.)

\textbf{The uniform gas.} Now take the uniform gas of Section 12.12 for dirac16complex: all momenta $\mathbf p$ (in $d$ space dimensions), particle occupations $f_+(E_p)=1/(e^{(E_p-\mu)/T}+1)$ and antiparticle occupations $f_-(E_p)=1/(e^{(E_p+\mu)/T}+1)$ (the holes of the sea). With the projectors $P_\pm(\mathbf p)=\tfrac12(1\pm h_{\mathbf p}/E_p)$ onto the positive and negative energies of $h_{\mathbf p}=-im\gamma^4-\gamma^4\gamma^lp_l$, the normal-ordered density matrix is $\rho=\int\frac{d^dp}{(2\pi)^d}\bigl[f_+P_+(\mathbf p)-f_-P_-(\mathbf p)\bigr]$. The number and scalar densities follow from $\mathrm{Tr}\,P_\pm=8$ and $\mathrm{Tr}(BC\,P_\pm)=\pm8m/E_p$:

\[
n=\mathrm{Tr}\,\rho=8\int\frac{d^dp}{(2\pi)^d}\bigl(f_+-f_-\bigr),\qquad
S=\mathrm{Tr}(BC\rho)=8\int\frac{d^dp}{(2\pi)^d}\,\frac m{E_p}\bigl(f_++f_-\bigr).
\]

Antiparticles lower $n$ and raise $S$. The exchange term needs the \textbf{kernel}

\[
K_{ab}(\mathbf p,\mathbf q)=\mathrm{Tr}\bigl(P_a(\mathbf p)\,BC\,P_b(\mathbf q)\,BC\bigr)=4\Bigl[1+ab\,\frac{m^2-\mathbf p\cdot\mathbf q}{E_pE_q}\Bigr],\qquad a,b=\pm1
\]

(\texttt{KS\_excha\allowbreak{}nge\_kern\allowbreak{}elPlusPl\allowbreak{}us}, \texttt{KS\_excha\allowbreak{}nge\_kern\allowbreak{}elPlusMi\allowbreak{}nus}; the Rust solver checks it with the full $16\times16$ projectors to $5\times10^{-15}$, \texttt{rust/spe\allowbreak{}ctrum/ex\allowbreak{}change-c\allowbreak{}heck.jso\allowbreak{}n}). Insert $\rho$ into $\mathrm{Tr}(BC\rho BC\rho)$: the four products of the two terms of $\rho$ give $f_+f_+'K_{++}-f_+f_-'K_{+-}-f_-f_+'K_{-+}+f_-f_-'K_{--}$ under the double integral (primes for the momentum $\mathbf q$). Collect the terms of the kernel:

\[
\mathrm{Tr}(BC\rho BC\rho)=4\!\int\!\!\int\Bigl[(f_+-f_-)(f_+'-f_-')+\frac{m^2-\mathbf p\cdot\mathbf q}{E_pE_q}(f_++f_-)(f_+'+f_-')\Bigr].
\]

For occupations that depend only on $|\mathbf p|$, the term with $\mathbf p\cdot\mathbf q$ integrates to zero, because the average of $\mathbf p\cdot\mathbf q$ over the directions of $\mathbf p$ vanishes (\texttt{KS\_excha\allowbreak{}nge\_angu\allowbreak{}larAvera\allowbreak{}geOfPdot\allowbreak{}QVanishe\allowbreak{}s}). The double integral then \textbf{factorizes} into $(n/8)^2$ and $(S/8)^2$:

\[
\mathrm{Tr}(BC\rho BC\rho)=4\Bigl[\bigl(\tfrac n8\bigr)^2+\bigl(\tfrac S8\bigr)^2\Bigr]=\frac{n^2+S^2}{16},\qquad
e_x(n,S)=-\frac\lambda{32}\bigl(n^2+S^2\bigr),\qquad e_H=\frac\lambda2S^2
\]

(erratum E4.7; \texttt{KS\_excha\allowbreak{}nge\_unif\allowbreak{}ormGasCl\allowbreak{}osedForm}). This closed form holds \textbf{exactly at every temperature} and in any number of space dimensions, since only the isotropy of the occupations was used; the temperature enters only through $n$ and $S$. A gas at rest ($k_F\to0$) has $S=n$ and $e_x=-\tfrac\lambda{16}n^2=-e_H/8$, the filled-shell value (\texttt{KS\_excha\allowbreak{}nge\_rest\allowbreak{}GasLimit}). The independent sympy checker tabulated the double integral by direct quadrature, and the closed form reproduces all 580 rows of that table to $5\times10^{-15}$ relative as evaluated by the checker (\texttt{python-t\allowbreak{}heory-re\allowbreak{}port.jso\allowbreak{}n}, entry \texttt{exchange\allowbreak{}Table}; the table is \texttt{artifact\allowbreak{}s/dirac1\allowbreak{}6complex\allowbreak{}/kohn-sh\allowbreak{}am/excha\allowbreak{}nge-tabl\allowbreak{}e.json}). Re-evaluated by the Rust solver from the tabulated columns in double precision, the largest relative deviation is $1.7\times10^{-14}$ for the 290 rows in $d=3$ and $5.3\times10^{-15}$ for the 290 rows in $d=4$ (checks \texttt{exchange\allowbreak{}\_table\_d\allowbreak{}3\_closed\allowbreak{}\_form} and \texttt{exchange\allowbreak{}\_table\_d\allowbreak{}4\_closed\allowbreak{}\_form} of \texttt{rust/spe\allowbreak{}ctrum/ex\allowbreak{}change-t\allowbreak{}able-che\allowbreak{}ck.json}); the right panel of the figure below shows these deviations row by row.

\begin{figure}[htbp]
\centering
\includegraphics[width=0.92\linewidth]{artifacts/dirac16complex/kohn-sham/figures/exchange_uniform_gas.png}
\caption{Left: the exchange energy density \texorpdfstring{$-e_x/\lambda$}{-e\_x/λ} of the uniform 8-fold gas against the number density at four temperatures, closed form (lines) and quadrature (markers). Right: the relative difference between quadrature and closed form for every row of the two tables, below \texorpdfstring{$10^{-12}$}{10\textasciicircum{}-12} everywhere (file exchange\_uniform\_gas.png in artifacts/dirac16complex/kohn-sham/figures).}
\end{figure}

\textbf{The local potentials.} Evaluated with the local proper densities $n_p(y)$ and $S_p(y)$, the closed form is the local density approximation of Section 12.12 for this system. Its two derivatives are the \textbf{exchange potentials}

\[
v_v=\frac{\partial e_x}{\partial n}=-\frac\lambda{16}\,n,\qquad v_s=\frac{\partial e_x}{\partial S}=-\frac\lambda{16}\,S
\]

(\texttt{KS\_excha\allowbreak{}nge\_ldaP\allowbreak{}otential\allowbreak{}s}). The vector potential $v_v$ multiplies the number density matrix and enters the equation as $\varepsilon\to\varepsilon-v_v(y)$; the scalar potential $v_s$ multiplies the scalar density matrix $BC$ and adds to the mass. The Kohn--Sham equation used in all runs (called KS-VS in the theory file) is the block equation of Section 13.5 with

\[
M_{\text{eff}}(y)=m+\lambda S_p(y)+v_s(y)=m+\tfrac{15}{16}\lambda S_p(y),\qquad v(y)=v_v(y)=-\tfrac1{16}\lambda\,n_p(y).
\]

The specification calls this pair $(M_{\text{eff}},v_v)$, together with the gravitational terms, the \textbf{Kohn--Sham fermion-gas thermodynamic pseudo-potential}.

\textbf{How good is the local form?} For the contact interaction the exact Fock term is local in $y$ as well (Section 12.6): with the $2\times2$ density matrix $\rho_\beta=\tfrac12(n_\beta+s_\beta\sigma_y+t_\beta\sigma_z+c_\beta\sigma_x)$ of each block $\beta$ (up to the signs $j$), $E_x^{HF}=-\tfrac\lambda4\int dV\sum_\beta\bigl(n_\beta^2+s_\beta^2-t_\beta^2-c_\beta^2\bigr)$ (\texttt{KS\_excha\allowbreak{}nge\_bloc\allowbreak{}kFormOfF\allowbreak{}ockTerm}). If the eight blocks carry equal densities $n/8$ and $S/8$ and no momentum current, this is $-\tfrac\lambda4\cdot8\bigl[(n/8)^2+(S/8)^2\bigr]=-\tfrac\lambda{32}(n^2+S^2)$, the closed form. Otherwise the two differ by the $t_\beta^2$ terms and by the unequal occupation of the blocks. The Kohn--Sham scheme of this chapter uses the closed form; that choice is the \textbf{approximation} of the exchange part.

\textbf{No correlation term.} Beyond Hartree--Fock a contact interaction with a coupling of mass dimension $-6$ is not renormalizable: its higher-order corrections cannot be made finite with a finite number of parameters (the standard power-counting argument of quantum field theory, which we quote and do not derive). A correlation energy would therefore depend on a cutoff and on renormalization choices that the project does not make. The scheme is \textbf{exchange-only} by choice, an assumption listed in Section 13.16: the uniform-gas Hartree--Fock energy written as a local functional of $n$ and $S$.

\subsection{The Mermin--Kohn--Sham functional and its equations}

\textbf{The functional.} Let the orbitals $\chi_n$ (each in one block, with a block type $j_n$, a momentum $k_n$ and a parity) be normalized, $\int\chi_n^\dagger\chi_n\,dy=1$, with occupations $f_n$, and let $dV=e^{6Hy}\ell^3dy$ be the proper volume element per unit coordinate extra-time volume. Mermin's free-energy functional (Section 12.14) is

\[
\begin{aligned}
F&=\sum_nf_n\langle\chi_n|h_0|\chi_n\rangle-T\,S_{\text{ent}}+E_H+E_x,\qquad h_0=j\Bigl[-i\sigma_x\frac d{dy}+m\sigma_y+\xi k\sigma_z\Bigr],\\
E_H&=\frac\lambda2\int S_p^2\,dV,\qquad E_x=-\frac\lambda{32}\int\bigl(n_p^2+S_p^2\bigr)\,dV,\qquad S_{\text{ent}}=-\sum_n\bigl[f_n\ln f_n+(1-f_n)\ln(1-f_n)\bigr],
\end{aligned}
\]

with $n_p=e^{-6Hy}\sum_nf_n\chi_n^\dagger\chi_n/\ell^3$ and $S_p=e^{-6Hy}\sum_nf_n\chi_n^\dagger(j_n\sigma_y)\chi_n/\ell^3$ (\texttt{kohn-sha\allowbreak{}m-theory\allowbreak{}.json}, key \texttt{function\allowbreak{}al}). (In the solver the sums run over the particle and sea levels with the weights $w_n$ of Section 13.7.)

\textbf{Stationarity with respect to the orbitals.} Vary $\chi_n^\dagger(y)$ with multipliers $f_n\varepsilon_n$ for the normalization. The derivative of $\int S_p^2\,dV$ is $2S_p(y)\cdot e^{6Hy}\ell^3\cdot e^{-6Hy}f_n(j_n\sigma_y)\chi_n(y)/\ell^3=2f_nS_p\,j_n\sigma_y\chi_n$: the factor $e^{6Hy}$ of the volume element cancels the $e^{-6Hy}$ of the density, so the equation is local with the flat measure $dy$. In the same way $E_x$ gives $-\tfrac\lambda{16}f_n\bigl(n_p+S_pj_n\sigma_y\bigr)\chi_n=f_n(v_v+v_sj_n\sigma_y)\chi_n$. Dividing by $f_n$:

\[
\Bigl[h_0+(\lambda S_p+v_s)\,j\sigma_y+v_v\Bigr]\chi_n=\varepsilon_n\chi_n
\quad\Longleftrightarrow\quad
j\Bigl[-i\sigma_x\frac d{dy}+M_{\text{eff}}\sigma_y+\xi k\sigma_z\Bigr]\chi_n+v_v\chi_n=\varepsilon_n\chi_n ,
\]

the Hamiltonian form of Section 13.5 with the potentials of Section 13.8 (\texttt{KS\_funct\allowbreak{}ional\_st\allowbreak{}ationari\allowbreak{}tyGivesK\allowbreak{}SEquatio\allowbreak{}n}).

\textbf{Stationarity with respect to the occupations.} Exactly as in Section 12.14, $\partial F/\partial f_n=\varepsilon_n+T\ln\frac{f_n}{1-f_n}-\mu=0$ gives the Fermi--Dirac occupations $f_n=1/(e^{(\varepsilon_n-\mu)/T}+1)$, with $\mu$ fixed by the particle number (\texttt{KS\_funct\allowbreak{}ional\_fe\allowbreak{}rmiDirac\allowbreak{}Occupati\allowbreak{}ons}).

\textbf{Total energy.} Both interaction functionals are quadratic in the densities, so $\sum_nf_n\varepsilon_n=\sum_nf_n\langle h_0\rangle_n+2E_H+2E_x$ (each quadratic functional contributes twice itself through its potential). Hence

\[
E=\sum_nf_n\varepsilon_n-E_H-E_x,\qquad F=E-TS_{\text{ent}}
\]

(\texttt{KS\_funct\allowbreak{}ional\_to\allowbreak{}talEnerg\allowbreak{}yDoubleC\allowbreak{}ounting}). \textbf{Worked example.} For $N=112$ at $\hat\lambda_1$ the solver reports $\sum w\varepsilon=61.1358649$, $E_H=0.0342610$ and $E_x=-0.0116081$ (\texttt{rust/scf\allowbreak{}/m1\_L3\_N\allowbreak{}112\_lamp\allowbreak{}1\_T0/run\allowbreak{}.json}, keys \texttt{\detokenize{ksSum}}, \texttt{hartreeE\allowbreak{}nergy}, \texttt{exchange\allowbreak{}Energy}). Then $E=61.1358649-0.0342610+0.0116081=61.1132120$, the value of the key \texttt{\detokenize{energy}}.

\textbf{Hellmann--Feynman relations.} Because $F$ is stationary, its derivative with respect to a parameter is the explicit partial derivative (the envelope theorem of Section 12.14): $dF/d\lambda=(E_H+E_x)/\lambda$, $dF/dm=\int S_p\,dV$ (the total scalar charge), $dF/dT=-S_{\text{ent}}$, and for each level $d\varepsilon_n/dk=j\int\xi\,\chi_n^\dagger\sigma_z\chi_n\,dy$ at fixed potentials. The checks of the family \texttt{KS\_funct\allowbreak{}ional} are named after these relations: \texttt{hellmann\allowbreak{}FeynmanL\allowbreak{}ambda}, \texttt{hellmann\allowbreak{}FeynmanM\allowbreak{}ass} and \texttt{hellmann\allowbreak{}FeynmanT\allowbreak{}emperatu\allowbreak{}re} in both exact reports; \texttt{hellmann\allowbreak{}FeynmanM\allowbreak{}omentum}, for $d\varepsilon_n/dk$, in \texttt{python-t\allowbreak{}heory-re\allowbreak{}port.jso\allowbreak{}n} only; and \texttt{merminSt\allowbreak{}ructure}, in both reports, for the thermodynamic relations of Section 12.14. The Rust solver checks $dE/dm=\int S_p\,dV$ by finite differences in two runs (checks \texttt{m1\_L3\_N1\allowbreak{}12\_lam0\_\allowbreak{}T0\_hellm\allowbreak{}ann\_feyn\allowbreak{}man\_dE\_d\allowbreak{}m} and \texttt{m1\_L3\_N1\allowbreak{}12\_lamp1\allowbreak{}\_T0\_hell\allowbreak{}mann\_fey\allowbreak{}nman\_dE\_\allowbreak{}dm} of \texttt{rust/scf\allowbreak{}/summary\allowbreak{}.json}). For $\lambda=0$ the scalar charge $\int\chi^\dagger j\sigma_y\chi\,dy$ of a single level is $d\varepsilon/dM$ (Hellmann--Feynman with $\partial h_j/\partial M=j\sigma_y$). At $k=0$ it is $M/\varepsilon$ for the massive levels of parity $+$, whose energy $\sqrt{M^2+(n\pi/L)^2}$ was found in Section 13.6 (for example $1/1.447972=0.6906$); it is positive for the other massive levels too, 0 for the zero modes, and \textbf{negative} for the brane band, because $c(M)$ of Section 13.6 decreases toward 1 as $M$ grows, so $d\varepsilon/dM\approx k\,dc/dM<0$ (the header of \texttt{src/emt.\allowbreak{}rs}; column \texttt{scalar\_c\allowbreak{}harge} of the \texttt{levels.c\allowbreak{}sv} files).

\subsection{How the Rust solver computes the Kohn--Sham states}

The solver \texttt{studies/\allowbreak{}dirac16c\allowbreak{}omplex\_k\allowbreak{}ohn\_sham} integrates the block equation as an ordinary differential equation in $y$ with the CVODE engine of Chapter 10; nothing is discretized into a matrix. Its module headers contain the derivations summarized here.

\textbf{The Prüfer angle.} Write the real solution as $a=r\cos\theta$, $b=-r\sin\theta$. From the real system of Section 13.5, $\theta'=(ba'-ab')/r^2$ and $(\ln r)'=(aa'+bb')/r^2$. Inserting $a'$ and $b'$, and using $ab=-\tfrac12r^2\sin2\theta$ and $a^2-b^2=r^2\cos2\theta$,

\[
\theta'=-j\bigl(\varepsilon-v(y)\bigr)+\xi(y)\,k\cos2\theta-M_{\text{eff}}(y)\sin2\theta,\qquad (\ln r)'=M_{\text{eff}}\cos2\theta+\xi k\sin2\theta .
\]

The first equation does not contain $r$: the angle can be integrated alone. The tip condition $b(-L)=0$ is $\theta(-L)=0$. At the brane, parity $+$ ($b(0)=0$) means $\sin\theta(0)=0$ and parity $-$ ($a(0)=0$) means $\cos\theta(0)=0$. The solver treats the block type $j=-1$ (its $s=+1$), for which $\partial\theta'/\partial\varepsilon=+1>0$. Then the end value $\Theta(\varepsilon)=\theta(0;\varepsilon)$ \textbf{increases strictly} with $\varepsilon$ (a solution with a larger $\varepsilon$ can never be overtaken by one with a smaller $\varepsilon$, since at a crossing point it would have the larger slope). It is also continuous in $\varepsilon$ (the solution of a differential equation depends continuously on a parameter in it, a standard fact that we quote), and it is unbounded in both directions. Indeed, for $j=-1$ the angle equation reads $\theta'=\varepsilon+\bigl[-v+\xi k\cos2\theta-M_{\text{eff}}\sin2\theta\bigr]$, and on $[-L,0]$ the bracket is bounded, $\lvert\,\cdot\,\rvert\le C$ with $C$ the largest value of $|v|+\xi|k|+|M_{\text{eff}}|$ there (note $\xi\le e^{HL-a_{4,0}}$). So $\varepsilon-C\le\theta'\le\varepsilon+C$, and integrating from $\theta(-L)=0$ gives $(\varepsilon-C)L\le\Theta(\varepsilon)\le(\varepsilon+C)L$: $\Theta$ grows without bound, roughly like $\varepsilon L$, as $\varepsilon\to\pm\infty$. A continuous, strictly increasing function that runs from $-\infty$ to $+\infty$ takes every real value exactly once, so

\[
\text{parity }+:\ \Theta(\varepsilon)=n\pi,\qquad \text{parity }-:\ \Theta(\varepsilon)=\tfrac\pi2+n\pi,\qquad n\in\mathbb Z,
\]

has exactly one solution for each integer $n$, the \textbf{Prüfer index} of the level. This is the oscillation theorem of the one-dimensional Dirac operator: every level in an energy window is found and none is missed. The block type $j=+1$ follows from $h_{+1}=-h_{-1}+2v$ (Section 13.5): its levels are those of the $j=-1$ problem with $\varepsilon\to-\varepsilon$ and $v\to-v$, with the same profiles.

\textbf{Finding a level and its profile.} For each momentum shell and parity, the solver brackets each target value of $\Theta$ by doubling the step in $\varepsilon$, then refines it by the Illinois variant of the regula falsi (a root finder that keeps the root bracketed). At the level it integrates the linear system for $(a,b)$ once more, together with the integrals $\int(a^2+b^2)$, $\int2ab$ and $\int\xi k(a^2-b^2)$, rescales the solution whenever $a^2+b^2$ leaves $[10^{-40},10^{40}]$, and normalizes it. The CVODE tolerances are $10^{-12}$ (relative) and $10^{-14}$ (absolute) with the largest step $0.05$ (\texttt{rust/spe\allowbreak{}ctrum/su\allowbreak{}mmary.js\allowbreak{}on}, key \texttt{toleranc\allowbreak{}es}). The potentials $M_{\text{eff}}(y)$ and $v(y)$ are tabulated on 301 equally spaced points of $[-L,0]$ and interpolated between them by a \textbf{natural cubic spline}: on each interval between two neighbouring points the interpolant is a polynomial of degree 3 that takes the tabulated values at both ends, and the pieces are chosen so that the first and the second derivative are continuous at every inner point, with the second derivative zero at the two ends (\texttt{src/spli\allowbreak{}ne.rs}); one run uses 601 points as a refinement (Section 13.15).

\textbf{The coupling strengths.} The proper densities near the tip grow like $e^{6HL}$ and with $N$, so one value of $\hat\lambda$ cannot give the same strength of interaction in every configuration. For each configuration $(m,L,N)$ the solver therefore measures, in the free ground state, the largest value of the pair $\bigl(\tfrac{15}{16}|S_p|,\ n_p/16\bigr)$ per unit $\hat\lambda$ (the \textbf{strength}, in units of $m$), and sets $\hat\lambda_1=0.1/\text{strength}$ and $\hat\lambda_2=1/\text{strength}$: at first order the pseudo-potential then reaches $0.1\,m$ and $1\,m$ (\texttt{coupling\allowbreak{}Rule} in \texttt{rust/spe\allowbreak{}ctrum/su\allowbreak{}mmary.js\allowbreak{}on}). For $m=1$, $L=3$ this gives $\hat\lambda_1=9.72989\times10^{-3}$ ($N=8$), $9.56776\times10^{-3}$ ($N=112$) and $1.57694\times10^{-4}$ ($N=1016$), and $\hat\lambda_2=10\,\hat\lambda_1$ in each case. Each configuration is solved for $\hat\lambda\in\{0,\pm\hat\lambda_1,\pm\hat\lambda_2\}$; $\hat\lambda>0$ is a repulsive and $\hat\lambda<0$ an attractive contact interaction.

\textbf{The self-consistent loop.} Starting from the free state, the loop of Section 12.11 is run with the densities $(n_c,S_c)$ as the iterated quantity and \textbf{Anderson mixing} (Section 12.13; mixing parameter 0.4, history of 6, at most 80 iterations; \texttt{paramete\allowbreak{}rs} in each \texttt{\detokenize{run.json}}). It stops when the largest change of $n_c$ and of $S_c$, divided by the largest of $|n_c|$ and $|S_c|$, is below $10^{-10}$. At $T>0$ the chemical potential is found by bisection at every iteration (Section 12.14).

\textbf{The ground state when the loop cannot settle} (erratum E4.8). In one canonical run ($N=1016$, $-\hat\lambda_2$) two levels cross at the Fermi level during the iteration, the integer occupations flip back and forth, and the loop stagnates (Section 12.13). The \textbf{ground state} is then \emph{defined} as the self-consistent solution reached by \textbf{continuation in the coupling} from $\lambda=0$: the solver solves at $\lambda/4$, $\lambda/2$, $3\lambda/4$ and $\lambda$, each step started from the previous solution with damped mixing, first with exact occupations and, where they slosh, with Fermi--Dirac occupation smearing of width $10^{-3}\,m$ (then $10^{-2}\,m$). That run converged with smearing $10^{-3}\,m$ at the physical temperature $T=0$ (so $F=E$); \texttt{\detokenize{run.json}} records the path (\texttt{zeroTemp\allowbreak{}eratureF\allowbreak{}allbackS\allowbreak{}tage} 1, \texttt{occupati\allowbreak{}onSmeari\allowbreak{}ng} 0.001, \texttt{exactZer\allowbreak{}oTempera\allowbreak{}tureOccu\allowbreak{}pations} false). Section 13.11 describes the resulting ensemble.

\textbf{The solver's own checks.} Every self-consistent run is checked for convergence, for the particle number ($\sum w=N$ and $\ell^3\int n_c\,dy=N$), for the energy identity $\int\rho\,dV=E$ of Section 13.14, and for the conservation of the energy--momentum tensor. With the checks specific to each subcommand (spectrum and theory agreement, excited states, thermodynamics, energy--momentum tensor), all 397 checks of the five \texttt{summary.\allowbreak{}json} files pass. A second run of the same build is byte-identical in all 330 output files (key \texttt{repeatFi\allowbreak{}lesCompa\allowbreak{}red}), and a run with tolerances divided by 10 and half the largest step agrees with the canonical one to $5.96\times10^{-8}$ relative in $E$ and $F$ (65 runs) and to $2.8\times10^{-8}$ absolute in 364829 eigenvalues (\texttt{rust/det\allowbreak{}erminism\allowbreak{}-report.\allowbreak{}json}).

\subsection{Ground states}

The table lists the self-consistent ground states for $m=1$, $L=3$ (\texttt{rust/scf\allowbreak{}/summary\allowbreak{}.json}; $E_0$, $\Delta_{KS}$ and the last column are also in \texttt{rust/exc\allowbreak{}ited/exc\allowbreak{}itations\allowbreak{}.csv}, whose column \texttt{\detokenize{mu}} equals $\varepsilon_{\text{HOMO}}$ except in the smeared last row, where $\mu=1.389567$). $E_0$ is the total energy of Section 13.9 relative to the filled Dirac sea (normal ordering), $\varepsilon_{\text{HOMO}}$ the highest occupied particle level, $\Delta_{KS}$ the Kohn--Sham gap, and the last column is the largest Hartree mass shift reached, $\max_y|\lambda S_p|/m$. All energies are in units of $m$.

\begingroup
\small
\begin{longtable}{@{}>{\raggedright\arraybackslash}p{0.133\linewidth}>{\raggedright\arraybackslash}p{0.133\linewidth}>{\raggedright\arraybackslash}p{0.133\linewidth}>{\raggedright\arraybackslash}p{0.133\linewidth}>{\raggedright\arraybackslash}p{0.133\linewidth}>{\raggedright\arraybackslash}p{0.133\linewidth}@{}}
\toprule
\textbf{$N$} & \textbf{$\hat\lambda/\hat\lambda_1$} & \textbf{$E_0$} & \textbf{$\varepsilon_{\text{HOMO}}$} & \textbf{$\Delta_{KS}$} & \textbf{$\max\lvert\lambda S_p\rvert/m$} \\
\midrule
\endfirsthead
\toprule
\textbf{$N$} & \textbf{$\hat\lambda/\hat\lambda_1$} & \textbf{$E_0$} & \textbf{$\varepsilon_{\text{HOMO}}$} & \textbf{$\Delta_{KS}$} & \textbf{$\max\lvert\lambda S_p\rvert/m$} \\
\midrule
\endhead
8 & 0 & 0 & 0 & 0.430734 & 0 \\
8 & +1 & $-0.000986833$ & $-0.000245560$ & 0.430944 & 0.0165 \\
8 & $-1$ & $+0.000986833$ & $+0.000245560$ & 0.430516 & 0.0165 \\
8 & +10 & $-0.00783039$ & $-0.00168579$ & 0.431899 & 0.647 \\
8 & $-10$ & $+0.00783039$ & $+0.00168579$ & 0.429137 & 0.647 \\
112 & 0 & 61.090723 & 0.704100 & 0.0955462 & 0 \\
112 & +1 & 61.113212 & 0.704476 & 0.0954841 & 0.0592 \\
112 & $-1$ & 61.068610 & 0.703736 & 0.0956051 & 0.0583 \\
112 & +10 & 61.328311 & 0.708384 & 0.0948132 & 1.358 \\
112 & $-10$ & 60.891305 & 0.700966 & 0.0960102 & 0.452 \\
1016 & 0 & 1127.136625 & 1.385919 & 0.0530077 & 0 \\
1016 & +1 & 1127.160980 & 1.385934 & 0.0532575 & 0.100 \\
1016 & $-1$ & 1127.111039 & 1.385904 & 0.0527200 & 0.114 \\
1016 & +10 & 1127.342203 & 1.386070 & 0.0544240 & 0.665 \\
1016 & $-10$ & 1126.858093 & 1.389461 & 0.0412551 & 1.372 \\
\bottomrule
\end{longtable}
\endgroup

($\hat\lambda/\hat\lambda_1=\pm10$ means $\pm\hat\lambda_2$. The last row is the smeared ensemble described below.) The largest exchange potential, $\max_y|v|/m$, is 0.591 for $N=8$ at $\pm\hat\lambda_2$, 0.893 for $N=112$ at $+\hat\lambda_2$ and 1.50 for $N=1016$ at $-\hat\lambda_2$: the chosen couplings put the pseudo-potential into the range from about $0.1\,m$ to $1.5\,m$, as intended.

\textbf{$N=8$: the brane zero modes.} Without interaction the ground state fills exactly the eight zero modes of Section 13.6; they have $\varepsilon=0$, so $E_0=0$, and their scalar density vanishes identically ($\chi_2\equiv0$ gives $\chi^\dagger\sigma_y\chi=0$). Their coordinate density is $n_c\propto e^{2my}$, so the fraction of the particles within the distance $1/H$ of the brane is $\int_{-1}^0e^{2y}dy\big/\int_{-3}^0e^{2y}dy=(1-e^{-2})/(1-e^{-6})=0.866813$; the solver reports $0.8668133$ (\texttt{braneFra\allowbreak{}ction\_wi\allowbreak{}thin\_1\_o\allowbreak{}ver\_H} in \texttt{rust/scf\allowbreak{}/m1\_L3\_N\allowbreak{}8\_lam0\_T\allowbreak{}0/run.js\allowbreak{}on}). With interaction, the exchange potential $v=-\lambda n_p/16$ shifts the zero modes to $\varepsilon\approx\langle v\rangle$, negative for $\lambda>0$. The table shows $E_0(-\lambda)=-E_0(\lambda)$ to all printed digits; Exercise 13.9 derives this exact property of the $N=8$ state.

\textbf{$N=112$ and $N=1016$: the brane band.} The larger states fill the zero modes and then the brane band shell by shell (Section 13.7). For $\lambda=0$ the occupied particle levels of $N=112$ are the zero modes and the band at $k=0.25$, $0.354$ and $0.433\,m$ ($\varepsilon=0.430734$, $0.587957$, $0.704100$); the lowest empty level is the band at $k=0.5\,m$ ($0.799646$), so $\Delta_{KS}=0.0955462$. For $N=1016$ the band is filled up to $k=0.935\,m$ (a shell of 192 states at $\varepsilon=1.385919$), and in addition the eight bulk states of parity $-$ at $k=0$ with $\varepsilon=1.292293$ (Section 13.6) are occupied (\texttt{rust/scf\allowbreak{}/m1\_L3\_N\allowbreak{}1016\_lam\allowbreak{}0\_T0/lev\allowbreak{}els.csv}). The lowest empty level of $N=1016$ is not a band level: it is the bulk level of parity $-$ at $k=0.25\,m$ ($\varepsilon=1.438927$), the continuation to $k=0.25\,m$ of that parity-minus $k=0$ level, so $\Delta_{KS}=1.438927-1.385919=0.0530077$ (Section 13.12). The interaction changes the energies by at most $0.3\,m$ out of 61 and 1127: the interaction is a perturbation, and the self-consistent energy shifts follow the first-order Hellmann--Feynman estimate $\lambda\,dE/d\lambda$ closely (figure \texttt{ground\_s\allowbreak{}tate\_ene\allowbreak{}rgy.png} in \texttt{artifact\allowbreak{}s/dirac1\allowbreak{}6complex\allowbreak{}/kohn-sh\allowbreak{}am/figur\allowbreak{}es}, not reproduced here). The fraction of particles within $1/H$ of the brane grows with $N$: 0.867 ($N=8$), 0.924 ($N=112$), 0.962 ($N=1016$) for $\lambda=0$; for $m=3$ it is 0.9975 ($N=8$) and 0.9990 ($N=112$), because a heavier zero mode $e^{My}$ is more strongly localized.

\begin{figure}[htbp]
\centering
\includegraphics[width=0.92\linewidth]{artifacts/dirac16complex/kohn-sham/figures/density_profiles.png}
\caption{Self-consistent ground states for \texorpdfstring{$m=1$}{m=1}, \texorpdfstring{$L=3$}{L=3}: the coordinate density \texorpdfstring{$n_c$}{n\_c} (particles per unit \texorpdfstring{$y$}{y}), the proper number density \texorpdfstring{$n_p$}{n\_p}, the proper scalar density \texorpdfstring{$S_p$}{S\_p}, the Hartree--exchange mass shift \texorpdfstring{$M_{\text{eff}}-m=\tfrac{15}{16}\lambda S_p$}{M\_eff-m=1516λS\_p}, the exchange potential \texorpdfstring{$v_x=-\lambda n_p/16$}{v\_x=-λn\_p/16}, and the highest occupied orbital against the notebook coordinate \texorpdfstring{$z=6Hx_0$}{z=6Hx\_0} (file density\_profiles.png in artifacts/dirac16complex/kohn-sham/figures).}
\end{figure}

\textbf{Where the interaction acts.} The figure shows the point of Section 13.7. The particles sit near the brane ($n_c$ largest at $y=0$), but the proper densities $n_p=e^{-6Hy}n_c$ are largest at the tip, so the pseudo-potential $(M_{\text{eff}}-m,\ v)$ is concentrated near $y=-L$. The scalar density is negative near the brane, where the brane band dominates, and for the larger $N$ it becomes positive near the tip, where the bulk levels contribute (Section 13.9).

\textbf{The level crossing at $N=1016$, $-\hat\lambda_2$} (erratum E4.8). The attraction pulls the 8-fold $k=0$ bulk level down from $1.447972$ to the Fermi level, where it meets the 192-fold band at $k=0.935\,m$. The converged smeared ensemble (Section 13.10) moves about 4.2 particles from the band (occupation 0.978059) into the $k=0$ level (occupation 0.526580, so $8\times0.526580=4.21$ particles), which lies $3.69\times10^{-3}\,m$ above the band (\texttt{rust/scf\allowbreak{}/m1\_L3\_N\allowbreak{}1016\_lam\allowbreak{}m2\_T0/le\allowbreak{}vels.csv}, column \texttt{\detokenize{f}}); its $\varepsilon_{\text{HOMO}}=1.389461$ is that fractionally occupied level, and its Kohn--Sham gap $0.0412551\,m$ is the distance to the parity-minus bulk level at $k=0.25\,m$ ($1.430716$; Section 13.12). This is an ensemble with fractional occupations, not a single determinant, and its numbers refer to that ensemble. An exploratory reference computation on twice finer grids converged, from its own start, to a different self-consistent state about $1.01\,m$ higher (\texttt{handoff/\allowbreak{}reviews/\allowbreak{}stage4\_r\allowbreak{}efgrid\_N\allowbreak{}0120\_exp\allowbreak{}loratory\allowbreak{}.log}): the nonlinear problem has more than one self-consistent solution here, and the continuation from $\lambda=0$ is what selects the one reported.

\begin{figure}[htbp]
\centering
\includegraphics[width=0.92\linewidth]{artifacts/dirac16complex/kohn-sham/figures/fermi_level_N1016_lamm2.png}
\caption{The level crossing at the Fermi level for \texorpdfstring{$N=1016$}{N=1016}: particle levels near the Fermi level against \texorpdfstring{$k$}{k} (marker area proportional to the number of states, colour the occupation), free (left) and in the smeared self-consistent ensemble at \texorpdfstring{$-\hat\lambda_2$}{-λ\_2} (right), where the 8-fold \texorpdfstring{$k=0$}{k=0} level and the 192-fold band at \texorpdfstring{$k=0.935\,m$}{k=0.935 m} share the particles (file fermi\_level\_N1016\_lamm2.png in artifacts/dirac16complex/kohn-sham/figures).}
\end{figure}

\begin{figure}[htbp]
\centering
\includegraphics[width=0.92\linewidth]{artifacts/dirac16complex/kohn-sham/figures/scf_history.png}
\caption{Convergence of the self-consistent loop: the density residual against the iteration for \texorpdfstring{$+\hat\lambda_2$}{+λ\_2} (left) and \texorpdfstring{$-\hat\lambda_2$}{-λ\_2} (right) and \texorpdfstring{$N=8$}{N=8}, 112, 1016; the tolerance is \texorpdfstring{$10^{-10}$}{10\textasciicircum{}-10}. The slow, irregular curve is the \texorpdfstring{$N=1016$}{N=1016} run at \texorpdfstring{$-\hat\lambda_2$}{-λ\_2} with its level crossing (file scf\_history.png in artifacts/dirac16complex/kohn-sham/figures).}
\end{figure}

\subsection{First excited states: Kohn--Sham gap, particle--hole pairs and Delta-SCF}

The three quantities of Section 12.15 are computed for every $T=0$ ground state of Section 13.11 (\texttt{rust/exc\allowbreak{}ited/}). The \textbf{particle--hole list} (\texttt{particle\allowbreak{}-hole.cs\allowbreak{}v} of each run) contains the differences $\varepsilon_{\text{particle}}-\varepsilon_{\text{hole}}$ in increasing order, where a level counts as a possible hole if its occupation is $f>10^{-12}$ and as a possible particle if $f<1-10^{-12}$ (erratum E4.9; at $T>0$ the thresholds are $f\ge\tfrac12$ and $f<\tfrac12$). The \textbf{Delta-SCF} state fixes the occupations by the identity of the levels (shell, parity, block type, Prüfer index), moves one particle from the HOMO group to the LUMO group, and is converged again; $\Delta_{\text{SCF}}=E_1-E_0$ (the header of \texttt{src/scf.\allowbreak{}rs}).

\begingroup
\small
\begin{longtable}{@{}>{\raggedright\arraybackslash}p{0.160\linewidth}>{\raggedright\arraybackslash}p{0.160\linewidth}>{\raggedright\arraybackslash}p{0.160\linewidth}>{\raggedright\arraybackslash}p{0.160\linewidth}>{\raggedright\arraybackslash}p{0.160\linewidth}@{}}
\toprule
\textbf{$N$} & \textbf{$\hat\lambda/\hat\lambda_1$} & \textbf{$\Delta_{KS}$} & \textbf{$\Delta_{\text{SCF}}$} & \textbf{$\Delta_{\text{SCF}}-\Delta_{KS}$} \\
\midrule
\endfirsthead
\toprule
\textbf{$N$} & \textbf{$\hat\lambda/\hat\lambda_1$} & \textbf{$\Delta_{KS}$} & \textbf{$\Delta_{\text{SCF}}$} & \textbf{$\Delta_{\text{SCF}}-\Delta_{KS}$} \\
\midrule
\endhead
8 & 0 & 0.430734 & 0.430734 & 0 \\
8 & +1 & 0.430944 & 0.430938 & $-5.79\times10^{-6}$ \\
8 & $-1$ & 0.430516 & 0.430523 & $+7.50\times10^{-6}$ \\
8 & +10 & 0.431899 & 0.431946 & $+4.72\times10^{-5}$ \\
8 & $-10$ & 0.429137 & 0.429161 & $+2.32\times10^{-5}$ \\
112 & 0 & 0.0955462 & 0.0955462 & 0 \\
112 & +1 & 0.0954841 & 0.0954842 & $+5.84\times10^{-8}$ \\
112 & $-1$ & 0.0956051 & 0.0956051 & $-5.18\times10^{-8}$ \\
112 & +10 & 0.0948132 & 0.0948140 & $+8.41\times10^{-7}$ \\
112 & $-10$ & 0.0960102 & 0.0960098 & $-3.28\times10^{-7}$ \\
1016 & 0 & 0.0530077 & 0.0530077 & $-3.7\times10^{-10}$ \\
1016 & +1 & 0.0532575 & 0.0532689 & $+1.14\times10^{-5}$ \\
1016 & $-1$ & 0.0527200 & 0.0527082 & $-1.18\times10^{-5}$ \\
1016 & +10 & 0.0544240 & 0.0545261 & $+1.02\times10^{-4}$ \\
1016 & $-10$ & 0.0412551 & 0.0436184 & $+2.36\times10^{-3}$ \\
\bottomrule
\end{longtable}
\endgroup

(Energies in units of $m$, $m=1$, $L=3$; \texttt{rust/exc\allowbreak{}ited/sum\allowbreak{}mary.jso\allowbreak{}n}. The last row belongs to the smeared ensemble.)

\textbf{What the first excited state is.} For $N=8$ the lowest particle--hole pair takes a particle out of a zero mode ($k=0$) and puts it into the first shell of the brane band ($k=0.25\,m$); the next pairs go to $k=0.354\,m$ (excitation $0.587957$) and so on (\texttt{rust/exc\allowbreak{}ited/m1\_\allowbreak{}L3\_N8\_la\allowbreak{}m0\_T0/pa\allowbreak{}rticle-h\allowbreak{}ole.csv}). For $N=112$ the first excitation moves a particle from the band shell at $k=0.433\,m$ to the next band shell, $k=0.5\,m$ (\texttt{rust/exc\allowbreak{}ited/m1\_\allowbreak{}L3\_N112\_\allowbreak{}lam0\_T0/\allowbreak{}particle\allowbreak{}-hole.cs\allowbreak{}v}). For $N=8$ and $N=112$ the first excited state is therefore an excitation \textbf{within the brane band}, and its energy is set by the spacing of the torus momenta ($\Delta k=0.25\,m$) along the band $\varepsilon\approx ck$. For $N=1016$ the lowest empty level is not a band level. It is the bulk level of parity $-$ at $k=0.25\,m$ ($\varepsilon=1.438927$ at $\lambda=0$), which has the parity, the block type ($j=-1$, Rust label $s=+1$) and the Prüfer index (column \texttt{\detokenize{index}}) of the $j=-1$ member of the 8-fold parity-minus $k=0$ level $1.292293$ of Section 13.6, and a positive scalar charge ($0.789$, column \texttt{scalar\_c\allowbreak{}harge} of \texttt{rust/scf\allowbreak{}/m1\_L3\_N\allowbreak{}1016\_lam\allowbreak{}0\_T0/lev\allowbreak{}els.csv}). So for $\hat\lambda=0$, $\pm\hat\lambda_1$ and $+\hat\lambda_2$ the first excitation moves a particle from the band at $k=0.935\,m$ into this bulk level ($0.0530077\,m$ at $\lambda=0$, \texttt{rust/exc\allowbreak{}ited/m1\_\allowbreak{}L3\_N1016\allowbreak{}\_lam0\_T0\allowbreak{}/particl\allowbreak{}e-hole.c\allowbreak{}sv}); its energy is not set by the band spacing. The ensemble at $-\hat\lambda_2$ is described below. Without interaction the levels do not move, so $\Delta_{\text{SCF}}=\Delta_{KS}$ exactly (the value $-3.7\times10^{-10}$ in the row $N=1016$, $\hat\lambda=0$ is at the level of the numerical accuracy of the two separate solutions). With interaction the orbital relaxation of Section 12.15 is at most $1.02\times10^{-4}\,m$ in all states with a single determinant; its sign depends on the case.

\textbf{The ensemble at $N=1016$, $-\hat\lambda_2$.} Here the particle--hole list starts with $6.6\times10^{-12}\,m$: this is not an excitation but the difference between the two numerically split halves of the fractionally occupied $k=0$ level (Section 13.11). Next comes $3.69\times10^{-3}\,m$ (a particle from the band at $k=0.935\,m$ into the $k=0$ level), then the Kohn--Sham gap $0.0412551\,m$ (from the $k=0$ level into the parity-minus bulk level at $k=0.25\,m$, $\varepsilon=1.430716$). The Delta-SCF excitation of this ensemble is $0.0436184\,m$ on the standard grid of 301 points and $0.0436194\,m$ on 601 points (run \texttt{m1\_L3\_N1\allowbreak{}016\_lamm\allowbreak{}2\_T0\_g60\allowbreak{}1}, check \texttt{excited\_\allowbreak{}grid\_ref\allowbreak{}inement\_\allowbreak{}delta\_sc\allowbreak{}f}).

\textbf{The cross-check of this number} (erratum E4.13 of \texttt{handoff/\allowbreak{}specs/ST\allowbreak{}AGE4\_SPE\allowbreak{}C.md}). This was the hardest number of the chapter to confirm. The reference solver first computed the state on its standard grids of 60, 120 and 240 points; its three Delta-SCF values $0.0430172$, $0.0434646$ and $0.0435817\,m$ were not yet close enough to their limit for an extrapolation in the grid spacing, which gave $0.0436218\,m$, and the checker failed with deviations up to $3.3\times10^{-6}\,m$ from the Rust values. On the finer grids of 120, 240 and 480 points the reference gives $0.0436199531\,m$. The Rust values are $0.0436184212\,m$ (301 points) and $0.0436193591\,m$ (601 points), and their extrapolation from 301 and 601 points is $0.0436196717\,m$. The checker compares the 301-point value with the reference: the difference, $1.53\times10^{-6}\,m$, lies inside the tolerance $1.94\times10^{-6}\,m$, which is $10^{-6}\,m$ plus the measured grid uncertainty of the Rust run, $\lvert\Delta_{\text{SCF}}(601)-\Delta_{\text{SCF}}(301)\rvert=9.4\times10^{-7}\,m$ (erratum E4.12); for the 601-point value the difference is $5.9\times10^{-7}\,m$ (check \texttt{canonica\allowbreak{}l\_deltaS\allowbreak{}CF}, true, and the entries \texttt{canonica\allowbreak{}l\_excite\allowbreak{}d\_m1\_L3\_\allowbreak{}N1016\_la\allowbreak{}mm2\_T0} and \texttt{canonica\allowbreak{}l\_excite\allowbreak{}d\_m1\_L3\_\allowbreak{}N1016\_la\allowbreak{}mm2\_T0\_g\allowbreak{}601} of \texttt{comparis\allowbreak{}ons} in \texttt{artifact\allowbreak{}s/dirac1\allowbreak{}6complex\allowbreak{}/kohn-sh\allowbreak{}am/pytho\allowbreak{}n-check-\allowbreak{}report.j\allowbreak{}son}). Quoted to the digits on which all these values agree, $\Delta_{\text{SCF}}=0.04362\,m$. The eigenvalues of two deep sea levels of the same 301-point run are the subject of the one failed check of the cross-check (Section 13.15).

\begin{figure}[htbp]
\centering
\includegraphics[width=0.92\linewidth]{artifacts/dirac16complex/kohn-sham/figures/ks_gap_delta_scf.png}
\caption{Left: the Kohn--Sham gap (filled) and the Delta-SCF excitation energy (open) of the ground states of \texorpdfstring{$N=8$}{N=8}, 112 and 1016 against \texorpdfstring{$\hat\lambda/\hat\lambda_1$}{λ/λ\_1}. Middle: the same quantities against \texorpdfstring{$N$}{N}, one line for each coupling. Right: the orbital relaxation \texorpdfstring{$\Delta_{\text{SCF}}-\Delta_{KS}$}{Δ\_SCF-Δ\_KS} on a signed logarithmic scale. The annotation in the left panel refers to the smeared ensemble of Section 13.11 (file ks\_gap\_delta\_scf.png in artifacts/dirac16complex/kohn-sham/figures).}
\end{figure}

\subsection{Finite temperature}

\textbf{The runs.} The subcommand \texttt{\detokenize{thermo}} solves the Mermin--Kohn--Sham equations of Section 13.9 at $T/m\in\{0,0.1,0.3,1\}$ for $N=8$ and $N=112$ ($m=1$, $L=3$) in three series (\texttt{rust/the\allowbreak{}rmo/summ\allowbreak{}ary.json}, key \texttt{seriesRu\allowbreak{}le}): \texttt{\detokenize{lam0}} without interaction; \texttt{\detokenize{lamp1}} with the $T=0$ coupling $\hat\lambda_1$, run only at the temperatures where its first-order pseudo-potential stays below $1\,m$; and \texttt{\detokenize{lamh}} with a smaller ``hot'' coupling, $\hat\lambda=2.368\times10^{-5}$, chosen so that the pseudo-potential is $0.1\,m$ at $T=m$ and the whole series belongs to one Hamiltonian. The \texttt{\detokenize{lamp1}} points at $T=m$ were not run: there the first-order pseudo-potential would be $41.1\,m$ ($N=8$) and $40.4\,m$ ($N=112$), because thermal pairs raise the largest scalar density of the free state to about 4504 (key \texttt{skippedR\allowbreak{}uns}). At $T>0$ both particle and sea levels have fractional weights, $w=f$ and $w=-(1-f)$ (Section 13.7): the gas contains \textbf{thermal particle--antiparticle pairs}, which do not change $N$.

\textbf{Results without interaction} (\texttt{rust/the\allowbreak{}rmo/ther\allowbreak{}modynami\allowbreak{}cs.csv}, series 0, column \texttt{S\_entrop\allowbreak{}y}; energies in units of $m$, the entropy $S_{\text{ent}}$ dimensionless):

\begingroup
\small
\begin{longtable}{@{}>{\raggedright\arraybackslash}p{0.133\linewidth}>{\raggedright\arraybackslash}p{0.133\linewidth}>{\raggedright\arraybackslash}p{0.133\linewidth}>{\raggedright\arraybackslash}p{0.133\linewidth}>{\raggedright\arraybackslash}p{0.133\linewidth}>{\raggedright\arraybackslash}p{0.133\linewidth}@{}}
\toprule
\textbf{$N$} & \textbf{$T/m$} & \textbf{$\mu$} & \textbf{$E$} & \textbf{$F$} & \textbf{$S_{\text{ent}}$} \\
\midrule
\endfirsthead
\toprule
\textbf{$N$} & \textbf{$T/m$} & \textbf{$\mu$} & \textbf{$E$} & \textbf{$F$} & \textbf{$S_{\text{ent}}$} \\
\midrule
\endhead
8 & 0.1 & 0.130419 & 1.031867 & $-0.393979$ & 14.25845 \\
8 & 0.3 & 0.013287 & 108.7346 & $-30.51150$ & 464.1537 \\
8 & 1 & 0.000376 & 46811.74 & $-11447.70$ & 58259.43 \\
112 & 0.1 & 0.697334 & 69.50938 & 53.97164 & 155.3773 \\
112 & 0.3 & 0.316932 & 167.7124 & $-12.41988$ & 600.4411 \\
112 & 1 & 0.010141 & 46813.93 & $-11447.15$ & 58261.08 \\
\bottomrule
\end{longtable}
\endgroup

At $T=0$ the values are those of Section 13.11 ($F=E=0$ for $N=8$, $61.090723$ for $N=112$, $S_{\text{ent}}=0$). \textbf{Worked check:} $F=E-TS_{\text{ent}}$ for $N=8$ at $T=0.1$ gives $1.031867-0.1\times14.25845=-0.393978$, the tabulated $F$ up to rounding in the last digit (with the unrounded $E=1.0318666$ and $S_{\text{ent}}=14.2584539$ one gets $-0.3939788$, which rounds to the tabulated $-0.393979$). The heat capacity $C_V=dE/dT$ is $33.618$, $1754.07$ and $2.383\times10^5$ for $N=8$ and $157.752$, $1465.74$ and $2.383\times10^5$ for $N=112$ at $T/m=0.1$, $0.3$ and $1$ (column \texttt{\detokenize{C_V}}: a self-consistent central difference with step $0.05\,T$ for $T\le0.3\,m$, the exact fixed-spectrum formula of Exercise 12.9 at $T=m$). Without interaction the exact fixed-spectrum value, $33.500$ and $1745.07$ for $N=8$, agrees with the central difference within the 2 percent allowed by the check \texttt{cv\_finit\allowbreak{}e\_differ\allowbreak{}ence\_vs\_\allowbreak{}exact} (the difference is the truncation error of the finite step).

\textbf{What the numbers say.} At $T=0.1\,m$ the gas is a slightly heated Fermi system near the brane. At $T=m$ it is a \textbf{plasma of thermal pairs}: the energy window of the solver then contains about 75000 levels (74945 for $N=8$, column \texttt{\detokenize{states}}), $E$ and $S_{\text{ent}}$ are almost the same for $N=8$ and $N=112$ (the eight or 112 extra particles are negligible against the pairs), and $\mu$ is close to 0 because the plasma is almost charge-symmetric. The free energy decreases with $T$ and the entropy is positive in every run (checks \texttt{free\_ene\allowbreak{}rgy\_decr\allowbreak{}eases}, \texttt{entropy\_\allowbreak{}positive}, \texttt{heat\_cap\allowbreak{}acity\_po\allowbreak{}sitive} of \texttt{rust/the\allowbreak{}rmo/summ\allowbreak{}ary.json}, 102 of 102 true). The Kohn--Sham gap at finite $T$ is read with the threshold $f\ge\tfrac12$ (Section 12.15): for $N=112$ at $T=0.1\,m$ the chemical potential $0.697$ lies just below the band level $0.704$, so the HOMO is the band at $0.588$ and the gap is $0.116143$; at $T=0.3\,m$, $\mu=0.317$ and the gap becomes the distance $0.430734$ between the zero modes and the first band shell.

\textbf{With interaction.} The hot coupling changes $F$ at $T=m$ from $-11447.70$ to $-11447.48$ ($N=8$), with $\max|\lambda S_p|/m=0.104$; the $T=0$ coupling $\hat\lambda_1$ changes $F$ at $T=0.3\,m$ from $-30.5115$ to $-30.5037$ ($N=8$) and from $-12.4199$ to $-12.4103$ ($N=112$). The interaction stays a small correction at all temperatures computed.

\begin{figure}[htbp]
\centering
\includegraphics[width=0.92\linewidth]{artifacts/dirac16complex/kohn-sham/figures/thermodynamics.png}
\caption{Mermin--Kohn--Sham thermodynamics for \texorpdfstring{$m=1$}{m=1}, \texorpdfstring{$L=3$}{L=3}: energy, free energy, entropy, heat capacity and Kohn--Sham gap against \texorpdfstring{$T$}{T} for \texorpdfstring{$N=8$}{N=8} and \texorpdfstring{$N=112$}{N=112} (filled: no interaction; open: with interaction), and the occupations of the \texorpdfstring{$N=8$}{N=8} levels at three temperatures, where the sea levels (\texorpdfstring{$\varepsilon<0$}{ε<0}) show their holes, the thermal antiparticles (file thermodynamics.png in artifacts/dirac16complex/kohn-sham/figures).}
\end{figure}

\subsection{The energy--momentum tensor and the Einstein source}

\textbf{The tensor of one orbital.} Insert the stationary orbital into the energy--momentum tensor of Chapter 7, $T_{\mu\nu}=-\tfrac14\bigl[\bar\Psi\gamma_\mu D_\nu\Psi+\bar\Psi\gamma_\nu D_\mu\Psi-(D_\mu\bar\Psi)\gamma_\nu\Psi-(D_\nu\bar\Psi)\gamma_\mu\Psi\bigr]+g_{\mu\nu}\mathcal L_s$, and use the expectation-value rule. The energy density is the simplest component. With $\gamma_4=-\gamma^4$, $D_4=\partial_4$ (the static connection has $\Omega_4=0$), $\partial_4\Psi=-i\varepsilon\Psi$ and $\partial_4\bar\Psi=+i\varepsilon\bar\Psi$, the bracket gives $-\tfrac14\bigl[2i\varepsilon\bar\Psi\gamma^4\Psi+2i\varepsilon\bar\Psi\gamma^4\Psi\bigr]=-i\varepsilon\Psi^\dagger C\gamma^4\Psi=\varepsilon\Psi^\dagger B\Psi$ (because $C\gamma^4=iB$), whose expectation is $\varepsilon\,u^\dagger B^2u=\varepsilon\,n$. With $g_{44}=-1$ the last term contributes $-\mathcal L_s$. On shell $\mathcal L_s=SU'(S)-U(S)$ (Chapter 7), and for $U=\tfrac\lambda2S^2$ this is $\tfrac\lambda2S^2$, whose expectation in the Kohn--Sham state is the interaction energy density of Section 13.8, $L_s=e_H+e_x=\tfrac\lambda2S_p^2-\tfrac\lambda{32}(n_p^2+S_p^2)$. Summed over the levels with the weights $w_n$, this gives $\rho=\sum_nw_n\varepsilon_nn_n-L_s$.

\textbf{The pressure along $y$.} For $\mu=\nu=y$ we have $g_{yy}=1$, $\gamma_y=\gamma^y=\gamma^0$ and $D_y=\partial_y$, because the spin connection of the static field has no $y$ component ($\Omega_0=\Omega_4=0$, Section 4.12). So

\[
T_{yy}=-\tfrac12\bigl[\bar\Psi\gamma^0\partial_y\Psi-(\partial_y\bar\Psi)\gamma^0\Psi\bigr]+\mathcal L_s .
\]

Insert $\Psi=e^{-i\varepsilon x_4}e^{i\mathbf k\cdot\mathbf x}e^{-3Hy}\chi$. The phases cancel between $\bar\Psi$ and $\Psi$, and $\partial_y(e^{-3Hy}\chi)=e^{-3Hy}(\chi'-3H\chi)$; the two terms with $-3H\chi$ give the same product $-3H\,e^{-6Hy}\chi^\dagger C\gamma^0\chi$ in both halves of the bracket, so they cancel in the difference. The expectation-value rule of Section 13.7 replaces $\Psi^\dagger$ by $u^\dagger B$, so the matrix between the two spinors becomes $BC\gamma^0=-i\gamma^4\gamma^0=i\gamma^0\gamma^4=iA_4$, which is $ij\sigma_x$ in the block (Section 13.4). With the factor $e^{-6Hy}/\ell^3$ of the proper densities, one orbital gives

\[
T^y{}_y-L_s=-\tfrac12\,ij\,\frac{e^{-6Hy}}{\ell^3}\bigl[\chi^\dagger\sigma_x\chi'-\chi'^\dagger\sigma_x\chi\bigr].
\]

Now use the block equation $\chi'=N\chi$ of Section 13.5 and its conjugate transpose $\chi'^\dagger=\chi^\dagger N^\dagger$: the bracket is $\chi^\dagger(\sigma_xN-N^\dagger\sigma_x)\chi$. With $N=M_{\text{eff}}\sigma_z-\xi k\sigma_y+ijE\sigma_x$, $N^\dagger=M_{\text{eff}}\sigma_z-\xi k\sigma_y-ijE\sigma_x$ and the products $\sigma_x\sigma_z=-i\sigma_y$, $\sigma_z\sigma_x=i\sigma_y$, $\sigma_x\sigma_y=i\sigma_z$, $\sigma_y\sigma_x=-i\sigma_z$ of Chapter 2,

\[
\sigma_xN-N^\dagger\sigma_x=-2iM_{\text{eff}}\sigma_y-2i\xi k\sigma_z+2ijE .
\]

Multiplying by $-\tfrac12ij$ and using $j^2=1$ and $i^2=-1$ turns the three terms into $-jM_{\text{eff}}\sigma_y$, $-j\xi k\sigma_z$ and $E$. With the densities of Section 13.7 ($n$ from $\chi^\dagger\chi$, $s$ from $j\chi^\dagger\sigma_y\chi$, $t$ from $j\chi^\dagger\sigma_z\chi$, each times $e^{-6Hy}/\ell^3$), one orbital therefore has

\[
T^y{}_y-L_s=(\varepsilon-v)\,n-M_{\text{eff}}\,s-\xi k\,t .
\]

Now sum over the levels with the weights $w_n$. At a given $y$ the potentials $M_{\text{eff}}-m=\tfrac{15}{16}\lambda S_p$ and $v=-\tfrac1{16}\lambda n_p$ of Section 13.8 are the same for every orbital, and $\sum_nw_ns_n=S_p$, $\sum_nw_nn_n=n_p$, so

\[
\sum_nw_n\bigl[(M_{\text{eff}}-m)\,s_n+v\,n_n\bigr]=\tfrac{15}{16}\lambda S_p^2-\tfrac1{16}\lambda n_p^2=\lambda S_p^2-\tfrac\lambda{16}\bigl(n_p^2+S_p^2\bigr)=2L_s .
\]

Writing $M_{\text{eff}}=m+(M_{\text{eff}}-m)$ in the one-orbital result and summing gives $p_y=T^y{}_y=\sum_nw_n\bigl[\varepsilon_nn_n-m\,s_n-\xi k_n\,t_n\bigr]-2L_s+L_s$, the formula below.

\textbf{The 3-space pressures.} Take $\mu=\nu=x_1$, the direction of $\mathbf k$. Here $\gamma_1=g_{11}\gamma^{x_1}=e^{Hy+a_{4,0}}\gamma^1$, and $D_1=\partial_1+\Omega_1$ with $\Omega_1$ a multiple of $\gamma^0\gamma^1$ (Section 4.12). The bracket of $T_{11}$ is $-\tfrac12\bigl[\bar\Psi\gamma_1\partial_1\Psi-(\partial_1\bar\Psi)\gamma_1\Psi\bigr]-\tfrac12\bar\Psi\{\gamma_1,\Omega_1\}\Psi$, and the connection term vanishes because $\{\gamma^1,\gamma^0\gamma^1\}=\gamma^1\gamma^0\gamma^1+\gamma^0\gamma^1\gamma^1=-\gamma^0+\gamma^0=0$. With $\partial_1\Psi=ik\Psi$ and $\partial_1\bar\Psi=-ik\bar\Psi$ the derivative terms give $-ik\,e^{Hy+a_{4,0}}\bar\Psi\gamma^1\Psi$, and raising the index with $g^{11}=e^{-2Hy-2a_{4,0}}$ gives

\[
T^1{}_1=-i\xi k\,\bar\Psi\gamma^1\Psi+\mathcal L_s .
\]

The matrix of the expectation-value rule is $BC\gamma^1=-i\gamma^4\gamma^1\to-i(-j\sigma_z)=ij\sigma_z$ (Section 13.4), so one orbital gives $T^1{}_1-L_s=-i\xi k\cdot ij\,\frac{e^{-6Hy}}{\ell^3}\chi^\dagger\sigma_z\chi=\xi k\,t$. For $\mu=\nu=x_2$ and $x_3$ nothing depends on the coordinate, and the connection term vanishes in the same way ($\{\gamma^2,\gamma^0\gamma^2\}=0$), so $T^2{}_2=T^3{}_3=L_s$ for this orbital. The sum $T^1{}_1+T^2{}_2+T^3{}_3$ does not change under rotations of 3-space (Section 13.3), so an orbital whose momentum points in another lattice direction contributes the same amount $\xi k\,t+3L_s$ to it. The average of the three 3-space pressures is therefore $p_3=\tfrac13\sum_nw_n\xi k_nt_n+L_s$.

\textbf{The extra times.} In the good sector nothing depends on $x_5,x_6,x_7$ (Section 13.3), and $\Omega_5$ is a multiple of $\gamma^0\gamma^5$, with $\{\gamma^5,\gamma^0\gamma^5\}=\gamma^5\gamma^0\gamma^5+\gamma^0(\gamma^5)^2=\gamma^0-\gamma^0=0$ (since $(\gamma^5)^2=-1$); the same holds for $x_6$ and $x_7$. So $T^5{}_5=T^6{}_6=T^7{}_7=L_s$, which is $p_t$.

Collecting, the results for the Kohn--Sham state, a sum over all levels $n$ with the weights $w_n$ and the proper densities $n_n,s_n,t_n$ of Section 13.7, are

\[
\begin{aligned}
\rho(y)&=\sum_nw_n\,\varepsilon_n\,n_n(y)-L_s(y),\qquad
p_y(y)=\sum_nw_n\bigl[\varepsilon_nn_n-m\,s_n-\xi k_n\,t_n\bigr]-L_s(y),\\
p_3(y)&=\tfrac13\sum_nw_n\,\xi k_n\,t_n+L_s(y),\qquad p_t(y)=L_s(y),
\end{aligned}
\]

where $p_y=T^y{}_y$ is the pressure along the hidden direction, $p_3$ the average of the three 3-space pressures and $p_t$ the pressure along each extra time (\texttt{KS\_emt\_r\allowbreak{}ho}, \texttt{KS\_emt\_p\allowbreak{}y}, \texttt{KS\_emt\_p\allowbreak{}1}, \texttt{KS\_emt\_p\allowbreak{}2p3}, \texttt{KS\_emt\_p\allowbreak{}t}; \texttt{kohn-sha\allowbreak{}m-theory\allowbreak{}.json}, keys \texttt{\detokenize{emt}}, \texttt{\detokenize{ksState}}). The off-diagonal components either are proportional to the $y$-current, which vanishes for eigenstates, or cancel over every closed shell $\{\mathbf k,-\mathbf k\}$ (\texttt{KS\_emt\_T\allowbreak{}y4Propor\allowbreak{}tionalTo\allowbreak{}Current}, \texttt{KS\_emt\_T\allowbreak{}y1Propor\allowbreak{}tionalTo\allowbreak{}Current}, \texttt{KS\_emt\_T\allowbreak{}41cancel\allowbreak{}sOverShe\allowbreak{}ll}, \texttt{KS\_emt\_o\allowbreak{}ffBlockC\allowbreak{}omponent\allowbreak{}sVanishP\allowbreak{}erOrbita\allowbreak{}l}). Without interaction $L_s=0$, so the extra times carry no pressure.

\textbf{Two identities that the solver checks.} (i) \textbf{Energy.} Each orbital has $\int n_n\,dV=\int e^{-6Hy}\chi_n^\dagger\chi_n\ell^{-3}\,e^{6Hy}\ell^3\,dy=1$, so $\int\rho\,dV=\sum_nw_n\varepsilon_n-E_H-E_x=E$, the total energy of Section 13.9 (check \texttt{energy\_f\allowbreak{}rom\_rho} of every run; for example $61.0907233$ both ways for $N=112$ at $\lambda=0$, \texttt{rust/emt\allowbreak{}/emt-sum\allowbreak{}mary.csv}). (ii) \textbf{Conservation.} For a static tensor that depends on $y$ only, $\nabla_\mu T^\mu{}_y=\partial_yT^y{}_y+\Gamma^\mu{}_{\mu y}T^y{}_y-\Gamma^\nu{}_{\mu y}T^\mu{}_\nu$ (summed over $\mu$ and $\nu$; we do not use the letter $\lambda$ as an index here, because it is the coupling). The Christoffel symbols of the warped metric are $\Gamma^\mu{}_{y\mu}=H$ on the six warped directions and $\Gamma^y{}_{\mu\mu}=-Hg_{\mu\mu}$ (\texttt{KS\_geome\allowbreak{}try\_chri\allowbreak{}stoffelC\allowbreak{}losedFor\allowbreak{}ms}), so $\Gamma^\mu{}_{\mu y}=6H$ and the last term is $H\bigl(3p_3+3p_t\bigr)$:

\[
p_y'+6H\,p_y-3H\,p_3-3H\,p_t=0 .
\]

The solver evaluates the left side for every self-consistent state (check \texttt{emt\_cons\allowbreak{}ervation}); the normalized residual is at most $2.1\times10^{-5}$ (column \texttt{conserva\allowbreak{}tion\_res\allowbreak{}idual} of \texttt{emt-summ\allowbreak{}ary.csv}).

\textbf{Averages and equations of state.} Proper-volume averages are $\langle X\rangle=\int_{-L}^0X\,e^{6Hy}dy\big/\int_{-L}^0e^{6Hy}dy$, and $w_y=\langle p_y\rangle/\langle\rho\rangle$, $w_3=\langle p_3\rangle/\langle\rho\rangle$, $w_t=\langle p_t\rangle/\langle\rho\rangle$. The comparison with the required source asks which $\kappa$ would satisfy $\rho_{\text{req}}=-21H^2/\kappa=\langle\rho\rangle$, that is $\kappa_{\text{needed}}=-21H^2/\langle\rho\rangle$ (\texttt{rust/emt\allowbreak{}/emt-sum\allowbreak{}mary.csv}; $m=1$, $L=3$, $T=0$):

\begingroup
\small
\begin{longtable}{@{}>{\raggedright\arraybackslash}p{0.133\linewidth}>{\raggedright\arraybackslash}p{0.133\linewidth}>{\raggedright\arraybackslash}p{0.133\linewidth}>{\raggedright\arraybackslash}p{0.133\linewidth}>{\raggedright\arraybackslash}p{0.133\linewidth}>{\raggedright\arraybackslash}p{0.133\linewidth}@{}}
\toprule
\textbf{$N$} & \textbf{$\hat\lambda/\hat\lambda_1$} & \textbf{$\langle\rho\rangle$} & \textbf{$w_y$} & \textbf{$w_3$} & \textbf{$\kappa_{\text{needed}}$} \\
\midrule
\endfirsthead
\toprule
\textbf{$N$} & \textbf{$\hat\lambda/\hat\lambda_1$} & \textbf{$\langle\rho\rangle$} & \textbf{$w_y$} & \textbf{$w_3$} & \textbf{$\kappa_{\text{needed}}$} \\
\midrule
\endhead
8 & 0 & 0 & --- & --- & --- \\
8 & +1 & $-3.73\times10^{-7}$ & 9.01 & 0.991 & $+5.63\times10^{7}$ \\
8 & $-1$ & $+3.73\times10^{-7}$ & 9.01 & 0.991 & $-5.63\times10^{7}$ \\
112 & 0 & 0.0230891 & 0.449 & 0.297 & $-909.5$ \\
112 & +1 & 0.0230976 & 0.452 & 0.297 & $-909.2$ \\
112 & $-1$ & 0.0230807 & 0.447 & 0.297 & $-909.8$ \\
1016 & 0 & 0.425999 & 0.336 & 0.290 & $-49.30$ \\
1016 & +1 & 0.426008 & 0.336 & 0.290 & $-49.29$ \\
1016 & $-1$ & 0.425989 & 0.336 & 0.290 & $-49.30$ \\
\bottomrule
\end{longtable}
\endgroup

The one finite-temperature state of the \texttt{\detokenize{emt}} set, $N=112$ at $+\hat\lambda_1$ and $T=0.3\,m$, has $\langle\rho\rangle=0.0633668$, $w_y=0.322$, $w_3=0.274$ and $\kappa_{\text{needed}}=-331.4$. (For $N=8$ at $\lambda=0$, $\langle\rho\rangle=0$ and the ratios are undefined.)

\textbf{Reading the table.} Every state with band or bulk levels has a \textbf{positive} average energy density and positive pressures along the hidden direction and in 3-space ($w_y$ between 0.32 and 0.45, $w_3$ close to $0.3$): it is a gas of fast, nearly massless fermions, as the brane band suggests. The field needs the opposite sign, $\rho_{\text{req}}<0$. A positive $\langle\rho\rangle$ can match $\rho_{\text{req}}$ only with a \textbf{negative} gravitational coupling, $\kappa\approx-909$ for $N=112$ and $\kappa\approx-49$ for $N=1016$; with the physical sign $\kappa>0$ these states cannot be the source. For $N=8$ the energy density vanishes without interaction (the zero modes have $\varepsilon=0$), and with $\hat\lambda_1>0$ it is tiny and negative, $-3.73\times10^{-7}$, which would need $\kappa=5.6\times10^7$; but then the mass condition below gives the opposite sign of $\kappa$. There is a second, structural, obstacle: the required source is the same at every $y$, whereas the computed $\rho(y)$ changes by several orders of magnitude between the tip and the brane (figure below), so these states cannot match it point by point, and the comparison of averages is the most favourable one. (The specification had expected a positive Kohn--Sham energy density in every run; the measurement shows that this fails for $N=8$, \texttt{positive\allowbreak{}RhoExpec\allowbreak{}tation} in \texttt{rust/emt\allowbreak{}/summary\allowbreak{}.json}.)

\textbf{The sourcing conditions of erratum E4.1.} For the homogeneous state $\Psi=u(x_4)$ of Stage 2 the exact conditions under which dirac16complex sources the static field are $mS=-36H^2/\kappa$ and $\lambda S^2=30H^2/\kappa$ (Chapter 9); together they require $\lambda S/m=-5/6$ and $mS<0$. The solver evaluates them on the average scalar density $\langle S_p\rangle$ as a diagnostic (\texttt{E41\_sour\allowbreak{}cingCond\allowbreak{}itions} in each \texttt{\detokenize{run.json}}). The ground states are dominated by the brane band, whose scalar charge is negative (Section 13.9), so $\langle S_p\rangle<0$: $-0.00788$ for $N=112$ and $-0.0872$ for $N=1016$ at $\lambda=0$. The mass condition alone then gives a positive $\kappa$ (4568 and 412.7), but the coupling condition needs $\lambda\langle S_p\rangle/m=-5/6$, whereas the computed states have $|\lambda\langle S_p\rangle/m|\le10^{-4}$ (the runs of \texttt{rust/emt\allowbreak{}/} use $\hat\lambda\in\{0,\pm\hat\lambda_1\}$). At first order the coupling would have to be $\hat\lambda\approx-\tfrac56m^7/\langle S_p\rangle=105.7$ ($N=112$) or $9.55$ ($N=1016$), about $10^3$ to $10^4$ times $\hat\lambda_2$, far outside the range in which this local scheme was set up. \textbf{In no computed run are the conditions met} (column \texttt{sourcing\allowbreak{}\_conditi\allowbreak{}ons\_met} of \texttt{emt-summ\allowbreak{}ary.csv} is 0 in every row). This is a statement about the computed states and couplings, not a proof that no Kohn--Sham state could ever source the field.

\begin{figure}[htbp]
\centering
\includegraphics[width=0.92\linewidth]{artifacts/dirac16complex/kohn-sham/figures/emt_profiles.png}
\caption{Energy--momentum tensor of the ground states at \texorpdfstring{$+\hat\lambda_1$}{+λ\_1} for \texorpdfstring{$N=112$}{N=112} and \texorpdfstring{$N=1016$}{N=1016}: the proper densities \texorpdfstring{$\rho$}{ρ}, \texorpdfstring{$p_y$}{p\_y}, \texorpdfstring{$p_3$}{p\_3}, \texorpdfstring{$p_t$}{p\_t} against \texorpdfstring{$y$}{y} (top), the local ratios \texorpdfstring{$w_y(y)$}{w\_y(y)} and \texorpdfstring{$w_3(y)$}{w\_3(y)} with the required value \texorpdfstring{$w_{\text{req}}=-5/7$}{w\_req=-5/7} (bottom left), and \texorpdfstring{$\rho\,e^{6Hy}$}{ρe\textasciicircum{}6Hy}, the energy per unit \texorpdfstring{$y$}{y} (bottom right) (file emt\_profiles.png in artifacts/dirac16complex/kohn-sham/figures).}
\end{figure}

\begin{figure}[htbp]
\centering
\includegraphics[width=0.92\linewidth]{artifacts/dirac16complex/kohn-sham/figures/einstein_source.png}
\caption{The Kohn--Sham states (the runs of \texttt{rust/emt\allowbreak{}/}) against the source the static field requires, from left to right: the proper-volume average \texorpdfstring{$\langle\rho\rangle$}{⟨ρ⟩} with \texorpdfstring{$\rho_{\text{req}}$}{ρ\_req} for \texorpdfstring{$\kappa=1$}{κ=1}; the ratios \texorpdfstring{$w_y$}{w\_y} and \texorpdfstring{$w_3$}{w\_3} with \texorpdfstring{$w_{\text{req}}=-5/7$}{w\_req=-5/7}; the \texorpdfstring{$\kappa$}{κ} obtained from \texorpdfstring{$\rho_{\text{req}}=\langle\rho\rangle$}{ρ\_req=⟨ρ⟩} and the \texorpdfstring{$\kappa$}{κ} obtained from the mass condition \texorpdfstring{$m\langle S_p\rangle=-36H^2/\kappa$}{m⟨S\_p⟩=-36H\textasciicircum{}2/κ}; and \texorpdfstring{$\lambda\langle S_p\rangle/m$}{λ⟨S\_p⟩/m} against the value \texorpdfstring{$-5/6$}{-5/6} that erratum E4.1 requires, on a signed logarithmic axis (file einstein\_source.png in artifacts/dirac16complex/kohn-sham/figures).}
\end{figure}

\textbf{The mirror sector: a Kohn--Sham map $(m,\lambda)\to(-m,\lambda)$.} At the level of the classical field, Chapter 15 proves that the chirality map $\Psi\to\gamma^8\Psi$ turns a solution with $(m,\lambda)$ into one with $(-m,-\lambda)$ and reverses the energy--momentum tensor (Theorem T1; rows L16 and L37 of the ledger in Chapter 0). For the quantum field the image $\gamma^8\Psi$ then carries the reversed Krein metric $-B$. In the Kohn--Sham setting of this chapter the universe of mass $-m$ is instead quantized on its own with the positive metric $B$, that is, with the standard expectation-value rule of Section 13.7. With that rule the block form of $\gamma^8$ (which maps the block $(j,s_2,s_3)$ onto $(-j,s_2,s_3)$, Section 13.4, and acts there as $\sigma_y$ up to a phase) keeps $\lambda$: it maps the problem $(m,\lambda)$ onto $(-m,\lambda)$ with the same energies, number densities and pressures (only the scalar density changes sign), so the energies of the two members add. This is the mirror-type pair of Theorem T2 (row L38), not the cancelling T1 pair (\texttt{artifact\allowbreak{}s/dirac1\allowbreak{}6complex\allowbreak{}/pair-cr\allowbreak{}eation/p\allowbreak{}airing-t\allowbreak{}heory.js\allowbreak{}on}, keys \texttt{T3.theor\allowbreak{}emStanda\allowbreak{}rdRule}, \texttt{T3.theor\allowbreak{}emImageR\allowbreak{}ule} and \texttt{T3.pairT\allowbreak{}otalsKS}). At the level of the $2\times2$ blocks this map has a simple exact form, which we derive here (it is stated in the header of \texttt{src/emt.\allowbreak{}rs}). With $\chi=(a,ib)$ we have $\sigma_y\chi=(b,ia)$: in the real system of Section 13.5 the map exchanges the two components and the block type, $\tilde a=b$, $\tilde b=a$, $\tilde j=-j$. Then

\[
\begin{aligned}
\tilde a'&=b'=(jE-\xi k)\,a-M_{\text{eff}}\,b=(jE-\xi k)\,\tilde b-M_{\text{eff}}\,\tilde a\\
&=(-M_{\text{eff}})\,\tilde a-(\tilde jE+\xi k)\,\tilde b,\\
\tilde b'&=a'=M_{\text{eff}}\,a-(jE+\xi k)\,b=M_{\text{eff}}\,\tilde b-(jE+\xi k)\,\tilde a\\
&=(\tilde jE-\xi k)\,\tilde a-(-M_{\text{eff}})\,\tilde b ,
\end{aligned}
\]

which is the real system with the mass $-M_{\text{eff}}$, the same $\varepsilon$, $k$ and $v$. So every level of the problem with mass function $M_{\text{eff}}$ is also a level, with the same energy, of the problem with $-M_{\text{eff}}$, provided the boundary conditions are transformed as well: the tip condition $b(-L)=0$ becomes $\tilde a(-L)=0$ (the other bag, $\theta=\pi$), and the brane parities $+$ and $-$ are exchanged. The densities transform as $n\to n$ (since $a^2+b^2$ is symmetric) and $s\to-s$ (since $2jab\to2\tilde j\tilde a\tilde b=-2jab$). The map is consistent with self-consistency: $M_{\text{eff}}=m+\tfrac{15}{16}\lambda S_p$ becomes $-M_{\text{eff}}=-m+\tfrac{15}{16}\lambda(-S_p)$ \textbf{with the same $\lambda$}, and $v=-\lambda n_p/16$ is unchanged. Hence the Kohn--Sham problem $(m,\lambda)$ with the tip bag $b(-L)=0$ is mapped exactly onto the problem $(-m,\lambda)$ with the bag $a(-L)=0$: the same energies, number densities and pressures, the opposite scalar density, and therefore the same $mS$, the same $\lambda S^2$ and the same verdict on the E4.1 conditions. This is the Kohn--Sham shadow of the pairing theorems: Chapter 15 proves it in full as Theorem T3 (Section 15.7), including how the normal ordering and the particle--sea classification are mapped, and reports its numerical test, which is partial (Section 15.10).

\subsection{Convergence, reproducibility and the status of the cross-check}

\textbf{Cutoff $L$.} The tip cutoff hardly matters once $L\ge3$: the Kohn--Sham gap of $N=8$ is $0.424302$, $0.430734$ and $0.430746\,m$ for $L=2$, 3 and 4, and the ground-state energy of $N=112$ is $60.725616$, $61.090723$ and $61.091032\,m$ (\texttt{rust/scf\allowbreak{}/summary\allowbreak{}.json}). Between $L=3$ and $L=4$ the energy changes by 5 parts in a million. The reason is not the first-order profile of Section 13.6: its slope $c=2/(1+e^{-L})$, whose integrand decays only like $e^{-(2M-H)|y|}$, changes by 3 percent between $L=3$ and $L=4$ ($1.905148$ and $1.964028$). The zero modes have $\varepsilon=0$ for every $L$. The band orbitals at $k\ne0$ are pushed toward the brane by the term $\xi(y)k=k\,e^{H|y|}$ (for $a_{4,0}=0$), which grows toward the tip, so they decay much faster than the first-order profile: the lowest band level (the $N=8$ gap) changes by $6.4\times10^{-3}\,m$ from $L=2$ to $L=3$ but only by $1.2\times10^{-5}\,m$ from $L=3$ to $L=4$, and the band level at $k=0.433\,m$ (the $\varepsilon_{\text{HOMO}}$ of $N=112$) by only $3\times10^{-8}\,m$ (key \texttt{\detokenize{epsLumo}} of the runs \texttt{m1\_L2\_N8\allowbreak{}\_lam0\_T0}, \texttt{m1\_L3\_N8\allowbreak{}\_lam0\_T0} and \texttt{m1\_L4\_N8\allowbreak{}\_lam0\_T0}, and key \texttt{\detokenize{epsHomo}} of the corresponding \texttt{\detokenize{N112}} runs, in \texttt{rust/scf\allowbreak{}/}).

\textbf{Grid.} The run on 601 instead of 301 points ($N=112$, $\hat\lambda_1$) gives the same ground-state energy $61.1132120\,m$ to all printed digits (check \texttt{grid\_ref\allowbreak{}inement\_\allowbreak{}energy}), and the hardest Delta-SCF changes from $0.0436184$ to $0.0436194\,m$ (Section 13.12).

\textbf{Torus spacing.} Halving $\Delta k$ to $0.125\,m$ at the same density means $N=8\cdot112=896$ (run \texttt{m1\_L3\_N8\allowbreak{}96\_lamp1\allowbreak{}\_T0\_dk0p\allowbreak{}125}). That value of $N$ is an \textbf{open} shell: 80 of the 192 states of the band shell at $k=0.468\,m$ are occupied, equally, so its Kohn--Sham gap is 0 by definition. Its energy per particle, $521.396546/896=0.581916\,m$, differs from the $N=112$ value $61.113212/112=0.545654\,m$ by 6.6 percent (erratum E4.11): the finite shell spacing of the torus still matters at this level of detail.

\textbf{The constant $a_{4,0}$.} The run with $a_{4,0}=0.5$ and its partner with $a_{4,0}=0$, $\Delta k\,e^{-0.5}=0.151633\,m$ and $\hat\lambda\,e^{1.5}$ (erratum E4.10; Exercise 13.12 explains the factors) have the same energy, $39.0088426\,m$, and the same levels (check \texttt{a4\_resca\allowbreak{}ling\_pai\allowbreak{}r\_exact}).

\textbf{Reproducibility.} Section 13.10 listed the determinism and refined-tolerance results. The canonical tree was produced with three release builds of the solver, which differ only in parts that do not change the other outputs; the solver's \texttt{README.m\allowbreak{}d} records their provenance (section ``Provenance of the canonical tree'').

\textbf{The independent cross-check (status).} The Stage-4 plan requires every number to be reproduced by an independent Python reference solver, which uses a matrix method on staggered grids instead of shooting and extrapolates in the grid spacing, and checked for the identities of this chapter (Hellmann--Feynman, current conservation, particle number, self-consistency) by \texttt{scripts/\allowbreak{}check\_di\allowbreak{}rac16com\allowbreak{}plex\_koh\allowbreak{}n\_sham.p\allowbreak{}y}. The final report of the checker, \texttt{artifact\allowbreak{}s/dirac1\allowbreak{}6complex\allowbreak{}/kohn-sh\allowbreak{}am/pytho\allowbreak{}n-check-\allowbreak{}report.j\allowbreak{}son}, compares the 56 reference runs under \texttt{artifact\allowbreak{}s/dirac1\allowbreak{}6complex\allowbreak{}/kohn-sh\allowbreak{}am/refer\allowbreak{}ence/} with 81 Rust runs (every Rust run has its reference partner; a refinement run such as \texttt{\detokenize{_g601}} is compared with the reference run of its base label) and has 63 checks, of which \textbf{62 are true}. The true checks include the internal identities of both solvers, the Hellmann--Feynman stationarity, the couplings, the ground-state energies, chemical potentials and gaps, the matching of every level with its branch and occupation, the particle--hole lists, the Delta-SCF energies (worst ratio of deviation to tolerance 0.79, for the smeared ensemble of Section 13.12), the profiles, the energy--momentum averages, the thermodynamics, the E4.1 sourcing record, the reproduction of one Rust run (\texttt{m1\_L2\_N8\allowbreak{}\_lamp1\_T\allowbreak{}0}) by the reference solver at exactly the Rust coupling, and the checker's own repetition of the two tests of Section 13.10 (check \texttt{rust\_rep\allowbreak{}eat\_byte\allowbreak{}\_identit\allowbreak{}y}, 335 files: the 330 output files and the five \texttt{summary.\allowbreak{}json}; check \texttt{rust\_ref\allowbreak{}ined\_con\allowbreak{}vergence}, 81 runs). \textbf{One check failed}, \texttt{canonica\allowbreak{}l\_eigenv\allowbreak{}alues}. It compares 58412 eigenvalues of the two solvers with the tolerance $10^{-6}\max(1,\lvert\varepsilon\rvert/m)\,m$ plus a small grid term, and its worst deviation is $2.195\times10^{-6}\,m$ against the tolerance $1.054\times10^{-6}\,m$ (ratio 2.08), at the level $\varepsilon=-1.05374908\,m$ of the 301-point run \texttt{m1\_L3\_N1\allowbreak{}016\_lamm\allowbreak{}2\_T0}, the smeared ensemble of Section 13.11. All its failures belong to this run: two deep levels of the Dirac sea at $k=0$, one of each parity, each in both block types (their free values are $-1.447972\,m$ and $-1.292293\,m$), counted once in the ground-state and once in the excited-state record of the run. In the 601-point run of the same ensemble (\texttt{m1\_L3\_N1\allowbreak{}016\_lamm\allowbreak{}2\_T0\_g60\allowbreak{}1}) all 79 compared levels agree with the reference to $1.83\times10^{-7}\,m$. That the failure is the grid error of the 301-point run for these two levels is a hypothesis to be tested, not a result (Section 18.7). One group of comparisons did not run: the self-tests of the reference solver are not in its committed summary (the entry \texttt{referenc\allowbreak{}e\_self\_t\allowbreak{}ests} of \texttt{comparis\allowbreak{}ons} has the status "not run"). The Stage-4 gate \texttt{scripts/\allowbreak{}verify\_s\allowbreak{}tage4\_ko\allowbreak{}hn\_sham.\allowbreak{}ps1} (or \texttt{\detokenize{.sh}}), which reruns the programs of the stage in a fresh copy of the repository (the reference solver in a reduced form) and requires the canonical Rust outputs to be reproduced byte for byte, also requires a fresh cross-check without a failed check and the PDFs of the two Stage-4 documents, which are not finished; it has not been run. The numbers of this chapter are therefore those of the Rust solver with its own checks, confirmed by the independent solver in 62 of the 63 checks; the eigenvalues of the two deep sea levels of the 301-point smeared run are confirmed only to about $2.2\times10^{-6}\,m$.

\subsection{What we proved and what we assumed}

\textbf{Proved} (derived in this chapter step by step, and verified by the named checks of the two exact reports): the warped form of the static primordial field, its curvature $R=-42H^2$ and $G^\mu{}_\nu=\mathrm{diag}(15,15,15,15,21,15,15,15)H^2$, the required source $\rho_{\text{req}}=-21H^2/\kappa$, $p_{\text{req}}=+15H^2/\kappa$, the extrinsic curvature and the Israel stress of the Z2 brane in the stated sign convention; the removal of the spin connection by the factor $W^{-3}$ and the reduced equation; the splitting of the 16-component equation into eight $2\times2$ blocks with the explicit block forms of $A_0,A_1,A_4,B,C,BC$ and the two block types; the Hamiltonian form, the relation $h_{-1}-v=-(h_{+1}-v)$, the conservation of the $y$-current, the boundary conditions and the self-adjointness they give; the exact $k=0$ spectrum and the first-order splitting $\pm ck$ of the brane zero modes; the Hartree--Fock energy $\tfrac\lambda2[S^2-\mathrm{Tr}(BC\rho BC\rho)]$, the filled-shell ratio $1/8$ and the exact closed form $e_x=-\tfrac\lambda{32}(n^2+S^2)$ of the uniform gas at every temperature; the Kohn--Sham equations and Fermi--Dirac occupations from Mermin's functional, the double-counting formula and the Hellmann--Feynman relations; the energy density and the pressures $p_y$, $p_3$, $p_t$ of an orbital and of the Kohn--Sham state, the identity $\int\rho\,dV=E$ and the form $p_y'+6Hp_y-3Hp_3-3Hp_t=0$ of the conservation law in this metric; the exact map of the Kohn--Sham problem $(m,\lambda)$ onto $(-m,\lambda)$ with exchanged boundary conditions (a mirror-type pair whose energies add, not the cancelling T1 pair); and (Exercise 13.9) the antisymmetry $E_0(-\lambda)=-E_0(\lambda)$ of the $N=8$ state.

\textbf{Computed} (by the Rust solver, with its own 397 checks, a byte-identical repeat and a refined-tolerance run, and compared with the independent reference solver in 63 checks, of which 62 are true; the failed one, on two deep sea levels of the 301-point smeared run, is open, and the Stage-4 gate has not been run; Section 13.15): the self-consistent ground states, the Kohn--Sham gaps, particle--hole lists and Delta-SCF energies, the finite-temperature thermodynamics, the energy--momentum tensor and its averages, the fact that the Kohn--Sham states satisfy the conservation law to a normalized residual of at most $2.1\times10^{-5}$ (check \texttt{emt\_cons\allowbreak{}ervation}), and the comparison with the required source: positive $\langle\rho\rangle$ in every state with band or bulk levels (so a negative $\kappa$ would be needed), and the conditions of erratum E4.1 met in no run.

\textbf{Assumed or chosen:} the static member $a_4=\text{const}$ of the primordial family; the good sector without extra-time dependence; the Z2 gluing at $y=0$ with its brane; the tip cutoff $L$ and the bag condition $\theta=0$ there; the torus with $\Delta k=0.25\,m$ and the closed-shell particle numbers; the expectation-value rule and normal ordering against the free Dirac sea, with the zero modes counted as particle levels and the particle--sea classification by continuity from $\lambda=0$; the local (exchange-only) functional built on the uniform-gas closed form, with no correlation term; a fixed metric (no back-reaction); the definition of the ground state by continuation in $\lambda$ and the smeared ensemble at the one level crossing; and the use of Delta-SCF as the approximation to the first excited state. \textbf{Not claimed:} that such a Kohn--Sham state existed in any universe, that it sources the primordial field, or anything about the creation of universes; Chapters 15 and 16 treat the pairing questions.

\subsection{Exercises}

\textbf{Exercise 13.1.} Using the two rules of Section 13.4, show that $J^2=1$ and $K_1^2=K_2^2=-1$, that $J$ commutes with $C$ and with $B$, and that $\gamma^8=\gamma^0\gamma^1\cdots\gamma^7$ anticommutes with $J$ and commutes with $K_1$ and $K_2$.

\textbf{Exercise 13.2.} Show that $A_0A_1A_4=-J$ and deduce $A_4\to j\sigma_x$ in the block basis.

\textbf{Exercise 13.3.} From $C=A_1K_1$ and $B=-iJK_1$ derive the block forms of $C$, $B$ and $BC$.

\textbf{Exercise 13.4.} Show that $\tilde\chi(y)=\sigma_z\chi(-y)$ solves the block equation with the mass function $-M(-y)$ if $\chi$ solves it with $M(y)$ (take $\xi$ and $v$ even), and conclude when the reflection is a symmetry.

\textbf{Exercise 13.5.} Check that $\chi=(e^{My},0)$ solves the block equation at $k=0$, $\varepsilon=0$, $v=0$ and satisfies the brane condition of parity $+$ and the tip condition. Compute its brane fraction (the fraction of $\int\chi^\dagger\chi\,dy$ within $|y|<1$) for $M=1$ and $L=2$ and $L=4$.

\textbf{Exercise 13.6.} For $M=1$, $L=3$, compute the two lowest positive levels of parity $+$ and the lowest positive level of parity $-$ at $k=0$.

\textbf{Exercise 13.7.} For $m=3$ ($H=1$) and $L=3$ compute $c$ and the first-order band energy $ck$ at the smallest torus momentum $k=0.25\,m=0.75\,H$. Compare with the Kohn--Sham gap $0.895833\,H$ of the free $N=8$ state for $m=3$ (\texttt{rust/scf\allowbreak{}/m3\_L3\_N\allowbreak{}8\_lam0\_T\allowbreak{}0/run.js\allowbreak{}on}).

\textbf{Exercise 13.8.} Prove $\frac d{dy}(\chi^\dagger\sigma_x\chi)=0$ for solutions of the block equation, and explain why a condition that kills the current at the tip makes it vanish everywhere.

\textbf{Exercise 13.9.} Show that for $N=8$ at $T=0$ the ground states at $\lambda$ and $-\lambda$ are related by moving every occupied orbital from block type $j$ to $-j$, and deduce $E_0(-\lambda)=-E_0(\lambda)$ and $\varepsilon_{\text{HOMO}}(-\lambda)=-\varepsilon_{\text{HOMO}}(\lambda)$. Why does the argument fail for $N=112$?

\textbf{Exercise 13.10.} Compute $S$ and $\mathrm{Tr}(BCP_+BCP_+)$ for the eight positive-energy states at rest, and the ratio $E_x/E_H$. Check that the uniform-gas closed form gives the same ratio when $S=n$.

\textbf{Exercise 13.11.} Derive $E=\sum w\varepsilon-E_H-E_x$ and evaluate it for $N=8$ at $+\hat\lambda_1$, where the solver reports $\sum w\varepsilon=-0.00196448147$, $E_H=5.06259\times10^{-6}$ and $E_x=-0.000982711$.

\textbf{Exercise 13.12.} Show that $a_{4,0}$ enters only through $k\,e^{-a_{4,0}}$, and that the problem with $a_{4,0}=0.5$ on the torus with $\Delta k=0.25\,m$ and coupling $\lambda$ is equivalent to the problem with $a_{4,0}=0$, $\Delta k'=0.25\,e^{-0.5}m$ and $\lambda'=\lambda\,e^{1.5}$.

\textbf{Exercise 13.13.} Compute the $\kappa$ that the free $N=1016$ state would need to be the source, from $\langle\rho\rangle=0.425999$, and explain why its sign rules it out.

\textbf{Exercise 13.14.} Repeat the Israel computation of Section 13.2 for the opposite gluing $W=e^{+H|y|}$ and find the sign of the brane energy density.

\textbf{Exercise 13.15.} Check that the swap $(a,b,j)\to(b,a,-j)$ of Section 13.14 maps the tip condition $b(-L)=0$ to $\tilde a(-L)=0$, exchanges the brane parities, keeps $n$ and reverses $s$.

\subsection{Answers to the exercises}

\textbf{Answer 13.1.} $J$ is a product of $k=3$ gammas: $J^2=(-1)^{3}(\gamma^0)^2(\gamma^1)^2(\gamma^4)^2=(-1)(1)(1)(-1)=1$. For $K_1$ ($k=2$): $(-1)^1\cdot1\cdot1=-1$; for $K_2$: $(-1)^1(-1)(-1)=-1$. $J$ commutes with its factors $\gamma^0,\gamma^1$ and anticommutes with $\gamma^2,\gamma^3$, so moving $J$ through $C=\gamma^0\gamma^1\gamma^2\gamma^3$ gives the sign $(+1)(+1)(-1)(-1)=+1$; $J$ commutes with $\gamma^4$, hence with $B=-iC\gamma^4$. Every $\gamma^a$ is a factor of $\gamma^8$ ($k=8$), so $\gamma^a\gamma^8=(-1)^7\gamma^8\gamma^a$: $\gamma^8$ anticommutes with every gamma, hence with a product of three ($J$) and commutes with a product of two ($K_1$, $K_2$).

\textbf{Answer 13.2.} $A_0A_1A_4=\gamma^0\gamma^0\gamma^1\gamma^0\gamma^4=\gamma^1\gamma^0\gamma^4=-\gamma^0\gamma^1\gamma^4=-J$ (using $(\gamma^0)^2=1$ and one swap). In the block $J=j$, $A_0=\sigma_z$, $A_1=-i\sigma_y$, and $\sigma_z(-i\sigma_y)=-\sigma_x$; so $-\sigma_xA_4=-j$, and multiplying by $\sigma_x$ (with $\sigma_x^2=1$) gives $A_4=j\sigma_x$.

\textbf{Answer 13.3.} $C=\gamma^0\gamma^1\gamma^2\gamma^3=A_1K_1\to(-i\sigma_y)(is_2)=s_2\sigma_y$. $B=-iC\gamma^4=-iJK_1\to-i\cdot j\cdot is_2=js_2$. $BC\to js_2\cdot s_2\sigma_y=j\sigma_y$ (since $s_2^2=1$).

\textbf{Answer 13.4.} $\tilde\chi'(y)=-\sigma_z\chi'(-y)=-\sigma_zN(-y)\sigma_z\,\tilde\chi(y)$. With $\sigma_z\sigma_z\sigma_z=\sigma_z$ and $\sigma_z\sigma_{x,y}\sigma_z=-\sigma_{x,y}$, $-\sigma_zN(-y)\sigma_z=-M(-y)\sigma_z-\xi(-y)k\sigma_y+ijE(-y)\sigma_x$, which for even $\xi$ and $E$ is $N(y)$ with the mass function $-M(-y)$. The reflection maps the problem onto itself exactly when $M(y)=-M(-y)$, an odd mass function.

\textbf{Answer 13.5.} With $k=\varepsilon=v=0$ the block equation reads $\chi_1'=M\chi_1$, $\chi_2'=-M\chi_2$, solved by $(e^{My},0)$; $\chi_2=0$ at $y=0$ and $y=-L$. The fraction is $\int_{-1}^0e^{2y}dy\big/\int_{-L}^0e^{2y}dy=(1-e^{-2})/(1-e^{-2L})$: $0.880797$ for $L=2$ and $0.864955$ for $L=4$. The solver reports $0.8808$ and $0.8650$ for the free $N=8$ states (\texttt{rust/scf\allowbreak{}/m1\_L2\_N\allowbreak{}8\_lam0\_T\allowbreak{}0} and \texttt{m1\_L4\_N8\allowbreak{}\_lam0\_T0}), whose particles are exactly these zero modes.

\textbf{Answer 13.6.} Parity $+$: $\sqrt{1+(\pi/3)^2}=1.447972$ and $\sqrt{1+(2\pi/3)^2}=2.320881$. Parity $-$: the smallest positive root of $\tan(3p)=-p$ is $p=0.818548$, so $\varepsilon=\sqrt{1+p^2}=1.292293$.

\textbf{Answer 13.7.} $c=\frac{2M}{2M-H}\cdot\frac{1-e^{-(2M-H)L}}{1-e^{-2ML}}=\frac65\cdot\frac{1-e^{-15}}{1-e^{-18}}=1.1999997$, and $ck=0.9000\,H$ at $k=0.75\,H$. The exact level is $0.895833\,H$, only 0.5 percent lower. For $m=1$ the first-order value was 10.6 percent too high at the same $k/m$: the heavier zero mode is more strongly localized at the brane, where $\xi\approx1$, so the linear approximation holds longer.

\textbf{Answer 13.8.} Section 13.5: $\frac d{dy}(\chi^\dagger\sigma_x\chi)=\chi^\dagger(N^\dagger\sigma_x+\sigma_xN)\chi=0$ because each Pauli matrix in $N$ anticommutes with $\sigma_x$ or, for the $\sigma_x$ term, appears with opposite signs in $N^\dagger$ and $N$. A quantity that is constant in $y$ and zero at $y=-L$ is zero everywhere, so the brane condition must, and does, also kill it; otherwise the eigenvalue problem would have no solution.

\textbf{Answer 13.9.} At $T=0$ only the eight zero-mode particle levels have nonzero weight (the sea is full, $w=-(1-f)=0$, and all other particle levels are empty). Move each occupied orbital $\chi$ unchanged from block type $j$ to $-j$ and replace $\lambda$ by $-\lambda$. Then $n_p$ is unchanged, $S_p\to-S_p$ (because $s=j\chi^\dagger\sigma_y\chi$), so $M_{\text{eff}}-m=\tfrac{15}{16}\lambda S_p$ is unchanged and $v=-\lambda n_p/16\to-v$. Writing $h_j=jD+v$ with $D$ built from $M_{\text{eff}}$, the new Hamiltonian is $-jD-v=-h_j$: the orbital is an eigenvector with the eigenvalue $-\varepsilon$, and the densities it produces are exactly those used, so the mapped state is self-consistent at $-\lambda$ (and it is the one reached by continuation, since both start from the same zero modes at $\lambda=0$; the zero modes are particle levels in both block types by convention). Then $\sum w\varepsilon\to-\sum w\varepsilon$, $E_H=\tfrac\lambda2\int S_p^2\to-E_H$ and $E_x\to-E_x$, so $E_0\to-E_0$, and $\varepsilon_{\text{HOMO}}\to-\varepsilon_{\text{HOMO}}$, as the table of Section 13.11 shows. For $N=112$ the occupied band levels have $\varepsilon>0$ and would be mapped to $\varepsilon<0$, which are sea levels; the mapped state is then not the ground state at $-\lambda$.

\textbf{Answer 13.10.} $S=\mathrm{Tr}(BCP_+)=8$ and $\mathrm{Tr}(BCP_+BCP_+)=\mathrm{Tr}(P_+)=8$ (Section 13.8), so $E_x/E_H=-\tfrac\lambda2\cdot8\big/\bigl(\tfrac\lambda2\cdot64\bigr)=-\tfrac18$. With $S=n$ the closed form gives $e_x=-\tfrac\lambda{32}\cdot2S^2=-\tfrac\lambda{16}S^2=-\tfrac18\cdot\tfrac\lambda2S^2$, the same ratio.

\textbf{Answer 13.11.} Section 13.9: the potentials of the two quadratic functionals contribute $2E_H+2E_x$ to $\sum w\varepsilon$. Numerically $E=-0.00196448147-0.00000506259+0.00098271090=-0.00098683316$, the reported $E_0=-9.86833\times10^{-4}$.

\textbf{Answer 13.12.} $a_{4,0}$ appears only in $\xi=e^{-Hy-a_{4,0}}$ (the volume factor $e^{6Hy}$ does not contain it), so $\xi k=e^{-Hy}(ke^{-a_{4,0}})$. The torus momenta $0.25\,m\,(n_1,n_2,n_3)$ become $0.25\,e^{-0.5}m\,(n_1,n_2,n_3)=0.151633\,m\,(n_1,n_2,n_3)$: the partner has $\Delta k'=0.151633\,m$ and $\ell'=\ell e^{0.5}$, with the same orbitals $\chi(y)$ and levels. Its proper densities $e^{-6Hy}\chi^\dagger\chi/\ell'^3$ are smaller by $e^{-1.5}$, so the potentials $\tfrac{15}{16}\lambda S_p$ and $-\lambda n_p/16$ agree only if $\lambda'=\lambda e^{1.5}$. The energies agree as well: $E_H\propto\lambda\int S_p^2e^{6Hy}\ell^3dy$ changes by $e^{1.5}\cdot e^{-3}\cdot e^{1.5}=1$, and likewise $E_x$.

\textbf{Answer 13.13.} $\kappa=-21H^2/\langle\rho\rangle=-21/0.425999=-49.30$. A negative $\kappa$ would reverse the sign of the gravitational coupling (a repulsive gravity for positive energy); with the physical sign $\kappa>0$ the required $\rho_{\text{req}}$ is negative, while this state has $\langle\rho\rangle>0$.

\textbf{Answer 13.14.} For $W=e^{+H|y|}$ the extrinsic curvature is $-H$ just below the brane and $+H$ just above, so $[K^\mu{}_\mu]=+2H$ on the six warped directions and $[K]=+12H$. Then $S^\mu{}_\mu=-(2H-12H)/\kappa=+10H/\kappa$ and $S^4{}_4=-(0-12H)/\kappa=+12H/\kappa$, so $\rho_{\text{brane}}=-S^4{}_4=-12H/\kappa<0$: this gluing would need a brane of negative energy. (It does not glue the notebook's patch, whose warp decreases toward the tip.)

\textbf{Answer 13.15.} The tip condition $b(-L)=0$ reads $\tilde a(-L)=0$ because $\tilde a=b$. At the brane, parity $+$ ($b(0)=0$) becomes $\tilde a(0)=0$, which is parity $-$ of the new system, and vice versa. $\tilde a^2+\tilde b^2=a^2+b^2$ keeps $n$, and $2\tilde j\tilde a\tilde b=-2jab$ reverses $s$.

\section{The Kohn--Sham approximation for dirac16complex00}

\subsection{What this chapter does, and the status of its results}

Chapter 13 built a Kohn--Sham (density-functional) approximation for dirac16complex, the field whose sixteen components anticommute, and computed its ground and first excited states in the static primordial field. This chapter does the same for the second field of the book, \textbf{dirac16complex00}, whose sixteen components are ordinary commuting complex numbers (Chapter 6). The request asked for such an approximation, with ground and first excited states, for each of the two fields (section 4 of \texttt{handoff/\allowbreak{}specs/ST\allowbreak{}AGE5\_SPE\allowbreak{}C.md}).

The main result fits in one sentence, and the chapter proves it step by step: \textbf{within the model defined in Section 14.2, the statistics of the components enters the mean-field energy only through the sign of the exchange term.} For dirac16complex the exchange energy density of the contact interaction is $e_x=-\tfrac\lambda{32}(n^2+S^2)$ (Section 13.8); for dirac16complex00 it is

\[
e_x^{00}=+\frac\lambda{32}\bigl(n^2+S^2\bigr),
\]

and consequently the Kohn--Sham effective mass and vector potential of dirac16complex00 are $M_{\text{eff}}^{00}=m+\tfrac{17}{16}\lambda S_p$ and $v^{00}=+\tfrac\lambda{16}n_p$, in place of $m+\tfrac{15}{16}\lambda S_p$ and $-\tfrac\lambda{16}n_p$.

The sentence contains a condition, ``within the model'', and a large part of the chapter is about it. A classical field has no occupation numbers, no Pauli principle and, as Section 14.2 shows, no lowest energy. A Kohn--Sham model of such a field must therefore \textbf{state what is being averaged}, and the answer changes the result: Section 14.5 exhibits one density matrix for which three different averaging rules give three different interaction energies. Section 14.8 states plainly what the model is and what it is not.

\textbf{Sources.} The binding specification is section 4 of \texttt{handoff/\allowbreak{}specs/ST\allowbreak{}AGE5\_SPE\allowbreak{}C.md}. The exact results come from two Wolfram packages, each with a verifier, and from two independent sympy checkers. The four reports, all in the folder \texttt{artifact\allowbreak{}s/dirac1\allowbreak{}6complex\allowbreak{}/pair-cr\allowbreak{}eation/} like every file named in this chapter without a folder, are:

\begin{itemize}
\item \texttt{wolfram-\allowbreak{}pairing-\allowbreak{}report.j\allowbreak{}son}, 141 of 141 checks true, from the package \texttt{wolfram/\allowbreak{}Dirac16C\allowbreak{}omplexPa\allowbreak{}iring.wl} and its verifier \texttt{scripts/\allowbreak{}verify\_d\allowbreak{}irac16co\allowbreak{}mplex\_pa\allowbreak{}iring.wl\allowbreak{}s};
\item \texttt{python-p\allowbreak{}airing-r\allowbreak{}eport.js\allowbreak{}on}, 172 of 172 checks true, written by the independent checker \texttt{scripts/\allowbreak{}check\_di\allowbreak{}rac16com\allowbreak{}plex\_pai\allowbreak{}ring.py};
\item \texttt{wolfram-\allowbreak{}dirac16c\allowbreak{}omplex00\allowbreak{}-report.\allowbreak{}json}, 46 of 46 checks true, from the package \texttt{wolfram/\allowbreak{}Dirac16C\allowbreak{}omplex00\allowbreak{}.wl} and its verifier \texttt{scripts/\allowbreak{}verify\_d\allowbreak{}irac16co\allowbreak{}mplex00.\allowbreak{}wls};
\item \texttt{python-d\allowbreak{}irac16co\allowbreak{}mplex00-\allowbreak{}report.j\allowbreak{}son}, 49 of 49 checks true, written by the independent checker \texttt{scripts/\allowbreak{}check\_di\allowbreak{}rac16com\allowbreak{}plex00.p\allowbreak{}y}.
\end{itemize}

The statistics sign is the check family whose names begin with \texttt{PAIR\_sta\allowbreak{}t\_} (Wolfram) and \texttt{\detokenize{S5_stat_}} (Python); the field theory of dirac16complex00 is the family \texttt{\detokenize{C00_}} (Wolfram) and \texttt{\detokenize{S5_}} (the dirac16complex00 checker). The exact formulas are collected in \texttt{pairing-\allowbreak{}theory.j\allowbreak{}son} (key \texttt{statisti\allowbreak{}cs}) and \texttt{dirac16c\allowbreak{}omplex00\allowbreak{}-theory.\allowbreak{}json} (keys \texttt{commutin\allowbreak{}gFieldSp\allowbreak{}ecifics} and \texttt{\detokenize{static}}). The numbers of Sections 14.9 and 14.10 come from the independent reference solver \texttt{scripts/\allowbreak{}ks\_refer\allowbreak{}ence\_sol\allowbreak{}ver.py}, driven by \texttt{scripts/\allowbreak{}ks\_refer\allowbreak{}ence\_pai\allowbreak{}rs.py}, whose outputs lie in \texttt{artifact\allowbreak{}s/dirac1\allowbreak{}6complex\allowbreak{}/pair-cr\allowbreak{}eation/r\allowbreak{}eference\allowbreak{}/} with the summary \texttt{referenc\allowbreak{}e-pairs-\allowbreak{}summary.\allowbreak{}json}. That summary records itself as incomplete (its key \texttt{\detokenize{complete}} is \texttt{\detokenize{false}}), and the committed Rust outputs of the Stage-5 subcommand \texttt{\detokenize{pairs}} in \texttt{artifact\allowbreak{}s/dirac1\allowbreak{}6complex\allowbreak{}/pair-cr\allowbreak{}eation/r\allowbreak{}ust/pair\allowbreak{}s/} contain no dirac16complex00 run at all. The Stage-5 documents planned in section 6 of \texttt{handoff/\allowbreak{}specs/ST\allowbreak{}AGE5\_SPE\allowbreak{}C.md} were not written when this chapter was written, and the Stage-5 gate had not been run; the chapter therefore cites the reports and outputs directly.

\textbf{Status of the results (read this first).}

\begin{itemize}
\item PROVED (derived here, and verified by the named exact checks): the Gaussian moments of Section 14.3; the four-amplitude averages for waves and for fermions (Section 14.4); the statistics sign of the exchange term (Theorem 14.4); the Hartree--Fock energy of the contact interaction for both fields, the filled-shell ratio $E_x=\pm E_H/8$ and the uniform-gas exchange $e_x=\pm\tfrac\lambda{32}(n^2+S^2)$ (Section 14.6); the Kohn--Sham potentials and equations of dirac16complex00 (Section 14.7); and the first two facts of Section 14.8 that limit the meaning of the model (the indefinite covariance of the model and the Krein weighting of its energies).
\item ASSUMED: the model itself (the prescriptions P1 to P4 of Section 14.2), together with everything that Chapter 13 assumes. The third point of Section 14.8, that the model is not a quantum theory, is a statement about how the model is defined (it is not obtained by quantizing the field); Pauli's spin--statistics theorem is quoted there, not proved, as the reason why no such quantization is attempted.
\item COMPUTED: the Kohn--Sham ground and first excited states of dirac16complex00 in one configuration ($m=3H$, $L=3/H$, $N=112$, $\hat\lambda\in\{0,+\hat\lambda_1\}$, $T=0$) and, without interaction, at one temperature; by one solver only, with no independent cross-check (Sections 14.9 and 14.10).
\item NOT COMPUTED: every other configuration of the planned run matrix, listed in Section 14.11.
\end{itemize}

\textbf{Units and conventions.} As in Chapter 13: $H=1$, coordinates $x_0,\dots,x_7$ with $x_4$ the time, $\eta=\mathrm{diag}(+1,+1,+1,+1,-1,-1,-1,-1)$, the notebook's gammas $\gamma^a$ (Chapter 2), $C=\gamma^0\gamma^1\gamma^2\gamma^3$, $\bar\Psi=\Psi^\dagger C$, $S=\bar\Psi\Psi$, $B=-iC\gamma^4$, $\gamma^8=\gamma^0\gamma^1\cdots\gamma^7$, and the Hermitian current $J^\mu=-i\bar\Psi\gamma^\mu\Psi$ with charge density $J^4=\Psi^\dagger B\Psi$ (Section 8.12). \textbf{The complex conjugate of a number $c$ is written $c^\ast$, never $\bar c$}, because the bar is reserved for the Dirac adjoint $\bar\Psi$. The numerical sections use $m=3$, so energies there are quoted in units of $H$ (not of $m$), with $m=3H$.

\subsection{A classical field has no occupation numbers: what must be averaged}

\textbf{The field and its Lagrangian.} At every point $x$ the field dirac16complex00 is a column $\Psi(x)=(\Psi_0(x),\dots,\Psi_{15}(x))^T$ of sixteen ordinary complex numbers that transform together as a Pin(4,4) spinor. Its Lagrangian is the same formula as for dirac16complex,

\[
\mathcal L=\sqrt{|g|}\,\Bigl[\tfrac12\bigl(\bar\Psi\gamma^\mu D_\mu\Psi-(D_\mu\bar\Psi)\gamma^\mu\Psi\bigr)-m\,\bar\Psi\Psi-\tfrac\lambda2(\bar\Psi\Psi)^2\Bigr],
\]

and so is its field equation $\gamma^\mu D_\mu\Psi=(m+\lambda S)\Psi$ (Chapters 6 and 7; checks \texttt{C00\_EL\_c\allowbreak{}ommuting\allowbreak{}Curved} and \texttt{C00\_EL\_i\allowbreak{}dentical\allowbreak{}FormBoth\allowbreak{}Statisti\allowbreak{}cs} of \texttt{wolfram-\allowbreak{}dirac16c\allowbreak{}omplex00\allowbreak{}-report.\allowbreak{}json}). For a given configuration of the field, $S(x)=\bar\Psi(x)\Psi(x)$ is a number at each point, and the interaction energy density is the number $\tfrac\lambda2S(x)^2$. Nothing is ``exchanged'': exchange, in Chapter 12, was a property of many-particle states.

\textbf{What Chapter 13 used.} The Kohn--Sham model of dirac16complex rested on four facts of the quantized fermion field: (i) each one-particle level is either empty or occupied once (the Pauli principle, Section 8.3); (ii) with the filled Dirac sea and normal ordering the energy is not negative (Section 8.10); (iii) the expectation-value rule $\langle\Psi^\dagger M\Psi\rangle=u^\dagger BMu$ makes the number density of every one-particle state positive (Section 8.12); and (iv) Wick's theorem computes the average of $S^2$ in a determinant or a thermal ensemble (Section 12.5). None of the four is available for a classical commuting field. Three exact results of Stage 5 show what goes wrong (Chapter 7 derives the first two in detail).

\textbf{Fact 1: the charge density has no sign.} The phase symmetry $\Psi\to e^{i\alpha}\Psi$ gives the conserved charge density $J^4=\Psi^\dagger B\Psi$ for commuting components too (check \texttt{C00\_curr\allowbreak{}ent\_cons\allowbreak{}ervation\allowbreak{}AndReali\allowbreak{}ty}). $B$ is Hermitian with $B^2=1$ and has the eigenvalues $+1$ and $-1$ eight times each (Section 2.12), so $J^4$ takes both signs. A worked example with the explicit matrix: the rows 6 and 15 of $B$ read $(Bu)_6=i\,u_{15}$ and $(Bu)_{15}=-i\,u_6$ (\texttt{dirac16c\allowbreak{}omplex00\allowbreak{}-theory.\allowbreak{}json}, key \texttt{commutin\allowbreak{}gFieldSp\allowbreak{}ecifics.\allowbreak{}B}). Write $e_n$ for the column with 1 in place $n$ (counting from 0) and take

\[
\Psi_+=i\,e_6+e_{15},\qquad \Psi_-=-i\,e_6+e_{15}.
\]

For $\Psi_+$: $(B\Psi_+)_6=i\cdot1=i=(\Psi_+)_6$ and $(B\Psi_+)_{15}=-i\cdot i=1=(\Psi_+)_{15}$, so $B\Psi_+=\Psi_+$ and $J^4=\Psi_+^\dagger\Psi_+=2$. For $\Psi_-$: $(B\Psi_-)_6=i=-(\Psi_-)_6$ and $(B\Psi_-)_{15}=-i(-i)=-1=-(\Psi_-)_{15}$, so $B\Psi_-=-\Psi_-$ and $J^4=-2$ (check \texttt{C00\_char\allowbreak{}ge\_indef\allowbreak{}inite}, which records these two vectors). The argument of Section 8.8 applies unchanged to commuting spinors: every charge density built from a Spin(4,4)-invariant form is indefinite. So the probability density of Dirac's 1928 theory has no analogue here, and dirac16complex00 is not a wave function with a probability interpretation.

\textbf{Fact 2: the classical energy has no lower bound.} In flat space with $\lambda=0$, a plane wave $\Psi=u\,e^{-i\omega x_4+i\mathbf k\cdot\mathbf x}$ that solves the field equation has the energy density $\rho=T_{44}=\omega\,J^4$ with $J^4=u^\dagger Bu$ (check \texttt{C00\_ener\allowbreak{}gy\_unbou\allowbreak{}ndedBelo\allowbreak{}w}). The report lists four exact solutions with $m=1$ and the wave numbers $(1,1,2,3,0,0,0)$ along $(x_0,x_1,x_2,x_3,x_5,x_6,x_7)$, so that $\omega^2=1+1+1+4+9=16$ and $\omega=\pm4$ (they are the entry \texttt{energyUn\allowbreak{}boundedB\allowbreak{}elow} under the key \texttt{commutin\allowbreak{}gFieldSp\allowbreak{}ecifics} of \texttt{dirac16c\allowbreak{}omplex00\allowbreak{}-theory.\allowbreak{}json}):

\begingroup
\small
\begin{longtable}{@{}>{\raggedright\arraybackslash}p{0.267\linewidth}>{\raggedright\arraybackslash}p{0.267\linewidth}>{\raggedright\arraybackslash}p{0.267\linewidth}@{}}
\toprule
\textbf{$\omega$} & \textbf{$J^4$} & \textbf{$\rho$} \\
\midrule
\endfirsthead
\toprule
\textbf{$\omega$} & \textbf{$J^4$} & \textbf{$\rho$} \\
\midrule
\endhead
$+4$ & $+224$ & $+896$ \\
$+4$ & $-224$ & $-896$ \\
$-4$ & $+32$ & $-128$ \\
$-4$ & $-32$ & $+128$ \\
\bottomrule
\end{longtable}
\endgroup

Multiplying a solution by a number $c$ gives a solution and multiplies $\rho$ by $|c|^2$, so the second solution can be made to carry as negative an energy as we like. The interaction does not help: with $m=1$ and $\lambda=-1$ the homogeneous states with $S=1$, $10$ and $100$ solve the nonlinear equation and have $\rho=mS+\tfrac\lambda2S^2=\tfrac12$, $10-50=-40$ and $100-5000=-4900$; for $\lambda=+1$ the report lists a family of self-consistent plane waves whose energy density decreases without bound as well. For the quantized Grassmann field the same indefinite matrix $B$ became the anticommutator, and the filled Dirac sea with normal ordering made the Hamiltonian non-negative in the good sector (Section 8.10). That mechanism needs the Pauli principle; a classical commuting field does not have it.

\textbf{Fact 3: no Pauli principle.} The amplitude of a mode of a classical field is any complex number. Nothing restricts the squared amplitude, the classical analogue of an occupation number, to 0 or 1.

\textbf{Consequence.} For dirac16complex00 the phrase ``ground state'' in the sense of Chapter 12 (the state of lowest energy, Section 12.2) has no meaning, because there is no lowest energy. A Kohn--Sham model of this field can therefore not be derived from a variational principle of the classical field. It must be \textbf{defined}, and the definition must say (a) over which collection of field configurations the densities $n$ and $S$ are averages, (b) which mean occupations the modes carry, and (c) which bilinear of the field is called the density.

\textbf{Why (a) matters: one mode, two answers.} Take one mode, $\Psi=c\,u$ with a fixed normalized spinor $u$ and a random complex amplitude $c$ whose squared modulus is $f$ on average. The scalar density is proportional to $|c|^2$, and the interaction energy to its square. (i) If every member of the collection has the same modulus, $|c|^2=f$, and only the phase of $c$ is random, then $\langle|c|^4\rangle=f^2=\langle|c|^2\rangle^2$. (ii) If $c$ is a Gaussian random number (Section 14.3), then $\langle|c|^4\rangle=2f^2$. Two collections with the same average density give average interaction energies that differ by a factor 2. The difference $\langle S^2\rangle-\langle S\rangle^2$ is exactly what the exchange term of a density functional must represent (Section 14.5), so a density functional must say which collection it describes.

\textbf{The Stage-5 model.} The Stage-5 specification (section 4 of \texttt{handoff/\allowbreak{}specs/ST\allowbreak{}AGE5\_SPE\allowbreak{}C.md}) answers (a) to (c) as follows, and \texttt{pairing-\allowbreak{}theory.j\allowbreak{}son} (key \texttt{statisti\allowbreak{}cs.posit\allowbreak{}ivity.st\allowbreak{}atus}) records the resulting status: the model is ``a DFT-motivated mean field of a random-phase (Gaussian) ensemble of classical modes with Fermi-Dirac occupations (Pauli filling imposed by prescription), not a quantum theory of commuting spinors''. In detail:

\begin{itemize}
\item \textbf{P1 (the ensemble).} The field is a sum of Kohn--Sham orbitals $u_n$ of Chapter 13 with independent random complex amplitudes $c_n$, and every average of products of amplitudes is computed with the rule of a circular Gaussian (Sections 14.3 and 14.4), with the mean squared amplitudes fixed by P2.
\item \textbf{P2 (the occupations).} The mean squared amplitudes are the weights $w_n$ of Chapter 13: the Fermi--Dirac occupation $w_n=f_n$ on a particle level and $w_n=-(1-f_n)$ on a sea level, with normal ordering against the free Dirac sea (Section 13.7). For a commuting field this ``Pauli filling'' is not a consequence of anything; it is imposed, as the user specified (``fermion-gas thermodynamics''). Note what the sea weights mean for P1: a sea level that is not completely full has $w_n<0$, and no random complex number has a negative mean squared modulus. So already this bookkeeping of particles and sea makes the ``ensemble'' formal: P1 fixes a rule for computing averages, the rule of a circular Gaussian with the numbers $w_n$ in place of $f$, and not a probability distribution of fields. (At $T=0$ every sea level is full and has $w_n=0$. Prescription P3 adds a second, independent reason why the ensemble is formal: the Krein signs of Section 14.8.)
\item \textbf{P3 (the densities).} The density of a bilinear $\Psi^\dagger M\Psi$ is computed with the expectation-value rule of the quantized fermion field, $u^\dagger BMu$ for a mode $u$ (Section 8.12), so that the two fields differ \textbf{only} in the statistics.
\item \textbf{P4 (the rest).} Exchange only, with no correlation term (Section 14.8), and everything else of Chapter 13: the static primordial field with the Z2 brane, the good sector, the eight $2\times2$ blocks, the boundary conditions, the torus, and the definition of the ground state by continuation in the coupling (erratum E4.8).
\end{itemize}

Section 14.8 examines what P3 costs.

\subsection{Random complex amplitudes from zero}

\textbf{Probability densities.} A random real number $x$ is described by a \textbf{probability density} $p(x)\ge0$ with $\int p(x)\,dx=1$: the probability that $x$ falls between $a$ and $b$ is $\int_a^bp(x)\,dx$, and the \textbf{average} (mean, expectation) of a function $g$ of $x$ is $\langle g\rangle=\int g(x)\,p(x)\,dx$. A random complex number $c=x+iy$ is a random point of the plane, described by a density $p(c)=p(x,y)$ with $\int\!\int p\,dx\,dy=1$ and $\langle g\rangle=\int\!\int g\,p\,dx\,dy$, which we write $\int g\,p\,d^2c$. Two random numbers $c_0,c_1$ are \textbf{independent} when their joint density is the product $p_0(c_0)\,p_1(c_1)$; then the double integral factorizes and $\langle g(c_0)h(c_1)\rangle=\langle g(c_0)\rangle\langle h(c_1)\rangle$.

\textbf{Polar coordinates in the plane.} Write $x=r\cos\varphi$ and $y=r\sin\varphi$, that is $c=r\,e^{i\varphi}$ (the polar form of Section 1.5), with $r\ge0$ and $0\le\varphi<2\pi$. The small region between the radii $r$ and $r+dr$ and the angles $\varphi$ and $\varphi+d\varphi$ is almost a rectangle with the sides $dr$ (along the radius) and $r\,d\varphi$ (along the circle of radius $r$), so its area is $r\,dr\,d\varphi$. Hence

\[
\int g\,d^2c=\int_0^\infty\!\!\int_0^{2\pi}g\bigl(re^{i\varphi}\bigr)\,r\,d\varphi\,dr .
\]

\textbf{The circular Gaussian.} For a number $f>0$ define

\[
p_f(c)=\frac1{\pi f}\,e^{-|c|^2/f}.
\]

It depends only on $|c|=r$, so all phases are equally likely: this is the precise meaning of a \textbf{random phase}. It is normalized: with the substitution $s=r^2/f$, $ds=2r\,dr/f$,

\[
\int p_f\,d^2c=\frac1{\pi f}\cdot2\pi\int_0^\infty e^{-r^2/f}\,r\,dr=\frac2f\cdot\frac f2\int_0^\infty e^{-s}\,ds=1 .
\]

(The density of the modulus alone, $(2r/f)\,e^{-r^2/f}$, is called the Rayleigh distribution.)

\textbf{A factorial integral.} For $k=0,1,2,\dots$ let $I_k=\int_0^\infty s^ke^{-s}\,ds$. Then $I_0=1$, and integrating by parts (Section 1.10) with $g'=e^{-s}$, $g=-e^{-s}$,

\[
I_k=\bigl[-s^ke^{-s}\bigr]_0^\infty+k\int_0^\infty s^{k-1}e^{-s}\,ds=k\,I_{k-1},
\]

because the boundary term vanishes at both ends ($e^{-s}$ decreases faster than any power grows, a fact of calculus used without proof). Hence $I_k=k(k-1)\cdots1=k!$.

\textbf{Proposition 14.1 (moments of one Gaussian amplitude).} For the circular Gaussian with $\langle|c|^2\rangle=f$ and integers $\alpha,\beta\ge0$,

\[
\bigl\langle(c^\ast)^\alpha c^\beta\bigr\rangle=\delta_{\alpha\beta}\,\alpha!\,f^\alpha .
\]

\emph{Proof.} In polar form $(c^\ast)^\alpha c^\beta=r^{\alpha+\beta}e^{i(\beta-\alpha)\varphi}$. The angular integral is $\int_0^{2\pi}e^{ik\varphi}d\varphi=2\pi$ for $k=0$ and $\bigl[e^{ik\varphi}/(ik)\bigr]_0^{2\pi}=0$ for every integer $k\ne0$, so only $\alpha=\beta$ survives. Then, with $s=r^2/f$, so that $r^{2\alpha}=f^\alpha s^\alpha$ and $r\,dr=\tfrac f2ds$,

\[
\bigl\langle|c|^{2\alpha}\bigr\rangle=\frac{2\pi}{\pi f}\int_0^\infty r^{2\alpha}e^{-r^2/f}\,r\,dr=\frac2f\cdot f^\alpha\cdot\frac f2\,I_\alpha=\alpha!\,f^\alpha .\qquad\square
\]

In particular $\langle|c|^2\rangle=f$, confirming the name of the parameter, and $\langle|c|^4\rangle=2f^2$: the squared modulus fluctuates as much as its mean, $\langle|c|^4\rangle-\langle|c|^2\rangle^2=f^2$.

\textbf{Worked example.} For $f=\tfrac12$: $\langle|c|^2\rangle=\tfrac12$, $\langle|c|^4\rangle=2\cdot\tfrac14=\tfrac12$, $\langle|c|^6\rangle=6\cdot\tfrac18=\tfrac34$, and $\langle c\rangle=\langle c^2\rangle=\langle(c^\ast)^2c\rangle=0$. For comparison, the \textbf{fixed-modulus} collection $c=\sqrt f\,e^{i\varphi}$ with a random phase alone has $\langle(c^\ast)^\alpha c^\beta\rangle=\delta_{\alpha\beta}f^\alpha$ (the same angular integral, no radial integral): $\langle|c|^4\rangle=f^2=\tfrac14$.

\textbf{The quantum versions.} A \textbf{fermion mode} (Section 8.3) is empty or occupied once. If it is occupied with probability $f$, its occupation number $k\in\{0,1\}$ has $\langle k\rangle=f$ and $\langle k(k-1)\rangle=0$, because $k(k-1)=0$ for $k=0$ and $k=1$. In operator language $b^\dagger b^\dagger bb=0$ because $(b^\dagger)^2=0$. A \textbf{boson mode} is described by an operator with the commutator $[b,b^\dagger]=1$ instead of the anticommutator; its number operator $N=b^\dagger b$ has the eigenvalues $k=0,1,2,\dots$ (the quantum harmonic oscillator, a standard fact that we quote), and $b^\dagger b^\dagger bb=b^\dagger(bb^\dagger-1)b=N^2-N=N(N-1)$. In the Gibbs state of Section 12.14 the probability of $k$ quanta in a level of energy $\varepsilon$ is proportional to $e^{-k(\varepsilon-\mu)/T}$, that is $p_k=(1-x)\,x^k$ with $x=e^{-(\varepsilon-\mu)/T}<1$ (normalized by the geometric series $\sum_kx^k=1/(1-x)$). Differentiating the geometric series once and twice, $\sum_kkx^k=x/(1-x)^2$ and $\sum_kk(k-1)x^k=2x^2/(1-x)^3$, so

\[
\langle k\rangle=\frac x{1-x}=:f,\qquad \langle k(k-1)\rangle=(1-x)\,\frac{2x^2}{(1-x)^3}=2f^2 .
\]

(The first formula is the Bose--Einstein occupation $f=1/(e^{(\varepsilon-\mu)/T}-1)$.) The four cases side by side:

\begingroup
\small
\begin{longtable}{@{}>{\raggedright\arraybackslash}p{0.267\linewidth}>{\raggedright\arraybackslash}p{0.267\linewidth}>{\raggedright\arraybackslash}p{0.267\linewidth}@{}}
\toprule
\textbf{one mode} & \textbf{mean number} & \textbf{mean number of ordered pairs} \\
\midrule
\endfirsthead
\toprule
\textbf{one mode} & \textbf{mean number} & \textbf{mean number of ordered pairs} \\
\midrule
\endhead
fermion, occupation probability $f$ & $f$ & $0$ \\
boson, thermal & $f$ & $2f^2$ \\
classical Gaussian amplitude & $f$ & $2f^2$ \\
classical fixed modulus, random phase & $f$ & $f^2$ \\
\bottomrule
\end{longtable}
\endgroup

The last column is $\langle k(k-1)\rangle$ for the quantum modes and $\langle|c|^4\rangle$ for the classical ones. The thermal boson and the classical Gaussian agree: both follow from $\langle(c^\ast)^\alpha c^\alpha\rangle=\alpha!f^\alpha$ and its operator analogue. The Wolfram verifier computes these moments exactly, $\langle b^\dagger b\rangle=f$, $\langle b^\dagger b^\dagger bb\rangle=2f^2$, $\langle b^\dagger b^\dagger b^\dagger bbb\rangle=6f^3$, $\langle|c|^4\rangle=2f^2$, $\langle|c|^6\rangle=6f^3$ and $\langle(c^\ast)^2c\rangle=0$ (check \texttt{PAIR\_sta\allowbreak{}t\_single\allowbreak{}ModeMome\allowbreak{}nts}, \texttt{wolfram-\allowbreak{}pairing-\allowbreak{}report.j\allowbreak{}son}).

\subsection{Four amplitudes: Isserlis' rule for waves and Wick's rule for fermions}

\textbf{Many modes.} Let $c_0,\dots,c_{M-1}$ be independent circular Gaussian amplitudes with $\langle|c_n|^2\rangle=f_n$. The exchange term of Section 14.5 needs the averages of products of two conjugated and two plain amplitudes.

\textbf{Proposition 14.2 (Isserlis' rule for four amplitudes).} For all mode labels $n,q,r,p$,

\[
\bigl\langle c_n^\ast\,c_q^\ast\,c_r\,c_p\bigr\rangle=f_nf_q\bigl(\delta_{np}\delta_{qr}+\delta_{nr}\delta_{qp}\bigr).
\]

\emph{Proof.} By independence the average is the product, over the modes, of the averages of the factors that belong to each mode, and by Proposition 14.1 a mode contributes zero unless it appears as many times conjugated as unconjugated. Two cases remain. (i) $n=q$: the mode $n$ appears twice conjugated, so both plain factors must belong to it, $p=r=n$, and the average is $\langle|c_n|^4\rangle=2f_n^2$. (ii) $n\ne q$: the modes $n$ and $q$ appear once conjugated each, so $\{p,r\}=\{n,q\}$, and either $p=n,r=q$ or $p=q,r=n$; each gives $\langle|c_n|^2\rangle\langle|c_q|^2\rangle=f_nf_q$. The formula reproduces both cases: for $n=q=p=r$ the bracket is $1+1$, and otherwise exactly the matching Kronecker product is 1. $\square$

\textbf{Proposition 14.3 (Wick's rule for four fermion operators).} In a quasi-free fermion state with the occupation probabilities $f_n$ of orthonormal modes (a Slater determinant when every $f_n$ is 0 or 1, a non-interacting thermal ensemble otherwise; Section 12.5),

\[
\bigl\langle a_n^\dagger\,a_q^\dagger\,a_r\,a_p\bigr\rangle=f_nf_q\bigl(\delta_{np}\delta_{qr}-\delta_{nr}\delta_{qp}\bigr).
\]

\emph{Proof.} Section 12.5 gives $\langle a_{p'}^\dagger a_{q'}^\dagger a_{s'}a_{r'}\rangle=\rho_{r'p'}\rho_{s'q'}-\rho_{s'p'}\rho_{r'q'}$ for any quasi-free state with the one-body density matrix $\rho_{ba}=\langle a_a^\dagger a_b\rangle$. Put $(p',q',s',r')=(n,q,r,p)$: the right side is $\rho_{pn}\rho_{rq}-\rho_{rn}\rho_{pq}$. In the basis of the modes $\rho$ is diagonal, $\rho_{pn}=f_n\delta_{pn}$, which gives the claim. $\square$

\textbf{Reading the two rules.} Both have two terms. The \textbf{direct} pairing joins $c_n^\ast$ with $c_p$ and $c_q^\ast$ with $c_r$; the \textbf{exchange} pairing joins $c_n^\ast$ with $c_r$ and $c_q^\ast$ with $c_p$. The direct term has the sign $+$ in both rules. The exchange term has the sign $+$ for Gaussian waves and $-$ for fermions: for fermions, reaching the exchange pairing requires one exchange of two anticommuting operators. A compact way to remember this: arrange the four two-point averages in the $2\times2$ matrix

\[
G=\begin{pmatrix}\langle c_n^\ast c_p\rangle&\langle c_n^\ast c_r\rangle\\ \langle c_q^\ast c_p\rangle&\langle c_q^\ast c_r\rangle\end{pmatrix}=\begin{pmatrix}g_{00}&g_{01}\\ g_{10}&g_{11}\end{pmatrix};
\]

the Gaussian average is its \textbf{permanent} $g_{00}g_{11}+g_{01}g_{10}$, and the fermion average is its \textbf{determinant} $g_{00}g_{11}-g_{01}g_{10}$. The general theorems, which we quote and do not use, say the same for products of any length: the average of $k$ conjugated and $k$ plain Gaussian amplitudes is the permanent of the $k\times k$ matrix of two-point averages (Isserlis, 1918), and the corresponding fermion average is its determinant (Wick, 1950).

\textbf{Worked examples.} With $n=q=p=r=0$ the Gaussian rule gives $2f_0^2$ and the fermion rule $0$ (the Pauli principle). With $n=p=0$ and $q=r=1$ both rules give $f_0f_1$ (direct pairing only). With $n=r=0$ and $q=p=1$ the Gaussian rule gives $+f_0f_1$ and the fermion rule $-f_0f_1$ (exchange pairing only).

\subsection{The statistics sign of the exchange term}

\textbf{The setting.} Let $\psi$ be a column of $d$ components (later $d=16$) built from $d$ orthonormal modes $u_0,\dots,u_{d-1}$, the columns of a unitary $d\times d$ matrix $U$:

\[
\psi=\sum_nc_n\,u_n=U\,c\quad(\text{waves}),\qquad \psi=\sum_na_n\,u_n=U\,a\quad(\text{fermions}).
\]

Define the \textbf{density matrix} $\rho=\sum_nf_n\,u_nu_n^\dagger=UFU^\dagger$ with $F=\mathrm{diag}(f_0,\dots,f_{d-1})$. Then the two-point average is

\[
\langle\psi_a^\ast\psi_b\rangle=\sum_{n,p}U_{an}^\ast U_{bp}\,\langle c_n^\ast c_p\rangle=\sum_nf_n\,U_{bn}U_{an}^\ast=\rho_{ba},
\]

using $\langle c_n^\ast c_p\rangle=f_n\delta_{np}$ (Proposition 14.1 and independence); for fermions the same computation with $\langle a_n^\dagger a_p\rangle=f_n\delta_{np}$ gives $\langle\psi_a^\dagger\psi_b\rangle=\rho_{ba}$. So both averages of a single bilinear are the same: $\langle\psi^\dagger X\psi\rangle=\sum_{a,b}X_{ab}\rho_{ba}=\mathrm{Tr}(X\rho)$ for every $d\times d$ matrix $X$.

\textbf{Theorem 14.4 (the statistics sign).} For all $d\times d$ matrices $X$ and $Y$,

\[
\bigl\langle(\psi^\dagger X\psi)(\psi^\dagger Y\psi)\bigr\rangle=\mathrm{Tr}(X\rho)\,\mathrm{Tr}(Y\rho)+\mathrm{sg}\,\mathrm{Tr}(X\rho\,Y\rho),
\]

with $\mathrm{sg}=+1$ for independent circular Gaussian amplitudes (a classical commuting field) and $\mathrm{sg}=-1$ for the normal-ordered product $:(\psi^\dagger X\psi)(\psi^\dagger Y\psi):$ of a fermion field in a quasi-free state. We call $\mathrm{sg}$ the \textbf{statistics sign}, with the name that the reports use (\texttt{pairing-\allowbreak{}theory.j\allowbreak{}son}, key \texttt{conventi\allowbreak{}ons.stat\allowbreak{}isticsSi\allowbreak{}gn}) and that Chapter 15 uses; Chapter 6 writes the same sign as $s$.

\emph{Proof for waves.} Write $\tilde X=U^\dagger XU$ and $\tilde Y=U^\dagger YU$ (the matrices in the basis of the modes). Then $\psi^\dagger X\psi=c^\dagger U^\dagger XUc=\sum_{n,p}c_n^\ast\tilde X_{np}c_p$, and, because amplitudes commute,

\[
\bigl\langle(\psi^\dagger X\psi)(\psi^\dagger Y\psi)\bigr\rangle=\sum_{n,p,q,r}\tilde X_{np}\tilde Y_{qr}\,\bigl\langle c_n^\ast c_q^\ast c_rc_p\bigr\rangle .
\]

Insert Proposition 14.2. The direct term sets $p=n$, $r=q$ and gives $\sum_nf_n\tilde X_{nn}\sum_qf_q\tilde Y_{qq}=\mathrm{Tr}(F\tilde X)\,\mathrm{Tr}(F\tilde Y)$; the exchange term sets $r=n$, $p=q$ and gives $\sum_{n,q}f_n\tilde X_{nq}f_q\tilde Y_{qn}=\mathrm{Tr}(F\tilde XF\tilde Y)$. Return to the original basis with the cyclic property of the trace and $UFU^\dagger=\rho$: $\mathrm{Tr}(F\tilde X)=\mathrm{Tr}(FU^\dagger XU)=\mathrm{Tr}(UFU^\dagger X)=\mathrm{Tr}(X\rho)$, and $\mathrm{Tr}(F\tilde XF\tilde Y)=\mathrm{Tr}(UFU^\dagger X\,UFU^\dagger Y)=\mathrm{Tr}(\rho X\rho Y)=\mathrm{Tr}(X\rho Y\rho)$.

\emph{Proof for fermions.} Now $\psi^\dagger X\psi=\sum_{n,p}a_n^\dagger\tilde X_{np}a_p$. Normal ordering moves every creation operator to the left of every annihilation operator, with the sign of the permutation, and drops the terms that the anticommutator would produce (Section 13.8): $:a_n^\dagger a_p\,a_q^\dagger a_r:=-a_n^\dagger a_q^\dagger a_pa_r=+a_n^\dagger a_q^\dagger a_ra_p$ (one exchange to move $a_q^\dagger$ past $a_p$, a second to exchange $a_p$ and $a_r$). Proposition 14.3 then gives the direct term with $+$ and the exchange term with $-$; the return to the original basis is the same. $\square$

\textbf{Remark 1: the full product.} Without normal ordering, $a_pa_q^\dagger=\delta_{pq}-a_q^\dagger a_p$ adds the term $\sum_{n,p,r}\tilde X_{np}\tilde Y_{pr}\langle a_n^\dagger a_r\rangle=\mathrm{Tr}(XY\rho)$, so the plain fermion product has the average $\mathrm{Tr}(X\rho)\mathrm{Tr}(Y\rho)-\mathrm{Tr}(X\rho Y\rho)+\mathrm{Tr}(XY\rho)$. A boson field ($b_pb_q^\dagger=\delta_{pq}+b_q^\dagger b_p$) gives the same with $+\mathrm{Tr}(X\rho Y\rho)$, and its normal-ordered average equals the classical Gaussian one. For classical amplitudes there is nothing to reorder and no $\mathrm{Tr}(XY\rho)$ term. The repository verifies all of this exactly (in exact arithmetic, with no rounding), and it is worth saying precisely what each check covers. (The prefix \texttt{\detokenize{PAIR_}} marks checks of \texttt{wolfram-\allowbreak{}pairing-\allowbreak{}report.j\allowbreak{}son}, the prefix \texttt{\detokenize{S5_}} checks of the independent \texttt{python-p\allowbreak{}airing-r\allowbreak{}eport.js\allowbreak{}on}.)

\begin{itemize}
\item \textbf{Fermions.} \texttt{PAIR\_sta\allowbreak{}t\_fermio\allowbreak{}nWickMin\allowbreak{}us} works on the complete 16-dimensional state space of four modes, in 12 quasi-free states (3 bases of modes times 4 sets of occupations, pure determinants and mixed states), with two fixed Hermitian $4\times4$ test matrices whose entries are complex numbers with rational real and imaginary parts; it checks both the normal-ordered and the full product. \texttt{S5\_stat\_\allowbreak{}fermionW\allowbreak{}ickMinus} does the same on its own 16-dimensional space of four modes, with its own 12 states and test matrices.
\item \textbf{Thermal bosons.} \texttt{PAIR\_sta\allowbreak{}t\_bosonT\allowbreak{}hermalWi\allowbreak{}ckPlus} checks the normal-ordered product for the same four occupation sets $\{1,1,0,0\}$, $\{1,0,1,1\}$, $\{\tfrac13,\tfrac25,0,1\}$ and $\{\tfrac12,\tfrac17,\tfrac34,\tfrac29\}$ of four modes in the same three bases; only its single-mode moments are computed symbolically in $f$, from the thermal series (the package also forms the full product, but this check does not test it). \texttt{S5\_stat\_\allowbreak{}bosonThe\allowbreak{}rmalWick\allowbreak{}Plus} checks the normal-ordered \textbf{and} the full product with symbolic occupations $f_0,f_1,f_2$ of three modes in one rational basis.
\item \textbf{Classical Gaussian amplitudes.} The moments are exact integrals over the density $p_f$. \texttt{PAIR\_sta\allowbreak{}t\_classi\allowbreak{}calGauss\allowbreak{}ianWickP\allowbreak{}lus} uses the four occupation sets and three bases above, and also checks that every four-amplitude moment of the Gaussian equals the corresponding normal-ordered moment of the thermal boson; \texttt{S5\_stat\_\allowbreak{}classica\allowbreak{}lGaussia\allowbreak{}nWickPlu\allowbreak{}s} uses symbolic occupations of three modes, and confirms that no $\mathrm{Tr}(XY\rho)$ term appears.
\end{itemize}

\textbf{Remark 2: fixed moduli.} If the amplitudes have fixed moduli, $c_n=\sqrt{f_n}\,e^{i\varphi_n}$ with independent random phases, only case (i) of Proposition 14.2 changes ($f_n^2$ instead of $2f_n^2$), so

\[
\bigl\langle(\psi^\dagger X\psi)(\psi^\dagger Y\psi)\bigr\rangle_{\text{fixed moduli}}=\mathrm{Tr}(X\rho)\,\mathrm{Tr}(Y\rho)+\mathrm{Tr}(X\rho Y\rho)-\sum_nf_n^2\,\tilde X_{nn}\tilde Y_{nn}
\]

(check \texttt{PAIR\_sta\allowbreak{}t\_fixedA\allowbreak{}mplitude\allowbreak{}PhasesDe\allowbreak{}viate}). The $+$ sign of the exchange term is therefore a property of \textbf{Gaussian} amplitudes, the amplitudes of thermal (chaotic) waves; it is not a property of every random-phase collection.

\textbf{Worked example: one density matrix, three answers.} Take $d=2$, the modes $u_0=(1,1)^T/\sqrt2$ and $u_1=(1,-1)^T/\sqrt2$, the occupations $f_0=1$, $f_1=\tfrac12$, and $X=Y=\sigma_x=\begin{pmatrix}0&1\\1&0\end{pmatrix}$, the first Pauli matrix of Section 2.3. Then

\[
\rho=f_0u_0u_0^\dagger+f_1u_1u_1^\dagger=\frac12\begin{pmatrix}f_0+f_1&f_0-f_1\\ f_0-f_1&f_0+f_1\end{pmatrix}=\begin{pmatrix}\tfrac34&\tfrac14\\ \tfrac14&\tfrac34\end{pmatrix},\qquad
\sigma_x\rho=\begin{pmatrix}\tfrac14&\tfrac34\\ \tfrac34&\tfrac14\end{pmatrix},
\]

so $\mathrm{Tr}(\sigma_x\rho)=\tfrac12$ and $\mathrm{Tr}(\sigma_x\rho\,\sigma_x\rho)=\tfrac1{16}+\tfrac9{16}+\tfrac9{16}+\tfrac1{16}=\tfrac54$. Since $\sigma_xu_0=u_0$ and $\sigma_xu_1=-u_1$, the mode-basis matrix is $\tilde X=\mathrm{diag}(1,-1)$. The three averages of $(\psi^\dagger\sigma_x\psi)^2$ are:

\begingroup
\small
\begin{longtable}{@{}>{\raggedright\arraybackslash}p{0.267\linewidth}>{\raggedright\arraybackslash}p{0.267\linewidth}>{\raggedright\arraybackslash}p{0.267\linewidth}@{}}
\toprule
\textbf{collection} & \textbf{formula} & \textbf{value} \\
\midrule
\endfirsthead
\toprule
\textbf{collection} & \textbf{formula} & \textbf{value} \\
\midrule
\endhead
Gaussian amplitudes ($\mathrm{sg}=+1$) & $\tfrac14+\tfrac54$ & $\tfrac32$ \\
fixed moduli, random phases & $\tfrac32-(f_0^2+f_1^2)$ & $\tfrac14$ \\
fermions, normal ordered ($\mathrm{sg}=-1$) & $\tfrac14-\tfrac54$ & $-1$ \\
\bottomrule
\end{longtable}
\endgroup

A direct check makes the three numbers transparent. In the mode basis $\psi^\dagger\sigma_x\psi=|c_0|^2-|c_1|^2$. With fixed moduli this is always $f_0-f_1=\tfrac12$, and its square is $\tfrac14$. With Gaussian amplitudes, $\langle(|c_0|^2-|c_1|^2)^2\rangle=2f_0^2-2f_0f_1+2f_1^2=2-1+\tfrac12=\tfrac32$. For fermions $\psi^\dagger\sigma_x\psi=\hat n_0-\hat n_1$, and $:\hat n_0^2:=a_0^\dagger a_0^\dagger a_0a_0=0$, so $\langle:(\hat n_0-\hat n_1)^2:\rangle=-2\langle\hat n_0\hat n_1\rangle=-2f_0f_1=-1$. All three collections have the same density matrix and the same average $\langle\psi^\dagger\sigma_x\psi\rangle=\tfrac12$. \textbf{The density matrix does not determine the average of a square; the statistics and the ensemble do.}

\textbf{Second example.} With the unrotated modes $u_n=e_n$ (so $\rho=\mathrm{diag}(f_0,f_1)$) and the same $X=Y=\sigma_x$, $\mathrm{Tr}(\sigma_x\rho)=0$ and $\mathrm{Tr}(\sigma_x\rho\sigma_x\rho)=2f_0f_1=1$: the Gaussian average is $+1$, the fermion average $-1$, and the fixed-modulus average $+1$ (the diagonal of $\sigma_x$ vanishes). With $X=Y=\mathrm{diag}(1,0)$ instead one gets $2f_0^2=2$ (Gaussian), $f_0^2=1$ (fixed moduli) and $0$ (fermions).

\textbf{What the sign means.} For fermions the exchange term removes the probability of finding two identical particles in the same state at the same place (Section 12.6: two fermions with the same internal label never meet at one point, and exchange removes exactly that part of the Hartree energy), so a repulsive contact interaction costs less than the Hartree estimate. Gaussian waves do the opposite: their intensity fluctuates, $\langle I^2\rangle=2\langle I\rangle^2$ for $I=|c|^2$, so high-intensity places are more likely than for a steady wave, and a repulsive contact interaction costs more. (In optics this ``bunching'' of thermal light is the effect observed by Hanbury Brown and Twiss in 1956; we quote it only as an illustration.)

\subsection{The Hartree--Fock energy of the contact interaction for both fields}

\textbf{The rule of the model.} For the quantized field dirac16complex, Section 13.8 showed that the expectation-value rule makes the two-point function $\langle\Psi_a^\dagger\Psi_b\rangle=(\rho B)_{ba}$, where $\rho=\sum_nw_nu_nu_n^\dagger$ is the Kohn--Sham density matrix of Section 13.7 (orbitals $u_n$ with weights $w_n$). Prescription P3 adopts the same two-point average for dirac16complex00: $\langle\Psi_a^\ast\Psi_b\rangle=K_{ba}$ with

\[
K=\rho B .
\]

In every state of the model each orbital lies in one of the eight blocks of Section 13.4, where $B$ is the number $\beta=js_2=\pm1$ (the \textbf{Krein sign} of the block). So $\rho$ commutes with $B$, and $K=\sum_nw_n\beta_n\,u_nu_n^\dagger$ is Hermitian. (The same holds for the uniform gas used below: its projectors $P_\pm(\mathbf p)$ are built from $\gamma^4$ and the products $\gamma^4\gamma^l$ with $l\le3$, which commute with $B$ by Section 8.11.) The model is therefore the Gaussian ensemble of Section 14.5 with the mean squared amplitude $w_n$ of each orbital replaced by $w_n\beta_n$, and it defines the average of a product of four amplitudes by the same pairing rule with these numbers. Both sides of Theorem 14.4 are polynomials in the occupations, and the proof used nothing but the pairing rule, so the theorem holds with $\rho$ replaced by $K$. (Section 14.8 returns to the fact that $w_n\beta_n$ can be negative, for two independent reasons: $\beta_n=-1$ in four of the eight blocks, and $w_n<0$ on a sea level that is not full.)

\textbf{The average of $S^2$.} Take $X=Y=C$, so that $\Psi^\dagger C\Psi=S$. The cyclic property of the trace gives

\[
\mathrm{Tr}(CK)=\mathrm{Tr}(C\rho B)=\mathrm{Tr}(BC\rho),\qquad \mathrm{Tr}(CKCK)=\mathrm{Tr}(C\rho B\,C\rho B)=\mathrm{Tr}(BC\rho\,BC\rho)
\]

(in the second identity the last factor $B$ is moved to the front). These two identities are check \texttt{PAIR\_sta\allowbreak{}t\_expect\allowbreak{}ationRul\allowbreak{}eTraces} (and \texttt{S5\_stat\_\allowbreak{}expectat\allowbreak{}ionRuleT\allowbreak{}races}), verified on an exact random Hermitian $16\times16$ matrix. With Theorem 14.4:

\textbf{Theorem 14.5 (Hartree--Fock energy for both statistics).} In the model, and for the quantized dirac16complex (Section 13.8),

\[
e_{HF}=\frac\lambda2\bigl\langle S^2\bigr\rangle=\underbrace{\frac\lambda2\,S^2}_{e_H}+\underbrace{\mathrm{sg}\,\frac\lambda2\,\mathrm{Tr}(BC\rho\,BC\rho)}_{e_x},\qquad S=\mathrm{Tr}(BC\rho),
\]

with $\mathrm{sg}=-1$ for dirac16complex and $\mathrm{sg}=+1$ for dirac16complex00. The Hartree term $e_H$ is the same for both fields; only the exchange term $e_x$ changes sign (\texttt{pairing-\allowbreak{}theory.j\allowbreak{}son}, key \texttt{statisti\allowbreak{}cs.hartr\allowbreak{}eeFock.f\allowbreak{}ormula}). For $\mathrm{sg}=-1$ this is the formula of Section 13.8, where the normal-ordered quantum average $\langle{:}S^2{:}\rangle$ appears; the repository also checks it on an exact Krein--Fock space of four moving modes with both Krein signs (check \texttt{S5\_stat\_\allowbreak{}expectat\allowbreak{}ionRuleK\allowbreak{}reinFock}).

\textbf{A filled shell: the ratio $+1/8$.} Fill the eight positive-energy states at rest, $\rho=P_+=\tfrac12(1+BC)$ (Section 13.8, where $BC=-i\gamma^4$ and $BC\,P_+=P_+$). Then $S=\mathrm{Tr}(BCP_+)=8$ and $\mathrm{Tr}(BCP_+BCP_+)=\mathrm{Tr}(P_+)=8=S^2/8$, so

\[
\begin{aligned}
&E_x=\mathrm{sg}\,\frac\lambda2\cdot8=\mathrm{sg}\,\frac18\cdot\frac\lambda2\cdot64=\mathrm{sg}\,\frac{E_H}8:\\
&E_x=-\frac{E_H}8\ \ (\text{dirac16complex}),\qquad E_x=+\frac{E_H}8\ \ (\text{dirac16complex00}).
\end{aligned}
\]

The ratio $1/8$ holds at every momentum (check \texttt{PAIR\_sta\allowbreak{}t\_filled\allowbreak{}ShellExc\allowbreak{}hangeRat\allowbreak{}io}, symbolic in $m$ and $\mathbf p$). \textbf{Worked example} at an exact point of the good sector recorded in \texttt{python-p\allowbreak{}airing-r\allowbreak{}eport.js\allowbreak{}on} (entry \texttt{S5\_stat.\allowbreak{}filledSh\allowbreak{}ellExcha\allowbreak{}ngeRatio}): $m=2$ and $\mathbf p=(0,1,2,4)$ along $(x_0,x_1,x_2,x_3)$, so $E_p^2=4+0+1+4+16=25$ and $E_p=5$. The report gives $\mathrm{Tr}(BCP_+)=\tfrac{16}5$, which is $8m/E_p=\tfrac{16}5$ (Section 13.8), and $\mathrm{Tr}(BCP_+BCP_+)=\tfrac{32}{25}$, which is $\bigl(\tfrac{16}5\bigr)^2/8=\tfrac{256}{200}=\tfrac{32}{25}$.

\textbf{The uniform gas.} Section 13.8 considered the uniform gas of particles and antiparticles with the density matrix $\rho=\int\frac{d^dp}{(2\pi)^d}\bigl[f_+P_+(\mathbf p)-f_-P_-(\mathbf p)\bigr]$, where $d$ is the number of \textbf{space dimensions} of the momenta $\mathbf p$ (as in Section 13.8; here $d$ is not the number of spinor components of Section 14.5), and with occupations that depend only on $|\mathbf p|$. It computed, for every $d$,

\[
\mathrm{Tr}(BC\rho\,BC\rho)=\frac{n^2+S^2}{16}.
\]

No step of that computation used the statistics: it is a statement about traces of projectors. (The Python checker notes that isotropy is more than needed: the term with $\mathbf p\cdot\mathbf q$ already vanishes when the occupations are unchanged under $\mathbf p\to-\mathbf p$; \texttt{S5\_stat.\allowbreak{}uniformG\allowbreak{}asExchan\allowbreak{}ge} in \texttt{python-p\allowbreak{}airing-r\allowbreak{}eport.js\allowbreak{}on}.) With Theorem 14.5:

\[
\begin{aligned}
&e_x=\mathrm{sg}\,\frac\lambda{32}\bigl(n^2+S^2\bigr):\\
&e_x=-\frac\lambda{32}\bigl(n^2+S^2\bigr)\ \ (\text{dirac16complex}),\qquad e_x^{00}=+\frac\lambda{32}\bigl(n^2+S^2\bigr)\ \ (\text{dirac16complex00}),
\end{aligned}
\]

exactly, at every temperature. This is the check \texttt{PAIR\_sta\allowbreak{}t\_unifor\allowbreak{}mGasExch\allowbreak{}ange} of the Wolfram report and the check \texttt{S5\_stat\_\allowbreak{}uniformG\allowbreak{}asExchan\allowbreak{}ge} of the Python report. The Python checker also verifies the trace identity on an exact \textbf{finite} gas: $m=2$, eleven momenta with four components along $(x_0,x_1,x_2,x_3)$ (so $d=4$), on the four shells $E=2,3,4,5$ (the shell $E=2$ is $\mathbf p=0$; each other momentum comes with $-\mathbf p$), with the occupations $f_+=1,\tfrac34,\tfrac13,\tfrac1{10}$ and $f_-=\tfrac17,\tfrac1{11},0,\tfrac1{13}$ on the four shells. It finds $n=\tfrac{510808}{15015}\approx34.0198$, $S=\tfrac{2403416}{75075}\approx32.0135$ and $\mathrm{Tr}(BC\rho BC\rho)=\tfrac{768720549416}{5636255625}\approx136.3885=(n^2+S^2)/16$ exactly (check \texttt{S5\_stat\_\allowbreak{}uniformG\allowbreak{}asExactF\allowbreak{}inite}; Exercise 14.6 recomputes $n$ and $S$ by hand). For dirac16complex00 this gas has $e_x^{00}=\tfrac\lambda2\cdot136.3885=68.194\,\lambda$.

\textbf{Two limits.} A gas at rest has $S=n$ (every particle contributes $m/E_p=1$), so $e_x=\mathrm{sg}\,\tfrac\lambda{16}n^2=\mathrm{sg}\,e_H/8$, the filled-shell value (the limit is part of check \texttt{PAIR\_sta\allowbreak{}t\_T0Tota\allowbreak{}lDerivat\allowbreak{}ive}). A fast gas has $|S|\ll n$, and then the exchange term $\pm\tfrac\lambda{32}n^2$ can be much larger than the Hartree term $\tfrac\lambda2S^2$; Section 14.10 meets exactly this situation.

\subsection{The Kohn--Sham equations of dirac16complex00}

\textbf{The frame is that of Chapter 13.} The static primordial field in the warped coordinate $y$ with the Z2 brane (Section 13.2), the good sector and the stationary ansatz with the factor $W^{-3}$ (Section 13.3), the eight $2\times2$ blocks (Section 13.4), the block equation, its Hamiltonian form, the conserved $y$-current and the boundary conditions (Section 13.5), and the torus and the bookkeeping of particles and sea (Section 13.7) are statements about the linear Dirac operator in the fixed field. They hold word for word for commuting components. The reduced equation and the block forms were nevertheless derived again for dirac16complex00 and verified in both of its reports: see the check \texttt{C00\_stat\allowbreak{}ic\_geome\allowbreak{}tryAndRe\allowbreak{}ducedEqu\allowbreak{}ation} and the check \texttt{C00\_stat\allowbreak{}ic\_block\allowbreak{}sFromSta\allowbreak{}ge4Theor\allowbreak{}y} in the Wolfram report, and the two checks with the same names after the prefix \texttt{\detokenize{S5_}} in the Python report.

\textbf{The local potentials.} Evaluate $e_x=\mathrm{sg}\,\tfrac\lambda{32}(n^2+S^2)$ with the local proper densities $n_p(y)$ and $S_p(y)$ (the local density approximation of Section 12.12). Its derivatives are the exchange potentials

\[
v_v=\frac{\partial e_x}{\partial n}=\mathrm{sg}\,\frac\lambda{16}\,n,\qquad v_s=\frac{\partial e_x}{\partial S}=\mathrm{sg}\,\frac\lambda{16}\,S,
\]

and the effective mass is $M_{\text{eff}}=m+\lambda S_p+v_s=m+\bigl(1+\tfrac{\mathrm{sg}}{16}\bigr)\lambda S_p$ (checks \texttt{PAIR\_sta\allowbreak{}t\_ldaPot\allowbreak{}entials}, \texttt{PAIR\_T3b\allowbreak{}lock\_sta\allowbreak{}tisticsC\allowbreak{}oefficie\allowbreak{}nts}, \texttt{S5\_stat\_\allowbreak{}ldaPoten\allowbreak{}tials}). Side by side:

\begingroup
\small
\begin{longtable}{@{}>{\raggedright\arraybackslash}p{0.267\linewidth}>{\raggedright\arraybackslash}p{0.267\linewidth}>{\raggedright\arraybackslash}p{0.267\linewidth}@{}}
\toprule
\textbf{quantity} & \textbf{dirac16complex ($\mathrm{sg}=-1$)} & \textbf{dirac16complex00 ($\mathrm{sg}=+1$)} \\
\midrule
\endfirsthead
\toprule
\textbf{quantity} & \textbf{dirac16complex ($\mathrm{sg}=-1$)} & \textbf{dirac16complex00 ($\mathrm{sg}=+1$)} \\
\midrule
\endhead
exchange energy density $e_x$ & $-\tfrac\lambda{32}(n^2+S^2)$ & $+\tfrac\lambda{32}(n^2+S^2)$ \\
filled shell $E_x/E_H$ & $-\tfrac18$ & $+\tfrac18$ \\
vector potential $v=v_v$ & $-\tfrac\lambda{16}n_p$ & $+\tfrac\lambda{16}n_p$ \\
scalar potential $v_s$ & $-\tfrac\lambda{16}S_p$ & $+\tfrac\lambda{16}S_p$ \\
effective mass $M_{\text{eff}}$ & $m+\tfrac{15}{16}\lambda S_p$ & $m+\tfrac{17}{16}\lambda S_p$ \\
\bottomrule
\end{longtable}
\endgroup

The pair $(M_{\text{eff}}^{00},v^{00})$ is the \textbf{Kohn--Sham fermion-gas thermodynamic pseudo-potential} of dirac16complex00: the name that section 3 of \texttt{handoff/\allowbreak{}specs/ST\allowbreak{}AGE4\_SPE\allowbreak{}C.md} gives to the pair $(M_{\text{eff}},v)$ and that Chapter 13 uses (the Stage-5 specification says ``effective potential'').

\textbf{The functional.} In the notation of Section 13.9 (orbitals $\chi_n$ with block type $j_n$ and momentum $k_n$, normalized by $\int\chi_n^\dagger\chi_n\,dy=1$, occupations $f_n$, proper volume element $dV=e^{6Hy}\ell^3dy$, the Pauli matrices $\sigma_x,\sigma_y,\sigma_z$ of Section 2.3, and the momentum weight $\xi(y)=e^{-Hy-a_{4,0}}$ of Section 13.3), the Mermin functional of dirac16complex00 is

\[
\begin{aligned}
F^{00}&=\sum_nf_n\langle\chi_n|h_0|\chi_n\rangle-T\,S_{\text{ent}}+E_H+E_x^{00},\qquad h_0=j\Bigl[-i\sigma_x\frac d{dy}+m\sigma_y+\xi k\sigma_z\Bigr],\\
E_H&=\frac\lambda2\int S_p^2\,dV,\qquad E_x^{00}=+\frac\lambda{32}\int\bigl(n_p^2+S_p^2\bigr)\,dV,
\end{aligned}
\]

with $n_p=e^{-6Hy}\sum_nf_n\chi_n^\dagger\chi_n/\ell^3$ and $S_p=e^{-6Hy}\sum_nf_n\chi_n^\dagger(j_n\sigma_y)\chi_n/\ell^3$. (As in Section 13.9, in the solver these sums, and the one-body sum of $F^{00}$, run over the particle and the sea levels with the weights $w_n$ of Section 13.7 in place of $f_n$: $w_n=f_n$ on a particle level and $w_n=-(1-f_n)$ on a sea level, prescription P2. This is why the total energy below is written with $w_n$.) It differs from the functional of Section 13.9 only in the sign of the exchange term. (The theory files and the Rust code call the momentum weight \texttt{\detokenize{kappa}}, and the check names write the Pauli matrices as sigma1, sigma2, sigma3; in this book $\kappa$ is the gravitational coupling of Section 13.2, as in Section 14.10.)

\textbf{Stationarity.} Vary $\chi_n^\dagger(y)$ as in Section 13.9. The derivative of $\int n_p^2\,dV$ is $2n_p(y)\cdot e^{6Hy}\ell^3\cdot e^{-6Hy}f_n\chi_n(y)/\ell^3=2f_nn_p\chi_n$: the factor $e^{6Hy}$ of the volume element cancels the factor $e^{-6Hy}$ of the density. In the same way $\int S_p^2\,dV$ gives $2f_nS_p\,j_n\sigma_y\chi_n$. Hence $E_x^{00}$ contributes $\tfrac\lambda{16}f_n\bigl(n_p+S_pj_n\sigma_y\bigr)\chi_n=f_n\bigl(v_v+v_sj_n\sigma_y\bigr)\chi_n$ and $E_H$ contributes $f_n\lambda S_pj_n\sigma_y\chi_n$. Dividing by $f_n$, the Kohn--Sham equation of dirac16complex00 is $h_j^{00}\chi_n=\varepsilon_n\chi_n$ with

\[
h_j^{00}=j\Bigl[-i\sigma_x\frac d{dy}+\Bigl(m+\tfrac{17}{16}\lambda S_p(y)\Bigr)\sigma_y+\xi(y)\,k\,\sigma_z\Bigr]+\frac\lambda{16}\,n_p(y),
\]

the Hamiltonian form $h_j$ of Section 13.5 with the new potentials (check \texttt{PAIR\_T3k\allowbreak{}s\_statio\allowbreak{}narityBo\allowbreak{}thStatis\allowbreak{}tics}, which carries the statistics sign as a symbol). In the real variables $\chi=(a,ib)$ of Section 13.5 it reads $a'=M_{\text{eff}}a-(jE+\xi k)b$ and $b'=(jE-\xi k)a-M_{\text{eff}}b$ with $M_{\text{eff}}=m+\tfrac{17}{16}\lambda S_p$ and $E=\varepsilon-\tfrac\lambda{16}n_p$. Stationarity with respect to the occupations gives the Fermi--Dirac occupations exactly as in Section 13.9: in this model they follow from the entropy term that prescription P2 puts into the functional, not from any property of the classical field.

\textbf{Total energy.} Both interaction functionals are quadratic in the densities, so the double-counting argument of Section 13.9 applies unchanged:

\[
E=\sum_nw_n\varepsilon_n-E_H-E_x^{00},\qquad F=E-TS_{\text{ent}},\qquad \frac{dF}{d\lambda}=\frac{E_H+E_x^{00}}\lambda .
\]

\textbf{Worked example.} To see the two sets of formulas at work, take at one point the numbers $\lambda=16$, $n=\tfrac12$ and $S=-\tfrac14$ (chosen for easy arithmetic, not taken from a run; $H=1$). Then $e_H=8\cdot\tfrac1{16}=\tfrac12$ and $\tfrac\lambda{32}(n^2+S^2)=\tfrac12\bigl(\tfrac14+\tfrac1{16}\bigr)=\tfrac5{32}$, so $e_x=-\tfrac5{32}$ for dirac16complex and $+\tfrac5{32}$ for dirac16complex00. The vector potentials are $\mp\tfrac{16}{16}\cdot\tfrac12=\mp\tfrac12$, the scalar potentials $\mp\tfrac{16}{16}\cdot(-\tfrac14)=\pm\tfrac14$, and the mass shifts $\lambda S+v_s=-4+\tfrac14=-\tfrac{15}4$ (dirac16complex, which is $\tfrac{15}{16}\cdot16\cdot(-\tfrac14)$) and $-4-\tfrac14=-\tfrac{17}4$ (dirac16complex00, which is $\tfrac{17}{16}\cdot16\cdot(-\tfrac14)$).

\textbf{The first-order effect of the sign.} By the Hellmann--Feynman relation just stated, $dE/d\lambda$ at $\lambda=0$ is the interaction energy per unit $\lambda$ evaluated with the densities of the free ($\lambda=0$) state (the envelope argument of Section 13.9: the change of the orbitals does not contribute at first order, because $F$ is stationary). Hence

\[
E(\lambda)=E(0)+E_H^{\text{free}}+\mathrm{sg}\,X^{\text{free}}+O(\lambda^2),\qquad X^{\text{free}}=\frac\lambda{32}\int\bigl(n_p^2+S_p^2\bigr)\,dV\Big|_{\lambda=0},
\]

with $X^{\text{free}}\ge0$ for $\lambda>0$. So at first order the energies of the two fields differ by $2X^{\text{free}}$, dirac16complex00 lying higher for a repulsive coupling. The orbital energies show the same sign: for $\lambda>0$ the exchange potential $v^{00}=+\tfrac\lambda{16}n_p$ is positive and shifts the particle levels up, whereas $v=-\tfrac\lambda{16}n_p$ shifts them down for dirac16complex; the mass shift, $\tfrac{17}{16}\lambda S_p$ against $\tfrac{15}{16}\lambda S_p$, is nearly the same for both fields. Section 14.10 confirms both statements numerically.

\subsection{What the model is, and what it is not}

The derivations of Sections 14.3 to 14.7 are exact. What they describe is fixed by the prescriptions P1 to P4. Two exact facts, proved below, limit what the resulting numbers mean, and a third point says what the model is not; it is a statement about how the model is defined, not a theorem.

\textbf{First fact: the covariance of the model is not positive.} For any collection of classical fields, whatever its probability distribution, the \textbf{covariance} matrix $K_{ab}=\langle\Psi_a\Psi_b^\ast\rangle$ is positive semidefinite: for every vector $v$ the number $v^\dagger\Psi=\sum_av_a^\ast\Psi_a$ has $|v^\dagger\Psi|^2=\sum_{a,b}v_a^\ast\Psi_a\Psi_b^\ast v_b$, whose average is

\[
\bigl\langle|v^\dagger\Psi|^2\bigr\rangle=\sum_{a,b}v_a^\ast K_{ab}v_b=v^\dagger Kv\ \ge\ 0,
\]

because the average of a non-negative quantity is not negative. Under prescription P3 the covariance is $K=\rho B$ (Section 14.6). Take the filled shell at rest, $\rho=P_+=\tfrac12(1-i\gamma^4)$. $P_+$ and $B$ are Hermitian and commute ($\gamma^4$ commutes with $B$, Section 8.11), so they have a common orthonormal basis of eigenvectors. On the eight-dimensional range of $P_+$ (the positive-energy states at rest) the Krein form $u^\dagger Bu$ has the signature (4,4) (Section 8.8): $B$ has the eigenvalue $+1$ four times and $-1$ four times there. Hence

\[
K=P_+B\ \text{ has the eigenvalues }\ +1\ (4\text{ times}),\quad -1\ (4\text{ times}),\quad 0\ (8\text{ times}),
\]

as the Wolfram verifier records (measurement \texttt{stat\_res\allowbreak{}tShellCo\allowbreak{}variance\allowbreak{}Eigenval\allowbreak{}ues} and check \texttt{PAIR\_sta\allowbreak{}t\_expect\allowbreak{}ationRul\allowbreak{}eCovaria\allowbreak{}nceIndef\allowbreak{}inite}, \texttt{wolfram-\allowbreak{}pairing-\allowbreak{}report.j\allowbreak{}son}). An explicit vector: the rest state $u_-=\tfrac12(e_0+e_4+i\,e_9-i\,e_{13})$ of Section 8.8 has $-i\gamma^4u_-=u_-$, so $P_+u_-=u_-$, and $Bu_-=-u_-$, so $Ku_-=-u_-$ and $u_-^\dagger Ku_-=-1$. The ``average of $|u_-^\dagger\Psi|^2$'' would be $-1$. \textbf{No probability distribution of classical fields has this covariance.} For the orbitals of the four blocks with Krein sign $js_2=-1$ the model is not an average over classical field configurations; it is a formal Gaussian functional whose covariance carries the Krein signs, which \texttt{pairing-\allowbreak{}theory.j\allowbreak{}son} (key \texttt{statisti\allowbreak{}cs.posit\allowbreak{}ivity.co\allowbreak{}nsequenc\allowbreak{}e}) calls a ``FORMAL (Krein-signed) Gaussian functional''.

\textbf{A second, independent source of negative variances.} In the model the mean squared amplitude of the orbital $u_n$ is $w_n\beta_n$, the eigenvalue of $K=\sum_nw_n\beta_nu_nu_n^\dagger$ on $u_n$ (Section 14.6). The Krein sign $\beta_n$ is one factor; the weight $w_n$ of prescription P2 is the other. On a sea level $w_n=-(1-f_n)$, which is negative whenever the level is not completely full, and at a temperature $T>0$ the Fermi--Dirac occupation $f_n=1/(e^{(\varepsilon_n-\mu)/T}+1)$ is smaller than 1 for every level, so every sea level carries a thermal hole of weight $1-f_n>0$ (tiny for deep levels, but not zero). The four possible cases are:

\begingroup
\small
\begin{longtable}{@{}>{\raggedright\arraybackslash}p{0.267\linewidth}>{\raggedright\arraybackslash}p{0.267\linewidth}>{\raggedright\arraybackslash}p{0.267\linewidth}@{}}
\toprule
\textbf{level} & \textbf{block with $\beta_n=+1$} & \textbf{block with $\beta_n=-1$} \\
\midrule
\endfirsthead
\toprule
\textbf{level} & \textbf{block with $\beta_n=+1$} & \textbf{block with $\beta_n=-1$} \\
\midrule
\endhead
particle, $w_n=f_n\ge0$ & $w_n\beta_n\ge0$ & $w_n\beta_n\le0$ \\
sea, $w_n=-(1-f_n)\le0$ & $w_n\beta_n\le0$ & $w_n\beta_n\ge0$ \\
\bottomrule
\end{longtable}
\endgroup

So at $T>0$ negative variances occur in the blocks of both Krein signs: from the occupied particle levels in the blocks with $\beta_n=-1$, and from the thermal holes of the sea in the blocks with $\beta_n=+1$. This applies to the runs at $T=0.3H$ of Section 14.10. At $T=0$ every sea level is full, $w_n=0$ there, and only the Krein sign remains as a source of negative variances.

\textbf{What the honest alternative gives.} An honest Gaussian ensemble of classical waves filling the rest shell has the covariance $K=\rho=P_+$. Its charge density is $\mathrm{Tr}(BP_+)=\tfrac12\mathrm{Tr}B-\tfrac i2\mathrm{Tr}(B\gamma^4)=0$, because $\mathrm{Tr}B=0$ and $B\gamma^4=-iC(\gamma^4)^2=iC$ with $\mathrm{Tr}C=0$; and its scalar density is $\mathrm{Tr}(CP_+)=\tfrac12\mathrm{Tr}C-\tfrac i2\mathrm{Tr}(C\gamma^4)=0$, because a product of four or of five different gammas is traceless (Chapter 2). The four modes with Krein sign $+1$ and the four with $-1$ cancel in both densities. The expectation-value rule instead gives $\mathrm{Tr}(B\,BP_+)=\mathrm{Tr}P_+=8$ and $\mathrm{Tr}(BC\,P_+)=8$. The Stage-5 model keeps the rule, so that both fields carry the same densities as functions of their orbitals; \texttt{pairing-\allowbreak{}theory.j\allowbreak{}son} records this as ``a modelling choice, stated as such'' (check \texttt{PAIR\_sta\allowbreak{}t\_expect\allowbreak{}ationRul\allowbreak{}eCovaria\allowbreak{}nceIndef\allowbreak{}inite} contains the two zero traces).

\textbf{Second fact: the model's energies are not the classical energies.} In a block of type $(j,s_2)$ the matrix $B$ is the number $js_2$ (Section 13.4). So a field $\Psi$ whose sixteen components lie in one block satisfies $B\Psi=js_2\Psi$ at every point, and, because $B$ is Hermitian, $\Psi^\dagger B=(B\Psi)^\dagger=js_2\,\Psi^\dagger$: the row $\Psi^\dagger B$ is the row $\Psi^\dagger$ multiplied by the number $js_2$. The expectation-value rule of the Kohn--Sham model replaces $\Psi^\dagger$ by $\Psi^\dagger B$ (Section 13.7), that is, for a single-block mode, by $js_2\,\Psi^\dagger$, and it does so everywhere, also inside derivatives of $\Psi^\dagger$. Every term of the one-body energy--momentum tensor (the part without the interaction, Section 13.14) contains exactly one factor built from $\Psi^\dagger$ (the adjoint $\bar\Psi=\Psi^\dagger C$ or its derivative $D_\mu\bar\Psi$) and one built from $\Psi$. So every term, and with it every component, is multiplied by the same number $js_2$, and since $(js_2)^2=1$,

\[
T_{\mu\nu}^{\text{classical}}=js_2\;T_{\mu\nu}^{\text{KS one-body}}
\]

for every single-block mode, component by component, and even for fields that do not solve the field equation. The energy density shows it most simply. By Section 13.14 its one-body part is $\varepsilon\,\Psi^\dagger B\Psi$. For a classical single-block field this is $js_2\,\varepsilon\,\Psi^\dagger\Psi$, while the Kohn--Sham orbital contributes $\varepsilon\,u^\dagger B^2u=\varepsilon\,n$. The Wolfram check of the general identity, \texttt{C00\_stat\allowbreak{}ic\_class\allowbreak{}icalMode\allowbreak{}IsKreinW\allowbreak{}eightedK\allowbreak{}S}, compares all 64 components of the tensor in each of the eight blocks. Its twin in the Python report is \texttt{S5\_stati\allowbreak{}c\_classi\allowbreak{}calModeI\allowbreak{}sKreinWe\allowbreak{}ightedKS}. The statement is recorded in \texttt{dirac16c\allowbreak{}omplex00\allowbreak{}-theory.\allowbreak{}json} under the key \texttt{static.c\allowbreak{}lassical\allowbreak{}VersusKo\allowbreak{}hnSham}.

\textbf{A worked example: two box modes.} The report computes one explicit mode in a block of each Krein sign (check \texttt{C00\_stat\allowbreak{}ic\_class\allowbreak{}icalEner\allowbreak{}gyKreinS\allowbreak{}igned}, and \texttt{S5\_stati\allowbreak{}c\_classi\allowbreak{}calEnerg\allowbreak{}yKreinSi\allowbreak{}gned} in the Python report; \texttt{dirac16c\allowbreak{}omplex00\allowbreak{}-theory.\allowbreak{}json}, key \texttt{static.b\allowbreak{}oxModes}). Take the free problem ($\lambda=0$, so $M_{\text{eff}}=m$ and $v=0$) with $m=3$ at zero momentum $k=0$ in a box of length $L=\pi/4$, and the first massive level of parity $+$ of Section 13.6: $p=\pi/L=4$, $\varepsilon=\sqrt{m^2+p^2}=\sqrt{9+16}=5$ and

\[
\chi_2=\sin4y,\qquad \chi_1=\frac{3\sin4y+4\cos4y}{5ij},
\]

with $\chi_2(0)=\sin0=0$ (parity $+$) and $\chi_2(-\pi/4)=\sin(-\pi)=0$ (the bag condition at the end of the box). The check puts this orbital into block 0 ($js_2=+1$) and into block 2 ($js_2=-1$) as the sixteen-component column $\chi=\chi_1v_++\chi_2v_-$, with the unnormalized basis vectors $v_\pm$ of Section 13.4 ($|v_\pm|^2=8$), verifies the full sixteen-component reduced equation $\gamma^0\chi'-5i\gamma^4\chi=3\chi$ and the two boundary conditions, and computes the classical energy density. At $\lambda=0$ the energy density is $\varepsilon\,\Psi^\dagger B\Psi=js_2\,\varepsilon\,\Psi^\dagger\Psi$, and $W^6\Psi^\dagger\Psi=\chi^\dagger\chi=8\bigl(|\chi_1|^2+|\chi_2|^2\bigr)$ (Section 13.3, with $W^6=\sqrt{|g|}$), so

\[
\rho\,W^6=js_2\cdot5\cdot8\Bigl[\sin^24y+\tfrac1{25}\bigl(3\sin4y+4\cos4y\bigr)^2\Bigr]=js_2\cdot\tfrac85\bigl(25-9\cos8y+12\sin8y\bigr)
\]

(Exercise 14.15 does the trigonometry). The report records exactly this profile for block 0 and its negative for block 2. Since $-9\cos8y+12\sin8y$ never goes below $-\sqrt{9^2+12^2}=-15$ (Exercise 14.15), the bracket is at least $10$: the classical energy density is positive everywhere in block 0 and negative everywhere in block 2. The Kohn--Sham model gives both orbitals the positive energy $\varepsilon=5$. For the four blocks of Krein sign $-1$ a positive Kohn--Sham orbital energy therefore corresponds to a classical wave of negative energy: the Kohn--Sham energies of dirac16complex00 are Krein-weighted energies, not the energies of classical field configurations. (The example also shows that the classical energy of dirac16complex00 has no lower bound in the static primordial field either: multiply the mode of block 2 by a large number.)

\textbf{Third point: it is not a quantum theory.} This is a statement about how the model is defined, not a theorem. One might hope to reinterpret the model as the quantization of dirac16complex00 with commutators instead of anticommutators. The model is not obtained that way: the Stage-5 specification does not attempt such a quantization, and neither does this chapter. The reason for not attempting it is a theorem that we quote and do not prove: in ordinary space-time with one time, Pauli's spin--statistics theorem (1940) shows that a spinor field quantized with commutators has either an energy without lower bound or states of negative norm. No such theorem is proved for signature (4,4) in this book (Chapter 8).

\textbf{The words ``ground state'' and ``excited state''.} In this chapter the \textbf{ground state} of dirac16complex00 means, by definition, the self-consistent solution of the Kohn--Sham equations of Section 14.7 reached by continuation in the coupling from $\lambda=0$ (erratum E4.8, the definition of Chapter 13). It is not the minimum of any energy of the classical field, which has none (Section 14.2). The \textbf{first excited state} means the three quantities of Section 12.15 computed in this model: the Kohn--Sham gap, the particle--hole list and the Delta-SCF energy.

\textbf{No correlation term.} For the quantized dirac16complex, Section 13.8 argued that a contact coupling of mass dimension $-6$ leaves no well-defined correlation energy to approximate. For the classical ensemble of the model the analogue of correlation would be the departure of the interacting ensemble from a Gaussian one; it is not computed either. The model is exchange-only for both fields.

\textbf{In one paragraph.} The dirac16complex00 Kohn--Sham model is a mean-field functional built with the averaging rule of a Gaussian (random-phase) ensemble of classical modes, with the weights of a fermion gas imposed (including the negative weights of sea levels that are not full) and with the Krein-signed densities of the quantized field adopted by prescription; because of these negative numbers it is a formal Gaussian functional, not an average over classical fields. It is built so that the two fields differ in exactly one respect: the sign of the exchange term, which Theorems 14.4 and 14.5 derive. Its ``ground state'' is a definition, not a minimum. Its energies are not the energies of classical field configurations. It is not a quantum theory of commuting spinors.

\subsection{How the model is solved, and what was computed}

\textbf{Two solvers.} The equations of Section 14.7 are those of Chapter 13 with two coefficients changed, $\tfrac{15}{16}\to\tfrac{17}{16}$ in the effective mass and $-\tfrac1{16}\to+\tfrac1{16}$ in the vector potential. Both Stage-4 solvers received an option for the statistics sign in Stage 5. The Rust solver \texttt{studies/\allowbreak{}dirac16c\allowbreak{}omplex\_k\allowbreak{}ohn\_sham}, which integrates the block equation by shooting (Section 13.10), has the subcommand \texttt{\detokenize{pairs}} (module \texttt{src/pair\allowbreak{}s.rs}, with the statistics \texttt{Commutin\allowbreak{}g} for dirac16complex00). The independent reference solver \texttt{scripts/\allowbreak{}ks\_refer\allowbreak{}ence\_sol\allowbreak{}ver.py}, driven by \texttt{scripts/\allowbreak{}ks\_refer\allowbreak{}ence\_pai\allowbreak{}rs.py} (field label \texttt{\detokenize{d16c00}}), uses a matrix method instead of shooting: it places one component of each block orbital on the nodes and the other on the half nodes of a uniform grid in $y$ (a staggered grid), turns the Hermitian block operator into a real symmetric matrix, diagonalizes it, solves every problem on three grids of $N_0$, $2N_0$ and $4N_0$ intervals ($N_0=120$ for $m=3$), and extrapolates to zero grid spacing by eliminating the error terms proportional to $h^2$ and $h^3$ (the header of \texttt{scripts/\allowbreak{}ks\_refer\allowbreak{}ence\_sol\allowbreak{}ver.py}). Without interaction the statistics does not enter at all, and the two fields are the same numerical problem: the reference runs of the two fields at $\lambda=0$ agree in every printed digit, and the Rust module states the same (``The two statistics at lambda = 0 are the same problem (bit for bit)'', header of \texttt{src/pair\allowbreak{}s.rs}).

\textbf{The couplings.} dirac16complex00 uses the same numerical $\hat\lambda=\lambda m^6$ as dirac16complex: the Stage-4 reference coupling of the configuration $(|m|,L,N)$, calibrated so that the first-order pseudo-potential of dirac16complex reaches $0.1\,m$ ($\hat\lambda_1$) or $1\,m$ ($\hat\lambda_2=10\hat\lambda_1$) (Section 13.10). For $(m,L,N)=(3,3,112)$ the file \texttt{artifact\allowbreak{}s/dirac1\allowbreak{}6complex\allowbreak{}/kohn-sh\allowbreak{}am/refer\allowbreak{}ence/cou\allowbreak{}plings/c\allowbreak{}oupling-\allowbreak{}m3\_L3\_N1\allowbreak{}12.json} gives $\hat\lambda_1=852.6397026$, that is $\lambda=\hat\lambda_1/m^6=852.6397026/729=1.169602$ in units of $H^{-6}$ (the coupling has the mass dimension $-6$, Section 13.8).

\textbf{What was planned.} The Stage-5 run matrix (the header of \texttt{scripts/\allowbreak{}ks\_refer\allowbreak{}ence\_pai\allowbreak{}rs.py}) is $(|m|,N)\in\{(1,8),(1,112),(3,8),(3,112)\}$ with $L=3$, $\Delta k=0.25\,|m|$, $\hat\lambda\in\{0,\pm\hat\lambda_1,\pm\hat\lambda_2\}$, $T=0$ for every configuration and $T=0.1\,|m|$ for $N=112$, both fields, and three universes each: $+M$ (the problem of Chapter 13), $-M$ with the transformed boundary conditions, and $-M$ with the untransformed ones (the two $-M$ universes belong to the pairing theorems of Chapter 15).

\textbf{What exists for dirac16complex00.} The committed summary \texttt{artifact\allowbreak{}s/dirac1\allowbreak{}6complex\allowbreak{}/pair-cr\allowbreak{}eation/r\allowbreak{}eference\allowbreak{}/referen\allowbreak{}ce-pairs\allowbreak{}-summary\allowbreak{}.json} (key \texttt{\detokenize{complete}}: \texttt{\detokenize{false}}) contains, for dirac16complex00 with $m=3$ and $N=112$, the following (``control'' is the $-M$ universe with the untransformed boundary conditions):

\begingroup
\small
\begin{longtable}{@{}>{\raggedright\arraybackslash}p{0.200\linewidth}>{\raggedright\arraybackslash}p{0.200\linewidth}>{\raggedright\arraybackslash}p{0.200\linewidth}>{\raggedright\arraybackslash}p{0.200\linewidth}@{}}
\toprule
\textbf{coupling, temperature} & \textbf{$+M$} & \textbf{$-M$} & \textbf{control} \\
\midrule
\endfirsthead
\toprule
\textbf{coupling, temperature} & \textbf{$+M$} & \textbf{$-M$} & \textbf{control} \\
\midrule
\endhead
$\lambda=0$, $T=0$ & converged & converged & converged \\
$\lambda=0$, $T=0.3H$ & converged & converged & converged \\
$+\hat\lambda_1$, $T=0$ & converged & pending & failed \\
$+\hat\lambda_2$, $T=0$ & pending & failed & failed \\
$-\hat\lambda_1$, $T=0$ & pending & pending & pending \\
$-\hat\lambda_2$, $T=0$ & pending & pending & pending \\
$\lambda\ne0$, $T=0.3H$ & not attempted & not attempted & not attempted \\
\bottomrule
\end{longtable}
\endgroup

For $m=3$, $N=8$ and for $m=1$ ($N=8$ and $N=112$) every dirac16complex00 run is pending. ``Pending'' means never run; the three failed runs stopped because the energy window of the self-consistent loop grew beyond the cap of $20\,\lvert m\rvert$ that the driver imposes. The keys \texttt{\detokenize{failed}}, \texttt{notAttem\allowbreak{}pted} and \texttt{\detokenize{pending}} of the summary list these runs, and the reason of each failure is the key \texttt{\detokenize{failed}} of its run record.

For the interacting runs at $T=0.3H$ the driver estimates the first-order pseudo-potential from the free thermal state, $|\hat\lambda_1|$ times its strength, and finds $84.1\,|m|$ (key \texttt{thermoFi\allowbreak{}rstOrder}): the free thermal gas at $T=0.1\,|m|$ has a much larger density near the tip than the brane-localized $T=0$ state that calibrates $\hat\lambda_1$, so the coupling calibrated at $T=0$ is far outside the window of Chapter 13 there. The committed Rust \texttt{\detokenize{pairs}} outputs in \texttt{artifact\allowbreak{}s/dirac1\allowbreak{}6complex\allowbreak{}/pair-cr\allowbreak{}eation/r\allowbreak{}ust/pair\allowbreak{}s/} contain only dirac16complex configurations and no \texttt{summary.\allowbreak{}json}; the checker \texttt{scripts/\allowbreak{}check\_di\allowbreak{}rac16com\allowbreak{}plex\_pai\allowbreak{}rs.py} has no committed report. \textbf{So every number of dirac16complex00 in this chapter comes from one solver, without an independent cross-check.}

\subsection{Ground and first excited states of both fields side by side}

The only configuration in which both fields are computed with interaction is $m=3H$, $L=3/H$, $N=112$, $a_{4,0}=0$, $\Delta k=0.25\,m=0.75H$, both brane parities, $T=0$. In the tables below, as in the labels of the runs, \textbf{d16c} stands for dirac16complex and \textbf{d16c00} for dirac16complex00. The numbers are those of the following run records of \texttt{referenc\allowbreak{}e-pairs-\allowbreak{}summary.\allowbreak{}json}:

\begin{Verbatim}[fontsize=\small]
d16c_m3_L3_N112_lam0_T0/plusM       the free state (d16c00 has the identical record)
d16c_m3_L3_N112_lamp1_T0/plusM      dirac16complex at +lambda_hat_1
d16c00_m3_L3_N112_lamp1_T0/plusM    dirac16complex00 at +lambda_hat_1
\end{Verbatim}

Energies are in units of $H$; $\hat\lambda/\hat\lambda_1=0$ is the free state, which is the same for both fields.

\textbf{Ground states} (keys \texttt{extrapol\allowbreak{}ated.tot\allowbreak{}al}, \texttt{\detokenize{ksSum}}, \texttt{\detokenize{hartree}}, \texttt{\detokenize{exchange}} and \texttt{\detokenize{mu}} of each run record; at $T=0$ the chemical potential is $\varepsilon_{\text{HOMO}}$):

\begingroup
\small
\begin{longtable}{@{}>{\raggedright\arraybackslash}p{0.114\linewidth}>{\raggedright\arraybackslash}p{0.114\linewidth}>{\raggedright\arraybackslash}p{0.114\linewidth}>{\raggedright\arraybackslash}p{0.114\linewidth}>{\raggedright\arraybackslash}p{0.114\linewidth}>{\raggedright\arraybackslash}p{0.114\linewidth}>{\raggedright\arraybackslash}p{0.114\linewidth}@{}}
\toprule
\textbf{field} & \textbf{$\hat\lambda/\hat\lambda_1$} & \textbf{$E_0$} & \textbf{$\sum w\varepsilon$} & \textbf{$E_H$} & \textbf{$E_x$} & \textbf{$\varepsilon_{\text{HOMO}}$} \\
\midrule
\endfirsthead
\toprule
\textbf{field} & \textbf{$\hat\lambda/\hat\lambda_1$} & \textbf{$E_0$} & \textbf{$\sum w\varepsilon$} & \textbf{$E_H$} & \textbf{$E_x$} & \textbf{$\varepsilon_{\text{HOMO}}$} \\
\midrule
\endhead
both & 0 & 131.448298 & 131.448298 & 0 & 0 & 1.542011 \\
d16c & +1 & 127.409033 & 123.547979 & 0.874212 & $-4.735266$ & 1.473938 \\
d16c00 & +1 & 136.882683 & 142.569624 & 0.947939 & $+4.739001$ & 1.644264 \\
\bottomrule
\end{longtable}
\endgroup

\textbf{First excited states} (keys \texttt{\detokenize{ksGap}} and \texttt{excited.\allowbreak{}deltaSCF}; $\varepsilon_{\text{LUMO}}$ is the particle energy of the first particle--hole pair):

\begingroup
\small
\begin{longtable}{@{}>{\raggedright\arraybackslash}p{0.133\linewidth}>{\raggedright\arraybackslash}p{0.133\linewidth}>{\raggedright\arraybackslash}p{0.133\linewidth}>{\raggedright\arraybackslash}p{0.133\linewidth}>{\raggedright\arraybackslash}p{0.133\linewidth}>{\raggedright\arraybackslash}p{0.133\linewidth}@{}}
\toprule
\textbf{field} & \textbf{$\hat\lambda/\hat\lambda_1$} & \textbf{$\varepsilon_{\text{LUMO}}$} & \textbf{$\Delta_{KS}$} & \textbf{$\Delta_{\text{SCF}}$} & \textbf{$\Delta_{\text{SCF}}-\Delta_{KS}$} \\
\midrule
\endfirsthead
\toprule
\textbf{field} & \textbf{$\hat\lambda/\hat\lambda_1$} & \textbf{$\varepsilon_{\text{LUMO}}$} & \textbf{$\Delta_{KS}$} & \textbf{$\Delta_{\text{SCF}}$} & \textbf{$\Delta_{\text{SCF}}-\Delta_{KS}$} \\
\midrule
\endhead
both & 0 & 1.776275 & 0.2342638 & 0.2342638 & $1\times10^{-15}$ \\
d16c & +1 & 1.708971 & 0.2350337 & 0.2350344 & $7.0\times10^{-7}$ \\
d16c00 & +1 & 1.879407 & 0.2351435 & 0.2351447 & $1.25\times10^{-6}$ \\
\bottomrule
\end{longtable}
\endgroup

For dirac16complex the Stage-4 Rust solver gives, for the same configuration, $E_0=127.409027$, $E_H=0.874212$, $E_x=-4.735270$ and $\Delta_{KS}=0.2350337$ at $+\hat\lambda_1$, and $E_0=131.448295$, $\Delta_{KS}=0.2342638$ at $\lambda=0$ (\texttt{artifact\allowbreak{}s/dirac1\allowbreak{}6complex\allowbreak{}/kohn-sh\allowbreak{}am/rust/\allowbreak{}scf/m3\_L\allowbreak{}3\_N112\_l\allowbreak{}amp1\_T0/\allowbreak{}run.json} and \texttt{m3\_L3\_N1\allowbreak{}12\_lam0\_\allowbreak{}T0/run.j\allowbreak{}son}; the Rust coupling $852.6397056$ is calibrated separately and differs from the reference one in the ninth digit). The two solvers agree to $6\times10^{-6}$ in $E_0$ and $6\times10^{-9}$ in the gap. (Section 13.15 describes how the Stage-4 cross-check works. Its final committed result is the report \texttt{artifact\allowbreak{}s/dirac1\allowbreak{}6complex\allowbreak{}/kohn-sh\allowbreak{}am/pytho\allowbreak{}n-check-\allowbreak{}report.j\allowbreak{}son}: 63 checks, of which 62 are true. The one that fails, \texttt{canonica\allowbreak{}l\_eigenv\allowbreak{}alues}, concerns a single run of Chapter 13, the smeared ground state for $m=1$, $L=3$, $N=1016$ at $-\hat\lambda_2$ computed by the Rust solver on 301 grid points. Of the 58412 level comparisons of all runs, the only ones outside their tolerance are four of the 79 compared levels of that run (in both of its records, the ground state and the excited state); the worst, the deep level $\varepsilon=-1.05374908\,m$, differs from the reference solver by $2.2\times10^{-6}$ against a tolerance of $1.05\times10^{-6}$ (key \texttt{canonica\allowbreak{}l\_eigenv\allowbreak{}alues\_de\allowbreak{}tail}), while the refined Rust run of the same state on 601 grid points agrees with the reference levels to $1.8\times10^{-7}$ (comparison \texttt{canonica\allowbreak{}l\_excite\allowbreak{}d\_m1\_L3\_\allowbreak{}N1016\_la\allowbreak{}mm2\_T0\_g\allowbreak{}601}). The Delta-SCF energy of the same run agrees within its tolerance after erratum E4.13 of \texttt{handoff/\allowbreak{}specs/ST\allowbreak{}AGE4\_SPE\allowbreak{}C.md} (check \texttt{canonica\allowbreak{}l\_deltaS\allowbreak{}CF} true). The Stage-4 gate had not been run.) There is no second solver for the dirac16complex00 column.

\textbf{Worked checks.} (i) The double-counting formula of Section 14.7 holds in both rows: $123.547979-0.874212+4.735266=127.409033$ and $142.569624-0.947939-4.739001=136.882684$ (the last digit is rounding; the unrounded record gives $136.8826835$). (ii) The statistics sign: the exchange energies have opposite signs and nearly equal magnitudes, $-4.735$ and $+4.739$. (iii) The first-order estimate of Section 14.7 says that the two energies differ by $2X^{\text{free}}$, twice the magnitude of the exchange energy of the free densities; each computed $\lvert E_x\rvert$ equals $X^{\text{free}}$ up to terms of second order in $\lambda$. Indeed the computed difference $136.882683-127.409033=9.473650$ is close to $E_x^{00}-E_x=4.739001+4.735266=9.474267$: the two numbers are equal at first order in $\lambda$, and here they differ by $6.2\times10^{-4}$. (iv) The levels move as predicted: the highest occupied level rises from $1.542$ to $1.644$ for dirac16complex00, whose exchange potential $+\tfrac\lambda{16}n_p$ is positive, and falls to $1.474$ for dirac16complex, whose exchange potential $-\tfrac\lambda{16}n_p$ is negative (the nearly equal mass shifts contribute to both). (v) This is a fast gas in the sense of Section 14.6: the exchange energies are five times the Hartree energies, because the brane band carries a small scalar density ($\lvert S\rvert\ll n$).

\textbf{What the first excited state is.} In all three rows the highest occupied level is the brane band (Section 13.6) on the torus shell $\nu=n_1^2+n_2^2+n_3^2=3$, that is $k=\sqrt3\cdot0.75H=1.299H$ ($0.433\,m$), with $4\cdot8=32$ states, and the lowest empty level is the band on the shell $\nu=4$, $k=1.5H$ ($0.5\,m$), with $4\cdot6=24$ states (reference level labels \texttt{3:+1:-1:\allowbreak{}0} and \texttt{4:+1:-1:\allowbreak{}0}, which name the shell $\nu$, the brane parity, the block type in the reference solver's convention and the index of the level, with the multiplicities \texttt{homoMult\allowbreak{}iplicity} 32 and \texttt{lumoMult\allowbreak{}iplicity} 24). The filling is the one of Section 13.7: $8+24+48+32=112$. The first excitation moves one particle from the shell $\nu=3$ to the shell $\nu=4$, an excitation \textbf{within the brane band}, exactly as for the $m=1$ states of Chapter 13. The next particle--hole energies of dirac16complex00 are $0.441136$, $0.516079$ and $0.626495$ (dirac16complex: $0.440969$, $0.515752$, $0.626300$; key \texttt{excited.\allowbreak{}particle\allowbreak{}Hole}). The interaction raises the gap by $7.7\times10^{-4}$ (dirac16complex) and $8.8\times10^{-4}$ (dirac16complex00), and the orbital relaxation $\Delta_{\text{SCF}}-\Delta_{KS}$ is of order $10^{-6}$ in both: at this coupling the difference between the two statistics shows in the total energy and in the absolute position of the levels, hardly in the excitation energy.

\textbf{Energy--momentum tensor.} The proper-volume averages of Section 13.14 (key \texttt{\detokenize{emt}} of each run; $\kappa_{\text{needed}}=-21H^2/\langle\rho\rangle$):

\begingroup
\small
\begin{longtable}{@{}>{\raggedright\arraybackslash}p{0.114\linewidth}>{\raggedright\arraybackslash}p{0.114\linewidth}>{\raggedright\arraybackslash}p{0.114\linewidth}>{\raggedright\arraybackslash}p{0.114\linewidth}>{\raggedright\arraybackslash}p{0.114\linewidth}>{\raggedright\arraybackslash}p{0.114\linewidth}>{\raggedright\arraybackslash}p{0.114\linewidth}@{}}
\toprule
\textbf{field} & \textbf{$\hat\lambda/\hat\lambda_1$} & \textbf{$\langle\rho\rangle$} & \textbf{$w_y$} & \textbf{$w_3$} & \textbf{$w_t$} & \textbf{$\kappa_{\text{needed}}$} \\
\midrule
\endfirsthead
\toprule
\textbf{field} & \textbf{$\hat\lambda/\hat\lambda_1$} & \textbf{$\langle\rho\rangle$} & \textbf{$w_y$} & \textbf{$w_3$} & \textbf{$w_t$} & \textbf{$\kappa_{\text{needed}}$} \\
\midrule
\endhead
both & 0 & 1.341376 & 0.1772 & 0.3291 & 0 & $-15.66$ \\
d16c & +1 & 1.300157 & 0.1753 & 0.3110 & $-0.0303$ & $-16.15$ \\
d16c00 & +1 & 1.396831 & 0.2371 & 0.3596 & $+0.0415$ & $-15.03$ \\
\bottomrule
\end{longtable}
\endgroup

The pressure along the extra times is $p_t=L_s=e_H+e_x$ (Section 13.14), so its sign follows the exchange sign: the average $\langle p_t\rangle$ is $-0.0394$ for dirac16complex and $+0.0580$ for dirac16complex00. In every row $\langle\rho\rangle>0$, while the static field requires $\rho_{\text{req}}=-21H^2/\kappa<0$; a negative $\kappa$ would be needed, and the sourcing conditions of erratum E4.1 are not met (for dirac16complex00 the run records $\lambda\langle S_p\rangle/m=-0.0347$ against the required $-\tfrac56$; key \texttt{emt.E41\_\allowbreak{}sourcing\allowbreak{}Conditio\allowbreak{}ns}). The verdict of Section 13.14 is unchanged for dirac16complex00 at this coupling.

\textbf{Finite temperature without interaction.} At $T=0.1\,m=0.3H$ and $\lambda=0$ both fields give $\mu=1.469481$, $E=161.026343$, $F=104.727753$, $S_{\text{ent}}=187.661967$ and $\Delta_{KS}=0.279429$ (runs \texttt{d16c\_m3\_\allowbreak{}L3\_N112\_\allowbreak{}lam0\_T0p\allowbreak{}1/plusM} and \texttt{d16c00\_m\allowbreak{}3\_L3\_N11\allowbreak{}2\_lam0\_T\allowbreak{}0p1/plus\allowbreak{}M}), and $F=E-TS_{\text{ent}}=161.026343-0.3\cdot187.661967=104.727753$ checks. The interacting thermal points were not attempted (Section 14.9), so no finite-temperature result distinguishes the two fields.

\textbf{The mirror universes, briefly.} Chapter 15 proves that the Kohn--Sham problem with mass $-m$, the same $\lambda$ and the transformed boundary conditions has exactly the same energies as the $+m$ problem, for both statistics (\texttt{pairing-\allowbreak{}theory.j\allowbreak{}son}, key \texttt{T3.theor\allowbreak{}emStanda\allowbreak{}rdRule}; checks \texttt{PAIR\_T3k\allowbreak{}s\_sigma2\allowbreak{}Function\allowbreak{}alInvari\allowbreak{}ant} and \texttt{PAIR\_T3k\allowbreak{}s\_sigma2\allowbreak{}KSOperat\allowbreak{}orEquiva\allowbreak{}riant}, which carry the statistics sign). For dirac16complex00 the reference solver computed this partner only at $\lambda=0$: $E_0=131.448295$ for $-M$ against $131.448298$ for $+M$, a difference of $2.7\times10^{-6}$ that comes from the staggered grid, which treats the two components of an orbital differently (\texttt{pairing-\allowbreak{}theory.j\allowbreak{}son}, key \texttt{T3.numer\allowbreak{}icsPresc\allowbreak{}ription.\allowbreak{}discreti\allowbreak{}sationNo\allowbreak{}te}). The interacting partner of dirac16complex00 is among the pending runs: \textbf{for the commuting field the pairing at $\lambda\ne0$ is proved (Chapter 15) but not demonstrated numerically.}

\subsection{What was not computed}

Stated plainly, for dirac16complex00:

\begin{itemize}
\item no Kohn--Sham state for $|m|=1$ (the configurations of Chapter 13) and none for $N=8$ ($N=1016$ was not part of the Stage-5 run matrix);
\item for $m=3$, $N=112$: no attractive coupling ($-\hat\lambda_1$, $-\hat\lambda_2$) and no strong coupling ($+\hat\lambda_2$) at $T=0$;
\item no interacting state at a finite temperature;
\item no run of the Rust solver, hence no comparison of two independent solvers, and no report of the pairs checker; the Stage-5 gate was not run;
\item no mirror ($-M$) partner at $\lambda\ne0$.
\end{itemize}

The missing runs were not completed because the Stage-5 numerics were stopped so that this book could be written (\texttt{HANDOFF.\allowbreak{}md}, section 2). The exact results of Sections 14.3 to 14.8 do not depend on them.

\subsection{What we proved and what we assumed}

\textbf{Proved} (derived step by step in this chapter and verified by the named exact checks): the normalization and the moments $\langle(c^\ast)^\alpha c^\beta\rangle=\delta_{\alpha\beta}\alpha!f^\alpha$ of the circular Gaussian, and the corresponding single-mode results for thermal bosons and for fermions; Isserlis' rule for four Gaussian amplitudes (Proposition 14.2) and Wick's rule for four fermion operators (Proposition 14.3); Theorem 14.4, $\langle(\psi^\dagger X\psi)(\psi^\dagger Y\psi)\rangle=\mathrm{Tr}(X\rho)\mathrm{Tr}(Y\rho)+\mathrm{sg}\,\mathrm{Tr}(X\rho Y\rho)$ with $\mathrm{sg}=+1$ for Gaussian waves and $\mathrm{sg}=-1$ for normal-ordered fermions, with the extra term $\mathrm{Tr}(XY\rho)$ of the full quantum products and the deviation of fixed-modulus waves; Theorem 14.5, the Hartree--Fock energy $\tfrac\lambda2[S^2+\mathrm{sg}\,\mathrm{Tr}(BC\rho BC\rho)]$ of the contact interaction for both fields; the filled-shell ratio $E_x=\mathrm{sg}\,E_H/8$; the exact uniform-gas exchange $e_x=\mathrm{sg}\,\tfrac\lambda{32}(n^2+S^2)$ at every temperature, hence $e_x^{00}=+\tfrac\lambda{32}(n^2+S^2)$; the potentials $v^{00}=+\tfrac\lambda{16}n_p$, $v_s^{00}=+\tfrac\lambda{16}S_p$, $M_{\text{eff}}^{00}=m+\tfrac{17}{16}\lambda S_p$; the Kohn--Sham equations of dirac16complex00 from its Mermin functional, and its double-counting formula; the indefinite charge density and the unbounded classical energy of dirac16complex00 (explicit solutions); the indefiniteness of the covariance $K=\rho B$ of the model and the vanishing charge and scalar density of an honest Gaussian filled shell; the Krein weighting $T^{\text{classical}}=js_2\,T^{\text{KS}}$ of single-block modes, component by component, with the explicit box modes whose energy densities are $\pm\tfrac85(25-9\cos8y+12\sin8y)$; and, directly from the definitions (no separate check is needed), the negative variances $w_n\beta_n$ of the model from the Krein signs and, at $T>0$, from the thermal holes of the sea.

\textbf{Computed} (reference solver only, no independent cross-check): the ground and first excited states of dirac16complex00 for $m=3H$, $L=3/H$, $N=112$ at $\lambda=0$ and $+\hat\lambda_1$ and $T=0$, their energy--momentum averages, and the free thermal state at $T=0.3H$ (Section 14.10).

\textbf{Assumed:} the model P1 to P4 of Section 14.2: the averaging rule of a Gaussian random-phase ensemble applied to the Kohn--Sham orbitals; Fermi--Dirac occupations and the particle/sea bookkeeping of Chapter 13 imposed on a commuting field, whose negative sea weights already make the ensemble formal; the expectation-value rule of the quantized fermion field adopted by prescription, which makes the covariance Krein-signed rather than positive; the definition of the model as a mean-field functional, not as a quantization of the field (the third point of Section 14.8); exchange only, with no correlation; and everything assumed in Chapter 13 (fixed static field with the Z2 brane, good sector, tip cutoff and bag condition, torus, normal ordering against the free sea, continuation in $\lambda$, no back-reaction). We quoted without proof the harmonic-oscillator spectrum of a boson mode, the general Isserlis and Wick theorems for products of any length (only the four-factor cases are used and proved), Pauli's spin--statistics theorem, and the Hanbury Brown--Twiss effect as an illustration.

\textbf{Not claimed:} that the model is a quantum theory of commuting spinors, that its ground state minimizes any energy of the classical field, that its energies are classical field energies, or anything about the creation of universes; Chapters 15 and 16 treat the pairing questions.

\subsection{Exercises}

\textbf{Exercise 14.1.} For the circular Gaussian with $f=2$, compute $\langle|c|^6\rangle$ from Proposition 14.1 and directly from the integral $\frac1{\pi f}\int r^6e^{-r^2/f}\,r\,dr\,d\varphi$. Show also that $\langle c^2\rangle=0$.

\textbf{Exercise 14.2.} Show that $\sum_{k\ge0}k(k-1)(k-2)\,x^k=6x^3/(1-x)^4$ for $0\le x<1$, and deduce that a thermal boson mode has $\langle k(k-1)(k-2)\rangle=6f^3$ with $f=x/(1-x)$. Compare with $\langle|c|^6\rangle$ of the Gaussian and with the fermion mode.

\textbf{Exercise 14.3.} Repeat the first worked example of Section 14.5 (the rotated modes, $X=Y=\sigma_x$) with $f_0=\tfrac12$ and $f_1=\tfrac14$: compute $\rho$, $\mathrm{Tr}(\sigma_x\rho)$, $\mathrm{Tr}(\sigma_x\rho\sigma_x\rho)$ and the three averages of $(\psi^\dagger\sigma_x\psi)^2$.

\textbf{Exercise 14.4.} One mode ($d=1$, $X=Y=1$). Compute (a) $\langle|c|^4\rangle$ for a Gaussian amplitude, (b) $\langle(a^\dagger a)^2\rangle$ for a fermion mode with occupation $f$, (c) $\langle(b^\dagger b)^2\rangle$ for a thermal boson mode, and check each against Theorem 14.4 and Remark 1 of Section 14.5.

\textbf{Exercise 14.5.} For the filled shell at the moving point of Section 14.6 ($\mathrm{Tr}(BCP_+)=\tfrac{16}5$, $\mathrm{Tr}(BCP_+BCP_+)=\tfrac{32}{25}$) and $\lambda=1$, compute $E_H$ and $E_x$ for both fields and the ratios $E_x/E_H$.

\textbf{Exercise 14.6.} Recompute the densities of the exact finite gas of Section 14.6. Each momentum contributes $8(f_+-f_-)$ to $n$ and $8\tfrac m{E}(f_++f_-)$ to $S$, with $m=2$; the shells $E=2,3,4,5$ contain $1,4,2,4$ momenta. Check that $(n^2+S^2)/16$ equals the trace $\tfrac{768720549416}{5636255625}$ recorded in \texttt{python-p\allowbreak{}airing-r\allowbreak{}eport.js\allowbreak{}on}.

\textbf{Exercise 14.7.} At a point with $\lambda=8$, $n=0.4$ and $S=0.2$ ($H=1$), compute $e_H$, $e_x$, $v_v$, $v_s$ and $M_{\text{eff}}-m$ for both fields, and check the coefficients $\tfrac{15}{16}$ and $\tfrac{17}{16}$.

\textbf{Exercise 14.8.} Along the uniform gas at $T=0$, adding particles fills the Fermi surface, where every particle has the energy $E_F$. Show that $dS/dn=m/E_F$ there, and deduce the total derivative $\frac{d e_x}{dn}=\mathrm{sg}\,\tfrac\lambda{16}\bigl(n+S\,\tfrac m{E_F}\bigr)$ (check \texttt{PAIR\_sta\allowbreak{}t\_T0Tota\allowbreak{}lDerivat\allowbreak{}ive}).

\textbf{Exercise 14.9.} (a) Prove that the covariance $K_{ab}=\langle\Psi_a\Psi_b^\ast\rangle$ of any collection of classical fields is Hermitian and positive semidefinite. (b) For $K=P_+B$ and the vector $u_+=\tfrac12(e_0-e_4-i\,e_9-i\,e_{13})$ of Section 8.8 compute $u_+^\dagger Ku_+$, and compare with the value for $u_-$ found in Section 14.8.

\textbf{Exercise 14.10.} In a block of type $(j,s_2)=(1,-1)$, show that $\Psi^\dagger B\Psi=-\Psi^\dagger\Psi$ and conclude the sign of the classical energy density of a stationary single-block wave with $\varepsilon>0$. What does the Kohn--Sham model assign to the same orbital?

\textbf{Exercise 14.11.} Using the table of ground states in Section 14.10: (a) verify the double-counting formula for both fields; (b) compute $E_0^{00}-E_0$ and $E_x^{00}-E_x$ and explain why they nearly agree; (c) compute $E_0-E_0(\lambda=0)$ for both fields and compare with $E_H+E_x$ of the same run.

\textbf{Exercise 14.12.} Show that the argument of Exercise 13.9 carries over to dirac16complex00: for $N=8$ at $T=0$ the ground states at $\lambda$ and $-\lambda$ are related by moving every occupied orbital from block type $j$ to $-j$, and $E_0^{00}(-\lambda)=-E_0^{00}(\lambda)$. For $\lambda>0$, do the zero modes of dirac16complex00 move up or down? (This state was not computed; the exercise asks for the exact argument.)

\textbf{Exercise 14.13.} Explain, from the functional of Section 14.7, why the two fields give identical Kohn--Sham results at $\lambda=0$ at every temperature, and why the reference runs at $T=0.3H$ therefore cannot distinguish them.

\textbf{Exercise 14.14.} A classical Gaussian filled rest shell with the positive covariance $K=P_+$: compute the average charge density $\langle\Psi^\dagger B\Psi\rangle$, the scalar density $\langle\Psi^\dagger C\Psi\rangle$ and, with Theorem 14.4, $\langle S^2\rangle$. Compare with the expectation-value rule of the model.

\textbf{Exercise 14.15.} The box mode of Section 14.8 ($m=3$, $k=0$, $\lambda=0$, $\varepsilon=5$, $\chi_2=\sin4y$, $\chi_1=(3\sin4y+4\cos4y)/(5ij)$ on $-\pi/4\le y\le0$). (a) Check that it solves the block equation of Section 13.5, which here reads $\chi_1'=3\chi_1+5ij\,\chi_2$ and $\chi_2'=-3\chi_2+5ij\,\chi_1$, and the two boundary conditions. (b) Show that $40\bigl[\sin^24y+\tfrac1{25}(3\sin4y+4\cos4y)^2\bigr]=\tfrac85\bigl(25-9\cos8y+12\sin8y\bigr)$. (c) Show that $a\cos\theta+b\sin\theta$ lies between $-\sqrt{a^2+b^2}$ and $+\sqrt{a^2+b^2}$ for all real $a,b,\theta$, and conclude that the profile of (b) lies between 16 and 64. (d) Integrate $\rho W^6$ over the box for the two blocks and compare with $js_2\,\varepsilon\int\chi^\dagger\chi\,dy$.

\subsection{Answers to the exercises}

\textbf{Answer 14.1.} Proposition 14.1 gives $\langle|c|^6\rangle=3!\,f^3=6\cdot8=48$. Directly: the angular integral gives $2\pi$, so $\langle|c|^6\rangle=\frac2f\int_0^\infty r^7e^{-r^2/f}dr$; with $s=r^2/f$, $r^6=f^3s^3$ and $r\,dr=\tfrac f2ds$, this is $\frac2f\cdot f^3\cdot\frac f2\int_0^\infty s^3e^{-s}ds=f^3\cdot3!=48$. $\langle c^2\rangle$ has $\alpha=0$, $\beta=2$; its angular integral $\int_0^{2\pi}e^{2i\varphi}d\varphi=0$, so it vanishes.

\textbf{Answer 14.2.} Differentiate $\sum_kx^k=(1-x)^{-1}$ three times: $\sum_kk(k-1)(k-2)x^{k-3}=6(1-x)^{-4}$; multiply by $x^3$. Then $\langle k(k-1)(k-2)\rangle=(1-x)\cdot6x^3/(1-x)^4=6\bigl(x/(1-x)\bigr)^3=6f^3$, equal to the Gaussian $\langle|c|^6\rangle=3!f^3$. For a fermion mode $k(k-1)(k-2)=0$ for $k\in\{0,1\}$, so the average is 0.

\textbf{Answer 14.3.} $\rho=\tfrac12\begin{pmatrix}f_0+f_1&f_0-f_1\\ f_0-f_1&f_0+f_1\end{pmatrix}=\begin{pmatrix}\tfrac38&\tfrac18\\ \tfrac18&\tfrac38\end{pmatrix}$. $\mathrm{Tr}(\sigma_x\rho)=2\cdot\tfrac18=\tfrac14=f_0-f_1$. $\sigma_x\rho=\begin{pmatrix}\tfrac18&\tfrac38\\ \tfrac38&\tfrac18\end{pmatrix}$, so $\mathrm{Tr}(\sigma_x\rho\sigma_x\rho)=\tfrac1{64}+\tfrac9{64}+\tfrac9{64}+\tfrac1{64}=\tfrac{20}{64}=\tfrac5{16}=f_0^2+f_1^2$. Gaussian: $\tfrac1{16}+\tfrac5{16}=\tfrac38$; fixed moduli: $\tfrac38-\tfrac5{16}=\tfrac1{16}=(f_0-f_1)^2$; fermions: $\tfrac1{16}-\tfrac5{16}=-\tfrac14=-2f_0f_1$.

\textbf{Answer 14.4.} (a) $2f^2$; Theorem 14.4 with $\rho=f$: $f\cdot f+f\cdot f=2f^2$. (b) $(a^\dagger a)^2=a^\dagger a$ (Section 8.3), so the average is $f$; Remark 1: $f^2-f^2+\mathrm{Tr}(1\cdot1\cdot f)=f$. (c) $\langle N^2\rangle=\langle N(N-1)\rangle+\langle N\rangle=2f^2+f$; the boson version of Remark 1: $f^2+f^2+f$. The classical Gaussian (a) equals the normal-ordered boson value $2f^2$, without the extra $f$.

\textbf{Answer 14.5.} $E_H=\tfrac12\bigl(\tfrac{16}5\bigr)^2=\tfrac{128}{25}=5.12$ for both fields. $E_x=\mathrm{sg}\cdot\tfrac12\cdot\tfrac{32}{25}=\mathrm{sg}\,\tfrac{16}{25}=\mp0.64$: $-0.64$ for dirac16complex and $+0.64$ for dirac16complex00. The ratios are $\mp\tfrac{16/25}{128/25}=\mp\tfrac18$.

\textbf{Answer 14.6.} Per shell, $n$ receives $8(f_+-f_-)\times$(number of momenta): $E=2$: $8\cdot\tfrac67=\tfrac{48}7$; $E=3$: $8\cdot\tfrac{29}{44}\cdot4=\tfrac{232}{11}$; $E=4$: $8\cdot\tfrac13\cdot2=\tfrac{16}3$; $E=5$: $8\cdot\tfrac3{130}\cdot4=\tfrac{48}{65}$. $S$ receives $8\tfrac2E(f_++f_-)\times$(number): $E=2$: $8\cdot\tfrac87=\tfrac{64}7$; $E=3$: $\tfrac{16}3\cdot\tfrac{37}{44}\cdot4=\tfrac{592}{33}$; $E=4$: $4\cdot\tfrac13\cdot2=\tfrac83$; $E=5$: $\tfrac{16}5\cdot\tfrac{23}{130}\cdot4=\tfrac{736}{325}$. Summing, $n=\tfrac{510808}{15015}\approx34.01985$ and $S=\tfrac{2403416}{75075}\approx32.01353$, and $(n^2+S^2)/16=\tfrac{768720549416}{5636255625}\approx136.38852$, the recorded trace (the common denominators are best handled with a computer or with patience).

\textbf{Answer 14.7.} $e_H=\tfrac82\cdot0.04=0.16$. $\tfrac\lambda{32}(n^2+S^2)=\tfrac14(0.16+0.04)=0.05$, so $e_x=-0.05$ (dirac16complex) and $+0.05$ (dirac16complex00). $v_v=\mp\tfrac8{16}\cdot0.4=\mp0.2$ and $v_s=\mp\tfrac8{16}\cdot0.2=\mp0.1$ (upper sign dirac16complex). $\lambda S=1.6$, so $M_{\text{eff}}-m=1.6-0.1=1.5=\tfrac{15}{16}\cdot1.6$ and $1.6+0.1=1.7=\tfrac{17}{16}\cdot1.6$.

\textbf{Answer 14.8.} At $T=0$ the added particles sit on the Fermi surface; each adds 1 to the number and $m/E_F$ to the scalar density (Section 13.8: a particle of energy $E_p$ contributes $m/E_p$ to $S$ per unit of $n$), so $dS/dn=m/E_F$. By the chain rule $\frac{de_x}{dn}=\frac{\partial e_x}{\partial n}+\frac{\partial e_x}{\partial S}\frac{dS}{dn}=\mathrm{sg}\,\tfrac\lambda{16}n+\mathrm{sg}\,\tfrac\lambda{16}S\,\tfrac m{E_F}$.

\textbf{Answer 14.9.} (a) $K_{ba}^\ast=\langle\Psi_b\Psi_a^\ast\rangle^\ast=\langle\Psi_b^\ast\Psi_a\rangle=K_{ab}$, so $K^\dagger=K$; and $v^\dagger Kv=\langle|v^\dagger\Psi|^2\rangle\ge0$ as shown in Section 14.8. (b) Section 8.8 shows $-i\gamma^4u_+=u_+$ and $Bu_+=u_+$, so $Ku_+=P_+Bu_+=P_+u_+=u_+$ and $u_+^\dagger Ku_+=u_+^\dagger u_+=1$: positive, as a covariance should be, while $u_-^\dagger Ku_-=-1$. Both vectors are positive-energy states at rest; they differ only in their Krein sign.

\textbf{Answer 14.10.} In the block $B$ acts as the number $js_2=-1$, so $B\Psi=-\Psi$ and $\Psi^\dagger B\Psi=-\Psi^\dagger\Psi\le0$. The energy density of the classical wave contains $\varepsilon\Psi^\dagger B\Psi=-\varepsilon\Psi^\dagger\Psi<0$ for $\varepsilon>0$: negative. The Kohn--Sham model assigns $\varepsilon\,u^\dagger B^2u=\varepsilon\,n>0$: the opposite sign, which is the factor $js_2$ of Section 14.8.

\textbf{Answer 14.11.} (a) Section 14.10, worked check (i). (b) $136.882683-127.409033=9.473650$ and $4.739001+4.735266=9.474267$. At first order both energies are $E(0)+E_H^{\text{free}}\pm X^{\text{free}}$, so their difference is $2X^{\text{free}}$. Each run's $\lvert E_x\rvert$ is $X^{\text{free}}$ up to terms of second order in $\lambda$ (its densities differ from the free ones at first order), so the sum of the two magnitudes is also $2X^{\text{free}}$ at first order. The two numbers therefore agree at first order in $\lambda$, and their difference, $6.2\times10^{-4}$, is of second order. (c) dirac16complex: $127.409033-131.448298=-4.039265$, against $E_H+E_x=0.874212-4.735266=-3.861054$; dirac16complex00: $136.882683-131.448298=5.434385$, against $0.947939+4.739001=5.686940$. The Hellmann--Feynman relation $dE/d\lambda=(E_H+E_x)/\lambda$ (Section 14.7, at $T=0$ where $F=E$) gives $E(\lambda)-E(0)=\int_0^\lambda\bigl(E_H+E_x\bigr)(\lambda')\,\frac{d\lambda'}{\lambda'}$, which equals $(E_H+E_x)(\lambda)$ at first order in $\lambda$; the computed differences, $-0.178$ and $-0.253$, are of second order and say that $(E_H+E_x)/\lambda$ changed along the continuation. Note that no minimum principle fixes the sign of this second-order term here: the one-body Dirac operator has no lower bound, and the particle levels are defined by continuity from $\lambda=0$ (Section 13.7), so the variational argument of Section 12.2 (the exact minimum lies below every trial state) cannot be used to predict it.

\textbf{Answer 14.12.} At $T=0$ only the eight zero-mode particle levels carry weight. Move each occupied orbital unchanged from $j$ to $-j$ and replace $\lambda$ by $-\lambda$. Then $n_p$ is unchanged and $S_p\to-S_p$ (because the scalar density of an orbital is $s=j\chi^\dagger\sigma_y\chi$ times $e^{-6Hy}/\ell^3$, Section 13.7), so $M_{\text{eff}}-m=\tfrac{17}{16}\lambda S_p$ is unchanged, and $v^{00}=\tfrac\lambda{16}n_p\to-v^{00}$. Writing $h_j^{00}=jD+v^{00}$ with $D=-i\sigma_x\frac d{dy}+M_{\text{eff}}\sigma_y+\xi k\sigma_z$ (the operator $D_k$ of Section 13.5, built from $M_{\text{eff}}$), the new operator is $-jD-v^{00}=-h_j^{00}$: each orbital is an eigenvector with eigenvalue $-\varepsilon$, and it reproduces the densities used, so the mapped state is self-consistent at $-\lambda$ (and it is reached by continuation, since both start from the same eight zero modes). Then $\sum w\varepsilon\to-\sum w\varepsilon$, $E_H=\tfrac\lambda2\int S_p^2dV\to-E_H$ and $E_x^{00}=\tfrac\lambda{32}\int(n_p^2+S_p^2)dV\to-E_x^{00}$, so $E_0^{00}\to-E_0^{00}$. The zero modes have $\varepsilon=0$ and $\chi_2\equiv0$, hence zero scalar density, so to first order they move by the average of $v^{00}=+\tfrac\lambda{16}n_p>0$: \textbf{up} for $\lambda>0$, the opposite of dirac16complex, whose zero modes move down (Section 13.11).

\textbf{Answer 14.13.} At $\lambda=0$ the terms $E_H$ and $E_x$ vanish, and the only place where $\mathrm{sg}$ appears (the exchange term and its potentials) is multiplied by $\lambda$. The functional, and with it the Kohn--Sham equations, the occupations at every $T$, the energies and the levels, are then identical for the two fields. The only thermal runs of Section 14.10 have $\lambda=0$, so they carry no information on the statistics.

\textbf{Answer 14.14.} With $K=P_+$: $\langle\Psi^\dagger B\Psi\rangle=\mathrm{Tr}(BP_+)=0$ and $\langle\Psi^\dagger C\Psi\rangle=\mathrm{Tr}(CP_+)=0$ (Section 14.8). Theorem 14.4 with $X=Y=C$ and $\rho\to P_+$ gives $\langle S^2\rangle=[\mathrm{Tr}(CP_+)]^2+\mathrm{Tr}(CP_+CP_+)=0+\mathrm{Tr}(CP_+CP_+)$; since $C$ commutes with $\gamma^4$, it commutes with $P_+$, so $CP_+CP_+=C^2P_+^2=P_+$ and $\langle S^2\rangle=\mathrm{Tr}P_+=8$. The honest ensemble thus has zero average scalar density but a nonzero average of $S^2$: its interaction energy is pure fluctuation. The rule of the model instead gives $\langle\Psi^\dagger B\Psi\rangle=8$, $S=8$ and $\langle S^2\rangle=64+8=72$, that is $E_H=\tfrac\lambda2\cdot64$ and $E_x=+\tfrac\lambda2\cdot8=E_H/8$.

\textbf{Answer 14.15.} (a) We use $(ij)^2=i^2j^2=-1$. First, $\chi_2'=4\cos4y$, and $-3\chi_2+5ij\,\chi_1=-3\sin4y+(3\sin4y+4\cos4y)=4\cos4y$. Second, $\chi_1'=(12\cos4y-16\sin4y)/(5ij)$; on the right, write $5ij\sin4y=(5ij)^2\sin4y/(5ij)=-25\sin4y/(5ij)$, so $3\chi_1+5ij\,\chi_2=(9\sin4y+12\cos4y-25\sin4y)/(5ij)=(12\cos4y-16\sin4y)/(5ij)$. Both equations hold. The boundary conditions: $\chi_2(0)=\sin0=0$ and $\chi_2(-\pi/4)=\sin(-\pi)=0$. (b) Write $s=\sin4y$ and $c=\cos4y$. Since $|5ij|=5$ and $3s+4c$ is real, $|\chi_1|^2=(3s+4c)^2/25$. Then $40\bigl[s^2+\tfrac1{25}(9s^2+24sc+16c^2)\bigr]=\tfrac85\bigl(34s^2+24sc+16c^2\bigr)$. The double-angle formulas with the angle $4y$ give $s^2=\tfrac12(1-\cos8y)$, $c^2=\tfrac12(1+\cos8y)$ and $sc=\tfrac12\sin8y$, so $34s^2+24sc+16c^2=17-17\cos8y+12\sin8y+8+8\cos8y=25-9\cos8y+12\sin8y$. (c) Let $R=\sqrt{a^2+b^2}>0$ (for $R=0$ there is nothing to show). The point $(a/R,b/R)$ lies on the unit circle, so $a=R\cos\varphi$ and $b=R\sin\varphi$ for some angle $\varphi$, and $a\cos\theta+b\sin\theta=R(\cos\varphi\cos\theta+\sin\varphi\sin\theta)=R\cos(\theta-\varphi)$, which lies between $-R$ and $R$. With $a=-9$, $b=12$ and $\theta=8y$: $R=\sqrt{81+144}=15$, so $25-9\cos8y+12\sin8y$ lies between $10$ and $40$, and the profile between $\tfrac85\cdot10=16$ and $\tfrac85\cdot40=64$. It is positive everywhere in block 0; in block 2 it is multiplied by $js_2=-1$ and lies between $-64$ and $-16$. (d) Over an interval of length $\pi/4$ the functions $\cos8y$ and $\sin8y$ run through exactly one period, so their integrals vanish: $\int_{-\pi/4}^0\cos8y\,dy=\bigl[\tfrac18\sin8y\bigr]_{-\pi/4}^0=0$ and $\int_{-\pi/4}^0\sin8y\,dy=\bigl[-\tfrac18\cos8y\bigr]_{-\pi/4}^0=0$. Hence $\int\rho W^6dy=js_2\cdot\tfrac85\cdot25\cdot\tfrac\pi4=js_2\cdot10\pi$. On the other hand $\int\chi^\dagger\chi\,dy=8\int\bigl(|\chi_1|^2+|\chi_2|^2\bigr)dy=\tfrac8{25}\int\bigl(25-9\cos8y+12\sin8y\bigr)dy=\tfrac8{25}\cdot\tfrac{25\pi}4=2\pi$, and $js_2\,\varepsilon\cdot2\pi=js_2\cdot10\pi$: the same. The classical energy of the mode (per unit coordinate volume of the six directions $x_1,x_2,x_3,x_5,x_6,x_7$, since $W^6dy=\sqrt{|g|}\,dy$ is the rest of the volume element) is $+10\pi$ in block 0 and $-10\pi$ in block 2, while the Kohn--Sham model assigns the orbital (normalized to $\int\chi^\dagger\chi\,dy=1$) the energy $+5$ in both blocks.

\section{The pairing theorems}

\subsection{What this chapter proves, and what it does not}

\textbf{The question.} The author's notebook opens with a hypothesis. Its cell 6 says that at the time $x_4=0$ "a pair of universes with MASSES ± M is created" (quoted from the survey \texttt{handoff/\allowbreak{}surveys/\allowbreak{}survey\_n\allowbreak{}otebook-\allowbreak{}physics.\allowbreak{}md}; Chapter 0 quotes the cells 6, 7 and 17 in full, and Chapter 16 discusses them). Before one can ask whether such a pair is created, one must know whether the equations of this book allow it at all. If a universe of mass $+M$ is a solution, is there a partner of mass $-M$? And what energy, momentum and charge do the two carry together? This chapter answers these two questions exactly, with three theorems called T1, T2 and T3 as in the Stage-5 specification \texttt{handoff/\allowbreak{}specs/ST\allowbreak{}AGE5\_SPE\allowbreak{}C.md} (its §5), and it reports the numerical test of the third one.

In this chapter a \textbf{universe of mass $+M$} means a configuration or solution of the field equations (from Section 15.7 on, a Kohn--Sham state of Chapter 13) whose mass parameter is $m=+M$, and its \textbf{partner of mass $-M$} is one with $m=-M$. The notebook writes its mass term with a number $M$, and the Lagrangian of this book has $m=-HM$ (Section 5.12). Every statement below relates $m$ to $-m$, so it does not matter which member of a pair is called the universe of mass $+M$.

\textbf{The three theorems in words.}

\begin{itemize}
\item \textbf{T1, the chirality map} (Sections 15.3 and 15.4). Multiply the field by the chirality $\gamma^8$ of Chapter 2. A configuration of the theory with mass $m$ and coupling $\lambda$ becomes a configuration of the theory with $-m$ and $-\lambda$, in every gravitational field. Its energy density, momentum, stresses and charge, computed from the bilinears of the field, all change sign; its scalar density $S=\bar\Psi\Psi$ does not. For the fields treated as \textbf{classical} fields, the pair made of a field and its image therefore carries \textbf{zero} total energy--momentum and zero total charge, at every point and at every time. For the quantized field no reading examined in the repository turns this into a cancellation between two independent universes; that question is OPEN (Section 15.4).
\item \textbf{T2, the mirror map} (Sections 15.5 and 15.6). Combine the field with the reflection of one space-like direction. A configuration with $(m,\lambda)$ becomes one with $(-m,\lambda)$, the \textbf{same} coupling. Energy, momentum and charge are unchanged (they are only reflected); the scalar density changes sign. In the mirror-extended primordial field of Chapter 9, the reflection of the hidden direction across the brane turns a universe of mass $+M$ on one side into a universe of mass $-M$ on the other side.
\item \textbf{T3, the Kohn--Sham level} (Section 15.7). In the Kohn--Sham approximation of Chapter 13 (and of Chapter 14 for dirac16complex00) the chirality map becomes a swap of the two components of every $2\times2$ block. With suitably transformed boundary conditions, the Kohn--Sham ground state and first excited state of the universe of mass $-M$ with the \textbf{same} $\lambda$ are the exact images of those of the universe of mass $+M$: the same energies, the same energy--momentum tensor, the opposite scalar density.
\end{itemize}

T1 and T3 both come from $\gamma^8$, yet T1 reverses the energy and needs $-\lambda$, while T3 keeps the energy and the same $\lambda$. Section 15.8 explains why: T1 is a statement about the classical bilinears of the field, T3 about the expectation values of the positive-norm quantum theory of Chapter 8, and the two differ by the matrix $B$, which $\gamma^8$ reverses. Section 15.9 shows what goes wrong when the boundary conditions are not transformed, and Section 15.10 reports the numerical tests.

\textbf{What is not proved.} This is the honesty rule of Chapter 0. The theorems show which partners exist and what they carry, and they show that the totals of a pair of the T1 type, treated as classical fields, are those of the empty state, so that the conservation of the totals and the gravitational field equations do not forbid such a pair (the separate conservation laws of the two members still do, unless each member is neutral; Section 15.4). They do \textbf{not} compute any process that creates a pair: no rate, no probability, no amplitude, and no dynamical big bang is derived. The Stage-1 document says it in one sentence that this book repeats: the pairing of the masses $\pm m$ is "a structural property of the equations, not a claim that universes of masses $\pm M$ are created in pairs" (\texttt{provenan\allowbreak{}ce/DIRAC\allowbreak{}16COMPLE\allowbreak{}X\_ARBITR\allowbreak{}ARY\_FIEL\allowbreak{}D.md}, §1.3, item 6). Section 15.11 lists precisely what is and what is not established. Chapter 16 takes up the question of pair creation, and Chapter 17 the question of matter and antimatter.

\textbf{Sources and their status.} The exact theory of the pairing theorems is the Wolfram package \texttt{wolfram/\allowbreak{}Dirac16C\allowbreak{}omplexPa\allowbreak{}iring.wl} with its verifier, and an independent sympy checker re-derives everything with its own code and compares the results. The programs and their reports:

\begin{Verbatim}[fontsize=\small]
wolfram/Dirac16ComplexPairing.wl, scripts/verify_dirac16complex_pairing.wls
scripts/check_dirac16complex_pairing.py          (independent, sympy)
artifacts/dirac16complex/pair-creation/
  wolfram-pairing-report.json           141 of 141 checks true
  python-pairing-report.json            172 of 172 checks true
  pairing-theory.json                   every map and formula, exact
  wolfram-dirac16complex00-report.json   46 of 46 checks true
  python-dirac16complex00-report.json    49 of 49 checks true
\end{Verbatim}

The Python pairing report includes the comparison with the Wolfram results (check \texttt{S5\_pairi\allowbreak{}ngAgrees\allowbreak{}WithWolf\allowbreak{}ram}). The last two reports belong to the Stage-5 field theory of dirac16complex00. Check names beginning with \texttt{\detokenize{PAIR_}} or \texttt{\detokenize{C00_}} are checks of the two Wolfram reports; the Python reports contain the same families under names beginning with \texttt{\detokenize{S5_}}. The matter--antimatter document \texttt{provenan\allowbreak{}ce/DIRAC\allowbreak{}16COMPLE\allowbreak{}X\_MATTER\allowbreak{}\_ANTIMAT\allowbreak{}TER.md} (its §7) uses the same results. Three things were \textbf{not finished} when this chapter was written: the numerical demonstration had been stopped part of the way through (Section 15.10 says exactly which runs exist), the Stage-5 cross-checker and the Stage-5 gate had not been run, and the Stage-5 documents \texttt{provenan\allowbreak{}ce/DIRAC\allowbreak{}16COMPLE\allowbreak{}X\_PAIR\_C\allowbreak{}REATION.\allowbreak{}md} and \texttt{provenan\allowbreak{}ce/DIRAC\allowbreak{}16COMPLE\allowbreak{}X00\_FIEL\allowbreak{}D\_THEORY\allowbreak{}.md} were not yet written. The exact results are therefore quoted as committed, and the chapter is built directly on the specification, the exact reports and the committed outputs.

\textbf{Conventions.} As everywhere in the book: counting from 0, coordinates $x_0,\dots,x_7$ with $x_4$ the time, $\eta=\mathrm{diag}(+1,+1,+1,+1,-1,-1,-1,-1)$, the notebook's gamma matrices $\gamma^0,\dots,\gamma^7$ (Section 2.8), $C=\gamma^0\gamma^1\gamma^2\gamma^3$, $\gamma^8=\gamma^0\gamma^1\cdots\gamma^7$, $B=-iC\gamma^4$, $\bar\Psi=\Psi^\dagger C$ and $D_\mu\Psi=\partial_\mu\Psi+\Omega_\mu\Psi$ with the canonical spin connection of Chapter 4. The two fields of Chapter 6, dirac16complex (anticommuting Grassmann components, quantized in Chapter 8) and dirac16complex00 (commuting components, a classical field), have the same Lagrangian,

\[
\mathcal L_{m,\lambda}=\sqrt{|g|}\,\Bigl[K-mS-\tfrac\lambda2S^2\Bigr],\qquad K=\tfrac12\bigl(\bar\Psi\gamma^\mu D_\mu\Psi-(D_\mu\bar\Psi)\gamma^\mu\Psi\bigr),\qquad S=\bar\Psi\Psi ,
\]

and the same field equations $E_{m,\lambda}[\Psi]=0$ and $\bar E_{m,\lambda}[\Psi]=0$ (Chapter 7), with

\[
E_{m,\lambda}[\Psi]:=\gamma^\mu D_\mu\Psi-(m+\lambda S)\Psi,\qquad \bar E_{m,\lambda}[\Psi]:=(D_\mu\bar\Psi)\gamma^\mu+(m+\lambda S)\bar\Psi .
\]

The energy--momentum tensor of Chapter 7 is

\[
T_{\mu\nu}=-\tfrac14\Bigl[\bar\Psi\gamma_\mu D_\nu\Psi+\bar\Psi\gamma_\nu D_\mu\Psi-(D_\mu\bar\Psi)\gamma_\nu\Psi-(D_\nu\bar\Psi)\gamma_\mu\Psi\Bigr]+g_{\mu\nu}\,\mathcal L_s,\qquad \mathcal L_s=\mathcal L/\sqrt{|g|},
\]

and we use it as an expression that can be evaluated on any configuration, a solution or not. The conserved current is $J^\mu=-i\bar\Psi\gamma^\mu\Psi$ (Section 8.12); in Gaussian normal gauge its time component is $J^4=\Psi^\dagger B\Psi$, and the charge of a slice is $Q=\int\sqrt{|g|}\,J^{x_4}\,d^7x$. We write $\mathcal L_{m,\lambda}[e,\Psi]$ or $T_{\mu\nu}[e,\Psi;m,\lambda]$ when the vielbein $e$ must be named.

\subsection{Tools: constant matrices, bilinears and parity}

Every proof of this chapter manipulates expressions $\Psi^\dagger X\Phi=\sum_{a,b}\Psi_a^\ast X_{ab}\Phi_b$, where $X$ is a $16\times16$ matrix of ordinary numbers (or ordinary functions of $x$) and $\Psi$, $\Phi$ are columns of field components or their derivatives. Chapter 6 calls them bilinears. This section collects the facts that we need about them.

\textbf{Six facts about $\gamma^8$, $C$ and $B$} (proved in Sections 2.9, 2.10, 2.12 and 8.6; checks \texttt{PAIR\_alg\allowbreak{}ebra\_gam\allowbreak{}ma8Prope\allowbreak{}rties}, \texttt{PAIR\_alg\allowbreak{}ebra\_CPr\allowbreak{}operties}, \texttt{PAIR\_alg\allowbreak{}ebra\_BPr\allowbreak{}operties}):

\begin{itemize}
\item (F1) $\gamma^8=\mathrm{diag}(-I_8,I_8)$ is real and diagonal, hence Hermitian, and $(\gamma^8)^2=1$.
\item (F2) $\gamma^8\gamma^a=-\gamma^a\gamma^8$ for every $a$. The curved gammas $\gamma^\mu=e_a{}^\mu\gamma^a$ and $\gamma_\mu=g_{\mu\nu}\gamma^\nu$ of any vielbein are combinations of the $\gamma^a$ with ordinary functions as coefficients, so $\gamma^8$ anticommutes with them too.
\item (F3) $\gamma^8$ commutes with $C$ (a product of four gammas) and with every $S^{ab}=\tfrac14[\gamma^a,\gamma^b]$ (two gammas), hence with every spinor connection $\Omega_\mu=\tfrac12\omega_{\mu ab}S^{ab}$, in every gravitational field.
\item (F4) $C$ is real and symmetric with $C^2=1$, and every $C\gamma^a$ is real and antisymmetric.
\item (F5) $B=-iC\gamma^4$ is Hermitian, $B^2=1$, and $BC=-i\gamma^4$.
\item (F6) $\gamma^8B\gamma^8=-B$.
\end{itemize}

\emph{Proof of (F6).} $\gamma^8B\gamma^8=-i\,\gamma^8C\gamma^4\gamma^8=-i\,C\gamma^8\gamma^4\gamma^8=-i\,C(-\gamma^4)(\gamma^8)^2=+iC\gamma^4=-B$, by (F3), (F2) and (F1). $\square$

In words: $\gamma^8$ reverses the matrix $B$. Chapter 8 showed that $B$ is the matrix of the canonical anticommutator, the Krein metric of the quantum theory (Section 8.7). Fact (F6) will decide the difference between T1 and T3.

\textbf{Lemma 15.1 (constant matrices inside bilinears).} Let $R$ and $X$ be constant $16\times16$ matrices of ordinary numbers, and let $\Psi$, $\Phi$ be columns of commuting components or of Grassmann components. Then

\[
(R\Psi)^\dagger X(R\Phi)=\Psi^\dagger\bigl(R^\dagger XR\bigr)\Phi ,
\]

and the same holds when $\Phi$ is replaced by a derivative $\partial_\mu\Phi$ or $D_\mu\Phi$.

\emph{Proof.} The components of $R\Phi$ are $(R\Phi)_b=\sum_cR_{bc}\Phi_c$, and those of $(R\Psi)^\ast$ are $(R\Psi)_a^\ast=\sum_dR_{ad}^\ast\Psi_d^\ast$. For Grassmann components this uses the rule $(c\theta)^\ast=c^\ast\theta^\ast$ for a number $c$ and one generator $\theta$ (Section 5.9). Hence

\[
\sum_{a,b}(R\Psi)_a^\ast X_{ab}(R\Phi)_b=\sum_{d,c}\Psi_d^\ast\Bigl(\sum_{a,b}R_{ad}^\ast X_{ab}R_{bc}\Bigr)\Phi_c=\sum_{d,c}\Psi_d^\ast\bigl(R^\dagger XR\bigr)_{dc}\Phi_c .
\]

Only ordinary numbers were moved; every $\Psi^\ast$ still stands to the left of every $\Phi$, and no two field components were exchanged. So the computation is the same for both statistics. A constant $R$ commutes with $\partial_\mu$. $\square$

\textbf{Lemma 15.2 (parity under $\gamma^8$).} If $P$ is a product of $k$ frame gamma matrices (repetitions allowed), then $\gamma^8P\gamma^8=(-1)^kP$. We call $P$ \textbf{even} if $k$ is even and \textbf{odd} if $k$ is odd. The parity of a product of matrices is the product of their parities.

\emph{Proof.} Move $\gamma^8$ from the left of $P$ to its right, one factor at a time. Each passage gives a factor $-1$ by (F2), so $\gamma^8P=(-1)^kP\gamma^8$; multiply on the right by $\gamma^8$ and use $(\gamma^8)^2=1$. If $\gamma^8X\gamma^8=\epsilon_XX$ and $\gamma^8Y\gamma^8=\epsilon_YY$, then $\gamma^8XY\gamma^8=\gamma^8X\gamma^8\gamma^8Y\gamma^8=\epsilon_X\epsilon_YXY$. $\square$

The matrices that occur in this book fall into the following classes (check \texttt{PAIR\_alg\allowbreak{}ebra\_bil\allowbreak{}inearPar\allowbreak{}ities} tests all 456 matrices $C\gamma^a$, $C\gamma^aS^{bc}$, $CS^{bc}\gamma^a$ and the matrix $C$):

\begingroup
\small
\begin{longtable}{@{}>{\raggedright\arraybackslash}p{0.200\linewidth}>{\raggedright\arraybackslash}p{0.200\linewidth}>{\raggedright\arraybackslash}p{0.200\linewidth}>{\raggedright\arraybackslash}p{0.200\linewidth}@{}}
\toprule
\textbf{matrix $X$} & \textbf{number of gamma factors} & \textbf{$\gamma^8X\gamma^8$} & \textbf{where it occurs} \\
\midrule
\endfirsthead
\toprule
\textbf{matrix $X$} & \textbf{number of gamma factors} & \textbf{$\gamma^8X\gamma^8$} & \textbf{where it occurs} \\
\midrule
\endhead
$C$ & 4 & $+C$ & $S=\bar\Psi\Psi$, the mass term, $U(S)$ \\
$C\gamma^a$, $C\gamma^aS^{bc}$, $CS^{bc}\gamma^a$ & 5, 7, 7 & $-X$ & kinetic term, $T_{\mu\nu}$, current $J^\mu$ \\
$B=-iC\gamma^4$ & 5 & $-B$ & charge density $J^4=\Psi^\dagger B\Psi$ \\
$BC=-i\gamma^4$ & 9 & $-BC$ & scalar density in the expectation-value rule \\
$BC\gamma^a$ & 10 & $+BC\gamma^a$ & energy and current in the expectation-value rule \\
\bottomrule
\end{longtable}
\endgroup

The last two rows show the rule that will matter in Section 15.8: \textbf{multiplying by $B$ reverses the parity}, because $B$ is odd.

\subsection{Theorem T1: the chirality map}

\textbf{Theorem 15.3 (T1, the chirality map).} Let $e$ be any vielbein (so any metric and its canonical spin connection), $m$ and $\lambda$ any real numbers, and $\Psi$ any configuration of commuting or of Grassmann components, not necessarily a solution. Put $\Psi_-:=\gamma^8\Psi$. Then:

\begin{enumerate}
\item $\bar\Psi_-=\bar\Psi\gamma^8$, $D_\mu\Psi_-=\gamma^8D_\mu\Psi$ and $D_\mu\bar\Psi_-=(D_\mu\bar\Psi)\gamma^8$;
\item $S[\Psi_-]=S[\Psi]$ and $K[\Psi_-]=-K[\Psi]$;
\item $\mathcal L_{m,\lambda}[\Psi_-]=-\mathcal L_{-m,-\lambda}[\Psi]$;
\item $E_{-m,-\lambda}[\Psi_-]=-\gamma^8E_{m,\lambda}[\Psi]$ and $\bar E_{-m,-\lambda}[\Psi_-]=-\bar E_{m,\lambda}[\Psi]\gamma^8$; hence $\Psi$ solves the field equations with $(m,\lambda)$ if and only if $\gamma^8\Psi$ solves them with $(-m,-\lambda)$;
\item $T_{\mu\nu}[\Psi_-;-m,-\lambda]=-T_{\mu\nu}[\Psi;m,\lambda]$ for all $\mu,\nu$;
\item $J^\mu[\Psi_-]=-J^\mu[\Psi]$, and the charge of every slice changes sign, $Q[\Psi_-]=-Q[\Psi]$.
\end{enumerate}

\emph{Proof.} Step 1 (adjoint and derivatives). By Lemma 15.1 with $R=\gamma^8$, which is Hermitian (F1), $\bar\Psi_-=(\gamma^8\Psi)^\dagger C=\Psi^\dagger\gamma^8C=\Psi^\dagger C\gamma^8=\bar\Psi\gamma^8$ by (F3). Since $\gamma^8$ is constant and commutes with $\Omega_\mu$ (F3), $D_\mu(\gamma^8\Psi)=\gamma^8\partial_\mu\Psi+\Omega_\mu\gamma^8\Psi=\gamma^8(\partial_\mu\Psi+\Omega_\mu\Psi)=\gamma^8D_\mu\Psi$, and in the same way $D_\mu\bar\Psi_-=\partial_\mu(\bar\Psi\gamma^8)-\bar\Psi\gamma^8\Omega_\mu=(D_\mu\bar\Psi)\gamma^8$.

Step 2 (scalar and kinetic term). $S[\Psi_-]=\bar\Psi\gamma^8\gamma^8\Psi=S[\Psi]$ by (F1). For the first half of $K$, $\bar\Psi_-\gamma^\mu D_\mu\Psi_-=\bar\Psi\gamma^8\gamma^\mu\gamma^8D_\mu\Psi=-\bar\Psi\gamma^\mu D_\mu\Psi$ by (F2) and (F1), and the second half changes sign in the same way. So $K[\Psi_-]=-K[\Psi]$. In the language of Lemma 15.2: the scalar has an even matrix, the kinetic term odd ones.

Step 3 (Lagrangian). With $U=\tfrac\lambda2S^2$, which is even in $S$,

\[
\mathcal L_{m,\lambda}[\Psi_-]=\sqrt{|g|}\Bigl[-K-mS-\tfrac\lambda2S^2\Bigr]=-\sqrt{|g|}\Bigl[K-(-m)S-\tfrac{(-\lambda)}2S^2\Bigr]=-\mathcal L_{-m,-\lambda}[\Psi].
\]

Step 4 (field equations). Using Step 1, (F2) and $S[\Psi_-]=S$,

\[
E_{-m,-\lambda}[\Psi_-]=\gamma^\mu\gamma^8D_\mu\Psi-(-m-\lambda S)\gamma^8\Psi=-\gamma^8\gamma^\mu D_\mu\Psi+\gamma^8(m+\lambda S)\Psi=-\gamma^8E_{m,\lambda}[\Psi],
\]

and $\bar E_{-m,-\lambda}[\Psi_-]=(D_\mu\bar\Psi)\gamma^8\gamma^\mu-(m+\lambda S)\bar\Psi\gamma^8=-\bigl[(D_\mu\bar\Psi)\gamma^\mu+(m+\lambda S)\bar\Psi\bigr]\gamma^8$. Since $\gamma^8$ is invertible, one side vanishes exactly when the other does.

Step 5 (energy--momentum tensor). Each of the four bilinears in the bracket of $T_{\mu\nu}$ contains exactly one curved gamma between an adjoint and a field, so each changes sign as in Step 2. The last term is $g_{\mu\nu}\mathcal L_s$, and by Step 3 (divided by $\sqrt{|g|}$) $\mathcal L_s[\Psi_-;-m,-\lambda]=-\mathcal L_s[\Psi;m,\lambda]$. So every term of $T_{\mu\nu}[\Psi_-;-m,-\lambda]$ is minus the corresponding term of $T_{\mu\nu}[\Psi;m,\lambda]$.

Step 6 (current and charge). $J^\mu[\Psi_-]=-i\bar\Psi\gamma^8\gamma^\mu\gamma^8\Psi=-J^\mu[\Psi]$. The charge is the integral of $\sqrt{|g|}\,J^{x_4}$ over the same slice with the same $\sqrt{|g|}$, so it changes sign.

Step 7 (both statistics). Every step multiplied matrices of ordinary numbers or used Lemma 15.1; nowhere were two field components exchanged. The proof therefore holds word for word for Grassmann components. (In the notebook's basis the map sends $\Psi_a\to-\Psi_a$ for $a\le7$ and $\Psi_a\to\Psi_a$ for $a\ge8$; for Grassmann generators this is a map of the algebra onto itself that respects all products, check \texttt{PAIR\_alg\allowbreak{}ebra\_chi\allowbreak{}ralBlock\allowbreak{}Structur\allowbreak{}e}.) $\square$

\textbf{Why $\lambda$ must change sign.} The potential $\tfrac\lambda2S^2$ is even, so it does not change sign when $K$ does. At a fixed coupling the map fails: $\mathcal L_{m,\lambda}[\gamma^8\Psi]+\mathcal L_{-m,\lambda}[\Psi]=\sqrt{|g|}\,[-K-mS-\tfrac\lambda2S^2+K+mS-\tfrac\lambda2S^2]=-\lambda\sqrt{|g|}\,S^2$, which is not zero when $\lambda\ne0$ and $S\ne0$. This is erratum E2 of \texttt{handoff/\allowbreak{}specs/CO\allowbreak{}NTRACT.m\allowbreak{}d} (§11); the verifiers confirm the failure (checks \texttt{PAIR\_T1g\allowbreak{}eneric\_n\allowbreak{}aiveFixe\allowbreak{}dLambdaF\allowbreak{}ails}, \texttt{PAIR\_T1g\allowbreak{}rassmann\allowbreak{}\_naiveFi\allowbreak{}xedLambd\allowbreak{}aFails}, \texttt{PAIR\_T1j\allowbreak{}ets\_naiv\allowbreak{}eFixedLa\allowbreak{}mbdaFail\allowbreak{}s}). For $\lambda=0$ the two members of the pair are the same free field with the masses $+m$ and $-m$.

\textbf{T1 relates two theories; it is not a symmetry of one.} Since $E_{m,\lambda}$ and $E_{-m,-\lambda}$ differ only in the sign of the mass term, $E_{m,\lambda}[\gamma^8\Psi]=E_{-m,-\lambda}[\gamma^8\Psi]-2(m+\lambda S)\gamma^8\Psi$. If $\Psi$ is a solution with $(m,\lambda)$, the first term vanishes by item 4, and $\gamma^8\Psi$ solves the original equations only where $(m+\lambda S)\Psi=0$. The image of a solution is a solution of the \emph{partner} theory (check \texttt{PAIR\_T1j\allowbreak{}ets\_onSh\allowbreak{}ellImage}, which also confirms that the image is not a solution with $(m,\lambda)$ or $(-m,\lambda)$ at the test points).

\textbf{A second view: reversing the frame.} Replace the vielbein $e_\mu{}^a$ by $-e_\mu{}^a$ and keep $\Psi$. The metric $g=e\,\eta\,e^T$ and $\sqrt{|g|}$ do not change; the curved gammas $\gamma^\mu=e_a{}^\mu\gamma^a$ change sign; the Christoffel symbols do not change, and neither does the spin connection $\omega_\mu{}^a{}_b=e_b{}^\nu(\Gamma^\rho{}_{\mu\nu}e_\rho{}^a-\partial_\mu e_\nu{}^a)$ of Section 4.11, because each of its terms contains two factors of the vielbein. Hence $K\to-K$, $S\to S$ and $\mathcal L_{m,\lambda}[-e,\Psi]=-\mathcal L_{-m,-\lambda}[e,\Psi]$: the same statement as item 3. Combining the two changes, $\mathcal L_{m,\lambda}[-e,\gamma^8\Psi]=\mathcal L_{m,\lambda}[e,\Psi]$. This is no accident: $\gamma^8=(\gamma^0\gamma^1)(\gamma^2\gamma^3)(\gamma^4\gamma^5)(\gamma^6\gamma^7)$ is the product of four rotations by $\pi$, $\exp(\pi S^{ab})=\gamma^a\gamma^b$ in the planes $(0,1)$, $(2,3)$, $(4,5)$ and $(6,7)$ (Exercise 2.14; the planes $(4,5)$ and $(6,7)$ are rotation planes because $\eta^{44}\eta^{55}=\eta^{66}\eta^{77}=+1$), and it turns every frame vector into its negative. So $(e,\Psi)\to(-e,\gamma^8\Psi)$ is a local frame rotation of spinor norm $+1$, under which the Lagrangian is invariant (Chapter 6; the checks are listed at the end of this section). The content of T1 can therefore be put in one sentence: \textbf{the theory with $(-m,-\lambda)$ is the theory with $(m,\lambda)$ whose action is multiplied by $-1$}, written in the relabelled field $\gamma^8\Psi$. Multiplying an action by $-1$ does not change its field equations, but it reverses everything that is computed from the action by a variation with respect to the metric (the energy--momentum tensor, Section 5.8) and the canonical structure of the quantum theory (Section 15.4).

\textbf{Worked example: a quantum at rest and its image.} Take flat space, $\lambda=0$ and the rest state of Section 8.8,

\[
u_+=\tfrac12\bigl(e_0-e_4-i\,e_9-i\,e_{13}\bigr),\qquad -i\gamma^4u_+=u_+,\qquad Cu_+=u_+,\qquad Bu_+=u_+ ,
\]

where $e_n$ is the column with 1 in place $n$. The field $\Psi_+=e^{-imx_4}u_+$ depends on $x_4$ only, so $\gamma^\mu D_\mu\Psi_+=\gamma^4\partial_4\Psi_+=-im\gamma^4\Psi_+=m\Psi_+$: it solves the free equation with mass $m$. Its image is $\Psi_-=\gamma^8\Psi_+=e^{-imx_4}\tilde u$ with $\tilde u:=\gamma^8u_+=\tfrac12(-e_0+e_4-i\,e_9-i\,e_{13})$, because $\gamma^8$ multiplies the components 0 to 7 by $-1$. With the rows of $\gamma^4$ and $B$ quoted in Section 11.4 and the rows of $C$ used in Section 8.8 ($(\gamma^4v)_0=-v_{13}$, $(\gamma^4v)_4=v_9$, $(\gamma^4v)_9=-v_4$, $(\gamma^4v)_{13}=v_0$; $(Cv)_0=-v_4$, $(Cv)_4=-v_0$, $(Cv)_9=v_{13}$, $(Cv)_{13}=v_9$; $(Bv)_0=iv_9$, $(Bv)_4=-iv_{13}$, $(Bv)_9=-iv_0$, $(Bv)_{13}=iv_4$) one finds $-i\gamma^4\tilde u=-\tilde u$, $C\tilde u=\tilde u$ and $B\tilde u=-\tilde u$ (Exercise 15.4). Hence $\gamma^4\partial_4\Psi_-=m(-i\gamma^4)\Psi_-=-m\Psi_-$: the image solves the free equation with mass $-m$, as item 4 says. The classical densities are:

\begingroup
\small
\begin{longtable}{@{}>{\raggedright\arraybackslash}p{0.200\linewidth}>{\raggedright\arraybackslash}p{0.200\linewidth}>{\raggedright\arraybackslash}p{0.200\linewidth}>{\raggedright\arraybackslash}p{0.200\linewidth}@{}}
\toprule
\textbf{quantity} & \textbf{$\Psi_+$ (mass $m$)} & \textbf{$\Psi_-=\gamma^8\Psi_+$ (mass $-m$)} & \textbf{sum} \\
\midrule
\endfirsthead
\toprule
\textbf{quantity} & \textbf{$\Psi_+$ (mass $m$)} & \textbf{$\Psi_-=\gamma^8\Psi_+$ (mass $-m$)} & \textbf{sum} \\
\midrule
\endhead
charge density $J^4=\Psi^\dagger B\Psi$ & $+1$ & $-1$ & 0 \\
scalar density $S=\Psi^\dagger C\Psi$ & $+1$ & $+1$ & 2 \\
energy density $\rho=m\,\Psi^\dagger B\Psi$ & $+m$ & $-m$ & 0 \\
\bottomrule
\end{longtable}
\endgroup

The energy density is the component $T_{44}$; for a field that oscillates like $e^{-i\varepsilon x_4}$ in flat space with $\lambda=0$ it equals $\varepsilon\,\Psi^\dagger B\Psi$ (the derivation of Section 13.14, with $\mathcal L_s=0$ on shell when $\lambda=0$), and here $\varepsilon=m$ for both members.

\textbf{How T1 was verified.} The identities of items 1 to 6 were checked in four independent ways, each in both pairing reports: as polynomial identities in completely generic independent symbols for the vielbein, the connection and the field, which is a check for every gravitational field (the Python checker also derives the Euler--Lagrange expressions from the Lagrangian); in a Grassmann algebra with 288 odd generators (Wolfram) and with 1440 odd generators (Python; measurement \texttt{S5\_T1gra\allowbreak{}ssmann} of \texttt{python-p\allowbreak{}airing-r\allowbreak{}eport.js\allowbreak{}on}); with exact jets of the general non-diagonal test vielbein G1 of Stage 1 at three points, including the on-shell images, the conservation of both energy--momentum tensors and the reversal of the frame (13 checks); and for generic fields of all eight coordinates in the primordial field of Chapter 9 with an arbitrary function $a_4$, all 64 components of $T_{\mu\nu}$. The Lagrangian part is also check \texttt{ALG\_gamm\allowbreak{}a8Map} of Stage 1.

\begin{Verbatim}[fontsize=\small]
PAIR_T1generic_lagrangian            PAIR_T1generic_fieldEquation
PAIR_T1generic_conjugateFieldEquation PAIR_T1generic_emtAll36
PAIR_T1generic_current               S5_T1generic_eulerLagrangeFromLagrangian
PAIR_T1grassmann_lagrangian          PAIR_T1grassmann_diracOperator
PAIR_T1grassmann_emt                 PAIR_T1grassmann_current
PAIR_T1jets_* (13 checks)            PAIR_T1primordial_* (6 checks)
PAIR_algebra_gamma8InIdentityComponent
PAIR_T1jets_vielbeinSignFlipGeometry PAIR_T1jets_vielbeinSignFlipIsT1
PAIR_T1jets_gamma8WithFrameSignIsSymmetry
\end{Verbatim}

\subsection{The chiral pair: zero total energy--momentum and charge}

\textbf{Corollary 15.4 (the chiral pair).} Assume:

\begin{itemize}
\item (a) both members live in the same gravitational field, that is, they are two fields on one spacetime with one vielbein $e$ and one metric $g$;
\item (b) the first member is a configuration $\Psi_+$ of the theory with $(m,\lambda)$, and the second is $\Psi_-=\gamma^8\Psi_+$, a configuration of the theory with $(-m,-\lambda)$ (for $\lambda=0$ both members belong to the same free field, with masses $+m$ and $-m$);
\item (c) the energy--momentum tensor, current and action of the pair are the sums of those of the members, each computed with its own parameters.
\end{itemize}

Then, at every point and every time $x_4$, in particular at $x_4=0$,

\[
\begin{aligned}
&T^{\mathrm{pair}}_{\mu\nu}=T_{\mu\nu}[\Psi_+;m,\lambda]+T_{\mu\nu}[\Psi_-;-m,-\lambda]=0,\qquad J^\mu_++J^\mu_-=0,\qquad Q_++Q_-=0,\\
&\mathcal L^{\mathrm{pair}}=\mathcal L_{m,\lambda}[\Psi_+]+\mathcal L_{-m,-\lambda}[\Psi_-]=0,\qquad S^{\mathrm{pair}}=2S_+ .
\end{aligned}
\]

If $\Psi_+$ solves its field equations, $\Psi_-$ solves its own. With the pair as the only source, the Einstein equations $G_{\mu\nu}=\kappa T^{\mathrm{pair}}_{\mu\nu}$, and equally the Einstein--Lovelock equations of Section 9.9, are the \textbf{source-free} equations.

\emph{Proof.} Items 5, 6, 3 and 2 of Theorem 15.3, applied to $\Psi=\Psi_+$, give the four sums and $S[\Psi_-]=S[\Psi_+]$; item 4 gives the statement on solutions. $\square$

(Checks \texttt{PAIR\_tot\allowbreak{}als\_fiel\allowbreak{}dLevelCh\allowbreak{}iralPair}, in the primordial field with an arbitrary $a_4$, all 64 components; \texttt{PAIR\_T1j\allowbreak{}ets\_pair\allowbreak{}EMTAndCu\allowbreak{}rrentVan\allowbreak{}ish}, on shell in the general field G1; in the matter--antimatter reports, \texttt{MA\_M4\_pa\allowbreak{}irEMTAnd\allowbreak{}CurrentG\allowbreak{}1\_commut\allowbreak{}ing} and \texttt{MA\_M4\_pa\allowbreak{}irEMTG1\_\allowbreak{}grassman\allowbreak{}n}.)

\textbf{What the corollary means.} Take the "empty" state $\Psi=0$: it has $T_{\mu\nu}=0$, no charge and no action. A T1 pair of classical fields has the same totals. The conservation of the \textbf{totals} (energy, momentum, charge) and the constraints of the gravitational field equations therefore do not forbid the appearance of such a pair in place of the empty state. This is the precise sense in which the pairing is \textbf{consistent} with the notebook's hypothesis. Four remarks keep it in proportion.

\begin{enumerate}
\item \emph{No mechanism is contained in it.} The field equations of each member contain the field in every term, so $\Psi=0$ is a solution. For the mode equations of Chapters 8 and 11, which are ordinary differential equations in $x_4$ in the good sector, the uniqueness theorem of Picard and Lindelöf (Section 10.2) says that data which vanish at one time stay zero at all times (the zero function is a solution; the set of times at which a solution vanishes is closed, by continuity, and open, by the uniqueness near each of its points, so it is the whole interval of existence): the classical equations do not create a pair out of the empty state. A creation process would need something beyond them, such as a quantum transition amplitude, a boundary condition at a singular moment or an interaction between the members; none is derived in this project.
\item \emph{The separate conservation laws still hold.} In the theory as built no term of the action contains both members, so each member obeys its own field equation, $E_{m,\lambda}[\Psi_+]=0$ and $E_{-m,-\lambda}[\Psi_-]=0$. The proof that the charge is conserved (Section 5.7 for the principle, Section 8.12 for the current; Theorem M1 of the matter--antimatter document \texttt{provenan\allowbreak{}ce/DIRAC\allowbreak{}16COMPLE\allowbreak{}X\_MATTER\allowbreak{}\_ANTIMAT\allowbreak{}TER.md}, its §4, gives the complete proof with its machine checks) uses only the field equation of the field whose current it is, so it applies to each member alone: $Q_+$ and $Q_-$ are conserved separately. In a gravitational field that does not depend on $x_4$, the shift of one member alone in $x_4$ is a symmetry as well, and Noether's theorem with the translations of Example 2 of Section 5.7 makes the energies $E_+$ and $E_-$ conserved separately. The empty state has $Q_+=0$ and $E_+=0$, and a conserved quantity that differs between two states forbids the transition between them. So a transition from the empty state to a T1 pair is \textbf{forbidden unless $Q_+=0$} (and, in such a static field, unless also $E_+=0$), although the totals are zero. Chapter 16 proves this as Proposition 16.5 (Section 16.9) and discusses what an interaction between the members, which the theory does not contain, would change (OPEN).
\item \emph{The corollary holds for every configuration, solution or not.} It is automatic, because the partner theory has the action of the original one multiplied by $-1$ (Section 15.3). It therefore selects no particular universe and says nothing about which universe is realized.
\item \emph{In eight-dimensional Einstein gravity a T1 pair cannot be the source of the primordial field.} Its total energy--momentum is zero, while the primordial field of Chapter 9 needs a nonzero, negative energy density in eight-dimensional Einstein gravity ($\rho_{\mathrm{req}}=-21H^2/\kappa$ for the static member, Section 9.15). If the pair is the only matter, the geometry must solve the source-free equations, which the primordial field does not in Einstein gravity (Section 9.10). For the Einstein--Lovelock equations with the terms of orders 2 and 3 the question is \textbf{OPEN}: these terms were not computed for this field, neither by the notebook nor by the project (Section 9.9; row L25 of the ledger of Section 0.9; Section 18.5). If some choice of the notebook's constants made the field a solution of the vacuum Einstein--Lovelock equations, a T1 pair would fit it exactly (Section 16.10). The pairing theorems do not change the sourcing analysis of Chapters 9 and 13.
\end{enumerate}

\textbf{The quantum field.} Corollary 15.4 is a statement about the fields treated as classical fields: two classical fields on one spacetime, the second in exactly the configuration $\gamma^8\Psi_+$ and with the opposite classical energy density. For the Grassmann field dirac16complex the energies and charges of the quantum theory are operators, and one must ask what the partner of mass $-M$ is as a quantum field. Chapter 8 found the canonical anticommutator $\{\Psi_a(x),\Psi_b^\dagger(y)\}=B_{ab}\,\delta^7(x-y)/\sqrt{|g|}$ in Gaussian normal gauge. For the image field $\Psi_-=\gamma^8\Psi$, built from the same operators, Lemma 15.1 and (F6) give

\[
\{\Psi_-,\Psi_-^\dagger\}=\gamma^8\{\Psi,\Psi^\dagger\}\gamma^8=\gamma^8B\gamma^8\,\frac{\delta^7}{\sqrt{|g|}}=-B\,\frac{\delta^7}{\sqrt{|g|}} :
\]

the image field carries the \textbf{reversed Krein metric} $-B$. This is exactly the canonical structure of the Lagrangian $-\mathcal L_{-m,-\lambda}$: the rule of Section 8.6 gives the anticommutator $i\mathcal K^{-1}\delta^7$, where $\mathcal K$ is the matrix in front of $\partial_4\Psi$ in the Lagrangian (Section 8.5 calls this matrix $K$; in this chapter the letter $K$ denotes the kinetic term of $\mathcal L$), and multiplying the Lagrangian by $-1$ multiplies $\mathcal K$ by $-1$ and hence the anticommutator by $-1$ (Exercise 15.3). Two different "universes of mass $-M$" can therefore be meant:

\begin{itemize}
\item \textbf{the image field} $\Psi_-=\gamma^8\Psi_+$, built from the same operators and acting on the same states, with the metric $-B$;
\item \textbf{an independently quantized theory} of mass $-m$: a second field with its own canonical anticommutator $+B$ and the positive Fock space of Section 8.10.
\end{itemize}

Neither of them is a second, independent universe whose energy and charge cancel those of the first. We take them in turn.

\textbf{The image reading cancels only formally.} The exact Fock-space model of the two Stage-5 verifiers (four rest modes with $m=1$, the filled Dirac sea, normal ordering as the subtraction of the sea value) confirms the operator identities $H[\Psi_-;-m]=-H[\Psi;m]$, $Q[\Psi_-]=-Q[\Psi]$ and $S[\Psi_-]=S[\Psi]$ for the energy operators $H=\Psi^\dagger Bh\Psi$ of these free modes, the charge and the scalar density, and the energy identity survives normal ordering, $:\!H[\Psi_-;-m]\!:\;=-:\!H[\Psi;m]\!:$ (checks \texttt{PAIR\_T1k\allowbreak{}rein\_ima\allowbreak{}geOperat\allowbreak{}orIdenti\allowbreak{}ties} and \texttt{S5\_T1kre\allowbreak{}in\_image\allowbreak{}Operator\allowbreak{}Identiti\allowbreak{}es}). The corresponding identity for the energy--momentum tensor, $:\!T_{\mu\nu}[\gamma^8\Psi;-m,-\lambda]\!:\;=-:\!T_{\mu\nu}[\Psi;m,\lambda]\!:$, is not a check of the Fock model: it follows from item 5 of Theorem 15.3, because $\gamma^8$ acts linearly and at each point separately, so that the subtraction of the sea value commutes with the map (\texttt{pairing-\allowbreak{}theory.j\allowbreak{}son}, key \texttt{\detokenize{T1krein}}, entry \texttt{imageFie\allowbreak{}ld}). Evaluated with the formulas of $\mathcal L_{-m,-\lambda}$, every quantum of the image has the energy $-|\varepsilon|$ and the opposite charge (check \texttt{PAIR\_T1k\allowbreak{}rein\_ima\allowbreak{}geExpect\allowbreak{}ationVal\allowbreak{}ues}).

These identities do \textbf{not} describe a second universe. The image field is the same quantum system as $\Psi_+$: the same operators, the same space of states and the same state. The sums $Q_++Q[\Psi_-]=0$ and $:\!H_+\!:+:\!H[\Psi_-;-m]\!:\;=0$ therefore have the form $X+(-X)=0$ and hold in every state; nothing in them is compensated by anything else. Measured with its own canonical structure, the image field has the energy and the charge of $\Psi_+$ itself. To see this, recall from Section 8.6 that for a field with the anticommutator $\{\Psi,\Psi^\dagger\}=A$ and a numerical matrix $X$ one has $i[\Psi^\dagger X\Psi,\Psi]=-iAX\Psi$. Take flat space and a mode of momentum $k$ with the mode Hamiltonian $h_k(m)$ of Section 8.9. The energy operator is $\Psi^\dagger Bh_k(m)\Psi$ (by the rule of Section 8.12 its expectation value in a normalized mode $u$ is $u^\dagger B^2h_ku=u^\dagger h_ku$). For $\Psi_+$, with $A=B$, the formula gives $-iB^2h_k(m)\Psi=-ih_k(m)\Psi$, which is the field equation $i\partial_4\Psi=h_k(m)\Psi$ of Section 8.9: $H_+$ generates the time evolution. For the image field, $A=-B$ and the energy matrix of $\mathcal L_{-m,-\lambda}$ is $Bh_k(-m)$, so $H[\Psi_-;-m]$ generates $-i(-B)Bh_k(-m)\Psi_-=+ih_k(-m)\Psi_-$. The actual evolution of the image is $\partial_4\Psi_-=\gamma^8\partial_4\Psi=-i\gamma^8h_k(m)\Psi=-ih_k(-m)\Psi_-$, because $\gamma^8h_k(m)\gamma^8=h_k(-m)$ (the mass term $-im\gamma^4$ is odd under $\gamma^8$ and the momentum terms $\gamma^4\gamma^j$ are even, Lemma 15.2). So $H[\Psi_-;-m]$ runs the time evolution backwards, and the image field's own generator of the time evolution is $-H[\Psi_-;-m]=+H_+$. In the same way, with $X=B$, the operator $Q[\Psi_-]$ generates the inverse phase rotation, and the image field's own charge is $-Q[\Psi_-]=+Q_+$. Its own source of gravity is the energy--momentum tensor of its own Lagrangian $-\mathcal L_{-m,-\lambda}$, obtained by varying that action with respect to the metric (Section 5.8; the remark at the end of Section 15.3), and this is $-T_{\mu\nu}[\Psi_-;-m,-\lambda]=+T_{\mu\nu}[\Psi_+;m,\lambda]$ by item 5 of Theorem 15.3. The values $-|\varepsilon|$ and $-1$ per quantum are expectation values of the formulas of $\mathcal L_{-m,-\lambda}$, and these are not the image field's own Hamiltonian and charge. The committed matter--antimatter report records exactly this (\texttt{artifact\allowbreak{}s/dirac1\allowbreak{}6complex\allowbreak{}/matter-\allowbreak{}antimatt\allowbreak{}er/wolfr\allowbreak{}am-matte\allowbreak{}r-antima\allowbreak{}tter-rep\allowbreak{}ort.json}, measurement \texttt{M5\_impli\allowbreak{}cation}, entry \texttt{kreinLev\allowbreak{}elCaveat}), with the matrix identities behind it in its check \texttt{MA\_M4\_kr\allowbreak{}einOnePa\allowbreak{}rticle}.

\textbf{The independent reading does not cancel.} Quantized on its own, with $+B$, the theory of mass $-m$ has the modes $\gamma^8u$ for the modes $u$ of the $+m$ theory, with the same frequencies, and each quantum has a positive energy, the charge $+1$ for a particle and the \textbf{opposite} scalar density (checks \texttt{PAIR\_T1k\allowbreak{}rein\_min\allowbreak{}usMModes}, \texttt{PAIR\_T1k\allowbreak{}rein\_ind\allowbreak{}ependent\allowbreak{}CARPlusB} and \texttt{PAIR\_T1k\allowbreak{}rein\_ind\allowbreak{}ependent\allowbreak{}Expectat\allowbreak{}ionValue\allowbreak{}s}). Two cases must be distinguished.

\begin{itemize}
\item With the parameters $(-m,\lambda)$, the same coupling, and in the state with the same mode occupations (the mapped state), the energy--momentum and the charge are those of the $+M$ universe: in the free Fock model this is the last check named, and with the interaction it is the Kohn--Sham statement T3 of Section 15.7, the Kohn--Sham form of the mirror pairing T2. \textbf{This pair doubles; it does not cancel.}
\item With the parameters $(-m,-\lambda)$ and $\lambda\ne0$, the kinetic and mass energies are again positive and add, while the self-interaction $\tfrac{(-\lambda)}2S^2$ is the negative of $\tfrac\lambda2S^2$ at the same value of $S^2$, which the mapped state has (its scalar density is $-S$). The interaction energy of the $-M$ universe then has the opposite sign, and the total is neither zero nor twice that of one universe.
\end{itemize}

In both cases the charges of the two universes are equal in the mapped state. They cancel only under an additional assumption on the state of the $-M$ universe (a state of the opposite charge, made of antiparticles), and even then the energies do not cancel (entry \texttt{kreinLev\allowbreak{}elCaveat}, part (b)).

The checks of the two readings, in the Fock model of the Wolfram pairing report (the Python report has the same checks with names beginning with \texttt{S5\_T1kre\allowbreak{}in\_}):

\begin{Verbatim}[fontsize=\small]
image field:        PAIR_T1krein_imageAnticommutatorMinusB
                    PAIR_T1krein_imageOperatorIdentities
                    PAIR_T1krein_imageExpectationValues
independent theory: PAIR_T1krein_minusMModes  PAIR_T1krein_independentCARPlusB
                    PAIR_T1krein_independentExpectationValues
\end{Verbatim}

The matter--antimatter document tabulates the energy $E$, charge $Q$ and scalar density $S$ of one quantum in three exact modes (its §7.2). The numbers are those of the measurement \texttt{M4.krein\allowbreak{}ModeFact\allowbreak{}s}, with the check \texttt{MA\_M4\_kr\allowbreak{}einModeF\allowbreak{}acts}, in the Python report \texttt{python-m\allowbreak{}atter-an\allowbreak{}timatter\allowbreak{}-report.\allowbreak{}json} of the same folder as the Wolfram report just named, and the Wolfram check \texttt{MA\_M4\_kr\allowbreak{}einOnePa\allowbreak{}rticle} verifies the matrix identities behind them. The column of the image field gives the expectation values of the formulas of $\mathcal L_{-m,-\lambda}$ under the anticommutator $-B$; by the discussion above they are not the image field's own energy and charge, which are those of the first column.

\begingroup
\small
\begin{longtable}{@{}>{\raggedright\arraybackslash}p{0.200\linewidth}>{\raggedright\arraybackslash}p{0.200\linewidth}>{\raggedright\arraybackslash}p{0.200\linewidth}>{\raggedright\arraybackslash}p{0.200\linewidth}@{}}
\toprule
\textbf{$m$, momentum $(k_0,k_1,k_2,k_3)$} & \textbf{$+m$ universe} & \textbf{image field, metric $-B$, formulas of $\mathcal L_{-m,-\lambda}$} & \textbf{independent $-m$ theory, metric $+B$} \\
\midrule
\endfirsthead
\toprule
\textbf{$m$, momentum $(k_0,k_1,k_2,k_3)$} & \textbf{$+m$ universe} & \textbf{image field, metric $-B$, formulas of $\mathcal L_{-m,-\lambda}$} & \textbf{independent $-m$ theory, metric $+B$} \\
\midrule
\endhead
1, $(1,1,2,3)$ & $(4,\,1,\,\tfrac14)$ & $(-4,\,-1,\,\tfrac14)$ & $(4,\,1,\,-\tfrac14)$ \\
3, $(1,1,1,2)$ & $(4,\,1,\,\tfrac34)$ & $(-4,\,-1,\,\tfrac34)$ & $(4,\,1,\,-\tfrac34)$ \\
$\tfrac35$, $(\tfrac45,0,0,0)$ & $(1,\,1,\,\tfrac35)$ & $(-1,\,-1,\,\tfrac35)$ & $(1,\,1,\,-\tfrac35)$ \\
\bottomrule
\end{longtable}
\endgroup

The energies are $E=\sqrt{m^2+k_0^2+k_1^2+k_2^2+k_3^2}$ and the scalar densities $m/E$. (Proof: write $h_k=mZ+P_k$ with $Z=-i\gamma^4=BC$ and the momentum part $P_k=-\gamma^4\sum_{j\ne4}k_j\gamma^j$; $Z^2=1$ and $ZP_k+P_kZ=0$ (Section 8.9), so $Zh_k+h_kZ=2m$, and for a normalized eigenvector $h_ku=Eu$ of $h_k$, which is Hermitian because there is no momentum along the extra times (Section 8.9), this gives $2E\,u^\dagger Zu=2m$.) For the first row $\sqrt{1+1+1+4+9}=4$.

\textbf{Conclusion.} The cancellation of Corollary 15.4 is a statement about \textbf{classical} fields: two independent classical fields with the Lagrangians $\mathcal L_{m,\lambda}$ and $\mathcal L_{-m,-\lambda}$, in the correlated configuration $\Psi_-=\gamma^8\Psi_+$, the second with the opposite classical energy density. Of the two quantum readings examined in the repository, the image field cancels only as an identity $X+(-X)=0$ within one quantum system, and an independently quantized partner does not cancel. \textbf{A quantum description of two independent universes whose energies and charges cancel is OPEN} (Section 16.6). The Fock-level results, moreover, are cited from Stage 5: all of their checks are true in the committed reports, but the matter--antimatter report labels its use of them PROVISIONAL until the Stage-5 gate has been run from a fresh clone (measurement \texttt{M4\_krein\allowbreak{}LevelSta\allowbreak{}tus} of the Wolfram matter--antimatter report).

For the \textbf{commuting} field dirac16complex00, which is a classical field, Corollary 15.4 is a statement about ordinary numbers and needs no hypothesis beyond (a) to (c); remark 2 above applies to it as well. Its classical charge density is itself indefinite and its classical energy is not bounded below (Chapter 7; checks \texttt{C00\_char\allowbreak{}ge\_indef\allowbreak{}inite} and \texttt{C00\_ener\allowbreak{}gy\_unbou\allowbreak{}ndedBelo\allowbreak{}w}), and the image of a configuration simply has the opposite energy and charge.

\subsection{Theorem T2: reflections of character minus one}

T1 reverses the energy. A different map relates $m$ and $-m$ without reversing the energy: a reflection. Chapter 6 derived it in flat space for the reflection of one coordinate; here we prove it in every gravitational field, where the reflection must act on the frame together with the field.

\textbf{Tools from Chapter 2.} A real vector $v=(v_0,\dots,v_7)$ with $n(v):=\eta(v,v)=\pm1$ gives the \textbf{unit vector} $u=\gamma(v)=\sum_av_a\gamma^a$ of Pin(4,4), with $u^2=n(u)$, $u^{-1}=n(u)\,u$, $\det u=1$ as a $16\times16$ matrix, and $u^\dagger=u^T$ because $u$ is real (Section 2.14). $u$ is \textbf{space-like} if $n(u)=+1$ and \textbf{time-like} if $n(u)=-1$. Two facts of Section 2.14 are used:

\begin{itemize}
\item the key identity $-u\,\gamma(w)\,u^{-1}=\gamma(R_uw)$ for every vector $w$, where $R_u$ is the reflection of vectors in the hyperplane orthogonal to $v$;
\item Theorem 2.8 with one factor, $u^TCu=-n(u)\,C$, hence $u^TC=-n(u)\,C\,u^{-1}$. The number $-n(u)$ is the Pin character of $u$; for a space-like $u$ it is $-1$.
\end{itemize}

\textbf{The reflected frame.} Let $e$ be a vielbein. We define a second vielbein $e_u$ in two steps. First, rotate the frame by the constant matrix $\Lambda$ that the spin matrix $u$ covers, $u^{-1}\gamma^au=\Lambda^a{}_b\gamma^b$; such a $\Lambda$ exists and lies in $\mathrm O(4,4)$ by Theorem 2.9. The local spin covariance theorem of Section 4.12 was stated for spin transformations, but its proof uses only two properties of the spin matrix, the covering relation and the determinant 1, and $u$ has both. So it applies to $u$ and gives, for the rotated frame, the spinor connection $u\Omega_\mu u^{-1}$ (the term $(\partial_\mu u)u^{-1}$ vanishes, $u$ being constant) and $D'_\mu(u\Psi)=u\,D_\mu\Psi$. The curved gammas of the rotated frame are $u\gamma^\mu u^{-1}$: the inverse vielbein of $e'_\mu{}^a=\Lambda^a{}_be_\mu{}^b$ is $e'_a{}^\mu=e_c{}^\mu(\Lambda^{-1})^c{}_a$, and the inverted covering relation $u\gamma^cu^{-1}=(\Lambda^{-1})^c{}_a\gamma^a$ gives $e'_a{}^\mu\gamma^a=e_c{}^\mu\,u\gamma^cu^{-1}$. Second, reverse all eight frame vectors; by Section 15.3 this changes the curved gammas by a sign and leaves the spinor connection unchanged. The result is a vielbein $e_u$ with

\[
\gamma_u^\mu=-u\,\gamma^\mu u^{-1},\qquad \Omega^{(u)}_\mu=u\,\Omega_\mu u^{-1},\qquad D^{(u)}_\mu(u\Psi)=u\,D_\mu\Psi .
\]

It describes the same metric, because $\{\gamma_u^\mu,\gamma_u^\nu\}=u\{\gamma^\mu,\gamma^\nu\}u^{-1}=2g^{\mu\nu}$. By the key identity, $\gamma_u^\mu=e_a{}^\mu\,\gamma(R_ue_a)$: every frame direction $a$ now carries the reflected vector $R_ue_a$ instead of $e_a$. So $e_u$ is the frame \textbf{reflected by $R_u$}; Stage 5 calls this the twisted frame action and writes $e\to e\,R_u$.

\textbf{Theorem 15.5 (T2, the mirror map).} Let $u$ be a unit vector, $e$ any vielbein, $e_u$ the reflected frame, and $\Psi$ any configuration of commuting or Grassmann components. Put $\Psi_u:=u\Psi$. Then

\[
S[e_u,\Psi_u]=-n(u)\,S[e,\Psi],\qquad K[e_u,\Psi_u]=n(u)\,K[e,\Psi],\qquad J^\mu[e_u,\Psi_u]=n(u)\,J^\mu[e,\Psi].
\]

For a \textbf{space-like} $u$ ($n(u)=+1$, Pin character $-1$):

\[
\begin{aligned}
&\mathcal L_{m,\lambda}[e_u,\Psi_u]=\mathcal L_{-m,\lambda}[e,\Psi],\qquad T_{\mu\nu}[e_u,\Psi_u;m,\lambda]=T_{\mu\nu}[e,\Psi;-m,\lambda],\\
&E_{-m,\lambda}[e_u,\Psi_u]=-u\,E_{m,\lambda}[e,\Psi],
\end{aligned}
\]

so $\Psi$ solves the field equations with $(m,\lambda)$ in the frame $e$ exactly when $u\Psi$ solves them with $(-m,\lambda)$, \textbf{the same coupling}, in the reflected frame; the energy--momentum tensor and the current are the same, and the scalar density is opposite. For a \textbf{time-like} $u$ ($n(u)=-1$) one gets instead $\mathcal L_{m,\lambda}[e_u,\Psi_u]=-\mathcal L_{-m,-\lambda}[e,\Psi]$, a map of the T1 type.

\emph{Proof.} Write $n=n(u)$. Step 1 (adjoint). By Lemma 15.1 and $u^\dagger=u^T$, $\bar\Psi_u=\Psi^\dagger u^TC=-n\,\Psi^\dagger C\,u^{-1}=-n\,\bar\Psi u^{-1}$. For the derivative, $D^{(u)}_\mu\bar\Psi_u=\partial_\mu\bar\Psi_u-\bar\Psi_u\,u\Omega_\mu u^{-1}=-n\bigl[(\partial_\mu\bar\Psi)u^{-1}-\bar\Psi\Omega_\mu u^{-1}\bigr]=-n\,(D_\mu\bar\Psi)\,u^{-1}$.

Step 2 (the bilinears). $S[e_u,\Psi_u]=(-n\bar\Psi u^{-1})(u\Psi)=-nS$. For the kinetic term,

\[
\bar\Psi_u\gamma_u^\mu D^{(u)}_\mu\Psi_u=\bigl(-n\bar\Psi u^{-1}\bigr)\bigl(-u\gamma^\mu u^{-1}\bigr)\bigl(uD_\mu\Psi\bigr)=n\,\bar\Psi\gamma^\mu D_\mu\Psi ,
\]

and in the same way $(D^{(u)}_\mu\bar\Psi_u)\gamma_u^\mu\Psi_u=n\,(D_\mu\bar\Psi)\gamma^\mu\Psi$, so $K\to nK$. The current has one gamma and no derivative: $J^\mu\to nJ^\mu$.

Step 3 (Lagrangian). For $n=+1$: $\mathcal L_{m,\lambda}[e_u,\Psi_u]=\sqrt{|g|}\,[K-m(-S)-\tfrac\lambda2S^2]=\mathcal L_{-m,\lambda}[e,\Psi]$. For $n=-1$: $\sqrt{|g|}\,[-K-mS-\tfrac\lambda2S^2]=-\mathcal L_{-m,-\lambda}[e,\Psi]$.

Step 4 (energy--momentum tensor). With $\gamma^u_\mu=g_{\mu\nu}\gamma_u^\nu=-u\gamma_\mu u^{-1}$, each bilinear of the bracket of $T_{\mu\nu}$ picks up the factor $n$ exactly as in Step 2, and $g_{\mu\nu}\mathcal L_s$ follows Step 3. For $n=+1$ every term of $T_{\mu\nu}[e_u,\Psi_u;m,\lambda]$ equals the corresponding term of $T_{\mu\nu}[e,\Psi;-m,\lambda]$.

Step 5 (field equation, $n=+1$). $E_{-m,\lambda}[e_u,\Psi_u]=(-u\gamma^\mu u^{-1})(uD_\mu\Psi)-(-m+\lambda(-S))u\Psi=-u\bigl[\gamma^\mu D_\mu\Psi-(m+\lambda S)\Psi\bigr]=-u\,E_{m,\lambda}[e,\Psi]$, and $u$ is invertible. As in Section 15.3 no two field components were exchanged, so the proof holds for both statistics. $\square$

\textbf{Remarks.}

\begin{itemize}
\item \emph{The other lift.} Without the reversal of the frame (the frame rotated only by $\Lambda$, curved gammas $u\gamma^\mu u^{-1}$), the kinetic term gets $-n$ instead of $n$, and a space-like $u$ gives $-\mathcal L_{m,-\lambda}$: this is erratum E3 of the contract and the "other lift" of Chapter 6 (check \texttt{PAIR\_T2f\allowbreak{}rame\_unt\allowbreak{}wistedSp\allowbreak{}acelikeC\allowbreak{}ontractE\allowbreak{}3}). Composing that lift with $\gamma^8$ gives back T2 (check \texttt{PAIR\_T2f\allowbreak{}rame\_gam\allowbreak{}ma8Times\allowbreak{}Untwiste\allowbreak{}dSpaceli\allowbreak{}ke}).
\item \emph{The canonical structure is preserved.} For $u=\gamma^0,\gamma^1,\gamma^2,\gamma^3$ (and $\gamma^4$) one has $uBu^\dagger=+B$, for $u=\gamma^5,\gamma^6,\gamma^7$ one has $uBu^\dagger=-B$ (check \texttt{PAIR\_alg\allowbreak{}ebra\_kre\allowbreak{}inUnderB\allowbreak{}asicRefl\allowbreak{}ections}; for example $\gamma^0B\gamma^0=-i\gamma^0C\gamma^4\gamma^0=-iC\gamma^4=B$, because $\gamma^0$ anticommutes with both $C$ and $\gamma^4$). The mirror map by a space-like basis direction therefore keeps the anticommutator $+B$: the mirror image of a state of the positive-norm quantum theory is again a state of that theory, with the same energies and charges.
\item \emph{The mirror pair doubles.} For a space-like $u$ the pair $\{(m,\lambda,e,\Psi),\ (-m,\lambda,e_u,u\Psi)\}$ has $T^{\mathrm{pair}}=2T$, $J^{\mathrm{pair}}=2J$, $\mathcal L^{\mathrm{pair}}=2\mathcal L$ and $S^{\mathrm{pair}}=0$ (check \texttt{PAIR\_tot\allowbreak{}als\_fiel\allowbreak{}dLevelMi\allowbreak{}rrorPair}). A T2 pair is a universe and its mirror image; it does not cancel.
\item \emph{For the Grassmann field without a reflection.} Chapter 6 found that the antilinear map $\Psi\to\gamma^8\Psi^\ast$ sends $\mathcal L_{m,\lambda}$ to $\mathcal L_{-m,\lambda}$ for Grassmann components; combined with a space-like reflection it becomes an exact symmetry that reverses the charge (the CP of Chapter 17).
\end{itemize}

\textbf{How T2 was verified.} For the eight basis vectors $\gamma^0,\dots,\gamma^7$ and for general rational space-like and time-like unit vectors, with exact jets of the general test vielbein G1: the character $-n(u)$ (\texttt{PAIR\_T2f\allowbreak{}rame\_cha\allowbreak{}racterIs\allowbreak{}MinusNor\allowbreak{}m}), the geometry of the reflected frames (\texttt{PAIR\_T2f\allowbreak{}rame\_fra\allowbreak{}meReflec\allowbreak{}tionsGeo\allowbreak{}metry}), $\mathcal L\to\mathcal L_{-m,\lambda}$ for space-like $u$ (\texttt{PAIR\_T2f\allowbreak{}rame\_twi\allowbreak{}stedSpac\allowbreak{}elikeMap\allowbreak{}sToMinus\allowbreak{}MSameLam\allowbreak{}bda}), the T1 type for time-like $u$ (\texttt{PAIR\_T2f\allowbreak{}rame\_twi\allowbreak{}stedTime\allowbreak{}like}), and the on-shell images (\texttt{PAIR\_T2f\allowbreak{}rame\_onS\allowbreak{}hellImag\allowbreak{}e}); the full table is \texttt{T2frameT\allowbreak{}able} in \texttt{pairing-\allowbreak{}theory.j\allowbreak{}son}.

\subsection{The mirror universe across the Z2 brane}

In the primordial field the mirror map becomes very concrete. Chapter 9 (Section 9.16) and Chapter 13 (Section 13.2) extend the static member of the notebook's field beyond the end $y=0$ of the notebook's chart by a mirror copy, glued along the \textbf{Z2 brane} $y=0$:

\[
ds^2=dy^2-dx_4^2+W(y)^2\Bigl[e^{2a_{4,0}}\bigl(dx_1^2+dx_2^2+dx_3^2\bigr)-e^{-2a_{4,0}}\bigl(dx_5^2+dx_6^2+dx_7^2\bigr)\Bigr],\qquad W(y)=e^{-H|y|}.
\]

Here $y$ is the proper hidden coordinate, the side $y<0$ is the notebook's patch and $y>0$ its mirror copy. We write $s_\mu=-1$ for $\mu=y$ and $s_\mu=+1$ for the other seven coordinates. The warp $W$ is even in $y$, and its derivative $W'$ is odd.

\textbf{The Dirac operator.} The vielbein is diagonal, $h=(1,\,We^{a_{4,0}}\,(\times3),\,1,\,We^{-a_{4,0}}\,(\times3))$ in the order $(y,x_1,\dots,x_7)$, so $\gamma^{x_y}=\gamma^0$, $\gamma^{x_l}=W^{-1}e^{-a_{4,0}}\gamma^l$ for $l=1,2,3$, $\gamma^{x_4}=\gamma^4$ and $\gamma^{x_j}=W^{-1}e^{a_{4,0}}\gamma^j$ for $j=5,6,7$. Every curved gamma is an even function of $y$ times the frame gamma with the same number. The diagonal-vielbein formula of Section 4.11 gives exactly six nonzero connection matrices, one for each warped direction $p\in\{1,2,3,5,6,7\}$:

\[
\Omega_p=c_p(y)\,S^{0p},\qquad c_l=-W'e^{a_{4,0}}\ (l=1,2,3),\qquad c_j=+W'e^{-a_{4,0}}\ (j=5,6,7),\qquad \Omega_y=\Omega_4=0 ,
\]

and every $c_p$ is an odd function of $y$. (For $y<0$ this is the connection of Section 9.13 written in the coordinate $y$; it gives $\gamma^\mu\Omega_\mu=3(W'/W)\gamma^0$, which is $+3H\gamma^0$ for $y<0$ and $-3H\gamma^0$ for $y>0$.)

\textbf{Lemma 15.6.} For $\Psi'(y,x):=\gamma^0\,\Psi(-y,x)$ (all other coordinates $x$ unchanged) and $y\ne0$:

\[
D_\mu\Psi'(y)=s_\mu\,\gamma^0\,(D_\mu\Psi)(-y)\quad\text{for every }\mu,\qquad \gamma^0\gamma^{x_\mu}(y)\gamma^0=-s_\mu\,\gamma^{x_\mu}(-y).
\]

\emph{Proof.} For $\mu=y$, the chain rule gives $\partial_y[\Psi(-y)]=-(\partial_y\Psi)(-y)$, and $\Omega_y=0$. For $\mu=4$ and the other coordinates, $\partial_\mu$ does not touch $y$. For a warped direction $p$ the connection term needs $S^{0p}\gamma^0=\tfrac12\gamma^0\gamma^p\gamma^0=-\tfrac12\gamma^p=-\gamma^0S^{0p}$, so that

\[
\begin{aligned}
\Omega_p(y)\,\gamma^0\Psi(-y)&=c_p(y)\,S^{0p}\gamma^0\Psi(-y)=-c_p(y)\,\gamma^0S^{0p}\Psi(-y)\\
&=\gamma^0\,c_p(-y)S^{0p}\Psi(-y)=\gamma^0\,\Omega_p(-y)\Psi(-y),
\end{aligned}
\]

using that $c_p$ is odd. Together, $D_p\Psi'(y)=\gamma^0(D_p\Psi)(-y)$. For the gammas, $\gamma^0\gamma^0\gamma^0=\gamma^0$ and $\gamma^0\gamma^a\gamma^0=-\gamma^a$ for $a\ne0$, that is $\gamma^0\gamma^a\gamma^0=-s_a\gamma^a$; the curved gammas are even functions of $y$ times frame gammas. $\square$

\textbf{Theorem 15.7 (T2 across the brane).} For $\Psi'(y)=\gamma^0\Psi(-y)$ and both statistics:

\[
\begin{aligned}
&\gamma^\mu D_\mu\Psi'(y)=-\gamma^0\,(\gamma^\mu D_\mu\Psi)(-y),\qquad S[\Psi'](y)=-S[\Psi](-y),\qquad K[\Psi'](y)=K[\Psi](-y),\\
&E_{-m,\lambda}[\Psi'](y)=-\gamma^0\,E_{m,\lambda}[\Psi](-y),\qquad \mathcal L_{-m,\lambda}[\Psi'](y)=\mathcal L_{m,\lambda}[\Psi](-y),\\
&T_{\mu\nu}[\Psi';-m,\lambda](y)=s_\mu s_\nu\,T_{\mu\nu}[\Psi;m,\lambda](-y),\qquad J^\mu[\Psi'](y)=s_\mu\,J^\mu[\Psi](-y).
\end{aligned}
\]

So $\Psi$ solves the field equations with $(m,\lambda)$ on the side $y<0$ exactly when $\Psi'$ solves them with $(-m,\lambda)$, the same $\lambda$, on the side $y>0$. The mirror universe has the same energy density $\rho=T_{44}$, the same pressures and the same charge density $J^4$ at mirror points; the scalar density and the flux along $y$ change sign.

\emph{Proof.} By Lemma 15.6, inserting $1=\gamma^0\gamma^0$ in front of each curved gamma,

\[
\begin{aligned}
\gamma^\mu D_\mu\Psi'(y)&=\sum_\mu\gamma^{x_\mu}(y)\,s_\mu\gamma^0(D_\mu\Psi)(-y)=\sum_\mu s_\mu\,\gamma^0\bigl(\gamma^0\gamma^{x_\mu}(y)\gamma^0\bigr)(D_\mu\Psi)(-y)\\
&=-\gamma^0\sum_\mu s_\mu^2\,\gamma^{x_\mu}(-y)(D_\mu\Psi)(-y)=-\gamma^0\,(\gamma^\mu D_\mu\Psi)(-y),
\end{aligned}
\]

because $s_\mu^2=1$. The adjoint is $\bar\Psi'(y)=\Psi(-y)^\dagger\gamma^0C=-\bar\Psi(-y)\gamma^0$, because $\gamma^0$ is real symmetric and anticommutes with $C$ (Section 2.9, (C1)). Hence $S[\Psi'](y)=-\bar\Psi\gamma^0\gamma^0\Psi=-S(-y)$, and for any bilinear with one curved gamma and one derivative,

\[
\bar\Psi'\gamma_\mu D_\nu\Psi'(y)=\bigl(-\bar\Psi\gamma^0\bigr)\gamma_\mu\bigl(s_\nu\gamma^0D_\nu\Psi\bigr)(-y)=-s_\nu\,\bar\Psi\bigl(\gamma^0\gamma_\mu\gamma^0\bigr)D_\nu\Psi(-y)=s_\mu s_\nu\,\bigl(\bar\Psi\gamma_\mu D_\nu\Psi\bigr)(-y).
\]

The same computation holds for the terms with $D_\mu\bar\Psi'$ (conjugate Lemma 15.6). The kinetic term consists of the terms with $\mu=\nu$ contracted with $\gamma^{x_\mu}$, where $s_\mu^2=1$, so $K[\Psi'](y)=K[\Psi](-y)$; the current has one gamma and no derivative, so $J^\mu[\Psi'](y)=s_\mu J^\mu[\Psi](-y)$. The field equation: $E_{-m,\lambda}[\Psi'](y)=-\gamma^0(\gamma^\mu D_\mu\Psi)(-y)-\bigl(-m-\lambda S(-y)\bigr)\gamma^0\Psi(-y)=-\gamma^0E_{m,\lambda}[\Psi](-y)$. The Lagrangian: $\sqrt{|g|}=W^6$ is even (the factors $e^{\pm a_{4,0}}$ cancel), and $K-(-m)(-S)-\tfrac\lambda2S^2=K-mS-\tfrac\lambda2S^2$. The energy--momentum tensor: the bracket terms were just computed, $g_{\mu\nu}$ is diagonal and even, and $\mathcal L_s$ maps as the Lagrangian. $\square$

(Checks \texttt{PAIR\_T2z\allowbreak{}2\_geomet\allowbreak{}ry}, \texttt{PAIR\_T2z\allowbreak{}2\_diracO\allowbreak{}perator}, \texttt{PAIR\_T2z\allowbreak{}2\_scalar\allowbreak{}Odd}, \texttt{PAIR\_T2z\allowbreak{}2\_lagran\allowbreak{}gianEven}, \texttt{PAIR\_T2z\allowbreak{}2\_emtPul\allowbreak{}lback} for all 64 components, \texttt{PAIR\_T2z\allowbreak{}2\_curren\allowbreak{}tPullbac\allowbreak{}k}; the Python checker adds \texttt{S5\_T2z2\_\allowbreak{}kineticE\allowbreak{}ven} and \texttt{S5\_T2z2\_\allowbreak{}fieldEqu\allowbreak{}ationMap\allowbreak{}sToMinus\allowbreak{}MSameLam\allowbreak{}bda}, and \texttt{PAIR\_T2z\allowbreak{}2\_sameMa\allowbreak{}ssFails} confirms that the image does \textbf{not} solve the equations with the same mass.)

\textbf{Corollary 15.8 (the mirror universe carries the mass $-M$).} Let the mass be a function $m(y)$. By the proof above, the image $\Psi'$ solves the equations with the mass function $y\mapsto-m(-y)$. A configuration that is \textbf{Z2-symmetric}, $\Psi(y)=\pm\gamma^0\Psi(-y)$, solves the field equations on both sides with one mass function $m(y)$ exactly when $\bigl(m(y)+m(-y)\bigr)\Psi(y)=0$, that is, where $\Psi\ne0$, when the mass function is \textbf{odd}: $m(-y)=-m(y)$. With $m=+M$ on the notebook's side, the mirror side carries $-M$. With the Hartree term the effective mass $m(y)+\lambda S(y)$ of a Z2-symmetric configuration is odd automatically, because $S$ is odd by Theorem 15.7. (Checks \texttt{PAIR\_T2z\allowbreak{}2\_massFu\allowbreak{}nctionMa\allowbreak{}p}, \texttt{PAIR\_T2z\allowbreak{}2\_symmet\allowbreak{}ricIffOd\allowbreak{}dMass}; this is erratum E4.6 of \texttt{handoff/\allowbreak{}specs/ST\allowbreak{}AGE4\_SPE\allowbreak{}C.md}, which Section 13.5 states.)

\emph{Proof.} The proof of Theorem 15.7 never used that $m$ is constant; with a mass function it says that $\Psi'(y)=\gamma^0\Psi(-y)$ solves the equations with the mass function $-m(-y)$ whenever $\Psi$ solves them with $m(y)$. Now let $\Psi$ solve the equations with $m(y)$ on both sides and be Z2-symmetric, so that $\Psi'=\pm\Psi$. Then $\Psi$ solves the equations with $m(y)$ and with $-m(-y)$; subtracting the two equations leaves $\bigl(m(y)+m(-y)\bigr)\Psi(y)=0$. Conversely, if $m$ is odd, take a solution on $y<0$ and define it on $y>0$ by $\Psi(y):=\pm\gamma^0\Psi(-y)$; by the same statement it solves the equations with $-m(-y)=m(y)$ there. $\square$

\textbf{Every Stage-4 Kohn--Sham state is already a mirror pair.} The Kohn--Sham problem of Chapter 13 imposes the parity conditions $\Psi(-y)=\pm\gamma^0\Psi(y)$ at the brane, that is, it is solved on $y\le0$ and continued to $y>0$ by the Z2 identification. By Corollary 15.8, every such state describes a universe of mass $+M$ on the side $y<0$ and a universe of mass $-M$ on the side $y>0$, with the same $\lambda$, the same energy density and pressures, the same charge density and the opposite scalar density, joined by the brane, which carries the Israel stress of Section 9.16 (\texttt{pairing-\allowbreak{}theory.j\allowbreak{}son}, key \texttt{T3.stage\allowbreak{}4Z2Pair}). This is a geometric realization of the notebook's "pair of universes of masses $\pm M$". It is a \textbf{model}: the Z2 gluing is a choice (Section 9.16), and nothing here says that the pair was created.

\textbf{Worked example: the brane zero mode.} At zero momentum, with a constant mass $M>0$ and no interaction, the block equation of Section 13.5 has on the side $y<0$ the zero mode $\chi=(e^{My},0)$ of parity $+$ (Section 13.6). In the block basis $\gamma^0$ acts as $\sigma_z$, so the Z2 continuation of parity $+$ is $\chi(y)=\sigma_z\chi(-y)=(e^{-My},0)$ for $y>0$. There the first block equation $\chi_1'=M_{\mathrm{eff}}\chi_1$ reads $-Me^{-My}=M_{\mathrm{eff}}e^{-My}$, so $M_{\mathrm{eff}}=-M$: the mirror side has the mass $-M$, and the mode $e^{-M|y|}$ is localized at the brane from both sides (Exercise 15.7).

\textbf{A remark on the even-mass symmetry.} For an even mass function, Stage 4 identified the reflection $P_B:\Psi'(y)=i\gamma^0\gamma^8\Psi(-y)$ as the symmetry (erratum E4.6, Section 13.5), without using it: the Kohn--Sham problem of Chapter 13 imposes the two parities as boundary conditions and assumes no reflection symmetry. At the level of the classical field it is, up to the factor $i$, the chirality map of T1 followed by the reflection of Theorem 15.7, so it maps $(m,\lambda)$ first to $(-m,-\lambda)$ and then to $(m,-\lambda)$, with $\mathcal L\to-\mathcal L$ (check \texttt{PAIR\_T2z\allowbreak{}2\_PBfiel\allowbreak{}dLevelFl\allowbreak{}ipsLambd\allowbreak{}a}). In the Kohn--Sham model it is a symmetry with the same $\lambda$, because the matrix $BC$ of the scalar density in the expectation-value rule is $P_B$-even while $C$ is $P_B$-odd (checks \texttt{PAIR\_T2z\allowbreak{}2\_rulePa\allowbreak{}rities}, \texttt{S5\_T2z2\_\allowbreak{}PBRuleMa\allowbreak{}trixEven\allowbreak{}COdd}; Exercise 15.16). This is a first instance of the effect explained in Section 15.8.

\subsection{Theorem T3: the Kohn--Sham level}

We now turn to the Kohn--Sham approximation of Chapter 13: a finite number $N$ of quanta in the static primordial field, described by orbitals in eight $2\times2$ blocks. The question is whether the Kohn--Sham ground state and first excited state of a universe of mass $+M$ have exact partners in a universe of mass $-M$.

\textbf{The Kohn--Sham problem, recalled.} (Chapter 13; for the statistics sign, Chapter 14.) The 16 components are split by the commuting labels $(j,s_2,s_3)$ of Section 13.4 into eight blocks of two components $\chi=(\chi_1,\chi_2)$. At a 3-momentum $\mathbf k=(k,0,0)$ an orbital of block type $j$ solves

\[
h_j(M_{\mathrm{eff}},v)\,\chi=\varepsilon\,\chi,\qquad h_j(M,v):=j\Bigl[-i\sigma_x\frac{d}{dy}+M(y)\,\sigma_y+\xi(y)\,k\,\sigma_z\Bigr]+v(y)
\]

on $-L\le y\le0$, with the momentum weight $\xi(y)=e^{-Hy-a_{4,0}}$ and the Pauli matrices $\sigma_x,\sigma_y,\sigma_z$, in the notation of Chapter 13 (the check names and \texttt{pairing-\allowbreak{}theory.j\allowbreak{}son} write the Pauli matrices as sigma1, sigma2, sigma3, and the theory files call the weight \texttt{\detokenize{kappa}}). The boundary conditions (Section 13.5) are a \textbf{parity} at the brane, $p=+$: $\chi_2(0)=0$ or $p=-$: $\chi_1(0)=0$, both sectors belonging to every state, and a \textbf{bag condition} at the tip, $(1-Q(\theta))\chi(-L)=0$ with $Q(\theta)=\cos\theta\,\sigma_z+\sin\theta\,\sigma_y$; Stage 4 uses $\theta=0$, that is $\chi_2(-L)=0$. The expectation-value rule of Section 8.12 gives, per orbital, the proper densities of Section 13.7,

\[
n=\frac{e^{-6Hy}}{\ell^3}\chi^\dagger\chi,\qquad s=\frac{e^{-6Hy}}{\ell^3}\,j\,\chi^\dagger\sigma_y\chi,\qquad t=\frac{e^{-6Hy}}{\ell^3}\,j\,\chi^\dagger\sigma_z\chi,\qquad c=\frac{e^{-6Hy}}{\ell^3}\,j\,\chi^\dagger\sigma_x\chi ,
\]

the number (charge) density, the scalar density, the density of the 3-momentum current and the current along $y$. Summed over the orbitals with the weights $w=f$ (particle levels) or $w=-(1-f)$ (sea levels), $f$ the Fermi--Dirac occupation at the temperature $T$ with the chemical potential $\mu$ fixed by $N$, they give $n_p(y)$ and $S_p(y)$. The Kohn--Sham potentials are

\[
M_{\mathrm{eff}}=m+\Bigl(1+\frac{\mathrm{sg}}{16}\Bigr)\lambda\,S_p,\qquad v=\mathrm{sg}\,\frac{\lambda}{16}\,n_p,
\]

where the \textbf{statistics sign} is $\mathrm{sg}=-1$ for dirac16complex (Section 13.8: $M_{\mathrm{eff}}=m+\tfrac{15}{16}\lambda S_p$, $v=-\tfrac{\lambda}{16}n_p$) and $\mathrm{sg}=+1$ for dirac16complex00, whose exchange term has the opposite sign (Chapter 14: $M_{\mathrm{eff}}=m+\tfrac{17}{16}\lambda S_p$, $v=+\tfrac{\lambda}{16}n_p$; checks \texttt{PAIR\_sta\allowbreak{}t\_unifor\allowbreak{}mGasExch\allowbreak{}ange}, \texttt{PAIR\_sta\allowbreak{}t\_ldaPot\allowbreak{}entials}). Chapter 6 writes this sign as $s$; here the letter $s$ is taken by the scalar density. The energy is $E=\sum w\,\varepsilon-E_H-E_x$ with $E_H+E_x=\lambda\int\bigl[\tfrac12S_p^2+\tfrac{\mathrm{sg}}{32}(n_p^2+S_p^2)\bigr]dV$ (Section 13.9), and the free energy is $F=E-TS_{\mathrm{ent}}$.

\textbf{Lemma 15.9 (the chirality in the reduced equation and in the blocks).}

\begin{enumerate}
\item In the 16-component reduced equation of Section 13.3, $\chi'=N(M)\chi$ with $N(M)=M\gamma^0-i\xi k\,\gamma^0\gamma^1+i(\varepsilon-v)\gamma^0\gamma^4$, one has $\gamma^8N(M)\gamma^8=N(-M)$.
\item $\gamma^8$ maps the block $(j,s_2,s_3)$ onto the block $(-j,s_2,s_3)$, and in the block bases it acts by $\tau\,\sigma_y$ with a number $\tau$ of modulus 1.
\end{enumerate}

\emph{Proof.} (1) $\gamma^0$ is odd, $\gamma^0\gamma^1$ and $\gamma^0\gamma^4$ are even (Lemma 15.2), so only the mass term changes sign. (2) Use the three commuting labels of Section 13.4, the matrices $J=\gamma^0\gamma^1\gamma^4$, $K_1=\gamma^2\gamma^3$ and $K_2=\gamma^5\gamma^6$ (the names are those of Section 13.4; here $J$ is not the current $J^\mu$, and $K_1$, $K_2$ have nothing to do with the kinetic term $K$). $J$ is odd, and $K_1$, $K_2$ are even. If a column $\phi$ has $J\phi=j\phi$, $K_1\phi=is_2\phi$ and $K_2\phi=is_3\phi$, then $J\gamma^8\phi=-\gamma^8J\phi=-j\,\gamma^8\phi$ while $K_{1,2}\gamma^8\phi=\gamma^8K_{1,2}\phi$: the image lies in the block $(-j,s_2,s_3)$. Write $\gamma^8[v_+\ v_-]=[w_+\ w_-]\,G$ with the bases $(v_+,v_-)$ of the first and $(w_+,w_-)$ of the second block and a $2\times2$ matrix $G$. In both blocks $A_0=\gamma^0$ acts as $\sigma_z$ and $A_1=\gamma^0\gamma^1$ as $-i\sigma_y$, while $A_4=\gamma^0\gamma^4$ acts as $j\sigma_x$ in the first and as $-j\sigma_x$ in the second block (Section 13.4). From $\gamma^8A_0=-A_0\gamma^8$ we get $G\sigma_z=-\sigma_zG$; from $\gamma^8A_4=A_4\gamma^8$ we get $G(j\sigma_x)=(-j\sigma_x)G$, that is $G\sigma_x=-\sigma_xG$. Write $G=g_0+g_1\sigma_x+g_2\sigma_y+g_3\sigma_z$ (every $2\times2$ matrix has this form). Anticommuting with $\sigma_z$ forces $g_0=g_3=0$, anticommuting with $\sigma_x$ forces $g_0=g_1=0$; so $G=\tau\sigma_y$ with $\tau:=g_2$. The bases are orthonormal and $\gamma^8$ is unitary (Hermitian with square 1), so $G$ is unitary and $|\tau|=1$. $\square$

The exact computation finds $\tau=\pm1$ for all eight blocks in the block basis of the Stage-4 theory file \texttt{kohn-sha\allowbreak{}m-theory\allowbreak{}.json} (check \texttt{PAIR\_T3b\allowbreak{}lock\_gam\allowbreak{}ma8IsSig\allowbreak{}ma2Betwe\allowbreak{}enPartne\allowbreak{}rBlocks}).

\textbf{Worked example.} Section 13.4 lists the basis of the block $(1,1,1)$; the same construction gives the basis of the block $(-1,1,1)$ (\texttt{kohn-sha\allowbreak{}m-theory\allowbreak{}.json}, block index 4). Multiplied by $\sqrt8$, the four vectors are

\[
\begin{aligned}
\sqrt8\,v_+&=(0,0,0,0,\,1,i,-1,i,\,0,0,0,0,\,1,i,-1,i),\\
\sqrt8\,v_-&=(i,-1,-i,-1,\,0,0,0,0,\,-i,1,i,1,\,0,0,0,0),\\
\sqrt8\,w_+&=(1,i,-1,i,\,0,0,0,0,\,1,i,-1,i,\,0,0,0,0),\\
\sqrt8\,w_-&=(0,0,0,0,\,i,-1,-i,-1,\,0,0,0,0,\,-i,1,i,1).
\end{aligned}
\]

Since $\gamma^8$ changes the sign of the components 0 to 7,

\[
\begin{aligned}
\sqrt8\,\gamma^8v_+&=(0,0,0,0,\,-1,-i,1,-i,\,0,0,0,0,\,1,i,-1,i)=\sqrt8\,i\,w_-,\\
\sqrt8\,\gamma^8v_-&=(-i,1,i,1,\,0,0,0,0,\,-i,1,i,1,\,0,0,0,0)=-\sqrt8\,i\,w_+,
\end{aligned}
\]

which is $\gamma^8[v_+\ v_-]=[w_+\ w_-]\,\sigma_y$, that is $\tau=1$ (Exercise 15.10).

\textbf{Lemma 15.10 (the block map).} For every block type $j$, mass function $M(y)$ and potential $v(y)$,

\[
\sigma_y\,h_j(M,v)\,\sigma_y=h_{-j}(-M,v),
\]

as differential operators. Hence if $\chi$ is an orbital of $h_j(M,v)$ with the level $\varepsilon$, then $\sigma_y\chi$ is an orbital of $h_{-j}(-M,v)$ with the \textbf{same} level $\varepsilon$ and the same 3-momentum.

\emph{Proof.} The Pauli matrices anticommute pairwise and square to 1, so $\sigma_y\sigma_x\sigma_y=-\sigma_x$, $\sigma_y\sigma_y\sigma_y=\sigma_y$, $\sigma_y\sigma_z\sigma_y=-\sigma_z$, and the constant $\sigma_y$ commutes with $d/dy$. Therefore $\sigma_yh_j(M,v)\sigma_y=j\bigl[+i\sigma_x\tfrac{d}{dy}+M\sigma_y-\xi k\sigma_z\bigr]+v=(-j)\bigl[-i\sigma_x\tfrac{d}{dy}+(-M)\sigma_y+\xi k\sigma_z\bigr]+v$. If $h_j\chi=\varepsilon\chi$, then $h_{-j}(-M,v)(\sigma_y\chi)=\sigma_yh_j(M,v)\sigma_y\sigma_y\chi=\sigma_yh_j\chi=\varepsilon\,\sigma_y\chi$. $\square$

(Checks \texttt{PAIR\_T3b\allowbreak{}lock\_sig\allowbreak{}ma2MapsB\allowbreak{}lockODE}, \texttt{PAIR\_T3b\allowbreak{}lock\_ham\allowbreak{}iltonian\allowbreak{}Sigma2}.) In the real form $\chi=(a,\,ib)$ of Section 13.5, $\sigma_y(a,\,ib)=(b,\,ia)$: the map is the swap $(a,b,j)\to(b,a,-j)$ that Section 13.14 found directly from the real system (check \texttt{PAIR\_T3b\allowbreak{}lock\_sig\allowbreak{}ma2IsRus\allowbreak{}tSwap}).

\textbf{Lemma 15.11 (the boundary conditions).} Under $\chi\to\sigma_y\chi$: the parity at the brane changes, $p\to-p$; the bag angle at the tip changes as $\theta\to\pi-\theta$; in particular the Stage-4 condition $\theta=0$ ($\chi_2(-L)=0$) becomes $\theta=\pi$ ($\chi_1(-L)=0$).

\emph{Proof.} $\sigma_y\chi=(-i\chi_2,\,i\chi_1)$, so $\chi_2(0)=0$ becomes $(\sigma_y\chi)_1(0)=0$ and conversely. For the bag, $\sigma_yQ(\theta)\sigma_y=\cos\theta\,\sigma_y\sigma_z\sigma_y+\sin\theta\,\sigma_y=-\cos\theta\,\sigma_z+\sin\theta\,\sigma_y=Q(\pi-\theta)$, since $\cos(\pi-\theta)=-\cos\theta$ and $\sin(\pi-\theta)=\sin\theta$. So $(1-Q(\theta))\chi=0$ is equivalent to $(1-Q(\pi-\theta))\sigma_y\chi=\sigma_y(1-Q(\theta))\chi=0$. For $\theta=0$: $Q(\pi)=-\sigma_z$ and $(1+\sigma_z)\sigma_y\chi(-L)=0$ says $(\sigma_y\chi)_1(-L)=0$. $\square$

(Checks \texttt{PAIR\_T3b\allowbreak{}lock\_par\allowbreak{}ityMap}, \texttt{PAIR\_T3b\allowbreak{}lock\_bag\allowbreak{}AngleMap}.) Since every Stage-4 Kohn--Sham state contains both parity sectors, the exchange $p\to-p$ only relabels the sectors. The one boundary condition that really changes is the tip bag.

\textbf{Lemma 15.12 (the densities).} Under $\chi\to\sigma_y\chi$ together with $j\to-j$, with the expectation-value rule,

\[
(n,\,s,\,t,\,c)\ \longrightarrow\ (n,\,-s,\,t,\,c).
\]

\emph{Proof.} $n\propto(\sigma_y\chi)^\dagger(\sigma_y\chi)=\chi^\dagger\chi$. $s\propto(-j)\,\chi^\dagger\sigma_y\sigma_y\sigma_y\chi=-j\,\chi^\dagger\sigma_y\chi$. $t\propto(-j)\,\chi^\dagger\sigma_y\sigma_z\sigma_y\chi=(-j)(-\chi^\dagger\sigma_z\chi)=j\,\chi^\dagger\sigma_z\chi$. $c\propto(-j)(-\chi^\dagger\sigma_x\chi)=j\,\chi^\dagger\sigma_x\chi$. The phase $\tau$ of Lemma 15.9 drops out, since $|\tau|^2=1$. $\square$

(Check \texttt{PAIR\_T3b\allowbreak{}lock\_den\allowbreak{}sityAndC\allowbreak{}urrentMa\allowbreak{}ps}; measurement \texttt{T3block\_\allowbreak{}densityS\allowbreak{}igns\_sig\allowbreak{}ma2} $=\{1,-1,1,1\}$.)

\textbf{Theorem 15.13 (T3, the Kohn--Sham level).} Let the \textbf{universe of mass $+M$} be the Kohn--Sham problem of Chapter 13 with the bare mass $m$, the coupling $\lambda$, the statistics sign $\mathrm{sg}$, both parity sectors, the tip bag angle $\theta$, $N$ quanta at the temperature $T$ (Fermi--Dirac occupations, normal ordering against the free Dirac sea, the particle and sea labels fixed by continuity from $\lambda=0$ with the zero modes counted as particles, Section 13.7). Let the \textbf{universe of mass $-M$} be the same problem with the bare mass $-m$, the \textbf{same} $\lambda$ and $\mathrm{sg}$, both parity sectors and the tip bag angle $\pi-\theta$. Then the map $\chi\to\sigma_y\chi$ in the partner block $(j,s_2,s_3)\to(-j,s_2,s_3)$, at the same momentum, sends the one problem exactly onto the other:

\begin{enumerate}
\item the self-consistent potentials correspond as $M_{\mathrm{eff}}\to-M_{\mathrm{eff}}$ and $v\to v$, and every self-consistent solution of the $+M$ problem is mapped onto a self-consistent solution of the $-M$ problem, and conversely;
\item the Kohn--Sham levels with their multiplicities, the occupations, the particle and sea labels, the weights, $\mu$, $E$, $F$ and $S_{\mathrm{ent}}$ are identical;
\item the profiles $n_p(y)$, $v(y)$, $t$, $c$ and the energy--momentum profiles $\rho(y)$, $p_y(y)$, $p_3(y)$, $p_t(y)$ are identical, while $S_p(y)\to-S_p(y)$ and $M_{\mathrm{eff}}(y)\to-M_{\mathrm{eff}}(y)$;
\item the ground state (defined by continuation in $\lambda$ from $\lambda=0$, erratum E4.8) and the first excited state (the Kohn--Sham gap, the particle--hole list, the Delta-SCF state) of the $+M$ universe are mapped exactly onto those of the $-M$ universe.
\end{enumerate}

With $-\lambda$ in place of $\lambda$ the map fails for $\lambda\ne0$.

\emph{Proof.} Step 1 (potentials). Suppose the $+M$ problem has the densities $(n_p,S_p)$. By Lemma 15.12 the mapped orbitals have the densities $(n_p,-S_p)$. With the bare mass $-m$ and the same $\lambda$ the potentials built from them are

\[
-m+\Bigl(1+\frac{\mathrm{sg}}{16}\Bigr)\lambda\,(-S_p)=-M_{\mathrm{eff}},\qquad \mathrm{sg}\,\frac{\lambda}{16}\,n_p=v ,
\]

exactly the potentials $(-M_{\mathrm{eff}},v)$ in which Lemma 15.10 places the mapped orbitals. With $-\lambda$ instead, the mass would be $-m+(1+\tfrac{\mathrm{sg}}{16})\lambda S_p\ne-M_{\mathrm{eff}}$ and the potential $-v\ne v$ whenever $\lambda S_p\ne0$ or $\lambda n_p\ne0$.

Step 2 (spectrum and boundary conditions). By Lemmas 15.10 and 15.11, $\chi$ is an eigenfunction of $h_j(M_{\mathrm{eff}},v)$ with the parity $p$ and the bag $\theta$ exactly when $\sigma_y\chi$ is an eigenfunction of $h_{-j}(-M_{\mathrm{eff}},v)$ with the parity $-p$ and the bag $\pi-\theta$, with the same $\varepsilon$. Since $\sigma_y^2=1$ the correspondence is one to one, so the two spectra coincide with multiplicities, level by level.

Step 3 (occupations and labels). The occupations depend only on $\varepsilon$, $\mu$ and $T$. At $\lambda=0$ the free problems are mapped onto each other by Step 2 (with $M_{\mathrm{eff}}=\pm m$, $v=0$), so the free partner of a mapped level is the mapped free partner, with the same $\varepsilon_{\mathrm{free}}$; hence the particle and sea labels, and the weights $f$ or $-(1-f)$, are the same. The zero modes $(e^{my},0)$ of the $+M$ problem go to the zero modes $\sigma_y(e^{my},0)=(0,\,ie^{my})$ of the $-M$ problem, again particle levels. The particle number $N=\sum w$ is the same function of $\mu$, so $\mu$ is the same.

Step 4 (self-consistency). Steps 1 to 3 say that the iteration map of the self-consistent loop (densities $\to$ potentials $\to$ orbitals and occupations $\to$ densities) of the $-M$ problem is the image of that of the $+M$ problem. Every iterate, and every fixed point, is mapped. The continuation path $\lambda/4,\lambda/2,3\lambda/4,\lambda$ that defines the ground state is the same path in both problems, because $\lambda$ is not changed, and so is every fallback with occupation smearing.

Step 5 (energies). $\sum w\varepsilon$ is the same. $E_H+E_x$ contains $S_p$ only through $S_p^2$ and $n_p$ through $n_p^2$, so it is the same. Hence $E$ is the same, and with the same occupations $S_{\mathrm{ent}}$ and $F$ are the same.

Step 6 (first excited state). The levels, the occupations and the thresholds for holes and particles are the same, so the Kohn--Sham gap and the particle--hole list are the same. The Delta-SCF state fixes the occupations by the identity of the levels (shell, parity, block type, Prüfer index); the mapped identities define the mapped constraint, and by Step 4 its self-consistent solution is the image, with the same energy $E_1$.

Step 7 (energy--momentum tensor). By Section 13.14 (with $L_s=\tfrac\lambda2S_p^2+\tfrac{\mathrm{sg}\lambda}{32}(n_p^2+S_p^2)$ for either statistics), $\rho=\sum w\,\varepsilon\,n-L_s$, $p_y=\sum w\,[\varepsilon n-m\,s-\xi k\,t]-L_s$, $p_3=\tfrac13\sum w\,\xi k\,t+L_s$ and $p_t=L_s$. Under the map $\varepsilon$, $w$, $n$, $t$ and $L_s$ are unchanged, and $m\,s\to(-m)(-s)=m\,s$. So all four profiles are unchanged. $\square$

The statement is the key \texttt{T3.theor\allowbreak{}emStanda\allowbreak{}rdRule} of \texttt{pairing-\allowbreak{}theory.j\allowbreak{}son}. Its checks, in both pairing reports (the negative control of the operator check confirms that $-\lambda$ does not map):

\begin{Verbatim}[fontsize=\small]
PAIR_T3block_potentialsStandardRule   PAIR_T3ks_stationarityBothStatistics
PAIR_T3ks_sigma2FunctionalInvariant   PAIR_T3ks_sigma2Densities
PAIR_T3ks_sigma2KSOperatorEquivariant PAIR_T3ks_zeroModeSplittingMapsExactly
PAIR_T3ks_massiveLevelsMap            PAIR_T3emt_standardRulePlusT
PAIR_T3emt_stage4CrossCheck
\end{Verbatim}

\textbf{The hypotheses, once more.} T3 is a theorem about the Kohn--Sham model, not about the exact many-body theory: it assumes the model of Chapter 13 (mean field with the local exchange of the uniform gas, no correlation term, a fixed metric, the good sector, the torus of momenta, the expectation-value rule and normal ordering against the free sea), and it needs the \textbf{transformed} tip condition. For dirac16complex00 the model is the one of Chapter 14: a DFT-motivated mean field in which the expectation-value rule and the Pauli filling are imposed by prescription. With the rule, the covariance of the occupied modes is indefinite (on a filled rest shell it has the eigenvalues $+1$ four times, $-1$ four times and $0$ eight times, check \texttt{PAIR\_sta\allowbreak{}t\_expect\allowbreak{}ationRul\allowbreak{}eCovaria\allowbreak{}nceIndef\allowbreak{}inite}), so for the commuting field the model is a formal (Krein-signed) Gaussian functional rather than the statistics of a classical random field (\texttt{pairing-\allowbreak{}theory.j\allowbreak{}son}, key \texttt{statisti\allowbreak{}cs.posit\allowbreak{}ivity}). The proof of T3 keeps $\mathrm{sg}$ as a symbol and holds for both values.

\textbf{Two remarks on notation and on a second route.} In the Rust solver the block label is $s=-j$ and the orbital is written $\chi=(a,\,ib)$; the map is then $(\text{shell},p,s,n)\to(\text{shell},-p,-s,-n)$ for (shell, parity, block type, Prüfer index), with the same $\varepsilon$ and the orbital $(a,b)\to\pm(b,a)$ (\texttt{\detokenize{levelMap}} in every \texttt{pairing.\allowbreak{}json}); the Stage-4 tip condition $b(-L)=0$ becomes $a(-L)=0$. The Python reference solver writes the block orbital as $(f,g)$, with parity $+$ meaning $g(0)=0$; in its notation the tip condition $g(-L)=0$ (called \texttt{\detokenize{g0}}) of the $+M$ universe becomes $f(-L)=0$ (called \texttt{\detokenize{f0}}), and the map is the swap $f\leftrightarrow g$ together with the block type (\texttt{pairing-\allowbreak{}theory.j\allowbreak{}son}, key \texttt{T3.numer\allowbreak{}icsPresc\allowbreak{}ription.\allowbreak{}componen\allowbreak{}tSwap}). The reflection $\gamma^1$ of the 3-space direction $x_1$ gives a second exact route: it acts as $\sigma_x$ in every block, keeps $j$, reverses $k$ and changes $\theta\to\theta+\pi$, with the same conclusions over closed shells $\{\mathbf k,-\mathbf k\}$ (checks \texttt{PAIR\_T3b\allowbreak{}lock\_gam\allowbreak{}ma1IsSig\allowbreak{}ma1InEve\allowbreak{}ryBlock}, \texttt{PAIR\_T3k\allowbreak{}s\_sigma1\allowbreak{}Function\allowbreak{}alInvari\allowbreak{}ant}).

\textbf{The pair totals at the Kohn--Sham level.} For the ordinary universe of mass $-M$ (the standard rule, the same $\lambda$) the pair of Theorem 15.13 is a \textbf{mirror pair}: $E^{\mathrm{pair}}=2E_+$, $F^{\mathrm{pair}}=2F_+$, charge $2N$, $S^{\mathrm{pair}}_p(y)=0$ and every energy--momentum profile doubled (check \texttt{PAIR\_tot\allowbreak{}als\_ksMi\allowbreak{}rrorPair}). It is the Kohn--Sham form of T2.

\subsection{Why T1 and T3 give different partners}

T1 (Section 15.3) says: the image of a universe of mass $+M$ has the coupling $-\lambda$ and the opposite energy. T3 (Section 15.7) says: the partner has the same $\lambda$ and the same energy. Both are built on $\gamma^8$. There is no contradiction, because the two theorems are about different quantities.

\textbf{Classical bilinears versus the expectation-value rule.} T1 is about the classical bilinears $\Psi^\dagger X\Psi$ of the field. The Kohn--Sham model uses instead the expectation values of the positive-norm quantum theory of Chapter 8. For one quantum in a normalized mode $u$ above the Dirac sea the rule of Section 8.12 is

\[
\bigl\langle\Psi^\dagger X\Psi\bigr\rangle=u^\dagger B\,X\,u .
\]

So every classical matrix $X$ is replaced by $BX$. By Lemma 15.2, $B$ is odd, and \textbf{multiplying by $B$ reverses the parity} under $\gamma^8$. The energy and the current, whose classical matrices are odd, become even; the scalar density, whose classical matrix $C$ is even, becomes odd. Under $u\to\gamma^8u$:

\begingroup
\small
\begin{longtable}{@{}>{\raggedright\arraybackslash}p{0.160\linewidth}>{\raggedright\arraybackslash}p{0.160\linewidth}>{\raggedright\arraybackslash}p{0.160\linewidth}>{\raggedright\arraybackslash}p{0.160\linewidth}>{\raggedright\arraybackslash}p{0.160\linewidth}@{}}
\toprule
\textbf{quantity} & \textbf{classical matrix $X$} & \textbf{classical bilinear under $\gamma^8$} & \textbf{rule matrix $BX$} & \textbf{expectation value under $\gamma^8$} \\
\midrule
\endfirsthead
\toprule
\textbf{quantity} & \textbf{classical matrix $X$} & \textbf{classical bilinear under $\gamma^8$} & \textbf{rule matrix $BX$} & \textbf{expectation value under $\gamma^8$} \\
\midrule
\endhead
charge density & $B$ & changes sign & $B^2=1$ & unchanged \\
scalar density & $C$ & unchanged & $BC=-i\gamma^4$ & changes sign \\
energy, momentum, stresses (the bracket of $T_{\mu\nu}$) & $C\gamma^a$ and $C\gamma^aS^{bc}$ & change sign & $BC\gamma^a$, $BC\gamma^aS^{bc}$ & unchanged \\
\bottomrule
\end{longtable}
\endgroup

The first two columns are T1: energy and charge change sign, the scalar density does not, and self-consistency of the mass $m+\lambda S$ then requires $\lambda\to-\lambda$. The last two columns are T3: energy and charge stay, the scalar density changes sign, and self-consistency of $m+\lambda S\to-(m+\lambda S)$ requires the same $\lambda$.

\textbf{The same statement for one block.} In the block $(j,s_2,s_3)$ the matrix $B$ acts as the number $js_2$ (Section 13.4), the \textbf{Krein sign} of the block. So for a mode that lives in one block, every classical bilinear is $js_2$ times the corresponding expectation value, and in particular the classical energy--momentum tensor of such a mode is $js_2$ times its Kohn--Sham one-body energy--momentum tensor (checks \texttt{C00\_stat\allowbreak{}ic\_class\allowbreak{}icalMode\allowbreak{}IsKreinW\allowbreak{}eightedK\allowbreak{}S} and \texttt{C00\_stat\allowbreak{}ic\_class\allowbreak{}icalEner\allowbreak{}gyKreinS\allowbreak{}igned} of \texttt{wolfram-\allowbreak{}dirac16c\allowbreak{}omplex00\allowbreak{}-report.\allowbreak{}json}; key \texttt{static.c\allowbreak{}lassical\allowbreak{}VersusKo\allowbreak{}hnSham} of \texttt{artifact\allowbreak{}s/dirac1\allowbreak{}6complex\allowbreak{}/pair-cr\allowbreak{}eation/d\allowbreak{}irac16co\allowbreak{}mplex00-\allowbreak{}theory.j\allowbreak{}son}). The block map $j\to-j$ reverses the Krein sign. The expectation values are unchanged (T3), so the classical bilinears change sign (T1).

\textbf{Worked example (continued).} The rest state $u_+$ of Section 15.3 and its image $\tilde u=\gamma^8u_+$. As a mode of the $-m$ theory, $\tilde u$ is a positive-energy mode: $h(-m)\tilde u=-m(-i\gamma^4)\tilde u=m\,\tilde u$ with $h(m)=-im\gamma^4$ (Section 8.9). Per quantum, with $BC=-i\gamma^4$:

\begingroup
\small
\begin{longtable}{@{}>{\raggedright\arraybackslash}p{0.160\linewidth}>{\raggedright\arraybackslash}p{0.160\linewidth}>{\raggedright\arraybackslash}p{0.160\linewidth}>{\raggedright\arraybackslash}p{0.160\linewidth}>{\raggedright\arraybackslash}p{0.160\linewidth}@{}}
\toprule
\textbf{quantity per quantum} & \textbf{$u_+$, mass $m$} & \textbf{$\tilde u$, classical bilinear} & \textbf{$\tilde u$, rule with $+B$ (T3)} & \textbf{$\tilde u$, rule with $-B$ (image field)} \\
\midrule
\endfirsthead
\toprule
\textbf{quantity per quantum} & \textbf{$u_+$, mass $m$} & \textbf{$\tilde u$, classical bilinear} & \textbf{$\tilde u$, rule with $+B$ (T3)} & \textbf{$\tilde u$, rule with $-B$ (image field)} \\
\midrule
\endhead
charge $\langle J^4\rangle$ & $u_+^\dagger B^2u_+=1$ & $\tilde u^\dagger B\tilde u=-1$ & $\tilde u^\dagger B^2\tilde u=1$ & $-1$ \\
scalar $\langle S\rangle$ & $u_+^\dagger(-i\gamma^4)u_+=1$ & $\tilde u^\dagger C\tilde u=1$ & $\tilde u^\dagger(-i\gamma^4)\tilde u=-1$ & $+1$ \\
energy & $u_+^\dagger h(m)u_+=m$ & $m\,\tilde u^\dagger B\tilde u=-m$ & $\tilde u^\dagger h(-m)\tilde u=m$ & $-m$ \\
\bottomrule
\end{longtable}
\endgroup

For $u_+$ the classical and the rule values coincide, because $u_+$ has the Krein sign $u_+^\dagger Bu_+=+1$. The last column evaluates the formulas of $\mathcal L_{-m,-\lambda}$ with the metric $-B$; by Section 15.4 the image field's own energy and charge, measured with its own canonical structure, are those of the column of $u_+$. (Exercise 15.4 recomputes the table by hand.)

\textbf{The Krein image at the Kohn--Sham level.} Section 15.4 showed that the image field $\gamma^8\Psi$ of the quantum theory carries the reversed metric $-B$, and that it is the same quantum system as $\Psi$, evaluated with the formulas of the partner theory. The Kohn--Sham construction below is of the same formal kind. The expectation-value rule of the image is $\langle\Psi_-^\dagger X\Psi_-\rangle=u^\dagger(-B)Xu$ with $u$ the mapped mode, and every one-body density gets one more sign: $(n,s,t,c)\to(-n,\,s,\,-t,\,-c)$ (check \texttt{PAIR\_T3b\allowbreak{}lock\_pot\allowbreak{}entialsI\allowbreak{}mageRule}). The potentials then match with $(m,\lambda)\to(-m,-\lambda)$: $-m+(1+\tfrac{\mathrm{sg}}{16})(-\lambda)S_p=-M_{\mathrm{eff}}$ and $\mathrm{sg}\tfrac{(-\lambda)}{16}(-n_p)=v$, while $E_H+E_x$ changes sign. The image has the same orbitals $\sigma_y\chi$, the same levels and the same occupations, and all its one-body densities and energies, including $T_{\mu\nu}$, are those of the ordinary $-M$ run with the opposite sign: $E\to-E$ and $T_{\mu\nu}\to-T_{\mu\nu}$ (checks \texttt{PAIR\_T3k\allowbreak{}s\_imageR\allowbreak{}uleEnerg\allowbreak{}yOdd}, \texttt{PAIR\_T3e\allowbreak{}mt\_image\allowbreak{}RuleMinu\allowbreak{}sT}). This is the Kohn--Sham form of T1. The image is not solved as a universe of its own: it is the ordinary $-M$ state with the sign of every one-body density and energy reversed (\texttt{pairing-\allowbreak{}theory.j\allowbreak{}son}, key \texttt{pairTota\allowbreak{}ls}, entry \texttt{kreinIma\allowbreak{}gePair}), that is, by T3, the $+M$ state measured with the formulas of the partner theory. One formal feature deserves mention: the derivatives of the image's energy with respect to the occupations (Janak's theorem, Section 12.15) are $-\varepsilon$, and since $f(\varepsilon;\mu,T)=f(-\varepsilon;-\mu,-T)$ for the Fermi function, the image is a stationary state of its free energy at the formal temperature $-T$ and chemical potential $-\mu$, with $F_-(-T)=-F_+(T)$ (check \texttt{PAIR\_T3k\allowbreak{}s\_occupa\allowbreak{}tionsAnd\allowbreak{}Temperat\allowbreak{}ure}). For the \textbf{Krein image pair}, formed as $(+M)-(-M)$ for every one-body density and energy: $E^{\mathrm{pair}}=0$, charge 0, $T^{\mathrm{pair}}_{\mu\nu}=0$ at every point, $S^{\mathrm{pair}}=2S_+$ (check \texttt{PAIR\_tot\allowbreak{}als\_ksKr\allowbreak{}einImage\allowbreak{}Pair}). This zero is of the $X+(-X)$ kind of Section 15.4: since the $-M$ state has the same energy and energy--momentum tensor as the $+M$ state (T3), subtracting the one from the other gives zero. It restates the equality $E(+M)=E(-M)$ of T3; it is not a cancellation between two independent universes.

\textbf{Summary of the four pairings.}

\begingroup
\small
\begin{longtable}{@{}>{\raggedright\arraybackslash}p{0.133\linewidth}>{\raggedright\arraybackslash}p{0.133\linewidth}>{\raggedright\arraybackslash}p{0.133\linewidth}>{\raggedright\arraybackslash}p{0.133\linewidth}>{\raggedright\arraybackslash}p{0.133\linewidth}>{\raggedright\arraybackslash}p{0.133\linewidth}@{}}
\toprule
\textbf{pairing} & \textbf{map} & \textbf{partner} & \textbf{$T_{\mu\nu}$} & \textbf{$Q$} & \textbf{$S$} \\
\midrule
\endfirsthead
\toprule
\textbf{pairing} & \textbf{map} & \textbf{partner} & \textbf{$T_{\mu\nu}$} & \textbf{$Q$} & \textbf{$S$} \\
\midrule
\endhead
T1, classical fields (Section 15.3) & $\Psi\to\gamma^8\Psi$ & $(-m,-\lambda)$ & $-T_{\mu\nu}$ & $-Q$ & $+S$ \\
T1, Krein image, formal (Sections 15.4, 15.8) & $\chi\to\sigma_y\chi$, metric $-B$ & $(-m,-\lambda)$ & $-T_{\mu\nu}$ in the formulas of the partner & $-Q$ in the formulas of the partner & $+S$ \\
T2, mirror (Sections 15.5, 15.6) & $\Psi\to u\Psi$, reflected frame & $(-m,\lambda)$ & $+T_{\mu\nu}$, reflected & $+Q$ & $-S$ \\
T3, Kohn--Sham (Section 15.7) & $\chi\to\sigma_y\chi$, metric $+B$ & $(-m,\lambda)$, bag $\pi-\theta$ & $+T_{\mu\nu}$ & $+Q$ & $-S$ \\
\bottomrule
\end{longtable}
\endgroup

Only the pairs of the first two rows have zero totals, and both need the coupling $-\lambda$ (or $\lambda=0$). The first row is a statement about classical fields. The second row is a formal construction: the same quantum system, measured with the metric $-B$ and the formulas of the partner theory, whose zero is of the kind $X+(-X)=0$ (Section 15.4); measured with its own canonical structure it has $+T_{\mu\nu}$ and $+Q$. No row describes a cancellation between two independent quantum universes; whether a quantum description of such a pair exists is OPEN.

\subsection{When the boundary conditions are not transformed: the control}

T3 needs the tip bag $\pi-\theta$ in the universe of mass $-M$. What happens if the $-M$ universe keeps the Stage-4 bag $\theta=0$ ($\chi_2(-L)=0$)? Write $P(m,\theta)$ for the Kohn--Sham problem with the bare mass $m$ and the tip bag angle $\theta$. By Theorem 15.13, $P(-m,\theta)$ is the image of $P(+m,\pi-\theta)$; so it is the image of $P(+m,\theta)$ only if the two bags give the same problem, which they do not for $\theta=0$, and the pairing fails. The exact theory makes this concrete at zero momentum, with a constant mass $M>0$ and no interaction (\texttt{pairing-\allowbreak{}theory.j\allowbreak{}son}, key \texttt{T3.untra\allowbreak{}nsformed\allowbreak{}BCContro\allowbreak{}l}). We derive the three differences.

\textbf{1. The zero mode moves to the tip.} With the mass $-M$ at $k=0$ and $\varepsilon=0$ the block equation reads $\chi_1'=-M\chi_1$, $\chi_2'=M\chi_2$. The parity-$+$ conditions $\chi_2(0)=\chi_2(-L)=0$ force $\chi_2\equiv0$, and then $\chi=(e^{-My},0)$: a zero mode that grows toward the tip $y=-L$, where the $+M$ zero mode $(e^{My},0)$ is smallest. The ratio of its density at the brane to that at the tip is $e^{-2ML}$, against $e^{+2ML}$ for $+M$ (check \texttt{PAIR\_T3k\allowbreak{}s\_zeroMo\allowbreak{}deUntran\allowbreak{}sformedC\allowbreak{}ontrol}). With the transformed bag the $-M$ zero mode is $(0,e^{My})$, localized at the brane like the $+M$ one (check \texttt{PAIR\_T3k\allowbreak{}s\_zeroMo\allowbreak{}deImage}).

\textbf{2. The brane band splits differently.} By the Hellmann--Feynman rule of Section 13.6, a zero mode $\chi$ moves at small $k$ as $\varepsilon\approx j\,c\,k$ with $c=\int\xi\,\chi^\dagger\sigma_z\chi\,dy\big/\int\chi^\dagger\chi\,dy$. For the tip zero mode $\chi^\dagger\sigma_z\chi=e^{-2My}$, and

\[
c_{\mathrm{ctrl}}=\frac{\int_{-L}^0e^{-Hy-a_{4,0}}e^{-2My}\,dy}{\int_{-L}^0e^{-2My}\,dy}=e^{-a_{4,0}}\,\frac{2M}{2M+H}\cdot\frac{e^{(2M+H)L}-1}{e^{2ML}-1},
\]

against $c(M)=e^{-a_{4,0}}\frac{2M}{2M-H}\cdot\frac{1-e^{-(2M-H)L}}{1-e^{-2ML}}$ for the paired problems (Section 13.6). For $M=H=1$, $a_{4,0}=0$ and $L=3$: $c=1.905148$ and $c_{\mathrm{ctrl}}=13.421975$. At $M=H$ the difference is $c_{\mathrm{ctrl}}-c=\tfrac23(Y-1)^2/(Y+1)>0$ with $Y=e^{HL}$ (Exercise 15.14; checks \texttt{PAIR\_T3k\allowbreak{}s\_contro\allowbreak{}lSplitti\allowbreak{}ngClosed\allowbreak{}Form}, \texttt{PAIR\_T3k\allowbreak{}s\_contro\allowbreak{}lDiffers\allowbreak{}FromPair\allowbreak{}edProble\allowbreak{}m}). The Rust solver measures the slopes $1.9051480$ (both paired universes) and $13.4219708$ (control) on its free spectra, against the closed forms $1.9051483$ and $13.4219752$ (\texttt{artifact\allowbreak{}s/dirac1\allowbreak{}6complex\allowbreak{}/pair-cr\allowbreak{}eation/r\allowbreak{}ust/pair\allowbreak{}s/free-s\allowbreak{}ection.j\allowbreak{}son}, key \texttt{splittin\allowbreak{}g}, relative differences $1.3\times10^{-7}$ and $3.3\times10^{-7}$).

\textbf{3. A bound state appears below the mass gap.} In the parity-$-$ sector ($\chi_1(0)=0$) with $\theta=0$ ($\chi_2(-L)=0$), eliminate $\chi_1$ as in Section 13.6: from $\chi_2'=-M_{\mathrm{eff}}\chi_2+ij\varepsilon\chi_1$ we get $\chi_1=(\chi_2'+M_{\mathrm{eff}}\chi_2)/(ij\varepsilon)$ and $\chi_2''=(M^2-\varepsilon^2)\chi_2$. The condition $\chi_1(0)=0$ becomes $\chi_2'(0)+M_{\mathrm{eff}}\chi_2(0)=0$. For $\varepsilon^2<M^2$ put $q=\sqrt{M^2-\varepsilon^2}$; the solution with $\chi_2(-L)=0$ is $\chi_2=\sinh(q(y+L))$, and the brane condition reads

\[
q\cosh(qL)+M_{\mathrm{eff}}\sinh(qL)=0 .
\]

For $M_{\mathrm{eff}}=+M$ both terms are positive and there is no solution: the $+M$ universe has no level with $|\varepsilon|<M$ in this sector (Section 13.6). For the control, $M_{\mathrm{eff}}=-M$, the condition is $\tanh(qL)=q/M$. The function $\tanh(qL)-q/M$ vanishes at $q=0$, has the slope $L-1/M$ there and is negative at $q=M$; since $\tanh$ is concave for positive arguments, there is exactly one root $q\in(0,M)$ when $ML>1$, and none when $ML\le1$. The level is

\[
\varepsilon_{\mathrm b}=\pm\sqrt{M^2-q^2}=\pm\frac{M}{\cosh(qL)}\approx\pm2M\,e^{-ML}\quad(ML\gg1),
\]

using $1-\tanh^2=1/\cosh^2$. For $M=1$, $L=3$: $q=0.9949015$ and $\varepsilon_{\mathrm b}=0.1008511214$ (Exercise 15.13); for $M=3$, $L=3$: $\varepsilon_{\mathrm b}=7.404590\times10^{-4}$. For $\varepsilon^2>M^2$ put $r=\sqrt{\varepsilon^2-M^2}$ (a wave number; we avoid the letter $p$, which is the parity here). Then $\chi_2''=-r^2\chi_2$, the solution with $\chi_2(-L)=0$ is $\chi_2=\sin(r(y+L))$, and the brane condition gives the level conditions $r\cos(rL)\pm M\sin(rL)=0$ ($+$ for the paired universes, $M_{\mathrm{eff}}=+M$; $-$ for the control, $M_{\mathrm{eff}}=-M$). They have no common solution: adding and subtracting them would give $r\cos(rL)=0$ and $M\sin(rL)=0$, which cannot both hold for $r>0$ and $M>0$, because $\cos$ and $\sin$ never vanish together. (Checks \texttt{PAIR\_T3k\allowbreak{}s\_mixedS\allowbreak{}ectorSol\allowbreak{}utions}, \texttt{PAIR\_T3k\allowbreak{}s\_contro\allowbreak{}lMixedSe\allowbreak{}ctorLeve\allowbreak{}lsDisjoi\allowbreak{}nt}, \texttt{PAIR\_T3k\allowbreak{}s\_contro\allowbreak{}lSubGapB\allowbreak{}oundStat\allowbreak{}e}.)

\textbf{Consequence.} With the untransformed bag the spectra differ already at $k=0$, so the occupations, energies, densities and the energy--momentum tensor of the $-M$ universe differ from those of the $+M$ universe. The numerical control of Section 15.10 shows exactly this, and item 3 predicts its Kohn--Sham gap for $N=8$ to twelve digits.

\subsection{The numerical demonstration}

\textbf{What was planned.} The Stage-5 specification (§5, T4) asks for the Kohn--Sham ground and first excited states of the $+M$ universe, of the transformed $-M$ universe and of the untransformed control, for both fields, $m\in\{1,3\}$ ($H=1$), $L=3$, $N\in\{8,112\}$, $\hat\lambda\in\{0,\pm\hat\lambda_1,\pm\hat\lambda_2\}$ with the couplings of Section 13.10, at $T=0$ and at $T=0.1\,|m|$, computed by the Rust solver (its \texttt{\detokenize{pairs}} subcommand, module \texttt{src/pair\allowbreak{}s.rs} of \texttt{studies/\allowbreak{}dirac16c\allowbreak{}omplex\_k\allowbreak{}ohn\_sham}, described in the solver's \texttt{README.m\allowbreak{}d}, section "Stage 5") and by the independent Python reference solver (\texttt{scripts/\allowbreak{}ks\_refer\allowbreak{}ence\_pai\allowbreak{}rs.py} with \texttt{scripts/\allowbreak{}ks\_refer\allowbreak{}ence\_sol\allowbreak{}ver.py}), and compared by \texttt{scripts/\allowbreak{}check\_di\allowbreak{}rac16com\allowbreak{}plex\_pai\allowbreak{}rs.py}.

\textbf{What exists} (committed outputs, state of 2026-09-30):

\begin{itemize}
\item \textbf{Rust, dirac16complex only}, in \texttt{artifact\allowbreak{}s/dirac1\allowbreak{}6complex\allowbreak{}/pair-cr\allowbreak{}eation/r\allowbreak{}ust/pair\allowbreak{}s/}, one folder per configuration with a \texttt{pairing.\allowbreak{}json} (the comparison of the three universes) and the run files of each universe: $m=1$, $N=8$ for all five couplings at $T=0$; $m=1$, $N=112$ for all five couplings at $T=0$ and at $T=0.1$; $m=3$, $N=8$ for $\hat\lambda\in\{0,\pm\hat\lambda_1\}$ at $T=0$. These 18 configurations are the ones described in this section. Seven more, each with its \texttt{pairing.\allowbreak{}json}, were committed after it was written: $m=3$, $N=112$ for $\hat\lambda\in\{0,+\hat\lambda_1\}$ at $T=0$ and at $T=0.1$ and for $-\hat\lambda_1$ at $T=0$, and $m=3$, $N=8$ for $\pm\hat\lambda_2$ at $T=0$ (Section 19.10). They show the same pairing (every level matched, no mismatch of multiplicities, total energies equal to at most $1.6\times10^{-8}$, level energies to at most $1.1\times10^{-9}$ except $1.8\times10^{-5}$ in \texttt{d16c\_m3\_\allowbreak{}L3\_N112\_\allowbreak{}lamp1\_T0\allowbreak{}p1}); in the two runs with $\pm\hat\lambda_2$ the excited-state iteration did not converge in either universe (field \texttt{excitedC\allowbreak{}onverged}). Still missing are the other five configurations of dirac16complex with $m=3$, $N=112$ and every dirac16complex00 configuration, and no summary file of the runs is committed. The untransformed control was solved only at $\lambda=0$; for $\lambda\ne0$ the solver refused to run it, for the following reason. The solver looks for the Kohn--Sham levels inside an energy window whose lower end, the \textbf{window floor}, is the fixed negative energy $-(2.5\,|m|+2\pi/L)$, that is $-4.594$ for $m=1$ and $L=3$ (function \texttt{window\_f\allowbreak{}loor} in \texttt{src/scf.\allowbreak{}rs} of the solver), and it labels each level as a particle or a sea level by following it from its free partner (Section 13.7). The free particle levels lie at $\varepsilon\ge0$, so this labelling is safe as long as the interaction cannot push a particle level below the floor. The interaction changes the block Hamiltonian $h_j$ of Section 15.7 by $j(M_{\mathrm{eff}}-m)\sigma_y+v$, and no level can move by more than the largest size of this change, $\max|M_{\mathrm{eff}}-m|+\max|v_x|$ (the solver's \textbf{shift bound}). The solver therefore requires, before its first self-consistent update, that the shift bound be at most the absolute value of the window floor; it calls this requirement its window premise. For the two paired universes the bound of the first update is $0.1$ (for $\hat\lambda_1$ by construction, Section 13.10); for the control it is $40.34$, against $4.594$, for $m=1$, $N=8$, $+\hat\lambda_1$ (\texttt{d16c\_m1\_\allowbreak{}L3\_N8\_la\allowbreak{}mp1\_T0/p\allowbreak{}airing.j\allowbreak{}son}, keys \texttt{firstUpd\allowbreak{}ateShift\allowbreak{}Bound\_ab\allowbreak{}sWindowF\allowbreak{}loor} and \texttt{controlO\allowbreak{}utcome}). The reason is that the untransformed tip condition puts the zero modes at the tip, where the proper densities, and with them $\lambda S_p$ and $v_x$, carry the factor $e^{6HL}\approx6.6\times10^7$ (Section 13.7). The \texttt{controlO\allowbreak{}utcome} of each interacting \texttt{pairing.\allowbreak{}json} records the same refusal with its own numbers.
\item \textbf{Reference solver}, in \texttt{artifact\allowbreak{}s/dirac1\allowbreak{}6complex\allowbreak{}/pair-cr\allowbreak{}eation/r\allowbreak{}eference\allowbreak{}/} with the summary \texttt{referenc\allowbreak{}e-pairs-\allowbreak{}summary.\allowbreak{}json}, which records itself as incomplete (\texttt{complete\allowbreak{}: false}, 131 runs pending): converged runs exist for $m=3$, $N=112$; for dirac16complex at $\lambda=0$ ($T=0$ and $T=0.1\,|m|$, all three universes) and at $\pm\hat\lambda_1$ ($T=0$, the $+M$ and $-M$ universes); for dirac16complex00 at $\lambda=0$ (both temperatures, all three universes) and the $+M$ universe alone at $+\hat\lambda_1$. The run at $+\hat\lambda_2$ did not converge (a collapse of the self-consistent loop was suspected). The interacting controls failed: in four of them the energy window of the reference solver had to grow to $[-64,\,64.59]$, beyond the solver's cap of $20\,|m|$, and the self-consistent loop was aborted, and in the fifth the level search did not terminate (entries \texttt{\detokenize{failed}} and \texttt{failedCo\allowbreak{}ntrol} of the summary). The interacting runs at $T>0$ (all three universes, both fields) were not attempted: with the couplings $\pm\hat\lambda_1$ and $\pm\hat\lambda_2$ of $T=0$, their first-order pseudo-potential in the free state at $T=0.3$ is $84\,|m|$ and $841\,|m|$, beyond the attempt limit $10\,|m|$ of the reference solver (entries \texttt{notAttem\allowbreak{}pted}).
\item \textbf{Not done}: the cross-check report of \texttt{scripts/\allowbreak{}check\_di\allowbreak{}rac16com\allowbreak{}plex\_pai\allowbreak{}rs.py} (no report is committed), the Stage-5 gate, and any figure of the pairs. No numerical test of T3 exists for dirac16complex00 at $\lambda\ne0$, where its exchange sign matters; at $\lambda=0$ the two fields are the same problem.
\end{itemize}

Two identity checks were completed: the $+M$ universe of every $T=0$ Rust configuration for which a Stage-4 run with the same parameters is committed (all except $m=3$, $N=8$ at $\pm\hat\lambda_2$) reproduces that Stage-4 run bit for bit (\texttt{stage4Re\allowbreak{}producti\allowbreak{}on} in each \texttt{pairing.\allowbreak{}json}), and the Stage-4 outputs of the Rust solver are unchanged by the Stage-5 additions (all nine checks of \texttt{artifact\allowbreak{}s/dirac1\allowbreak{}6complex\allowbreak{}/pair-cr\allowbreak{}eation/r\allowbreak{}ust/stag\allowbreak{}e4-ident\allowbreak{}ity-repo\allowbreak{}rt.json} true).

\textbf{Rust results at $T=0$} (dirac16complex, $L=3$, $H=1$; energies in units of $H$, which is $m$ for $m=1$; the run label gives $m$, $N$ and $\hat\lambda/\hat\lambda_1$, where $\pm10$ means $\pm\hat\lambda_2$; numbers from the \texttt{pairing.\allowbreak{}json} of each configuration). The columns give the ground-state energy of the $+M$ universe, its difference from that of the transformed $-M$ universe, and the total scalar charges $S_{\mathrm{tot}}=\ell^3\int S_c\,dy$ of both:

\begingroup
\small
\begin{longtable}{@{}>{\raggedright\arraybackslash}p{0.160\linewidth}>{\raggedright\arraybackslash}p{0.160\linewidth}>{\raggedright\arraybackslash}p{0.160\linewidth}>{\raggedright\arraybackslash}p{0.160\linewidth}>{\raggedright\arraybackslash}p{0.160\linewidth}@{}}
\toprule
\textbf{$m$, $N$, $\hat\lambda/\hat\lambda_1$} & \textbf{$E_0(+M)$} & \textbf{$\lvert E_0(+M)-E_0(-M)\rvert$} & \textbf{$S_{\mathrm{tot}}(+M)$} & \textbf{$S_{\mathrm{tot}}(-M)$} \\
\midrule
\endfirsthead
\toprule
\textbf{$m$, $N$, $\hat\lambda/\hat\lambda_1$} & \textbf{$E_0(+M)$} & \textbf{$\lvert E_0(+M)-E_0(-M)\rvert$} & \textbf{$S_{\mathrm{tot}}(+M)$} & \textbf{$S_{\mathrm{tot}}(-M)$} \\
\midrule
\endhead
1, 8, 0 & 0 & 0 & 0 & 0 \\
1, 8, +1 & $-0.000986833163$ & $1.3\times10^{-12}$ & $+0.00790445514$ & $-0.00790445514$ \\
1, 8, $-1$ & $+0.000986833163$ & $1.3\times10^{-12}$ & $-0.00790445514$ & $+0.00790445514$ \\
1, 8, +10 & $-0.00783039195$ & $1.7\times10^{-10}$ & $+0.0506051452$ & $-0.0506051448$ \\
1, 8, $-10$ & $+0.00783039195$ & $1.7\times10^{-10}$ & $-0.0506051452$ & $+0.0506051448$ \\
1, 112, 0 & 61.0907233 & $7.3\times10^{-11}$ & $-20.853342$ & $+20.853342$ \\
1, 112, +1 & 61.1132120 & $4.9\times10^{-10}$ & $-20.9453233$ & $+20.9453233$ \\
1, 112, $-1$ & 61.0686105 & $1.2\times10^{-9}$ & $-20.7616373$ & $+20.7616373$ \\
1, 112, +10 & 61.3283111 & $7.0\times10^{-9}$ & $-21.7686201$ & $+21.7686201$ \\
1, 112, $-10$ & 60.8913051 & $8.1\times10^{-9}$ & $-20.0394132$ & $+20.0394132$ \\
3, 8, 0 & 0 & 0 & 0 & 0 \\
3, 8, +1 & $-1.19999991$ & $3.7\times10^{-12}$ & below $10^{-12}$ & below $10^{-12}$ \\
3, 8, $-1$ & $+1.19999991$ & $3.7\times10^{-12}$ & below $10^{-12}$ & below $10^{-12}$ \\
\bottomrule
\end{longtable}
\endgroup

The $+M$ energies reproduce the committed Stage-4 runs bit for bit (key \texttt{stage4Re\allowbreak{}producti\allowbreak{}on} of each \texttt{pairing.\allowbreak{}json}). For $m=1$ they are the values of the table of Section 13.11; the $m=3$ runs, which that table does not list, are \texttt{artifact\allowbreak{}s/dirac1\allowbreak{}6complex\allowbreak{}/kohn-sh\allowbreak{}am/rust/\allowbreak{}scf/m3\_L\allowbreak{}3\_N8\_lam\allowbreak{}0\_T0/run\allowbreak{}.json} and the two runs \texttt{m3\_L3\_N8\allowbreak{}\_lamp1\_T\allowbreak{}0} and \texttt{m3\_L3\_N8\allowbreak{}\_lamm1\_T\allowbreak{}0} next to it ($E_0=\mp1.1999999136735815$ at $\pm\hat\lambda_1$). For $m=3$, $N=8$ the eight occupied orbitals are zero modes with $\chi_2\equiv0$, whose scalar density vanishes (Section 13.11), so only the exchange potential acts. The first excited state:

\begingroup
\small
\begin{longtable}{@{}>{\raggedright\arraybackslash}p{0.133\linewidth}>{\raggedright\arraybackslash}p{0.133\linewidth}>{\raggedright\arraybackslash}p{0.133\linewidth}>{\raggedright\arraybackslash}p{0.133\linewidth}>{\raggedright\arraybackslash}p{0.133\linewidth}>{\raggedright\arraybackslash}p{0.133\linewidth}@{}}
\toprule
\textbf{$m$, $N$, $\hat\lambda/\hat\lambda_1$} & \textbf{$\Delta_{KS}(+M)$} & \textbf{difference} & \textbf{$\Delta_{\mathrm{SCF}}(+M)$} & \textbf{difference} & \textbf{largest $\lvert\Delta\varepsilon\rvert$ over all levels} \\
\midrule
\endfirsthead
\toprule
\textbf{$m$, $N$, $\hat\lambda/\hat\lambda_1$} & \textbf{$\Delta_{KS}(+M)$} & \textbf{difference} & \textbf{$\Delta_{\mathrm{SCF}}(+M)$} & \textbf{difference} & \textbf{largest $\lvert\Delta\varepsilon\rvert$ over all levels} \\
\midrule
\endhead
1, 8, 0 & 0.430733679 & $3.0\times10^{-12}$ & 0.430733679 & $3.0\times10^{-12}$ & $3.3\times10^{-10}$ \\
1, 8, +1 & 0.430943911 & $4.3\times10^{-12}$ & 0.430938122 & $8.2\times10^{-12}$ & $2.9\times10^{-10}$ \\
1, 8, $-1$ & 0.430515902 & $7.4\times10^{-12}$ & 0.430523397 & $2.6\times10^{-11}$ & $3.0\times10^{-10}$ \\
1, 8, +10 & 0.431898697 & $5.9\times10^{-13}$ & 0.431945853 & $1.9\times10^{-10}$ & $4.9\times10^{-9}$ \\
1, 8, $-10$ & 0.429137447 & $1.6\times10^{-11}$ & 0.429160673 & $1.0\times10^{-10}$ & $6.6\times10^{-9}$ \\
1, 112, 0 & 0.0955462022 & $3.7\times10^{-12}$ & 0.0955462022 & $3.7\times10^{-12}$ & $3.3\times10^{-10}$ \\
1, 112, +1 & 0.0954841338 & $9.5\times10^{-12}$ & 0.0954841923 & $9.4\times10^{-10}$ & $3.2\times10^{-10}$ \\
1, 112, $-1$ & 0.0956051341 & $2.5\times10^{-11}$ & 0.0956050822 & $1.3\times10^{-9}$ & $3.1\times10^{-10}$ \\
1, 112, +10 & 0.0948132008 & $6.2\times10^{-11}$ & 0.0948140420 & $2.0\times10^{-9}$ & $1.6\times10^{-9}$ \\
1, 112, $-10$ & 0.0960101550 & $1.5\times10^{-11}$ & 0.0960098274 & $8.9\times10^{-10}$ & $3.8\times10^{-9}$ \\
3, 8, 0 & 0.895832758 & $1.7\times10^{-11}$ & 0.895832758 & $1.7\times10^{-11}$ & $6.0\times10^{-10}$ \\
3, 8, +1 & 0.895832758 & $5.6\times10^{-12}$ & 0.897264439 & $8.2\times10^{-10}$ & $6.4\times10^{-10}$ \\
3, 8, $-1$ & 0.895832758 & $2.4\times10^{-12}$ & 0.894861136 & $4.6\times10^{-10}$ & $6.2\times10^{-10}$ \\
\bottomrule
\end{longtable}
\endgroup

In every configuration of this section, at $T=0$ and at $T=0.1\,m$, each level of the $+M$ universe (between 662 and 977 levels per configuration in the solver's energy window) has its partner under the map $(\text{shell},p,s,n)\to(\text{shell},-p,-s,-n)$ with no mismatch of multiplicity, the orbitals agree after the swap $(a,b)\to(b,a)$ to at most $1.5\times10^{-8}$, the profiles $n_c$, $n_p$, $v$, $\rho$, $p_y$, $p_3$, $p_t$ agree and $S_c$, $S_p$ are opposite to at most $2.0\times10^{-12}$ in the solver's normalized norms, and $M_{\mathrm{eff}}$ is opposite to at most $1.2\times10^{-7}$ relative (keys \texttt{pairingD\allowbreak{}eviation\allowbreak{}\_plusM\_v\allowbreak{}s\_minusM} of the \texttt{pairing.\allowbreak{}json} files).

\textbf{Rust results at $T=0.1\,m$} ($m=1$, $N=112$): the Mermin states of Section 13.13 and their partners.

\begingroup
\small
\begin{longtable}{@{}>{\raggedright\arraybackslash}p{0.133\linewidth}>{\raggedright\arraybackslash}p{0.133\linewidth}>{\raggedright\arraybackslash}p{0.133\linewidth}>{\raggedright\arraybackslash}p{0.133\linewidth}>{\raggedright\arraybackslash}p{0.133\linewidth}>{\raggedright\arraybackslash}p{0.133\linewidth}@{}}
\toprule
\textbf{$\hat\lambda/\hat\lambda_1$} & \textbf{$E$} & \textbf{$F$} & \textbf{$S_{\mathrm{ent}}$} & \textbf{$\mu$} & \textbf{largest difference of these four} \\
\midrule
\endfirsthead
\toprule
\textbf{$\hat\lambda/\hat\lambda_1$} & \textbf{$E$} & \textbf{$F$} & \textbf{$S_{\mathrm{ent}}$} & \textbf{$\mu$} & \textbf{largest difference of these four} \\
\midrule
\endhead
0 & 69.5093780 & 53.9716429 & 155.377351 & 0.697334281 & $1.7\times10^{-9}$ \\
+1 & 69.5379866 & 53.9885813 & 155.494053 & 0.697612981 & $1.6\times10^{-9}$ \\
$-1$ & 69.4815823 & 53.9546932 & 155.268891 & 0.697058519 & $6.9\times10^{-10}$ \\
+10 & 69.8065526 & 54.1412640 & 156.652886 & 0.700242844 & $1.5\times10^{-8}$ \\
$-10$ & 69.2714394 & 53.8109715 & 154.604679 & 0.694812714 & $1.5\times10^{-8}$ \\
\bottomrule
\end{longtable}
\endgroup

The row $\lambda=0$ is the Stage-4 thermodynamics of Section 13.13 ($E=69.50938$, $F=53.97164$).

\textbf{How good is the agreement?} The pairing holds, level by level, to between $10^{-12}$ and $10^{-8}$. That is the precision of the solver itself: its tolerances are $10^{-12}$ (relative) and $10^{-14}$ (absolute) in the integration and $10^{-10}$ in the self-consistent loop (\texttt{paramete\allowbreak{}rs} of each \texttt{\detokenize{run.json}}), and the Stage-4 refined-tolerance run changes energies by up to $5.96\times10^{-8}$ relative (Section 13.10). The two universes are computed as independent problems, with different boundary conditions and different signs of the mass, so this agreement is a genuine numerical test of Theorem 15.13, not a copy.

\textbf{The pair totals.} Each \texttt{pairing.\allowbreak{}json} also forms the two pair totals of Section 15.8 from the two runs: the mirror pair ($+M$ plus $-M$) and the formal Krein image pair ($+M$ minus $-M$ for every one-body density and energy). For $m=1$, $N=112$, $+\hat\lambda_2$, $T=0$:

\begingroup
\small
\begin{longtable}{@{}>{\raggedright\arraybackslash}p{0.267\linewidth}>{\raggedright\arraybackslash}p{0.267\linewidth}>{\raggedright\arraybackslash}p{0.267\linewidth}@{}}
\toprule
\textbf{quantity} & \textbf{mirror pair, $(+M)+(-M)$} & \textbf{Krein image pair, $(+M)-(-M)$} \\
\midrule
\endfirsthead
\toprule
\textbf{quantity} & \textbf{mirror pair, $(+M)+(-M)$} & \textbf{Krein image pair, $(+M)-(-M)$} \\
\midrule
\endhead
energy $E$ & 122.6566221 & $-7.0\times10^{-9}$ \\
charge (number of quanta) & 224 & 0 \\
total scalar charge $S_{\mathrm{tot}}$ & $-5.3\times10^{-9}$ & $-43.5372402$ \\
$\langle\rho\rangle$ & 0.0463577818 & $-2.7\times10^{-12}$ \\
$\langle p_y\rangle$ & 0.0220049755 & $-2.6\times10^{-13}$ \\
$\langle p_3\rangle$ & 0.0137914530 & $-1.8\times10^{-13}$ \\
\bottomrule
\end{longtable}
\endgroup

The mirror pair has twice the energy, charge and pressures of one universe and zero scalar charge; the Krein image pair has zero energy, charge and pressures and twice the scalar charge, $2\times(-21.7686201)$. Both are what Sections 15.7 and 15.8 predict. The two columns carry different information. The mirror column adds two universes that were computed independently. The Krein image column is the difference $(+M)-(-M)$ of the same two runs (the \texttt{definiti\allowbreak{}on} of \texttt{kreinIma\allowbreak{}gePair} in the file): its zero energy, charge and pressures restate, to the precision of the solver, the equality $E(+M)=E(-M)$ of T3. They are a zero of the $X+(-X)$ kind (Section 15.8), not a computed cancellation between two independent universes.

\textbf{The untransformed control} (the $-M$ universe with the Stage-4 bag $\theta=0$), at $\lambda=0$:

\begingroup
\small
\begin{longtable}{@{}>{\raggedright\arraybackslash}p{0.160\linewidth}>{\raggedright\arraybackslash}p{0.160\linewidth}>{\raggedright\arraybackslash}p{0.160\linewidth}>{\raggedright\arraybackslash}p{0.160\linewidth}>{\raggedright\arraybackslash}p{0.160\linewidth}@{}}
\toprule
\textbf{configuration} & \textbf{quantity} & \textbf{$+M$ universe} & \textbf{$-M$, transformed} & \textbf{$-M$, untransformed} \\
\midrule
\endfirsthead
\toprule
\textbf{configuration} & \textbf{quantity} & \textbf{$+M$ universe} & \textbf{$-M$, transformed} & \textbf{$-M$, untransformed} \\
\midrule
\endhead
$m=1$, $N=8$, $T=0$ & $\Delta_{KS}$ & 0.4307336786 & 0.4307336786 & 0.1008511214 \\
$m=3$, $N=8$, $T=0$ & $\Delta_{KS}$ & 0.8958327583 & 0.8958327583 & 0.0007404590 \\
$m=1$, $N=112$, $T=0$ & $E_0$ & 61.0907233 & 61.0907233 & 56.2654962 \\
$m=1$, $N=112$, $T=0.1$ & $F$ & 53.9716429 & 53.9716429 & 49.2723994 \\
\bottomrule
\end{longtable}
\endgroup

The control does not pair. For $N=8$ its eight occupied orbitals are the tip zero modes at $\varepsilon=0$, and its lowest excitation goes to the sub-gap bound state of Section 15.9: the Rust gaps $0.10085112137571$ ($m=1$) and $0.00074045901682$ ($m=3$) agree with $\varepsilon_{\mathrm b}=M/\cosh(qL)=0.10085112137514$ and $0.00074045901623$ to $6\times10^{-13}$. For $N=112$ the control's ground state is a different state altogether, $4.8\,m$ lower in energy, in which a level of 32 states at the Fermi level is filled to three quarters, so that its Kohn--Sham gap is 0 (\texttt{rust/pai\allowbreak{}rs/d16c\_\allowbreak{}m1\_L3\_N1\allowbreak{}12\_lam0\_\allowbreak{}T0/minus\allowbreak{}M\_contro\allowbreak{}l/levels\allowbreak{}.csv}, column \texttt{\detokenize{f}}).

\textbf{The reference solver} ($m=3$, $N=112$, $L=3$, $H=1$; grids $N_0=120$, 240, 480 with the extrapolation of the reference solver; \texttt{referenc\allowbreak{}e-pairs-\allowbreak{}summary.\allowbreak{}json}):

\begingroup
\small
\begin{longtable}{@{}>{\raggedright\arraybackslash}p{0.160\linewidth}>{\raggedright\arraybackslash}p{0.160\linewidth}>{\raggedright\arraybackslash}p{0.160\linewidth}>{\raggedright\arraybackslash}p{0.160\linewidth}>{\raggedright\arraybackslash}p{0.160\linewidth}@{}}
\toprule
\textbf{$\hat\lambda/\hat\lambda_1$, $T$} & \textbf{$E(+M)$} & \textbf{$E(-M)$} & \textbf{$S_{\mathrm{tot}}(+M)$} & \textbf{$S_{\mathrm{tot}}(-M)$} \\
\midrule
\endfirsthead
\toprule
\textbf{$\hat\lambda/\hat\lambda_1$, $T$} & \textbf{$E(+M)$} & \textbf{$E(-M)$} & \textbf{$S_{\mathrm{tot}}(+M)$} & \textbf{$S_{\mathrm{tot}}(-M)$} \\
\midrule
\endhead
0, 0 & 131.448298 & 131.448295 & $-7.206331$ & $+7.206334$ \\
0, 0.3 & 161.026343 & 161.026341 & $-7.971180$ & $+7.971183$ \\
+1, 0 & 127.409033 & 127.409029 & $-8.466395$ & $+8.466399$ \\
$-1$, 0 & 135.701559 & 135.701558 & $-6.388991$ & $+6.388992$ \\
\bottomrule
\end{longtable}
\endgroup

These are dirac16complex runs; $T=0.3$ is $T=0.1\,|m|$, where the free energies are $104.727753$ and $104.727750$. The Kohn--Sham gaps of the two universes are

\begingroup
\small
\begin{longtable}{@{}>{\raggedright\arraybackslash}p{0.267\linewidth}>{\raggedright\arraybackslash}p{0.267\linewidth}>{\raggedright\arraybackslash}p{0.267\linewidth}@{}}
\toprule
\textbf{$\hat\lambda/\hat\lambda_1$, $T$} & \textbf{$\Delta_{KS}(+M)$} & \textbf{$\Delta_{KS}(-M)$} \\
\midrule
\endfirsthead
\toprule
\textbf{$\hat\lambda/\hat\lambda_1$, $T$} & \textbf{$\Delta_{KS}(+M)$} & \textbf{$\Delta_{KS}(-M)$} \\
\midrule
\endhead
0, 0 & 0.2342637881 & 0.2342637838 \\
0, 0.3 & 0.2794285932 & 0.2794285878 \\
+1, 0 & 0.2350337045 & 0.2350336992 \\
$-1$, 0 & 0.2336532576 & 0.2336532545 \\
\bottomrule
\end{longtable}
\endgroup

The dirac16complex00 runs at $\lambda=0$ give the same numbers to all digits, as they must, since the statistics enters only through the exchange term, which is proportional to $\lambda$. The reference controls at $\lambda=0$ give $E_0=119.118131$ at $T=0$ and $F=93.121665$ at $T=0.3$, far from the paired values. The pairing defect of the reference solver, about $3\times10^{-6}$ in $E$, is much larger than that of the Rust solver. This too was predicted (\texttt{pairing-\allowbreak{}theory.j\allowbreak{}son}, \texttt{T3.numer\allowbreak{}icsPresc\allowbreak{}ription.\allowbreak{}discreti\allowbreak{}sationNo\allowbreak{}te}): the reference solver puts the two components of $\chi$ on two staggered grids, and the swap $\chi\to\sigma_y\chi$ exchanges the roles of the two grids, so the two universes are discretized differently and pair only up to the discretization error. Both values agree with the Stage-4 Rust energy $131.4482953$ of the same $+M$ problem (\texttt{artifact\allowbreak{}s/dirac1\allowbreak{}6complex\allowbreak{}/kohn-sh\allowbreak{}am/rust/\allowbreak{}scf/m3\_L\allowbreak{}3\_N112\_l\allowbreak{}am0\_T0/r\allowbreak{}un.json}) to within $3\times10^{-6}$.

\textbf{Reproduction.} The exact verifier and the independent checker run from the repository root; written into \texttt{\detokenize{build/}}, they leave the committed reports untouched. In Git Bash:

\begin{Verbatim}[fontsize=\small]
export PYTHONUTF8=1
wolframscript -file scripts/verify_dirac16complex_pairing.wls \
    build/pairing/wolfram-pairing-report.json
python scripts/check_dirac16complex_pairing.py \
    --output build/pairing/python-pairing-report.json
\end{Verbatim}

In PowerShell set \texttt{\$env:PYT\allowbreak{}HONUTF8 \allowbreak{}= "1"} and replace each line-ending backslash by a backtick. Each program prints one line per check and ends with the lines \texttt{check\_co\allowbreak{}unt} and \texttt{failed\_c\allowbreak{}heck\_cou\allowbreak{}nt}. For the Kohn--Sham pair runs, the solver's \texttt{README.m\allowbreak{}d} (section "Stage 5") gives the Rust command, and the reference script takes an output folder; in Git Bash, writing into \texttt{\detokenize{build/}} (the Rust subcommand writes into the subfolder \texttt{\detokenize{pairs/}} of the folder given):

\begin{Verbatim}[fontsize=\small]
cargo build --release --manifest-path studies/dirac16complex_kohn_sham/Cargo.toml
./studies/dirac16complex_kohn_sham/target/release/dirac16complex_kohn_sham \
    pairs --output build/pairs-rust
python scripts/ks_reference_pairs.py --output build/pairs-reference
\end{Verbatim}

Section 19.10 records which of these runs are committed.

\subsection{What the pairing theorems establish, and what they do not}

\textbf{Established (proved, with exact machine checks):}

\begin{itemize}
\item For both fields and in every gravitational field, the chirality map sends the theory $(m,\lambda)$ to the theory $(-m,-\lambda)$, whose action is minus the original one, and every configuration to a configuration; the energy--momentum tensor and the current, computed from the bilinears of the field, change sign (T1).
\item A T1 pair of classical fields has zero total energy--momentum, charge and action at every point and every time, so the gravitational field equations with the pair as source are the source-free equations, and the conservation of the totals does not forbid its appearance from the empty state (Corollary 15.4). In the theory as built the two members do not interact, so the charge of each member, and in a static gravitational field its energy, is conserved on its own; this forbids the appearance of a T1 pair from the empty state unless each member is neutral (and, in such a field, has zero energy) (Section 15.4, remark 2; Proposition 16.5).
\item For the quantized field the image $\gamma^8\Psi$ carries the anticommutator $-B$, and in the exact Fock model its operators satisfy $H[\Psi_-;-m]=-H[\Psi;m]$, $Q[\Psi_-]=-Q[\Psi]$ and $S[\Psi_-]=S[\Psi]$ (cited from Stage 5, PROVISIONAL until the Stage-5 gate). These are identities $X+(-X)=0$ within one quantum system: measured with its own canonical structure the image field has $+H_+$, $+Q_+$ and $+T_{\mu\nu}[\Psi_+]$ (Section 15.4). An independently quantized $-M$ field has positive energies.
\item A reflection of a space-like direction, acting on the field and the frame, sends $(m,\lambda)$ to $(-m,\lambda)$ with the same energy--momentum and charge (T2); across the Z2 brane of the static primordial field this is the mirror universe of mass $-M$, and every Stage-4 Kohn--Sham state is such a mirror pair.
\item In the Kohn--Sham model of Chapters 13 and 14, for both statistics, the universe of mass $-M$ with the same $\lambda$ and the transformed tip condition has exactly the same ground state, first excited state, thermodynamics and energy--momentum tensor as the universe of mass $+M$, and the opposite scalar density (T3); with the untransformed tip condition the pairing fails.
\end{itemize}

\textbf{Computed:} the Rust and reference results of Section 15.10, which confirm T3 to the precision of each solver for dirac16complex in the configurations listed there, and confirm the failure of the untransformed control.

\textbf{Not established:}

\begin{itemize}
\item No process that creates a pair is derived: no rate, probability or amplitude, no wave function of the universe, no dynamical big bang. The theorems show that the totals of a cancelling pair of classical fields are consistent with the conservation of the totals and with the gravitational constraints, not that such a pair occurs.
\item The cancelling (T1) pair is established for classical fields only, with the coupling $-\lambda$ (or $\lambda=0$) and the exact correlation $\Psi_-=\gamma^8\Psi_+$ for the second member. For the quantized field no reading examined in the repository gives two independent universes whose energies and charges cancel: the image field with the metric $-B$ is the same quantum system, and an independently quantized $-M$ universe has positive energies, which add. A quantum description of two independent universes whose energies and charges cancel is \textbf{OPEN} (Section 15.4; Section 16.6).
\item Nothing in the theory selects the cancelling pair over the mirror pair of T2 and T3, whose energy doubles, and no interaction between the members, which would be needed before only the totals are conserved, is contained in the theory (OPEN; Section 16.9).
\item The Kohn--Sham statements concern a mean-field model (no correlation, fixed metric, the good sector); for dirac16complex00 the model is a formal functional with Pauli filling imposed by prescription.
\item The numerical demonstration is incomplete (Section 15.10): no Rust run of dirac16complex00, no run with $m=3$, $N=112$ in Rust, no cross-check report and no Stage-5 gate.
\item The primordial field still requires a negative energy density in eight-dimensional Einstein gravity; the pairing does not supply it. Whether the field solves the vacuum Einstein--Lovelock equations for some couplings, so that a pair with zero totals could accompany it, is OPEN, because the Lovelock terms of orders 2 and 3 were not computed (row L25 of Section 0.9; Section 18.5).
\end{itemize}

Whether the big bang creates universes in pairs is therefore \textbf{not} answered by these theorems; it remains the notebook's hypothesis, and Chapter 16 discusses it. What the opposite charges of a T1 pair mean for the matter--antimatter asymmetry is the subject of Chapter 17.

\subsection{What we proved and what we assumed}

\textbf{What we proved.} With complete derivations in this chapter: how a constant matrix passes through a bilinear, for commuting and for Grassmann components (Lemma 15.1); the parity of products of gammas under $\gamma^8$ and the fact that $\gamma^8$ reverses $B$ (Lemma 15.2, fact (F6)); Theorem T1 in every gravitational field, for both statistics, with the Lagrangian, field equations, energy--momentum tensor and current (Theorem 15.3), the failure at fixed $\lambda$, the equivalence with a reversal of the frame, and the fact that the partner theory has the negated action; the vanishing totals of the chiral pair of classical fields and the source-free gravitational equations (Corollary 15.4), together with the separate conservation of the charge (and, in a static field, of the energy) of each member, which forbids the appearance of a pair from the empty state unless each member is neutral; the anticommutator $-B$ of the image field, and the fact that the image field, measured with its own canonical structure, has the generators $+H_+$ and $+Q_+$ and the source $+T_{\mu\nu}[\Psi_+]$, so that its Fock-level cancellation is an identity $X+(-X)=0$ within one quantum system; Theorem T2 for every unit vector in every gravitational field, with the reflected frame built from the covariance theorem of Section 4.12 (Theorem 15.5); its explicit form across the Z2 brane, the pull-back of the energy--momentum tensor and the current, and the odd mass function of a Z2-symmetric configuration (Lemma 15.6, Theorem 15.7, Corollary 15.8); the block form $\tau\,\sigma_y$ of $\gamma^8$ (Lemma 15.9), the map $\sigma_yh_j(M,v)\sigma_y=h_{-j}(-M,v)$ (Lemma 15.10), the transformation of the parity and bag conditions (Lemma 15.11) and of the densities (Lemma 15.12); Theorem T3 for both statistics (Theorem 15.13) and its formal Krein-image form with $-\lambda$, whose zero pair totals are of the $X+(-X)$ kind; the explanation of the difference between T1 and T3 by the parity of $B$; and, for the untransformed control, the tip zero mode, the splitting constant $c_{\mathrm{ctrl}}$ and the sub-gap bound state with $\tanh(qL)=q/M$.

\textbf{What we took from elsewhere.} The algebra of the gammas, $C$, $B$, the Pin group and its reflections (Chapter 2); the vielbein, the spin connection and the covariance theorem (Chapter 4); the Lagrangian, field equations and energy--momentum tensor (Chapters 6 and 7); the canonical anticommutator, the positive Fock space and the expectation-value rule (Chapter 8); the static primordial field and the Z2 gluing (Chapter 9); the block reduction, boundary conditions, densities, Kohn--Sham functional and energy--momentum formulas (Chapter 13) and the statistics sign of the exchange (Chapter 14); the Fock-level model of the image field (provisional until the Stage-5 gate) and the numbers of the one-particle table, from the Stage-5 and matter--antimatter reports; the proof of charge conservation, Theorem M1 of the matter--antimatter document, and the separate conservation laws of a pair of fields that do not interact (Proposition 16.5); the uniqueness theorem of Picard and Lindelöf (Chapter 10). The exact checks at test points confirm the general derivations; they do not replace them.

\textbf{What we assumed.} The Z2 gluing of the primordial field is a choice of model (Section 9.16). The Kohn--Sham model is an approximation with the assumptions listed after Theorem 15.13, and for dirac16complex00 it is defined by prescription. That the two members of a T1 pair live on one spacetime and that their totals add is hypothesis (a) to (c) of Corollary 15.4; that the second member is exactly $\gamma^8\Psi_+$ is part of the definition of the pair, not a result. The cancellation is established for classical fields; whether a quantum description of two independent universes whose energies and charges cancel exists is OPEN. Which of the pairings, if any, describes the notebook's "pair of universes" is not decided by the equations; it is an interpretation, and no statement of this chapter makes it.

\subsection{Exercises}

\textbf{Exercise 15.1.} Decide whether $\gamma^8X\gamma^8=+X$ or $-X$ for $X=C$, $C\gamma^3$, $C\gamma^0\gamma^1$, $B$, $BC$, $BC\gamma^1$ and $\gamma^0\gamma^4$.

\textbf{Exercise 15.2.} Show that $E_{m,\lambda}[\gamma^8\Psi]=-\gamma^8E_{m,\lambda}[\Psi]-2(m+\lambda S)\gamma^8\Psi$. Conclude that the image of a solution with $(m,\lambda)$ is not a solution with $(m,\lambda)$ wherever $\Psi\ne0$ and $m+\lambda S\ne0$.

\textbf{Exercise 15.3.} In flat space and Gaussian normal gauge, the Lagrangian $\mathcal L_{m,\lambda}$ has, after a total derivative, the form $\Psi^\dagger\mathcal K\partial_4\Psi-\mathcal H$ with $\mathcal K=C\gamma^4$ (Section 8.5, where this matrix is called $K$; here $K$ is the kinetic term). Use the rule $\{\Psi,\Psi^\dagger\}=i\mathcal K^{-1}\delta^7$ of Section 8.6 to find the anticommutator of the Lagrangian $-\mathcal L_{-m,-\lambda}$, and compare it with $\{\gamma^8\Psi,(\gamma^8\Psi)^\dagger\}$ computed from $\{\Psi,\Psi^\dagger\}=B\,\delta^7$. What are then the image field's own generator of the time evolution and its own charge, in terms of $H_+$ and $Q_+$?

\textbf{Exercise 15.4.} Let $\tilde u=\gamma^8u_+=\tfrac12(-e_0+e_4-ie_9-ie_{13})$. With the rows of $\gamma^4$, $C$ and $B$ quoted in Section 15.3, show $-i\gamma^4\tilde u=-\tilde u$, $C\tilde u=\tilde u$ and $B\tilde u=-\tilde u$, and recompute the table of Section 15.8.

\textbf{Exercise 15.5.} For a time-like unit vector $u$ ($n(u)=-1$) go through the proof of Theorem 15.5 and show $\mathcal L_{m,\lambda}[e_u,u\Psi]=-\mathcal L_{-m,-\lambda}[e,\Psi]$. Then compute $uBu^\dagger$ for $u=\gamma^4$ and for $u=\gamma^5$.

\textbf{Exercise 15.6.} Show that $S^{0p}\gamma^0=-\gamma^0S^{0p}$ for $p\ne0$, and that the coefficient $c_l=-W'e^{a_{4,0}}$ of $\Omega_l$ is odd in $y$ for $W=e^{-H|y|}$. Deduce $D_l\Psi'(y)=\gamma^0(D_l\Psi)(-y)$ for $\Psi'(y)=\gamma^0\Psi(-y)$.

\textbf{Exercise 15.7.} Check that $\chi(y)=(e^{-M|y|},0)$ solves the block equation of Section 13.5 at $k=0$, $\varepsilon=0$, $v=0$ on both sides of the brane with the mass function $M(y)=-M\,\mathrm{sgn}(y)$, that it satisfies $\chi(y)=\sigma_z\chi(-y)$, and that its number density is even in $y$.

\textbf{Exercise 15.8.} Using only $\sigma_a\sigma_b=-\sigma_b\sigma_a$ ($a\ne b$) and $\sigma_a^2=1$, show $\sigma_y\sigma_x\sigma_y=-\sigma_x$ and $\sigma_y\sigma_z\sigma_y=-\sigma_z$, and verify $\sigma_yh_j(M,v)\sigma_y=h_{-j}(-M,v)$ term by term.

\textbf{Exercise 15.9.} Show $\sigma_yQ(\theta)\sigma_y=Q(\pi-\theta)$. Write the bag conditions for $\theta=0$, $\theta=\pi$ and $\theta=\pi/2$ in components, and say which of them the map sends into which.

\textbf{Exercise 15.10.} With the vectors $v_\pm$ and $w_\pm$ of Section 15.7, verify $\gamma^8v_+=i\,w_-$ and $\gamma^8v_-=-i\,w_+$ (up to the common factor $1/\sqrt8$), and write the $2\times2$ matrix $G$ of Lemma 15.9.

\textbf{Exercise 15.11.} (a) Transform $(n,s,t,c)$ with the image rule, in which the metric is $-B$. (b) For dirac16complex00 ($M_{\mathrm{eff}}=m+\tfrac{17}{16}\lambda S_p$, $v=\tfrac{\lambda}{16}n_p$) check that the ordinary map needs $(m,\lambda)\to(-m,\lambda)$ and the image map $(m,\lambda)\to(-m,-\lambda)$.

\textbf{Exercise 15.12.} At $\lambda=0$ and $k=0$ the level $\varepsilon=\sqrt{M^2+(\pi/L)^2}$ of parity $+$ has the scalar charge $\int\chi^\dagger j\sigma_y\chi\,dy=d\varepsilon/dM$ (Section 13.9). Show that the image level in the $-M$ problem has the opposite scalar charge, in agreement with Lemma 15.12, and evaluate both for $M=1$, $L=3$.

\textbf{Exercise 15.13.} For the control with $M=1$, $L=3$: bracket the root of $f(q)=\tanh(3q)-q$ in $(0.9,1)$, halve the bracket six times, and compute $\varepsilon_{\mathrm b}=1/\cosh(3q)$. Compare with the Rust gap $0.1008511214$ of the control of Section 15.10. Estimate $\varepsilon_{\mathrm b}$ for $M=3$, $L=3$ with the formula $2Me^{-ML}$.

\textbf{Exercise 15.14.} Show that at $M=H$ and $a_{4,0}=0$, $c(M)=2Y/(Y+1)$ and $c_{\mathrm{ctrl}}=\tfrac23(Y^2+Y+1)/(Y+1)$ with $Y=e^{HL}$, hence $c_{\mathrm{ctrl}}-c=\tfrac23(Y-1)^2/(Y+1)$. Evaluate for $L=3$.

\textbf{Exercise 15.15.} From the tables of Section 15.10 compute, for $m=1$, $N=112$, $+\hat\lambda_2$, $T=0$, the energies and total scalar charges of the mirror pair and of the Krein image pair, and compare with the pair table. Which pair has zero energy and charge, which quantity doubles in each, and what kind of zero is the zero energy of the Krein image pair?

\textbf{Exercise 15.16.} Show that $P_B=i\gamma^0\gamma^8$ is Hermitian with $P_B^2=1$, that $P_BCP_B=-C$, $P_BBP_B=-B$ and therefore $P_B(BC)P_B=BC$. Explain why $P_B$ reverses $\lambda$ for the classical field but keeps it in the Kohn--Sham model.

\subsection{Answers to the exercises}

\textbf{Answer 15.1.} Count the gamma factors (Lemma 15.2): $C$ has 4, even, $+C$; $C\gamma^3$ has 5, odd, $-X$; $C\gamma^0\gamma^1$ has 6, even, $+X$; $B$ has 5, odd, $-B$; $BC$ has 9, odd, $-BC$ (equivalently $BC=-i\gamma^4$ has 1); $BC\gamma^1$ has 10, even, $+X$; $\gamma^0\gamma^4$ has 2, even, $+X$.

\textbf{Answer 15.2.} $E_{m,\lambda}[\Phi]-E_{-m,-\lambda}[\Phi]=-(m+\lambda S)\Phi-(m+\lambda S)\Phi=-2(m+\lambda S)\Phi$ for any $\Phi$, because the two expressions differ only in the sign of $m+\lambda S$ and $S[\gamma^8\Psi]=S[\Psi]$. With $\Phi=\gamma^8\Psi$ and item 4 of Theorem 15.3, $E_{m,\lambda}[\gamma^8\Psi]=-\gamma^8E_{m,\lambda}[\Psi]-2(m+\lambda S)\gamma^8\Psi$. If $\Psi$ solves $E_{m,\lambda}=0$, what remains is $-2(m+\lambda S)\gamma^8\Psi$, which vanishes only where $(m+\lambda S)\Psi=0$, since $\gamma^8$ is invertible.

\textbf{Answer 15.3.} For $\mathcal L_{m,\lambda}$, $\mathcal K=C\gamma^4$. Since $\gamma^4$ commutes with $C$, $(C\gamma^4)(-\gamma^4C)=-C(\gamma^4)^2C=C^2=1$, so $\mathcal K^{-1}=-\gamma^4C=-C\gamma^4$ (Section 8.5 with $g^{44}=-1$) and $i\mathcal K^{-1}=-iC\gamma^4=B$. For $-\mathcal L_{-m,-\lambda}$ the term with $\partial_4\Psi$ does not depend on $m$ or $\lambda$, and the overall sign gives $\mathcal K=-C\gamma^4$, hence $i\mathcal K^{-1}=-B$. On the other hand $\{\gamma^8\Psi,(\gamma^8\Psi)^\dagger\}=\gamma^8\{\Psi,\Psi^\dagger\}\gamma^8=\gamma^8B\gamma^8\,\delta^7=-B\,\delta^7$. The two agree: the image field is canonically a field of $-\mathcal L_{-m,-\lambda}$. Its own Hamiltonian and charge are therefore those of the Lagrangian $-\mathcal L_{-m,-\lambda}$, that is minus the expressions $H[\Psi_-;-m]$ and $Q[\Psi_-]$ of $\mathcal L_{-m,-\lambda}$: $-H[\Psi_-;-m]=+H_+$ and $-Q[\Psi_-]=+Q_+$, by the operator identities of Section 15.4. The same follows from the Heisenberg equation: with $A=-B$, $i[H[\Psi_-;-m],\Psi_-]=-i(-B)Bh_k(-m)\Psi_-=+ih_k(-m)\Psi_-$, the reversed evolution, so the evolution $\partial_4\Psi_-=-ih_k(-m)\Psi_-$ is generated by $-H[\Psi_-;-m]$. The image field is the same quantum system as $\Psi_+$, with the same energy and charge.

\textbf{Answer 15.4.} The components of $\tilde u$ are $u_0=-\tfrac12$, $u_4=\tfrac12$, $u_9=-\tfrac i2$, $u_{13}=-\tfrac i2$. $(\gamma^4u)_0=-u_{13}=\tfrac i2$, so $(-i\gamma^4u)_0=\tfrac12=-u_0$; $(\gamma^4u)_4=u_9=-\tfrac i2$, so $(-i\gamma^4u)_4=-\tfrac12=-u_4$; $(\gamma^4u)_9=-u_4=-\tfrac12$, so $(-i\gamma^4u)_9=\tfrac i2=-u_9$; $(\gamma^4u)_{13}=u_0=-\tfrac12$, so $(-i\gamma^4u)_{13}=\tfrac i2=-u_{13}$. Hence $-i\gamma^4\tilde u=-\tilde u$. $(Cu)_0=-u_4=-\tfrac12=u_0$, $(Cu)_4=-u_0=\tfrac12=u_4$, $(Cu)_9=u_{13}=u_9$, $(Cu)_{13}=u_9=u_{13}$: $C\tilde u=\tilde u$. $(Bu)_0=iu_9=\tfrac12=-u_0$, $(Bu)_4=-iu_{13}=-\tfrac12=-u_4$, $(Bu)_9=-iu_0=\tfrac i2=-u_9$, $(Bu)_{13}=iu_4=\tfrac i2=-u_{13}$: $B\tilde u=-\tilde u$. Since $\tilde u^\dagger \tilde u=4\cdot\tfrac14=1$: classically $\tilde u^\dagger B\tilde u=-1$, $\tilde u^\dagger C\tilde u=1$, energy $m\cdot(-1)=-m$. With the rule and $+B$: charge $u^\dagger B^2u=1$, scalar $u^\dagger(-i\gamma^4)u=-1$, energy $u^\dagger h(-m)u=u^\dagger(im\gamma^4)u=-m\,u^\dagger(-i\gamma^4)u=m$. With $-B$ every rule value changes sign: $-1$, $+1$, $-m$. For $u_+$ the values of Section 8.8 ($-i\gamma^4u_+=u_+$, $Cu_+=u_+$, $Bu_+=u_+$) give 1, 1 and $m$ in both columns.

\textbf{Answer 15.5.} With $n=-1$: $\bar\Psi_u=+\bar\Psi u^{-1}$, so $S\to S$; the kinetic bilinear becomes $(\bar\Psi u^{-1})(-u\gamma^\mu u^{-1})(uD_\mu\Psi)=-\bar\Psi\gamma^\mu D_\mu\Psi$, so $K\to-K$. Then $\mathcal L\to\sqrt{|g|}[-K-mS-\tfrac\lambda2S^2]=-\mathcal L_{-m,-\lambda}$. For $u=\gamma^4$: $u^\dagger=(\gamma^4)^T=-\gamma^4$, and $\gamma^4B(-\gamma^4)=i\gamma^4C\gamma^4\gamma^4=-i\gamma^4C=-iC\gamma^4=B$ (since $\gamma^4$ commutes with $C$). For $u=\gamma^5$: $u^\dagger=-\gamma^5$, and $\gamma^5B(-\gamma^5)=i\gamma^5C\gamma^4\gamma^5=iC\gamma^5\gamma^4\gamma^5=-iC\gamma^4(\gamma^5)^2=iC\gamma^4=-B$ (here $\gamma^5$ commutes with $C$ and anticommutes with $\gamma^4$).

\textbf{Answer 15.6.} For $p\ne0$, $S^{0p}=\tfrac12\gamma^0\gamma^p$, so $S^{0p}\gamma^0=\tfrac12\gamma^0\gamma^p\gamma^0=-\tfrac12\gamma^0\gamma^0\gamma^p=-\tfrac12\gamma^p$ and $\gamma^0S^{0p}=\tfrac12\gamma^0\gamma^0\gamma^p=\tfrac12\gamma^p$: they are opposite. $W=e^{-H|y|}$ is even, so $W'(-y)=-W'(y)$ and $c_l(-y)=-c_l(y)$. Then $D_l[\gamma^0\Psi(-y)]=\gamma^0(\partial_l\Psi)(-y)+c_l(y)S^{0l}\gamma^0\Psi(-y)=\gamma^0(\partial_l\Psi)(-y)-c_l(y)\gamma^0S^{0l}\Psi(-y)=\gamma^0\bigl[\partial_l\Psi+c_l(-y)S^{0l}\Psi\bigr](-y)=\gamma^0(D_l\Psi)(-y)$.

\textbf{Answer 15.7.} At $k=0$, $\varepsilon=v=0$ the block equation is $\chi_1'=M(y)\chi_1$, $\chi_2'=-M(y)\chi_2$. With $\chi_2=0$ and $\chi_1=e^{-M|y|}$: for $y<0$, $\chi_1=e^{My}$ and $\chi_1'=M\chi_1$ with $M(y)=+M$; for $y>0$, $\chi_1=e^{-My}$ and $\chi_1'=-M\chi_1$ with $M(y)=-M$. So $M(y)=-M\,\mathrm{sgn}(y)$, an odd function. $\sigma_z\chi(-y)=(e^{-M|y|},0)=\chi(y)$. The density $\chi^\dagger\chi=e^{-2M|y|}$ is even, and the scalar density $j\chi^\dagger\sigma_y\chi=0$ because $\chi_2=0$.

\textbf{Answer 15.8.} $\sigma_y\sigma_x\sigma_y=-\sigma_x\sigma_y\sigma_y=-\sigma_x$ and $\sigma_y\sigma_z\sigma_y=-\sigma_z\sigma_y\sigma_y=-\sigma_z$. Term by term, with the constant $\sigma_y$ commuting with $d/dy$:

\[
\begin{aligned}
&\sigma_y\bigl(-ij\sigma_x\tfrac{d}{dy}\bigr)\sigma_y=+ij\sigma_x\tfrac{d}{dy}=-i(-j)\sigma_x\tfrac{d}{dy},\qquad \sigma_y\bigl(jM\sigma_y\bigr)\sigma_y=jM\sigma_y=(-j)(-M)\sigma_y,\\
&\sigma_y\bigl(j\xi k\sigma_z\bigr)\sigma_y=-j\xi k\sigma_z=(-j)\xi k\sigma_z,\qquad \sigma_yv\sigma_y=v .
\end{aligned}
\]

Together these are the four terms of $h_{-j}(-M,v)$.

\textbf{Answer 15.9.} $\sigma_yQ(\theta)\sigma_y=\cos\theta\,(-\sigma_z)+\sin\theta\,\sigma_y=\cos(\pi-\theta)\sigma_z+\sin(\pi-\theta)\sigma_y$. $\theta=0$: $Q=\sigma_z$, $(1-\sigma_z)\chi=(0,2\chi_2)$, the condition $\chi_2(-L)=0$. $\theta=\pi$: $Q=-\sigma_z$, the condition $\chi_1(-L)=0$. $\theta=\pi/2$: $Q=\sigma_y$, $(1-\sigma_y)\chi=(\chi_1+i\chi_2,\ \chi_2-i\chi_1)=0$, that is $\chi_1=-i\chi_2$. The map sends $\theta=0$ to $\pi$, $\pi$ to $0$, and $\pi/2$ to itself.

\textbf{Answer 15.10.} $\gamma^8$ multiplies the components 0 to 7 by $-1$. Multiply every vector by $\sqrt8$. The nonzero components of $\sqrt8\,w_-$ are $(i,-1,-i,-1)$ in places 4 to 7 and $(-i,1,i,1)$ in places 12 to 15; times $i$ they become $(-1,-i,1,-i)$ and $(1,i,-1,i)$. The nonzero components of $\sqrt8\,v_+$ are $(1,i,-1,i)$ in places 4 to 7 and the same in places 12 to 15; $\gamma^8$ turns the first group into $(-1,-i,1,-i)$ and keeps the second. So $\gamma^8v_+=i\,w_-$. In the same way $\sqrt8\,v_-$ has $(i,-1,-i,-1)$ in places 0 to 3 and $(-i,1,i,1)$ in places 8 to 11; $\gamma^8$ gives $(-i,1,i,1)$ and $(-i,1,i,1)$; and $-i$ times $\sqrt8\,w_+$, which has $(1,i,-1,i)$ in both groups, is $(-i,1,i,1)$ in both groups. So $\gamma^8v_-=-i\,w_+$. Hence $\gamma^8[v_+\ v_-]=[w_+\ w_-]\,G$ with

\[
G=\begin{pmatrix}0&-i\\ i&0\end{pmatrix}=\sigma_y :
\]

the first column of $G$ holds the coefficients $(0,i)$ of the image of $v_+$, the second those of $v_-$, $(-i,0)$.

\textbf{Answer 15.11.} (a) With $-B$ each expectation value gets one more sign: $(n,s,t,c)\to(-n,\,s,\,-t,\,-c)$. (b) Ordinary map, $(n_p,S_p)\to(n_p,-S_p)$: $m'+\tfrac{17}{16}\lambda'(-S_p)$ must equal $-m-\tfrac{17}{16}\lambda S_p$ for all $S_p$, so $m'=-m$ and $\lambda'=\lambda$; then $v'=\tfrac{\lambda}{16}n_p=v$, as Lemma 15.10 requires. Image map, $(n_p,S_p)\to(-n_p,S_p)$: $m'+\tfrac{17}{16}\lambda'S_p=-m-\tfrac{17}{16}\lambda S_p$ gives $m'=-m$ and $\lambda'=-\lambda$, and then $v'=\tfrac{-\lambda}{16}(-n_p)=v$. The same holds with $\tfrac{15}{16}$ and $-\tfrac{1}{16}$ for dirac16complex.

\textbf{Answer 15.12.} For the $+M$ problem $\varepsilon=\sqrt{M^2+q^2}$ with $q=\pi/L$, so $d\varepsilon/dM=M/\varepsilon$. The $-M$ problem has the mass $M'=-M$ and the same level $\varepsilon=\sqrt{M'^2+q^2}$ (Lemma 15.10 and check \texttt{PAIR\_T3k\allowbreak{}s\_massiv\allowbreak{}eLevelsM\allowbreak{}ap}), so its scalar charge is $d\varepsilon/dM'=M'/\varepsilon=-M/\varepsilon$: opposite, as Lemma 15.12 says. For $M=1$, $L=3$: $q=1.047198$, $\varepsilon=1.447972$ and the scalar charges are $\pm0.690621$.

\textbf{Answer 15.13.} $f(0.9)=+0.0910$ and $f(1)=-0.0049$, so a root lies in $(0.9,1)$. Halving: $f(0.95)=+0.0433$, bracket $(0.95,1)$; $f(0.975)=+0.0193$, $(0.975,1)$; $f(0.9875)=+0.0072$, $(0.9875,1)$; $f(0.99375)=+0.0011$, $(0.99375,1)$; $f(0.996875)=-0.0019$, $(0.99375,0.996875)$; $f(0.9953125)=-0.0004$, $(0.99375,0.9953125)$. Continuing gives $q=0.9949015$. Then $\cosh(3q)=\cosh(2.984705)=9.915606$ and $\varepsilon_{\mathrm b}=0.1008511$, the Rust gap $0.1008511214$ of the control to all digits shown. For $M=3$, $L=3$: $2Me^{-ML}=6e^{-9}=7.4046\times10^{-4}$, close to the exact $7.40459\times10^{-4}$ and to the Rust gap $0.0007404590$.

\textbf{Answer 15.14.} At $M=H=1$: $c=\frac{2}{1}\cdot\frac{1-e^{-L}}{1-e^{-2L}}=\frac{2}{1+e^{-L}}=\frac{2Y}{Y+1}$, and $c_{\mathrm{ctrl}}=\frac23\cdot\frac{e^{3L}-1}{e^{2L}-1}=\frac23\cdot\frac{(Y-1)(Y^2+Y+1)}{(Y-1)(Y+1)}=\frac23\cdot\frac{Y^2+Y+1}{Y+1}$. The difference is $\frac{\frac23(Y^2+Y+1)-2Y}{Y+1}=\frac23\cdot\frac{Y^2-2Y+1}{Y+1}=\frac23\cdot\frac{(Y-1)^2}{Y+1}$, positive for $Y>1$. For $L=3$, $Y=20.085537$ and the difference is $11.516827=13.421975-1.905148$.

\textbf{Answer 15.15.} Mirror pair: $E=61.3283111+61.3283111=122.6566221$, charge $112+112=224$, $S_{\mathrm{tot}}=-21.7686201+21.7686201=0$ (the file gives $-5.3\times10^{-9}$). Krein image pair: the image has the $-M$ values with the opposite sign, so $E=61.3283111-61.3283111=0$ ($-7.0\times10^{-9}$ in the file), charge $112-112=0$, $S_{\mathrm{tot}}=-21.7686201-21.7686201=-43.5372402$. The Krein image pair cancels in energy, charge and pressures and doubles the scalar charge; the mirror pair doubles energy, charge and pressures and cancels the scalar charge. The two cancellations are of different kinds. The mirror pair's zero scalar charge compares two universes computed independently. The Krein image pair is formed as $(+M)-(-M)$ from the same two runs, so its zero energy only restates $E(+M)=E(-M)$, the content of T3; it is a zero of the $X+(-X)$ kind (Section 15.8), not a cancellation between two independent universes.

\textbf{Answer 15.16.} $P_B^\dagger=-i\gamma^8\gamma^0=i\gamma^0\gamma^8=P_B$ (both factors are real symmetric and they anticommute), and $P_B^2=-\gamma^0\gamma^8\gamma^0\gamma^8=\gamma^0\gamma^0\gamma^8\gamma^8=1$. $P_BCP_B=-\gamma^0\gamma^8C\gamma^0\gamma^8=-\gamma^0C\gamma^8\gamma^0\gamma^8=\gamma^0C\gamma^0\gamma^8\gamma^8=\gamma^0C\gamma^0=-C$, because $C$ commutes with $\gamma^8$ and anticommutes with $\gamma^0$. $P_BBP_B$: $\gamma^0$ commutes with $B$ ($\gamma^0B\gamma^0=B$, Section 15.5) and $\gamma^8$ anticommutes with it (F6), so $P_BBP_B=-\gamma^0\gamma^8B\gamma^0\gamma^8=\gamma^0\gamma^8B\gamma^8\gamma^0=-\gamma^0B\gamma^0=-B$. Hence $P_B(BC)P_B=(P_BBP_B)(P_BCP_B)=(-B)(-C)=BC$. For the classical field, $P_B$ reverses the kinetic term (through its factor $\gamma^8$, as in T1) and the scalar density $\bar\Psi\Psi$, whose matrix $C$ is $P_B$-odd; so $\mathcal L_{m,\lambda}\to\sqrt{|g|}\,[-K+mS-\tfrac\lambda2S^2]=-\mathcal L_{m,-\lambda}$: the same mass, the opposite coupling. In the Kohn--Sham model the scalar density has the matrix $BC$, which $P_B$ keeps, and $P_B$ keeps the number-density matrix $B^2=1$ as well (check \texttt{PAIR\_T2z\allowbreak{}2\_rulePa\allowbreak{}rities}); so the densities, the potentials and hence an even mass function with the same $\lambda$ are mapped onto themselves.

\section{Does the big bang create universes in pairs?}

\subsection{The question, and the answer in brief}

The request that this book answers asks, among other things, for a proof "that the 'big bang' creates universes in pairs". The author's notebook states the same idea as a hypothesis: at the time $x_4=0$ a universe of mass $+M$ and a universe of mass $-M$ are created together. This chapter examines that statement with everything the earlier chapters have built, and it says exactly how much of it the repository establishes.

The answer, in five sentences, each with its status word from Section 0.3:

\begin{enumerate}
\item \textbf{PROVED.} The field equations of dirac16complex (anticommuting components) and of dirac16complex00 (commuting components) have exact \textbf{pairing symmetries} (Chapter 15): the chirality map $\Psi\mapsto\gamma^8\Psi$ turns every solution with mass and coupling $(m,\lambda)$ into a solution with $(-m,-\lambda)$ (theorem T1); a reflection of a space-like frame direction turns every solution with $(m,\lambda)$ into one with $(-m,\lambda)$, and across the brane of the primordial field it produces a mirror universe of mass $-M$ (theorem T2); and the same two statements hold exactly for the Kohn--Sham ground and first excited states (theorem T3).
\item \textbf{PROVED (a consequence of T1), with a limit that is also PROVED.} A pair made of a universe $\Psi_+$ with $(m,\lambda)$ and its image $\Psi_-=\gamma^8\Psi_+$ with $(-m,-\lambda)$ carries zero total energy, momentum, stress and charge at every point of every gravitational field and at every time $x_4$. So the \textbf{totals} of such a pair are those of the empty state, and Einstein's equations with the pair as their only source are the equations without a source. This holds for the fields treated as classical fields, with the partner's coupling $-\lambda$ and the partner in exactly the configuration $\gamma^8\Psi_+$; the partner then has the opposite classical energy density, $-\rho_+$ (negative wherever $\rho_+>0$, as in Worked example 16.B). The limit: in the theory as built the two members do not interact, so the charge of each member, and in a gravitational field that does not change with time also the energy of each member, is conserved \textbf{on its own}. These separate conservation laws forbid the appearance of a cancelling pair out of the empty state unless each member has zero charge (and, in such a static field, zero energy), exactly as for the mirror pair of item 3 (Section 16.9). Only an interaction between the members, which the theory does not contain, could leave the totals as the only conserved quantities (OPEN). For the quantized field no reading examined in the repository gives two independent universes whose energies and charges cancel (Section 16.6): at the quantum level the cancellation is OPEN.
\item \textbf{PROVED.} A pair with zero total energy and momentum is invisible to gravity, so it can only accompany a gravitational field that needs no source. In eight-dimensional Einstein gravity the primordial field of the notebook needs a nonzero source (Chapter 9), so such a pair cannot be what produces it; with the Lovelock terms of orders 2 and 3 this is OPEN (Section 16.10). The mirror pair of T2, which can live on the two sides of the brane, has equal, not opposite, energies.
\item \textbf{OPEN (not derived).} No creation process, no creation rate, no probability amplitude, no wave function of the universe and no dynamical big bang is derived anywhere in the project.
\item \textbf{HYPOTHESIS.} "The big bang creates universes in pairs" therefore remains what the notebook calls it, a hypothesis. Stage 1 of the project put it in one sentence, which this chapter repeats: the pairing of the masses $\pm m$ is "a structural property of the equations, not a claim that universes of masses $\pm M$ are created in pairs" (Stage-1 document \texttt{provenan\allowbreak{}ce/DIRAC\allowbreak{}16COMPLE\allowbreak{}X\_ARBITR\allowbreak{}ARY\_FIEL\allowbreak{}D.md}, §1.3, item 6).
\end{enumerate}

\textbf{The word "consistent".} In this chapter "a process is consistent with a law" means "the law does not forbid it". It is a statement of logic, and it can be proved: to prove that energy conservation does not forbid a process, one shows that the energy before equals the energy after. It is much weaker than "the process happens". Section 16.4 shows with an everyday example of physics how large the difference is.

\textbf{Plan.} Section 16.2 quotes the hypothesis and reads it word by word. Section 16.3 separates three different questions hidden in it. Sections 16.4 and 16.5 show, with two examples from ordinary physics and from Chapter 11, what "creation" means when it is actually computed. Sections 16.6 and 16.7 state what is proved, with the key steps of the proofs (the complete proofs are in Chapter 15), and Section 16.8 what is computed. Section 16.9 derives what is consistent with the conservation laws and what the separate conservation laws of the two members forbid, and Section 16.10 confronts the zero total source of a cancelling pair with the nonzero source that the primordial field needs. Sections 16.11 and 16.12 list what is not derived and what a calculation of pair creation would require. Section 16.13 is the ledger of the chapter.

\textbf{Sources.} The exact Stage-5 theory consists of two Wolfram packages with their verifiers and two independent sympy checkers. Their reports and the collected exact statements are these files, with the programs that write them:

\begin{Verbatim}[fontsize=\small]
artifacts/dirac16complex/pair-creation/
  wolfram-pairing-report.json            141 of 141 checks true
    (scripts/verify_dirac16complex_pairing.wls, wolfram/Dirac16ComplexPairing.wl)
  python-pairing-report.json             172 of 172 checks true
    (scripts/check_dirac16complex_pairing.py)
  wolfram-dirac16complex00-report.json    46 of 46 checks true
    (scripts/verify_dirac16complex00.wls, wolfram/Dirac16Complex00.wl)
  python-dirac16complex00-report.json     49 of 49 checks true
    (scripts/check_dirac16complex00.py)
  pairing-theory.json                    the exact pairing statements
  dirac16complex00-theory.json           the exact statements on dirac16complex00
\end{Verbatim}

Check names beginning with \texttt{\detokenize{PAIR_}} belong to the Wolfram pairing report; the Python report repeats each family under names beginning with \texttt{\detokenize{S5_}}. The numerical pair runs are in the subfolders \texttt{\detokenize{rust/}} and \texttt{referenc\allowbreak{}e/} of the same folder. Further sources are the matter--antimatter analysis \texttt{provenan\allowbreak{}ce/DIRAC\allowbreak{}16COMPLE\allowbreak{}X\_MATTER\allowbreak{}\_ANTIMAT\allowbreak{}TER.md}, the Stage-1 and Stage-2 documents, the notebook survey \texttt{handoff/\allowbreak{}surveys/\allowbreak{}survey\_n\allowbreak{}otebook-\allowbreak{}physics.\allowbreak{}md}, and Chapters 8, 9, 11 and 13 of this book.

\textbf{State of the work when this chapter was written.} Stages 1 to 3 of the project are complete and verified by their gates. In Stage 4 the Rust Kohn--Sham solver of Chapter 13 is complete, its final cross-check against the independent reference solver has 63 checks of which 1 failed (\texttt{artifact\allowbreak{}s/dirac1\allowbreak{}6complex\allowbreak{}/kohn-sh\allowbreak{}am/pytho\allowbreak{}n-check-\allowbreak{}report.j\allowbreak{}son}, the check \texttt{canonica\allowbreak{}l\_eigenv\allowbreak{}alues}), and the Stage-4 gate has not been run; this chapter uses Stage-4 outputs in Sections 16.8 and 16.10, and they carry that status. The exact Stage-5 theory is complete in the sense that its four reports are committed and every check in them is true. The Stage-5 numerics are partial (Section 16.8), the Stage-5 documents have not been written, and no Stage-5 gate exists. The chapter quotes only committed files and says at each number where it comes from.

\subsection{The hypothesis in the author's own words}

The notebook is a working file, and its statements about pairs of universes are spread over a few text cells. The survey of the notebook quotes them (\texttt{handoff/\allowbreak{}surveys/\allowbreak{}survey\_n\allowbreak{}otebook-\allowbreak{}physics.\allowbreak{}md}, item 3 and the list of section titles). Cell 6 says:

"HYPOTHESIS: If, employing the Einstein eqs (or Einstein-Lovelock eqs), superluminal inflation/deflation exists, then at time x4 = 0 ... a pair of universes with MASSES ± M is created".

Cell 7, the title of a section, asks:

"Bigger Bang: Question: Are Universe (s) of masses ± M created in pairs at time x4 = 0 (before the particles of the standard model exist) ?"

Cell 17 records the task:

"M is the mass of the Ψ16 field"; "TODO: prove Universe(s) of masses ±M are created in pairs!"

Cell 23, in the middle of a derivation of the source of the metric, adds a hope: that the source tensor of the spinor field be $\Lambda g^{\mu\nu}$, "and H = some function of M, where Universe(s) of masses ± M created in pairs at time x4 = 0".

The project documents have responded to the TODO of cell 17 twice, in the same careful way, without claiming to have carried it out. Stage 1: the pairing of the masses by the chirality map "is a structural property of the equations, not a claim that universes of masses $\pm M$ are created in pairs" (Stage-1 document, §1.3, item 6). Stage 2 (the Stage-2 document, §16): the pairing "echoes the notebook's hypothesis", but it is "structural only", and "nothing about a creation process at $x_4=0$, about a pair of universes, or about a physical preference between the two signs follows from it".

\textbf{Reading the hypothesis word by word.} Each word needs a precise meaning before anything can be proved about it.

\begin{itemize}
\item \emph{"A universe of mass $M$."} Cell 17 says that $M$ is the mass of the field $\Psi_{16}$. A "universe of mass $M$" is therefore a configuration of the dirac16complex field (or of dirac16complex00) whose mass parameter in the Lagrangian is $M$, living in some gravitational field. It is \textbf{not} the total gravitational mass or energy of that universe; Section 16.6 shows that a universe with a negative mass parameter can have positive energy. The dictionary between the notebook and this book is $m=-HM$ (Sections 5.12 and 9.14): the notebook writes the mass term as $+HM\,\Psi^TC\Psi$, this book as $-m\,\bar\Psi\Psi$. So the notebook's universe of mass $+M$ has the book's mass $m=-HM$. The sign of the mass is a convention at this point, and Section 16.7 shows that it is even less than that: it is tied to the orientation of the frame.
\item \emph{"At time $x_4=0$."} $x_4$ is the time in which everything evolves (Section 0.6). The gravitational field is the notebook's metric of Chapter 9, which is \textbf{prescribed}, not computed (status ASSUMED); its free function $a_4(t)$, $t=Hx_4$, is not determined by any equation of the notebook or of the project.
\item \emph{"Before the particles of the standard model exist."} The theory contains no particles of the Standard Model at all; this part of the sentence is outside what the theory can address.
\item \emph{"Big bang."} In cosmology the big bang is the hot and dense early phase from which the present expansion of the universe started; in the classical solutions of Einstein's equations that describe it, the scale factor (the number that measures the size of every region of space, Section 9.4) goes to zero at a finite time in the past, and the equations break down there (a \textbf{singularity}). The notebook's section title speaks of a "Bigger Bang", and the notebook uses the moment $x_4=0$ for it. Section 16.11 shows that the notebook's metric has no singularity at $x_4=0$.
\item \emph{"Created in pairs."} This is the claim itself. Section 16.3 separates three things it can mean.
\item \emph{The premise "if ... superluminal inflation/deflation exists".} Chapter 9 has examined it for Einstein gravity: every member of the notebook's family of metrics needs a source with negative energy density, $\rho_{\mathrm{req}}=-3H^2(7+a_4'^2)/\kappa<0$ (Section 9.10, PROVED), and for linear $a_4$ a dirac16complex condensate that does not depend on $x_0$ is an exact source of that kind (Proposition 9.4, PROVED in the mean-field reading). So the premise can be met in Einstein gravity, but only with a negative energy density. With the Lovelock terms of orders 2 and 3, which neither the notebook nor the project computes, the question is OPEN.
\end{itemize}

\textbf{Status.} The hypothesis of cells 6, 7 and 17 is a HYPOTHESIS (the honesty ledger of Section 0.9 records it so). The hope of cell 23 that the source be $\Lambda g$ cannot be realized in eight-dimensional Einstein gravity (Section 9.10, PROVED), and no relation between $H$ and $M$ is derived anywhere (OPEN).

\subsection{Three different questions}

The sentence "the big bang creates universes in pairs" packs three questions of very different difficulty into one.

\textbf{Q1. Is it allowed?} Do the conservation laws of the theory (energy, momentum, charge) and the equations of gravity permit a transition from a state without universes to a state with a pair? This is a question about \textbf{constraints}: it asks whether something forbids the transition. It can be answered by comparing conserved quantities before and after.

\textbf{Q2. Does it happen in a classical solution?} Is there a solution of the coupled equations (the field equations of the two universes together with the gravitational field equations) that contains no universes before some moment and the pair after it? This is a question about \textbf{dynamics} in the classical theory.

\textbf{Q3. How likely is it?} In the quantum theory, what is the probability amplitude for the transition, and what is the probability, or the number of pairs created per unit time? This is a question about \textbf{quantum dynamics}; its answer is a number.

The repository answers Q1 for the fields treated as classical fields (Section 16.9), and the answer has two parts. The totals of a cancelling pair (energy, momentum, charge) are those of the empty state, so the conservation of the totals does not forbid its appearance. But in the theory as built the two members do not interact, so the charge of each member (and, in a gravitational field that does not change with time, the energy of each member) is conserved on its own, and these separate laws forbid the appearance of any pair, cancelling or mirror, whose members carry a nonzero charge (or, in such a field, a nonzero energy). At the quantum level Q1 is still open (Section 16.6). The repository does not answer Q2 or Q3 (Sections 16.11 and 16.12): both are OPEN.

\textbf{Symmetry is not realization.} A pairing theorem says: for every solution with mass $+M$ there is a solution with mass $-M$. It does not say that whenever the first is present, the second is present too. An everyday comparison makes the difference plain. The laws of mechanics do not change when everything is reflected in a mirror; so for every right hand that the laws allow, a left hand is allowed as well. That does not mean that every right hand comes with a left hand, still less that hands are created in right--left pairs. To go from "the partner is allowed" to "the partner is created together with the original", one needs an answer to Q2 or Q3.

\subsection{A lesson from ordinary physics: when can a photon make a pair?}

The difference between "allowed" and "happens" is best seen in a process that is completely understood: the creation of an electron and a positron (its antiparticle) from light. It also shows that the conservation laws can forbid a process that one might naively expect.

\textbf{Energy, momentum and mass.} We use ordinary space with three directions and one time, and units in which the speed of light is 1. A particle of mass $m$ with momentum $\mathbf p=(p_1,p_2,p_3)$ has the energy $E=\sqrt{m^2+|\mathbf p|^2}$, where $|\mathbf p|^2=p_1^2+p_2^2+p_3^2$. This is the relation $E^2=k_1^2+k_2^2+k_3^2+m^2$ of Section 2.4, which the Dirac equation was built to reproduce. A photon has $m=0$, so $E=|\mathbf p|$. For a collection of particles, the total energy $E_{\mathrm{tot}}$ and the total momentum $\mathbf p_{\mathrm{tot}}$ are the sums over the particles, and we define its \textbf{invariant mass} $\mu$ by

\[
\mu^2:=E_{\mathrm{tot}}^2-|\mathbf p_{\mathrm{tot}}|^2 .
\]

For a single particle $\mu=m$. The conservation of energy and momentum says that $E_{\mathrm{tot}}$ and $\mathbf p_{\mathrm{tot}}$ are the same before and after any process; hence $\mu^2$ is the same before and after as well. (These conservation laws, and the relation between $E$, $\mathbf p$ and $m$, are the standard physics of special relativity; we use them as given.)

\textbf{Lemma 16.1.} Two bodies with invariant masses $\mu_1,\mu_2\ge0$ (each a particle or a collection of particles, with energies $E_i=\sqrt{\mu_i^2+|\mathbf p_i|^2}$) have together an invariant mass of at least $\mu_1+\mu_2$:

\[
(E_1+E_2)^2-|\mathbf p_1+\mathbf p_2|^2\ \ge\ (\mu_1+\mu_2)^2 .
\]

\emph{Proof.} Write $a=|\mathbf p_1|$ and $b=|\mathbf p_2|$. Multiplying out,

\[
(E_1+E_2)^2-|\mathbf p_1+\mathbf p_2|^2=E_1^2-a^2+E_2^2-b^2+2\bigl(E_1E_2-\mathbf p_1\cdot\mathbf p_2\bigr)=\mu_1^2+\mu_2^2+2\bigl(E_1E_2-\mathbf p_1\cdot\mathbf p_2\bigr).
\]

The dot product satisfies $\mathbf p_1\cdot\mathbf p_2=ab\cos\theta\le ab$, where $\theta$ is the angle between the two vectors. So it is enough to show $E_1E_2\ge\mu_1\mu_2+ab$. Both sides are non-negative, so we may compare their squares:

\[
E_1^2E_2^2=(\mu_1^2+a^2)(\mu_2^2+b^2)=\mu_1^2\mu_2^2+\mu_1^2b^2+\mu_2^2a^2+a^2b^2,
\]

\[
(\mu_1\mu_2+ab)^2=\mu_1^2\mu_2^2+2\mu_1\mu_2ab+a^2b^2 .
\]

The difference is $\mu_1^2b^2+\mu_2^2a^2-2\mu_1\mu_2ab=(\mu_1b-\mu_2a)^2\ge0$. Hence $E_1E_2-\mathbf p_1\cdot\mathbf p_2\ge\mu_1\mu_2$, and the left-hand side is at least $\mu_1^2+\mu_2^2+2\mu_1\mu_2=(\mu_1+\mu_2)^2$. $\square$

\textbf{Proposition 16.2 (a single photon cannot make a pair).} A single photon in empty space cannot turn into an electron and a positron.

\emph{Proof.} Before: one photon, $\mu^2=E^2-|\mathbf p|^2=0$. After: two particles of mass $m>0$ each, so by Lemma 16.1 $\mu^2\ge(2m)^2>0$. Since $\mu^2$ is conserved, $0\ge4m^2$, which is false. $\square$

Charge does not forbid the process (the electron has charge $-1$ and the positron $+1$, total 0, the same as the photon's), and neither does energy on its own (the photon can have as much energy as we like). It is the combination of energy and momentum conservation that forbids it.

\textbf{Proposition 16.3 (with a nucleus it is allowed).} A photon of energy $E_\gamma$ that hits a nucleus of mass $M_N$ at rest can produce an electron--positron pair (the nucleus remaining in the final state) only if

\[
E_\gamma\ \ge\ 2m\Bigl(1+\frac{m}{M_N}\Bigr),
\]

and the conservation laws allow it for every energy above this threshold.

\emph{Proof.} Before: total energy $E_\gamma+M_N$, total momentum that of the photon, of size $E_\gamma$. So $\mu^2=(E_\gamma+M_N)^2-E_\gamma^2=M_N^2+2E_\gamma M_N$. After: the pair has, by Lemma 16.1, an invariant mass of at least $2m$, and applying the lemma once more to the pair and the nucleus, the final state has $\mu\ge2m+M_N$. Conservation requires $M_N^2+2E_\gamma M_N\ge(2m+M_N)^2=M_N^2+4mM_N+4m^2$, that is $E_\gamma\ge2m+2m^2/M_N$. Conversely, if the inequality holds, there are final momenta of the three bodies with exactly the initial total energy and momentum (Exercise 16.4 constructs them), so no conservation law is violated. $\square$

\textbf{Worked example.} In units of the electron mass, $m=1$. For a nucleus with $M_N=4$ the threshold is $2(1+\tfrac14)=2.5$. For $M_N=1000$ it is $2.002$, and for $M_N=1836$ (about the mass of a proton in units of the electron mass) it is $2.00109$. For a heavy nucleus the photon must bring the two rest energies, $2m$, and a tiny bit more: the nucleus takes up the momentum but almost no energy. With the standard value of the electron's rest energy, about 0.511 MeV, the threshold near a heavy nucleus is about 1.022 MeV.

\textbf{Two photons.} Two photons of energies $E_a$ and $E_b$ moving head-on have the momenta $E_a\mathbf n$ and $-E_b\mathbf n$ for a unit vector $\mathbf n$, so

\[
\mu^2=(E_a+E_b)^2-(E_a-E_b)^2=4E_aE_b,
\]

so the conservation laws require $4E_aE_b\ge4m^2$, that is $E_aE_b\ge m^2$; above this value they allow a pair (the construction of Exercise 16.4 works here as well). Here no third body is needed.

\textbf{What the example teaches.} Three lessons carry over to universes.

\begin{enumerate}
\item Conservation laws are \textbf{necessary} conditions. They can rule a process out (Proposition 16.2); they never rule a process in.
\item Whether an allowed process happens, and how often, is computed from the \textbf{dynamics}. For light near a nucleus this is quantum electrodynamics, the quantum theory of electrons and light, which gives the probability of pair creation as a \textbf{cross section}: an effective target area $\sigma$ attributed to each nucleus, such that a photon aimed at random into an area $A$ around one nucleus (an area much larger than $\sigma$) makes a pair with the probability $\sigma/A$. The calculation goes back to Bethe and Heitler (H. Bethe and W. Heitler, Proc. R. Soc. Lond. A 146, 83 (1934)), and pair creation by two photons was first computed by Breit and Wheeler (G. Breit and J. A. Wheeler, Phys. Rev. 46, 1087 (1934)). This book quotes these results; it does not derive them.
\item For a pair of universes of masses $\pm M$ related by T1, treated as classical fields, the \textbf{totals} of the pair (energy, momentum, charge) are zero, the same as for the empty state (Section 16.9). Here the analogy with light breaks down at two points. First, the electron and the positron are two quanta of one and the same field, the electron field, whose charge is the sum of their charges, $-1+1=0$; the two universes of a cancelling pair are two different fields, each with a charge of its own. Second, light and the electron field interact (the interaction of quantum electrodynamics is what turns light into matter), so the energy of the light alone is not conserved: the photon hands its energy to the pair, and only the total energy is conserved. The two universes do not interact, so the charge of each universe, and in a gravitational field that does not change with time the energy of each universe, is conserved on its own. These separate laws forbid the appearance of a pair out of the empty state unless each universe has zero charge (and, in such a field, zero energy). The analogue of lesson 2, a dynamics that produces the pair and a computed probability, does not exist in the project either (Sections 16.11 and 16.12).
\end{enumerate}

\textbf{Status.} Lemma 16.1 and Propositions 16.2 and 16.3 are PROVED here from the conservation of energy and momentum and the relation $E^2=m^2+|\mathbf p|^2$, which are ASSUMED as standard physics. The existence and size of the pair-creation cross sections are quoted from the literature, not computed in this book.

\subsection{What creation looks like when it is computed}

The project itself contains one complete calculation of pair creation: not of universes, but of particle--antiparticle pairs of the dirac16complex field in an expanding universe. It is experiment EXP-4b of Stage 3, described in Section 11.11. Recalling its structure shows what a calculation of the creation of universes would have to provide.

\textbf{The ingredients of EXP-4b.}

\begin{enumerate}
\item \textbf{A quantum field with a vacuum.} The field is quantized in the good sector (no momentum along the extra times, Section 8.10); its vacuum is the filled Dirac sea.
\item \textbf{A time-dependent background.} Three-dimensional space expands with a prescribed scale factor $a(t)$: de Sitter inflation glued to a radiation era at $t=0$.
\item \textbf{The dynamics.} Every momentum mode obeys its equation $i\,\dot u=h(t)\,u$ with a Hamiltonian $h(t)$ that changes with time because the momentum is redshifted by the expansion.
\item \textbf{The answer: a number.} A mode that starts in its negative-energy level ends with a weight $|\beta_k|^2$ in the positive-energy level; that is a created particle, and the hole it leaves is a created antiparticle. Summed over the modes, the number of created quanta per comoving volume is $na^3$.
\end{enumerate}

\textbf{The result} (COMPUTED; \texttt{artifact\allowbreak{}s/dirac1\allowbreak{}6complex\allowbreak{}/numeric\allowbreak{}s/exp4/s\allowbreak{}ummary.j\allowbreak{}son}, key \texttt{\detokenize{pair}}, list \texttt{\detokenize{masses}}, entry $m=1$, value \texttt{\detokenize{nA3}}): for $m=H_{\mathrm{inf}}$ the sudden end of inflation creates $na^3=4.412\times10^{-3}\,H_{\mathrm{inf}}^3$ quanta, and for $m=0$ none (Section 11.11 explains why: the massless equation does not notice an expansion of this type). Gravitational particle creation of this kind was discovered by Parker (L. Parker, Phys. Rev. Lett. 21, 562 (1968)).

\textbf{What the analogous calculation for universes would need.} A quantum theory in which a whole universe plays the role of a quantum; a "vacuum" that describes the absence of universes; a dynamics, either an interaction or a time-dependent background of some kind, that connects this state with states containing universes; and the resulting amplitude. None of these four ingredients exists in the notebook or in the project. Section 16.12 lists them in detail. The pairing theorems, to which we now turn, are statements of a different kind: they compare solutions, they do not produce them.

\subsection{What is proved (1): the chirality map and the cancelling pair}

\textbf{The setting.} Both fields of the book, dirac16complex (anticommuting components) and dirac16complex00 (commuting components), have the same Lagrangian density (Chapters 6 and 8)

\[
\mathcal L_{m,\lambda}=\sqrt{|g|}\,\Bigl[K-mS-\tfrac\lambda2S^2\Bigr],\qquad K=\tfrac12\bigl(\bar\Psi\gamma^\mu D_\mu\Psi-(D_\mu\bar\Psi)\gamma^\mu\Psi\bigr),\qquad S=\bar\Psi\Psi,\qquad \bar\Psi=\Psi^\dagger C,
\]

with $D_\mu\Psi=\partial_\mu\Psi+\Omega_\mu\Psi$, $D_\mu\bar\Psi=\partial_\mu\bar\Psi-\bar\Psi\Omega_\mu$, $\Omega_\mu=\tfrac12\omega_{\mu ab}S^{ab}$ and $\gamma^\mu=e_a{}^\mu\gamma^a$. The subscripts record the mass $m$ and the coupling $\lambda$ of the potential $U=\tfrac\lambda2S^2$. The field equation is $E_{m,\lambda}[\Psi]:=\gamma^\mu D_\mu\Psi-(m+\lambda S)\Psi=0$. The conserved Hermitian current is $J^\mu=-i\bar\Psi\gamma^\mu\Psi$, whose time component in Gaussian normal gauge is the charge density $J^4=\Psi^\dagger B\Psi$ with $B=-iC\gamma^4$ (Section 8.12). The energy--momentum tensor is (Chapter 7; Sections 9.12 and 13.14 use the same formula)

\[
T_{\mu\nu}=-\tfrac14\Bigl[\bar\Psi\gamma_\mu D_\nu\Psi+\bar\Psi\gamma_\nu D_\mu\Psi-(D_\mu\bar\Psi)\gamma_\nu\Psi-(D_\nu\bar\Psi)\gamma_\mu\Psi\Bigr]+g_{\mu\nu}\,\mathcal L_s,\qquad \mathcal L_s=\mathcal L/\sqrt{|g|},
\]

with $\gamma_\mu=g_{\mu\nu}\gamma^\nu$ and the energy density $\rho=T_{44}$.

\textbf{Four matrix facts} (Chapter 2, Sections 2.10 and 2.12). (i) $\gamma^8=\gamma^0\gamma^1\cdots\gamma^7=\mathrm{diag}(-I_8,I_8)$ is real and diagonal, hence Hermitian, and $(\gamma^8)^2=1$. (ii) $\gamma^8$ anticommutes with every $\gamma^a$, hence with every curved $\gamma^\mu=e_a{}^\mu\gamma^a$ and every $\gamma_\mu$ (the coefficients are ordinary functions). (iii) $\gamma^8$ commutes with $C$ and with every $S^{ab}$, hence with every $\Omega_\mu$ and with $D_\mu$. (iv) $\gamma^8B\gamma^8=-B$: indeed $\gamma^8B\gamma^8=-iC\gamma^8\gamma^4\gamma^8=-iC(-\gamma^4)(\gamma^8)^2=iC\gamma^4=-B$, using (iii) to move $\gamma^8$ past $C$ and (ii) to move it past $\gamma^4$.

\textbf{The bilinear rule.} Put $\Psi_-:=\gamma^8\Psi$. By (i), $\Psi_-^\dagger=\Psi^\dagger\gamma^8$, and by (iii), $\bar\Psi_-=\Psi^\dagger\gamma^8C=\Psi^\dagger C\gamma^8=\bar\Psi\gamma^8$. Also $D_\mu\Psi_-=\gamma^8D_\mu\Psi$ and $D_\mu\bar\Psi_-=(D_\mu\bar\Psi)\gamma^8$, because $\gamma^8$ is constant and commutes with $\Omega_\mu$. Therefore every bilinear $\bar\Psi X\Psi$ (with or without derivatives on either side) goes over into $\bar\Psi\gamma^8X\gamma^8\Psi$. If $X$ contains exactly one gamma matrix, possibly multiplied by spin generators $S^{ab}$ and ordinary functions, then (ii) and (iii) give $\gamma^8X\gamma^8=-X$; if $X=1$, then $\gamma^8X\gamma^8=1$. So

\[
K[\Psi_-]=-K[\Psi],\qquad S[\Psi_-]=S[\Psi],\qquad J^\mu[\Psi_-]=-J^\mu[\Psi],
\]

and every kinetic bilinear inside $T_{\mu\nu}$ changes sign.

\textbf{Theorem T1 (chirality map; PROVED).} For every gravitational field (every vielbein), every mass $m$, every coupling $\lambda$, every field configuration $\Psi$ (solution or not) and both statistics:

\[
\begin{aligned}
&\mathcal L_{m,\lambda}[\gamma^8\Psi]=-\mathcal L_{-m,-\lambda}[\Psi],\qquad E_{-m,-\lambda}[\gamma^8\Psi]=-\gamma^8E_{m,\lambda}[\Psi],\\
&T_{\mu\nu}[\gamma^8\Psi;-m,-\lambda]=-T_{\mu\nu}[\Psi;m,\lambda],
\end{aligned}
\]

and $J^\mu[\gamma^8\Psi]=-J^\mu[\Psi]$. In particular $\Psi$ solves the field equations with $(m,\lambda)$ if and only if $\gamma^8\Psi$ solves them with $(-m,-\lambda)$.

\emph{Proof.} Lagrangian: by the bilinear rule, $\mathcal L_{m,\lambda}[\gamma^8\Psi]=\sqrt{|g|}\,[-K-mS-\tfrac\lambda2S^2]=-\sqrt{|g|}\,[K-(-m)S-\tfrac{(-\lambda)}2S^2]=-\mathcal L_{-m,-\lambda}[\Psi]$. Field equation: $\gamma^\mu D_\mu(\gamma^8\Psi)=\gamma^\mu\gamma^8D_\mu\Psi=-\gamma^8\gamma^\mu D_\mu\Psi$ by (ii), and $S[\gamma^8\Psi]=S[\Psi]$, so

\[
E_{-m,-\lambda}[\gamma^8\Psi]=-\gamma^8\gamma^\mu D_\mu\Psi-(-m-\lambda S)\gamma^8\Psi=-\gamma^8\bigl[\gamma^\mu D_\mu\Psi-(m+\lambda S)\Psi\bigr]=-\gamma^8E_{m,\lambda}[\Psi].
\]

Since $\gamma^8$ is invertible, one side vanishes exactly when the other does. Energy--momentum tensor: the bracket consists of kinetic bilinears, which change sign, and $\mathcal L_s[\gamma^8\Psi;-m,-\lambda]=-\mathcal L_s[\Psi;m,\lambda]$ is the Lagrangian identity with $(m,\lambda)$ replaced by $(-m,-\lambda)$ and divided by $\sqrt{|g|}$. Current: bilinear rule. Statistics: the map is linear and keeps every $\Psi^\dagger$ to the left of every $\Psi$, so no two anticommuting factors are ever exchanged; the proof holds word for word for Grassmann components (in the notebook basis the map simply multiplies $\Psi_0,\dots,\Psi_7$ by $-1$). $\square$ (Chapter 6 proves the Lagrangian identity in the same way, and Section 15.3 gives every step of the whole theorem.)

\textbf{The coupling must change sign.} The potential $\tfrac\lambda2S^2$ is even in $S$ while the kinetic term is odd under the map. At a fixed coupling the identity fails: $\mathcal L_{m,\lambda}[\gamma^8\Psi]+\mathcal L_{-m,\lambda}[\Psi]=\sqrt{|g|}\,[-K-mS-\tfrac\lambda2S^2+K+mS-\tfrac\lambda2S^2]=-\lambda\sqrt{|g|}\,S^2\ne0$ (erratum E2 of \texttt{handoff/\allowbreak{}specs/CO\allowbreak{}NTRACT.m\allowbreak{}d}; check \texttt{PAIR\_T1g\allowbreak{}eneric\_n\allowbreak{}aiveFixe\allowbreak{}dLambdaF\allowbreak{}ails}). So the partner universe of T1 has mass $-m$ \textbf{and} coupling $-\lambda$; only for $\lambda=0$ are the two members the same free field with the masses $+m$ and $-m$.

\textbf{The same map, read as a reversed frame.} Replacing the vielbein $e$ by $-e$ leaves the metric $g_{\mu\nu}=e_\mu{}^a\eta_{ab}e_\nu{}^b$, $\sqrt{|g|}$ and the spin connection unchanged (in $\omega_\mu{}^a{}_b$ the vielbein and its inverse appear once each, so the two signs cancel) and reverses every $\gamma^\mu$. Hence $K\to-K$, $S\to S$, and $\mathcal L_{m,\lambda}[-e,\Psi]=-\mathcal L_{-m,-\lambda}[e,\Psi]$: the partner of T1 is also the original field described with the reversed frame (checks \texttt{PAIR\_T1j\allowbreak{}ets\_viel\allowbreak{}beinSign\allowbreak{}FlipIsT1} and \texttt{PAIR\_T1j\allowbreak{}ets\_gamm\allowbreak{}a8WithFr\allowbreak{}ameSignI\allowbreak{}sSymmetr\allowbreak{}y}). The signs of $m$ and $\lambda$, together with the overall sign of the Lagrangian, are therefore tied to a choice of frame; Section 16.7 meets a cleaner form of the same fact.

\textbf{How T1 is verified.} Stage 1 checked the Lagrangian identity with the exact matrices (check \texttt{ALG\_gamm\allowbreak{}a8Map} in both algebra reports of Stage 1). Stage 5 checked every identity of T1 as an exact polynomial identity in completely generic symbols, in a genuine Grassmann algebra with 288 odd generators, with exact jets of a general non-diagonal vielbein, and in the notebook's primordial field with an arbitrary function $a_4$ and fields that depend on all eight coordinates. The checks, all true in \texttt{wolfram-\allowbreak{}pairing-\allowbreak{}report.j\allowbreak{}son}, are

\begin{Verbatim}[fontsize=\small]
PAIR_T1generic_lagrangian       PAIR_T1generic_fieldEquation
PAIR_T1generic_conjugateFieldEquation  PAIR_T1generic_emtAll36
PAIR_T1generic_current          PAIR_T1grassmann_lagrangian
PAIR_T1grassmann_diracOperator  PAIR_T1grassmann_emt   PAIR_T1grassmann_current
PAIR_T1jets_lagrangian          PAIR_T1jets_emt        PAIR_T1jets_fieldEquations
PAIR_T1primordial_lagrangian    PAIR_T1primordial_diracOperator
PAIR_T1primordial_emt64         PAIR_T1primordial_current
\end{Verbatim}

and the independent Python checker repeats the four families as \texttt{S5\_T1gen\allowbreak{}eric}, \texttt{S5\_T1gra\allowbreak{}ssmann}, \texttt{S5\_T1jet\allowbreak{}s} and \texttt{S5\_T1pri\allowbreak{}mordial}. The theorem is about infinitely many fields, so the proof is the derivation above; the checks confirm it exactly.

\textbf{Worked example 16.A (T1 with the actual matrices).} Take the constant spinor $\Psi=e_0+e_4-i\,e_9$, where $e_k$ is the column with 1 in position $k$ (counting from 0) and 0 elsewhere. From the table of $C$ and $B/i$ in Section 2.9: $(Cu)_0=-u_4$, $(Cu)_4=-u_0$, $(Cu)_9=u_{13}$; $(Bu)_0=i\,u_9$, $(Bu)_4=-i\,u_{13}$, $(Bu)_9=-i\,u_0$. Only the components 0, 4 and 9 of $\Psi$ are nonzero, so only these rows are needed:

\[
\begin{aligned}
&(C\Psi)_0=-\Psi_4=-1,\qquad (C\Psi)_4=-\Psi_0=-1,\qquad (C\Psi)_9=\Psi_{13}=0,\\
&(B\Psi)_0=i\Psi_9=1,\qquad (B\Psi)_4=-i\Psi_{13}=0,\qquad (B\Psi)_9=-i\Psi_0=-i .
\end{aligned}
\]

Hence $S=\Psi^\dagger C\Psi=1\cdot(-1)+1\cdot(-1)+i\cdot0=-2$ and $J^4=\Psi^\dagger B\Psi=1\cdot1+1\cdot0+i\cdot(-i)=2$ (remember $\Psi_9^\ast=+i$). Now $\gamma^8=\mathrm{diag}(-I_8,I_8)$ reverses the components 0 to 7, so $\gamma^8\Psi=-e_0-e_4-i\,e_9$. Repeating: $(C\gamma^8\Psi)_0=1$, $(C\gamma^8\Psi)_4=1$, so $S[\gamma^8\Psi]=(-1)(1)+(-1)(1)=-2$; and $(B\gamma^8\Psi)_0=i(-i)=1$, $(B\gamma^8\Psi)_9=-i(-1)=i$, so $J^4[\gamma^8\Psi]=(-1)(1)+(i)(i)=-2$. The scalar density is unchanged and the charge density has changed sign, exactly as T1 says. (The same numbers come out of the exact fixture \texttt{artifact\allowbreak{}s/dirac1\allowbreak{}6complex\allowbreak{}/arbitra\allowbreak{}ry-field\allowbreak{}/algebra\allowbreak{}-fixture\allowbreak{}.json}.)

\textbf{Corollary 16.4 (the cancelling pair; PROVED).} Let $\Psi_+$ be a solution of the theory $\mathcal L_{m,\lambda}$, and let a second field $\Psi_-$ of the theory $\mathcal L_{-m,-\lambda}$, in the same gravitational field, be in the configuration $\Psi_-=\gamma^8\Psi_+$ (by T1 this is a solution). Then at every point and at every time $x_4$, in particular at $x_4=0$,

\[
T^{\mathrm{pair}}_{\mu\nu}:=T_{\mu\nu}[\Psi_+;m,\lambda]+T_{\mu\nu}[\Psi_-;-m,-\lambda]=0,\qquad J^\mu_++J^\mu_-=0,\qquad \mathcal L^{\mathrm{pair}}=0,\qquad S^{\mathrm{pair}}=2S[\Psi_+].
\]

The pair carries no net energy density, momentum density, stress or charge density anywhere, and hence no total energy, momentum or charge.

\emph{Proof.} Add the identities of T1 for $\Psi=\Psi_+$. $\square$ (Corollary 15.4 states the same result with its hypotheses listed one by one.)

The Stage-5 checks are \texttt{PAIR\_T1j\allowbreak{}ets\_pair\allowbreak{}EMTAndCu\allowbreak{}rrentVan\allowbreak{}ish} and \texttt{PAIR\_tot\allowbreak{}als\_fiel\allowbreak{}dLevelCh\allowbreak{}iralPair}. The matter--antimatter analysis proves the same statement independently, as its Corollary M4.1, with the checks

\begin{Verbatim}[fontsize=\small]
artifacts/dirac16complex/matter-antimatter/wolfram-matter-antimatter-report.json
  MA_M4_pairEMTAndCurrentG1_commuting   MA_M4_pairEMTG1_grassmann
\end{Verbatim}

\textbf{Worked example 16.B (a pair at rest in flat space).} In flat space ($g=\eta$, $\Omega_\mu=0$) with $\lambda=0$, take a constant column $u$ and $\Psi_+=u\,e^{-i\varepsilon x_4}$, a field that depends only on the time. We first compute its energy density for any such stationary solution, treating the components as commuting numbers (the classical field dirac16complex00). With $\gamma_4=\eta_{44}\gamma^4=-\gamma^4$ and $D_4=\partial_4$, the $44$ component of $T_{\mu\nu}$ is

\[
\rho=T_{44}=-\tfrac14\bigl[2\bar\Psi\gamma_4\partial_4\Psi-2(\partial_4\bar\Psi)\gamma_4\Psi\bigr]+\eta_{44}\mathcal L_s=\tfrac12\bigl[\bar\Psi\gamma^4\partial_4\Psi-(\partial_4\bar\Psi)\gamma^4\Psi\bigr]-\mathcal L_s .
\]

Now $\partial_4\Psi=-i\varepsilon\Psi$ and, since $\bar\Psi=u^\dagger C\,e^{+i\varepsilon x_4}$, $\partial_4\bar\Psi=+i\varepsilon\bar\Psi$. The bracket is $-i\varepsilon\bar\Psi\gamma^4\Psi-i\varepsilon\bar\Psi\gamma^4\Psi$, so the first term is $-i\varepsilon\,\Psi^\dagger C\gamma^4\Psi=\varepsilon\,\Psi^\dagger B\Psi$. On a solution, $\mathcal L_s=SU'(S)-U(S)$ (Chapter 7: the field equations turn $K$ into $(m+U')S$), which is 0 for $U=0$. Hence

\[
\rho=\varepsilon\,\Psi^\dagger B\Psi=\varepsilon\,u^\dagger Bu .
\]

Take $m>0$ and the rest state $u_+=\tfrac12(e_0-e_4-i\,e_9-i\,e_{13})$ of Section 8.8, which satisfies $-i\gamma^4u_+=u_+$, $Cu_+=u_+$ and $Bu_+=u_+$. The field equation $\gamma^4\partial_4\Psi_+=m\Psi_+$ reads $-i\varepsilon\gamma^4u_+=mu_+$, which holds with $\varepsilon=m$. So $\Psi_+=u_+e^{-imx_4}$ is a solution with mass $m$, and

\[
\rho_+=m\,u_+^\dagger Bu_+=m,\qquad J^4_+=u_+^\dagger Bu_+=1,\qquad S_+=u_+^\dagger Cu_+=1 .
\]

The partner is $\Psi_-=\gamma^8u_+e^{-imx_4}$ with $\gamma^8u_+=\tfrac12(-e_0+e_4-i\,e_9-i\,e_{13})=:w$. It solves the equation with mass $-m$: $-im\gamma^4w=-im\gamma^4\gamma^8u_+=im\gamma^8\gamma^4u_+=-m\gamma^8(-i\gamma^4u_+)=-m\,w$, which is $\gamma^4\partial_4\Psi_-=-m\Psi_-$; its frequency is again $\varepsilon=m$. With $\gamma^8B\gamma^8=-B$ and $\gamma^8C\gamma^8=C$,

\[
\rho_-=m\,w^\dagger Bw=m\,u_+^\dagger\gamma^8B\gamma^8u_+=-m,\qquad J^4_-=-1,\qquad S_-=u_+^\dagger\gamma^8C\gamma^8u_+=+1 .
\]

(Check by hand with the tables of Section 2.9: $(Bw)_0=iw_9=\tfrac12$, $(Bw)_4=-iw_{13}=-\tfrac12$, $(Bw)_9=-iw_0=\tfrac i2$, $(Bw)_{13}=iw_4=\tfrac i2$, and $w^\dagger Bw=-\tfrac14-\tfrac14-\tfrac14-\tfrac14=-1$.) Both energy densities agree with the formula $\rho=mS+U$ that Chapter 7 derives for fields that depend only on the time: $\rho_+=m\,S_+=m$ and $\rho_-=(-m)\,S_-=-m$. The pair has $\rho_++\rho_-=0$ and $J^4_++J^4_-=0$, while its scalar density is $2$: this is Corollary 16.4 in the smallest possible example. The cancellation works because the partner has the \textbf{opposite} classical energy density, here the negative value $-m$.

\textbf{The classical energy has no fixed sign anyway.} The other rest state of Section 8.8, $u_-=\tfrac12(e_0+e_4+i\,e_9-i\,e_{13})$, also has $-i\gamma^4u_-=u_-$, hence also frequency $m$ for mass $m$, but $u_-^\dagger Bu_-=-1$: a universe of mass $+m$ in the state $u_-e^{-imx_4}$ has the classical energy density $-m$. The classical energy of the commuting field is not bounded below and its charge density has no fixed sign; Stage 5 proves both statements exactly for dirac16complex00 (\texttt{artifact\allowbreak{}s/dirac1\allowbreak{}6complex\allowbreak{}/pair-cr\allowbreak{}eation/d\allowbreak{}irac16co\allowbreak{}mplex00-\allowbreak{}theory.j\allowbreak{}son}, key \texttt{\detokenize{fields}}, and the reports \texttt{wolfram-\allowbreak{}dirac16c\allowbreak{}omplex00\allowbreak{}-report.\allowbreak{}json} and \texttt{python-d\allowbreak{}irac16co\allowbreak{}mplex00-\allowbreak{}report.j\allowbreak{}son}). Exercise 16.7 even builds a single classical universe with zero energy density and zero charge density. So, for the classical field, zero total energy is not a special property of pairs.

\textbf{The quantum field: two readings of the partner.} For the quantized field dirac16complex, energies and charges are not the classical bilinears but the expectation values of Section 8.12: a quantum in the normalized positive-energy mode $u$ has $\langle\Psi^\dagger X\Psi\rangle=u^\dagger BXu$, with $X=Bh$ for the energy ($h$ the mode Hamiltonian), $X=B$ for the charge and $X=C$ for the scalar density. At rest, $h=-im\gamma^4$ for mass $m$ and $BC=-i\gamma^4$. The partner universe can be given a quantum structure in two ways. The Stage-5 exact theory computes both (\texttt{pairing-\allowbreak{}theory.j\allowbreak{}son}, key \texttt{\detokenize{T1krein}}), with the checks

\begin{Verbatim}[fontsize=\small]
PAIR_T1krein_imageAnticommutatorMinusB   PAIR_T1krein_imageOperatorIdentities
PAIR_T1krein_imageExpectationValues      PAIR_T1krein_independentCARPlusB
PAIR_T1krein_independentExpectationValues
\end{Verbatim}

\begin{itemize}
\item \textbf{The image field.} The partner is the operator $\Psi_-=\gamma^8\Psi_+$ acting on the same space of states. Its canonical anticommutator is $\gamma^8B\gamma^8=-B$ (fact (iv)): it carries the \textbf{reversed Krein metric}. This is the anticommutator that the Lagrangian $-\mathcal L_{-m,-\lambda}$, which by T1 equals $\mathcal L_{m,\lambda}[\gamma^8\,\cdot\,]$, produces (in Chapter 8 the anticommutator comes from the kinetic term, and an overall minus sign reverses it). Its expectation rule is $\langle X\rangle=w^\dagger(-B)Xw$ with $w=\gamma^8u$. Every operator of this partner is $\gamma^8$ times the corresponding operator of the $+M$ universe.
\item \textbf{The independent field.} The theory $\mathcal L_{-m,-\lambda}$ is quantized on its own, with the positive structure $+B$ of Chapter 8. Its modes are the $w=\gamma^8u$ (same energies $\varepsilon$), and its rule is $\langle X\rangle=w^\dagger BXw$.
\end{itemize}

For the rest quantum of Worked example 16.B the three descriptions give (energy $E$, charge $Q$, scalar density $S$):

\begingroup
\small
\begin{longtable}{@{}>{\raggedright\arraybackslash}p{0.200\linewidth}>{\raggedright\arraybackslash}p{0.200\linewidth}>{\raggedright\arraybackslash}p{0.200\linewidth}>{\raggedright\arraybackslash}p{0.200\linewidth}@{}}
\toprule
\textbf{universe} & \textbf{$E$} & \textbf{$Q$} & \textbf{$S$} \\
\midrule
\endfirsthead
\toprule
\textbf{universe} & \textbf{$E$} & \textbf{$Q$} & \textbf{$S$} \\
\midrule
\endhead
$+M$ universe, mode $u_+$, rule $u^\dagger BXu$ & $u_+^\dagger h(m)u_+=m$ & $u_+^\dagger u_+=1$ & $u_+^\dagger(-i\gamma^4)u_+=1$ \\
image partner, mode $w$, rule $w^\dagger(-B)Xw$ & $-w^\dagger h(-m)w=-m$ & $-w^\dagger w=-1$ & $-w^\dagger(-i\gamma^4)w=1$ \\
independent partner, mode $w$, rule $w^\dagger BXw$ & $w^\dagger h(-m)w=m$ & $w^\dagger w=1$ & $w^\dagger(-i\gamma^4)w=-1$ \\
\bottomrule
\end{longtable}
\endgroup

Each entry is one line of algebra: $B^2=1$; $h(-m)w=im\gamma^4w=m\,w$ because $-i\gamma^4w=-w$ (from $\gamma^4\gamma^8=-\gamma^8\gamma^4$); for the image energy $w^\dagger(-B)(Bh(-m))w=-w^\dagger h(-m)w$; and $w^\dagger(-i\gamma^4)w=-w^\dagger w=-1$. So the image partner has $(E,Q,S)\to(-E,-Q,S)$ and the independent partner $(E,Q,S)\to(E,Q,-S)$. (Section 15.8 explains the reason in general: the expectation-value rule multiplies every classical matrix by $B$, and $\gamma^8$ reverses $B$.) This is the general rule found by the exact theory (\texttt{pairing-\allowbreak{}theory.j\allowbreak{}son}, \texttt{\detokenize{T1krein}}, entries \texttt{imageFie\allowbreak{}ld} and \texttt{independ\allowbreak{}entQuant\allowbreak{}isation}) and by the matter--antimatter analysis (the matter--antimatter document, §7.2, with three further exact examples).

\textbf{What the two readings give (PROVED; the Fock-level part cited from Stage 5).} Two facts decide the question.

\begin{enumerate}
\item \emph{The image reading cancels only formally.} With the image rule the partner's $(E,Q,S)$ are $(-E,-Q,S)$, and at the level of operators Stage 5 proves ${:}T_{\mu\nu}[\gamma^8\Psi;-m,-\lambda]{:}=-{:}T_{\mu\nu}[\Psi;m,\lambda]{:}$ for the normal-ordered energy--momentum tensor. But the image field is built from the same operators, on the same states, as $\Psi_+$. The identities $Q_++Q[\Psi_-]=0$ and $H_++H[\Psi_-;-m]=0$ therefore have the form $X+(-X)=0$ and hold in every state; they are not a compensation by a second universe. Relative to its own canonical structure (the anticommutator $-B$, that is the Lagrangian $-\mathcal L_{-m,-\lambda}$), the image field's generator of time evolution is $-H[\Psi_-;-m]=+H_+$, its generator of phases is $-Q[\Psi_-]=+Q_+$, and its source of gravity, the energy--momentum tensor of its own Lagrangian (Section 5.8 defines it by varying the field's own action), is $-T_{\mu\nu}[\Psi_-;-m,-\lambda]=+T_{\mu\nu}[\Psi_+;m,\lambda]$. It is the same quantum system as the $+M$ universe, with the same energy and charge. The values $-|\varepsilon|$ and $-1$ per quantum in the table are expectation values of the formulas of $\mathcal L_{-m,-\lambda}$, not of that field's own Hamiltonian and charge. This observation is recorded in the committed matter--antimatter report \texttt{artifact\allowbreak{}s/dirac1\allowbreak{}6complex\allowbreak{}/matter-\allowbreak{}antimatt\allowbreak{}er/wolfr\allowbreak{}am-matte\allowbreak{}r-antima\allowbreak{}tter-rep\allowbreak{}ort.json} (measurement \texttt{M5\_impli\allowbreak{}cation}, entry \texttt{kreinLev\allowbreak{}elCaveat}, with the matrix identities of the check \texttt{MA\_M4\_kr\allowbreak{}einOnePa\allowbreak{}rticle}), and Chapter 17 discusses it as well.
\item \emph{The independent reading adds.} Quantized with its own positive structure, the $-M$ universe has positive energies, so the energies of the pair add; the charges cancel only under the additional assumption that the $-M$ universe is in a state of opposite charge; and for $\lambda\ne0$ its interaction energy has the opposite sign. This is the mirror pairing of Section 16.7.
\end{enumerate}

\textbf{Conclusion.} The cancellation of Corollary 16.4 is a statement about \textbf{classical} fields: two independent classical fields, with the Lagrangians $\mathcal L_{m,\lambda}$ and $\mathcal L_{-m,-\lambda}$, in the correlated configuration $\Psi_-=\gamma^8\Psi_+$, the second with the opposite classical energy density $-\rho_+$ (for dirac16complex00 this is the physical content; for dirac16complex it is a statement about its classical Grassmann-valued field). Of the two quantum readings examined in the repository, neither describes two independent, consistently quantized universes whose energies and charges cancel. Whether any quantum description of a physical universe of mass $-M$ realizes the classical cancellation is \textbf{OPEN}. (The matter--antimatter report also records the Fock-level statements taken from Stage 5 as provisional until the Stage-5 gate has passed; measurement \texttt{M4\_krein\allowbreak{}LevelSta\allowbreak{}tus}.)

\subsection{What is proved (2): mirror universes and Kohn--Sham pairs}

\textbf{Theorem T2 (mirror map; PROVED).} Let $u=\sum_av_a\gamma^a$ be a real unit vector of the Clifford algebra that is \textbf{space-like}, $n(u)=\eta(v,v)=+1$ (for example $u=\gamma^0$, $\gamma^1$, $\gamma^2$ or $\gamma^3$). Replace the field by $u\Psi$ and, at the same time, the vielbein by the reflected vielbein, in which the frame direction $v$ is reversed (the metric does not change). Then

\[
\mathcal L_{m,\lambda}[\text{reflected frame},\,u\Psi]=\mathcal L_{-m,\lambda}[\text{original frame},\,\Psi],
\]

the energy--momentum tensor is \textbf{unchanged}, $T\to+T$, the current is unchanged, $J\to+J$, and the scalar density changes sign, $S\to-S$. Solutions with $(m,\lambda)$ go over into solutions with $(-m,\lambda)$: the \textbf{same} coupling.

\emph{Key step.} The scalar density picks up the Pin character of Section 2.14: by Theorem 2.8 with one factor, $u^TCu=-n(u)\,C$, and since $u$ is real, $u^\dagger=u^T$; hence $S=\Psi^\dagger C\Psi\to\Psi^\dagger u^TCu\Psi=-n(u)\,S=-S$ for a space-like $u$. The kinetic term, on the other hand, keeps its sign, because the reflection of the frame supplies a second minus sign that cancels the one from the character (Chapter 6 shows this in flat space, where the frame reflection is the reflection of the coordinate $x_b$; Section 15.5 carries it out in every gravitational field; \texttt{pairing-\allowbreak{}theory.j\allowbreak{}son}, key \texttt{\detokenize{T2}}, entry \texttt{\detokenize{proof}}). With $K\to K$ and $S\to-S$, $K-mS-\tfrac\lambda2S^2\to K+mS-\tfrac\lambda2S^2$, which is the Lagrangian with mass $-m$ and the same $\lambda$. For a time-like $u$ the same construction gives a map of the T1 type instead (a table of all eight basis vectors and two general unit vectors is \texttt{T2frameT\allowbreak{}able} in \texttt{pairing-\allowbreak{}theory.j\allowbreak{}son}). Checks: \texttt{PAIR\_T2f\allowbreak{}rame\_twi\allowbreak{}stedSpac\allowbreak{}elikeMap\allowbreak{}sToMinus\allowbreak{}MSameLam\allowbreak{}bda}, \texttt{PAIR\_T2f\allowbreak{}rame\_sca\allowbreak{}larAndCu\allowbreak{}rrentSig\allowbreak{}ns}, \texttt{PAIR\_T2f\allowbreak{}rame\_cha\allowbreak{}racterIs\allowbreak{}MinusNor\allowbreak{}m}, \texttt{PAIR\_T2f\allowbreak{}rame\_onS\allowbreak{}hellImag\allowbreak{}e} and the Python family \texttt{S5\_T2fra\allowbreak{}me}.

\textbf{What T2 says about the sign of the mass.} The theory with mass $-m$ in one frame is the theory with mass $+m$ in the reflected frame, with the same energy--momentum tensor and the same charge. Only the scalar density, which itself changes sign under the reflection, distinguishes them. So the sign of the mass has meaning only relative to an orientation of the frame (the exact theory records this reading in \texttt{pairing-\allowbreak{}theory.j\allowbreak{}son}, key \texttt{\detokenize{T2z2}}, entry \texttt{interpre\allowbreak{}tation}). A "universe of mass $-M$" in the sense of T2 is the mirror image of a universe of mass $+M$.

\textbf{The mirror universe across the brane} (PROVED in Section 15.6; the checks are listed at the end of this paragraph). In the Z2 geometry of Section 9.16, the static primordial field on $y<0$ glued to its mirror image on $y>0$, the reflection $y\to-y$ is an \textbf{isometry}, that is, a map of spacetime to itself that leaves the metric unchanged (here because the metric depends on $y$ only through $|y|$), and it reflects the frame direction 0. Take $u=\gamma^0$ and define $\Psi'(y,x)=\gamma^0\Psi(-y,x)$, where $x$ stands for the other seven coordinates. Then $\Psi$ solves the field equation with $(m,\lambda)$ on $y<0$ exactly when $\Psi'$ solves it with $(-m,\lambda)$ on $y>0$; the energy density and all pressures of $\Psi'$ at $y$ equal those of $\Psi$ at $-y$; the charge density is the same on both sides; the scalar density is opposite. A configuration with $\Psi(-y)=\pm\gamma^0\Psi(y)$ obeys the field equation with one mass function on both sides exactly when that mass function is odd, $m(-y)=-m(y)$: the mirror universe carries the mass $-M$ (Section 13.5 and erratum E4.6 of \texttt{handoff/\allowbreak{}specs/ST\allowbreak{}AGE4\_SPE\allowbreak{}C.md}). The parity conditions of the Kohn--Sham problem of Chapter 13 are of this form, so \textbf{every Kohn--Sham state of Chapter 13 is already a pair of universes of masses $+M$ and $-M$ in this mirror sense} (\texttt{pairing-\allowbreak{}theory.j\allowbreak{}son}, key \texttt{\detokenize{T3}}, entry \texttt{stage4Z2\allowbreak{}Pair}). The checks of this paragraph are

\begin{Verbatim}[fontsize=\small]
PAIR_T2z2_diracOperator    PAIR_T2z2_emtPullback     PAIR_T2z2_currentPullback
PAIR_T2z2_massFunctionMap  PAIR_T2z2_symmetricIffOddMass
\end{Verbatim}

\textbf{Theorem T3 (Kohn--Sham level; PROVED in Section 15.7).} In the static primordial field, with the Kohn--Sham model of Chapter 13 or its counterpart for the commuting field in Chapter 14 (the proof covers both statistics), the $-M$ universe with the \textbf{same} coupling $\lambda$ and the tip condition $a(-L)=0$ instead of $b(-L)=0$ (in the bag family of Section 13.5, the angle $\theta=\pi$ instead of $\theta=0$) is the exact image of the $+M$ universe under the block map $\chi\to\sigma_y\chi$, $j\to-j$, which in the real variables of Section 13.14 is the swap $(a,b,j)\to(b,a,-j)$. Its spectra (with multiplicities), occupations, chemical potential, energy, free energy, entropy, Kohn--Sham gap, particle--hole list and Delta-SCF energy are identical to those of the $+M$ universe; its number density and all components of its energy--momentum tensor are identical; its scalar density and effective mass are reversed, $S_p\to-S_p$, $M_{\mathrm{eff}}\to-M_{\mathrm{eff}}$. The self-consistent iteration maps step by step, so the ground state defined by continuation in $\lambda$ and the first excited state map exactly. The image of the $+M$ state under T1 (coupling $-\lambda$, Krein metric $-B$) has the same orbitals and levels but reversed one-body densities and energies: $n\to-n$, $S\to S$, $E\to-E$, $T_{\mu\nu}\to-T_{\mu\nu}$. Checks: the families \texttt{PAIR\_T3b\allowbreak{}lock\_*} and \texttt{PAIR\_T3k\allowbreak{}s\_*} (for example \texttt{PAIR\_T3k\allowbreak{}s\_sigma2\allowbreak{}KSOperat\allowbreak{}orEquiva\allowbreak{}riant}, \texttt{PAIR\_T3k\allowbreak{}s\_sigma2\allowbreak{}Function\allowbreak{}alInvari\allowbreak{}ant}, \texttt{PAIR\_T3k\allowbreak{}s\_imageR\allowbreak{}uleEnerg\allowbreak{}yOdd}) and the Python families \texttt{S5\_T3blo\allowbreak{}ck} and \texttt{\detokenize{S5_T3ks}}.

\textbf{The boundary condition matters (PROVED; Section 15.9 derives the details).} If the $-M$ universe keeps the untransformed tip condition, the pairing fails. The exact theory shows it on the brane zero modes of Section 13.6: with the transformed condition the $-M$ zero mode has the same first-order band splitting $c(M)$ as the $+M$ one, while with the untransformed condition the constant becomes

\[
c_{\mathrm{ctrl}}=e^{-a_{4,0}}\,\frac{2M}{2M+H}\cdot\frac{e^{(2M+H)L}-1}{e^{2ML}-1},\qquad c(M)=e^{-a_{4,0}}\,\frac{2M}{2M-H}\cdot\frac{1-e^{-(2M-H)L}}{1-e^{-2ML}} .
\]

For $M=H=1$, $L=3$, $a_{4,0}=0$: $c(M)=2(1-e^{-3})/(1-e^{-6})=1.905148$ and $c_{\mathrm{ctrl}}=2(e^9-1)/\bigl(3(e^6-1)\bigr)=13.421975$ (\texttt{pairing-\allowbreak{}theory.j\allowbreak{}son}, key \texttt{\detokenize{T3}}, entry \texttt{untransf\allowbreak{}ormedBCC\allowbreak{}ontrol}; checks \texttt{PAIR\_T3k\allowbreak{}s\_zeroMo\allowbreak{}deUntran\allowbreak{}sformedC\allowbreak{}ontrol}, \texttt{PAIR\_T3k\allowbreak{}s\_contro\allowbreak{}lSplitti\allowbreak{}ngClosed\allowbreak{}Form}).

\textbf{The pair totals at the Kohn--Sham level} (PROVED; checks \texttt{PAIR\_tot\allowbreak{}als\_ksMi\allowbreak{}rrorPair} and \texttt{PAIR\_tot\allowbreak{}als\_ksKr\allowbreak{}einImage\allowbreak{}Pair}). For a $+M$ state with $N$ particles above the sea, energy $E_+$ and energy--momentum tensor $T_+$:

\begin{itemize}
\item the \textbf{mirror pair} ($+M$ universe and the ordinary $-M$ universe of T3): $E_{\mathrm{pair}}=2E_+$, charge $2N$, total scalar density 0, $T_{\mathrm{pair}}=2T_+$;
\item the \textbf{Krein image pair} ($+M$ universe and its T1 image with $-\lambda$ and $-B$): $E_{\mathrm{pair}}=0$, charge 0, $T_{\mathrm{pair}}=0$ at every point, total scalar density $2S_+$. The image is obtained from the ordinary $-M$ state by reversing the sign of every one-body density and energy (\texttt{pairing-\allowbreak{}theory.j\allowbreak{}son}, key \texttt{pairTota\allowbreak{}ls}), so this zero is of the $X+(-X)$ kind discussed in Section 16.6.
\end{itemize}

\textbf{Two kinds of pair.} The theorems therefore produce two different kinds of $\pm M$ pair, and it matters which one is meant:

\begingroup
\small
\begin{longtable}{@{}>{\raggedright\arraybackslash}p{0.267\linewidth}>{\raggedright\arraybackslash}p{0.267\linewidth}>{\raggedright\arraybackslash}p{0.267\linewidth}@{}}
\toprule
\textbf{property} & \textbf{cancelling pair (T1)} & \textbf{mirror pair (T2, T3)} \\
\midrule
\endfirsthead
\toprule
\textbf{property} & \textbf{cancelling pair (T1)} & \textbf{mirror pair (T2, T3)} \\
\midrule
\endhead
partner of $\Psi_+$ & $\gamma^8\Psi_+$ & $u\Psi_+$ with the reflected frame; $\gamma^0\Psi_+(-y)$ across the brane \\
parameters of the partner & $(-m,-\lambda)$ & $(-m,\lambda)$ \\
where the partner lives & at the same points of the same spacetime & reflected; on the other side of the brane \\
energy--momentum of the partner & $-T_+$ & $+T_+$ (reflected) \\
charge of the partner & $-Q_+$ & $+Q_+$ \\
scalar density of the partner & $+S_+$ & $-S_+$ \\
totals of the pair & $T=0$, $Q=0$ & $T=2T_+$, $Q=2Q_+$ \\
quantum reading & image field (metric $-B$): cancels only as $X+(-X)$ within one system; two cancelling universes OPEN & independent field (metric $+B$): energies add \\
\bottomrule
\end{longtable}
\endgroup

\subsection{What is computed}

The numerical side of Stage 5 (called T4 in \texttt{handoff/\allowbreak{}specs/ST\allowbreak{}AGE5\_SPE\allowbreak{}C.md}) solves the Kohn--Sham equations of the $+M$ universe, of the $-M$ universe with the transformed tip condition, and of a \textbf{control}: the $-M$ universe with the untransformed tip condition. When this chapter was written it was partial. What exists is the following; Section 15.10 gives the complete tables.

\textbf{The Rust runs (COMPUTED).} The Rust solver of Chapter 13 has a subcommand \texttt{\detokenize{pairs}}, whose outputs are committed under \texttt{artifact\allowbreak{}s/dirac1\allowbreak{}6complex\allowbreak{}/pair-cr\allowbreak{}eation/r\allowbreak{}ust/pair\allowbreak{}s/} for 25 configurations, all for dirac16complex. The 18 that were committed when this section was written are $m=1$, $L=3$, $N=8$ at $T=0$ with the five couplings $\hat\lambda\in\{0,\pm\hat\lambda_1,\pm\hat\lambda_2\}$; $m=1$, $L=3$, $N=112$ with the same five couplings at $T=0$ and at $T=0.1\,m$; and $m=3$, $L=3$, $N=8$ at $T=0$ with $\hat\lambda\in\{0,\pm\hat\lambda_1\}$. Each run writes a file \texttt{pairing.\allowbreak{}json}. In all 18 runs every level of the $+M$ universe is matched by a level of the $-M$ universe under the level map of T3, with no mismatch of multiplicities; the level energies agree to at most $6.6\times10^{-9}$, and the total energies to at most $1.5\times10^{-8}$ (key \texttt{pairingD\allowbreak{}eviation\allowbreak{}\_plusM\_v\allowbreak{}s\_minusM}). Every $+M$ run at $T=0$ reproduces the corresponding Stage-4 run of Chapter 13 bit for bit (key \texttt{stage4Re\allowbreak{}producti\allowbreak{}on}), and the Stage-4 outputs were left unchanged by the new subcommand (\texttt{artifact\allowbreak{}s/dirac1\allowbreak{}6complex\allowbreak{}/pair-cr\allowbreak{}eation/r\allowbreak{}ust/stag\allowbreak{}e4-ident\allowbreak{}ity-repo\allowbreak{}rt.json}, 9 of 9 checks true). The seven configurations committed later ($m=3$, $N=112$ with $\hat\lambda\in\{0,+\hat\lambda_1\}$ at $T=0$ and at $T=0.1\,m$ and with $-\hat\lambda_1$ at $T=0$; $m=3$, $N=8$ with $\pm\hat\lambda_2$ at $T=0$; Section 19.10 lists all 25) show the same pairing: every level is matched with no mismatch of multiplicities, the total energies agree to at most $1.6\times10^{-8}$, and the level energies to at most $1.1\times10^{-9}$, except in \texttt{d16c\_m3\_\allowbreak{}L3\_N112\_\allowbreak{}lamp1\_T0\allowbreak{}p1} ($1.8\times10^{-5}$, at $T=0.1\,m$). In the two runs with $\pm\hat\lambda_2$ the iteration for the excited state did not converge in either universe (field \texttt{excitedC\allowbreak{}onverged} under the key \texttt{universe\allowbreak{}s}), and no Stage-4 run with their parameters exists to compare with (key \texttt{stage4Re\allowbreak{}producti\allowbreak{}on}). Some pair totals (key \texttt{pairTota\allowbreak{}ls}; energies in units of $m$, $m=1$):

\begingroup
\small
\begin{longtable}{@{}>{\raggedright\arraybackslash}p{0.160\linewidth}>{\raggedright\arraybackslash}p{0.160\linewidth}>{\raggedright\arraybackslash}p{0.160\linewidth}>{\raggedright\arraybackslash}p{0.160\linewidth}>{\raggedright\arraybackslash}p{0.160\linewidth}@{}}
\toprule
\textbf{run} & \textbf{$E_+$} & \textbf{mirror pair $E$} & \textbf{mirror pair charge} & \textbf{Krein image pair $E$} \\
\midrule
\endfirsthead
\toprule
\textbf{run} & \textbf{$E_+$} & \textbf{mirror pair $E$} & \textbf{mirror pair charge} & \textbf{Krein image pair $E$} \\
\midrule
\endhead
\texttt{d16c\_m1\_\allowbreak{}L3\_N8\_la\allowbreak{}m0\_T0} & 0 & 0 & 16 & 0 \\
\texttt{d16c\_m1\_\allowbreak{}L3\_N8\_la\allowbreak{}mp1\_T0} & $-9.868332\times10^{-4}$ & $-1.973666\times10^{-3}$ & 16 & $1.3\times10^{-12}$ \\
\texttt{d16c\_m1\_\allowbreak{}L3\_N112\_\allowbreak{}lam0\_T0} & 61.090723287 & 122.181446574 & 224 & $-7.3\times10^{-11}$ \\
\texttt{d16c\_m1\_\allowbreak{}L3\_N112\_\allowbreak{}lamp2\_T0\allowbreak{}p1} & 69.806552563 & 139.613105141 & 224 & $-1.5\times10^{-8}$ \\
\bottomrule
\end{longtable}
\endgroup

The mirror-pair energy recorded in the file is the sum $E_++E_-$ of the two universes, each computed separately; it equals $2E_+$ up to the small difference between $E_-$ and $E_+$. For the last row, $2\times69.806552563=139.613105126$, which differs from the recorded sum by $1.5\times10^{-8}$. The mirror pair doubles the energy and the charge; the Krein image pair, computed from the $-M$ run by reversing signs (Section 16.7), has zero energy up to the numerical error (at most $1.5\times10^{-8}$ in absolute value over the 18 runs) and zero charge (at most $2.2\times10^{-12}$). The last row is at the temperature $T=0.1\,m$.

\textbf{The control (COMPUTED).} Without interaction the control converges and differs from the $+M$ universe, as the exact theory predicts: for $N=8$ its Kohn--Sham gap is $0.1008511$ instead of $0.4307337$ (it is the energy of a bound state that exists only with the untransformed condition, Section 15.9), and for $N=112$ its energy is $56.265496$ instead of $61.090723$ (\texttt{d16c\_m1\_\allowbreak{}L3\_N8\_la\allowbreak{}m0\_T0/pa\allowbreak{}iring.js\allowbreak{}on} and \texttt{d16c\_m1\_\allowbreak{}L3\_N112\_\allowbreak{}lam0\_T0/\allowbreak{}pairing.\allowbreak{}json}, key \texttt{universe\allowbreak{}s}, entry \texttt{minusM\_c\allowbreak{}ontrol}). The controls without interaction at $T=0.1\,m$ and for $m=3$ converge as well and also differ from the $+M$ universe (the files \texttt{d16c\_m1\_\allowbreak{}L3\_N112\_\allowbreak{}lam0\_T0p\allowbreak{}1/pairin\allowbreak{}g.json} and \texttt{d16c\_m3\_\allowbreak{}L3\_N8\_la\allowbreak{}m0\_T0/pa\allowbreak{}iring.js\allowbreak{}on}, key \texttt{controlO\allowbreak{}utcome}). With interaction the control was \textbf{not run}. Before the first self-consistent update the solver checks that the shift of the levels which that update can cause stays inside its energy window. With the untransformed tip condition the zero modes sit at the tip, where the proper densities are enhanced by $e^{6HL}$, and the bound on the shift exceeds the window in every interacting control (for example 40.34 against 4.59 in \texttt{d16c\_m1\_\allowbreak{}L3\_N8\_la\allowbreak{}mp1\_T0/p\allowbreak{}airing.j\allowbreak{}son}, key \texttt{universe\allowbreak{}s}, entry \texttt{minusM\_c\allowbreak{}ontrol}, field \texttt{\detokenize{error}}, which begins "not run: the first SCF update violates the window premise of the solver"). So the numerical failure of the pairing with the untransformed condition is shown only in the runs without interaction. That the transformed boundary condition is not a detail, and that without it there is no pairing, rests on the exact control of Section 16.7 (Section 15.9) together with these runs.

\textbf{What is not computed.} No Kohn--Sham pair runs of the commuting field dirac16complex00 were made with the Rust solver; the reference solver computed a few of them for $m=3$, $N=112$, and the pairing of that field with interaction is not demonstrated numerically (Chapter 14 lists what exists). The independent reference solver has computed only part of its pair matrix: its summary \texttt{artifact\allowbreak{}s/dirac1\allowbreak{}6complex\allowbreak{}/pair-cr\allowbreak{}eation/r\allowbreak{}eference\allowbreak{}/referen\allowbreak{}ce-pairs\allowbreak{}-summary\allowbreak{}.json} records \texttt{\detokenize{complete}} as false, with 131 runs still pending, and the runs it has finished ($m=3$, $N=112$) do not overlap with the 18 Rust configurations, so no pair run has yet been computed by both solvers. The checker that compares the two solvers, \texttt{scripts/\allowbreak{}check\_di\allowbreak{}rac16com\allowbreak{}plex\_pai\allowbreak{}rs.py}, exists, but no report of it is committed; the Stage-5 gate does not exist, and the Stage-5 documents have not been written. The numerical pairing is therefore COMPUTED by one solver, with its own checks; its independent cross-check is OPEN. The exact statements of Sections 16.6 and 16.7 do not depend on these numbers.

\subsection{What is consistent: the totals of a cancelling pair are those of the empty state}

\textbf{When a conservation law forbids a transition.} A \textbf{conserved quantity} is a number, computed from the state of the system, that does not change in time. It forbids a transition from a state A to a state B exactly when its values in A and B differ. A system can have several conserved quantities, and each of them forbids on its own: a transition is allowed by the conservation laws only if \textbf{every} conserved quantity has the same value in A and in B. In that case the conservation laws say nothing: they neither forbid the transition nor make it happen. Section 16.4 was an example: charge alone allows a photon to become an electron--positron pair; energy and momentum together forbid it.

\textbf{The empty state.} Classically, the empty state is $\Psi_+=\Psi_-=0$ everywhere; every bilinear of either field vanishes, so $T_{\mu\nu}=0$ and $J^\mu=0$ for each member separately, and every conserved quantity of each member is zero. In the quantum theory of the good sector it is the vacuum of Section 8.10, in which the normal-ordered energy, momentum and charge are zero.

\textbf{The conserved quantities of one field.} (a) The \textbf{charge} $Q=\int\sqrt{|g|}\,J^{x_4}\,d^7x$ over a slice $x_4=\text{const}$ does not depend on $x_4$, in every gravitational field and for both statistics, because the Lagrangian is invariant under $\Psi\to e^{i\alpha}\Psi$ (Section 5.7 for the principle, Section 8.12 for the current; Theorem M1 of the matter--antimatter document, §4.1, gives the complete proof with its machine checks). Among these checks, all true in the Wolfram matter--antimatter report named in Section 16.6, are the identities \texttt{MA\_M1\_no\allowbreak{}etherIde\allowbreak{}ntity\_co\allowbreak{}mmuting\_\allowbreak{}G1} and \texttt{MA\_M1\_no\allowbreak{}etherIde\allowbreak{}ntity\_gr\allowbreak{}assmann\_\allowbreak{}G1} and the conservation on solutions \texttt{MA\_M1\_on\allowbreak{}ShellCon\allowbreak{}servatio\allowbreak{}n\_G1}. (b) The \textbf{energy--momentum tensor} is conserved in the local sense, $\nabla_\mu T^\mu{}_\nu=0$ on every solution (Chapter 7); Einstein's equations require exactly this of their source (Section 9.9). (c) If the gravitational field does not depend on $x_4$, as the static primordial field does not, the Lagrangian density does not depend on $x_4$ explicitly, and Noether's theorem with the translation of Example 2 of Section 5.7 makes the total \textbf{energy} a conserved quantity; the same holds for the momentum along every coordinate on which nothing depends. We call a gravitational field that does not depend on $x_4$ \textbf{static}.

\textbf{The conserved quantities of a pair of fields that do not interact.} A pair consists of two fields, $\Psi_+$ with the Lagrangian $\mathcal L_{m,\lambda}$ and $\Psi_-$ with the Lagrangian $\mathcal L_{-m,-\lambda}$, and its action is the sum

\[
S_{\mathrm{pair}}=\int\Bigl(\mathcal L_{m,\lambda}[\Psi_+]+\mathcal L_{-m,-\lambda}[\Psi_-]\Bigr)\,d^8x .
\]

In the theory as built no term contains both fields (Section 16.11, item 1). Four steps follow.

\begin{enumerate}
\item \emph{Each member obeys its own equation.} The field equation of $\Psi_+$ is obtained by varying $\Psi_+$ alone (Section 5.3). The second term does not contain $\Psi_+$, so it contributes nothing, and the equation is $E_{m,\lambda}[\Psi_+]=0$, the equation of a single universe, whatever $\Psi_-$ does. In the same way $\Psi_-$ obeys $E_{-m,-\lambda}[\Psi_-]=0$.
\item \emph{Each charge is conserved on its own.} The proof of (a) (Theorem M1, item 4: $\partial_\mu(\sqrt{|g|}\,J^\mu)=0$ on every solution) uses nothing but the field equation of the field whose current it is. Applied to $\Psi_+$ it gives $\partial_\mu(\sqrt{|g|}\,J^\mu_+)=0$, so the charge $Q_+$ of $\Psi_+$ does not depend on $x_4$; applied to $\Psi_-$ it gives the same for $Q_-$. Noether's theorem says the same thing in another way: the rotation $\Psi_+\to e^{i\alpha}\Psi_+$ with $\Psi_-$ left alone does not change the first term of $S_{\mathrm{pair}}$ (by (a)) and does not touch the second, so it is a symmetry of its own, with its own conserved charge $Q_+$.
\item \emph{In a static field each energy is conserved on its own.} If nothing depends on $x_4$, the shift of $\Psi_+$ alone in $x_4$, with $\Psi_-$ left alone, is also a symmetry of $S_{\mathrm{pair}}$, so by (c) the energy $E_+$ of $\Psi_+$ is conserved on its own; the same holds for $E_-$ and for the momenta of each member along every coordinate on which nothing depends.
\item \emph{This does not depend on how the gravitational field is found.} Steps 1 and 2 hold in every gravitational field, so they hold whether the metric is prescribed (as in this book) or computed from Einstein's equations with the pair as source; step 3 holds in every static one.
\end{enumerate}

\textbf{Proposition 16.5 (what the conservation laws say about a cancelling pair; PROVED for classical fields).} Let $\Psi_+$ be a classical field of the theory $\mathcal L_{m,\lambda}$ and $\Psi_-$ a second, independent classical field of the theory $\mathcal L_{-m,-\lambda}$, in the correlated configuration $\Psi_-=\gamma^8\Psi_+$ (the cancelling pair of Corollary 16.4), in the theory as built, in which no term couples the two members. In any gravitational field and at any time $x_4$, in particular at $x_4=0$:

\begin{enumerate}
\item \emph{(the totals)} every density built from $T_{\mu\nu}$ or $J^\mu$ vanishes for the pair at every point, so the pair's total charge $Q_++Q_-$, and in a static field its total energy and its total momenta, are zero, the values of the empty state; the \textbf{conservation of the totals} therefore does not forbid a transition from the empty state to the pair;
\item \emph{(gravity)} the gravitational field equations with the pair as their only source, $G^\mu{}_\nu=\kappa\,T^{\mathrm{pair}\,\mu}{}_\nu$ (or the Einstein--Lovelock equations of Section 4.9 with the same right-hand side), are the equations \textbf{without any source}: whatever geometry is possible without the pair is also possible with it, and the pair changes nothing in it;
\item \emph{(the separate laws)} the charges $Q_+$ and $Q_-=-Q_+$ of the two members are conserved separately, and so, in a static field, are the energies $E_+$ and $E_-=-E_+$ and the momenta along every coordinate on which nothing depends. The empty state has $Q_+=0$ and $E_+=0$. Therefore a transition from the empty state to a cancelling pair is \textbf{forbidden unless $Q_+=0$}, and in a static field also unless $E_+=0$ and the momenta of $\Psi_+$ vanish: the same conditions as for the mirror pair at the end of this section.
\end{enumerate}

\emph{Proof.} Item 1: by Corollary 16.4 the densities of the two members cancel at every point, so the total charge, energy and momenta, which are integrals of these densities over a slice, are zero. Item 2: the right-hand side is $\kappa\cdot0=0$. Item 3: by steps 2 and 3 above, $Q_+$ and $Q_-$ (and in a static field $E_+$, $E_-$ and the momenta) are conserved quantities of the pair, each on its own; by Corollary 16.4, $Q_-=-Q_+$ and $E_-=-E_+$. The empty state has the value 0 for each of them. A conserved quantity whose values in the two states differ forbids the transition, so the pair can be reached from the empty state only if $Q_+=0$ (and then $Q_-=0$ as well), and, in a static field, only if also $E_+=0$ and the momenta of $\Psi_+$ are zero. $\square$

\textbf{Worked example 16.C.} Take the pair of Worked example 16.B in a periodic box of 7-volume $V$ (the slice coordinates are identified periodically, so every integral over a slice is finite; this is one of the cases of Theorem M1, item 4). In flat space $\sqrt{|g|}=1$, and $J^4_+=1$, $\rho_+=m$ at every point, so $Q_+=V$ and $E_+=mV$, while $Q_-=-V$ and $E_-=-mV$. The totals are $Q_++Q_-=0$ and $E_++E_-=0$, the values of the empty state, as item 1 says. But $Q_+=V\ne0$ and, since flat space is static, $E_+=mV\ne0$: by item 3 this pair cannot appear from the empty state, although its totals are zero. The conditions of item 3 on the charge and the energy can be met: Exercise 16.7 builds a configuration of the commuting field whose charge density and energy density vanish at every point, so that $Q_+=0$ and $E_+=0$ for it and for its partner. Item 3 does not forbid a pair of that kind (provided also its momenta vanish); but Proposition 16.7 below shows that no classical field of this kind can grow out of zero fields at all, whatever the values of its conserved quantities.

\textbf{What the phrase "consistent with the conservation laws" means here.} The Stage-5 plan of the project (\texttt{handoff/\allowbreak{}specs/ST\allowbreak{}AGE5\_SPE\allowbreak{}C.md}, the corollary of T1) says that a pair of this type "carries no net energy-momentum and no net charge, so its creation from the field-free state is consistent with every conservation law and constraint". Proposition 16.5 gives the exact extent of this sentence. It is true of the totals of the pair and of the gravitational equations (items 1 and 2). It is not true of the separate conservation laws of the two members (item 3), which in the theory as built are exact and forbid the appearance of a pair with $Q_+\ne0$. Two conditions are built into items 1 and 2: the partner has the coupling $-\lambda$, and its classical energy density is the opposite, $-\rho_+$, of that of $\Psi_+$ (negative wherever $\rho_+>0$, as in Worked example 16.B).

\textbf{When would only the totals be conserved? (HYPOTHESIS; OPEN).} A cancelling pair with $Q_+\ne0$ would be allowed by the conservation laws only if the separate laws of item 3 were broken while the totals stayed conserved. That needs an \textbf{interaction between the two members}: a term of the action that contains both fields and that is unchanged by the common rotation $(\Psi_+,\Psi_-)\to(e^{i\alpha}\Psi_+,e^{i\alpha}\Psi_-)$ but not by the separate rotations. To see what such a term looks like, take a term containing the product $\bar\Psi_+\Psi_-=\Psi_+^\dagger C\Psi_-$ (together with its complex conjugate). Under $\Psi_+\to e^{i\alpha}\Psi_+$ and $\Psi_-\to e^{i\beta}\Psi_-$ one has $\Psi_+^\dagger\to e^{-i\alpha}\Psi_+^\dagger$, so $\bar\Psi_+\Psi_-\to e^{i(\beta-\alpha)}\,\bar\Psi_+\Psi_-$, which is unchanged exactly when $\beta=\alpha$. With such a term only the charge of the common rotation, $Q_++Q_-$, would be conserved, and in a static field only the total energy. Such a term would also change the field equations and add an energy of its own, so the configuration of Corollary 16.4 would in general no longer be a solution, and even the zero of the totals would have to be examined again. This is an illustration of the kind of term that would be needed, not a proposal: the theory as built contains no such term, and whether one could be added consistently (with the symmetries of the theory, and without making the empty state unstable, Section 16.12, item 5), and whether it would then create cancelling pairs, has not been studied (OPEN). Even then, for classical fields, a pair could not grow out of zero fields (Proposition 16.7), because such a term, like every term of the field equations, vanishes when both fields vanish.

\textbf{At the quantum level it is OPEN.} For the quantized field the corresponding statement would need a quantum description of two independent universes whose energies and charges cancel. Section 16.6 showed that neither reading examined in the repository provides one: the image reading cancels only as $X+(-X)=0$ within a single quantum system, and in the independent reading the energies add. Whether the quantum theory admits a cancelling pair of independent universes at all is OPEN. In both readings the charge of the $+M$ universe is still conserved on its own (in the good sector it is the operator $Q=\sum(b^\ast b-d^\ast d)$ of Section 8.12, and Theorem M1 holds for both statistics), so the limit of item 3 does not disappear in the quantum theory.

\textbf{Three remarks.}

\begin{enumerate}
\item \emph{The idea is old.} That a universe with zero total energy could appear "from nothing" without violating energy conservation was proposed by Tryon (E. P. Tryon, Nature 246, 396 (1973)), and the phrase "creation of universes from nothing" is the title of a paper of Vilenkin (A. Vilenkin, Phys. Lett. B 117, 25 (1982)). What T1 adds is an exact, verified mechanism for a zero of the totals: an explicit partner whose energy--momentum tensor is minus that of the original at every point. In the theory as built this zero is not enough, because the charge of each member, and in a static field its energy, is conserved on its own (item 3). Zero total energy has never, by itself, produced a rate: that is always a separate calculation.
\item \emph{Consistency is not existence.} Section 16.4 showed conservation laws allowing a process whose probability must still be computed. Items 1 and 2 of Proposition 16.5 are statements of that kind; item 3 shows that for a pair with $Q_+\ne0$ not even the allowing holds in the theory as built.
\item \emph{The partner must be exact.} The cancellation holds only for $\Psi_-=\gamma^8\Psi_+$ at every point, with the coupling $-\lambda$. Any other configuration of mass $-M$ leaves a remainder. A creation process, if one existed, would have to produce this exact correlation between the two members.
\end{enumerate}

\textbf{The mirror pair is not consistent with the conservation laws unless it is neutral (PROVED).} For the mirror pair of T2 and T3 the partner has the same charge as the original (Section 16.7), so the pair has the total charge $2Q_+$. The charge is exactly conserved and the empty state has charge 0; therefore a mirror pair can appear from the empty state only if $Q_+=0$. In a static field the same argument with the energy requires $E_+=0$. For the mirror pair already the totals forbid this; the separate laws of item 3 give the same conditions, and they are the same conditions as for the cancelling pair. Every Kohn--Sham state of Chapter 13 has a nonzero charge, its particle number $N$ (8, 112, 896 or 1016 in the committed runs), so no computed Kohn--Sham pair, mirror or Krein image, could appear from the empty state: the free state with $N=8$ has $E_+=0$, but its mirror pair has the charge 16 (Section 16.8), and its $+M$ member alone has the charge 8. An independently quantized partner in a state of opposite charge (made of antiparticles) would cancel the total charge, but that is an additional assumption about its state, and its energy would still add (the matter--antimatter document, §7.3, which also concludes, in agreement with the caveat recorded in the matter--antimatter report, measurement \texttt{M5\_impli\allowbreak{}cation}, entry \texttt{kreinLev\allowbreak{}elCaveat}, and with Section 16.6, that for the image-field pair the cancellation is only an identity $X+(-X)=0$ within one quantum system, so that the cancellation of a pair is a statement about classical fields only). In a gravitational field that changes with time, the energy is not conserved (the expanding universes of Chapter 11 create particle pairs of positive energy, Section 11.11), and only the charges constrain a mirror pair or a cancelling pair.

\subsection{Zero total source against the source that the primordial field needs}

A cancelling pair is invisible to gravity (Proposition 16.5, item 2). The notebook, however, wants the pair to be created in its own primordial gravitational field. The two statements have to be put side by side.

\textbf{Proposition 16.6 (a cancelling pair cannot produce the primordial field; PROVED).} Suppose that a T1 pair is the only matter and that gravity obeys the eight-dimensional Einstein equations, with or without a cosmological constant $\Lambda$. Then no member of the notebook's primordial family of metrics is a solution, whatever the function $a_4$.

\emph{Proof.} By Corollary 16.4 the equations become $G^\mu{}_\nu+\Lambda\,\delta^\mu{}_\nu=0$ (Section 4.8 writes $G_{\mu\nu}+\Lambda g_{\mu\nu}=\kappa T_{\mu\nu}$; raising one index gives this mixed form). For $\Lambda=0$ this fails because $G^4{}_4=3H^2(7+a_4'^2)\ge21H^2>0$ (Section 9.8). For $\Lambda\ne0$ all eight diagonal components of $G^\mu{}_\nu$ would have to be equal to $-\Lambda$; already $G^0{}_0=G^4{}_4$ means $-3H^2(a_4'^2-5)=3H^2(7+a_4'^2)$, that is $15-3a_4'^2=21+3a_4'^2$, that is $a_4'^2=-1$, impossible for a real function. For the static member ($a_4'=a_4''=0$), $G^\mu{}_\nu=\mathrm{diag}(15,15,15,15,21,15,15,15)\,H^2$ (Section 9.15), and $15H^2+\Lambda=0$ and $21H^2+\Lambda=0$ cannot both hold. $\square$

So a pair with zero total energy--momentum is compatible only with a gravitational field that needs no source, while the primordial field of the notebook needs one: in Einstein gravity $\rho_{\mathrm{req}}=-3H^2(7+a_4'^2)/\kappa<0$, and for the static member $\rho_{\mathrm{req}}=-21H^2/\kappa$, $p_{\mathrm{req}}=+15H^2/\kappa$ (Sections 9.10 and 9.15). There are two ways to read the notebook's picture, and each must be stated honestly.

\textbf{Reading A: the cancelling pair lives in a geometry that something else produces.} The pair then adds nothing to the source, and the geometry is whatever the other matter makes it. For the static primordial field, the candidates for that other matter are:

\begin{itemize}
\item an $x_0$-independent condensate of dirac16complex, which is an exact source if and only if its three-gamma bilinears vanish, $mS=-36H^2/\kappa$ and $\lambda S^2=30H^2/\kappa$ (Proposition 9.4 and Section 9.15; PROVED in the mean-field reading); its energy density is negative, and whether such a state is realized or stable is OPEN;
\item the Kohn--Sham states of Chapter 13, which do \textbf{not} qualify: every state with band or bulk levels has a positive average energy density, for example $\langle\rho\rangle=0.0230891$ for $N=112$ without interaction, which would need the negative coupling $\kappa=-909.5$ (COMPUTED; \texttt{artifact\allowbreak{}s/dirac1\allowbreak{}6complex\allowbreak{}/kohn-sh\allowbreak{}am/rust/\allowbreak{}emt/emt-\allowbreak{}summary.\allowbreak{}csv}; Section 13.14), and the sourcing conditions are met in no computed run;
\item the Lovelock terms of orders 2 and 3: if some choice of the notebook's constants $w_1,w_2,w_3,\Lambda$ made the field a solution of the \textbf{vacuum} Einstein--Lovelock equations, which is what the notebook's cell 14 writes down, then a cancelling pair would fit it exactly. Nobody has computed these terms for this field (Section 4.9), so this is OPEN.
\end{itemize}

In reading A the pair adds nothing to the source, so it plays no role in producing the universe it lives in; and by Proposition 16.5, item 3, it could appear from the empty state only if each member had zero charge (and, in the static field, zero energy).

\textbf{Reading B: the two universes are the two halves of the Z2 geometry.} This is the reading that Stage 4 built as a model of the notebook's picture (Section 9.16, status ASSUMED). The partner is the mirror universe of T2 on the other side of the brane. Nothing cancels: the energy density of the $-M$ half at $y$ equals that of the $+M$ half at $-y$, the charges are equal, and each half needs the negative bulk source $\rho_{\mathrm{req}}=-21H^2/\kappa$, while the brane needs the positive surface energy density $12H/\kappa$ and the surface pressure $-10H/\kappa$ (Section 9.16, PROVED as requirements of the geometry). No computed state supplies the bulk source (Chapter 13), and the matter on the brane is not derived from anything (OPEN).

\begingroup
\small
\begin{longtable}{@{}>{\raggedright\arraybackslash}p{0.267\linewidth}>{\raggedright\arraybackslash}p{0.267\linewidth}>{\raggedright\arraybackslash}p{0.267\linewidth}@{}}
\toprule
\textbf{property} & \textbf{reading A (cancelling pair)} & \textbf{reading B (mirror pair across the brane)} \\
\midrule
\endfirsthead
\toprule
\textbf{property} & \textbf{reading A (cancelling pair)} & \textbf{reading B (mirror pair across the brane)} \\
\midrule
\endhead
pairing theorem & T1 (with $-\lambda$; classical fields) & T2 and T3 (same $\lambda$) \\
total energy--momentum of the pair & 0 at every point & twice that of one universe \\
total charge of the pair & 0 & twice that of one universe \\
appearance from the empty state & not forbidden by the totals, but forbidden by the separate conservation of each member's charge unless each universe is neutral (and, in a static field, of zero energy); allowed for a charged pair only with a coupling of the members that the theory lacks (OPEN); quantum level OPEN & forbidden unless each universe is neutral (and, in a static field, of zero energy); here already the totals forbid it \\
what the gravitational field needs & a separate source, or a vacuum Einstein--Lovelock solution (OPEN) & $\rho_{\mathrm{req}}<0$ in each half and a brane stress, not supplied by any computed state \\
status of the picture & totals consistent with the empty state, appearance of a charged pair forbidden in the theory as built, nothing derived & a chosen model (ASSUMED), its sources OPEN \\
\bottomrule
\end{longtable}
\endgroup

In neither reading is the notebook's field produced by a pair of universes. In reading A the totals of the pair are those of the empty state, but the pair does not act on the geometry, and in the theory as built the separate conservation laws of its two members forbid its appearance unless each member is neutral. In reading B the pair is the mirror structure of the geometry itself, and that structure needs sources that the theory has not supplied.

\subsection{What is not derived}

The following items are not derived anywhere in the notebook or the project. Each is OPEN unless stated otherwise.

\begin{enumerate}
\item \textbf{No creation process.} The field equations of the $+M$ and $-M$ members are separate equations; no term of the theory couples them, and gravity, in reading A, sees only their sum, which is zero. Because nothing couples them, the charge of each member, and in a static field its energy, is conserved on its own (Proposition 16.5, item 3). There is no equation whose solution describes a transition from a state without universes to a state with a pair. For classical fields such a transition is even excluded within every setting in which the book can decide it (Proposition 16.7 below).
\item \textbf{No rate, probability or amplitude.} Section 16.12 lists what their calculation would need.
\item \textbf{No wave function of the universe.} The notebook's file name speaks of a "wave function of the universe", and a section title of the notebook (cell 308, according to the survey) introduces "the wave function, Ψ16, for this Universe": the notebook uses the phrase for the classical field $\Psi_{16}$ itself. In quantum cosmology the phrase means something else: a quantum state of the geometry together with the matter fields, as in the proposal of Hartle and Hawking (J. B. Hartle and S. W. Hawking, Phys. Rev. D 28, 2960 (1983)). According to its survey, the notebook contains no field quantization (only classical fields and single-particle wave functions), and the project quantizes only the matter field in a prescribed geometry (Chapter 8). A wave function of the universe in the quantum-cosmology sense exists in neither.
\item \textbf{No dynamical big bang.} The metric is prescribed (ASSUMED), its function $a_4$ is free, and the moment $x_4=0$ is regular (Proposition 16.8 below). Nothing in the field marks a beginning.
\item \textbf{No selection of pairs.} T1 to T3 hold for every solution. They say that every universe of mass $+M$ has an allowed partner of mass $-M$; they do not say that the partner is present (Section 16.3).
\item \textbf{Nothing singles out $x_4=0$.} T1 holds at every time $x_4$. In the static field, used by Chapters 13 and 15, nothing depends on $x_4$ at all; for the notebook's choice $a_4=t$, the moment $x_4=0$ is only the moment at which $a_4=0$.
\item \textbf{No value of $M$, and no relation between $H$ and $M$.} The mass is a free parameter, and the relation "H = some function of M" hoped for in cell 23 is not derived.
\item \textbf{No solution of the matter--antimatter problem.} The pairing theorems do not create charge. The charge of each universe is exactly conserved in every gravitational field (Theorem M1), and each field has an exact symmetry that reverses the charge, charge conjugation for dirac16complex00 (in every gravitational field) and a CP-type symmetry for dirac16complex (in flat space and in gravitational fields in which the reflection of a space-like coordinate is an isometry) (Theorem M2), so two of Sakharov's three conditions for producing an excess of matter fail; the theory contains no baryons of the Standard Model, and no departure from equilibrium is computed (PROVED and stated in Chapter 17, which follows the matter--antimatter analysis \texttt{provenan\allowbreak{}ce/DIRAC\allowbreak{}16COMPLE\allowbreak{}X\_MATTER\allowbreak{}\_ANTIMAT\allowbreak{}TER.md}). A cancelling pair with opposite charges $Q_+=-Q_-\ne0$ has zero total charge, but since each charge is conserved on its own, such a pair cannot arise from a neutral state in the theory as built: its charges can only be assumed as initial data (the matter--antimatter document, §8.3, item 1). That it explains the observed excess of matter is not established.
\end{enumerate}

\textbf{Proposition 16.7 (a classical field cannot appear from nothing; PROVED in the setting stated).} Consider a classical field of the good sector described by finitely many mode amplitudes, collected in a column $u(x_4)$, that obeys

\[
i\,\frac{du}{dx_4}=F(u,x_4),
\]

where $F$ is a smooth function of $u$ with $F(0,x_4)=0$ for every $x_4$. This covers one mode in flat space, $F=h_ku$ with the constant matrix of Section 8.9; the modes in the expanding universes of Chapter 11, $F=h(x_4)u$ with a matrix that changes with time; and any finite set of coupled modes with the interaction $\lambda S\Psi$, which is of third order in $u$. If $u=0$ at one time, then $u=0$ at all times.

\emph{Proof.} The constant function $u(x_4)=0$ is a solution, because $F(0,x_4)=0$. Let $u$ be any solution, defined on an interval $I$ of times, with $u(t_0)=0$ at some time $t_0$ in $I$. The theorem of Picard and Lindelöf (Section 10.2, quoted there) says only that near each initial time there is exactly one solution with a given initial value; we have to carry this over to the whole interval. Suppose that $u(t_2)\ne0$ at some time $t_2>t_0$ in $I$. Consider the times $s$ between $t_0$ and $t_2$ such that $u=0$ on the whole interval from $t_0$ to $s$, and let $s^\ast$ be their least upper bound (the smallest number that none of them exceeds). Since $u$ is continuous, $u(s^\ast)=0$, so $s^\ast\ne t_2$ and $s^\ast<t_2$. Now apply the theorem with the initial time $s^\ast$ and the initial value 0: near $s^\ast$ the zero function is the only solution with this value, so $u=0$ on an interval from $s^\ast$ to some $s^\ast+\delta$ with $\delta>0$. Then $u=0$ on the whole interval from $t_0$ to $s^\ast+\delta$, which contradicts the choice of $s^\ast$. So no such $t_2$ exists, and the same argument with the direction of time reversed treats the times before $t_0$. Hence $u=0$ on all of $I$. (This is the argument of Section 15.4, remark 1: the set of times at which $u$ vanishes is closed, by continuity, and open, by uniqueness near each of its points, so it is the whole interval.) $\square$ (For a constant matrix the solution is explicit, $u(x_4)=e^{-ih_kx_4}u(0)$, Section 8.9, and $u(0)=0$ gives $u=0$ at once.)

So, within the classical theory, a pair cannot arise from the empty configuration by the field equations. It must either be present in the initial data (and then it is not created), or the equations must break down at an initial singular moment (which the notebook's field does not have, by Proposition 16.8), or the equations must be changed. Two limits of this statement: it concerns classical fields, and quantum fields behave differently (the vacuum of EXP-4b creates pairs, Section 16.5); and for the full field equation in 4+4 dimensions, whose initial-value problem is not well posed outside the good sector (Section 8.9), no uniqueness theorem is proved in this book.

\textbf{Proposition 16.8 (the notebook's field is regular at $x_4=0$; PROVED).} For every smooth function $a_4$, on the notebook's chart $0<z<\pi/2$, the metric of Section 9.2 is smooth and non-degenerate at every finite time $x_4$, in particular at $x_4=0$: its eight diagonal entries are finite and nonzero, $\det g=+\cos^2z\ne0$, and its Ricci scalar $R=6H^2(a_4'^2-7)$ is finite.

\emph{Proof.} The entries are $\cot^2z$, $s^{1/3}e^{\pm2a_4}$ and $-1$ with $s=\sin z$ (Section 9.2); on the chart $\cot z$ and $s$ are positive and finite, and $e^{\pm2a_4}$ is positive and finite whenever $a_4$ is finite. The determinant is $\cos^2z$ (Section 9.3) and the Ricci scalar is $6H^2(a_4'^2-7)$ (Section 9.8), finite whenever $a_4'$ is. $\square$

For the notebook's example $a_4=t$ the scale factor of 3-space, $s^{1/6}e^{Hx_4}$, tends to zero only as $x_4\to-\infty$, never at a finite time. The field therefore has no big-bang singularity at $x_4=0$ or at any other finite time on its chart. (Its chart ends at $z=\pi/2$, where only the coordinate $x_0$ degenerates, and at the tip $z\to0$, which lies at infinite proper distance; Sections 9.5 and 9.16.)

\subsection{What a calculation of pair creation would require}

Questions Q2 and Q3 of Section 16.3 can be attacked in two ways. Both are OPEN; this section lists what each would need, so that a student can see where the missing work lies (Chapter 18 collects the open problems).

\textbf{The classical route (Q2).} One would need a solution of the coupled system: the field equations of both members together with Einstein's (or the Einstein--Lovelock) equations, whose source is the total energy--momentum tensor, so that the metric is \textbf{computed} rather than prescribed. By Proposition 16.7 a classical pair cannot grow out of zero fields, so the solution would have to start from a genuinely singular moment, at which the geometry degenerates and the equations lose their meaning, together with a rule for what emerges from it. Such a rule is an initial condition, not a derivation. The notebook's metric has no such moment (Proposition 16.8). The classical route can therefore describe at most how a pair that is already there evolves.

\textbf{The quantum route (Q3).} Here the analogue of Section 16.5 is needed, one level up:

\begin{enumerate}
\item \textbf{A space of states with different numbers of universes}, containing a state with no universe and states with one universe or with a pair. In quantum cosmology the corresponding object is a wave function of the geometry and the matter fields, which obeys the Wheeler--DeWitt equation (B. S. DeWitt, Phys. Rev. 160, 1113 (1967)); the proposals of Hartle and Hawking and of Vilenkin cited in Sections 16.9 and 16.11 are of this kind. Nothing of this exists in the project, which quantizes only the matter field in a prescribed geometry.
\item \textbf{A probability interpretation.} Probabilities need a positive inner product. The canonical state space of dirac16complex is a Krein space, and a positive Fock space exists only in the good sector (Chapter 8). For the cancelling pair no quantum description with two independent universes is known (Section 16.6), so this item is not a formality.
\item \textbf{A dynamics that connects the states:} an interaction term, or a time-dependent background, whose matrix element between the empty state and the pair state is not zero. In the theory as built nothing couples the two members, and gravity sees only their total, which vanishes. For a pair whose members carry the charges $\pm Q_+\ne0$ the interaction would also have to break the separate conservation of the two charges while keeping their sum conserved (Proposition 16.5, item 3, and the paragraph after it).
\item \textbf{The amplitude and what follows from it:} the amplitude $\langle\text{pair}|\,U\,|\text{empty}\rangle$ of the evolution operator $U$, the probability $|\langle\text{pair}|U|\text{empty}\rangle|^2$, and a rate, which could then be compared with the number $na^3=4.412\times10^{-3}H_{\mathrm{inf}}^3$ that EXP-4b computes for particle pairs.
\item \textbf{An answer to the stability question.} If a quantum theory contained an independent universe of mass $-M$ with negative energies, so that pairs of zero total energy could be created by some interaction, energy conservation would not limit their number, since every configuration of the $+M$ member would have a partner with the opposite energy. In ordinary quantum field theory, quanta of negative energy that interact with quanta of positive energy are known to make the empty state unstable (J. M. Cline, S. Jeon and G. D. Moore, Phys. Rev. D 70, 043543 (2004)). A calculation of pair creation in the cancelling reading would have to face this question; it has not been examined here.
\item \textbf{Numbers.} The mass $M$, the length scale $1/H$, the gravitational coupling $\kappa$ and the coupling $\lambda$ would have to be fixed. The theory as built does not fix any of them.
\end{enumerate}

\textbf{The photon analogy, completed.}

\begingroup
\small
\begin{longtable}{@{}>{\raggedright\arraybackslash}p{0.267\linewidth}>{\raggedright\arraybackslash}p{0.267\linewidth}>{\raggedright\arraybackslash}p{0.267\linewidth}@{}}
\toprule
\textbf{ingredient} & \textbf{photon near a nucleus} & \textbf{pair of universes of masses $\pm M$} \\
\midrule
\endfirsthead
\toprule
\textbf{ingredient} & \textbf{photon near a nucleus} & \textbf{pair of universes of masses $\pm M$} \\
\midrule
\endhead
conservation of the totals satisfied? & yes, above the threshold of Proposition 16.3 & yes, for a cancelling pair of classical fields (Proposition 16.5, item 1); quantum level OPEN \\
something to take up momentum or energy & the nucleus & not needed: all totals are zero \\
a coupling that leaves only the totals conserved & yes: light and the electron field interact, so the energy of the light alone is not conserved & none: the members do not interact, so each member's charge (and in a static field its energy) is conserved on its own, which forbids a pair with $Q_+\ne0$ (Proposition 16.5, item 3) \\
dynamics connecting initial and final states & quantum electrodynamics & none in the theory \\
computed probability & the Bethe--Heitler cross section (literature) & none (OPEN) \\
\bottomrule
\end{longtable}
\endgroup

The first two lines are where the pairing theorems stand. The last three are the missing steps, and the third line shows that the missing coupling is needed even before any rate is computed.

\subsection{The ledger of this chapter}

\begingroup
\small
\begin{longtable}{@{}>{\raggedright\arraybackslash}p{0.267\linewidth}>{\raggedright\arraybackslash}p{0.267\linewidth}>{\raggedright\arraybackslash}p{0.267\linewidth}@{}}
\toprule
\textbf{statement} & \textbf{status} & \textbf{where verified} \\
\midrule
\endfirsthead
\toprule
\textbf{statement} & \textbf{status} & \textbf{where verified} \\
\midrule
\endhead
"At $x_4=0$ a pair of universes of masses $\pm M$ is created" (notebook, cells 6, 7 and 17) & HYPOTHESIS & \texttt{handoff/\allowbreak{}surveys/\allowbreak{}survey\_n\allowbreak{}otebook-\allowbreak{}physics.\allowbreak{}md}; Section 16.2 \\
The primordial metric and its function $a_4$ & ASSUMED (prescribed) & Chapter 9 \\
A single photon cannot create a pair; the threshold near a nucleus & PROVED (from standard kinematics) & Section 16.4 \\
An expanding universe creates particle pairs: $na^3=4.412\times10^{-3}H_{\mathrm{inf}}^3$ for $m=H_{\mathrm{inf}}$ & COMPUTED & \texttt{artifact\allowbreak{}s/dirac1\allowbreak{}6complex\allowbreak{}/numeric\allowbreak{}s/exp4/s\allowbreak{}ummary.j\allowbreak{}son}; Section 11.11 \\
T1: $\mathcal L_{m,\lambda}[\gamma^8\Psi]=-\mathcal L_{-m,-\lambda}[\Psi]$, $T\to-T$, $J\to-J$, both statistics, every gravitational field & PROVED & \texttt{PAIR\_T1g\allowbreak{}eneric\_*}, \texttt{PAIR\_T1g\allowbreak{}rassmann\allowbreak{}\_*}, \texttt{PAIR\_T1j\allowbreak{}ets\_*}, \texttt{PAIR\_T1p\allowbreak{}rimordia\allowbreak{}l\_*}; \texttt{ALG\_gamm\allowbreak{}a8Map}; Section 16.6 \\
Corollary 16.4: the cancelling pair has $T=0$, $J=0$ at every point and every $x_4$ & PROVED & \texttt{PAIR\_T1j\allowbreak{}ets\_pair\allowbreak{}EMTAndCu\allowbreak{}rrentVan\allowbreak{}ish}, \texttt{PAIR\_tot\allowbreak{}als\_fiel\allowbreak{}dLevelCh\allowbreak{}iralPair}; Section 16.6 \\
Quantum level: the image field cancels only as $X+(-X)$ within one quantum system; an independently quantized partner adds & PROVED (Fock level cited from Stage 5, provisional until its gate) & \texttt{PAIR\_T1k\allowbreak{}rein\_*}; \texttt{MA\_M4\_kr\allowbreak{}einOnePa\allowbreak{}rticle}; Section 16.6 \\
A quantum description of two independent universes whose energies and charges cancel & OPEN & Section 16.6 \\
T2: space-like reflection gives $(m,\lambda)\to(-m,\lambda)$ with $T\to+T$, $J\to+J$; the mirror universe across the brane & PROVED & \texttt{PAIR\_T2f\allowbreak{}rame\_*}, \texttt{PAIR\_T2z\allowbreak{}2\_*}; Section 16.7 \\
T3: the Kohn--Sham states of the $+M$ and $-M$ universes coincide with the transformed tip condition; the pairing fails without it & PROVED & \texttt{PAIR\_T3b\allowbreak{}lock\_*}, \texttt{PAIR\_T3k\allowbreak{}s\_*}; Section 16.7 \\
Kohn--Sham pair totals: mirror pair $2E_+$, charge $2N$; Krein image pair 0 & PROVED; COMPUTED in 25 Rust runs & \texttt{PAIR\_tot\allowbreak{}als\_ksMi\allowbreak{}rrorPair}, \texttt{PAIR\_tot\allowbreak{}als\_ksKr\allowbreak{}einImage\allowbreak{}Pair}; \texttt{rust/pai\allowbreak{}rs/*/pai\allowbreak{}ring.jso\allowbreak{}n}; Section 16.8 \\
Independent cross-check of the pair numerics; Kohn--Sham pair runs of dirac16complex00 & OPEN & \texttt{referenc\allowbreak{}e/refere\allowbreak{}nce-pair\allowbreak{}s-summar\allowbreak{}y.json}; Section 16.8 \\
The totals of a cancelling pair of classical fields (energy, momentum, charge) are those of the empty state, and the gravitational equations with the pair as only source are source-free (Proposition 16.5, items 1 and 2) & PROVED (partner with $-\lambda$ and the opposite classical energy density $-\rho_+$) & \texttt{PAIR\_T1j\allowbreak{}ets\_pair\allowbreak{}EMTAndCu\allowbreak{}rrentVan\allowbreak{}ish}, \texttt{PAIR\_tot\allowbreak{}als\_fiel\allowbreak{}dLevelCh\allowbreak{}iralPair} (Corollary 16.4); Section 16.9 \\
In the theory as built each member's charge, and in a static field its energy and momenta, are conserved separately, so a cancelling pair with $Q_+\ne0$ (or, in a static field, $E_+\ne0$) cannot appear from the empty state (Proposition 16.5, item 3) & PROVED & Theorem M1 (\texttt{\detokenize{MA_M1_*}}, \texttt{MA\_M5\_im\allowbreak{}plicatio\allowbreak{}n}); Section 16.9 \\
A coupling between the members that would leave only the totals conserved, and so allow a charged cancelling pair & OPEN (not in the theory; a HYPOTHESIS if assumed) & Section 16.9 \\
A mirror pair with nonzero charge cannot appear from the empty state & PROVED & Section 16.9 \\
A cancelling pair cannot produce the primordial field in Einstein gravity (Proposition 16.6) & PROVED & Section 16.10 \\
The primordial field solves the vacuum Einstein--Lovelock equations for some couplings & OPEN & Sections 4.9 and 16.10 \\
A classical field of the good sector cannot appear from zero (Proposition 16.7) & PROVED & Section 16.11 \\
The notebook's field is regular at $x_4=0$ (Proposition 16.8) & PROVED & Section 16.11 \\
A creation process, rate, amplitude or wave function of the universe; a dynamical big bang & OPEN & Sections 16.11 and 16.12; Chapter 18 \\
"The big bang creates universes in pairs" & HYPOTHESIS, not derived: the totals of a classical cancelling pair are those of the empty state, but in the theory as built the separate conservation laws forbid the appearance of a pair whose members are charged & this chapter \\
\bottomrule
\end{longtable}
\endgroup

\subsection{What we proved and what we assumed}

\textbf{Proved in this chapter} (with the complete argument given here, or with the key steps given here and the complete proof in Chapter 15, and in every case with the exact machine checks named): the kinematics of pair creation by light, namely that a single photon in empty space cannot make an electron--positron pair, that near a nucleus of mass $M_N$ the conservation laws allow it exactly when $E_\gamma\ge2m(1+m/M_N)$, and that two head-on photons need $E_aE_b\ge m^2$ (Lemma 16.1, Propositions 16.2 and 16.3); theorem T1, $\mathcal L_{m,\lambda}[\gamma^8\Psi]=-\mathcal L_{-m,-\lambda}[\Psi]$ with $E\to-\gamma^8E$, $T_{\mu\nu}\to-T_{\mu\nu}$ and $J^\mu\to-J^\mu$, for both statistics in every gravitational field, and its reading as a reversal of the frame; Corollary 16.4, the cancelling pair with $T^{\mathrm{pair}}_{\mu\nu}=0$ and $J^{\mathrm{pair}}=0$ at every point and every $x_4$; the rest-frame worked examples, which show the cancellation and show that the classical energy of dirac16complex00 has no fixed sign; the two quantum readings of the partner, of which the image field with the Krein metric $-B$ cancels only as an identity $X+(-X)=0$ within one quantum system and the independently quantized field adds; theorem T2, the mirror map $(m,\lambda)\to(-m,\lambda)$ with $T\to+T$ and $J\to+J$, and the mirror universe of mass $-M$ across the brane; theorem T3 and the Kohn--Sham pair totals (the mirror pair doubles energy and charge; the Krein image pair has zero energy, charge and energy--momentum, a zero of the $X+(-X)$ kind); that the totals of the cancelling pair of classical fields are those of the empty state and that the gravitational equations with the pair as only source are the source-free equations (Proposition 16.5, items 1 and 2); that in the theory as built, where the members do not interact, the charge of each member, and in a static field its energy and momenta, are conserved separately, so that a cancelling pair with $Q_+\ne0$ (or, in a static field, $E_+\ne0$) cannot appear from the empty state (Proposition 16.5, item 3, with Worked example 16.C); that a mirror pair with nonzero charge cannot appear from the empty state; that a cancelling pair cannot produce any member of the primordial family in Einstein gravity, with or without a cosmological constant (Proposition 16.6); that a classical field of the good sector described by finitely many modes cannot appear from zero (Proposition 16.7); and that the notebook's field is regular at $x_4=0$ and at every finite time (Proposition 16.8).

\textbf{Computed:} the particle pairs created by an expanding universe (EXP-4b, Section 11.11), which serve as the model of a creation calculation; the Kohn--Sham pair runs of the Rust solver for dirac16complex (25 configurations), which reproduce the $+M$ universe in the $-M$ universe to at most $1.6\times10^{-8}$ in the energy, give the mirror and Krein image pair totals, and show, in the runs without interaction, that the pairing fails with the untransformed tip condition (the interacting controls were not run, because the solver's energy-window test refused them); the average energy densities of the Kohn--Sham states of Chapter 13, which are positive (Stage-4 outputs, with the cross-check status stated in Section 16.1).

\textbf{Assumed or quoted:} the conservation of energy and momentum and the relation $E^2=m^2+|\mathbf p|^2$ of special relativity; the literature on pair-creation cross sections, on gravitational particle creation, on zero-energy universes, on the wave function of the universe and on the instability caused by negative-energy quanta, all cited and none used as a result; the theorem of Picard and Lindelöf (Section 10.2); the prescribed primordial metric and its free function $a_4$ (Chapter 9); the Z2 gluing of Stage 4 as a model of the notebook's picture (Section 9.16); the Kohn--Sham model of Chapter 13 with its approximations; the mean-field reading of the bilinears in the source analysis of Chapter 9; and the defining assumptions of the cancelling pair: two independent, non-interacting classical fields in the correlated configuration $\Psi_-=\gamma^8\Psi_+$, the partner with the coupling $-\lambda$ and with the opposite classical energy density $-\rho_+$ (negative wherever $\rho_+>0$).

\textbf{Open:} a creation process, a rate, a probability amplitude, a wave function of the universe in the sense of quantum cosmology, a dynamical big bang; a coupling between the two members that would break their separate conservation laws and leave only the totals conserved (the theory as built has none); a quantum description of two independent universes of masses $\pm M$ whose energies and charges cancel; whether the primordial field solves the vacuum Einstein--Lovelock equations for some couplings; the sources of the Z2 geometry (bulk and brane); the independent cross-check of the Kohn--Sham pair numerics and the pair runs of dirac16complex00.

\textbf{Hypothesis:} that the big bang creates universes of masses $\pm M$ in pairs. The pairing theorems give a cancelling pair of classical fields whose totals are those of the empty state; in the theory as built the separate conservation laws of its two members still forbid its appearance unless each member is neutral (and, in a static field, of zero energy), and nothing derives its creation. In the words of Stage 1, the pairing is "a structural property of the equations, not a claim that universes of masses $\pm M$ are created in pairs".

\subsection{Exercises}

\textbf{Exercise 16.1.} Give each statement one of the five status words of Section 0.3, with one sentence of reason. (a) For a T1 pair in the static primordial field, $T^{\mathrm{pair}}_{\mu\nu}=0$ at $x_4=0$. (b) At $x_4=0$ a pair of universes of masses $\pm M$ is created. (c) The probability per unit time that a pair of universes is created. (d) In the Rust run \texttt{d16c\_m1\_\allowbreak{}L3\_N112\_\allowbreak{}lam0\_T0} the mirror pair has the energy 122.1814466. (e) The two halves of the Z2 geometry are glued at $y=0$.

\textbf{Exercise 16.2.} Two electrons ($m=1$) have the momenta $\mathbf p_1=(3,0,0)$ and $\mathbf p_2=(0,4,0)$. Compute $E_1$, $E_2$ and the invariant mass squared of the pair, and check Lemma 16.1.

\textbf{Exercise 16.3.} Show that a single photon cannot turn into one neutral particle of mass $\mu>0$, and that such a particle, at rest, cannot decay into a single photon.

\textbf{Exercise 16.4.} Complete the proof of Proposition 16.3: for $E_\gamma\ge2m(1+m/M_N)$, construct momenta of the electron, the positron and the nucleus, all along the direction $\mathbf n$ of the photon, with the initial total energy $E_\gamma+M_N$ and the initial total momentum $E_\gamma\mathbf n$. (Hint: give the electron and the positron the momenta $(q+r)\mathbf n$ and $(q-r)\mathbf n$, the nucleus $(E_\gamma-2q)\mathbf n$, choose $q=mE_\gamma/(2m+M_N)$, and vary $r\ge0$.)

\textbf{Exercise 16.5.} In units $m=1$: (a) compute the threshold of Proposition 16.3 for $M_N=4$ and for $M_N=1836$. (b) Two photons collide head-on, one with $E_a=3$. How large must $E_b$ be at least? (c) If the two photons have equal energies, how large must each be?

\textbf{Exercise 16.6.} Using the tables of Section 2.9 ($C$: row 1 reads $-5$, row 5 reads $-1$, row 8 reads $+12$, row 12 reads $+8$; $B/i$: row 1 reads $-8$, row 5 reads $+12$, row 8 reads $+1$, row 12 reads $-5$), compute $S=\Psi^\dagger C\Psi$ and $J^4=\Psi^\dagger B\Psi$ for $\Psi=e_1+e_5+i\,e_8$ and for $\gamma^8\Psi$.

\textbf{Exercise 16.7.} In flat space with $\lambda=0$ and $m>0$, let $\Psi=\tfrac1{\sqrt2}(u_++u_-)\,e^{-imx_4}$ with the rest states $u_\pm$ of Section 8.8. Show that $\Psi$ solves the field equation with mass $m$, and compute its classical energy density $\rho$, its charge density $J^4$ and its scalar density $S$. What does the result say about "zero total energy" as a property of pairs?

\textbf{Exercise 16.8.} The file

\begin{Verbatim}[fontsize=\small]
artifacts/dirac16complex/pair-creation/rust/pairs/d16c_m1_L3_N112_lam0_T0/pairing.json
\end{Verbatim}

gives $E_+=61.090723287$ for the $+M$ universe with $N=112$ particles (key \texttt{universe\allowbreak{}s}, entry \texttt{\detokenize{plusM}}, field \texttt{\detokenize{E}}). (a) What are the energy and the charge of the mirror pair and of the Krein image pair? (b) Compare these totals with those of the empty state. (c) Now use also the separate conservation of the charge and of the energy of each universe (Proposition 16.5, item 3): may either pair appear from the empty state in the static field, as far as the conservation laws are concerned?

\textbf{Exercise 16.9.} Show directly that the static primordial field cannot be produced by a T1 pair together with a cosmological constant: which value of $\Lambda$ would the component $G^0{}_0$ require, and which the component $G^4{}_4$?

\textbf{Exercise 16.10.} For the notebook's choice $a_4=t$, at the point with $\sin z=3/5$ and with $H=1$: write down the diagonal of the metric at $x_4=0$ and at $x_4=-10$, and the Ricci scalar. When does the scale factor of 3-space vanish?

\textbf{Exercise 16.11.} Let $h(t)=t\,\sigma_x$ with the Pauli matrix $\sigma_x$ of Section 2.3. (a) Show that $u(t)=\bigl(\cos\tfrac{t^2}2,\,-i\sin\tfrac{t^2}2\bigr)^T$ solves $i\,\dot u=h(t)u$ with $u(0)=(1,0)^T$. (b) Show that $u^\dagger u$ is constant. (c) Which solution has $u(0)=0$, and why is it the only one?

\textbf{Exercise 16.12.} (a) Assuming $K\to K$ and $S\to-S$ (the map of T2), show that $\mathcal L_{m,\lambda}$ goes over into $\mathcal L_{-m,\lambda}$. (b) Assuming $K\to-K$ and $S\to S$ (the map of T1), show that it goes over into $-\mathcal L_{-m,-\lambda}$. (c) Explain why no map with $K\to-K$ and $S\to S$ can turn $\mathcal L_{m,\lambda}$ into $-\mathcal L_{-m,\lambda}$ when $\lambda\ne0$.

\textbf{Exercise 16.13.} Repeat the table of Section 16.6 for a quantum in the rest mode $u_-$ of Section 8.8 (Krein sign $-1$): compute $(E,Q,S)$ in the $+M$ universe, for the image partner and for the independent partner. Compare the energy with the classical energy density $m\,u_-^\dagger Bu_-$.

\textbf{Exercise 16.14.} Let $\Psi_+\to e^{i\alpha}\Psi_+$ and $\Psi_-\to e^{i\beta}\Psi_-$ with constant real $\alpha$ and $\beta$. (a) By which factor are $S_+=\bar\Psi_+\Psi_+$, $S_-=\bar\Psi_-\Psi_-$ and $\bar\Psi_+\Psi_-=\Psi_+^\dagger C\Psi_-$ multiplied? (b) For which pairs $(\alpha,\beta)$ is a term containing $\bar\Psi_+\Psi_-$ unchanged, and which combination of the two charges would a Lagrangian with such a term still conserve? (c) The partner of Worked example 16.B has $J^4_-=-1$ in a periodic box of 7-volume $V$. Using (b), explain why in the theory as built (without such a term) the conservation of charge forbids this pair to appear from the empty state, and why with such a term the conservation of charge would no longer forbid it.

\subsection{Answers to the exercises}

\textbf{Answer 16.1.} (a) PROVED: it is Corollary 16.4, which holds for the two classical fields of a cancelling pair in every gravitational field at every $x_4$. (b) HYPOTHESIS: it is the notebook's claim; nothing in the project derives it. (c) OPEN: no calculation of such a probability exists (Section 16.12). (d) COMPUTED: it is a floating-point result of the Rust solver, recorded in the file's key \texttt{pairTota\allowbreak{}ls}. (e) ASSUMED: the Z2 gluing is a construction chosen by Stage 4 to model the notebook's picture, not a result of any equation (Section 9.16).

\textbf{Answer 16.2.} $E_1=\sqrt{1+9}=\sqrt{10}\approx3.1623$ and $E_2=\sqrt{1+16}=\sqrt{17}\approx4.1231$. The total momentum is $(3,4,0)$ with $|\mathbf p_1+\mathbf p_2|^2=25$. So $\mu^2=(\sqrt{10}+\sqrt{17})^2-25=10+17+2\sqrt{170}-25=2+2\sqrt{170}\approx28.08$. Lemma 16.1 requires $\mu^2\ge(1+1)^2=4$, which holds. By the proof of Lemma 16.1, equality needs both $\mathbf p_1\cdot\mathbf p_2=ab$ (the two momenta point in the same direction, or one of them is zero) and $\mu_1b=\mu_2a$ (the sizes of the momenta are proportional to the masses); for two electrons, $\mu_1=\mu_2$, this means $\mathbf p_1=\mathbf p_2$. Here $\mathbf p_1\cdot\mathbf p_2=0<ab=12$ and the sizes 3 and 4 differ, so the inequality is strict.

\textbf{Answer 16.3.} A photon has $\mu^2=E^2-|\mathbf p|^2=0$; one particle of mass $\mu$ has $\mu^2>0$. Since the invariant mass squared is conserved, $0=\mu^2$ is impossible. The decay is the same statement read backwards: the particle at rest has $E=\mu$ and $\mathbf p=0$, so $E^2-|\mathbf p|^2=\mu^2>0$, while one photon has 0.

\textbf{Answer 16.4.} Write $P=E_\gamma$ and $\mu_0=2m+M_N$. The total momentum is $(q+r)+(q-r)+(P-2q)=P$ for every $q$ and $r$. The total energy is the continuous function

\[
f(r)=\sqrt{m^2+(q+r)^2}+\sqrt{m^2+(q-r)^2}+\sqrt{M_N^2+(P-2q)^2}.
\]

With $q=mP/\mu_0$ the nucleus has the momentum $P-2q=M_NP/\mu_0$, so every body has momentum proportional to its mass, and $f(0)=\frac{m}{\mu_0}\sqrt{\mu_0^2+P^2}+\frac{m}{\mu_0}\sqrt{\mu_0^2+P^2}+\frac{M_N}{\mu_0}\sqrt{\mu_0^2+P^2}=\sqrt{\mu_0^2+P^2}$. As $r$ grows, $f(r)$ grows without bound. By the intermediate value theorem $f$ takes every value between $f(0)$ and infinity, so a value $r\ge0$ with $f(r)=E_\gamma+M_N$ exists provided $\sqrt{\mu_0^2+E_\gamma^2}\le E_\gamma+M_N$. Squaring, this is $(2m+M_N)^2\le M_N^2+2E_\gamma M_N$, which is exactly the threshold condition $E_\gamma\ge2m(1+m/M_N)$. For two photons the same construction with the pair alone ($\mu_0=2m$, $P=E_a-E_b$, $q=P/2$) works exactly when $4m^2+(E_a-E_b)^2\le(E_a+E_b)^2$, that is $E_aE_b\ge m^2$.

\textbf{Answer 16.5.} (a) $2(1+\tfrac14)=2.5$ and $2(1+\tfrac1{1836})\approx2.00109$. (b) $E_b\ge m^2/E_a=\tfrac13$. (c) $E^2\ge1$, so each photon needs at least $E=1=m$: two photons of equal energy must each bring one rest energy.

\textbf{Answer 16.6.} For $\Psi$: $(C\Psi)_1=-\Psi_5=-1$, $(C\Psi)_5=-\Psi_1=-1$, $(C\Psi)_8=\Psi_{12}=0$, so $S=1\cdot(-1)+1\cdot(-1)+(-i)\cdot0=-2$. $(B\Psi)_1=-i\Psi_8=-i\cdot i=1$, $(B\Psi)_5=i\Psi_{12}=0$, $(B\Psi)_8=i\Psi_1=i$, so $J^4=1\cdot1+1\cdot0+(-i)\cdot i=1+1=2$ (with $\Psi_8^\ast=-i$). For $\gamma^8\Psi=-e_1-e_5+i\,e_8$ (the components 1 and 5 lie in the upper half and change sign): $(C\gamma^8\Psi)_1=1$, $(C\gamma^8\Psi)_5=1$, so $S=(-1)(1)+(-1)(1)=-2$; $(B\gamma^8\Psi)_1=-i\cdot i=1$, $(B\gamma^8\Psi)_8=i\cdot(-1)=-i$, so $J^4=(-1)(1)+(-i)(-i)=-1-1=-2$. The scalar density is unchanged and the charge density changes sign, as T1 requires.

\textbf{Answer 16.7.} Both $u_+$ and $u_-$ satisfy $-i\gamma^4u_\pm=u_\pm$ (Section 8.8), so $\gamma^4\partial_4\Psi=-im\gamma^4\Psi=m\Psi$: $\Psi$ solves the field equation with mass $m$ and has frequency $\varepsilon=m$. From Section 8.8, $Bu_+=u_+$, $Bu_-=-u_-$, $Cu_+=u_+$, $Cu_-=-u_-$, and $u_\pm^\dagger u_\pm=1$; one also checks $u_+^\dagger u_-=\tfrac14\bigl(1\cdot1+(-1)\cdot1+(i)(i)+(i)(-i)\bigr)=0$. Hence $B(u_++u_-)=u_+-u_-$ and $(u_++u_-)^\dagger(u_+-u_-)=1-1+u_-^\dagger u_+-u_+^\dagger u_-=0$, so $J^4=0$ and, by the formula of Worked example 16.B, $\rho=m\,\Psi^\dagger B\Psi=0$. In the same way $C(u_++u_-)=u_+-u_-$ gives $S=0$. A single classical universe of mass $m$ can have zero energy density and zero charge density. For the classical field, zero total energy is therefore not a property that singles out pairs; what T1 adds is the exact cancellation of every component of the energy--momentum tensor between a configuration and its partner.

\textbf{Answer 16.8.} (a) Mirror pair: $E=2\times61.090723287=122.181446574$ and charge $2\times112=224$. The file's \texttt{pairTota\allowbreak{}ls.mirro\allowbreak{}rPair.to\allowbreak{}tals} records $E=122.18144657377286$, which is the sum $E_++E_-$ of the two separately computed universes; it differs from $2E_+$ by $7.3\times10^{-11}$, the numerical difference between $E_-$ and $E_+$. Rounded to six decimals it is 122.181447. Krein image pair: $E=0$ (the file records $-7.3\times10^{-11}$, which is $E_+-E_-$, the same numerical difference) and charge 0. (b) The empty state has energy 0 and charge 0. The mirror pair differs from it in both totals; the Krein image pair agrees with it in both. (c) The mirror pair is forbidden already by its totals: charge conservation and, in the static field, energy conservation. The Krein image pair is forbidden too. Its $+M$ universe alone has the charge 112 and the energy 61.090723287, and in the theory as built each universe's charge and, in the static field, its energy are conserved on their own (Proposition 16.5, item 3), while the empty state has the value 0 for each. Moreover its zero is of the form $X+(-X)$: the image is the same quantum system as the $+M$ universe, with its own charge $+112$ relative to its own canonical structure (Section 16.6). So \textbf{neither} pair can appear from the empty state, as far as the conservation laws are concerned, and whether two independent quantum universes can cancel at all is OPEN.

\textbf{Answer 16.9.} With a T1 pair as the only matter, the equations are $G^\mu{}_\nu+\Lambda\delta^\mu{}_\nu=0$. For the static field $G^0{}_0=15H^2$ requires $\Lambda=-15H^2$, and $G^4{}_4=21H^2$ requires $\Lambda=-21H^2$. One constant cannot have both values.

\textbf{Answer 16.10.} At $\sin z=3/5$: $\cot^2z=16/9$ and $s^{1/3}=(3/5)^{1/3}\approx0.8434$. At $x_4=0$, $a_4=0$ and the diagonal is $(16/9,\ 0.8434,\ 0.8434,\ 0.8434,\ -1,\ -0.8434,\ -0.8434,\ -0.8434)$. At $x_4=-10$, $a_4=-10$, $e^{2a_4}=e^{-20}\approx2.061\times10^{-9}$ and $e^{-2a_4}\approx4.852\times10^{8}$, so the diagonal is $(16/9,\ 1.738\times10^{-9}\ (\times3),\ -1,\ -4.092\times10^{8}\ (\times3))$: very different, but finite and nonzero. The Ricci scalar is $6(1-7)=-36$ at every time. The scale factor of 3-space, $s^{1/6}e^{x_4}$, is positive at every finite time and tends to 0 only as $x_4\to-\infty$: there is no moment at which it vanishes.

\textbf{Answer 16.11.} (a) $\dot u=\bigl(-t\sin\tfrac{t^2}2,\ -it\cos\tfrac{t^2}2\bigr)^T$, so $i\dot u=\bigl(-it\sin\tfrac{t^2}2,\ t\cos\tfrac{t^2}2\bigr)^T$. And $h(t)u=t\,\sigma_xu$ exchanges the two components: $t\bigl(-i\sin\tfrac{t^2}2,\ \cos\tfrac{t^2}2\bigr)^T$, the same. At $t=0$, $u=(1,0)^T$. (b) $u^\dagger u=\cos^2\tfrac{t^2}2+\sin^2\tfrac{t^2}2=1$. (c) $u(t)=0$ for all $t$ solves the equation and has $u(0)=0$. By the theorem of Picard and Lindelöf (Section 10.2) any solution with $u(0)=0$ agrees with it near $t=0$; and the set of times at which a solution with $u(0)=0$ vanishes is closed (by continuity) and open (by the same theorem applied at each of its points), so it is the whole interval of times, exactly as in the proof of Proposition 16.7. So the zero function is the only solution with $u(0)=0$. Nothing can grow out of zero, even though $h$ changes with time: this is Proposition 16.7 for one mode.

\textbf{Answer 16.12.} (a) $\mathcal L_{m,\lambda}/\sqrt{|g|}=K-mS-\tfrac\lambda2S^2\to K-m(-S)-\tfrac\lambda2(-S)^2=K-(-m)S-\tfrac\lambda2S^2=\mathcal L_{-m,\lambda}/\sqrt{|g|}$. (b) $K-mS-\tfrac\lambda2S^2\to-K-mS-\tfrac\lambda2S^2=-\bigl[K-(-m)S-\tfrac{(-\lambda)}2S^2\bigr]=-\mathcal L_{-m,-\lambda}/\sqrt{|g|}$. (c) Under such a map the image is $-K-mS-\tfrac\lambda2S^2$, while $-\mathcal L_{-m,\lambda}/\sqrt{|g|}=-K-mS+\tfrac\lambda2S^2$. They differ by $-\lambda S^2$, which is not zero for $\lambda\ne0$ (and for a field with $S\ne0$). The even potential forces $\lambda\to-\lambda$ in every map that reverses the kinetic term.

\textbf{Answer 16.13.} $h(m)u_-=-im\gamma^4u_-=m\,u_-$, so in the $+M$ universe $E=u_-^\dagger h(m)u_-=m$, $Q=u_-^\dagger u_-=1$ and $S=u_-^\dagger(-i\gamma^4)u_-=u_-^\dagger u_-=1$. For $w=\gamma^8u_-$ one has $-i\gamma^4w=-w$ (because $\gamma^4$ anticommutes with $\gamma^8$), hence $h(-m)w=m\,w$. The image partner gives $E=-w^\dagger h(-m)w=-m$, $Q=-w^\dagger w=-1$ and $S=-w^\dagger(-i\gamma^4)w=1$; the independent partner gives $E=m$, $Q=1$ and $S=-1$. These are the same numbers as for $u_+$: with the expectation rule of Section 8.12 a quantum in a positive-energy mode has positive energy whatever the Krein sign of the mode. The classical energy density of the same mode is $m\,u_-^\dagger Bu_-=-m$. The classical bilinear and the quantum expectation value can therefore have opposite signs, which is why the classical cancellation of Corollary 16.4 does not carry over automatically to the quantum field (Section 16.6).

\textbf{Answer 16.14.} (a) The phases are ordinary numbers, so they can be moved to the front of every product, for commuting and for Grassmann components alike. With $\Psi_+^\dagger\to e^{-i\alpha}\Psi_+^\dagger$: $S_+=\Psi_+^\dagger C\Psi_+\to e^{-i\alpha}e^{i\alpha}S_+=S_+$, the factor 1; in the same way $S_-$ gets the factor 1; and $\bar\Psi_+\Psi_-=\Psi_+^\dagger C\Psi_-\to e^{-i\alpha}e^{i\beta}\,\bar\Psi_+\Psi_-=e^{i(\beta-\alpha)}\,\bar\Psi_+\Psi_-$. (b) The factor $e^{i(\beta-\alpha)}$ equals 1 exactly when $\beta-\alpha$ is a whole multiple of $2\pi$, that is, for the common rotation $\beta=\alpha$ (up to the harmless multiples of $2\pi$). A Lagrangian with such a term is therefore still unchanged by the common rotation, but no longer by the rotation of $\Psi_+$ alone or of $\Psi_-$ alone. Noether's theorem for the common rotation gives the current from the Noether formula of Theorem M1, summed over both fields; the new term contains no derivatives, so it adds nothing to that formula, and the conserved charge is $Q_++Q_-$. The separate charges $Q_+$ and $Q_-$ are no longer conserved. (c) The pair has $Q_+=V$ and $Q_-=-V$. Without such a term $Q_+$ is conserved on its own (Proposition 16.5, item 3), and the empty state has $Q_+=0\ne V$, so charge conservation forbids the transition. With such a term only $Q_++Q_-$ is conserved, and its value is $V-V=0$ both for the pair and for the empty state, so charge conservation no longer forbids it. That would not yet make the transition possible: the new term changes the field equations and the energy, so the energy balance and the pair itself would have to be examined again, and for classical fields Proposition 16.7 still excludes growth from zero fields (Section 16.9).

\section{Matter and antimatter}

\subsection{What this chapter answers}

The request that this book answers ends with the words "this theory solves matter anti-matter mysteries" (Section 0.2). This chapter explains what those mysteries are, starting from zero, and then states exactly what the dirac16complex theory contributes to them. The honest answer comes first, because everything else in the chapter supports it.

\textbf{The honest answer.} The theory as built does not solve the matter--antimatter problem, and within the theory the central step of such a solution is impossible. The Lagrangian of both fields of this book, dirac16complex (anticommuting components) and dirac16complex00 (commuting components), does not change when the field is multiplied by a constant phase $e^{i\alpha}$. By Noether's theorem the associated charge $Q$ is conserved in every gravitational field, as long as no charge flows in or out through the boundary of space (Theorem M1, Section 17.9), so no process of the theory can create a net charge inside one universe. The theory also has exact symmetries that turn every solution with charge $Q$ into a solution with charge $-Q$ without reversing the time $x_4$: for dirac16complex00 the charge conjugation C, in every gravitational field; for dirac16complex a combination CP of charge conjugation with a spatial reflection, in flat space and in every gravitational field in which the reflection is an isometry (leaves the metric unchanged), for example the primordial field of Chapter 9 with a reflection of $x_1$, $x_2$ or $x_3$ (Theorem M2, Sections 17.10 and 17.11). A generic gravitational field breaks this CP, but that opens no way to an asymmetry, because the charge stays conserved. Nothing in the repository computes a departure from thermal equilibrium. The first two facts negate the first two of the three conditions that Sakharov found in 1967 to be necessary for creating a matter--antimatter asymmetry (Section 17.5); the third condition is not addressed by any computation, and since the first fails exactly, a departure from equilibrium could not help. What the theory does provide is exact but limited: as classical fields, a universe of mass $+M$ and its chirality image of mass $-M$ carry opposite charges and zero total charge (Theorem M4, Section 17.13). At the quantum level the image is the same quantum system with the Krein metric $-B$, and no cancellation between two independently quantised universes follows (Section 17.13). Under three hypotheses that nothing in the repository derives, this gives a picture in which the whole pair is symmetric while each member is asymmetric (Section 17.14). The picture predicts no number; in particular it does not predict the observed baryon-to-photon ratio $\eta\approx6\times10^{-10}$.

\textbf{The plan.} Sections 17.2 to 17.7 explain the problem from zero: antiparticles and conserved numbers, the size of the observed asymmetry, why the universe is not a patchwork of matter and antimatter regions, Sakharov's three conditions with a derivation in plain words, why the Standard Model of particle physics does not suffice, and the class of "universe and anti-universe" ideas. Sections 17.8 to 17.16 carry out the exact analysis of the theory: the theorems M1 to M4, the conditional scenario M5 and the scorecard M6, followed by the list of what would have to be added to the theory. Section 17.17 lists the machine checks and the commands that rerun them.

\textbf{Sources.} The exact analysis of this chapter is that of the matter-antimatter document, whose title is "Matter and antimatter in the dirac16complex theory: what can be proved", written under its own specification; this chapter follows it and explains each step for a beginner. Its machine checks are in two independent reports, one written by a Wolfram Language verifier (44 of 44 checks true) and one by an independent Python checker (75 of 75 checks true); the Wolfram verifier also writes a file of exact results. The committed files are not all of the same age, and the chapter says where this matters. On two points it follows the wording of the committed Wolfram report and exact-results file, which were rewritten after the committed text of the document: the scope of CP in a gravitational field (Section 17.11) and the reading of the pair of universes at the quantum level (Sections 17.13 and 17.14). The committed Python report was written by an earlier version of the Python checker; the checker in the same commit has 77 checks, two more than that report, and a rerun of it gives 77 of 77 true (Section 17.17). The quantum statements about the pair of universes use the exact pairing results of Stage 5 (141 of 141 checks true). The files are

\begin{Verbatim}[fontsize=\small]
provenance/DIRAC16COMPLEX_MATTER_ANTIMATTER.md        the matter-antimatter document
handoff/specs/MATTER_ANTIMATTER_SPEC.md               its specification
artifacts/dirac16complex/matter-antimatter/
    wolfram-matter-antimatter-report.json             Wolfram report (44 checks)
    python-matter-antimatter-report.json              Python report (75 checks, earlier
                                                      version of the checker)
    matter-antimatter-theory.json                     exact results (Wolfram)
wolfram/Dirac16ComplexMatterAntimatter.wl             the Wolfram package
scripts/verify_dirac16complex_matter_antimatter.wls   the Wolfram verifier
scripts/check_dirac16complex_matter_antimatter.py     the Python checker
artifacts/dirac16complex/pair-creation/
    pairing-theory.json                               Stage-5 exact pairing results
    wolfram-pairing-report.json                       Stage-5 report (141 checks)
\end{Verbatim}

Check names in typewriter type that begin with \texttt{\detokenize{MA_}} belong to the two matter-antimatter reports (two of them, marked where they appear, only to the Python checker of this commit and not yet to its committed report), and those that begin with \texttt{\detokenize{PAIR_}} to the Stage-5 pairing report. The literature cited in Sections 17.2 to 17.7 is listed in Section 17.18.

\textbf{Labels.} As everywhere in this book (Section 0.3): PROVED means a complete derivation is given here, and where the repository has an exact machine check it is named; COMPUTED means floating-point data read from a committed file; ASSUMED marks a convention, an input or an approximation; HYPOTHESIS marks a claim that nothing in the repository derives; OPEN marks a question that the project has not answered. The conventions are those of the whole book: counting from 0, coordinates $x_0,\dots,x_7$ with $x_4$ the time, $\eta=\mathrm{diag}(+1,+1,+1,+1,-1,-1,-1,-1)$, the notebook's gamma matrices $\gamma^0,\dots,\gamma^7$ (Chapter 2), $C=\gamma^0\gamma^1\gamma^2\gamma^3$, the chirality $\gamma^8=\gamma^0\gamma^1\cdots\gamma^7=\mathrm{diag}(-I_8,I_8)$, $\bar\Psi=\Psi^\dagger C$ and $B=-iC\gamma^4$.

\subsection{Particles, antiparticles and conserved numbers}

\textbf{Antiparticles.} In 1928 Dirac wrote down his equation for the electron (Chapter 2). Its solutions come with negative as well as positive energies, and Dirac concluded that a particle with the mass of the electron and the opposite electric charge must exist. Such a particle, the positron, was found in cosmic rays in 1932 (C. D. Anderson, Phys. Rev. 43, 491 (1933)). Every known particle has an \textbf{antiparticle} with the same mass and the opposite values of all its charges: the electron has the positron, the proton the antiproton, the neutron the antineutron (both neutral, but with opposite baryon number, defined below). A few neutral particles, such as the photon, are their own antiparticles.

\textbf{Annihilation and pair creation.} A particle and its antiparticle can \textbf{annihilate}: an electron and a positron that meet turn into two (or more) photons. Conversely, energy can turn into a particle--antiparticle pair: a photon of enough energy that passes near an atomic nucleus, which takes up momentum, can turn into an electron and a positron. (A single free photon cannot do this: energy and momentum cannot both be conserved.) In both processes the total electric charge is unchanged, because the two members of a pair carry opposite charges. Chapter 11 met a gravitational version of pair creation: in an expanding universe the change of the mode Hamiltonian with time lifts quanta out of the Dirac sea and leaves antiparticles behind (Section 11.11).

\textbf{Three conserved numbers.} Three kinds of charge matter here.

\begin{itemize}
\item The \textbf{electric charge}, measured in units of the proton charge. It is exactly conserved in every process ever observed.
\item The \textbf{baryon number} $B$. Protons, neutrons and their heavier relatives are \textbf{baryons} and have $B=+1$; their antiparticles have $B=-1$; each quark carries $B=\tfrac13$ and each antiquark $-\tfrac13$. Particles made of a quark and an antiquark, such as the pions $\pi^+,\pi^-,\pi^0$, have $B=0$, and so do electrons, neutrinos and photons.
\item The \textbf{lepton number} $L$. Electrons, muons, tau leptons and neutrinos have $L=+1$ and their antiparticles $L=-1$.
\end{itemize}

"Matter" in cosmology means baryons: the atoms of stars, planets and gas owe almost all their mass to protons and neutrons. The \textbf{matter--antimatter problem} is the question why the universe contains baryons but essentially no antibaryons.

\textbf{Worked example: bookkeeping.} A number is conserved in a process if its total before equals its total after. Count the three numbers in three observed processes (a bar denotes an antiparticle, $\bar\nu_e$ is the electron antineutrino, $\gamma$ a photon):

\begingroup
\small
\begin{longtable}{@{}>{\raggedright\arraybackslash}p{0.200\linewidth}>{\raggedright\arraybackslash}p{0.200\linewidth}>{\raggedright\arraybackslash}p{0.200\linewidth}>{\raggedright\arraybackslash}p{0.200\linewidth}@{}}
\toprule
\textbf{Process} & \textbf{electric charge} & \textbf{baryon number $B$} & \textbf{lepton number $L$} \\
\midrule
\endfirsthead
\toprule
\textbf{Process} & \textbf{electric charge} & \textbf{baryon number $B$} & \textbf{lepton number $L$} \\
\midrule
\endhead
$n\to p+e^-+\bar\nu_e$ & $0\to+1-1+0=0$ & $1\to1+0+0=1$ & $0\to0+1-1=0$ \\
$e^++e^-\to\gamma+\gamma$ & $+1-1=0\to0$ & $0\to0$ & $-1+1=0\to0$ \\
$p+\bar p\to\pi^++\pi^-+\pi^0$ & $+1-1=0\to+1-1+0=0$ & $1-1=0\to0$ & $0\to0$ \\
\bottomrule
\end{longtable}
\endgroup

All three numbers are conserved in all three processes. The first line, the decay of a free neutron, turns a neutron into a proton; $B$ is unchanged because both are baryons.

\textbf{Conservation and symmetry.} A number is \textbf{conserved} when no process can change it. Noether's theorem (Section 5.7) connects conservation with symmetry: if the action of a field theory does not change under a continuous transformation that is the same at every point, then there is a current $j^\mu$ whose divergence vanishes on every solution, and the integral of its time component over space, the \textbf{charge}, does not change with time. The phase symmetry $\psi\to e^{i\alpha}\psi$ of Example 1 of Section 5.7 is the simplest case. Theorem M1 of Section 17.9 is the same theorem for the field of this book, and it is the reason why the theory, as built, cannot produce an asymmetry of its own charge.

\subsection{What is observed: the baryon-to-photon ratio}

\textbf{The definition.} The universe is filled with the photons of the \textbf{cosmic microwave background}, the thermal radiation left over from the hot early universe, now at the temperature $T_0=2.7255$ K (D. J. Fixsen, Astrophys. J. 707, 916 (2009)). Because both the baryons and these photons are diluted in the same way by the expansion (after the first seconds of the universe their numbers per comoving volume stay constant), their ratio is a fixed number that characterizes our universe, the \textbf{baryon-to-photon ratio}

\[
\eta=\frac{n_B-n_{\bar B}}{n_\gamma},
\]

where $n_B$, $n_{\bar B}$ and $n_\gamma$ are the numbers of baryons, antibaryons and photons per unit volume. Today $n_{\bar B}$ is negligible, so $\eta$ is simply the number of baryons per photon.

\textbf{Two independent measurements.} The value of $\eta$ is measured in two independent ways. First, the nuclei of the light elements (deuterium, helium-3, helium-4, lithium-7) were made in the first minutes of the universe (\textbf{big-bang nucleosynthesis}), and how much of each was made depends on $\eta$; the observed abundances fix $\eta$ (R. H. Cyburt, B. D. Fields, K. A. Olive and T.-H. Yeh, Rev. Mod. Phys. 88, 015004 (2016)). Second, the pattern of hot and cold spots in the microwave background (its "acoustic peaks") depends on the density of baryons at the time the radiation was released. The Planck satellite measured it; the Planck 2018 cosmological parameters give

\[
\Omega_bh^2=0.0224\pm0.0001
\]

(Planck Collaboration, Astron. Astrophys. 641, A6 (2020)). The two determinations agree.

\textbf{What $\Omega_bh^2$ means.} $H_0$ is today's Hubble rate (Section 11.2), written $H_0=100\,h$ km/s/Mpc with a pure number $h$ (1 Mpc, a megaparsec, is $3.0857\times10^{22}$ m). The \textbf{critical density} $\rho_{\mathrm{crit}}=3H_0^2/(8\pi G)$, with Newton's constant $G$, is the total density of a spatially flat universe with the Hubble rate $H_0$ (it is the first Friedmann equation of Section 11.2, $H^2=\tfrac{8\pi G}3\rho$, solved for $\rho$). $\Omega_b$ is the mass density of baryons divided by $\rho_{\mathrm{crit}}$. The combination $\Omega_bh^2$ is what the microwave background measures directly: it is proportional to the baryon mass density itself, because $\rho_{\mathrm{crit}}\propto H_0^2\propto h^2$.

\textbf{From $\Omega_bh^2$ to $\eta$, step by step.} We use the standard values of the constants (E. Tiesinga, P. J. Mohr, D. B. Newell and B. N. Taylor, Rev. Mod. Phys. 93, 025010 (2021)): $G=6.6743\times10^{-11}\ \mathrm{m^3\,kg^{-1}\,s^{-2}}$, Boltzmann's constant $k_B=1.380649\times10^{-23}$ J/K, Planck's constant divided by $2\pi$, $\hbar=1.054572\times10^{-34}$ J s, the speed of light $c=2.99792458\times10^8$ m/s and the proton mass $m_p=1.672622\times10^{-27}$ kg.

\begin{enumerate}
\item \emph{The unit of the Hubble rate.} $100$ km/s/Mpc $=10^5\ \mathrm{m\,s^{-1}}/(3.0857\times10^{22}\ \mathrm m)=3.2408\times10^{-18}\ \mathrm{s^{-1}}$.
\item \emph{The critical density.} $\rho_{\mathrm{crit}}=h^2\cdot3\,(3.2408\times10^{-18})^2/(8\pi\cdot6.6743\times10^{-11})=1.8783\times10^{-26}\,h^2$ kg/m$^3$. Hence the baryon mass density is $\rho_b=\Omega_b\rho_{\mathrm{crit}}=1.8783\times10^{-26}\,(\Omega_bh^2)$ kg/m$^3$.
\item \emph{The number of baryons.} Almost all baryons are protons and neutrons, whose masses differ by 0.14 percent, so the number density is $n_B\approx\rho_b/m_p=11.230\,(\Omega_bh^2)$ per m$^3$. (A correction we do not apply: the fraction $Y_p\approx0.245$ of the baryonic mass is bound in helium-4 nuclei (the observed helium mass fraction; Cyburt et al., cited above). Masses of particles are usually quoted as rest energies $mc^2$ in MeV (mega-electronvolts, $1\ \mathrm{MeV}=1.602177\times10^{-13}$ J). A helium-4 nucleus has $3727.38$ MeV (Tiesinga et al., cited above), that is $931.84$ MeV per nucleon, less than the proton's $938.27$ MeV. Counting the helium nucleons with their own mass multiplies $n_B$ by $(1-Y_p)+Y_p\cdot938.27/931.84=0.755+0.245\times1.00690=1.0017$. This raises the result by about 0.2 percent, and the coefficient of step 5 from $2.7342\times10^{-8}$ to about $2.739\times10^{-8}$. The change is below the 0.45 percent precision of $\Omega_bh^2=0.0224\pm0.0001$ used here, so we keep the simpler value.)
\item \emph{The number of photons.} The microwave background has the spectrum of a black body. For a black body of temperature $T$ the number of photons per unit volume is $n_\gamma=\frac{2\zeta(3)}{\pi^2}\bigl(\frac{k_BT}{\hbar c}\bigr)^3$, where $\zeta(3)=\sum_{n\ge1}n^{-3}=1.2020569$; this standard formula of statistical physics is quoted, not derived, in this book. With $k_BT_0/(\hbar c)=1190.23$ per m and $2\zeta(3)/\pi^2=0.243588$: $n_\gamma=0.243588\times(1190.23)^3=4.1073\times10^8$ per m$^3$, that is about 411 photons per cubic centimetre.
\item \emph{The ratio.} $\eta=11.230\,(\Omega_bh^2)/(4.1073\times10^8)=2.7342\times10^{-8}\,\Omega_bh^2$. For $\Omega_bh^2=0.0224\pm0.0001$ this gives
\end{enumerate}

\[
\eta=6.12\times10^{-10},\qquad\text{between }6.10\times10^{-10}\text{ and }6.15\times10^{-10}.
\]

In words: there are about six baryons for every ten billion photons, or $1/\eta\approx1.6\times10^9$ photons per baryon.

\textbf{Why this is a problem.} Run the expansion backwards. At temperatures at which the thermal energy far exceeds the rest energy of a proton, baryons, antibaryons and photons were all abundant and turned into each other all the time. If the early universe had contained exactly as many baryons as antibaryons, then, as it cooled, almost all of them would have annihilated with each other, and far less matter than we observe would be left (L. Canetti, M. Drewes and M. Shaposhnikov, New J. Phys. 14, 095012 (2012)). The present baryons are therefore the remnant of a small excess of baryons over antibaryons in the early universe. One could simply put this excess into the initial conditions of the universe. The standard view, reviewed in the article just cited, is that it should instead be \textbf{generated by physical processes} from a state without asymmetry; explaining how is the matter--antimatter problem, also called \textbf{baryogenesis}.

\subsection{No antimatter domains}

Could the universe be symmetric on the whole, with some regions made of matter and others of antimatter? Galaxies made of antimatter would shine exactly like galaxies of matter, so their light alone would not reveal them. But wherever a region of matter touches a region of antimatter, particles and antiparticles annihilate. Cohen, De Rújula and Glashow showed that after the universe became transparent, annihilation at the boundaries of such regions cannot be avoided, and that the gamma rays it produces would exceed the observed diffuse gamma-ray background unless the region of matter in which we live is essentially the whole visible universe (A. G. Cohen, A. De Rújula and S. L. Glashow, Astrophys. J. 495, 539 (1998)). A universe made of domains of matter and antimatter, symmetric on the whole, is therefore excluded by observation.

A second observation matters below. Antimatter responds to gravity like matter, as far as tested: antihydrogen atoms released in the ALPHA-g experiment at CERN move downward in the gravity of the Earth, and repulsive "antigravity" is ruled out for them (ALPHA Collaboration, Nature 621, 716 (2023)). The member of mass $-M$ of the pair of universes of Section 17.13 is therefore \textbf{not} antimatter in our universe: its negative mass is a parameter of a second field configuration, and nothing in this book says that antimatter has negative mass or falls upward.

\subsection{Sakharov's three conditions, derived in plain words}

In 1967 Andrei Sakharov identified three ingredients that any process which creates the baryon asymmetry from a state without asymmetry must contain (A. D. Sakharov, JETP Lett. 5, 24 (1967)). Each condition follows from a short argument, which we give now. Throughout, "the system" is the whole universe, it starts in a state with baryon number $B=0$ and without any asymmetry, and we ask what is needed for $B\ne0$ to appear later.

\textbf{Condition 1: baryon-number violation.}

\emph{Proposition 17.1.} If every process conserves $B$, then $B(t)=B(0)$ at every time.

\emph{Proof.} Conservation means $dB/dt=0$, so $B$ is constant. $\square$

Starting from $B=0$, the universe would have $B=0$ forever. So some process must change $B$. In the dirac16complex theory the role of $B$ could only be played by the charge $Q$, and Theorem M1 shows that $Q$ is conserved: this condition fails.

\textbf{Condition 2: C and CP violation.} \textbf{Charge conjugation} C exchanges every particle with its antiparticle. \textbf{Parity} P reflects space, $\mathbf x\to-\mathbf x$. \textbf{CP} does both. A transformation is a \textbf{symmetry of the dynamics} if it turns every possible history of the system into another possible history, running in the same direction of time. The next proposition says that a symmetry that reverses $B$ forbids the creation of an average asymmetry.

\emph{Proposition 17.2.} Let $\mathcal S$ be a symmetry of the dynamics that reverses $B$: if a history $h$ has the baryon number $B[h](t)$ at time $t$, then its image $\mathcal Sh$ has $B[\mathcal Sh](t)=-B[h](t)$. Suppose the initial state is an ensemble (a list of possible initial configurations $c$ with probabilities $p(c)$) in which every configuration and its image are equally likely, $p(\mathcal Sc)=p(c)$. Then the average baryon number $\langle B(t)\rangle$ is zero at every time.

\emph{Proof.} Let $h_c$ be the history that starts from $c$. Because $\mathcal S$ is a symmetry, the history that starts from $\mathcal Sc$ is $\mathcal Sh_c$. The map $c\to\mathcal Sc$ is one-to-one, so a sum over all $c$ may be rewritten as a sum over all $\mathcal Sc$:

\[
\langle B(t)\rangle=\sum_cp(c)\,B[h_c](t)=\sum_cp(\mathcal Sc)\,B[\mathcal Sh_c](t)=\sum_cp(c)\bigl(-B[h_c](t)\bigr)=-\langle B(t)\rangle ,
\]

so $\langle B(t)\rangle=0$. $\square$

Apply this with $\mathcal S=\mathrm C$: an initial state without asymmetry contains particles and antiparticles in equal numbers, so it is C-symmetric, and if C were a symmetry of the dynamics no average asymmetry could ever develop. Apply it with $\mathcal S=\mathrm{CP}$: a homogeneous and isotropic initial state is also symmetric under the reflection P, and P does not change $B$, so CP reverses $B$ as well; if CP were a symmetry, again no asymmetry could develop. Hence \textbf{both C and CP must be violated}. In the dirac16complex theory one of them is exact for each field (Theorem M2): C for dirac16complex00, in every gravitational field, and CP for dirac16complex, in flat space and in every gravitational field in which the reflection is an isometry (a map of the points that leaves the metric unchanged, Section 17.10). The latter include the fields that are homogeneous and isotropic in the space-like directions, the natural setting of the CP step above. So this condition fails too. A generic gravitational field breaks the CP of dirac16complex (Section 17.11), but that cannot produce an asymmetry, because the first condition fails exactly in every gravitational field.

\emph{Worked example (a toy model).} Consider a heavy particle $X$ with two ways to decay: into two quarks ($B=\tfrac23$) with probability $r$, or into an antiquark and an antilepton ($B=-\tfrac13$) with probability $1-r$. Its antiparticle $\bar X$ decays into two antiquarks ($B=-\tfrac23$) with probability $\bar r$, or into a quark and a lepton ($B=\tfrac13$) with probability $1-\bar r$. The decays themselves violate $B$. The average baryon number produced by one $X$ is $\tfrac23r-\tfrac13(1-r)=r-\tfrac13$, by one $\bar X$ it is $-\tfrac23\bar r+\tfrac13(1-\bar r)=\tfrac13-\bar r$, and a population of $N$ particles $X$ and $N$ antiparticles $\bar X$ produces

\[
\Delta B=N\bigl(r-\bar r\bigr).
\]

For $N=1000$, $r=0.51$ and $\bar r=0.49$ this is $\Delta B=20$. If C were a symmetry, the decay $X\to qq$ and its C-image $\bar X\to\bar q\bar q$ would have the same rate, and since $X$ and $\bar X$ have the same total decay rate (by the CPT theorem mentioned below), $r=\bar r$ and $\Delta B=0$; if CP were a symmetry, the same would hold after summing over the directions of the decay products. Only if both are violated can $r\ne\bar r$.

\textbf{Condition 3: departure from thermal equilibrium.} In thermal equilibrium at temperature $T$ each quantum state of energy $E$ is occupied with the Fermi--Dirac probability $f(E)=1/(e^{(E-\mu)/T}+1)$ (Section 12.14; Boltzmann's constant is set to 1), where the chemical potential $\mu$ belongs to a conserved number. For a particle of baryon number $+1$ in a state of energy $E$ the occupation is $f_+=1/(e^{(E-\mu_B)/T}+1)$, and for its antiparticle in the state of the same energy it is $f_-=1/(e^{(E+\mu_B)/T}+1)$: the antiparticle carries the opposite baryon number, so it sees $-\mu_B$. The two energies are equal because particles and antiparticles have the same mass, a consequence of the CPT theorem of ordinary relativistic quantum field theory (quoted, not proved in this book). If reactions that change $B$ run back and forth in equilibrium, $B$ is not a conserved number, and the equilibrium state has no chemical potential for it: the Gibbs state of Section 12.14 is built from the Hamiltonian and the conserved numbers only. Then $\mu_B=0$, $f_+=f_-$ for every energy, and the average baryon number vanishes. An asymmetry can therefore only be produced while the universe is \textbf{out of equilibrium}, for example during a phase transition or in the decay of heavy particles that happens too late for the inverse processes to keep up.

\emph{Worked example.} Take a state with $E=2T$. At $\mu_B=0$, $f_+=f_-=1/(e^2+1)=0.119203$. A chemical potential $\mu_B=0.1\,T$ would give $f_+=1/(e^{1.9}+1)=0.130108$ and $f_-=1/(e^{2.1}+1)=0.109097$, a surplus of $0.021012$ particles per state; but in equilibrium with $B$-violating reactions $\mu_B$ must vanish, and the surplus with it.

These three conditions are the yardstick of the scorecard of Section 17.15.

\subsection{Why the Standard Model is not enough}

The Standard Model of particle physics contains all three ingredients in principle. Baryon number is violated at high temperature by the electroweak \textbf{sphaleron} processes (V. A. Kuzmin, V. A. Rubakov and M. E. Shaposhnikov, Phys. Lett. B 155, 36 (1985)). C is violated maximally by the weak interaction, and CP is violated by the complex phase of the quark mixing matrix. A departure from equilibrium could occur at the electroweak transition in the early universe. The consensus of the literature, reviewed by Canetti, Drewes and Shaposhnikov (New J. Phys. 14, 095012 (2012)), is nevertheless that the Standard Model cannot produce the observed asymmetry, for two reasons.

\begin{enumerate}
\item Lattice computations showed that for Higgs masses of the order of the W-boson mass and above, the electroweak transition is not a first-order phase transition and provides no strong departure from equilibrium (K. Kajantie, M. Laine, K. Rummukainen and M. E. Shaposhnikov, Phys. Rev. Lett. 77, 2887 (1996)). The Higgs boson discovered at the Large Hadron Collider (ATLAS Collaboration, Phys. Lett. B 716, 1 (2012); CMS Collaboration, Phys. Lett. B 716, 30 (2012)) is heavier than the end point found there.
\item The CP violation of the quark mixing matrix produces an asymmetry far too small (M. B. Gavela, P. Hernández, J. Orloff and O. Pène, Mod. Phys. Lett. A 9, 795 (1994); P. Huet and E. Sather, Phys. Rev. D 51, 379 (1995)).
\end{enumerate}

Physics beyond the Standard Model is therefore needed. This section reports the literature; nothing in it is computed in this book.

\subsection{Pairs of universes: global symmetry, local asymmetry}

One class of ideas replaces the question "why is there more matter than antimatter?" by a global symmetry: the universe as a whole is symmetric, but it consists of two parts, each of which is asymmetric in the opposite way. The known published example is the \textbf{CPT-symmetric universe} of Boyle, Finn and Turok (L. Boyle, K. Finn and N. Turok, Phys. Rev. Lett. 121, 251301 (2018)). They propose that the universe after the big bang is the CPT image of the universe before it, so that the epochs before and after the bang form a universe--antiuniverse pair, and they argue that CPT symmetry selects a vacuum state and gives a new interpretation of the cosmological baryon asymmetry. This book neither depends on nor tests that proposal; it is cited as the known example of the class.

The pair of universes of this theory (Section 17.13) is of a different kind. Both members exist at the same time $x_4$, and they are related by the chirality map $\Psi\to\gamma^8\Psi$ together with the change of the mass $m\to-m$, not by a reflection of time through a bang. The resemblance is only structural: a symmetric whole made of two asymmetric parts. Such a picture, even when it is consistent, does not explain the size of the asymmetry in either part, and each part must separately agree with every observation of Sections 17.3 and 17.4.

\subsection{The theory in this chapter: two fields, one current}

\textbf{The two fields.} Both fields of this book are columns $\Psi=(\Psi_0,\dots,\Psi_{15})^T$ of 16 complex components at every point, which carry the irreducible Clifford module of Pin(4,4) (Chapter 2). For \textbf{dirac16complex} the components are anticommuting Grassmann numbers (Section 5.9): $\Psi_a\Psi_b=-\Psi_b\Psi_a$, $\Psi_a\Psi_b^\ast=-\Psi_b^\ast\Psi_a$, and complex conjugation reverses the order of a product, $(\theta_1\theta_2)^\ast=\theta_2^\ast\theta_1^\ast$. For \textbf{dirac16complex00} the components are ordinary commuting complex numbers (Chapter 6). We write

\[
s=-1\ \text{for anticommuting components},\qquad s=+1\ \text{for commuting components},
\]

and call $s$ the \textbf{statistics sign}. The two theories have the same Lagrangian and differ only in $s$.

\textbf{The Lagrangian} (Chapter 6), for both fields:

\[
\mathcal L=\sqrt{|g|}\,\bigl[K-mS-U(S)\bigr],\qquad S=\bar\Psi\Psi,\qquad K=\tfrac12\bigl(\bar\Psi\gamma^\mu D_\mu\Psi-(D_\mu\bar\Psi)\gamma^\mu\Psi\bigr),
\]

with the curved gammas $\gamma^\mu=e_a{}^\mu\gamma^a$, the spinor covariant derivatives $D_\mu\Psi=\partial_\mu\Psi+\Omega_\mu\Psi$ and $D_\mu\bar\Psi=\partial_\mu\bar\Psi-\bar\Psi\Omega_\mu$, and the canonical spinor connection $\Omega_\mu=\tfrac12\omega_{\mu ab}S^{ab}$ (Chapter 4). For anticommuting components $U$ must be a polynomial in $S$ of degree at most 16, because $S^{17}=0$ (Section 5.9); for commuting components $U$ may be any function. We write $\mathcal L_{m,\lambda}$ for the Lagrangian with mass $m$ and $U=\tfrac\lambda2S^2$.

\textbf{The field equations} (Chapter 7) are $E=0$ and $\bar E=0$ with

\[
E:=\gamma^\mu D_\mu\Psi-\bigl(m+U'(S)\bigr)\Psi,\qquad \bar E:=(D_\mu\bar\Psi)\gamma^\mu+\bigl(m+U'(S)\bigr)\bar\Psi .
\]

\textbf{The current and the charge.} The current of the phase symmetry is

\[
j^\mu=\bar\Psi\gamma^\mu\Psi=\Psi^\dagger C\gamma^\mu\Psi .
\]

Each matrix $C\gamma^a$ is real and antisymmetric (expression [1], Section 2.9), hence anti-Hermitian, so $j^\mu$ is anti-Hermitian; for commuting components it is a purely imaginary number (fact 4 of Section 1.5). The Hermitian current of this book is $J^\mu=-ij^\mu=-i\bar\Psi\gamma^\mu\Psi$ (Section 8.12). The matter-antimatter document states its theorems for $j^\mu$; the factor $-i$ changes nothing in them. The \textbf{charge} of the slice $\Sigma_{x_4}=\{x_4=\text{const}\}$ is

\[
Q(x_4)=\int_{\Sigma_{x_4}}\sqrt{|g|}\,J^{x_4}\,d^7x ,
\]

where $J^{x_4}$ is the component of $J^\mu$ along the coordinate $x_4$ and $d^7x=dx_0\,dx_1\,dx_2\,dx_3\,dx_5\,dx_6\,dx_7$. In Gaussian normal gauge ($g^{44}=-1$ and $\gamma^{x_4}=\gamma^4$, Section 8.6) the charge density is

\[
J^4=-i\Psi^\dagger C\gamma^4\Psi=\Psi^\dagger B\Psi,\qquad B=-iC\gamma^4 .
\]

$B$ is Hermitian with $B^2=1$ and has eight eigenvalues $+1$ and eight eigenvalues $-1$ (Section 2.12), so for a classical field $Q$ is an \textbf{indefinite} quadratic form: it can be positive, negative or zero. After quantization $B$ is the Krein metric of the canonical anticommutator $\{\Psi_a,\Psi_b^\dagger\}=B_{ab}\,\delta^7/\sqrt{|g|}$ (Section 8.6), and in the good sector the normal-ordered charge is $Q=\sum\bigl(b^\ast b-d^\ast d\bigr)$: particles carry $Q=+1$ and antiparticles $Q=-1$ (Section 8.12).

\textbf{Test geometries.} The machine checks of this chapter use three exactly given gravitational fields: G1, the generic non-diagonal polynomial vielbein of Stage 1, at three rational points (Wolfram checks); and, built independently from exact vielbein jets (values and derivatives) at one point, G\_A, a non-diagonal vielbein with space--space, space--time and time--time mixing, and G\_B, a diagonal vielbein that depends on $x_0$ and $x_4$ only and has reflection symmetries (Python checks). As always, a check at test points confirms a general proof; the proofs below are the general ones (Section 0.3).

\subsection{Theorem M1: exact U(1) symmetry and charge conservation}

\textbf{Theorem M1.} For both statistics, every admissible potential $U$ and every vielbein:

\begin{enumerate}
\item (global phase) $\mathcal L[e^{i\alpha}\Psi]=\mathcal L[\Psi]$ for every constant real $\alpha$;
\item (local phase and Noether current) for a phase $\alpha(x)$ that depends on the point, $\mathcal L[e^{i\alpha}\Psi]=\mathcal L[\Psi]+i\sqrt{|g|}\,(\partial_\mu\alpha)\,j^\mu$, so the Noether current of the phase symmetry is $N^\mu=i\sqrt{|g|}\,j^\mu$;
\item (off-shell identity) for every configuration, solution or not, $\partial_\mu\bigl(\sqrt{|g|}\,j^\mu\bigr)=\sqrt{|g|}\,\bigl(\bar E\,\Psi+\bar\Psi\,E\bigr)$;
\item (conservation) on every solution $\nabla_\mu j^\mu=0$, and $Q(x_4)$ does not depend on $x_4$ whenever the flux of $\sqrt{|g|}\,J^\mu$ through the boundary of the slices vanishes (fields of compact support, sufficient fall-off at infinity, or slice coordinates identified periodically);
\item (charge density) in Gaussian normal gauge $J^4=\Psi^\dagger B\Psi$, with $B$ Hermitian, $B^2=1$, eight eigenvalues $+1$, eight eigenvalues $-1$ and trace 0.
\end{enumerate}

\textbf{Consequence.} No solution of the field equations of either field, for any $U$ and in any gravitational field with $g^{44}\ne0$, changes $Q$, as long as the boundary flux of item 4 vanishes. The charge of a universe is then fixed by its initial data at $x_4=0$.

\textbf{Proof of item 1.} For commuting components the substitution multiplies $\Psi$ by $e^{i\alpha}$, $\Psi^\dagger$ by $e^{-i\alpha}$ and hence $\bar\Psi=\Psi^\dagger C$ by $e^{-i\alpha}$ ($C$ is real). Every term of $\mathcal L$ is either a bilinear $\Psi^\dagger X\Psi$ with a matrix $X$ built from the geometry (the kinetic term and $S$), which the phase does not touch, or a function of $S$; each bilinear is multiplied by $e^{-i\alpha}e^{i\alpha}=1$. For anticommuting components complex conjugation is antilinear, $(c\theta)^\ast=c^\ast\theta^\ast$, so replacing every generator $\Psi_a$ by $e^{i\alpha}\Psi_a$ and every $\Psi_a^\ast$ by $e^{-i\alpha}\Psi_a^\ast$, and every product of generators by the product of the replaced generators in the same order, is a consistent rule for the whole Grassmann algebra (an automorphism). Every monomial of $\mathcal L$ contains as many factors of the kind $\Psi$ as of the kind $\Psi^\ast$: this is clear for $K$ and $S$, and $S^k$ consists of $\binom{16}k$ monomials with $k$ factors of each kind (Section 5.9). So every monomial is multiplied by $(e^{i\alpha}e^{-i\alpha})^k=1$, and every admissible $U$, a polynomial of degree at most 16 in $S$, is invariant. $\square$

\textbf{Proof of item 2.} With $\alpha=\alpha(x)$ the product rule gives

\[
D_\mu\bigl(e^{i\alpha}\Psi\bigr)=e^{i\alpha}\bigl(D_\mu\Psi+i(\partial_\mu\alpha)\Psi\bigr),\qquad D_\mu\bigl(e^{-i\alpha}\bar\Psi\bigr)=e^{-i\alpha}\bigl(D_\mu\bar\Psi-i(\partial_\mu\alpha)\bar\Psi\bigr),
\]

while $S$ and $U(S)$ do not change. Inserting this into $K$, the phases $e^{\pm i\alpha}$ cancel and

\[
K\ \to\ \tfrac12\Bigl(\bar\Psi\gamma^\mu\bigl(D_\mu\Psi+i\partial_\mu\alpha\,\Psi\bigr)-\bigl(D_\mu\bar\Psi-i\partial_\mu\alpha\,\bar\Psi\bigr)\gamma^\mu\Psi\Bigr)=K+i(\partial_\mu\alpha)\,\bar\Psi\gamma^\mu\Psi=K+i(\partial_\mu\alpha)\,j^\mu .
\]

So $\mathcal L$ changes by $i\sqrt{|g|}\,(\partial_\mu\alpha)j^\mu$, and the coefficient of $\partial_\mu\alpha$ is the Noether current. (It agrees with the Noether formula of Section 5.7, applied with the Grassmann derivative conventions of Section 5.10.) $\square$

For commuting components item 2 already gives conservation by the argument of Section 5.7. If $\alpha$ vanishes near the boundary, $\Psi\to e^{i\alpha}\Psi$ is an allowed variation of the field, so on a solution the change of the action vanishes to first order in $\alpha$: $\int i\sqrt{|g|}\,(\partial_\mu\alpha)j^\mu\,d^8x=0$. Integrating by parts, $\int i\,\alpha\,\partial_\mu(\sqrt{|g|}\,j^\mu)\,d^8x=0$ for every such $\alpha$, and the fundamental lemma of Section 5.2 gives $\partial_\mu(\sqrt{|g|}\,j^\mu)=0$. Item 3 gives more: an identity that holds for every configuration and needs no variational argument, so it holds verbatim for anticommuting components.

\textbf{Proof of item 3.} The product rule (the derivative $\partial_\mu$ is even and needs no Grassmann signs) gives

\[
\partial_\mu\bigl(\sqrt{|g|}\,\bar\Psi\gamma^\mu\Psi\bigr)=\sqrt{|g|}\,(\partial_\mu\bar\Psi)\gamma^\mu\Psi+\bar\Psi\,\partial_\mu\bigl(\sqrt{|g|}\,\gamma^\mu\bigr)\Psi+\sqrt{|g|}\,\bar\Psi\gamma^\mu\partial_\mu\Psi .
\]

The only geometric input is the divergence identity of Section 4.12, a consequence of the covariant constancy of the gammas:

\[
\partial_\mu\bigl(\sqrt{|g|}\,\gamma^\mu\bigr)=\sqrt{|g|}\,\bigl(\gamma^\mu\Omega_\mu-\Omega_\mu\gamma^\mu\bigr).
\]

Inserting it and collecting the terms,

\[
\begin{aligned}
\partial_\mu\bigl(\sqrt{|g|}\,j^\mu\bigr)&=\sqrt{|g|}\,\bigl[(\partial_\mu\bar\Psi-\bar\Psi\Omega_\mu)\gamma^\mu\Psi+\bar\Psi\gamma^\mu(\partial_\mu\Psi+\Omega_\mu\Psi)\bigr]\\
&=\sqrt{|g|}\,\bigl[(D_\mu\bar\Psi)\gamma^\mu\Psi+\bar\Psi\gamma^\mu D_\mu\Psi\bigr].
\end{aligned}
\]

On the other hand, by the definitions of $E$ and $\bar E$,

\[
\bar E\,\Psi+\bar\Psi\,E=(D_\mu\bar\Psi)\gamma^\mu\Psi+\bigl(m+U'(S)\bigr)S+\bar\Psi\gamma^\mu D_\mu\Psi-\bigl(m+U'(S)\bigr)S ,
\]

because $U'(S)$ is an even element (a polynomial in $S$), which commutes with every component, so that $\bar\Psi\,U'(S)\Psi=U'(S)\,S$. The two potential terms cancel, and the two expressions agree. No factor was reordered anywhere, so the argument holds for both statistics. $\square$

\textbf{Proof of item 4.} On a solution $E=\bar E=0$, so $\partial_\mu(\sqrt{|g|}\,j^\mu)=0$, and $\nabla_\mu j^\mu=|g|^{-1/2}\partial_\mu(\sqrt{|g|}\,j^\mu)=0$ by the divergence formula of Section 4.5. Multiply by $-i$: $\partial_\mu(\sqrt{|g|}\,J^\mu)=0$. Now integrate this over the slab between the slices $x_4=t_1$ and $x_4=t_2$, and over a box in the seven other coordinates. By the fundamental theorem of calculus in each variable (Section 1.10), the term with $\mu=4$ gives $Q(t_2)-Q(t_1)$, and the seven other terms give the flux of $\sqrt{|g|}\,J^i$ through the side faces of the box:

\[
Q(t_2)-Q(t_1)=-\int_{t_1}^{t_2}dx_4\oint\sqrt{|g|}\,J^i\,dS_i .
\]

This is pure calculus and holds in any signature. The flux vanishes when the fields vanish near the side faces, when they fall off fast enough as the box grows, or when opposite faces are identified (a torus), because then the contributions of opposite faces cancel. Then $Q(t_2)=Q(t_1)$. That $J^4$ is indefinite plays no role: $Q$ is conserved although it is not positive. $\square$

\textbf{Proof of item 5.} $\sqrt{|g|}\,J^{x_4}=-i\sqrt{|g|}\,\Psi^\dagger C\gamma^{x_4}\Psi$, and in Gaussian normal gauge $\gamma^{x_4}=\gamma^4$ and $-iC\gamma^4=B$. The properties of $B$ are proved in Section 2.12. $\square$

\textbf{The quantized field.} The identity of item 3 is algebraic, and in it every $\Psi^\dagger$ stands to the left of every $\Psi$; it therefore holds for the operator-valued field as well. The Heisenberg equations of the quantized theory are the field equations (Section 8.6), and normal ordering changes $Q$ by a constant. Hence the normal-ordered charge $\sum(b^\ast b-d^\ast d)$ is conserved in the quantized theory too. (This step is derived; it is not a separate machine check.)

\textbf{Worked example: a charge density that can be negative.} Take flat space, $\lambda=0$, and the two rest states of Section 8.8,

\[
u_+=\tfrac12\bigl(e_0-e_4-i\,e_9-i\,e_{13}\bigr),\qquad u_-=\tfrac12\bigl(e_0+e_4+i\,e_9-i\,e_{13}\bigr),
\]

which satisfy $-i\gamma^4u_\pm=u_\pm$, $Bu_+=u_+$ and $Bu_-=-u_-$. Their scalar product is $u_+^\dagger u_-=\tfrac14\bigl(1\cdot1+(-1)\cdot1+i\cdot i+i\cdot(-i)\bigr)=\tfrac14(1-1-1+1)=0$ (the components of $u_+^\dagger$ are the complex conjugates $\tfrac12,-\tfrac12,\tfrac i2,\tfrac i2$ at the places 0, 4, 9, 13). The commuting field $\Psi=(a\,u_++b\,u_-)\,e^{-imx_4}$ with complex numbers $a,b$ solves the free field equation, because $\gamma^4\partial_4\Psi=\gamma^4(-im)\Psi=m(-i\gamma^4)\Psi=m\Psi$. Its charge density is

\[
J^4=\Psi^\dagger B\Psi=|a|^2u_+^\dagger u_+-|b|^2u_-^\dagger u_-=|a|^2-|b|^2 ,
\]

since the cross terms contain $u_+^\dagger Bu_-=-u_+^\dagger u_-=0$. For $a=\tfrac35$, $b=\tfrac45$ it is $\tfrac9{25}-\tfrac{16}{25}=-\tfrac7{25}$, although the field has positive frequency (both modes are eigenvectors of the rest Hamiltonian $h_0=m(-i\gamma^4)$ of Section 8.9 with the positive eigenvalue $m$) and Hilbert density $|a|^2+|b|^2=1$. Its classical energy density is negative too: for a free plane wave of frequency $\omega$ it is $\rho=\omega J^4$ (Section 7.11), here $\rho=mJ^4=-\tfrac7{25}m$. The charge density is constant in $x_4$, as Theorem M1 requires. In the quantized theory, by contrast, a one-particle state built on any normalized positive-energy mode $u$ has the charge $u^\dagger B\,B\,u=u^\dagger u=1$ (the expectation-value rule of Section 8.12).

\textbf{The Kohn--Sham states have a fixed particle number (COMPUTED).} The Kohn--Sham states of Chapters 13 and 14 are computed at a fixed particle number $N$: the Mermin functional contains the constraint $\sum_nf_n=N$ with the chemical potential $\mu$ as its Lagrange multiplier (Section 12.14). For the Stage-4 states of dirac16complex (Chapter 13), in the no-sea convention of Stage 4 the net occupation, particle occupations minus sea holes, is the Kohn--Sham image of the charge $Q$. For the Stage-5 states of dirac16complex00 (Chapter 14) no more is claimed than that they are equilibrium (Mermin) states at fixed $N$: their densities are adopted by prescription, and an honest Gaussian ensemble of classical waves filling the rest shell would have zero charge density (Section 14.8). As a record of data, not a proof, the Wolfram verifier reads the committed Stage-4 runs (check \texttt{MA\_M1\_ks\allowbreak{}FixedNet\allowbreak{}NumberRe\allowbreak{}corded}, measurement \texttt{M1\_ksFix\allowbreak{}edNetNum\allowbreak{}ber} of \texttt{wolfram-\allowbreak{}matter-a\allowbreak{}ntimatte\allowbreak{}r-report\allowbreak{}.json}): 56 reference-solver run files with 168 levels and 33 Rust run files all have the net occupation $N$ to $10^{-9}$ relative, with largest absolute deviations $3.5073\times10^{-10}$ (reference solver) and $1.012\times10^{-11}$ (Rust). This is the only check of the two matter-antimatter reports that uses floating-point numbers.

\textbf{How M1 was verified.} The Wolfram package represents anticommuting components in an exact Grassmann algebra with 1440 odd generators (the values and first and second derivatives of $\Psi$ and $\Psi^\ast$; compare Section 5.10) and commuting components as exact symbols. With an undefined function $U$ the argument of $U$ is identically invariant, so $U(S)$ is invariant for every function (commuting components, flat space and the three G1 points). For anticommuting components at the three G1 points, with $U=\tfrac\lambda2S^2+\tfrac{c_3}3S^3$, the Lagrangian density has 1848 monomials, all of U(1) charge 0, and the explicit phase automorphism leaves it unchanged; the powers $S^k$, $k=1,\dots,17$, have 16, 120, 560, 1820, 4368, 8008, 11440, 12870, 11440, 8008, 4368, 1820, 560, 120, 16, 1 and 0 monomials. The off-shell identity of item 3 holds exactly at the three G1 points (1088 monomials on its left side for anticommuting components). As a \textbf{negative control}, with the notebook's contraction of the spin connection (Section 4.14) the identity fails at all three points, as it must, since the divergence identity fails for that contraction. On shell, with exact solutions of the field equations at the three G1 points for $(m,\lambda)=(3/7,5/11)$, $(-2/5,7/13)$ and $(5/9,-3/8)$, the divergence of the current vanishes, while for generic off-shell data it does not. The independent Python checker verifies, in G\_A and in G\_B and for both statistics, the U(1) invariance with the exact phase $e^{i\alpha}=(3+4i)/5$, the Noether current, the Euler--Lagrange expressions and the off-shell identity as exact polynomial identities in fully generic field jets (Lagrangian densities of 1364 and 1636 monomials in G\_A and 952 and 1224 in G\_B for anticommuting and commuting components), and it proves the divergence identity for arbitrary first vielbein jets: both sides are linear in the 512 numbers $\partial_\lambda e_\nu{}^a$ at a point, and all 512 directions are verified exactly (measurements \texttt{\detokenize{M1}} of \texttt{python-m\allowbreak{}atter-an\allowbreak{}timatter\allowbreak{}-report.\allowbreak{}json}). The main check names are

\begin{Verbatim}[fontsize=\small]
wolfram-matter-antimatter-report.json
  MA_M1_u1InvarianceCommutingGenericU_flat  MA_M1_u1InvarianceGrassmann_G1
  MA_M1_noetherCurrentLocalPhase_grassmann_G1  MA_M1_noetherIdentity_grassmann_G1
  MA_M1_noetherIdentity_commuting_G1  MA_M1_negativeControlNotebookConnection
  MA_M1_onShellConservation_G1  MA_M1_chargeDensityMatrix  MA_M1_stage1ChecksCited
python-matter-antimatter-report.json
  MA_M1_noetherIdentity_<X>_<G>  (X = grassmann, commuting; G = G_A..., G_B...)
  MA_M1_divergenceIdentityAllFirstJets  MA_M1
\end{Verbatim}

\subsection{Theorem M2, part 1: the maps and the transformation rule}

Sakharov's second condition asks whether the theory has symmetries that exchange matter and antimatter. Chapter 6 met the complex conjugations and the reflections of a space-like direction in flat space. This section derives one transformation rule for all candidates of this kind and for both fields. In flat space the rule is a statement about one theory; in a gravitational field it holds in a frame form, which the end of the section explains. Section 17.11 then finds the exact symmetries among the candidates, and its Lemma D shows that no other constant matrix gives one.

\textbf{The maps.} For a subset $R\subseteq\{0,\dots,7\}$ of directions let $r_a=-1$ for $a\in R$ and $r_a=+1$ otherwise, and let $x\mapsto Rx$ be the \textbf{reflection} $x_a\mapsto r_ax_a$ of the coordinates in $R$. Write $s_R=|R\cap\{0,1,2,3\}|$ for the number of reflected space-like directions and $\sigma_R=(-1)^{s_R}$. For a set $A=\{a_1<\dots<a_k\}$ let $\Gamma_A=\gamma^{a_1}\gamma^{a_2}\cdots\gamma^{a_k}$ be its Clifford monomial (Section 2.5), with $\Gamma_\varnothing=1$, and write $R^c$ for the complement of $R$. For a constant invertible $16\times16$ matrix $M$ we consider the \textbf{linear} and the \textbf{antilinear} maps

\[
T^{\mathrm{lin}}_{M,R}:\ \Psi'(x)=M\,\Psi(Rx),\qquad T^{\mathrm{anti}}_{M,R}:\ \Psi'(x)=M\,\Psi^\ast(Rx).
\]

An antilinear map contains a complex conjugation: it turns $c\Psi$ into $c^\ast\Psi'$. For $R=\varnothing$ the antilinear maps are the candidates for \textbf{charge conjugation} C. For $R\ne\varnothing$ the maps with a Pin(4,4) matrix $M$ (Section 2.14) are reflections, called P-type if $R\subseteq\{0,1,2,3\}$ (space is reflected) and T-type if $4\in R$ (the time $x_4$ is reflected), and their combinations with C. The textbook form $\Psi'=M\bar\Psi^T$ of a charge conjugation is included: since $C^T=C$, $\bar\Psi^T=(\Psi^\dagger C)^T=C\Psi^\ast$, so $M\bar\Psi^T=(MC)\Psi^\ast$.

\textbf{Lemma A (the only charge conjugations).} Let $\theta$ be a sign, $+1$ or $-1$ (it has nothing to do with the metric $\eta$). The solutions of $\gamma^aM=\theta\,M(\gamma^a)^\ast$ for all $a$ are the multiples of $1$ for $\theta=+1$ and the multiples of $\gamma^8$ for $\theta=-1$. The solutions of $\gamma^aM=\zeta M(\gamma^a)^T$ for all $a$ are the multiples of $C$ for $\zeta=-1$ and of $\gamma^8C$ for $\zeta=+1$. Hence every constant charge conjugation compatible with the Dirac operator is, up to a factor, one of

\[
\mathrm C_0:\ \Psi\to\Psi^\ast,\qquad \mathrm C_8:\ \Psi\to\gamma^8\Psi^\ast .
\]

(Why this condition: the conjugate of $\gamma^a\partial_a\Psi$ is $(\gamma^a)^\ast\partial_a\Psi^\ast$, and $\gamma^a\partial_a(M\Psi^\ast)$ is a multiple of $M$ times it for every field only if $\gamma^aM=\pm M(\gamma^a)^\ast$.)

\emph{Proof.} The gammas are real, so the first condition reads $\gamma^aM=\theta M\gamma^a$. For $\theta=+1$, $M$ commutes with every $\gamma^a$, and by Corollary 2.3 it is a multiple of $1$. For $\theta=-1$, $\gamma^8$ anticommutes with every $\gamma^a$ (Section 2.10), so $\gamma^8M$ commutes with every $\gamma^a$: $\gamma^a\gamma^8M=-\gamma^8\gamma^aM=\gamma^8M\gamma^a$. Hence $\gamma^8M=c\cdot1$ and $M=c\,\gamma^8$. For the transposed condition, $C$ is symmetric and $C\gamma^a$ antisymmetric, so $(\gamma^a)^TC=(C\gamma^a)^T=-C\gamma^a$, that is $(\gamma^a)^T=-C\gamma^aC^{-1}$. The condition becomes $\gamma^aM=-\zeta MC\gamma^aC^{-1}$, that is $\gamma^a(MC)=-\zeta(MC)\gamma^a$, and the first part, applied to $MC$, gives $MC\propto1$ for $\zeta=-1$ and $MC\propto\gamma^8$ for $\zeta=+1$. With $C^{-1}=C$ and $\gamma^8C=C\gamma^8$ this is the claim. $\square$

\textbf{Lemma B (sign patterns).} For every set $R$, the equation $r_a\,M^{-1}\gamma^aM=\varepsilon\,\gamma^a$ for all $a$, with one common sign $\varepsilon$, has up to a factor exactly two solutions: $M=\Gamma_R$ with $\varepsilon=(-1)^{|R|}$, and $M=\Gamma_{R^c}$ with $\varepsilon=-(-1)^{|R|}$. Both are real orthogonal matrices ($M^\dagger=M^T=M^{-1}$), and

\[
M^\dagger CM=\sigma_R\,C,\qquad \sigma_R=(-1)^{s_R}.
\]

\emph{Proof.} By rule (R3) of Section 2.5, for a monomial $\Gamma_A$ of degree $k=|A|$: $\Gamma_A^{-1}\gamma^a\Gamma_A=(-1)^k\gamma^a$ if $a\notin A$ and $(-1)^{k-1}\gamma^a$ if $a\in A$. For $A=R$ this is $(-1)^{|R|}r_a\gamma^a$, so $r_a\Gamma_R^{-1}\gamma^a\Gamma_R=(-1)^{|R|}\gamma^a$ (use $r_a^2=1$). For $A=R^c$, whose size $8-|R|$ has the parity of $|R|$, the roles of "in" and "not in" are exchanged, and $r_a\Gamma_{R^c}^{-1}\gamma^a\Gamma_{R^c}=-(-1)^{|R|}\gamma^a$. If $M$ and $M'$ solve the same equation with the same sign, then $M^{-1}\gamma^aM=M'^{-1}\gamma^aM'$, so $M'M^{-1}$ commutes with every $\gamma^a$ and is a multiple of $1$ (Corollary 2.3). Every monomial is a product of signed permutation matrices, hence real and orthogonal. Finally, $M^\dagger CM=M^{-1}\gamma^0\gamma^1\gamma^2\gamma^3M=\prod_{a=0}^{3}\bigl(M^{-1}\gamma^aM\bigr)=\prod_{a=0}^{3}\bigl(\varepsilon r_a\gamma^a\bigr)=\varepsilon^4\,r_0r_1r_2r_3\,C=\sigma_RC$. $\square$

\textbf{Lemma C (the statistics sign).} For every matrix $Y$, $\Psi^TY\Psi^\ast=s\,\Psi^\dagger Y^T\Psi$, and the same holds with a derivative on either factor.

\emph{Proof.} $\Psi^TY\Psi^\ast=\sum_{i,j}\Psi_iY_{ij}\Psi^\ast_j=s\sum_{i,j}\Psi_j^\ast Y_{ij}\Psi_i=s\,\Psi^\dagger Y^T\Psi$: exchanging two odd factors costs a sign, exchanging two commuting factors does not. $\square$

\textbf{Theorem M2 (the transformation rule).} Let $M$ be one of the two solutions of Lemma B for the set $R$, with its sign $\varepsilon$ (a factor of modulus 1 in front of $M$ changes nothing). Then, for $U=\tfrac\lambda2S^2$,

\[
\mathcal L_{m,\lambda}[T\Psi](x)=\kappa\,\mathcal L_{\sigma\kappa m,\,\kappa\lambda}[\Psi](Rx),\qquad j'^a(x)=c_j\,r_a\,j^a(Rx),
\]

where

\[
\begin{aligned}
&\text{linear maps:}&&\kappa=\sigma_R\varepsilon,\qquad \sigma=\sigma_R,\qquad c_j=\kappa,\\
&\text{antilinear maps:}&&\kappa=s\,\sigma_R\varepsilon,\qquad \sigma=s\,\sigma_R,\qquad c_j=-\kappa .
\end{aligned}
\]

In words: the kinetic term is multiplied by $\kappa$, the scalar $S$ by $\sigma$, and the current by $c_j$ in addition to the sign $r_a$ that every vector component gets under the reflection.

\emph{Proof for linear maps (flat space).} Now $\partial_a\Psi'(x)=r_aM(\partial_a\Psi)(Rx)$ and $\Psi'^\dagger=\Psi^\dagger M^\dagger$. By Lemma B, $M^\dagger CM=\sigma_RC$ and $M^\dagger C\gamma^aM=(M^{-1}CM)(M^{-1}\gamma^aM)=\sigma_R\varepsilon\,r_a\,C\gamma^a$. Hence $S\to\sigma_RS$ and $j^a\to\sigma_R\varepsilon\,r_aj^a$. In $K$ every term carries one matrix $C\gamma^a$ and one derivative $\partial_a$, so it acquires $\sigma_R\varepsilon\,r_a\cdot r_a=\sigma_R\varepsilon=\kappa$. Therefore, using $\kappa^2=1$ and $(\sigma S)^2=S^2$,

\[
\kappa K-m\sigma S-\tfrac\lambda2S^2=\kappa\Bigl[K-(\sigma\kappa m)S-\tfrac{\kappa\lambda}2S^2\Bigr],
\]

which is $\kappa\,\mathcal L_{\sigma\kappa m,\kappa\lambda}$ (divided by $\sqrt{|g|}=1$). $\square$

\emph{Proof for antilinear maps (flat space).} Now $\Psi'^\dagger=\Psi^TM^\dagger$ (the conjugate of $M\Psi^\ast$ is $M\Psi$ because $M$ is real). By Lemma C and Lemma B,

\[
S'=\Psi^T\bigl(M^\dagger CM\bigr)\Psi^\ast=\sigma_R\,\Psi^TC\Psi^\ast=s\sigma_R\,\Psi^\dagger C^T\Psi=s\sigma_R\,S .
\]

For the kinetic term put $Y=M^\dagger C\gamma^aM=\sigma_R\varepsilon\,r_aC\gamma^a$ and use $(C\gamma^a)^T=-C\gamma^a$:

\[
\begin{aligned}
\Psi'^\dagger C\gamma^a\partial_a\Psi'&=r_a\,\Psi^TY\partial_a\Psi^\ast=r_a\,s\,(\partial_a\Psi^\dagger)Y^T\Psi=-s\sigma_R\varepsilon\,(\partial_a\Psi^\dagger)C\gamma^a\Psi,\\
-(\partial_a\Psi'^\dagger)C\gamma^a\Psi'&=-r_a\,(\partial_a\Psi^T)Y\Psi^\ast=-r_a\,s\,\Psi^\dagger Y^T\partial_a\Psi=s\sigma_R\varepsilon\,\Psi^\dagger C\gamma^a\partial_a\Psi .
\end{aligned}
\]

Adding and halving, $K\to s\sigma_R\varepsilon\,K$. In the same way $j'^a=\Psi^TY\Psi^\ast=s\Psi^\dagger Y^T\Psi=-s\sigma_R\varepsilon\,r_aj^a$. The Lagrangian follows as in the linear case. $\square$

\emph{Curved fields: the frame form.} In a gravitational field the maps are used in \textbf{frame form}: the points are not moved, $\Psi'(x)=M\Psi(x)$ (or $M\Psi^\ast(x)$), and at the same time the frame is changed by the constant signs, $e'_\mu{}^a=r_a\,e_\mu{}^a$ (no sum over $a$). This leaves the metric unchanged, $g'_{\mu\nu}=\sum_ar_a^2\,e_\mu{}^a\eta_{aa}e_\nu{}^a=g_{\mu\nu}$, because $\eta$ is diagonal and $r_a^2=1$, and hence also the Christoffel symbols. The inverse frame changes in the same way, $e'_a{}^\mu=r_a\,e_a{}^\mu$, because $\sum_\mu r_ae_a{}^\mu\,r_be_\mu{}^b=r_ar_b\,\delta_a{}^b=\delta_a{}^b$. So $\gamma'^\mu=\sum_ar_a\,e_a{}^\mu\gamma^a$, and the formula of Section 4.11 for the spin connection gives

\[
\omega'_\mu{}^a{}_b=e'_b{}^\nu\bigl(\Gamma^\rho{}_{\mu\nu}\,e'_\rho{}^a-\partial_\mu e'_\nu{}^a\bigr)=r_b\,e_b{}^\nu\bigl(\Gamma^\rho{}_{\mu\nu}\,r_ae_\rho{}^a-r_a\,\partial_\mu e_\nu{}^a\bigr)=r_ar_b\,\omega_\mu{}^a{}_b ,
\]

and after lowering the index with the diagonal $\eta$, $\omega'_{\mu ab}=r_ar_b\,\omega_{\mu ab}$. Multiplying the equation $r_aM^{-1}\gamma^aM=\varepsilon\gamma^a$ of Lemma B by $M$ from the left and by $M^{-1}$ from the right gives $r_a\gamma^a=\varepsilon M\gamma^aM^{-1}$, that is $M\gamma^aM^{-1}=\varepsilon r_a\gamma^a$ (multiply by $\varepsilon$ and use $\varepsilon^2=1$), hence $MS^{ab}M^{-1}=\tfrac14[\varepsilon r_a\gamma^a,\varepsilon r_b\gamma^b]=r_ar_bS^{ab}$, $\Omega'_\mu=\tfrac12\omega'_{\mu ab}S^{ab}=M\Omega_\mu M^{-1}$ and $D'_\mu(M\Psi)=\partial_\mu(M\Psi)+M\Omega_\mu M^{-1}M\Psi=M\,D_\mu\Psi$. For antilinear maps one also needs $(D_\mu\Psi)^\ast=D_\mu\Psi^\ast$, which holds because $\Omega_\mu$ is real. With these replacements the flat computation goes through unchanged: $\partial_a$ is replaced by $e_a{}^\mu D_\mu$, and the sign $r_a$ that the reflection of the points supplied in flat space is now supplied by the frame change. So the rule holds in every gravitational field in the frame form

\[
\mathcal L_{m,\lambda}[T\Psi;\,e\,\mathrm{diag}(r)](x)=\kappa\,\mathcal L_{\sigma\kappa m,\,\kappa\lambda}[\Psi;\,e](x),
\]

where the second argument names the frame that is used. Read this carefully: it relates the theory on the frame $e\,\mathrm{diag}(r)$ to the theory on the frame $e$. For $R=\varnothing$ the two frames are the same, and the rule is a statement about one theory in every gravitational field. For $R\ne\varnothing$ it is a statement about one theory only when the frame change can be undone; Section 17.11 examines when this happens.

\emph{The coordinate form.} Suppose now that the reflection of the points carries the frame into the reflected frame, $e_\mu{}^a(x)=r_\mu r_a\,e_\mu{}^a(Rx)$ for all $\mu$ and $a$ (no sums). This holds in flat space, and for example for a diagonal frame whose entries do not depend on the reflected coordinates; it makes the reflection an \textbf{isometry}, a map of the points that leaves the metric unchanged, $g_{\mu\nu}(x)=r_\mu r_\nu\,g_{\mu\nu}(Rx)$. Then the inverse frame obeys the same rule, the Christoffel symbols obey $\Gamma^\rho{}_{\mu\nu}(x)=r_\rho r_\mu r_\nu\,\Gamma^\rho{}_{\mu\nu}(Rx)$ (each index of a derivative or of the metric brings its sign), and the formula of Section 4.11 gives $\omega_\mu{}^a{}_b(x)=r_\mu\,r_ar_b\,\omega_\mu{}^a{}_b(Rx)$. For the reflected field $\Psi'(x)=M\Psi(Rx)$ every derivative and every connection term with index $\mu$ brings the sign $r_\mu$, every frame factor $e_a{}^\mu$ brings $r_\mu r_a$, and the two signs $r_\mu$ cancel. What remains is exactly the frame form at the point $Rx$. So in such a field the coordinate form holds for the one theory on the frame $e$: $\mathcal L_{m,\lambda}[T\Psi](x)=\kappa\,\mathcal L_{\sigma\kappa m,\kappa\lambda}[\Psi](Rx)$. $\square$

\subsection{Theorem M2, part 2: C, P and CP for the two fields}

Theorem M2 treats the two matrices of Lemma B. The next lemma shows that no other constant matrix can give an exact symmetry, so that nothing is missed. Here, and in Corollary M2.1, a map is an \textbf{exact symmetry} if $\mathcal L_{m,\lambda}[T\Psi](x)=\mathcal L_{m,\lambda}[\Psi](Rx)$ for every field $\Psi$, in flat space or in a gravitational field in which the reflection is an isometry (Section 17.10).

\textbf{Lemma D (no other matrices).} Let $M$ be any constant invertible $16\times16$ matrix and $R$ any set, and suppose that the linear or the antilinear map with $M$ and $R$ is an exact symmetry in flat space for some $m\ne0$ and some $\lambda$. Then $r_aM^{-1}\gamma^aM=\gamma^a$ for all $a$, so that $M=c\,M_0$, where $M_0$ is the solution of Lemma B with $\varepsilon=+1$ and $c$ a number. Moreover $|c|^2\sigma_R=1$ for the linear map and $|c|^2\sigma_R=s$ for the antilinear map; hence $|c|=1$, and $\sigma_R=1$, respectively $\sigma_R=s$.

\emph{Proof.} Both sides of the defining relation are polynomials in the values $\Psi_i$, $\Psi_i^\ast$ and the first derivatives $\partial_a\Psi_i$, $\partial_a\Psi_i^\ast$ at one point, which can be chosen independently of each other. For anticommuting components different monomials are independent elements of the Grassmann algebra (Section 5.9); for commuting components a polynomial in the real and imaginary parts of complex numbers that vanishes for all their values has only zero coefficients. So the coefficient of each monomial must be the same on both sides. The coefficient of $\Psi_i^\ast\Psi_j$ (no derivative) comes from the mass term, and the coefficient of $\Psi_i^\ast\,\partial_a\Psi_j$ from the first half of $K$, which is $\tfrac12(C\gamma^a)_{ij}$ on the right-hand side.

\emph{Linear map.} $\Psi'^\dagger=\Psi^\dagger M^\dagger$ and $\partial_a\Psi'(x)=r_aM(\partial_a\Psi)(Rx)$. The mass terms give $m\,(M^\dagger CM)_{ij}=m\,C_{ij}$, so $M^\dagger CM=C$ because $m\ne0$. The kinetic terms give $\tfrac12r_a(M^\dagger C\gamma^aM)_{ij}=\tfrac12(C\gamma^a)_{ij}$, so $r_aM^\dagger C\gamma^aM=C\gamma^a$. Multiplying $M^\dagger CM=C$ by $M^{-1}$ from the right gives $M^\dagger C=CM^{-1}$; inserted into the second equation, $r_a\,CM^{-1}\gamma^aM=C\gamma^a$, and multiplying by $C^{-1}=C$ from the left, $r_aM^{-1}\gamma^aM=\gamma^a$.

\emph{Antilinear map.} $\Psi'^\dagger=\Psi^TM^\dagger$. By Lemma C, $S'=\Psi^T(M^\dagger CM)\Psi^\ast=s\,\Psi^\dagger(M^\dagger CM)^T\Psi$, so the mass terms give $s(M^\dagger CM)^T=C$, that is $M^\dagger CM=sC$ ($C$ is symmetric). In $K'$ only the second half, $-(\partial_a\Psi'^\dagger)C\gamma^a\Psi'=-r_a(\partial_a\Psi^T)Y\Psi^\ast$ with $Y=M^\dagger C\gamma^aM$, contains monomials of the kind $\Psi^\ast\,\partial_a\Psi$; by Lemma C it equals $-r_as\,\Psi^\dagger Y^T\partial_a\Psi$. So $-\tfrac12r_as\,Y^T=\tfrac12C\gamma^a$, and transposing with $(C\gamma^a)^T=-C\gamma^a$, $Y=s\,r_a\,C\gamma^a$. With $M^\dagger C=sCM^{-1}$ (from $M^\dagger CM=sC$) this reads $s\,CM^{-1}\gamma^aM=s\,r_aC\gamma^a$, hence again $r_aM^{-1}\gamma^aM=\gamma^a$.

In both cases $M$ solves the equation of Lemma B with $\varepsilon=+1$, whose solutions are the multiples of $M_0$. Finally $M^\dagger CM=|c|^2M_0^\dagger CM_0=|c|^2\sigma_RC$ by Lemma B, and this must equal $C$ (linear map) or $sC$ (antilinear map). $\square$

The same comparison with the signs kept shows a little more (derived here, not a separate machine check): if a map of this kind with $m\ne0$ satisfies $\mathcal L_{m,\lambda}[T\Psi](x)=\kappa\,\mathcal L_{m',\lambda'}[\Psi](Rx)$ with a sign $\kappa$ and some $m',\lambda'$, then $M^\dagger CM=\sigma'C$ with $\sigma'=\kappa m'/m$ for the linear map ($\sigma'\ne0$, because $M^\dagger CM$ is invertible), the kinetic terms give $r_aM^{-1}\gamma^aM=(\kappa/\sigma')\gamma^a$, and squaring both sides ($(M^{-1}\gamma^aM)^2=(\gamma^a)^2$) gives $(\kappa/\sigma')^2=1$. For the antilinear map the same steps give $M^\dagger CM=s\sigma'C$ and again $r_aM^{-1}\gamma^aM=(\kappa/\sigma')\gamma^a$. So $M$ is again a solution of Lemma B, with $\varepsilon=\kappa/\sigma'$, and comparing $M^\dagger CM$ with Lemma B as above shows that the factor in front of it has modulus 1. The 1024 maps counted below therefore contain, up to such factors, every map of this kind whose Lagrangian is $\pm\mathcal L_{m',\lambda'}$.

\textbf{Corollary M2.1 (exact symmetries).} A map of Section 17.10 is an exact symmetry for all $m$ and $\lambda$ if and only if $\kappa=\sigma=1$. It is then a symmetry for every potential $U$, because $K\to K$ and $S\to S$. Explicitly:

\begin{itemize}
\item linear maps: an exact symmetry with the reflected set $R$ exists if and only if $s_R$ is even, for any number of reflected time-like directions; it is the monomial among $\Gamma_R$, $\Gamma_{R^c}$ that has $\varepsilon=+1$;
\item antilinear maps, commuting components: if and only if $s_R$ is even;
\item antilinear maps, anticommuting components: if and only if $s_R$ is odd.
\end{itemize}

\emph{Proof.} By Lemma D an exact symmetry of flat space, whatever its constant matrix, is up to a factor of modulus 1 one of the maps of Theorem M2, so it suffices to examine these (in a gravitational field in which the reflection is an isometry they obey the same rule, Section 17.10); for them the rule gives $\mathcal L_{m,\lambda}\to\kappa\mathcal L_{\sigma\kappa m,\kappa\lambda}$, which equals $\mathcal L_{m,\lambda}$ for all $m$ and $\lambda$ exactly when $\kappa=\sigma=1$. For linear maps, $\sigma=\sigma_R=1$ means $s_R$ even, and then $\kappa=\varepsilon=1$ picks one of the two monomials, which have opposite signs $\varepsilon$ (Lemma B). For antilinear maps, $\sigma=s\sigma_R=1$ means $\sigma_R=s$ ($s_R$ even for $s=+1$, odd for $s=-1$), and then $\kappa=\varepsilon=1$ again picks one monomial. So for each allowed set $R$ there is, up to a factor of modulus 1, exactly one exact map of each type. The Lagrangian relation equates the densities at the points $x$ and $Rx$; a reflection does not change volumes, so the action is unchanged, and solutions are mapped to solutions. $\square$

No map of this kind sends $\mathcal L_{m,\lambda}$ to $-\mathcal L_{m,\lambda}$ when $\lambda\ne0$: that would need $\kappa=-1$ together with $\sigma\kappa=1$ and $\kappa\lambda=\lambda$, which contradict each other. For $\lambda=0$ the maps with $\kappa=\sigma=-1$ send $\mathcal L_{m,0}$ to $-\mathcal L_{m,0}$, and they are then symmetries of the field equations, which are the same for $\mathcal L$ and $-\mathcal L$.

\textbf{Corollary M2.2 (the charge).} For an exact symmetry $c_j=+1$ if the map is linear and $c_j=-1$ if it is antilinear. Hence, if $x_4$ is not reflected ($r_4=+1$), an exact linear symmetry preserves $J^4$ and $Q$, and an exact antilinear symmetry reverses them, $Q\to-Q$.

\emph{Proof.} $c_j=\kappa=1$ for linear and $c_j=-\kappa=-1$ for antilinear exact maps. With $r_4=+1$ the rule gives $j'^4(x)=c_j\,j^4(Rx)$, and integrating over a slice, whose coordinates are merely reflected, gives $Q'=c_jQ$. $\square$

\textbf{Counting.} For each statistics there are $2^8=256$ sets $R$, 2 matrices and 2 types, so 1024 maps. Their Lagrangian images fall into the four classes $\mathcal L_{m,\lambda}$, $\mathcal L_{-m,\lambda}$, $-\mathcal L_{m,-\lambda}$ and $-\mathcal L_{-m,-\lambda}$ with 256 maps each. The exact symmetries number 256 for each statistics: 128 linear ($R\cap\{0,1,2,3\}$ of even size, 8 choices, times any subset of $\{4,5,6,7\}$, 16 choices) and 128 antilinear. The exact symmetries that reverse $Q$ without reflecting $x_4$ are the exact antilinear maps with $4\notin R$: 8 choices of $R\cap\{0,1,2,3\}$ with the right parity times 8 subsets of $\{5,6,7\}$, that is \textbf{64 for each statistics}.

\textbf{Worked example: three rows from the rule.} (i) $\mathrm C_0$ on commuting components: $R=\varnothing$, $M=1=\Gamma_\varnothing$, so $\varepsilon=(-1)^0=+1$, $\sigma_R=+1$ and $s=+1$; hence $\kappa=\sigma=+1$ and $c_j=-1$. The map is exact and reverses the charge. (ii) The same map on anticommuting components: now $s=-1$, so $\kappa=\sigma=-1$ and $\mathcal L_{m,\lambda}\to-\mathcal L_{(-1)(-1)m,\,-\lambda}=-\mathcal L_{m,-\lambda}$: not a symmetry. (iii) $\mathrm C_8\mathrm P_1$ on anticommuting components, the map $\Psi'(x)=\Gamma_{\{0,2,3,4,5,6,7\}}\Psi^\ast(R_1x)$ that reflects $x_1$: $R=\{1\}$ and $M=\Gamma_{R^c}$, so $\varepsilon=-(-1)^1=+1$, $s_R=1$, $\sigma_R=-1$ and $s=-1$; hence $\kappa=(-1)(-1)(+1)=+1$, $\sigma=(-1)(-1)=+1$ and $c_j=-1$. The map is exact and reverses the charge. (In $\gamma^8\gamma^1=\gamma^0\gamma^1\gamma^2\cdots\gamma^7\gamma^1$ move the last factor to the left past the six factors $\gamma^7,\dots,\gamma^2$, which gives the sign $(-1)^6=+1$, and use $(\gamma^1)^2=1$: so $\gamma^8\gamma^1=\Gamma_{\{0,2,3,4,5,6,7\}}$, and the map is $\mathrm C_8$ applied after the reflection $\Psi\to\gamma^1\Psi(R_1x)$.)

A direct way to see (i): the gammas are real and $S=\Psi^\dagger C\Psi$ is real for commuting components, so the complex conjugate of the field equation $\gamma^\mu D_\mu\Psi=(m+\lambda S)\Psi$ is the same equation for $\Psi^\ast$. Every solution $\Psi$ has the partner solution $\Psi^\ast$, and its charge density is $(\Psi^\ast)^\dagger B\Psi^\ast=\Psi^TB\Psi^\ast=\Psi^\dagger B^T\Psi=-\Psi^\dagger B\Psi$, because $B^T=-B$ ($B=-iC\gamma^4$ with $C\gamma^4$ antisymmetric).

\textbf{C, P, CP, T and CPT for the two fields.} The table applies the rule to the maps with names. $R_Ax$ is $x$ with the coordinates $x_a$, $a\in A$, reversed. "exact" means $\mathcal L\to\mathcal L$; the effect on $Q$ is given for the maps that do not reflect $x_4$ and whose effect on the charge the matter-antimatter document records. $\mathrm P_b$ is the Pin(4,4) lift $\gamma^b$ of the reflection of $x_b$ (Section 2.14).

\begingroup
\small
\begin{longtable}{@{}>{\raggedright\arraybackslash}p{0.200\linewidth}>{\raggedright\arraybackslash}p{0.200\linewidth}>{\raggedright\arraybackslash}p{0.200\linewidth}>{\raggedright\arraybackslash}p{0.200\linewidth}@{}}
\toprule
\textbf{Map} & \textbf{Definition} & \textbf{dirac16complex00 (commuting)} & \textbf{dirac16complex (anticommuting)} \\
\midrule
\endfirsthead
\toprule
\textbf{Map} & \textbf{Definition} & \textbf{dirac16complex00 (commuting)} & \textbf{dirac16complex (anticommuting)} \\
\midrule
\endhead
$\mathrm C_0$ & $\Psi^\ast(x)$ & exact, $Q\to-Q$ & $-\mathcal L_{m,-\lambda}$, $Q\to Q$ \\
$\mathrm C_8$ & $\gamma^8\Psi^\ast(x)$ & $-\mathcal L_{-m,-\lambda}$, $Q\to Q$ & $\mathcal L_{-m,\lambda}$, $Q\to-Q$ \\
chirality & $\gamma^8\Psi(x)$ & $-\mathcal L_{-m,-\lambda}$, $Q\to-Q$ & $-\mathcal L_{-m,-\lambda}$, $Q\to-Q$ \\
$\mathrm P_b$, $b\le3$ & $\gamma^b\Psi(R_bx)$ & $\mathcal L_{-m,\lambda}$ & $\mathcal L_{-m,\lambda}$ \\
$\mathrm C_0\mathrm P_b$ & $\gamma^b\Psi^\ast(R_bx)$ & $\mathcal L_{-m,\lambda}$ & $-\mathcal L_{-m,-\lambda}$ \\
$\mathrm C_8\mathrm P_b$ & $\Gamma_{\{b\}^c}\Psi^\ast(R_bx)$ & $-\mathcal L_{m,-\lambda}$ & exact, $Q\to-Q$ \\
$\mathrm C_8\mathrm P_{123}$ & $\Gamma_{\{0,4,5,6,7\}}\Psi^\ast(R_{123}x)$ & $-\mathcal L_{m,-\lambda}$ & exact, $Q\to-Q$ \\
$\mathrm P_{0123}$ & $C\Psi(R_{0123}x)$ & exact, $Q\to Q$ & exact, $Q\to Q$ \\
$\mathrm C_0\mathrm P_{0123}$ & $C\Psi^\ast(R_{0123}x)$ & exact, $Q\to-Q$ & $-\mathcal L_{m,-\lambda}$ \\
T, linear & $\Gamma_{\{4\}^c}\Psi(R_4x)$ & exact & exact \\
T, antilinear & $\Gamma_{\{4\}^c}\Psi^\ast(R_4x)$ & exact & $-\mathcal L_{m,-\lambda}$ \\
total inversion & $\gamma^8\Psi(-x)$ & exact & exact \\
CPT & $\gamma^8\Psi^\ast(-x)$ & exact & $-\mathcal L_{m,-\lambda}$ \\
$\mathrm C\mathrm P_{123}\mathrm T$ & $\gamma^1\gamma^2\gamma^3\gamma^4\Psi^\ast(R_{1234}x)$ & $-\mathcal L_{m,-\lambda}$ & exact \\
\bottomrule
\end{longtable}
\endgroup

Reading of the table:

\begin{itemize}
\item \textbf{dirac16complex00 (commuting field).} C is exact and reverses the charge. P and CP with one (or three) reflected space-like directions are not exact for $m\ne0$; P is exact at $m=0$. CP with an even number of reflected space-like directions, for example $\mathrm C_0\mathrm P_{0123}$, is exact. T (linear and antilinear) and the full antilinear inversion CPT are exact.
\item \textbf{dirac16complex (anticommuting field).} No constant C is exact for $m\ne0$: $\mathrm C_0$ gives $-\mathcal L_{m,-\lambda}$, and $\mathrm C_8$ gives $\mathcal L_{-m,\lambda}$, the theory with the opposite mass (so C is exact at $m=0$). CP with an odd number of reflected space-like directions is exact and reverses the charge: $\mathrm C_8\mathrm P_b$ for $b=0,1,2,3$ and $\mathrm C_8\mathrm P_{123}$. T is exact as a linear map (after quantization it is of antiunitary type, see below). The antilinear total inversion is not exact, while $\mathrm C\mathrm P_{123}\mathrm T$ is.
\end{itemize}

\textbf{Which entries hold in a gravitational field.} The table is computed in flat space. The rows with $R=\varnothing$, the maps $\mathrm C_0$, $\mathrm C_8$ and the chirality map, move neither the points nor the frame, so by Section 17.10 their entries hold for one theory in every gravitational field; in particular $\mathrm C_0$ is an exact, charge-reversing symmetry of dirac16complex00 in every gravitational field. Every other row reflects points. In a gravitational field in which the reflection is an isometry, such as the primordial field of Chapter 9 for the reflections of $x_1$, $x_2$ and $x_3$ (its metric is diagonal and depends only on $x_0$ and $x_4$, Section 9.2), the coordinate form of Section 17.10 holds and the entry is unchanged. In a generic gravitational field only the frame form holds, and it relates the theory on the frame $e$ to the theory on the frame $e\,\mathrm{diag}(r)$. For the CP maps of the anticommuting field these are really two different theories, for two reasons.

\begin{enumerate}
\item \emph{No frame rotation undoes the frame change.} $\mathcal L$ is invariant under the frame rotations of $\mathrm{Spin}_0(4,4)$ (Section 2.14), whose vector matrices $\Lambda$ are products of exponentials and can therefore be joined to the unit matrix by letting all angles go to 0. Write $\Lambda$ in blocks, with $A$ the $4\times4$ block of the space-like directions and $G$ the block that maps them to the time-like ones. The space-like block of the equation $\Lambda^T\eta\Lambda=\eta$ reads $A^TA-G^TG=I_4$, so $A^TA=I_4+G^TG$, whose eigenvalues are at least 1, and $(\det A)^2=\det(A^TA)\ge1$. So $\det A$ never lies between $-1$ and $1$; it equals 1 at the unit matrix and changes continuously with the angles, so it is at least 1 for every such $\Lambda$. The matrix $\mathrm{diag}(r)$ of a CP map has $\det A=(-1)^{s_R}=-1$ ($s_R$ is odd), so it is not a frame rotation of this kind.
\item \emph{No other exact symmetry undoes it either.} The only matrices that turn the gammas as $\mathrm{diag}(r)$ does are its two linear Pin lifts $\Gamma_R$ and $\Gamma_{R^c}$ (Lemma B), and by Corollary M2.1 neither is an exact symmetry when $s_R$ is odd.
\end{enumerate}

At a fixed frame, the internal part $\Psi\to\Gamma_{R^c}\Psi^\ast$ of such a map multiplies the kinetic terms of the different frame directions by different signs, and nothing compensates them: it turns $\mathcal L$ into none of $\pm\mathcal L_{\pm m,\pm\lambda}$. \textbf{So CP of dirac16complex is exact in flat space and in every gravitational field in which the reflection is an isometry, and a generic gravitational field breaks it.} This opens no way to an asymmetry: Theorem M1 holds in every gravitational field, so the charge stays conserved and Sakharov's first condition fails exactly. The committed exact-results file states this scope (\texttt{matter-a\allowbreak{}ntimatte\allowbreak{}r-theory\allowbreak{}.json}, key \texttt{\detokenize{M2}}, entries \texttt{answer.g\allowbreak{}rassmann} and \texttt{answer.s\allowbreak{}akharov2}), and the committed Python checker verifies it in the check \texttt{MA\_M2\_cp\allowbreak{}ScopeInC\allowbreak{}urvedFie\allowbreak{}lds}: at a fixed frame of the generic field G\_A the internal parts of $\mathrm C_8\mathrm P_0,\dots,\mathrm C_8\mathrm P_3$ and $\mathrm C_8\mathrm P_{123}$ give none of $\pm\mathcal L_{\pm m,\pm\lambda}$, while with the reflection the same matrices give the exact, charge-reversing maps; all $8\times16=128$ sign patterns with an odd number of reflected space-like directions have $\det A=-1$; and neither linear Pin lift of any of them is exact. (This check is not in the committed Python report, which was written by an earlier version of the checker; Section 17.17.) The same limitation applies to every table row that reflects points: it holds in flat space and wherever the reflection is an isometry.

\textbf{For both fields an exact symmetry exists that preserves $x_4$ and reverses $Q$:} $\mathrm C_0$ for the commuting field, in every gravitational field, and $\mathrm C_8\mathrm P_b$ for the anticommuting field, in flat space and in gravitational fields in which the reflection of $x_b$ is an isometry. By Proposition 17.2, such a symmetry alone forbids the creation of an average charge from a symmetric state.

\textbf{The quantized anticommuting field.} After quantization (Chapter 8), a symmetry must be carried by an operator on the space of states: a \textbf{unitary} operator preserves inner products (Section 8.2); an \textbf{antiunitary} operator preserves them up to complex conjugation and complex-conjugates every number it passes, which is how time reversal is implemented in ordinary quantum mechanics. A map of the field operators can be implemented by such an operator only if it preserves the canonical anticommutator $\{\Psi_a,\Psi_b^\dagger\}=B_{ab}\,\delta^7$ (flat space, Gaussian normal gauge). A unitary operator leaves numbers alone, so the transformed operators must reproduce the number matrix $B$; an antiunitary operator complex-conjugates numbers, so they must reproduce $B^\ast=-B$. We call a map of \textbf{unitary type} if it passes the first test (it is then a linear automorphism of the anticommutation relations) and of \textbf{antiunitary type} if it passes the second (an antilinear automorphism). Passing the test is necessary for an implementation, not sufficient; whether an actual operator exists is the OPEN question at the end of this paragraph. Lemma B gives $M\gamma^aM^{-1}=\varepsilon r_a\gamma^a$ (multiply $r_aM^{-1}\gamma^aM=\varepsilon\gamma^a$ by $\varepsilon r_aM$ from the left and by $M^{-1}$ from the right) and hence, multiplying four of these, $MCM^{-1}=\sigma_RC$, so that $MBM^{-1}=\sigma_R\varepsilon\,r_4\,B$; and $B^T=-B$. For a linear map $\Psi\to M\Psi$ and for the antilinear map $\Psi\to M\Psi^{\dagger T}$ (the operator version of $M\Psi^\ast$; its adjoint is $M\Psi$ because $M$ is real):

\[
\begin{aligned}
&\{M\Psi,(M\Psi)^\dagger\}=MBM^\dagger=\sigma_R\varepsilon\,r_4\,B,\\
&\{(M\Psi^{\dagger T})_a,\ (M\Psi)_b\}=\sum_{c,d}M_{ac}M_{bd}\{\Psi^\dagger_c,\Psi_d\}=\bigl(MB^TM^\dagger\bigr)_{ab}=-\sigma_R\varepsilon\,r_4\,B_{ab} .
\end{aligned}
\]

For an exact symmetry of the anticommuting field ($\sigma_R\varepsilon=1$ for linear maps; $\sigma_R=-1$ and $\varepsilon=1$ for antilinear maps) both right-hand sides equal $r_4B$. Hence every exact symmetry is compatible with the canonical structure: it is of unitary type if it preserves $x_4$ and of antiunitary type if it reverses $x_4$ (derived; the Wolfram verifier confirms it for all 256 exact maps, check \texttt{MA\_M2\_ca\allowbreak{}nonicalS\allowbreak{}tructure}). The charge conjugation of unitary type is $\mathrm C_8$ ($\gamma^8B^T\gamma^8=-\gamma^8B\gamma^8=B$), whereas $\mathrm C_0$ is of antiunitary type although it preserves $x_4$ ($1\cdot B^T\cdot1=-B$). For the CP of unitary type, $\mathrm C_8\mathrm P_b$ with $M=\Gamma_{\{b\}^c}$, the charge density becomes

\[
J'^4=\sum_{a,d}\bigl(M\Psi\bigr)_aB_{ad}\bigl(M\Psi^{\dagger T}\bigr)_d=\sum_{c,e}\Psi_c\,\bigl(M^TBM\bigr)_{ce}\,\Psi^\dagger_e .
\]

Reordering the operators with the anticommutator, $\Psi_c\Psi^\dagger_e=-\Psi^\dagger_e\Psi_c+B_{ce}\delta^7$, and using $(M^TBM)^T=M^TB^TM=-M^TBM$, gives

\[
J'^4=-\Psi^\dagger(M^TBM)^T\Psi+(\text{a number})=\Psi^\dagger\,M^TBM\,\Psi+(\text{a number}).
\]

By Lemma B, $M^TBM=M^{-1}BM=\sigma_R\varepsilon\,r_4\,B=-B$ for this map. So $J'^4=-J^4$ up to a constant, and the normal-ordered charge is reversed, $:\!Q\!:\ \to-:\!Q\!:$ (derived). \textbf{OPEN:} whether each of these maps of unitary or antiunitary type is actually implemented by a unitary or antiunitary operator on the positive ($J=B$) Fock space of the good sector (Section 8.10) is not decided in the repository (the committed reports say the same: \texttt{MA\_M2\_ca\allowbreak{}nonicalS\allowbreak{}tructure} of the Wolfram report and \texttt{MA\_M2\_ca\allowbreak{}nonicalS\allowbreak{}tructure\allowbreak{}OfExactG\allowbreak{}rassmann\allowbreak{}Symmetri\allowbreak{}es} of the Python report).

\textbf{A remark on the slicing in signature (4,4).} The group $\mathrm{Spin}_0(4,4)$ of products of exponentials $\exp(\theta S^{ab})$ (Section 2.14), under which $\mathcal L$ is invariant, contains the element

\[
\exp(\pi S^{45})=\cos\tfrac\pi2+\gamma^4\gamma^5\sin\tfrac\pi2=\gamma^4\gamma^5 ,
\]

because $(\gamma^4\gamma^5)^2=-(\gamma^4)^2(\gamma^5)^2=-1$ makes the pair $(4,5)$ of two time-like directions a "rotation" pair (Section 2.11). It rotates the plane of the times $x_4$ and $x_5$ by the angle $\pi$, $(x_4,x_5)\to(-x_4,-x_5)$. In flat space, as a linear map with $R=\{4,5\}$ and $M=\Gamma_R$, it has $\varepsilon=(-1)^2=+1$ and $s_R=0$, so it is exact, and $c_j=+1$ with $r_4=-1$ gives $J'^4(x)=-J^4(Rx)$: it maps every classical solution of charge $Q$ into one of charge $-Q$. The product of four such rotations, in the planes $(0,1)$, $(2,3)$, $(4,5)$ and $(6,7)$, is $\gamma^8$, the total inversion, which is exact as well. In a world with one time the direction of time cannot be reversed by a continuous rotation; with two or more times it can. In signature (4,4) the time orientation of the slicing $x_4=\text{const}$ can therefore be reversed continuously, and the sign of the \textbf{classical} charge $Q$ refers to a chosen slicing.

After quantization of the anticommuting field this changes. Both maps reverse $x_4$, so they are of antiunitary type: for $M=\gamma^4\gamma^5$ the anticommutator gives $MBM^\dagger=\sigma_R\varepsilon\,r_4\,B=-B$, and the same holds for $M=\gamma^8$ ($R=\{0,\dots,7\}$, $\varepsilon=(-1)^8=1$, $s_R=4$, $r_4=-1$). An antiunitary operator $A$ with $A\Psi A^{-1}=M\Psi(Rx)$ complex-conjugates the numbers $B_{cd}$ in $J^4=\sum_{c,d}\Psi^\dagger_cB_{cd}\Psi_d$, so

\[
AJ^4(x)A^{-1}=\Psi^\dagger M^TB^\ast M\Psi\,(Rx)=+J^4(Rx),
\]

because $B^\ast=-B$ and $M^TBM=M^{-1}BM=\sigma_R\varepsilon\,r_4\,B=-B$ (Lemma B, as in the computation of $MBM^{-1}$ above). The quantum charge is preserved, as it is under time reversal in ordinary quantum mechanics (derived; the committed Python checker records $M^TB^\ast M=+B$ for both maps in the entry \texttt{facts.qu\allowbreak{}antisedG\allowbreak{}rassmann\allowbreak{}Field} of its measurement for \texttt{MA\_M2\_sp\allowbreak{}in0Conta\allowbreak{}insCharg\allowbreak{}eReversi\allowbreak{}ngTimeRo\allowbreak{}tation}, an entry that the committed report, written by an earlier version of the checker, does not contain). These maps reverse $x_4$; they are not C or CP in Sakharov's sense, and the scorecard of Section 17.15 does not use them.

\textbf{How M2 was verified.} The Wolfram verifier finds the intertwiner spaces of Lemma A one-dimensional each (bases $1$, $\gamma^8$, $C$, $\gamma^8C$); the 256 monomials realise the 256 sign patterns bijectively; the statistics sign $s$ is computed with an exact random matrix $Y$; for all 1024 maps per statistics the Lagrangian density computed from the transformed field jets (commuting symbols, respectively the Grassmann algebra) equals $\kappa\mathcal L_{\sigma\kappa m,\kappa\lambda}$ as the rule predicts, with the eight current signs, without a single mismatch; the frame form in the curved field G1 is checked for 31 frame reflections for both statistics; and all 256 exact maps of the anticommuting field preserve the canonical structure. The independent Python checker verifies the solution spaces of Lemma A, the internal maps on curved jets in G\_A, named reflections, three generic non-axis unit vectors, reflections in the curved field G\_B, the rotation $\exp(\pi S^{45})$ and the total inversion, and the complete character table of the 1024 maps per statistics (4 classes of 256, 256 exact symmetries, 64 exact charge-reversing symmetries without $x_4$ reversal). The two implementations agree on all 2048 classification rows and on all 22 items of the C, P, CP, T and CPT summary (check \texttt{MA\_agree\allowbreak{}sWithWol\allowbreak{}fram}). The Python checker of the same commit adds the check \texttt{MA\_M2\_cp\allowbreak{}ScopeInC\allowbreak{}urvedFie\allowbreak{}lds} described above, which the committed Python report does not yet contain (Section 17.17). Lemma D and the remark after it are derived in this chapter and are not separate machine checks. The main check names are

\begin{Verbatim}[fontsize=\small]
wolfram-matter-antimatter-report.json
  MA_M2_conjugationIntertwiners  MA_M2_signPatternClassification  MA_M2_statisticsSign
  MA_M2_lagrangianFlatCommuting_all  MA_M2_lagrangianFlatGrassmann_all
  MA_M2_frameLevelG1_commuting  MA_M2_frameLevelG1_grassmann
  MA_M2_symmetrySummaryAndChargeReversal  MA_M2_canonicalStructure
python-matter-antimatter-report.json
  MA_M2_chargeConjugationSolutionSpaces  MA_M2_spin0ContainsChargeReversingTimeRotation
  MA_M2_discreteGroupCharacterTable  MA_M2_C_and_CP_status
  MA_M2_canonicalStructureOfExactGrassmannSymmetries  MA_M2
checker of this commit only (not in the committed Python report)
  MA_M2_cpScopeInCurvedFields
\end{Verbatim}

\subsection{Theorem M3: which charge-violating terms the symmetry allows}

Sakharov's first condition needs an interaction that changes the charge. The simplest candidates are \textbf{Majorana-type} terms, built with $\Psi^T$ in place of $\Psi^\dagger$, such as $\Psi^TM\Psi$ with a constant matrix $M$ (named after Ettore Majorana, who used such terms for neutral fermions). Under $\Psi\to e^{i\alpha}\Psi$ both factors are multiplied by $e^{i\alpha}$, so the term is multiplied by $e^{2i\alpha}$: it has \textbf{U(1) charge 2} and violates the symmetry of Theorem M1. Which of them does the local Lorentz symmetry of the theory allow? Theorem M3 answers this. It is a classification of what the symmetry permits; it is not a claim that any such term is present in the theory, or natural.

\textbf{The condition for invariance.} Under a spin transformation $\Psi\to R\Psi$ (Section 2.14) the term becomes $\Psi^TR^TMR\Psi$. For $R=\exp(\theta S^{ab})=1+\theta S^{ab}+O(\theta^2)$,

\[
R^TMR=M+\theta\bigl((S^{ab})^TM+MS^{ab}\bigr)+O(\theta^2).
\]

So invariance under all of $\mathrm{Spin}_0(4,4)$ requires $(S^{ab})^TM+MS^{ab}=0$ for all 28 generators. Conversely this condition suffices: with $R(\theta)=\exp(\theta S)$, $dR/d\theta=SR=RS$, so $\frac d{d\theta}\bigl(R^TMR\bigr)=R^T\bigl(S^TM+MS\bigr)R=0$, hence $R^TMR=M$ for every $\theta$, and for every product of such exponentials.

\textbf{Theorem M3.}

\begin{enumerate}
\item (invariant forms) The constant matrices $M$ with $(S^{ab})^TM+MS^{ab}=0$ for all 28 generators are exactly $M=C(\alpha P_-+\beta P_+)$ with complex numbers $\alpha,\beta$, where $P_\mp=\tfrac12(1\mp\gamma^8)$. This space has dimension 2, all its elements are symmetric, and it contains no nonzero antisymmetric matrix.
\item (the full group) The elements of $\mathrm{Spin}(4,4)$ of spinor norm $-1$ (Section 2.14), such as $g=\gamma^a\gamma^b$ with $a$ space-like and $b$ time-like, act by $g^T(CP_\pm)g=-CP_\pm$. No nonzero form is strictly invariant under the full $\mathrm{Spin}(4,4)$; both basis forms carry the spinor-norm character.
\item (Pin characters) For every unit vector $u=u_a\gamma^a$ with $u^2=n(u)=\pm1$: $u^TCu=-n(u)\,C$ and $u^T(C\gamma^8)u=+n(u)\,C\gamma^8$, while $u^T(CP_\pm)u=-n(u)\,CP_\mp$. So $C$ carries the character $-n(u)$ (the character of $\bar\Psi\Psi$), $C\gamma^8$ the character $+n(u)$, and the chiral forms $CP_\pm$ are $\mathrm{Spin}_0$-invariant but not Pin-covariant, because $u$ exchanges the two chiralities.
\item (derivative terms) The invariant derivative bilinears $\Psi^TM\gamma^a\partial_a\Psi$ have $M$ in the same two-dimensional space. $M\gamma^a$ is symmetric for every $a$ exactly when $M\propto C\gamma^8$, and antisymmetric for every $a$ exactly when $M\propto C$.
\item (what survives) For \textbf{anticommuting} components every invariant mass-type term $\Psi^TM\Psi$ vanishes identically: \textbf{no Majorana mass term exists}. The derivative term $\sqrt{|g|}\,\Psi^TC\gamma^\mu D_\mu\Psi$ is a total divergence, while $\sqrt{|g|}\,\Psi^TC\gamma^8\gamma^\mu D_\mu\Psi$ survives, with the flat-space Euler--Lagrange expression $2C\gamma^8\gamma^a\partial_a\Psi$. For \textbf{commuting} components the two chiral mass terms $\Psi^TCP_\pm\Psi$ survive (equivalently $\Psi^TC\Psi$ and $\Psi^TC\gamma^8\Psi$), and so does the derivative term $\sqrt{|g|}\,\Psi^TC\gamma^\mu D_\mu\Psi$, which has the form of the notebook's Lagrangian Lg[]; $\sqrt{|g|}\,\Psi^TC\gamma^8\gamma^\mu D_\mu\Psi$ is then a total divergence.
\item (charge) Every term of items 1 to 5 has U(1) charge 2, and none of them is contained in $\mathcal L$.
\end{enumerate}

\emph{Proof of item 1.} From $(CS^{ab})^T=-CS^{ab}$ (Section 2.11) and $C^T=C$ we get $(S^{ab})^TC=-CS^{ab}$, that is $(S^{ab})^T=-CS^{ab}C^{-1}$. The condition becomes $-CS^{ab}C^{-1}M+MS^{ab}=0$; multiplying by $C^{-1}$ from the left, $S^{ab}(C^{-1}M)=(C^{-1}M)S^{ab}$: the matrix $C^{-1}M$ commutes with all 28 generators. By Theorem 2.12 and the remark after it (it holds for $\mathrm{Spin}_0(4,4)$), this commutant is spanned by $P_-$ and $P_+$. So $M=C(\alpha P_-+\beta P_+)$. The chirality $\gamma^8=\mathrm{diag}(-I_8,I_8)$ is symmetric and commutes with $C$, so $(CP_\pm)^T=P_\pm^TC^T=P_\pm C=CP_\pm$: both basis matrices are symmetric, and so is every combination. $\square$

\emph{Proof of item 3.} From $(\gamma^a)^TC=-C\gamma^a$ and linearity, $u^TC=-Cu$, so $u^TCu=-Cu\,u=-n(u)\,C$. Since $u$ anticommutes with $\gamma^8$: $u^TC\gamma^8u=-Cu\gamma^8u=C\gamma^8u\,u=n(u)\,C\gamma^8$. Since $uP_\pm=P_\mp u$: $u^TCP_\pm u=-CuP_\pm u=-CP_\mp u\,u=-n(u)\,CP_\mp$. $\square$

\emph{Proof of item 2.} Apply item 3 twice with $u=\gamma^a$ ($n=+1$) and $u=\gamma^b$ ($n=-1$): $(\gamma^a)^T(CP_\pm)\gamma^a=-n(\gamma^a)\,CP_\mp=-CP_\mp$, and then $g^T(CP_\pm)g=(\gamma^b)^T(-CP_\mp)\gamma^b=-\bigl(-n(\gamma^b)\bigr)CP_\pm=-CP_\pm$. A form $\alpha CP_-+\beta CP_+$ is therefore mapped to its negative by $g$, and it is invariant only if it is zero. $\square$

\emph{Proof of item 4.} For $M=C\gamma^8$: $(C\gamma^8\gamma^a)^T=(\gamma^a)^T\gamma^8C=(\gamma^a)^TC\gamma^8=-C\gamma^a\gamma^8=C\gamma^8\gamma^a$, symmetric. For $M=C$: $(C\gamma^a)^T=-C\gamma^a$, antisymmetric. A combination $\alpha CP_-+\beta CP_+$ with $\beta=-\alpha$ is proportional to $C\gamma^8$, with $\beta=\alpha$ proportional to $C$, and otherwise has neither property. That invariance puts $M$ into this space follows as for item 1, because a global spin transformation $R$ with its vector matrix $\Lambda$ ($R^{-1}\gamma^aR=\Lambda^a{}_b\gamma^b$) turns $\Psi^TM\gamma^a\partial_a\Psi$ into $\Psi^TR^TMR\,\gamma^b\partial_b\Psi$; the verifiers find the same two-dimensional space for the derivative forms (check \texttt{MA\_M3\_ki\allowbreak{}neticInv\allowbreak{}ariantFo\allowbreak{}rms}). $\square$

\emph{Proof of item 5.} For anticommuting components $\Psi_i\Psi_j=-\Psi_j\Psi_i$, so $\Psi^TM\Psi=\Psi^TM_A\Psi$ with the antisymmetric part $M_A=\tfrac12(M-M^T)$ (Lemma 5.4). The map $M\mapsto(S^{ab})^TM+MS^{ab}$ sends symmetric matrices to symmetric ones and antisymmetric ones to antisymmetric ones (transpose it). Moreover a bilinear $\Psi^TX\Psi=2\sum_{i<j}X_{ij}\Psi_i\Psi_j$ with antisymmetric $X$ vanishes only if $X=0$, because the monomials $\Psi_i\Psi_j$ ($i<j$) are independent. So the invariance of $\Psi^TM_A\Psi$ is exactly the condition of item 1 for $M_A$, and item 1 makes $M_A$ symmetric and antisymmetric at once: $M_A=0$, and every invariant mass-type term vanishes. For derivative terms: if $A$ is a constant antisymmetric matrix, Lemma 5.5 gives $\Psi^TA\,\partial_a\Psi=\tfrac12\partial_a(\Psi^TA\Psi)$, a total derivative, which is the case $A=C\gamma^a$; in curved space, with $\sqrt{|g|}$ and the canonical connection, this is the theorem of Section 5.13 that Lg[] is a pure divergence for a Grassmann field. If instead $X$ is constant and symmetric, the term $\Psi^TX\partial_a\Psi$ has the left derivatives $\partial_L/\partial\Psi_c=(X\partial_a\Psi)_c$ and $\partial_L/\partial(\partial_a\Psi_c)=-(X\Psi)_c$ (move $\partial_a\Psi_c$ to the left past the odd $\Psi_i$), so its Euler--Lagrange expression is $(X\partial_a\Psi)_c+\partial_a(X\Psi)_c=2(X\partial_a\Psi)_c\ne0$; this is the case $X=C\gamma^8\gamma^a$. For commuting components only the symmetric part of a matrix counts in $\Psi^TM\Psi$ (Section 1.8), $CP_\pm$ are symmetric and nonzero, and the roles of symmetric and antisymmetric in the derivative terms are exchanged. $\square$

\emph{Proof of item 6.} $\Psi^T\to e^{i\alpha}\Psi^T$ and $\Psi\to e^{i\alpha}\Psi$, so every term $\Psi^TX\Psi$ is multiplied by $e^{2i\alpha}$; $\mathcal L$ contains only bilinears $\Psi^\dagger\cdots\Psi$ and functions of $S$ (Section 17.8). $\square$

\textbf{Worked example.} In the notebook's basis $C=\mathrm{diag}(-\sigma,\sigma)$ with $\sigma=\begin{pmatrix}0&I_4\\ I_4&0\end{pmatrix}$ (Section 2.9) and $P_-=\mathrm{diag}(I_8,0)$, so $CP_-=\mathrm{diag}(-\sigma,0)$: its only nonzero entries are $-1$ in the places $(0,4),(1,5),(2,6),(3,7)$ and their mirror images. For commuting components

\[
\Psi^TCP_-\Psi=-2\bigl(\Psi_0\Psi_4+\Psi_1\Psi_5+\Psi_2\Psi_6+\Psi_3\Psi_7\bigr).
\]

With $\Psi_0=1$, $\Psi_4=2$ and all other components 0 this is $-4$; after the phase $\Psi\to e^{i\alpha}\Psi$ it is $-4e^{2i\alpha}$, which for $\alpha=\pi/2$ is $+4$. A term $g\,\Psi^TCP_-\Psi$ plus its complex conjugate in the Lagrangian would therefore break the U(1) of Theorem M1 down to the two phases with $e^{2i\alpha}=1$, namely $\Psi\to\pm\Psi$. For anticommuting components the same expression is $-(\Psi_0\Psi_4+\Psi_4\Psi_0)-\dots=0$: the Majorana mass term does not exist.

\textbf{Pin characters of the Majorana-type terms.} Under $\Psi\to u\Psi$ with $u=\gamma^b$ ($n=\eta_{bb}$), $\Psi^TC\Psi$ picks up the factor $-n$ and $\Psi^TC\gamma^8\Psi$ the factor $+n$ (item 3); for the derivative terms the factor also depends on the lift of the reflection (the twisted and untwisted lifts of Section 2.14), as recorded in check \texttt{MA\_M3\_pi\allowbreak{}nCharact\allowbreak{}ersMajor\allowbreak{}anaTerms}.

\textbf{The notebook's Lagrangian.} Lg[] (Section 5.12) uses Transpose[Psi16], not ConjugateTranspose: its kinetic matrix $\sigma_{16}\,\mathrm T16^a=C\gamma^a$ and its mass matrix $\sigma_{16}=C$ belong to the family of item 1, with the Pin character $-n(u)$. It is therefore of Majorana type. For a complex commuting field it has U(1) charge 2; for a real commuting field it is non-trivial; for a Grassmann field it is a total derivative with a vanishing mass term (Chapter 5). dirac16complex00 uses the charge-0 Lagrangian $\mathcal L$ of Section 17.8 instead (Chapter 6).

\textbf{Quartic examples (not a classification).} For anticommuting components invariant charge-violating terms of higher order do exist. The matrices $C\gamma^a\gamma^b$ with $a\ne b$ are antisymmetric (transpose and use $(\gamma^a)^TC=-C\gamma^a$ twice), so the bilinears $T^{ab}=\Psi^TC\gamma^a\gamma^b\Psi$ do not vanish for Grassmann components, and they transform as an antisymmetric tensor. Their full contractions are $\mathrm{Spin}_0$-invariant:

\[
\begin{aligned}
&Q_4=\sum_{a<b}\eta_{aa}\eta_{bb}\,\bigl(\Psi^TC\gamma^a\gamma^b\Psi\bigr)^2 &&\text{(40 monomials, charge 4)},\\
&Q_2=\sum_{a<b}\eta_{aa}\eta_{bb}\,\bigl(\Psi^TC\gamma^a\gamma^b\Psi\bigr)\bigl(\Psi^\dagger C\gamma^a\gamma^b\Psi\bigr) &&\text{(160 monomials, charge 2)}.
\end{aligned}
\]

These are examples only: higher-order charge-violating terms are not classified (OPEN), and nothing here says that such a term is present, natural, or sufficient for generating an asymmetry.

\textbf{How M3 was verified.} Wolfram: the null space of the 28 conditions has dimension 2 with basis $CP_-$, $CP_+$, a symmetric subspace of dimension 2 and an antisymmetric subspace of dimension 0; the character table over the four characters has the dimensions 0, 1, 1, 0 (trivial, $-n$, $+n$, determinant); all four mass-type forms vanish for Grassmann components (a control with an antisymmetric matrix has 8 nonzero monomials), while for commuting components $CP_-$, $CP_+$, $C$ and $C\gamma^8$ give forms of rank 8, 8, 16 and 16; the kinetic Euler--Lagrange expressions, the charge 2, the Pin characters, the identification of Lg[] and $Q_4$. Python, independently: the null space as the exact integer basis $\tfrac12(C\gamma^8-C)$, $\tfrac12(C\gamma^8+C)$; the survival table with monomial counts ($\Psi^TC\Psi$ and $\Psi^TC\gamma^8\Psi$ have 0 monomials for Grassmann and 8 each for commuting components, while the charge-0 bilinears $\Psi^\dagger C\Psi$ and $\Psi^\dagger C\gamma^8\Psi$ have 16); the derivative terms in the curved field G\_A; $Q_4$ with 40 and $Q_2$ with 160 monomials, and the agreement of $Q_4$ with 8 times the chiral form of the Wolfram package. The main check names are

\begin{Verbatim}[fontsize=\small]
wolfram-matter-antimatter-report.json
  MA_M3_spinInvariantForms  MA_M3_pinCharacterForms  MA_M3_kineticInvariantForms
  MA_M3_grassmannSurvival  MA_M3_commutingSurvival  MA_M3_u1Charge
  MA_M3_notebookLgIsMajoranaType  MA_M3_extraGrassmannQuarticCharge4
python-matter-antimatter-report.json
  MA_M3_invariantFormsSpan_C_Cgamma8  MA_M3_allInvariantFormsSymmetric
  MA_M3_massTypeSurvivalAndCharge  MA_M3_derivativeTypeSurvivalCurved
  MA_M3_quarticChargeViolatingExamplesGrassmann  MA_M3
\end{Verbatim}

\subsection{Theorem M4: the chirality pair of universes}

\textbf{Theorem M4 (field level; the pairing theorem T1 of Stage 5).} For every vielbein, every $m$ and $\lambda$, both statistics and every configuration $\Psi$ (solution or not), with $\Psi_-:=\gamma^8\Psi$:

\[
\begin{aligned}
&\bar\Psi_-=\bar\Psi\gamma^8,\qquad S[\Psi_-]=S[\Psi],\qquad K[\Psi_-]=-K[\Psi],\qquad j^\mu[\Psi_-]=-j^\mu[\Psi],\\
&\mathcal L_{m,\lambda}[\gamma^8\Psi]=-\mathcal L_{-m,-\lambda}[\Psi],\qquad E_{-m,-\lambda}[\gamma^8\Psi]=-\gamma^8E_{m,\lambda}[\Psi],\\
&T_{\mu\nu}[\gamma^8\Psi;-m,-\lambda]=-T_{\mu\nu}[\Psi;m,\lambda].
\end{aligned}
\]

Hence $\Psi$ solves the field equations with $(m,\lambda)$ if and only if $\gamma^8\Psi$ solves them with $(-m,-\lambda)$.

\emph{Proof.} $\gamma^8$ is Hermitian, $(\gamma^8)^2=1$, it anticommutes with every $\gamma^a$, and it commutes with $C$ and with every $S^{ab}$ (both are even, Section 2.10), hence with every $\Omega_\mu$ and with $D_\mu$. Therefore $\bar\Psi_-=\Psi^\dagger\gamma^8C=\Psi^\dagger C\gamma^8=\bar\Psi\gamma^8$, and a bilinear $\bar\Psi X\Psi$ goes to $\bar\Psi\gamma^8X\gamma^8\Psi$. This is $-\bar\Psi X\Psi$ when $X$ contains one gamma matrix (the current, the kinetic term, the kinetic part of $T_{\mu\nu}$) and $+\bar\Psi X\Psi$ for $X=1$ (the scalar $S$). Consequently

\[
\mathcal L_{m,\lambda}[\gamma^8\Psi]=\sqrt{|g|}\bigl[-K-mS-\tfrac\lambda2S^2\bigr]=-\sqrt{|g|}\Bigl[K-(-m)S-\tfrac{(-\lambda)}2S^2\Bigr]=-\mathcal L_{-m,-\lambda}[\Psi].
\]

For the field equation, $\gamma^\mu D_\mu(\gamma^8\Psi)=-\gamma^8\gamma^\mu D_\mu\Psi$ and $(-m-\lambda S)\gamma^8\Psi=-\gamma^8(m+\lambda S)\Psi$, so $E_{-m,-\lambda}[\gamma^8\Psi]=-\gamma^8\gamma^\mu D_\mu\Psi+\gamma^8(m+\lambda S)\Psi=-\gamma^8E_{m,\lambda}[\Psi]$. The energy--momentum tensor (Chapter 7) is $T_{\mu\nu}=-\tfrac14[\text{kinetic bilinears}]+g_{\mu\nu}\mathcal L_s$: the kinetic bilinears change sign, and $\mathcal L_s[\gamma^8\Psi;-m,-\lambda]=-\mathcal L_s[\Psi;m,\lambda]$ is the Lagrangian identity with $(m,\lambda)$ replaced by $(-m,-\lambda)$. The map is linear and keeps every $\Psi^\dagger$ to the left of every $\Psi$, so the proof holds for both statistics. $\square$

\textbf{Why $\lambda$ must change sign.} The potential $U=\tfrac\lambda2S^2$ is even in $S$, and $S$ does not change, so at a fixed $\lambda$ the identity fails: $\mathcal L_{m,\lambda}[\gamma^8\Psi]+\mathcal L_{-m,\lambda}[\Psi]=\sqrt{|g|}\bigl[-K-mS-\tfrac\lambda2S^2+K+mS-\tfrac\lambda2S^2\bigr]=-\lambda\sqrt{|g|}\,S^2\ne0$ (the erratum E2 of the project contract).

\textbf{Corollary M4.1 (the pair).} Let $\Psi_+$ carry $(m,\lambda)$ and $\Psi_-=\gamma^8\Psi_+$ carry $(-m,-\lambda)$ in the same gravitational field. Then at every point and every $x_4$, in particular at $x_4=0$,

\[
j^\mu_++j^\mu_-=0,\qquad Q_++Q_-=0,\qquad T^{\mathrm{pair}}_{\mu\nu}=T_{\mu\nu}[\Psi_+;m,\lambda]+T_{\mu\nu}[\Psi_-;-m,-\lambda]=0 .
\]

The pair carries no net charge and no net energy, momentum or stress, so the Einstein (or Einstein--Lovelock) equations with the pair as their source are the source-free equations. The energies cancel because the classical energy density of the member of mass $-M$ is minus that of the member of mass $+M$. This is a statement about conservation laws and constraints; it is not a computed creation rate or amplitude (Chapter 16).

\textbf{Worked example.} The rest state $u_+$ of Section 17.9 has $Bu_+=u_+$. Its image $v=\gamma^8u_+$ flips the signs of the upper eight components: $v=\tfrac12(-e_0+e_4-i\,e_9-i\,e_{13})$. With the rows of $B$ (Section 2.9: $(Bw)_0=iw_9$, $(Bw)_4=-iw_{13}$, $(Bw)_9=-iw_0$, $(Bw)_{13}=iw_4$) one finds $(Bv)_0=i(-\tfrac i2)=\tfrac12=-v_0$, and likewise for the other three places, so $Bv=-v$ and $v^\dagger Bv=-1$: the charge densities of $\Psi_+=u_+e^{-imx_4}$ and of its partner add to $1+(-1)=0$. The partner solves the field equation with the mass $-m$: $h_0(-m)=im\gamma^4=\gamma^8(-im\gamma^4)\gamma^8$, so $h_0(-m)v=\gamma^8h_0(m)u_+=m\,v$.

\textbf{One-particle (Krein) level.} In the good sector in flat space, with no momenta along $x_5,x_6,x_7$, the mode Hamiltonian $h_k(m)=-im\gamma^4-\gamma^4\sum_{j\ne4}k_j\gamma^j$ is Hermitian, commutes with $B$ and satisfies $h_k(m)^2=(m^2+k^2)\,1$ (Sections 8.9 and 8.10). Since $\gamma^8\gamma^4\gamma^8=-\gamma^4$ and $\gamma^8\gamma^4\gamma^j\gamma^8=\gamma^4\gamma^j$ (two anticommutations), and $B$ is $-i$ times a product of five gammas,

\[
\gamma^8h_k(m)\gamma^8=h_k(-m),\qquad \gamma^8B\gamma^8=-B .
\]

So $\gamma^8$ maps the positive-energy (negative-energy) eigenspace of $h_k(m)$ onto the positive-energy (negative-energy) eigenspace of $h_k(-m)$ with the same energy: if $h_k(m)u=Eu$, then $h_k(-m)\gamma^8u=\gamma^8h_k(m)u=E\gamma^8u$. It preserves the Hilbert norm $u^\dagger u$ ($\gamma^8$ is orthogonal) and reverses the Krein norm $u^\dagger Bu$. At rest the $B$-form on each positive-energy space has signature (4,4) (Section 8.8), and the Gram matrix of the image is minus the original one.

Which energy and charge an image mode carries depends on the canonical structure given to the image field. With the expectation-value rule of Section 8.12, $\langle\Psi^\dagger X\Psi\rangle=u^\dagger BXu$ for a Hilbert-normalised positive-energy mode $u$, the energy is $X=Bh$, the charge $X=B$ and the scalar density $X=C$. For the image field $\Psi_-=\gamma^8\Psi$ the canonical anticommutator is $\gamma^8B\gamma^8=-B$, and the rule reads $u_-^\dagger(-B)Xu_-$ with $u_-=\gamma^8u$. For an independently quantised theory of mass $-m$ with its own positive structure the rule is $u_-^\dagger BXu_-$. The table gives $(E,Q,S)$ per quantum for three exact positive-energy modes $u$ with mass $m$ and momentum $k=(k_0,k_1,k_2,k_3)$ (measurement \texttt{M4.krein\allowbreak{}ModeFact\allowbreak{}s} of \texttt{python-m\allowbreak{}atter-an\allowbreak{}timatter\allowbreak{}-report.\allowbreak{}json}):

\begingroup
\small
\begin{longtable}{@{}>{\raggedright\arraybackslash}p{0.200\linewidth}>{\raggedright\arraybackslash}p{0.200\linewidth}>{\raggedright\arraybackslash}p{0.200\linewidth}>{\raggedright\arraybackslash}p{0.200\linewidth}@{}}
\toprule
\textbf{$m$, $k$} & \textbf{universe of mass $+m$} & \textbf{image field, metric $-B$} & \textbf{independent, metric $+B$} \\
\midrule
\endfirsthead
\toprule
\textbf{$m$, $k$} & \textbf{universe of mass $+m$} & \textbf{image field, metric $-B$} & \textbf{independent, metric $+B$} \\
\midrule
\endhead
$1$, $(1,1,2,3)$ & $(4,1,1/4)$ & $(-4,-1,1/4)$ & $(4,1,-1/4)$ \\
$3$, $(1,1,1,2)$ & $(4,1,3/4)$ & $(-4,-1,3/4)$ & $(4,1,-3/4)$ \\
$3/5$, $(4/5,0,0,0)$ & $(1,1,3/5)$ & $(-1,-1,3/5)$ & $(1,1,-3/5)$ \\
\bottomrule
\end{longtable}
\endgroup

In the first column $E=\sqrt{m^2+k^2}$ (for example $\sqrt{1+1+1+4+9}=4$), $Q=u^\dagger BBu=u^\dagger u=1$ and $S=m/E$. The last value is the flat-space case $M_{\mathrm{eff}}=m$ of the result $s=M_{\mathrm{eff}}/E$ of Section 11.4, and it takes two lines: $BC=-i\gamma^4$ (Section 2.12), so $S=u^\dagger BCu=u^\dagger Zu$ with $Z=-i\gamma^4$; and $h_k(m)=mZ+X$ with $X=-\gamma^4\sum_jk_j\gamma^j$, where $Z^2=1$ and $ZX+XZ=0$ (each $\gamma^4\gamma^j$ with $j\ne4$ anticommutes with $\gamma^4$), so $Zh+hZ=2m$. Sandwiching this between $u^\dagger$ and $u$ with $hu=Eu$, $u^\dagger h=Eu^\dagger$ and $u^\dagger u=1$ gives $2E\,u^\dagger Zu=2m$, that is $S=m/E$. The general pattern follows from the two matrix identities. For the image field: $u_-^\dagger(-B)(Bh(-m))u_-=-u^\dagger\gamma^8h(-m)\gamma^8u=-u^\dagger h(m)u=-E$, $u_-^\dagger(-B)Bu_-=-u^\dagger u=-1$ and $u_-^\dagger(-B)Cu_-=-u^\dagger(\gamma^8B\gamma^8)(\gamma^8C\gamma^8)u=u^\dagger BCu=S$: so $(E,Q,S)\to(-E,-Q,S)$. For the independent field: $u_-^\dagger h(-m)u_-=E$, $u_-^\dagger u_-=1$ and $u_-^\dagger BCu_-=-S$: so $(E,Q,S)\to(E,Q,-S)$.

\textbf{Fock level (from Stage 5; PROVISIONAL).} The Stage-5 Wolfram pairing report has all 141 of its checks true, and its exact Fock-level result (the checks \texttt{PAIR\_T1k\allowbreak{}rein\_*}, in a Fock model of four rest modes with $m=1$, the sea formed by the two negative-energy modes, and normal ordering as the subtraction of the sea value; key \texttt{\detokenize{T1krein}} of \texttt{pairing-\allowbreak{}theory.j\allowbreak{}son}) is as follows.

\begin{itemize}
\item The image field $\Psi_-=\gamma^8\Psi$, on the same Fock space and in the same state, has the anticommutator $-B$ on its modes; the operator identities $H[\Psi_-;-m]=-H[\Psi;m]$, $Q[\Psi_-]=-Q[\Psi]$ and $S[\Psi_-]=S[\Psi]$ hold, normal ordering commutes with the map, and $:\!T_{\mu\nu}[\gamma^8\Psi;-m,-\lambda]\!:\ =-:\!T_{\mu\nu}[\Psi;m,\lambda]\!:$ as operators. Stage 5 describes the image as a universe with the Krein metric $-B$, energy $-|\epsilon|$ and charge $-1$ per quantum; these numbers are the expectation values of the formulas of $\mathcal L_{-m,-\lambda}$, not the image field's own energy and charge (the caveat below).
\item An independently quantised theory of mass $-m$ with the positive ($J=B$) structure has the modes $w_n=\gamma^8u_n$ (same energy $\epsilon$, Krein sign $-\beta_n$). A quantum in $w_n$ has energy $+|\epsilon|$, charge $+1$ and the opposite scalar density; with the same mode occupations such a universe has, in this free Fock model, the same energy--momentum and the same charge as the universe of mass $+m$, so the pair totals add instead of cancelling (Stage 5 calls this the T2-type pairing; its Kohn--Sham form is the mirror map of Section 13.14, treated in Chapter 15). For $\lambda\ne0$ the energies do not simply add: with the parameters $(-m,-\lambda)$ the interaction energy of the $-M$ universe is $-\tfrac\lambda2S'^2$ with $S'=-S_+$, that is $-\tfrac\lambda2S_+^2$, the opposite of that of the $+M$ universe (entry \texttt{kreinLev\allowbreak{}elCaveat} of the Wolfram matter-antimatter report). Only the kinetic and mass energies add.
\end{itemize}

The Stage-5 files carry no finality flag: the Wolfram matter-antimatter report records the status of these statements as provisional until the Stage-5 gate has passed (measurement \texttt{M4\_krein\allowbreak{}LevelSta\allowbreak{}tus} of \texttt{wolfram-\allowbreak{}matter-a\allowbreak{}ntimatte\allowbreak{}r-report\allowbreak{}.json}).

\textbf{How particles and antiparticles map} (derived from these operator identities, as in the matter-antimatter document). In the universe of mass $+m$ a particle has $(E,Q)=(+|\epsilon|,+1)$ and an antiparticle, a hole in the sea, has $(+|\epsilon|,-1)$. The image field describes the same particles and antiparticles, with the same energies and charges when these are measured with its own generators $-H[\Psi_-;-m]$ and $-Q[\Psi_-]$; only the formulas of $\mathcal L_{-m,-\lambda}$, evaluated on the same states, give $(-|\epsilon|,-1)$ and $(-|\epsilon|,+1)$. In the independently quantised $-M$ universe, particles and antiparticles have the same energies and charges as in the $+M$ universe and the opposite scalar density. For the modes of positive and negative Krein norm: a mode $u_n$ with Krein sign $\beta_n=u_n^\dagger Bu_n$ goes to $w_n=\gamma^8u_n$ with $w_n^\dagger Bw_n=-\beta_n$. Measured with the image field's own metric $-B$ the sign is $\beta_n$ again; measured with $+B$ (independent quantization) every positive-norm mode becomes a negative-norm mode and conversely.

So the cancellation of Corollary M4.1 does not carry over to the quantum level as a cancellation between two universes. If the $-M$ member is the image field with the Krein metric $-B$, the operator totals vanish, but only as identities of the form $X+(-X)=0$ within one quantum system (the caveat that follows). If the $-M$ universe is an independently quantised field with positive energies, the charges would cancel only by an additional assumption about the state of the $-M$ universe, and the energies would not cancel (the kinetic and mass energies add, and for $\lambda\ne0$ the interaction energies have opposite signs).

\textbf{A caveat on what the operator cancellation means.} The committed Wolfram matter-antimatter report records a further point (measurement \texttt{M5\_impli\allowbreak{}cation}, entry \texttt{kreinLev\allowbreak{}elCaveat}, and the matrix checks of \texttt{M4\_krein\allowbreak{}OneParti\allowbreak{}cle}). In the image-field reading, $\Psi_-=\gamma^8\Psi_+$ is built from the same operators, on the same Fock space and in the same state, as $\Psi_+$. The identities $Q_++Q[\Psi_-]=0$ and $H_++H[\Psi_-;-m]=0$ are therefore of the form $X+(-X)=0$ and hold in every state; they are not a compensation by a second, independent universe. Relative to its own canonical structure (anticommutator $-B$, Lagrangian $-\mathcal L_{-m,-\lambda}$), the image field's generator of time evolution is $-H[\Psi_-;-m]=+H_+$ and its generator of phases is $-Q[\Psi_-]=+Q_+$: it is the same quantum system as $\Psi_+$, with the same energy and charge, and the values $-|\epsilon|$ and $-1$ per quantum are expectation values of the $\mathcal L_{-m,-\lambda}$ formulas rather than of that field's own Hamiltonian and charge. In the independent reading the kinetic and mass energies add, the charges cancel only by an additional assumption on the state of the $-M$ universe, and for $\lambda\ne0$ its interaction energy has the opposite sign. The committed report sums this up in one sentence: at the quantum level no reading of H1 gives a cancellation between two independent, consistently quantised universes. \textbf{Of the two readings examined in the repository, neither describes two independent, consistently quantised universes whose charges and energies cancel.} Whether any quantum description of a physical universe of mass $-M$ gives the classical cancellation of Corollary M4.1 is \textbf{OPEN}.

\textbf{How M4 was verified.} Wolfram: the matrix facts (the current matrices are odd under $\gamma^8$, the mass matrix is even, $\bar\Psi\to\bar\Psi\gamma^8$, $B\to-B$); the current flip in G1 for Grassmann components at three points; for commuting components at the three G1 points the field jets of $T^{\mathrm{pair}}_{\mu\nu}$ and of the pair current vanish while the individual ones do not, and the image of an exact solution solves the $(-m,-\lambda)$ equations; for Grassmann components all 64 components of $T^{\mathrm{pair}}_{\mu\nu}$ vanish at the point p1, while each member has 64 nonzero components; and the one-particle facts above. Python, independently: the $\gamma^8$ matrix facts, including $(\gamma^8)^TC\gamma^a\gamma^8=-C\gamma^a$ and the commutation with $\Omega_\mu$ in G\_A; in G\_A for both statistics the current flip, $S$ invariant, $K\to-K$, the Lagrangian identity, the Euler--Lagrange pairing and the pairing of all 36 independent components of $T_{\mu\nu}$, as exact polynomial identities; the three Krein mode samples of the table; and the citation of the Stage-5 report with all twelve \texttt{PAIR\_T1k\allowbreak{}rein} checks true. The Python checker of the same commit adds the check \texttt{MA\_M4\_im\allowbreak{}ageField\allowbreak{}FockMode\allowbreak{}l}, which the committed Python report does not yet contain (Section 17.17): in an exact Fock model of four modes of the good sector it builds the image field as a genuine operator on the same Fock space and verifies the caveat above, namely that the image field's own generators of $x_4$-evolution and of phases are $-H[\Psi_-;-m]=+H_+$ and $-Q[\Psi_-]=+Q_+$, that $Q_++Q[\Psi_-]=0$ in every Fock state, and that an independently quantised $-m$ field with the same occupations has $H'=H_+$, $Q'=Q_+$ and $S'=-S_+$. The main check names are

\begin{Verbatim}[fontsize=\small]
wolfram-matter-antimatter-report.json
  MA_M4_currentFlipMatrix  MA_M4_currentFlipG1_grassmann
  MA_M4_pairEMTAndCurrentG1_commuting  MA_M4_pairEMTG1_grassmann  MA_M4_kreinOneParticle
python-matter-antimatter-report.json
  MA_M4_gamma8MatrixFacts  MA_M4_gamma8ChargeFlipCurved  MA_M4_gamma8EulerLagrangePairing
  MA_M4_gamma8EMTPairing  MA_M4_kreinModeFacts  MA_M4
checker of this commit only (not in the committed Python report)
  MA_M4_imageFieldFockModel
wolfram-pairing-report.json (Stage 5, cited)
  PAIR_T1krein_imageAnticommutatorMinusB  PAIR_T1krein_imageOperatorIdentities
  PAIR_T1krein_imageExpectationValues  PAIR_T1krein_independentCARPlusB
  PAIR_T1krein_independentExpectationValues  (and seven more PAIR_T1krein_ checks)
\end{Verbatim}

\subsection{M5: the conditional scenario, stated as a hypothesis}

The pair of Section 17.13 suggests a picture of the kind described in Section 17.7: our universe has an excess of charge, and a partner universe has the opposite excess. This section states that picture as precisely as possible, separates what is proved from what is assumed, and says what the picture does not do.

\textbf{The hypotheses.} The following three statements are \textbf{hypotheses}. None of them is derived anywhere in the repository.

\begin{itemize}
\item \textbf{H1 (HYPOTHESIS, not derived).} Our universe is one member of a chirality pair created together: two independent classical fields in the same gravitational field, $\Psi_+$ governed by the Lagrangian $\mathcal L_{m,\lambda}$ and $\Psi_-$ governed by the Lagrangian $+\mathcal L_{-m,-\lambda}$, in the correlated configuration $\Psi_-=\gamma^8\Psi_+$ (the correlation is part of the assumption). The sign in front of the second Lagrangian matters. The field equations of $+\mathcal L_{-m,-\lambda}$ and of $-\mathcal L_{-m,-\lambda}$ are the same, but the Noether current and the energy--momentum tensor are computed from the Lagrangian itself: with $+\mathcal L_{-m,-\lambda}$ the Noether current of $\Psi_-$ is $i\sqrt{|g|}\,j^\mu[\Psi_-]$ (Theorem M1 holds for every $m$ and $\lambda$), and its energy--momentum tensor is $T_{\mu\nu}[\Psi_-;-m,-\lambda]=-T_{\mu\nu}[\Psi_+;m,\lambda]$ (Theorem M4), so the second member has the classical energy $-E_+$. This is the form of H1 recorded in the committed Wolfram report (measurement \texttt{M5\_impli\allowbreak{}cation}, entry \texttt{hypothes\allowbreak{}es}). At the quantum level H1 has no reading in which two independent, consistently quantised universes cancel each other's charge and energy (the caveat of Section 17.13). No creation process, rate or amplitude is computed anywhere in the repository.
\item \textbf{H2 (HYPOTHESIS, not derived).} The creation assigns $Q_+=-Q_-\ne0$. Given H1, the relation $Q_-=-Q_+$ follows from M4; the content of H2 is $Q_+\ne0$, and nothing in the repository computes the value or the sign of $Q_+$.
\item \textbf{H3 (HYPOTHESIS, not derivable within the theory).} The dirac16complex U(1) charge is identified with the baryon number $B$ (or with $B-L$). The theory contains no Standard-Model baryons, quarks or leptons, so this identification cannot be derived from it.
\end{itemize}

\textbf{Proposition M5 (classical level).} If H1, H2 and H3 hold, then $Q_+(x_4)+Q_-(x_4)=0$ at every $x_4$, each of $Q_+$ and $Q_-$ is separately conserved (when the boundary flux vanishes), and the baryon excess $B_+=Q_+\ne0$ seen in one member is compensated exactly by $B_-=-Q_+$ in the other. The pair also carries $T^{\mathrm{pair}}_{\mu\nu}=0$ for the classical bilinears: the member of mass $-M$ has the classical energy $-E_+$.

\emph{Proof.} By M4, $\sqrt{|g|}\,j^{x_4}[\gamma^8\Psi_+]=-\sqrt{|g|}\,j^{x_4}[\Psi_+]$ at every point. By H1 the charge of $\Psi_-$ is the Noether charge of $+\mathcal L_{-m,-\lambda}$, which is built from the same current $j^\mu[\Psi_-]$, so $Q_-=-Q_+$ at every $x_4$, and $Q_++Q_-=0$. By M1 each charge is conserved, $dQ_\pm/dx_4=0$. By H3, $B_\pm=Q_\pm$, and by H2, $B_+\ne0$. The statement about $T^{\mathrm{pair}}_{\mu\nu}$ is Corollary M4.1, with $T_{\mu\nu}[\Psi_-;-m,-\lambda]$ the energy--momentum tensor of $+\mathcal L_{-m,-\lambda}$. $\square$

The Wolfram verifier checks the implication symbolically: $Q_++Q_-$ simplifies to 0, and $dQ_-/dx_4$ simplifies to 0 given $dQ_+/dx_4=0$ (check \texttt{MA\_M5\_im\allowbreak{}plicatio\allowbreak{}n}, with the ingredients computed in the checks of M1 and M4). The committed report records its status in these words: the classical-level implication is proved; the scenario itself is a hypothesis, not a result. The implication is elementary; its substance lies entirely in the hypotheses.

\textbf{At the quantum level (PROVISIONAL, and no cancellation between two universes).} Proposition M5 is not proved for the quantized fields. The operator identity $:\!T_{\mu\nu}[\gamma^8\Psi;-m,-\lambda]\!:\ =-:\!T_{\mu\nu}[\Psi;m,\lambda]\!:$ of Section 17.13 is cited from the Stage-5 pairing results, which are provisional until the Stage-5 gate has passed (measurement \texttt{M4\_krein\allowbreak{}LevelSta\allowbreak{}tus}). And even taken as given, by the caveat of Section 17.13 it is an identity of the form $X+(-X)=0$ within one quantum system, not a compensation by a second universe; for an independently quantised $-M$ universe the energies do not cancel and the charges cancel only by an additional assumption.

\textbf{What M5 does not do.}

\begin{enumerate}
\item It does not produce an asymmetry. By M1, $Q_+$ is constant, so a nonzero $Q_+$ today is the same as a nonzero $Q_+$ at $x_4=0$: the asymmetry of our universe is an initial condition, which H2 puts in by hand.
\item It does not predict $\eta$. The observed $\eta\approx6\times10^{-10}$ is \textbf{not predicted}: nothing here computes the magnitude or the sign of $Q_+$, nor the photon content of the universe to which $\eta$ refers.
\item It does not satisfy Sakharov's conditions; it replaces them by a global symmetry of the pair together with an initial condition in each member.
\item It depends on the reading of H1. At the classical level the cancellation holds for two correlated classical fields, the second governed by $+\mathcal L_{-m,-\lambda}$, and the member of mass $-M$ then has the classical energy $-E_+$. At the quantum level, when the $-M$ member is the image field $\gamma^8\Psi_+$ with the Krein metric $-B$, the totals vanish only as an operator identity of the form $X+(-X)=0$ within one quantum system (the caveat of Section 17.13), not as a cancellation between two universes; if the $-M$ universe is instead quantised independently with its own positive structure, its quanta carry charge $+1$ and energy $+|\epsilon|$, its energies add to ours, and its charge cancels ours only under an additional assumption about its state. Which reading, if any, describes a physical universe of mass $-M$ is OPEN.
\item It does not relate the dirac16complex charge to the baryons of the Standard Model (H3), and it does not show that pairs are created (H1). Whether the big bang creates universes in pairs is the subject of Chapter 16; no creation process, rate or amplitude is derived anywhere in the repository.
\item It is not a statement about the observed world. The model lives in signature (4,4), with four time-like directions, while observed spacetime has one time (row L43 of the ledger of Chapter 0); applying the scenario to "our universe" is a further assumption, and nothing in it is fitted to or compared with data. Its quantum statements are formal: the canonical state space is a Krein space (Section 8.7), a positive Fock space exists only in the good sector (Section 8.10), the extra-time sector is ill-posed (Section 8.13), and dirac16complex00 is treated as a classical field.
\end{enumerate}

\subsection{M6: the Sakharov scorecard, and the answer}

The three conditions of Section 17.5, applied to the theory as built:

\begingroup
\small
\begin{longtable}{@{}>{\raggedright\arraybackslash}p{0.267\linewidth}>{\raggedright\arraybackslash}p{0.267\linewidth}>{\raggedright\arraybackslash}p{0.267\linewidth}@{}}
\toprule
\textbf{Sakharov condition} & \textbf{Status in the theory as built} & \textbf{What would have to be added} \\
\midrule
\endfirsthead
\toprule
\textbf{Sakharov condition} & \textbf{Status in the theory as built} & \textbf{What would have to be added} \\
\midrule
\endhead
1. Violation of the conserved number (baryon number; here the charge $Q$) & Fails, by M1: $Q$ is exactly conserved in every gravitational field (when the boundary flux vanishes), for both statistics and every potential & A U(1)-violating interaction from the list of M3 (Section 17.16, item 1) \\
2. C and CP violation & Fails, by M2: C is exact for the commuting field in every gravitational field, and CP is exact for the anticommuting field in flat space and in gravitational fields in which the reflection is an isometry; each reverses $Q$ and preserves $x_4$. A generic gravitational field breaks this CP, which does not help, because condition 1 fails & Terms that break all 64 exact charge-reversing symmetries without $x_4$ reversal (Section 17.16, item 2) \\
3. Departure from thermal equilibrium & Not addressed by any computation; the Kohn--Sham states of Chapters 13 and 14 are equilibrium states at fixed $N$ & A non-equilibrium history and a computation of rates (Section 17.16, item 3) \\
\bottomrule
\end{longtable}
\endgroup

Since condition 1 fails exactly, no departure from equilibrium could create a net charge inside one universe of this theory. The computed checks behind the rows are

\begin{Verbatim}[fontsize=\small]
row 1  wolfram-matter-antimatter-report.json
         MA_M1_noetherIdentity_grassmann_G1  MA_M1_noetherIdentity_commuting_G1
       python-matter-antimatter-report.json
         MA_M1_noetherIdentity_<X>_<G>  (both statistics, both geometries)
row 2  wolfram-matter-antimatter-report.json
         MA_M2_symmetrySummaryAndChargeReversal
       python-matter-antimatter-report.json
         MA_M2_discreteGroupCharacterTable  MA_M2_C_and_CP_status
       checker of this commit only (not in the committed Python report)
         MA_M2_cpScopeInCurvedFields
row 3  none; nothing was computed
\end{Verbatim}

The same three rows, in the checker's own words, are stored in the committed Python report as the measurement \texttt{M6\_sakha\allowbreak{}rovScore\allowbreak{}card}, and that report records the honest answer as the measurement \texttt{honestAn\allowbreak{}swer}. The committed report was written by an earlier version of the checker (Section 17.17). The checker of the same commit words the rows more precisely: row 1 adds the vanishing boundary flux, and row 2 adds the scope of CP given in the table above and lists \texttt{MA\_M2\_cp\allowbreak{}ScopeInC\allowbreak{}urvedFie\allowbreak{}lds}. Its honest answer adds that at the quantum level no reading gives a cancellation between two independent universes, because the image field is the same quantum system (Section 17.13).

\textbf{The answer.} The dirac16complex theory as built does \textbf{not} solve the matter--antimatter problem. Its charge is exactly conserved (M1), it has exact charge-reversing symmetries (M2: C for dirac16complex00 in every gravitational field, CP for dirac16complex in flat space and wherever the reflection is an isometry), nothing computes a departure from equilibrium, it contains no baryons of the Standard Model, and it predicts no value of $\eta$. What it contributes is exact but structural: a pair of universes of masses $+M$ and $-M$ with opposite charges and zero total charge (M4), which, under the three hypotheses of M5, would realise the idea of a symmetric whole made of two asymmetric parts. The statement "this theory solves the matter--antimatter mysteries" is therefore not established; it cannot be proved within the theory as built, because its central part, the creation of a net charge inside one universe, is excluded by Theorem M1.

\subsection{What would have to be added to the theory}

The following list is a consequence of M1 to M3. It describes necessary ingredients; none of them is present in the theory, none is claimed to be natural, and none has been shown to be sufficient.

\begin{enumerate}
\item \textbf{A charge-violating interaction.} By M3, for the anticommuting field no Majorana-type mass term exists. The lowest-order candidates are the derivative term $\sqrt{|g|}\,\Psi^TC\gamma^8\gamma^\mu D_\mu\Psi$ plus its Hermitian conjugate (charge 2), and quartic terms such as $Q_4$ (charge 4) and $Q_2$ (charge 2) plus their conjugates. For the commuting field the mass terms $\Psi^TCP_\pm\Psi$ and the derivative term $\sqrt{|g|}\,\Psi^TC\gamma^\mu D_\mu\Psi$ (the form of the notebook's Lg[]), each plus its complex conjugate, have charge 2. Each such term breaks the U(1) of M1 down to a discrete set of phases (derived): $e^{i\alpha}$ with $e^{2i\alpha}=1$ for charge-2 terms, and with $e^{4i\alpha}=1$ for charge-4 terms.
\item \textbf{Breaking of every charge-reversing symmetry.} The added terms must break every exact symmetry of Section 17.11 that reverses $Q$ without reversing $x_4$: $\mathrm C_0$ and the 63 others for the commuting field ($\mathrm C_0$ holds in every gravitational field), and the CP family for the anticommuting field (64 maps, among them $\mathrm C_8\mathrm P_b$ for $b=0,1,2,3$ and $\mathrm C_8\mathrm P_{123}$), which hold in flat space and wherever the reflection is an isometry. A generic gravitational field already breaks the CP family (Section 17.11), but that is of use only together with item 1: as long as the charge is conserved, no symmetry breaking can create a net charge. A remark that is derived but not machine-checked: the U(1) rotation $\Psi\to e^{i\alpha}\Psi$ is a symmetry of $\mathcal L$ and multiplies the coefficient $g$ of a charge-2 term by $e^{2i\alpha}$, so the phase of a single such coefficient can be removed by redefining the field; only relative phases between several charge-violating terms can be physical. Whether a given combination breaks the CP family is not computed (OPEN).
\item \textbf{A departure from equilibrium.} A dynamical history in which the charge-violating processes are out of equilibrium, for example during the evolution of the primordial field of Chapter 9, and a computation of their rates. Nothing of this kind has been computed.
\item \textbf{A link to baryons.} The dirac16complex charge would have to be connected to the baryon number of the Standard Model (H3), which requires couplings to Standard-Model fields that the theory does not contain.
\item \textbf{A computation of $\eta$.} Only with items 1 to 4 could a value of $\eta$ be computed and compared with the observed $\eta\approx6\times10^{-10}$ of Section 17.3.
\end{enumerate}

\subsection{How the analysis was verified, and how to rerun it}

\textbf{The reports.} Two independent exact implementations were written for the analysis, and one Stage-5 report is used as an input:

\begin{Verbatim}[fontsize=\small]
artifacts/dirac16complex/matter-antimatter/wolfram-matter-antimatter-report.json
  44 of 44 checks true; producer scripts/verify_dirac16complex_matter_antimatter.wls
  with wolfram/Dirac16ComplexMatterAntimatter.wl; also writes
  artifacts/dirac16complex/matter-antimatter/matter-antimatter-theory.json
artifacts/dirac16complex/matter-antimatter/python-matter-antimatter-report.json
  75 of 75 checks true; producer scripts/check_dirac16complex_matter_antimatter.py
  with scripts/grassmann_algebra.py, in an earlier version; the report records
  for the checker the sha256
  20914546ed5894ef223be8a1ae7ffecb62b308fdbf13d88f35a939a3c27bd526
  while the committed checker file has the sha256
  ef1a00dc02e72186cd6c78d83ad3d930c372ccd9440983a307e8e04cae622231
  and 77 checks (all true when rerun)
artifacts/dirac16complex/pair-creation/wolfram-pairing-report.json   (Stage-5 input)
  141 of 141 checks true; producer scripts/verify_dirac16complex_pairing.wls
\end{Verbatim}

Both matter-antimatter implementations compute exactly: integers, fractions, Gaussian rationals (complex numbers with rational parts), exact symbolic algebra and exact Grassmann algebras. No floating-point number decides any check, with the one labelled exception \texttt{MA\_M1\_ks\allowbreak{}FixedNet\allowbreak{}NumberRe\allowbreak{}corded} (Section 17.9). Each starts from the gamma matrices built from their definitions and compares them with the exact fixture \texttt{artifact\allowbreak{}s/dirac1\allowbreak{}6complex\allowbreak{}/arbitra\allowbreak{}ry-field\allowbreak{}/algebra\allowbreak{}-fixture\allowbreak{}.json}. The Python checker uses nothing produced by Wolfram as truth: it compares its own results with the Wolfram theory file only in its agreement check \texttt{MA\_agree\allowbreak{}sWithWol\allowbreak{}fram} (2048 classification rows, 22 summary items, the 64 charge-reversing symmetries, the invariant forms, the survival table, $Q_4$, the Krein facts and the $\gamma^8$ facts, with no disagreement). Two things about the committed Python report are older than the files next to it. First, it records, as \texttt{inputSha\allowbreak{}256}, the sha256 of the Wolfram theory file it compared with, \texttt{a3a85c9a\allowbreak{}...}; the committed theory file has the sha256 \texttt{1505a396\allowbreak{}...} (the Wolfram verifier was rerun afterwards), so the agreement check of that report refers to the earlier version. Second, it was written by an earlier version of the Python checker (the two sha256 values are in the list above). The checker file in the same commit has 77 checks, two more than the report: \texttt{MA\_M2\_cp\allowbreak{}ScopeInC\allowbreak{}urvedFie\allowbreak{}lds}, which verifies the scope of CP in a gravitational field (Section 17.11), and \texttt{MA\_M4\_im\allowbreak{}ageField\allowbreak{}FockMode\allowbreak{}l}, which verifies in an exact Fock model that the image field is the same quantum system as $\Psi_+$ (Sections 17.13 and 17.14). It also words some measurements more precisely: the rows of \texttt{M6\_sakha\allowbreak{}rovScore\allowbreak{}card}, the hypothesis H1 of \texttt{M5\_condi\allowbreak{}tionalSc\allowbreak{}enario} (the second member with the Lagrangian $+\mathcal L_{-m,-\lambda}$) and \texttt{honestAn\allowbreak{}swer}. A rerun of this checker against the committed theory file gives 77 of 77 checks true, and the 75 checks of the committed report are among them, all true (Section 19.10 records such a rerun from a fresh clone). Regenerating the committed report is unfinished work of the matter--antimatter analysis. The Stage-5 input files carry no finality flag; they are cited by their sha256 values in both matter-antimatter reports, and the Fock-level statements taken from them are provisional (Section 17.13).

\textbf{Counting the checks yourself.} With the method of Section 0.8, start Python in the repository root:

\begin{Verbatim}[fontsize=\small]
>>> import json
>>> folder = "artifacts/dirac16complex/matter-antimatter/"
>>> for name in ["wolfram-matter-antimatter-report.json",
...              "python-matter-antimatter-report.json"]:
...     checks = json.load(open(folder + name, encoding="utf-8"))["checks"]
...     print(name, sum(checks.values()), len(checks))
...
\end{Verbatim}

(The last line, an empty continuation line, ends the loop.) It prints the two file names followed by \texttt{\detokenize{44 44}} and \texttt{\detokenize{75 75}}: all checks of the committed reports are true.

\textbf{Rerunning the verifiers without touching the committed files.} Run from the repository root; the output goes to the folder \texttt{build/ma\allowbreak{}/}, which Git ignores. The Wolfram step takes about 11 to 13 minutes and the Python checker about a minute (681 and 53 seconds in the rerun of Section 19.10; the matter-antimatter document gives about 13 minutes and half a minute). In Git Bash:

\begin{Verbatim}[fontsize=\small]
export PYTHONUTF8=1
wolframscript -file scripts/verify_dirac16complex_matter_antimatter.wls \
    build/ma/wolfram-report.json build/ma/theory.json
python scripts/check_dirac16complex_matter_antimatter.py \
    --theory build/ma/theory.json --output build/ma/python-report.json
\end{Verbatim}

In PowerShell (the line continuation is a backtick):

\begin{Verbatim}[fontsize=\small]
$env:PYTHONUTF8 = "1"
wolframscript -file scripts/verify_dirac16complex_matter_antimatter.wls `
    build/ma/wolfram-report.json build/ma/theory.json
python scripts/check_dirac16complex_matter_antimatter.py `
    --theory build/ma/theory.json --output build/ma/python-report.json
\end{Verbatim}

Each checker prints one line \texttt{check\_<n\allowbreak{}ame>=tru\allowbreak{}e} (or \texttt{\detokenize{false}}) per check and ends with \texttt{check\_co\allowbreak{}unt} and \texttt{failed\_c\allowbreak{}heck\_cou\allowbreak{}nt}. Expect \texttt{check\_co\allowbreak{}unt=44} and \texttt{failed\_c\allowbreak{}heck\_cou\allowbreak{}nt=0} from the Wolfram verifier, and \texttt{check\_co\allowbreak{}unt=77} and \texttt{failed\_c\allowbreak{}heck\_cou\allowbreak{}nt=0} from the Python checker: 77, not the 75 of the committed report, for the reason given above. The Wolfram report path is a plain positional argument, and an optional second positional argument sets the path of the theory file (the matter-antimatter document explains why no double hyphen is used). The typeset matter-antimatter document is \texttt{provenan\allowbreak{}ce/DIRAC\allowbreak{}16COMPLE\allowbreak{}X\_MATTER\allowbreak{}\_ANTIMAT\allowbreak{}TER.pdf}; Chapter 19 collects the commands of all stages.

\subsection{Literature cited in this chapter}

\begin{itemize}
\item C. D. Anderson, "The positive electron", Phys. Rev. 43, 491 (1933).
\item A. D. Sakharov, "Violation of CP invariance, C asymmetry, and baryon asymmetry of the universe", JETP Lett. 5, 24 (1967).
\item Planck Collaboration (N. Aghanim et al.), "Planck 2018 results. VI. Cosmological parameters", Astron. Astrophys. 641, A6 (2020), arXiv:1807.06209; erratum Astron. Astrophys. 652, C4 (2021).
\item D. J. Fixsen, "The temperature of the cosmic microwave background", Astrophys. J. 707, 916 (2009).
\item E. Tiesinga, P. J. Mohr, D. B. Newell and B. N. Taylor, "CODATA recommended values of the fundamental physical constants: 2018", Rev. Mod. Phys. 93, 025010 (2021).
\item R. H. Cyburt, B. D. Fields, K. A. Olive and T.-H. Yeh, "Big bang nucleosynthesis: Present status", Rev. Mod. Phys. 88, 015004 (2016).
\item L. Canetti, M. Drewes and M. Shaposhnikov, "Matter and antimatter in the universe", New J. Phys. 14, 095012 (2012), arXiv:1204.4186.
\item A. G. Cohen, A. De Rújula and S. L. Glashow, "A matter-antimatter universe?", Astrophys. J. 495, 539 (1998).
\item ALPHA Collaboration (E. K. Anderson et al.), "Observation of the effect of gravity on the motion of antimatter", Nature 621, 716 (2023).
\item V. A. Kuzmin, V. A. Rubakov and M. E. Shaposhnikov, "On anomalous electroweak baryon-number non-conservation in the early universe", Phys. Lett. B 155, 36 (1985).
\item K. Kajantie, M. Laine, K. Rummukainen and M. E. Shaposhnikov, "Is there a hot electroweak phase transition at $m_H\gtrsim m_W$?", Phys. Rev. Lett. 77, 2887 (1996), arXiv:hep-ph/9605288.
\item ATLAS Collaboration, "Observation of a new particle in the search for the Standard Model Higgs boson with the ATLAS detector at the LHC", Phys. Lett. B 716, 1 (2012).
\item CMS Collaboration, "Observation of a new boson at a mass of 125 GeV with the CMS experiment at the LHC", Phys. Lett. B 716, 30 (2012).
\item M. B. Gavela, P. Hernández, J. Orloff and O. Pène, "Standard Model CP-violation and baryon asymmetry", Mod. Phys. Lett. A 9, 795 (1994).
\item P. Huet and E. Sather, "Electroweak baryogenesis and standard model CP violation", Phys. Rev. D 51, 379 (1995), arXiv:hep-ph/9404302.
\item L. Boyle, K. Finn and N. Turok, "CPT-Symmetric Universe", Phys. Rev. Lett. 121, 251301 (2018).
\end{itemize}

\subsection{What we proved and what we assumed}

\textbf{Proved in this chapter} (complete derivations, with the exact machine checks named where the repository has them): the three Sakharov conditions in the form of Propositions 17.1 and 17.2 and of the equilibrium argument of Section 17.5, each under its stated premises; the conversion $\eta=2.7342\times10^{-8}\,\Omega_bh^2$, given the quoted constants and the quoted black-body photon density, and the size (about 0.2 percent) of the helium correction it leaves out; \textbf{Theorem M1}: the exact U(1) invariance of $\mathcal L$ for both statistics and every admissible potential, the Noether current, the off-shell identity $\partial_\mu(\sqrt{|g|}\,j^\mu)=\sqrt{|g|}(\bar E\Psi+\bar\Psi E)$ in every gravitational field, the conservation of $Q$ when the boundary flux vanishes, the indefinite charge density $\Psi^\dagger B\Psi$, and (derived) the conservation of the normal-ordered charge after quantization; \textbf{Theorem M2}: the lemmas A, B and C; the transformation rule for all linear and antilinear maps built from a constant matrix and a reflection of frame directions, in flat space, in every gravitational field in the frame form (which relates the theories on the frames $e$ and $e\,\mathrm{diag}(r)$), and in the coordinate form wherever the reflection is an isometry; Lemma D (derived here), which shows that no other constant matrix gives an exact symmetry; the classification of the exact symmetries; the counts (1024 maps, four classes of 256, 256 exact symmetries, 64 exact charge-reversing symmetries without $x_4$ reversal for each statistics); the table of C, P, CP, T and CPT, with its scope: C exact for dirac16complex00 in every gravitational field; no constant C but an exact CP for dirac16complex with $m\ne0$, in flat space and wherever the reflection is an isometry, broken by a generic gravitational field (no frame rotation of $\mathrm{Spin}_0(4,4)$ and no other exact symmetry undoes its frame change); the compatibility of the exact symmetries with the canonical anticommutator (of unitary type if they preserve $x_4$, of antiunitary type if they reverse it), and (derived) the reversal of the normal-ordered charge by the CP of unitary type; the continuous reversal of $x_4$ and of the classical charge by $\exp(\pi S^{45})=\gamma^4\gamma^5$ in signature (4,4), and (derived) the preservation of the quantum charge by these maps, which are of antiunitary type; \textbf{Theorem M3}: the invariant bilinear forms are $C(\alpha P_-+\beta P_+)$, all symmetric, so the anticommuting field has no Majorana mass term, with the Pin characters, the derivative terms and the charge 2 of every such term; \textbf{Theorem M4}: the chirality map turns every solution with $(m,\lambda)$ into one with $(-m,-\lambda)$, reverses the current and the energy--momentum tensor, so the pair has zero total charge and zero total energy--momentum in every gravitational field, and at the one-particle level $\gamma^8$ preserves energies and Hilbert norms and reverses Krein norms; the classical-level implication of \textbf{M5} (Proposition M5), for two classical fields governed by $\mathcal L_{m,\lambda}$ and $+\mathcal L_{-m,-\lambda}$.

\textbf{Computed (floating-point data):} only the record that the committed Kohn--Sham runs have the net occupation $N$ to $10^{-9}$ relative.

\textbf{Cited and provisional:} the Fock-level statements of Section 17.13, taken from the Stage-5 pairing report (\texttt{PAIR\_T1k\allowbreak{}rein}, all 141 checks of that report true), provisional until the Stage-5 gate has passed, including the operator identity $:\!T^{\mathrm{pair}}_{\mu\nu}\!:\ =0$ in the image-field reading, which is not a part of Proposition M5; and the caveat on their meaning recorded in the Wolfram matter-antimatter report (the image field is the same quantum system; in the independent reading, in the free Fock model, the kinetic and mass energies add, and for $\lambda\ne0$ the interaction energies have opposite signs). The two checks \texttt{MA\_M2\_cp\allowbreak{}ScopeInC\allowbreak{}urvedFie\allowbreak{}lds} and \texttt{MA\_M4\_im\allowbreak{}ageField\allowbreak{}FockMode\allowbreak{}l} belong to the checker of this commit and are not in the committed Python report (Section 17.17).

\textbf{Quoted from the literature, not derived:} every observational statement (the value of $\Omega_bh^2$, the temperature of the microwave background, the agreement with nucleosynthesis, the exclusion of antimatter domains, the ALPHA-g result), the analysis of the Standard Model, the black-body photon density, the values of the constants, the CPT theorem of ordinary quantum field theory, and the description of the proposal of Boyle, Finn and Turok.

\textbf{Assumed:} the Lagrangian of Section 17.8 for both statistics (Chapter 6), with anticommuting components for the fermion field dirac16complex; the signature (4,4) as the setting of the model, which is not the signature of observed spacetime (row L43 of Chapter 0), so that applying the scenario of M5 to our universe is a further assumption; the canonical quantization with respect to $x_4$ and the restriction of every quantum statement to the good sector (Chapter 8), where alone a positive Fock space exists, so that the quantum statements are formal (Krein space, ill-posed extra-time sector) and dirac16complex00 is treated as a classical field; Gaussian normal gauge wherever $B$ appears; the vanishing of the boundary flux for charge conservation; a fixed gravitational field as the background of the discrete symmetries of M2; and the class of maps examined in M2 (constant matrices, complex conjugation, reflections of frame directions).

\textbf{Hypotheses (not derived):} H1, H2 and H3 of Section 17.14.

\textbf{OPEN:} a creation process, rate or amplitude for a pair of universes (Chapter 16); which quantum reading, if any, gives a physical universe of mass $-M$ whose charge and energy cancel those of ours (of the two readings examined, neither does); whether the exact symmetries of unitary and antiunitary type are implemented by actual operators on the positive Fock space of the good sector; whether any combination of charge-violating terms breaks the CP family; higher-order charge-violating terms; and everything that a computation of $\eta$ would need (Section 17.16).

\textbf{The conclusion, stated plainly:} the theory as built does not solve the matter--antimatter problem. It contains no baryon-number-violating interaction (its charge is exactly conserved in every gravitational field when no charge flows through the boundary), it has an exact charge-reversing symmetry for each of its two fields (C for dirac16complex00 in every gravitational field; CP for dirac16complex in flat space and wherever the reflection is an isometry, and where a gravitational field breaks this CP the charge is still conserved), it contains no departure-from-equilibrium computation and no Standard-Model baryons, and it predicts no baryon-to-photon ratio.

\subsection{Exercises}

\textbf{Exercise 17.1.} For each process decide whether the electric charge, $B$, $L$ and $B-L$ are conserved: (a) $p\to e^++\pi^0$ (a decay that has never been observed; some extensions of the Standard Model predict it); (b) $\mu^-\to e^-+\bar\nu_e+\nu_\mu$ (the muon $\mu^-$ and the muon neutrino $\nu_\mu$ have $L=+1$); (c) $\bar p+p\to\pi^++\pi^-$.

\textbf{Exercise 17.2.} With $\eta=2.7342\times10^{-8}\,\Omega_bh^2$ (Section 17.3), compute $\eta$ and the number of photons per baryon for $\Omega_bh^2=0.0220$ and $0.0230$. Why does the coefficient not depend on $h$?

\textbf{Exercise 17.3.} A quantum state has the energy $E=T$. Compute the particle and antiparticle occupations $f_\pm=1/(e^{(E\mp\mu_B)/T}+1)$ for $\mu_B=0$ and for $\mu_B=0.1\,T$, and the surplus per state. Why must $\mu_B$ vanish when reactions that change $B$ are in equilibrium?

\textbf{Exercise 17.4.} In the toy model of Section 17.5 take $N=5000$, $r=0.52$ and $\bar r=0.48$. (a) Compute $\Delta B$. (b) Which Sakharov conditions does the toy meet by construction, and which one is hidden in the assumption that the decays run in one direction only? (c) What is $\Delta B$ if C is an exact symmetry?

\textbf{Exercise 17.5.} For commuting components and a constant $\alpha$, find the factor that multiplies each of the following under $\Psi\to e^{i\alpha}\Psi$: (a) $S=\bar\Psi\Psi$, (b) $j^\mu$, (c) $\Psi^TC\Psi$, (d) $S^2$, (e) $S\,\Psi^TC\gamma^8\Psi$. Which of them may appear in a Lagrangian with the U(1) symmetry of Theorem M1?

\textbf{Exercise 17.6.} For the commuting field $\Psi=(a\,u_++b\,u_-)\,e^{-imx_4}$ of Section 17.9 compute $J^4$, $\Psi^\dagger\Psi$ and the classical energy density $\rho$ for $(a,b)=(1,0)$, $(0,1)$, $(1/\sqrt2,1/\sqrt2)$ and $(1/\sqrt2,i/\sqrt2)$. Which charge does one quantum in the mode $u_-$ carry in the quantized theory?

\textbf{Exercise 17.7.} Take two components and the matrix $Y$ whose only nonzero entry is $Y_{01}=1$. Compute $\Psi^TY\Psi^\ast$ and $\Psi^\dagger Y^T\Psi$ for commuting and for anticommuting components, and read off the statistics sign $s$ of Lemma C.

\textbf{Exercise 17.8.} Use the rule of Theorem M2 to find $\kappa$, $\sigma$ and $c_j$ for $\mathrm C_0\mathrm P_{0123}$ ($M=C$, $R=\{0,1,2,3\}$, antilinear) and for $\mathrm P_b$ ($M=\gamma^b$, $R=\{b\}$ with $b\le3$, linear), for both statistics, and compare with the table of Section 17.11.

\textbf{Exercise 17.9.} (a) How many subsets of $\{0,1,2,3\}$ have an even number of elements, and how many an odd number? (b) Derive the number 64 of exact charge-reversing symmetries without $x_4$ reversal. (c) Why is there exactly one exact map for each allowed set $R$?

\textbf{Exercise 17.10.} (a) Show that $(CP_+)^T=CP_+$. (b) For commuting components compute $\Psi^TCP_+\Psi$ for $\Psi_8=1$, $\Psi_{12}=3$ and all other components 0, and its value after $\Psi\to e^{i\pi/2}\Psi$. (c) Show that $\Psi^TCP_+\Psi=0$ for anticommuting components.

\textbf{Exercise 17.11.} For the unit vector $u=\gamma^4$ compute $u^TCu$ and $u^T(C\gamma^8)u$ directly, and compare with item 3 of Theorem M3.

\textbf{Exercise 17.12.} Compute $\gamma^8u_-$ for the rest state $u_-$ of Section 17.9 and its Krein norm $(\gamma^8u_-)^\dagger B(\gamma^8u_-)$. What is the total charge density of the pair made of $u_-e^{-imx_4}$ and its image?

\textbf{Exercise 17.13.} Show that $\mathcal L_{m,\lambda}[\gamma^8\Psi]+\mathcal L_{-m,-\lambda}[\Psi]=0$ and $\mathcal L_{m,\lambda}[\gamma^8\Psi]+\mathcal L_{-m,\lambda}[\Psi]=-\lambda\sqrt{|g|}\,S^2$.

\textbf{Exercise 17.14.} (a) Show that $(\gamma^4\gamma^5)^2=-1$ and $\exp(\pi S^{45})=\gamma^4\gamma^5$. (b) Show that conjugation by $R=\gamma^4\gamma^5$ maps $\gamma^4\to-\gamma^4$ and $\gamma^5\to-\gamma^5$ and leaves $\gamma^0$ unchanged. (c) Why is this exact symmetry not a candidate for C or CP in Sakharov's argument?

\textbf{Exercise 17.15.} Show that $\gamma^8B^T\gamma^8=B$ and $B^T=-B$. Which of the charge conjugations $\mathrm C_0$ and $\mathrm C_8$ is of unitary type in the quantized anticommuting theory (Section 17.11), and what is still open about it?

\textbf{Exercise 17.16.} (a) Suppose H1 and H3 hold, and $Q_+=0$ at $x_4=0$. What is the baryon excess $B_+$ of our universe today? (b) Explain why the scenario M5 cannot predict $\eta$, even if H1 to H3 hold.

\textbf{Exercise 17.17.} The term $g\sqrt{|g|}\,\Psi^TC\gamma^8\gamma^\mu D_\mu\Psi$ plus its Hermitian conjugate is added to the Lagrangian of the anticommuting field, with a complex number $g\ne0$. (a) Which constant phases $e^{i\alpha}$ remain symmetries? (b) Show that the phase of $g$ can be changed at will by redefining the field. (c) Which of Sakharov's conditions would this term address, and what would still be missing?

\textbf{Exercise 17.18.} (a) Write down $\mathrm{diag}(r)$ for the CP map $\mathrm C_8\mathrm P_1$ ($R=\{1\}$) and for the map $\mathrm C_0\mathrm P_{0123}$ ($R=\{0,1,2,3\}$), and the determinant of the $4\times4$ block of the space-like directions in each case. (b) Which of the two frame changes is a frame rotation of $\mathrm{Spin}_0(4,4)$? (Hint: the rotation $\exp(\pi S^{01})$ turns the plane of $x_0$ and $x_1$ by the angle $\pi$.) (c) Is $\mathrm C_8\mathrm P_1$ a symmetry of dirac16complex in the primordial field of Chapter 9? In the generic test field G\_A?

\subsection{Answers to the exercises}

\textbf{Answer 17.1.} (a) Charge: $+1\to+1+0$, conserved. $B$: $1\to0$, violated. $L$: $0\to-1$ (the positron is an antilepton), violated. $B-L$: $1\to0-(-1)=1$, conserved. (b) Charge: $-1\to-1+0+0$, conserved; $B$: $0\to0$; $L$: $1\to1-1+1=1$, conserved; hence also $B-L$. (c) Charge: $-1+1=0\to+1-1=0$; $B$: $-1+1=0\to0$; $L$: $0\to0$: all conserved.

\textbf{Answer 17.2.} $\eta=2.7342\times10^{-8}\times0.0220=6.015\times10^{-10}$ and $2.7342\times10^{-8}\times0.0230=6.289\times10^{-10}$; the numbers of photons per baryon are $1/\eta=1.662\times10^9$ and $1.590\times10^9$. The critical density is proportional to $H_0^2$, hence to $h^2$, so the baryon mass density $\Omega_b\rho_{\mathrm{crit}}$ depends on $\Omega_b$ and $h$ only through the product $\Omega_bh^2$; this is also why the microwave background measures this product.

\textbf{Answer 17.3.} For $\mu_B=0$: $f_+=f_-=1/(e+1)=0.268941$, no surplus. For $\mu_B=0.1\,T$: $f_+=1/(e^{0.9}+1)=0.289050$ and $f_-=1/(e^{1.1}+1)=0.249740$, a surplus of $0.039311$ particles per state. A chemical potential belongs to a conserved number: the Gibbs state is built from the Hamiltonian and the conserved numbers. If reactions that change $B$ are in equilibrium, $B$ is not conserved, the equilibrium state has no chemical potential for it, $\mu_B=0$, and particles and antiparticles are equally populated.

\textbf{Answer 17.4.} (a) $\Delta B=N(r-\bar r)=5000\times0.04=200$. (b) The decays violate $B$ (condition 1), and $r\ne\bar r$ requires C and CP violation (condition 2). The departure from equilibrium (condition 3) is hidden in the assumption that the inverse processes, such as $q+q\to X$, do not run; in equilibrium they would, and by the argument of Section 17.5 the average asymmetry would be erased. (c) An exact C gives $r=\bar r$ and $\Delta B=0$.

\textbf{Answer 17.5.} (a) $e^{-i\alpha}e^{i\alpha}=1$; (b) 1; (c) $e^{2i\alpha}$; (d) 1; (e) $1\cdot e^{2i\alpha}=e^{2i\alpha}$. Only (a), (b), (d) and functions of invariant quantities may appear; (c) and (e) carry U(1) charge 2.

\textbf{Answer 17.6.} $J^4=|a|^2-|b|^2$ and $\Psi^\dagger\Psi=|a|^2+|b|^2$ (Section 17.9). For $(1,0)$: $J^4=1$; for $(0,1)$: $J^4=-1$; for $(1/\sqrt2,1/\sqrt2)$ and $(1/\sqrt2,i/\sqrt2)$: $J^4=\tfrac12-\tfrac12=0$. In all four cases $\Psi^\dagger\Psi=1$. All four fields have the frequency $\omega=m$, so $\rho=mJ^4$ (Section 7.11): $\rho=m$, $-m$, $0$ and $0$. The second field has positive frequency and negative classical energy. In the quantized theory one quantum in the positive-energy mode $u_-$ has the charge $u_-^\dagger B\,B\,u_-=u_-^\dagger u_-=1$, like every particle.

\textbf{Answer 17.7.} $\Psi^TY\Psi^\ast=\Psi_0\Psi_1^\ast$. The only nonzero entry of $Y^T$ is $(Y^T)_{10}=1$, so $\Psi^\dagger Y^T\Psi=\Psi_1^\ast\Psi_0$. For commuting components the two are equal: $s=+1$. For anticommuting components $\Psi_0\Psi_1^\ast=-\Psi_1^\ast\Psi_0$: $s=-1$.

\textbf{Answer 17.8.} $\mathrm C_0\mathrm P_{0123}$: $M=C=\Gamma_R$ with $|R|=4$, so $\varepsilon=(-1)^4=+1$; $s_R=4$, $\sigma_R=+1$. Then $\kappa=s$, $\sigma=s$ and $c_j=-s$. Commuting ($s=1$): exact, $c_j=-1$, so $Q\to-Q$. Anticommuting ($s=-1$): $\kappa=\sigma=-1$, $\mathcal L_{m,\lambda}\to-\mathcal L_{m,-\lambda}$. $\mathrm P_b$: $M=\gamma^b=\Gamma_R$ with $|R|=1$, so $\varepsilon=-1$; $s_R=1$, $\sigma_R=-1$. Then $\kappa=\sigma_R\varepsilon=+1$, $\sigma=-1$, $c_j=+1$: $\mathcal L_{m,\lambda}\to\mathcal L_{-m,\lambda}$ for both statistics. All four entries agree with the table.

\textbf{Answer 17.9.} (a) Even: $\binom40+\binom42+\binom44=1+6+1=8$; odd: $\binom41+\binom43=4+4=8$. (b) The charge-reversing exact symmetries that keep $x_4$ are the exact antilinear maps with $4\notin R$. The part $R\cap\{0,1,2,3\}$ must have even size (commuting) or odd size (anticommuting): 8 choices; the part $R\cap\{5,6,7\}$ is arbitrary: $2^3=8$ choices. In total $8\times8=64$. (c) By Lemma D, an exact map with the set $R$ has, up to a factor of modulus 1, a matrix that solves the equation of Lemma B with $\varepsilon=+1$; no other constant matrix is possible. The two solutions $\Gamma_R$ and $\Gamma_{R^c}$ of Lemma B have opposite signs $\varepsilon$, so exactly one of them has $\varepsilon=+1$.

\textbf{Answer 17.10.} (a) $(CP_+)^T=P_+^TC^T=P_+C=CP_+$, because $P_+=\mathrm{diag}(0,I_8)$ is symmetric and commutes with $C$ ($\gamma^8$ does). (b) $CP_+=\mathrm{diag}(0,\sigma)$ has the entries $+1$ at $(8,12),(9,13),(10,14),(11,15)$ and at the mirror places, so $\Psi^TCP_+\Psi=2(\Psi_8\Psi_{12}+\Psi_9\Psi_{13}+\Psi_{10}\Psi_{14}+\Psi_{11}\Psi_{15})=2\cdot1\cdot3=6$. After the phase it is $e^{i\pi}\cdot6=-6$. (c) Each pair of mirror entries gives, for example, $\Psi_8\Psi_{12}+\Psi_{12}\Psi_8=0$ for anticommuting components.

\textbf{Answer 17.11.} $\gamma^4$ is antisymmetric, $(\gamma^4)^T=-\gamma^4$, and commutes with $C$ (it anticommutes with all four factors of $C$). So $u^TCu=-\gamma^4C\gamma^4=-C(\gamma^4)^2=C$. With $n(u)=-1$ item 3 predicts $-n(u)C=C$. Next, $u^TC\gamma^8u=-\gamma^4C\gamma^8\gamma^4=-C\gamma^4\gamma^8\gamma^4=C\gamma^8\gamma^4\gamma^4=-C\gamma^8$, and item 3 predicts $n(u)C\gamma^8=-C\gamma^8$.

\textbf{Answer 17.12.} $u_-=\tfrac12(e_0+e_4+i\,e_9-i\,e_{13})$, and $\gamma^8$ reverses the upper eight components: $v=\gamma^8u_-=\tfrac12(-e_0-e_4+i\,e_9-i\,e_{13})$. With the rows of $B$: $(Bv)_0=iv_9=-\tfrac12=v_0$, $(Bv)_4=-iv_{13}=-\tfrac12=v_4$, $(Bv)_9=-iv_0=\tfrac i2=v_9$, $(Bv)_{13}=iv_4=-\tfrac i2=v_{13}$. So $Bv=v$ and the Krein norm is $+1$, as $\gamma^8B\gamma^8=-B$ requires ($u_-^\dagger Bu_-=-1$). The total charge density of the pair is $-1+1=0$.

\textbf{Answer 17.13.} With $S[\gamma^8\Psi]=S$ and $K[\gamma^8\Psi]=-K$: $\mathcal L_{m,\lambda}[\gamma^8\Psi]=\sqrt{|g|}(-K-mS-\tfrac\lambda2S^2)$. Adding $\mathcal L_{-m,-\lambda}[\Psi]=\sqrt{|g|}(K+mS+\tfrac\lambda2S^2)$ gives 0. Adding $\mathcal L_{-m,\lambda}[\Psi]=\sqrt{|g|}(K+mS-\tfrac\lambda2S^2)$ gives $-\lambda\sqrt{|g|}\,S^2$.

\textbf{Answer 17.14.} (a) $\gamma^4\gamma^5\gamma^4\gamma^5=-\gamma^4\gamma^4\gamma^5\gamma^5=-(-1)(-1)=-1$. With $S^{45}=\tfrac12\gamma^4\gamma^5$ and $J=\gamma^4\gamma^5$, $J^2=-1$, Section 2.11 gives $\exp(\pi S^{45})=\exp(\tfrac\pi2J)=\cos\tfrac\pi2+J\sin\tfrac\pi2=J$. (b) $R^{-1}=-\gamma^4\gamma^5$. Then $R\gamma^4R^{-1}=-\gamma^4\gamma^5\gamma^4\gamma^4\gamma^5=\gamma^4\gamma^5\gamma^5=-\gamma^4$ and $R\gamma^5R^{-1}=-\gamma^4\gamma^5\gamma^5\gamma^4\gamma^5=\gamma^4\gamma^4\gamma^5=-\gamma^5$. $\gamma^0$ anticommutes with both $\gamma^4$ and $\gamma^5$, so it commutes with $R$ and is unchanged. (c) It reverses $x_4$: it maps a history running forward in $x_4$ to one running backward. Proposition 17.2 needs a symmetry that maps histories to histories in the same direction of time, which C and CP do and this rotation does not.

\textbf{Answer 17.15.} $C$ is symmetric, $\gamma^4$ antisymmetric, and they commute, so $(C\gamma^4)^T=(\gamma^4)^TC=-\gamma^4C=-C\gamma^4$ and $B^T=-i(C\gamma^4)^T=iC\gamma^4=-B$. Then $\gamma^8B^T\gamma^8=-\gamma^8B\gamma^8=-(-B)=B$. By Section 17.11 a map $\Psi\to M\Psi^{\dagger T}$ gives the anticommutator $MB^TM^\dagger$: for $\mathrm C_8$ ($M=\gamma^8$) this is $B$, so $\mathrm C_8$ is of unitary type; for $\mathrm C_0$ ($M=1$) it is $-B$, so $\mathrm C_0$ is of antiunitary type although it does not reverse $x_4$. Being of unitary type is necessary for an implementation by a unitary operator, not sufficient: whether a unitary operator on the positive Fock space of the good sector implements $\mathrm C_8$ is OPEN.

\textbf{Answer 17.16.} (a) By M1 the charge is conserved, so $Q_+=0$ at all times, and by H3 $B_+=0$: no baryon excess. The excess of our universe must be assumed as an initial condition, which is the content of H2. (b) By M1, $Q_+$ never changes, so its value is an initial condition put in by H2; nothing in the repository computes it. Moreover $\eta$ counts the photons of our universe, and the theory contains no photons of the Standard Model, so even the denominator of $\eta$ is outside the theory.

\textbf{Answer 17.17.} (a) The added term is multiplied by $e^{2i\alpha}$ and its conjugate by $e^{-2i\alpha}$, so only phases with $e^{2i\alpha}=1$ remain: $\alpha\in\{0,\pi\}$, that is $\Psi\to\pm\Psi$. (b) Redefine $\Psi=e^{i\beta}\Psi'$. The original Lagrangian is unchanged (Theorem M1), and the new term becomes $ge^{2i\beta}\sqrt{|g|}\,\Psi'^TC\gamma^8\gamma^\mu D_\mu\Psi'$; choosing $2\beta=-\arg g$ makes the coefficient real and positive. Only relative phases between several charge-violating terms could have a physical meaning. (c) It addresses condition 1: the charge would no longer be conserved. For condition 2 every exact charge-reversing symmetry of the CP family would have to be broken as well, and whether this term does so is not computed. Condition 3 needs a non-equilibrium history and a computation of rates, and a link to the baryons of the Standard Model and a computation of $\eta$ would still be missing (Section 17.16).

\textbf{Answer 17.18.} (a) The two frame changes are

\[
\mathrm C_8\mathrm P_1:\ \mathrm{diag}(1,-1,1,1,1,1,1,1),\qquad \mathrm C_0\mathrm P_{0123}:\ \mathrm{diag}(-1,-1,-1,-1,1,1,1,1).
\]

The space-like block of the first is $\mathrm{diag}(1,-1,1,1)$, with determinant $-1$; that of the second is $-I_4$, with determinant $(-1)^4=+1$. (b) The second. The rotation by $\pi$ in the plane of $x_0,x_1$ reverses $x_0$ and $x_1$ and leaves the other six coordinates alone, the one in the plane of $x_2,x_3$ reverses $x_2$ and $x_3$, and their product is the frame change of $\mathrm C_0\mathrm P_{0123}$. The first cannot be a frame rotation of $\mathrm{Spin}_0(4,4)$, because every such rotation has a space-like block with determinant at least 1 (Section 17.11, "Which entries hold in a gravitational field"). (c) The primordial field has a diagonal metric that depends only on $x_0$ and $x_4$ (Section 9.2), so the reflection of $x_1$ is an isometry, and $\mathrm C_8\mathrm P_1$ in its coordinate form is an exact, charge-reversing symmetry there. G\_A is a generic non-diagonal field without this isometry: only the frame form holds, which relates the theories on the frames $e$ and $e\,\mathrm{diag}(r)$, and at a fixed frame the internal part of $\mathrm C_8\mathrm P_1$ is not a symmetry (check \texttt{MA\_M2\_cp\allowbreak{}ScopeInC\allowbreak{}urvedFie\allowbreak{}lds}). The charge is conserved in both fields (Theorem M1).

\section{Open problems and how a student could attack them}

\subsection{What this chapter is for}

A textbook usually ends where its theory is complete. This one cannot, because the theory it teaches is not complete, and the honesty rule of Section 0.2 forbids pretending otherwise. The earlier chapters marked every unanswered question with the word OPEN (Section 0.3). This chapter collects all of them: the open questions of this book and those listed in the project documents. For each one it gives five things:

\begin{enumerate}
\item \textbf{the question}, stated as precisely as we can;
\item \textbf{why it is hard}, that is, which obstacle is known and which step nobody has taken;
\item \textbf{what is known}, with the status words of Section 0.3 and the file or check where each statement is verified;
\item \textbf{a first step} that a student who has worked through this book could take, often carried out on a small example in this chapter;
\item \textbf{where to start} in the repository: the files, programs and reports that already exist.
\end{enumerate}

Nothing in this chapter is claimed to be solved. Two of the problems below, a creation amplitude for pairs of universes (Section 18.9) and baryogenesis (Section 18.10), are the two parts of the author's request that this project has not established (Section 0.2). They appear here because they are open, and the chapter says exactly why.

\textbf{A label used only in this chapter.} Several first steps below contain a short calculation of our own. Such a calculation is marked \textbf{derived here}. It is complete as written and uses only results proved earlier in the book, but no program of the repository has checked it. In the language of Section 0.3 its status is that of a derivation without a machine check, and the first task of a student who wants to build on it is to redo it and then to verify it exactly, by machine, in two independent programs, as the project does for every exact statement (Section 18.2).

\textbf{The state of the project on which the problems stand} (2026-09-30, from \texttt{README.m\allowbreak{}d}, \texttt{HANDOFF.\allowbreak{}md} section 2 and the committed reports). Stages 1 to 3 are complete, and their gates passed from fresh public clones of the repository. In Stage 4 the exact theory is complete (125 of 125 Wolfram and 157 of 157 sympy checks) and so is the Rust Kohn--Sham solver, but the final cross-check against the independent reference solver has one failed check out of 63 and the Stage-4 gate has not been run (Section 18.7). In Stage 5 the exact theory is committed and checked (Wolfram 46 of 46 and 141 of 141 checks, Python 49 of 49 and 172 of 172), the Kohn--Sham pair numerics are partial, and the Stage-5 documents are not written (Section 18.8). The exact theorems of the matter--antimatter analysis are verified (44 of 44 checks in \texttt{wolfram-\allowbreak{}matter-a\allowbreak{}ntimatte\allowbreak{}r-report\allowbreak{}.json} and 77 of 77 in \texttt{python-m\allowbreak{}atter-an\allowbreak{}timatter\allowbreak{}-report.\allowbreak{}json}, both in \texttt{artifact\allowbreak{}s/dirac1\allowbreak{}6complex\allowbreak{}/matter-\allowbreak{}antimatt\allowbreak{}er/}), and its document \texttt{provenan\allowbreak{}ce/DIRAC\allowbreak{}16COMPLE\allowbreak{}X\_MATTER\allowbreak{}\_ANTIMAT\allowbreak{}TER.md} is committed; \texttt{HANDOFF.\allowbreak{}md} still lists the review and the fixing of that document as remaining, so the analysis is not called complete here. Its answer is negative (Chapter 17).

\textbf{The list.} The table gives every problem, where it is treated and its status.

\begingroup
\small
\begin{longtable}{@{}>{\raggedright\arraybackslash}p{0.200\linewidth}>{\raggedright\arraybackslash}p{0.200\linewidth}>{\raggedright\arraybackslash}p{0.200\linewidth}>{\raggedright\arraybackslash}p{0.200\linewidth}@{}}
\toprule
\textbf{Problem} & \textbf{Section} & \textbf{Status} & \textbf{Where it arose} \\
\midrule
\endfirsthead
\toprule
\textbf{Problem} & \textbf{Section} & \textbf{Status} & \textbf{Where it arose} \\
\midrule
\endhead
An interacting quantum field theory of dirac16complex, and renormalisation & 18.3 & OPEN & Chapters 8, 12 and 13 \\
The extra-time instability, and signature (4,4) against the observed world & 18.4 & OPEN & Chapters 5, 8 and 11 \\
What determines $a_4$; Einstein--Lovelock gravity with an acceptable source & 18.5 & OPEN & Chapters 4 and 9 \\
Can a Kohn--Sham state be the source of the primordial field? & 18.6 & OPEN; no computed state is & Chapter 13 \\
The last Stage-4 cross-check and the Stage-4 gate & 18.7 & OPEN; 62 of 63 cross-checks pass & Chapter 13 \\
The Stage-5 numerics & 18.8 & OPEN; partly computed & Chapters 14 and 15 \\
A creation amplitude for pairs of universes & 18.9 & OPEN; pair creation itself is a HYPOTHESIS & Chapters 15 and 16 \\
Baryogenesis in this framework & 18.10 & OPEN; the theory as built fails Sakharov's first two conditions & Chapter 17 \\
Smaller open items of the book and the project documents & 18.11 & OPEN & various \\
\bottomrule
\end{longtable}
\endgroup

The problems are not independent. The instability of Section 18.4 is one of the obstacles to the quantum theory of Section 18.3; the gravitational field equations of Section 18.5 decide both the sourcing question of Section 18.6 and any creation amplitude of Section 18.9; and baryogenesis (Section 18.10) needs a quantum theory, a non-equilibrium history and new interactions at once. A student is well advised to start with the problems whose first step is a finite calculation: Sections 18.5, 18.6 and 18.7.

\subsection{How to work on an open problem in this repository}

The project answered every question it did answer by the same procedure, and a student who attacks an open problem should follow it, because it is what makes an answer trustworthy. The steps are these.

\begin{enumerate}
\item \textbf{Write the question down precisely}, together with what would count as an answer and what would not. The project's specifications are examples: \texttt{handoff/\allowbreak{}specs/ST\allowbreak{}AGE4\_SPE\allowbreak{}C.md}, \texttt{handoff/\allowbreak{}specs/ST\allowbreak{}AGE5\_SPE\allowbreak{}C.md} and \texttt{handoff/\allowbreak{}specs/MA\allowbreak{}TTER\_ANT\allowbreak{}IMATTER\_\allowbreak{}SPEC.md}. When a measurement contradicts the text of a specification, the specification receives an \textbf{erratum} that records the correction and the evidence (the Stage-4 specification has thirteen, E4.1 to E4.13).
\item \textbf{Derive the result by hand}, step by step, as this book does.
\item \textbf{Verify every exact statement in two independent programs}: a Wolfram Language package with its verifier and an independent Python checker built on sympy, each writing a report whose entry \texttt{\detokenize{checks}} maps check names to true or false (Section 0.8). Include \textbf{negative controls}, tests that must fail for deliberately wrong input (Section 10.11).
\item \textbf{Compute every number with two independent solvers} and compare them with a checker whose tolerances are fixed before the comparison. A tolerance is changed only on the basis of a recorded measurement, and the change is itself an erratum (E4.12 and E4.13 are examples). Repeat each run to show that it is deterministic, and repeat it with tighter tolerances to show that it has converged (Section 10.11).
\item \textbf{Never weaken a check to make it pass.} A check that fails is information; the procedure is to find out why.
\item \textbf{Write a document and a gate}: the document states what is proved, what is computed, what is assumed and what is not claimed; the gate reruns everything from a fresh copy of the repository (Chapter 19).
\end{enumerate}

The tools that already exist, and that the first steps below use, are the following.

\begingroup
\small
\begin{longtable}{@{}>{\raggedright\arraybackslash}p{0.267\linewidth}>{\raggedright\arraybackslash}p{0.267\linewidth}>{\raggedright\arraybackslash}p{0.267\linewidth}@{}}
\toprule
\textbf{Tool} & \textbf{What it does} & \textbf{Examples in the repository} \\
\midrule
\endfirsthead
\toprule
\textbf{Tool} & \textbf{What it does} & \textbf{Examples in the repository} \\
\midrule
\endhead
exact Wolfram Language packages and verifiers & prove identities with exact integers, fractions and symbols, and write a report & \texttt{wolfram/\allowbreak{}*.wl} with \texttt{scripts/\allowbreak{}verify\_*\allowbreak{}.wls} \\
independent exact Python checkers & redo the same statements with sympy, without Wolfram code & \texttt{scripts/\allowbreak{}check\_di\allowbreak{}rac16com\allowbreak{}plex\_*.p\allowbreak{}y} \\
an exact Grassmann algebra & anticommuting variables as computer objects & \texttt{scripts/\allowbreak{}grassman\allowbreak{}n\_algebr\allowbreak{}a.py} \\
Rust solvers with the CVODE engine & ordinary differential equations in $x_4$ (Stage 3) and shooting in $y$ (Stages 4 and 5) & \texttt{studies/\allowbreak{}dirac16c\allowbreak{}omplex\_c\allowbreak{}osmology}, \texttt{studies/\allowbreak{}dirac16c\allowbreak{}omplex\_k\allowbreak{}ohn\_sham} \\
independent numpy reference solvers & a second numerical method for the Kohn--Sham problems & \texttt{scripts/\allowbreak{}ks\_refer\allowbreak{}ence\_sol\allowbreak{}ver.py}, \texttt{scripts/\allowbreak{}ks\_refer\allowbreak{}ence\_pai\allowbreak{}rs.py} \\
checkers of numerical results & compare two solvers, test identities, tolerances fixed in advance & \texttt{scripts/\allowbreak{}check\_di\allowbreak{}rac16com\allowbreak{}plex\_koh\allowbreak{}n\_sham.p\allowbreak{}y}, \texttt{scripts/\allowbreak{}check\_di\allowbreak{}rac16com\allowbreak{}plex\_pai\allowbreak{}rs.py} \\
gates & rerun a whole stage from a fresh clone & \texttt{scripts/\allowbreak{}verify\_s\allowbreak{}tage1\_ar\allowbreak{}bitrary\_\allowbreak{}field.ps\allowbreak{}1} to \texttt{scripts/\allowbreak{}verify\_s\allowbreak{}tage4\_ko\allowbreak{}hn\_sham.\allowbreak{}ps1}, each with a \texttt{\detokenize{.sh}} twin \\
the document builder & Markdown to LaTeX to a reproducible PDF & \texttt{scripts/\allowbreak{}build\_pr\allowbreak{}ovenance\allowbreak{}\_pdf.py} \\
\bottomrule
\end{longtable}
\endgroup

Chapter 19 lists the commands that run each of them in PowerShell and in Bash.

\subsection{Problem 1: an interacting quantum field theory, and renormalisation}

\textbf{The question.} Chapter 8 quantized dirac16complex canonically, but only as a free field: the anticommutator, the Krein space, the Fock space of the good sector and the expectation-value rule are statements about the theory with $U=0$. The interaction $U(S)=\tfrac\lambda2S^2$ entered afterwards only through mean-field approximations: the homogeneous condensates of Chapter 11 and the Hartree--Fock and exchange-only Kohn--Sham functionals of Chapter 13. The open question is whether there is a quantum theory of the \textbf{interacting} field: a space of states with positive probabilities, a Hamiltonian that is bounded below, and predictions that are finite and do not depend on how a calculation is cut off at high energy. The ledger of Section 0.9 lists it as OPEN, and the Stage-1 document says the same (§1.3, item 4: no interacting Fock space or renormalization is constructed).

\textbf{Why it is hard.} There are four separate obstacles, and each would be serious on its own.

\emph{(1) The indefinite inner product.} No positive Hilbert space carries the canonical anticommutator (Theorem 8.1). The state space is a Krein space, and every Spin(4,4)-invariant charge density is indefinite (Sections 8.7 and 8.8). Probabilities in quantum theory are squared lengths in a positive space (Section 8.2). The free theory escapes in the good sector, where a positive Fock space exists (Section 8.10), but nobody has shown that an interaction keeps a positive subspace of states invariant.

\emph{(2) Power counting.} Measure every quantity in powers of a mass, with $\hbar=c=1$ (Section 0.6). A length has mass dimension $-1$, so the volume element $d^Dx$ has dimension $-D$, and because the action $\int\mathcal L\,d^Dx$ is a pure number, the Lagrangian density has dimension $D$. The kinetic term $\bar\Psi\gamma^\mu\partial_\mu\Psi$ contains two fields and one derivative (dimension 1), so $2[\Psi]+1=D$, that is

\[
[\Psi]=\frac{D-1}2 .
\]

The mass term $m\bar\Psi\Psi$ then gives $[m]+2[\Psi]=D$, so $[m]=1$, as it should be. The interaction $\lambda(\bar\Psi\Psi)^2$ contains four fields: $[\lambda]+4[\Psi]=D$, so

\[
[\lambda]=D-2(D-1)=2-D,\qquad [\lambda]=-6\ \text{ for }D=8 .
\]

This is the value found in Section 13.8 (check \texttt{KS\_excha\allowbreak{}nge\_coup\allowbreak{}lingDime\allowbreak{}nsion}). A coupling of negative mass dimension makes the dimensionless strength of the interaction at an energy $E$, the combination $\lambda E^{D-2}=\lambda E^6$, grow without bound. The standard power-counting argument of quantum field theory (chapter 10 of M. E. Peskin and D. V. Schroeder, \emph{An Introduction to Quantum Field Theory}, Addison-Wesley 1995; quoted, not derived here) concludes that every higher order of perturbation theory in $\lambda$ needs new kinds of counterterms. (A \textbf{counterterm} is a new term added to the Lagrangian, with a coefficient chosen to absorb a correction that depends on the cutoff, the largest momentum kept in the calculation. A theory is renormalisable when finitely many kinds of counterterms suffice, whose effect can be absorbed into the constants already present, such as $m$ and $\lambda$.) So the theory is \textbf{not renormalisable}. Section 13.8 used the same fact to explain why the Kohn--Sham functional has no correlation term. The four-fermion theory of the weak interaction in four dimensions has the same problem ($[\lambda]=2-4=-2$); it is understood today as the low-energy limit of a renormalisable theory. Whether dirac16complex has such a completion is not known.

\emph{(3) The interaction does not respect the good sector} (derived here). On a slice $x_4=\text{const}$ take the seven slice coordinates periodic, so that every field is a sum of plane waves $e^{i\mathbf k\cdot\mathbf x}$ with allowed wave numbers $\mathbf k=(k_0,k_1,k_2,k_3,k_5,k_6,k_7)$. A term of $\int(\bar\Psi\Psi)^2\,d^7x$ that removes two quanta with wave numbers $\mathbf k_1,\mathbf k_2$ and creates two with $\mathbf k_3,\mathbf k_4$ contains the integral of $e^{i(\mathbf k_1+\mathbf k_2-\mathbf k_3-\mathbf k_4)\cdot\mathbf x}$ over the slice. On a period of length $\ell$ the integral $\int_0^\ell e^{2\pi inx/\ell}dx$ is zero for every integer $n\ne0$, so the term vanishes unless $\mathbf k_1+\mathbf k_2=\mathbf k_3+\mathbf k_4$ in each of the seven directions: momentum is conserved. It is conserved, but it does not forbid two quanta without extra-time momentum ($k_5=0$) from turning into two quanta with $k_5=+q$ and $k_5=-q$. Energy does not forbid it either: by the dispersion relation $E^2=m^2+k_0^2+\dots+k_3^2-k_5^2-k_6^2-k_7^2$ of Section 8.9, two quanta at rest have the energy $2m$, and two quanta with the wave numbers $(k_1,k_5)=(p,q)$ and $(-p,-q)$ have $2\sqrt{m^2+p^2-q^2}$, which equals $2m$ whenever $p^2=q^2$. The final quanta carry extra-time momentum, so their mode Hamiltonian is not Hermitian (Section 8.9) and they lie outside the positive Fock space of Section 8.10. This is a kinematic statement, not a rate: whether the matrix element of such a process vanishes for another reason, for example because of the Krein structure, is part of the open problem.

\emph{(4) No Euclidean shortcut.} The usual tool of quantum field theory for defining and computing loop corrections replaces the time by an imaginary time, which turns the one minus sign of the four-dimensional metric into a plus and makes the metric positive. Here the same step applied to $x_4$ removes only one of the four minus signs of $\eta$: the signature becomes (5,3), still indefinite, and the standard Euclidean methods do not apply directly.

\textbf{What is known.}

\begin{itemize}
\item PROVED: the Krein structure is forced, $B=-iC\gamma^4$ is Hermitian with eight eigenvalues of each sign, and every invariant density is indefinite (Theorem 8.1, Section 8.8; checks \texttt{ALG\_char\allowbreak{}geFormB}, \texttt{ALG\_inva\allowbreak{}riantFor\allowbreak{}ms} and \texttt{QNT\_krei\allowbreak{}nSignatu\allowbreak{}re} in both algebra reports of \texttt{artifact\allowbreak{}s/dirac1\allowbreak{}6complex\allowbreak{}/arbitra\allowbreak{}ry-field\allowbreak{}/}).
\item PROVED by derivation, without a machine check: a positive Fock space and a normal-ordered Hamiltonian $H\ge0$ in the good sector (Section 8.10; Stage-1 document, §10.7).
\item PROVED: the first order of perturbation theory in $\lambda$ for quasi-free states, the Hartree--Fock energy $\tfrac\lambda2\bigl[S^2-\mathrm{Tr}(BC\rho BC\rho)\bigr]$ and its exactly local uniform-gas form $e_x=-\tfrac\lambda{32}(n^2+S^2)$ (Section 13.8; checks \texttt{KS\_excha\allowbreak{}nge\_wick\allowbreak{}TheoremH\allowbreak{}F} and \texttt{KS\_excha\allowbreak{}nge\_unif\allowbreak{}ormGasCl\allowbreak{}osedForm}). For commuting components the exchange term changes sign (Chapter 14). The checks of the sign are \texttt{PAIR\_sta\allowbreak{}t\_fermio\allowbreak{}nWickMin\allowbreak{}us} and \texttt{PAIR\_sta\allowbreak{}t\_classi\allowbreak{}calGauss\allowbreak{}ianWickP\allowbreak{}lus} of the Stage-5 Wolfram pairing report and their twins with the prefix \texttt{\detokenize{S5_}} in the independent Python pairing report.
\item PROVED: an exact finite Fock model with Krein anticommutators: four modes of the rest sector with $m=1$, a Dirac sea made of the two negative-energy modes, and normal ordering as the subtraction of the sea value (checks \texttt{PAIR\_T1k\allowbreak{}rein\_*} in the Wolfram pairing report and \texttt{S5\_T1kre\allowbreak{}in\_*} in the Python one). It is a free model.
\item PROVED: the classical energy of the commuting field dirac16complex00 is not bounded below (check \texttt{C00\_ener\allowbreak{}gy\_unbou\allowbreak{}ndedBelo\allowbreak{}w} in \texttt{wolfram-\allowbreak{}dirac16c\allowbreak{}omplex00\allowbreak{}-report.\allowbreak{}json} and \texttt{S5\_energ\allowbreak{}y\_unboun\allowbreak{}dedBelow} in \texttt{python-d\allowbreak{}irac16co\allowbreak{}mplex00-\allowbreak{}report.j\allowbreak{}son}, both in the Stage-5 folder \texttt{artifact\allowbreak{}s/dirac1\allowbreak{}6complex\allowbreak{}/pair-cr\allowbreak{}eation/}). For the anticommuting field the positivity of the good sector comes from the Pauli principle and the filled Dirac sea (Section 8.10), a mechanism that commuting components do not have (\texttt{dirac16c\allowbreak{}omplex00\allowbreak{}-theory.\allowbreak{}json}, key \texttt{commutin\allowbreak{}gFieldSp\allowbreak{}ecifics}). No spin--statistics theorem for signature (4,4) is proved (Section 8.15).
\item ASSUMED: the Kohn--Sham functional of Chapter 13 is exchange-only, without correlation (Section 13.8).
\end{itemize}

\textbf{Worked example: how coupling two modes of opposite Krein sign can make frequencies complex} (derived here). Take two modes with the Krein form $B=\mathrm{diag}(1,-1)$, the two-mode toy of Section 8.7. A mode Hamiltonian $h$ conserves the Krein norm $u^\dagger Bu$ exactly when $h^\dagger B=Bh$ (the computation of Section 8.9). Write $h=\begin{pmatrix}a&b\\c&d\end{pmatrix}$ with complex entries. Then

\[
h^\dagger B=\begin{pmatrix}a^\ast&-c^\ast\\b^\ast&-d^\ast\end{pmatrix},\qquad Bh=\begin{pmatrix}a&b\\-c&-d\end{pmatrix},
\]

and the two agree exactly when $a$ and $d$ are real and $c=-b^\ast$. So the most general Krein-norm-conserving $h$ is

\[
h=\begin{pmatrix}\omega_0&g\\-g^\ast&\omega_1\end{pmatrix},
\]

where $\omega_0$ and $\omega_1$ are the frequencies of the two uncoupled modes and $g$ couples them. Its eigenvalues solve $(\omega_0-x)(\omega_1-x)+\lvert g\rvert^2=0$:

\[
x=\frac{\omega_0+\omega_1}2\pm\sqrt{\Bigl(\frac{\omega_0-\omega_1}2\Bigr)^2-\lvert g\rvert^2}.
\]

They are real when $\lvert g\rvert\le\lvert\omega_0-\omega_1\rvert/2$. For a stronger coupling they form a complex pair, and the solution $e^{-ixt}$ of the eigenvalue with positive imaginary part grows like $e^{\varkappa t}$ with $\varkappa=\sqrt{\lvert g\rvert^2-(\omega_0-\omega_1)^2/4}$. Compare two modes of the same Krein sign ($B=1$): then the condition is $h^\dagger=h$, that is $c=b^\ast$, the eigenvalues are $\tfrac{\omega_0+\omega_1}2\pm\sqrt{(\omega_0-\omega_1)^2/4+\lvert g\rvert^2}$, and they are real for every coupling. With numbers: $\omega_0=3$, $\omega_1=1$, so the threshold is $\lvert g\rvert=1$. For $g=\tfrac12$ the eigenvalues are $2\pm\sqrt{3}/2=2.866025$ and $1.133975$; for $g=2$ they are $2\pm i\sqrt3$, and the growth rate is $\varkappa=\sqrt3=1.732051$. Exercise 8.11 was the case $\omega_0=\omega_1=0$, $g=1$.

The toy explains the extra-time instability of Chapter 8 exactly. At rest the mode Hamiltonian is $h=-im\gamma^4-k_5\gamma^4\gamma^5$ for a wave number $k_5$ along $x_5$. Its first term commutes with $B$ and has the eigenvalues $\pm m$. Its second term anticommutes with $B$ (Section 8.9) and with $\gamma^4$, because $\gamma^4(\gamma^4\gamma^5)=-\gamma^5=-(\gamma^4\gamma^5)\gamma^4$; so it connects a mode of energy $+m$ and one Krein sign with a mode of energy $-m$ and the other. The toy with $\omega_0=m$, $\omega_1=-m$ and $\lvert g\rvert=\lvert k_5\rvert$ has the eigenvalues $\pm\sqrt{m^2-k_5^2}$, exactly the frequencies $\pm E$ of Section 8.9 at zero 3-momentum, and its frequencies become complex exactly when $k_5^2>m^2$ (Exercise 18.3). At $k_5^2=m^2$ the two eigenvalues coincide at 0, the matrix squares to zero and the mode grows linearly, as for the sample $m=1$, $k_5=1$ of Section 8.9 (Exercise 8.6); part (c) of Exercise 18.3 treats this case in the toy, and the case $g=1$ of Exercise 18.2 is the same situation with the double eigenvalue 2. The lesson for an interacting theory is that an interaction which couples modes of opposite Krein sign can destabilise the system once the coupling exceeds half the difference of the frequencies. Point (3) above showed only that the contact interaction may scatter quanta of the good sector into modes with extra-time momentum, which is a kinematic statement; it said nothing about Krein signs. Whether the contact interaction, through those extra-time modes, couples good-sector quanta to modes of opposite Krein sign is exactly what steps (b) and (c) of the first step below test.

\textbf{A first step.} (a) Redo the power counting and the toy (Exercises 18.1 to 18.3). (b) Take the exact four-mode Krein Fock model of the Stage-5 pairing package (\texttt{wolfram/\allowbreak{}Dirac16C\allowbreak{}omplexPa\allowbreak{}iring.wl}, function \texttt{checkT1K\allowbreak{}rein}), add the normal-ordered contact interaction restricted to its four modes, and compute the spectrum of the resulting finite matrix exactly as a function of $\lambda$. In the rest sector the matrix $BC=-i\gamma^4$ of the scalar density commutes with $B$ (Section 8.6), so one expects no mixing of Krein signs there; verify this. (c) Add two modes with extra-time wave numbers $+q$ and $-q$ (their mode Hamiltonian is that of Section 8.9) and the part of the contact interaction that scatters two rest quanta into them (point 3), and follow the eigenvalues as $\lambda$ and $q$ vary: do complex pairs appear, as in the toy? Every step is a finite exact computation and can be checked in sympy as well. (d) For renormalisation, compute the second-order (correlation) energy of the uniform gas of Section 13.8 with an upper cutoff $\Lambda$ on all momenta, and determine how it depends on $\Lambda$; power counting predicts that the dependence cannot be absorbed into $m$ and $\lambda$ (quoted, not derived).

\textbf{Where to start.} The Stage-1 document, §10, for the canonical quantization, and these files:

\begin{Verbatim}[fontsize=\small]
artifacts/dirac16complex/arbitrary-field/wolfram-algebra-report.json
artifacts/dirac16complex/arbitrary-field/python-algebra-report.json
wolfram/Dirac16ComplexAlgebra.wl          with scripts/verify_dirac16complex_algebra.wls
scripts/check_dirac16complex_algebra.py   the independent algebra checker
wolfram/Dirac16ComplexPairing.wl          the four-mode Krein Fock model (checkT1Krein)
scripts/check_dirac16complex_pairing.py   its independent twin
artifacts/dirac16complex/pair-creation/wolfram-pairing-report.json
artifacts/dirac16complex/pair-creation/python-pairing-report.json
wolfram/Dirac16ComplexKohnSham.wl         Wick's theorem on an exact Fock space
scripts/grassmann_algebra.py              anticommuting variables
\end{Verbatim}

\subsection{Problem 2: the extra times, their instability, and the observed world}

\textbf{The question.} Two linked questions. (a) The world we observe has one time and three space directions. The theory lives in four space-like and four time-like directions, signature (4,4) (Section 1.12), and the repository says plainly that this is not observed spacetime (\texttt{README.m\allowbreak{}d}, "Scientific boundary"). How, if at all, does such a theory describe the observed world? (b) Waves with momentum along the extra times $x_5,x_6,x_7$ do not oscillate but grow (Section 8.9, experiment EXP-5), and every quantum statement and every Kohn--Sham result of this book is restricted \textbf{by hand} to the good sector, where nothing depends on the extra times (Section 8.13). Is there a mechanism that suppresses, removes or freezes these modes, and a formulation of the theory whose initial-value problem is well posed, that is, whose solutions exist, are unique and depend continuously on the initial data?

\textbf{Why it is hard.}

\begin{itemize}
\item \emph{The equation is ultrahyperbolic.} The slices $x_4=\text{const}$ contain three time-like directions. For $E^2<0$ the growth rate $\varkappa=\sqrt{k_5^2+k_6^2+k_7^2-m^2-k_0^2-\dots-k_3^2}$ of Section 8.9 is unbounded as the extra-time momentum grows, so an arbitrarily small change of the initial data with a short enough wavelength grows arbitrarily fast: the solution does not depend continuously on its data. This is the classical argument against several time directions (see, for example, M. Tegmark, "On the dimensionality of spacetime", Class. Quantum Grav. 14, L69 (1997)).
\item \emph{The energy is not bounded below}: for a scalar field in 4+4 dimensions (Section 5.7), for the commuting field dirac16complex00 (check \texttt{C00\_ener\allowbreak{}gy\_unbou\allowbreak{}ndedBelo\allowbreak{}w}), and for dirac16complex outside the good sector.
\item \emph{The notebook's background makes it worse.} When 3-space inflates, the extra times deflate, the physical extra-time momentum grows like $e^{a_4}$, and every mode with extra-time momentum eventually passes the onset of growth (Section 8.13).
\item \emph{The extra dimensions are not stabilised.} In EXP-2, 8-dimensional Einstein gravity with this matter ends with all seven directions other than $x_4$ expanding equally, so that an observer in 3-space sees the effective equation of state $w_{\mathrm{eff}}\to4/3$; EXP-3 and EXP-4 simply assume static extra dimensions (Section 11.13).
\item \emph{The brane does not remove the extra times.} The metric induced on the brane of the Z2 construction (Section 9.16) is $h=\mathrm{diag}(e^{2a_{4,0}},e^{2a_{4,0}},e^{2a_{4,0}},-1,-e^{-2a_{4,0}},-e^{-2a_{4,0}},-e^{-2a_{4,0}})$: three positive and four negative entries, signature (3,4). The Kohn--Sham particles do gather near the brane (between 86.7 and 96.2 percent of them lie within $1/H$ of it for $m=1$, Section 13.11), which hides the direction $x_0$; but a world on the brane still has four time-like directions.
\end{itemize}

\textbf{What is known.}

\begin{itemize}
\item PROVED: the flat-space mode Hamiltonian satisfies $h_k^2=E^2$, it is Hermitian exactly when there is no extra-time momentum, and it always conserves the Krein norm (Section 8.9; check \texttt{QNT\_flat\allowbreak{}ModeHami\allowbreak{}ltonian} in both algebra reports).
\item COMPUTED (EXP-5, Section 8.13; the numbers are in \texttt{artifact\allowbreak{}s/dirac1\allowbreak{}6complex\allowbreak{}/numeric\allowbreak{}s/exp5/s\allowbreak{}ummary.j\allowbreak{}son}): for the comoving extra-time momenta $q=0.05$ and $0.1$ ($m=1$) the growth starts at $t_\ast=\ln(m/q)=2.99573$ and $2.30259$; the Hilbert norm grows to $1.258\times10^{16}$ and $1.313\times10^{16}$ by $x_4=t_\ast+3$, while the Krein norm drifts by at most $1.08\times10^{-9}$ relative.
\item ASSUMED: the restriction to the good sector (Section 8.13).
\item COMPUTED: the EXP-2 result on the extra dimensions quoted above (Section 11.13).
\end{itemize}

\textbf{Worked example 1: which modes keep oscillating} (derived here, with the frozen-coefficient estimate of Section 8.13). In the notebook's field a mode with 3-momentum $k_1$ along $x_1$ and extra-time momentum $k_5$ along $x_5$ has, with frozen coefficients,

\[
E^2=M_{\mathrm{eff}}^2+K^2+\bigl(k_1e^{-H\zeta-a_4}\bigr)^2-\bigl(k_5e^{-H\zeta+a_4}\bigr)^2
\]

(Section 8.13), with further terms of the same two kinds for $k_2,k_3$ and $k_6,k_7$. Three histories of $a_4$ behave differently.

\begin{enumerate}
\item \emph{$a_4$ grows without bound} (the notebook's $a_4=t$). The factor $e^{a_4}$ grows and $e^{-a_4}$ shrinks. If any extra-time momentum is nonzero, the last term eventually dominates and $E^2<0$: the mode grows from then on. If all extra-time momenta vanish, $E^2\ge M_{\mathrm{eff}}^2+K^2\ge0$ at every time. So in this history \textbf{the good sector is exactly the set of modes that oscillate forever}.
\item \emph{$a_4$ constant} (the static member of Chapter 13). Then $E^2$ does not change with time at a fixed position $\zeta$, and a mode that oscillates at the start oscillates forever. The set of oscillating modes is larger than the good sector: it contains every mode with $E^2\ge0$.
\item \emph{$a_4$ decreasing without bound.} Now the extra-time terms shrink, and every mode eventually oscillates.
\end{enumerate}

The good sector is therefore singled out dynamically only by the notebook's inflating history, and even there generic initial data contain the other modes; the calculation identifies the good sector, it does not remove the rest.

\textbf{Worked example 2: compact extra times do not help} (derived here). One might hope that rolling up the extra times into small circles, as is done with extra space dimensions, would remove the dangerous modes. Let $x_5$ be a circle of circumference $2\pi R$. A field on it is periodic, so $k_5=n/R$ with an integer $n$. In flat space

\[
E^2=m^2+\lvert\mathbf k\rvert^2-\frac{n^2}{R^2},
\]

where $\lvert\mathbf k\rvert^2=k_0^2+\dots+k_3^2$. For every radius $R$ and every $\mathbf k$ the modes with $\lvert n\rvert>R\sqrt{m^2+\lvert\mathbf k\rvert^2}$ have $E^2<0$ and grow. A small circle makes the first growing mode appear at a smaller $\lvert n\rvert$, not later: with $m=1$, $\mathbf k=0$ and $R=2$, the modes $n=\pm1$ oscillate ($E^2=\tfrac34$), $n=\pm2$ grow linearly ($E^2=0$) and every $\lvert n\rvert\ge3$ grows exponentially (Exercise 18.5).

\textbf{A known idea from the literature.} For the constant-coefficient ultrahyperbolic wave equation, W. Craig and S. Weinstein showed that the initial-value problem becomes well posed when the initial data satisfy a nonlocal constraint ("On determinism and well-posedness in multiple time dimensions", Proc. R. Soc. A 465, 3023 (2009)). For the mode equation of Section 8.9 the natural constraint of this kind keeps only the modes with $E^2\ge0$. Worked example 1 shows what that means in the notebook's background: with an inflating $a_4$, the only modes that satisfy such a constraint for all times are those of the good sector.

\textbf{A first step.} (a) Extend the reduced Kohn--Sham equation of Section 13.3 by an extra-time momentum and study its spectrum. With the ansatz factor $e^{ik_5x_5}$ and the curved gamma $\gamma^{x_5}=e^{-Hy+a_{4,0}}\gamma^5$ of Section 13.3, the reduced equation acquires the term $-i\,e^{-Hy+a_{4,0}}k_5\,\gamma^0\gamma^5\chi$ on its right-hand side. By the two rules of Section 13.4, $\gamma^0\gamma^5$ anticommutes with $J=\gamma^0\gamma^1\gamma^4$ and with $K_2=\gamma^5\gamma^6$ and commutes with $K_1=\gamma^2\gamma^3$ (derived here; Exercise 18.4 asks for the details). It therefore maps the block $(j,s_2,s_3)$ onto $(-j,s_2,-s_3)$: the eight $2\times2$ blocks couple in pairs into four $4\times4$ systems, and the two-component methods of Chapter 13 (the Prüfer angle of Section 13.10) no longer apply. A matrix method such as that of \texttt{scripts/\allowbreak{}ks\_refer\allowbreak{}ence\_sol\allowbreak{}ver.py} shows complex eigenvalues directly; the question is whether they appear and where the corresponding modes live in $y$. (b) Reproduce the EXP-5 onset times for other momenta with the committed program. (c) For question (a) of the problem, study the literature on theories with several times, and ask which structure could turn the three extra times into something unobservable; the observations above (the brane has signature (3,4), compactification does not help) rule out the two simplest answers.

\textbf{Where to start.} The Stage-1 document, §10.10, the Stage-2 document, §14.3, and these files:

\begin{Verbatim}[fontsize=\small]
artifacts/dirac16complex/numerics/exp5/            the EXP-5 outputs
studies/dirac16complex_cosmology/src/exp5.rs       the EXP-5 program
scripts/check_dirac16complex_exp5.py               its independent checker
studies/dirac16complex_cosmology/src/exp2.rs       EXP-2 (8-dimensional Einstein gravity)
studies/dirac16complex_kohn_sham/src/blocks.rs     the 2x2 block reduction
studies/dirac16complex_kohn_sham/src/shooting.rs   the shooting method
scripts/ks_reference_solver.py                     the matrix method
\end{Verbatim}

\subsection{Problem 3: what determines \texorpdfstring{$a_4$}{a\_4}, and Einstein--Lovelock gravity}

\textbf{The question.} The primordial field contains one arbitrary function, $a_4(t)$ (Section 9.2). The notebook intends it to be fixed by the Einstein--Lovelock equations of eight dimensions, with the Lovelock terms of orders 2 and 3 (cell 14; Sections 4.9 and 9.9), but no code of the notebook defines those terms and the project does not compute them either. In 8-dimensional Einstein gravity the field needs the energy density $\rho_{\mathrm{req}}=-3H^2(7+a_4'^2)/\kappa\le-21H^2/\kappa<0$ for every $a_4$ (Section 9.10), and the weak, null and dominant energy conditions fail (Section 9.11). The open questions: (a) do the Einstein--Lovelock equations with the orders 2 and 3 admit the notebook's field as a vacuum solution, or with a source of positive energy? (b) Which equation fixes $a_4$? (c) Can the slopes $a_4'=\tfrac23(M\mp1)$ that the notebook lists in cell 150 without derivation (Section 9.10) be derived?

\textbf{Why it is hard.} The Lovelock tensors of orders 2 and 3 are quadratic and cubic in the curvature; written with the generalized Kronecker delta (Section 4.9) they contain determinants with five and seven rows, and for a time-dependent $a_4$ they give nonlinear equations in $a_4'$ and $a_4''$. At the brane of the Z2 construction the Israel junction condition of Section 9.16 holds only for Einstein gravity; with a Gauss--Bonnet term it acquires further terms (derived in the literature, for example S. C. Davis, "Generalized Israel junction conditions for a Gauss--Bonnet brane world", Phys. Rev. D 67, 024030 (2003)). What "physically acceptable" means in signature (4,4) is itself unsettled, since the energy conditions of Section 9.11 are conditions for one observer and one choice of null vectors, not a complete classification. And the normalisation that the notebook intends for its Lovelock1, Lovelock2 and Lovelock3 is not recorded.

\textbf{What is known.}

\begin{itemize}
\item PROVED (Chapter 9, Stage-2 checks \texttt{P\_einste\allowbreak{}in\_requi\allowbreak{}redSourc\allowbreak{}e}, \texttt{P\_einste\allowbreak{}in\_rhoRe\allowbreak{}quiredNe\allowbreak{}gative}): the Einstein tensor for every $a_4$ and the negative required energy density.
\item PROVED (Proposition 9.2; check \texttt{P\_source\allowbreak{}\_einstei\allowbreak{}nTransve\allowbreak{}rseDiffe\allowbreak{}renceIs2\allowbreak{}H2a4pp}): in Einstein gravity no dirac16complex state that depends on $x_0$ and $x_4$ only can be the source where $a_4''\ne0$. Read the other way round (derived here, one line): with such a source the field equations \textbf{force} $a_4''=0$, so $a_4$ is linear, $a_4=ct$ plus a constant; and by condition (3) of Proposition 9.4, $\lambda S^2=2H^2(15-3c^2)/\kappa$, the slope is fixed by the source, $c^2=5-\kappa\lambda S^2/(6H^2)$. In Einstein gravity with a homogeneous dirac16complex source, $a_4$ is therefore determined up to the sign of its slope and an additive constant, but only for a source of negative energy density.
\item PROVED (Section 4.9): in eight dimensions only the Lovelock orders $k=0,1,2,3$ exist.
\item PROVED (Section 9.15; checks \texttt{KS\_geome\allowbreak{}try\_cons\allowbreak{}tantCurv\allowbreak{}atureSev\allowbreak{}enSpace} and \texttt{KS\_geome\allowbreak{}try\_kret\allowbreak{}schmannC\allowbreak{}onstant} in \texttt{wolfram-\allowbreak{}kohn-sha\allowbreak{}m-report\allowbreak{}.json}): the static member ($a_4$ constant) is the product of the flat time line $x_4$ with a seven-dimensional space of constant curvature $-H^2$ in the directions $y,x_1,x_2,x_3,x_5,x_6,x_7$.
\end{itemize}

\textbf{Worked example: the Lovelock tensors of the static member} (derived here). No program of the repository computes the Lovelock terms of orders 2 and 3, and Chapter 16 records them as not computed. The derivation below is a first step toward changing that status, and the status changes only when the result has been checked exactly by machine. The last fact of the list above makes every Lovelock tensor of the static member computable by counting. Call the seven directions $y,x_1,x_2,x_3,x_5,x_6,x_7$ the set $\Sigma$ ($n=7$ elements) and put $K=-H^2$.

\emph{Step 1: the curvature.} Constant curvature $K$ means $R_{rsmn}=K(g_{rm}g_{sn}-g_{rn}g_{sm})$ for indices in $\Sigma$ (Section 9.15). Raising indices as in Section 4.9, $R^{\alpha\beta}{}_{\mu\nu}=g^{\beta\lambda}R^\alpha{}_{\lambda\mu\nu}=K\bigl(\delta^\alpha_\mu\delta^\beta_\nu-\delta^\alpha_\nu\delta^\beta_\mu\bigr)=K\,\delta^{\alpha\beta}_{\mu\nu}$ for $\alpha,\beta,\mu,\nu\in\Sigma$. Every component with an index 4 vanishes, because $x_4$ is a flat factor.

\emph{Step 2: the contraction.} In $E^{(k)\mu}{}_\nu=-\frac1{2^{k+1}}\,\delta^{\mu\,\alpha_1\beta_1\cdots\alpha_k\beta_k}_{\nu\,\gamma_1\delta_1\cdots\gamma_k\delta_k}R^{\gamma_1\delta_1}{}_{\alpha_1\beta_1}\cdots R^{\gamma_k\delta_k}{}_{\alpha_k\beta_k}$ (Section 9.9; the letters $\alpha_i,\beta_i,\gamma_i,\delta_i$ are summed indices here) insert $R^{\gamma_i\delta_i}{}_{\alpha_i\beta_i}=K\bigl(\delta^{\gamma_i}_{\alpha_i}\delta^{\delta_i}_{\beta_i}-\delta^{\gamma_i}_{\beta_i}\delta^{\delta_i}_{\alpha_i}\bigr)$. Summing over $\gamma_i,\delta_i$ replaces them in the generalized delta by $\alpha_i,\beta_i$, once directly and once in exchanged order; the exchange of two lower indices flips the sign of the determinant, so the two terms add. Each factor gives $2K$, and

\[
E^{(k)\mu}{}_\nu=-\frac{K^k}2\,\delta^{\mu\,\alpha_1\beta_1\cdots\alpha_k\beta_k}_{\nu\,\alpha_1\beta_1\cdots\alpha_k\beta_k},
\]

where the $2k$ repeated indices are summed over $\Sigma$ only.

\emph{Step 3: counting.} The generalized delta with the same list of indices above and below is $+1$ when all its indices are different (the determinant of the unit matrix) and 0 when two coincide (two equal rows). For $\mu\ne\nu$ the lower list would have to be a rearrangement of the upper one, which needs $\nu=\mu$, so all off-diagonal components vanish. For $\mu=\nu$ the sum counts the ordered lists of $p=2k$ different elements of $\Sigma$ that are also different from $\mu$. If $\mu\in\Sigma$ there are $n-1=6$ elements to choose from, and the number of such lists is $6\cdot5\cdots(7-p)=6!/(6-p)!$. If $\mu=4$, which is not in $\Sigma$, there are $n=7$, and the number is $7!/(7-p)!$. Hence, with no sum over $\mu$,

\[
E^{(k)\mu}{}_\mu=-\frac{K^k}2\cdot\frac{6!}{(6-2k)!}\ \ (\mu\in\Sigma),\qquad E^{(k)4}{}_4=-\frac{K^k}2\cdot\frac{7!}{(7-2k)!}.
\]

With $K=-H^2$:

\begingroup
\small
\begin{longtable}{@{}>{\raggedright\arraybackslash}p{0.160\linewidth}>{\raggedright\arraybackslash}p{0.160\linewidth}>{\raggedright\arraybackslash}p{0.160\linewidth}>{\raggedright\arraybackslash}p{0.160\linewidth}>{\raggedright\arraybackslash}p{0.160\linewidth}@{}}
\toprule
\textbf{$k$} & \textbf{$6!/(6-2k)!$} & \textbf{$7!/(7-2k)!$} & \textbf{$E^{(k)\mu}{}_\mu$ for each $\mu\in\Sigma$} & \textbf{$E^{(k)4}{}_4$} \\
\midrule
\endfirsthead
\toprule
\textbf{$k$} & \textbf{$6!/(6-2k)!$} & \textbf{$7!/(7-2k)!$} & \textbf{$E^{(k)\mu}{}_\mu$ for each $\mu\in\Sigma$} & \textbf{$E^{(k)4}{}_4$} \\
\midrule
\endhead
0 & 1 & 1 & $-\tfrac12$ & $-\tfrac12$ \\
1 & 30 & 42 & $15H^2$ & $21H^2$ \\
2 & 360 & 840 & $-180H^4$ & $-420H^4$ \\
3 & 720 & 5040 & $360H^6$ & $2520H^6$ \\
\bottomrule
\end{longtable}
\endgroup

The row $k=1$ is the Einstein tensor $G^\mu{}_\nu=\mathrm{diag}(15,15,15,15,21,15,15,15)H^2$ of Section 9.15, verified there by \texttt{KS\_geome\allowbreak{}try\_eins\allowbreak{}teinMixe\allowbreak{}dDiag}: the method passes its first test.

\emph{Step 4: the field equations.} Write the constant coefficients of the Einstein--Lovelock equations as $c_k$, as Section 9.9 does (Chapter 4 calls them $\alpha_k$; we do not, because the letters $\alpha_1,\beta_1,\dots$ are summed indices in the formula of Step 2). The equations $\sum_kc_kE^{(k)\mu}{}_\nu=\kappa T^\mu{}_\nu$ (Sections 4.9 and 9.9) with $\rho=-T^4{}_4$ and the pressure $p=T^\mu{}_\mu$ (no sum) in each of the seven directions become two equations:

\[
\begin{aligned}
\kappa\,\rho_{\mathrm{req}}&=\tfrac12c_0-21c_1H^2+420c_2H^4-2520c_3H^6,\\
\kappa\,p_{\mathrm{req}}&=-\tfrac12c_0+15c_1H^2-180c_2H^4+360c_3H^6 .
\end{aligned}
\]

For Einstein gravity ($c_1=1$, all others 0) they give back $\rho_{\mathrm{req}}=-21H^2/\kappa$ and $p_{\mathrm{req}}=15H^2/\kappa$ (Section 9.15). Two consequences follow.

\begin{itemize}
\item \emph{The sign of the required energy can change.} With $c_1=1$, $c_0=c_3=0$ and a Gauss--Bonnet coefficient $c_2$: $\kappa\rho_{\mathrm{req}}=H^2(420c_2H^2-21)$, which is positive when $c_2>1/(20H^2)$. For example $c_2=1/(10H^2)$ gives $\rho_{\mathrm{req}}=21H^2/\kappa>0$ and $p_{\mathrm{req}}=-3H^2/\kappa$ (Exercise 18.7).
\item \emph{The static member can be a vacuum solution.} Both right-hand sides vanish exactly when their sum and the second one vanish. The sum is $\kappa(\rho_{\mathrm{req}}+p_{\mathrm{req}})=-6c_1H^2+240c_2H^4-2160c_3H^6=-6H^2\bigl(c_1-40c_2H^2+360c_3H^4\bigr)$, so the conditions are
\end{itemize}

\[
c_1-40c_2H^2+360c_3H^4=0,\qquad c_0=2\bigl(15c_1H^2-180c_2H^4+360c_3H^6\bigr).
\]

For $c_1=1$ and $c_3=0$: $c_2=1/(40H^2)$ and $c_0=21H^2$, that is a negative cosmological constant $\Lambda=-\tfrac12c_0=-\tfrac{21}2H^2$ in the convention $c_0=-2\Lambda$ of Section 9.9 ($\alpha_0=-2\Lambda$ in Section 4.9). Read the other way, given the couplings $c_k$ the vacuum condition fixes $H$.

\textbf{What the worked example does not settle.} It is derived here and has not been checked by any program. It covers only the static member, in the bulk $y<0$; the notebook's inflating $a_4$ needs $E^{(2)}$ and $E^{(3)}$ for a time-dependent $a_4$, which are not computed. The brane at $y=0$ needs the Gauss--Bonnet junction conditions, which are not computed. Whether couplings of this size give a physically acceptable theory (for example, whether small perturbations of this background grow) is not examined. And the notebook's normalisation of its Lovelock terms is unknown, so nothing here says which values of its numbers $w_1,w_2,w_3$ and $\Lambda$ (cell 14) this corresponds to.

\textbf{A first step.} (a) Verify the table exactly by machine, twice. The Wolfram package \texttt{wolfram/\allowbreak{}Dirac16C\allowbreak{}omplexKo\allowbreak{}hnSham.w\allowbreak{}l} already computes the Riemann tensor of the static member (it is used by \texttt{KS\_geome\allowbreak{}try\_cons\allowbreak{}tantCurv\allowbreak{}atureSev\allowbreak{}enSpace}); add the generalized-delta contractions for $k=2,3$. Do the same independently in sympy, starting from the geometry engine \texttt{scripts/\allowbreak{}d16c\_geo\allowbreak{}metry\_sy\allowbreak{}mpy.py}. Include a negative control, for example a metric that is not of constant curvature. (b) Compute $E^{(2)}$ and $E^{(3)}$ for the notebook's metric with a general $a_4(t)$; the Stage-2 package \texttt{wolfram/\allowbreak{}Dirac16C\allowbreak{}omplexPr\allowbreak{}imordial\allowbreak{}.wl} computes its Riemann tensor exactly. Write the Einstein--Lovelock equations as ordinary differential equations for $a_4$, in vacuum and with the homogeneous source of Proposition 9.4, and ask whether the slopes of cell 150 appear. (c) Derive the junction conditions at the brane in the Gauss--Bonnet case.

\textbf{Where to start.} Chapter 9; the Stage-2 document (§15 for the source, §5 to §7 for the curvature); the Stage-4 specification, §1; D. Lovelock, J. Math. Phys. 12, 498 (1971), cited in Section 4.9; and these files:

\begin{Verbatim}[fontsize=\small]
wolfram/Dirac16ComplexPrimordial.wl       the primordial field, any a4 (Stage 2)
scripts/check_dirac16complex_primordial.py
artifacts/dirac16complex/primordial-field/wolfram-primordial-report.json
wolfram/Dirac16ComplexKohnSham.wl         the static member and its brane (Stage 4)
scripts/check_dirac16complex_kohn_sham_theory.py
artifacts/dirac16complex/kohn-sham/wolfram-kohn-sham-report.json
scripts/d16c_geometry_sympy.py            the independent sympy geometry engine
handoff/specs/STAGE4_SPEC.md
\end{Verbatim}

\subsection{Problem 4: can a Kohn--Sham state be the source of the primordial field?}

\textbf{The question.} In 8-dimensional Einstein gravity the static member of the primordial field needs the source $\rho_{\mathrm{req}}=-21H^2/\kappa$ and $p_{\mathrm{req}}=+15H^2/\kappa$ in all seven directions other than $x_4$ (Section 9.15). A homogeneous state $\Psi=u(x_4)$ is an exact source exactly when its three-gamma bilinears vanish, $mS=-36H^2/\kappa$ and $\lambda S^2=30H^2/\kappa$ (Proposition 9.4 with $c=0$; erratum E4.1 of \texttt{handoff/\allowbreak{}specs/ST\allowbreak{}AGE4\_SPE\allowbreak{}C.md}). Chapter 13 asked whether a Kohn--Sham state of finitely many quanta can meet these conditions and found that no computed state does (Section 13.14). The open question is whether \textbf{any} many-particle state of dirac16complex, or of dirac16complex00, can be the source of the static field: in the fixed field, or self-consistently, with the metric computed from the state (back-reaction).

\textbf{Why it is hard.}

\begin{itemize}
\item \emph{Sign.} The field needs a negative energy density; every computed state with band or bulk levels has a positive proper-volume average $\langle\rho\rangle$, so the match would need a negative gravitational coupling (Section 13.14). In the good sector the free, normal-ordered energy is not negative (Section 8.10), so a negative energy must come from the interaction or from a negative-energy sector.
\item \emph{Strength.} The two conditions together need $\lambda\langle S_p\rangle/m=-\tfrac56$; at first order this means $\hat\lambda\approx105.7$ for $N=112$ and $9.55$ for $N=1016$, $10^3$ to $10^4$ times the largest coupling used, far outside the range for which the local exchange-only scheme was set up (Section 13.14; column \texttt{lambda\_h\allowbreak{}at\_neede\allowbreak{}d\_first\_\allowbreak{}order} of \texttt{artifact\allowbreak{}s/dirac1\allowbreak{}6complex\allowbreak{}/kohn-sh\allowbreak{}am/rust/\allowbreak{}emt/emt-\allowbreak{}summary.\allowbreak{}csv}).
\item \emph{Shape.} The required source is the same at every $y$, while the computed energy density varies by orders of magnitude: between 0.00865 and 3.495 for $N=112$ and between 0.181 and 4768 for $N=1016$ (both at $\lambda=0$, $m=H=1$; columns \texttt{\detokenize{rho_min}} and \texttt{\detokenize{rho_max}} of the same file).
\item \emph{Back-reaction.} If the metric is computed from the state, the orbitals and the metric must be found together: a coupled nonlinear boundary-value problem in $y$ that nobody has set up.
\end{itemize}

\textbf{What is known.}

\begin{itemize}
\item PROVED (Proposition 9.4): exact homogeneous sources exist in the c-number (mean-field) reading of the bilinears, always with negative energy density. The Stage-2 report checks them in \texttt{P\_source\allowbreak{}\_x0Indep\allowbreak{}endentSo\allowbreak{}urceCond\allowbreak{}itions} and \texttt{P\_source\allowbreak{}\_x0Indep\allowbreak{}endentEx\allowbreak{}actExamp\allowbreak{}les}.
\item PROVED for the commuting field: an exact classical solution of dirac16complex00 that is the source of the static field, with $m=-30$, $\lambda=\tfrac{125}6$, $S=\tfrac65$, $M_{\mathrm{eff}}=m+\lambda S=-5$, $\rho=-21$, $p=15$, $w=-\tfrac57$ in units $H=\kappa=1$ (\texttt{artifact\allowbreak{}s/dirac1\allowbreak{}6complex\allowbreak{}/pair-cr\allowbreak{}eation/d\allowbreak{}irac16co\allowbreak{}mplex00-\allowbreak{}theory.j\allowbreak{}son}, key \texttt{primordi\allowbreak{}al.stati\allowbreak{}cSourceE\allowbreak{}xample}; checks \texttt{C00\_prim\allowbreak{}ordial\_s\allowbreak{}taticFie\allowbreak{}ldSource\allowbreak{}dExactly} and \texttt{S5\_primo\allowbreak{}rdial\_st\allowbreak{}aticFiel\allowbreak{}dSourced\allowbreak{}Exactly}). Indeed $mS=-36$, $\lambda S^2=\tfrac{125}6\cdot\tfrac{36}{25}=30$, $\rho=mS+\tfrac\lambda2S^2=-36+15=-21$ and $p=\tfrac\lambda2S^2=15$. Its negative energy density is possible because the classical energy of dirac16complex00 is not bounded below (Section 18.3).
\item COMPUTED: no Kohn--Sham state meets the conditions (the column \texttt{sourcing\allowbreak{}\_conditi\allowbreak{}ons\_met} of \texttt{rust/emt\allowbreak{}/emt-sum\allowbreak{}mary.csv} is 0 in every row; Section 13.14).
\item PROVED (Stage 5, Theorem T3; Section 13.14): the $-M$ universe with transformed boundary conditions has the same energies and pressures as the $+M$ universe and the opposite scalar density, so $mS$ and $\lambda S^2$ are unchanged: the mirror universe does not help.
\item PROVED (Stage 5; \texttt{artifact\allowbreak{}s/dirac1\allowbreak{}6complex\allowbreak{}/pair-cr\allowbreak{}eation/p\allowbreak{}airing-t\allowbreak{}heory.js\allowbreak{}on}, key \texttt{T3.theor\allowbreak{}emImageR\allowbreak{}ule}; check \texttt{PAIR\_tot\allowbreak{}als\_ksKr\allowbreak{}einImage\allowbreak{}Pair}): the Krein image of a Kohn--Sham state (the $\gamma^8$ image with the Krein metric $-B$ and the coupling $-\lambda$) has the same orbitals and eigenvalues, the number density reversed ($n\to-n$), the scalar density unchanged ($S\to S$), and every energy and every component of the energy--momentum tensor reversed.
\end{itemize}

\textbf{Worked example 1: the fixed static field needs a uniform and isotropic source} (derived here). By the table of Section 18.5 every Lovelock tensor $E^{(k)}$ of the static member is diagonal and constant, with the same value on all seven directions $y,x_1,x_2,x_3,x_5,x_6,x_7$. Whatever the coefficients $c_k$, the Einstein--Lovelock equations for the fixed static member therefore demand a source whose energy density is the same at every point and whose pressure is one and the same constant in all seven directions. The Kohn--Sham states are neither uniform (see "Shape" above) nor isotropic: for $N=112$ at $\lambda=0$ the proper-volume averages are $\langle p_y\rangle=0.0103747$, $\langle p_3\rangle=0.0068653$ and $\langle p_t\rangle=0$ (\texttt{rust/emt\allowbreak{}/emt-sum\allowbreak{}mary.csv}); without interaction the extra times carry no pressure at all (Section 13.14). So in the \textbf{fixed} static field these states are excluded in every Einstein--Lovelock theory, not only in Einstein gravity. The exact homogeneous sources above are isotropic by construction: their pressure is $SU'(S)-U(S)$ in all seven directions (Proposition 9.4). A Kohn--Sham state held near a brane is a different kind of object, and the remaining possibility is to let the metric respond.

\textbf{Worked example 2: what back-reaction demands} (derived here). Call a metric \textbf{ultrastatic} if it is diagonal, $g_{44}=-1$, and no entry depends on $x_4$. The static member is ultrastatic, and so is every static warp $W(y)$ of Section 9.16. For such a metric every Christoffel symbol with an index 4 vanishes, by the rules of Section 9.6: $\Gamma^p{}_{4p}=\tfrac12\partial_4\ln\lvert g_{pp}\rvert=0$ and $\Gamma^4{}_{\mu4}=\tfrac12\partial_\mu\ln\lvert g_{44}\rvert=0$ by (C1), $\Gamma^4{}_{pp}=-\partial_4g_{pp}/(2g_{44})=0$ and $\Gamma^\rho{}_{44}=-\partial_\rho g_{44}/(2g_{\rho\rho})=0$ by (C2), and those with three different indices by (C3). Every term of $R_{44}=\partial_\rho\Gamma^\rho{}_{44}-\partial_4\Gamma^\rho{}_{\rho4}+\Gamma^\rho{}_{\rho\lambda}\Gamma^\lambda{}_{44}-\Gamma^\rho{}_{4\lambda}\Gamma^\lambda{}_{\rho4}$ contains such a symbol, so $R_{44}=0$. Section 9.11 derived from the trace-reversed Einstein equations that $5\rho+\sum_{\mu\ne4}p_{(\mu)}=6R_{44}/\kappa$. Hence, in 8-dimensional Einstein gravity, every source of an ultrastatic metric must satisfy at every point

\[
5\rho+p_y+3p_3+3p_t=0,
\]

with $p_3$ the average 3-space pressure and $p_t$ the pressure along each extra time (or, in general, the sum of the seven pressures in place of $p_y+3p_3+3p_t$). Matter with $\rho>0$ and non-negative pressures can never do this: an ultrastatic geometry needs matter that exactly saturates the strong energy condition along $x_4$. (The same computation in four dimensions gives $\rho+3p=0$, the condition that made Einstein add a cosmological constant to obtain his static universe of 1917; Exercise 18.9.) Both exact sources above pass the test: for Proposition 9.4 with $c=0$, $5(mS+U)+7(SU'-U)=5mS+6\lambda S^2=5(-36)+6\cdot30=0$ in units $H=\kappa=1$, and for the dirac16complex00 example $5(-21)+7\cdot15=0$ (Exercise 18.8).

For a Kohn--Sham state in the static member insert the energy--momentum tensor of Section 13.14, $\rho=\sum_nw_n\varepsilon_nn_n-L_s$, $p_y=\sum_nw_n(\varepsilon_nn_n-ms_n-\xi k_nt_n)-L_s$, $p_3=\tfrac13\sum_nw_n\xi k_nt_n+L_s$ and $p_t=L_s$, with the momentum weight $\xi(y)=e^{-Hy-a_{4,0}}$ of Section 13.3 (the Greek letter xi; it is not the gravitational coupling $\kappa$ of the equation above). The interaction terms cancel ($-5-1+3+3=0$) and so do the momentum terms:

\[
5\rho+p_y+3p_3+3p_t=6\sum_nw_n\varepsilon_nn_n(y)-m\,S_p(y).
\]

A Kohn--Sham state of the static member could therefore be the source of an ultrastatic metric only if $6\sum_nw_n\varepsilon_nn_n=mS_p$ at every $y$. (For a general warp the Kohn--Sham tensor has to be derived again, and whether the identity keeps this form is part of the first step below.) The committed states are far from it. Evaluated on the committed profiles (\texttt{rust/emt\allowbreak{}/<run>/p\allowbreak{}rofiles.\allowbreak{}csv}, columns \texttt{\detokenize{rho}}, \texttt{\detokenize{p_y}}, \texttt{\detokenize{p_3}}, \texttt{\detokenize{p_t}}), the left-hand side lies between 0.0519 and 27.07 for $N=112$ and between 1.085 and 28609 for $N=1016$ (both at $\lambda=0$), and in those files the identity with $6(\rho+L_s)-mS_p$ on the right holds to rounding (Exercise 18.10). The free state with $N=8$ satisfies the condition trivially, because its whole energy--momentum tensor vanishes (the eight zero modes have $\varepsilon=0$ and $S_p=0$, Section 13.11); but a vanishing tensor is the source only of a vacuum metric, and the static member is not a vacuum solution of Einstein gravity.

\textbf{Worked example 3: the Krein image of a computed state.} Take $N=112$ at $\lambda=0$ (\texttt{rust/emt\allowbreak{}/emt-sum\allowbreak{}mary.csv}, $m=H=1$): $\langle\rho\rangle=0.0230891$ and $\langle S_p\rangle=-0.00788147$. Its Krein image has $\langle\rho\rangle=-0.0230891$, negative as the field needs, and matching $\rho_{\mathrm{req}}=-21/\kappa$ on average gives $\kappa=21/0.0230891=909.5>0$. But the image carries the mass $-m=-1$ and the unchanged $\langle S_p\rangle$, so the mass condition, used on averages as a diagnostic exactly as in Section 13.14, reads $(-1)(-0.00788147)=-36/\kappa$ and gives $\kappa=-4568<0$. The two conditions contradict each other, the coupling condition $\lambda S^2=30/\kappa$ cannot hold at $\lambda=0$ at all, and the image inherits the anisotropy of worked example 1. The image is not a source either.

\textbf{A first step.} (a) Use the necessary condition $6\sum_nw_n\varepsilon_nn_n=mS_p$ as a cheap test for every candidate state; it needs only the committed profiles. (b) Generalise worked example 2 to static metrics in which $g_{44}$ is warped as well, $g_{44}=-e^{2f_4(y)}$: then $R_{44}\ne0$, and Lemma 9.1, with $x_4$ treated as one more fibre direction over the one-dimensional base $y$ (the proof of the lemma works unchanged for a base with one coordinate), gives the new condition. (c) Allow different warps for 3-space and the extra times, $g_{pp}\propto e^{2f_3(y)}$ for $p=1,2,3$ and $e^{2f_t(y)}$ for $p=5,6,7$. Lemma 9.1 with the base $(y,x_4)$, on which nothing depends on $x_4$, gives $R^p{}_p=-(f_p''+\Theta f_p')$ with $\Theta=3f_3'+3f_t'$, hence (derived here)

\[
\kappa\,(p_3-p_t)=G^1{}_1-G^5{}_5=-(f_3''-f_t'')-3(f_3'+f_t')(f_3'-f_t'),
\]

so the pressure difference of the Kohn--Sham states needs $f_3'\ne f_t'$ somewhere, that is a difference $f_3-f_t=2a_4(y)$ that depends on $y$. A constant difference is not enough: the static member has $f_3-f_t=2a_{4,0}$, which is nonzero whenever $a_{4,0}\ne0$, but $f_3'=f_t'=H$ (Section 9.15), and the right-hand side vanishes. Then set up the self-consistent problem: the unknown warps, the Kohn--Sham orbitals in that metric, and Einstein's equations as ordinary differential equations in $y$. (d) Combine with Section 18.5: with a Gauss--Bonnet term the required energy density of the static member can be positive, but by worked example 1 the fixed static metric still needs a uniform isotropic source, so back-reaction is needed there as well.

\textbf{Where to start.} Section 13.14; the Stage-2 document, §15.5; erratum E4.1 of the Stage-4 specification; and these files:

\begin{Verbatim}[fontsize=\small]
artifacts/dirac16complex/kohn-sham/rust/emt/emt-summary.csv    averages of every state
artifacts/dirac16complex/kohn-sham/rust/emt/<run>/profiles.csv the profiles in y
studies/dirac16complex_kohn_sham/src/emt.rs                    the energy-momentum tensor
studies/dirac16complex_kohn_sham/src/shooting.rs               the shooting method
artifacts/dirac16complex/pair-creation/dirac16complex00-theory.json
artifacts/dirac16complex/pair-creation/pairing-theory.json     key T3: the pairing maps
handoff/specs/STAGE4_SPEC.md
\end{Verbatim}

\subsection{Problem 5: the last Stage-4 cross-check and the Stage-4 gate}

\textbf{The question.} Every number of Chapter 13 was computed by the Rust solver and is to be confirmed by the independent Python reference solver, which uses a matrix method on staggered grids and extrapolates in the grid spacing (Section 13.15). The report of the final cross-check has 63 checks, of which 62 are true:

\begin{Verbatim}[fontsize=\small]
artifacts/dirac16complex/kohn-sham/python-check-report.json   the report
scripts/check_dirac16complex_kohn_sham.py                      its producer
\end{Verbatim}

The problem is to resolve the one failed check and then to run the Stage-4 gate to the end.

\textbf{What is known.} The failed check is \texttt{canonica\allowbreak{}l\_eigenv\allowbreak{}alues}. It compares 58412 eigenvalues of the two solvers with the tolerance $10^{-6}\max(1,\lvert\varepsilon\rvert/m)\,m$ plus a small grid term, and its worst deviation is $2.195\times10^{-6}$ against the tolerance $1.054\times10^{-6}$ (ratio 2.08). All its failures belong to the run \texttt{m1\_L3\_N1\allowbreak{}016\_lamm\allowbreak{}2\_T0}, the smeared ensemble at the attractive coupling $-\hat\lambda_2$ of Sections 13.11 and 13.12: two deep levels of the Dirac sea at $k=0$, one in each parity sector, each in both block types, and each appearing in the ground-state and in the excited-state record of that run. Everything else about the run agrees: the ground-state energy (Rust 1126.8580930, reference 1126.8580892, deviation $3.8\times10^{-6}$ against a tolerance of $1.1\times10^{-3}$) and the Kohn--Sham gap (deviation $1.8\times10^{-8}$), both from the \texttt{comparis\allowbreak{}ons} block of the report. The Delta-SCF of this run, the one open number of Section 13.12, now agrees: the reference on the finer grids 120/240/480 gives 0.0436199531, against the Rust values 0.0436193591 (601 points) and 0.0436196717 (extrapolated from 301 and 601 points) (erratum E4.13 of \texttt{handoff/\allowbreak{}specs/ST\allowbreak{}AGE4\_SPE\allowbreak{}C.md}).

The worse of the two failing levels is the $k=0$ sea level of parity $+$ whose free value is $-1.447972$, the negative-energy partner of the $k=0$ bulk level that the attraction pulls down to the Fermi level (Section 13.11); the other is the $k=0$ sea level of parity $-$ whose free value is $-1.292293$ (Section 13.6). The interaction moves both up, and three committed files give their eigenvalues:

\begingroup
\small
\begin{longtable}{@{}>{\raggedright\arraybackslash}p{0.200\linewidth}>{\raggedright\arraybackslash}p{0.200\linewidth}>{\raggedright\arraybackslash}p{0.200\linewidth}>{\raggedright\arraybackslash}p{0.200\linewidth}@{}}
\toprule
\textbf{Source} & \textbf{File (under \texttt{artifact\allowbreak{}s/dirac1\allowbreak{}6complex\allowbreak{}/kohn-sh\allowbreak{}am/})} & \textbf{$\varepsilon$, parity $+$} & \textbf{$\varepsilon$, parity $-$} \\
\midrule
\endfirsthead
\toprule
\textbf{Source} & \textbf{File (under \texttt{artifact\allowbreak{}s/dirac1\allowbreak{}6complex\allowbreak{}/kohn-sh\allowbreak{}am/})} & \textbf{$\varepsilon$, parity $+$} & \textbf{$\varepsilon$, parity $-$} \\
\midrule
\endhead
Rust, 301 points (canonical) & \texttt{rust/scf\allowbreak{}/m1\_L3\_N\allowbreak{}1016\_lam\allowbreak{}m2\_T0/le\allowbreak{}vels.csv} & $-1.05374908333$ & $-0.98844655130$ \\
Rust, 601 points & \texttt{rust/exc\allowbreak{}ited/m1\_\allowbreak{}L3\_N1016\allowbreak{}\_lamm2\_T\allowbreak{}0\_g601/l\allowbreak{}evels.cs\allowbreak{}v} & $-1.05374707094$ & $-0.98844473928$ \\
reference, grids 120/240/480, extrapolated & \texttt{referenc\allowbreak{}e/m1\_L3\_\allowbreak{}N1016\_la\allowbreak{}mm2\_T0/s\allowbreak{}pectrum.\allowbreak{}csv} & $-1.05374688831$ & $-0.98844457574$ \\
\bottomrule
\end{longtable}
\endgroup

For the level of parity $+$ the differences are: Rust (301) minus reference $=-2.195\times10^{-6}$, the failing deviation; Rust (601) minus reference $=-1.83\times10^{-7}$, well inside the tolerance; Rust (301) minus Rust (601) $=-2.012\times10^{-6}$. The level of parity $-$ shows the same pattern ($-1.976\times10^{-6}$, $-1.64\times10^{-7}$ and $-1.812\times10^{-6}$). The committed numbers therefore suggest that both deviations are the grid error of the 301-point Rust run for these two deep sea levels. That is a hypothesis to test, not a result. The checker already adds the measured Rust grid uncertainty $\lvert X(601)-X(301)\rvert$ to the tolerances of the ground-state energy, the chemical potential, the Kohn--Sham gap, the Delta-SCF energy and the lowest particle--hole energy of a run with a 601-point partner (erratum E4.12; function \texttt{rust\_gri\allowbreak{}d\_uncert\allowbreak{}ainties} of the checker), but not to single eigenvalues.

\textbf{Why it is hard.} Not conceptually; the difficulty is discipline. This run is the hardest of Stage 4: two levels cross at the Fermi level, the ground state is an ensemble with smeared occupations, and the nonlinear problem has more than one self-consistent solution (an exploratory reference run on finer grids converged to a state about $1.01\,m$ higher; Section 13.11). A tolerance may be changed only on the basis of a measurement, and the change must be recorded.

\textbf{A first step.} (1) Test the hypothesis: solve the run on a finer $y$ grid (1201 points, say) with the Rust solver, and check that its values of the two levels move toward the reference values in the way Richardson's rule of Section 10.11 predicts for a grid error. (2) If it does, decide, with that measurement, whether the rule of erratum E4.12 should be extended to single eigenvalues of runs with a finer-grid partner (Exercise 18.11 shows that the check would then pass), and record the decision as a new erratum of the Stage-4 specification. Enlarging the tolerance without such a measurement would be a weakened check, which the project does not allow. (3) Rerun the checker. (4) Run the gate: \texttt{scripts/\allowbreak{}verify\_s\allowbreak{}tage4\_ko\allowbreak{}hn\_sham.\allowbreak{}ps1} in PowerShell or \texttt{bash scr\allowbreak{}ipts/ver\allowbreak{}ify\_stag\allowbreak{}e4\_kohn\_\allowbreak{}sham.sh}, first with \texttt{\detokenize{-DryRun}} (PowerShell) or \texttt{--dry-ru\allowbreak{}n} (Bash), which prints every step and its expected wall time without running anything. The gate requires a fresh checker report without a failed check, and it rebuilds and verifies the two Stage-4 PDFs, which do not exist yet: the Stage-4 document exists only as the Markdown draft \texttt{provenan\allowbreak{}ce/DIRAC\allowbreak{}16COMPLE\allowbreak{}X\_KOHN\_S\allowbreak{}HAM\_PRIM\allowbreak{}ORDIAL.m\allowbreak{}d}, and the Stage-4 student guide is not written. A full run takes hours; the header of the PowerShell gate estimates about three hours for the Rust step alone. Chapter 19 describes the gates in general.

\textbf{Where to start.} The entries \texttt{\detokenize{checks}}, \texttt{measurem\allowbreak{}ents} and \texttt{comparis\allowbreak{}ons} of the cross-check report, the Stage-4 row of \texttt{HANDOFF.\allowbreak{}md}, and these files:

\begin{Verbatim}[fontsize=\small]
artifacts/dirac16complex/kohn-sham/python-check-report.json
scripts/check_dirac16complex_kohn_sham.py        the checker (rust_grid_uncertainties)
scripts/ks_reference_solver.py                   the reference (its list canonical_runs)
studies/dirac16complex_kohn_sham/README.md       the Rust solver
scripts/verify_stage4_kohn_sham.ps1              the gate (PowerShell)
scripts/verify_stage4_kohn_sham.sh               the gate (Bash)
\end{Verbatim}

\subsection{Problem 6: the Stage-5 numerics}

\textbf{The question.} Item T4 of \texttt{handoff/\allowbreak{}specs/ST\allowbreak{}AGE5\_SPE\allowbreak{}C.md} (§5) is a numerical demonstration, not a theorem: the Kohn--Sham pairing theorem T3 (Chapter 15) is proved exactly, and T4 asks to show it at work in computed states, in the sense of the word COMPUTED of Section 0.3. The specification asks for the ground and first excited states of the $+M$ and $-M$ universes and for the untransformed-boundary-condition control, for both fields, with $m=1$ and $3$, $L=3$, $N=8$ and $112$, $\lambda\in\{0,\pm\hat\lambda_1,\pm\hat\lambda_2\}$, and "$T=0$ and one $T>0$", computed by both solvers, with the pairing verified to solver precision and the pair totals reported. So there are three universes per configuration: the $+M$ universe, the $-M$ universe with the transformed boundary conditions, and the $-M$ universe with the untransformed ones (the control, which must not pair). The specification also asks for a checker, the two Stage-5 documents and the Stage-5 gate. The concrete matrix that the reference driver plans is recorded under the key \texttt{\detokenize{matrix}} of \texttt{artifact\allowbreak{}s/dirac1\allowbreak{}6complex\allowbreak{}/pair-cr\allowbreak{}eation/r\allowbreak{}eference\allowbreak{}/referen\allowbreak{}ce-pairs\allowbreak{}-summary\allowbreak{}.json}: the temperature $T=0.1\lvert m\rvert$ (entry \texttt{T\_over\_a\allowbreak{}bsM}), used only for $N=112$ (entry \texttt{\detokenize{thermoN}}). Counting from that entry gives 60 configurations (two fields; four pairs $(\lvert m\rvert,N)$ with five couplings at $T=0$; the two pairs with $N=112$ with five couplings at $T=0.1\lvert m\rvert$), that is 180 runs of single universes.

\textbf{What is committed} (2026-09-30).

\begin{itemize}
\item The exact theory: \texttt{wolfram-\allowbreak{}dirac16c\allowbreak{}omplex00\allowbreak{}-report.\allowbreak{}json} (46 of 46 checks), \texttt{python-d\allowbreak{}irac16co\allowbreak{}mplex00-\allowbreak{}report.j\allowbreak{}son} (49 of 49), \texttt{wolfram-\allowbreak{}pairing-\allowbreak{}report.j\allowbreak{}son} (141 of 141) and \texttt{python-p\allowbreak{}airing-r\allowbreak{}eport.js\allowbreak{}on} (172 of 172), all in \texttt{artifact\allowbreak{}s/dirac1\allowbreak{}6complex\allowbreak{}/pair-cr\allowbreak{}eation/}.
\item The Rust subcommand \texttt{\detokenize{pairs}} (\texttt{studies/\allowbreak{}dirac16c\allowbreak{}omplex\_k\allowbreak{}ohn\_sham\allowbreak{}/src/pai\allowbreak{}rs.rs}) with outputs under \texttt{artifact\allowbreak{}s/dirac1\allowbreak{}6complex\allowbreak{}/pair-cr\allowbreak{}eation/r\allowbreak{}ust/pair\allowbreak{}s/}: 25 of the 60 configurations, all for dirac16complex, each with a \texttt{pairing.\allowbreak{}json}: $m=1$, $N=8$ (five couplings, $T=0$); $m=1$, $N=112$ (five couplings at $T=0$ and at $T=0.1$); $m=3$, $N=8$ (five couplings, $T=0$); $m=3$, $N=112$ ($\lambda=0$ and $+\hat\lambda_1$ at $T=0$ and at $T=0.1$, $-\hat\lambda_1$ at $T=0$). The list is that of commit \texttt{\detokenize{1f2dd69}} (Section 19.10). There is no run for dirac16complex00.
\item In every committed run the $-M$ universe with the transformed boundary conditions reproduces the $+M$ universe: every level is matched with no mismatch of multiplicities, the largest difference of an eigenvalue is $1.8\times10^{-5}$ (run \texttt{d16c\_m3\_\allowbreak{}L3\_N112\_\allowbreak{}lamp1\_T0\allowbreak{}p1}, at $T=0.1\,\lvert m\rvert$) and at most $6.6\times10^{-9}$ in every other run (the largest in \texttt{d16c\_m1\_\allowbreak{}L3\_N8\_la\allowbreak{}mm2\_T0}), the largest difference of an energy is $1.6\times10^{-8}$, and the scalar totals are opposite (key \texttt{pairingD\allowbreak{}eviation\allowbreak{}\_plusM\_v\allowbreak{}s\_minusM} of each \texttt{pairing.\allowbreak{}json}). In the two runs with $m=3$, $N=8$ at $\pm\hat\lambda_2$ the iteration for the excited state did not converge in either universe (field \texttt{excitedC\allowbreak{}onverged}). Every $+M$ run at $T=0$ for which a Stage-4 run with the same parameters is committed (all except those two) is bit for bit that Stage-4 run (key \texttt{stage4Re\allowbreak{}producti\allowbreak{}on}).
\item The controls ran only at $\lambda=0$ (six runs), and they differ as the exact theory predicts. For $m=1$, $N=8$ the control's Kohn--Sham gap is 0.100851 against 0.430734 for the pair. This value is the sub-gap bound state predicted by the exact theory (\texttt{pairing-\allowbreak{}theory.j\allowbreak{}son}, key \texttt{T3.untra\allowbreak{}nsformed\allowbreak{}BCContro\allowbreak{}l}; check \texttt{PAIR\_T3k\allowbreak{}s\_contro\allowbreak{}lSubGapB\allowbreak{}oundStat\allowbreak{}e}): $\varepsilon=\sqrt{M^2-q^2}$ with $\tanh(qL)=q/M$ (Section 15.9), which for $M=1$, $L=3$ gives $q=0.994902$ and $\varepsilon=0.100851$ (the analytic value listed in \texttt{rust/pai\allowbreak{}rs/free-\allowbreak{}section.\allowbreak{}json}). For $m=3$, $N=8$ the control gap is 0.000740 against 0.895833, for $m=1$, $N=112$ the control's ground energy is 56.26550 against 61.09072, and for $m=3$, $N=112$ it is 119.11813 against 131.44830. The interacting controls were not run: their first self-consistent update would violate the energy-window premise of the solver, because the untransformed tip bag puts the zero modes at the tip, where the proper-density factor is $e^{6HL}$ (the entry \texttt{minusM\_c\allowbreak{}ontrol} under the key \texttt{universe\allowbreak{}s} of each \texttt{pairing.\allowbreak{}json} records this).
\item The pair totals (key \texttt{pairTota\allowbreak{}ls}): for \texttt{d16c\_m1\_\allowbreak{}L3\_N112\_\allowbreak{}lam0\_T0} the mirror pair (the $+M$ universe and the ordinary $-M$ universe) has $E=122.181447=2\times61.090723$ and charge 224, and the Krein-image pair has $E=-7.3\times10^{-11}$, zero to solver precision, and charge 0, as Theorem T3 predicts.
\item The reference (\texttt{scripts/\allowbreak{}ks\_refer\allowbreak{}ence\_pai\allowbreak{}rs.py}; \texttt{artifact\allowbreak{}s/dirac1\allowbreak{}6complex\allowbreak{}/pair-cr\allowbreak{}eation/r\allowbreak{}eference\allowbreak{}/referen\allowbreak{}ce-pairs\allowbreak{}-summary\allowbreak{}.json}, whose flag \texttt{\detokenize{complete}} is false): of the 180 runs, 49 are recorded, all for $\lvert m\rvert=3$, $N=112$, and 17 of them converged; 131 are pending, among them every run with $\lvert m\rvert=1$ and every run with $N=8$, for both fields; 6 failed (five controls and one $-M$ run of dirac16complex00; five of them because their energy window exceeded the solver's cap of $20\lvert m\rvert$, one with the recorded error "tail: window never exhausted"); 26 did not converge, 24 of them interacting runs at $T=0.1\lvert m\rvert$ that were not attempted, and two runs of dirac16complex at $+\hat\lambda_2$.
\item Not committed: a report of the pairs checker, the two Stage-5 documents and the Stage-5 gate. Their names are listed below.
\end{itemize}

\begin{Verbatim}[fontsize=\small]
scripts/check_dirac16complex_pairs.py          the pairs checker (no committed report)
provenance/DIRAC16COMPLEX00_FIELD_THEORY       planned document (not written)
provenance/DIRAC16COMPLEX_PAIR_CREATION        planned document (not written)
scripts/verify_stage5_pair_creation            planned gate (not written)
\end{Verbatim}

\textbf{Why it is hard.} The window premise stops the interacting controls. The couplings of the $m=3$ configurations are large ($\hat\lambda_1=42863$ for $N=8$ and $852.6$ for $N=112$, from the Stage-4 reference couplings listed in \texttt{referenc\allowbreak{}e-pairs-\allowbreak{}summary.\allowbreak{}json}), because the rule of Section 13.10 scales them to a fixed pseudo-potential. The commuting field has the opposite exchange sign, $M_{\mathrm{eff}}=m+\tfrac{17}{16}\lambda S_p$ and $v_x=+\tfrac1{16}\lambda n_p$ (Chapter 14), so attraction and repulsion exchange roles, and its Kohn--Sham model is a formal functional with the Pauli filling imposed by prescription (\texttt{pairing-\allowbreak{}theory.j\allowbreak{}son}, key \texttt{notProve\allowbreak{}d}). The staggered grid of the reference solver treats the two components of a block differently, so there the pairing holds only up to the discretisation error and must be judged after extrapolation (\texttt{pairing-\allowbreak{}theory.j\allowbreak{}son}, key \texttt{T3.numer\allowbreak{}icsPresc\allowbreak{}ription}). And the reference runs take hours.

\textbf{A first step.} (1) Check one pair by hand: in \texttt{rust/pai\allowbreak{}rs/free-\allowbreak{}section.\allowbreak{}json} compare the $k=0$ levels of the $+M$ and $-M$ universes with each other and with the exact values of Section 13.6 ($0$ and $\pm\sqrt{1+(n\pi/3)^2}$, for example $\pm1.447972$, in the parity sector with the zero mode; $\pm1.292293$ in the other), and find the control's sub-gap level. (2) Run the pairs checker on what exists, writing its report into the git-ignored folder \texttt{\detokenize{build/}} (first command below), and read which comparisons ran and which were recorded as not run. Expect failed checks as well: on the committed files the checker currently ends with a nonzero exit status and a list of lines beginning with \texttt{\detokenize{FAILED}}. Some of these failures come from the incomplete matrix; for example, one check demands that all 180 reference runs be present. Others come from tolerances applied to quantities that vanish; for example, a relative comparison of $\lambda\langle S_p\rangle/m$ between the $+M$ and $-M$ universes of $m=3$, $N=8$, where the committed \texttt{pairing.\allowbreak{}json} of \texttt{d16c\_m3\_\allowbreak{}L3\_N8\_la\allowbreak{}mm1\_T0} gives $\langle S_p\rangle=8.3\times10^{-15}$ and $-8.4\times10^{-15}$, that is zero to rounding. The other failures are not analysed in any committed file. Sorting every failure into one of these kinds, or finding a real defect, is part of this problem. No committed report of the checker exists, so none of its numbers is a result of this book, and a failed check here does not contradict the pairing deviations quoted above, which are read from the committed \texttt{pairing.\allowbreak{}json} files. (3) Run one pending reference configuration, again into \texttt{\detokenize{build/}} (second command; the option \texttt{\detokenize{--list}} prints all 180 labels). (4) Compute the missing dirac16complex00 configurations with the Rust \texttt{\detokenize{pairs}} subcommand, whose parameters are described in the section "Stage 5" of the crate's \texttt{README.m\allowbreak{}d}.

\begin{Verbatim}[fontsize=\small]
python scripts/check_dirac16complex_pairs.py --report build/pairs-check.json
python scripts/ks_reference_pairs.py --output build/refpairs \
    --runs d16c_m1_L3_N8_lam0_T0/plusM,d16c_m1_L3_N8_lam0_T0/minusM
\end{Verbatim}

(In PowerShell the second command is written on one line, without the backslash.)

\textbf{Where to start.} The Stage-5 specification, its §4 and §5, its document outline, and these files:

\begin{Verbatim}[fontsize=\small]
handoff/specs/STAGE5_SPEC.md
handoff/specs/STAGE5_DOC_OUTLINE.md
artifacts/dirac16complex/pair-creation/            the exact reports and theory files
artifacts/dirac16complex/pair-creation/rust/pairs/ the Rust pair runs
artifacts/dirac16complex/pair-creation/reference/  the reference pair runs
studies/dirac16complex_kohn_sham/src/pairs.rs      the Rust subcommand pairs
scripts/ks_reference_pairs.py                      the reference driver
scripts/check_dirac16complex_pairs.py              the checker
\end{Verbatim}

\subsection{Problem 7: a creation amplitude for pairs of universes}

\textbf{The question.} The notebook states as a hypothesis that at the time $x_4=0$ a pair of universes with the masses $\pm M$ is created (cell 6), asks whether such universes are created in pairs (cell 7), and records the task "TODO: prove Universe(s) of masses ±M are created in pairs!" (cell 17; all three are quoted in Section 0.2). The project proves exact pairing theorems (Chapter 15), and Chapter 16 states exactly what they establish and what they do not; the creation itself remains a HYPOTHESIS. The open problem is to find a dynamical theory in which the creation of such a pair is a process with a computable amplitude, probability or rate, to compute it, and so to decide whether pairs are created.

\textbf{What is known.}

\begin{itemize}
\item PROVED (Theorem T1 of Chapter 15; checks \texttt{ALG\_gamm\allowbreak{}a8Map} of Stage 1 and \texttt{PAIR\_tot\allowbreak{}als\_fiel\allowbreak{}dLevelCh\allowbreak{}iralPair}): a universe $\Psi_+$ with $(m,\lambda)$ and its chirality image $\gamma^8\Psi_+$ with $(-m,-\lambda)$, in the same gravitational field, have zero total energy--momentum tensor, zero total current and even zero total Lagrangian density, at every point and every $x_4$, in every gravitational field (\texttt{pairing-\allowbreak{}theory.j\allowbreak{}son}, key \texttt{pairTota\allowbreak{}ls.field\allowbreak{}Level}). As classical fields, such a pair therefore has the totals of the field-free state: the conservation of the totals and the constraints of the field equations do not forbid its appearance, and at the level of the classical action it costs nothing. In the theory as built, however, the two members do not interact, so each member's charge (and, in a static field, its energy) is conserved on its own, and a pair with $Q_+\ne0$ cannot appear from the field-free state (Proposition 16.5, item 3).
\item PROVED (Theorem T2; \texttt{PAIR\_tot\allowbreak{}als\_fiel\allowbreak{}dLevelMi\allowbreak{}rrorPair}): the mirror pair, related by a Pin(4,4) reflection combined with $\gamma^8$ and with the same $\lambda$, has the same energy--momentum in both members, so its totals double instead of cancelling. In the Z2 construction every Kohn--Sham state of Stage 4 continues across the brane to a universe of mass $-M$ with the same $\lambda$ (\texttt{pairing-\allowbreak{}theory.j\allowbreak{}son}, key \texttt{T3.stage\allowbreak{}4Z2Pair}; checks \texttt{PAIR\_T2z\allowbreak{}2\_*}).
\item PROVED at the level of a Fock space (checks \texttt{PAIR\_T1k\allowbreak{}rein\_*}, cited from Stage 5 and provisional until its gate): it matters which field is called "the $-M$ universe". The image field $\gamma^8\Psi$ on the same Fock space carries the Krein metric $-B$; the formulas of $\mathcal L_{-m,-\lambda}$ give each of its quanta the energy $-\lvert\varepsilon\rvert$ and the charge $-1$, and the operator totals of the pair vanish. But the image is the same quantum system as $\Psi_+$ (the same operators, the same Fock space and the same state), so these totals vanish as identities $X+(-X)=0$, not by a compensation between two universes; measured with its own generators, the image has the energy and the charge of $\Psi_+$ (Section 16.6). An independently quantized $-M$ field with its own positive structure has the energy $+\lvert\varepsilon\rvert$ and the charge $+1$ per quantum: the energies add, and the charges cancel only under an additional assumption about its state (matter--antimatter document, §7.3). So no reading gives a cancellation between two independent, consistently quantized universes; whether any quantum description of a universe of mass $-M$ realizes the classical cancellation is OPEN.
\item COMPUTED (EXP-4b, Sections 11.11 and 11.13): the creation of \textbf{quanta} in particle--antiparticle pairs by a changing background. It was computed with the \textbf{Bogoliubov method}, the name of the computation of Section 11.11: follow each mode through the changing background and measure how much of a mode that starts in its positive-energy level ends on the negative-energy level. That weight is $\lvert\beta_k\rvert^2$ of Section 11.11, and the numbers $\beta_k$ are called \textbf{Bogoliubov coefficients}. L. Parker introduced the method for expanding universes (Phys. Rev. 183, 1057 (1969)). For example $na^3=4.412\times10^{-3}H_{\mathrm{inf}}^3$ for $m=H_{\mathrm{inf}}$ with the sudden end of inflation. It is the only creation process computed in the repository, and it creates quanta inside one universe, not universes.
\item Stage 1 said it in one sentence: the pairing of the masses $\pm m$ is "a structural property of the equations, not a claim that universes of masses $\pm M$ are created in pairs" (Stage-1 document, §1.3, item 6).
\end{itemize}

\textbf{Why it is hard.}

\begin{itemize}
\item \emph{The metric has no dynamics here.} The primordial field is prescribed, and what determines $a_4$ is open (Section 18.5). Without a dynamical geometry there is no process "creation of a universe".
\item \emph{A universe is not a quantum in a fixed background.} Creating one needs a quantum theory of the geometry itself: a wave function of the universe (the notebook's file name speaks of one) or a sum over geometries. The known proposals, the canonical equation of B. S. DeWitt (Phys. Rev. 160, 1113 (1967)), the creation "from nothing" of A. Vilenkin (Phys. Lett. B 117, 25 (1982)) and the no-boundary wave function of J. B. Hartle and S. W. Hawking (Phys. Rev. D 28, 2960 (1983)), are debated in four dimensions, and none has been formulated for this 8-dimensional theory.
\item \emph{A classical pair cannot appear from nothing.} Chapter 16 shows from the uniqueness of the solutions of the field equations in the good sector that a classical field which vanishes at one time vanishes at all times. A creation process must therefore be quantum, or start at a singular moment, which the notebook's field does not have (Chapter 16).
\item \emph{The matter side inherits Problems 1 and 2}: no interacting quantum theory, a Krein space, and the extra-time instability. In addition, if a $-M$ member with negative energies interacted with the $+M$ member, pairs of zero total energy could be produced without limit; Chapter 16 describes this stability question.
\item \emph{The cancelling pair has special parameters}: it needs $\lambda\to-\lambda$ (or $\lambda=0$), and as a quantum field it has no known reading as two independent universes: with the Krein metric $-B$ its $-M$ member is the same quantum system as the $+M$ member (Section 16.6).
\item \emph{"At $x_4=0$" is not well defined yet.} The notebook's field singles out no moment ($a_4$ is an arbitrary function, and the static member does not depend on $x_4$ at all), and in signature (4,4) a rotation by $\pi$ in the plane of two time-like directions, which belongs to the identity component of the symmetry group, reverses $x_4$ and the charge (matter--antimatter document, §5.7): "before" and "after" refer to a chosen slicing.
\end{itemize}

\textbf{Worked example: a reduced model of the primordial family} (derived here). A standard first step of quantum cosmology is to keep only a few variables of the metric and to quantize them. The family of Section 9.5 has one variable, $a_4$. Add a \textbf{lapse} $N(x_4)>0$ through $g_{44}=-N^2$, so that $d\tau=N\,dx_4$ is proper time along $x_4$.

\emph{The reduced action.} In the coordinates $(\zeta,\tau)$ the metric has the form of Lemma 9.1 with $f_p=H\zeta+\sigma_pa_4$, so the computation of Section 9.8 applies with the $t$-derivative replaced by the $\tau$-derivative: $R=6(da_4/d\tau)^2-42H^2=6\dot a_4^2/N^2-42H^2$, where $\dot a_4=da_4/dx_4$. No second derivative of $a_4$ appears, because the 7-volume is constant ($\Theta_4=0$ and $\sum_p\sigma_p=0$ in Section 9.8), so no boundary term is needed. With $\sqrt{\lvert g\rvert}=N\,e^{6H\zeta}$, the Einstein--Hilbert action $\tfrac1{2\kappa}\int\sqrt{\lvert g\rvert}\,R\,d^8x$ of Section 9.9, taken over a region of the slice whose proper 7-volume is $V=\int e^{6H\zeta}\,d\zeta\,d^6x$ (the same for every $a_4$ and every $x_4$, Section 9.4), becomes $\int L\,dx_4$ with

\[
L=\frac V{2\kappa}\Bigl(\frac{6\dot a_4^2}N-42H^2N\Bigr)-NV\rho .
\]

The last term is a homogeneous source with a constant energy density $\rho$, such as the condensate of Proposition 9.4 ($\rho=mS+U$). It is chosen so that its derivative with respect to $N$ at $N=1$ is $-VT_{44}=-V\rho$, which is what the definition $T_{\mu\nu}=-(2/\sqrt{\lvert g\rvert})\,\delta I_{\mathrm{matter}}/\delta g^{\mu\nu}$ of Section 9.9 requires (with $g^{44}=-1/N^2$).

\emph{The equations.} The momentum of $a_4$ is $p=\partial L/\partial\dot a_4=6V\dot a_4/(\kappa N)$, and the energy function $p\dot a_4-L$ of Section 5.2 is

\[
\mathcal H=N\Bigl[\frac{\kappa p^2}{12V}+\frac{21H^2V}\kappa+V\rho\Bigr]
\]

(insert $\dot a_4=\kappa pN/(6V)$: $p\dot a_4=\kappa p^2N/(6V)$ and $L=\kappa p^2N/(12V)-21H^2VN/\kappa-NV\rho$). The variable $N$ appears in $L$ without a time derivative, so its Euler--Lagrange equation is $\partial L/\partial N=0$. At $N=1$ this reads $-\tfrac V\kappa(3\dot a_4^2+21H^2)-V\rho=0$, that is

\[
\rho=-\frac{3\dot a_4^2+21H^2}\kappa=-\frac{3H^2\bigl(7+a_4'^2\bigr)}\kappa \qquad(\dot a_4=Ha_4'),
\]

exactly the requirement $\rho_{\mathrm{req}}$ of Section 9.10. The Euler--Lagrange equation of $a_4$ is $\frac{d}{dx_4}\bigl(6V\dot a_4/\kappa\bigr)=0$, so $a_4$ is linear, as found in Section 18.5 from Proposition 9.2. The reduced model agrees with Chapter 9 on both counts. (The other Einstein equations, those of the directions $\zeta$ and of the transverse trace, are not contained in it; they give condition (3) of Proposition 9.4.)

\emph{Quantization.} On wave functions $\psi(a_4)$ the operator $p=-i\,d/da_4$ satisfies the rule $[a_4,p]=i$ of Section 8.2: $a_4(-i\psi')+i(a_4\psi)'=i\psi$. The constraint $\mathcal H=0$ becomes the one-variable \textbf{Wheeler--DeWitt equation}

\[
-\frac\kappa{12V}\,\psi''+V\Bigl(\frac{21H^2}\kappa+\rho\Bigr)\psi=0,\qquad\text{that is}\qquad \psi''=\frac{12V^2}\kappa\Bigl(\frac{21H^2}\kappa+\rho\Bigr)\psi .
\]

If $\rho>-21H^2/\kappa$, in particular in vacuum, for every source of positive energy, and for a chiral pair, whose total energy--momentum vanishes, the right-hand side is a positive multiple of $\psi$ and the solutions are real exponentials $e^{\pm\ell a_4}$ with $\ell=V\sqrt{12(21H^2/\kappa+\rho)/\kappa}$: no oscillating, "classically allowed" solution exists, which is the quantum counterpart of the fact that no $a_4$ solves Einstein's equations with such a source. If $\rho<-21H^2/\kappa$ the solutions are $e^{\pm i\ell'a_4}$ with $\ell'=V\sqrt{12(-\rho-21H^2/\kappa)/\kappa}$, and the momentum $p=\pm\ell'$ belongs to the classical solution $a_4=ct$ with $3H^2c^2=-\kappa\rho-21H^2$; the notebook's $a_4=t$ belongs to $\rho=-24H^2/\kappa$ (Exercise 18.12).

\textbf{What the reduced model says, and what it does not.} In 8-dimensional Einstein gravity a chiral pair, whose total energy--momentum is zero, is invisible to the constraint: adding it changes nothing, and the geometry it lives in must be a vacuum solution, which the notebook's field is not. By Section 18.5 (derived there, not yet checked by machine) the static member is a vacuum solution of Einstein--Gauss--Bonnet gravity with tuned couplings, so there, as far as the bulk field equations are concerned, a chiral pair has a consistent home. The reduced model leaves out the Einstein equations of the other directions, the Lovelock terms, the brane, the matter modes and their instability, and the Krein structure; it has no boundary condition for $\psi$ (the no-boundary proposal and the tunnelling proposal are two choices from the literature); and it contains no process in which the number of universes changes. A creation amplitude needs a framework in which universes themselves can appear, and none is formulated in this repository.

\textbf{A first step.} (a) Redo the reduction (Exercise 18.12) and verify it exactly by machine: the variation with respect to $N$ must give $G^4{}_4$ of Section 9.8, and that with respect to $a_4$ the combination $G^1{}_1-G^5{}_5=2H^2a_4''$. (b) Add the Gauss--Bonnet term; its reduced form needs $E^{(2)}$ for a time-dependent $a_4$, the first step of Problem 3. (c) Formulate the matter side of a pair: the exact Fock model of \texttt{wolfram/\allowbreak{}Dirac16C\allowbreak{}omplexPa\allowbreak{}iring.wl} contains both readings, the image field with $-B$ and the independent quantization. (d) Compare with the one creation process that is computed, the creation of quanta in EXP-4b, measured by the Bogoliubov coefficients $\beta_k$ of Section 11.11.

\textbf{Where to start.} Chapters 15 and 16; the matter--antimatter document, §7; the four papers cited above; and these files (the keys \texttt{\detokenize{T1}}, \texttt{\detokenize{T1krein}}, \texttt{\detokenize{T2}}, \texttt{\detokenize{T3}}, \texttt{pairTota\allowbreak{}ls}, \texttt{notProve\allowbreak{}d} and \texttt{notebook\allowbreak{}Hypothes\allowbreak{}is} of the first one):

\begin{Verbatim}[fontsize=\small]
artifacts/dirac16complex/pair-creation/pairing-theory.json
wolfram/Dirac16ComplexPairing.wl                 the exact pairing theory, Fock model
studies/dirac16complex_cosmology/src/exp4.rs     EXP-4: creation of quanta
artifacts/dirac16complex/numerics/exp4/          its outputs
wolfram/Dirac16ComplexPrimordial.wl              the metric family and its curvature
\end{Verbatim}

\subsection{Problem 8: baryogenesis in this framework}

\textbf{The question.} The observed universe contains matter and almost no antimatter; the excess is measured by the baryon-to-photon ratio $\eta\approx6\times10^{-10}$ (matter--antimatter document, §2.2, from the Planck 2018 results and big-bang nucleosynthesis). The author asks that this theory solve the matter--antimatter mysteries. Chapter 17 shows that the theory as built does not: its charge is exactly conserved (Theorem M1), it has exact symmetries that reverse the charge without reversing $x_4$ (Theorem M2: charge conjugation $\Psi\to\Psi^\ast$ for dirac16complex00, a CP for dirac16complex), nothing computes a departure from equilibrium, and it contains no baryons of the Standard Model. The open problem is whether an extension of the framework could generate an asymmetry from a symmetric state, and what it would predict.

\textbf{Why it is hard.} Each of Sakharov's three conditions fails or is not addressed in the theory as built (the scorecard M6 of Chapter 17). In addition: identifying the charge of dirac16complex with baryon number (hypothesis H3 of Chapter 17) would need couplings to Standard-Model fields that the theory does not contain; the pair scenario M5 puts the asymmetry in as an initial condition and is not a mechanism; the quantum reading of the pair is unsettled (Section 18.9); the sign of the charge refers to a chosen slicing in (4,4) (matter--antimatter document, §5.7); and every quantum computation meets Problems 1 and 2.

\textbf{What is known.}

\begin{itemize}
\item PROVED (Chapter 17): M1, M2, and the classification M3 of the charge-violating terms allowed by the symmetry. For the anticommuting field no Majorana-type mass term exists; the derivative term $\Psi^TC\gamma^8\gamma^\mu D_\mu\Psi$ and quartic terms such as $Q_4$ and $Q_2$ survive. For the commuting field the two chiral mass terms $\Psi^TCP_\pm\Psi$ and the derivative term $\Psi^TC\gamma^\mu D_\mu\Psi$ (the form of the notebook's Lg[]) survive (checks \texttt{\detokenize{MA_M3_*}} of both matter--antimatter reports).
\item Derived in the matter--antimatter document (§10, item 2), not machine-checked: the phase of the coefficient of a single charge-2 term can be removed by redefining the field, so only relative phases between several such terms can be physical.
\item OPEN: whether a given combination of charge-violating terms breaks all 64 charge-reversing symmetries of each field (matter--antimatter document, §10); whether the CP of dirac16complex is implemented on the positive Fock space of the good sector (§5.6); charge-violating terms of higher order than the quartic examples are not classified (§6.3).
\item HYPOTHESIS: the conditional scenario M5 (Chapter 17), which predicts no value of $\eta$.
\end{itemize}

\textbf{Worked example 1: the first two Sakharov conditions for one complex variable} (derived here). Take a complex variable $z(t)$ with

\[
L=\dot z^\ast\dot z-\omega^2z^\ast z-\epsilon\bigl(z^2+z^{\ast2}\bigr),\qquad \epsilon\ \text{real}.
\]

Without the last term $L$ does not change under $z\to e^{i\alpha}z$, and Noether's theorem (Section 5.7, computed as in its Example 1) gives the conserved charge $Q=i(z^\ast\dot z-z\dot z^\ast)$, up to an overall sign, which is a convention; for $z=e^{-i\omega t}$ it is $Q=2\omega>0$. The last term has charge 2: it is multiplied by $e^{2i\alpha}$. The Euler--Lagrange equation of $z^\ast$ is $\ddot z=-\omega^2z-2\epsilon z^\ast$, and therefore

\[
\frac{dQ}{dt}=i\bigl(z^\ast\ddot z-z\ddot z^\ast\bigr)=i\bigl(-2\epsilon z^{\ast2}+2\epsilon z^2\bigr)=-4\epsilon\,\mathrm{Im}(z^2).
\]

The charge now changes: \textbf{the first condition} is met. But $L$ is unchanged by $z\to z^\ast$, which maps every solution $z(t)$ to a solution $z^\ast(t)$ with the charge $-Q(t)$. Starting from any collection of initial conditions that contains, with each one, its complex conjugate, the average charge stays zero at all times: this map is the "C" of the toy, and as long as it is a symmetry \textbf{the second condition} fails. A complex coefficient does not help by itself: the term $-(\epsilon z^2+\epsilon^\ast z^{\ast2})$ with $\epsilon=\lvert\epsilon\rvert e^{i\varphi}$ becomes $-\lvert\epsilon\rvert(w^2+w^{\ast2})$ for $w=e^{i\varphi/2}z$, and $w\to w^\ast$, that is $z\to e^{-i\varphi}z^\ast$, is again a charge-reversing symmetry. Two charge-violating terms whose phases cannot be removed together are needed (Exercise 18.14). The toy contains nothing that could meet \textbf{the third condition}: it has no out-of-equilibrium history.

\textbf{Worked example 2: the same computation for the commuting field} (derived here). Add to the Lagrangian of dirac16complex00 a charge-2 mass term of M3, $-\sqrt{\lvert g\rvert}\,(gX+g^\ast X^\ast)$ with $X=\Psi^TM\Psi$ and $M=CP_+$ or $CP_-$, which are real and symmetric because $C$ and $\gamma^8$ are real, symmetric and commute (Section 8.8, Chapter 17). The Euler--Lagrange expression of the old Lagrangian with respect to $\Psi^\ast$ is $\sqrt{\lvert g\rvert}\,CE$ with $E=\gamma^\mu D_\mu\Psi-(m+U')\Psi$, as in Chapter 17 (\texttt{dirac16c\allowbreak{}omplex00\allowbreak{}-theory.\allowbreak{}json}, key \texttt{fieldEqu\allowbreak{}ations}). Since $X^\ast=\Psi^\dagger M\Psi^\ast$ and $M$ is symmetric, the derivative of the new term with respect to $\Psi^\ast$ is $-2g^\ast\sqrt{\lvert g\rvert}\,M\Psi^\ast$, and the field equation becomes $CE=2g^\ast M\Psi^\ast$, that is $E=2g^\ast CM\Psi^\ast$ (because $C^2=1$). For commuting components the conjugate expression satisfies $\bar E=-E^\dagger C$ (use $(\gamma^a)^TC=-C\gamma^a$ and $(S^{ab})^TC=-CS^{ab}$, Sections 5.12 and 8.8), so $\bar E=-2g\,\Psi^TM$. The identity $\partial_\mu\bigl(\sqrt{\lvert g\rvert}\,j^\mu\bigr)=\sqrt{\lvert g\rvert}\,(\bar E\Psi+\bar\Psi E)$ of Theorem M1 holds for every configuration, so

\[
\partial_\mu\bigl(\sqrt{\lvert g\rvert}\,j^\mu\bigr)=\sqrt{\lvert g\rvert}\,\bigl(-2gX+2g^\ast X^\ast\bigr)=-4i\sqrt{\lvert g\rvert}\,\mathrm{Im}(gX),
\]

and with the real current $J^\mu=-ij^\mu$ of Section 8.12, when no charge flows out through the boundary of the slice,

\[
\frac{dQ}{dx_4}=-4\int\sqrt{\lvert g\rvert}\,\mathrm{Im}\bigl(g\,\Psi^TM\Psi\bigr)\,d^7x ,
\]

the field version of the toy. With one term the phase of $g$ can be removed and the charge conjugation $\Psi\to\Psi^\ast$, combined with a constant phase, remains a symmetry, exactly as in the toy.

\textbf{A first step.} (a) The toy and its two-term version (Exercises 18.13 and 18.14). (b) Worked example 2 numerically: the homogeneous state $\Psi=u(x_4)$ of dirac16complex00 is an exact solution for every potential (\texttt{dirac16c\allowbreak{}omplex00\allowbreak{}-theory.\allowbreak{}json}, key \texttt{primordi\allowbreak{}al.stage\allowbreak{}2Homogen\allowbreak{}eous}); add the term $gX+g^\ast X^\ast$, integrate the 16 complex equations in $x_4$ (with the fourth-order Runge--Kutta method of Section 10.11, or the CVODE driver of Stage 3), and test the formula for $dQ/dx_4$ along the solution, as the checkers of Section 10.11 test conservation laws. (c) Take both chiral mass terms $g_+X_++g_-X_-$ with a relative phase and test, with the transformation rule of Theorem M2 (Chapter 17), all 64 charge-reversing symmetries of the commuting field: a finite exact computation. (d) Only then add a departure from equilibrium, for example a time-dependent background as in EXP-4b, and compute the charge produced from a symmetric initial ensemble. (e) Whether Theorem T1 still pairs solutions after such a term is added must be derived again: since $\gamma^8$ is diagonal (hence symmetric) and commutes with $C$ and with $P_\pm$, the term $X=\Psi^TM\Psi$ becomes $\Psi^T\gamma^8M\gamma^8\Psi=X$ under $\Psi\to\gamma^8\Psi$, so the map now needs $g\to-g$ in addition to $m\to-m$ and $\lambda\to-\lambda$ (derived here; check it).

\textbf{Where to start.} Chapter 17 and the matter--antimatter document (§6 for M3, §10 for what would have to be added); A. D. Sakharov, JETP Lett. 5, 24 (1967); the review of L. Canetti, M. Drewes and M. Shaposhnikov, New J. Phys. 14, 095012 (2012); as the known published example of a pair-of-universes idea, L. Boyle, K. Finn and N. Turok, Phys. Rev. Lett. 121, 251301 (2018), on which nothing in this repository depends; and these files, of which the first three contain the transformation rule of M2 and the classification of all 1024 maps per field:

\begin{Verbatim}[fontsize=\small]
wolfram/Dirac16ComplexMatterAntimatter.wl
scripts/verify_dirac16complex_matter_antimatter.wls
scripts/check_dirac16complex_matter_antimatter.py
wolfram/Dirac16Complex00.wl       dirac16complex00, including the homogeneous states
provenance/DIRAC16COMPLEX_MATTER_ANTIMATTER.md
\end{Verbatim}

\subsection{Smaller open items}

The remaining open items of the book and of the project documents are smaller, or are parts of the problems above. Each has a first step.

\begingroup
\small
\begin{longtable}{@{}>{\raggedright\arraybackslash}p{0.267\linewidth}>{\raggedright\arraybackslash}p{0.267\linewidth}>{\raggedright\arraybackslash}p{0.267\linewidth}@{}}
\toprule
\textbf{Item} & \textbf{Where it arises} & \textbf{A first step} \\
\midrule
\endfirsthead
\toprule
\textbf{Item} & \textbf{Where it arises} & \textbf{A first step} \\
\midrule
\endhead
The $q$ term of the notebook is attributed to the rule of cell 1058 by a reconstruction (HYPOTHESIS) & Section 9.14; Stage-2 document, §11.3 & rerun the stored chain of cells 1058 to 1137 in an older Mathematica kernel, like the session that produced the stored outputs \\
Stability of the classical solutions under small perturbations is not studied & Stage-1 document, §1.3, item 3 & linearize the field equation around the homogeneous condensate of Proposition 9.4 and find the frequencies of the perturbations with the mode methods of Section 8.9 \\
Stability and quantum (Krein-space) status of the negative-energy source of Proposition 9.4 & Stage-2 document, §20, item 2 & the same linearization, and the expectation-value rule of Section 8.12 applied to the condensate \\
Sources that depend on the transverse coordinates, and plane waves with $\mathrm{Im}\,K=-3H$, $\mathrm{Re}\,K\ne0$ & Stage-2 document, §20, item 2 & extend Section 9.12; the off-diagonal components of $T_{\mu\nu}$ must vanish as well \\
A spin--statistics theorem for signature (4,4) & Section 8.15 & compare the commuting and the anticommuting quantization of the good sector; the unbounded classical energy of dirac16complex00 (Section 18.3) is the starting point \\
Is the CP of dirac16complex implemented on the positive Fock space of the good sector? & matter--antimatter document, §5.6 & test with the fixture whether its matrix preserves the good sector and is compatible with $J=B$ (Section 8.11 does the same for rotations) \\
Charge-violating terms of higher order than the quartic examples & matter--antimatter document, §6.3 & classify with the null-space method of M3 \\
Is the Kohn--Sham state reached by continuation in $\lambda$ the lowest self-consistent state when several exist? & Sections 13.10 and 13.11 & a search from many starting densities for $N=1016$ at $-\hat\lambda_2$ \\
The continuum limit of the 3-torus (6.6 percent change of $E/N$ at $\Delta k=0.125\,m$) & Section 13.15; erratum E4.11 & closed-shell runs at smaller $\Delta k$ and the same density \\
Correlation beyond the exchange-only functional & Section 13.8 & part of Problem 1 (Section 18.3) \\
The meaning of the Kohn--Sham model of dirac16complex00, with the Pauli filling imposed by prescription & Chapter 14; \texttt{pairing-\allowbreak{}theory.j\allowbreak{}son}, key \texttt{notProve\allowbreak{}d} & a statistical mechanics of the classical commuting field with Gaussian ensembles (the Wick sign of Section 18.3) \\
Dark matter: abundance, darkness, stabilised extra dimensions, structure formation, the hidden-space momentum $k_0$ & Section 11.13 & turn the EXP-4b yields into a relic density for physical values of $m$ and $H_{\mathrm{inf}}$ \\
Dark energy: other potentials, inhomogeneous states, a mechanism that stabilises the extra dimensions & Section 11.14 & repeat EXP-3 with a different $U(S)$ \\
The Fermi-degeneracy pressure of the condensate & Section 11.13; Stage-3 document, §4.6 & compute it with the uniform-gas formulas of Section 13.8 \\
Kohn--Sham states in a time-dependent primordial field ($a_4'\ne0$) & Chapter 13 uses only the static member & a time-dependent mean field in the good sector of the notebook's $a_4=t$ background \\
\bottomrule
\end{longtable}
\endgroup

\subsection{What we proved and what we assumed}

\textbf{Derived here}, with complete arguments in this chapter but without a check by any program of the repository: the power counting $[\Psi]=(D-1)/2$ and $[\lambda]=2-D$; the kinematic possibility that the contact interaction scatters quanta of the good sector into modes with extra-time momentum; the two-mode Krein toy, its threshold $\lvert g\rvert=\lvert\omega_0-\omega_1\rvert/2$ and its agreement with the extra-time dispersion relation at rest, including the linear growth at $k_5^2=m^2$ (whether the contact interaction actually couples modes of opposite Krein sign is left open); the frozen-coefficient classification of the modes that keep oscillating for three histories of $a_4$, and the failure of compact extra times to remove the growing modes; the coupling of the $2\times2$ blocks in pairs by an extra-time momentum; the Lovelock tensors of all orders for the static member, the required source they imply, a positive required energy density for $c_2>1/(20H^2)$, and the vacuum solution of Einstein--Gauss--Bonnet gravity with $c_2=1/(40H^2)$ and $c_0=21H^2$; the linearity of $a_4$ forced by a homogeneous source in Einstein gravity, with its slope; the uniform and isotropic source that the fixed static metric needs in every Einstein--Lovelock theory; the condition $5\rho+\sum_{\mu\ne4}p_{(\mu)}=0$ for every source of an ultrastatic metric in Einstein gravity and its Kohn--Sham form $6\sum_nw_n\varepsilon_nn_n=mS_p$; the pressure difference produced by different warps of 3-space and the extra times, which needs $f_3'\ne f_t'$, a difference of the warps that depends on $y$; the reduced model of the primordial family, its constraint, its agreement with Chapter 9 and its Wheeler--DeWitt equation; the charge-violation toy, the removability of a single phase and the condition for two terms; the formula $dQ/dx_4=-4\int\sqrt{\lvert g\rvert}\,\mathrm{Im}(g\Psi^TM\Psi)\,d^7x$ for the commuting field with a charge-2 mass term; and the extended pairing map with $g\to-g$. Each of these is a first step for a student, and the first task is to verify it exactly, by machine, in two independent programs.

\textbf{Computed for this chapter from committed files by a short calculation} (Exercise 18.10 shows one of them): the ranges of $5\rho+p_y+3p_3+3p_t$ on the committed Kohn--Sham profiles; the differences between the Rust and the reference eigenvalues of Section 18.7; the Krein-image diagnostics of Section 18.6; the counts and largest deviations of the Stage-5 numerics in Section 18.8. The files are named where the numbers are used.

\textbf{Quoted from earlier chapters and from the committed reports}, with the status given there: every statement listed under "What is known" in Sections 18.3 to 18.10.

\textbf{Quoted from the literature and not derived}: the power-counting argument for non-renormalisability; the argument against several time directions; the well-posedness of the ultrahyperbolic problem under a nonlocal constraint; Lovelock's theorem; the junction conditions of Gauss--Bonnet gravity; the proposals of DeWitt, Vilenkin and Hartle and Hawking; the attribution of the Bogoliubov method of particle creation to Parker (the method itself is derived in Section 11.11); Sakharov's conditions; the observed value of $\eta$.

\textbf{Assumed}: the frozen-coefficient estimate of Section 18.4; the matter term $-NV\rho$ with a constant $\rho$ in the reduced model of Section 18.9; the quantization rule $p=-i\,d/da_4$ for a variable of the metric.

\textbf{Not claimed}: that any problem of this chapter is solved; that universes are created in pairs, by the big bang or otherwise; that this theory explains the matter--antimatter asymmetry; that the Einstein--Gauss--Bonnet couplings of Section 18.5 describe an acceptable theory; that any Kohn--Sham state is the source of the primordial field.

\subsection{Exercises}

\textbf{Exercise 18.1.} (a) Compute $[\Psi]$ and $[\lambda]$ for $D=2,4,5,8$. For which $D$ is the contact coupling dimensionless? (b) Two quanta of mass $m=1$ at rest turn into two quanta with the wave numbers $(k_1,k_5)=(p,q)$ and $(-p,-q)$. Show that energy conservation requires $p^2=q^2$, and compute the energy of each final quantum for $p=q=2$.

\textbf{Exercise 18.2.} For the Krein toy of Section 18.3 with $\omega_0=3$ and $\omega_1=1$, compute the eigenvalues for $g=\tfrac12$, $1$ and $2$. For $g=1$ show that $h-2$ squares to zero and describe the growth of the solutions. Compare with two modes of the same Krein sign and $g=2$.

\textbf{Exercise 18.3.} (a) Show that $h=\begin{pmatrix}m&g\\-g^\ast&-m\end{pmatrix}$ has the eigenvalues $\pm\sqrt{m^2-\lvert g\rvert^2}$. With $\lvert g\rvert=\lvert k_5\rvert$, compare with $E^2$ of Section 8.9 at zero 3-momentum. (b) Give the growth rate for $m=1$, $k_5=2$. Which sample of the table of Section 8.9 has the same $E^2$? (c) Take $m=1$ and $g=k_5=1$. Show that $h^2=0$, write the solution of $i\dot u=hu$ with $u(0)=(1,0)^T$, and describe how it grows.

\textbf{Exercise 18.4.} With the two rules of Section 13.4 show that $\gamma^0\gamma^5$ anticommutes with $J=\gamma^0\gamma^1\gamma^4$ and with $K_2=\gamma^5\gamma^6$ and commutes with $K_1=\gamma^2\gamma^3$. Compute $(\gamma^0\gamma^5)^2$ and compare it with $A_1^2$ and $A_4^2$ of Section 13.4.

\textbf{Exercise 18.5.} Let $x_5$ be a circle of radius $R$, $m=1$ and $\mathbf k=0$. (a) For $R=2$ compute $E^2$ for $n=0,\pm1,\pm2,\pm3$ and say how each mode behaves; give the growth rate for $n=3$. (b) Which modes grow for $R=\tfrac12$?

\textbf{Exercise 18.6.} (a) Compute $6!/(6-2k)!$ and $7!/(7-2k)!$ for $k=0,1,2,3$ and the entries of $E^{(2)}$ and $E^{(3)}$ of the static member. (b) Check that the row $k=1$ is the Einstein tensor of Section 9.15. (c) What does the counting of Section 18.5 give for $k=4$, and why does this agree with Section 4.9?

\textbf{Exercise 18.7.} Use the coefficients $c_k$ of Section 18.5 (called $\alpha_k$ in Chapter 4). (a) Derive the vacuum couplings $c_2=1/(40H^2)$ and $c_0=21H^2$ for $c_1=1$, $c_3=0$. (b) With $c_1=1$, $c_2=1/(10H^2)$ and $c_0=c_3=0$ compute $\rho_{\mathrm{req}}$, $p_{\mathrm{req}}$ and $w$. (c) Show that Einstein gravity with a cosmological constant alone ($c_2=c_3=0$) has no static vacuum of this form.

\textbf{Exercise 18.8.} For the dirac16complex00 source of Section 18.6 ($m=-30$, $\lambda=\tfrac{125}6$, $S=\tfrac65$, $H=\kappa=1$) check $mS=-36$, $\lambda S^2=30$, $M_{\mathrm{eff}}$, $\rho$, $p$ and $w$, the kinetic and potential energies $\mathrm{KE}_L=\tfrac12S(m+U')$ and $\mathrm{PE}_L=\rho-\mathrm{KE}_L$ of a homogeneous state (Chapter 7; \texttt{dirac16c\allowbreak{}omplex00\allowbreak{}-theory.\allowbreak{}json}, key \texttt{energySp\allowbreak{}lits}), and the condition $5\rho+7p=0$.

\textbf{Exercise 18.9.} In $D$ dimensions write $t$ for the index of the time, the direction in which the metric is static (in Section 18.6 it is $x_4$, with the index 4). An ultrastatic metric is then diagonal, has $g_{tt}=-1$ and does not depend on the time. Use the trace-reversed Einstein equations $R_{\mu\nu}=\kappa\bigl(T_{\mu\nu}-g_{\mu\nu}T/(D-2)\bigr)$ of Section 9.9 to show that every source of an ultrastatic metric satisfies $(D-3)\rho+\sum p=0$, the sum running over the pressures $p_{(\mu)}=T^\mu{}_\mu$ (no sum) of the $D-1$ directions other than the time. For $D=4$ and a source made of dust of density $\rho_m$ and a cosmological constant $\Lambda$, show that $\Lambda=\kappa\rho_m/2$.

\textbf{Exercise 18.10.} (a) Repeat worked example 3 of Section 18.6 for $N=1016$ at $\lambda=0$, where \texttt{rust/emt\allowbreak{}/emt-sum\allowbreak{}mary.csv} gives $\langle\rho\rangle=0.425999$ and $\langle S_p\rangle=-0.0872325$. (b) Write the Python lines that compute the smallest and the largest value of $5\rho+p_y+3p_3+3p_t$ on the committed profile of this run, and check the identity of worked example 2 on it.

\textbf{Exercise 18.11.} From the table of Section 18.7 compute, for both levels, the differences Rust (301) minus reference, Rust (601) minus reference and Rust (301) minus Rust (601). Test the rule "deviation at most the tolerance plus the measured grid uncertainty $\lvert\varepsilon(601)-\varepsilon(301)\rvert$", with the tolerance $1.054\times10^{-6}$ for the level of parity $+$ and at least $1.0\times10^{-6}$ for the level of parity $-$. Why does this arithmetic alone not justify changing the checker?

\textbf{Exercise 18.12.} (a) Derive the momentum $p$ and the function $\mathcal H$ of the reduced model of Section 18.9. (b) For $\rho=-33H^2/\kappa$ find the slope $c$ of the classical solution and compare with Exercise 9.5. (c) For a chiral pair ($\rho=0$) write the general solution of the Wheeler--DeWitt equation.

\textbf{Exercise 18.13.} For the toy of Section 18.10 with $\omega=1$ and $\epsilon=\tfrac1{10}$ write $z=x+iy$. (a) Show that $Q=2(y\dot x-x\dot y)$ and that $x$ and $y$ oscillate with the frequencies $\omega_1=\sqrt{\omega^2+2\epsilon}$ and $\omega_2=\sqrt{\omega^2-2\epsilon}$. (b) Show that real initial data, for example $z(0)=1$, $\dot z(0)=\tfrac3{10}$, give $Q=0$ at all times. (c) For $z(0)=1$, $\dot z(0)=-i$ compute $Q(t)$ and check $Q(0)=2$ and $dQ/dt=-4\epsilon\,\mathrm{Im}(z^2)$ at $t=0$.

\textbf{Exercise 18.14.} Add to the toy of Section 18.10 the terms $-(\epsilon_1z^2+\epsilon_1^\ast z^{\ast2})-(\epsilon_2z^4+\epsilon_2^\ast z^{\ast4})$ with $\epsilon_k=\lvert\epsilon_k\rvert e^{i\varphi_k}$. Show that a map $z\to e^{i\beta}z^\ast$ is a symmetry for some $\beta$ exactly when $2\varphi_1-\varphi_2$ is a multiple of $\pi$, and that the combination $2\varphi_1-\varphi_2$ does not change under $z\to e^{i\alpha}z$. What would still be missing for an asymmetry?

\subsection{Answers to the exercises}

\textbf{Answer 18.1.} (a) $[\Psi]=(D-1)/2$ and $[\lambda]=2-D$: for $D=2$, $\tfrac12$ and $0$; for $D=4$, $\tfrac32$ and $-2$; for $D=5$, $2$ and $-3$; for $D=8$, $\tfrac72$ and $-6$. The coupling is dimensionless only for $D=2$. (b) Each final quantum has $E=\sqrt{1+p^2-q^2}$, and $2E=2$ requires $p^2=q^2$. For $p=q=2$: $E=\sqrt{1+4-4}=1$ each, total 2, the energy of the two quanta at rest.

\textbf{Answer 18.2.} The threshold is $\lvert\omega_0-\omega_1\rvert/2=1$. $g=\tfrac12$: $2\pm\sqrt{1-\tfrac14}=2.866025$ and $1.133975$. $g=1$: the double eigenvalue 2. $g=2$: $2\pm i\sqrt3$, growth rate $\sqrt3=1.732051$. For $g=1$, $h=\begin{pmatrix}3&1\\-1&1\end{pmatrix}$ and $h-2=\begin{pmatrix}1&1\\-1&-1\end{pmatrix}$, whose square is $\begin{pmatrix}1-1&1-1\\-1+1&-1+1\end{pmatrix}=0$. So $e^{-iht}=e^{-2it}\bigl(1-i(h-2)t\bigr)$ (the series of the exponential of the nilpotent part stops after two terms): the solutions grow linearly, as in the case $E^2=0$ of Section 8.9. With the same Krein sign, $h=\begin{pmatrix}3&2\\2&1\end{pmatrix}$ has the real eigenvalues $2\pm\sqrt5=4.236068$ and $-0.236068$.

\textbf{Answer 18.3.} (a) $\det(h-x)=(m-x)(-m-x)+\lvert g\rvert^2=x^2-m^2+\lvert g\rvert^2$, so $x=\pm\sqrt{m^2-\lvert g\rvert^2}$. With $\lvert g\rvert=\lvert k_5\rvert$ this is $\pm E$ with $E^2=m^2-k_5^2$, the relation of Section 8.9 at zero 3-momentum. (b) For $m=1$, $k_5=2$: $E^2=-3$ and the growth rate is $\sqrt3=1.732051$; the table of Section 8.9 has the same $E^2=-3$ for the sample $m=1$, $k_7=2$. (c) $h=\begin{pmatrix}1&1\\-1&-1\end{pmatrix}$ and $h^2=\begin{pmatrix}1-1&1-1\\-1+1&-1+1\end{pmatrix}=0$: both eigenvalues are 0, and $E^2=1-1=0$. The series of the exponential stops after two terms, $e^{-iht}=1-iht$, so $u(t)=(1-iht)u(0)=\bigl(1-it,\ it\bigr)^T$. Its squared Hilbert length is $u^\dagger u=(1+t^2)+t^2=1+2t^2$, which grows without bound (the components grow linearly in $t$); with $B=\mathrm{diag}(1,-1)$ of Section 18.3 its Krein norm $u^\dagger Bu=(1+t^2)-t^2=1$ stays constant, as it must for $h^\dagger B=Bh$. This is the linear growth of the sample $m=1$, $k_5=1$ of Section 8.9 (Exercise 8.6), and the same matrix as $h-2$ in Answer 18.2.

\textbf{Answer 18.4.} $\gamma^0$ is a factor of $J$ ($k=3$) and commutes with it; $\gamma^5$ is not a factor and anticommutes with it; so $\gamma^0\gamma^5J=-J\gamma^0\gamma^5$. $\gamma^0$ is not a factor of $K_2$ ($k=2$) and commutes with it; $\gamma^5$ is a factor and anticommutes; so $\gamma^0\gamma^5$ anticommutes with $K_2$. Neither $\gamma^0$ nor $\gamma^5$ is a factor of $K_1$, so both commute with it. $(\gamma^0\gamma^5)^2=(-1)^1(\gamma^0)^2(\gamma^5)^2=-(1)(-1)=+1$, the same as $A_4^2=(\gamma^0\gamma^4)^2=+1$ and opposite to $A_1^2=-1$: an extra-time momentum enters the reduced equation like the frequency term, not like a 3-momentum.

\textbf{Answer 18.5.} (a) $E^2=1-n^2/4$: $n=0$ gives 1 and $n=\pm1$ gives $\tfrac34$ (oscillation); $n=\pm2$ gives 0 (linear growth); $n=\pm3$ gives $-\tfrac54$, exponential growth with $\varkappa=\sqrt5/2=1.118034$, and every larger $\lvert n\rvert$ grows faster. (b) $E^2=1-4n^2<0$ for every $n\ne0$: all modes except $n=0$ grow. A smaller circle makes things worse.

\textbf{Answer 18.6.} (a) $6!/(6-2k)!=1,30,360,720$ and $7!/(7-2k)!=1,42,840,5040$. With $K=-H^2$: $E^{(2)}=-\tfrac{H^4}2\cdot360=-180H^4$ on the seven directions and $-\tfrac{H^4}2\cdot840=-420H^4$ on $x_4$; $E^{(3)}=-\tfrac{-H^6}2\cdot720=360H^6$ and $-\tfrac{-H^6}2\cdot5040=2520H^6$. (b) $k=1$: $-\tfrac{-H^2}2\cdot30=15H^2$ and $-\tfrac{-H^2}2\cdot42=21H^2$, the entries of $G^\mu{}_\nu=\mathrm{diag}(15,15,15,15,21,15,15,15)H^2$. (c) For $k=4$ one needs lists of 8 different elements of $\Sigma$, which has only 7 elements: there are none, and $E^{(4)}=0$. Section 4.9 found the same: the generalized delta of $E^{(4)}$ has $2k+1=9>8$ upper indices and vanishes.

\textbf{Answer 18.7.} (a) With $c_3=0$ the first condition of Section 18.5 is $c_1-40c_2H^2=0$, so $c_2=1/(40H^2)$; the second gives $c_0=2(15H^2-180H^4/(40H^2))=2(15-4.5)H^2=21H^2$. Check on $x_4$: $-\tfrac{21}2+21-\tfrac{420}{40}=0$. (b) $\kappa\rho_{\mathrm{req}}=-21H^2+420H^4/(10H^2)=21H^2$ and $\kappa p_{\mathrm{req}}=15H^2-180H^4/(10H^2)=-3H^2$, so $\rho_{\mathrm{req}}=21H^2/\kappa>0$, $p_{\mathrm{req}}=-3H^2/\kappa$ and $w=-\tfrac17$. (c) The first condition becomes $c_1=0$, which contradicts $c_1=1$: no value of $\Lambda$ makes both components vanish, because $E^{(1)}$ has different entries on $x_4$ (21) and on the seven directions (15), while $E^{(0)}$ has equal ones.

\textbf{Answer 18.8.} $mS=-30\cdot\tfrac65=-36$; $\lambda S^2=\tfrac{125}6\cdot\tfrac{36}{25}=30$; $M_{\mathrm{eff}}=-30+\tfrac{125}6\cdot\tfrac65=-30+25=-5$; $\rho=mS+\tfrac\lambda2S^2=-36+15=-21$; $p=\tfrac\lambda2S^2=15$; $w=-\tfrac57$. $\mathrm{KE}_L=\tfrac12\cdot\tfrac65\cdot(-30+25)=-3$ and $\mathrm{PE}_L=-21+3=-18$, the values recorded in \texttt{dirac16c\allowbreak{}omplex00\allowbreak{}-theory.\allowbreak{}json}. $5\rho+7p=-105+105=0$.

\textbf{Answer 18.9.} Label the time direction by $t$ (in this book's eight dimensions $t$ is $x_4$, and the index 0 belongs to the space-like coordinate $x_0$, so it is not used for the time here). For an ultrastatic metric $R_{tt}=0$ (the argument of Section 18.6 for $R_{44}$ works in every dimension). With $T_{tt}=\rho$, $g_{tt}=-1$ and $T=g^{tt}T_{tt}+\sum p=-\rho+\sum p$: $R_{tt}=\kappa\bigl(T_{tt}-g_{tt}T/(D-2)\bigr)=\kappa\bigl(\rho+(-\rho+\sum p)/(D-2)\bigr)=\kappa\bigl((D-3)\rho+\sum p\bigr)/(D-2)=0$. For $D=8$ this is $5\rho+\sum_{\mu\ne4}p_{(\mu)}=0$, the condition of Section 18.6. For $D=4$: $\rho+3p=0$ with $p$ the (equal) pressure. A cosmological constant, moved to the right-hand side of $G^\mu{}_\nu+\Lambda\delta^\mu{}_\nu=\kappa T^\mu{}_\nu$, acts like a source with $\rho_\Lambda=\Lambda/\kappa$ and $p_\Lambda=-\Lambda/\kappa$. With dust: $\rho_m+\Lambda/\kappa-3\Lambda/\kappa=0$, so $\Lambda=\kappa\rho_m/2$.

\textbf{Answer 18.10.} (a) The image has $\langle\rho\rangle=-0.425999$, so $\kappa=21/0.425999=49.30>0$ from the energy density; the mass condition $(-1)(-0.0872325)=-36/\kappa$ gives $\kappa=-412.7<0$ (the $+M$ state gave $+412.7$, Section 13.14). The two conditions contradict each other, and the coupling condition cannot hold at $\lambda=0$. (b) In the repository root:

\begin{Verbatim}[fontsize=\small]
import csv
path = "artifacts/dirac16complex/kohn-sham/rust/emt/m1_L3_N1016_lam0_T0/profiles.csv"
rows = list(csv.DictReader(open(path)))
def comb(r):
    return (5 * float(r["rho"]) + float(r["p_y"])
            + 3 * float(r["p_3"]) + 3 * float(r["p_t"]))
def ks(r):
    return 6 * (float(r["rho"]) + float(r["L_s"])) - 1.0 * float(r["S_p"])
values = [comb(r) for r in rows]
print(min(values), max(values))
print(max(abs(comb(r) - ks(r)) / max(1.0, abs(comb(r))) for r in rows))
\end{Verbatim}

The first line prints about 1.0852 and 28609.2, the second a number of the order of $10^{-16}$: the identity holds to rounding. (Here $m=1$, and $\sum_nw_n\varepsilon_nn_n=\rho+L_s$ by Section 13.14.)

\textbf{Answer 18.11.} Parity $+$: $-1.05374908333+1.05374688831=-2.195\times10^{-6}$; $-1.05374707094+1.05374688831=-1.83\times10^{-7}$; $-1.05374908333+1.05374707094=-2.012\times10^{-6}$. Parity $-$: $-1.976\times10^{-6}$, $-1.64\times10^{-7}$ and $-1.812\times10^{-6}$. The extended rule passes: $2.195\times10^{-6}\le1.054\times10^{-6}+2.012\times10^{-6}=3.066\times10^{-6}$ and $1.976\times10^{-6}\le1.0\times10^{-6}+1.812\times10^{-6}=2.812\times10^{-6}$. But a rule chosen after seeing that it passes proves nothing. The change is justified only if an independent measurement, such as a finer-grid Rust run converging toward the reference, shows that the difference between 301 and 601 points is a real grid error of the Rust value; the rule must then be applied to every run with such a partner and recorded as an erratum.

\textbf{Answer 18.12.} (a) $p=\partial L/\partial\dot a_4=\tfrac V{2\kappa}\cdot\tfrac{12\dot a_4}N=\tfrac{6V\dot a_4}{\kappa N}$, and $p\dot a_4-L=N\bigl[\kappa p^2/(12V)+21H^2V/\kappa+V\rho\bigr]$ as in Section 18.9. (b) $3H^2c^2=-\kappa\rho-21H^2=33H^2-21H^2=12H^2$, so $c=\pm2$; Exercise 9.5 found $\rho_{\mathrm{req}}=-33H^2/\kappa$ for $a_4=2t$. (c) $\psi''=\ell^2\psi$ with $\ell=V\sqrt{12\cdot21}\,H/\kappa=6\sqrt7\,VH/\kappa=15.8745\,VH/\kappa$, so $\psi=A\,e^{\ell a_4}+B\,e^{-\ell a_4}$: no oscillating solution.

\textbf{Answer 18.13.} (a) $z^\ast\dot z=(x-iy)(\dot x+i\dot y)=x\dot x+y\dot y+i(x\dot y-y\dot x)$ and $z\dot z^\ast$ is its complex conjugate, so $Q=i\cdot2i(x\dot y-y\dot x)=2(y\dot x-x\dot y)$. Since $z^2+z^{\ast2}=2(x^2-y^2)$, the Lagrangian is $\dot x^2+\dot y^2-\omega^2(x^2+y^2)-2\epsilon(x^2-y^2)$, with the equations $\ddot x=-(\omega^2+2\epsilon)x$ and $\ddot y=-(\omega^2-2\epsilon)y$. (b) Real data have $y(0)=\dot y(0)=0$, so $y=0$ at all times and $Q=0$. (c) $x=\cos\omega_1t$ and $y=-\sin(\omega_2t)/\omega_2$ with $\omega_1=\sqrt{1.2}=1.095445$ and $\omega_2=\sqrt{0.8}=0.894427$, and $Q(t)=2\bigl[(\omega_1/\omega_2)\sin\omega_1t\,\sin\omega_2t+\cos\omega_1t\,\cos\omega_2t\bigr]$. At $t=0$, $Q=2$. Since $z^2=x^2-y^2+2ixy$, $\mathrm{Im}(z^2)=2xy=0$ at $t=0$, so $dQ/dt=0$ there; differentiating $Q(t)$ gives the same, because each term of $dQ/dt$ contains $\sin\omega_1t$ or $\sin\omega_2t$. For later times $Q$ oscillates, while the conjugate solution carries $-Q(t)$.

\textbf{Answer 18.14.} The kinetic and mass terms do not change under $z\to e^{i\beta}z^\ast$. The term $\epsilon_1z^2$ becomes $\epsilon_1e^{2i\beta}z^{\ast2}$, which must equal the old coefficient of $z^{\ast2}$, $\epsilon_1^\ast$: $e^{2i\beta}=e^{-2i\varphi_1}$. The term $\epsilon_2z^4$ requires $e^{4i\beta}=e^{-2i\varphi_2}$. The first gives $e^{4i\beta}=e^{-4i\varphi_1}$, so both hold exactly when $e^{-4i\varphi_1}=e^{-2i\varphi_2}$, that is when $2\varphi_1-\varphi_2$ is a multiple of $\pi$ (and then $\beta=-\varphi_1$ works). Under $z\to e^{i\alpha}z$ the coefficients become $\epsilon_1e^{2i\alpha}$ and $\epsilon_2e^{4i\alpha}$, and $2(\varphi_1+2\alpha)-(\varphi_2+4\alpha)=2\varphi_1-\varphi_2$: the combination is physical. When it is not a multiple of $\pi$, no map of this form reverses the charge; still missing are a check of all other charge-reversing maps, a departure from equilibrium, and a link to baryons.

\section{Reproducing everything}

This appendix shows how to rerun the programs of this book on your own computer, starting from nothing but an internet connection. Every machine check that the book names, and every number that it takes from a program, comes from a committed file of the repository, and the text names the file; derivations marked "derived here" or "not machine-checked" (in Chapters 11, 17 and 18, for example) have no machine check at all. This appendix regenerates the committed files and checks. It gives every command twice, once for PowerShell (the command shell of Windows) and once for Bash (the shell of Git Bash on Windows, of macOS and of Linux), it says what each command prints when it succeeds and how long it took on the test computer, and it states plainly which stages can be checked completely today and which cannot yet.

\textbf{Which commit, and what was measured.} The commands of this appendix were run for this edition from fresh clones of the public repository at commit \texttt{\detokenize{1f2dd69}} (2026-09-30), the newest commit when the appendix was revised, on the computer described in Section 19.2, unless a paragraph says otherwise. That commit contains drafts of the chapters of this book, not their final text; the final files of the book are added by a later commit. To work with exactly the files that were measured, check that commit out after cloning (\texttt{git chec\allowbreak{}kout 1f2\allowbreak{}dd69}, Section 19.4); at a later commit, \texttt{git diff\allowbreak{} --stat \allowbreak{}1f2dd69 \allowbreak{}HEAD} lists every file that has changed since. The wall times, the numbers of unit tests and the printed lines quoted in this appendix are measurements of these runs. Their logs are not committed, so, unlike the other numbers of the book, they cannot be looked up in the repository: a reader checks them by running the same commands. Where a run takes hours, the appendix does not repeat it; it quotes the wall time recorded in the committed gate logs, in the headers of the gate scripts, in \texttt{README.m\allowbreak{}d} or in \texttt{HANDOFF.\allowbreak{}md}, and names the source.

\subsection{What reproducing means, and the state of each stage}

\textbf{Three levels.} Section 0.8 described three levels of checking a statement: following the derivation, reading the committed report of the machine check, and running the programs again. This appendix is about the third level. Running a program again means: the program recomputes a result from its inputs, writes a new report, and the new report is compared with the committed one. Two kinds of agreement are used. For exact computations (integers, fractions, symbols) and for the numerical outputs of a fixed program on a fixed engine, the new file must be \textbf{byte-identical} to the committed one: the two files are the same sequence of bytes. For numerical results computed in another way (a finer grid, tighter tolerances, a different solver) the new numbers must agree with the old ones within a stated tolerance.

\textbf{Fingerprints.} Byte identity is tested with a fingerprint of the file. The \textbf{sha256} of a file is a number computed from all of its bytes and written with 64 \textbf{hexadecimal} digits. Hexadecimal digits are the digits of base 16: the ten digits 0 to 9, followed by the letters a, b, c, d, e and f, which stand for the values 10 to 15 (the upper-case letters A to F denote the same digits). For example, the hexadecimal number \texttt{\detokenize{1f}} is $1\cdot16+15=31$. Two files have the same sha256 exactly when they are byte-identical, for every practical purpose: changing a single byte changes the fingerprint completely. The names of the commits of the repository (Section 19.4) are hexadecimal numbers of the same kind.

\textbf{Gates.} For each finished stage the repository contains a \textbf{gate}: a script that runs every program of the stage in a fixed order, stops at the first failure, compares every output with the committed file, and ends with one line that says \texttt{\detokenize{OK}} or \texttt{\detokenize{FAILED}}. Each gate exists twice, as a PowerShell script (\texttt{\detokenize{.ps1}}) and as a Bash script (\texttt{\detokenize{.sh}}); the two are called twins, run the same steps and print the same final line.

\textbf{The state at commit \texttt{\detokenize{1f2dd69}}.} The table lists what can be checked. The column "Gate" names the scripts in the folder \texttt{\detokenize{scripts/}}.

\begingroup
\small
\begin{longtable}{@{}>{\raggedright\arraybackslash}p{0.200\linewidth}>{\raggedright\arraybackslash}p{0.200\linewidth}>{\raggedright\arraybackslash}p{0.200\linewidth}>{\raggedright\arraybackslash}p{0.200\linewidth}@{}}
\toprule
\textbf{Stage} & \textbf{Gate} & \textbf{Final line of a successful run} & \textbf{State at commit \texttt{\detokenize{1f2dd69}}} \\
\midrule
\endfirsthead
\toprule
\textbf{Stage} & \textbf{Gate} & \textbf{Final line of a successful run} & \textbf{State at commit \texttt{\detokenize{1f2dd69}}} \\
\midrule
\endhead
1: the field in an arbitrary gravitational field & \texttt{verify\_s\allowbreak{}tage1\_ar\allowbreak{}bitrary\_\allowbreak{}field} & \texttt{stage1\_a\allowbreak{}rbitrary\allowbreak{}\_field\_v\allowbreak{}erificat\allowbreak{}ion=OK} & complete; both twins passed from fresh public clones, and the Bash twin passed again for this edition (Section 19.6) \\
2: the primordial field & \texttt{verify\_s\allowbreak{}tage2\_pr\allowbreak{}imordial\allowbreak{}\_field} & \texttt{stage2\_p\allowbreak{}rimordia\allowbreak{}l\_field\_\allowbreak{}verifica\allowbreak{}tion=OK} & complete; both twins passed from fresh public clones, and the Bash twin passed again for this edition (Section 19.7) \\
3: dark-sector numerics & \texttt{verify\_s\allowbreak{}tage3\_da\allowbreak{}rk\_secto\allowbreak{}r} & \texttt{stage3\_d\allowbreak{}ark\_sect\allowbreak{}or\_verif\allowbreak{}ication=\allowbreak{}OK} & complete; both twins passed from fresh public clones of earlier commits; at \texttt{\detokenize{1f2dd69}} the Bash twin passes steps 00 to 30 and stops in step 31, all unit tests, because four tests of Stage 4 fail (Sections 19.8 and 19.11) \\
4: Kohn--Sham states & \texttt{verify\_s\allowbreak{}tage4\_ko\allowbreak{}hn\_sham} & \texttt{stage4\_k\allowbreak{}ohn\_sham\allowbreak{}\_verific\allowbreak{}ation=OK} & not finished: the gate has never been run to the end and cannot pass yet (Section 19.9) \\
5: dirac16complex00 and the pairing theorems & none yet & none & exact theory committed, every check of its four reports true, one of its results cited as provisional (below); Kohn--Sham pair numerics partial; documents and gate not written (Section 19.10) \\
matter and antimatter & none & none & exact theory committed and checked; document built and registered (36 pages); all 49 tests of its two test files pass (Sections 19.10 and 19.11) \\
this textbook & the assembler and the PDF builder & \texttt{textbook\allowbreak{}\_assembl\allowbreak{}y=OK}, \texttt{provenan\allowbreak{}ce\_pdf=O\allowbreak{}K} & Section 19.12 \\
\bottomrule
\end{longtable}
\endgroup

One Stage-5 result needs a warning. The exact Fock-level checks \texttt{PAIR\_T1k\allowbreak{}rein\_*} and \texttt{S5\_T1kre\allowbreak{}in\_*} are all true, but the text that the Stage-5 files attach to them, which calls $-|\varepsilon|$ and $-1$ the energy and the charge of an "image universe", is to be corrected (Section 15.4 explains why these are not the image field's own energy and charge; \texttt{HANDOFF.\allowbreak{}md}, section 0.4, item C, records the correction to be made in \texttt{wolfram/\allowbreak{}Dirac16C\allowbreak{}omplexPa\allowbreak{}iring.wl} and \texttt{scripts/\allowbreak{}check\_di\allowbreak{}rac16com\allowbreak{}plex\_pai\allowbreak{}ring.py}). Until the Stage-5 gate has passed, the matter--antimatter reports cite this result as PROVISIONAL, and the correction may change the names and the number of the checks of the Stage-5 pairing reports.

The sources of the table are \texttt{README.m\allowbreak{}d} (section "Stages and documents"), \texttt{HANDOFF.\allowbreak{}md} (sections 0.4 and 2), the committed gate logs in \texttt{handoff/\allowbreak{}reviews/}, the committed reports named in Sections 19.6 to 19.12 and the runs of this edition. The rest of this appendix goes through the stages one by one.

\subsection{The computer and the tools}

\textbf{The test computer.} The committed outputs and every timing of this appendix come from one computer: an Intel Core Ultra 9 275HX processor with 24 logical processors, running Windows 11 Pro for Workstations (build 10.0.26200). While the commands of this edition were timed, other jobs of the project were running on the same computer and kept its processor busy (Windows reported a load of 100 percent), and three groups of the runs of this appendix ran at the same time, each in its own fresh clone: the six exact programs of Section 19.10; the Stage-3 gate of Section 19.8, whose step 31 is the unit-test run of Section 19.11; and the remaining commands of Sections 19.3 to 19.12. Two short runs were made afterwards in these clones, as Sections 19.10 and 19.11 say. The timings are therefore upper values, and an otherwise idle computer is faster. The Stage-3 programs were also tested on Ubuntu 24.04 running inside Windows (student guide \texttt{provenan\allowbreak{}ce/DIRAC\allowbreak{}16COMPLE\allowbreak{}X\_STUDEN\allowbreak{}T\_GUIDE.\allowbreak{}md}, §2).

\textbf{The tools.} Six programs are needed. The versions are the ones that produced the committed files.

\begingroup
\small
\begin{longtable}{@{}>{\raggedright\arraybackslash}p{0.267\linewidth}>{\raggedright\arraybackslash}p{0.267\linewidth}>{\raggedright\arraybackslash}p{0.267\linewidth}@{}}
\toprule
\textbf{Tool} & \textbf{Version used} & \textbf{Needed for} \\
\midrule
\endfirsthead
\toprule
\textbf{Tool} & \textbf{Version used} & \textbf{Needed for} \\
\midrule
\endhead
Git & 2.51.2 & cloning the repository and the numerical engine; all stages \\
Python, with numpy, sympy, matplotlib, nbformat, nbclient, nbconvert, ipykernel & Python 3.14.5; numpy 2.4.6, sympy 1.14.0, matplotlib 3.11.0, nbformat 5.10.4, nbclient 0.10.2, nbconvert 7.16.6, ipykernel 7.1.0 & every stage and the textbook \\
WolframScript with the Wolfram Language (a Wolfram Engine or Mathematica) & WolframScript 1.14.0, Wolfram Language 15.0.1 & the exact Wolfram verifiers of Stages 1, 2, 4 and 5 and of the matter--antimatter analysis; optional in Stages 3 and 4 \\
A TeX distribution with pdflatex & MiKTeX 26.5 (pdfTeX 4.27) & every PDF \\
Rust with cargo, rustfmt and clippy & rustc and cargo 1.91.1 & Stages 3, 4 and 5 (the numerical programs) \\
PowerShell 7 (on Windows) & 7.6.6 & the \texttt{\detokenize{.ps1}} gates \\
\bottomrule
\end{longtable}
\endgroup

On macOS and Linux the Bash twins are used and PowerShell is not needed; there, TeX Live takes the place of MiKTeX. Python, Git, Rust and a TeX distribution are free programs; WolframScript needs an installed and activated Wolfram Engine or Mathematica. Without WolframScript the Stage-3 and Stage-4 gates skip their Wolfram steps and say so in their output; the Stage-1 and Stage-2 gates stop with a message that WolframScript was not found.

\textbf{Installing the Python packages.} The file \texttt{requirem\allowbreak{}ents-sta\allowbreak{}ge3.txt} in the repository root pins six of the packages; the seventh, nbconvert, is needed by the notebook steps of the Stage-3 and Stage-4 gates. In either shell, from the repository root:

\begin{Verbatim}[fontsize=\small]
python -m pip install -r requirements-stage3.txt
python -m pip install nbconvert==7.16.6
\end{Verbatim}

With other versions of the packages every check still passes, but the last digits of some files written by Python (for example the EXP-3 fits and the executed notebooks) can differ from the committed ones, and then the byte-for-byte comparisons of the Stage-3 gate fail (student guide, §3.5 and §14). On macOS and Linux, install the packages into a virtual environment, as the student guide explains in its §3.5.

\textbf{Checking the tools.} Each tool reports its version. In either shell:

\begin{Verbatim}[fontsize=\small]
git --version
python --version
python -c "import numpy, sympy; print(numpy.__version__, sympy.__version__)"
wolframscript -version
pdflatex --version
cargo --version
\end{Verbatim}

On the test computer these printed, in order, \texttt{git vers\allowbreak{}ion 2.51\allowbreak{}.2.windo\allowbreak{}ws.1}, \texttt{Python 3\allowbreak{}.14.5}, \texttt{2.4.6 1.\allowbreak{}14.0}, \texttt{WolframS\allowbreak{}cript 1.\allowbreak{}14.0 for\allowbreak{} Microso\allowbreak{}ft Windo\allowbreak{}ws (64-b\allowbreak{}it)}, a first line containing \texttt{MiKTeX-p\allowbreak{}dfTeX 4.\allowbreak{}27 (MiKT\allowbreak{}eX 26.5)} and \texttt{cargo 1.\allowbreak{}91.1 (ea\allowbreak{}2d97820 \allowbreak{}2025-10-\allowbreak{}10)}. Before its version, MiKTeX may print a reminder that you have not checked for updates; it does not affect anything in this appendix. In PowerShell, \texttt{\$PSVersi\allowbreak{}onTable.\allowbreak{}PSVersio\allowbreak{}n} shows the version of PowerShell itself (here 7.6.6).

\subsection{Two shells: PowerShell and Bash}

\textbf{What a shell is.} A shell is a program in which you type commands, one per line; pressing Enter runs the command. The shell always works in one folder, the \textbf{current folder}; a file named without a folder, such as \texttt{README.m\allowbreak{}d}, means the file of that name in the current folder, and a \textbf{relative path} such as \texttt{scripts/\allowbreak{}setup\_so\allowbreak{}lver.sh} means the file \texttt{setup\_so\allowbreak{}lver.sh} in the subfolder \texttt{\detokenize{scripts}} of the current folder. Every command of this appendix is to be typed in the \textbf{repository root}, the folder that contains \texttt{README.m\allowbreak{}d}. A program ends with an \textbf{exit code}, a whole number: 0 means success, anything else means failure. PowerShell shows the exit code of the last program when you type \texttt{\$LASTEXI\allowbreak{}TCODE}; Bash shows it when you type \texttt{\detokenize{echo $?}}.

\textbf{Six differences that matter here.}

\begin{enumerate}
\item \emph{Running a script.} In PowerShell a script in the current tree is run by its path, \texttt{.\textbackslash{}script\allowbreak{}s\textbackslash{}setup\_\allowbreak{}solver.p\allowbreak{}s1}, or, for the gates, by PowerShell 7 explicitly: \texttt{pwsh -No\allowbreak{}Profile \allowbreak{}-File sc\allowbreak{}ripts/ve\allowbreak{}rify\_sta\allowbreak{}ge2\_prim\allowbreak{}ordial\_f\allowbreak{}ield.ps1}. In Bash a script is given to the Bash program: \texttt{bash scr\allowbreak{}ipts/set\allowbreak{}up\_solve\allowbreak{}r.sh}.
\item \emph{Environment variables.} An environment variable is a named setting that programs started from the shell can read. PowerShell sets one with \texttt{\$env:PYT\allowbreak{}HONUTF8 \allowbreak{}= "1"}, Bash with \texttt{export P\allowbreak{}YTHONUTF\allowbreak{}8=1}. The setting lasts until the shell is closed. \texttt{PYTHONUT\allowbreak{}F8=1} makes Python read and write text as UTF-8, so that Greek letters in the output of the checkers survive; the gates set it themselves, and for single commands you set it once per shell.
\item \emph{Continuing a long command.} A command that does not fit on one line can be continued: in PowerShell the line ends with a backtick character (the grave accent) after a space, in Bash with a backslash. Every command box of this appendix whose lines end in one of these two characters is one single command.
\item \emph{Program names.} On Windows a compiled program ends in \texttt{\detokenize{.exe}}, and PowerShell runs a program whose path is stored in a variable with the call operator \texttt{\detokenize{&}}, as in \texttt{\& \$bin e\allowbreak{}xp5}. In Bash the \texttt{\detokenize{.exe}} may be left out in Git Bash and does not exist on macOS and Linux.
\item \emph{Paths.} Python, cargo and WolframScript accept forward slashes on every system, so the paths inside the commands are written with \texttt{\detokenize{/}} in both shells.
\item \emph{Windows PowerShell 5.1.} Windows also contains an older PowerShell, version 5.1, started by the command \texttt{powershe\allowbreak{}ll}. The gates need PowerShell 7 (the command \texttt{\detokenize{pwsh}}); started from version 5.1, each gate starts itself again under PowerShell 7 and passes on its exit code, or stops with a message if PowerShell 7 is not installed (header of each \texttt{\detokenize{.ps1}} gate).
\end{enumerate}

\textbf{Worked example.} Setting the environment variable and counting the checks of a report, first in PowerShell:

\begin{Verbatim}[fontsize=\small]
$env:PYTHONUTF8 = "1"
python -c "
import json, sys
c = json.load(open(sys.argv[1], encoding='utf-8'))['checks']
print(sum(c.values()), len(c))
" artifacts/dirac16complex/pair-creation/wolfram-pairing-report.json
\end{Verbatim}

The lines from \texttt{python -\allowbreak{}c "} to the closing \texttt{\detokenize{"}} and the path after it form one command: the three lines between the double quotes are a small Python program, which Python receives as one piece of text, and the path after the closing quote is handed to that program as \texttt{sys.argv\allowbreak{}[1]} (the first word after the program). Both shells allow a quoted text to continue over several lines, so in Bash only the first line differs:

\begin{Verbatim}[fontsize=\small]
export PYTHONUTF8=1
\end{Verbatim}

Both shells print the one line \texttt{\detokenize{141 141}}: all 141 checks of the Stage-5 pairing report are true. This is the method of Section 0.8 written as one command. Replace the path to count the checks of any other report whose checks are stored in a top-level entry \texttt{\detokenize{checks}}. The EXP-3 fit analysis \texttt{artifact\allowbreak{}s/dirac1\allowbreak{}6complex\allowbreak{}/numeric\allowbreak{}s/exp3/f\allowbreak{}its.json} is an exception: it stores its checks one level down, in its entry \texttt{validati\allowbreak{}on}, and for it the command stops with the error message \texttt{KeyError\allowbreak{}: 'check\allowbreak{}s'}.

\subsection{A fresh clone and the numerical engine}

\textbf{Git, repositories and commits.} \textbf{Git} is a program that records the history of a folder of files. The folder together with its whole recorded history is a \textbf{repository}. The history is a sequence of \textbf{commits}: a commit is a recorded state of all the files of the repository, together with a short message, and it is named by a 40-digit hexadecimal fingerprint of its content; the first seven digits, such as \texttt{\detokenize{1f2dd69}}, are enough to name a commit of this repository. The files as they lie on your disk are the \textbf{working tree}; the command \texttt{git stat\allowbreak{}us --sho\allowbreak{}rt} lists every file of the working tree that is changed or new compared with the last commit (files in folders that Git is told to ignore excepted), and prints nothing when there is none. A \textbf{clone} is a complete copy of a repository, with its whole history, made by Git; a \textbf{fresh} clone is one made just before a test, so that no file of an earlier run and no local change can enter the test.

\textbf{Cloning.} Choose a folder for the clone. On Windows, choose a folder with a short path, such as \texttt{C:\textbackslash{}work}: by default a Windows program cannot open a file whose full path is longer than 260 characters, and some files of the repository lie more than 100 characters deep inside the clone. For this appendix the first clones were made 168 characters deep, and they failed in two ways: the Rust linker stopped with the error \texttt{\detokenize{LNK1104}} (cannot open file), and WolframScript did not see two Stage-4 run files whose paths then had 271 and 272 characters, so that one report counted them wrong (Section 19.10). Then type, in either shell:

\begin{Verbatim}[fontsize=\small]
git clone https://github.com/once-ere/Dirac_claude.git
cd Dirac_claude
git log --oneline -1
\end{Verbatim}

The first command downloads the repository into the new folder \texttt{Dirac\_cl\allowbreak{}aude} (about 440 megabytes; 11 seconds on the test computer), the second makes that folder the current folder, and the third prints the commit you have, its short hexadecimal name followed by its message. All later commands are typed in this folder. To work with the files of the commit at which this appendix was measured, type

\begin{Verbatim}[fontsize=\small]
git checkout 1f2dd69
\end{Verbatim}

Git then reports that you are in "detached HEAD" state, which only means that the working tree shows an older commit than the newest one (\texttt{\detokenize{HEAD}} is Git's name for the commit you have); \texttt{git chec\allowbreak{}kout mai\allowbreak{}n} returns to the newest commit. Some folders are deliberately missing from a clone: \texttt{\detokenize{build/}} (scratch output), \texttt{\detokenize{vendor/}} (the numerical engine, fetched below) and the private reference folder \texttt{dirac-ma\allowbreak{}in/} of Stage 1 (see Section 19.6). The file \texttt{.gitigno\allowbreak{}re} lists them, and Git never records their contents.

\textbf{The numerical engine.} The Rust programs of Stages 3, 4 and 5 use CVODE, the solver of the pure-Rust translation of SUNDIALS 7.8.0 (Chapter 10). It is not stored in this repository; a setup script downloads it at a fixed commit (a \textbf{pin}) into \texttt{vendor/r\allowbreak{}ustSolve\allowbreak{}It}. PowerShell:

\begin{Verbatim}[fontsize=\small]
.\scripts\setup_solver.ps1 -Platform win11
\end{Verbatim}

Bash (Git Bash, macOS, Linux):

\begin{Verbatim}[fontsize=\small]
bash scripts/setup_solver.sh win11
\end{Verbatim}

Expected output of the first run (6 seconds on the test computer):

\begin{Verbatim}[fontsize=\small]
solver_platform=win11
solver_commit=a8fdff459adfe181573d7924b18bffbdf378fdb3
solver_setup=OK
\end{Verbatim}

A second run finds the engine and prints \texttt{solver\_s\allowbreak{}etup=ALR\allowbreak{}EADY-PRE\allowbreak{}SENT} instead of \texttt{solver\_s\allowbreak{}etup=OK}. If the folder holds a different commit, the script refuses to continue and asks you to remove \texttt{vendor/r\allowbreak{}ustSolve\allowbreak{}It} and run it again; an interrupted download is reported in the same way. If PowerShell refuses to run the script with the message that running scripts is disabled on this system, run it as \texttt{pwsh -No\allowbreak{}Profile \allowbreak{}-Executi\allowbreak{}onPolicy\allowbreak{} Bypass \allowbreak{}-File sc\allowbreak{}ripts/se\allowbreak{}tup\_solv\allowbreak{}er.ps1 -\allowbreak{}Platform\allowbreak{} win11}; the option \texttt{-Executi\allowbreak{}onPolicy\allowbreak{} Bypass} applies to this one command only.

\textbf{Which engine.} There are three engines, \texttt{\detokenize{win11}}, \texttt{\detokenize{macos}} and \texttt{\detokenize{linux}}, in three repositories (Section 10.9 lists their pinned commits). The committed numerical files were produced with the \texttt{\detokenize{win11}} engine, and that engine reproduces them byte for byte on Windows and on Linux; the \texttt{\detokenize{macos}} and \texttt{\detokenize{linux}} engines contain a different mathematical library, so with them all checks pass but some files differ in their last digits (Section 10.10; \texttt{README.m\allowbreak{}d}, section "Reproducing"). This is why every command of this appendix, and every gate, uses \texttt{\detokenize{win11}}, also on macOS and Linux. Without an argument the Bash script chooses the engine from the operating system, which is not what the byte comparisons need.

\subsection{Reading what a verifier prints, and rerunning one without touching committed files}

\textbf{Verifiers and checkers.} A \textbf{verifier} in this book is a program that recomputes exact statements and writes a report; the Wolfram Language ones are called verifiers (file names \texttt{scripts/\allowbreak{}verify\_*\allowbreak{}.wls}), the independent Python ones checkers (\texttt{scripts/\allowbreak{}check\_*.\allowbreak{}py}). Every one of them prints a line for each check and for each measurement (some also print progress lines with the elapsed time) and ends with two summary lines:

\begin{Verbatim}[fontsize=\small]
check_ALG_clifford=true
check_ALG_gammaTransposeSymmetry=true
...
measurement_ALG_cliffordPairsVerified=64
measurement_ALG_chargeMatrixSignature=[8,8,0]
...
check_count=20
failed_check_count=0
\end{Verbatim}

(This is the Python algebra checker of Stage 1; \texttt{\detokenize{...}} marks omitted lines in every output box of this appendix.) A check is a statement that is either true or false; a \textbf{measurement} is a recorded value, such as a count or a matrix; here 64 ordered pairs of gamma matrices were tested and the charge matrix has 8 positive, 8 negative and 0 zero eigenvalues. The program exits with code 0 when every check is true and with a nonzero code otherwise (some verifiers also refuse, with a nonzero code, when the number of checks differs from the number they expect). The same checks and measurements are written into the report, a JSON file whose entry \texttt{\detokenize{checks}} holds the true or false values (Section 0.8).

\textbf{What a gate prints.} A gate prints the tools it found (lines \texttt{stageN\_t\allowbreak{}ool\_...=}), then for each step a line \texttt{stageN\_s\allowbreak{}tep=<nam\allowbreak{}e>} when the step starts and \texttt{stageN\_s\allowbreak{}tep\_ok=<\allowbreak{}name>} when it has succeeded, with the full output of the step in between. After the steps it audits the reports and prints their check counts and the sha256 of every output file (Section 19.1). The last line is the verdict. When a step fails, the gate stops and prints instead, for example in Stage 2,

\begin{Verbatim}[fontsize=\small]
stage2_failed_step=stage2-02-check-primordial
stage2_failed_log=build/logs/stage2-02-check-primordial-bash.log
stage2_primordial_field_verification=FAILED
\end{Verbatim}

and exits with a nonzero code; the other gates print the same three lines with their own stage number and name. Every step writes its complete output into a \textbf{log file} in \texttt{build/lo\allowbreak{}gs/} (the Bash twins add \texttt{\detokenize{-bash}} to the file name so that the logs of the two twins do not overwrite each other); the log begins with the start time and the command and ends with the finish time and the exit code.

\textbf{Rerunning one program into a scratch folder.} Many verifiers write their report into the committed folder by default. To check a single statement without changing any committed file, write the report into the folder \texttt{\detokenize{build/}}, which Git ignores. Two rules make this work:

\begin{enumerate}
\item The Wolfram verifiers take the report path as a plain argument after the script name, never after a separator \texttt{\detokenize{--}}: with WolframScript 1.14, \texttt{wolframs\allowbreak{}cript -f\allowbreak{}ile s.wl\allowbreak{}s -- r.j\allowbreak{}son} loses the path, while \texttt{wolframs\allowbreak{}cript -f\allowbreak{}ile s.wl\allowbreak{}s r.json} passes it (the headers of all gates record this test). Some Wolfram verifiers also write a theory file next to the report; with a report path in \texttt{\detokenize{build/}} that file goes to \texttt{\detokenize{build/}} as well.
\item The Python checkers take the report path from an option, \texttt{\detokenize{--output}} for most of them. Two Stage-1 programs write to a fixed path; for them the helper \texttt{scripts/\allowbreak{}run\_with\allowbreak{}\_report\_\allowbreak{}path.py} redirects the report (Section 19.6).
\end{enumerate}

After any rerun, \texttt{git stat\allowbreak{}us --sho\allowbreak{}rt} lists every committed file that has changed (Section 19.4); it prints nothing when all committed files are unchanged. To put the committed version of the artifacts back, type \texttt{git rest\allowbreak{}ore arti\allowbreak{}facts}.

\textbf{Worked example: one Stage-1 checker.} PowerShell:

\begin{Verbatim}[fontsize=\small]
$env:PYTHONUTF8 = "1"
python scripts/check_dirac16complex_algebra.py --wolfram-report= `
    --output build/textbook/python-algebra-report.json
\end{Verbatim}

Bash:

\begin{Verbatim}[fontsize=\small]
export PYTHONUTF8=1
python scripts/check_dirac16complex_algebra.py --wolfram-report= \
    --output build/textbook/python-algebra-report.json
\end{Verbatim}

Both print 20 check lines and end with \texttt{check\_co\allowbreak{}unt=20}, \texttt{failed\_c\allowbreak{}heck\_cou\allowbreak{}nt=0} and the line \texttt{report=b\allowbreak{}uild/tex\allowbreak{}tbook/py\allowbreak{}thon-alg\allowbreak{}ebra-rep\allowbreak{}ort.json}; the committed report has 21 checks, because the empty option \texttt{--wolfra\allowbreak{}m-report\allowbreak{}=} switches off the comparison with the Wolfram report (Section 2.16). On the test computer the run took 4 seconds in Bash and 1 second in PowerShell, and \texttt{git stat\allowbreak{}us --sho\allowbreak{}rt} printed nothing afterwards.

\subsection{Stage 1: the field in an arbitrary gravitational field}

\textbf{What the gate does.} The Stage-1 gate runs thirteen steps (header of \texttt{scripts/\allowbreak{}verify\_s\allowbreak{}tage1\_ar\allowbreak{}bitrary\_\allowbreak{}field.sh}); each writes its log into \texttt{build/lo\allowbreak{}gs/}. The times in the table are those of the run described below, read from the start and finish times in the logs:

\begingroup
\small
\begin{longtable}{@{}>{\raggedright\arraybackslash}p{0.267\linewidth}>{\raggedright\arraybackslash}p{0.267\linewidth}>{\raggedright\arraybackslash}p{0.267\linewidth}@{}}
\toprule
\textbf{Step} & \textbf{Command (from the log of the run described below)} & \textbf{Time} \\
\midrule
\endfirsthead
\toprule
\textbf{Step} & \textbf{Command (from the log of the run described below)} & \textbf{Time} \\
\midrule
\endhead
\texttt{stage1-0\allowbreak{}1-build-\allowbreak{}fixture} & \texttt{python s\allowbreak{}cripts/b\allowbreak{}uild\_dir\allowbreak{}ac16comp\allowbreak{}lex\_fixt\allowbreak{}ure.py -\allowbreak{}-output \allowbreak{}build/st\allowbreak{}age1/alg\allowbreak{}ebra-fix\allowbreak{}ture.jso\allowbreak{}n} & under 1 s \\
\texttt{stage1-0\allowbreak{}2-check-\allowbreak{}algebra} & \texttt{python s\allowbreak{}cripts/c\allowbreak{}heck\_dir\allowbreak{}ac16comp\allowbreak{}lex\_alge\allowbreak{}bra.py -\allowbreak{}-wolfram\allowbreak{}-report=\allowbreak{} --outpu\allowbreak{}t build/\allowbreak{}stage1/p\allowbreak{}ython-al\allowbreak{}gebra-re\allowbreak{}port.jso\allowbreak{}n} & 3 s \\
\texttt{stage1-0\allowbreak{}3-wolfra\allowbreak{}m-algebr\allowbreak{}a} & \texttt{wolframs\allowbreak{}cript -f\allowbreak{}ile scri\allowbreak{}pts/veri\allowbreak{}fy\_dirac\allowbreak{}16comple\allowbreak{}x\_algebr\allowbreak{}a.wls bu\allowbreak{}ild/stag\allowbreak{}e1/wolfr\allowbreak{}am-algeb\allowbreak{}ra-repor\allowbreak{}t.json} & 23 s \\
\texttt{stage1-0\allowbreak{}4-wolfra\allowbreak{}m-geomet\allowbreak{}ry} & \texttt{wolframs\allowbreak{}cript -f\allowbreak{}ile scri\allowbreak{}pts/veri\allowbreak{}fy\_dirac\allowbreak{}16comple\allowbreak{}x\_geomet\allowbreak{}ry.wls b\allowbreak{}uild/sta\allowbreak{}ge1/wolf\allowbreak{}ram-geom\allowbreak{}etry-rep\allowbreak{}ort.json} & 409 s \\
\texttt{stage1-0\allowbreak{}5-check-\allowbreak{}geometry} & \texttt{python s\allowbreak{}cripts/r\allowbreak{}un\_with\_\allowbreak{}report\_p\allowbreak{}ath.py b\allowbreak{}uild/sta\allowbreak{}ge1/pyth\allowbreak{}on-geome\allowbreak{}try-repo\allowbreak{}rt.json \allowbreak{}scripts/\allowbreak{}check\_di\allowbreak{}rac16com\allowbreak{}plex\_geo\allowbreak{}metry.py\allowbreak{} --wolfr\allowbreak{}am-repor\allowbreak{}t build/\allowbreak{}stage1/w\allowbreak{}olfram-g\allowbreak{}eometry-\allowbreak{}report.j\allowbreak{}son} & 405 s \\
\texttt{stage1-0\allowbreak{}6-grassm\allowbreak{}ann-demo} & \texttt{python s\allowbreak{}cripts/r\allowbreak{}un\_with\_\allowbreak{}report\_p\allowbreak{}ath.py b\allowbreak{}uild/sta\allowbreak{}ge1/gras\allowbreak{}smann-de\allowbreak{}mo-repor\allowbreak{}t.json s\allowbreak{}cripts/d\allowbreak{}emo\_gras\allowbreak{}smann\_la\allowbreak{}grangian\allowbreak{}s.py} & 57 s \\
\texttt{stage1-0\allowbreak{}7-check-\allowbreak{}algebra-\allowbreak{}crossche\allowbreak{}ck} & \texttt{python s\allowbreak{}cripts/c\allowbreak{}heck\_dir\allowbreak{}ac16comp\allowbreak{}lex\_alge\allowbreak{}bra.py -\allowbreak{}-wolfram\allowbreak{}-report \allowbreak{}build/st\allowbreak{}age1/wol\allowbreak{}fram-alg\allowbreak{}ebra-rep\allowbreak{}ort.json\allowbreak{} --outpu\allowbreak{}t build/\allowbreak{}stage1/p\allowbreak{}ython-al\allowbreak{}gebra-re\allowbreak{}port.jso\allowbreak{}n} & 2 s \\
\texttt{stage1-0\allowbreak{}8-python\allowbreak{}-tests} & \texttt{python -\allowbreak{}m unitte\allowbreak{}st disco\allowbreak{}ver -s t\allowbreak{}ests -p \allowbreak{}"test\_d1\allowbreak{}6c\_[ag]*\allowbreak{}.py" -v} (99 tests) & 54 s \\
\texttt{stage1-0\allowbreak{}9-public\allowbreak{}ation-te\allowbreak{}sts} & \texttt{python -\allowbreak{}m unitte\allowbreak{}st disco\allowbreak{}ver -s t\allowbreak{}ests -p \allowbreak{}test\_pub\allowbreak{}lication\allowbreak{}\_tooling\allowbreak{}.py -v} (68 tests) & 94 s \\
\texttt{stage1-1\allowbreak{}0-summar\allowbreak{}y} & \texttt{python s\allowbreak{}cripts/b\allowbreak{}uild\_sta\allowbreak{}ge1\_summ\allowbreak{}ary.py -\allowbreak{}-output \allowbreak{}build/st\allowbreak{}age1/sta\allowbreak{}ge1-summ\allowbreak{}ary.json\allowbreak{} --repor\allowbreak{}ts-direc\allowbreak{}tory bui\allowbreak{}ld/stage\allowbreak{}1} & under 1 s \\
\texttt{stage1-1\allowbreak{}1-proven\allowbreak{}ance-pdf} & \texttt{python s\allowbreak{}cripts/b\allowbreak{}uild\_pro\allowbreak{}venance\_\allowbreak{}pdf.py p\allowbreak{}rovenanc\allowbreak{}e/DIRAC1\allowbreak{}6COMPLEX\allowbreak{}\_ARBITRA\allowbreak{}RY\_FIELD\allowbreak{}.md} & 16 s \\
\texttt{stage1-1\allowbreak{}2-public\allowbreak{}ation-re\allowbreak{}check} & \texttt{python -\allowbreak{}m unitte\allowbreak{}st disco\allowbreak{}ver -s t\allowbreak{}ests -p \allowbreak{}test\_d16\allowbreak{}c\_arbitr\allowbreak{}ary\_fiel\allowbreak{}d\_public\allowbreak{}ation.py\allowbreak{} -v} (27 tests) & 1 s \\
\texttt{stage1-1\allowbreak{}3-public\allowbreak{}-clone-a\allowbreak{}udit} & \texttt{python s\allowbreak{}cripts/v\allowbreak{}erify\_st\allowbreak{}age1\_pub\allowbreak{}lic\_clon\allowbreak{}e\_audit.\allowbreak{}py --com\allowbreak{}mitted a\allowbreak{}rtifacts\allowbreak{}/dirac16\allowbreak{}complex/\allowbreak{}arbitrar\allowbreak{}y-field \allowbreak{}--regene\allowbreak{}rated bu\allowbreak{}ild/stag\allowbreak{}e1} & 1 s \\
\bottomrule
\end{longtable}
\endgroup

Steps 02 and 07 run the same Python checker twice: first alone, then again with the comparison against the Wolfram algebra report that step 03 has just written (check \texttt{ALG\_wolf\allowbreak{}ramAgree\allowbreak{}ment}). Step 05 compares the Python geometry with the Wolfram geometry of step 04 (check \texttt{GEO\_wolf\allowbreak{}ramAgree\allowbreak{}ment}). The helper \texttt{scripts/\allowbreak{}run\_with\allowbreak{}\_report\_\allowbreak{}path.py} is needed because the geometry checker and the Grassmann demonstration write their reports to a fixed path and have no output option; the helper runs the committed code unchanged and only redirects the report (its header explains why no option was added: the sha256 of the two scripts is recorded in the committed reports). Steps 08 and 09 each report one skipped test, the two tests that need the folder \texttt{dirac-ma\allowbreak{}in/} described next.

\textbf{Two modes: with and without dirac-main.} Three exact data files of the separately published reference implementation dirac-main (https://github.com/once-ere/dirac) are compared with the Stage-1 matrices (Chapter 3). The folder \texttt{dirac-ma\allowbreak{}in/} that would hold them is not part of this repository, so a fresh clone lacks it. The gate first prints \texttt{stage1\_d\allowbreak{}irac\_mai\allowbreak{}n=presen\allowbreak{}t} or \texttt{stage1\_d\allowbreak{}irac\_mai\allowbreak{}n=absent}. With \texttt{\detokenize{absent}} it runs in the \textbf{public-clone mode}: every regenerated file goes into \texttt{build/st\allowbreak{}age1/} and the committed reports stay untouched; the verifiers record the dirac-main comparisons as "not-run"; and the extra step 13 requires every regenerated file to equal the committed one, byte for byte or value for value, except for exactly the entries that the missing folder changes (its log lists each allowed difference in a line \texttt{stage1\_a\allowbreak{}udit\_all\allowbreak{}owed=}). With \texttt{\detokenize{present}} the gate rewrites the committed reports in place, runs the dirac-main comparisons as well, and needs no step 13. Both modes were run from fresh public clones earlier on 2026-09-30 and printed the OK line: the Bash and PowerShell twins without dirac-main, and the Bash twin with a copy of the folder (\texttt{handoff/\allowbreak{}reviews/\allowbreak{}stage1\_g\allowbreak{}ate\_2026\allowbreak{}-09-30\_p\allowbreak{}ublic\_sh\allowbreak{}.log}, \texttt{...\_publ\allowbreak{}ic\_ps1.l\allowbreak{}og} and \texttt{...\_priv\allowbreak{}ate\_sh.l\allowbreak{}og}).

\textbf{Running it.} PowerShell:

\begin{Verbatim}[fontsize=\small]
pwsh -NoProfile -File scripts/verify_stage1_arbitrary_field.ps1
\end{Verbatim}

Bash:

\begin{Verbatim}[fontsize=\small]
bash scripts/verify_stage1_arbitrary_field.sh
\end{Verbatim}

The last lines of a successful run in the public-clone mode (sha256 values cut after 16 digits):

\begin{Verbatim}[fontsize=\small]
stage1_report_checks=build/stage1/wolfram-algebra-report.json 21 true
stage1_report_checks=build/stage1/wolfram-geometry-report.json 43 true
stage1_report_checks=build/stage1/python-algebra-report.json 21 true
stage1_report_checks=build/stage1/python-geometry-report.json 52 true
stage1_report_checks=build/stage1/grassmann-demo-report.json 16 true
stage1_sha256=39cc390b53266bda...  build/stage1/wolfram-algebra-report.json
...
stage1_sha256=83252f8478aa1f05...  provenance/DIRAC16COMPLEX_ARBITRARY_FIELD.pdf
stage1_skipped=dirac-main cross-checks (dirac-main/ is a git-ignored reference ...
stage1_arbitrary_field_verification=OK
\end{Verbatim}

The eight sha256 lines cover the five reports, \texttt{stage1-s\allowbreak{}ummary.j\allowbreak{}son} and the \texttt{\detokenize{.tex}} and \texttt{\detokenize{.pdf}} of the Stage-1 document; the \texttt{stage1\_s\allowbreak{}kipped} line is longer on the screen and names the log in which the audit lists the comparisons that were not run. The five reports hold $21+43+21+52+16=153$ checks, the 153 of 153 of Chapter 0. Five of the eight fingerprints equal those of the committed files: the Wolfram geometry report, the Python geometry report, the Grassmann report, and the \texttt{\detokenize{.tex}} and the \texttt{\detokenize{.pdf}} of the document, which the gate rebuilds byte for byte. Three differ, and that is expected: the Wolfram algebra report (\texttt{39cc390b\allowbreak{}...} against the committed \texttt{d43adeeb\allowbreak{}...}) and the Python algebra report (\texttt{2be5b910\allowbreak{}...} against \texttt{3275b2af\allowbreak{}...}) record the dirac-main comparisons as "not-run", and \texttt{stage1-s\allowbreak{}ummary.j\allowbreak{}son} (\texttt{03165f0f\allowbreak{}...} against \texttt{a65972b5\allowbreak{}...}) collects their entries. These are exactly the differences that step 13 allows (its log lists them).

\textbf{Wall time.} The Bash twin, run for this edition from a fresh clone (public-clone mode), took 1095 seconds, about 18 minutes; the table above gives the time of each step. The two exact geometry programs dominate. In the committed runs of 2026-09-30 the Python geometry checker took 190 seconds (Bash) and 227 seconds (PowerShell) without dirac-main and 621 seconds with it (the line \texttt{runtime\_\allowbreak{}seconds=} of each log), so the whole gate needs roughly 10 to 20 minutes, depending on the load of the computer.

\textbf{Single programs.} Each row of the table is an ordinary command and can be typed on its own, in either shell (after setting \texttt{PYTHONUT\allowbreak{}F8} as in Section 19.3); the gate's folder \texttt{build/st\allowbreak{}age1/} can be replaced by any folder below \texttt{\detokenize{build/}}. The Wolfram verifiers print progress lines, their checks and measurements, and end with \texttt{check\_co\allowbreak{}unt=21} (algebra) and \texttt{check\_co\allowbreak{}unt=43} (geometry); the Python geometry checker ends with \texttt{check\_co\allowbreak{}unt=52}, the Grassmann demonstration with \texttt{check\_co\allowbreak{}unt=16}, all with \texttt{failed\_c\allowbreak{}heck\_cou\allowbreak{}nt=0}. Run step 04 before step 05 and step 03 before step 07, because the later steps read the reports of the earlier ones.

\subsection{Stage 2: the primordial field}

\textbf{What the gate does.} The Stage-2 gate has four steps (header of \texttt{scripts/\allowbreak{}verify\_s\allowbreak{}tage2\_pr\allowbreak{}imordial\allowbreak{}\_field.s\allowbreak{}h}):

\begingroup
\small
\begin{longtable}{@{}>{\raggedright\arraybackslash}p{0.267\linewidth}>{\raggedright\arraybackslash}p{0.267\linewidth}>{\raggedright\arraybackslash}p{0.267\linewidth}@{}}
\toprule
\textbf{Step} & \textbf{Program} & \textbf{What it does} \\
\midrule
\endfirsthead
\toprule
\textbf{Step} & \textbf{Program} & \textbf{What it does} \\
\midrule
\endhead
\texttt{stage2-0\allowbreak{}1-wolfra\allowbreak{}m-primor\allowbreak{}dial} & \texttt{scripts/\allowbreak{}verify\_d\allowbreak{}irac16co\allowbreak{}mplex\_pr\allowbreak{}imordial\allowbreak{}.wls} & the exact Wolfram verifier of Chapter 9; writes \texttt{wolfram-\allowbreak{}primordi\allowbreak{}al-repor\allowbreak{}t.json} and \texttt{primordi\allowbreak{}al-compo\allowbreak{}nents.js\allowbreak{}on} \\
\texttt{stage2-0\allowbreak{}2-check-\allowbreak{}primordi\allowbreak{}al} & \texttt{scripts/\allowbreak{}check\_di\allowbreak{}rac16com\allowbreak{}plex\_pri\allowbreak{}mordial.\allowbreak{}py} & the independent sympy checker; compares its field equations with the components written by step 1 \\
\texttt{stage2-0\allowbreak{}3-python\allowbreak{}-tests} & \texttt{python -\allowbreak{}m unitte\allowbreak{}st} with the files \texttt{tests/te\allowbreak{}st\_d16c\_\allowbreak{}primordi\allowbreak{}al*.py} & the unit tests of Stage 2 \\
\texttt{stage2-0\allowbreak{}4-proven\allowbreak{}ance-pdf} & \texttt{scripts/\allowbreak{}build\_pr\allowbreak{}ovenance\allowbreak{}\_pdf.py} & rebuilds \texttt{provenan\allowbreak{}ce/DIRAC\allowbreak{}16COMPLE\allowbreak{}X\_PRIMOR\allowbreak{}DIAL\_FIE\allowbreak{}LD.pdf} and compares it with the registered edition \\
\bottomrule
\end{longtable}
\endgroup

The gate writes its reports over the committed ones in \texttt{artifact\allowbreak{}s/dirac1\allowbreak{}6complex\allowbreak{}/primord\allowbreak{}ial-fiel\allowbreak{}d/}; when everything agrees the new files are byte-identical, so \texttt{git stat\allowbreak{}us --sho\allowbreak{}rt} shows nothing afterwards. It then requires every check to be true and requires the Python report to have compared its equations with the components of this very run.

\textbf{Running it.} PowerShell:

\begin{Verbatim}[fontsize=\small]
pwsh -NoProfile -File scripts/verify_stage2_primordial_field.ps1
\end{Verbatim}

Bash:

\begin{Verbatim}[fontsize=\small]
bash scripts/verify_stage2_primordial_field.sh
\end{Verbatim}

The last lines of a successful run (the folder \texttt{artifact\allowbreak{}s/dirac1\allowbreak{}6complex\allowbreak{}/primord\allowbreak{}ial-fiel\allowbreak{}d/} is shortened to \texttt{\detokenize{.../}} here, and each sha256 line ends with the file name):

\begin{Verbatim}[fontsize=\small]
stage2_step_ok=stage2-04-provenance-pdf
stage2_report_checks=.../wolfram-primordial-report.json 126 true
stage2_report_checks=.../python-primordial-report.json 16 true, wolframAgreement=compared
stage2_sha256=a0164273df62f2e1...  .../wolfram-primordial-report.json
stage2_sha256=a5d8caa8cb226082...  .../primordial-components.json
stage2_sha256=2d9cc593bd7812be...  .../python-primordial-report.json
stage2_sha256=c1cd3773ef2f15c4...  provenance/DIRAC16COMPLEX_PRIMORDIAL_FIELD.tex
stage2_sha256=a6f6d7b813591412...  provenance/DIRAC16COMPLEX_PRIMORDIAL_FIELD.pdf
stage2_primordial_field_verification=OK
\end{Verbatim}

(the sha256 values are cut after 16 digits here; the gate prints all 64).

\textbf{Wall time.} The Bash twin, run for this edition from a fresh clone, took 165 seconds: 93 for the Wolfram verifier, 16 for the Python checker, 31 for the 53 unit tests and 13 for the PDF, the rest for the audits. The committed logs of the earlier fresh-clone runs (\texttt{handoff/\allowbreak{}reviews/\allowbreak{}stage2\_g\allowbreak{}ate\_2026\allowbreak{}-09-30\_f\allowbreak{}resh\_clo\allowbreak{}ne\_OK.lo\allowbreak{}g} for Bash and \texttt{stage2\_g\allowbreak{}ate\_ps1\_\allowbreak{}2026-09-\allowbreak{}30\_fresh\allowbreak{}\_clone\_O\allowbreak{}K.log} for PowerShell, both containing the OK line) record 88 and 121 seconds for the Wolfram step and 32 seconds for the unit tests.

\textbf{Single programs.} To rerun the two verifiers into the scratch folder \texttt{build/te\allowbreak{}xtbook/s\allowbreak{}tage2/}, in PowerShell:

\begin{Verbatim}[fontsize=\small]
$env:PYTHONUTF8 = "1"
$s = "build/textbook/stage2"
wolframscript -file scripts/verify_dirac16complex_primordial.wls `
    "$s/wolfram-primordial-report.json"
python scripts/check_dirac16complex_primordial.py `
    --output "$s/python-primordial-report.json" `
    --wolfram-components "$s/primordial-components.json"
\end{Verbatim}

and in Bash:

\begin{Verbatim}[fontsize=\small]
export PYTHONUTF8=1
s=build/textbook/stage2
wolframscript -file scripts/verify_dirac16complex_primordial.wls \
    "$s/wolfram-primordial-report.json"
python scripts/check_dirac16complex_primordial.py \
    --output "$s/python-primordial-report.json" \
    --wolfram-components "$s/primordial-components.json"
\end{Verbatim}

The variable \texttt{\detokenize{s}} only abbreviates the folder. The Wolfram verifier first prints one progress line with the clock time for each group of checks, from \texttt{zero-tes\allowbreak{}t sanity} to a line \texttt{done in \allowbreak{}... s}; then its checks and measurements; then \texttt{check\_co\allowbreak{}unt=126}, \texttt{failed\_c\allowbreak{}heck\_cou\allowbreak{}nt=0} and \texttt{elapsed\_\allowbreak{}seconds=}; and last two lines \texttt{\detokenize{report=}} and \texttt{componen\allowbreak{}ts=} with the full paths of the report and of \texttt{primordi\allowbreak{}al-compo\allowbreak{}nents.js\allowbreak{}on}, which it writes next to the report. It took 82 seconds in Bash and 46 seconds in PowerShell. The Python checker ends with \texttt{check\_co\allowbreak{}unt=16}, \texttt{failed\_c\allowbreak{}heck\_cou\allowbreak{}nt=0}, sixteen lines \texttt{timing\_.\allowbreak{}..=} and the line \texttt{\detokenize{report=}}, after 6 and 4 seconds. The two Wolfram files were byte-identical to the committed ones. The Python report differs from the committed one in a single entry, as it should: its entry \texttt{inputSha\allowbreak{}256} records the sha256 of the components file it compared with under that file's path, here \texttt{build/te\allowbreak{}xtbook/s\allowbreak{}tage2/pr\allowbreak{}imordial\allowbreak{}-compone\allowbreak{}nts.json} instead of the committed path; the recorded sha256 itself is the same.

\subsection{Stage 3: the dark-sector experiments}

\textbf{What the gate does.} The Stage-3 gate runs 33 steps, numbered 00 to 32 (header of \texttt{scripts/\allowbreak{}verify\_s\allowbreak{}tage3\_da\allowbreak{}rk\_secto\allowbreak{}r.sh}). In groups:

\begingroup
\small
\begin{longtable}{@{}>{\raggedright\arraybackslash}p{0.400\linewidth}>{\raggedright\arraybackslash}p{0.400\linewidth}@{}}
\toprule
\textbf{Steps} & \textbf{What they do} \\
\midrule
\endfirsthead
\toprule
\textbf{Steps} & \textbf{What they do} \\
\midrule
\endhead
00 & copy the committed Stage-3 files into \texttt{build/st\allowbreak{}age3/sna\allowbreak{}pshot}, to prove at the end that they did not change \\
01, 02 & install the \texttt{\detokenize{win11}} engine (Section 19.4) and check that it is exactly the pinned commit, unmodified \\
03 to 06 & \texttt{cargo fm\allowbreak{}t --chec\allowbreak{}k}, \texttt{cargo cl\allowbreak{}ippy} with every warning treated as an error, \texttt{cargo te\allowbreak{}st --rel\allowbreak{}ease}, \texttt{cargo bu\allowbreak{}ild --re\allowbreak{}lease} for the crate \texttt{studies/\allowbreak{}dirac16c\allowbreak{}omplex\_c\allowbreak{}osmology} \\
07, 08 & run all five experiments twice, into \texttt{build/st\allowbreak{}age3/run\allowbreak{}-a} and \texttt{build/st\allowbreak{}age3/run\allowbreak{}-b}; each run must end with \texttt{\detokenize{SUCCESS}} \\
09 & the 62 output files of the experiments must be byte-identical in both runs and in the committed folders \texttt{artifact\allowbreak{}s/dirac1\allowbreak{}6complex\allowbreak{}/numeric\allowbreak{}s/exp1/} to \texttt{\detokenize{exp5/}} \\
10, 11 & the EXP-3 fit analysis \texttt{scripts/\allowbreak{}analyze\_\allowbreak{}dirac16c\allowbreak{}omplex\_e\allowbreak{}xp3.py}, compared with the committed fit files \\
12 to 17 & the five independent checkers \texttt{scripts/\allowbreak{}check\_di\allowbreak{}rac16com\allowbreak{}plex\_exp\allowbreak{}1.py} to \texttt{...\_exp5\allowbreak{}.py}, each with a repeat run and a run with tightened tolerances, and an audit of their reports \\
18, 19 & the collected summary \texttt{numerics\allowbreak{}-summary\allowbreak{}.json}, compared with the committed one \\
20 to 25 & the Jupyter notebook, executed twice without any stored output, audited, and required to rewrite the committed figures byte for byte \\
26, 27 & the Mathematica notebook (skipped without WolframScript) \\
28, 29 & the two Stage-3 PDFs, rebuilt and compared with their registered editions \\
30 to 32 & every snapshotted file unchanged; all unit tests of the repository; every output written during this run \\
\bottomrule
\end{longtable}
\endgroup

\textbf{Running it.} It needs everything of Section 19.2 except that WolframScript is optional. PowerShell:

\begin{Verbatim}[fontsize=\small]
pwsh -NoProfile -File scripts/verify_stage3_dark_sector.ps1
\end{Verbatim}

Bash:

\begin{Verbatim}[fontsize=\small]
bash scripts/verify_stage3_dark_sector.sh
\end{Verbatim}

A successful run ends with

\begin{Verbatim}[fontsize=\small]
stage3_audit_fresh=OK
stage3_step_ok=stage3-32-fresh-outputs
stage3_dark_sector_verification=OK
\end{Verbatim}

Without WolframScript, two more lines appear: \texttt{stage3\_s\allowbreak{}kipped\_s\allowbreak{}tep=stag\allowbreak{}e3-26-ma\allowbreak{}thematic\allowbreak{}a-notebo\allowbreak{}ok,...} where step 26 would run, and \texttt{stage3\_m\allowbreak{}athemati\allowbreak{}ca=SKIPP\allowbreak{}ED (...)} just before the final line, which then still reads OK, because every other step passed.

\textbf{The committed OK runs.} The Bash twin took 28 minutes from a fresh clone of commit \texttt{\detokenize{1fb83c3}} on 2026-09-30, from 06:09 to 06:37, including 316 seconds for the 323 unit tests of that commit (\texttt{handoff/\allowbreak{}reviews/\allowbreak{}stage3\_g\allowbreak{}ate\_2026\allowbreak{}-09-30\_f\allowbreak{}resh\_clo\allowbreak{}ne\_OK.lo\allowbreak{}g}, which contains the OK line). The PowerShell twin also printed OK, from a fresh clone of commit \texttt{\detokenize{348c2e1}} (\texttt{handoff/\allowbreak{}reviews/\allowbreak{}stage3\_g\allowbreak{}ate\_ps1\_\allowbreak{}2026-09-\allowbreak{}30\_fresh\allowbreak{}\_clone\_O\allowbreak{}K.log}; that run was suspended for about 40 minutes at the user's request, so its duration is not a measurement). To repeat one of these runs exactly, check out its commit after cloning (\texttt{git chec\allowbreak{}kout 1fb\allowbreak{}83c3}, Section 19.4). The two complete runs of the experiments (steps 07 and 08) and the notebook executions take most of the time; \texttt{README.m\allowbreak{}d} summarizes it as "tens of minutes".

\textbf{The run at this edition's commit.} For this edition the Bash twin was run from a fresh clone of commit \texttt{\detokenize{1f2dd69}}. It passed steps 00 to 30: among them the two complete runs of the five experiments, the byte comparison of their 62 files, the five checkers with their repeat and refined runs, both executions of the Jupyter notebook, the Mathematica notebook (WolframScript was installed, so no step was skipped) and the two PDFs. It then stopped in step 31, after 2302 seconds (38 minutes), with

\begin{Verbatim}[fontsize=\small]
stage3_failed_step=stage3-31-unit-tests
stage3_failed_log=build/logs/stage3-31-unit-tests-bash.log
stage3_dark_sector_verification=FAILED
\end{Verbatim}

Step 31 runs every unit test of the repository, and four of the 579 tests of this commit fail. They test the committed report of the Stage-4 Jupyter notebook, not Stage 3 (see Section 19.11); every step that concerns Stage 3 passed.

\textbf{What changed for Stage 3 since the OK runs.} The following command, in either shell, lists every change between the commit of the Bash OK run and commit \texttt{\detokenize{1f2dd69}} in the files that the Stage-3 gate uses: the crate and its outputs, the gate, its audit program and its helpers, the engine setup, the EXP-3 analysis and the five checkers, the two notebooks with their runner and auditor, the Mathematica verifier, the PDF builder and its converter, the two Stage-3 documents, the build settings and the package pins:

\begin{Verbatim}[fontsize=\small]
git diff --stat 1fb83c3 1f2dd69 -- studies/dirac16complex_cosmology \
    artifacts/dirac16complex/numerics scripts/verify_stage3_dark_sector.sh \
    scripts/verify_stage3_dark_sector.ps1 scripts/verify_stage3_dark_sector_audit.py \
    scripts/run_logged.sh scripts/run_logged.ps1 scripts/setup_solver.sh \
    scripts/setup_solver.ps1 scripts/analyze_dirac16complex_exp3.py \
    "scripts/check_dirac16complex_exp*.py" notebooks/dirac16complex_dark_sector.ipynb \
    notebooks/Dirac16ComplexDarkSector.nb notebooks/run_notebook.py \
    notebooks/check_notebook.py scripts/verify_dirac16complex_mathematica_notebook.wls \
    scripts/build_provenance_pdf.py scripts/build_dissertation_tex.py \
    "provenance/DIRAC16COMPLEX_DARK_SECTOR_NUMERICS.*" \
    "provenance/DIRAC16COMPLEX_STUDENT_GUIDE.*" .cargo/config.toml \
    requirements-stage3.txt
\end{Verbatim}

(In PowerShell, end the continued lines with a backtick instead of the backslash.) It prints

\begin{Verbatim}[fontsize=\small]
 scripts/build_dissertation_tex.py | 14 ++++++++++++++
 scripts/build_provenance_pdf.py   |  7 +++++++
 2 files changed, 21 insertions(+)
\end{Verbatim}

so only the two PDF tools changed: they gained the option \texttt{--number\allowbreak{}-section\allowbreak{}s-from-z\allowbreak{}ero} for this book (Section 19.12), which the Stage-3 PDFs do not use. The run at \texttt{\detokenize{1f2dd69}} described above rebuilt both Stage-3 PDFs byte for byte (steps 28 and 29), and every committed Stage-3 file was unchanged at its end (step 30).

\textbf{Single programs.} The student guide (\texttt{provenan\allowbreak{}ce/DIRAC\allowbreak{}16COMPLE\allowbreak{}X\_STUDEN\allowbreak{}T\_GUIDE.\allowbreak{}md}) explains the Stage-3 programs step by step and was itself tested by replaying all of its command boxes; here are the essential commands, run for this edition from a fresh clone after Section 19.4. Build the program (in either shell):

\begin{Verbatim}[fontsize=\small]
cd studies/dirac16complex_cosmology
cargo build --release
cd ../..
\end{Verbatim}

(32 seconds on the test computer; the last line of the output begins with \texttt{\detokenize{Finished}}, and any compiler warning would stop the build, because the crate forbids warnings). Then run EXP-1 and EXP-5 into a scratch folder and check them, in PowerShell:

\begin{Verbatim}[fontsize=\small]
$env:PYTHONUTF8 = "1"
$bin = ".\studies\dirac16complex_cosmology\target\release\dirac16complex_cosmology.exe"
& $bin exp1 --output build/textbook
& $bin exp5 --output build/textbook
python scripts/check_dirac16complex_exp1.py --output build/textbook
python scripts/check_dirac16complex_exp5.py --output build/textbook
\end{Verbatim}

and in Bash:

\begin{Verbatim}[fontsize=\small]
export PYTHONUTF8=1
bin=./studies/dirac16complex_cosmology/target/release/dirac16complex_cosmology
$bin exp1 --output build/textbook
$bin exp5 --output build/textbook
python scripts/check_dirac16complex_exp1.py --output build/textbook
python scripts/check_dirac16complex_exp5.py --output build/textbook
\end{Verbatim}

The program prints one line \texttt{PASS - n\allowbreak{}ame: det\allowbreak{}ail} per self-check (10 for EXP-1) and ends with

\begin{Verbatim}[fontsize=\small]
exp1: solver_steps=33866 rhs_evaluations=34950 files=29 verdict=SUCCESS
SUCCESS
\end{Verbatim}

for EXP-1 (1 second) and with \texttt{exp5: so\allowbreak{}lver\_ste\allowbreak{}ps=2548 \allowbreak{}rhs\_eval\allowbreak{}uations=\allowbreak{}4800 fil\allowbreak{}es=5 ver\allowbreak{}dict=SUC\allowbreak{}CESS} and \texttt{\detokenize{SUCCESS}} for EXP-5 (under 1 second). The 29 and 5 files it wrote into \texttt{build/te\allowbreak{}xtbook/e\allowbreak{}xp1/} and \texttt{build/te\allowbreak{}xtbook/e\allowbreak{}xp5/} were byte-identical to the committed ones. The checkers end with \texttt{check\_co\allowbreak{}unt=23} (EXP-1, 2 seconds) and \texttt{check\_co\allowbreak{}unt=19} (EXP-5, 4 seconds), each with \texttt{failed\_c\allowbreak{}heck\_cou\allowbreak{}nt=0}. The committed checker reports have 25 and 21 checks: the two missing checks, \texttt{repeatBy\allowbreak{}teIdenti\allowbreak{}ty} and \texttt{refinedC\allowbreak{}onvergen\allowbreak{}ce}, need a repeat run and a run with tightened tolerances, which the checker performs itself when it is given the options \texttt{\detokenize{--repeat}} and \texttt{--refine\allowbreak{}d}, each followed by a scratch folder (student guide, §9). All five experiments at once:

\begin{Verbatim}[fontsize=\small]
$bin all --output build/textbook/all
\end{Verbatim}

(in PowerShell \texttt{\& \$bin a\allowbreak{}ll --out\allowbreak{}put buil\allowbreak{}d/textbo\allowbreak{}ok/all}) ran 167 seconds on the busy test computer and wrote the 62 output files of the five experiments, every one byte-identical to the committed file of the same name; the student guide measured 80 to 90 seconds for the same run (its §7.4).

\subsection{Stage 4: the Kohn--Sham states, and why its gate cannot pass yet}

\textbf{What exists and what it gives.} Stage 4 (Chapter 13) has five layers. The state of each at commit \texttt{\detokenize{1f2dd69}}, read from the committed files:

\begingroup
\small
\begin{longtable}{@{}>{\raggedright\arraybackslash}p{0.267\linewidth}>{\raggedright\arraybackslash}p{0.267\linewidth}>{\raggedright\arraybackslash}p{0.267\linewidth}@{}}
\toprule
\textbf{Layer} & \textbf{Program} & \textbf{Committed result} \\
\midrule
\endfirsthead
\toprule
\textbf{Layer} & \textbf{Program} & \textbf{Committed result} \\
\midrule
\endhead
exact theory, Wolfram & \texttt{scripts/\allowbreak{}verify\_d\allowbreak{}irac16co\allowbreak{}mplex\_ko\allowbreak{}hn\_sham.\allowbreak{}wls} & \texttt{wolfram-\allowbreak{}kohn-sha\allowbreak{}m-report\allowbreak{}.json}: 125 of 125 checks true \\
exact theory, sympy & \texttt{scripts/\allowbreak{}check\_di\allowbreak{}rac16com\allowbreak{}plex\_koh\allowbreak{}n\_sham\_t\allowbreak{}heory.py} & \texttt{python-t\allowbreak{}heory-re\allowbreak{}port.jso\allowbreak{}n}: 157 of 157 checks true \\
Rust Kohn--Sham solver & \texttt{studies/\allowbreak{}dirac16c\allowbreak{}omplex\_k\allowbreak{}ohn\_sham} & own checks all true: 33 (\texttt{\detokenize{spectrum}}), 137 (\texttt{\detokenize{scf}}), 65 (\texttt{\detokenize{excited}}), 102 (\texttt{\detokenize{thermo}}), 60 (\texttt{\detokenize{emt}}), in the \texttt{summary.\allowbreak{}json} of each folder of \texttt{\detokenize{rust/}}; a repeat run byte-identical in all 330 files compared, and a run with tolerances divided by 10 agreeing to $5.96\times10^{-8}$ relative in the energies of 65 runs and to $2.8\times10^{-8}$ in 364829 levels (\texttt{rust/det\allowbreak{}erminism\allowbreak{}-report.\allowbreak{}json}, 4 of 4 checks true) \\
independent reference solver & \texttt{scripts/\allowbreak{}ks\_refer\allowbreak{}ence\_sol\allowbreak{}ver.py} & 56 runs in \texttt{referenc\allowbreak{}e/} (\texttt{referenc\allowbreak{}e/refere\allowbreak{}nce-summ\allowbreak{}ary.json}) \\
cross-checker, Rust against reference & \texttt{scripts/\allowbreak{}check\_di\allowbreak{}rac16com\allowbreak{}plex\_koh\allowbreak{}n\_sham.p\allowbreak{}y} & \texttt{python-c\allowbreak{}heck-rep\allowbreak{}ort.json}: 63 checks, 62 true, 1 false \\
\bottomrule
\end{longtable}
\endgroup

All files of the table lie in \texttt{artifact\allowbreak{}s/dirac1\allowbreak{}6complex\allowbreak{}/kohn-sh\allowbreak{}am/}. The one false check is \texttt{canonica\allowbreak{}l\_eigenv\allowbreak{}alues}. Its measurement \texttt{canonica\allowbreak{}l\_eigenv\allowbreak{}alues\_de\allowbreak{}tail} records 58412 comparisons of Kohn--Sham levels between the two solvers; the worst one, a deep level at $\varepsilon=-1.0537\,m$ of the smeared $N=1016$ ensemble at $-\hat\lambda_2$ (Section 13.11), differs by $2.195\times10^{-6}\,m$, while the tolerance there is $1.054\times10^{-6}\,m$ (ratio 2.08). All eight failing comparisons belong to this one run (four in its ground-state record and four in its excited-state record); every other run passes. The Delta-SCF excitation energy of that ensemble, which was the open comparison when Chapter 13 was written, now agrees within its tolerance (check \texttt{canonica\allowbreak{}l\_deltaS\allowbreak{}CF}, worst ratio 0.79), after the reference was run on twice finer grids for this one case (erratum E4.13 of \texttt{handoff/\allowbreak{}specs/ST\allowbreak{}AGE4\_SPE\allowbreak{}C.md}). The report also records one comparison that did not run, the reference solver's self-tests (\texttt{comparis\allowbreak{}onsNotRu\allowbreak{}n}).

Two more Stage-4 reports belong to the notebooks: the Mathematica notebook's \texttt{mathemat\allowbreak{}ica-repo\allowbreak{}rt.json} has 94 of 94 checks true, and the Jupyter notebook's \texttt{notebook\allowbreak{}-report.\allowbreak{}json} has 68 of 74 checks true and the verdict FAILURE. Its six false checks (their names begin with \texttt{gauntlet\allowbreak{}\_} and \texttt{\detokenize{r07_}}) repeat the cross-check comparisons as they stood when the notebook was last executed, before the reruns that closed the Delta-SCF comparison and before the final cross-check report existed (the notebook report itself records that it found \texttt{python-c\allowbreak{}heck-rep\allowbreak{}ort.json} absent); the notebook has not been executed again since.

\textbf{Why the gate cannot pass yet.} The gate \texttt{scripts/\allowbreak{}verify\_s\allowbreak{}tage4\_ko\allowbreak{}hn\_sham.\allowbreak{}\{ps1,sh\}} exists and was tested in parts, but it has never been run to the end. It cannot end with OK at present, for four reasons that the gate itself detects:

\begin{itemize}
\item Documents are missing. Step 00 copies every committed Stage-4 file that the gate must leave unchanged, among them \texttt{provenan\allowbreak{}ce/DIRAC\allowbreak{}16COMPLE\allowbreak{}X\_KOHN\_S\allowbreak{}HAM\_PRIM\allowbreak{}ORDIAL.t\allowbreak{}ex} and \texttt{\detokenize{.pdf}} and the three files of the Stage-4 student guide \texttt{provenan\allowbreak{}ce/DIRAC\allowbreak{}16COMPLE\allowbreak{}X\_KOHN\_S\allowbreak{}HAM\_STUD\allowbreak{}ENT\_GUID\allowbreak{}E}, and fails because these do not exist yet (only \texttt{provenan\allowbreak{}ce/DIRAC\allowbreak{}16COMPLE\allowbreak{}X\_KOHN\_S\allowbreak{}HAM\_PRIM\allowbreak{}ORDIAL.m\allowbreak{}d} is committed; the partial run below); steps 32 and 33 build the two PDFs and compare them with registered editions, and no Stage-4 edition is registered in \texttt{provenan\allowbreak{}ce/pdf-s\allowbreak{}pecifica\allowbreak{}tions.js\allowbreak{}on}.
\item The collected summary \texttt{artifact\allowbreak{}s/dirac1\allowbreak{}6complex\allowbreak{}/kohn-sh\allowbreak{}am/kohn-\allowbreak{}sham-sum\allowbreak{}mary.jso\allowbreak{}n}, with which step 23 compares its fresh summary, does not exist yet.
\item Step 21 requires every check of the fresh cross-check report to be true, and \texttt{canonica\allowbreak{}l\_eigenv\allowbreak{}alues} is false.
\item Step 35 runs the unit tests of Stage 4, and four of them, which test the committed report of the Jupyter notebook, fail (Section 19.11).
\end{itemize}

The numbers of Chapter 13 are therefore those of the Rust solver with its own checks and with the cross-check status just described.

\textbf{The dry run.} The option \texttt{\detokenize{-DryRun}} (PowerShell) or \texttt{--dry-ru\allowbreak{}n} (Bash) prints every step with its expected wall time and its command, runs nothing, and ends with \texttt{stage4\_k\allowbreak{}ohn\_sham\allowbreak{}\_verific\allowbreak{}ation=DR\allowbreak{}Y-RUN}:

\begin{Verbatim}[fontsize=\small]
pwsh -NoProfile -File scripts/verify_stage4_kohn_sham.ps1 -DryRun
\end{Verbatim}

\begin{Verbatim}[fontsize=\small]
bash scripts/verify_stage4_kohn_sham.sh --dry-run
\end{Verbatim}

It lists 37 steps, \texttt{stage4-0\allowbreak{}0-snapsh\allowbreak{}ot} to \texttt{stage4-3\allowbreak{}6-fresh-\allowbreak{}outputs}, and took 4 seconds in Bash and under a second in PowerShell. The header of the PowerShell twin explains each step and records the measured time of each part (its paragraph TIMING): about 4 hours in total, most of it in step 14, which recomputes the complete canonical Rust tree (about 3 hours, dominated by the \texttt{\detokenize{thermo}} subcommand; an estimate, since this step has not yet been run as a whole).

\textbf{A partial run.} The option \texttt{\detokenize{-Steps}} (PowerShell) or \texttt{\detokenize{--steps}} (Bash) runs only the listed steps and then ends with \texttt{stage4\_k\allowbreak{}ohn\_sham\allowbreak{}\_verific\allowbreak{}ation=PA\allowbreak{}RTIAL} and the exit code 3, never with OK. Steps 01 to 13 (engine, exact theory, generated constants, the four cargo steps and the configuration check) need no long Rust run:

\begin{Verbatim}[fontsize=\small]
pwsh -NoProfile -File scripts/verify_stage4_kohn_sham.ps1 `
    -Steps 01,02,03,04,05,06,07,08,09,10,11,12,13
\end{Verbatim}

\begin{Verbatim}[fontsize=\small]
bash scripts/verify_stage4_kohn_sham.sh \
    --steps 01,02,03,04,05,06,07,08,09,10,11,12,13
\end{Verbatim}

For this edition the Bash command was run from a fresh clone. All thirteen selected steps passed in 151 seconds, and the run ended with

\begin{Verbatim}[fontsize=\small]
stage4_selected_steps=01,02,03,04,05,06,07,08,09,10,11,12,13
stage4_kohn_sham_verification=PARTIAL
\end{Verbatim}

and the exit code 3. The exact Wolfram theory (step 03) took 19 seconds and the sympy theory (step 04) 54 seconds; step 05 found their four output files byte-identical to the committed ones; steps 06 to 08 regenerated the Rust file of exact constants, \texttt{studies/\allowbreak{}dirac16c\allowbreak{}omplex\_k\allowbreak{}ohn\_sham\allowbreak{}/src/gen\allowbreak{}erated.r\allowbreak{}s}, and its report \texttt{rust/gen\allowbreak{}erator-r\allowbreak{}eport.js\allowbreak{}on} byte for byte; and step 11 ran the 49 Rust tests of the solver in 65 seconds (47 passed, and the 2 timing probes that are marked to be skipped in a normal run were skipped). The PowerShell command, run afterwards in the same clone, also passed all thirteen steps and ended with the same two lines, after 121 seconds (step 11: 55 seconds). In both runs the log of step 11 contains one line that begins with \texttt{\detokenize{[ERROR]}} and ends with \texttt{the righ\allowbreak{}t-hand s\allowbreak{}ide rout\allowbreak{}ine fail\allowbreak{}ed in an\allowbreak{} unrecov\allowbreak{}erable m\allowbreak{}anner}. The engine prints such a line whenever one integration stops with an error; the solver is written to handle such errors (for example, the function \texttt{\detokenize{profile}} in \texttt{studies/\allowbreak{}dirac16c\allowbreak{}omplex\_k\allowbreak{}ohn\_sham\allowbreak{}/src/sho\allowbreak{}oting.rs} repeats a failed integration with other settings), and the line \texttt{test res\allowbreak{}ult: ok.\allowbreak{} 47 pass\allowbreak{}ed; 0 fa\allowbreak{}iled; 2 \allowbreak{}ignored} shows that no test failed.

Step 00 is left out of the list because at this commit it fails. With the list \texttt{00,01,..\allowbreak{}.,13} the same run stopped after 1 second. It first named the missing file, in the line \texttt{stage4\_a\allowbreak{}udit\_pro\allowbreak{}blem=can\allowbreak{}not snap\allowbreak{}shot mis\allowbreak{}sing fil\allowbreak{}e} followed by the path \texttt{provenan\allowbreak{}ce/DIRAC\allowbreak{}16COMPLE\allowbreak{}X\_KOHN\_S\allowbreak{}HAM\_PRIM\allowbreak{}ORDIAL.t\allowbreak{}ex}, and then printed

\begin{Verbatim}[fontsize=\small]
stage4_audit_snapshot=FAILED
stage4_failed_step=stage4-00-snapshot
stage4_failed_log=build/logs/stage4-00-snapshot-bash.log
stage4_kohn_sham_verification=FAILED
\end{Verbatim}

The snapshot step copies, among other files, the LaTeX and PDF files of the Stage-4 document and the Markdown, LaTeX and PDF files of the Stage-4 student guide, so that the gate can prove at the end that they did not change; none of these five files exists yet (the first reason above).

\textbf{Single programs.} The exact theory into a scratch folder, in PowerShell:

\begin{Verbatim}[fontsize=\small]
$env:PYTHONUTF8 = "1"
$k = "build/textbook/stage4"
wolframscript -file scripts/verify_dirac16complex_kohn_sham.wls `
    "$k/wolfram-kohn-sham-report.json"
python scripts/check_dirac16complex_kohn_sham_theory.py `
    --output "$k/python-theory-report.json" --table "$k/exchange-table.json"
\end{Verbatim}

and in Bash:

\begin{Verbatim}[fontsize=\small]
export PYTHONUTF8=1
k=build/textbook/stage4
wolframscript -file scripts/verify_dirac16complex_kohn_sham.wls \
    "$k/wolfram-kohn-sham-report.json"
python scripts/check_dirac16complex_kohn_sham_theory.py \
    --output "$k/python-theory-report.json" --table "$k/exchange-table.json"
\end{Verbatim}

These are the programs of steps 03 and 04 of the gate, with a scratch folder instead of \texttt{build/st\allowbreak{}age4/the\allowbreak{}ory}. In the partial run above they printed \texttt{check\_co\allowbreak{}unt=125} (Wolfram, 19 seconds) and \texttt{check\_co\allowbreak{}unt=157} (sympy, 54 seconds), each with \texttt{failed\_c\allowbreak{}heck\_cou\allowbreak{}nt=0}. The Wolfram verifier writes the theory file \texttt{kohn-sha\allowbreak{}m-theory\allowbreak{}.json} next to its report, and the sympy checker writes the exchange table given after \texttt{\detokenize{--table}}; all four files were byte-identical to the committed ones in \texttt{artifact\allowbreak{}s/dirac1\allowbreak{}6complex\allowbreak{}/kohn-sh\allowbreak{}am/}.

The Rust solver is built and run like the Stage-3 program (in either shell; in PowerShell call the program as \texttt{\detokenize{& $ks}} with the path ending in \texttt{\detokenize{.exe}}):

\begin{Verbatim}[fontsize=\small]
cargo build --release --manifest-path studies/dirac16complex_kohn_sham/Cargo.toml
ks=./studies/dirac16complex_kohn_sham/target/release/dirac16complex_kohn_sham
$ks print-config
$ks spectrum --output build/textbook/ks
\end{Verbatim}

The line \texttt{\detokenize{ks=...}} is Bash; in PowerShell write \texttt{\$ks = ".\allowbreak{}\textbackslash{}studies\allowbreak{}\textbackslash{}dirac16\allowbreak{}complex\_\allowbreak{}kohn\_sha\allowbreak{}m\textbackslash{}target\allowbreak{}\textbackslash{}release\allowbreak{}\textbackslash{}dirac16\allowbreak{}complex\_\allowbreak{}kohn\_sha\allowbreak{}m.exe"}. For this edition the build took 1 second, because step 12 of the partial run above had already compiled the engine and the solver; \texttt{print-co\allowbreak{}nfig} ended with \texttt{\detokenize{SUCCESS}}, and \texttt{\detokenize{spectrum}} ran 37 seconds and ended with

\begin{Verbatim}[fontsize=\small]
spectrum: solver_steps=39152872 rhs_evaluations=47587645 files=15 verdict=SUCCESS
SUCCESS
\end{Verbatim}

Its 15 files in \texttt{build/te\allowbreak{}xtbook/k\allowbreak{}s/spectr\allowbreak{}um/} were byte-identical to the committed ones in \texttt{artifact\allowbreak{}s/dirac1\allowbreak{}6complex\allowbreak{}/kohn-sh\allowbreak{}am/rust/\allowbreak{}spectrum\allowbreak{}/}. The other subcommands, \texttt{\detokenize{scf}}, \texttt{\detokenize{excited}}, \texttt{\detokenize{thermo}} and \texttt{\detokenize{emt}}, take from minutes to hours (the TIMING paragraph named above); \texttt{\detokenize{all}} runs the five of them.

\subsection{Stage 5 and the matter--antimatter analysis}

\textbf{What exists.} Stage 5 (Chapters 6, 7, 14 and 15) and the matter--antimatter analysis (Chapter 17) have no gate yet. Their exact theory is committed, each part in two independent implementations, and every check of the committed reports is true:

\begingroup
\small
\begin{longtable}{@{}>{\raggedright\arraybackslash}p{0.267\linewidth}>{\raggedright\arraybackslash}p{0.267\linewidth}>{\raggedright\arraybackslash}p{0.267\linewidth}@{}}
\toprule
\textbf{Part} & \textbf{Program} & \textbf{Committed report (checks true)} \\
\midrule
\endfirsthead
\toprule
\textbf{Part} & \textbf{Program} & \textbf{Committed report (checks true)} \\
\midrule
\endhead
dirac16complex00, Wolfram & \texttt{scripts/\allowbreak{}verify\_d\allowbreak{}irac16co\allowbreak{}mplex00.\allowbreak{}wls} with \texttt{wolfram/\allowbreak{}Dirac16C\allowbreak{}omplex00\allowbreak{}.wl} & \texttt{wolfram-\allowbreak{}dirac16c\allowbreak{}omplex00\allowbreak{}-report.\allowbreak{}json} (46 of 46) and \texttt{dirac16c\allowbreak{}omplex00\allowbreak{}-theory.\allowbreak{}json} \\
dirac16complex00, sympy & \texttt{scripts/\allowbreak{}check\_di\allowbreak{}rac16com\allowbreak{}plex00.p\allowbreak{}y} & \texttt{python-d\allowbreak{}irac16co\allowbreak{}mplex00-\allowbreak{}report.j\allowbreak{}son} (49 of 49) \\
pairing theorems, Wolfram & \texttt{scripts/\allowbreak{}verify\_d\allowbreak{}irac16co\allowbreak{}mplex\_pa\allowbreak{}iring.wl\allowbreak{}s} with \texttt{wolfram/\allowbreak{}Dirac16C\allowbreak{}omplexPa\allowbreak{}iring.wl} & \texttt{wolfram-\allowbreak{}pairing-\allowbreak{}report.j\allowbreak{}son} (141 of 141) and \texttt{pairing-\allowbreak{}theory.j\allowbreak{}son} \\
pairing theorems, sympy & \texttt{scripts/\allowbreak{}check\_di\allowbreak{}rac16com\allowbreak{}plex\_pai\allowbreak{}ring.py} & \texttt{python-p\allowbreak{}airing-r\allowbreak{}eport.js\allowbreak{}on} (172 of 172) \\
matter and antimatter, Wolfram & \texttt{scripts/\allowbreak{}verify\_d\allowbreak{}irac16co\allowbreak{}mplex\_ma\allowbreak{}tter\_ant\allowbreak{}imatter.\allowbreak{}wls} with \texttt{wolfram/\allowbreak{}Dirac16C\allowbreak{}omplexMa\allowbreak{}tterAnti\allowbreak{}matter.w\allowbreak{}l} & \texttt{wolfram-\allowbreak{}matter-a\allowbreak{}ntimatte\allowbreak{}r-report\allowbreak{}.json} (44 of 44) and \texttt{matter-a\allowbreak{}ntimatte\allowbreak{}r-theory\allowbreak{}.json} \\
matter and antimatter, sympy & \texttt{scripts/\allowbreak{}check\_di\allowbreak{}rac16com\allowbreak{}plex\_mat\allowbreak{}ter\_anti\allowbreak{}matter.p\allowbreak{}y} & \texttt{python-m\allowbreak{}atter-an\allowbreak{}timatter\allowbreak{}-report.\allowbreak{}json} (77 of 77) \\
\bottomrule
\end{longtable}
\endgroup

The Stage-5 reports lie in \texttt{artifact\allowbreak{}s/dirac1\allowbreak{}6complex\allowbreak{}/pair-cr\allowbreak{}eation/}, the matter--antimatter reports in \texttt{artifact\allowbreak{}s/dirac1\allowbreak{}6complex\allowbreak{}/matter-\allowbreak{}antimatt\allowbreak{}er/}. The matter--antimatter document \texttt{provenan\allowbreak{}ce/DIRAC\allowbreak{}16COMPLE\allowbreak{}X\_MATTER\allowbreak{}\_ANTIMAT\allowbreak{}TER.md} is written, and its PDF is registered in \texttt{provenan\allowbreak{}ce/pdf-s\allowbreak{}pecifica\allowbreak{}tions.js\allowbreak{}on} (edition \texttt{dirac16c\allowbreak{}omplex-m\allowbreak{}atter-an\allowbreak{}timatter}, 36 pages at commit \texttt{\detokenize{1f2dd69}}; the document records in its §1 and §11.4 what it proves and what it cites).

\textbf{A Stage-5 result cited as provisional.} All twelve Wolfram checks \texttt{PAIR\_T1k\allowbreak{}rein\_*} and all thirteen sympy checks \texttt{S5\_T1kre\allowbreak{}in\_*} of the exact Fock-level model are true. Their physical reading in the Stage-5 files is not: the entry \texttt{imageFie\allowbreak{}ld} of the key \texttt{\detokenize{T1krein}} of \texttt{pairing-\allowbreak{}theory.j\allowbreak{}son} calls the values $-|\varepsilon|$ and $-1$ per quantum the energy and the charge of an "image universe", while they are the values of the formulas of $\mathcal L_{-m,-\lambda}$ evaluated on the image field, whose own energy and charge are $+H_+$ and $+Q_+$ (Section 15.4; matter--antimatter document, §7.3; check \texttt{MA\_M4\_im\allowbreak{}ageField\allowbreak{}FockMode\allowbreak{}l}). \texttt{HANDOFF.\allowbreak{}md} (section 0.4, item C) records the correction to be made in \texttt{wolfram/\allowbreak{}Dirac16C\allowbreak{}omplexPa\allowbreak{}iring.wl} and \texttt{scripts/\allowbreak{}check\_di\allowbreak{}rac16com\allowbreak{}plex\_pai\allowbreak{}ring.py} before the Stage-5 documents are written, and both matter--antimatter reports cite the Stage-5 result as PROVISIONAL until the Stage-5 gate has passed (measurement \texttt{M4\_krein\allowbreak{}LevelSta\allowbreak{}tus} of the Wolfram report). The correction may change the names and the number of the checks of the two pairing reports, and with them the counts of this section and the rows of Section 20.14.

\textbf{What does not exist yet.} The two Stage-5 documents \texttt{provenan\allowbreak{}ce/DIRAC\allowbreak{}16COMPLE\allowbreak{}X00\_FIEL\allowbreak{}D\_THEORY} and \texttt{provenan\allowbreak{}ce/DIRAC\allowbreak{}16COMPLE\allowbreak{}X\_PAIR\_C\allowbreak{}REATION} and the gate \texttt{scripts/\allowbreak{}verify\_s\allowbreak{}tage5\_pa\allowbreak{}ir\_creat\allowbreak{}ion.\{ps1\allowbreak{},sh\}} planned in \texttt{handoff/\allowbreak{}specs/ST\allowbreak{}AGE5\_SPE\allowbreak{}C.md} (its §6) are not written. The Kohn--Sham demonstration of the pairing (Chapter 15) is \textbf{partial}:

\begin{itemize}
\item The Rust solver has the new subcommand \texttt{\detokenize{pairs}} (file \texttt{studies/\allowbreak{}dirac16c\allowbreak{}omplex\_k\allowbreak{}ohn\_sham\allowbreak{}/src/pai\allowbreak{}rs.rs}). At commit \texttt{\detokenize{1f2dd69}} its committed outputs, in \texttt{artifact\allowbreak{}s/dirac1\allowbreak{}6complex\allowbreak{}/pair-cr\allowbreak{}eation/r\allowbreak{}ust/pair\allowbreak{}s/}, cover 25 parameter sets of dirac16complex, each with the three universes $+M$, $-M$ and the $-M$ control and a file \texttt{pairing.\allowbreak{}json}: mass $m=1$ with $N=8$ at the five couplings $0,\pm\hat\lambda_1,\pm\hat\lambda_2$ ($T=0$) and with $N=112$ at the same five couplings at $T=0$ and at $T=0.1\,m$ (15 sets); mass $m=3$ with $N=8$ at the same five couplings at $T=0$ (5 sets); and mass $m=3$ with $N=112$ at the coupling 0 and at $+\hat\lambda_1$, each at $T=0$ and $T=0.1\,m$, and at $-\hat\lambda_1$ at $T=0$ (5 sets). The folder has no collected summary, and there are no committed Rust runs of dirac16complex00. (Chapters 15 and 16 describe in detail the 18 sets that were committed when they were written, that is all of these except the five sets with $m=3$, $N=112$ and the two with $m=3$, $N=8$ at $\pm\hat\lambda_2$, and summarize these seven later sets in a few sentences; Chapter 18 lists all 25. Further runs are added as they finish, so a later commit can have more.)
\item The reference runs of the independent solver (\texttt{scripts/\allowbreak{}ks\_refer\allowbreak{}ence\_pai\allowbreak{}rs.py}, outputs in \texttt{artifact\allowbreak{}s/dirac1\allowbreak{}6complex\allowbreak{}/pair-cr\allowbreak{}eation/r\allowbreak{}eference\allowbreak{}/}) are marked incomplete in their summary \texttt{referenc\allowbreak{}e-pairs-\allowbreak{}summary.\allowbreak{}json} (entry \texttt{\detokenize{complete}} is false): of the 49 runs recorded there, 17 converged (10 of dirac16complex, 7 of dirac16complex00), 2 did not converge, 6 failed and 24 were not attempted, and 131 runs of the planned matrix are listed as pending.
\item The checker \texttt{scripts/\allowbreak{}check\_di\allowbreak{}rac16com\allowbreak{}plex\_pai\allowbreak{}rs.py} (Rust against reference, pairing identities, pair totals) exists, but no report of it is committed.
\item One Stage-5 numerical report is complete: \texttt{artifact\allowbreak{}s/dirac1\allowbreak{}6complex\allowbreak{}/pair-cr\allowbreak{}eation/r\allowbreak{}ust/stag\allowbreak{}e4-ident\allowbreak{}ity-repo\allowbreak{}rt.json}, 9 of 9 checks true, which shows that adding the subcommand \texttt{\detokenize{pairs}} left the Stage-4 outputs byte-identical (a fresh \texttt{\detokenize{spectrum}} run and quick \texttt{\detokenize{scf}} and \texttt{\detokenize{excited}} runs).
\end{itemize}

\textbf{Rerunning the exact theory.} Into the scratch folder \texttt{build/te\allowbreak{}xtbook/s\allowbreak{}5}, in PowerShell:

\begin{Verbatim}[fontsize=\small]
$env:PYTHONUTF8 = "1"
$o = "build/textbook/s5"
python scripts/check_dirac16complex00.py `
    --output "$o/python-dirac16complex00-report.json"
python scripts/check_dirac16complex_pairing.py --output "$o/python-pairing-report.json"
python scripts/check_dirac16complex_matter_antimatter.py `
    --output "$o/python-matter-antimatter-report.json"
wolframscript -file scripts/verify_dirac16complex00.wls `
    "$o/wolfram-dirac16complex00-report.json"
wolframscript -file scripts/verify_dirac16complex_pairing.wls `
    "$o/wolfram-pairing-report.json"
wolframscript -file scripts/verify_dirac16complex_matter_antimatter.wls `
    "$o/wolfram-matter-antimatter-report.json"
\end{Verbatim}

and in Bash:

\begin{Verbatim}[fontsize=\small]
export PYTHONUTF8=1
o=build/textbook/s5
python scripts/check_dirac16complex00.py \
    --output "$o/python-dirac16complex00-report.json"
python scripts/check_dirac16complex_pairing.py --output "$o/python-pairing-report.json"
python scripts/check_dirac16complex_matter_antimatter.py \
    --output "$o/python-matter-antimatter-report.json"
wolframscript -file scripts/verify_dirac16complex00.wls \
    "$o/wolfram-dirac16complex00-report.json"
wolframscript -file scripts/verify_dirac16complex_pairing.wls \
    "$o/wolfram-pairing-report.json"
wolframscript -file scripts/verify_dirac16complex_matter_antimatter.wls \
    "$o/wolfram-matter-antimatter-report.json"
\end{Verbatim}

The three Python checkers compare their own results with the committed Wolfram theory files (they never use them as truth, only in their agreement checks), so they can run first. Each Wolfram verifier writes its theory file next to its report, here into \texttt{build/te\allowbreak{}xtbook/s\allowbreak{}5/}. For this edition the six Bash commands were run from a fresh clone, one after the other:

\begingroup
\small
\begin{longtable}{@{}>{\raggedright\arraybackslash}p{0.200\linewidth}>{\raggedright\arraybackslash}p{0.200\linewidth}>{\raggedright\arraybackslash}p{0.200\linewidth}>{\raggedright\arraybackslash}p{0.200\linewidth}@{}}
\toprule
\textbf{Command} & \textbf{Wall time} & \textbf{Final lines} & \textbf{New files against the committed ones} \\
\midrule
\endfirsthead
\toprule
\textbf{Command} & \textbf{Wall time} & \textbf{Final lines} & \textbf{New files against the committed ones} \\
\midrule
\endhead
\texttt{check\_di\allowbreak{}rac16com\allowbreak{}plex00.p\allowbreak{}y} & 372 s & \texttt{check\_co\allowbreak{}unt=49}, \texttt{failed\_c\allowbreak{}heck\_cou\allowbreak{}nt=0} & report byte-identical \\
\texttt{check\_di\allowbreak{}rac16com\allowbreak{}plex\_pai\allowbreak{}ring.py} & 144 s & \texttt{check\_co\allowbreak{}unt=172}, \texttt{failed\_c\allowbreak{}heck\_cou\allowbreak{}nt=0} & report byte-identical \\
\texttt{check\_di\allowbreak{}rac16com\allowbreak{}plex\_mat\allowbreak{}ter\_anti\allowbreak{}matter.p\allowbreak{}y} & 44 s & \texttt{check\_co\allowbreak{}unt=77}, \texttt{failed\_c\allowbreak{}heck\_cou\allowbreak{}nt=0} & report byte-identical \\
\texttt{verify\_d\allowbreak{}irac16co\allowbreak{}mplex00.\allowbreak{}wls} & 369 s & \texttt{check\_co\allowbreak{}unt=46}, \texttt{failed\_c\allowbreak{}heck\_cou\allowbreak{}nt=0} & report and theory file byte-identical \\
\texttt{verify\_d\allowbreak{}irac16co\allowbreak{}mplex\_pa\allowbreak{}iring.wl\allowbreak{}s} & 338 s & \texttt{check\_co\allowbreak{}unt=141}, \texttt{failed\_c\allowbreak{}heck\_cou\allowbreak{}nt=0} & report and theory file byte-identical \\
\texttt{verify\_d\allowbreak{}irac16co\allowbreak{}mplex\_ma\allowbreak{}tter\_ant\allowbreak{}imatter.\allowbreak{}wls} & 453 s; 326 s from a clone with a short path & \texttt{check\_co\allowbreak{}unt=44}, \texttt{failed\_c\allowbreak{}heck\_cou\allowbreak{}nt=0} & theory file byte-identical; report byte-identical from a clone with a short path (below) \\
\bottomrule
\end{longtable}
\endgroup

Together they took 29 minutes on the busy test computer, and every check of every report was true. This clone, however, lay 168 characters deep (Section 19.4), and there the Wolfram matter--antimatter report differed from the committed one in three counts of its measurement \texttt{M1\_ksFix\allowbreak{}edNetNum\allowbreak{}ber}: 55 instead of 56 reference run files, 165 instead of 168 reference levels and 32 instead of 33 Rust run files. The verifier finds these files of Stage 4 by listing folders, and WolframScript did not list the two files \texttt{m1\_L3\_N1\allowbreak{}12\_lamp1\allowbreak{}rescaled\allowbreak{}\_T0\_dk0p\allowbreak{}15163266\allowbreak{}49281583\allowbreak{}6/run.js\allowbreak{}on} (one in \texttt{kohn-sha\allowbreak{}m/refere\allowbreak{}nce/}, one in \texttt{kohn-sha\allowbreak{}m/rust/s\allowbreak{}cf/}), whose full paths there had 272 and 271 characters, more than the 260 that Windows allows by default. The check \texttt{MA\_M1\_ks\allowbreak{}FixedNet\allowbreak{}NumberRe\allowbreak{}corded} stayed true, because every file that was read passed. Run again from the clone of Sections 19.3 to 19.9, which lay 83 characters deep, the same command wrote a report and a theory file byte-identical to the committed ones. So all twelve new files reproduce the committed files byte for byte when the clone has a short path.

\textbf{The matter--antimatter PDF.} In either shell:

\begin{Verbatim}[fontsize=\small]
python scripts/build_provenance_pdf.py provenance/DIRAC16COMPLEX_MATTER_ANTIMATTER.md
\end{Verbatim}

The builder converts the Markdown twice, runs pdflatex three times on each copy, requires both logs to be free of warnings and the two PDFs to be byte-identical, and compares the result with the registered edition (Section 19.12 describes it in detail). It ran 5 seconds and ended with

\begin{Verbatim}[fontsize=\small]
measurement_pageCount=36
...
check_count=14
failed_check_count=0
provenance_pdf=OK
\end{Verbatim}

and the rebuilt PDF was byte-identical to the one committed at \texttt{\detokenize{1f2dd69}} (sha256 \texttt{f4d6cb51\allowbreak{}02849614\allowbreak{}...}).

\subsection{The unit tests}

\textbf{What they test.} Besides the verifiers, the repository contains \textbf{unit tests} in the folder \texttt{\detokenize{tests/}}: small programs that test the tools (the Markdown converter, the PDF builder, the textbook assembler, the gates) and pin the committed evidence (for example, that the check counts, numbers and sha256 values quoted in a document agree with the reports). Python's module \texttt{\detokenize{unittest}} finds and runs them. One file of tests, in either shell:

\begin{Verbatim}[fontsize=\small]
python -m unittest discover -s tests -p "test_d16c_textbook_assembler.py" -v
\end{Verbatim}

This runs the 34 tests of the textbook assembler (1 second) and ends with \texttt{\detokenize{OK}}. The pattern after \texttt{\detokenize{-p}} chooses the files; the gates use their own patterns (Sections 19.6 to 19.8). All tests of the repository:

\begin{Verbatim}[fontsize=\small]
python -m unittest discover -s tests -v
\end{Verbatim}

Step 31 of the Stage-3 gate runs exactly this command. In the run of Section 19.8, from a fresh clone of commit \texttt{\detokenize{1f2dd69}}, it ran 579 tests in 505 seconds and ended with

\begin{Verbatim}[fontsize=\small]
Ran 579 tests in 505.024s

FAILED (failures=4, skipped=2)
\end{Verbatim}

The two skipped tests need the folder \texttt{dirac-ma\allowbreak{}in/} (Section 19.6). The four failures are tests of the file \texttt{tests/te\allowbreak{}st\_d16c\_\allowbreak{}kohn\_sha\allowbreak{}m\_notebo\allowbreak{}ok.py}, which has 30 tests. All four test the committed report of the Stage-4 Jupyter notebook, a committed record of unfinished work, not a tool or a result of Stages 1 to 3:

\begin{itemize}
\item \texttt{test\_ver\allowbreak{}dict\_suc\allowbreak{}cess}: the report's verdict is FAILURE, with the six false checks of Section 19.9;
\item \texttt{test\_rep\allowbreak{}ort\_is\_c\allowbreak{}urrent}: the summary \texttt{referenc\allowbreak{}e-summar\allowbreak{}y.json} of the Stage-4 reference runs (Section 19.9), on which the report depends, has changed since the report was written, so the sha256 that the report records for it is out of date;
\item \texttt{test\_rea\allowbreak{}udit\_of\_\allowbreak{}the\_comm\allowbreak{}itted\_no\allowbreak{}tebook\_r\allowbreak{}eproduce\allowbreak{}s\_the\_fi\allowbreak{}rst\_exec\allowbreak{}ution}: for the same reason, a new audit of the committed notebook records other hashes than the report;
\item \texttt{test\_two\allowbreak{}\_executi\allowbreak{}ons\_with\allowbreak{}\_identic\allowbreak{}al\_resul\allowbreak{}ts}: the report names an older way of executing the notebook (through a Jupyter kernel) than the one the test expects (\texttt{notebook\allowbreak{}s/run\_no\allowbreak{}tebook.p\allowbreak{}y}).
\end{itemize}

The notebook has to be executed again once the cross-check is settled. The 49 tests of the two files \texttt{tests/te\allowbreak{}st\_d16c\_\allowbreak{}matter\_a\allowbreak{}ntimatte\allowbreak{}r.py} and \texttt{tests/te\allowbreak{}st\_d16c\_\allowbreak{}matter\_a\allowbreak{}ntimatte\allowbreak{}r\_public\allowbreak{}ation.py}, run on their own with the pattern \texttt{"test\_d1\allowbreak{}6c\_matte\allowbreak{}r\_antima\allowbreak{}tter*.py\allowbreak{}"} in the same clone after the gate had finished, all passed (35 seconds).

Every test that the gates of Stages 1 and 2 run passed (Sections 19.6 and 19.7). The Stage-3 gate, however, runs all 579 tests in its step 31 and therefore cannot end with OK at this commit (Section 19.8), and the Stage-4 gate runs the four failing notebook tests in its step 35.

\subsection{Building this textbook}

\textbf{The two programs.} The book is written as 21 chapter files, one per chapter, in the folder of the first line below. The assembler (second line) checks them and joins them, with the title and the abstract, into the one Markdown file of the third line; the PDF builder (fourth line) turns that file into LaTeX and the LaTeX into the PDF:

\begin{Verbatim}[fontsize=\small]
provenance/textbook/chapters/     00-how-to-read.md ... 20-glossary-and-check-index.md
scripts/build_textbook.py         the assembler
provenance/DIRAC16COMPLEX_TEXTBOOK.md   the assembled book (.tex and .pdf beside it)
scripts/build_provenance_pdf.py   the PDF builder
\end{Verbatim}

In PowerShell:

\begin{Verbatim}[fontsize=\small]
python scripts/build_textbook.py
python scripts/build_provenance_pdf.py provenance/DIRAC16COMPLEX_TEXTBOOK.md `
    --developer-layout --number-sections-from-zero
\end{Verbatim}

In Bash:

\begin{Verbatim}[fontsize=\small]
python scripts/build_textbook.py
python scripts/build_provenance_pdf.py provenance/DIRAC16COMPLEX_TEXTBOOK.md \
    --developer-layout --number-sections-from-zero
\end{Verbatim}

The option \texttt{--develo\allowbreak{}per-layo\allowbreak{}ut} chooses the layout of the student guides (smaller tables with ragged right margins, code that may break inside long names), and \texttt{--number\allowbreak{}-section\allowbreak{}s-from-z\allowbreak{}ero} makes LaTeX number the chapters from 0, as the book does; without it every chapter would be printed with a number one too high (the header of \texttt{scripts/\allowbreak{}build\_di\allowbreak{}ssertati\allowbreak{}on\_tex.p\allowbreak{}y}).

\textbf{What the assembler checks.} Before it writes anything, the assembler runs the checks listed in its header and prints each as \texttt{check\_<n\allowbreak{}ame>=tru\allowbreak{}e|false}: the chapter file names and numbers (00 to 20, no gaps, no duplicates), UTF-8 text with line feeds only, exactly one heading \texttt{\#\# N. Ti\allowbreak{}tle} per chapter and headings \texttt{\#\#\# N.M \allowbreak{}Title} numbered 1, 2, 3, ... in order, no section number twice, every internal reference of the form "Chapter N" or "Section N.M" pointing to an existing chapter or section, and the Markdown convertible by the LaTeX converter (tables with equal numbers of cells, code lines of at most 89 characters, figures that exist). The header lists fourteen checks; a run makes twelve of them with \texttt{\detokenize{--check}}, thirteen when it also writes the output or compares it, and ten with \texttt{--check \allowbreak{}--allow-\allowbreak{}missing}. Only when all of them pass does it write the output; it then prints the output path, its size and its sha256, and ends with \texttt{textbook\allowbreak{}\_assembl\allowbreak{}y=OK}. The output depends only on the chapter files and on the assembler, so two runs give the same bytes. Three options change what it does:

\begin{itemize}
\item \texttt{\detokenize{--check}} runs the checks and writes nothing;
\item \texttt{--check \allowbreak{}--allow-\allowbreak{}missing} accepts a book whose planned chapters are not all written yet: it reports each missing chapter, reports the references into missing chapters as "pending" instead of failing, and (without \texttt{\detokenize{--check}}) writes a draft into \texttt{build/te\allowbreak{}xtbook/} instead of \texttt{provenan\allowbreak{}ce/};
\item \texttt{--check \allowbreak{}--verify\allowbreak{}-output} also requires the existing \texttt{provenan\allowbreak{}ce/DIRAC\allowbreak{}16COMPLE\allowbreak{}X\_TEXTBO\allowbreak{}OK.md} to equal the assembly byte for byte.
\end{itemize}

For this edition the assembler ran from the fresh clone of commit \texttt{\detokenize{1f2dd69}}, which contains the 21 chapter files in their draft form, with \texttt{\detokenize{--check}} in about a second; the run ended with \texttt{check\_co\allowbreak{}unt=12}, \texttt{failed\_c\allowbreak{}heck\_cou\allowbreak{}nt=0} and \texttt{textbook\allowbreak{}\_assembl\allowbreak{}y=OK}. Without \texttt{\detokenize{--check}} it wrote \texttt{provenan\allowbreak{}ce/DIRAC\allowbreak{}16COMPLE\allowbreak{}X\_TEXTBO\allowbreak{}OK.md} (2101864 bytes; 13 checks, all true).

\textbf{What the PDF builder checks.} For a Markdown file \texttt{\detokenize{D/X.md}} the builder (its header lists the steps) converts the file twice, into \texttt{\detokenize{D/X.tex}} and a second copy, with two different settings of Python's hash seed, and requires the two to be byte-identical; runs pdflatex three times on each copy, in two fresh folders below \texttt{\detokenize{build/X/}}; searches both final logs for warnings (lines containing \texttt{LaTeX Wa\allowbreak{}rning}, a package warning, \texttt{\detokenize{Overfull}}, \texttt{Underful\allowbreak{}l} or \texttt{Undefine\allowbreak{}d contro\allowbreak{}l sequen\allowbreak{}ce}, or beginning with the error mark \texttt{\detokenize{!}}) and fails on any; requires the two PDFs to be byte-identical and to have US-letter pages; and, in its default \textbf{verify mode}, compares the page count and the sha256 of the PDF with the edition registered for this file in \texttt{provenan\allowbreak{}ce/pdf-s\allowbreak{}pecifica\allowbreak{}tions.js\allowbreak{}on}. Only then does it copy the PDF to \texttt{\detokenize{D/X.pdf}} and print \texttt{provenan\allowbreak{}ce\_pdf=O\allowbreak{}K}. The PDF is reproducible because the LaTeX preamble written by the converter switches off everything that would make two runs differ: the date, the random file identifier and the compression of the PDF objects.

\textbf{Registering an edition.} After an edit of the book the maintainers rebuild it with the option \texttt{--regist\allowbreak{}er}, which writes the new page count and sha256 into the registry instead of comparing with it, and then update the sha256 pins of the book's publication test (\texttt{tests/te\allowbreak{}st\_d16c\_\allowbreak{}textbook\allowbreak{}\_publica\allowbreak{}tion.py}, planned by \texttt{handoff/\allowbreak{}specs/TE\allowbreak{}XTBOOK\_S\allowbreak{}PEC.md}, §2). A reader never needs \texttt{--regist\allowbreak{}er}: in verify mode a rebuilt book must reproduce the registered edition byte for byte, and that is the check. Until the textbook edition is registered, the verify mode reports its registry checks as failed while all other checks can pass.

For this edition the same command was run in that clone on the draft just assembled. It took 32 seconds and produced a PDF of 656 pages; both pdflatex logs were free of warnings and the two PDFs were byte-identical. As expected for a file without a registered edition (commit \texttt{\detokenize{1f2dd69}} registers none for the book), it ended with

\begin{Verbatim}[fontsize=\small]
check_count=14
failed_check_count=5
failed_checks=editionRegistered,registeredPath,registeredPageCount,registeredSha256,...
\end{Verbatim}

(the fifth failed check, \texttt{provenan\allowbreak{}cePdfCop\allowbreak{}y}, means that the PDF is not copied to its final place when a check fails), and with the exit code 1. The page count and the sha256 of the book you read are those of its registered edition, not of this draft.

\subsection{A complete session from a fresh clone}

The commands below collect, in order, the gates and checks of this appendix. They assume the tools of Section 19.2 and a short folder path (Section 19.4); to reproduce exactly the results stated below, add \texttt{git chec\allowbreak{}kout 1f2\allowbreak{}dd69} after the second line. PowerShell:

\begin{Verbatim}[fontsize=\small]
git clone https://github.com/once-ere/Dirac_claude.git
cd Dirac_claude
python -m pip install -r requirements-stage3.txt
python -m pip install nbconvert==7.16.6
.\scripts\setup_solver.ps1 -Platform win11
pwsh -NoProfile -File scripts/verify_stage1_arbitrary_field.ps1
pwsh -NoProfile -File scripts/verify_stage2_primordial_field.ps1
pwsh -NoProfile -File scripts/verify_stage3_dark_sector.ps1
pwsh -NoProfile -File scripts/verify_stage4_kohn_sham.ps1 -DryRun
python scripts/build_textbook.py --check
git status --short
\end{Verbatim}

Bash (Git Bash, macOS, Linux):

\begin{Verbatim}[fontsize=\small]
git clone https://github.com/once-ere/Dirac_claude.git
cd Dirac_claude
python -m pip install -r requirements-stage3.txt
python -m pip install nbconvert==7.16.6
bash scripts/setup_solver.sh win11
bash scripts/verify_stage1_arbitrary_field.sh
bash scripts/verify_stage2_primordial_field.sh
bash scripts/verify_stage3_dark_sector.sh
bash scripts/verify_stage4_kohn_sham.sh --dry-run
python scripts/build_textbook.py --check
git status --short
\end{Verbatim}

Then run the partial Stage-4 gate and the single programs of Sections 19.9 and 19.10 for Stages 4 and 5 and the matter--antimatter analysis. Expected results at commit \texttt{\detokenize{1f2dd69}}: the Stage-1 and Stage-2 gates end with their OK lines after about 18 and 3 minutes; the Stage-3 gate runs for about 40 minutes and stops in its step 31 with \texttt{stage3\_d\allowbreak{}ark\_sect\allowbreak{}or\_verif\allowbreak{}ication=\allowbreak{}FAILED}, because of the four failing unit tests of Section 19.11 (at the commits of its committed logs it ended with OK); the Stage-4 dry run ends with \texttt{\detokenize{DRY-RUN}}; the assembler ends with \texttt{textbook\allowbreak{}\_assembl\allowbreak{}y=OK}; and the last command prints nothing, because the gates write their work into the git-ignored folder \texttt{\detokenize{build/}}, and every committed file that a gate rewrites in place (the Stage-2 gate rewrites its reports and its PDF) is rewritten with the same bytes. The full Stage-4 gate, which would take about 4 hours, cannot pass yet (Section 19.9), and Stage 5 has no gate yet (Section 19.10). At a later commit, \texttt{git diff\allowbreak{} --stat \allowbreak{}1f2dd69 \allowbreak{}HEAD} lists every file that has changed since, so that you can see which of these results may differ.

\subsection{What we proved and what we assumed}

This appendix proves no theorem. It records \textbf{measurements of reproducibility}, made for this edition from fresh public clones of commit \texttt{\detokenize{1f2dd69}} on the test computer of Section 19.2, or quoted from a committed log or script header that is named where it is used. The measurements of this edition's own runs (the wall times, the test counts and the printed lines) come from runs whose logs are not committed; a reader checks them by repeating the commands. The measurements are: the Bash twins of the Stage-1 and Stage-2 gates ended with their OK lines (after 1095 and 165 seconds); the Bash twin of the Stage-3 gate passed its steps 00 to 30 and failed in its step 31, the unit tests; the Stage-2 verifiers, the Stage-3 programs and checkers, steps 01 to 13 of the Stage-4 gate in both shells and its \texttt{\detokenize{spectrum}} subcommand, the six exact programs of Stage 5 and the matter--antimatter analysis, and the two PDFs of Sections 19.10 and 19.12 were rerun, and every file that should be byte-identical was byte-identical once the clone had a short path, while the files that differ (the three Stage-1 files of the public-clone mode and the Stage-2 Python report written into a scratch folder) differ only in the entries explained in Sections 19.6 and 19.7; and the complete unit-test suite ran 579 tests, of which 4 failed and 2 were skipped (Section 19.11). It also records, from the committed files and these runs, the \textbf{state of the stages}: Stages 1 and 2 are complete, and the Bash twins of their gates passed at this commit; Stage 3 is complete and its gate passed at earlier commits, but at this commit its unit-test step fails because of four tests of Stage 4; the Stage-4 gate cannot pass yet, for the four reasons of Section 19.9, among them one false check of the cross-checker; Stage 5 has its exact theory, with its Fock-level result cited as provisional until a correction of its reading and the Stage-5 gate (Section 19.10), but only part of its Kohn--Sham numerics, no documents and no gate; and the matter--antimatter analysis has its two reports, its registered document and 49 passing tests.

It \textbf{assumes}: that the tools behave as their versions say (a different version of Python's packages, of pdflatex or of the Rust compiler can change the last digits or bytes of some files without changing any check, as Sections 19.2 and 19.4 explain); that the \texttt{\detokenize{win11}} engine is used for every byte comparison; that the reader's computer has the processor feature FMA (Section 10.10); that on Windows the clone lies in a folder with a short path (Section 19.4); and that the wall times, measured on a busy computer with several of these runs at the same time, are upper values for that computer and only indications for another. A byte-identical rerun shows that a program reproduces its committed output; it does not show that the output is correct. Correctness rests on the derivations of the chapters and on the independent checks named there.

\subsection{Exercises}

\textbf{Exercise 19.1.} Rewrite this Bash command for PowerShell:

\begin{Verbatim}[fontsize=\small]
export PYTHONUTF8=1
python scripts/check_dirac16complex_exp3.py \
    --output build/textbook
\end{Verbatim}

\textbf{Exercise 19.2.} A run of the Stage-3 gate ends with

\begin{Verbatim}[fontsize=\small]
stage3_failed_step=stage3-09-compare-outputs
stage3_failed_log=build/logs/stage3-09-compare-outputs-bash.log
stage3_dark_sector_verification=FAILED
\end{Verbatim}

What did the failed step compare, what is the most likely cause when every earlier step passed, and is any number of Chapter 11 wrong because of it?

\textbf{Exercise 19.3.} Use the one-command counter of Section 19.3 on the five Stage-1 reports of Section 19.6 and add the results. Which number of Chapter 0 do you obtain?

\textbf{Exercise 19.4.} Why does the Stage-1 gate write its reports into \texttt{build/st\allowbreak{}age1/} when the folder \texttt{dirac-ma\allowbreak{}in/} is missing, instead of into the committed folder as in the other mode?

\textbf{Exercise 19.5.} The Python algebra checker of Section 19.5 printed \texttt{check\_co\allowbreak{}unt=20}, the committed report has 21 checks, and the EXP-5 checker of Section 19.8 printed \texttt{check\_co\allowbreak{}unt=19} against 21. Name the missing checks in each case and say what has to be added to the command line to obtain them.

\textbf{Exercise 19.6.} Compute the sha256 of \texttt{provenan\allowbreak{}ce/DIRAC\allowbreak{}16COMPLE\allowbreak{}X\_MATTER\allowbreak{}\_ANTIMAT\allowbreak{}TER.pdf} in both shells and compare it with the value registered for the edition \texttt{dirac16c\allowbreak{}omplex-m\allowbreak{}atter-an\allowbreak{}timatter} in \texttt{provenan\allowbreak{}ce/pdf-s\allowbreak{}pecifica\allowbreak{}tions.js\allowbreak{}on} of the same clone.

\textbf{Exercise 19.7.} Add the step times of the Stage-1 run in the table of Section 19.6, counting the steps that took under a second as 0, and compare the sum with the total of 1095 seconds. Where did the rest of the time go?

\textbf{Exercise 19.8.} (a) The false Stage-4 check \texttt{canonica\allowbreak{}l\_eigenv\allowbreak{}alues} compares a difference of $2.195\times10^{-6}\,m$ with a tolerance of $1.054\times10^{-6}\,m$. Compute the ratio. (b) List the four reasons why the Stage-4 gate cannot end with OK today, and for each say what would have to be done.

\textbf{Exercise 19.9.} You changed one sentence of Chapter 7 and want to know, within a second and without building a PDF, whether every "Section N.M" reference of the book still resolves. Which command do you run, and which line of its output answers the question?

\textbf{Exercise 19.10.} A friend clones the repository, checks out the older commit \texttt{\detokenize{4cd47fe}} (\texttt{git chec\allowbreak{}kout 4cd\allowbreak{}47fe}, Section 19.4), reruns the matter--antimatter checker into a scratch folder as in Section 19.10 and reports: "77 checks, all true, but the committed report has 75: the repository is broken." Explain what the two numbers mean and whether the friend's conclusion is right. What does the same rerun give at commit \texttt{\detokenize{1f2dd69}}?

\textbf{Exercise 19.11.} (a) The hexadecimal number \texttt{\detokenize{ff}} is the largest number with two hexadecimal digits. Which decimal number is it? (b) How many different values can a sha256 have? Write the answer as a power of 2, using that one hexadecimal digit carries exactly as much information as four binary digits.

\subsection{Answers to the exercises}

\textbf{Answer 19.1.} The environment variable is set with \texttt{\detokenize{$env:}}, and the backslash at the end of a continued line becomes a backtick:

\begin{Verbatim}[fontsize=\small]
$env:PYTHONUTF8 = "1"
python scripts/check_dirac16complex_exp3.py `
    --output build/textbook
\end{Verbatim}

(The command is short enough to be written on one line as well.) Before it, run EXP-3 into the same folder, for example \texttt{\& \$bin e\allowbreak{}xp3 --ou\allowbreak{}tput bui\allowbreak{}ld/textb\allowbreak{}ook} with \texttt{\detokenize{$bin}} from Section 19.8, and the EXP-3 analysis \texttt{python s\allowbreak{}cripts/a\allowbreak{}nalyze\_d\allowbreak{}irac16co\allowbreak{}mplex\_ex\allowbreak{}p3.py --\allowbreak{}output b\allowbreak{}uild/tex\allowbreak{}tbook}, because the checker reads their files (student guide, §9.3). In this order the three commands printed \texttt{\detokenize{SUCCESS}}, \texttt{check\_co\allowbreak{}unt=10} and \texttt{check\_co\allowbreak{}unt=30} for this edition.

\textbf{Answer 19.2.} Step 09 compares the 62 output files of the five experiments, from the two complete runs of steps 07 and 08, with each other and with the committed files, byte for byte; its log names every file that differs. Steps 01 and 02 have already made sure that the pinned \texttt{\detokenize{win11}} engine is installed and unmodified, and steps 07 and 08 ended with \texttt{\detokenize{SUCCESS}}, so the program passed all its self-checks. The pinned engine is known to reproduce the committed files byte for byte on Windows 11 and on Ubuntu 24.04 (Section 10.10). A difference therefore points to one of three things: a system on which byte identity has not been established (for example macOS), a build that did not use the FMA setting of \texttt{.cargo/c\allowbreak{}onfig.to\allowbreak{}ml} (Section 10.10), or a committed file that was changed in the working tree (\texttt{git stat\allowbreak{}us --sho\allowbreak{}rt} shows it). No number of Chapter 11 becomes wrong through such a difference: the self-checks and the independent checkers test the results against exact solutions and conservation laws, not against bytes, and a difference in the last digits is far below every tolerance. The byte comparison tests reproducibility, not correctness.

\textbf{Answer 19.3.} The counter prints one line per run. For the five reports in the folder \texttt{artifact\allowbreak{}s/dirac1\allowbreak{}6complex\allowbreak{}/arbitra\allowbreak{}ry-field\allowbreak{}/} it printed:

\begingroup
\small
\begin{longtable}{@{}>{\raggedright\arraybackslash}p{0.400\linewidth}>{\raggedright\arraybackslash}p{0.400\linewidth}@{}}
\toprule
\textbf{Report given to the counter} & \textbf{Printed line} \\
\midrule
\endfirsthead
\toprule
\textbf{Report given to the counter} & \textbf{Printed line} \\
\midrule
\endhead
\texttt{wolfram-\allowbreak{}algebra-\allowbreak{}report.j\allowbreak{}son} & \texttt{\detokenize{21 21}} \\
\texttt{wolfram-\allowbreak{}geometry\allowbreak{}-report.\allowbreak{}json} & \texttt{\detokenize{43 43}} \\
\texttt{python-a\allowbreak{}lgebra-r\allowbreak{}eport.js\allowbreak{}on} & \texttt{\detokenize{21 21}} \\
\texttt{python-g\allowbreak{}eometry-\allowbreak{}report.j\allowbreak{}son} & \texttt{\detokenize{52 52}} \\
\texttt{grassman\allowbreak{}n-demo-r\allowbreak{}eport.js\allowbreak{}on} & \texttt{\detokenize{16 16}} \\
\bottomrule
\end{longtable}
\endgroup

The sum is $21+43+21+52+16=153$: the 153 of 153 exact Stage-1 checks of the table in Section 0.5.

\textbf{Answer 19.4.} Without \texttt{dirac-ma\allowbreak{}in/} the verifiers cannot run the comparisons with the three dirac-main files and record them as "not-run". Their reports therefore differ from the committed ones in those entries and in the hashes that depend on them. Writing them over the committed reports would change committed evidence; writing them into \texttt{build/st\allowbreak{}age1/} keeps the committed files unchanged, and step 13 then proves that the regenerated reports agree with the committed ones everywhere except in exactly the entries that the missing folder explains.

\textbf{Answer 19.5.} The algebra checker lacks \texttt{ALG\_wolf\allowbreak{}ramAgree\allowbreak{}ment}, its comparison with the Wolfram algebra report; the empty option \texttt{--wolfra\allowbreak{}m-report\allowbreak{}=} switches it off, and the option \texttt{--wolfra\allowbreak{}m-report} followed by the path of a Wolfram report switches it on (step 07 of the Stage-1 gate). The EXP-5 checker lacks \texttt{repeatBy\allowbreak{}teIdenti\allowbreak{}ty} and \texttt{refinedC\allowbreak{}onvergen\allowbreak{}ce}. With the option \texttt{\detokenize{--repeat}} followed by a folder below \texttt{\detokenize{build/}}, the checker runs the program once more into that folder and requires every file to be byte-identical; with \texttt{--refine\allowbreak{}d} followed by another folder, it runs the program with ten times tighter tolerances and requires the errors to shrink (student guide, §9.1). For example \texttt{python s\allowbreak{}cripts/c\allowbreak{}heck\_dir\allowbreak{}ac16comp\allowbreak{}lex\_exp5\allowbreak{}.py --ou\allowbreak{}tput bui\allowbreak{}ld/textb\allowbreak{}ook --re\allowbreak{}peat bui\allowbreak{}ld/repea\allowbreak{}t --refi\allowbreak{}ned buil\allowbreak{}d/refine\allowbreak{}d}, written with a line continuation in either shell, printed \texttt{check\_re\allowbreak{}peatByte\allowbreak{}Identity\allowbreak{}=true}, \texttt{check\_re\allowbreak{}finedCon\allowbreak{}vergence\allowbreak{}=true} and \texttt{check\_co\allowbreak{}unt=21} after 8 seconds for this edition.

\textbf{Answer 19.6.} In PowerShell the command

\begin{Verbatim}[fontsize=\small]
(Get-FileHash -Algorithm SHA256 provenance/DIRAC16COMPLEX_MATTER_ANTIMATTER.pdf).Hash
\end{Verbatim}

prints the fingerprint in upper-case letters (without \texttt{(...).Ha\allowbreak{}sh}, \texttt{Get-File\allowbreak{}Hash} prints a small table with the same value in its column \texttt{\detokenize{Hash}}):

\begin{Verbatim}[fontsize=\small]
F4D6CB5102849614D7E643D868963F958713C99AE8A890CEE55715AEFB4AAD6C
\end{Verbatim}

In Bash, \texttt{sha256su\allowbreak{}m proven\allowbreak{}ance/DIR\allowbreak{}AC16COMP\allowbreak{}LEX\_MATT\allowbreak{}ER\_ANTIM\allowbreak{}ATTER.pd\allowbreak{}f} prints the same 64 digits in lower case, followed by the file name (Git Bash puts a \texttt{\detokenize{*}} in front of the name; on macOS use \texttt{shasum -\allowbreak{}a 256} instead of \texttt{sha256su\allowbreak{}m}):

\begin{Verbatim}[fontsize=\small]
f4d6cb5102849614d7e643d868963f958713c99ae8a890cee55715aefb4aad6c
\end{Verbatim}

The registry entry holds this lower-case value and 36 pages: the same number, since upper-case and lower-case letters denote the same hexadecimal digits (Section 19.1). These are the values at commit \texttt{\detokenize{1f2dd69}}; a later edition of the document has another sha256, and the registry of the same clone always holds the one that the PDF of that clone must have.

\textbf{Answer 19.7.} $0+3+23+409+405+57+2+54+94+0+16+1+1=1065$ seconds. The remaining 30 seconds lie between the steps: the gate looks up the tools, removes and creates \texttt{build/st\allowbreak{}age1/}, starts a new process for every step, audits the five reports with short Python programs and computes the eight sha256 values.

\textbf{Answer 19.8.} (a) $2.195/1.054=2.08$, the ratio recorded in the measurement \texttt{canonica\allowbreak{}l\_eigenv\allowbreak{}alues\_de\allowbreak{}tail}. (b) First, the Stage-4 document has no LaTeX file, no PDF and no registered edition, and the Stage-4 student guide does not exist (steps 00, 32 and 33): both documents have to be finished, built and registered. Second, \texttt{kohn-sha\allowbreak{}m-summar\allowbreak{}y.json} does not exist (step 23): it has to be generated from the final results and committed. Third, the fresh cross-check report must have all checks true (step 21), and \texttt{canonica\allowbreak{}l\_eigenv\allowbreak{}alues} is false for one level of the smeared $N=1016$ ensemble: the disagreement between the two solvers has to be resolved, by a finer computation that brings them together or by a justified tolerance for this case. Fourth, the Jupyter notebook has to be executed again, so that its committed report is current and all its checks are true, which the unit tests of step 35 require. Only after that can the gate run through step 14, the reproduction of the complete Rust tree (about 3 hours), to the OK line.

\textbf{Answer 19.9.} \texttt{python s\allowbreak{}cripts/b\allowbreak{}uild\_tex\allowbreak{}tbook.py\allowbreak{} --check} (add \texttt{--allow-\allowbreak{}missing} while chapters are missing). The line \texttt{check\_cr\allowbreak{}ossRefer\allowbreak{}encesRes\allowbreak{}olve=tru\allowbreak{}e} answers it; any unresolved reference is printed before it as a line beginning with \texttt{\detokenize{problem=}}. The option \texttt{--list-r\allowbreak{}eference\allowbreak{}s} prints every reference with its classification.

\textbf{Answer 19.10.} At commit \texttt{\detokenize{4cd47fe}} the committed report had been written by an earlier version of the checker, with 75 checks: the report records for its checker the sha256 \texttt{20914546\allowbreak{}...}, while the checker file of that commit has the sha256 \texttt{ef1a00dc\allowbreak{}...}. The checker of that commit has two more checks, \texttt{MA\_M2\_cp\allowbreak{}ScopeInC\allowbreak{}urvedFie\allowbreak{}lds} and \texttt{MA\_M4\_im\allowbreak{}ageField\allowbreak{}FockMode\allowbreak{}l}, and all 77 are true. Nothing is broken in the sense of a false check: every committed check is reproduced as true, and two more checks pass. What the friend has found is that the committed report of that commit was not the newest one; regenerating it changes its check count and its recorded hashes, not any conclusion. The report was regenerated in commit \texttt{\detokenize{27794e8}} and has not changed since: at commit \texttt{\detokenize{1f2dd69}} the committed report has 77 checks, and the rerun of Section 19.10 wrote a byte-identical file.

\textbf{Answer 19.11.} (a) $\mathtt{ff}=15\cdot16+15=255$. (b) Each hexadecimal digit takes one of $16=2^4$ values, so 64 digits take $16^{64}=(2^4)^{64}=2^{256}$ values, which is where the name sha256 comes from. Two different files therefore have the same sha256 only by an accident whose probability is far below anything that could matter in practice.

\section{Glossary and index of verifier checks}

This last chapter has two parts. The glossary (Sections 20.2 to 20.8) lists the technical terms of the book in alphabetical order, each with a short definition and the chapter in which the term is introduced; symbols and notation follow in their own section. The index of verifier checks (Sections 20.9 to 20.15) lists every machine check that the book names, with the committed report that contains it, its value there and the chapters that cite it. The chapter ends, like every chapter, with what it proves and assumes, exercises and answers.

\subsection{How to use this chapter}

\textbf{The glossary.} Each entry has the form "\textbf{term} (Chapter N)", followed by one or two sentences. The chapter is the one in which the term is first defined or first explained; where two chapters are named, the first gives a short definition and the second the full treatment, or the term has two uses. A term that returns later is not listed again under the later chapters. The definitions are short reminders, not replacements for the chapters: the full definition, the worked examples and the proofs are in the chapter named. The entries are sorted letter by letter, as if each term were written as one word: spaces, hyphens, dashes, apostrophes, brackets, upper and lower case, accents and the markup of mathematical symbols are ignored (so "line element" comes after "linearly independent", and "$q$ vector" is sorted as "qvector"). Digits come before letters ("dirac16complex" before "Dirac adjoint", "Q1" before "$q$ vector"), except that a term beginning with a digit is sorted as if that digit were written as a word ("3-space" under T). An entry that names several terms is sorted by the first of them, the part before the first comma (so "double, double precision" comes before "double cover"). A term that is usually abbreviated is listed under its abbreviation with a cross-reference ("BDF: see backward differentiation formula"). Where a term has a special meaning in this project that differs from its everyday meaning (for example "gate", "check", "good sector"), the entry gives the project's meaning.

\textbf{The index of checks.} A check is a named statement that a verifier has decided to be true or false and recorded in a report (Section 0.8). Section 20.9 explains how the index was made and verified.

\subsection{Glossary: A to C}

\begin{itemize}
\item \textbf{absolute tolerance, atol} (Chapter 10). The error that CVODE accepts in a component of the state whose size is near zero; together with the relative tolerance it sets the weight of each component in the error test.
\item \textbf{action} (Chapter 5). The integral of the Lagrangian over time (for a particle) or of the Lagrangian density over spacetime (for a field); the equations of motion are the conditions under which it does not change to first order.
\item \textbf{Adams methods} (Chapter 10). Linear multistep methods that integrate a polynomial through past values of the right-hand side; Adams--Bashforth is explicit, Adams--Moulton implicit. CVODE's Adams--Moulton family was used for EXP-2 to EXP-5.
\item \textbf{adiabatic dressing} (Chapter 11). The small admixture of the other energy level that a slowly changing mode carries at every instant and that disappears when the change stops; it is not a created particle.
\item \textbf{adiabatic theorem, adiabatic expansion} (Chapter 11). A state that starts in a level of a slowly changing Hamiltonian stays in that level; the expansion in the rate of change gives the created amplitude when the change is not slow.
\item \textbf{adjoint} (Chapter 8). For an operator $O$, the operator $O^\dagger$ with $\langle f,Oh\rangle=\langle O^\dagger f,h\rangle$; for the standard inner product it is the conjugate transpose. See also Hermitian conjugate, Krein adjoint, Dirac adjoint.
\item \textbf{algebra} (Chapter 3). A vector space with a bilinear product; it may have a unit and may be associative or commutative.
\item \textbf{alternative} (Chapter 3). An algebra in which $x(xy)=(xx)y$ and $(yx)x=y(xx)$ for all $x,y$; the split octonions are alternative although not associative.
\item \textbf{amplification factor} (Chapter 10). The number $R(z)$, $z=h\lambda$, by which a one-step method multiplies the solution of $y'=\lambda y$ in each step; its modulus decides stability.
\item \textbf{Anderson mixing, Pulay mixing} (Chapter 12). A way to find the self-consistent density that combines several previous inputs and their residuals so as to make the residual smallest; the Rust Kohn--Sham solver uses it.
\item \textbf{annihilate, annihilation operator} (Chapter 8). The operator $b$ (or $a_p$) that removes one fermion from a mode and gives zero on an empty mode; in particle physics a particle and its antiparticle also annihilate each other (Chapter 17).
\item \textbf{anticommutation relations} (Chapter 12). The rules $\{a_p,a_q^\dagger\}=\delta_{pq}$ and $\{a_p,a_q\}=\{a_p^\dagger,a_q^\dagger\}=0$ of fermion creation and annihilation operators.
\item \textbf{anticommutator} (Chapter 1). $\{A,B\}=AB+BA$; two matrices anticommute when it is zero.
\item \textbf{anti-Hermitian} (Chapter 1). A square matrix with $A^\dagger=-A$; a real antisymmetric matrix is anti-Hermitian, and $i$ times an anti-Hermitian matrix is Hermitian.
\item \textbf{antilinear} (Chapter 17). A map $A$ with $A(cv)=c^\ast A(v)$ for complex numbers $c$, such as complex conjugation; an antiunitary map is an antilinear map that preserves the modulus of inner products.
\item \textbf{antiparticle} (Chapter 8). In the good-sector Fock space, a hole in the filled Dirac sea: it has positive energy and the opposite U(1) charge of a particle.
\item \textbf{antisymmetric matrix} (Chapter 1). A square matrix with $A^T=-A$; its diagonal is zero.
\item \textbf{associative} (Chapter 1). A product with $(AB)C=A(BC)$; matrix multiplication is associative, the product of the split octonions is not.
\item \textbf{associator} (Chapter 3). $(xy)z-x(yz)$; it measures the failure of associativity, for example $(e_1e_2)e_4-e_1(e_2e_4)=2e_7$ for the split octonions.
\item \textbf{ASSUMED} (Chapter 0). The status of a statement that is a starting point and not derived: a convention, a physical input or an approximation.
\item \textbf{atol}: see absolute tolerance.
\item \textbf{atomic units} (Chapter 12). Units in which $\hbar$, the electron mass and the electron charge are 1.
\item \textbf{aufbau rule} (Chapter 12). Occupy the lowest available one-particle levels (German "building up").
\item \textbf{average of a representation, $K_X$} (Chapter 3). The sum $\sum_A\gamma'_AX\gamma_A^{-1}$ over all 256 monomials; for a suitable matrix $X$ it gives a nonzero intertwiner between two sets of gamma matrices.
\item \textbf{back-reaction} (Chapter 18). The response of the metric to the matter in it, computed self-consistently instead of keeping the gravitational field fixed; not computed anywhere in the repository.
\item \textbf{backward differentiation formula, BDF} (Chapter 10). An implicit multistep method that differentiates a polynomial through the unknown and past values of the solution; stable for stiff problems. Used with Newton iteration for EXP-1.
\item \textbf{bag condition} (Chapter 13). The boundary condition at the tip cutoff $y=-L$ of the Kohn--Sham problem, $\chi_2(-L)=0$, the member $\theta=0$ of a family $(1-Q(\theta))\chi(-L)=0$ that kills the current.
\item \textbf{baryogenesis} (Chapter 17). Any process in the early universe that creates the observed excess of baryons over antibaryons; it needs the three conditions of Sakharov.
\item \textbf{baryon, baryon number} (Chapter 17). Baryons are particles made of three quarks, such as protons and neutrons; the baryon number counts baryons minus antibaryons and is conserved by every process observed so far.
\item \textbf{baryon-to-photon ratio} (Chapter 17). $\eta$, the number of baryons per photon of the cosmic microwave background, $\eta=6.12\times10^{-10}$ from the quoted measurements; the measured size of the matter--antimatter asymmetry, which this theory does not predict.
\item \textbf{Bash}: see command shell.
\item \textbf{basis} (Chapter 1). A list of linearly independent vectors of which every vector is a combination; all bases of $\mathbb R^n$ have $n$ vectors.
\item \textbf{BDF}: see backward differentiation formula.
\item \textbf{Bianchi identities} (Chapter 4). The first, $R^\rho{}_{\sigma\mu\nu}+R^\rho{}_{\mu\nu\sigma}+R^\rho{}_{\nu\sigma\mu}=0$, and the second, the vanishing cyclic sum of covariant derivatives of the Riemann tensor; the contracted form gives $\nabla_\mu G^\mu{}_\nu=0$.
\item \textbf{big-bang nucleosynthesis} (Chapter 17). The formation of the light nuclei in the first minutes of the universe; the observed abundances fix the baryon-to-photon ratio.
\item \textbf{bilinear} (Chapters 3 and 5). Linear in each of two arguments; in field theory a bilinear is an expression $\Psi^\dagger M\Psi$ or $\bar\Psi M\Psi$ quadratic in the field.
\item \textbf{binomial coefficient, factorial, binomial theorem} (Chapter 2). $\binom nk=\frac{n!}{k!(n-k)!}$, the number of subsets with $k$ elements of a set with $n$ elements (with the factorial $n!=1\cdot2\cdots n$, $0!=1$); the binomial theorem is $(x+y)^n=\sum_k\binom nkx^ky^{n-k}$.
\item \textbf{bisection} (Chapter 10). A root finder that halves a bracketing interval in every step, keeping the half on which the function changes sign.
\item \textbf{block matrix} (Chapter 1). A matrix cut into smaller matrices (blocks), multiplied block by block with the order of factors kept.
\item \textbf{block type} (Chapter 13). The label $j=\pm1$ (the eigenvalue of $J=\gamma^0\gamma^1\gamma^4$) of the eight $2\times2$ blocks of the reduced Kohn--Sham equation; the two types carry opposite spectra of $h-v$.
\item \textbf{Bogoliubov method, Bogoliubov coefficients} (Chapter 18; the computation is that of Chapter 11). Following each mode through a changing background and measuring the weight $|\beta_k|^2$ with which a mode that starts in its positive-energy level ends on the negative-energy level; the $\beta_k$ are the Bogoliubov coefficients. It is the only creation process computed in the repository, for particles in an expanding 3-space, not for universes.
\item \textbf{bonding orbital} (Chapter 12). The lower orbital $\varphi_b=(\mathrm L+\mathrm R)/\sqrt2$ of the two-site model, the symmetric combination of the two site orbitals.
\item \textbf{boost} (Chapter 2). A spin transformation $\exp(\theta S^{ab})$ with one space-like and one time-like index; it mixes the two directions hyperbolically. A boost pair of frame indices is such a pair (Chapter 4).
\item \textbf{bosons} (Chapter 12). Particles whose many-body wave function is unchanged when two of them are exchanged.
\item \textbf{bounce} (Chapter 11). The point of an EXP-3 run where the expansion rate reaches zero ($E^2=0$); it lies where the mean field is no longer valid.
\item \textbf{bracket (of a root)} (Chapter 10). An interval at whose ends a function has opposite signs, so that it contains a zero; the Stage-4 solver first brackets each level and then refines it.
\item \textbf{brane} (Chapter 9). A surface that carries energy and momentum of its own; here the surface $y=0$ where two mirror copies of the static primordial field are glued.
\item \textbf{brane band} (Chapter 13). The Kohn--Sham levels $\varepsilon=\pm ck$ (to first order in the momentum $k$) into which the eight brane zero modes split at nonzero momentum; for $N=112$ and $N=1016$ the particles fill it up to a closed shell.
\item \textbf{brane fraction} (Chapter 13). The fraction of the norm $\int\chi^\dagger\chi\,dy$ of an orbital, or of the particles of a state, within $|y|<1$, that is within the distance $1/H$ of the brane ($H=1$).
\item \textbf{byte-identical} (Chapter 10). Two files are byte-identical when they are the same sequence of bytes; used as the test of reproducibility with a fixed program and engine.
\item \textbf{canonical energy--momentum tensor} (Chapter 5). The Noether current of translations, $\Theta^\mu{}_\nu=\frac{\partial\mathcal L}{\partial(\partial_\mu\phi_A)}\partial_\nu\phi_A-\delta^\mu{}_\nu\mathcal L$.
\item \textbf{canonical Hartree--Fock equations} (Chapter 12). $\hat F\varphi_a=\varepsilon_a\varphi_a$, the Hartree--Fock equations after the occupied orbitals have been rotated so that the matrix of Lagrange multipliers is diagonal.
\item \textbf{canonical quantization} (Chapter 8). Turning the fields into operators whose commutators or anticommutators at equal time follow from the Hamiltonian form of the Lagrangian; here with respect to the time $x_4$.
\item \textbf{canonical spin connection} (Chapter 4). The unique solution $\omega_\mu{}^a{}_b$ of the vielbein postulate, lowered with $\eta$ to the antisymmetric $\omega_{\mu ab}$; the only connection used by the project.
\item \textbf{Cartan--Dieudonné theorem} (Chapter 2). Every element of O(4,4) is a product of at most eight reflections; quoted without proof.
\item \textbf{chain rule} (Chapter 1). The derivative of a composed function is the sum over the intermediate variables of the products of partial derivatives.
\item \textbf{characteristic polynomial} (Chapter 1). $p(x)=\det(xI-A)$; its roots are the eigenvalues of $A$.
\item \textbf{charge} (Chapters 5 and 7). The conserved quantity of a continuous symmetry, $Q=\int j^4\,d^7x$ over a slice; for dirac16complex the U(1) charge $Q=\int\sqrt{|g|}\,J^4\,d^7x$ of the phase symmetry, which in the good sector counts particles minus antiparticles (Chapter 8).
\item \textbf{charge conjugation, C} (Chapters 6 and 17). A map that exchanges particles and antiparticles; for the fields of this book an antilinear map $\Psi\to M\Psi^\ast$ with a constant matrix, such as $C_0:\Psi\to\Psi^\ast$ and $C_8:\Psi\to\gamma^8\Psi^\ast$. For the commuting field $C_0$ is an exact symmetry that reverses the charge; for the Grassmann field no constant charge conjugation is a symmetry when $m\ne0$.
\item \textbf{charge matrix, $C$} (Chapter 2). $C=\gamma^0\gamma^1\gamma^2\gamma^3$, the notebook's $\sigma_{16}$; real, symmetric, $C^2=1$, signature (8,8); it defines the Dirac adjoint.
\item \textbf{charge sloshing} (Chapter 12). The failure of a plain self-consistency iteration in which the density jumps back and forth between two regions.
\item \textbf{chart} (Chapter 4). A part of a manifold together with a one-to-one assignment of coordinates.
\item \textbf{check} (Chapter 0). A named statement decided by a verifier and recorded in its report as \texttt{\detokenize{true}} or \texttt{\detokenize{false}}.
\item \textbf{checker} (Chapter 2). An independent Python program (\texttt{scripts/\allowbreak{}check\_*.\allowbreak{}py}, usually built on sympy) that recomputes the exact or numerical results of a Wolfram verifier or of a Rust program and writes its own report.
\item \textbf{chemical potential} (Chapter 12). The Lagrange multiplier $\mu$ that fixes the particle number; the change of the minimal energy per added particle.
\item \textbf{chirality, $\gamma^8$} (Chapter 2). The product $\gamma^0\gamma^1\cdots\gamma^7$; Hermitian, $(\gamma^8)^2=1$, anticommuting with every $\gamma^a$, equal to $\mathrm{diag}(-I_8,I_8)$ in the notebook basis. The chirality projectors $P_\mp=\tfrac12(1\mp\gamma^8)$ select the two halves.
\item \textbf{chirality map, T1} (Chapter 15). The map $\Psi\to\gamma^8\Psi$; it turns the theory with $(m,\lambda)$ into the one with $(-m,-\lambda)$ and reverses the sign of the Lagrangian, the current and the energy--momentum tensor.
\item \textbf{Christoffel symbols} (Chapter 4). The coefficients $\Gamma^\rho{}_{\mu\nu}$ of the Levi-Civita connection, computed from the metric and its first derivatives.
\item \textbf{Clifford algebra, $\mathrm{Cl}(p,q)$} (Chapter 2). The algebra generated by elements $\gamma^a$ with $\gamma^a\gamma^b+\gamma^b\gamma^a=2\eta^{ab}$; $\mathrm{Cl}(4,4)$ is the algebra of all real $16\times16$ matrices.
\item \textbf{Clifford picture, tensor-product picture} (Chapter 2). The construction of real $16\times16$ gammas as Kronecker products of $2\times2$ matrices, the "Clifford picture" of the reference implementation dirac-main.
\item \textbf{Clifford relation} (Chapter 2). $\gamma^a\gamma^b+\gamma^b\gamma^a=2\eta^{ab}I$.
\item \textbf{Clifford vector} (Chapter 2). A combination $v=\sum_cv_c\gamma^c$ of the gamma matrices.
\item \textbf{clone, fresh clone} (Chapter 19). A complete copy of a repository, with its whole history, made by the command \texttt{git clon\allowbreak{}e}; a fresh clone is one made just before a test, so that no file of an earlier run and no local change can enter the test.
\item \textbf{closed shell} (Chapter 13). A particle number that fills complete Kohn--Sham levels (for example $N=8$, 112, 1016); for it the Kohn--Sham gap is well defined.
\item \textbf{collisionless} (Chapter 11). A gas whose particles do not interact after the start, so that each momentum mode keeps its occupation.
\item \textbf{command shell, PowerShell, Bash} (Chapter 19). A program in which commands are typed, one per line, and run in a current folder; the book gives every command for PowerShell (the shell of Windows) and for Bash (the shell of Git Bash on Windows, of macOS and of Linux). Not to be confused with the shells of Kohn--Sham levels (Chapter 13).
\item \textbf{commit} (Chapter 19). A recorded state of all the files of a repository, with a message, named by a 40-digit hexadecimal fingerprint of its content; a short beginning of the name, such as \texttt{\detokenize{1f2dd69}}, is enough to name it. The command \texttt{git chec\allowbreak{}kout} followed by the name restores the files of that state.
\item \textbf{commutant} (Chapter 2). The set of matrices that commute with every matrix of a representation; its dimension decides irreducibility and equivalence.
\item \textbf{commutator} (Chapter 1). $[A,B]=AB-BA$; two matrices commute when it is zero.
\item \textbf{comoving} (Chapter 11). Measured in coordinates that expand with the universe; a comoving momentum $k$ corresponds to the physical momentum $k/a$.
\item \textbf{completeness relation} (Chapter 8). $\sum_s(u_s)_a(u_s)_b^\ast+\sum_s(v_s)_a(v_s)_b^\ast=\delta_{ab}$ for an orthonormal basis of positive- and negative-energy columns; it expresses that the basis spans all of $\mathbb C^{16}$.
\item \textbf{complex conjugate} (Chapter 1). $z^\ast=a-ib$ for $z=a+ib$; for a matrix, conjugation of every entry.
\item \textbf{complexification} (Chapter 6). The Lagrangian of dirac16complex00 viewed as two copies of the notebook's real Lagrangian, one for the real and one for the imaginary part of the field, coupled only through the interaction $U$.
\item \textbf{complex number} (Chapter 1). A number $a+ib$ with real $a,b$ and $i^2=-1$.
\item \textbf{components} (Chapter 1). The numbers $v_0,\dots,v_{n-1}$ of a vector (counted from 0); in Chapter 4 the numbers of a tensor in a chosen chart or frame.
\item \textbf{composition algebra} (Chapter 3). An algebra with unit and a nondegenerate quadratic norm with $N(xy)=N(x)N(y)$.
\item \textbf{COMPUTED} (Chapter 0). The status of a numerical result of a program, reproducible and checked but not a proof.
\item \textbf{conjugation (Grassmann)} (Chapter 5). The rule $F\to F^\ast$ on a Grassmann algebra; it is antilinear and reverses products, $(FG)^\ast=G^\ast F^\ast$.
\item \textbf{connection coefficients} (Chapter 4). The numbers $\Gamma^\mu{}_{\nu\lambda}$ added to partial derivatives to make the covariant derivative a tensor.
\item \textbf{conserved current} (Chapter 5). A current $j^\mu$ with $\partial_\mu j^\mu=0$ on solutions; its time component integrates to a conserved charge.
\item \textbf{conserved quantity} (Chapter 16). A number computed from the state that does not change in time; it forbids a transition exactly when its values before and after differ, so conservation laws are necessary, never sufficient, conditions for a process.
\item \textbf{constraint} (Chapters 5 and 11). A relation that the variables must satisfy at every time: a momentum that is a function of the coordinates (Chapter 5), or an Einstein equation without second time derivatives (Chapter 11).
\item \textbf{continuation in the coupling} (Chapter 13). Defining the Kohn--Sham ground state as the self-consistent solution reached from $\lambda=0$ through intermediate couplings.
\item \textbf{continuous symmetry} (Chapter 5). A family of field changes under which the Lagrangian density changes only by a total divergence.
\item \textbf{contraction} (Chapter 13). The number (times the unit operator) that an equal-time anticommutator of two field operators gives; Wick's theorem writes expectation values as sums of products of contractions.
\item \textbf{convergence (of a numerical method)} (Chapter 10). A method converges if its error tends to 0 as the step size $h$ tends to 0.
\item \textbf{correlation energy} (Chapter 12). The part of the ground-state energy that no single Slater determinant captures, $E_0-E_{HF}$.
\item \textbf{cosmic microwave background} (Chapter 17). The thermal radiation left from the early universe; its temperature and pattern measure the baryon density.
\item \textbf{cosmological constant} (Chapter 4). A constant $\Lambda$ in the Einstein equations; as a fluid it has $w=-1$ (Chapter 11).
\item \textbf{counterterm, renormalization} (Chapter 18). Terms added to a Lagrangian to absorb the infinities of perturbation theory; a theory is renormalizable if finitely many kinds suffice. The contact interaction of dirac16complex in eight dimensions is not renormalizable.
\item \textbf{covariance} (Chapter 14). The matrix $K_{ab}=\langle\Psi_a\Psi_b^\ast\rangle$ of a collection of classical fields; it is positive semidefinite for every probability distribution, while the expectation-value rule of the model makes it indefinite.
\item \textbf{covariant derivative} (Chapter 4). A derivative corrected by connection coefficients so that it maps tensors to tensors; for spinors $D_\mu=\partial_\mu+\Omega_\mu$.
\item \textbf{covector} (Chapter 4). An object with lower-index components that transform with $\partial x/\partial x'$, such as a gradient.
\item \textbf{CP, CPT} (Chapter 17). CP combines charge conjugation with a spatial reflection (parity); CPT combines charge conjugation, parity and time reversal.
\item \textbf{CPL form} (Chapter 11). The Chevallier--Polarski--Linder parametrization $w(a)=w_0+w_a(1-a)$ of a slowly changing equation of state.
\item \textbf{CPT-symmetric universe} (Chapter 17). The proposal of L. Boyle, K. Finn and N. Turok (2018) that the universe after the big bang is the CPT image of the universe before it; cited as the known published example of a universe that is symmetric as a whole but made of two oppositely asymmetric parts.
\item \textbf{creation operator} (Chapters 8 and 12). The operator $b^\dagger$ (or $a_p^\dagger$) that adds one fermion to an empty mode and gives zero on an occupied mode.
\item \textbf{critical density} (Chapter 11). $\rho_{\mathrm{crit}}=3H_0^2/\kappa_4$, the energy density of a spatially flat universe with today's Hubble rate.
\item \textbf{cross-check, cross-checker} (Chapter 13). The comparison of the Stage-4 Rust solver with the independent Python reference solver \texttt{scripts/\allowbreak{}ks\_refer\allowbreak{}ence\_sol\allowbreak{}ver.py}, made by the checker \texttt{scripts/\allowbreak{}check\_di\allowbreak{}rac16com\allowbreak{}plex\_koh\allowbreak{}n\_sham.p\allowbreak{}y} with tolerances fixed in advance.
\item \textbf{cross section} (Chapter 16). The effective target area $\sigma$ that quantum electrodynamics assigns to a nucleus for pair creation by a photon: aimed at random into an area $A$, the photon makes a pair with probability $\sigma/A$; the example of a computed creation probability.
\item \textbf{curvature} (Chapter 4). The failure of parallel transport around small loops, measured by the Riemann tensor; for spinors by $F_{\mu\nu}$.
\item \textbf{curved gamma matrices} (Chapter 4). $\gamma^\mu=e_a{}^\mu\gamma^a$, the gammas with a coordinate index, which satisfy $\gamma^\mu\gamma^\nu+\gamma^\nu\gamma^\mu=2g^{\mu\nu}$.
\item \textbf{CVODE} (Chapter 10). The initial-value solver of SUNDIALS (variable step and order, Adams and BDF families), used here in its pure-Rust translation.
\item \textbf{cyclic property} (Chapter 2). $\mathrm{tr}(AB)=\mathrm{tr}(BA)$.
\end{itemize}

\subsection{Glossary: D to F}

\begin{itemize}
\item \textbf{d16c, d16c00} (Chapter 14). Abbreviations of dirac16complex and dirac16complex00 in tables and in the names of numerical runs, such as \texttt{d16c\_m3\_\allowbreak{}L3\_N112\_\allowbreak{}lam0\_T0}.
\item \textbf{dark energy} (Chapter 11). Whatever makes the expansion of the observed universe speed up today; it needs a large negative pressure, and in the standard model of cosmology it is a cosmological constant. EXP-3 tests whether a condensate of dirac16complex can supply it (the answer of Chapter 11 is no).
\item \textbf{dark matter} (Chapter 11). Gravitating matter that cannot be seen; it clusters like ordinary matter and has almost no pressure. EXP-4 tests whether the quanta of dirac16complex behave like it (the answer of Chapter 11 is a qualified yes).
\item \textbf{DEC}: see dominant energy condition.
\item \textbf{deceleration parameter} (Chapter 11). $q=-\ddot aa/\dot a^2$, positive when the expansion slows down and negative when it speeds up.
\item \textbf{delta function, Dirac delta} (Chapters 8 and 9). The generalized function $\delta(x)$ with $\int\delta(x)f(x)\,dx=f(0)$; it describes a source concentrated on a point or, in Chapter 9, on the brane.
\item \textbf{Delta-SCF} (Chapter 12). An excitation energy obtained as the difference of two self-consistent energies, of the excited and of the ground state, $E_1-E_0$; unlike the Kohn--Sham gap it includes the relaxation of the orbitals.
\item \textbf{dense direct solver} (Chapter 10). The solution of a linear system by Gaussian elimination on the full matrix; used by CVODE's Newton iteration in EXP-1.
\item \textbf{density} (Chapter 12). The expected number of particles per unit volume at a point, $n(\mathbf r)$; for orbitals with occupations $f_a$ it is $\sum_af_a|\varphi_a(\mathbf r)|^2$. The basic variable of density functional theory.
\item \textbf{density functional theory, DFT} (Chapter 12). The method that computes the ground-state energy as a functional of the density instead of the many-body wave function; in practice through the Kohn--Sham equations.
\item \textbf{density matrix} (Chapter 14). $\rho=\sum_nf_n\,u_nu_n^\dagger$, the matrix of two-point averages $\langle\psi_a^\ast\psi_b\rangle=\rho_{ba}$; in Chapter 12 the one-body density matrix of a many-body state.
\item \textbf{density operator} (Chapter 12). The operator $\hat\rho=\sum_ip_i|\Psi_i\rangle\langle\Psi_i|$ of a statistical ensemble; at a temperature $T$ it is $e^{-(\hat H-\mu\hat N)/T}$ divided by its trace.
\item \textbf{density parameters} (Chapter 11). $\Omega_i=\rho_i/\rho_{\mathrm{crit}}$, the energy densities of the components of the universe today in units of the critical density.
\item \textbf{derived here} (Chapter 18). The mark of a short calculation that is complete as written but has not been checked by any program of the repository.
\item \textbf{determinant} (Chapter 1). The number $\det A=\sum_\pi\mathrm{sgn}(\pi)\prod_iA_{i\pi(i)}$; it is nonzero exactly when $A$ is invertible, and $\det(AB)=\det A\det B$.
\item \textbf{DFT}: see density functional theory.
\item \textbf{diagonal matrix} (Chapter 1). A square matrix whose entries off the main diagonal are zero.
\item \textbf{difference quotient} (Chapter 10). An approximate derivative $(f(x+h)-f(x))/h$; CVODE uses difference quotients to build the Jacobian matrix.
\item \textbf{dimension} (Chapter 1). The number of vectors in a basis of a vector space.
\item \textbf{dirac16complex} (Chapter 0). The field theory of this book: a 16-component complex spinor field with anticommuting (Grassmann) components in eight dimensions of signature (4,4), with a mass term, a scalar interaction and the canonical spin connection.
\item \textbf{dirac16complex00} (Chapter 0). The same Lagrangian for a field with ordinary commuting complex components (a classical field); its real restriction is the notebook's Lagrangian (Chapter 6).
\item \textbf{Dirac adjoint} (Chapters 2 and 6). The row $\bar\Psi=\Psi^\dagger C$ built with the charge matrix $C$; the bilinears $\bar\Psi M\Psi$ with suitable $M$ do not change under spin transformations.
\item \textbf{Dirac bracket} (Chapter 8). The bracket of Hamiltonian mechanics with second-class constraints, built from the inverse of the matrix of the constraint brackets; its quantized form gives the canonical anticommutator.
\item \textbf{dirac-main} (Chapter 2). The reference implementation of the gamma matrices and the split octonions with which Stage 1 compares (its tensor matrices are the Clifford picture); its files are kept in the git-ignored folder \texttt{dirac-ma\allowbreak{}in/}, so that a public clone runs the gates without those comparisons (Section 19.6).
\item \textbf{Dirac operator} (Chapter 4). The operator $\gamma^\mu D_\mu$ of the field equation in a curved space; its square is given by the Lichnerowicz formula.
\item \textbf{Dirac sea} (Chapter 8). The state in which every negative-energy mode is filled; the vacuum of the good sector, in which a hole is an antiparticle with positive energy.
\item \textbf{Dirac's exchange energy} (Chapter 12). The exchange energy per volume of the uniform electron gas, $-\tfrac34(3/\pi)^{1/3}n^{4/3}$ in atomic units; the local exchange of the LDA.
\item \textbf{direct and exchange terms} (Chapter 14). The two ways of pairing the four amplitudes of an average $\langle c_n^\ast c_q^\ast c_rc_p\rangle$; the exchange term has the sign $+$ for Gaussian waves and $-$ for fermions.
\item \textbf{distance modulus} (Chapter 11). $\mu=5\log_{10}(d_L/10\,\mathrm{pc})$, the logarithm of the luminosity distance used in supernova tables.
\item \textbf{dominant energy condition, DEC} (Chapter 9). The weak energy condition together with the requirement that energy does not flow faster than light; it needs $\rho\ge0$ in particular.
\item \textbf{dot product, cross product} (Chapter 3). For vectors of three components $u\cdot r=u_1r_1+u_2r_2+u_3r_3$ (a symmetric number) and $u\times r$ (an antisymmetric vector); the quaternion product contains both.
\item \textbf{double, double precision} (Chapter 10). The standard 64-bit floating-point number, with about 16 significant decimal digits.
\item \textbf{double cover} (Chapter 2). A map from a group onto another that sends exactly two elements (here $R$ and $-R$) to each image; Pin(4,4) and Spin(4,4) are double covers of O(4,4) and SO(4,4).
\item \textbf{dry run, partial run} (Chapter 19). Two modes of the Stage-4 gate: a dry run lists every step with its command and its expected time and runs nothing; a partial run runs only the listed steps and ends with \texttt{\detokenize{PARTIAL}}, never with OK.
\item \textbf{dual (of an antisymmetric $4\times4$ matrix)} (Chapter 3). $({\ast}A)_{pq}=\tfrac12\sum_{r,s}\epsilon_{pqrs}A_{rs}$; the matrix is self-dual if ${\ast}A=A$ and anti-self-dual if ${\ast}A=-A$.
\item \textbf{dummy index} (Chapter 1). An index that is summed over; its name can be changed without changing the expression.
\item \textbf{dynamics} (Chapter 16). The equations that say how a state changes in time and how often a transition happens; conservation laws alone do not decide this.
\item \textbf{edition, registered edition, verify mode} (Chapter 19). The page count and the sha256 of a document's PDF as recorded in \texttt{provenan\allowbreak{}ce/pdf-s\allowbreak{}pecifica\allowbreak{}tions.js\allowbreak{}on}. In its verify mode, the default, the PDF builder requires a rebuilt PDF to reproduce the registered edition byte for byte; with the option \texttt{--regist\allowbreak{}er} it records a new edition instead.
\item \textbf{eigenvalue, eigenvector, eigenspace} (Chapters 1 and 2). A number $\lambda$ and a nonzero vector $v$ with $Av=\lambda v$; the eigenspace of $\lambda$ is the subspace of all such $v$ together with 0.
\item \textbf{eigenvalue problem} (Chapter 10). The task of finding the values of a parameter (for example the Kohn--Sham level $\varepsilon$) for which a differential equation with boundary conditions has a nonzero solution.
\item \textbf{Einstein equations} (Chapter 4). $G^\mu{}_\nu=\kappa T^\mu{}_\nu$ (in this book with $\kappa$ of eight dimensions), relating the curvature of spacetime to its energy--momentum source.
\item \textbf{Einstein--Lovelock equations} (Chapter 4). The Einstein equations supplemented by the Lovelock tensors of orders 2 and 3, the most general equations of this kind in eight dimensions; the notebook names them, but neither the notebook nor the project computes the higher orders.
\item \textbf{Einstein summation convention} (Chapter 1). An index that appears once up and once down in a product is summed over.
\item \textbf{Einstein tensor} (Chapter 4). $G_{\mu\nu}=R_{\mu\nu}-\tfrac12g_{\mu\nu}R$; its covariant divergence vanishes identically.
\item \textbf{electric charge} (Chapter 17). The conserved charge of electromagnetism; particles and antiparticles carry opposite electric charges.
\item \textbf{elementary matrix} (Chapter 2). $E_{ij}$, the matrix with a 1 in row $i$ and column $j$ and 0 everywhere else; the $E_{ij}$ form a basis of all matrices.
\item \textbf{energy condition} (Chapter 9). An inequality for the energy density and the pressures of a source, such as the weak, null, strong and dominant conditions; the primordial field violates some of them.
\item \textbf{energy density} (Chapter 5). The energy per volume: for the canonical tensor of Chapter 5 the component $\Theta^4{}_4$; for the metric energy--momentum tensor and an observer moving along $x_4$ with $g_{44}=-1$, $\rho=T_{44}=-T^4{}_4$ (Chapters 7 and 9).
\item \textbf{energy--momentum tensor} (Chapter 5). The tensor $T^\mu{}_\nu$ that contains energy density, momentum density, pressures and stresses; in Chapter 7 the metric (Hilbert) energy--momentum tensor from the variation of the action with respect to the metric, the source of the Einstein equations.
\item \textbf{energy projectors} (Chapter 8). $\Lambda_\pm=\tfrac12(1\pm h_k/E)$, which split the modes of a momentum into those with energy $+E$ and $-E$.
\item \textbf{ensemble} (Chapter 12). A statistical mixture of many-body states with probabilities; at a temperature $T$ the grand-canonical ensemble.
\item \textbf{entropy} (Chapter 12). $S=-\mathrm{tr}(\hat\rho\ln\hat\rho)$, for independent fermions $-\sum_n[f_n\ln f_n+(1-f_n)\ln(1-f_n)]$.
\item \textbf{envelope theorem} (Chapter 12). The derivative of the minimum of a function with respect to a parameter equals the partial derivative at the minimizer; it underlies the Hellmann--Feynman relations.
\item \textbf{environment variable} (Chapter 19). A named setting of the command shell that the programs started from it can read, such as \texttt{PYTHONUT\allowbreak{}F8=1}.
\item \textbf{equation of state, equation-of-state parameter} (Chapter 5). A relation between pressure and energy density; the parameter is $w=p/\rho$ ($w=0$ for dust, $\tfrac13$ for radiation, $-1$ for a cosmological constant).
\item \textbf{equivalent representations} (Chapter 2). Two representations $\rho,\rho'$ with an invertible intertwiner, $\rho'(g)=X\rho(g)X^{-1}$.
\item \textbf{erratum} (Chapter 18). A recorded correction of a specification or of the notebook, with its evidence, such as the errata E4.1 to E4.13 of the Stage-4 specification; errata are cited from Chapter 3 on.
\item \textbf{Euler--Lagrange equations} (Chapter 5). The conditions for the action to be stationary, $\frac{d}{dt}\frac{\partial L}{\partial\dot q}=\frac{\partial L}{\partial q}$, and their field version with $\partial_\mu$.
\item \textbf{Euler--Lagrange expression} (Chapter 5). The left side $E^A$ of the Euler--Lagrange field equation of the component $A$; it vanishes exactly on solutions and is unchanged when a total divergence is added to $\mathcal L$.
\item \textbf{Euler's method, explicit and implicit Euler} (Chapter 10). The simplest method for $y'=f(t,y)$: $y_{n+1}=y_n+hf(t_n,y_n)$ (explicit); the implicit (backward) form evaluates $f$ at the new point.
\item \textbf{even, odd} (Chapters 2 and 5). A Clifford monomial or a Grassmann monomial of even or odd degree; even Grassmann elements commute with everything, odd ones anticommute with each other.
\item \textbf{evolution equations} (Chapter 11). The Einstein equations with second time derivatives, which advance the metric in time, as opposed to the constraint.
\item \textbf{exact symmetry} (Chapter 17). A map $T$ with $\mathcal L_{m,\lambda}[T\Psi](x)=\mathcal L_{m,\lambda}[\Psi](Rx)$ for every field, in flat space or in a gravitational field in which the reflection $R$ is an isometry.
\item \textbf{exchange--correlation energy} (Chapter 12). The unknown part $E_{xc}[n]$ of the Kohn--Sham energy functional that contains everything beyond the non-interacting kinetic energy, the external energy and the Hartree energy; it must be approximated.
\item \textbf{exchange energy} (Chapter 12). The Fock term $E_x=-\tfrac12\int\!\int|\rho(x,x')|^2w(x,x')\,dx\,dx'$ of the energy of a Slater determinant; it has no classical analogue, comes from the antisymmetry of the wave function and lowers the energy for a repulsive interaction.
\item \textbf{exchange-only} (Chapter 13). A Kohn--Sham calculation that keeps the exact exchange of the interaction and no correlation energy; Chapter 13 is of this kind.
\item \textbf{exchange potentials} (Chapter 13). The derivatives $v_v=\partial e_x/\partial n=-\tfrac\lambda{16}n$ and $v_s=\partial e_x/\partial S=-\tfrac\lambda{16}S$ of the local exchange energy density; $v_v$ shifts the energy, $v_s$ adds to the mass.
\item \textbf{exit code} (Chapter 19). The whole number with which a program ends: 0 for success, any other number for failure.
\item \textbf{EXP-1 to EXP-5} (Chapter 11). The five numerical experiments of Stage 3: the field in the primordial field of the notebook (EXP-1), a homogeneous self-gravitating universe in eight dimensions (EXP-2), a condensate as the dark energy of the late universe (EXP-3), the quanta of an expanding 3-space as dark matter and their creation (EXP-4), and the deflating extra times (EXP-5).
\item \textbf{expectation value} (Chapter 8). $\langle f,Of\rangle$ for a normalized state $f$, the average of many measurements of $O$.
\item \textbf{expectation-value rule} (Chapter 11). The rule that replaces a bilinear of the quantum field by its average in the state (normal ordered against the sea); it defines the mean-field sources.
\item \textbf{explicit, implicit method} (Chapter 10). An explicit method computes the new value from known values; an implicit one solves an equation that contains the new value, and is needed for stiff problems.
\item \textbf{exponential wall} (Chapter 12). The growth of the size of the many-body wave function exponentially with the number of particles, which makes a direct computation impossible.
\item \textbf{extra times} (Chapter 9). The three coordinates $x_5,x_6,x_7$, time-like like $x_4$; in the primordial field they deflate.
\item \textbf{extrinsic curvature} (Chapter 9). The rate of change of the induced metric of a surface along its unit normal; its jump across the brane is fixed by the Israel junction condition.
\item \textbf{Fermi--Dirac distribution} (Chapter 12). $f(\varepsilon)=1/(e^{(\varepsilon-\mu)/T}+1)$, the occupation of a fermion level at temperature $T$.
\item \textbf{Fermi momentum} (Chapter 12). The largest momentum occupied in the ground state of a uniform gas, $k_F=(3\pi^2n)^{1/3}$ for two spin states.
\item \textbf{fermion mode, boson mode} (Chapter 14). A fermion mode is empty or occupied once; a boson mode has a commutator $[b,b^\dagger]=1$ and can be occupied any number of times.
\item \textbf{fermions} (Chapter 5). Particles described by anticommuting variables or operators; their many-body wave functions change sign when two of them are exchanged, and no mode holds more than one.
\item \textbf{fibre, base} (Chapter 9). In a warped metric the coordinates split into base coordinates (in Lemma 9.1 the two coordinates $x_0$ and $x_4$) and fibre coordinates (the six others), whose scale factors depend only on the base coordinates.
\item \textbf{field} (Chapter 5). A quantity that has a value at every point of spacetime, such as a scalar field $\phi(x)$ or the spinor field $\Psi(x)$.
\item \textbf{field equations} (Chapter 5). The Euler--Lagrange equations of a field; for the two fields of this book derived in Chapter 7.
\item \textbf{first excited state} (Chapter 12). The state of lowest energy above the ground state; in Chapter 13 obtained by moving one particle from the highest occupied to the lowest empty level.
\item \textbf{first-order adiabatic basis} (Chapter 11). The basis of instantaneous levels corrected to first order in the rate of change, in which EXP-4 measures the created weight $|\beta_k|^2$ so that the adiabatic dressing is not counted as creation.
\item \textbf{first-order form} (Chapter 10). A system of first derivatives equivalent to a higher-order equation, obtained by treating derivatives as new unknowns.
\item \textbf{first quantization} (Chapter 6). The description of a single particle by a wave function.
\item \textbf{fixed-modulus} (Chapter 14). A random amplitude $c=\sqrt f\,e^{i\varphi}$ with a random phase and a fixed modulus; unlike Gaussian amplitudes it has $\langle|c|^4\rangle=f^2$.
\item \textbf{fixed point} (Chapter 12). A density that the self-consistency map returns unchanged; the Kohn--Sham solution.
\item \textbf{fixed-point iteration} (Chapter 10). Solving $x=g(x)$ by repeating $x\leftarrow g(x)$; CVODE solves the implicit equation of each Adams step this way.
\item \textbf{FMA} (Chapter 10). Fused multiply-add, a processor instruction that computes $ab+c$ with a single rounding; the committed numerical files are byte-identical only when the program is built with it.
\item \textbf{Fock operator} (Chapter 12). The one-particle operator of the Hartree--Fock equations, kinetic plus external plus Hartree plus exchange operator.
\item \textbf{Fock space} (Chapter 8). The space of states with any number of particles and antiparticles, built from the vacuum by creation operators; positive in the good sector.
\item \textbf{frame, vielbein} (Chapter 4). At every point a set of eight orthonormal directions, given by $e_\mu{}^a$ with $g_{\mu\nu}=e_\mu{}^a\eta_{ab}e_\nu{}^b$; spinor components are measured in it.
\item \textbf{frame form} (Chapter 17). The way a reflection is applied in a gravitational field: the points are not moved, $\Psi'(x)=M\Psi(x)$, and the frame directions are multiplied by the signs $r_a$, which leaves the metric unchanged.
\item \textbf{Friedmann equations} (Chapter 11). The Einstein equations of a homogeneous and isotropic four-dimensional universe, relating the Hubble rate and its derivative to the energy density and the pressure.
\item \textbf{functional} (Chapter 5). A rule that assigns a number to a whole function, such as the action to a path or the energy to a density.
\item \textbf{functional derivative} (Chapter 12). $\delta F/\delta n(\mathbf r)$, defined by $\delta F=\int\frac{\delta F}{\delta n}\,\delta n\,d^3r$ to first order.
\item \textbf{fundamental gap} (Chapter 12). $E(N+1)+E(N-1)-2E(N)$, the difference of ionization energy and electron affinity; not equal to the Kohn--Sham gap in general.
\item \textbf{fundamental symmetry} (Chapter 8). An operator $J$ with $J=J^\dagger=J^{-1}$ that turns the indefinite form of a Krein space into a positive inner product $[f,Jh]$.
\end{itemize}

\subsection{Glossary: G to K}

\begin{itemize}
\item \textbf{G1, G2} (Chapter 4). The two test geometries of Stage 1: G1 is a generic non-diagonal vielbein $e_\mu{}^a=\delta_\mu{}^a+P_\mu{}^a(x)$ with polynomial entries, tested at three rational points; G2 is the primordial field. Check names end in \texttt{\detokenize{_G1}} or \texttt{\detokenize{_G2}} accordingly.
\item \textbf{gamma matrices} (Chapter 2). Sixteen-by-sixteen matrices $\gamma^0,\dots,\gamma^7$ with $\gamma^a\gamma^b+\gamma^b\gamma^a=2\eta^{ab}$; the book uses the real matrices of the notebook.
\item \textbf{gate} (Chapter 0). A script that runs every program of a stage in a fixed order, stops at the first failure, compares the outputs with the committed files and ends with one line \texttt{...\_veri\allowbreak{}fication\allowbreak{}=OK} or \texttt{\detokenize{FAILED}}; each exists as a PowerShell and a Bash twin (Chapter 19).
\item \textbf{Gauss--Bonnet combination} (Chapter 9). The Lovelock Lagrangian of order 2, $R^2-4R_{\mu\nu}R^{\mu\nu}+R_{\mu\nu\rho\sigma}R^{\mu\nu\rho\sigma}$; in eight dimensions it would add a term to the Einstein equations that neither the notebook nor the project computes.
\item \textbf{Gaussian amplitudes} (Chapter 14). Random complex amplitudes with the density $e^{-|c|^2/f}/(\pi f)$, the amplitudes of thermal (chaotic) waves; for them the exchange term of an average has the sign $+$.
\item \textbf{Gaussian normal} (Chapter 7). A time coordinate with $g_{44}=-1$ and $g_{4i}=0$ for $i\ne4$, with the frame's time direction along $x_4$; the observer of the energy density and the pressures.
\item \textbf{Gauss-Legendre quadrature} (Chapter 11). A rule $\sum_nw_nF(k_n)$ with nodes and weights chosen so that it integrates every polynomial up to degree $2n-1$ exactly; EXP-4 uses 48 nodes.
\item \textbf{generalized Kronecker delta} (Chapter 4). $\delta^{\mu_1\cdots\mu_k}_{\nu_1\cdots\nu_k}$, the determinant of the $k\times k$ matrix of ordinary Kronecker deltas; the Lovelock tensors are built with it.
\item \textbf{generators} (Chapter 2). Elements whose products and combinations give a whole algebra or group: the $\gamma^a$ generate the Clifford algebra, and the $S^{ab}$ generate the part of Spin(4,4) connected to 1.
\item \textbf{geodesic} (Chapter 4). A curve that is as straight as the geometry allows: its velocity is parallel transported along itself; freely falling particles move on geodesics.
\item \textbf{Git} (Chapter 19). The program that records the history of a folder of files as a sequence of commits and copies it between computers; the repository of this book is kept with it.
\item \textbf{global error} (Chapter 10). The difference between the computed and the exact solution after many steps; for a method of order $p$ it is proportional to $h^p$.
\item \textbf{good sector} (Chapter 8). The modes of the field that do not depend on the three extra times $x_5,x_6,x_7$; in it the Fock space is positive and the Hamiltonian not negative. The numerical chapters are restricted to it, by assumption.
\item \textbf{grand potential} (Chapter 12). $\Omega=\mathrm{Tr}[\hat\rho(\hat H-\mu\hat N)]+T\,\mathrm{Tr}[\hat\rho\ln\hat\rho]$; the equilibrium state at temperature $T$ minimizes it.
\item \textbf{Grassmann numbers, Grassmann algebra} (Chapter 5). Numbers generated by anticommuting symbols $\theta_i$ with $\theta_i\theta_j=-\theta_j\theta_i$, so $\theta_i^2=0$; the classical description of fermion fields. The components of dirac16complex are Grassmann-odd.
\item \textbf{gravitational pair creation} (Chapter 11). The appearance of particle--antiparticle pairs out of the vacuum through the expansion of space alone; EXP-4 computes it for an inflation followed by radiation.
\item \textbf{ground state} (Chapter 12). The state of lowest energy; its energy is the minimum of the energy expectation value over all normalized states.
\item \textbf{group} (Chapter 2). A set with an associative product, a unit and inverses; a subgroup is a subset that is itself a group.
\item \textbf{Hamiltonian} (Chapter 5). The energy as a function of coordinates and momenta, $H=p\dot q-L$; it generates the time evolution through Hamilton's equations (Chapter 8) and, as an operator, the Schrödinger equation.
\item \textbf{Hamilton's equations} (Chapter 8). $\dot q=\partial H/\partial p$ and $\dot p=-\partial H/\partial q$, the first-order form of the equations of motion.
\item \textbf{Hartree approximation} (Chapter 12). The mean-field approximation that keeps the Hartree potential and drops the exchange; it keeps the self-interaction of each particle.
\item \textbf{Hartree energy, Hartree potential} (Chapter 12). $E_H=\tfrac12\int\!\int n(x)w(x,x')n(x')$, the classical interaction energy of the density with itself, and its functional derivative $v_H=\int w\,n$.
\item \textbf{Hartree--Fock approximation} (Chapter 12). The best single Slater determinant, found by minimizing the energy over all choices of orthonormal orbitals; restricted if both spins share one spatial orbital, unrestricted otherwise.
\item \textbf{Hartree--Fock energy density} (Chapter 13). For the contact interaction, $e_{HF}=\tfrac\lambda2(S^2-\mathrm{Tr}(BC\rho BC\rho))$ with $S=\mathrm{Tr}(BC\rho)$: the Hartree term $\tfrac\lambda2S^2$ and the exchange (Fock) term.
\item \textbf{heat capacity} (Chapter 12). $C_V=dE/dT$ at fixed particle number; for the Kohn--Sham gas of Chapter 13 $C_V=T\,dS_s/dT$.
\item \textbf{Heisenberg equation} (Chapter 8). $\dot O=i[H,O]$ for an operator without explicit time dependence; the quantum form of Hamilton's equations.
\item \textbf{Hellmann--Feynman rule} (Chapter 13). For a normalized eigenvector $\chi$ of $h(k)$ with boundary conditions that do not depend on $k$, $d\varepsilon/dk=\langle\chi|(\partial_kh)\chi\rangle$; the Stage-4 checks whose names begin with \texttt{KS\_funct\allowbreak{}ional\_he\allowbreak{}llmannFe\allowbreak{}ynman} test the analogous relations for the derivatives of the energy functional with respect to the coupling, the mass and the temperature.
\item \textbf{Hermitian} (Chapter 1). A square matrix with $A^\dagger=A$; its eigenvalues are real and it has an orthonormal basis of eigenvectors.
\item \textbf{Hermitian conjugate, conjugate transpose} (Chapter 1). $A^\dagger=(A^\ast)^T$, the conjugate transpose.
\item \textbf{Hermitian form} (Chapter 8). A rule $[f,h]=f^\dagger Gh$ with a Hermitian matrix $G$; it is non-degenerate if $Gf=0$ only for $f=0$ and indefinite if some $[f,f]$ are negative.
\item \textbf{hexadecimal} (Chapter 19). Writing whole numbers in base 16 with the sixteen digits 0 to 9 and a to f (a stands for 10, b for 11, and so on up to f for 15); the upper-case letters A to F denote the same digits. The names of commits and the sha256 fingerprints are written in hexadecimal.
\item \textbf{hidden space} (Chapter 9). The coordinate $x_0$, space-like like $x_1,x_2,x_3$; in the static member of the primordial field the warp factor depends on it.
\item \textbf{Hilbert adjoint, Hilbert norm} (Chapters 8 and 11). The adjoint and the norm $u^\dagger u$ with respect to the positive inner product, as opposed to the Krein adjoint and the Krein norm $u^\dagger Bu$.
\item \textbf{Hilbert space} (Chapter 8). A complex vector space with a positive inner product; it gives probabilities.
\item \textbf{Hohenberg--Kohn theorems} (Chapter 12). The ground-state density determines the external potential up to a constant, and the ground-state energy is the minimum of a functional of the density; the foundation of DFT.
\item \textbf{hole} (Chapter 12). An empty level below the highest occupied one, created by removing a particle; in Chapter 8 a hole in the Dirac sea is an antiparticle.
\item \textbf{homogeneous frames} (Chapter 4). Diagonal frames of the metrics $ds^2=-dt^2+\sum_{j\ne4}\eta_{jj}h_j(t)^2dx_j^2$ with $t=x_4$, whose scale factors $h_j$ depend only on the time; the metrics of cosmology in this book.
\item \textbf{homomorphism} (Chapter 2). A map between groups that respects the products; its kernel is the set of elements sent to the unit.
\item \textbf{honesty rule} (Chapter 0). The rule of this book that every statement carries its status (PROVED, COMPUTED, ASSUMED, HYPOTHESIS or OPEN) and that no claim goes beyond what the committed files support.
\item \textbf{Hubble length} (Chapter 11). $c/H$, the distance light travels in one expansion time; a wave is inside the horizon when its physical wavelength is much shorter.
\item \textbf{Hubble rate} (Chapter 9). The logarithmic rate of change of a scale factor, $H=\dot a/a$; today's value in ordinary cosmology is $H_0$ (Chapter 11).
\item \textbf{Hurwitz's theorem} (Chapter 3). The only composition algebras with a positive norm have dimensions 1, 2, 4 and 8 (the real numbers, complex numbers, quaternions and octonions); quoted, not proved.
\item \textbf{HYPOTHESIS} (Chapter 0). The status of a statement that is proposed but neither proved nor computed, such as the notebook's claim that universes are created in pairs.
\item \textbf{identity matrix} (Chapter 1). The matrix $I$ with ones on the diagonal and zeros elsewhere, $IA=AI=A$.
\item \textbf{Illinois variant} (Chapter 10). The regula falsi with the improvement that, when the same end is replaced twice in a row, the function value at the other end is halved; the level finder of the Stage-4 solver.
\item \textbf{image field} (Chapter 15). The quantum field $\Psi_-=\gamma^8\Psi_+$, built from the same operators as $\Psi_+$ and acting on the same states; its anticommutator is the reversed Krein metric $-B$. It is the same quantum system as $\Psi_+$: its own generator of the evolution in $x_4$, its own charge and its own source of gravity are $+H_+$, $+Q_+$ and $+T_{\mu\nu}[\Psi_+]$. The identities $H_++H[\Psi_-;-m]=0$ and $Q_++Q[\Psi_-]=0$, in which $H[\Psi_-;-m]$ and $Q[\Psi_-]$ are the formulas of $\mathcal L_{-m,-\lambda}$ evaluated on the image, have the form $X+(-X)=0$; they are not a cancellation by a second universe (matter--antimatter document, §7.3).
\item \textbf{image (of a matrix)} (Chapter 2). The set of all columns $Mu$, a subspace; its dimension is the rank.
\item \textbf{indefinite} (Chapter 8). A form or inner product that takes both positive and negative values, such as $f^\dagger Bf$.
\item \textbf{initial-value problem} (Chapter 10). A differential equation together with the value of the solution at a starting time.
\item \textbf{inner product} (Chapter 8). A rule $\langle f,h\rangle$, linear in $h$ and antilinear in $f$, with $\langle h,f\rangle=\langle f,h\rangle^\ast$; positive if $\langle f,f\rangle>0$ for $f\ne0$.
\item \textbf{inside the horizon} (Chapter 11). A wave whose physical wavelength is much shorter than the Hubble length, $K\gg H$.
\item \textbf{interacting quantum field theory} (Chapter 18). A quantum theory of the field with its interaction, with positive probabilities, a Hamiltonian bounded below and finite predictions; for dirac16complex it is OPEN.
\item \textbf{interaction, $U(S)$} (Chapter 6). The term $-\sqrt{|g|}\,U(S)$ of the Lagrangian, a function of the scalar density $S=\bar\Psi\Psi$; for the contact interaction of Chapters 11 and 13, $U=\tfrac\lambda2S^2$.
\item \textbf{interpolation} (Chapter 10). Computing the solution between two steps from a polynomial through computed values; CVODE returns output at requested times this way.
\item \textbf{intertwiner} (Chapter 2). A matrix $X$ with $X\rho(g)=\rho'(g)X$ for all $g$; an invertible one makes the two representations equivalent.
\item \textbf{invariant bilinear form} (Chapters 6 and 8). A matrix $G$ with $(S^{ab})^TG+GS^{ab}=0$ for all $a,b$, so that $\Psi^TG\Phi$ does not change under spin rotations; Chapter 8 shows that no such form gives a positive charge density in signature (4,4), and Theorem M3 of Chapter 17 lists all of them.
\item \textbf{invariant mass} (Chapter 16). $\mu^2=E_{\mathrm{tot}}^2-|\mathbf p_{\mathrm{tot}}|^2$ of a group of particles; conserved in every process, so a single photon in empty space cannot make a pair.
\item \textbf{inversion (of a permutation)} (Chapter 1). A pair of places $i<j$ with $\sigma(i)>\sigma(j)$; the sign of the permutation is $(-1)$ to the number of inversions.
\item \textbf{invertible} (Chapter 1). A square matrix $A$ with an inverse $A^{-1}$, $AA^{-1}=A^{-1}A=I$; equivalently $\det A\ne0$.
\item \textbf{irreducible} (Chapter 2). A representation with no invariant subspace other than zero and the whole space; by Schur's lemma its commutant consists of multiples of the identity.
\item \textbf{isometry} (Chapter 16). A map of spacetime to itself that leaves the metric unchanged, such as the reflection $y\to-y$ of the Z2 geometry.
\item \textbf{Israel junction condition} (Chapter 9). The relation between the jump of the extrinsic curvature across a surface and the energy--momentum on the surface; it fixes the source that the brane must carry.
\item \textbf{Jacobian matrix} (Chapter 10). The matrix of partial derivatives $\partial f_i/\partial y_j$ of the right-hand side; implicit methods need it for Newton's method. (Chapter 4 uses the Jacobian matrix $\partial x'/\partial x$ of a change of coordinates.)
\item \textbf{Jacobi's formula} (Chapter 1). $\partial\det A=\det A\,\mathrm{tr}(A^{-1}\partial A)$.
\item \textbf{Janak's theorem} (Chapter 12). The derivative of the Kohn--Sham energy with respect to the occupation of a level equals the level's energy, $\partial E/\partial f_b=\varepsilon_b$.
\item \textbf{jets} (Chapter 5). The field, its first derivatives and higher derivatives treated as independent variables at a point; the Euler--Lagrange expression is a function of them.
\item \textbf{JSON} (Chapter 0). A plain-text format for lists and tables (JavaScript Object Notation) that Python reads directly; every report of the repository is a JSON file.
\item \textbf{Kasner exponents} (Chapter 11). The powers $p^{\mathrm K}_i$ in $h_i\propto(t-t_s)^{p^{\mathrm K}_i}$ near a singular time of a homogeneous metric (EXP-2).
\item \textbf{kernel} (Chapter 2). The elements a homomorphism sends to the unit; in Chapter 13 the exchange kernel of the contact interaction.
\item \textbf{ket} (Chapter 12). Dirac's notation $|\Psi\rangle$ for a state vector; the bra $\langle\Psi|$ is its adjoint.
\item \textbf{kinetic and potential energy} (Chapter 7). Two splits of the energy density, both definitions: the Lagrangian split $\mathrm{KE}_L=\tfrac12K_4$, $\mathrm{PE}_L=\rho-\mathrm{KE}_L$, and the Hamiltonian split $\mathrm{KE}_H=-K_\perp$, $\mathrm{PE}_H=mS+U(S)$.
\item \textbf{kinetic term} (Chapter 6). The part $K$ of the Lagrangian with one derivative of the field, symmetrized between $\Psi$ and $\bar\Psi$.
\item \textbf{kink formula} (Chapter 11). The leading large-momentum estimate $|\beta_k|^2\approx(mkH_{\mathrm{inf}}^2/4E^4)^2$ of the particles created when the second derivative of the scale factor jumps.
\item \textbf{Klein's inequality} (Chapter 12). $\mathrm{Tr}[\hat\rho(\ln\hat\rho-\ln\hat\sigma)]\ge0$ for density operators; it proves that the Gibbs state minimizes the grand potential.
\item \textbf{Kohn--Sham equations} (Chapter 12). One-particle equations $[-\tfrac12\nabla^2+v_s]\varphi_a=\varepsilon_a\varphi_a$ with $v_s=v+v_H+v_{xc}$, whose occupied orbitals give the density of the interacting system if $E_{xc}$ is exact.
\item \textbf{Kohn--Sham fermion-gas thermodynamic pseudo-potential} (Chapter 13). The Stage-4 specification's name for the pair $M_{\mathrm{eff}}=m+\tfrac{15}{16}\lambda S_p$ and $v_v=-\tfrac1{16}\lambda n_p$ of the Kohn--Sham equation, together with the gravitational terms: the local density approximation built from the exact exchange of the contact interaction in the uniform gas, used as an approximation in the primordial field.
\item \textbf{Kohn--Sham gap} (Chapter 12). $\Delta_{KS}=\varepsilon_{\mathrm{LUMO}}-\varepsilon_{\mathrm{HOMO}}$, the smallest difference in the particle--hole spectrum; the first estimate of the lowest excitation energy.
\item \textbf{Kohn--Sham potential} (Chapter 12). The local potential $v_s$ of the Kohn--Sham equations.
\item \textbf{Koopmans' theorem} (Chapter 12). In Hartree--Fock theory the energy to remove a particle from an orbital, with all other orbitals frozen, is minus the orbital energy.
\item \textbf{Krein adjoint} (Chapter 8). The adjoint with respect to the indefinite form $f^\dagger Bh$; the canonical $\Psi^\dagger=\chi B$ is the Krein adjoint of $\Psi$, while $\chi$ is its Hilbert adjoint.
\item \textbf{Krein image pair} (Chapters 15 and 16). A formal construction at the Kohn--Sham level: a $+M$ state together with its T1 image with $-\lambda$ and the reversed metric $-B$, formed as $(+M)-(-M)$ for every one-body density and energy, where $-M$ is the ordinary $-M$ state of T3 (check \texttt{PAIR\_tot\allowbreak{}als\_ksKr\allowbreak{}einImage\allowbreak{}Pair}). Its total energy, charge and energy--momentum tensor are zero at every point, but the zero is of the kind $X+(-X)=0$: the image is the $+M$ state measured with the formulas of the partner theory, so the zero restates the equality $E(+M)=E(-M)$ of T3. It is not a pair of two universes whose energies and charges cancel.
\item \textbf{Krein norm} (Chapter 8). $u^\dagger Bu$, which can be positive, negative or zero and is conserved for every mode, in and out of the good sector; Chapter 11 uses it in EXP-5.
\item \textbf{Krein sign} (Chapters 7 and 14). The value $\beta=js_2=\pm1$ by which the matrix $B$ acts on one of the eight $2\times2$ blocks of the Kohn--Sham problem, the sign of the form $u^\dagger Bu$ on that block; four blocks have $+1$ and four $-1$.
\item \textbf{Krein space} (Chapter 8). A complex vector space with a non-degenerate indefinite Hermitian form and a fundamental symmetry; the natural space of one-particle states of dirac16complex.
\item \textbf{Krein-unitary} (Chapter 8). A matrix $R$ with $R^\dagger BR=B$; it preserves the canonical anticommutator.
\item \textbf{Kronecker delta} (Chapter 1). $\delta_{ij}=1$ for $i=j$ and 0 otherwise; the entries of the identity matrix.
\end{itemize}

\subsection{Glossary: L to O}

\begin{itemize}
\item \textbf{Lagrange multiplier} (Chapter 12). A number $\lambda$ added with a constraint $c(x)=c_0$ to a function being minimized; at the constrained minimum it equals the change of the minimum per unit change of $c_0$ (the chemical potential is one).
\item \textbf{Lagrangian} (Chapter 5). The function $L(q,\dot q)$ whose time integral is the action; for a particle usually kinetic minus potential energy.
\item \textbf{Lagrangian density} (Chapter 5). The function $\mathcal L$ of the fields and their derivatives whose integral over spacetime is the action of a field theory; for the two fields of this book given in Chapter 6.
\item \textbf{lapse} (Chapter 18). The factor $N(x_4)$ in $g_{44}=-N^2$ that relates the coordinate time to proper time; the reduced model of Chapter 18 needs it to obtain its constraint.
\item \textbf{LDA}: see local density approximation.
\item \textbf{left derivative, right derivative} (Chapter 5). The two ways of differentiating a product of Grassmann variables, removing the variable after moving it to the left or to the right end; they differ by a sign for odd products.
\item \textbf{Leibniz rule} (Chapter 1). The product rule $\partial(fg)=(\partial f)g+f\,\partial g$; for Grassmann variables and derivatives it acquires signs (Chapter 5).
\item \textbf{lepton number} (Chapter 17). $L$: electrons, muons, tau leptons and neutrinos have $L=+1$, their antiparticles $L=-1$.
\item \textbf{Levi-Civita connection} (Chapter 4). The unique connection that is torsion-free and preserves the metric; its coefficients are the Christoffel symbols. The project uses no other.
\item \textbf{Levy--Lieb constrained search} (Chapter 12). The definition of the universal functional as a minimum over all wave functions with a given density.
\item \textbf{Lichnerowicz formula} (Chapter 4). $(\gamma^\mu D_\mu)^2\Psi=g^{\mu\nu}(D_\mu D_\nu\Psi-\Gamma^\lambda{}_{\mu\nu}D_\lambda\Psi)-\tfrac14R\Psi$, the square of the Dirac operator in a curved space.
\item \textbf{light-cone basis, light-cone coordinates} (Chapter 3). Coordinates such as $x_0\pm x_7$, sums and differences of a space-like and a time-like coordinate, in which the norm of signature (4,4) becomes $y_0y_4+y_1y_5+y_2y_6+y_3y_7$; the notebook's $\tau_a$ are octonion multiplications written in such a basis.
\item \textbf{linear} (Chapter 1). A map with $f(ax+by)=af(x)+bf(y)$; matrices are the linear maps of $\mathbb R^n$ or $\mathbb C^n$.
\item \textbf{linear combination} (Chapter 1). A sum $\sum_ic_iv_i$ of vectors multiplied by numbers.
\item \textbf{linearly independent, dependent} (Chapter 1). Vectors are independent if no nontrivial combination of them is zero, dependent otherwise.
\item \textbf{linear mixing} (Chapter 12). The self-consistency update $n_{\mathrm{in}}\leftarrow(1-\beta)n_{\mathrm{in}}+\beta n_{\mathrm{out}}$ with a fraction $0<\beta<1$; it cures charge sloshing in simple cases.
\item \textbf{linear multistep method} (Chapter 10). A method that computes the next value from several previous values and right-hand sides; Adams methods and BDF are of this kind.
\item \textbf{line element} (Chapter 1). $ds^2=g_{\mu\nu}dx^\mu dx^\nu$, the squared length of a small displacement.
\item \textbf{local density approximation, LDA} (Chapter 12). The approximation that treats each small volume as a piece of uniform gas with the local density, $E_{xc}\approx\int e_{xc}(n(\mathbf r))\,d^3r$; Chapter 13 uses its analogue for the exchange of the contact interaction.
\item \textbf{local error} (Chapter 10). The error made in one step that starts from the exact solution; for Euler's method $\tfrac12h^2y''$, proportional to $h^2$, while the global error is proportional to $h$.
\item \textbf{localized at the brane} (Chapter 13). A Kohn--Sham orbital whose weight grows toward $y=0$, such as the zero mode $\chi=(e^{My},0)^T$.
\item \textbf{log file} (Chapter 19). The file in \texttt{build/lo\allowbreak{}gs/} into which a gate writes the complete output of one step, beginning with the start time and the command and ending with the finish time and the exit code.
\item \textbf{Lorentz transformation, Lorentz boost, Lorentz invariance} (Chapters 2 and 4). In this book a Lorentz transformation is a change of the frame directions that preserves the metric $\eta$, an element of O(4,4); a boost mixes one space-like and one time-like direction hyperbolically (Chapter 2). A local Lorentz transformation turns the frame differently at every point (Chapter 4), and a term of the Lagrangian is Lorentz-invariant if it does not change under the spin transformations that accompany such a turning (local spin invariance, Section 4.12).
\item \textbf{Lovelock gravity, Lovelock orders} (Chapter 4). The family of gravitational equations with at most second derivatives of the metric, built from powers of the curvature: order 1 is Einstein's, order 2 the Gauss--Bonnet term; in eight dimensions orders up to 3 exist (Lovelock's theorem, quoted).
\item \textbf{lowering, raising an index} (Chapter 1). Contracting with $g_{\mu\nu}$ or $g^{\mu\nu}$ (or $\eta$ for frame indices) to turn an upper into a lower index or back.
\item \textbf{machine epsilon} (Chapter 10). The spacing of double-precision numbers near 1, $2^{-52}\approx2.2\times10^{-16}$; the limit of every relative accuracy.
\item \textbf{Magnus method} (Chapter 10). An integrator for $i\dot u=h(t)u$ that advances each step by the exponential of $-i$ times a Hermitian matrix, a unitary matrix, so that $u^\dagger u$ is kept exactly; the EXP-4 checker uses a fourth-order Magnus method.
\item \textbf{Majorana-type term} (Chapter 17). A term built with $\Psi^T$ instead of $\Psi^\dagger$, such as $\Psi^TM\Psi$; it has U(1) charge 2 and would violate charge conservation. Theorem M3 shows that the anticommuting field has no Lorentz-invariant Majorana mass term.
\item \textbf{manifold} (Chapter 4). A space that is covered by charts, each of which looks like a piece of $\mathbb R^n$, with smooth changes of coordinates where charts overlap.
\item \textbf{mass dimension} (Chapter 13). The power of a mass unit that a quantity carries in units with $\hbar=c=1$ (a length has mass dimension $-1$); it measures the size of the coupling $\lambda$.
\item \textbf{mass term} (Chapter 6). The part $-m\sqrt{|g|}\,\bar\Psi\Psi$ of the Lagrangian, linear in the mass and quadratic in the components.
\item \textbf{matrix} (Chapter 1). A rectangular array of numbers; matrix multiplication is associative and distributive but not commutative.
\item \textbf{matter--antimatter problem} (Chapter 17). The question why the universe contains baryons but essentially no antibaryons. Chapter 17 shows that the theory of this book, as built, does not solve it. It contains no baryons and no baryon-number-violating interaction; its U(1) charge is exactly conserved (Theorem M1), so Sakharov's first condition fails; each field has an exact symmetry that reverses the charge (C for dirac16complex00, CP for dirac16complex; Theorem M2), so the second fails for the Lagrangian; and no departure from equilibrium is computed.
\item \textbf{max\_step} (Chapter 10). The largest step size allowed to CVODE in an experiment.
\item \textbf{mean field} (Chapter 11). The approximation in which each quantum moves in the average field of all the others, with a bilinear such as $S$ replaced by its expectation value; the Hartree idea (Chapter 12).
\item \textbf{measurement} (Chapters 0 and 19). A value that a verifier records in its report next to the checks, such as a count, a number, a list or a matrix (entry \texttt{measurem\allowbreak{}ents}; printed as a line \texttt{measurem\allowbreak{}ent\_<nam\allowbreak{}e>=}); unlike a check it is neither true nor false.
\item \textbf{Mermin's functional, Mermin's Kohn--Sham functional} (Chapter 12). The free energy functional of the Kohn--Sham ensemble at temperature $T$, $F=\sum_af_a\langle\varphi_a|-\tfrac12\nabla^2|\varphi_a\rangle-TS_s+\int vn+E_H+F_{xc}$; its stationary point has Fermi--Dirac occupations.
\item \textbf{Mermin's theorem} (Chapter 12). The finite-temperature analogue of the Hohenberg--Kohn theorems (1965): the equilibrium density determines the potential, and the grand potential is minimal at it.
\item \textbf{metric} (Chapter 1). A symmetric matrix $g_{\mu\nu}$ (in Chapter 4 a field of such matrices) that defines lengths and angles through the line element; of signature (4,4) in this book.
\item \textbf{mirror map, T2} (Chapter 15). The combination of the field with the reflection of one space-like direction; it maps the theory with $(m,\lambda)$ to $(-m,\lambda)$ and leaves energy, momentum and charge unchanged (only reflected).
\item \textbf{mirror pair} (Chapter 15). A $+M$ universe with its ordinary $-M$ partner of T3 (the same $\lambda$): its pair energy is $2E_+$, its charge $2N$ and its scalar density 0.
\item \textbf{mismatch} (Chapter 10). In a shooting method, the amount by which the solution started from one boundary misses the condition at the other; a level is where it vanishes.
\item \textbf{mixed-product rule} (Chapter 1). $(A\otimes B)(C\otimes D)=(AC)\otimes(BD)$ for tensor products.
\item \textbf{mode amplitude} (Chapter 11). The column $u(t)$ of 16 complex numbers in the plane-wave ansatz $\Psi=V^{-1/2}e^{i\sum_jk_jx_j}u(t)$.
\item \textbf{mode Hamiltonian} (Chapter 8). The $16\times16$ matrix $h_k$ with $i\partial_4u=h_ku$ for a plane wave of momentum $k$; its square is $E^2$ times the identity.
\item \textbf{mode operators} (Chapter 8). The 16 operators $\Psi_k$ obtained from the field by a Fourier transform over the slice; they annihilate and create the quanta of momentum $k$.
\item \textbf{modulus} (Chapter 1). $|z|=\sqrt{a^2+b^2}$ for $z=a+ib$.
\item \textbf{momentum} (Chapter 5). $p=\partial L/\partial\dot q$, the canonical momentum; for a first-order Lagrangian it is a function of the coordinates (a constraint).
\item \textbf{monomial, degree} (Chapters 2 and 5). A product of generators, such as $\gamma_A=\gamma^{a_1}\cdots\gamma^{a_k}$ with increasing indices (there are 256 for eight gammas) or a product of distinct Grassmann generators; the number $k$ of factors is its degree, and it is even or odd with $k$.
\item \textbf{natural cubic spline} (Chapter 13). Piecewise cubic interpolation with continuous first and second derivatives and zero second derivatives at the ends; the Stage-4 solver interpolates its potentials this way.
\item \textbf{NEC}: see null energy condition.
\item \textbf{negative control} (Chapter 17). A computation that must fail, run to show that a check can detect the failure; for example the notebook's contraction of the spin connection, for which a divergence identity must fail.
\item \textbf{Newton's constant} (Chapter 11). The coupling $G$ of gravity; in Chapter 11 a bound on its variation in time rules out one variant of the model.
\item \textbf{Newton's method} (Chapter 10). Solving $g(x)=0$ by the iteration $x\leftarrow x-g(x)/g'(x)$, the zero of the tangent line; CVODE uses it (with the Jacobian matrix) for the implicit equations of BDF.
\item \textbf{nilpotent} (Chapter 8). A matrix with a vanishing power, such as $h_k^2=0$; its exponential is a polynomial, and a mode with such a Hamiltonian grows linearly.
\item \textbf{Noether's theorem, Noether current} (Chapter 5). Theorem 5.2: every continuous symmetry of a Lagrangian gives a current, its Noether current, that is conserved on the solutions of the field equations. The phase symmetry gives the U(1) current and its charge (Chapter 7; Theorem M1 of Chapter 17), the translations the canonical energy--momentum tensor.
\item \textbf{non-degenerate} (Chapter 8). A Hermitian form $f^\dagger Gh$ with $Gf=0$ only for $f=0$.
\item \textbf{non-interacting $v$-representable} (Chapter 12). A density that is the ground-state density of some non-interacting system in some potential; the Kohn--Sham scheme assumes it.
\item \textbf{nonlinear} (Chapter 12). Equations whose operator depends on the unknowns, like the Hartree--Fock and Kohn--Sham equations; they are solved by iteration.
\item \textbf{norm} (Chapters 2 and 3). For a vector $u$ of $\mathbb R^{4,4}$ the number $n(u)=\eta(u,u)$, which can be negative ($u$ is a unit vector if $n(u)=\pm1$); for a split octonion the norm $N(x)=x\bar x$ of signature (4,4).
\item \textbf{normalized} (Chapter 12). A state with $\langle\Psi|\Psi\rangle=1$.
\item \textbf{normal ordering} (Chapters 7 and 8). Writing products of creation and annihilation operators with the annihilation operators on the right (with the sign of each fermion exchange), which removes the constant contribution of the filled Dirac sea.
\item \textbf{notebook, the} (Chapter 0). The Mathematica notebook \texttt{Pair\_Cre\allowbreak{}ation\_of\allowbreak{}\_Univers\allowbreak{}es\_WaveF\allowbreak{}unctionO\allowbreak{}fUnivers\allowbreak{}e-4+4-Ei\allowbreak{}nstein-L\allowbreak{}ovelock-\allowbreak{}Nash.nb} by Patrick L. Nash in the repository root, the starting point of the project.
\item \textbf{N-representable} (Chapter 12). A density that comes from some antisymmetric $N$-particle wave function.
\item \textbf{null} (Chapter 1). A nonzero vector of zero length, $g(v,v)=0$.
\item \textbf{null energy condition, NEC} (Chapter 9). $T(k,k)\ge0$ for every null vector $k$; for $k=e_4+e_0$ it reads $\rho+p_{(0)}\ge0$. In four dimensions it is $\rho+p\ge0$, which a phantom fluid violates (Chapter 11).
\item \textbf{number operator} (Chapter 8). $N=b^\dagger b$, whose eigenvalues 0 and 1 count the fermions in a mode.
\item \textbf{occupation numbers, occupation probabilities} (Chapter 12). The numbers $n_p\in\{0,1\}$ of particles in the orbitals of a determinant, and at a temperature the probabilities $f_a\in[0,1]$.
\item \textbf{octonions} (Chapter 3). The eight-dimensional composition algebra that is neither commutative nor associative but alternative; the split octonions are its form with a norm of signature (4,4).
\item \textbf{ODE}: see ordinary differential equation.
\item \textbf{off shell, on shell} (Chapter 7). A statement holds off shell if it holds for every field configuration, and on shell if it holds only for solutions of the field equations.
\item \textbf{one-body density matrix} (Chapter 12). $\rho=\sum_{a\in O}|\varphi_a\rangle\langle\varphi_a|$, with $\rho_{qp}=\langle a_p^\dagger a_q\rangle$; every expectation value in a Slater determinant is built from it (Wick's theorem).
\item \textbf{onto} (Chapter 2). A map is onto if every element of the target is the image of some element; $\Lambda_{\mathrm t}$ maps Pin(4,4) onto O(4,4) (by the quoted Cartan--Dieudonné theorem).
\item \textbf{OPEN} (Chapter 0). The status of a question that the project has not answered.
\item \textbf{open interval} (Chapter 1). $(a,b)$, the set of real numbers $x$ with $a<x<b$.
\item \textbf{operator} (Chapter 8). A rule that maps states to states, linear unless said otherwise.
\item \textbf{orbital} (Chapters 7 and 12). A one-particle wave function, one of the factors of a Slater determinant; in the Kohn--Sham chapters one solution $\chi$ of the block equation, occupied by one quantum.
\item \textbf{orbital relaxation} (Chapter 12). The change of the Kohn--Sham levels when a particle is moved; it is the difference between the Delta-SCF energy and the Kohn--Sham gap.
\item \textbf{order} (Chapter 10). Of a differential equation, its highest derivative; of a numerical method, the power $p$ with a global error proportional to $h^p$.
\item \textbf{ordinary differential equation} (Chapter 10). An equation for an unknown function of one variable and its derivatives.
\item \textbf{orthogonal} (Chapters 1 and 2). Two vectors are orthogonal if their scalar product is zero; a real matrix $O$ is orthogonal if $O^TO=I$.
\item \textbf{out of equilibrium} (Chapter 17). A system whose state is not the thermal equilibrium state; Sakharov's third condition, since in equilibrium the average baryon number vanishes.
\end{itemize}

\subsection{Glossary: P to R}

\begin{itemize}
\item \textbf{pairing symmetries} (Chapter 16). The exact maps of Chapter 15 that relate solutions with the masses $m$ and $-m$: the chirality map T1 and the mirror map T2.
\item \textbf{pairing theorems T1, T2, T3} (Chapter 0; proved in Chapter 15). T1: the chirality map turns every configuration with $(m,\lambda)$ into one with $(-m,-\lambda)$ and reverses energy--momentum and charge; T2: the mirror map turns $(m,\lambda)$ into $(-m,\lambda)$ and keeps them; T3: the Kohn--Sham form of the pairing, a swap of the two components of every $2\times2$ block with transformed boundary conditions.
\item \textbf{parity} (Chapters 2, 13 and 17). Of an element of Pin(4,4), whether it is a product of an even or an odd number of unit gammas; of a Kohn--Sham orbital, its behaviour under the brane reflection $y\to-y$ (the two parity sectors are the two boundary conditions at $y=0$); in particle physics (P), the reflection of space $\mathbf x\to-\mathbf x$.
\item \textbf{part connected to 1} (Chapter 8). The part of a group that products of exponentials $\exp(\theta S^{ab})$ reach; for Spin(4,4) it is smaller than the whole group.
\item \textbf{partial derivative} (Chapter 1). The derivative with respect to one variable while the others are kept fixed, $\partial_\mu=\partial/\partial x^\mu$ (written $\partial/\partial x_\mu$ in Chapter 9, which keeps the notebook's lower-index names; Section 1.8).
\item \textbf{particle, hole, particle--hole spectrum} (Chapter 12). An excitation moves a particle from an occupied level (leaving a hole) to an empty one; the list of the differences $\varepsilon_a-\varepsilon_i$ is the particle--hole spectrum, and in Chapter 13 the file \texttt{particle\allowbreak{}-hole.cs\allowbreak{}v} of each run.
\item \textbf{particle level, sea level} (Chapter 13). The classification of the Kohn--Sham levels of Chapter 13 by continuity from the free problem: a level whose free partner has $\varepsilon_{\mathrm{free}}\ge0$ is a particle level, otherwise a sea level; the exact zero modes count as particle levels by convention.
\item \textbf{partner of mass $-M$} (Chapter 15). For a universe of mass $+M$, a configuration or Kohn--Sham state with the mass parameter $m=-M$.
\item \textbf{path} (Chapter 5). A function $q(t)$ on a time interval; the action assigns a number to every path.
\item \textbf{Pauli matrices} (Chapter 2). The Hermitian, traceless $2\times2$ matrices $\sigma_x=\begin{pmatrix}0&1\\1&0\end{pmatrix}$, $\sigma_y=\begin{pmatrix}0&-i\\i&0\end{pmatrix}$, $\sigma_z=\mathrm{diag}(1,-1)$; Chapters 4 and 6 use the numbered names $\sigma_1=\sigma_x$, $\sigma_2=\sigma_y$, $\sigma_3=\sigma_z$ (Section 2.3).
\item \textbf{periodic boundary conditions} (Chapter 12). The requirement that wave functions repeat from one face of a box to the opposite one (a torus); the allowed momenta are then $2\pi/\ell$ times integers.
\item \textbf{permanent} (Chapter 14). The determinant without signs, $g_{00}g_{11}+g_{01}g_{10}$ for a $2\times2$ matrix; Gaussian averages are permanents, fermion averages determinants.
\item \textbf{permutation} (Chapter 1). A rearrangement of $0,\dots,n-1$; its sign is $(-1)^{\text{number of inversions}}$.
\item \textbf{phantom} (Chapter 11). A fluid with $\rho>0$ and $w<-1$, that is $\rho+p<0$; it violates the null energy condition. The phantom epoch of EXP-3 is an artefact of the mean field.
\item \textbf{phase} (Chapter 1). The angle $\varphi$ in the polar form $z=|z|e^{i\varphi}$; a phase factor is a number $e^{i\alpha}$.
\item \textbf{physical momentum} (Chapter 11). $K_j=k_j/h_j$, the comoving momentum divided by the scale factor of its direction.
\item \textbf{Pin(4,4), Spin(4,4)} (Chapter 2). The groups of products of unit gammas $\gamma(u_1)\cdots\gamma(u_k)$, and of products with an even number of factors; they act on spinors and are double covers of O(4,4) and SO(4,4).
\item \textbf{plasma of thermal pairs} (Chapter 13). The Kohn--Sham gas at $T=m$, dominated by thermal particle--antiparticle pairs: the energy window of the solver then contains about 75000 levels, and the few added particles hardly change the energy and the entropy.
\item \textbf{Poisson bracket} (Chapter 8). $\{F,G\}=\frac{\partial F}{\partial q}\frac{\partial G}{\partial p}-\frac{\partial F}{\partial p}\frac{\partial G}{\partial q}$; with it $\dot F=\{F,H\}$.
\item \textbf{polar form} (Chapter 1). $z=|z|e^{i\varphi}$.
\item \textbf{PowerShell}: see command shell.
\item \textbf{prescribed background} (Chapter 9). A metric that is taken as given rather than solved for; the primordial field is used this way.
\item \textbf{pressure} (Chapter 5). The force per area exerted by a fluid, the spatial diagonal components $T^i{}_i$ of the energy--momentum tensor for an observer at rest; in eight dimensions there are seven principal pressures $p_{(i)}$.
\item \textbf{primary constraints} (Chapter 8). The relations between momenta and fields that follow directly from a first-order Lagrangian, such as $\Pi_a-\sqrt{|g|}(\Psi^\dagger C\gamma^{x_4})_a\approx0$.
\item \textbf{primary gamma matrices} (Chapter 2). The notebook's gammas $\gamma^a$ in block form, built from eight real $8\times8$ matrices $\tau_a$ and their partners $\bar\tau_a$; the gammas of the whole project.
\item \textbf{primitive integer normalization, row-major order} (Chapter 3). The rule that fixes an intertwiner completely: its entries are scaled to integers without a common factor, with the first nonzero entry in row-major order (row 0 from left to right, then row 1, and so on) positive.
\item \textbf{primordial field} (Chapter 9). The eight-dimensional metric of the notebook, in which 3-space inflates with $e^{a_4}$ and the three extra times deflate with $e^{-a_4}$; its static member (constant $a_4$) with a Z2 brane is the background of Chapters 13 to 15.
\item \textbf{primordial window} (Chapter 11). The smooth switch-on and switch-off of the expansion used in EXP-1, $a_4'(t)=\tfrac A4(1+\tanh\frac{t-t_1}{\Delta})(1-\tanh\frac{t-t_2}{\Delta})$.
\item \textbf{probability density} (Chapter 14). A function $p\ge0$ with $\int p=1$ that describes a random number; averages are $\langle g\rangle=\int g\,p$.
\item \textbf{projector} (Chapter 1). A matrix with $P^2=P$; for example the chirality projectors $P_\mp=\tfrac12(1\mp\gamma^8)$.
\item \textbf{proper densities, coordinate densities} (Chapter 13). Densities per unit proper volume, such as $n_p$ and $S_p$, and per unit $y$ and coordinate 3-volume, $n_c=e^{6Hy}n_p$; the total particle number is $N=\ell^3\int n_c\,dy$.
\item \textbf{PROVED} (Chapter 0). The status of a statement derived in the book from stated assumptions, usually with an exact machine check.
\item \textbf{Prüfer angle, Prüfer index} (Chapters 10 and 13). The angle $\theta$ in $a=r\cos\theta$, $b=-r\sin\theta$ of the two real components of an orbital; its end value increases with the energy, and the integer $n$ in the condition it must meet at $y=0$ numbers the levels (the Prüfer index), so that none is missed.
\item \textbf{Pulay mixing}: see Anderson mixing.
\item \textbf{Q1, Q2, Q3} (Chapter 16). The three questions into which Chapter 16 splits the hypothesis of pair creation: is it allowed by the conservation laws, does it happen in a classical solution, and how likely is it.
\item \textbf{quantization} (Chapter 8). The passage from classical fields to operators on a space of states; in this book canonical and formal (no regularization or renormalization).
\item \textbf{quantum dynamics} (Chapter 16). The quantum theory that gives amplitudes and rates of transitions; it would be needed to say how likely the creation of a pair of universes is, and it is not computed.
\item \textbf{quantum field} (Chapter 6). A field whose components are operators that create and annihilate particles and anticommute (for fermions).
\item \textbf{quaternions} (Chapter 3). The four-dimensional associative but non-commutative algebra with $i^2=j^2=k^2=ijk=-1$.
\item \textbf{$q$ vector} (Chapter 9). The Stage-2 document's name for the difference between the notebook's stored expressions and their Clifford-consistent rebuild (check \texttt{P\_notebo\allowbreak{}okCompar\allowbreak{}e\_qTermF\allowbreak{}romNonCl\allowbreak{}iffordEx\allowbreak{}traTimeG\allowbreak{}ammas}).
\item \textbf{random phase} (Chapter 14). A random complex amplitude whose distribution depends only on its modulus, so that all phases are equally likely.
\item \textbf{rank} (Chapter 2). The largest number of independent elements in a list of matrices or rows, the dimension of their span.
\item \textbf{real restriction} (Chapter 6). Setting the imaginary part of the commuting field of dirac16complex00 to zero; it gives the notebook's real Lagrangian $\mathrm{Lg}[\,]$ with the canonical connection.
\item \textbf{redshift} (Chapter 11). $z$ with $1+z=1/a$, the stretching of wavelengths of light emitted when the scale factor was $a$.
\item \textbf{reduced frequency} (Chapter 11). $\mu=M_{\mathrm{eff}}(a=1)/H_0$, a small stand-in (3 or 7) for the enormous ratio of a real fermion mass to the Hubble rate, used in EXP-3.
\item \textbf{refined convergence} (Chapter 11). A checker test that reruns a program with ten times tighter tolerances and requires the errors to shrink or stay at an identified floor; the check \texttt{refinedC\allowbreak{}onvergen\allowbreak{}ce}.
\item \textbf{reflection} (Chapters 2 and 17). $R_u(v)=v-2\frac{\eta(u,v)}{\eta(u,u)}u$ for a unit vector $u$, which reverses $u$ and keeps its orthogonal hyperplane; in Chapter 17 a reflection $x_a\to r_ax_a$ of a set of coordinates.
\item \textbf{regula falsi} (Chapter 10). The secant rule applied to the two ends of a bracket, keeping the part in which the sign changes; the Stage-4 solver uses its Illinois variant.
\item \textbf{regulator} (Chapter 12). A small auxiliary number, such as $\kappa$ in $e^{-\kappa R}/R$, that makes an integral converge during a calculation and is sent to zero at the end.
\item \textbf{relative residual, truncation estimate} (Chapter 10). The quantities with which the Mathematica notebook tests stored solutions: the largest difference between a derivative computed from the samples and the right-hand side, relative to the size of the right-hand side, and an estimate of the error of that derivative.
\item \textbf{relative tolerance, rtol} (Chapter 10). The relative error per step that CVODE accepts; with the absolute tolerance it defines the weights of the error norm.
\item \textbf{repeat-run byte identity} (Chapter 11). A checker test that reruns a program into a fresh folder and requires every file to be byte-identical; the check \texttt{repeatBy\allowbreak{}teIdenti\allowbreak{}ty}.
\item \textbf{report} (Chapter 0). A JSON file written by a verifier, with the named checks and their values \texttt{\detokenize{true}} or \texttt{\detokenize{false}}, the measurements and the hashes of the inputs.
\item \textbf{repository} (Chapter 19). A folder of files together with its complete recorded history, the sequence of its commits; the repository of this book is https://github.com/once-ere/Dirac\_claude.
\item \textbf{representation} (Chapter 2). A homomorphism $\rho$ from a group to invertible matrices, $\rho(gh)=\rho(g)\rho(h)$; the group then acts on the columns.
\item \textbf{representation theorem} (Chapter 2). Under Pin(4,4) the 16 components of a spinor form one irreducible block; under Spin(4,4) they split into two inequivalent irreducible halves of 8.
\item \textbf{restricted, unrestricted Hartree--Fock}: see Hartree--Fock approximation.
\item \textbf{reversed Krein metric} (Chapter 15). The matrix $-B$ in the anticommutator of the image field $\gamma^8\Psi$; it is the canonical structure of the Lagrangian $-\mathcal L_{-m,-\lambda}$.
\item \textbf{Ricci tensor, Ricci scalar} (Chapter 4). $R_{\mu\nu}=R^\rho{}_{\mu\rho\nu}$ and $R=g^{\mu\nu}R_{\mu\nu}$, contractions of the Riemann tensor.
\item \textbf{Richardson's rule} (Chapter 10). An estimate of the error of a result computed with step $h$ from the difference to the result with step $h/2$, for a method of known order.
\item \textbf{Riemann tensor} (Chapter 4). $R^\rho{}_{\sigma\mu\nu}$, the curvature: the change of a vector transported around a small loop.
\item \textbf{right-hand side} (Chapter 10). The function $f$ in $Y'=f(t,Y)$; the cost of a solver is counted in its evaluations.
\item \textbf{RK4}: see Runge--Kutta method.
\item \textbf{rtol}: see relative tolerance.
\item \textbf{Runge--Kutta method, RK4} (Chapter 10). The classical one-step method of order 4 with four evaluations of the right-hand side per step; used by the independent checkers of EXP-4 and EXP-5.
\item \textbf{rustSolveIt} (Chapter 10). The repositories of the author's Rust translation of SUNDIALS (\texttt{sundials\allowbreak{}\_rs}), which the setup scripts clone at a pinned commit into \texttt{vendor/r\allowbreak{}ustSolve\allowbreak{}It}; with the \texttt{\detokenize{win11}} platform it is the numerical engine of Stages 3 to 5.
\end{itemize}

\subsection{Glossary: S to Z}

\begin{itemize}
\item \textbf{Sakharov's conditions} (Chapter 17). The three conditions for creating a baryon excess from a symmetric start: a process that violates baryon number, violation of C and of CP, and a departure from thermal equilibrium. The theory of this book, which contains no baryons and no baryon-number-violating interaction, fails the first, because its U(1) charge is exactly conserved (Theorem M1); it fails the second for its Lagrangian, because each field has an exact symmetry that reverses the charge: C for dirac16complex00, in every gravitational field, and CP for dirac16complex, in flat space and in gravitational fields with the corresponding reflection isometry (Theorem M2); and no computation addresses the third (the scorecard M6 of Chapter 17).
\item \textbf{scalar} (Chapter 1). A single number, as opposed to a vector or matrix; in geometry a quantity that does not change under changes of coordinates or frames.
\item \textbf{scalar density, $S$} (Chapter 6). $S=\bar\Psi\Psi=\Psi^\dagger C\Psi$, the frame-independent bilinear on which the mass term and the interaction depend.
\item \textbf{scalar field} (Chapter 5). A field with one component that does not change under changes of frame, such as $\phi(x)$ with the Klein--Gordon Lagrangian.
\item \textbf{scalar product} (Chapter 1). $\langle u,v\rangle=\sum_iu_i^\ast v_i$ for complex columns (the dot product for real ones).
\item \textbf{scale factor} (Chapter 9). The number that multiplies a coordinate interval to give its proper length, for example $s^{1/6}e^{a_4}$ for the directions of 3-space in the primordial field.
\item \textbf{Schrödinger equation} (Chapter 8). $i\partial_t\psi=H\psi$, the time evolution of a quantum state.
\item \textbf{Schur's lemma} (Chapter 2). The commutant of an irreducible complex representation consists of multiples of the identity; used to prove that the gammas and the two halves are irreducible and inequivalent.
\item \textbf{SEC}: see strong energy condition.
\item \textbf{secant rule} (Chapter 10). A root finder that replaces the function by the straight line through the last two points and takes its zero as the next point.
\item \textbf{second class} (Chapter 8). Constraints whose brackets with each other form an invertible matrix; they are solved with the Dirac bracket.
\item \textbf{second quantization} (Chapter 6). Turning a classical field into operators that create and annihilate particles; in Chapter 12 the occupation-number description of many fermions.
\item \textbf{self-consistency} (Chapter 12). The condition that the orbitals computed in a potential give back the density or density matrix from which the potential was built; reached by iteration.
\item \textbf{set, subset, empty set} (Chapter 2). A set is a collection of distinct elements; a subset contains only elements of the set; the empty set has none. The subsets $A$ of $\{0,\dots,7\}$ label the 256 gamma monomials $\gamma_A$.
\item \textbf{sha256} (Chapter 19). A 64-digit hexadecimal fingerprint computed from the bytes of a file; two files have the same sha256 exactly when they are byte-identical, for every practical purpose.
\item \textbf{shells} (Chapter 13). The groups of lattice momenta with the same $|\mathbf k|^2=\Delta k^2\,\nu$, $\nu=n_1^2+n_2^2+n_3^2$, whose Kohn--Sham levels coincide.
\item \textbf{shift bound, window floor} (Chapter 15). The Stage-4 solver searches for levels above a fixed negative energy, the window floor ($-4.594$ for $m=1$); the shift bound $\max|M_{\mathrm{eff}}-m|+\max|v_x|$ limits how far the interaction can move a level, and the solver requires it to be at most the absolute value of the floor (its window premise).
\item \textbf{shooting method} (Chapter 10). Solving a boundary-value or eigenvalue problem by integrating from one end with trial values and adjusting them until the condition at the other end holds.
\item \textbf{signature} (Chapter 1). The numbers of positive and negative eigenvalues of a symmetric matrix, such as (4,4) for $\eta$; by Sylvester's law of inertia it does not depend on the basis.
\item \textbf{signed permutation matrix} (Chapter 2). A matrix with exactly one nonzero entry $\pm1$ in every row and every column; it is orthogonal. The notebook's gammas are of this kind.
\item \textbf{singularity} (Chapter 16). A time at which the volume of space goes to zero and the equations break down; the notebook's metric has none at $x_4=0$.
\item \textbf{Slater determinant} (Chapter 12). The antisymmetric wave function $\det[\varphi_a(x_i)]/\sqrt{N!}$ of $N$ fermions in $N$ orthonormal orbitals.
\item \textbf{slices} (Chapter 8). The seven-dimensional surfaces $x_4=\mathrm{const}$ on which the canonical quantization is set up.
\item \textbf{smearing} (Chapter 12). Replacing integer occupations by Fermi--Dirac occupations of a small width to make a self-consistency loop converge; the result is an ensemble. One run of Chapter 13 (smeared $N=1016$) needs it.
\item \textbf{smooth} (Chapter 1). Having continuous derivatives of every order.
\item \textbf{sound speed} (Chapter 11). $c_s^2=dp/d\rho$; a negative $c_s^2$ makes small perturbations grow instead of oscillate.
\item \textbf{source-free} (Chapter 15). Gravitational field equations with zero right-hand side. A chirality pair of two independent classical fields in the correlated configuration $\Psi_+$, $\Psi_-=\gamma^8\Psi_+$ (the second with the Lagrangian $\mathcal L_{-m,-\lambda}$) has zero total energy--momentum, so as the only source it leaves the source-free equations. This holds for classical fields only: for the quantized field no reading of the pair examined in the repository gives two universes whose energy--momentum cancels (Chapter 15).
\item \textbf{space-like, time-like} (Chapter 1). A vector with positive, respectively negative, squared length; in signature (4,4) the frame directions 0 to 3 are space-like and 4 to 7 time-like.
\item \textbf{span} (Chapter 2). The set of all linear combinations of a list of vectors or matrices.
\item \textbf{spectral theorem} (Chapter 1). A Hermitian matrix has real eigenvalues and an orthonormal basis of eigenvectors.
\item \textbf{sphaleron} (Chapter 17). A process of the electroweak theory that violates baryon number at high temperature.
\item \textbf{spin connection} (Chapter 4). The antisymmetric $\omega_{\mu ab}$ that makes the frame covariantly constant; see canonical spin connection. The notebook's contraction of it is not the canonical one (Chapter 4).
\item \textbf{spin generators} (Chapter 2). $S^{ab}=\tfrac14[\gamma^a,\gamma^b]$; their exponentials rotate spinors, and they commute with $\gamma^8$.
\item \textbf{spinor} (Chapter 2). A column of 16 components on which the gammas and the spin transformations act; a spinor in eight dimensions has at least 16 components.
\item \textbf{spinor connection, $\Omega_\mu$} (Chapter 4). $\Omega_\mu=\tfrac12\omega_{\mu ab}S^{ab}=\tfrac18\omega_{\mu ab}[\gamma^a,\gamma^b]$, the term in $D_\mu=\partial_\mu+\Omega_\mu$ that makes the derivative of a spinor covariant; its curvature is $F_{\mu\nu}$.
\item \textbf{spinor field} (Chapter 6). A field that attaches to every point a column of 16 components measured in the local frame, changing by a spin transformation when the frame is turned.
\item \textbf{spinor norm} (Chapter 2). $N(g)=n(u_1)\cdots n(u_k)=\pm1$ for $g=\gamma(u_1)\cdots\gamma(u_k)$ in Pin(4,4), from $g^TCg=(-1)^kN(g)C$.
\item \textbf{split octonions} (Chapter 3). The eight-dimensional composition algebra with a norm of signature (4,4); the notebook's $\tau_a$ are multiplications by split octonions.
\item \textbf{square matrix} (Chapter 1). A matrix with as many rows as columns.
\item \textbf{stability region, stable} (Chapter 10). A method is stable for $z=h\lambda$ if its amplification factor has $|R(z)|\le1$; the set of such $z$ is its stability region.
\item \textbf{standard basis} (Chapter 1). The vectors $e_0,\dots,e_{n-1}$ of $\mathbb R^n$, $e_k$ having the component 1 in place $k$ and 0 elsewhere.
\item \textbf{standard model of cosmology, $\Lambda$CDM} (Chapter 11). The model with radiation, ordinary matter, cold dark matter and a cosmological constant ($w=-1$), the reference against which the experiments of Chapter 11 compare a varying $w$.
\item \textbf{start value} (Chapter 10). The value of the solution at the initial time of an initial-value problem.
\item \textbf{state vector} (Chapter 10). The column $Y(t)$ of all unknowns of a system of ordinary differential equations.
\item \textbf{static member} (Chapter 13). The member of the primordial family with constant $a_4=a_{4,0}$, a metric that does not change in time; the background of the Kohn--Sham states.
\item \textbf{stationary} (Chapter 5). An action is stationary on a path if it does not change to first order under every variation with fixed end points.
\item \textbf{stationary state} (Chapter 12). An eigenstate of the Hamiltonian; its time dependence is only a phase.
\item \textbf{statistics sign, $s$} (Chapter 6). $s=+1$ for commuting and $s=-1$ for Grassmann components, produced by reordering factors, as in $\Psi^TY\Phi^\ast=s\,\Phi^\dagger Y^T\Psi$; it decides the sign of exchange (Chapter 14) and which charge conjugations are symmetries (Chapter 17).
\item \textbf{step size} (Chapter 10). The distance $h$ between two consecutive points of a numerical solution; CVODE chooses it adaptively.
\item \textbf{stiff} (Chapter 10). A problem with components that decay much faster than the solution of interest changes; explicit methods then need steps as small as the fastest decay time.
\item \textbf{stop time} (Chapter 10). A time beyond which CVODE must not integrate; every integration of the project sets it to its last output time.
\item \textbf{strength} (Chapter 13). For a configuration $(m,L,N)$, the largest first-order size of the pseudo-potential per unit coupling in the free ground state; the couplings are $\hat\lambda_1=0.1/\text{strength}$ and $\hat\lambda_2=1/\text{strength}$.
\item \textbf{strong energy condition, SEC} (Chapter 9). $(T_{\mu\nu}-g_{\mu\nu}T/(D-2))u^\mu u^\nu\ge0$ for time-like $u$, equivalent through the Einstein equations to $R_{\mu\nu}u^\mu u^\nu\ge0$.
\item \textbf{subgroup} (Chapter 2). A subset of a group that is itself a group under the same product.
\item \textbf{subspace} (Chapter 2). A set of vectors or matrices that contains every linear combination of its elements; the span of any list is one.
\item \textbf{SUNDIALS} (Chapter 10). A library of solvers for differential equations from the Lawrence Livermore National Laboratory; CVODE is its initial-value solver, used here in the Rust translation \texttt{sundials\allowbreak{}\_rs}.
\item \textbf{Sylvester's law of inertia} (Chapter 1). The numbers of positive, negative and zero eigenvalues of a symmetric matrix do not change under $A\to P^TAP$ with invertible $P$.
\item \textbf{symmetric} (Chapter 1). A square matrix with $A^T=A$.
\item \textbf{symmetric difference} (Chapter 2). $A\triangle B$, the set of indices in exactly one of $A$ and $B$; $\gamma_A\gamma_B=\pm\gamma_{A\triangle B}$.
\item \textbf{symmetrized} (Chapter 6). The kinetic term with the derivative acting once on $\Psi$ and once on $\bar\Psi$ with opposite signs, which makes the Lagrangian real without a factor $i$.
\item \textbf{symmetry of the dynamics} (Chapter 17). A transformation that turns every possible history into another possible history running in the same direction of time.
\item \textbf{system} (Chapter 10). Several ordinary differential equations for several unknown functions, collected in a state vector.
\item \textbf{T1, T2, T3}: see pairing theorems, chirality map, mirror map.
\item \textbf{tangent CPL parameters} (Chapter 11). The values $w_0=w(1)$ and $w_a=-dw/da$ at $a=1$ of a model's equation of state, compared with the CPL fits of supernova data.
\item \textbf{tensor} (Chapter 4). An object with upper and lower indices whose components transform with one factor $\partial x'/\partial x$ or $\partial x/\partial x'$ per index.
\item \textbf{tensor product, Kronecker product} (Chapter 1). $A\otimes B$, the block matrix with blocks $a_{ij}B$.
\item \textbf{Theorems M1 to M6} (Chapter 17). The results of the matter--antimatter analysis: M1 exact U(1) symmetry and charge conservation for both fields; M2 the classification of C, P, CP and their combinations; M3 the charge-violating terms that the symmetry allows (no Majorana mass term for the anticommuting field); M4 the chirality pair, for two independent classical fields in the correlated configuration $\Psi_-=\gamma^8\Psi_+$, with zero total charge and energy--momentum (for classical fields only: at the quantum level no reading gives a cancellation between two independent universes); M5 a conditional scenario stated as a hypothesis, which predicts no value of the asymmetry; M6 the Sakharov scorecard (the first and the second condition fail, the third is not addressed).
\item \textbf{thermal particle--antiparticle pairs} (Chapter 13). At $T>0$ particle and sea levels have fractional weights; the gas contains pairs that do not change $N$.
\item \textbf{Thomas--Fermi model} (Chapter 12). The 1927 approximation of the kinetic energy by that of a uniform gas, $\int C_Fn^{5/3}$; too crude, which is why Kohn--Sham keeps the kinetic energy exact.
\item \textbf{3-space, time, extra times, hidden space} (Chapter 9). The coordinates $x_1,x_2,x_3$ (ordinary space), $x_4$ (the time of evolution), $x_5,x_6,x_7$ (the extra times) and $x_0$ (the hidden space).
\item \textbf{tip} (Chapter 9). The end $\zeta\to-\infty$ of the hidden direction of the static primordial field, where the warp factor goes to 0; the Kohn--Sham problem cuts it off at $y=-L$ (Chapter 13).
\item \textbf{total derivative, total divergence} (Chapter 5). A term $\frac{d}{dt}F$ or $\partial_\mu V^\mu$; added to a Lagrangian it does not change the field equations.
\item \textbf{totally antisymmetric part} (Chapter 6). The part of $\omega_{cab}$ antisymmetric in all three indices; only it enters the Lagrangian, so a diagonal vielbein contributes no spin connection to it.
\item \textbf{trace} (Chapter 1). $\mathrm{tr}A=\sum_iA_{ii}$; $\mathrm{tr}(AB)=\mathrm{tr}(BA)$.
\item \textbf{traceless} (Chapter 2). A matrix with trace zero; every gamma monomial except the unit is traceless.
\item \textbf{transformation law of the metric} (Chapter 1). $g'_{\alpha\beta}=\frac{\partial x^\mu}{\partial x'^\alpha}\frac{\partial x^\nu}{\partial x'^\beta}g_{\mu\nu}$, in matrix form $g'=J^TgJ$, under a change of coordinates.
\item \textbf{transformed boundary condition} (Chapter 15). The tip condition of the $-M$ universe with the bag angle $\pi-\theta$ instead of $\theta$, needed for T3; with the untransformed condition the $-M$ universe is the control.
\item \textbf{transition state} (Chapter 12). Slater's estimate of an excitation energy from the Kohn--Sham levels with half a particle moved.
\item \textbf{transpose} (Chapter 1). $A^T$, the matrix with rows and columns exchanged.
\item \textbf{trapezoidal rule} (Chapter 10). The implicit method $Y_{n+1}=Y_n+\tfrac h2(f_n+f_{n+1})$, with $R(z)=(1+z/2)/(1-z/2)$.
\item \textbf{triality} (Chapter 2). The symmetry of Spin(4,4) that permutes the vector representation and the two spinor halves; quoted, not used.
\item \textbf{twin} (Chapters 6 and 19). The PowerShell or the Bash version of a gate, which run the same steps and print the same final line (Chapter 19); for checks, the sympy check \texttt{\detokenize{S5_x}} that repeats the Wolfram check \texttt{\detokenize{C00_x}} or \texttt{\detokenize{PAIR_x}} with the same ending (Chapter 6).
\item \textbf{twisted adjoint action} (Chapter 2). The map $\Lambda_{\mathrm t}$ from Pin(4,4) to O(4,4) defined by $\alpha(g)\gamma(v)g^{-1}=\gamma(\Lambda_{\mathrm t}(g)v)$ with $\alpha(g)=(-1)^{|g|}g$; the untwisted map $\Lambda_{\mathrm u}$ omits $\alpha$. The kernel of $\Lambda_{\mathrm t}$ is $\{\pm1\}$.
\item \textbf{U(1) charge 2} (Chapter 17). The property of a term such as $\Psi^TM\Psi$ that is multiplied by $e^{2i\alpha}$ under $\Psi\to e^{i\alpha}\Psi$; such a term violates charge conservation.
\item \textbf{U(1) symmetry, phase symmetry} (Chapter 7). The invariance of the Lagrangian under $\Psi\to e^{i\alpha}\Psi$ with a constant $\alpha$; its Noether current is the conserved current and its charge the U(1) charge (Theorem M1 of Chapter 17).
\item \textbf{ulp} (Chapter 10). The unit in the last place, the gap between a double and the next one; $\mathrm{ulp}(1)=2.2\times10^{-16}$.
\item \textbf{ultrahyperbolic} (Chapter 5). A wave equation with more than one time-like direction, such as the Klein--Gordon equation in signature (4,4); its initial-value problem is not well-posed.
\item \textbf{ultrastatic} (Chapter 18). A diagonal metric with $g_{44}=-1$ none of whose entries depends on $x_4$; the static member of the primordial field is one.
\item \textbf{uniform gas} (Chapter 12). Infinitely many particles at constant density; its exchange and correlation energies give the local density approximation.
\item \textbf{unit} (Chapter 3). An element $1$ with $1x=x1=x$ for every $x$ of an algebra.
\item \textbf{unitary} (Chapter 1). A complex matrix with $U^\dagger U=I$; it preserves scalar products. An antiunitary operator does so up to complex conjugation (Chapter 17).
\item \textbf{unitary type, antiunitary type} (Chapter 17). A map of the field is of unitary type if it preserves the canonical anticommutation relations linearly and of antiunitary type if it does so antilinearly; passing the test is necessary, not sufficient, for an operator that implements the map.
\item \textbf{unit normal} (Chapter 9). The unit vector $n=\partial_y$ orthogonal to a surface $y=\mathrm{const}$.
\item \textbf{unit test} (Chapter 19). A small program in the folder \texttt{\detokenize{tests/}} that tests one tool of the repository or pins one piece of committed evidence, for example a check count quoted in a document; Python's module \texttt{\detokenize{unittest}} finds and runs them.
\item \textbf{unit vector} (Chapter 2). A vector $u$ with $\eta(u,u)=\pm1$; its gamma $\gamma(u)$ is invertible.
\item \textbf{universal functional} (Chapter 12). $F[n]=\min_{\Psi\to n}\langle\Psi|\hat T+\hat W|\Psi\rangle$, the same for every external potential.
\item \textbf{universe of mass $+M$} (Chapter 15). A configuration, solution or Kohn--Sham state with the mass parameter $m=+M$; its partner of mass $-M$ has $m=-M$.
\item \textbf{untransformed control} (Chapter 15). The $-M$ universe computed with the Stage-4 bag condition $\theta=0$ instead of the transformed one; it shows that the transformation of the boundary condition matters.
\item \textbf{vacuum} (Chapter 8). The state annihilated by all annihilation operators; in the good sector the filled Dirac sea with no particles or antiparticles.
\item \textbf{variation} (Chapter 5). A small change $\epsilon\xi(t)$ of a path (or of a field) that vanishes at the end points.
\item \textbf{vector} (Chapter 1). A column of numbers, or an element of a vector space; in geometry an object with an upper index.
\item \textbf{verifier} (Chapters 0 and 19). In this book an exact program in the Wolfram Language (\texttt{scripts/\allowbreak{}verify\_*\allowbreak{}.wls}) that recomputes statements and writes a report; the independent Python programs that do the same are called checkers.
\item \textbf{vielbein}: see frame.
\item \textbf{vielbein postulate} (Chapter 4). The condition that the frame is covariantly constant, $\partial_\mu e_\nu{}^a-\Gamma^\lambda{}_{\mu\nu}e_\lambda{}^a+\omega_\mu{}^a{}_be_\nu{}^b=0$; it determines the canonical spin connection.
\item \textbf{volume element} (Chapter 1). $\sqrt{|g|}\,d^8x$, the proper volume of a small coordinate box.
\item \textbf{warped, warp factor} (Chapter 9). A metric in which several directions carry a common factor that depends on another coordinate, such as $e^{2H\zeta}$ in the static primordial field.
\item \textbf{wave function} (Chapter 6). A complex function whose squared modulus gives probabilities; in Chapter 12 the many-body wave function $\Psi(x_0,\dots,x_{N-1})$.
\item \textbf{weak energy condition, WEC} (Chapter 9). $T(u,u)\ge0$ for every time-like $u$; for $u=e_4$ it is $\rho\ge0$.
\item \textbf{WEC}: see weak energy condition.
\item \textbf{weighted root-mean-square norm} (Chapter 10). $\|e\|=\sqrt{\frac1n\sum_c(e_c/(\mathrm{rtol}|Y_c|+\mathrm{atol}))^2}$; CVODE accepts a step if it is at most 1.
\item \textbf{well-posed} (Chapter 8). An initial-value problem whose solution exists, is unique and depends continuously on the data (Hadamard); the full dirac16complex problem is not.
\item \textbf{Wheeler--DeWitt equation} (Chapters 16 and 18). The equation of quantum cosmology for a wave function of the geometry and the matter fields; Chapter 18 derives a one-variable version for a reduced model of the primordial family, and no program of the repository solves it.
\item \textbf{Wick's theorem} (Chapter 12). In a Slater determinant or a non-interacting thermal ensemble every expectation value is a sum of products of the one-body density matrix.
\item \textbf{win11} (Chapter 10). The platform name with which \texttt{scripts/\allowbreak{}setup\_so\allowbreak{}lver} installs the pinned engine; the byte identity of the committed numerical files is established with it on Windows 11 and on Ubuntu 24.04.
\item \textbf{window} (Chapter 8). The interval of times, away from the turning point, over which EXP-5 compares the growth of a mode with the WKB estimate.
\item \textbf{WKB approximation} (Chapter 8). The approximation (after Wentzel, Kramers and Brillouin) that treats a slowly changing coefficient as constant over the local growth or oscillation time; it estimates the growth of the extra-time modes of EXP-5.
\item \textbf{working tree} (Chapter 19). The files of a repository as they are on the disk. The command \texttt{git stat\allowbreak{}us --sho\allowbreak{}rt} lists every file of the working tree that is changed or new compared with the last commit (files in ignored folders such as \texttt{\detokenize{build/}} excepted), and prints nothing when there is none.
\item \textbf{Z2 brane} (Chapter 9). The surface $y=0$ at which two mirror copies of the static primordial patch are glued; a choice of model, not a derived result, with a kink in the warp factor that requires a source on the brane.
\item \textbf{Z2-symmetric} (Chapter 15). A configuration with $\Psi(y)=\pm\gamma^0\Psi(-y)$; it solves the field equations on both sides of the brane with one mass function exactly when the mass function is odd, $m(-y)=-m(y)$, where $\Psi\ne0$.
\item \textbf{zero mode} (Chapters 8 and 13). In Chapter 8, the part of the field that does not depend on $x_5,x_6,x_7$; in Chapter 13, the Kohn--Sham level $\varepsilon=0$ with $\chi=(e^{My},0)^T$ at zero momentum, localized at the brane.
\item \textbf{zero vector} (Chapter 1). The vector with all components zero.
\end{itemize}

\subsection{Symbols and notation}

The table lists the symbols that are used in more than one chapter, with the chapter that introduces them. Indices are counted from 0 throughout the book (Chapter 0): coordinates $x_0,\dots,x_7$, frame directions $a,b=0,\dots,7$, spinor components $0,\dots,15$.

\begingroup
\small
\begin{longtable}{@{}>{\raggedright\arraybackslash}p{0.267\linewidth}>{\raggedright\arraybackslash}p{0.267\linewidth}>{\raggedright\arraybackslash}p{0.267\linewidth}@{}}
\toprule
\textbf{Symbol} & \textbf{Meaning} & \textbf{Chapter} \\
\midrule
\endfirsthead
\toprule
\textbf{Symbol} & \textbf{Meaning} & \textbf{Chapter} \\
\midrule
\endhead
$\eta_{ab}$ & the flat metric of signature (4,4): $\eta_{aa}=+1$ for $a\le3$ (space-like), $\eta_{aa}=-1$ for $a\ge4$ (time-like), zero off the diagonal & 1 \\
$x_0$; $x_1,x_2,x_3$; $x_4$; $x_5,x_6,x_7$ & the hidden space; ordinary 3-space; the time of evolution; the three extra times & 9 \\
$\gamma^a$, $a=0,\dots,7$ & the primary gamma matrices, real $16\times16$, $\{\gamma^a,\gamma^b\}=2\eta^{ab}$ & 2 \\
$\gamma^8=\gamma^0\gamma^1\cdots\gamma^7$ & the chirality, $\mathrm{diag}(-I_8,I_8)$ in the notebook basis & 2 \\
$P_\mp=\tfrac12(1\mp\gamma^8)$ & the projectors on the two halves of chirality $\mp1$ & 2 \\
$C=\gamma^0\gamma^1\gamma^2\gamma^3$ & the charge matrix, the notebook's $\sigma_{16}$; real, symmetric, $C^2=1$ & 2 \\
$B=-iC\gamma^4$ & Hermitian, $B^2=1$, signature (8,8); the matrix of the canonical anticommutator and of the Krein form & 2, 8 \\
$S^{ab}=\tfrac14[\gamma^a,\gamma^b]$ & the spin generators & 2 \\
$\Psi$, $\bar\Psi=\Psi^\dagger C$ & the 16-component field and its Dirac adjoint & 2, 6 \\
$s=\pm1$ & the statistics sign: $+1$ for commuting (dirac16complex00), $-1$ for Grassmann (dirac16complex) components & 6 \\
$S=\bar\Psi\Psi$ & the scalar density & 6 \\
$m$, $M$ & the mass of the Lagrangian and the notebook's mass number, $m=-HM$ & 5 \\
$\lambda$, $U(S)=\tfrac\lambda2S^2$ & the coupling and the contact interaction & 6, 13 \\
$M_{\mathrm{eff}}=m+\lambda S$ & the effective mass of the mean field & 11 \\
$\hat\lambda=\lambda m^6$, $\hat\lambda_1$, $\hat\lambda_2$ & the dimensionless coupling of Stage 4 and its two strengths & 13 \\
$e_\mu{}^a$, $g_{\mu\nu}$, $\sqrt{\lvert g\rvert}$ & the vielbein, the metric $g_{\mu\nu}=e_\mu{}^a\eta_{ab}e_\nu{}^b$ and the volume factor & 4 \\
$\Gamma^\rho{}_{\mu\nu}$, $\omega_{\mu ab}$, $\Omega_\mu$, $D_\mu$ & the Christoffel symbols, the spin connection, the spinor connection and the spinor covariant derivative $D_\mu=\partial_\mu+\Omega_\mu$ & 4 \\
$R^\rho{}_{\sigma\mu\nu}$, $R_{\mu\nu}$, $R$, $G_{\mu\nu}$ & the Riemann tensor, the Ricci tensor, the Ricci scalar and the Einstein tensor & 4 \\
$\kappa$ & the gravitational coupling of the eight-dimensional Einstein equations $G^\mu{}_\nu=\kappa T^\mu{}_\nu$ ($\kappa_4$ in four dimensions, Chapter 11) & 4 \\
$\Theta^\mu{}_\nu$, $T^\mu{}_\nu$ & the canonical and the metric energy--momentum tensors & 5, 7 \\
$\rho$, $p_{(i)}$, $w$ & the energy density, the pressures and the equation-of-state parameter & 5 \\
$j^\mu$, $Q$ & a conserved current and its charge & 5, 7 \\
$J^\mu=-i\bar\Psi\gamma^\mu\Psi$, $J^4=\Psi^\dagger B\Psi$ & the Hermitian U(1) current of the book's fields and its charge density; $Q=\int\sqrt{\lvert g\rvert}\,J^4\,d^7x$ & 7 \\
$H$, $z=6Hx_0$, $t=Hx_4$, $a_4(t)$ & the notebook's inverse length, its hidden angle and dimensionless time, and the free function of the primordial field & 9 \\
$\zeta$, $y$ & the proper hidden coordinate of the static primordial field (called $y$ from Chapter 13 on) & 9, 13 \\
$L$, $\theta$ & the tip cutoff $y=-L$ and the bag angle of the tip condition & 13 \\
$h_k$, $E$ & the mode Hamiltonian of a plane wave and its energy, $h_k^2=E^2$ & 8 \\
$\varkappa$ & a growth rate (not the coupling $\kappa$): $\varkappa=\sqrt{-E^2}$ for a growing extra-time mode, $\varkappa=\sqrt{Q^2-m^2}$ in EXP-5 & 8, 11 \\
$j=\pm1$, $\beta=\pm1$ & the block type of a Kohn--Sham block and its Krein sign & 13, 14 \\
$N$, $T$, $\mu$, $f$ & the particle number, the temperature, the chemical potential and the Fermi--Dirac occupation & 12 \\
$n_p$, $S_p$; $n_c$, $S_c$ & the proper and the coordinate number and scalar densities & 13 \\
$E_H$, $E_x$, $\Delta_{KS}$ & the Hartree energy, the exchange energy and the Kohn--Sham gap & 12 \\
$h$, rtol, atol & the step size and the relative and absolute tolerances of a numerical integration & 10 \\
$a$, $z$, $H_0$, $w_0$, $w_a$ & in cosmology: the scale factor, the redshift, today's Hubble rate and the CPL parameters (the redshift $z$ is not the angle $z$ of Chapter 9) & 11 \\
T1, T2, T3 & the three pairing theorems & 15 \\
M1 to M6, H1 to H3 & the results of the matter--antimatter analysis and the three hypotheses of the conditional scenario M5 & 17 \\
$\eta$ (in Chapters 17 and 18) & the baryon-to-photon ratio, not the flat metric & 17 \\
\bottomrule
\end{longtable}
\endgroup

\subsection{How the index of checks was made, and how to look up a check yourself}

\textbf{What is indexed.} A verifier writes its report as a JSON file with an entry \texttt{\detokenize{checks}}, which maps every check name to \texttt{\detokenize{true}} or \texttt{\detokenize{false}} (Section 0.8); the Rust programs write the same kind of entry into the \texttt{summary.\allowbreak{}json} of each run folder, and two Rust files of Stage 4, \texttt{rust/spe\allowbreak{}ctrum/ex\allowbreak{}change-t\allowbreak{}able-che\allowbreak{}ck.json} and \texttt{rust/spe\allowbreak{}ctrum/th\allowbreak{}eory-agr\allowbreak{}eement.j\allowbreak{}son}, store for each check a small record whose entry \texttt{\detokenize{passed}} is \texttt{\detokenize{true}} or \texttt{\detokenize{false}}, next to a text \texttt{\detokenize{detail}}. Sections 20.10 to 20.14 list every check name that Chapters 0 to 19 cite, in running text, in code or in a listing, together with the committed files that record it (the rule follows), its value there, and the chapters that cite it. Three kinds of name are indexed:

\begin{itemize}
\item a single check, such as \texttt{ALG\_clif\allowbreak{}ford};
\item a family written with a final \texttt{\detokenize{*}} or with placeholders in angle brackets: \texttt{PAIR\_T1k\allowbreak{}rein\_*} means every check whose name begins with \texttt{PAIR\_T1k\allowbreak{}rein\_}, \texttt{MA\_M1\_no\allowbreak{}etherIde\allowbreak{}ntity\_<X\allowbreak{}>\_<G>} (Chapter 17) every check whose name begins with \texttt{MA\_M1\_no\allowbreak{}etherIde\allowbreak{}ntity\_}, and \texttt{*\_energy\allowbreak{}\_from\_rh\allowbreak{}o} every Rust self-check whose name ends with \texttt{\_energy\_\allowbreak{}from\_rho} (one per run); the column "Value" then gives the number of checks in the family and whether all of them are true;
\item a name whose ending the text writes separately, such as \texttt{GEO\_viel\allowbreak{}beinPost\allowbreak{}ulate\_G1} and \texttt{\detokenize{_G2}}, or \texttt{hellmann\allowbreak{}FeynmanM\allowbreak{}ass} after the family name \texttt{KS\_funct\allowbreak{}ional} (Chapter 13), is indexed under its full names.
\end{itemize}

Names of \textbf{measurements} (numbers and tables that a report records next to its checks, such as \texttt{ALG\_tauE\allowbreak{}qualsLef\allowbreak{}tMultipl\allowbreak{}icationP\allowbreak{}erIndex} in Chapter 3 or \texttt{canonica\allowbreak{}l\_eigenv\allowbreak{}alues\_de\allowbreak{}tail} in Chapter 19) are not checks and are not indexed; nor are the names of keys of theory files, such as \texttt{fieldEqu\allowbreak{}ations}. A name that the text uses only as a prefix of a whole stage, such as \texttt{\detokenize{S5_}} in "the \texttt{\detokenize{S5_}} twins", is not indexed either; the prefixes are explained at the head of each section.

\textbf{Which files the column "Files" names.} Three kinds of committed file, and no others:

\begin{enumerate}
\item every committed JSON file under \texttt{artifact\allowbreak{}s/} whose \textbf{top-level} entry \texttt{\detokenize{checks}} contains the name: the report of the program that decided the check (for the two Rust files with records, the value shown is the entry \texttt{\detokenize{passed}} of the record);
\item the EXP-3 fit analysis \texttt{numerics\allowbreak{}/exp3/fi\allowbreak{}ts.json}, which has no top-level entry \texttt{\detokenize{checks}} and stores its own checks one level down, in its entry \texttt{validati\allowbreak{}on};
\item the Wolfram reports W00 and WMA (keys below), where they \textbf{cite} a check of Stage 1, 2 or 4: each records the name and the value of every check it cites in an entry \texttt{\detokenize{checks}} inside its measurements (\texttt{citation\allowbreak{}s} in W00, \texttt{M1\_stage\allowbreak{}1ChecksC\allowbreak{}ited} in WMA). All these recorded values are true, as in the cited reports.
\end{enumerate}

Other entries named \texttt{\detokenize{checks}} that lie deeper inside a file are copies or records, not checks decided by the program that wrote the file, and are not indexed: the collected summary \texttt{numerics\allowbreak{}/numeric\allowbreak{}s-summar\allowbreak{}y.json} repeats the checks of the five experiments; the theory files \texttt{pair-cre\allowbreak{}ation/pa\allowbreak{}iring-th\allowbreak{}eory.jso\allowbreak{}n} and \texttt{matter-a\allowbreak{}ntimatte\allowbreak{}r/matter\allowbreak{}-antimat\allowbreak{}ter-theo\allowbreak{}ry.json} list, next to each theorem, the checks that verify it; and the entries \texttt{engine.c\allowbreak{}hecks} of the two Mathematica reports and \texttt{M1\_charg\allowbreak{}eDensity\allowbreak{}.checks} of WMA record sub-steps that no chapter cites. The collected Stage-1 summary \texttt{arbitrar\allowbreak{}y-field/\allowbreak{}stage1-s\allowbreak{}ummary.j\allowbreak{}son} has a top-level entry \texttt{\detokenize{checks}}, but it writes the name of the report in front of every check (such as \texttt{wolfram-\allowbreak{}algebra:\allowbreak{}ALG\_clif\allowbreak{}ford}), so no cited name occurs in it.

\textbf{How it was verified.} The index was produced by a short program that reads every committed JSON file under \texttt{artifact\allowbreak{}s/} of the commit named in Section 20.15, collects the names by the rule just stated, and then searches the chapter files for these names; a second, separately written program read the finished tables back and compared every row, file, value and citation with the committed files and the chapters. Every row of Sections 20.10 to 20.14 therefore names a check that exists, with the value shown, in the committed files shown. The first program also listed the names in the chapters that look like check names but occur in no indexed entry \texttt{\detokenize{checks}}; each of them was read in its sentence and found to be a measurement, a comparison recorded inside a report or a name inside a program, and none of them is called a check by the text.

\textbf{The keys of the files.} To keep the tables of Sections 20.10 to 20.14 short, their column "Files" names each committed file by a key:

\begingroup
\small
\begin{longtable}{@{}>{\raggedright\arraybackslash}p{0.267\linewidth}>{\raggedright\arraybackslash}p{0.267\linewidth}>{\raggedright\arraybackslash}p{0.267\linewidth}@{}}
\toprule
\textbf{Key} & \textbf{Committed file (in \texttt{artifact\allowbreak{}s/dirac1\allowbreak{}6complex\allowbreak{}/})} & \textbf{Written by} \\
\midrule
\endfirsthead
\toprule
\textbf{Key} & \textbf{Committed file (in \texttt{artifact\allowbreak{}s/dirac1\allowbreak{}6complex\allowbreak{}/})} & \textbf{Written by} \\
\midrule
\endhead
WA & \texttt{arbitrar\allowbreak{}y-field/\allowbreak{}wolfram-\allowbreak{}algebra-\allowbreak{}report.j\allowbreak{}son} & Stage 1, Wolfram algebra verifier \\
PA & \texttt{arbitrar\allowbreak{}y-field/\allowbreak{}python-a\allowbreak{}lgebra-r\allowbreak{}eport.js\allowbreak{}on} & Stage 1, Python algebra checker \\
WG & \texttt{arbitrar\allowbreak{}y-field/\allowbreak{}wolfram-\allowbreak{}geometry\allowbreak{}-report.\allowbreak{}json} & Stage 1, Wolfram geometry verifier \\
PG & \texttt{arbitrar\allowbreak{}y-field/\allowbreak{}python-g\allowbreak{}eometry-\allowbreak{}report.j\allowbreak{}son} & Stage 1, Python geometry checker \\
GD & \texttt{arbitrar\allowbreak{}y-field/\allowbreak{}grassman\allowbreak{}n-demo-r\allowbreak{}eport.js\allowbreak{}on} & Stage 1, Grassmann demonstration \\
WP & \texttt{primordi\allowbreak{}al-field\allowbreak{}/wolfram\allowbreak{}-primord\allowbreak{}ial-repo\allowbreak{}rt.json} & Stage 2, Wolfram verifier \\
PP & \texttt{primordi\allowbreak{}al-field\allowbreak{}/python-\allowbreak{}primordi\allowbreak{}al-repor\allowbreak{}t.json} & Stage 2, Python checker \\
E1S & \texttt{numerics\allowbreak{}/exp1/su\allowbreak{}mmary.js\allowbreak{}on} & EXP-1, self-checks of the Rust program \\
E1C & \texttt{numerics\allowbreak{}/exp1/py\allowbreak{}thon-che\allowbreak{}ck-repor\allowbreak{}t.json} & EXP-1, independent checker \\
E2S & \texttt{numerics\allowbreak{}/exp2/su\allowbreak{}mmary.js\allowbreak{}on} & EXP-2, self-checks of the Rust program \\
E2C & \texttt{numerics\allowbreak{}/exp2/py\allowbreak{}thon-che\allowbreak{}ck-repor\allowbreak{}t.json} & EXP-2, independent checker \\
E3S & \texttt{numerics\allowbreak{}/exp3/su\allowbreak{}mmary.js\allowbreak{}on} & EXP-3, self-checks of the Rust program \\
E3C & \texttt{numerics\allowbreak{}/exp3/py\allowbreak{}thon-che\allowbreak{}ck-repor\allowbreak{}t.json} & EXP-3, independent checker \\
E3F & \texttt{numerics\allowbreak{}/exp3/fi\allowbreak{}ts.json} & EXP-3, fit analysis (checks in its entry \texttt{validati\allowbreak{}on}) \\
E4S & \texttt{numerics\allowbreak{}/exp4/su\allowbreak{}mmary.js\allowbreak{}on} & EXP-4, self-checks of the Rust program \\
E4C & \texttt{numerics\allowbreak{}/exp4/py\allowbreak{}thon-che\allowbreak{}ck-repor\allowbreak{}t.json} & EXP-4, independent checker \\
E5S & \texttt{numerics\allowbreak{}/exp5/su\allowbreak{}mmary.js\allowbreak{}on} & EXP-5, self-checks of the Rust program \\
E5C & \texttt{numerics\allowbreak{}/exp5/py\allowbreak{}thon-che\allowbreak{}ck-repor\allowbreak{}t.json} & EXP-5, independent checker \\
WK & \texttt{kohn-sha\allowbreak{}m/wolfra\allowbreak{}m-kohn-s\allowbreak{}ham-repo\allowbreak{}rt.json} & Stage 4, Wolfram exact theory \\
PT & \texttt{kohn-sha\allowbreak{}m/python\allowbreak{}-theory-\allowbreak{}report.j\allowbreak{}son} & Stage 4, sympy exact theory \\
PC & \texttt{kohn-sha\allowbreak{}m/python\allowbreak{}-check-r\allowbreak{}eport.js\allowbreak{}on} & Stage 4, cross-checker (Rust against reference) \\
NB & \texttt{kohn-sha\allowbreak{}m/notebo\allowbreak{}ok-repor\allowbreak{}t.json} & Stage 4, Jupyter notebook \\
RSP & \texttt{kohn-sha\allowbreak{}m/rust/s\allowbreak{}pectrum/\allowbreak{}summary.\allowbreak{}json} & Stage 4, Rust subcommand spectrum \\
RXT & \texttt{kohn-sha\allowbreak{}m/rust/s\allowbreak{}pectrum/\allowbreak{}exchange\allowbreak{}-table-c\allowbreak{}heck.jso\allowbreak{}n} & Stage 4, Rust check of the exchange table (records with \texttt{\detokenize{passed}}) \\
RTA & \texttt{kohn-sha\allowbreak{}m/rust/s\allowbreak{}pectrum/\allowbreak{}theory-a\allowbreak{}greement\allowbreak{}.json} & Stage 4, Rust agreement with the exact theory (records with \texttt{\detokenize{passed}}) \\
RSC & \texttt{kohn-sha\allowbreak{}m/rust/s\allowbreak{}cf/summa\allowbreak{}ry.json} & Stage 4, Rust subcommand scf \\
REX & \texttt{kohn-sha\allowbreak{}m/rust/e\allowbreak{}xcited/s\allowbreak{}ummary.j\allowbreak{}son} & Stage 4, Rust subcommand excited \\
RTH & \texttt{kohn-sha\allowbreak{}m/rust/t\allowbreak{}hermo/su\allowbreak{}mmary.js\allowbreak{}on} & Stage 4, Rust subcommand thermo \\
REM & \texttt{kohn-sha\allowbreak{}m/rust/e\allowbreak{}mt/summa\allowbreak{}ry.json} & Stage 4, Rust subcommand emt \\
W00 & \texttt{pair-cre\allowbreak{}ation/wo\allowbreak{}lfram-di\allowbreak{}rac16com\allowbreak{}plex00-r\allowbreak{}eport.js\allowbreak{}on} & Stage 5, Wolfram, dirac16complex00 (with its citation records) \\
P00 & \texttt{pair-cre\allowbreak{}ation/py\allowbreak{}thon-dir\allowbreak{}ac16comp\allowbreak{}lex00-re\allowbreak{}port.jso\allowbreak{}n} & Stage 5, sympy, dirac16complex00 \\
WPR & \texttt{pair-cre\allowbreak{}ation/wo\allowbreak{}lfram-pa\allowbreak{}iring-re\allowbreak{}port.jso\allowbreak{}n} & Stage 5, Wolfram, pairing theorems \\
PPR & \texttt{pair-cre\allowbreak{}ation/py\allowbreak{}thon-pai\allowbreak{}ring-rep\allowbreak{}ort.json} & Stage 5, sympy, pairing theorems \\
WMA & \texttt{matter-a\allowbreak{}ntimatte\allowbreak{}r/wolfra\allowbreak{}m-matter\allowbreak{}-antimat\allowbreak{}ter-repo\allowbreak{}rt.json} & matter and antimatter, Wolfram (with its citation records) \\
PMA & \texttt{matter-a\allowbreak{}ntimatte\allowbreak{}r/python\allowbreak{}-matter-\allowbreak{}antimatt\allowbreak{}er-repor\allowbreak{}t.json} & matter and antimatter, sympy \\
\bottomrule
\end{longtable}
\endgroup

\textbf{Looking up a check yourself.} The following command prints every committed file whose entry \texttt{\detokenize{checks}} contains a given name, with its value. In PowerShell and in Bash alike (only the quoting of Section 19.3 matters, and the lines between the double quotes are the same):

\begin{Verbatim}[fontsize=\small]
python -c "
import json, pathlib, sys
for p in sorted(pathlib.Path('artifacts').rglob('*.json')):
    d = json.loads(p.read_text(encoding='utf-8'))
    c = d.get('checks') if isinstance(d, dict) else None
    if isinstance(c, dict) and sys.argv[1] in c:
        print(p.as_posix(), c[sys.argv[1]])
" ALG_clifford
\end{Verbatim}

Run from the root of a fresh clone of commit \texttt{\detokenize{1f2dd69}} (Section 19.4), it printed

\begin{Verbatim}[fontsize=\small]
artifacts/dirac16complex/arbitrary-field/python-algebra-report.json True
artifacts/dirac16complex/arbitrary-field/wolfram-algebra-report.json True
\end{Verbatim}

in both shells within a second. With \texttt{canonica\allowbreak{}l\_eigenv\allowbreak{}alues} instead of \texttt{ALG\_clif\allowbreak{}ford} it printed the one line \texttt{artifact\allowbreak{}s/dirac1\allowbreak{}6complex\allowbreak{}/kohn-sh\allowbreak{}am/pytho\allowbreak{}n-check-\allowbreak{}report.j\allowbreak{}son Fals\allowbreak{}e}. With \texttt{theory\_z\allowbreak{}ero\_mode\allowbreak{}\_splitti\allowbreak{}ng} it printed two lines: \texttt{artifact\allowbreak{}s/dirac1\allowbreak{}6complex\allowbreak{}/kohn-sh\allowbreak{}am/rust/\allowbreak{}spectrum\allowbreak{}/summary\allowbreak{}.json Tr\allowbreak{}ue}, and for \texttt{rust/spe\allowbreak{}ctrum/th\allowbreak{}eory-agr\allowbreak{}eement.j\allowbreak{}son} the whole record, \texttt{\{'passed\allowbreak{}': True,\allowbreak{} 'detail\allowbreak{}': ...\}} with a long text in place of the dots, whose entry \texttt{\detokenize{passed}} is the value. The short command looks only at top-level entries \texttt{\detokenize{checks}}, so it does not find the checks of \texttt{numerics\allowbreak{}/exp3/fi\allowbreak{}ts.json} (item 2 above) or the citation records of W00 and WMA (item 3); the index includes them.

\subsection{Index of checks: Stage 1, the field in an arbitrary gravitational field}

The prefixes are \texttt{\detokenize{ALG_}} (algebra), \texttt{\detokenize{GEO_}} (geometry), \texttt{\detokenize{LAG_}} (Lagrangian), \texttt{\detokenize{EMT_}} (energy--momentum tensor), \texttt{\detokenize{QNT_}} (quantization), \texttt{\detokenize{GR_}} (Grassmann demonstration) and \texttt{\detokenize{NEG_}} (negative controls); the endings \texttt{\detokenize{_G1}} and \texttt{\detokenize{_G2}} name the two test geometries (Chapter 4). The Wolfram reports of Stage 5 (W00) and of the matter--antimatter analysis (WMA) record some of these checks, under the same names, in their citation records (Section 20.9, item 3); the column "Files" then lists those reports too. The table has 90 rows; the keys in the column "Files" are those of Section 20.9.

\begingroup
\small
\begin{longtable}{@{}>{\raggedright\arraybackslash}p{0.200\linewidth}>{\raggedright\arraybackslash}p{0.200\linewidth}>{\raggedright\arraybackslash}p{0.200\linewidth}>{\raggedright\arraybackslash}p{0.200\linewidth}@{}}
\toprule
\textbf{Check} & \textbf{Files} & \textbf{Value} & \textbf{Cited in chapters} \\
\midrule
\endfirsthead
\toprule
\textbf{Check} & \textbf{Files} & \textbf{Value} & \textbf{Cited in chapters} \\
\midrule
\endhead
\texttt{ALG\_CAnt\allowbreak{}icommuta\allowbreak{}torGamma\allowbreak{}SSymmetr\allowbreak{}ic} & PG & true & 5 \\
\texttt{ALG\_CCom\allowbreak{}mutatorG\allowbreak{}ammaSAnt\allowbreak{}isymmetr\allowbreak{}ic} & PG & true & 5 \\
\texttt{ALG\_char\allowbreak{}geFormB} & WA, PA, W00, WMA & true & 2, 8, 18 \\
\texttt{ALG\_char\allowbreak{}geMatrix} & WA, PA, WG & true & 0, 2, 3, 5 \\
\texttt{ALG\_chir\allowbreak{}ality} & WA, PA & true & 0, 2, 3 \\
\texttt{ALG\_clif\allowbreak{}ford} & WA, PA & true & 0, 2, 3, 5 \\
\texttt{ALG\_clif\allowbreak{}fordPict\allowbreak{}ureInter\allowbreak{}twiner} & WA, PA & true & 0, 2, 3 \\
\texttt{ALG\_expr\allowbreak{}ession1} & WA, PA & true & 0, 2, 3, 5 \\
\texttt{ALG\_fait\allowbreak{}hful} & WA, PA & true & 2, 4 \\
\texttt{ALG\_fixt\allowbreak{}ureAgree\allowbreak{}ment} & WA, PA & true & 2, 3 \\
\texttt{ALG\_gamm\allowbreak{}a8Map} & WA, PA, W00, WMA & true & 0, 6, 15, 16, 18 \\
\texttt{ALG\_gamm\allowbreak{}aTranspo\allowbreak{}seSymmet\allowbreak{}ry} & WA, PA & true & 2, 5 \\
\texttt{ALG\_gras\allowbreak{}smannLem\allowbreak{}mas} & WG & true & 5 \\
\texttt{ALG\_inva\allowbreak{}riantFor\allowbreak{}ms} & WA, PA, W00, WMA & true & 0, 8, 18 \\
\texttt{ALG\_octo\allowbreak{}nionPict\allowbreak{}ureInter\allowbreak{}twiner} & WA, PA & true & 0, 3 \\
\texttt{ALG\_pinI\allowbreak{}rreducib\allowbreak{}leComple\allowbreak{}x} & WA, PA, W00 & true & 0, 2, 4 \\
\texttt{ALG\_pinL\allowbreak{}iftChara\allowbreak{}cter} & WA, PA, W00, WMA & true & 2, 6 \\
\texttt{ALG\_SabG\allowbreak{}ammaComm\allowbreak{}utator} & PG & true & 2, 4 \\
\texttt{ALG\_spin\allowbreak{}Decompos\allowbreak{}ition} & WA, PA, W00 & true & 0, 2 \\
\texttt{ALG\_spin\allowbreak{}Transpos\allowbreak{}ePropert\allowbreak{}ies} & WA, PA & true & 2, 4, 5 \\
\texttt{ALG\_wolf\allowbreak{}ramAgree\allowbreak{}ment} & PA & true & 2, 3, 19 \\
\texttt{EMT\_cons\allowbreak{}ervation\allowbreak{}\_G1} & WG, PG, W00, WMA & true & 0, 7 \\
\texttt{EMT\_cons\allowbreak{}ervation\allowbreak{}\_G2} & WG, PG, W00, WMA & true & 7 \\
\texttt{EMT\_homo\allowbreak{}geneousR\allowbreak{}eduction} & PG & true & 0, 7, 11 \\
\texttt{EMT\_homo\allowbreak{}geneousR\allowbreak{}eduction\allowbreak{}\_G3} & WG, W00 & true & 0, 7, 11 \\
\texttt{EMT\_homo\allowbreak{}geneousT\allowbreak{}imeSpace\allowbreak{}Vanishes\allowbreak{}OnShell\_\allowbreak{}G3} & PG & true & 7 \\
\texttt{EMT\_onsh\allowbreak{}ellLagra\allowbreak{}ngianSUp\allowbreak{}rimeMinu\allowbreak{}sU\_G1} & PG & true & 7 \\
\texttt{EMT\_onsh\allowbreak{}ellLagra\allowbreak{}ngianSUp\allowbreak{}rimeMinu\allowbreak{}sU\_G2} & PG & true & 7 \\
\texttt{EMT\_symm\allowbreak{}etricHer\allowbreak{}mitian\_G\allowbreak{}1} & WG, W00 & true & 0, 7 \\
\texttt{EMT\_symm\allowbreak{}etricHer\allowbreak{}mitian\_G\allowbreak{}2} & WG, W00 & true & 7 \\
\texttt{EMT\_trac\allowbreak{}e\_G1} & WG, PG, W00 & true & 0, 7 \\
\texttt{EMT\_trac\allowbreak{}e\_G2} & WG, PG, W00 & true & 7 \\
\texttt{EMT\_trac\allowbreak{}eOffShel\allowbreak{}lIdentit\allowbreak{}y\_G1} & PG & true & 7 \\
\texttt{EMT\_vari\allowbreak{}ation} & PG & true & 7 \\
\texttt{EMT\_vari\allowbreak{}ation\_G3} & WG, W00 & true & 7 \\
\texttt{GEO\_anti\allowbreak{}commutat\allowbreak{}orGammaO\allowbreak{}megaVani\allowbreak{}shesDiag\allowbreak{}onal\_G2} & PG & true & 4, 5, 6 \\
\texttt{GEO\_curv\allowbreak{}ature\_G1} & WG, PG & true & 4 \\
\texttt{GEO\_curv\allowbreak{}ature\_G2} & WG, PG & true & 4 \\
\texttt{GEO\_diag\allowbreak{}onalSlas\allowbreak{}hFormula\allowbreak{}\_G2} & WG, PG & true & 4 \\
\texttt{GEO\_diag\allowbreak{}onalSlas\allowbreak{}hFormula\allowbreak{}\_G3} & PG & true & 4 \\
\texttt{GEO\_dive\allowbreak{}rgenceId\allowbreak{}entity\_G\allowbreak{}1} & WG, PG, WMA & true & 4, 5 \\
\texttt{GEO\_dive\allowbreak{}rgenceId\allowbreak{}entity\_G\allowbreak{}2} & WG, PG, WMA & true & 4, 5 \\
\texttt{GEO\_fram\allowbreak{}eNondege\allowbreak{}nerate\_G\allowbreak{}1} & WG, PG & true & 4 \\
\texttt{GEO\_gamm\allowbreak{}aCovaria\allowbreak{}ntConsta\allowbreak{}ncy\_G1} & WG, PG & true & 0, 4, 5 \\
\texttt{GEO\_gamm\allowbreak{}aCovaria\allowbreak{}ntConsta\allowbreak{}ncy\_G2} & WG, PG & true & 0, 4, 5 \\
\texttt{GEO\_lich\allowbreak{}nerowicz\allowbreak{}Constant\allowbreak{}SameG1G2} & PG & true & 4 \\
\texttt{GEO\_lich\allowbreak{}nerowicz\allowbreak{}\_G1} & WG, PG & true & 4, 7 \\
\texttt{GEO\_lich\allowbreak{}nerowicz\allowbreak{}\_G2} & WG, PG & true & 4, 7 \\
\texttt{GEO\_note\allowbreak{}bookCont\allowbreak{}ractionF\allowbreak{}ails\_G1} & WG, PG & true & 0, 4 \\
\texttt{GEO\_note\allowbreak{}bookCont\allowbreak{}ractionF\allowbreak{}ails\_G2} & WG, PG & true & 0, 4 \\
\texttt{GEO\_omeg\allowbreak{}aAntisym\allowbreak{}metry\_G1} & WG, PG & true & 4 \\
\texttt{GEO\_omeg\allowbreak{}aAntisym\allowbreak{}metry\_G2} & WG, PG & true & 4 \\
\texttt{GEO\_prim\allowbreak{}ordialIn\allowbreak{}variants\allowbreak{}\_G2} & WG, PG & true & 4 \\
\texttt{GEO\_sqrt\allowbreak{}gSquared\allowbreak{}EqualsDe\allowbreak{}tg\_G1} & PG & true & 4 \\
\texttt{GEO\_viel\allowbreak{}beinPost\allowbreak{}ulate\_G1} & WG, PG & true & 4 \\
\texttt{GEO\_viel\allowbreak{}beinPost\allowbreak{}ulate\_G2} & WG, PG & true & 4 \\
\texttt{GEO\_wolf\allowbreak{}ramAgree\allowbreak{}ment} & PG & true & 19 \\
\texttt{GR\_bilin\allowbreak{}earOnlyA\allowbreak{}ntisymme\allowbreak{}tricPart\allowbreak{}Survives} & GD & true & 5 \\
\texttt{GR\_compl\allowbreak{}exPsiEqu\allowbreak{}ation} & GD & true & 7 \\
\texttt{GR\_compl\allowbreak{}exQuarti\allowbreak{}cEL} & GD & true & 0, 7 \\
\texttt{GR\_conju\allowbreak{}gationRu\allowbreak{}les} & GD & true & 5 \\
\texttt{GR\_curre\allowbreak{}ntHermit\allowbreak{}ian} & GD, WMA & true & 7, 8 \\
\texttt{GR\_emtHe\allowbreak{}rmitian} & GD & true & 7 \\
\texttt{GR\_kinet\allowbreak{}icSymmet\allowbreak{}ricMatri\allowbreak{}xContras\allowbreak{}t} & GD & true & 5 \\
\texttt{GR\_kinet\allowbreak{}icTotalD\allowbreak{}erivativ\allowbreak{}eReal} & GD & true & 5 \\
\texttt{GR\_lagra\allowbreak{}ngianHer\allowbreak{}mitian} & GD & true & 6 \\
\texttt{GR\_massT\allowbreak{}ermVanis\allowbreak{}hesReal} & GD & true & 5 \\
\texttt{GR\_noteb\allowbreak{}ookLgELT\allowbreak{}rivial} & GD & true & 5 \\
\texttt{GR\_noteb\allowbreak{}ookLgPur\allowbreak{}eDiverge\allowbreak{}nce} & GD & true & 0, 5 \\
\texttt{GR\_quart\allowbreak{}icTermPo\allowbreak{}lynomial} & GD & true & 5 \\
\texttt{GR\_scala\allowbreak{}rBilinea\allowbreak{}rHermiti\allowbreak{}an} & GD & true & 5 \\
\texttt{GR\_unsym\allowbreak{}metrized\allowbreak{}KineticN\allowbreak{}otHermit\allowbreak{}ian} & GD & true & 6 \\
\texttt{LAG\_eule\allowbreak{}rLagrang\allowbreak{}ePsibar\_\allowbreak{}G1} & WG, PG, W00, WMA & true & 0, 7 \\
\texttt{LAG\_eule\allowbreak{}rLagrang\allowbreak{}ePsibar\_\allowbreak{}G2} & WG, PG, W00, WMA & true & 7 \\
\texttt{LAG\_eule\allowbreak{}rLagrang\allowbreak{}ePsi\_G1} & WG, PG, W00, WMA & true & 0, 7 \\
\texttt{LAG\_eule\allowbreak{}rLagrang\allowbreak{}ePsi\_G2} & WG, PG, W00, WMA & true & 7 \\
\texttt{LAG\_herm\allowbreak{}iticity\_\allowbreak{}G1} & WG, W00 & true & 6 \\
\texttt{LAG\_herm\allowbreak{}iticity\_\allowbreak{}G2} & WG, W00 & true & 6 \\
\texttt{LAG\_loca\allowbreak{}lSpinInv\allowbreak{}ariance\_\allowbreak{}G1} & WG, W00, WMA & true & 4, 6 \\
\texttt{LAG\_loca\allowbreak{}lSpinInv\allowbreak{}ariance\_\allowbreak{}G2} & WG, W00, WMA & true & 4, 6 \\
\texttt{LAG\_note\allowbreak{}bookLgGr\allowbreak{}assmannT\allowbreak{}rivial\_G\allowbreak{}1} & WG, W00 & true & 0, 5 \\
\texttt{LAG\_note\allowbreak{}bookLgGr\allowbreak{}assmannT\allowbreak{}rivial\_G\allowbreak{}2} & WG, W00 & true & 0, 5 \\
\texttt{NEG\_note\allowbreak{}bookConn\allowbreak{}ectionDe\allowbreak{}tected\_G\allowbreak{}2} & WG & true & 4 \\
\texttt{QNT\_cano\allowbreak{}nicalMom\allowbreak{}entum\_G1} & WG, W00 & true & 8 \\
\texttt{QNT\_cano\allowbreak{}nicalMom\allowbreak{}entum\_G2} & WG, W00 & true & 8 \\
\texttt{QNT\_curr\allowbreak{}entHermi\allowbreak{}ticity} & WA, PA, WMA & true & 7, 8 \\
\texttt{QNT\_curv\allowbreak{}edAntico\allowbreak{}mmutator\allowbreak{}Matrix} & WA & true & 8 \\
\texttt{QNT\_flat\allowbreak{}ModeHami\allowbreak{}ltonian} & WA, PA, W00 & true & 0, 8, 18 \\
\texttt{QNT\_krei\allowbreak{}nSignatu\allowbreak{}re} & WA, PA, W00 & true & 0, 7, 8, 18 \\
\texttt{QNT\_unit\allowbreak{}aryAndKr\allowbreak{}einSubgr\allowbreak{}oups} & WA, PA & true & 2, 8 \\
\bottomrule
\end{longtable}
\endgroup

\subsection{Index of checks: Stage 2, the primordial field}

Every name begins with \texttt{\detokenize{P_}}. Where the column "Files" names W00, the Stage-5 report records the check in its citation records (Section 20.9, item 3). The table has 70 rows; the keys in the column "Files" are those of Section 20.9.

\begingroup
\small
\begin{longtable}{@{}>{\raggedright\arraybackslash}p{0.200\linewidth}>{\raggedright\arraybackslash}p{0.200\linewidth}>{\raggedright\arraybackslash}p{0.200\linewidth}>{\raggedright\arraybackslash}p{0.200\linewidth}@{}}
\toprule
\textbf{Check} & \textbf{Files} & \textbf{Value} & \textbf{Cited in chapters} \\
\midrule
\endfirsthead
\toprule
\textbf{Check} & \textbf{Files} & \textbf{Value} & \textbf{Cited in chapters} \\
\midrule
\endhead
\texttt{P\_a4line\allowbreak{}ar} & PP & true & 4 \\
\texttt{P\_a4line\allowbreak{}ar\_cell1\allowbreak{}50Altern\allowbreak{}ativesRh\allowbreak{}oNegativ\allowbreak{}e} & WP & true & 9 \\
\texttt{P\_a4line\allowbreak{}ar\_value\allowbreak{}s} & WP & true & 9 \\
\texttt{P\_blocks\allowbreak{}\_fourBlo\allowbreak{}cksOfFou\allowbreak{}r} & WP & true & 9 \\
\texttt{P\_blocks\allowbreak{}\_matchNo\allowbreak{}tebookSe\allowbreak{}ts} & WP & true & 9 \\
\texttt{P\_christ\allowbreak{}offel} & PP & true & 4 \\
\texttt{P\_christ\allowbreak{}offel\_cl\allowbreak{}osedForm\allowbreak{}s512} & WP & true & 4, 9 \\
\texttt{P\_christ\allowbreak{}offel\_co\allowbreak{}unt} & WP & true & 4, 9 \\
\texttt{P\_einste\allowbreak{}in} & PP & true & 0, 4 \\
\texttt{P\_einste\allowbreak{}in\_energ\allowbreak{}yConditi\allowbreak{}onForms} & WP & true & 9 \\
\texttt{P\_einste\allowbreak{}in\_Gmixe\allowbreak{}dClosedF\allowbreak{}orms} & WP & true & 4, 9 \\
\texttt{P\_einste\allowbreak{}in\_noteb\allowbreak{}ookCell5\allowbreak{}83} & WP & true & 4, 9 \\
\texttt{P\_einste\allowbreak{}in\_noteb\allowbreak{}ookCell5\allowbreak{}84} & WP & true & 4, 9 \\
\texttt{P\_einste\allowbreak{}in\_offDi\allowbreak{}agonalZe\allowbreak{}ro} & WP & true & 4, 9 \\
\texttt{P\_einste\allowbreak{}in\_R44} & WP & true & 4, 9 \\
\texttt{P\_einste\allowbreak{}in\_requi\allowbreak{}redSourc\allowbreak{}e} & WP & true & 0, 9, 18 \\
\texttt{P\_einste\allowbreak{}in\_rhoRe\allowbreak{}quiredNe\allowbreak{}gative} & WP & true & 0, 4, 9, 18 \\
\texttt{P\_einste\allowbreak{}in\_ricci\allowbreak{}Scalar} & WP & true & 4, 9 \\
\texttt{P\_gammaC\allowbreak{}onst\_div\allowbreak{}ergenceI\allowbreak{}dentity} & WP & true & 4, 9 \\
\texttt{P\_gammaC\allowbreak{}onst\_div\allowbreak{}ergenceI\allowbreak{}dentityF\allowbreak{}ailsNote\allowbreak{}bookCont\allowbreak{}raction} & WP & true & 4 \\
\texttt{P\_gammaC\allowbreak{}onst\_Dmu\allowbreak{}GammaNuZ\allowbreak{}ero64} & WP & true & 4, 9 \\
\texttt{P\_gammaC\allowbreak{}onst\_not\allowbreak{}ebookCon\allowbreak{}traction\allowbreak{}ClosedFo\allowbreak{}rms} & WP & true & 4, 9 \\
\texttt{P\_gammaC\allowbreak{}onst\_not\allowbreak{}ebookCon\allowbreak{}traction\allowbreak{}Fails} & WP & true & 4, 9 \\
\texttt{\detokenize{P_metric}} & PP & true & 0 \\
\texttt{P\_metric\allowbreak{}\_detG\_eq\allowbreak{}uals\_plu\allowbreak{}s\_cos2z} & WP & true & 0, 1, 4, 9 \\
\texttt{P\_metric\allowbreak{}\_noteboo\allowbreak{}kCell106\allowbreak{}0Det} & WP & true & 9 \\
\texttt{P\_metric\allowbreak{}\_signatu\allowbreak{}re44} & WP & true & 4, 9 \\
\texttt{P\_metric\allowbreak{}\_sqrtAbs\allowbreak{}DetG\_cos\allowbreak{}z} & WP & true & 4, 9 \\
\texttt{P\_metric\allowbreak{}\_vielbei\allowbreak{}nProduct} & WP & true & 4, 9 \\
\texttt{P\_modes\_\allowbreak{}exactRed\allowbreak{}uction} & WP & true & 11 \\
\texttt{P\_notebo\allowbreak{}okCompar\allowbreak{}e\_*} & WP & all 16 true & 9 \\
\texttt{P\_notebo\allowbreak{}okCompar\allowbreak{}e\_cell11\allowbreak{}37Reprod\allowbreak{}uced16of\allowbreak{}16} & WP & true & 0, 9 \\
\texttt{P\_notebo\allowbreak{}okCompar\allowbreak{}e\_commut\allowbreak{}ingQ1Dro\allowbreak{}psOutCli\allowbreak{}ffordGam\allowbreak{}mas} & WP & true & 9 \\
\texttt{P\_notebo\allowbreak{}okCompar\allowbreak{}e\_correc\allowbreak{}tVsStore\allowbreak{}dDifferO\allowbreak{}nlyByQ} & WP & true & 0, 9 \\
\texttt{P\_notebo\allowbreak{}okCompar\allowbreak{}e\_eLaztC\allowbreak{}ell1096} & WP & true & 9 \\
\texttt{P\_notebo\allowbreak{}okCompar\allowbreak{}e\_litera\allowbreak{}lRebuild\allowbreak{}Residual\allowbreak{}IsExactl\allowbreak{}yQTerms} & WP & true & 9 \\
\texttt{P\_notebo\allowbreak{}okCompar\allowbreak{}e\_qOnlyI\allowbreak{}nYZ0to7} & WP & true & 9 \\
\texttt{P\_notebo\allowbreak{}okCompar\allowbreak{}e\_qTermF\allowbreak{}romNonCl\allowbreak{}iffordEx\allowbreak{}traTimeG\allowbreak{}ammas} & WP & true & 9 \\
\texttt{P\_notebo\allowbreak{}okCompar\allowbreak{}e\_qVecto\allowbreak{}rIs2HqAt\allowbreak{}Cell1079} & WP & true & 9 \\
\texttt{P\_notebo\allowbreak{}okCompar\allowbreak{}e\_recons\allowbreak{}tructedG\allowbreak{}amma5Not\allowbreak{}Clifford} & WP & true & 9 \\
\texttt{P\_notebo\allowbreak{}okCompar\allowbreak{}e\_recons\allowbreak{}tructedG\allowbreak{}ammaForm} & WP & true & 9 \\
\texttt{P\_notebo\allowbreak{}okCompar\allowbreak{}e\_recons\allowbreak{}truction\allowbreak{}Reproduc\allowbreak{}esStored\allowbreak{}Ela} & WP & true & 0, 9 \\
\texttt{P\_notebo\allowbreak{}okCompar\allowbreak{}e\_relabe\allowbreak{}lCell111\allowbreak{}1} & WP & true & 9 \\
\texttt{P\_notebo\allowbreak{}okCompar\allowbreak{}e\_stored\allowbreak{}Ela16} & WP & true & 9 \\
\texttt{\detokenize{P_Omega}} & PP & true & 0 \\
\texttt{P\_Omega\_\allowbreak{}closedFo\allowbreak{}rms} & WP & true & 4, 9 \\
\texttt{P\_Omega\_\allowbreak{}contract\allowbreak{}Diagonal\allowbreak{}Formula} & WP & true & 4 \\
\texttt{P\_Omega\_\allowbreak{}gammaSla\allowbreak{}sh3Hgamm\allowbreak{}a0} & WP & true & 0, 4, 9 \\
\texttt{P\_Omega\_\allowbreak{}OmegaGam\allowbreak{}maAndAnt\allowbreak{}icommuta\allowbreak{}tor} & WP & true & 9 \\
\texttt{P\_Omega\_\allowbreak{}slashA4I\allowbreak{}ndepende\allowbreak{}nt} & WP & true & 9 \\
\texttt{\detokenize{P_source}} & PP & true & 0 \\
\texttt{P\_source\allowbreak{}\_einstei\allowbreak{}nTransve\allowbreak{}rseDiffe\allowbreak{}renceIs2\allowbreak{}H2a4pp} & WP & true & 0, 9, 18 \\
\texttt{P\_source\allowbreak{}\_realKDi\allowbreak{}agonalIs\allowbreak{}C1OverSP\allowbreak{}lusC2Ove\allowbreak{}rS2} & WP & true & 9 \\
\texttt{P\_source\allowbreak{}\_realKPl\allowbreak{}aneWaveC\allowbreak{}annotSou\allowbreak{}rce} & WP & true & 9 \\
\texttt{P\_source\allowbreak{}\_transve\allowbreak{}rsePress\allowbreak{}uresEqua\allowbreak{}lForEver\allowbreak{}yX0X4Sta\allowbreak{}te} & WP & true & 9 \\
\texttt{P\_source\allowbreak{}\_x0Indep\allowbreak{}endentDi\allowbreak{}agonalOn\allowbreak{}Shell} & WP, W00 & true & 9 \\
\texttt{P\_source\allowbreak{}\_x0Indep\allowbreak{}endentEx\allowbreak{}actExamp\allowbreak{}les} & WP, W00 & true & 0, 9, 18 \\
\texttt{P\_source\allowbreak{}\_x0Indep\allowbreak{}endentOf\allowbreak{}fDiagona\allowbreak{}lAre15Bi\allowbreak{}linears} & WP & true & 9 \\
\texttt{P\_source\allowbreak{}\_x0Indep\allowbreak{}endentSo\allowbreak{}urceCond\allowbreak{}itions} & WP & true & 0, 9, 18 \\
\texttt{P\_source\allowbreak{}\_x0Indep\allowbreak{}endentSt\allowbreak{}ateSolve\allowbreak{}sDiracEx\allowbreak{}actly} & WP & true & 9 \\
\texttt{P\_spinco\allowbreak{}nn} & PP & true & 4 \\
\texttt{P\_spinco\allowbreak{}nn\_antis\allowbreak{}ymmetry} & WP & true & 4, 9 \\
\texttt{P\_spinco\allowbreak{}nn\_close\allowbreak{}dForms} & WP & true & 4, 9 \\
\texttt{P\_spinco\allowbreak{}nn\_count\allowbreak{}24} & WP & true & 4, 9 \\
\texttt{P\_spinco\allowbreak{}nn\_noteb\allowbreak{}ookCell5\allowbreak{}01OmegaM\allowbreak{}uIJEqual\allowbreak{}sMixedOm\allowbreak{}ega} & WP & true & 4, 9 \\
\texttt{P\_spinco\allowbreak{}nn\_vielb\allowbreak{}einPostu\allowbreak{}late512} & WP & true & 4, 9 \\
\texttt{P\_zeta\_g\allowbreak{}ZetaZeta\allowbreak{}IsOne} & WP & true & 9 \\
\texttt{P\_zeta\_r\allowbreak{}ange} & WP & true & 9 \\
\texttt{P\_zeta\_w\allowbreak{}arpedMet\allowbreak{}ric} & WP & true & 4, 9 \\
\texttt{P\_zeta\_w\allowbreak{}arpFacto\allowbreak{}r} & WP & true & 9 \\
\bottomrule
\end{longtable}
\endgroup

\subsection{Index of checks: Stage 3, the dark-sector experiments}

Names with underscores between lower-case words (such as \texttt{seven\_vo\allowbreak{}lume\_con\allowbreak{}stant}) are self-checks of the Rust program, stored in the \texttt{summary.\allowbreak{}json} of each experiment; names in camel case (such as \texttt{eigenmod\allowbreak{}eLaws}) are checks of the independent Python checkers. The table has 52 rows; the keys in the column "Files" are those of Section 20.9.

\begingroup
\small
\begin{longtable}{@{}>{\raggedright\arraybackslash}p{0.200\linewidth}>{\raggedright\arraybackslash}p{0.200\linewidth}>{\raggedright\arraybackslash}p{0.200\linewidth}>{\raggedright\arraybackslash}p{0.200\linewidth}@{}}
\toprule
\textbf{Check} & \textbf{Files} & \textbf{Value} & \textbf{Cited in chapters} \\
\midrule
\endfirsthead
\toprule
\textbf{Check} & \textbf{Files} & \textbf{Value} & \textbf{Cited in chapters} \\
\midrule
\endhead
\texttt{a4\_profi\allowbreak{}le\_indep\allowbreak{}endence} & E1S & true & 11 \\
\texttt{backward\allowbreak{}\_kasner\_\allowbreak{}exponent\allowbreak{}s} & E2S & true & 11 \\
\texttt{bounceAn\allowbreak{}dStop} & E3C & true & 11 \\
\texttt{boundNon\allowbreak{}negative} & E2C & true & 11 \\
\texttt{closedFo\allowbreak{}rmVerifi\allowbreak{}ed} & E2C & true & 11 \\
\texttt{constrai\allowbreak{}ntPreser\allowbreak{}ved} & E2C & true & 11 \\
\texttt{constrai\allowbreak{}nt\_prese\allowbreak{}rved} & E2S & true & 11 \\
\texttt{decelera\allowbreak{}tionRoot\allowbreak{}s} & E3C & true & 11 \\
\texttt{diracEqu\allowbreak{}ationFd} & E2C & true & 11 \\
\texttt{eigenmod\allowbreak{}eLaws} & E1C & true & 11 \\
\texttt{einstein\allowbreak{}SourceNe\allowbreak{}gative} & E1C & true & 11 \\
\texttt{einstein\allowbreak{}\_source\_\allowbreak{}negative\allowbreak{}\_energy} & E1S & true & 11 \\
\texttt{einstein\allowbreak{}TensorFr\allowbreak{}omMetric} & E2C & true & 11 \\
\texttt{energy\_s\allowbreak{}quared\_c\allowbreak{}hanges\_s\allowbreak{}ign\_at\_t\allowbreak{}star} & E5S & true & 11 \\
\texttt{exactSol\allowbreak{}ution} & E1C, E2C & true & 11 \\
\texttt{exact\_so\allowbreak{}lution} & E2S & true & 11 \\
\texttt{exact\_so\allowbreak{}lution\_a\allowbreak{}ll\_runs} & E1S & true & 11 \\
\texttt{extra\_ti\allowbreak{}mes\_turn\allowbreak{}\_to\_expa\allowbreak{}nsion} & E2S & true & 11 \\
\texttt{gammaVar\allowbreak{}iantRepr\allowbreak{}oducesUn\allowbreak{}iteTange\allowbreak{}nt} & E3F & true & 11 \\
\texttt{growth\_m\allowbreak{}atches\_w\allowbreak{}kb\_first\allowbreak{}\_order} & E5S & true & 8 \\
\texttt{growth\_m\allowbreak{}atches\_w\allowbreak{}kb\_leadi\allowbreak{}ng\_order} & E5S & true & 8, 11 \\
\texttt{hilbert\_\allowbreak{}norm\_sup\allowbreak{}erexpone\allowbreak{}ntial\_gr\allowbreak{}owth} & E5S & true & 11 \\
\texttt{hubble\_p\allowbreak{}ositive\_\allowbreak{}backward\allowbreak{}\_stop\_at\allowbreak{}\_predict\allowbreak{}ed\_bounc\allowbreak{}e} & E3S & true & 11 \\
\texttt{kasnerEx\allowbreak{}ponents} & E2C & true & 11 \\
\texttt{kreinCon\allowbreak{}servedNo\allowbreak{}rmalized} & E5C & true & 11 \\
\texttt{mixedSta\allowbreak{}teOscill\allowbreak{}ationMat\allowbreak{}chesExac\allowbreak{}t} & E1C & true & 11 \\
\texttt{muIndepe\allowbreak{}ndence} & E3C & true & 11 \\
\texttt{offDiago\allowbreak{}nalStres\allowbreak{}sVanishe\allowbreak{}s} & E2C, E3C & true & 11 \\
\texttt{pairPaul\allowbreak{}i} & E4C & true & 11 \\
\texttt{pairSmoo\allowbreak{}thTransi\allowbreak{}tionRefe\allowbreak{}rence} & E4C & true & 11 \\
\texttt{pairSudd\allowbreak{}enSpectr\allowbreak{}umMagnus\allowbreak{}Referenc\allowbreak{}e} & E4C & true & 11 \\
\texttt{pairTail\allowbreak{}MatchesK\allowbreak{}inkTheor\allowbreak{}y} & E4C & true & 11 \\
\texttt{pair\_tai\allowbreak{}l\_matche\allowbreak{}s\_kink\_t\allowbreak{}heory} & E4S & true & 11 \\
\texttt{phantomC\allowbreak{}rossing} & E3C & true & 11 \\
\texttt{phantom\_\allowbreak{}crossing\allowbreak{}\_at\_cube\allowbreak{}\_root\_2a\allowbreak{}bs\_x0} & E3S & true & 11 \\
\texttt{phantom\_\allowbreak{}iff\_nega\allowbreak{}tive\_kin\allowbreak{}etic\_ene\allowbreak{}rgy} & E2S & true & 11 \\
\texttt{phantomS\allowbreak{}tructure} & E2C & true & 11 \\
\texttt{phaseDri\allowbreak{}ftIsTime\allowbreak{}Rounding} & E2C & true & 11 \\
\texttt{refinedC\allowbreak{}onvergen\allowbreak{}ce} & E1C, E2C, E3C, E4C, E5C & true & 11, 19 \\
\texttt{repeatBy\allowbreak{}teIdenti\allowbreak{}ty} & E1C, E2C, E3C, E4C, E5C & true & 10, 11, 19 \\
\texttt{rho\_froz\allowbreak{}en\_all\_r\allowbreak{}uns} & E1S & true & 11 \\
\texttt{rho\_p\_w\_\allowbreak{}match\_cl\allowbreak{}osed\_for\allowbreak{}m} & E3S & true & 11 \\
\texttt{rho\_p\_w\_\allowbreak{}mu\_indep\allowbreak{}endent} & E3S & true & 11 \\
\texttt{rhoZeroL\allowbreak{}ocation} & E3C & true & 11 \\
\texttt{seven\_vo\allowbreak{}lume\_con\allowbreak{}stant} & E1S & true & 11 \\
\texttt{sigma\_fr\allowbreak{}om\_spino\allowbreak{}r\_equals\allowbreak{}\_a\_minus\allowbreak{}\_3} & E3S & true & 11 \\
\texttt{spinorPh\allowbreak{}aseWithi\allowbreak{}nTimeRou\allowbreak{}nding} & E2C & true & 11 \\
\texttt{S\_times\_\allowbreak{}V\_consta\allowbreak{}nt} & E2S & true & 11 \\
\texttt{tangentC\allowbreak{}PL} & E3C & true & 11 \\
\texttt{wkbFirst\allowbreak{}Order} & E5C & true & 8, 11 \\
\texttt{wkbLateR\allowbreak{}ate} & E5C & true & 8, 11 \\
\texttt{wkbLeadi\allowbreak{}ng} & E5C & true & 8, 11 \\
\bottomrule
\end{longtable}
\endgroup

\subsection{Index of checks: Stage 4, the Kohn--Sham states}

Names beginning with \texttt{\detokenize{KS_}} are exact checks of the Wolfram verifier and of the sympy checker; names in lower case with underscores are self-checks of the Rust solver or checks of the cross-checker. A Rust self-check that is made for every run carries the name of the run in front, such as \texttt{m1\_L3\_N1\allowbreak{}12\_lam0\_\allowbreak{}T0\_energ\allowbreak{}y\_from\_r\allowbreak{}ho}, and is indexed as a family, here \texttt{*\_energy\allowbreak{}\_from\_rh\allowbreak{}o}. Where the column "Files" names W00, the Stage-5 report records the check in its citation records (Section 20.9, item 3). The table has 102 rows; the keys in the column "Files" are those of Section 20.9.

\begingroup
\small
\begin{longtable}{@{}>{\raggedright\arraybackslash}p{0.200\linewidth}>{\raggedright\arraybackslash}p{0.200\linewidth}>{\raggedright\arraybackslash}p{0.200\linewidth}>{\raggedright\arraybackslash}p{0.200\linewidth}@{}}
\toprule
\textbf{Check} & \textbf{Files} & \textbf{Value} & \textbf{Cited in chapters} \\
\midrule
\endfirsthead
\toprule
\textbf{Check} & \textbf{Files} & \textbf{Value} & \textbf{Cited in chapters} \\
\midrule
\endhead
\texttt{a4\_resca\allowbreak{}ling\_pai\allowbreak{}r\_exact} & RSC & true & 13 \\
\texttt{canonica\allowbreak{}l\_deltaS\allowbreak{}CF} & PC & true & 14, 19 \\
\texttt{canonica\allowbreak{}l\_eigenv\allowbreak{}alues} & PC & false & 14, 16, 18, 19 \\
\texttt{*\_cv\_fin\allowbreak{}ite\_diff\allowbreak{}erence\_v\allowbreak{}s\_exact} & RTH & all 4 true & 13 \\
\texttt{*\_emt\_co\allowbreak{}nservati\allowbreak{}on} & RSC, REM & all 43 true & 13 \\
\texttt{*\_energy\allowbreak{}\_from\_rh\allowbreak{}o} & RSC, REM & all 43 true & 13 \\
\texttt{*\_entrop\allowbreak{}y\_positi\allowbreak{}ve} & RTH & all 16 true & 13 \\
\texttt{exchange\allowbreak{}\_table\_d\allowbreak{}3\_closed\allowbreak{}\_form} & RSP, RXT & true & 13 \\
\texttt{exchange\allowbreak{}\_table\_d\allowbreak{}4\_closed\allowbreak{}\_form} & RSP, RXT & true & 13 \\
\texttt{excited\_\allowbreak{}grid\_ref\allowbreak{}inement\_\allowbreak{}delta\_sc\allowbreak{}f} & REX & true & 13 \\
\texttt{*\_free\_e\allowbreak{}nergy\_de\allowbreak{}creases} & RTH & all 16 true & 13 \\
\texttt{gauntlet\allowbreak{}\_*} & NB & 56 of 60 true & 19 \\
\texttt{grid\_ref\allowbreak{}inement\_\allowbreak{}energy} & RSC & true & 13 \\
\texttt{*\_heat\_c\allowbreak{}apacity\_\allowbreak{}positive} & RTH & all 16 true & 13 \\
\texttt{KS\_bound\allowbreak{}ary\_bagF\allowbreak{}amily} & WK, PT & true & 13 \\
\texttt{KS\_bound\allowbreak{}ary\_bagT\allowbreak{}hetaZero\allowbreak{}IsEvenPa\allowbreak{}rity} & WK, PT & true & 13 \\
\texttt{KS\_bound\allowbreak{}ary\_curr\allowbreak{}entBlock\allowbreak{}Form} & WK, PT & true & 13 \\
\texttt{KS\_bound\allowbreak{}ary\_curr\allowbreak{}entConse\allowbreak{}rvedAlon\allowbreak{}gY} & WK, PT & true & 13 \\
\texttt{KS\_bound\allowbreak{}ary\_curr\allowbreak{}entMatri\allowbreak{}x} & WK, PT & true & 13 \\
\texttt{KS\_bound\allowbreak{}ary\_hilb\allowbreak{}ertNormN\allowbreak{}otConser\allowbreak{}vedAlong\allowbreak{}Y} & WK, PT & true & 13 \\
\texttt{KS\_bound\allowbreak{}ary\_pari\allowbreak{}tyA\_symm\allowbreak{}etryIffM\allowbreak{}assOdd} & WK, PT, W00 & true & 13 \\
\texttt{KS\_bound\allowbreak{}ary\_pari\allowbreak{}tyB\_coup\allowbreak{}lesJBloc\allowbreak{}ks} & WK, PT & true & 13 \\
\texttt{KS\_bound\allowbreak{}ary\_pari\allowbreak{}tyB\_symm\allowbreak{}etryForE\allowbreak{}venMass} & WK, PT & true & 13 \\
\texttt{KS\_bound\allowbreak{}ary\_pari\allowbreak{}tyKillsC\allowbreak{}urrent} & WK, PT & true & 13 \\
\texttt{KS\_bound\allowbreak{}ary\_pari\allowbreak{}tyProjec\allowbreak{}torsBloc\allowbreak{}kForm} & WK, PT & true & 13 \\
\texttt{KS\_bound\allowbreak{}ary\_self\allowbreak{}AdjointB\allowbreak{}oundaryT\allowbreak{}erm} & WK, PT & true & 13 \\
\texttt{KS\_emt\_o\allowbreak{}ffBlockC\allowbreak{}omponent\allowbreak{}sVanishP\allowbreak{}erOrbita\allowbreak{}l} & WK, PT & true & 13 \\
\texttt{KS\_emt\_p\allowbreak{}1} & WK, PT & true & 13 \\
\texttt{KS\_emt\_p\allowbreak{}2p3} & WK, PT & true & 13 \\
\texttt{KS\_emt\_p\allowbreak{}t} & WK, PT & true & 13 \\
\texttt{KS\_emt\_p\allowbreak{}y} & WK, PT & true & 13 \\
\texttt{KS\_emt\_r\allowbreak{}ho} & WK, PT, W00 & true & 13 \\
\texttt{KS\_emt\_T\allowbreak{}41cancel\allowbreak{}sOverShe\allowbreak{}ll} & WK, PT & true & 13 \\
\texttt{KS\_emt\_T\allowbreak{}y1Propor\allowbreak{}tionalTo\allowbreak{}Current} & WK, PT & true & 13 \\
\texttt{KS\_emt\_T\allowbreak{}y4Propor\allowbreak{}tionalTo\allowbreak{}Current} & WK, PT & true & 13 \\
\texttt{KS\_excha\allowbreak{}nge\_angu\allowbreak{}larAvera\allowbreak{}geOfPdot\allowbreak{}QVanishe\allowbreak{}s} & WK, PT & true & 0, 13 \\
\texttt{KS\_excha\allowbreak{}nge\_bloc\allowbreak{}kFormOfF\allowbreak{}ockTerm} & WK, PT & true & 0, 13 \\
\texttt{KS\_excha\allowbreak{}nge\_coup\allowbreak{}lingDime\allowbreak{}nsion} & WK, PT & true & 13, 18 \\
\texttt{KS\_excha\allowbreak{}nge\_fill\allowbreak{}edShellO\allowbreak{}neEighth} & WK, PT & true & 13 \\
\texttt{KS\_excha\allowbreak{}nge\_fill\allowbreak{}edShellS\allowbreak{}calarDen\allowbreak{}sity} & WK, PT & true & 13 \\
\texttt{KS\_excha\allowbreak{}nge\_kern\allowbreak{}elPlusMi\allowbreak{}nus} & WK & true & 13 \\
\texttt{KS\_excha\allowbreak{}nge\_kern\allowbreak{}elPlusPl\allowbreak{}us} & WK & true & 13 \\
\texttt{KS\_excha\allowbreak{}nge\_ldaP\allowbreak{}otential\allowbreak{}s} & PT & true & 13 \\
\texttt{KS\_excha\allowbreak{}nge\_rest\allowbreak{}GasLimit} & WK, PT & true & 13 \\
\texttt{KS\_excha\allowbreak{}nge\_unif\allowbreak{}ormGasCl\allowbreak{}osedForm} & WK, PT & true & 0, 13, 18 \\
\texttt{KS\_excha\allowbreak{}nge\_wick\allowbreak{}TheoremH\allowbreak{}F} & WK, PT & true & 13, 18 \\
\texttt{KS\_funct\allowbreak{}ional} & PT & true & 13 \\
\texttt{KS\_funct\allowbreak{}ional\_fe\allowbreak{}rmiDirac\allowbreak{}Occupati\allowbreak{}ons} & WK, PT & true & 13 \\
\texttt{KS\_funct\allowbreak{}ional\_he\allowbreak{}llmannFe\allowbreak{}ynmanLam\allowbreak{}bda} & WK, PT & true & 13 \\
\texttt{KS\_funct\allowbreak{}ional\_he\allowbreak{}llmannFe\allowbreak{}ynmanMas\allowbreak{}s} & WK, PT & true & 13 \\
\texttt{KS\_funct\allowbreak{}ional\_he\allowbreak{}llmannFe\allowbreak{}ynmanMom\allowbreak{}entum} & PT & true & 13 \\
\texttt{KS\_funct\allowbreak{}ional\_he\allowbreak{}llmannFe\allowbreak{}ynmanTem\allowbreak{}perature} & WK, PT & true & 13 \\
\texttt{KS\_funct\allowbreak{}ional\_me\allowbreak{}rminStru\allowbreak{}cture} & WK, PT & true & 13 \\
\texttt{KS\_funct\allowbreak{}ional\_st\allowbreak{}ationari\allowbreak{}tyGivesK\allowbreak{}SEquatio\allowbreak{}n} & WK, PT & true & 13 \\
\texttt{KS\_funct\allowbreak{}ional\_to\allowbreak{}talEnerg\allowbreak{}yDoubleC\allowbreak{}ounting} & WK, PT & true & 13 \\
\texttt{KS\_geome\allowbreak{}try\_bran\allowbreak{}eEnergyP\allowbreak{}ositive} & WK, PT & true & 9, 13 \\
\texttt{KS\_geome\allowbreak{}try\_chri\allowbreak{}stoffelC\allowbreak{}losedFor\allowbreak{}ms} & WK, PT & true & 13 \\
\texttt{KS\_geome\allowbreak{}try\_chri\allowbreak{}stoffelC\allowbreak{}ount18} & WK, PT & true & 9 \\
\texttt{KS\_geome\allowbreak{}try\_cons\allowbreak{}tantCurv\allowbreak{}atureSev\allowbreak{}enSpace} & WK, PT & true & 9, 13, 18 \\
\texttt{KS\_geome\allowbreak{}try\_eins\allowbreak{}teinMixe\allowbreak{}dDiag} & WK, PT & true & 4, 9, 13, 18 \\
\texttt{KS\_geome\allowbreak{}try\_extr\allowbreak{}insicCur\allowbreak{}vature} & WK, PT & true & 9, 13 \\
\texttt{KS\_geome\allowbreak{}try\_indu\allowbreak{}cedMetri\allowbreak{}cFlat} & WK & true & 9 \\
\texttt{KS\_geome\allowbreak{}try\_isra\allowbreak{}elConven\allowbreak{}tionRand\allowbreak{}allSundr\allowbreak{}um} & PT & true & 9, 13 \\
\texttt{KS\_geome\allowbreak{}try\_isra\allowbreak{}elJump} & WK, PT & true & 9, 13 \\
\texttt{KS\_geome\allowbreak{}try\_isra\allowbreak{}elStress} & WK, PT & true & 0, 9, 13 \\
\texttt{KS\_geome\allowbreak{}try\_kret\allowbreak{}schmannC\allowbreak{}onstant} & WK & true & 9, 18 \\
\texttt{KS\_geome\allowbreak{}try\_note\allowbreak{}bookChar\allowbreak{}t} & WK & true & 9 \\
\texttt{KS\_geome\allowbreak{}try\_requ\allowbreak{}iredSour\allowbreak{}ce} & WK, PT & true & 0, 4, 9, 13 \\
\texttt{KS\_geome\allowbreak{}try\_rhoR\allowbreak{}equiredN\allowbreak{}egative} & WK, PT & true & 9, 13 \\
\texttt{KS\_geome\allowbreak{}try\_ricc\allowbreak{}iMixed} & WK, PT & true & 9 \\
\texttt{KS\_geome\allowbreak{}try\_ricc\allowbreak{}iScalarM\allowbreak{}inus42H2} & WK, PT & true & 0, 4, 9, 13 \\
\texttt{KS\_geome\allowbreak{}try\_sign\allowbreak{}ature44} & WK, PT & true & 9 \\
\texttt{KS\_geome\allowbreak{}try\_sqrt\allowbreak{}DetG\_W6} & WK, PT & true & 9, 13 \\
\texttt{KS\_reduc\allowbreak{}tion\_a4I\allowbreak{}sMomentu\allowbreak{}mRescali\allowbreak{}ng} & WK, PT & true & 9, 13 \\
\texttt{KS\_reduc\allowbreak{}tion\_alg\allowbreak{}ebraDim8} & WK, PT & true & 13 \\
\texttt{KS\_reduc\allowbreak{}tion\_ans\allowbreak{}atzRemov\allowbreak{}es3H} & WK, PT, W00 & true & 13 \\
\texttt{KS\_reduc\allowbreak{}tion\_bas\allowbreak{}isUnitar\allowbreak{}y} & WK, PT & true & 13 \\
\texttt{KS\_reduc\allowbreak{}tion\_blo\allowbreak{}ckHamilt\allowbreak{}onianEqu\allowbreak{}ivalentT\allowbreak{}oODE} & WK, PT & true & 13 \\
\texttt{KS\_reduc\allowbreak{}tion\_blo\allowbreak{}ckODEMat\allowbreak{}rix} & WK, PT & true & 0, 13 \\
\texttt{KS\_reduc\allowbreak{}tion\_blo\allowbreak{}cksA0A1A\allowbreak{}4} & WK, PT & true & 13 \\
\texttt{KS\_reduc\allowbreak{}tion\_blo\allowbreak{}cksBC} & WK, PT, W00 & true & 13 \\
\texttt{KS\_reduc\allowbreak{}tion\_blo\allowbreak{}ckTypes} & WK, PT & true & 13 \\
\texttt{KS\_reduc\allowbreak{}tion\_fiv\allowbreak{}eMatrice\allowbreak{}sBlockDi\allowbreak{}agonal} & WK, PT & true & 13 \\
\texttt{KS\_reduc\allowbreak{}tion\_fla\allowbreak{}tMeasure} & WK, PT & true & 13 \\
\texttt{KS\_reduc\allowbreak{}tion\_gam\allowbreak{}ma8Swaps\allowbreak{}J} & WK, PT & true & 13 \\
\texttt{KS\_reduc\allowbreak{}tion\_gam\allowbreak{}maSlashO\allowbreak{}mega3Hga\allowbreak{}mma0} & WK & true & 13 \\
\texttt{KS\_reduc\allowbreak{}tion\_hMi\allowbreak{}nusEqual\allowbreak{}sMinusHP\allowbreak{}lus} & WK, PT & true & 0, 13 \\
\texttt{KS\_reduc\allowbreak{}tion\_JK1\allowbreak{}K2commut\allowbreak{}e} & WK, PT & true & 13 \\
\texttt{KS\_reduc\allowbreak{}tion\_k0M\allowbreak{}assiveLe\allowbreak{}vels} & WK, PT & true & 13 \\
\texttt{KS\_reduc\allowbreak{}tion\_k0Z\allowbreak{}eroMode} & WK, PT & true & 13 \\
\texttt{KS\_reduc\allowbreak{}tion\_mir\allowbreak{}rorPatch\allowbreak{}SameForm} & WK & true & 13 \\
\texttt{KS\_reduc\allowbreak{}tion\_pro\allowbreak{}jectorsR\allowbreak{}ank2} & WK, PT & true & 13 \\
\texttt{KS\_reduc\allowbreak{}tion\_rec\allowbreak{}onstruct\allowbreak{}FromBloc\allowbreak{}ks} & WK, PT & true & 13 \\
\texttt{KS\_reduc\allowbreak{}tion\_red\allowbreak{}ucedEqua\allowbreak{}tionODEF\allowbreak{}orm} & WK, PT & true & 13 \\
\texttt{KS\_reduc\allowbreak{}tion\_rot\allowbreak{}ationalS\allowbreak{}ymmetry} & WK, PT & true & 13 \\
\texttt{KS\_reduc\allowbreak{}tion\_sig\allowbreak{}ma3Conju\allowbreak{}gationFl\allowbreak{}ipsK} & WK, PT & true & 13 \\
\texttt{KS\_reduc\allowbreak{}tion\_wit\allowbreak{}houtW3th\allowbreak{}e3HTermS\allowbreak{}urvives} & WK, PT & true & 13 \\
\texttt{KS\_reduc\allowbreak{}tion\_zer\allowbreak{}oModeSpl\allowbreak{}itting} & WK, PT & true & 13 \\
\texttt{m1\_L3\_N1\allowbreak{}12\_lam0\_\allowbreak{}T0\_hellm\allowbreak{}ann\_feyn\allowbreak{}man\_dE\_d\allowbreak{}m} & RSC & true & 13 \\
\texttt{m1\_L3\_N1\allowbreak{}12\_lamp1\allowbreak{}\_T0\_hell\allowbreak{}mann\_fey\allowbreak{}nman\_dE\_\allowbreak{}dm} & RSC & true & 13 \\
\texttt{\detokenize{r07_*}} & NB & 0 of 2 true & 19 \\
\texttt{theory\_z\allowbreak{}ero\_mode\allowbreak{}\_splitti\allowbreak{}ng} & RSP, RTA & true & 13 \\
\bottomrule
\end{longtable}
\endgroup

\subsection{Index of checks: Stage 5 and the matter--antimatter analysis}

The prefixes are \texttt{\detokenize{C00_}} (dirac16complex00, Wolfram), \texttt{\detokenize{S5_}} (the two sympy checkers of Stage 5), \texttt{\detokenize{PAIR_}} (pairing theorems, Wolfram) and \texttt{\detokenize{MA_}} (matter and antimatter, both programs). The table has 269 rows; the keys in the column "Files" are those of Section 20.9.

\begingroup
\small
\begin{longtable}{@{}>{\raggedright\arraybackslash}p{0.200\linewidth}>{\raggedright\arraybackslash}p{0.200\linewidth}>{\raggedright\arraybackslash}p{0.200\linewidth}>{\raggedright\arraybackslash}p{0.200\linewidth}@{}}
\toprule
\textbf{Check} & \textbf{Files} & \textbf{Value} & \textbf{Cited in chapters} \\
\midrule
\endfirsthead
\toprule
\textbf{Check} & \textbf{Files} & \textbf{Value} & \textbf{Cited in chapters} \\
\midrule
\endhead
\texttt{C00\_alge\allowbreak{}bra\_anti\allowbreak{}commutat\allowbreak{}orTotall\allowbreak{}yAntisym\allowbreak{}metric} & W00 & true & 6 \\
\texttt{C00\_alge\allowbreak{}bra\_lead\allowbreak{}Facts} & W00 & true & 6 \\
\texttt{C00\_alge\allowbreak{}bra\_pinM\allowbreak{}oduleAnd\allowbreak{}SpinInva\allowbreak{}riance} & W00 & true & 6 \\
\texttt{C00\_char\allowbreak{}ge\_indef\allowbreak{}inite} & W00 & true & 7, 14, 15 \\
\texttt{C00\_conn\allowbreak{}ection\_n\allowbreak{}onTrivia\allowbreak{}lCouplin\allowbreak{}g} & W00 & true & 6, 7 \\
\texttt{C00\_conn\allowbreak{}ection\_O\allowbreak{}megaTerm\allowbreak{}InFieldE\allowbreak{}quation} & W00 & true & 6 \\
\texttt{C00\_curr\allowbreak{}ent\_cons\allowbreak{}ervation\allowbreak{}AndReali\allowbreak{}ty} & W00 & true & 7, 14 \\
\texttt{C00\_EL\_c\allowbreak{}ommuting\allowbreak{}Curved} & W00 & true & 7, 14 \\
\texttt{C00\_EL\_c\allowbreak{}ommuting\allowbreak{}Flat} & W00 & true & 7 \\
\texttt{C00\_EL\_c\allowbreak{}ommuting\allowbreak{}GeneralS\allowbreak{}moothU} & W00 & true & 6, 7 \\
\texttt{C00\_EL\_i\allowbreak{}dentical\allowbreak{}FormBoth\allowbreak{}Statisti\allowbreak{}cs} & W00 & true & 7, 14 \\
\texttt{C00\_EMT\_\allowbreak{}conserva\allowbreak{}tionOnSh\allowbreak{}ell} & W00 & true & 7 \\
\texttt{C00\_EMT\_\allowbreak{}generalS\allowbreak{}moothU} & W00 & true & 6, 7 \\
\texttt{C00\_EMT\_\allowbreak{}homogene\allowbreak{}ousEquat\allowbreak{}ionsOfSt\allowbreak{}ate} & W00 & true & 7 \\
\texttt{C00\_EMT\_\allowbreak{}observer\allowbreak{}SplitGau\allowbreak{}ssianNor\allowbreak{}mal} & W00 & true & 7 \\
\texttt{C00\_EMT\_\allowbreak{}symmetri\allowbreak{}cAndReal} & W00 & true & 7 \\
\texttt{C00\_EMT\_\allowbreak{}traceOnS\allowbreak{}hell} & W00 & true & 7 \\
\texttt{C00\_EMT\_\allowbreak{}vielbein\allowbreak{}Variatio\allowbreak{}n} & W00 & true & 7 \\
\texttt{C00\_ener\allowbreak{}gy\_unbou\allowbreak{}ndedBelo\allowbreak{}w} & W00 & true & 7, 14, 15, 18 \\
\texttt{C00\_lagr\allowbreak{}angian\_c\allowbreak{}oefficie\allowbreak{}ntHermit\allowbreak{}icity} & W00 & true & 6 \\
\texttt{C00\_lagr\allowbreak{}angian\_g\allowbreak{}rassmann\allowbreak{}FlatHerm\allowbreak{}itianAnd\allowbreak{}EL} & W00 & true & 6 \\
\texttt{C00\_lagr\allowbreak{}angian\_r\allowbreak{}ealCommu\allowbreak{}ting} & W00 & true & 6 \\
\texttt{C00\_mass\allowbreak{}Term\_com\allowbreak{}mutingEx\allowbreak{}plicitSp\allowbreak{}inors} & W00 & true & 6 \\
\texttt{C00\_mass\allowbreak{}Term\_dis\allowbreak{}persionF\allowbreak{}lat} & W00 & true & 6 \\
\texttt{C00\_mass\allowbreak{}Term\_gra\allowbreak{}ssmannNo\allowbreak{}nzero} & W00 & true & 6 \\
\texttt{C00\_prim\allowbreak{}ordial\_E\allowbreak{}Lcommuti\allowbreak{}ng} & W00 & true & 7 \\
\texttt{C00\_prim\allowbreak{}ordial\_g\allowbreak{}eometryM\allowbreak{}atchesSt\allowbreak{}age2} & W00 & true & 7 \\
\texttt{C00\_prim\allowbreak{}ordial\_h\allowbreak{}omogeneo\allowbreak{}usState} & W00 & true & 7 \\
\texttt{C00\_prim\allowbreak{}ordial\_s\allowbreak{}taticFie\allowbreak{}ldSource\allowbreak{}dExactly} & W00 & true & 7, 18 \\
\texttt{C00\_real\allowbreak{}Restrict\allowbreak{}ion\_comm\allowbreak{}utingFla\allowbreak{}tDecompo\allowbreak{}sition} & W00 & true & 6 \\
\texttt{C00\_real\allowbreak{}Restrict\allowbreak{}ion\_comm\allowbreak{}utingNon\allowbreak{}Trivial} & W00 & true & 6 \\
\texttt{C00\_real\allowbreak{}Restrict\allowbreak{}ion\_comm\allowbreak{}utingNot\allowbreak{}ebookCon\allowbreak{}traction} & W00 & true & 6 \\
\texttt{C00\_real\allowbreak{}Restrict\allowbreak{}ion\_gras\allowbreak{}smannCom\allowbreak{}plexDeco\allowbreak{}mpositio\allowbreak{}n} & W00 & true & 6 \\
\texttt{C00\_real\allowbreak{}Restrict\allowbreak{}ion\_gras\allowbreak{}smannCur\allowbreak{}vedTrivi\allowbreak{}al} & W00 & true & 6 \\
\texttt{C00\_real\allowbreak{}Restrict\allowbreak{}ion\_gras\allowbreak{}smannFla\allowbreak{}tTrivial} & W00 & true & 6 \\
\texttt{C00\_real\allowbreak{}Restrict\allowbreak{}ion\_rela\allowbreak{}tionToL1} & W00 & true & 6 \\
\texttt{C00\_stat\allowbreak{}ic\_block\allowbreak{}sFromSta\allowbreak{}ge4Theor\allowbreak{}y} & W00 & true & 14 \\
\texttt{C00\_stat\allowbreak{}ic\_class\allowbreak{}icalEner\allowbreak{}gyKreinS\allowbreak{}igned} & W00 & true & 7, 14, 15 \\
\texttt{C00\_stat\allowbreak{}ic\_class\allowbreak{}icalMode\allowbreak{}IsKreinW\allowbreak{}eightedK\allowbreak{}S} & W00 & true & 7, 14, 15 \\
\texttt{C00\_stat\allowbreak{}ic\_geome\allowbreak{}tryAndRe\allowbreak{}ducedEqu\allowbreak{}ation} & W00 & true & 7, 14 \\
\texttt{MA\_agree\allowbreak{}sWithWol\allowbreak{}fram} & PMA & true & 17 \\
\texttt{\detokenize{MA_M1}} & PMA & true & 17 \\
\texttt{\detokenize{MA_M1_*}} & WMA, PMA & all 46 true & 16 \\
\texttt{MA\_M1\_ch\allowbreak{}argeDens\allowbreak{}ityMatri\allowbreak{}x} & WMA & true & 17 \\
\texttt{MA\_M1\_di\allowbreak{}vergence\allowbreak{}Identity\allowbreak{}AllFirst\allowbreak{}Jets} & PMA & true & 17 \\
\texttt{MA\_M1\_ks\allowbreak{}FixedNet\allowbreak{}NumberRe\allowbreak{}corded} & WMA & true & 17, 19 \\
\texttt{MA\_M1\_ne\allowbreak{}gativeCo\allowbreak{}ntrolNot\allowbreak{}ebookCon\allowbreak{}nection} & WMA & true & 7, 17 \\
\texttt{MA\_M1\_no\allowbreak{}etherCur\allowbreak{}rentLoca\allowbreak{}lPhase\_g\allowbreak{}rassmann\allowbreak{}\_G1} & WMA & true & 17 \\
\texttt{MA\_M1\_no\allowbreak{}etherIde\allowbreak{}ntity\_*} & WMA, PMA & all 6 true & 17 \\
\texttt{MA\_M1\_no\allowbreak{}etherIde\allowbreak{}ntity\_co\allowbreak{}mmuting\_\allowbreak{}G1} & WMA & true & 7, 16, 17 \\
\texttt{MA\_M1\_no\allowbreak{}etherIde\allowbreak{}ntity\_gr\allowbreak{}assmann\_\allowbreak{}G1} & WMA & true & 7, 16, 17 \\
\texttt{MA\_M1\_on\allowbreak{}ShellCon\allowbreak{}servatio\allowbreak{}n\_G1} & WMA & true & 16, 17 \\
\texttt{MA\_M1\_st\allowbreak{}age1Chec\allowbreak{}ksCited} & WMA & true & 17 \\
\texttt{MA\_M1\_u1\allowbreak{}Invarian\allowbreak{}ceCommut\allowbreak{}ingGener\allowbreak{}icU\_flat} & WMA & true & 17 \\
\texttt{MA\_M1\_u1\allowbreak{}Invarian\allowbreak{}ceCommut\allowbreak{}ingGener\allowbreak{}icU\_G1} & WMA & true & 6 \\
\texttt{MA\_M1\_u1\allowbreak{}Invarian\allowbreak{}ceGrassm\allowbreak{}ann\_G1} & WMA & true & 6, 17 \\
\texttt{\detokenize{MA_M2}} & PMA & true & 17 \\
\texttt{\detokenize{MA_M2_*}} & WMA, PMA & all 26 true & 6 \\
\texttt{MA\_M2\_C\_\allowbreak{}and\_CP\_s\allowbreak{}tatus} & PMA & true & 6, 17 \\
\texttt{MA\_M2\_ca\allowbreak{}nonicalS\allowbreak{}tructure} & WMA & true & 17 \\
\texttt{MA\_M2\_ca\allowbreak{}nonicalS\allowbreak{}tructure\allowbreak{}OfExactG\allowbreak{}rassmann\allowbreak{}Symmetri\allowbreak{}es} & PMA & true & 17 \\
\texttt{MA\_M2\_ch\allowbreak{}argeConj\allowbreak{}ugationS\allowbreak{}olutionS\allowbreak{}paces} & PMA & true & 17 \\
\texttt{MA\_M2\_co\allowbreak{}njugatio\allowbreak{}nIntertw\allowbreak{}iners} & WMA & true & 17 \\
\texttt{MA\_M2\_cp\allowbreak{}ScopeInC\allowbreak{}urvedFie\allowbreak{}lds} & PMA & true & 17, 19 \\
\texttt{MA\_M2\_di\allowbreak{}screteGr\allowbreak{}oupChara\allowbreak{}cterTabl\allowbreak{}e} & PMA & true & 17 \\
\texttt{MA\_M2\_fr\allowbreak{}ameLevel\allowbreak{}G1\_commu\allowbreak{}ting} & WMA & true & 17 \\
\texttt{MA\_M2\_fr\allowbreak{}ameLevel\allowbreak{}G1\_grass\allowbreak{}mann} & WMA & true & 17 \\
\texttt{MA\_M2\_in\allowbreak{}ternalMa\allowbreak{}psCurved\allowbreak{}Jets} & PMA & true & 6 \\
\texttt{MA\_M2\_la\allowbreak{}grangian\allowbreak{}FlatComm\allowbreak{}uting\_al\allowbreak{}l} & WMA & true & 6, 17 \\
\texttt{MA\_M2\_la\allowbreak{}grangian\allowbreak{}FlatGras\allowbreak{}smann\_al\allowbreak{}l} & WMA & true & 6, 17 \\
\texttt{MA\_M2\_na\allowbreak{}medRefle\allowbreak{}ctionsFl\allowbreak{}atJets} & PMA & true & 6 \\
\texttt{MA\_M2\_si\allowbreak{}gnPatter\allowbreak{}nClassif\allowbreak{}ication} & WMA & true & 17 \\
\texttt{MA\_M2\_sp\allowbreak{}in0Conta\allowbreak{}insCharg\allowbreak{}eReversi\allowbreak{}ngTimeRo\allowbreak{}tation} & PMA & true & 17 \\
\texttt{MA\_M2\_st\allowbreak{}atistics\allowbreak{}Sign} & WMA & true & 17 \\
\texttt{MA\_M2\_sy\allowbreak{}mmetrySu\allowbreak{}mmaryAnd\allowbreak{}ChargeRe\allowbreak{}versal} & WMA & true & 17 \\
\texttt{\detokenize{MA_M3}} & PMA & true & 17 \\
\texttt{\detokenize{MA_M3_*}} & WMA, PMA & all 17 true & 18 \\
\texttt{MA\_M3\_al\allowbreak{}lInvaria\allowbreak{}ntFormsS\allowbreak{}ymmetric} & PMA & true & 17 \\
\texttt{MA\_M3\_co\allowbreak{}mmutingS\allowbreak{}urvival} & WMA & true & 17 \\
\texttt{MA\_M3\_de\allowbreak{}rivative\allowbreak{}TypeSurv\allowbreak{}ivalCurv\allowbreak{}ed} & PMA & true & 17 \\
\texttt{MA\_M3\_ex\allowbreak{}traGrass\allowbreak{}mannQuar\allowbreak{}ticCharg\allowbreak{}e4} & WMA & true & 17 \\
\texttt{MA\_M3\_gr\allowbreak{}assmannS\allowbreak{}urvival} & WMA & true & 17 \\
\texttt{MA\_M3\_in\allowbreak{}variantF\allowbreak{}ormsSpan\allowbreak{}\_C\_Cgamm\allowbreak{}a8} & PMA & true & 17 \\
\texttt{MA\_M3\_ki\allowbreak{}neticInv\allowbreak{}ariantFo\allowbreak{}rms} & WMA & true & 17 \\
\texttt{MA\_M3\_ma\allowbreak{}ssTypeSu\allowbreak{}rvivalAn\allowbreak{}dCharge} & PMA & true & 17 \\
\texttt{MA\_M3\_no\allowbreak{}tebookLg\allowbreak{}IsMajora\allowbreak{}naType} & WMA & true & 17 \\
\texttt{MA\_M3\_pi\allowbreak{}nCharact\allowbreak{}erForms} & WMA & true & 17 \\
\texttt{MA\_M3\_pi\allowbreak{}nCharact\allowbreak{}ersMajor\allowbreak{}anaTerms} & WMA & true & 17 \\
\texttt{MA\_M3\_qu\allowbreak{}articCha\allowbreak{}rgeViola\allowbreak{}tingExam\allowbreak{}plesGras\allowbreak{}smann} & PMA & true & 17 \\
\texttt{MA\_M3\_sp\allowbreak{}inInvari\allowbreak{}antForms} & WMA & true & 17 \\
\texttt{MA\_M3\_u1\allowbreak{}Charge} & WMA & true & 17 \\
\texttt{\detokenize{MA_M4}} & PMA & true & 17 \\
\texttt{MA\_M4\_cu\allowbreak{}rrentFli\allowbreak{}pG1\_gras\allowbreak{}smann} & WMA & true & 17 \\
\texttt{MA\_M4\_cu\allowbreak{}rrentFli\allowbreak{}pMatrix} & WMA & true & 17 \\
\texttt{MA\_M4\_ga\allowbreak{}mma8Char\allowbreak{}geFlipCu\allowbreak{}rved} & PMA & true & 17 \\
\texttt{MA\_M4\_ga\allowbreak{}mma8EMTP\allowbreak{}airing} & PMA & true & 17 \\
\texttt{MA\_M4\_ga\allowbreak{}mma8Eule\allowbreak{}rLagrang\allowbreak{}ePairing} & PMA & true & 17 \\
\texttt{MA\_M4\_ga\allowbreak{}mma8Matr\allowbreak{}ixFacts} & PMA & true & 17 \\
\texttt{MA\_M4\_im\allowbreak{}ageField\allowbreak{}FockMode\allowbreak{}l} & PMA & true & 17, 19 \\
\texttt{MA\_M4\_kr\allowbreak{}einModeF\allowbreak{}acts} & PMA & true & 15, 17 \\
\texttt{MA\_M4\_kr\allowbreak{}einOnePa\allowbreak{}rticle} & WMA & true & 15, 16, 17 \\
\texttt{MA\_M4\_pa\allowbreak{}irEMTAnd\allowbreak{}CurrentG\allowbreak{}1\_commut\allowbreak{}ing} & WMA & true & 15, 16, 17 \\
\texttt{MA\_M4\_pa\allowbreak{}irEMTG1\_\allowbreak{}grassman\allowbreak{}n} & WMA & true & 15, 16, 17 \\
\texttt{MA\_M5\_im\allowbreak{}plicatio\allowbreak{}n} & WMA & true & 16, 17 \\
\texttt{PAIR\_alg\allowbreak{}ebra\_bil\allowbreak{}inearPar\allowbreak{}ities} & WPR & true & 15 \\
\texttt{PAIR\_alg\allowbreak{}ebra\_BPr\allowbreak{}operties} & WPR & true & 15 \\
\texttt{PAIR\_alg\allowbreak{}ebra\_chi\allowbreak{}ralBlock\allowbreak{}Structur\allowbreak{}e} & WPR & true & 15 \\
\texttt{PAIR\_alg\allowbreak{}ebra\_CPr\allowbreak{}operties} & WPR & true & 15 \\
\texttt{PAIR\_alg\allowbreak{}ebra\_gam\allowbreak{}ma8InIde\allowbreak{}ntityCom\allowbreak{}ponent} & WPR & true & 6, 15 \\
\texttt{PAIR\_alg\allowbreak{}ebra\_gam\allowbreak{}ma8Prope\allowbreak{}rties} & WPR & true & 6, 15 \\
\texttt{PAIR\_alg\allowbreak{}ebra\_kre\allowbreak{}inUnderB\allowbreak{}asicRefl\allowbreak{}ections} & WPR & true & 15 \\
\texttt{PAIR\_sta\allowbreak{}t\_*} & WPR & all 13 true & 14 \\
\texttt{PAIR\_sta\allowbreak{}t\_bosonT\allowbreak{}hermalWi\allowbreak{}ckPlus} & WPR & true & 14 \\
\texttt{PAIR\_sta\allowbreak{}t\_classi\allowbreak{}calGauss\allowbreak{}ianWickP\allowbreak{}lus} & WPR & true & 14, 18 \\
\texttt{PAIR\_sta\allowbreak{}t\_expect\allowbreak{}ationRul\allowbreak{}eCovaria\allowbreak{}nceIndef\allowbreak{}inite} & WPR & true & 14, 15 \\
\texttt{PAIR\_sta\allowbreak{}t\_expect\allowbreak{}ationRul\allowbreak{}eTraces} & WPR & true & 14 \\
\texttt{PAIR\_sta\allowbreak{}t\_fermio\allowbreak{}nWickMin\allowbreak{}us} & WPR & true & 14, 18 \\
\texttt{PAIR\_sta\allowbreak{}t\_filled\allowbreak{}ShellExc\allowbreak{}hangeRat\allowbreak{}io} & WPR & true & 14 \\
\texttt{PAIR\_sta\allowbreak{}t\_fixedA\allowbreak{}mplitude\allowbreak{}PhasesDe\allowbreak{}viate} & WPR & true & 14 \\
\texttt{PAIR\_sta\allowbreak{}t\_ldaPot\allowbreak{}entials} & WPR & true & 14, 15 \\
\texttt{PAIR\_sta\allowbreak{}t\_single\allowbreak{}ModeMome\allowbreak{}nts} & WPR & true & 14 \\
\texttt{PAIR\_sta\allowbreak{}t\_T0Tota\allowbreak{}lDerivat\allowbreak{}ive} & WPR & true & 14 \\
\texttt{PAIR\_sta\allowbreak{}t\_unifor\allowbreak{}mGasExch\allowbreak{}ange} & WPR & true & 14, 15 \\
\texttt{PAIR\_T1g\allowbreak{}eneric\_*} & WPR & all 9 true & 7, 16 \\
\texttt{PAIR\_T1g\allowbreak{}eneric\_c\allowbreak{}onjugate\allowbreak{}FieldEqu\allowbreak{}ation} & WPR & true & 7, 15, 16 \\
\texttt{PAIR\_T1g\allowbreak{}eneric\_c\allowbreak{}urrent} & WPR & true & 7, 15, 16 \\
\texttt{PAIR\_T1g\allowbreak{}eneric\_e\allowbreak{}mtAll36} & WPR & true & 7, 15, 16 \\
\texttt{PAIR\_T1g\allowbreak{}eneric\_f\allowbreak{}ieldEqua\allowbreak{}tion} & WPR & true & 7, 15, 16 \\
\texttt{PAIR\_T1g\allowbreak{}eneric\_l\allowbreak{}agrangia\allowbreak{}n} & WPR & true & 6, 15, 16 \\
\texttt{PAIR\_T1g\allowbreak{}eneric\_n\allowbreak{}aiveFixe\allowbreak{}dLambdaF\allowbreak{}ails} & WPR & true & 6, 15, 16 \\
\texttt{PAIR\_T1g\allowbreak{}rassmann\allowbreak{}\_*} & WPR & all 7 true & 7, 16 \\
\texttt{PAIR\_T1g\allowbreak{}rassmann\allowbreak{}\_current} & WPR & true & 7, 15, 16 \\
\texttt{PAIR\_T1g\allowbreak{}rassmann\allowbreak{}\_diracOp\allowbreak{}erator} & WPR & true & 7, 15, 16 \\
\texttt{PAIR\_T1g\allowbreak{}rassmann\allowbreak{}\_emt} & WPR & true & 7, 15, 16 \\
\texttt{PAIR\_T1g\allowbreak{}rassmann\allowbreak{}\_lagrang\allowbreak{}ian} & WPR & true & 6, 15, 16 \\
\texttt{PAIR\_T1g\allowbreak{}rassmann\allowbreak{}\_naiveFi\allowbreak{}xedLambd\allowbreak{}aFails} & WPR & true & 15 \\
\texttt{PAIR\_T1j\allowbreak{}ets\_*} & WPR & all 13 true & 7, 15, 16 \\
\texttt{PAIR\_T1j\allowbreak{}ets\_cons\allowbreak{}ervation\allowbreak{}Both} & WPR & true & 7 \\
\texttt{PAIR\_T1j\allowbreak{}ets\_emt} & WPR & true & 7, 16 \\
\texttt{PAIR\_T1j\allowbreak{}ets\_fiel\allowbreak{}dEquatio\allowbreak{}ns} & WPR & true & 7, 16 \\
\texttt{PAIR\_T1j\allowbreak{}ets\_gamm\allowbreak{}a8WithFr\allowbreak{}ameSignI\allowbreak{}sSymmetr\allowbreak{}y} & WPR & true & 6, 15, 16 \\
\texttt{PAIR\_T1j\allowbreak{}ets\_lagr\allowbreak{}angian} & WPR & true & 6, 16 \\
\texttt{PAIR\_T1j\allowbreak{}ets\_naiv\allowbreak{}eFixedLa\allowbreak{}mbdaFail\allowbreak{}s} & WPR & true & 15 \\
\texttt{PAIR\_T1j\allowbreak{}ets\_onSh\allowbreak{}ellImage} & WPR & true & 7, 15 \\
\texttt{PAIR\_T1j\allowbreak{}ets\_pair\allowbreak{}EMTAndCu\allowbreak{}rrentVan\allowbreak{}ish} & WPR & true & 15, 16 \\
\texttt{PAIR\_T1j\allowbreak{}ets\_viel\allowbreak{}beinSign\allowbreak{}FlipGeom\allowbreak{}etry} & WPR & true & 6, 15 \\
\texttt{PAIR\_T1j\allowbreak{}ets\_viel\allowbreak{}beinSign\allowbreak{}FlipIsT1} & WPR & true & 6, 15, 16 \\
\texttt{PAIR\_T1k\allowbreak{}rein\_*} & WPR & all 12 true & 16, 17, 18, 19 \\
\texttt{PAIR\_T1k\allowbreak{}rein\_ima\allowbreak{}geAntico\allowbreak{}mmutator\allowbreak{}MinusB} & WPR & true & 15, 16, 17 \\
\texttt{PAIR\_T1k\allowbreak{}rein\_ima\allowbreak{}geExpect\allowbreak{}ationVal\allowbreak{}ues} & WPR & true & 15, 16, 17 \\
\texttt{PAIR\_T1k\allowbreak{}rein\_ima\allowbreak{}geOperat\allowbreak{}orIdenti\allowbreak{}ties} & WPR & true & 15, 16, 17 \\
\texttt{PAIR\_T1k\allowbreak{}rein\_ind\allowbreak{}ependent\allowbreak{}CARPlusB} & WPR & true & 15, 16, 17 \\
\texttt{PAIR\_T1k\allowbreak{}rein\_ind\allowbreak{}ependent\allowbreak{}Expectat\allowbreak{}ionValue\allowbreak{}s} & WPR & true & 15, 16, 17 \\
\texttt{PAIR\_T1k\allowbreak{}rein\_min\allowbreak{}usMModes} & WPR & true & 15 \\
\texttt{PAIR\_T1p\allowbreak{}rimordia\allowbreak{}l\_*} & WPR & all 6 true & 15, 16 \\
\texttt{PAIR\_T1p\allowbreak{}rimordia\allowbreak{}l\_curren\allowbreak{}t} & WPR & true & 16 \\
\texttt{PAIR\_T1p\allowbreak{}rimordia\allowbreak{}l\_diracO\allowbreak{}perator} & WPR & true & 16 \\
\texttt{PAIR\_T1p\allowbreak{}rimordia\allowbreak{}l\_emt64} & WPR & true & 7, 16 \\
\texttt{PAIR\_T1p\allowbreak{}rimordia\allowbreak{}l\_lagran\allowbreak{}gian} & WPR & true & 6, 16 \\
\texttt{PAIR\_T2f\allowbreak{}rame\_*} & WPR & all 12 true & 6, 16 \\
\texttt{PAIR\_T2f\allowbreak{}rame\_cha\allowbreak{}racterIs\allowbreak{}MinusNor\allowbreak{}m} & WPR & true & 15, 16 \\
\texttt{PAIR\_T2f\allowbreak{}rame\_fra\allowbreak{}meReflec\allowbreak{}tionsGeo\allowbreak{}metry} & WPR & true & 15 \\
\texttt{PAIR\_T2f\allowbreak{}rame\_gam\allowbreak{}ma8Times\allowbreak{}Untwiste\allowbreak{}dSpaceli\allowbreak{}ke} & WPR & true & 15 \\
\texttt{PAIR\_T2f\allowbreak{}rame\_onS\allowbreak{}hellImag\allowbreak{}e} & WPR & true & 15, 16 \\
\texttt{PAIR\_T2f\allowbreak{}rame\_sca\allowbreak{}larAndCu\allowbreak{}rrentSig\allowbreak{}ns} & WPR & true & 16 \\
\texttt{PAIR\_T2f\allowbreak{}rame\_twi\allowbreak{}stedSpac\allowbreak{}elikeMap\allowbreak{}sToMinus\allowbreak{}MSameLam\allowbreak{}bda} & WPR & true & 6, 15, 16 \\
\texttt{PAIR\_T2f\allowbreak{}rame\_twi\allowbreak{}stedTime\allowbreak{}like} & WPR & true & 6, 15 \\
\texttt{PAIR\_T2f\allowbreak{}rame\_unt\allowbreak{}wistedSp\allowbreak{}acelikeC\allowbreak{}ontractE\allowbreak{}3} & WPR & true & 6, 15 \\
\texttt{PAIR\_T2z\allowbreak{}2\_*} & WPR & all 11 true & 16, 18 \\
\texttt{PAIR\_T2z\allowbreak{}2\_curren\allowbreak{}tPullbac\allowbreak{}k} & WPR & true & 15, 16 \\
\texttt{PAIR\_T2z\allowbreak{}2\_diracO\allowbreak{}perator} & WPR & true & 15, 16 \\
\texttt{PAIR\_T2z\allowbreak{}2\_emtPul\allowbreak{}lback} & WPR & true & 15, 16 \\
\texttt{PAIR\_T2z\allowbreak{}2\_geomet\allowbreak{}ry} & WPR & true & 15 \\
\texttt{PAIR\_T2z\allowbreak{}2\_lagran\allowbreak{}gianEven} & WPR & true & 15 \\
\texttt{PAIR\_T2z\allowbreak{}2\_massFu\allowbreak{}nctionMa\allowbreak{}p} & WPR & true & 15, 16 \\
\texttt{PAIR\_T2z\allowbreak{}2\_PBfiel\allowbreak{}dLevelFl\allowbreak{}ipsLambd\allowbreak{}a} & WPR & true & 15 \\
\texttt{PAIR\_T2z\allowbreak{}2\_rulePa\allowbreak{}rities} & WPR & true & 15 \\
\texttt{PAIR\_T2z\allowbreak{}2\_sameMa\allowbreak{}ssFails} & WPR & true & 15 \\
\texttt{PAIR\_T2z\allowbreak{}2\_scalar\allowbreak{}Odd} & WPR & true & 15 \\
\texttt{PAIR\_T2z\allowbreak{}2\_symmet\allowbreak{}ricIffOd\allowbreak{}dMass} & WPR & true & 15, 16 \\
\texttt{PAIR\_T3b\allowbreak{}lock\_*} & WPR & all 18 true & 16 \\
\texttt{PAIR\_T3b\allowbreak{}lock\_bag\allowbreak{}AngleMap} & WPR & true & 15 \\
\texttt{PAIR\_T3b\allowbreak{}lock\_den\allowbreak{}sityAndC\allowbreak{}urrentMa\allowbreak{}ps} & WPR & true & 15 \\
\texttt{PAIR\_T3b\allowbreak{}lock\_gam\allowbreak{}ma1IsSig\allowbreak{}ma1InEve\allowbreak{}ryBlock} & WPR & true & 15 \\
\texttt{PAIR\_T3b\allowbreak{}lock\_gam\allowbreak{}ma8IsSig\allowbreak{}ma2Betwe\allowbreak{}enPartne\allowbreak{}rBlocks} & WPR & true & 15 \\
\texttt{PAIR\_T3b\allowbreak{}lock\_ham\allowbreak{}iltonian\allowbreak{}Sigma2} & WPR & true & 15 \\
\texttt{PAIR\_T3b\allowbreak{}lock\_par\allowbreak{}ityMap} & WPR & true & 15 \\
\texttt{PAIR\_T3b\allowbreak{}lock\_pot\allowbreak{}entialsI\allowbreak{}mageRule} & WPR & true & 15 \\
\texttt{PAIR\_T3b\allowbreak{}lock\_pot\allowbreak{}entialsS\allowbreak{}tandardR\allowbreak{}ule} & WPR & true & 15 \\
\texttt{PAIR\_T3b\allowbreak{}lock\_sig\allowbreak{}ma2IsRus\allowbreak{}tSwap} & WPR & true & 15 \\
\texttt{PAIR\_T3b\allowbreak{}lock\_sig\allowbreak{}ma2MapsB\allowbreak{}lockODE} & WPR & true & 15 \\
\texttt{PAIR\_T3b\allowbreak{}lock\_sta\allowbreak{}tisticsC\allowbreak{}oefficie\allowbreak{}nts} & WPR & true & 14 \\
\texttt{PAIR\_T3e\allowbreak{}mt\_image\allowbreak{}RuleMinu\allowbreak{}sT} & WPR & true & 15 \\
\texttt{PAIR\_T3e\allowbreak{}mt\_stage\allowbreak{}4CrossCh\allowbreak{}eck} & WPR & true & 15 \\
\texttt{PAIR\_T3e\allowbreak{}mt\_stand\allowbreak{}ardRuleP\allowbreak{}lusT} & WPR & true & 15 \\
\texttt{PAIR\_T3k\allowbreak{}s\_*} & WPR & all 19 true & 16 \\
\texttt{PAIR\_T3k\allowbreak{}s\_contro\allowbreak{}lDiffers\allowbreak{}FromPair\allowbreak{}edProble\allowbreak{}m} & WPR & true & 15 \\
\texttt{PAIR\_T3k\allowbreak{}s\_contro\allowbreak{}lMixedSe\allowbreak{}ctorLeve\allowbreak{}lsDisjoi\allowbreak{}nt} & WPR & true & 15 \\
\texttt{PAIR\_T3k\allowbreak{}s\_contro\allowbreak{}lSplitti\allowbreak{}ngClosed\allowbreak{}Form} & WPR & true & 15, 16 \\
\texttt{PAIR\_T3k\allowbreak{}s\_contro\allowbreak{}lSubGapB\allowbreak{}oundStat\allowbreak{}e} & WPR & true & 15, 18 \\
\texttt{PAIR\_T3k\allowbreak{}s\_imageR\allowbreak{}uleEnerg\allowbreak{}yOdd} & WPR & true & 15, 16 \\
\texttt{PAIR\_T3k\allowbreak{}s\_massiv\allowbreak{}eLevelsM\allowbreak{}ap} & WPR & true & 15 \\
\texttt{PAIR\_T3k\allowbreak{}s\_mixedS\allowbreak{}ectorSol\allowbreak{}utions} & WPR & true & 15 \\
\texttt{PAIR\_T3k\allowbreak{}s\_occupa\allowbreak{}tionsAnd\allowbreak{}Temperat\allowbreak{}ure} & WPR & true & 15 \\
\texttt{PAIR\_T3k\allowbreak{}s\_sigma1\allowbreak{}Function\allowbreak{}alInvari\allowbreak{}ant} & WPR & true & 15 \\
\texttt{PAIR\_T3k\allowbreak{}s\_sigma2\allowbreak{}Densitie\allowbreak{}s} & WPR & true & 15 \\
\texttt{PAIR\_T3k\allowbreak{}s\_sigma2\allowbreak{}Function\allowbreak{}alInvari\allowbreak{}ant} & WPR & true & 14, 15, 16 \\
\texttt{PAIR\_T3k\allowbreak{}s\_sigma2\allowbreak{}KSOperat\allowbreak{}orEquiva\allowbreak{}riant} & WPR & true & 14, 15, 16 \\
\texttt{PAIR\_T3k\allowbreak{}s\_statio\allowbreak{}narityBo\allowbreak{}thStatis\allowbreak{}tics} & WPR & true & 14, 15 \\
\texttt{PAIR\_T3k\allowbreak{}s\_zeroMo\allowbreak{}deImage} & WPR & true & 15 \\
\texttt{PAIR\_T3k\allowbreak{}s\_zeroMo\allowbreak{}deSplitt\allowbreak{}ingMapsE\allowbreak{}xactly} & WPR & true & 15 \\
\texttt{PAIR\_T3k\allowbreak{}s\_zeroMo\allowbreak{}deUntran\allowbreak{}sformedC\allowbreak{}ontrol} & WPR & true & 15, 16 \\
\texttt{PAIR\_tot\allowbreak{}als\_fiel\allowbreak{}dLevelCh\allowbreak{}iralPair} & WPR & true & 15, 16, 18 \\
\texttt{PAIR\_tot\allowbreak{}als\_fiel\allowbreak{}dLevelMi\allowbreak{}rrorPair} & WPR & true & 15, 18 \\
\texttt{PAIR\_tot\allowbreak{}als\_ksKr\allowbreak{}einImage\allowbreak{}Pair} & WPR & true & 15, 16, 18 \\
\texttt{PAIR\_tot\allowbreak{}als\_ksMi\allowbreak{}rrorPair} & WPR & true & 15, 16 \\
\texttt{S5\_agree\allowbreak{}sWithWol\allowbreak{}fram} & P00 & true & 6 \\
\texttt{S5\_algeb\allowbreak{}ra\_antic\allowbreak{}ommutato\allowbreak{}rTotally\allowbreak{}Antisymm\allowbreak{}etric} & P00 & true & 6 \\
\texttt{S5\_algeb\allowbreak{}ra\_pinMo\allowbreak{}duleAndS\allowbreak{}pinInvar\allowbreak{}iance} & P00 & true & 6 \\
\texttt{S5\_conne\allowbreak{}ction\_no\allowbreak{}nTrivial\allowbreak{}Coupling} & P00 & true & 6 \\
\texttt{S5\_conne\allowbreak{}ction\_Om\allowbreak{}egaTermI\allowbreak{}nFieldEq\allowbreak{}uation} & P00 & true & 6 \\
\texttt{S5\_EL\_gr\allowbreak{}assmannC\allowbreak{}urved} & P00 & true & 7 \\
\texttt{S5\_EL\_id\allowbreak{}enticalF\allowbreak{}ormBothS\allowbreak{}tatistic\allowbreak{}s} & P00 & true & 7 \\
\texttt{S5\_EMT\_c\allowbreak{}onservat\allowbreak{}ionOnShe\allowbreak{}ll} & P00 & true & 7 \\
\texttt{S5\_EMT\_h\allowbreak{}omogeneo\allowbreak{}usEquati\allowbreak{}onsOfSta\allowbreak{}te} & P00 & true & 7 \\
\texttt{S5\_EMT\_s\allowbreak{}ymmetric\allowbreak{}AndReal} & P00 & true & 7 \\
\texttt{S5\_EMT\_t\allowbreak{}raceOnSh\allowbreak{}ell} & P00 & true & 7 \\
\texttt{S5\_EMT\_v\allowbreak{}ielbeinV\allowbreak{}ariation} & P00 & true & 7 \\
\texttt{S5\_energ\allowbreak{}y\_unboun\allowbreak{}dedBelow} & P00 & true & 18 \\
\texttt{S5\_lagra\allowbreak{}ngian\_co\allowbreak{}efficien\allowbreak{}tHermiti\allowbreak{}city} & P00 & true & 6 \\
\texttt{S5\_lagra\allowbreak{}ngian\_re\allowbreak{}alCommut\allowbreak{}ing} & P00 & true & 6 \\
\texttt{S5\_massT\allowbreak{}erm\_comm\allowbreak{}utingExp\allowbreak{}licitSpi\allowbreak{}nors} & P00 & true & 6 \\
\texttt{S5\_massT\allowbreak{}erm\_disp\allowbreak{}ersionFl\allowbreak{}at} & P00 & true & 6 \\
\texttt{S5\_massT\allowbreak{}erm\_gras\allowbreak{}smannNon\allowbreak{}zero} & P00 & true & 6 \\
\texttt{S5\_pairi\allowbreak{}ngAgrees\allowbreak{}WithWolf\allowbreak{}ram} & PPR & true & 15 \\
\texttt{S5\_primo\allowbreak{}rdial\_ho\allowbreak{}mogeneou\allowbreak{}sState} & P00 & true & 7 \\
\texttt{S5\_primo\allowbreak{}rdial\_st\allowbreak{}aticFiel\allowbreak{}dSourced\allowbreak{}Exactly} & P00 & true & 7, 18 \\
\texttt{S5\_realR\allowbreak{}estricti\allowbreak{}on\_commu\allowbreak{}tingNonT\allowbreak{}rivial} & P00 & true & 6 \\
\texttt{S5\_realR\allowbreak{}estricti\allowbreak{}on\_commu\allowbreak{}tingNote\allowbreak{}bookCont\allowbreak{}raction} & P00 & true & 6 \\
\texttt{S5\_realR\allowbreak{}estricti\allowbreak{}on\_grass\allowbreak{}mannComp\allowbreak{}lexDecom\allowbreak{}position} & P00 & true & 6 \\
\texttt{\detokenize{S5_stat}} & PPR & true & 14 \\
\texttt{S5\_stat\_\allowbreak{}*} & PPR & all 15 true & 14 \\
\texttt{S5\_stat\_\allowbreak{}bosonThe\allowbreak{}rmalWick\allowbreak{}Plus} & PPR & true & 14 \\
\texttt{S5\_stat\_\allowbreak{}classica\allowbreak{}lGaussia\allowbreak{}nWickPlu\allowbreak{}s} & PPR & true & 14 \\
\texttt{S5\_stat\_\allowbreak{}expectat\allowbreak{}ionRuleK\allowbreak{}reinFock} & PPR & true & 14 \\
\texttt{S5\_stat\_\allowbreak{}expectat\allowbreak{}ionRuleT\allowbreak{}races} & PPR & true & 14 \\
\texttt{S5\_stat\_\allowbreak{}fermionW\allowbreak{}ickMinus} & PPR & true & 14 \\
\texttt{S5\_stati\allowbreak{}c\_classi\allowbreak{}calEnerg\allowbreak{}yKreinSi\allowbreak{}gned} & P00 & true & 14 \\
\texttt{S5\_stati\allowbreak{}c\_classi\allowbreak{}calModeI\allowbreak{}sKreinWe\allowbreak{}ightedKS} & P00 & true & 14 \\
\texttt{S5\_stati\allowbreak{}c\_geomet\allowbreak{}ryAndRed\allowbreak{}ucedEqua\allowbreak{}tion} & P00 & true & 7 \\
\texttt{S5\_stati\allowbreak{}c\_source\allowbreak{}ProperCh\allowbreak{}art} & P00 & true & 7 \\
\texttt{S5\_stat\_\allowbreak{}ldaPoten\allowbreak{}tials} & PPR & true & 14 \\
\texttt{S5\_stat\_\allowbreak{}uniformG\allowbreak{}asExactF\allowbreak{}inite} & PPR & true & 14 \\
\texttt{S5\_stat\_\allowbreak{}uniformG\allowbreak{}asExchan\allowbreak{}ge} & PPR & true & 14 \\
\texttt{S5\_T1gen\allowbreak{}eric} & PPR & true & 16 \\
\texttt{S5\_T1gen\allowbreak{}eric\_eul\allowbreak{}erLagran\allowbreak{}geFromLa\allowbreak{}grangian} & PPR & true & 15 \\
\texttt{S5\_T1gen\allowbreak{}eric\_fie\allowbreak{}ldEquati\allowbreak{}on} & PPR & true & 7 \\
\texttt{S5\_T1gen\allowbreak{}eric\_lag\allowbreak{}rangian} & PPR & true & 6 \\
\texttt{S5\_T1gra\allowbreak{}ssmann} & PPR & true & 15, 16 \\
\texttt{S5\_T1jet\allowbreak{}s} & PPR & true & 16 \\
\texttt{S5\_T1kre\allowbreak{}in\_*} & PPR & all 13 true & 15, 18, 19 \\
\texttt{S5\_T1kre\allowbreak{}in\_image\allowbreak{}Operator\allowbreak{}Identiti\allowbreak{}es} & PPR & true & 15 \\
\texttt{S5\_T1pri\allowbreak{}mordial} & PPR & true & 16 \\
\texttt{S5\_T2fra\allowbreak{}me} & PPR & true & 16 \\
\texttt{S5\_T2z2\_\allowbreak{}fieldEqu\allowbreak{}ationMap\allowbreak{}sToMinus\allowbreak{}MSameLam\allowbreak{}bda} & PPR & true & 15 \\
\texttt{S5\_T2z2\_\allowbreak{}kineticE\allowbreak{}ven} & PPR & true & 15 \\
\texttt{S5\_T2z2\_\allowbreak{}PBRuleMa\allowbreak{}trixEven\allowbreak{}COdd} & PPR & true & 15 \\
\texttt{S5\_T3blo\allowbreak{}ck} & PPR & true & 16 \\
\texttt{\detokenize{S5_T3ks}} & PPR & true & 16 \\
\bottomrule
\end{longtable}
\endgroup

\subsection{Checks of the tools, and the state of the index}

\textbf{Checks that no committed report stores.} Chapter 19 also names checks of the two programs that build this book. They are printed as lines \texttt{check\_<n\allowbreak{}ame>=tru\allowbreak{}e} or \texttt{\detokenize{false}} when the program runs and are not stored in a committed file:

\begingroup
\small
\begin{longtable}{@{}>{\raggedright\arraybackslash}p{0.200\linewidth}>{\raggedright\arraybackslash}p{0.200\linewidth}>{\raggedright\arraybackslash}p{0.200\linewidth}>{\raggedright\arraybackslash}p{0.200\linewidth}@{}}
\toprule
\textbf{Check} & \textbf{Program that prints it} & \textbf{What it requires} & \textbf{Cited in chapters} \\
\midrule
\endfirsthead
\toprule
\textbf{Check} & \textbf{Program that prints it} & \textbf{What it requires} & \textbf{Cited in chapters} \\
\midrule
\endhead
\texttt{crossRef\allowbreak{}erencesR\allowbreak{}esolve} & \texttt{scripts/\allowbreak{}build\_te\allowbreak{}xtbook.p\allowbreak{}y} & every reference "Chapter N" and "Section N.M" of the book names an existing chapter or section & 19 \\
\texttt{editionR\allowbreak{}egistere\allowbreak{}d} & \texttt{scripts/\allowbreak{}build\_pr\allowbreak{}ovenance\allowbreak{}\_pdf.py} & the document has an edition in \texttt{provenan\allowbreak{}ce/pdf-s\allowbreak{}pecifica\allowbreak{}tions.js\allowbreak{}on} & 19 \\
\texttt{register\allowbreak{}edPath} & \texttt{scripts/\allowbreak{}build\_pr\allowbreak{}ovenance\allowbreak{}\_pdf.py} & the registered edition names this PDF file & 19 \\
\texttt{register\allowbreak{}edPageCo\allowbreak{}unt} & \texttt{scripts/\allowbreak{}build\_pr\allowbreak{}ovenance\allowbreak{}\_pdf.py} & the PDF has the registered number of pages & 19 \\
\texttt{register\allowbreak{}edSha256} & \texttt{scripts/\allowbreak{}build\_pr\allowbreak{}ovenance\allowbreak{}\_pdf.py} & the PDF has the registered sha256 & 19 \\
\texttt{provenan\allowbreak{}cePdfCop\allowbreak{}y} & \texttt{scripts/\allowbreak{}build\_pr\allowbreak{}ovenance\allowbreak{}\_pdf.py} & the verified PDF was copied to its final place & 19 \\
\bottomrule
\end{longtable}
\endgroup

\textbf{The commit of the index.} The index of Sections 20.10 to 20.14 was made from the committed reports of commit \texttt{\detokenize{1f2dd69}} (2026-09-30), the commit from which Chapter 19 ran its commands, and from the chapter files of this edition. Two later changes of the reports are known in advance. The Rust runs of the subcommand \texttt{\detokenize{pairs}} of Stage 5 are still being added; their files contain no entry \texttt{\detokenize{checks}} and do not enter the index. And the physical reading of the Stage-5 Fock-level result \texttt{\detokenize{T1krein}} is to be corrected in the Stage-5 programs (\texttt{HANDOFF.\allowbreak{}md}, section 0.4, item C; Section 19.10), which can change the names and the number of the checks \texttt{PAIR\_T1k\allowbreak{}rein\_*} and \texttt{S5\_T1kre\allowbreak{}in\_*} and of the matter--antimatter checks that cite them.

\textbf{Counting.} The five tables list 583 rows: 558 single checks and 25 families. Every single check listed is true in every committed file shown, except \texttt{canonica\allowbreak{}l\_eigenv\allowbreak{}alues} (Section 19.9); every family is entirely true, except the two families \texttt{gauntlet\allowbreak{}\_*} and \texttt{\detokenize{r07_*}} of the Jupyter notebook's Stage-4 report, whose six false checks Section 19.9 explains.

\subsection{What we proved and what we assumed}

This chapter proves nothing new. The glossary repeats, in one or two sentences each, definitions made in Chapters 0 to 19; where an entry and a chapter seem to differ, the chapter is authoritative, and the entry names it. In particular the entries on the pairs of universes (image field, Krein image pair, source-free, Theorems M1 to M6) repeat the precise statements of Chapters 15 to 17: the cancellation of energy, momentum and charge in a chirality pair is proved for two classical fields only, and every zero at the quantum or Kohn--Sham level is of the kind $X+(-X)=0$. The index of checks is a \textbf{computed} list: every row was produced by a program from the committed reports and the chapter files of this edition, by the rule of Section 20.9, and checked by a second program; the value shown is the value in the committed files shown. It assumes that the chapter files and the reports named in Section 20.15 are the ones of the edition you read; a later commit can add checks, add citations or change a value, and then the lookup command of Section 20.9 gives the current answer. That a check is listed as true says only that the program that wrote the report found it true; what the check establishes, and under which assumptions, is stated in the chapter that cites it.

\subsection{Exercises}

\textbf{Exercise 20.1.} Use the glossary to explain the difference between the Hilbert norm and the Krein norm of a mode amplitude. Which of the two is conserved for a mode with momentum along an extra time, and why does that not make such a mode harmless?

\textbf{Exercise 20.2.} Give the status word of Chapter 0 (PROVED, COMPUTED, ASSUMED, HYPOTHESIS or OPEN) for each statement: (a) the chirality map turns every solution with $(m,\lambda)$ into one with $(-m,-\lambda)$; (b) the Kohn--Sham gap of the free ground state with $N=8$; (c) the restriction of the numerical chapters to the good sector; (d) universes of masses $+M$ and $-M$ are created in pairs at $x_4=0$; (e) there is a consistent interacting quantum theory of dirac16complex.

\textbf{Exercise 20.3.} The symbols $\eta$ and $z$ each have two meanings in this book. Name them and say how the reader can tell them apart.

\textbf{Exercise 20.4.} Find in the index the one single check that is false in its committed report. Which programs are compared in it, and in which chapters is it discussed?

\textbf{Exercise 20.5.} The index lists the family \texttt{*\_energy\allowbreak{}\_from\_rh\allowbreak{}o} with 43 checks in two files. Use the table of Stage 4 and the lookup command of Section 20.9 to explain the number.

\textbf{Exercise 20.6.} Change the lookup command of Section 20.9 so that it counts, in every committed file, the checks whose names begin with a given text, and how many of them are true. Run it for \texttt{PAIR\_T1k\allowbreak{}rein\_} and compare the result with the index and with Chapter 17.

\textbf{Exercise 20.7.} Why does the statistics sign decide whether complex conjugation is a symmetry of the Lagrangian? Answer with the glossary entries for statistics sign and charge conjugation, and name the chapter where the proof is.

\textbf{Exercise 20.8.} A classmate writes: "The index shows that \texttt{PAIR\_T1k\allowbreak{}rein\_*} is all true, so the pair of universes has zero energy at the quantum level." Use the glossary entries for the image field, the reversed Krein metric and the Krein image pair to explain what the checks establish and what they do not.

\subsection{Answers to the exercises}

\textbf{Answer 20.1.} The Hilbert norm is $u^\dagger u$, the norm of the positive inner product; the Krein norm is $u^\dagger Bu$ with the indefinite matrix $B=-iC\gamma^4$ (Chapter 8; Chapter 11 uses it for EXP-5). For every mode $h^\dagger B=Bh$, so the Krein norm is conserved always; the Hilbert norm is conserved only when the mode Hamiltonian is Hermitian, which holds in the good sector. A mode with enough momentum along an extra time has $E^2<0$ and grows exponentially (Chapter 8): in the committed EXP-5 run quoted there its Hilbert norm grows to $1.258\times10^{16}$, while its Krein norm keeps its initial value, up to a drift below $10^{-8}$ times the Hilbert norm, that is up to rounding (check \texttt{kreinCon\allowbreak{}servedNo\allowbreak{}rmalized} of the EXP-5 checker). A conserved quantity that can be negative does not bound the size of the solution, because large positive and negative contributions can cancel in it; so its conservation does not make the mode harmless.

\textbf{Answer 20.2.} (a) PROVED (Theorem T1, Chapter 15, with exact checks). (b) COMPUTED (Chapter 13; a number of the Rust solver with its own checks and the cross-check status of Section 19.9). (c) ASSUMED: the restriction is imposed, not derived (Chapter 8). (d) HYPOTHESIS: the notebook's claim, which is not derived, and no creation process is computed. Chapter 16 shows what the conservation laws say about it (Proposition 16.5, for classical fields): a T1 pair has the totals of the empty state, so the total energy--momentum, the total charge and the gravitational constraints allow it; but each member's charge is conserved separately, so in the theory as built the pair cannot appear from the empty state unless $Q_+=0$. (e) OPEN (Chapter 18).

\textbf{Answer 20.3.} $\eta$ is the flat metric $\mathrm{diag}(+1,+1,+1,+1,-1,-1,-1,-1)$ everywhere except in Chapters 17 and 18, where $\eta$ also denotes the baryon-to-photon ratio ($6.12\times10^{-10}$ in Chapter 17, $\eta\approx6\times10^{-10}$ in Chapter 18); the metric always carries indices ($\eta_{ab}$, $\eta^{ab}$) or appears as the matrix $\eta$, the ratio is a single number. $z=6Hx_0$ is the hidden angle of the primordial field in Chapter 9, and in Chapter 11 $z$ is the redshift, $1+z=1/a$; the chapter and the context (geometry of the notebook, or cosmology of the late universe) decide.

\textbf{Answer 20.4.} \texttt{canonica\allowbreak{}l\_eigenv\allowbreak{}alues} in \texttt{kohn-sha\allowbreak{}m/python\allowbreak{}-check-r\allowbreak{}eport.js\allowbreak{}on} is false (Section 20.13). It compares the Kohn--Sham levels of the Rust solver with those of the independent Python reference solver; one deep level of the smeared $N=1016$ ensemble differs by more than the tolerance (Section 19.9). The index lists Chapters 14, 16, 18 and 19 as citing it.

\textbf{Answer 20.5.} The Rust solver writes one self-check \texttt{<run>\_en\allowbreak{}ergy\_fro\allowbreak{}m\_rho} for every run of the subcommand \texttt{\detokenize{scf}} (33 runs) and of the subcommand \texttt{\detokenize{emt}} (10 runs), in \texttt{rust/scf\allowbreak{}/summary\allowbreak{}.json} and \texttt{rust/emt\allowbreak{}/summary\allowbreak{}.json}: $33+10=43$. The check requires that the integral of the energy density equals the total energy of the run (Chapter 13). Since the names begin with the name of the run, the lookup command of Section 20.9 must be given a full name, such as \texttt{m1\_L3\_N1\allowbreak{}12\_lamp1\allowbreak{}\_T0\_ener\allowbreak{}gy\_from\_\allowbreak{}rho}; it then prints both files, because this run appears in both.

\textbf{Answer 20.6.} Replace the last two lines of the program:

\begin{Verbatim}[fontsize=\small]
python -c "
import json, pathlib, sys
for p in sorted(pathlib.Path('artifacts').rglob('*.json')):
    d = json.loads(p.read_text(encoding='utf-8'))
    c = d.get('checks') if isinstance(d, dict) else None
    if isinstance(c, dict):
        hits = [v for k, v in c.items() if k.startswith(sys.argv[1])]
        if hits:
            print(p.as_posix(), len(hits), 'true:', hits.count(True))
" PAIR_T1krein_
\end{Verbatim}

In both shells it printed

\begin{Verbatim}[fontsize=\small]
artifacts/dirac16complex/pair-creation/wolfram-pairing-report.json 12 true: 12
\end{Verbatim}

that is twelve checks, all true, in one file, as the index says and as Chapter 17 states ("all twelve \texttt{PAIR\_T1k\allowbreak{}rein} checks true").

\textbf{Answer 20.7.} Complex conjugation $C_0:\Psi\to\Psi^\ast$ has to move a conjugated factor past an unconjugated one, and the statistics sign $s$ records the sign of that reordering: $+1$ for commuting components, $-1$ for Grassmann components. Theorem 6.10 shows that $C_0$ sends the kinetic term $K\to sK$ and the scalar density $S\to sS$. For the commuting field ($s=+1$) the Lagrangian is unchanged, so $C_0$ is an exact symmetry that reverses the charge: the charge conjugation of dirac16complex00. For the Grassmann field ($s=-1$) it reverses the sign of the Lagrangian and of $\lambda$, and no constant charge conjugation is a symmetry when $m\ne0$ (Chapter 6; the classification is Theorem M2 of Chapter 17).

\textbf{Answer 20.8.} The twelve checks \texttt{PAIR\_T1k\allowbreak{}rein\_*}, all true in \texttt{pair-cre\allowbreak{}ation/wo\allowbreak{}lfram-pa\allowbreak{}iring-re\allowbreak{}port.jso\allowbreak{}n}, verify in an exact Fock-space model with four rest modes of mass $m=1$ both readings of the $-M$ universe (Section 15.4). For the image field $\Psi_-=\gamma^8\Psi_+$ on the same state space they verify its anticommutator $-B$ (the reversed Krein metric), the operator identities $H[\Psi_-;-m]=-H[\Psi;m]$, $Q[\Psi_-]=-Q[\Psi]$ and $S[\Psi_-]=S[\Psi]$, and the expectation values of the formulas of $\mathcal L_{-m,-\lambda}$, which are $-|\varepsilon|$ for the energy and $-1$ for the charge of each quantum. For the $-M$ theory quantized independently they verify its anticommutator $+B$ and its expectation values. The classmate's conclusion does not follow, for four reasons.

\begin{enumerate}
\item \emph{The image field is the same quantum system as $\Psi_+$.} Measured with its own canonical structure, its energy and its charge are $+H_+$ and $+Q_+$, not $-H_+$ and $-Q_+$, because $H[\Psi_-;-m]$ generates the evolution in $x_4$ backwards and $Q[\Psi_-]$ the inverse phase rotation (Section 15.4). This is proved in an independent exact Fock model by the check \texttt{MA\_M4\_im\allowbreak{}ageField\allowbreak{}FockMode\allowbreak{}l} of \texttt{matter-a\allowbreak{}ntimatte\allowbreak{}r/python\allowbreak{}-matter-\allowbreak{}antimatt\allowbreak{}er-repor\allowbreak{}t.json}. The zero sums $H_++H[\Psi_-;-m]=0$ and $Q_++Q[\Psi_-]=0$ have the form $X+(-X)=0$ and hold in every state; they do not describe two universes of zero total energy.
\item \emph{The independently quantized $-M$ universe does not cancel.} Its kinetic and mass energies are positive and add to those of the $+M$ universe, and its charge cancels that of the $+M$ universe only under an additional assumption on its state (the entry \texttt{kreinLev\allowbreak{}elCaveat} quoted in Chapters 15 and 17).
\item \emph{The Krein image pair of the glossary is a different object.} It is the Kohn--Sham construction $(+M)-(-M)$ of Section 15.8 (check \texttt{PAIR\_tot\allowbreak{}als\_ksKr\allowbreak{}einImage\allowbreak{}Pair}), not the Fock model of \texttt{PAIR\_T1k\allowbreak{}rein\_*}, and its zero energy is again of the kind $X+(-X)=0$. Neither object is a quantum description of two independent universes whose energies and charges cancel; whether one exists is OPEN (Section 16.6), and no creation process is computed anywhere in the repository.
\item \emph{The Stage-5 citation is provisional.} The matter--antimatter analysis cites the Fock-level result of Stage 5 as PROVISIONAL until the Stage-5 gate has passed (measurement \texttt{M4\_krein\allowbreak{}LevelSta\allowbreak{}tus} of \texttt{wolfram-\allowbreak{}matter-a\allowbreak{}ntimatte\allowbreak{}r-report\allowbreak{}.json}), and the physical reading written into the Stage-5 files, the entry \texttt{imageFie\allowbreak{}ld} of the key \texttt{\detokenize{T1krein}} of \texttt{pairing-\allowbreak{}theory.j\allowbreak{}son}, which calls $-|\varepsilon|$ and $-1$ the energy and charge of an "image universe", is to be corrected (\texttt{HANDOFF.\allowbreak{}md}, section 0.4, item C). An all-true family of checks confirms identities; it does not confirm a reading attached to them.
\end{enumerate}

\end{document}
